\documentclass[10pt,oneside]{book}

\usepackage[T1]{fontenc}
\usepackage[utf8]{inputenc}
\usepackage[french]{babel}
\usepackage{amsmath,amssymb,mathrsfs}
\usepackage[all]{xy}
\usepackage{tikz-cd}
\usepackage[a4paper,margin=22mm]{geometry}
\usepackage[
  colorlinks=true,
  linkcolor=blue,
  hypertexnames=false,
  unicode=true,
  pdfencoding=auto
]{hyperref}
\hypersetup{
  pdftitle={Éléments de géométrie algébrique -- édition française cumulative},
  pdfauthor={Alexander Grothendieck et Jean Dieudonné},
  pdfsubject={Texte diplomatique français cumulatif depuis l'autorité NUMDAM}
}

\renewcommand{\thesection}{\arabic{section}}
\renewcommand{\thesubsection}{\thesection.\arabic{subsection}}
\newcommand{\oldpage}[2][]{%
  \setcounter{footnote}{0}%
  \marginpar{\scriptsize\textbf{#1}~|~#2}\ignorespaces}
\newenvironment{env}[1][]{\par\medskip\noindent\textbf{(#1)}\ }{\par}
\newenvironment{definition}[1][]{%
  \par\medskip\noindent\itshape Définition (#1). ---\ }{\par}
\newenvironment{proposition}[1][]{%
  \par\medskip\noindent\itshape Proposition (#1). ---\ }{\par}
\newenvironment{theorem}[1][]{%
  \par\medskip\noindent\itshape Théorème (#1). ---\ }{\par}
\newenvironment{corollary}[1][]{%
  \par\medskip\noindent\itshape Corollaire (#1). ---\ }{\par}
\newenvironment{remark}[1][]{%
  \par\medskip\noindent\itshape Remarque (#1). ---\ }{\par}
\newenvironment{remarks}[1][]{%
  \par\medskip\noindent\itshape Remarques (#1). ---\ }{\par}
\newenvironment{scholium}[1][]{%
  \par\medskip\noindent\textit{Scholie (#1). ---\ }}{\par}
\newenvironment{lemma}[1][]{%
  \par\medskip\noindent\itshape Lemme (#1). ---\ }{\par}
\newenvironment{example}[1][]{%
  \par\medskip\noindent\textit{Exemple (#1). ---\ }}{\par}
\newenvironment{examples}[1][]{%
  \par\medskip\noindent\textit{Exemples (#1). ---\ }}{\par}
\newenvironment{notation}[1][]{%
  \par\medskip\noindent\textit{Notations (#1). ---\ }}{\par}
\newenvironment{proof}{%
  \par\medskip\noindent\textit{Démonstration. ---\ }}{\par}
\newenvironment{namedtheorem}[2][]{%
  \par\medskip\noindent\itshape #2 (#1). ---\ }{\par}
\newenvironment{namedlemma}[2][]{%
  \par\medskip\noindent\itshape #2 (#1). ---\ }{\par}

\begin{document}
\input{ega1/frontmatter-fr.tex}
\input{ega1/intro-fr.tex}
\input{ega1/ega0-1-fr.tex}

\clearpage
\setcounter{section}{0}
\setcounter{subsection}{0}
\setcounter{equation}{0}
\setcounter{footnote}{0}

\input{ega1/chapter1-frontmatter-fr.tex}
\input{ega1/ega1-1-fr.tex}
\input{ega1/ega1-2-fr.tex}
\input{ega1/ega1-3-fr.tex}
\input{ega1/ega1-4-fr.tex}
\input{ega1/ega1-5-fr.tex}
\section{Conditions de finitude}
\label{section:I.6-fr}

\subsection{Préschémas noethériens et localement noethériens.}
\label{subsection:I.6.1-fr}

\begin{definition}[6.1.1]
\label{I.6.1.1-fr}
On dit qu'un préschéma $X$ est \emph{noethérien} (resp. \emph{localement
noethérien}) s'il est réunion finie (resp. réunion) d'ouverts affines
$V_\alpha$ tels que l'anneau de chacun des schémas induits sur les $V_\alpha$
soit noethérien.
\end{definition}

Il résulte aussitôt de (\hyperref[I.1.5.2-fr]{1.5.2}) que si $X$ est
localement noethérien, le faisceau structural $\mathscr{O}_X$ est un
\emph{faisceau cohérent d'anneaux}, la question étant locale. \emph{Tout
sous-$\mathscr{O}_X$-}

\oldpage[I]{141}
\emph{Module (resp. tout $\mathscr{O}_X$-Module quotient) quasi-cohérent}
d'un $\mathscr{O}_X$-Module \emph{cohérent} $\mathscr{F}$ est alors
\emph{cohérent}, car la question est de nouveau locale, et il suffit
d'appliquer (\hyperref[I.1.5.1-fr]{1.5.1}),
(\hyperref[I.1.4.1-fr]{1.4.1}) et
(\hyperref[I.1.3.10-fr]{1.3.10}), joints au fait qu'un sous-module (resp.
module quotient) d'un module de type fini sur un anneau noethérien est de
type fini. Plus particulièrement, \emph{tout faisceau d'idéaux
quasi-cohérent} de $\mathscr{O}_X$ est \emph{cohérent}.

Si un préschéma $X$ est réunion finie (resp. réunion) d'ouverts $W_\lambda$
tels que les préschémas induits sur les $W_\lambda$ soient noethériens (resp.
localement noethériens), il est clair que $X$ est noethérien (resp.
localement noethérien).

\begin{proposition}[6.1.2]
\label{I.6.1.2-fr}
Pour qu'un préschéma $X$ soit noethérien, il faut et il suffit qu'il soit
localement noethérien et que son espace sous-jacent soit quasi-compact.
L'espace sous-jacent de $X$ est alors noethérien.
\end{proposition}

\begin{proof}
La première assertion découle aussitôt des définitions et de
(\hyperref[I.1.1.10-fr]{1.1.10 (ii)}). La seconde résulte de
(\hyperref[I.1.1.6-fr]{1.1.6}) et du fait que tout espace réunion finie de
sous-espaces noethériens est noethérien
(\hyperref[0.2.2.3-fr]{0, 2.2.3}).
\end{proof}

\begin{proposition}[6.1.3]
\label{I.6.1.3-fr}
Soit $X$ un schéma affine d'anneau $A$. Les conditions suivantes sont
équivalentes~: a) $X$ est noethérien~; b) $X$ est localement noethérien~;
c) $A$ est noethérien.
\end{proposition}

\begin{proof}
L'équivalence de a) et b) résulte de (\hyperref[I.6.1.2-fr]{6.1.2}) et de ce
que tout schéma affine a un espace sous-jacent quasi-compact
(\hyperref[I.1.1.10-fr]{1.1.10})~; il est clair en outre que c) entraîne a).
Pour voir que a) entraîne c), remarquons qu'il y a un recouvrement fini
$(V_i)$ de $X$ par des ouverts affines tels que l'anneau $A_i$ du préschéma
induit sur $V_i$ soit noethérien. Soit alors $(\mathfrak{a}_n)$ une suite
croissante d'idéaux de $A$~; il lui correspond canoniquement de façon
biunivoque (\hyperref[I.1.3.7-fr]{1.3.7}) une suite croissante
$(\widetilde{\mathfrak{a}}_n)$ de faisceaux d'idéaux dans
$\widetilde{A}\equiv\mathscr{O}_X$~; pour voir que la suite
$(\mathfrak{a}_n)$ est stationnaire, il suffit de prouver que la suite
$(\widetilde{\mathfrak{a}}_n)$ l'est. Or, la restriction
$\widetilde{\mathfrak{a}}_n|V_i$ est un faisceau quasi-cohérent d'idéaux dans
$\mathscr{O}_X|V_i$, étant l'image réciproque de
$\widetilde{\mathfrak{a}}_n$ par l'injection canonique $V_i\to X$
(\hyperref[0.5.1.4-fr]{0, 5.1.4})~; $\widetilde{\mathfrak{a}}_n|V_i$ est
donc de la forme $\widetilde{\mathfrak{a}}_{ni}$, où $\mathfrak{a}_{ni}$ est
un idéal de $A_i$ (\hyperref[I.1.3.7-fr]{1.3.7}). Comme $A_i$ est noethérien,
la suite $(\mathfrak{a}_{ni})$ est stationnaire pour tout $i$, d'où la
proposition.
\end{proof}

On notera que le raisonnement précédent prouve aussi que \emph{si $X$ est un
préschéma noethérien, toute suite croissante de faisceaux cohérents d'idéaux
de $\mathscr{O}_X$ est stationnaire}.

\begin{proposition}[6.1.4]
\label{I.6.1.4-fr}
Tout sous-préschéma d'un préschéma noethérien (resp. localement noethérien)
est noethérien (resp. localement noethérien).
\end{proposition}

\begin{proof}
Il suffit de faire la démonstration pour un préschéma noethérien $X$~; en
outre, on est aussitôt ramené par la définition
(\hyperref[I.6.1.1-fr]{6.1.1}) au cas où $X$ est un schéma affine. Comme tout
sous-préschéma de $X$ est un sous-préschéma fermé d'un préschéma induit sur
un ouvert (\hyperref[I.4.1.3-fr]{4.1.3}), on peut se borner à considérer le
cas d'un sous-préschéma $Y$ fermé ou induit sur un ouvert de $X$. Le cas où
$Y$ est fermé est immédiat, car si $A$ est l'anneau de $X$, on sait que $Y$
est un schéma affine d'anneau $A/\mathfrak{J}$, où $\mathfrak{J}$ est un
idéal de $A$ (\hyperref[I.4.2.3-fr]{4.2.3})~; comme $A$ est noethérien
(\hyperref[I.6.1.3-fr]{6.1.3}), il en est de même de $A/\mathfrak{J}$.

Supposons maintenant $Y$ ouvert dans $X$~; l'espace sous-jacent $Y$ est
noethérien (\hyperref[I.6.1.2-fr]{6.1.2}), donc quasi-compact, et par suite
réunion finie d'ouverts $D(f_i)$ ($f_i\in A$)~;
\oldpage[I]{142}
tout revient à démontrer la proposition lorsque $Y=D(f)$ avec $f\in A$. Mais
alors $Y$ est un schéma affine dont l'anneau est isomorphe à $A_f$
(\hyperref[I.1.3.6-fr]{1.3.6})~; comme $A$ est noethérien
(\hyperref[I.6.1.3-fr]{6.1.3}), il en est de même de $A_f$.
\end{proof}

\begin{env}[6.1.5]
\label{I.6.1.5-fr}
On notera que le \emph{produit} de deux $S$-préschémas noethériens n'est pas
nécessairement noethérien, même si ces préschémas sont affines, car le produit
tensoriel de deux algèbres noethériennes n'est pas nécessairement un anneau
noethérien (cf. (\hyperref[I.6.3.8-fr]{6.3.8})).
\end{env}

\begin{proposition}[6.1.6]
\label{I.6.1.6-fr}
Si $X$ est un préschéma noethérien, le Nilradical $\mathscr{N}_X$ de
$\mathscr{O}_X$ est nilpotent.
\end{proposition}

On peut en effet recouvrir $X$ par un nombre fini d'ouverts affines $U_i$ et il
suffit de prouver qu'il existe des entiers $n_i$ tels que
$(\mathscr{N}_X|U_i)^{n_i}=0$~; si $n$ est le plus grand des $n_i$, on aura
alors $\mathscr{N}_X^n=0$. On est donc ramené au cas où
$X=\operatorname{Spec}(A)$ est affine, $A$ étant un anneau noethérien~; en
vertu de (\hyperref[I.5.1.1-fr]{5.1.1}) et de
(\hyperref[I.1.3.13-fr]{1.3.13}), il suffit d'observer que le nilradical de
$A$ est nilpotent (\cite{I-11}, p.~127, cor.~4).

\begin{corollary}[6.1.7]
\label{I.6.1.7-fr}
Soit $X$ un préschéma noethérien~; pour que $X$ soit un schéma affine, il faut
et il suffit que $X_{\mathrm{red}}$ le soit.
\end{corollary}

Cela résulte de (\hyperref[I.6.1.6-fr]{6.1.6}) et
(\hyperref[I.5.1.10-fr]{5.1.10}).

\begin{lemma}[6.1.8]
\label{I.6.1.8-fr}
Soient $X$ un espace topologique, $x$ un point de $X$, $U$ un voisinage ouvert
de $x$ n'ayant qu'un nombre fini de composantes irréductibles. Alors il existe
un voisinage $V$ de $x$ tel que tout voisinage ouvert de $x$ contenu dans $V$
soit connexe.
\end{lemma}

Soient $U_i$ ($1\leq i\leq m$) les composantes irréductibles de $U$ ne
contenant pas $x$~; le complémentaire dans $U$ de la réunion des $U_i$ est un
voisinage ouvert $V$ de $x$ dans $U$, donc aussi dans $X$~; c'est d'ailleurs
le complémentaire dans $X$ de la réunion des composantes irréductibles de $X$
qui ne contiennent pas $x$ (\hyperref[0.2.1.6-fr]{0, 2.1.6}). Soit alors $W$
un voisinage ouvert de $x$ contenu dans $V$. Les composantes irréductibles de
$W$ sont les traces sur $W$ des composantes irréductibles de $U$ qui
rencontrent $W$ (\hyperref[0.2.1.6-fr]{0, 2.1.6}), donc ces composantes
contiennent $x$~; comme elles sont connexes, il en est de même de $W$.

\begin{corollary}[6.1.9]
\label{I.6.1.9-fr}
Un espace topologique localement noethérien est localement connexe (ce qui
entraîne entre autres que ses composantes connexes sont ouvertes).
\end{corollary}

\begin{proposition}[6.1.10]
\label{I.6.1.10-fr}
Soit $X$ un espace topologique localement noethérien. Les conditions suivantes
sont équivalentes~:
\begin{enumerate}
  \item[a)] Les composantes irréductibles de $X$ sont ouvertes.
  \item[b)] Les composantes irréductibles de $X$ sont identiques à ses
    composantes connexes.
  \item[c)] Les composantes connexes de $X$ sont irréductibles.
  \item[d)] Deux composantes irréductibles distinctes de $X$ ne se rencontrent
    pas.
\end{enumerate}

Enfin, si $X$ est un préschéma, ces conditions équivalent aussi à~:
\begin{enumerate}
  \item[e)] Pour tout $x\in X$, $\operatorname{Spec}(\mathscr{O}_x)$ est
    irréductible (autrement dit le nilradical de $\mathscr{O}_x$ est premier).
\end{enumerate}
\end{proposition}

Il est immédiat que a) entraîne b), car un espace irréductible est connexe, et
a) entraîne que les composantes irréductibles de $X$ sont des ensembles à la
fois ouverts et fermés. Il est trivial que b) entraîne c)~; inversement, un
ensemble fermé $F$ contenant
\oldpage[I]{143}
une composante connexe $C$ de $X$ et distinct de $C$ ne peut être
irréductible, car cet ensemble n'étant pas connexe est réunion de deux
ensembles non vides disjoints à la fois ouverts et fermés dans $F$, donc
fermés dans $X$~; par suite c) entraîne b). On en conclut aussitôt que c)
entraîne d), deux composantes connexes distinctes étant sans point commun.

Nous n'avons pas utilisé jusqu'ici le fait que $X$ est localement noethérien.
Supposons maintenant cette hypothèse réalisée et montrons que d) entraîne
a)~: en vertu de (\hyperref[0.2.1.6-fr]{0, 2.1.6}), on peut se borner au cas
où l'espace $X$ est noethérien, donc n'a qu'un nombre fini de composantes
irréductibles. Comme celles-ci sont fermées et deux à deux disjointes, elles
sont ouvertes.

Enfin l'équivalence de d) et e) est valable sans supposer que l'espace
sous-jacent au préschéma $X$ soit localement noethérien. On peut en effet se
ramener au cas où $X=\operatorname{Spec}(A)$ est affine en vertu de
(\hyperref[0.2.1.6-fr]{0, 2.1.6})~; dire que $x$ n'est contenu que dans une
seule composante irréductible de $X$ signifie alors que $\mathfrak{j}_x$ ne
contient qu'un seul idéal minimal de $A$ (\hyperref[I.1.1.14-fr]{1.1.14}), ce
qui équivaut à dire que $\mathfrak{j}_x\mathscr{O}_x$ ne contient qu'un seul
idéal minimal de $\mathscr{O}_x$, d'où la conclusion.

\begin{corollary}[6.1.11]
\label{I.6.1.11-fr}
Soit $X$ un espace localement noethérien. Pour que $X$ soit irréductible, il
faut et il suffit que $X$ soit connexe et non vide, et que deux composantes
irréductibles distinctes de $X$ ne se rencontrent pas. Si $X$ est un
préschéma, cette dernière condition équivaut à ce que
$\operatorname{Spec}(\mathscr{O}_x)$ est irréductible pour tout $x\in X$.
\end{corollary}

La dernière partie a été vue dans (\hyperref[I.6.1.10-fr]{6.1.10})~; il n'y a
donc à prouver que la suffisance des conditions de la première assertion.
Mais d'après (\hyperref[I.6.1.10-fr]{6.1.10}), ces conditions entraînent que
les composantes irréductibles de $X$ sont ses composantes connexes, et comme
$X$ est connexe et non vide, il est irréductible.

\begin{corollary}[6.1.12]
\label{I.6.1.12-fr}
Soit $X$ un préschéma localement noethérien. Pour que $X$ soit intègre, il
faut et il suffit que $X$ soit connexe et que $\mathscr{O}_x$ soit intègre
pour tout $x\in X$.
\end{corollary}

\begin{proposition}[6.1.13]
\label{I.6.1.13-fr}
Soit $X$ un préschéma localement noethérien, et soit $x\in X$ un point tel que
le nilradical $\mathscr{N}_x$ de $\mathscr{O}_x$ soit premier (resp. que
$\mathscr{O}_x$ soit réduit, resp. intègre)~; alors il existe un voisinage
ouvert $U$ de $x$ qui est irréductible (resp. réduit, resp. intègre).
\end{proposition}

Il suffit de considérer les deux cas où $\mathscr{N}_x$ est premier et où
$\mathscr{N}_x=0$, la troisième hypothèse étant conjonction des deux
premières. Si $\mathscr{N}_x$ est premier, $x$ n'appartient qu'à une seule
composante irréductible $Y$ de $X$ (\hyperref[I.6.1.10-fr]{6.1.10})~; la
réunion des composantes irréductibles de $X$ ne contenant pas $x$ est fermée
(l'ensemble de ces composantes étant localement fini), et le complémentaire
$U$ de cette réunion est donc ouvert et contenu dans $Y$, donc irréductible
(\hyperref[0.2.1.6-fr]{0, 2.1.6}). Si $\mathscr{N}_x=0$, on a aussi
$\mathscr{N}_y=0$ pour tout $y$ dans un voisinage de $x$, car $\mathscr{N}$
est quasi-cohérent (\hyperref[I.5.1.1-fr]{5.1.1}), et par suite cohérent
puisque $X$ est localement noethérien, et la conclusion résulte de
(\hyperref[0.5.2.2-fr]{0, 5.2.2}).

\subsection{Préschémas artiniens.}
\label{subsection:I.6.2-fr}

\begin{definition}[6.2.1]
\label{I.6.2.1-fr}
On dit qu'un préschéma est \emph{artinien} s'il est affine et si son anneau
est artinien.
\end{definition}

\oldpage[I]{144}
\begin{proposition}[6.2.2]
\label{I.6.2.2-fr}
Étant donné un préschéma $X$, les conditions suivantes sont équivalentes~:
\begin{enumerate}
  \item[a)] $X$ est un schéma artinien~;
  \item[b)] $X$ est noethérien et son espace sous-jacent est discret~;
  \item[c)] $X$ est noethérien et les points de son espace sous-jacent sont
    fermés (condition $\mathrm{T}_1$).
\end{enumerate}

Lorsqu'il en est ainsi, l'espace sous-jacent à $X$ est fini, et l'anneau $A$
de $X$ est composé direct des anneaux locaux (artiniens) des points de $X$.
\end{proposition}

On sait que a) entraîne la dernière assertion (\cite{I-13}, p.~205, th.~3),
tout idéal premier de $A$ est alors maximal et est l'image réciproque de
l'idéal maximal de l'un des composants locaux de $A$, donc l'espace $X$ est
fini et discret~; a) entraîne donc b) et b) entraîne évidemment c). Pour voir
que c) entraîne a), montrons d'abord que $X$ est alors fini~; on peut en effet
se ramener au cas où $X$ est affine, et on sait qu'un anneau noethérien dont
tous les idéaux premiers sont maximaux est artinien (\cite{I-13}, p.~203),
d'où notre assertion. L'espace sous-jacent $X$ est alors discret, somme
topologique d'un nombre fini de points $x_i$, et les anneaux locaux
$\mathscr{O}_{x_i}=A_i$ sont artiniens~; il est clair que $X$ est isomorphe au
schéma affine spectre premier de l'anneau $A$ composé direct des $A_i$
(\hyperref[I.1.7.3-fr]{1.7.3}).

\subsection{Morphismes de type fini.}
\label{subsection:I.6.3-fr}

\begin{definition}[6.3.1]
\label{I.6.3.1-fr}
On dit qu'un morphisme $f:X\to Y$ est \emph{de type fini} si $Y$ est réunion
d'une famille $(V_\alpha)$ d'ouverts affines ayant la propriété suivante~:
\begin{enumerate}
  \item[(P)] $f^{-1}(V_\alpha)$ est réunion finie d'ouverts affines
    $U_{\alpha i}$ tels que chacun des anneaux $A(U_{\alpha i})$ soit une
    algèbre de type fini sur $A(V_\alpha)$.
\end{enumerate}
On dit alors aussi que $X$ est un préschéma de type fini sur $Y$, ou un
$Y$-préschéma de type fini.
\end{definition}

\begin{proposition}[6.3.2]
\label{I.6.3.2-fr}
Si $f:X\to Y$ est un morphisme de type fini, tout ouvert affine $W$ de $Y$
possède la propriété (P) de la déf. (\hyperref[I.6.3.1-fr]{6.3.1}).
\end{proposition}

Nous démontrerons d'abord le

\begin{lemma}[6.3.2.1]
\label{I.6.3.2.1-fr}
Si $T\subset Y$ est un ouvert affine, possédant la propriété (P), alors, pour
tout $g\in A(T)$, $D(g)$ possède la propriété (P).
\end{lemma}

En effet, par hypothèse, $f^{-1}(T)$ est réunion finie d'ouverts affines
$Z_j$ tels que $A(Z_j)$ soit une algèbre de type fini sur $A(T)$~; soit
$\varphi_j:A(T)\to A(Z_j)$ l'homomorphisme d'anneaux correspondant à la
restriction de $f$ à $Z_j$ (\hyperref[I.2.2.4-fr]{2.2.4}), et posons
$g_j=\varphi_j(g)$~; on a alors
$f^{-1}(D(g))\cap Z_j=D(g_j)$
(\hyperref[I.1.2.2.2-fr]{1.2.2.2}). Or,
$A(D(g_j))=A(Z_j)_{g_j}=A(Z_j)[1/g_j]$ est de type fini sur $A(Z_j)$ et
\emph{a fortiori} sur $A(T)$ en vertu de l'hypothèse, donc aussi sur
$A(D(g))=A(T)[1/g]$, ce qui démontre le lemme.

Ce lemme étant établi, comme $W$ est quasi-compact
(\hyperref[I.1.1.10-fr]{1.1.10}), il existe un recouvrement fini de $W$ par
des ensembles de la forme $D(g_i)$, où chaque $g_i$ appartient à un anneau
$A(V_{\alpha(i)})$. Chaque $D(g_i)$, étant quasi-compact, est réunion finie
d'ensembles $D(h_{ik})$ où $h_{ik}\in A(W)$~; si
$\varphi_i:A(W)\to A(D(g_i))$ est l'application canonique, on a
$D(h_{ik})=D(\varphi_i(h_{ik}))$ en vertu de
(\hyperref[I.1.2.2.2-fr]{1.2.2.2}). En vertu de
(\hyperref[I.6.3.2.1-fr]{6.3.2.1}), chacun des
$f^{-1}(D(h_{ik}))$ admet un recouvrement fini par des ouverts affines
$U_{ijk}$ tels que $A(U_{ijk})$ soit une algèbre de type fini sur
$A(D(h_{ik}))=A(W)[1/h_{ik}]$, d'où la proposition.

\oldpage[I]{145}
On peut donc dire que la notion de préschéma de type fini sur $Y$ est
\emph{locale sur $Y$}.

\begin{proposition}[6.3.3]
\label{I.6.3.3-fr}
Soient $X$, $Y$ deux schémas affines~; pour que $X$ soit de type fini sur
$Y$, il faut et il suffit que $A(X)$ soit une algèbre de type fini sur
$A(Y)$.
\end{proposition}

La condition étant évidemment suffisante, prouvons qu'elle est nécessaire.
Posons $A=A(Y)$, $B=A(X)$~; en vertu de
(\hyperref[I.6.3.2-fr]{6.3.2}), il existe un recouvrement ouvert affine fini
$(V_i)$ de $X$ tel que chacun des anneaux $A(V_i)$ soit une $A$-algèbre de
type fini. En outre, les $V_i$ étant quasi-compacts, on peut recouvrir chacun
d'eux par un nombre fini d'ouverts de la forme $D(g_{ij})\subset V_i$, où
$g_{ij}\in B$~; si $\varphi_i$ est l'homomorphisme $B\to A(V_i)$ correspondant
à l'injection canonique $V_i\to X$, on a
\[
  B_{g_{ij}}=(A(V_i))_{\varphi_i(g_{ij})}
  =A(V_i)[1/\varphi_i(g_{ij})],
\]
donc $B_{g_{ij}}$ est une $A$-algèbre de type fini. On peut donc se ramener
au cas où $V_i=D(g_i)$ avec $g_i\in B$. Par hypothèse, il existe une partie
finie $F_i$ de $B$ et un entier $n_i\geq 0$ tels que $B_{g_i}$ soit l'algèbre
engendrée sur $A$ par les éléments $b_i/g_i^{n_i}$, où $b_i$ parcourt $F_i$.
Comme les $g_i$ sont en nombre fini, on peut d'ailleurs supposer tous les
$n_i$ égaux à un même entier $n$. En outre, comme les $D(g_i)$ forment un
recouvrement de $X$, l'idéal engendré dans $B$ par les $g_i$ est égal à $B$,
autrement dit, il existe des $h_i\in B$ tels que
$\sum_i h_i g_i=1$. Soit alors $F$ la partie finie de $B$, réunion des $F_i$,
de l'ensemble des $g_i$ et de l'ensemble des $h_i$~; montrons que le
sous-anneau $B'=A[F]$ de $B$ est égal à $B$. Par hypothèse, pour tout $b\in B$
et tout $i$, l'image canonique de $b$ dans $B_{g_i}$ est de la forme
$b'_i/g_i^{m_i}$, où $b'_i\in B'$~; en multipliant les $b'_i$ par des
puissances convenables des $g_i$, on peut encore supposer tous les $m_i$
égaux à un même entier $m$. Par définition des anneaux de fractions, il y a
donc un entier $N$ (dépendant de $b$) tel que $N\geq m$ et $g_i^N b\in B'$ pour
tout $i$~; or, \emph{dans l'anneau $B'$}, les $g_i^N$ engendrent l'idéal $B'$,
puisqu'il en est ainsi des $g_i$ (les $h_i$ appartenant à $B'$)~; il y a donc
des $c_i\in B'$ tels que $\sum_i c_i g_i^N=1$, d'où
$b=\sum_i c_i g_i^N b\in B'$, C.Q.F.D.

\begin{proposition}[6.3.4]
\label{I.6.3.4-fr}
\begin{enumerate}
  \item[(i)] Toute immersion fermée est de type fini.
  \item[(ii)] Le composé de deux morphismes de type fini est de type fini.
  \item[(iii)] Si $f:X\to X'$, $g:Y\to Y'$ sont deux $S$-morphismes de type
    fini, $f\times_S g$ est de type fini.
  \item[(iv)] Si $f:X\to Y$ est un $S$-morphisme de type fini, $f_{(S')}$ est
    de type fini pour toute extension $g:S'\to S$ du préschéma de base.
  \item[(v)] Si le composé $g\circ f$ de deux morphismes est de type fini, et
    si $g$ est séparé, $f$ est de type fini.
  \item[(vi)] Si un morphisme $f$ est de type fini, il en est de même de
    $f_{\mathrm{red}}$.
\end{enumerate}
\end{proposition}

En vertu de (\hyperref[I.5.5.12-fr]{5.5.12}), il suffit de démontrer (i),
(ii) et (iv).

Pour établir (i), on peut se borner au cas d'une injection canonique
$X\to Y$, $X$ étant un sous-préschéma fermé de $Y$~; en outre
(\hyperref[I.6.3.2-fr]{6.3.2}), on peut supposer $Y$ affine, auquel cas $X$
est aussi affine (\hyperref[I.4.2.3-fr]{4.2.3}) et son anneau est isomorphe à
un anneau quotient $A/\mathfrak{J}$, où $A$ est l'anneau de $Y$ et
$\mathfrak{J}$ un idéal de $A$~; comme $A/\mathfrak{J}$ est de type fini sur
$A$, la conclusion en résulte.

Démontrons maintenant (ii). Soient $f:X\to Y$, $g:Y\to Z$ deux morphismes de
type fini, et soit $U$ un ouvert affine de $Z$~; $g^{-1}(U)$ admet un
recouvrement fini par des ouverts affines $V_i$ tels que $A(V_i)$ soit une
algèbre de type fini sur $A(U)$ (\hyperref[I.6.3.2-fr]{6.3.2})~; de même,

\oldpage[I]{146}
chacun des $f^{-1}(V_i)$ admet un recouvrement fini par des ouverts affines
$W_{ij}$ tels que $A(W_{ij})$ soit une algèbre de type fini sur $A(V_i)$, et
par suite aussi une algèbre de type fini sur $A(U)$~; d'où la conclusion.

Enfin, pour démontrer (iv), on peut se borner au cas où $S=Y$~; en effet,
$f_{(S')}$ est aussi égal à $f_{(Y_{(S')})}$, $f$ étant considéré comme un
$Y$-morphisme, et l'extension de la base étant $Y_{(S')}\to Y$
(\hyperref[I.3.3.9-fr]{3.3.9}). Soient alors $p$, $q$ les projections
$X_{(S')}\to X$ et $X_{(S')}\to S'$. Soit $V$ un ouvert affine dans $S$~;
$f^{-1}(V)$ est réunion finie d'ouverts affines $W_i$ dont chacun est tel que
$A(W_i)$ soit une algèbre de type fini sur $A(V)$
(\hyperref[I.6.3.2-fr]{6.3.2}). Soit $V'$ un ouvert affine de $S'$ contenu
dans $g^{-1}(V)$~; comme $f\circ p=g\circ q$, $q^{-1}(V')$ est contenu dans la
réunion des $p^{-1}(W_i)$~; d'autre part, l'intersection
$p^{-1}(W_i)\cap q^{-1}(V')$ s'identifie au produit $W_i\times_V V'$
(\hyperref[I.3.2.7-fr]{3.2.7}), qui est un schéma affine d'anneau isomorphe à
$A(W_i)\otimes_{A(V)}A(V')$ (\hyperref[I.3.2.2-fr]{3.2.2})~; ce dernier étant
par hypothèse une algèbre de type fini sur $A(V')$, la proposition est
démontrée.

\begin{corollary}[6.3.5]
\label{I.6.3.5-fr}
Soit $f:X\to Y$ un morphisme d'immersion. Si l'espace sous-jacent à $Y$
(resp. $X$) est localement noethérien (resp. noethérien), $f$ est de type fini.
\end{corollary}

On peut toujours supposer $Y$ affine (\hyperref[I.6.3.2-fr]{6.3.2})~; si
l'espace sous-jacent à $Y$ est localement noethérien, on peut en outre le
supposer noethérien et alors l'espace sous-jacent à $X$, qui en est un
sous-espace, est noethérien. Autrement dit, on peut supposer $Y$ affine et
l'espace sous-jacent à $X$ noethérien~; il existe alors un recouvrement de $X$
par un nombre fini d'ouverts affines $D(g_i)\subset Y$, où $g_i\in A(Y)$, tels
que $X\cap D(g_i)$ soit fermé dans $D(g_i)$ (donc un schéma affine
(\hyperref[I.4.2.3-fr]{4.2.3})), puisque $X$ est localement fermé dans $Y$
(\hyperref[I.4.1.3-fr]{4.1.3}). Alors $A(X\cap D(g_i))$ est une algèbre de type
fini sur $A(D(g_i))$, d'après (\hyperref[I.6.3.4-fr]{6.3.4, (i)}) et
(\hyperref[I.6.3.3-fr]{6.3.3}), et
$A(D(g_i))=A(Y)_{g_i}=A(Y)[1/g_i]$ est de type fini sur $A(Y)$, ce qui achève
la démonstration.

\begin{corollary}[6.3.6]
\label{I.6.3.6-fr}
Soient $f:X\to Y$, $g:Y\to Z$ deux morphismes. Si $g\circ f$ est de type fini,
et si $X$ est noethérien, ou $X\times_Z Y$ localement noethérien, $f$ est de
type fini.
\end{corollary}

Cela résulte aussitôt de la démonstration de
(\hyperref[I.5.5.12-fr]{5.5.12}) et de (\hyperref[I.6.3.5-fr]{6.3.5}) appliquée
au morphisme d'immersion $\Gamma_f$.

\begin{proposition}[6.3.7]
\label{I.6.3.7-fr}
Soit $f:X\to Y$ un morphisme de type fini~; si $Y$ est noethérien (resp.
localement noethérien), $X$ est noethérien (resp. localement noethérien).
\end{proposition}

On peut se borner à faire la démonstration lorsque $Y$ est noethérien. Alors $Y$
est réunion finie d'ouverts affines $V_i$ tels que les $A(V_i)$ soient des
anneaux noethériens. En vertu de (\hyperref[I.6.3.2-fr]{6.3.2}), chacun des
$f^{-1}(V_i)$ est réunion d'un nombre fini d'ouverts affines $W_{ij}$ tels que
les $A(W_{ij})$ soient des algèbres de type fini sur $A(V_i)$, donc des anneaux
noethériens~; cela prouve que $X$ est noethérien.

\begin{corollary}[6.3.8]
\label{I.6.3.8-fr}
Soit $X$ un préschéma de type fini sur $S$. Pour toute extension de la base
$S'\to S$ telle que $S'$ soit noethérien (resp. localement noethérien),
$X_{(S')}$ est noethérien (resp. localement noethérien).
\end{corollary}

Cela résulte de (\hyperref[I.6.3.7-fr]{6.3.7}), $X_{(S')}$ étant de type fini
sur $S'$ en vertu de (\hyperref[I.6.3.4-fr]{6.3.4, (iv)}).

On peut encore dire que dans un produit $X\times_S Y$ de $S$-préschémas, si
\emph{l'un} des
\oldpage[I]{147}
facteurs $X$, $Y$ est \emph{de type fini} sur $S$ et
\emph{l'autre noethérien} (resp. \emph{localement noethérien}), alors
$X\times_S Y$ est \emph{noethérien} (resp. \emph{localement noethérien}).

\begin{corollary}[6.3.9]
\label{I.6.3.9-fr}
Soit $X$ un préschéma de type fini sur un préschéma localement noethérien $S$.
Alors tout $S$-morphisme $f:X\to Y$ est de type fini.
\end{corollary}

En effet, on peut supposer $S$ noethérien~; si $\varphi:X\to S$,
$\psi:Y\to S$ sont les morphismes structuraux, on a $\varphi=\psi\circ f$, et
$X$ est noethérien en vertu de (\hyperref[I.6.3.7-fr]{6.3.7})~; $f$ est donc
de type fini en vertu de (\hyperref[I.6.3.6-fr]{6.3.6}).

\begin{proposition}[6.3.10]
\label{I.6.3.10-fr}
Soit $f:X\to Y$ un morphisme de type fini. Pour que $f$ soit surjectif, il faut
et il suffit que, pour tout corps algébriquement clos $\Omega$, l'application
$X(\Omega)\to Y(\Omega)$ correspondant à $f$
(\hyperref[I.3.4.1-fr]{3.4.1}) soit surjective.
\end{proposition}

La condition est suffisante, comme on le voit en considérant pour tout
$y\in Y$, une extension algébriquement close $\Omega$ de $k(y)$, et le
diagramme commutatif
\[
  \xymatrix{
    X\ar[dd]_f\\
    & \operatorname{Spec}(\Omega)\ar[ul]\ar[dl]\\
    Y
  }
\]
(\emph{cf.} (\hyperref[I.3.5.3-fr]{3.5.3})). Inversement, supposons $f$
surjective, et soit $g:\{\xi\}=\operatorname{Spec}(\Omega)\to Y$ un morphisme,
$\Omega$ étant un corps algébriquement clos. Si on considère le diagramme
\[
  \xymatrix{
    X\ar[d]_f &
    X_{(\Omega)}\ar[l]\ar[d]^{f_{(\Omega)}}\\
    Y &
    \operatorname{Spec}(\Omega)\ar[l]
  }
\]
il suffit donc de montrer qu'il existe dans $X_{(\Omega)}$ un
\emph{point rationnel sur $\Omega$}
(\hyperref[I.3.3.14-fr]{3.3.14}, \hyperref[I.3.4.3-fr]{3.4.3} et
\hyperref[I.3.4.4-fr]{3.4.4}). Comme $f$ est surjectif, $X_{(\Omega)}$ n'est
pas vide (\hyperref[I.3.5.10-fr]{3.5.10}) et comme $f$ est de type fini, il en
est de même de $f_{(\Omega)}$ (\hyperref[I.6.3.4-fr]{6.3.4, (iv)})~; donc
$X_{(\Omega)}$ contient un ouvert affine non vide $Z$ tel que $A(Z)$ soit une
algèbre de type fini non nulle sur $\Omega$. En vertu du th. des zéros de
Hilbert \cite{I-21}, il existe un $\Omega$-homomorphisme $A(Z)\to\Omega$, donc
une section de $X_{(\Omega)}$ au-dessus de $\operatorname{Spec}(\Omega)$, ce
qui démontre la proposition.

\subsection{Préschémas algébriques.}
\label{subsection:I.6.4-fr}

\begin{definition}[6.4.1]
\label{I.6.4.1-fr}
Étant donné un corps $K$, on appelle \emph{$K$-préschéma algébrique} un
préschéma $X$ de type fini sur $K$~; $K$ est appelé le \emph{corps de base}
de $X$. Si en outre $X$ est un schéma (ou, ce qui revient au même
(\hyperref[I.5.5.8-fr]{5.5.8}), si $X$ est un $K$-schéma), on dit aussi que
$X$ est un \emph{$K$-schéma algébrique}.
\end{definition}

Tout $K$-préschéma algébrique est \emph{noethérien}
(\hyperref[I.6.3.7-fr]{6.3.7}).

\begin{proposition}[6.4.2]
\label{I.6.4.2-fr}
Soit $X$ un $K$-préschéma algébrique. Pour qu'un point $x\in X$ soit fermé,
il faut et il suffit que $k(x)$ soit une extension algébrique de $K$, de degré
fini.
\end{proposition}

On peut supposer $X$ affine, l'anneau $A$ de $X$ étant une $K$-algèbre de type
fini. En effet, les ouverts affines $U$ de $X$ tels que $A(U)$ soit une
$K$-algèbre de type fini forment un recouvrement de $X$
(\hyperref[I.6.3.1-fr]{6.3.1}). Les points fermés de $X$ sont alors les points
tels que $\mathfrak{j}_x$ soit
\oldpage[I]{148}
un idéal maximal de $A$, autrement dit tels que $A/\mathfrak{j}_x$ soit un
corps (nécessairement égal à $k(x)$). Comme $A/\mathfrak{j}_x$ est une
$K$-algèbre de type fini, on voit que si $x$ est fermé, $k(x)$ est un corps
qui est une algèbre de type fini sur $K$, donc nécessairement une $K$-algèbre
de \emph{rang fini} \cite{I-21}. Inversement, si $k(x)$ est de rang fini sur
$K$, il en est de même de $A/\mathfrak{j}_x\subset k(x)$ et comme tout anneau
intègre qui est une $K$-algèbre de rang fini est un corps, on a
$A/\mathfrak{j}_x=k(x)$, donc $x$ est fermé.

\begin{corollary}[6.4.3]
\label{I.6.4.3-fr}
Soient $K$ un corps algébriquement clos, $X$ un $K$-préschéma algébrique~; les
points fermés de $X$ sont alors les points rationnels sur $K$
(\hyperref[I.3.4.4-fr]{3.4.4}) et s'identifient canoniquement aux points de
$X$ à valeurs dans $K$.
\end{corollary}

\begin{proposition}[6.4.4]
\label{I.6.4.4-fr}
Soit $X$ un préschéma algébrique sur un corps $K$. Les propriétés suivantes
sont équivalentes~:
\begin{enumerate}
  \item[\rm(a)] $X$ est artinien.
  \item[\rm(b)] L'espace sous-jacent à $X$ est discret.
  \item[\rm(c)] L'espace sous-jacent à $X$ n'a qu'un nombre fini de points
    fermés.
  \item[\rm(c')] L'espace sous-jacent à $X$ est fini.
  \item[\rm(d)] Les points de $X$ sont fermés.
  \item[\rm(e)] $X$ est isomorphe à $\operatorname{Spec}(A)$, où $A$ est une
    $K$-algèbre de rang fini.
\end{enumerate}
\end{proposition}

Comme $X$ est noethérien, il résulte de (\hyperref[I.6.2.2-fr]{6.2.2}) que
les conditions a), b), d) sont équivalentes et entraînent c) et c')~; par
ailleurs, il est clair que e) entraîne a). Reste à voir que c) entraîne d) et
e)~; on peut se borner au cas où $X$ est affine. Alors $A(X)$ est une
$K$-algèbre de type fini (\hyperref[I.6.3.3-fr]{6.3.3}), donc un anneau de
Jacobson (\cite[p.~3-11 et 3-12]{I-1}), dans lequel il n'y a par hypothèse
qu'un nombre fini d'idéaux maximaux. Comme une intersection finie d'idéaux
premiers ne peut être un idéal premier que si elle est égale à l'un d'eux,
tout idéal premier de $A(X)$ est donc maximal, d'où d). En outre, on sait
alors (\hyperref[I.6.2.2-fr]{6.2.2}) que $A(X)$ est une $K$-algèbre
artinienne de type fini, donc nécessairement de \emph{rang fini} \cite{I-21}.

\begin{env}[6.4.5]
\label{I.6.4.5-fr}
Lorsque les conditions de (\hyperref[I.6.4.4-fr]{6.4.4}) sont satisfaites,
on dit que $X$ est un schéma \emph{fini sur $K$}
(\emph{cf.} (\hyperref[II.6.1.1-fr]{II, 6.1.1})), ou un
\emph{$K$-schéma fini}, de \emph{rang} $[A:K]$ que l'on note aussi
$\operatorname{rg}_K(X)$~; si $X$, $Y$ sont deux schémas finis sur $K$, on a
\[
  \operatorname{rg}_K(X\amalg Y)
  =\operatorname{rg}_K(X)+\operatorname{rg}_K(Y)
  \tag{6.4.5.1}\label{I.6.4.5.1-fr}
\]
\[
  \operatorname{rg}_K(X\times_K Y)
  =\operatorname{rg}_K(X)\operatorname{rg}_K(Y)
  \tag{6.4.5.2}\label{I.6.4.5.2-fr}
\]
comme il résulte de (\hyperref[I.3.2.2-fr]{3.2.2}).
\end{env}

\begin{corollary}[6.4.6]
\label{I.6.4.6-fr}
Soit $X$ un schéma fini sur un corps $K$. Pour toute extension $K'$ de $K$,
$X\otimes_K K'$ est un schéma fini sur $K'$, et son rang sur $K'$ est égal au
rang de $X$ sur $K$.
\end{corollary}

En effet, si $A=A(X)$, on a $[A\otimes_K K':K']=[A:K]$.

\begin{corollary}[6.4.7]
\label{I.6.4.7-fr}
Soit $X$ un schéma fini sur un corps $K$~; on pose
$n=\sum_{x\in X}[k(x):K]_s$
\emph{(on rappelle que si $K'$ est une extension de $K$, $[K':K]_s$ est le
\emph{rang séparable} de $K'$ sur $K$, rang de la plus grande extension
algébrique séparable de $K$ contenue dans $K'$)}~;
\oldpage[I]{149}
alors, pour toute extension algébriquement close $\Omega$ de $K$, l'espace
sous-jacent à $X\otimes_K\Omega$ a exactement $n$ points, qui s'identifient
aux points de $X$ à valeurs dans $\Omega$.
\end{corollary}

On peut évidemment se borner au cas où l'anneau $A=A(X)$ est local
(\hyperref[I.6.2.2-fr]{6.2.2})~; soient $\mathfrak m$ son idéal maximal,
$L=A/\mathfrak m$ son corps résiduel, extension algébrique de $K$. Les points
de $X$ à valeurs dans $\Omega$ correspondent alors biunivoquement aux
$\Omega$-sections de $X\otimes_K\Omega$
(\hyperref[I.3.4.1-fr]{3.4.1} et \hyperref[I.3.3.14-fr]{3.3.14}), et aussi
aux $K$-homomorphismes de $L$ dans $\Omega$
(\hyperref[I.1.7.3-fr]{1.7.3}), d'où la proposition (Bourbaki, \emph{Alg.},
chap.~V, \textsection7, n\textsuperscript{o}~5, prop.~8), compte tenu de
(\hyperref[I.6.4.3-fr]{6.4.3}).

\begin{env}[6.4.8]
\label{I.6.4.8-fr}
Le nombre $n$ défini dans (\hyperref[I.6.4.7-fr]{6.4.7}) s'appelle le
\emph{rang séparable} de $A$ (ou de $X$) sur $K$, ou aussi le
\emph{nombre géométrique de points de $X$}~; il est donc égal au nombre
d'éléments de $X(\Omega)_K$. Il résulte aussitôt de cette définition que,
pour toute extension $K'$ de $K$, $X\otimes_K K'$ a même nombre géométrique
de points que $X$. Si on désigne ce nombre par $n(X)$, il est clair que si
$X$, $Y$ sont deux schémas finis sur $K$, on a
\[
  n(X\amalg Y)=n(X)+n(Y).
  \tag{6.4.8.1}\label{I.6.4.8.1-fr}
\]

Sous les mêmes hypothèses, on a aussi
\[
  n(X\times_K Y)=n(X)n(Y)
  \tag{6.4.8.2}\label{I.6.4.8.2-fr}
\]
comme il résulte aussitôt de l'interprétation de $n(X)$ comme nombre
d'éléments de $X(\Omega)_K$ et de la formule
(\hyperref[I.3.4.3.1-fr]{3.4.3.1}).
\end{env}

\begin{proposition}[6.4.9]
\label{I.6.4.9-fr}
Soient $K$ un corps, $X$, $Y$ deux $K$-préschémas algébriques, $f:X\to Y$ un
$K$-morphisme, $\Omega$ une extension algébriquement close de $K$, de degré
de transcendance infini sur $K$. Pour que $f$ soit surjectif, il faut et il
suffit que l'application $X(\Omega)_K\to Y(\Omega)_K$ correspondant à $f$
(\hyperref[I.3.4.1-fr]{3.4.1}) soit surjective.
\end{proposition}

La nécessité résulte de (\hyperref[I.6.3.10-fr]{6.3.10}), en remarquant que
$f$ est nécessairement de type fini (\hyperref[I.6.3.9-fr]{6.3.9}). Pour voir
que la condition est suffisante, on raisonne comme dans
(\hyperref[I.6.3.10-fr]{6.3.10}), en remarquant que pour tout $y\in Y$,
$k(y)$ est une extension de $K$ de type fini, et par suite est $K$-isomorphe
à un sous-corps de $\Omega$.

\begin{remark}[6.4.10]
\label{I.6.4.10-fr}
Nous verrons au chap.~IV que la conclusion de
(\hyperref[I.6.4.9-fr]{6.4.9}) est encore valable sans hypothèse relative au
degré de transcendance de $\Omega$ sur $K$.
\end{remark}

\begin{proposition}[6.4.11]
\label{I.6.4.11-fr}
Si $f:X\to Y$ est un morphisme de type fini, pour tout $y\in Y$, la fibre
$f^{-1}(y)$ est un préschéma algébrique sur le corps résiduel $k(y)$ et pour
tout $x\in f^{-1}(y)$, $k(x)$ est une extension de type fini de $k(y)$.
\end{proposition}

Comme $f^{-1}(y)=X\otimes_Y k(y)$
(\hyperref[I.3.6.3-fr]{3.6.3}), la proposition résulte de
(\hyperref[I.6.3.4-fr]{6.3.4, (iv)}) et de
(\hyperref[I.6.3.3-fr]{6.3.3}).

\begin{proposition}[6.4.12]
\label{I.6.4.12-fr}
Soient $f:X\to Y$, $g:Y'\to Y$ deux morphismes~; posons
$X'=X\times_Y Y'$ et soit $f'=f_{(Y')}:X'\to Y'$. Soit $y'\in Y'$,
$y=g(y')$~; si la fibre $f^{-1}(y)$ est un schéma algébrique fini sur $k(y)$,
alors la fibre $f'^{-1}(y')$ est un schéma algébrique fini sur $k(y')$, ayant
même rang et même nombre géométrique de points que $f^{-1}(y)$.
\end{proposition}

Compte tenu de la transitivité des fibres
(\hyperref[I.3.6.5-fr]{3.6.5}), cela résulte aussitôt de
(\hyperref[I.6.4.6-fr]{6.4.6}) et
(\hyperref[I.6.4.8-fr]{6.4.8}).

\begin{env}[6.4.13]
\label{I.6.4.13-fr}
La prop. (\hyperref[I.6.4.11-fr]{6.4.11}) montre que les morphismes de type
fini correspondent intuitivement aux «~familles algébriques de variétés
algébriques~», les points de $Y$ jouant
\oldpage[I]{150}
le rôle de «~paramètres~», ce qui donne à ces morphismes une signification
«~géométrique~». Les morphismes qui ne sont pas de type fini interviendront
surtout par la suite dans les questions de «~changement du préschéma de
base~», par exemple par localisation ou complétion.
\end{env}

\subsection{Détermination locale d'un morphisme.}
\label{subsection:I.6.5-fr}

\begin{proposition}[6.5.1]
\label{I.6.5.1-fr}
Soient $X$, $Y$ deux $S$-préschémas, $Y$ étant de type fini sur $S$~; soient
$x\in X$, $y\in Y$ au-dessus d'un même point $s\in S$.
\begin{enumerate}
  \item[(i)] Si deux $S$-morphismes $f=(\psi,\theta)$,
    $f'=(\psi',\theta')$ de $X$ dans $Y$ sont tels que
    $\psi(x)=\psi'(x)=y$, et que les $\mathscr{O}_s$-homomorphismes (locaux)
    $\theta_x^\sharp$ et ${\theta'}_x^\sharp$ de $\mathscr{O}_y$ dans
    $\mathscr{O}_x$ soient identiques, $f$ et $f'$ coïncident dans un
    voisinage ouvert de $x$.
  \item[(ii)] Supposons en outre $S$ localement noethérien. Pour tout
    $\mathscr{O}_s$-homomorphisme local
    $\varphi:\mathscr{O}_y\to\mathscr{O}_x$, il existe un voisinage ouvert
    $U$ de $x$ dans $X$ et un $S$-morphisme $f=(\psi,\theta)$ de $U$ dans $Y$
    tels que $\psi(x)=y$ et $\theta_x^\sharp=\varphi$.
\end{enumerate}
\end{proposition}

\begin{enumerate}
  \item[(i)] La question étant locale sur $S$, $X$ et $Y$, on peut supposer
    $S$, $X$, $Y$ affines d'anneaux respectifs $A$, $B$, $C$, $f$ et $f'$
    étant de la forme $({}^a\varphi,\widetilde{\varphi})$ et
    $({}^a\varphi',\widetilde{\varphi}')$ respectivement, où $\varphi$ et
    $\varphi'$ sont deux $A$-homomorphismes de $C$ dans $B$ tels que
    $\varphi^{-1}(\mathfrak{j}_x)={\varphi'}^{-1}(\mathfrak{j}_x)
    =\mathfrak{j}_y$, et les homomorphismes $\varphi_x$ et $\varphi'_x$ de
    $C_y$ dans $B_x$, déduits de $\varphi$ et $\varphi'$, sont identiques~;
    on peut en outre supposer que $C$ est une $A$-algèbre de type fini. Soient
    $c_i$ ($1\leq i\leq n$) des générateurs de la $A$-algèbre $C$, et posons
    $b_i=\varphi(c_i)$, $b'_i=\varphi'(c_i)$~; par hypothèse, on a
    $b_i/1=b'_i/1$ dans l'anneau de fractions $B_x$ ($1\leq i\leq n$). Cela
    signifie qu'il existe des éléments $s_i\in B-\mathfrak{j}_x$ tels que
    $s_i(b_i-b'_i)=0$ pour $1\leq i\leq n$, et on peut évidemment supposer
    tous les $s_i$ égaux à un même élément $g\in B-\mathfrak{j}_x$. On en
    conclut que l'on a $b_i/1=b'_i/1$ pour $1\leq i\leq n$ dans l'anneau de
    fractions $B_g$~; si $i_g$ est l'homomorphisme canonique $B\to B_g$, on a
    par suite $i_g\circ\varphi=i_g\circ\varphi'$~; donc les restrictions de
    $f$ et $f'$ à $D(g)$ sont identiques.
  \item[(ii)] On peut se ramener à la même situation que dans (i), et supposer
    en outre que l'anneau $A$ est noethérien. Soient $c_i$
    ($1\leq i\leq n$) des générateurs de la $A$-algèbre $C$, et soit
    $\alpha:A[X_1,\ldots,X_n]\to C$ l'homomorphisme de l'algèbre de polynômes
    $A[X_1,\ldots,X_n]$ sur $C$ transformant $X_i$ en $c_i$ pour
    $1\leq i\leq n$. Soit d'autre part $i_y$ l'homomorphisme canonique
    $C\to C_y$, et considérons l'homomorphisme composé
    \[
      \beta:A[X_1,\ldots,X_n]
      \xrightarrow{\alpha} C
      \xrightarrow{i_y} C_y
      \xrightarrow{\varphi} B_x.
    \]
    Désignons par $\mathfrak{a}$ le noyau de $\beta$~; comme $A$ est
    noethérien, il en est de même de $A[X_1,\ldots,X_n]$, et par suite
    $\mathfrak{a}$ admet un système fini de générateurs
    $Q_j(X_1,\ldots,X_n)$ ($1\leq j\leq m$). D'autre part, chacun des
    éléments $\varphi(i_y(c_i))$ peut s'écrire $b_i/s_i$, où $b_i\in B$ et
    $s_i\notin\mathfrak{j}_x$~; on peut en outre supposer tous les $s_i$ égaux
    à un même élément $g\in B-\mathfrak{j}_x$. Cela étant, on a par hypothèse
    $Q_j(b_1/g,\ldots,b_n/g)=0$ dans $B_x$~; posons
    \[
      Q_j(X_1/T,\ldots,X_n/T)
      =P_j(X_1,\ldots,X_n,T)/T^{k_j}
    \]
    où $P_j$ est homogène de degré $k_j$. Soit alors
    $d_j=P_j(b_1,\ldots,b_n,g)\in B$. Par hypothèse, on a $t_jd_j=0$ pour un
    $t_j\in B-\mathfrak{j}_x$ ($1\leq j\leq m$), et on peut évidemment
    supposer tous les $t_j$
\oldpage[I]{151}
    égaux à un même élément $h\in B-\mathfrak{j}_x$~; on en conclut que
    $P_j(hb_1,\ldots,hb_n,hg)=0$ pour $1\leq j\leq m$. Cela étant, considérons
    l'homomorphisme $\rho$ de $A[X_1,\ldots,X_n]$ dans l'anneau de fractions
    $B_{hg}$ qui applique $X_i$ sur $hb_i/hg$ ($1\leq i\leq n$)~; l'image de
    $\mathfrak{a}$ par cet homomorphisme est 0, et il en est \emph{a fortiori}
    de même de l'image par $\rho$ du noyau $\alpha^{-1}(0)$. Donc $\rho$ se
    factorise en
    $A[X_1,\ldots,X_n]\xrightarrow{\alpha}C\xrightarrow{\gamma}B_{hg}$, avec
    $\gamma(c_i)=hb_i/hg$, et il est clair que si $i_x$ est l'homomorphisme
    canonique $B_{hg}\to B_x$, le diagramme
    \[
      \xymatrix{
        C\ar[r]^\gamma\ar[d]_{i_y} & B_{hg}\ar[d]^{i_x}\\
        C_y\ar[r]_\varphi & B_x
      }
      \tag{6.5.1.1}\label{I.6.5.1.1-fr}
    \]
    est commutatif~; on a donc $\varphi=\gamma_x$, et comme $\varphi$ est un
    homomorphisme local, ${}^a\gamma(x)=y$~; $f=({}^a\gamma,\widetilde{\gamma})$
    est donc un $S$-morphisme du voisinage $D(hg)$ de $x$ dans $Y$ qui répond à
    la question.
\end{enumerate}

\begin{corollary}[6.5.2]
\label{I.6.5.2-fr}
Sous les hypothèses de (\hyperref[I.6.5.1-fr]{6.5.1, (ii)}), si en outre $X$
est de type fini sur $S$, on peut supposer le morphisme $f$ de type fini.
\end{corollary}

Cela résulte de (\hyperref[I.6.3.6-fr]{6.3.6}).

\begin{corollary}[6.5.3]
\label{I.6.5.3-fr}
Supposons vérifiées les hypothèses de
(\hyperref[I.6.5.1-fr]{6.5.1, (ii)}), et supposons en outre que $Y$ soit
intègre, et $\varphi$ un homomorphisme injectif. Alors on peut supposer que
$f=({}^a\gamma,\widetilde{\gamma})$, où $\gamma$ est injectif.
\end{corollary}

En effet, on peut supposer $C$ intègre
(\hyperref[I.5.1.4-fr]{5.1.4}), donc $i_y$ injectif~; il résulte alors du
diagramme (\hyperref[I.6.5.1.1-fr]{6.5.1.1}) que $\gamma$ est injectif.

\begin{proposition}[6.5.4]
\label{I.6.5.4-fr}
Soient $f=(\psi,\theta):X\to Y$ un morphisme de type fini, $x$ un point de
$X$, $y=\psi(x)$.
\begin{enumerate}
  \item[(i)] Pour que $f$ soit une immersion locale au point $x$
    (\hyperref[I.4.5.1-fr]{4.5.1}), il faut et il suffit que
    $\theta_x^\sharp:\mathscr{O}_y\to\mathscr{O}_x$ soit surjectif.
  \item[(ii)] On suppose en outre $Y$ localement noethérien. Pour que $f$ soit
    un isomorphisme local au point $x$
    (\hyperref[I.4.5.2-fr]{4.5.2}), il faut et il suffit que
    $\theta_x^\sharp$ soit un isomorphisme.
\end{enumerate}
\end{proposition}

\begin{enumerate}
  \item[(ii)] En vertu de (\hyperref[I.6.5.1-fr]{6.5.1}), il existe alors un
    voisinage ouvert $V$ de $y$ et un morphisme $g:V\to X$ tels que $g\circ f$
    (resp. $f\circ g$) soit défini et coïncide avec l'identité dans un
    voisinage de $x$ (resp. $y$), d'où on tire aisément que $f$ est un
    isomorphisme local.
  \item[(i)] La question étant locale sur $X$ et $Y$, on peut supposer $X$ et
    $Y$ affines, d'anneaux respectifs $A$, $B$~; on a
    $f=({}^a\varphi,\widetilde{\varphi})$, où $\varphi$ est un homomorphisme
    d'anneaux $B\to A$ qui fait de $A$ une $B$-algèbre de type fini~; on a
    $\varphi^{-1}(\mathfrak{j}_x)=\mathfrak{j}_y$ et l'homomorphisme
    $\varphi_x:B_y\to A_x$ déduit de $\varphi$ est surjectif. Soit $(t_i)$
    ($1\leq i\leq n$) un système de générateurs de la $B$-algèbre $A$~;
    l'hypothèse sur $\varphi_x$ implique qu'il existe des $b_i\in B$ et un
    $c\in B-\mathfrak{j}_y$ tels que, dans l'anneau de fractions $A_x$, on ait
    $t_i/1=\varphi(b_i)/\varphi(c)$ pour $1\leq i\leq n$. Par suite
    (\hyperref[I.1.3.3-fr]{1.3.3}), il existe $a\in A-\mathfrak{j}_x$ tel que,
    si l'on pose $g=a\varphi(c)$, on ait aussi
    $t_i/1=a\varphi(b_i)/g$ \emph{dans l'anneau de fractions $A_g$}. Cela étant,
    il existe par hypothèse un polynôme $Q(X_1,\ldots,X_n)$ à coefficients dans
    l'anneau $\varphi(B)$,
\oldpage[I]{152}
    tel que $a=Q(t_1,\ldots,t_n)$~; posons
    \[
      Q(X_1/T,\ldots,X_n/T)
      =P(X_1,\ldots,X_n,T)/T^m,
    \]
    où $P$ est homogène de degré $m$. Dans l'anneau $A_g$, on a
    \[
      a/1
      =a^mP(\varphi(b_1),\ldots,\varphi(b_n),\varphi(c))/g^m
      =a^m\varphi(d)/g^m
    \]
    où $d\in B$. Comme, dans $A_g$,
    $g/1=(a/1)(\varphi(c)/1)$ est inversible par définition, il en est de même
    de $a/1$ et de $\varphi(c)/1$, et on peut donc écrire
    \[
      a/1=(\varphi(d)/1)(\varphi(c)/1)^{-m}.
    \]
    On en conclut que $\varphi(d)/1$ est aussi inversible dans $A_g$. Posons
    alors $h=cd$~; comme $\varphi(h)/1$ est inversible dans $A_g$,
    l'homomorphisme composé
    \[
      B\xrightarrow{\varphi}A\longrightarrow A_g
    \]
    se factorise en
    \[
      B\longrightarrow B_h\xrightarrow{\gamma}A_g
    \]
    (\hyperref[0.1.2.4-fr]{0, 1.2.4}). Montrons que $\gamma$ est
    \emph{surjectif}~; il suffit de vérifier que l'image de $B_h$ dans $A_g$
    contient les $t_i/1$ et $(g/1)^{-1}$. Or, on a
    \[
      (g/1)^{-1}
      =(\varphi(c)/1)^{m-1}(\varphi(d)/1)^{-1}
      =\gamma(c^m/h),
    \]
    et
    \[
      a/1=\gamma(d^{m+1}/h^m),
    \]
    donc
    \[
      (a\varphi(b_i))/1=\gamma(b_i d^{m+1}/h^m),
    \]
    et comme
    $t_i/1=(a\varphi(b_i)/1)(g/1)^{-1}$, notre assertion est démontrée. Le
    choix de $h$ implique que $\psi(D(g))\subset D(h)$, et la restriction de
    $f$ à $D(g)$ est égale à
    $({}^a\gamma,\widetilde{\gamma})$~; comme $\gamma$ est surjectif, cette
    restriction est une immersion fermée de $D(g)$ dans $D(h)$
    (\hyperref[I.4.2.3-fr]{4.2.3}).
\end{enumerate}

\begin{corollary}[6.5.5]
\label{I.6.5.5-fr}
Soit $f=(\psi,\theta):X\to Y$ un morphisme de type fini. On suppose $X$
irréductible, on désigne par $x$ son point générique, et on pose $y=\psi(x)$.
\begin{enumerate}
  \item[(i)] Pour que $f$ soit une immersion locale en un point de $X$, il
    faut et il suffit que
    $\theta_x^\sharp:\mathscr{O}_y\to\mathscr{O}_x$ soit surjectif.
  \item[(ii)] On suppose de plus $Y$ irréductible et localement noethérien.
    Pour que $f$ soit un isomorphisme local en un point de $X$, il faut et il
    suffit que $y$ soit le point générique de $Y$ (ou, ce qui revient au même
    (\hyperref[0.2.1.4-fr]{0, 2.1.4}) que $f$ soit un morphisme
    \emph{dominant}) et que $\theta_x^\sharp$ soit un isomorphisme (autrement
    dit, que $f$ soit \emph{birationnel}
    (\hyperref[I.2.2.9-fr]{2.2.9})).
\end{enumerate}
\end{corollary}

Il est clair que (i) résulte de
(\hyperref[I.6.5.4-fr]{6.5.4, (i)}) compte tenu de ce que tout ouvert non vide
de $X$ contient $x$~; de même (ii) résulte de
(\hyperref[I.6.5.4-fr]{6.5.4, (ii)}).

\subsection{Morphismes quasi-compacts et morphismes localement de type fini.}
\label{subsection:I.6.6-fr}

\begin{definition}[6.6.1]
\label{I.6.6.1-fr}
On dit qu'un morphisme $f:X\to Y$ est \emph{quasi-compact} si l'image
réciproque par $f$ de tout ouvert quasi-compact de $Y$ est quasi-compacte.
\end{definition}

Soit $\mathfrak{B}$ une base de la topologie de $Y$ formée d'ouverts
quasi-compacts (par exemple d'ouverts affines)~; pour que $f$ soit
quasi-compact, il faut et il suffit que l'image réciproque par $f$ de tout
ensemble de $\mathfrak{B}$ soit quasi-compacte (ou, ce qui revient au même,
réunion \emph{finie} d'ouverts affines), car tout ouvert quasi-compact de $Y$
est réunion finie d'ensembles de $\mathfrak{B}$. Par exemple, si $X$ est
\emph{quasi-compact} et $Y$ \emph{affine}, alors \emph{tout} morphisme
$f:X\to Y$ est quasi-compact~: en effet, $X$ est réunion finie d'ensembles
ouverts affines $U_i$, et pour tout ouvert affine $V$ de $Y$,
$U_i\cap f^{-1}(V)$ est affine
(\hyperref[I.5.5.10-fr]{5.5.10}), donc quasi-compact.

Si $f:X\to Y$ est un morphisme quasi-compact, il est clair que pour tout
ouvert $V$ de $Y$, la restriction de $f$ à $f^{-1}(V)$ est un morphisme
quasi-compact $f^{-1}(V)\to V$. Inversement, si $(U_\alpha)$ est un
recouvrement ouvert de $Y$ et $f:X\to Y$ un morphisme tel que les restrictions
$f^{-1}(U_\alpha)\to U_\alpha$ soient quasi-compactes, alors $f$ est
quasi-compact.

\begin{definition}[6.6.2]
\label{I.6.6.2-fr}
On dit qu'un morphisme $f:X\to Y$ est \emph{localement de type fini} si, pour
tout $x\in X$, il existe un voisinage ouvert $U$ de $x$ et un voisinage ouvert
$V\supset f(U)$ de $y$ tels que la
\oldpage[I]{153}
restriction de $f$ à $U$ soit un morphisme de type fini de $U$ dans $V$. On
dit alors aussi que $X$ est un préschéma localement de type fini sur $Y$, ou
un $Y$-préschéma localement de type fini.
\end{definition}

Il résulte aussitôt de (\hyperref[I.6.3.2-fr]{6.3.2}) que si $f$ est
localement de type fini, alors, pour tout ouvert $W$ de $Y$, la restriction
de $f$ à $f^{-1}(W)$ est un morphisme $f^{-1}(W)\to W$ qui est localement de
type fini.

Si $Y$ est localement noethérien et si $X$ est localement de type fini sur
$Y$, $X$ est localement noethérien en vertu de
(\hyperref[I.6.3.7-fr]{6.3.7}).

\begin{proposition}[6.6.3]
\label{I.6.6.3-fr}
Pour qu'un morphisme $f:X\to Y$ soit de type fini, il faut et il suffit qu'il
soit quasi-compact et localement de type fini.
\end{proposition}

La nécessité des conditions est immédiate, vu
(\hyperref[I.6.3.1-fr]{6.3.1}) et la remarque suivant
(\hyperref[I.6.6.1-fr]{6.6.1}). Inversement, supposons ces conditions
satisfaites et soit $U$ un ouvert affine de $Y$, d'anneau $A$~; pour tout
$x\in f^{-1}(U)$, il y a par hypothèse un voisinage
$V(x)\subset f^{-1}(U)$ de $x$ et un voisinage $W(x)\subset U$ de
$y=f(x)$, contenant $f(V(x))$ et tel que la restriction de $f$ à $V(x)$ soit
un morphisme $V(x)\to W(x)$ qui est de type fini. En remplaçant $W(x)$ par
un voisinage $W_1(x)\subset W(x)$ de $x$ de la forme $D(g)$ (avec $g\in A$),
et $V(x)$ par $V(x)\cap f^{-1}(W_1(x))$, on peut supposer que $W(x)$ est de
la forme $D(g)$, donc de type fini sur $U$ (puisque son anneau s'écrit
$A[1/g]$)~; par suite $V(x)$ est de type fini sur $U$. En outre
$f^{-1}(U)$ est quasi-compact par hypothèse, donc réunion d'un nombre fini
d'ouverts $V(x_i)$, ce qui achève la démonstration.

\begin{proposition}[6.6.4]
\label{I.6.6.4-fr}
\begin{enumerate}
  \item[(i)] Une immersion $X\to Y$ est quasi-compacte si elle est fermée,
    ou si l'espace sous-jacent à $Y$ est localement noethérien ou si l'espace
    sous-jacent à $X$ est noethérien.
  \item[(ii)] Le composé de deux morphismes quasi-compacts est quasi-compact.
  \item[(iii)] Si $f:X\to Y$ est un $S$-morphisme quasi-compact, il en est de
    même de $f_{(S')}:X_{(S')}\to Y_{(S')}$ pour toute extension $g:S'\to S$
    du préschéma de base.
  \item[(iv)] Si $f:X\to X'$ et $g:Y\to Y'$ sont deux $S$-morphismes
    quasi-compacts, $f\times_S g$ est quasi-compact.
  \item[(v)] Si le composé de deux morphismes $f:X\to Y$, $g:Y\to Z$ est
    quasi-compact et si $g$ est séparé, ou l'espace sous-jacent à $X$
    localement noethérien, $f$ est quasi-compact.
  \item[(vi)] Pour qu'un morphisme $f$ soit quasi-compact, il faut et il
    suffit que $f_{\mathrm{red}}$ le soit.
\end{enumerate}
\end{proposition}

On notera que (vi) est évident puisque la propriété d'être quasi-compact pour
un morphisme ne dépend que de l'application continue correspondante des
espaces sous-jacents. Démontrons de même la partie de (v) correspondant au cas
où l'espace sous-jacent $X$ est supposé localement noethérien. Posons
$h=g\circ f$, et soit $U$ un ouvert quasi-compact dans $Y$~; $g(U)$ est
quasi-compact (non nécessairement ouvert) dans $Z$, donc contenu dans une
réunion finie d'ouverts quasi-compacts $V_j$
(\hyperref[I.2.1.3-fr]{2.1.3}), et $f^{-1}(U)$ est par suite contenu dans la
réunion des $h^{-1}(V_j)$, qui sont des sous-espaces quasi-compacts de $X$,
donc des sous-espaces noethériens. On en conclut
(\hyperref[0.2.2.3-fr]{0, 2.2.3}) que $f^{-1}(U)$ est un espace noethérien,
et \emph{a fortiori} quasi-compact.

Pour prouver les autres assertions, il suffit de démontrer (i), (ii) et (iii)
(\hyperref[I.5.5.12-fr]{5.5.12}). Or, (ii) est évidente, et (i) résulte de
(\hyperref[I.6.3.5-fr]{6.3.5}) lorsque l'espace $Y$ est localement noethérien
ou l'espace $X$ noethérien et est évidente pour une immersion fermée. Pour
établir (iii),
\oldpage[I]{154}
on peut se borner au cas où $S=Y$ (\hyperref[I.3.3.11-fr]{3.3.11})~; posons
$f'=f_{(S')}$, et soit $U'$ un ouvert quasi-compact dans $S'$. Pour tout
$s'\in U'$, soit $T$ un voisinage ouvert affine de $g(s')$ dans $S$, et soit
$W$ un voisinage ouvert affine de $s'$ contenu dans
$U'\cap g^{-1}(T)$~; il suffira de montrer que $f'^{-1}(W)$ est
quasi-compact~; autrement dit, on peut se ramener à prouver que lorsque $S$
et $S'$ sont affines, l'espace sous-jacent à $X\times_S S'$ est
quasi-compact. Mais comme $X$ est alors par hypothèse réunion finie d'ouverts
affines $V_j$, $X\times_S S'$ est réunion des espaces sous-jacents aux
schémas affines $V_j\times_S S'$
(\hyperref[I.3.2.2-fr]{3.2.2} et \hyperref[I.3.2.7-fr]{3.2.7}), ce qui achève
de prouver la proposition.

On notera aussi que si $X=X'\amalg X''$ est somme de deux préschémas, un
morphisme $f:X\to Y$ est quasi-compact si et seulement si ses restrictions à
$X'$ et $X''$ le sont.

\begin{proposition}[6.6.5]
\label{I.6.6.5-fr}
Soit $f:X\to Y$ un morphisme quasi-compact. Pour que $f$ soit dominant, il
faut et il suffit que pour tout point générique $y$ d'une composante
irréductible de $Y$, $f^{-1}(y)$ contienne le point générique d'une composante
irréductible de $X$.
\end{proposition}

Il est immédiat que la condition est suffisante (sans supposer $f$
quasi-compact). Pour voir qu'elle est nécessaire, considérons un voisinage
ouvert affine $U$ de $y$~; $f^{-1}(U)$ est quasi-compact, donc réunion
\emph{finie} d'ouverts affines $V_i$, et l'hypothèse que $f$ est dominant
implique que $y$ appartient à l'adhérence dans $U$ d'un des $f(V_i)$. On peut
évidemment supposer $X$ et $Y$ réduits~; comme l'adhérence dans $X$ d'une
composante irréductible de $V_i$ est une composante irréductible de $X$
(\hyperref[0.2.1.6-fr]{0, 2.1.6}), on peut remplacer $X$ par $V_i$, $Y$ par
le sous-préschéma fermé réduit de $U$ ayant $\overline{f(V_i)}\cap U$ pour
espace sous-jacent (\hyperref[I.5.2.1-fr]{5.2.1}), et on est ainsi ramené à
démontrer la proposition lorsque $X=\operatorname{Spec}(A)$ et
$Y=\operatorname{Spec}(B)$ sont affines et réduits. Comme $f$ est dominant,
$B$ est alors un sous-anneau de $A$ (\hyperref[I.1.2.7-fr]{1.2.7}), et la
proposition résulte alors du fait que tout idéal premier minimal de $B$ est
l'intersection de $B$ et d'un idéal premier minimal de $A$
(\hyperref[0.1.5.8-fr]{0, 1.5.8}).

\begin{proposition}[6.6.6]
\label{I.6.6.6-fr}
\begin{enumerate}
  \item[(i)] Toute immersion locale est localement de type fini.
  \item[(ii)] Si deux morphismes $f:X\to Y$, $g:Y\to Z$ sont localement de
    type fini, il en est de même de $g\circ f$.
  \item[(iii)] Si $f:X\to Y$ est un $S$-morphisme localement de type fini,
    $f_{(S')}:X_{(S')}\to Y_{(S')}$ est localement de type fini pour toute
    extension $S'\to S$ du préschéma de base.
  \item[(iv)] Si $f:X\to X'$ et $g:Y\to Y'$ sont deux $S$-morphismes
    localement de type fini, $f\times_S g$ est localement de type fini.
  \item[(v)] Si le composé $g\circ f$ de deux morphismes est localement de
    type fini, $f$ est localement de type fini.
  \item[(vi)] Si un morphisme $f$ est localement de type fini, il en est de
    même de $f_{\mathrm{red}}$.
\end{enumerate}
\end{proposition}

En vertu de (\hyperref[I.5.5.12-fr]{5.5.12}), il suffit de démontrer (i), (ii)
et (iii). Si $j:X\to Y$ est une immersion locale, pour tout $x\in X$, il y a
un voisinage ouvert $V$ de $j(x)$ dans $Y$ et un voisinage ouvert $U$ de $x$
dans $X$ tels que la restriction de $j$ à $U$ soit une immersion fermée
$U\to V$ (\hyperref[I.4.5.1-fr]{4.5.1}), donc cette restriction est de type
fini. Pour établir (ii), considérons un point $x\in X$~; il y a par hypothèse
un voisinage ouvert $W$ de $g(f(x))$ et un voisinage ouvert $V$ de $f(x)$
tels que $g(V)\subset W$ et que $V$ soit de type fini sur $W$~; en outre
$f^{-1}(V)$ est localement de type fini sur $V$
(\hyperref[I.6.6.2-fr]{6.6.2}), donc il y a un voisinage ouvert $U$ de $x$ qui
\oldpage[I]{155}
est contenu dans $f^{-1}(V)$ et de type fini sur $V$~; par suite on a
$g(f(U))\subset W$ et $U$ est de type fini sur $W$
(\hyperref[I.6.3.4-fr]{6.3.4, (ii)}). Enfin, pour démontrer (iii), on peut se
borner au cas où $Y=S$ (\hyperref[I.3.3.11-fr]{3.3.11})~; pour tout
$x'\in X'=X_{(S')}$, soient $x$ l'image de $x'$ dans $X$, $s$ l'image de $x$
dans $S$, $T$ un voisinage ouvert de $s$, $T'$ son image réciproque dans $S'$,
$U$ un voisinage ouvert de $x$ dont l'image est contenue dans $T$ et qui est de
type fini sur $T$~; alors $U\times_S T'=U\times_T T'$ est un voisinage ouvert
de $x'$ (\hyperref[I.3.2.7-fr]{3.2.7}) qui est de type fini sur $T'$
(\hyperref[I.6.3.4-fr]{6.3.4, (iv)}).

\begin{corollary}[6.6.7]
\label{I.6.6.7-fr}
Soient $X$, $Y$ deux $S$-préschémas qui sont localement de type fini sur $S$.
Si $S$ est localement noethérien, $X\times_S Y$ est localement noethérien.
\end{corollary}

En effet, $X$ étant localement de type fini sur $S$, est localement
noethérien, et $X\times_S Y$ est localement de type fini sur $X$, donc est
aussi localement noethérien.

\begin{remark}[6.6.8]
\label{I.6.6.8-fr}
La proposition (\hyperref[I.6.3.10-fr]{6.3.10}) et sa démonstration s'étendent
aussitôt au cas où l'on suppose seulement que le morphisme $f$ est localement
de type fini. De même, les propositions
(\hyperref[I.6.4.2-fr]{6.4.2}) et (\hyperref[I.6.4.9-fr]{6.4.9}) restent
valables lorsqu'on suppose que les préschémas $X$, $Y$ qui figurent dans leur
énoncé sont seulement localement de type fini sur le corps $K$.
\end{remark}

\section{Applications rationnelles}
\label{section:I.7-fr}

\subsection{Applications rationnelles et fonctions rationnelles.}
\label{subsection:I.7.1-fr}

\begin{env}[7.1.1]
\label{I.7.1.1-fr}
Soient $X$, $Y$ deux préschémas, $U$, $V$, deux ouverts denses dans $X$, $f$
(resp. $g$) un morphisme de $U$ (resp. $V$) dans $Y$~; nous dirons que $f$ et
$g$ sont
\emph{équivalents} s'ils coïncident dans un ouvert dense dans $U\cap V$.
Comme une intersection finie d'ouverts denses dans $X$ est un ouvert dense
dans $X$, il est clair que cette relation est une
\emph{relation d'équivalence}.
\end{env}

\begin{definition}[7.1.2]
\label{I.7.1.2-fr}
Étant donnés deux préschémas $X$, $Y$, on appelle
\emph{application rationnelle de $X$ dans $Y$} une classe d'équivalence de
morphismes de parties ouvertes denses de $X$ dans $Y$, pour la relation
définie dans (\hyperref[I.7.1.1-fr]{7.1.1}). Si $X$ et $Y$ sont des
$S$-préschémas, on dit qu'une telle classe est une
\emph{$S$-application rationnelle} s'il existe un représentant de cette classe
qui est un $S$-morphisme. On appelle
\emph{$S$-section rationnelle de $X$} toute $S$-application rationnelle de $S$
dans $X$. On appelle
\emph{fonction rationnelle sur un préschéma $X$} toute $X$-section rationnelle
du $X$-préschéma $X\otimes_{\mathbf{Z}}\mathbf{Z}[T]$ (où $T$ est une
indéterminée).

Par abus de langage, lorsqu'il ne sera question que de $S$-préschémas, on dira
«~application rationnelle~» au lieu de «~$S$-application rationnelle~» si
aucune confusion ne peut en résulter.

Soient $f$ une application rationnelle de $X$ dans $Y$, $U$ un ouvert de
$X$~; si $f_1$, $f_2$ sont des morphismes appartenant à la classe $f$, définis
respectivement dans des ouverts denses $V$, $W$ de $X$, les restrictions
$f_1|(U\cap V)$, $f_2|(U\cap W)$ coïncident dans $U\cap V\cap W$ qui est dense
dans $U$~; la classe de morphismes $f$ définit donc une application rationnelle
de $U$ dans $Y$, appelée
\emph{restriction de $f$ à $U$} et notée $f|U$.

Si, à tout $S$-morphisme $f:X\to Y$, on fait correspondre la $S$-application
rationnelle
\oldpage[I]{156}
à laquelle appartient $f$, on définit une application canonique de
$\operatorname{Hom}_S(X,Y)$ dans l'ensemble des $S$-applications rationnelles
de $X$ dans $Y$. On désigne par $\Gamma_{\mathrm{rat}}(X/Y)$ l'ensemble des
$Y$-sections rationnelles de $X$, et on a donc une application canonique
$\Gamma(X/Y)\to\Gamma_{\mathrm{rat}}(X/Y)$. Il est clair en outre que si $X$ et
$Y$ sont deux $S$-préschémas, l'ensemble des $S$-applications rationnelles de
$X$ dans $Y$ s'identifie canoniquement à
$\Gamma_{\mathrm{rat}}((X\times_S Y)/X)$
(\hyperref[I.3.3.14-fr]{3.3.14}).
\end{definition}

\begin{env}[7.1.3]
\label{I.7.1.3-fr}
Il résulte aussitôt de (\hyperref[I.7.1.2-fr]{7.1.2}) et de
(\hyperref[I.3.3.14-fr]{3.3.14}) que les fonctions rationnelles sur $X$
s'identifient canoniquement aux
\emph{classes d'équivalence des sections du faisceau structural
$\mathscr{O}_X$} au-dessus d'ouverts partout denses de $X$, deux telles
sections étant équivalentes si elles coïncident dans un ouvert partout dense
contenu dans l'intersection de leurs ensembles de définition. Il en résulte en
particulier que les fonctions rationnelles sur $X$ forment un
\emph{anneau} $R(X)$.
\end{env}

\begin{env}[7.1.4]
\label{I.7.1.4-fr}
Lorsque $X$ est un préschéma \emph{irréductible}, tout ouvert non vide est dense
dans $X$~; on peut encore dire que les ouverts non vides de $X$ sont les
\emph{voisinages ouverts du point générique $x$} de $X$. Dire que deux
morphismes de parties ouvertes non vides de $X$ dans $Y$ sont équivalents
signifie donc dans ce cas qu'ils ont \emph{même germe} au point $x$. Autrement
dit, les applications rationnelles (resp. $S$-applications rationnelles)
$X\to Y$ s'identifient aux \emph{germes de morphismes} (resp. de
$S$-morphismes) de parties ouvertes non vides de $X$ dans $Y$ au point
générique $x$ de $X$. En particulier~:
\end{env}

\begin{proposition}[7.1.5]
\label{I.7.1.5-fr}
Si $X$ est un préschéma irréductible, l'anneau $R(X)$ des fonctions rationnelles
sur $X$ s'identifie canoniquement à l'anneau local $\mathscr{O}_x$ du point
générique $x$ de $X$. C'est un anneau local de dimension 0, et par suite un
anneau local artinien lorsque $X$ est noethérien~; c'est un corps lorsque $X$
est intègre, et il s'identifie au corps des fractions de $A(X)$ lorsqu'en outre
$X$ est un schéma affine.
\end{proposition}

Vu ce qui précède et l'identification des fonctions rationnelles aux sections
de $\mathscr{O}_X$ au-dessus d'un ouvert partout dense, la première assertion
n'est autre que la définition de la fibre d'un faisceau en un point. Pour les
autres assertions, on peut se borner au cas où $X$ est affine d'anneau $A$~;
alors $\mathfrak{j}_x$ est le nilradical de $A$, et $\mathscr{O}_x$ est donc de
dimension 0~; si $A$ est intègre, $\mathfrak{j}_x=(0)$, et $\mathscr{O}_x$ est
donc le corps des fractions de $A$. Enfin, si $A$ est noethérien, on sait
(\cite{I-11}, p.~127, cor.~4) que $\mathfrak{j}_x$ est nilpotent et
$\mathscr{O}_x=A_{\mathfrak{j}_x}$ artinien.

Si $X$ est \emph{intègre}, l'anneau $\mathscr{O}_z$ est intègre pour tout
$z\in X$~; tout ouvert affine $U$ contenant $z$ contient aussi $x$, et $R(U)$,
égal au corps des fractions de $A(U)$, s'identifie donc à $R(X)$~; on en
conclut que $R(X)$ s'identifie aussi au
\emph{corps des fractions de $\mathscr{O}_z$}~: l'identification canonique de
$\mathscr{O}_z$ à un sous-anneau de $R(X)$ consiste à faire correspondre à tout
germe de section $s\in\mathscr{O}_z$ l'unique fonction rationnelle sur $X$,
classe d'une section de $\mathscr{O}_X$ (nécessairement définie sur un ouvert
partout dense) ayant pour germe $s$ au point $z$.

\begin{env}[7.1.6]
\label{I.7.1.6-fr}
Supposons maintenant que $X$ ait un nombre \emph{fini} de composantes
irréductibles $X_i$ ($1\leq i\leq n$) (ce qui est le cas lorsque l'espace
sous-jacent à $X$ est \emph{noethérien})~; soit $X'_i$ l'ouvert de $X$,
complémentaire par rapport à $X_i$ de la réunion des $X_j\cap X_i$ pour
$j\neq i$~; $X'_i$ est irréductible, son point générique $x_i$ est le point
générique de $X_i$, et les $X'_i$ sont deux à deux disjoints, leur réunion
étant dense dans $X$ (\hyperref[0.2.1.6-fr]{0, 2.1.6}). Pour
\oldpage[I]{157}
tout ouvert partout dense $U$ de $X$, $U_i=U\cap X'_i$ est un ouvert non vide
dense dans $X'_i$, les $U_i$ étant deux à deux sans point commun, donc
$U'=\bigcup_i U_i$ est dense dans $X$. Se donner un morphisme de $U'$ dans $Y$
revient à se donner (arbitrairement) un morphisme de chacun des $U_i$ dans $Y$.
Donc~:
\end{env}

\begin{proposition}[7.1.7]
\label{I.7.1.7-fr}
Soient $X$, $Y$ deux préschémas (resp. $S$-préschémas) tels que $X$ ait un
nombre fini de composantes irréductibles $X_i$, de points génériques $x_i$
($1\leq i\leq n$). Si $R_i$ est l'ensemble des germes de morphismes (resp.
$S$-morphismes) de parties ouvertes de $X$ dans $Y$ au point $x_i$, l'ensemble
des applications rationnelles (resp. $S$-applications rationnelles) de $X$
dans $Y$ s'identifie au produit des $R_i$ ($1\leq i\leq n$).
\end{proposition}

\begin{corollary}[7.1.8]
\label{I.7.1.8-fr}
Soit $X$ un préschéma noethérien. L'anneau des fonctions rationnelles sur $X$
est un anneau artinien, dont les composants locaux sont les anneaux
$\mathscr{O}_{x_i}$ des points génériques $x_i$ des composantes irréductibles
de $X$.
\end{corollary}

\begin{corollary}[7.1.9]
\label{I.7.1.9-fr}
Soient $A$ un anneau noethérien, et soit $X=\operatorname{Spec}(A)$. Si $Q$ est
le complémentaire de la réunion des idéaux premiers minimaux de $A$, l'anneau
des fonctions rationnelles sur $X$ s'identifie canoniquement à l'anneau de
fractions $Q^{-1}A$.
\end{corollary}

Cela résultera du lemme suivant~:

\begin{lemma}[7.1.9.1]
\label{I.7.1.9.1-fr}
Pour qu'un élément $f\in A$ soit tel que $D(f)$ soit partout dense dans $X$, il
faut et il suffit que $f\in Q$~; tout ouvert dense dans $X$ contient un ouvert
de la forme $D(f)$, où $f\in Q$.
\end{lemma}

En effet, supposons ce lemme démontré~; l'anneau des sections
$\Gamma(D(f),\mathscr{O}_X)$ s'identifiant à $A_f$
(\hyperref[I.1.3.6-fr]{1.3.6} et \hyperref[I.1.3.7-fr]{1.3.7}), il résulte du
fait que les $D(f)$ avec $f\in Q$ forment un ensemble cofinal dans l'ensemble
ordonné (pour $\supset$) des ouverts denses dans $X$, et de la déf.
(\hyperref[I.7.1.1-fr]{7.1.1}) que l'anneau des fonctions rationnelles sur $X$
s'identifie à la limite inductive des $A_f$ pour $f\in Q$ (pour la relation de
préordre «~$g$ est multiple de $f$~»), c'est-à-dire à $Q^{-1}A$
(\hyperref[0.1.4.5-fr]{0, 1.4.5}).

Pour démontrer (\hyperref[I.7.1.9.1-fr]{7.1.9.1}), désignons encore par $X_i$
les composantes irréductibles de $X$ ($1\leq i\leq n$)~; si $D(f)$ est dense
dans $X$, $D(f)\cap X_i\neq\varnothing$ pour $1\leq i\leq n$ et
réciproquement~; mais cela signifie que $f\notin\mathfrak{p}_i$ pour
$1\leq i\leq n$, en posant $\mathfrak{p}_i=\mathfrak{j}(X_i)$, et comme les
$\mathfrak{p}_i$ sont les idéaux premiers minimaux de $A$
(\hyperref[I.1.1.14-fr]{1.1.14}), les relations
$f\notin\mathfrak{p}_i$ ($1\leq i\leq n$) équivalent à $f\in Q$, d'où la
première assertion du lemme. D'autre part, si $U$ est un ouvert dense dans $X$,
le complémentaire de $U$ est un ensemble de la forme $V(\mathfrak{a})$, où
$\mathfrak{a}$ est un idéal qui n'est contenu dans aucun des $\mathfrak{p}_i$~;
il n'est donc pas contenu dans leur réunion (\cite{I-10}, p.~13), et il existe
donc un $f\in\mathfrak{a}$ appartenant à $Q$~; d'où $D(f)\subset U$, ce qui
achève la démonstration.

\begin{env}[7.1.10]
\label{I.7.1.10-fr}
Supposons de nouveau $X$ irréductible, de point générique $x$. Comme tout
ouvert non vide $U$ de $X$ contient $x$, et par suite contient aussi tout
$z\in X$ tel que $x\in\overline{\{z\}}$, tout morphisme $U\to Y$ peut se
composer avec le morphisme canonique $\operatorname{Spec}(\mathscr{O}_x)\to X$
(\hyperref[I.2.4.1-fr]{2.4.1})~; et deux morphismes dans $Y$ de deux parties
ouvertes non vides de $X$, qui coïncident dans une partie ouverte non vide de
$X$ donnent par composition le même morphisme
$\operatorname{Spec}(\mathscr{O}_x)\to Y$. Autrement dit, à toute application
rationnelle de $X$ dans $Y$ correspond ainsi un morphisme
$\operatorname{Spec}(\mathscr{O}_x)\to Y$ bien déterminé.
\end{env}

\oldpage[I]{158}
\begin{proposition}[7.1.11]
\label{I.7.1.11-fr}
Soient $X$, $Y$ deux $S$-préschémas~; on suppose $X$ irréductible, de point
générique $x$, et $Y$ de type fini sur $S$. Deux $S$-applications rationnelles
de $X$ dans $Y$, auxquelles correspond le même $S$-morphisme
$\operatorname{Spec}(\mathscr{O}_x)\to Y$ sont identiques. Si on suppose de
plus $S$ localement noethérien, tout $S$-morphisme de
$\operatorname{Spec}(\mathscr{O}_x)$ dans $Y$ correspond à une $S$-application
rationnelle (et une seule) de $X$ dans $Y$.
\end{proposition}

Compte tenu de ce que tout ouvert non vide de $X$ est partout dense, cela
résulte aussitôt de (\hyperref[I.6.5.1-fr]{6.5.1}).

\begin{corollary}[7.1.12]
\label{I.7.1.12-fr}
On suppose $S$ localement noethérien, et les autres hypothèses de
(\hyperref[I.7.1.11-fr]{7.1.11}) satisfaites. Les $S$-applications rationnelles
de $X$ dans $Y$ s'identifient alors aux points du $S$-préschéma $Y$, à valeurs
dans le $S$-préschéma $\operatorname{Spec}(\mathscr{O}_x)$.
\end{corollary}

Cela n'est autre que (\hyperref[I.7.1.11-fr]{7.1.11}), avec la terminologie
introduite dans (\hyperref[I.3.4.1-fr]{3.4.1}).

\begin{corollary}[7.1.13]
\label{I.7.1.13-fr}
On suppose remplies les conditions de (\hyperref[I.7.1.12-fr]{7.1.12}). Soit
$s$ l'image de $x$ dans $S$. La donnée d'une $S$-application rationnelle de
$X$ dans $Y$ équivaut à la donnée d'un point $y$ de $Y$ au-dessus de $s$, et
d'un $\mathscr{O}_s$-homomorphisme local
$\mathscr{O}_y\to\mathscr{O}_x=R(X)$.
\end{corollary}

Cela résulte de (\hyperref[I.7.1.11-fr]{7.1.11}) et de
(\hyperref[I.2.4.4-fr]{2.4.4}).

En particulier~:

\begin{corollary}[7.1.14]
\label{I.7.1.14-fr}
Sous les conditions de (\hyperref[I.7.1.12-fr]{7.1.12}), les $S$-applications
rationnelles de $X$ dans $Y$ ne dépendent (pour $Y$ donné) que du
$S$-préschéma $\operatorname{Spec}(\mathscr{O}_x)$ et en particulier restent
les mêmes lorsqu'on remplace $X$ par
$\operatorname{Spec}(\mathscr{O}_z)$, pour tout $z\in X$.
\end{corollary}

En effet, comme $z\in\overline{\{x\}}$, $x$ est le point générique de
$Z=\operatorname{Spec}(\mathscr{O}_z)$, et
$\mathscr{O}_{X,x}=\mathscr{O}_{Z,x}$.

Lorsque $X$ est intègre, $R(X)=\mathscr{O}_x=k(x)$ est un corps
(\hyperref[I.7.1.5-fr]{7.1.5})~; les corollaires précédents se spécialisent
alors en~:

\begin{corollary}[7.1.15]
\label{I.7.1.15-fr}
On suppose vérifiées les conditions de (\hyperref[I.7.1.12-fr]{7.1.12}) et en
outre que $X$ est intègre. Soit $s$ l'image de $x$ dans $S$. Alors les
$S$-applications rationnelles de $X$ dans $Y$ s'identifient aux points
géométriques de $Y\otimes_S k(s)$ à valeurs dans l'extension $R(X)$ de $k(s)$,
autrement dit chacune d'elles équivaut à la donnée d'un point $y\in Y$
au-dessus de $s$ et d'un $k(s)$-monomorphisme de $k(y)$ dans
$k(x)=R(X)$.
\end{corollary}

Les points de $Y$ au-dessus de $s$ s'identifient en effet à ceux de
$Y\otimes_S k(s)$ (\hyperref[I.3.6.3-fr]{3.6.3}) et les
$\mathscr{O}_s$-homomorphismes locaux $\mathscr{O}_y\to R(X)$ aux
$k(s)$-monomorphismes $k(y)\to R(X)$.

Plus particulièrement~:

\begin{corollary}[7.1.16]
\label{I.7.1.16-fr}
Soient $k$ un corps, $X$, $Y$ deux préschémas algébriques
(\hyperref[I.6.4.1-fr]{6.4.1}) sur $k$~; on suppose en outre $X$ intègre. Alors
les $k$-applications rationnelles de $X$ dans $Y$ s'identifient aux points
géométriques de $Y$ à valeurs dans l'extension $R(X)$ de $k$
(\hyperref[I.3.4.4-fr]{3.4.4}).
\end{corollary}

\subsection{Domaine de définition d'une application rationnelle.}
\label{subsection:I.7.2-fr}

\begin{env}[7.2.1]
\label{I.7.2.1-fr}
Soient $X$, $Y$ deux préschémas, $f$ une application rationnelle de $X$ dans
$Y$. On dit que $f$ est \emph{définie en un point $x\in X$} s'il existe un
ensemble ouvert partout dense $U$ contenant $x$ et un morphisme $U\to Y$
appartenant à la classe d'équivalence $f$. L'ensemble des points $x\in X$ où
$f$ est définie est appelé le \emph{domaine de définition de $f$}~; il est
clair que c'est un ouvert partout dense dans $X$.
\end{env}

\oldpage[I]{159}
\begin{proposition}[7.2.2]
\label{I.7.2.2-fr}
Soient $X$, $Y$ deux $S$-préschémas, tels que $X$ soit réduit et $Y$ séparé
sur $S$. Soit $f$ une $S$-application rationnelle de $X$ dans $Y$, $U_0$ son
domaine de définition. Il existe alors un $S$-morphisme et un seul $U_0\to Y$
appartenant à la classe $f$.
\end{proposition}

Comme pour tout morphisme $U\to Y$ appartenant à la classe $f$, on a
nécessairement $U\subset U_0$, il est clair que la proposition sera
conséquence du

\begin{lemma}[7.2.2.1]
\label{I.7.2.2.1-fr}
Sous les hypothèses de (\hyperref[I.7.2.2-fr]{7.2.2}), soient $U_1$, $U_2$
deux ouverts partout denses de $X$, $f_i:U_i\to Y$ ($i=1,2$) deux
$S$-morphismes tels qu'il existe un ouvert $V\subset U_1\cap U_2$, dense dans
$X$ et dans lequel $f_1$ et $f_2$ coïncident. Alors $f_1$ et $f_2$ coïncident
dans $U_1\cap U_2$.
\end{lemma}

On peut évidemment se borner au cas où $X=U_1=U_2$. Comme $X$ (donc $V$)
est réduit, $X$ est le plus petit sous-préschéma fermé de $X$ majorant $V$
(\hyperref[I.5.2.2-fr]{5.2.2}). Soit
$g=(f_1,f_2)_S:X\to Y\times_S Y$~; comme par hypothèse la diagonale
$T=\Delta_Y(Y)$ est un sous-préschéma fermé de $Y\times_S Y$,
$Z=g^{-1}(T)$ est un sous-préschéma fermé de $X$
(\hyperref[I.4.4.1-fr]{4.4.1}). Si $h:V\to Y$ est la restriction commune de
$f_1$ et $f_2$ à $V$, la restriction de $g$ à $V$ est $g'=(h,h)_S$, qui se
factorise en $g'=\Delta_Y\circ h$~; comme $\Delta_Y^{-1}(T)=Y$, on a
$g'^{-1}(T)=V$, et par suite $Z$ est un sous-préschéma fermé de $X$ induisant
$V$, donc majorant $V$, ce qui entraîne $Z=X$. De la relation
$g^{-1}(T)=X$, on déduit (\hyperref[I.4.4.1-fr]{4.4.1}) que $g$ se factorise
en $\Delta_Y\circ f$, où $f$ est un morphisme $X\to Y$, ce qui entraîne par
définition du morphisme diagonal que $f_1=f_2=f$.

Il est clair que le morphisme $U_0\to Y$ défini dans
(\hyperref[I.7.2.2-fr]{7.2.2}) est l'unique morphisme de la classe $f$ qui
\emph{ne puisse être prolongé} à un morphisme d'une partie ouverte de $X$
contenant strictement $U_0$. \emph{Sous les hypothèses de
(\hyperref[I.7.2.2-fr]{7.2.2})}, on peut donc \emph{identifier} les
applications rationnelles de $X$ dans $Y$ aux morphismes \emph{non
prolongeables} (à des ouverts strictement plus grands) d'ouverts partout
denses de $X$ dans $Y$. Avec cette identification, la prop.
(\hyperref[I.7.2.2-fr]{7.2.2}) entraîne~:

\begin{corollary}[7.2.3]
\label{I.7.2.3-fr}
Les hypothèses sur $X$ et $Y$ étant celles de
(\hyperref[I.7.2.2-fr]{7.2.2}), soit $U$ un ouvert partout dense de $X$. Il
existe une correspondance biunivoque canonique entre les $S$-morphismes de
$U$ dans $Y$ et les $S$-applications rationnelles de $X$ dans $Y$ définies
en tous les points de $U$.
\end{corollary}

En vertu de (\hyperref[I.7.2.2-fr]{7.2.2}), pour tout $S$-morphisme $f$ de
$U$ dans $Y$, il existe en effet une $S$-application rationnelle et une seule
$\overline f$ de $X$ dans $Y$ qui prolonge $f$.

\begin{corollary}[7.2.4]
\label{I.7.2.4-fr}
Soient $S$ un schéma, $X$ un $S$-préschéma réduit, $Y$ un $S$-schéma,
$f:U\to Y$ un $S$-morphisme d'un ouvert dense $U$ de $X$ dans $Y$. Si
$\overline f$ est la $\mathbf{Z}$-application rationnelle de $X$ dans $Y$ qui
prolonge $f$, $\overline f$ est un $S$-morphisme (et est par suite la
$S$-application rationnelle de $X$ dans $Y$ prolongeant $f$).
\end{corollary}

En effet, si $\varphi:X\to S$, $\psi:Y\to S$ sont les morphismes structuraux,
$U_0$ le domaine de définition de $\overline f$, $j$ l'injection $U_0\to X$,
il suffit de prouver que $\psi\circ\overline f=\varphi\circ j$, ce qui résulte
aussitôt de (\hyperref[I.7.2.2.1-fr]{7.2.2.1}), puisque $f$ est un
$S$-morphisme.

\begin{corollary}[7.2.5]
\label{I.7.2.5-fr}
Soient $X$, $Y$ deux $S$-préschémas~; on suppose $X$ réduit, $X$ et $Y$
séparés sur $S$. Soient $p:Y\to X$ un $S$-morphisme (faisant de $Y$ un
$X$-préschéma), $U$ un ouvert partout dense de $X$, $f$ une $U$-section de
$Y$~; alors l'application rationnelle $\overline f$ de $X$ dans $Y$
prolongeant $f$ est une $X$-section rationnelle de $Y$.
\end{corollary}

\oldpage[I]{160}
Il faut prouver que $p\circ\overline f$ est l'identité dans le domaine de
définition de $\overline f$~; puisque $X$ est séparé sur $S$, cela résulte
encore de (\hyperref[I.7.2.2.1-fr]{7.2.2.1}).

\begin{corollary}[7.2.6]
\label{I.7.2.6-fr}
Soient $X$ un préschéma réduit, $U$ un ouvert partout dense de $X$. Il y a
correspondance biunivoque canonique entre les sections de $\mathscr{O}_X$ au-dessus
de $U$ et les fonctions rationnelles $f$ sur $X$ définies en tout point de $U$.
\end{corollary}

Compte tenu de (\hyperref[I.7.2.3-fr]{7.2.3}),
(\hyperref[I.7.1.2-fr]{7.1.2}) et (\hyperref[I.7.1.3-fr]{7.1.3}), il suffit de
remarquer que le $X$-préschéma $X\otimes_{\mathbf Z}\mathbf Z[T]$ est séparé
au-dessus de $X$ (\hyperref[I.5.5.1-fr]{5.5.1, (iv)}).

\begin{corollary}[7.2.7]
\label{I.7.2.7-fr}
Soient $Y$ un préschéma réduit, $f:X\to Y$ un morphisme séparé, $U$ un ouvert
partout dense de $Y$, $g:U\to f^{-1}(U)$ une $U$-section de $f^{-1}(U)$, $Z$
le sous-préschéma réduit de $X$ ayant $\overline{g(U)}$ pour espace sous-jacent
(\hyperref[I.5.2.1-fr]{5.2.1}). Pour que $g$ soit restriction d'une
$Y$-section de $X$ (autrement dit (\hyperref[I.7.2.5-fr]{7.2.5}) pour que
l'application rationnelle de $Y$ dans $X$ prolongeant $g$ soit partout
définie), il faut et il suffit que la restriction de $f$ à $Z$ soit un
isomorphisme de $Z$ sur $Y$.
\end{corollary}

La restriction de $f$ à $f^{-1}(U)$ est un morphisme séparé
(\hyperref[I.5.5.1-fr]{5.5.1, (i)}), donc $g$ est une immersion fermée
(\hyperref[I.5.4.6-fr]{5.4.6}), et par suite
$g(U)=Z\cap f^{-1}(U)$ et le sous-préschéma induit par $Z$ sur l'ouvert
$g(U)$ de $Z$ est identique au sous-préschéma fermé de $f^{-1}(U)$ associé à
$g$ (\hyperref[I.5.2.1-fr]{5.2.1}). Il est clair alors que la condition de
l'énoncé est suffisante, car si elle est remplie et si $f_Z:Z\to Y$ est la
restriction de $f$ à $Z$ et $\overline g:Y\to Z$ l'isomorphisme réciproque,
$\overline g$ prolonge $g$. Inversement, si $g$ est restriction à $U$ d'une
$Y$-section $h$ de $X$, $h$ est une immersion fermée
(\hyperref[I.5.4.6-fr]{5.4.6}), donc $h(Y)$ est fermé, et comme il est contenu
dans $Z$, il est égal à $Z$, et il résulte de
(\hyperref[I.5.2.1-fr]{5.2.1}) que $h$ est nécessairement un isomorphisme de
$Y$ sur le sous-préschéma fermé $Z$ de $X$.

\begin{env}[7.2.8]
\label{I.7.2.8-fr}
Soient $X$, $Y$ deux $S$-préschémas, $X$ étant supposé réduit et $Y$ séparé
sur $S$. Soit $f$ une $S$-application rationnelle de $X$ dans $Y$, et soit
$x$ un point de $X$~; on peut composer $f$ avec le $S$-morphisme canonique
$\operatorname{Spec}(\mathscr{O}_x)\to X$
(\hyperref[I.2.4.1-fr]{2.4.1}) pourvu que la trace sur
$\operatorname{Spec}(\mathscr{O}_x)$ du domaine de définition de $f$ soit
dense dans $\operatorname{Spec}(\mathscr{O}_x)$ (identifié à l'ensemble des
$z\in X$ tels que $x\in\overline{\{z\}}$
(\hyperref[I.2.4.2-fr]{2.4.2})). Ceci aura lieu dans les cas suivants~:
\begin{enumerate}
  \item[1\textsuperscript{o}] $X$ est irréductible (donc intègre), car alors
    le point générique $\xi$ de $X$ est le point générique de
    $\operatorname{Spec}(\mathscr{O}_x)$~; comme le domaine de définition $U$
    de $f$ contient $\xi$, $U\cap\operatorname{Spec}(\mathscr{O}_x)$ contient
    $\xi$, donc est dense dans $\operatorname{Spec}(\mathscr{O}_x)$.
  \item[2\textsuperscript{o}] $X$ est localement noethérien~; notre assertion
    résulte en effet alors du
\end{enumerate}
\end{env}

\begin{lemma}[7.2.8.1]
\label{I.7.2.8.1-fr}
Soient $X$ un préschéma dont l'espace sous-jacent est localement noethérien,
$x$ un point de $X$. Les composantes irréductibles de
$\operatorname{Spec}(\mathscr{O}_x)$ sont les traces sur
$\operatorname{Spec}(\mathscr{O}_x)$ des composantes irréductibles de $X$
contenant $x$. Pour qu'un ouvert $U\subset X$ soit tel que
$U\cap\operatorname{Spec}(\mathscr{O}_x)$ soit dense dans
$\operatorname{Spec}(\mathscr{O}_x)$, il faut et il suffit qu'il rencontre les
composantes irréductibles de $X$ contenant $x$ (ce qui a lieu en particulier
si $U$ est dense dans $X$).
\end{lemma}

La seconde assertion résulte évidemment de la première, et il suffit donc de
démontrer celle-ci. Comme $\operatorname{Spec}(\mathscr{O}_x)$ est contenu
dans tout ouvert affine $U$ contenant $x$, et que les composantes
irréductibles de $U$ contenant $x$ sont les traces sur $U$ des composantes
irréductibles de $X$ contenant $x$
(\hyperref[0.2.1.6-fr]{0, 2.1.6}), on peut supposer $X$ affine d'anneau $A$.
Comme les idéaux premiers de $A_x$ correspondent biunivoquement aux idéaux
premiers
\oldpage[I]{161}
de $A$ contenus dans $\mathfrak{j}_x$
(\hyperref[0.1.2.6-fr]{0, 1.2.6}), les idéaux premiers minimaux de $A_x$
correspondent aux idéaux premiers minimaux de $A$ contenus dans
$\mathfrak{j}_x$, d'où le lemme (\hyperref[I.1.1.14-fr]{1.1.14}).

Cela étant, supposons que l'on soit dans l'un des deux cas précités. Si $U$
est le domaine de définition de la $S$-application rationnelle $f$, désignons
par $f'$ l'application rationnelle de $\operatorname{Spec}(\mathscr{O}_x)$
dans $Y$ qui coïncide (compte tenu de (\hyperref[I.2.4.2-fr]{2.4.2})) avec
$f$ dans $U\cap\operatorname{Spec}(\mathscr{O}_x)$~; nous dirons que cette
application rationnelle est \emph{induite par $f$}.

\begin{proposition}[7.2.9]
\label{I.7.2.9-fr}
Soient $S$ un préschéma localement noethérien, $X$ un $S$-préschéma réduit,
$Y$ un $S$-schéma de type fini. On suppose en outre $X$ irréductible ou
localement noethérien. Soient alors $f$ une $S$-application rationnelle de
$X$ dans $Y$, $x$ un point de $X$. Pour que $f$ soit définie au point $x$,
il faut et il suffit que l'application rationnelle $f'$ de
$\operatorname{Spec}(\mathscr{O}_x)$ dans $Y$, induite par $f$
(\hyperref[I.7.2.8-fr]{7.2.8}), soit un morphisme.
\end{proposition}

La condition étant évidemment nécessaire (puisque
$\operatorname{Spec}(\mathscr{O}_x)$ est contenu dans tout ouvert contenant
$x$), prouvons qu'elle est suffisante. En vertu de
(\hyperref[I.6.5.1-fr]{6.5.1}), il existe un voisinage ouvert $U$ de $x$ dans
$X$ et un $S$-morphisme $g$ de $U$ dans $Y$, induisant $f'$ sur
$\operatorname{Spec}(\mathscr{O}_x)$. Si $X$ est irréductible, $U$ est dense
dans $X$, et en vertu de (\hyperref[I.7.2.3-fr]{7.2.3}) on peut supposer que
$g$ est une $S$-application rationnelle. En outre, le point générique de $X$
appartient à $\operatorname{Spec}(\mathscr{O}_x)$ et au domaine de définition
de $f$, donc $f$ et $g$ coïncident en ce point, et par suite dans un ensemble
ouvert non vide de $X$ (\hyperref[I.6.5.1-fr]{6.5.1}). Mais comme $f$ et $g$
sont des $S$-applications rationnelles, elles sont identiques
(\hyperref[I.7.2.3-fr]{7.2.3}), donc $f$ est définie en $x$.

Si maintenant on suppose $X$ localement noethérien, on peut supposer $U$
noethérien~; il n'y a alors qu'un nombre fini de composantes irréductibles
$X_i$ de $X$ contenant $x$ (\hyperref[I.7.2.8.1-fr]{7.2.8.1}), et on peut
supposer que ce sont les seules rencontrant $U$, en remplaçant au besoin $U$
par un ouvert plus petit (puisqu'il n'y a qu'un nombre fini de composantes
irréductibles de $X$ rencontrant $U$, $U$ étant noethérien). On voit alors
comme ci-dessus que $f$ et $g$ coïncident dans un ouvert non vide de chacun
des $X_i$. Tenant compte du fait que chacun des $X_i$ est contenu dans
$\overline U$, considérons alors le morphisme $f_1$, défini dans un ouvert
dense de $U\cup(X-\overline U)$, égal à $g$ dans $U$ et à $f$ dans
l'intersection de $X-\overline U$ et du domaine de définition de $f$. Comme
$U\cup(X-\overline U)$ est dense dans $X$, $f_1$ et $f$ coïncident dans un
ouvert dense de $X$, et comme $f$ est une application rationnelle, $f$ est
une extension de $f_1$ (\hyperref[I.7.2.3-fr]{7.2.3}), donc est définie au
point $x$.

\subsection{Faisceau des fonctions rationnelles.}
\label{subsection:I.7.3-fr}

\begin{env}[7.3.1]
\label{I.7.3.1-fr}
Soit $X$ un préschéma. Pour tout ouvert $U\subset X$, désignons par $R(U)$
l'anneau des fonctions rationnelles sur $U$
(\hyperref[I.7.1.3-fr]{7.1.3})~; c'est une
$\Gamma(U,\mathscr{O}_X)$-algèbre. En outre, si $V\subset U$ est un second
ouvert de $X$, toute section de $\mathscr{O}_X$ au-dessus d'une partie
ouverte partout dense de $U$ donne par restriction à $V$ une section
au-dessus d'une partie ouverte partout dense de $V$, et si deux sections
coïncident au-dessus d'une partie ouverte partout dense de $U$, leurs
restrictions à $V$ coïncident au-dessus d'une partie ouverte partout dense de
$V$. On définit donc ainsi un di-homomorphisme d'algèbres
$\rho_{V,U}:R(U)\to R(V)$, et il est clair que si $U\supset V\supset W$ sont
trois ouverts de $X$, on a
$\rho_{W,U}=\rho_{W,V}\circ\rho_{V,U}$~; les $R(U)$ définissent donc un
\emph{préfaisceau} d'algèbres sur $X$.
\end{env}

\oldpage[I]{162}
\begin{definition}[7.3.2]
\label{I.7.3.2-fr}
On appelle \emph{faisceau des fonctions rationnelles} sur un préschéma $X$
et on désigne par $\mathscr{R}(X)$ la $\mathscr{O}_X$-Algèbre associée au
préfaisceau formé des $R(U)$.
\end{definition}

Pour tout préschéma $X$ et tout ouvert $U\subset X$, il est clair que le
faisceau induit $\mathscr{R}(X)|U$ n'est autre que $\mathscr{R}(U)$.

\begin{proposition}[7.3.3]
\label{I.7.3.3-fr}
Soit $X$ un préschéma tel que la famille $(X_\lambda)$ de ses composantes
irréductibles soit localement finie (ce qui est en particulier le cas lorsque
l'espace sous-jacent à $X$ est localement noethérien). Alors le
$\mathscr{O}_X$-Module $\mathscr{R}(X)$ est quasi-cohérent, et pour tout
ouvert $U$ de $X$ ne rencontrant qu'un nombre fini de composantes $X_\lambda$,
$R(U)$ est égal à $\Gamma(U,\mathscr{R}(X))$ et s'identifie canoniquement au
composé direct des anneaux locaux des points génériques des $X_\lambda$ telles
que $U\cap X_\lambda\neq\varnothing$.
\end{proposition}

On peut évidemment se limiter au cas où $X$ n'a qu'un nombre fini de
composantes irréductibles $X_i$, de points génériques $x_i$
($1\leq i\leq n$). Le fait que $R(U)$ est canoniquement identifié au composé
direct des $\mathscr{O}_{x_i}=R(X_i)$ tels que
$U\cap X_i\neq\varnothing$ résulte alors de
(\hyperref[I.7.1.7-fr]{7.1.7}). Montrons en outre que le préfaisceau
$U\to R(U)$ vérifie les axiomes des faisceaux, ce qui prouvera que
$R(U)=\Gamma(U,\mathscr{R}(X))$. En effet, il vérifie (F 1) d'après ce qui
précède. Pour voir qu'il satisfait à (F 2), considérons un recouvrement d'un
ouvert $U$ de $X$ par des ouverts $V_\alpha\subset U$~; si les
$s_\alpha\in R(V_\alpha)$ sont telles que les restrictions de $s_\alpha$ et
$s_\beta$ à $V_\alpha\cap V_\beta$ coïncident pour tout couple d'indices, on
en conclut que pour tout indice $i$ tel que $U\cap X_i\neq\varnothing$, les
composantes dans $R(X_i)$ de toutes les $s_\alpha$ telles que
$V_\alpha\cap X_i\neq\varnothing$ sont les mêmes~; désignant par $t_i$ cette
composante, il est clair que l'élément de $R(U)$ ayant les $t_i$ pour
composantes a pour restriction $s_\alpha$ à chaque $V_\alpha$. Enfin, pour
voir que $\mathscr{R}(X)$ est quasi-cohérent, on peut se limiter au cas où
$X=\operatorname{Spec}(A)$ est affine~; en prenant pour $U$ les ouverts
affines de la forme $D(f)$, où $f\in A$, il résulte de ce qui précède et de la
définition (\hyperref[I.1.3.4-fr]{1.3.4}) que l'on a
$\mathscr{R}(X)=\widetilde M$, où $M$ est somme directe des $A$-modules
$A_{x_i}$.

\begin{corollary}[7.3.4]
\label{I.7.3.4-fr}
Soit $X$ un préschéma réduit n'ayant qu'un nombre fini de composantes
irréductibles, et soient $X_i$ ($1\leq i\leq n$) les sous-préschémas fermés
réduits de $X$ ayant pour espaces sous-jacents les composantes irréductibles de
$X$ (\hyperref[I.5.2.1-fr]{5.2.1}). Si $h_i$ est l'injection canonique
$X_i\to X$, $\mathscr{R}(X)$ est alors composée directe des
$\mathscr{O}_X$-Algèbres $(h_i)_*(\mathscr{R}(X_i))$.
\end{corollary}

\begin{corollary}[7.3.5]
\label{I.7.3.5-fr}
Si $X$ est irréductible, tout $\mathscr{R}(X)$-Module quasi-cohérent
$\mathscr{F}$ est un faisceau simple.
\end{corollary}

Il suffit de montrer que tout $x\in X$ admet un voisinage $U$ tel que
$\mathscr{F}|U$ soit un faisceau simple
(\hyperref[0.3.6.2-fr]{0, 3.6.2}), autrement dit on est ramené au cas où $X$
est affine~; on peut en outre supposer que $\mathscr{F}$ est le conoyau d'un
homomorphisme
$(\mathscr{R}(X))^{(I)}\to(\mathscr{R}(X))^{(J)}$
(\hyperref[0.5.1.3-fr]{0, 5.1.3}), et tout revient à voir que
$\mathscr{R}(X)$ est un faisceau simple~; mais cela est évident puisque
$\Gamma(U,\mathscr{R}(X))=R(X)$ pour tout ouvert non vide $U$, $U$ contenant
le point générique de $X$.

\begin{corollary}[7.3.6]
\label{I.7.3.6-fr}
Si $X$ est irréductible, pour tout $\mathscr{O}_X$-Module quasi-cohérent
$\mathscr{F}$, $\mathscr{F}\otimes_{\mathscr{O}_X}\mathscr{R}(X)$ est un
faisceau simple~; si en outre $X$ est réduit (donc intègre),
$\mathscr{F}\otimes_{\mathscr{O}_X}\mathscr{R}(X)$ est isomorphe à un
faisceau de la forme $(\mathscr{R}(X))^{(I)}$.
\end{corollary}

La seconde assertion résulte de ce que $R(X)$ est alors un corps.

\begin{proposition}[7.3.7]
\label{I.7.3.7-fr}
Supposons que le préschéma $X$ soit localement intègre ou localement
\oldpage[I]{163}
noethérien. Alors $\mathscr{R}(X)$ est une $\mathscr{O}_X$-Algèbre
quasi-cohérente~; si en outre $X$ est réduit (ce qui est le cas lorsque $X$
est localement intègre), l'homomorphisme canonique
$\mathscr{O}_X\to\mathscr{R}(X)$ est injectif.
\end{proposition}

La question étant locale, la première assertion résulte de
(\hyperref[I.7.3.3-fr]{7.3.3})~; la seconde résulte aussitôt de
(\hyperref[I.7.2.3-fr]{7.2.3}).

\begin{env}[7.3.8]
\label{I.7.3.8-fr}
Soient $X$, $Y$ deux préschémas ayant chacun un nombre fini de composantes
irréductibles, et soit $f:X\to Y$ un morphisme dont la restriction à
l'ensemble des points génériques des composantes irréductibles de $X$ est une
surjection sur l'ensemble des points génériques des composantes irréductibles
de $Y$. Alors on a
\[
  f^*(\mathscr{R}(Y))=\mathscr{R}(X).
  \tag{7.3.8.1}\label{I.7.3.8.1-fr}
\]
\end{env}

En effet, on est ramené (en vertu de
(\hyperref[I.7.3.3-fr]{7.3.3})) au cas où $X$ et $Y$ sont irréductibles, de
points génériques $x$, $y$, avec $f(x)=y$~; donc
\[
  (f^*(\mathscr{R}(Y)))_x
  =\mathscr{O}_y\otimes_{\mathscr{O}_y}\mathscr{O}_x
  =\mathscr{O}_x
\]
(\hyperref[0.4.3.1-fr]{0, 4.3.1}), ce qui démontre (7.3.8.1) en vertu de
(\hyperref[I.7.3.5-fr]{7.3.5}).

\subsection{Faisceaux de torsion et faisceaux sans torsion.}
\label{subsection:I.7.4-fr}

\begin{env}[7.4.1]
\label{I.7.4.1-fr}
Soit $X$ un préschéma \emph{intègre}. Pour tout $\mathscr{O}_X$-Module
$\mathscr{F}$, l'homomorphisme canonique
$\mathscr{O}_X\to\mathscr{R}(X)$ définit par tensorisation un homomorphisme
(dit encore \emph{canonique})
\[
  \mathscr{F}\to
  \mathscr{F}\otimes_{\mathscr{O}_X}\mathscr{R}(X)
\]
qui, sur chaque fibre, n'est autre que l'homomorphisme $z\to z\otimes 1$ de
$\mathscr{F}_x$ dans
$\mathscr{F}_x\otimes_{\mathscr{O}_x}R(X)$. Le \emph{noyau $\mathscr{T}$}
de cet homomorphisme est un sous-$\mathscr{O}_X$-Module de $\mathscr{F}$,
appelé \emph{faisceau de torsion de $\mathscr{F}$}~; il est quasi-cohérent si
$\mathscr{F}$ est quasi-cohérent
(\hyperref[I.4.1.1-fr]{4.1.1} et \hyperref[I.7.3.6-fr]{7.3.6}). On dit que
$\mathscr{F}$ est \emph{sans torsion} si $\mathscr{T}=0$ et que
$\mathscr{F}$ est un \emph{faisceau de torsion} si
$\mathscr{T}=\mathscr{F}$. Pour tout $\mathscr{O}_X$-Module $\mathscr{F}$,
$\mathscr{F}/\mathscr{T}$ est sans torsion. On déduit de
(\hyperref[I.7.3.5-fr]{7.3.5}) que~:
\end{env}

\begin{proposition}[7.4.2]
\label{I.7.4.2-fr}
Si $X$ est un préschéma intègre, tout $\mathscr{O}_X$-Module quasi-cohérent
sans torsion $\mathscr{F}$ est isomorphe à un sous-faisceau $\mathscr{G}$
d'un faisceau simple de la forme $(\mathscr{R}(X))^{(I)}$, engendré (en tant
que $\mathscr{R}(X)$-Module) par $\mathscr{G}$.
\end{proposition}

Le cardinal de $I$ est appelé le \emph{rang de $\mathscr{F}$}~; pour tout
ouvert affine non vide $U$ de $X$, le rang de $\mathscr{F}$ est égal au rang
de $\Gamma(U,\mathscr{F})$ en tant que $\Gamma(U,\mathscr{O}_X)$-module,
comme on le voit aussitôt en considérant le point générique de $X$, contenu
dans $U$. En particulier~:

\begin{corollary}[7.4.3]
\label{I.7.4.3-fr}
Sur un préschéma intègre $X$, tout $\mathscr{O}_X$-Module quasi-cohérent sans
torsion et de rang $1$ (en particulier tout $\mathscr{O}_X$-Module
inversible) est isomorphe à un sous-$\mathscr{O}_X$-Module de
$\mathscr{R}(X)$, et réciproquement.
\end{corollary}

\begin{corollary}[7.4.4]
\label{I.7.4.4-fr}
Soient $X$ un préschéma intègre, $\mathscr{L}$, $\mathscr{L}'$ deux
$\mathscr{O}_X$-Modules sans torsion, $f$ (resp. $f'$) une section de
$\mathscr{L}$ (resp. $\mathscr{L}'$) au-dessus de $X$. Pour que
$f\otimes f'=0$, il faut et il suffit que l'une des sections $f$, $f'$ soit
nulle.
\end{corollary}

Soit $x$ le point générique de $X$~; on a par hypothèse
$(f\otimes f')_x=f_x\otimes f'_x=0$. Comme $\mathscr{L}_x$ et
$\mathscr{L}'_x$ s'identifient à des sous-$\mathscr{O}_x$-Modules du corps
$\mathscr{O}_x$, la relation précédente entraîne $f_x=0$ ou $f'_x=0$, et par
suite $f=0$ ou $f'=0$ puisque $\mathscr{L}$ et $\mathscr{L}'$ sont sans
torsion (\hyperref[I.7.3.5-fr]{7.3.5}).

\begin{proposition}[7.4.5]
\label{I.7.4.5-fr}
Soient $X$, $Y$ deux préschémas intègres, $f:X\to Y$ un morphisme dominant.
Pour tout $\mathscr{O}_X$-Module quasi-cohérent sans torsion $\mathscr{F}$,
$f_*(\mathscr{F})$ est un $\mathscr{O}_Y$-Module sans torsion.
\end{proposition}

\oldpage[I]{164}
\begin{proof}
Comme $f_*$ est exact à gauche
(\hyperref[0.4.2.1-fr]{0, 4.2.1}), il suffit, en vertu de
(\hyperref[I.7.4.2-fr]{7.4.2}), de prouver la proposition lorsque
$\mathscr{F}=(\mathscr{R}(X))^{(I)}$. Or, tout ouvert non vide $U$ de $Y$
contient le point générique de $Y$, donc $f^{-1}(U)$ contient le point
générique de $X$ (\hyperref[0.2.1.5-fr]{0, 2.1.5}), donc on a alors
\[
  \Gamma(U,f_*(\mathscr{F}))
  =\Gamma(f^{-1}(U),\mathscr{F})=(R(X))^{(I)}~;
\]
autrement dit, $f_*(\mathscr{F})$ est le faisceau simple de fibre
$(R(X))^{(I)}$, considéré comme $\mathscr{R}(Y)$-Module, et il est évidemment
sans torsion.
\end{proof}

\begin{proposition}[7.4.6]
\label{I.7.4.6-fr}
Soient $X$ un préschéma intègre, $x$ son point générique. Pour tout
$\mathscr{O}_X$-Module quasi-cohérent de type fini $\mathscr{F}$, les
conditions suivantes sont équivalentes~: a) $\mathscr{F}$ est un faisceau de
torsion~; b) $\mathscr{F}_x=0$~; c)
$\operatorname{Supp}(\mathscr{F})\neq X$.
\end{proposition}

En vertu de (\hyperref[I.7.3.5-fr]{7.3.5}) et de
(\hyperref[I.7.4.1-fr]{7.4.1}), les relations $\mathscr{F}_x=0$ et
$\mathscr{F}\otimes_{\mathscr{O}_X}\mathscr{R}(X)=0$ sont équivalentes, donc
a) et b) sont équivalentes~; d'autre part,
$\operatorname{Supp}(\mathscr{F})$ est fermé dans $X$
(\hyperref[0.5.2.2-fr]{0, 5.2.2}), et comme tout ouvert non vide de $X$
contient $x$, b) et c) sont équivalentes.

\begin{env}[7.4.7]
\label{I.7.4.7-fr}
On étend (par abus de langage) les définitions de
(\hyperref[I.7.4.1-fr]{7.4.1}) au cas où $X$ est un préschéma
\emph{réduit} n'ayant qu'un nombre \emph{fini} de composantes irréductibles~;
il résulte alors de (\hyperref[I.7.3.4-fr]{7.3.4}) que l'équivalence de a) et
c) dans (\hyperref[I.7.4.6-fr]{7.4.6}) est encore valable pour un tel
préschéma.
\end{env}

\section{Les schémas de Chevalley}
\label{section:I.8-fr}

\subsection{Anneaux locaux apparentés.}
\label{subsection:I.8.1-fr}

Pour tout anneau local $A$, nous noterons $\mathfrak m(A)$ l'idéal maximal de
$A$.

\begin{lemma}[8.1.1]
\label{I.8.1.1-fr}
Soient $A$, $B$ deux anneaux locaux tels que $A\subset B$~; les conditions
suivantes sont équivalentes~:
\begin{enumerate}
  \item[(i)] $\mathfrak m(B)\cap A=\mathfrak m(A)$~;
  \item[(ii)] $\mathfrak m(A)\subset\mathfrak m(B)$~;
  \item[(iii)] $1$ n'appartient pas à l'idéal de $B$ engendré par
    $\mathfrak m(A)$.
\end{enumerate}
\end{lemma}

Il est évident que (i) entraîne (ii) et (ii) entraîne (iii)~; enfin, si (iii)
est vérifiée, $\mathfrak m(B)\cap A$ contient $\mathfrak m(A)$ et ne contient
pas $1$, donc est égal à $\mathfrak m(A)$.

Lorsque les conditions équivalentes de
(\hyperref[I.8.1.1-fr]{8.1.1}) sont remplies, on dit que $B$ \emph{domine}
$A$~; il revient au même de dire que l'injection $A\to B$ est un
homomorphisme \emph{local}. Il est clair que, dans l'ensemble des sous-anneaux
locaux d'un anneau $R$, la relation de domination est une relation d'ordre.

\begin{env}[8.1.2]
\label{I.8.1.2-fr}
Considérons maintenant un \emph{corps} $R$. Pour tout sous-anneau $A$ de $R$,
nous désignerons par $L(A)$ l'ensemble des anneaux locaux $A_\mathfrak p$, où
$\mathfrak p$ parcourt le spectre premier de $A$~; ils sont identifiés à des
sous-anneaux de $R$ contenant $A$. Comme
$\mathfrak p=(\mathfrak p A_\mathfrak p)\cap A$, l'application
$\mathfrak p\mapsto A_\mathfrak p$ de $\operatorname{Spec}(A)$ dans $L(A)$
est bijective.
\end{env}

\begin{lemma}[8.1.3]
\label{I.8.1.3-fr}
Soient $R$ un corps, $A$ un sous-anneau de $R$. Pour qu'un sous-anneau local
$M$ de $R$ domine un anneau $A_\mathfrak p\in L(A)$, il faut et il suffit que
$A\subset M$~; l'anneau local $A_\mathfrak p$ dominé par $M$ est alors unique
et correspond à $\mathfrak p=\mathfrak m(M)\cap A$.
\end{lemma}

En effet, si $M$ domine $A_\mathfrak p$, on a
$\mathfrak m(M)\cap A_\mathfrak p=\mathfrak p A_\mathfrak p$ d'après
(\hyperref[I.8.1.1-fr]{8.1.1}), d'où l'unicité de $\mathfrak p$~; d'autre
part, si $A\subset M$, $\mathfrak m(M)\cap A=\mathfrak p$ est premier dans
$A$, et comme $A-\mathfrak p\subset M$, on a $A_\mathfrak p\subset M$ et
$\mathfrak p A_\mathfrak p\subset\mathfrak m(M)$, donc $M$ domine
$A_\mathfrak p$.

\oldpage[I]{165}
\begin{lemma}[8.1.4]
\label{I.8.1.4-fr}
Soient $R$ un corps, $M$, $N$ deux sous-anneaux locaux de $R$, $P$ le
sous-anneau de $R$ engendré par $M\cup N$. Les conditions suivantes sont
équivalentes~:
\begin{enumerate}
  \item[(i)] Il existe un idéal premier $\mathfrak p$ de $P$ tel que
    $\mathfrak m(M)=\mathfrak p\cap M$,
    $\mathfrak m(N)=\mathfrak p\cap N$.
  \item[(ii)] L'idéal $\mathfrak a$ engendré dans $P$ par
    $\mathfrak m(M)\cup\mathfrak m(N)$ est distinct de $P$.
  \item[(iii)] Il existe un sous-anneau local $Q$ de $R$ dominant à la fois
    $M$ et $N$.
\end{enumerate}
\end{lemma}

Il est clair que (i) implique (ii)~; inversement, si $\mathfrak a\ne P$,
$\mathfrak a$ est contenu dans un idéal maximal $\mathfrak n$ de $P$ et comme
$1\notin\mathfrak n$, $\mathfrak n\cap M$ contient $\mathfrak m(M)$ et est
distinct de $M$, donc $\mathfrak n\cap M=\mathfrak m(M)$ et de même
$\mathfrak n\cap N=\mathfrak m(N)$. Il est clair que si $Q$ domine $M$ et
$N$, on a $P\subset Q$, et
\[
  \mathfrak m(M)=\mathfrak m(Q)\cap M
  =(\mathfrak m(Q)\cap P)\cap M,\qquad
  \mathfrak m(N)=(\mathfrak m(Q)\cap P)\cap N,
\]
donc (iii) entraîne (i)~; la réciproque est évidente en prenant
$Q=P_\mathfrak p$.

Lorsque les conditions de (\hyperref[I.8.1.4-fr]{8.1.4}) sont satisfaites,
on dit, avec C.~Chevalley, que les anneaux locaux $M$ et $N$ sont
\emph{apparentés}.

\begin{proposition}[8.1.5]
\label{I.8.1.5-fr}
Soient $A$, $B$ deux sous-anneaux d'un corps $R$, $C$ le sous-anneau de $R$
engendré par $A\cup B$. Les conditions suivantes sont équivalentes~:
\begin{enumerate}
  \item[(i)] Pour tout anneau local $Q$ contenant $A$ et $B$, on a
    $A_\mathfrak p=B_\mathfrak q$, en posant
    $\mathfrak p=\mathfrak m(Q)\cap A$,
    $\mathfrak q=\mathfrak m(Q)\cap B$.
  \item[(ii)] Pour tout idéal premier $\mathfrak r$ de $C$, on a
    $A_\mathfrak p=B_\mathfrak q$, en posant
    $\mathfrak p=\mathfrak r\cap A$, $\mathfrak q=\mathfrak r\cap B$.
  \item[(iii)] Si $M\in L(A)$ et $N\in L(B)$ sont apparentés, ils sont
    identiques.
  \item[(iv)] On a $L(A)\cap L(B)=L(C)$.
\end{enumerate}
\end{proposition}

Les lemmes (\hyperref[I.8.1.3-fr]{8.1.3}) et
(\hyperref[I.8.1.4-fr]{8.1.4}) prouvent que (i) et (iii) sont équivalentes~;
il est clair que (i) entraîne (ii) en l'appliquant à $Q=C_\mathfrak r$~;
inversement, (ii) entraîne (i), car si $Q$ contient $A\cup B$, il contient
$C$, et si $\mathfrak r=\mathfrak m(Q)\cap C$, on a
$\mathfrak p=\mathfrak r\cap A$ et $\mathfrak q=\mathfrak r\cap B$ d'après
(\hyperref[I.8.1.3-fr]{8.1.3}). Il est immédiat que (iv) implique (i), car si
$Q$ contient $A\cup B$, il domine un anneau local
$C_\mathfrak r\in L(C)$ par (\hyperref[I.8.1.3-fr]{8.1.3})~; on a par
hypothèse $C_\mathfrak r\in L(A)\cap L(B)$, et
(\hyperref[I.8.1.1-fr]{8.1.1}) et (\hyperref[I.8.1.3-fr]{8.1.3}) prouvent que
$C_\mathfrak r=A_\mathfrak p=B_\mathfrak q$. Prouvons enfin que (iii)
entraîne (iv). Soit $Q\in L(C)$~; $Q$ domine un $M\in L(A)$ et un
$N\in L(B)$ (\hyperref[I.8.1.3-fr]{8.1.3}), donc $M$ et $N$, étant
apparentés, sont identiques par hypothèse. Comme on a alors $C\subset M$,
$M$ domine un $Q'\in L(C)$ (\hyperref[I.8.1.3-fr]{8.1.3}), donc $Q$ domine
$Q'$, ce qui (\hyperref[I.8.1.3-fr]{8.1.3}) entraîne nécessairement
$Q=Q'=M$, donc $Q\in L(A)\cap L(B)$. Inversement, si
$Q\in L(A)\cap L(B)$, on a $C\subset Q$, donc
(\hyperref[I.8.1.3-fr]{8.1.3}) $Q$ domine un
$Q''\in L(C)\subset L(A)\cap L(B)$~; $Q$ et $Q''$ étant apparentés sont
identiques, donc $Q''=Q\in L(C)$, ce qui achève la démonstration.

\subsection{Anneaux locaux d'un schéma intègre.}
\label{subsection:I.8.2-fr}

\begin{env}[8.2.1]
\label{I.8.2.1-fr}
Soient $X$ un préschéma \emph{intègre}, $R$ son corps des fonctions
rationnelles, identique à l'anneau local du point générique $a$ de $X$~; pour
tout $x\in X$, on sait que $\mathscr{O}_x$ s'identifie canoniquement à un
sous-anneau de $R$ (\hyperref[I.7.1.5-fr]{7.1.5}), et pour toute fonction
rationnelle $f\in R$, le domaine de définition $\delta(f)$ de $f$ est
l'ensemble ouvert des $x\in X$ tels que $f\in\mathscr{O}_x$. Il en résulte
(\hyperref[I.7.2.6-fr]{7.2.6}) que pour tout ouvert $U\subset X$, on a
\begin{equation}
  \Gamma(U,\mathscr{O}_X)=\bigcap_{x\in U}\mathscr{O}_x.
  \tag{8.2.1.1}\label{I.8.2.1.1-fr}
\end{equation}
\end{env}

\oldpage[I]{166}
\begin{proposition}[8.2.2]
\label{I.8.2.2-fr}
Soient $X$ un préschéma intègre, $R$ son corps des fonctions rationnelles.
Pour que $X$ soit un schéma, il faut et il suffit que la relation
«~$\mathscr{O}_x$ et $\mathscr{O}_y$ sont apparentés~»
(\hyperref[I.8.1.4-fr]{8.1.4}) entre points $x$, $y$ de $X$ implique $x=y$.
\end{proposition}

Supposons cette condition vérifiée, et montrons que $X$ est séparé. Soient
$U$ et $V$ deux ouverts affines distincts de $X$, $A$ et $B$ leurs anneaux,
identifiés à des sous-anneaux de $R$~; $U$ (resp. $V$) s'identifie donc
(\hyperref[I.8.1.2-fr]{8.1.2}) à $L(A)$ (resp. $L(B)$), et l'hypothèse
entraîne (\hyperref[I.8.1.5-fr]{8.1.5}) que si $C$ est le sous-anneau de $R$
engendré par $A\cup B$, $W=U\cap V$ s'identifie à
$L(A)\cap L(B)=L(C)$. En outre, on sait
(\cite{I-1}, p.~5-03, prop.~4 \emph{bis}) que tout sous-anneau $E$ de $R$ est
égal à l'intersection des anneaux locaux appartenant à $L(E)$~; $C$
s'identifie donc à l'intersection des anneaux $\mathscr{O}_z$ pour $z\in W$,
autrement dit (\hyperref[I.8.2.1.1-fr]{8.2.1.1}) à
$\Gamma(W,\mathscr{O}_X)$. Considérons alors le sous-préschéma induit par $X$
sur $W$~; à l'homomorphisme identique
$\varphi:C\to\Gamma(W,\mathscr{O}_X)$ correspond
(\hyperref[I.2.2.4-fr]{2.2.4}) un morphisme
$\Phi=(\psi,\theta):W\to\operatorname{Spec}(C)$~; nous allons voir que $\Phi$
est un \emph{isomorphisme} de préschémas, d'où résultera que $W$ est un ouvert
\emph{affine}. L'identification de $W$ à
$L(C)=\operatorname{Spec}(C)$ montre que $\psi$ est \emph{bijective}. D'autre
part, pour tout $x\in W$, $\theta_x^\sharp$ est l'injection
$C_\mathfrak r\to\mathscr{O}_x$, si
$\mathfrak r=\mathfrak m_x\cap C$, et par définition $C_\mathfrak r$ est
identifié à $\mathscr{O}_x$, donc $\theta_x^\sharp$ est bijective. Il reste
donc à voir que $\psi$ est un \emph{homéomorphisme}, autrement dit que pour
toute partie fermée $F\subset W$, $\psi(F)$ est fermé dans
$\operatorname{Spec}(C)$. Or, $F$ est la trace sur $W$ d'une partie fermée de
$U$, de la forme $V(\mathfrak a)$, où $\mathfrak a$ est un idéal de $A$~;
montrons que $\psi(F)=V(\mathfrak a C)$, ce qui prouvera notre assertion. En
effet, les idéaux premiers de $C$ contenant $\mathfrak a C$ sont les idéaux
premiers de $C$ contenant $\mathfrak a$, donc les idéaux de la forme
$\psi(x)=\mathfrak m_x\cap C$ où
$\mathfrak a\subset\mathfrak m_x$ et $x\in W$~; comme
$\mathfrak a\subset\mathfrak m_x$ équivaut à
$x\in V(\mathfrak a)=W\cap F$ pour $x\in U$, on a bien
$\psi(F)=V(\mathfrak a C)$.

Il s'ensuit que $X$ est séparé, car $U\cap V$ est affine et son anneau $C$ est
engendré par la réunion $A\cup B$ des anneaux de $U$ et $V$
(\hyperref[I.5.5.6-fr]{5.5.6}).

Inversement, supposons $X$ séparé, et soient $x$, $y$ deux points de $X$ tels
que $\mathscr{O}_x$ et $\mathscr{O}_y$ soient apparentés. Soit $U$ (resp.
$V$) un ouvert affine contenant $x$ (resp. $y$), d'anneau $A$ (resp. $B$)~;
on sait alors que $U\cap V$ est affine et que son anneau $C$ est engendré par
$A\cup B$ (\hyperref[I.5.5.6-fr]{5.5.6}). Si
$\mathfrak p=\mathfrak m_x\cap A$, $\mathfrak q=\mathfrak m_y\cap B$, on a
$A_\mathfrak p=\mathscr{O}_x$, $B_\mathfrak q=\mathscr{O}_y$, et comme
$A_\mathfrak p$ et $B_\mathfrak q$ sont apparentés, il existe un idéal premier
$\mathfrak r$ de $C$ tel que $\mathfrak p=\mathfrak r\cap A$,
$\mathfrak q=\mathfrak r\cap B$ (\hyperref[I.8.1.4-fr]{8.1.4}). Mais alors il
existe un point $z\in U\cap V$ tel que
$\mathfrak r=\mathfrak m_z\cap C$ puisque $U\cap V$ est affine, et on a
évidemment $x=z$, et $y=z$, d'où $x=y$.

\begin{corollary}[8.2.3]
\label{I.8.2.3-fr}
Soient $X$ un schéma intègre, $x$, $y$ deux points de $X$. Pour que
$x\in\overline{\{y\}}$, il faut et il suffit que
$\mathscr{O}_x\subset\mathscr{O}_y$, autrement dit que toute fonction
rationnelle définie en $x$ soit définie en $y$.
\end{corollary}

La condition est évidemment nécessaire puisque le domaine de définition
$\delta(f)$ d'une fonction rationnelle $f\in R$ est ouvert~; montrons qu'elle
est suffisante. Si $\mathscr{O}_x\subset\mathscr{O}_y$, il existe un idéal
premier $\mathfrak p$ de $\mathscr{O}_x$ tel que $\mathscr{O}_y$ domine
$(\mathscr{O}_x)_\mathfrak p$ (\hyperref[I.8.1.3-fr]{8.1.3})~; or
(\hyperref[I.2.4.2-fr]{2.4.2}) il existe $z\in X$ tel que
$x\in\overline{\{z\}}$ et que
$\mathscr{O}_z=(\mathscr{O}_x)_\mathfrak p$~; comme $\mathscr{O}_z$ et
$\mathscr{O}_y$ sont apparentés, on a $z=y$ par
(\hyperref[I.8.2.2-fr]{8.2.2}), d'où le corollaire.

\begin{corollary}[8.2.4]
\label{I.8.2.4-fr}
Si $X$ est un schéma intègre, l'application $x\mapsto\mathscr{O}_x$ est
injective~; autrement dit, si $x$, $y$ sont deux points distincts de $X$, il
existe une fonction rationnelle définie en l'un de ces points et non en
l'autre.
\end{corollary}

\oldpage[I]{167}
Cela résulte de (\hyperref[I.8.2.3-fr]{8.2.3}) et de l'axiome
$(T_0)$ (\hyperref[I.2.1.4-fr]{2.1.4}).

\begin{corollary}[8.2.5]
\label{I.8.2.5-fr}
Soit $X$ un schéma intègre dont l'espace sous-jacent est noethérien~;
lorsque $f$ parcourt le corps $R$ des fonctions rationnelles sur $X$, les
ensembles $\delta(f)$ engendrent la topologie de $X$.
\end{corollary}

En effet, toute partie fermée de $X$ est alors réunion finie d'ensembles
fermés irréductibles, c'est-à-dire de la forme $\overline{\{y\}}$
(\hyperref[I.2.1.5-fr]{2.1.5}). Or, si
$x\notin\overline{\{y\}}$, il existe une fonction rationnelle $f$ définie en
$x$ et non en $y$ (\hyperref[I.8.2.3-fr]{8.2.3}), autrement dit, on a
$x\in\delta(f)$ et $\delta(f)$ ne rencontre pas $\overline{\{y\}}$. Le
complémentaire de $\overline{\{y\}}$ est par suite réunion d'ensembles de la
forme $\delta(f)$, et en vertu de la première remarque, tout ouvert de $X$ est
réunion d'intersections finies d'ouverts de la forme $\delta(f)$.

\begin{env}[8.2.6]
\label{I.8.2.6-fr}
Le cor. (\hyperref[I.8.2.5-fr]{8.2.5}) montre que la topologie de $X$ est
entièrement caractérisée par la donnée de la famille d'anneaux locaux
$(\mathscr{O}_x)_{x\in X}$ ayant $R$ pour corps des fractions. Il revient au
même d'ailleurs de dire que les parties fermées de $X$ sont définies de la
façon suivante~: étant donnée une partie finie
$\{x_1,\ldots,x_n\}$ de $X$, on considère l'ensemble des $y\in X$ tels que
$\mathscr{O}_y\subset\mathscr{O}_{x_i}$ pour un indice $i$ au moins, et ces
ensembles (pour tous les choix de $\{x_1,\ldots,x_n\}$) sont les ensembles
fermés de $X$. En outre, une fois connue la topologie de $X$, le faisceau
structural $\mathscr{O}_X$ est aussi bien déterminé par la famille des
$\mathscr{O}_x$, puisque
$\Gamma(U,\mathscr{O}_X)=\bigcap_{x\in U}\mathscr{O}_x$
(\hyperref[I.8.2.1.1-fr]{8.2.1.1}). La famille
$(\mathscr{O}_x)_{x\in X}$ détermine donc complètement le préschéma $X$
lorsque $X$ est un schéma intègre dont l'espace sous-jacent est noethérien.
\end{env}

\begin{proposition}[8.2.7]
\label{I.8.2.7-fr}
Soient $X$, $Y$ deux schémas intègres, $f:X\to Y$ un morphisme dominant
(\hyperref[I.2.2.6-fr]{2.2.6}), $K$ (resp. $L$) le corps des fonctions
rationnelles sur $X$ (resp. $Y$). Alors $L$ s'identifie à un sous-corps de
$K$, et pour tout $x\in X$, $\mathscr{O}_{f(x)}$ est l'unique anneau local de
$Y$ dominé par $\mathscr{O}_x$.
\end{proposition}

En effet, si $f=(\psi,\theta)$ et si $a$ est le point générique de $X$,
$\psi(a)$ est le point générique de $Y$
(\hyperref[0.2.1.5-fr]{0, 2.1.5})~; $\theta_a^\sharp$ est par suite un
monomorphisme du corps $L=\mathscr{O}_{\psi(a)}$ dans le corps
$K=\mathscr{O}_a$. Comme tout ouvert affine non vide $U$ de $Y$ contient
$\psi(a)$, il résulte de (\hyperref[I.2.2.4-fr]{2.2.4}) que l'homomorphisme
$\Gamma(U,\mathscr{O}_Y)\to\Gamma(\psi^{-1}(U),\mathscr{O}_X)$ correspondant
à $f$ est la restriction de $\theta_a^\sharp$ à
$\Gamma(U,\mathscr{O}_Y)$. Donc, pour tout $x\in X$,
$\theta_x^\sharp$ est la restriction à $\mathscr{O}_{\psi(x)}$ de
$\theta_a^\sharp$, et est par suite un monomorphisme. On sait en outre que
$\theta_x^\sharp$ est un homomorphisme local, donc, si on identifie $L$ à un
sous-corps de $K$ par $\theta_a^\sharp$, $\mathscr{O}_{\psi(x)}$ est dominé
par $\mathscr{O}_x$ (\hyperref[I.8.1.1-fr]{8.1.1})~; c'est d'ailleurs le seul
anneau local de $Y$ dominé par $\mathscr{O}_x$ puisque deux anneaux locaux de
$Y$ qui sont apparentés sont identiques
(\hyperref[I.8.2.2-fr]{8.2.2}).

\begin{proposition}[8.2.8]
\label{I.8.2.8-fr}
Soient $X$ un préschéma irréductible, $f:X\to Y$ une immersion locale
(resp. un isomorphisme local)~; on suppose en outre le morphisme $f$ séparé.
Alors $f$ est une immersion (resp. une immersion ouverte).
\end{proposition}

Soit $f=(\psi,\theta)$~; il suffit, dans les deux cas, de prouver que $\psi$
est un \emph{homéomorphisme} de $X$ sur $\psi(X)$
(\hyperref[I.4.5.3-fr]{4.5.3}). Remplaçant $f$ par $f_{\mathrm{red}}$
(\hyperref[I.5.1.6-fr]{5.1.6} et
\hyperref[I.5.5.1-fr]{5.5.1, (vi)}), on peut supposer $X$ et $Y$ réduits. Si
$Y'$ est le sous-préschéma fermé réduit de $Y$ ayant pour espace sous-jacent
$\overline{\psi(X)}$, $f$ se factorise en
$X\xrightarrow{f'}Y'\xrightarrow{j}Y$, où $j$ est l'injection canonique
(\hyperref[I.5.2.2-fr]{5.2.2}). Il résulte de
(\hyperref[I.5.5.1-fr]{5.5.1, (v)}) que $f'$ est encore un morphisme séparé~;
en outre, $f'$ est encore

\oldpage[I]{168}
une immersion locale (resp. un isomorphisme local), car la question étant
locale sur $X$ et $Y$, on peut se borner pour le démontrer au cas où $f$ est
une immersion fermée (resp. une immersion ouverte) et notre assertion découle
alors aussitôt de (\hyperref[I.4.2.2-fr]{4.2.2}).

On peut donc supposer que $f$ est un morphisme \emph{dominant}, ce qui
entraîne que $Y$ est, lui aussi, irréductible
(\hyperref[0.2.1.5-fr]{0, 2.1.5}), donc que $X$ et $Y$ sont tous deux
\emph{intègres}. En outre, la question étant locale sur $Y$, on peut supposer
que $Y$ est un schéma affine~; comme $f$ est séparé, $X$ est un schéma
(\hyperref[I.5.5.1-fr]{5.5.1, (ii)}), et on est finalement dans les
hypothèses de (\hyperref[I.8.2.7-fr]{8.2.7}). Alors, pour tout $x\in X$,
$\theta_x^\sharp$ est injectif~; mais l'hypothèse que $f$ est une immersion
locale implique que $\theta_x^\sharp$ est surjectif
(\hyperref[I.4.2.2-fr]{4.2.2}), donc $\theta_x^\sharp$ est bijectif,
autrement dit (avec l'identification de
(\hyperref[I.8.2.7-fr]{8.2.7})) on a
$\mathscr{O}_{\psi(x)}=\mathscr{O}_x$. Cela implique par
(\hyperref[I.8.2.4-fr]{8.2.4}) que $\psi$ est une application
\emph{injective}, ce qui démontre déjà la proposition lorsque $f$ est un
isomorphisme local (\hyperref[I.4.5.3-fr]{4.5.3}). Lorsqu'on suppose seulement
que $f$ est une immersion locale, pour tout $x\in X$ il existe un voisinage
ouvert $U$ de $x$ dans $X$ et un voisinage ouvert $V$ de $\psi(x)$ dans $Y$
tels que la restriction de $\psi$ à $U$ soit un homéomorphisme de $U$ sur une
partie \emph{fermée} de $V$. Or, $U$ est dense dans $X$, donc $\psi(U)$ est
dense dans $Y$ et \emph{a fortiori} dans $V$, ce qui prouve que
$\psi(U)=V$~; comme $\psi$ est injective, $\psi^{-1}(V)=U$ et ceci achève de
montrer que $\psi$ est un homéomorphisme de $X$ sur $\psi(X)$.

\subsection{Les schémas de Chevalley.}
\label{subsection:I.8.3-fr}

\begin{env}[8.3.1]
\label{I.8.3.1-fr}
Soient $X$ un schéma intègre \emph{noethérien}, $R$ son corps des fonctions
rationnelles~; désignons par $X'$ l'ensemble des sous-anneaux locaux
$\mathscr{O}_x\subset R$, où $x$ parcourt $X$. L'ensemble $X'$ vérifie les
trois conditions suivantes~:
\begin{enumerate}
  \item[(Sch. 1)] \emph{Pour tout $M\in X'$, $R$ est le corps des fractions
    de $M$.}
  \item[(Sch. 2)] \emph{Il existe un ensemble fini de sous-anneaux
    noethériens $A_i$ de $R$ tels que $X'=\bigcup_i L(A_i)$ et que, pour tout
    couple d'indices $i$, $j$, le sous-anneau $A_{ij}$ de $R$ engendré par
    $A_i\cup A_j$ soit une algèbre de type fini sur $A_i$.}
  \item[(Sch. 3)] \emph{Deux éléments $M$, $N$ de $X'$ qui sont apparentés
    sont identiques.}
\end{enumerate}
\end{env}

On a en effet vu dans (\hyperref[I.8.2.1-fr]{8.2.1}) que (Sch. 1) est
satisfaite, et (Sch. 3) résulte de (\hyperref[I.8.2.2-fr]{8.2.2}). Pour
démontrer (Sch. 2), il suffit de recouvrir $X$ par un nombre fini d'ouverts
affines $U_i$, d'anneaux noethériens, et de prendre
$A_i=\Gamma(U_i,\mathscr{O}_X)$~; l'hypothèse que $X$ est un schéma entraîne
que $U_i\cap U_j$ est affine et que
$\Gamma(U_i\cap U_j,\mathscr{O}_X)=A_{ij}$
(\hyperref[I.5.5.6-fr]{5.5.6})~; en outre, comme l'espace $U_i$ est
noethérien, l'immersion $U_i\cap U_j\to U_i$ est de type fini
(\hyperref[I.6.3.5-fr]{6.3.5}), donc $A_{ij}$ est une $A_i$-algèbre de type
fini (\hyperref[I.6.3.3-fr]{6.3.3}).

\begin{env}[8.3.2]
\label{I.8.3.2-fr}
Les structures dont les axiomes sont (Sch. 1), (Sch. 2) et (Sch. 3)
généralisent les «~schémas~» au sens de C.~Chevalley, qui suppose en outre que
$R$ est une extension de type fini d'un corps $K$ et que les $A_i$ sont des
$K$-algèbres de type fini (ce qui rend inutile une partie de (Sch. 2))
\cite{I-1}. Inversement, si on a une telle structure sur un ensemble $X'$,
on peut lui associer un schéma intègre $X$ en utilisant les remarques de
(\hyperref[I.8.2.6-fr]{8.2.6})~: l'espace sous-jacent de $X$ est égal à $X'$
muni de la topologie définie dans (\hyperref[I.8.2.6-fr]{8.2.6}), et du
faisceau $\mathscr{O}_X$

\oldpage[I]{169}
tel que $\Gamma(U,\mathscr{O}_X)=\bigcap_{x\in U}\mathscr{O}_x$ pour tout
ouvert $U\subset X$, avec une définition évidente des homomorphismes de
restriction. Nous laissons au lecteur le soin de vérifier qu'on obtient bien
ainsi un schéma intègre, dont les anneaux locaux sont les éléments de $X'$~;
nous n'utiliserons pas ce résultat par la suite.
\end{env}

\section{Compléments sur les faisceaux quasi-cohérents}
\label{section:I.9-fr}

\subsection{Produit tensoriel de faisceaux quasi-cohérents.}
\label{subsection:I.9.1-fr}

\begin{proposition}[9.1.1]
\label{I.9.1.1-fr}
Soit $X$ un préschéma (resp. un préschéma localement noethérien). Soient
$\mathscr{F}$ et $\mathscr{G}$ deux $\mathscr{O}_X$-Modules quasi-cohérents
(resp. cohérents)~; alors
$\mathscr{F}\otimes_{\mathscr{O}_X}\mathscr{G}$ est quasi-cohérent (resp.
cohérent) et de type fini si $\mathscr{F}$ et $\mathscr{G}$ sont de type
fini. Si $\mathscr{F}$ admet une présentation finie et si $\mathscr{G}$ est
quasi-cohérent (resp. cohérent), alors
$\mathscr{H}\!om_{\mathscr{O}_X}(\mathscr{F},\mathscr{G})$ est quasi-cohérent
(resp. cohérent).
\end{proposition}

La question étant locale, on peut supposer $X$ affine (resp. affine
noethérien)~; en outre, si $\mathscr{F}$ est cohérent, on peut supposer qu'il
est le conoyau d'un homomorphisme
$\mathscr{O}_X^m\to\mathscr{O}_X^n$. Les assertions relatives aux faisceaux
quasi-cohérents résultent alors de
(\hyperref[I.1.3.12-fr]{1.3.12}) et
(\hyperref[I.1.3.9-fr]{1.3.9})~; les assertions relatives aux faisceaux
cohérents résultent de (\hyperref[I.1.5.1-fr]{1.5.1}) et du fait que, si $M$
et $N$ sont des modules de type fini sur un anneau noethérien $A$,
$M\otimes_A N$ et $\operatorname{Hom}_A(M,N)$ sont des $A$-modules de type
fini.

\begin{definition}[9.1.2]
\label{I.9.1.2-fr}
Soient $X$, $Y$ deux $S$-préschémas, $p$, $q$ les projections de
$X\times_S Y$, $\mathscr{F}$ (resp. $\mathscr{G}$) un
$\mathscr{O}_X$-Module (resp. un $\mathscr{O}_Y$-Module) quasi-cohérent. On
appelle produit tensoriel de $\mathscr{F}$ et $\mathscr{G}$ sur
$\mathscr{O}_S$ (ou sur $S$) et on note
$\mathscr{F}\otimes_{\mathscr{O}_S}\mathscr{G}$ (ou
$\mathscr{F}\otimes_S\mathscr{G}$) le produit tensoriel
$p^*(\mathscr{F})\otimes_{\mathscr{O}_{X\times_S Y}}q^*(\mathscr{G})$ sur le
préschéma $X\times_S Y$.
\end{definition}

Si $X_i$ ($1\leq i\leq n$) sont des $S$-préschémas, $\mathscr{F}_i$ un
$\mathscr{O}_{X_i}$-Module quasi-cohérent ($1\leq i\leq n$), on définit de
la même manière le produit tensoriel
$\mathscr{F}_1\otimes_S\mathscr{F}_2\otimes_S\cdots\otimes_S\mathscr{F}_n$
sur le préschéma
$Z=X_1\times_S X_2\times_S\cdots\times_S X_n$~; c'est un
$\mathscr{O}_Z$-module \emph{quasi-cohérent} en vertu de
(\hyperref[I.9.1.1-fr]{9.1.1}) et de
(\hyperref[0.5.1.4-fr]{0, 5.1.4})~; il est \emph{cohérent} si les
$\mathscr{F}_i$ le sont et si $Z$ est \emph{localement noethérien} en vertu
de (\hyperref[I.9.1.1-fr]{9.1.1}),
(\hyperref[0.5.3.11-fr]{0, 5.3.11}) et
(\hyperref[I.6.1.1-fr]{6.1.1}).

On notera que si on prend $X=Y=S$, la définition
(\hyperref[I.9.1.2-fr]{9.1.2}) redonne le produit tensoriel de
$\mathscr{O}_S$-Modules. En outre, comme
$q^*(\mathscr{O}_Y)=\mathscr{O}_{X\times_S Y}$
(\hyperref[0.4.3.4-fr]{0, 4.3.4}), le produit
$\mathscr{F}\otimes_S\mathscr{O}_Y$ s'identifie canoniquement à
$p^*(\mathscr{F})$, et de même
$\mathscr{O}_X\otimes_S\mathscr{G}$ s'identifie canoniquement à
$q^*(\mathscr{G})$. Plus particulièrement, si on prend $Y=S$ et qu'on
désigne par $f$ le morphisme structural $X\to Y$, on a
$\mathscr{O}_X\otimes_Y\mathscr{G}=f^*(\mathscr{G})$~: le produit tensoriel
ordinaire et l'image réciproque apparaissent donc comme des cas particuliers
du produit tensoriel général.

La déf. (\hyperref[I.9.1.2-fr]{9.1.2}) entraîne immédiatement que, pour $X$
et $Y$ fixés, $\mathscr{F}\otimes_S\mathscr{G}$ est un
\emph{bifoncteur covariant additif et exact à droite} en $\mathscr{F}$ et
$\mathscr{G}$.

\begin{proposition}[9.1.3]
\label{I.9.1.3-fr}
Soient $S$, $X$, $Y$ trois schémas affines d'anneaux respectifs, $A$, $B$,
$C$, $B$ et $C$ étant des $A$-algèbres. Soit $M$ (resp. $N$) un $B$-module
(resp. $C$-module), $\mathscr{F}=\widetilde{M}$ (resp.
$\mathscr{G}=\widetilde{N}$) le faisceau quasi-cohérent associé~; alors
$\mathscr{F}\otimes_S\mathscr{G}$ est canoniquement isomorphe au faisceau
associé au $(B\otimes_A C)$-module $M\otimes_A N$.
\end{proposition}

\begin{proof}
\oldpage[I]{170}
En effet, en vertu de (\hyperref[I.1.6.5-fr]{1.6.5}),
$\mathscr{F}\otimes_S\mathscr{G}$ est canoniquement isomorphe au faisceau
associé au $(B\otimes_A C)$-module
\[
  \bigl(M\otimes_B(B\otimes_A C)\bigr)
  \otimes_{B\otimes_A C}
  \bigl((B\otimes_A C)\otimes_C N\bigr)
\]
et en raison des isomorphismes canoniques entre produits tensoriels, ce
dernier est isomorphe à
\[
  M\otimes_B(B\otimes_A C)\otimes_C N
  =(M\otimes_B B)\otimes_A(C\otimes_C N)=M\otimes_A N.
\]
\end{proof}

\begin{proposition}[9.1.4]
\label{I.9.1.4-fr}
Soient $f:T\to X$, $g:T\to Y$ deux $S$-morphismes, $\mathscr{F}$ (resp.
$\mathscr{G}$) un $\mathscr{O}_X$-Module (resp. $\mathscr{O}_Y$-Module)
quasi-cohérent. On a alors
\[
  (f,g)_S^*(\mathscr{F}\otimes_S\mathscr{G})
  =f^*(\mathscr{F})\otimes_{\mathscr{O}_T}g^*(\mathscr{G}).
\]
\end{proposition}

\begin{proof}
Si $p$, $q$ sont les projections de $X\times_S Y$, la formule résulte en
effet des relations $(f,g)_S^*\circ p^*=f^*$ et
$(f,g)_S^*\circ q^*=g^*$ (\hyperref[0.3.5.5-fr]{0, 3.5.5}), et du fait
qu'une image réciproque d'un produit tensoriel de faisceaux algébriques est
le produit tensoriel de leurs images réciproques
(\hyperref[0.4.3.3-fr]{0, 4.3.3}).
\end{proof}

\begin{corollary}[9.1.5]
\label{I.9.1.5-fr}
Soient $f:X\to X'$, $g:Y\to Y'$ deux $S$-morphismes, $\mathscr{F}'$ (resp.
$\mathscr{G}'$) un $\mathscr{O}_{X'}$-Module (resp.
$\mathscr{O}_{Y'}$-Module) quasi-cohérent. On a alors
\[
  (f\times_S g)^*(\mathscr{F}'\otimes_S\mathscr{G}')
  =f^*(\mathscr{F}')\otimes_S g^*(\mathscr{G}').
\]
\end{corollary}

\begin{proof}
Cela résulte de (\hyperref[I.9.1.4-fr]{9.1.4}) et du fait que
$f\times_S g=(f\circ p,g\circ q)_S$, $p$ et $q$ étant les projections de
$X\times_S Y$.
\end{proof}

\begin{corollary}[9.1.6]
\label{I.9.1.6-fr}
Soient $X$, $Y$, $Z$ trois $S$-préschémas, $\mathscr{F}$ (resp.
$\mathscr{G}$, $\mathscr{H}$) un $\mathscr{O}_X$-Module (resp.
$\mathscr{O}_Y$-Module, $\mathscr{O}_Z$-Module) quasi-cohérent~; le faisceau
$\mathscr{F}\otimes_S\mathscr{G}\otimes_S\mathscr{H}$ est l'image
réciproque de
$(\mathscr{F}\otimes_S\mathscr{G})\otimes_S\mathscr{H}$ par l'isomorphisme
canonique de $X\times_S Y\times_S Z$ sur
$(X\times_S Y)\times_S Z$.
\end{corollary}

\begin{proof}
En effet, cet isomorphisme s'écrit $(p_1,p_2)_S\times_S p_3$, en désignant
par $p_1$, $p_2$, $p_3$ les projections de $X\times_S Y\times_S Z$.

De même, l'image réciproque de $\mathscr{G}\otimes_S\mathscr{F}$ par
l'isomorphisme canonique de $X\times_S Y$ sur $Y\times_S X$ est
$\mathscr{F}\otimes_S\mathscr{G}$.
\end{proof}

\begin{corollary}[9.1.7]
\label{I.9.1.7-fr}
Si $X$ est un $S$-préschéma, tout $\mathscr{O}_X$-Module quasi-cohérent
$\mathscr{F}$ est l'image réciproque de
$\mathscr{F}\otimes_S\mathscr{O}_S$ par l'isomorphisme canonique de $X$ sur
$X\times_S S$ (\hyperref[I.3.3.3-fr]{3.3.3}).
\end{corollary}

\begin{proof}
En effet, cet isomorphisme est $(1_X,\varphi)_S$, où $\varphi$ est le
morphisme structural $X\to S$, et le corollaire résulte de
(\hyperref[I.9.1.4-fr]{9.1.4}) et du fait que
$\varphi^*(\mathscr{O}_S)=\mathscr{O}_X$.
\end{proof}

\begin{env}[9.1.8]
\label{I.9.1.8-fr}
Soient $X$ un $S$-préschéma, $\mathscr{F}$ un $\mathscr{O}_X$-Module
quasi-cohérent, $\varphi:S'\to S$ un morphisme~; on désigne par
$\mathscr{F}_{(\varphi)}$ ou $\mathscr{F}_{(S')}$ le faisceau
quasi-cohérent $\mathscr{F}\otimes_S\mathscr{O}_{S'}$, sur
$X\times_S S'=X_{(\varphi)}=X_{(S')}$~; donc
$\mathscr{F}_{(S')}=p^*(\mathscr{F})$, où $p$ est la projection
$X_{(S')}\to X$.
\end{env}

\begin{proposition}[9.1.9]
\label{I.9.1.9-fr}
Soit $\varphi':S''\to S'$ un morphisme. Pour tout
$\mathscr{O}_X$-Module quasi-cohérent $\mathscr{F}$ sur le $S$-préschéma
$X$, $(\mathscr{F}_{(\varphi)})_{(\varphi')}$ est l'image réciproque de
$\mathscr{F}_{(\varphi\circ\varphi')}$ par l'isomorphisme canonique
$(X_{(\varphi)})_{(\varphi')}\xrightarrow{\sim}
X_{(\varphi\circ\varphi')}$ (\hyperref[I.3.3.9-fr]{3.3.9}).
\end{proposition}

\begin{proof}
Cela résulte aussitôt des définitions et de
(\hyperref[I.3.3.9-fr]{3.3.9}), et s'écrit encore
\[
  (\mathscr{F}\otimes_S\mathscr{O}_{S'})
  \otimes_{S'}\mathscr{O}_{S''}
  =\mathscr{F}\otimes_S\mathscr{O}_{S''}.
  \tag{9.1.9.1}\label{I.9.1.9.1-fr}
\]
\end{proof}

\begin{proposition}[9.1.10]
\label{I.9.1.10-fr}
Soient $Y$ un $S$-préschéma, $f:X\to Y$ un $S$-morphisme. Pour tout
$\mathscr{O}_Y$-Module quasi-cohérent $\mathscr{G}$ et tout morphisme
$S'\to S$, on a
\[
  (f_{(S')})^*(\mathscr{G}_{(S')})
  =(f^*(\mathscr{G}))_{(S')}.
\]
\end{proposition}

\begin{proof}
Cela résulte aussitôt de la commutativité du diagramme
\oldpage[I]{171}
\[
  \xymatrix{
    X_{(S')}\ar[r]^{f_{(S')}}\ar[d] &
    Y_{(S')}\ar[d]\\
    X\ar[r]_f &
    Y
  }
\]
\end{proof}

\begin{corollary}[9.1.11]
\label{I.9.1.11-fr}
Soient $X$, $Y$ deux $S$-préschémas, $\mathscr{F}$ (resp. $\mathscr{G}$) un
$\mathscr{O}_X$-Module (resp. $\mathscr{O}_Y$-Module) quasi-cohérent.
L'image réciproque du faisceau
$(\mathscr{F}_{(S')}\otimes_{S'}\mathscr{G}_{(S')})$ par l'isomorphisme
canonique
$(X\times_S Y)_{(S')}\xrightarrow{\sim}
(X_{(S')})\times_{S'}(Y_{(S')})$
(\hyperref[I.3.3.10-fr]{3.3.10}) est égale à
$(\mathscr{F}\otimes_S\mathscr{G})_{(S')}$.
\end{corollary}

\begin{proof}
Si $p$, $q$ sont les projections de $X\times_S Y$, l'isomorphisme en
question n'est autre que $(p_{(S')},q_{(S')})_{S'}$~; le corollaire résulte
des prop. (\hyperref[I.9.1.4-fr]{9.1.4}) et
(\hyperref[I.9.1.10-fr]{9.1.10}).
\end{proof}

\begin{proposition}[9.1.12]
\label{I.9.1.12-fr}
Avec les notations de (\hyperref[I.9.1.2-fr]{9.1.2}) soient $z$ un point de
$X\times_S Y$, $x=p(z)$, $y=q(z)$~; la fibre
$(\mathscr{F}\otimes_S\mathscr{G})_z$ est isomorphe à
$(\mathscr{F}_x\otimes_{\mathscr{O}_x}\mathscr{O}_z)
\otimes_{\mathscr{O}_z}
(\mathscr{G}_y\otimes_{\mathscr{O}_y}\mathscr{O}_z)
=\mathscr{F}_x\otimes_{\mathscr{O}_x}\mathscr{O}_z
\otimes_{\mathscr{O}_y}\mathscr{G}_y$.
\end{proposition}

\begin{proof}
Comme on peut se ramener au cas affine, la proposition résulte de la formule
(\hyperref[I.1.6.5.1-fr]{1.6.5.1}).
\end{proof}

\begin{corollary}[9.1.13]
\label{I.9.1.13-fr}
Si $\mathscr{F}$ et $\mathscr{G}$ sont de type fini, on a
\[
  \operatorname{Supp}(\mathscr{F}\otimes_S\mathscr{G})
  =p^{-1}(\operatorname{Supp}(\mathscr{F}))
  \cap q^{-1}(\operatorname{Supp}(\mathscr{G})).
\]
\end{corollary}

\begin{proof}
Comme $p^*(\mathscr{F})$ et $q^*(\mathscr{G})$ sont de type fini sur
$\mathscr{O}_{X\times_S Y}$, on est ramené, par
(\hyperref[I.9.1.12-fr]{9.1.12}) et
(\hyperref[0.1.7.5-fr]{0, 1.7.5}), au cas particulier où
$\mathscr{G}=\mathscr{O}_Y$, c'est-à-dire à démontrer la formule
\[
  \operatorname{Supp}(p^{-1}(\mathscr{F}))
  =p^{-1}(\operatorname{Supp}(\mathscr{F})).
  \tag{9.1.13.1}\label{I.9.1.13.1-fr}
\]

Le même raisonnement que dans (\hyperref[0.1.7.5-fr]{0, 1.7.5}) ramène à
vérifier que l'on a, pour tout $z\in X\times_S Y$,
$\mathscr{O}_z/\mathfrak{m}_x\mathscr{O}_z\neq0$ (avec $x=p(z)$), ce qui
découle du fait que l'homomorphisme
$\mathscr{O}_x\to\mathscr{O}_z$ est \emph{local} par hypothèse.
\end{proof}

Nous laissons au lecteur le soin d'étendre à un produit d'un nombre
quelconque de facteurs les résultats démontrés dans ce numéro pour deux
facteurs.

\subsection{Image directe d'un faisceau quasi-cohérent.}
\label{subsection:I.9.2-fr}

\begin{proposition}[9.2.1]
\label{I.9.2.1-fr}
Soit $f:X\to Y$ un morphisme de préschémas. On suppose qu'il existe un
recouvrement $(Y_\alpha)$ de $Y$ par des ouverts affines ayant la propriété
suivante~: chacun des $f^{-1}(Y_\alpha)$ admet un recouvrement fini
$(X_{\alpha i})$ par des ouverts affines contenus dans
$f^{-1}(Y_\alpha)$, tel que chacune des intersections
$X_{\alpha i}\cap X_{\alpha j}$ soit elle-même réunion finie d'ouverts
affines. Dans ces conditions, pour tout $\mathscr{O}_X$-Module
quasi-cohérent $\mathscr{F}$, $f_*(\mathscr{F})$ est un
$\mathscr{O}_Y$-Module quasi-cohérent.
\end{proposition}

\begin{proof}
La question étant locale sur $Y$, on peut supposer $Y$ égal à l'un des
$Y_\alpha$, donc supprimer les indices $\alpha$.

\begin{enumerate}
\item[a)] Supposons d'abord que les $X_i\cap X_j$ soient eux-mêmes des
ouverts \emph{affines}. Posons
$\mathscr{F}_i=\mathscr{F}|X_i$,
$\mathscr{F}_{ij}=\mathscr{F}|(X_i\cap X_j)$, et soient
$\mathscr{F}'_i$ et $\mathscr{F}'_{ij}$ les images de $\mathscr{F}_i$ et
$\mathscr{F}_{ij}$ respectivement par les restrictions de $f$ à $X_i$ et
$X_i\cap X_j$~; on sait que les $\mathscr{F}'_i$ et
$\mathscr{F}'_{ij}$ sont quasi-cohérents
(\hyperref[I.1.6.3-fr]{1.6.3}). Posons
$\mathscr{G}=\bigoplus_i\mathscr{F}'_i$,
$\mathscr{H}=\bigoplus_{i,j}\mathscr{F}'_{ij}$~; $\mathscr{G}$ et
$\mathscr{H}$ sont des $\mathscr{O}_Y$-Modules quasi-cohérents~; nous allons
définir un homomorphisme $u:\mathscr{G}\to\mathscr{H}$ tel que
$f_*(\mathscr{F})$ soit le \emph{noyau} de $u$~; il en résultera que
$f_*(\mathscr{F})$ est quasi-cohérent
(\hyperref[I.1.3.9-fr]{1.3.9}). Il suffit de définir $u$
\oldpage[I]{172}
comme homomorphisme de préfaisceaux~; tenant compte des définitions de
$\mathscr{G}$ et $\mathscr{H}$, il suffit donc, pour tout ensemble ouvert
$W\subset Y$, de définir un homomorphisme
\[
  u_W:\bigoplus_i\Gamma(f^{-1}(W)\cap X_i,\mathscr{F})
  \longrightarrow
  \bigoplus_{i,j}\Gamma(f^{-1}(W)\cap X_i\cap X_j,\mathscr{F})
\]
de façon à satisfaire aux conditions de compatibilité usuelles lorsque $W$
varie. Si, pour toute section
$s_i\in\Gamma(f^{-1}(W)\cap X_i,\mathscr{F})$, on désigne par $s_{i|j}$ sa
restriction à $f^{-1}(W)\cap X_i\cap X_j$, on posera
\[
  u_W((s_i))=(s_{i|j}-s_{j|i})
\]
et les conditions de compatibilité sont remplies de façon évidente. Pour
prouver que le noyau $\mathscr{R}$ de $u$ est $f_*(\mathscr{F})$, définissons
un homomorphisme de $f_*(\mathscr{F})$ dans $\mathscr{R}$ en faisant
correspondre à toute section $s\in\Gamma(f^{-1}(W),\mathscr{F})$ la famille
$(s_i)$, où $s_i$ est la restriction de $s$ à $f^{-1}(W)\cap X_i$~; les
axiomes (F 1) et (F 2) des faisceaux (G, II, 1.1) entraînent que cet
homomorphisme est \emph{bijectif}, ce qui termine la démonstration dans ce cas.

\item[b)] Dans le cas général, le même raisonnement s'applique une fois que
l'on a établi que les $\mathscr{F}'_{ij}$ sont quasi-cohérents. Or, par
hypothèse, $X_i\cap X_j$ est réunion finie d'ouverts affines $X_{ijk}$~; et
comme les $X_{ijk}$ sont des ouverts affines \emph{dans un schéma},
l'intersection de deux quelconques d'entre eux est encore un ouvert affine
(\hyperref[I.5.5.6-fr]{5.5.6}). On est donc ramené au premier cas, et
(9.2.1) est donc démontrée.
\end{enumerate}
\end{proof}

\begin{corollary}[9.2.2]
\label{I.9.2.2-fr}
La conclusion de (9.2.1) est valable dans chacun des cas suivants~:
\begin{enumerate}
\item[a)] $f$ est séparé et quasi-compact.
\item[b)] $f$ est séparé et de type fini.
\item[c)] $f$ est quasi-compact et l'espace sous-jacent à $X$ est localement
noethérien.
\end{enumerate}
\end{corollary}

\begin{proof}
Dans le cas a), les $X_{\alpha i}\cap X_{\alpha j}$ sont affines
(\hyperref[I.5.5.6-fr]{5.5.6}). Le cas b) est un cas particulier de a)
(\hyperref[I.6.6.3-fr]{6.6.3}). Enfin, dans le cas c), on peut se ramener au
cas où $Y$ est affine et l'espace sous-jacent à $X$ noethérien~; alors $X$
admet un recouvrement ouvert affine fini $(X_i)$, et les $X_i\cap X_j$, étant
quasi-compacts, sont réunions finies d'ouverts affines
(\hyperref[I.2.1.3-fr]{2.1.3}).
\end{proof}

\subsection{Prolongement des sections de faisceaux quasi-cohérents.}
\label{subsection:I.9.3-fr}

\begin{theorem}[9.3.1]
\label{I.9.3.1-fr}
Soit $X$ un préschéma dont l'espace sous-jacent est noethérien, ou un schéma
dont l'espace sous-jacent est quasi-compact. Soient $\mathscr{L}$ un
$\mathscr{O}_X$-Module inversible
(\hyperref[0.5.4.1-fr]{0, 5.4.1}), $f$ une section de $\mathscr{L}$ au-dessus
de $X$, $X_f$ l'ensemble ouvert des $x\in X$ où $f(x)\neq0$
(\hyperref[0.5.5.1-fr]{0, 5.5.1}), $\mathscr{F}$ un
$\mathscr{O}_X$-Module quasi-cohérent.
\begin{enumerate}
\item[(i)] Si $s\in\Gamma(X,\mathscr{F})$ est telle que $s|X_f=0$, il existe
un entier $n>0$ tel que $s\otimes f^{\otimes n}=0$.
\item[(ii)] Pour toute section $s\in\Gamma(X_f,\mathscr{F})$, il existe un
entier $n>0$ tel que $s\otimes f^{\otimes n}$ se prolonge en une section de
$\mathscr{F}\otimes\mathscr{L}^{\otimes n}$ au-dessus de $X$.
\end{enumerate}
\end{theorem}

\begin{proof}
\begin{enumerate}
\item[(i)] Comme l'espace sous-jacent à $X$ est quasi-compact, donc réunion
finie d'ouverts affines $U_i$ tels que $\mathscr{L}|U_i$ soit isomorphe à
$\mathscr{O}_X|U_i$, on est ramené au cas où $X$ est affine et
$\mathscr{L}=\mathscr{O}_X$. Dans ce cas, $f$ s'identifie à un élément de
$A(X)$ et on a $X_f=D(f)$~; $s$ s'identifie à un élément d'un
$A(X)$-module $M$ et $s|X_f$ à l'élément correspondant de $M_f$, et le
résultat est trivial, compte tenu de la définition d'un module de fractions.
\oldpage[I]{173}

\item[(ii)] De nouveau $X$ est réunion finie d'ouverts affines $U_i$
($1\leq i\leq r$) tels que $\mathscr{L}|U_i\simeq\mathscr{O}_X|U_i$, et pour
chaque $i$,
$(s\otimes f^{\otimes n})|(U_i\cap X_f)$ s'identifie par l'isomorphisme
précédent à
$(f|(U_i\cap X_f))^n(s|(U_i\cap X_f))$. On sait alors
(\hyperref[I.1.4.1-fr]{1.4.1}) qu'il existe un entier $n>0$ tel que pour
chaque $i$, $(s\otimes f^{\otimes n})|(U_i\cap X_f)$ se prolonge en une
section $s_i$ de $\mathscr{F}\otimes\mathscr{L}^{\otimes n}$ au-dessus de
$U_i$. Soit $s_{i|j}$ la restriction de $s_i$ à $U_i\cap U_j$~; on a par
définition $s_{i|j}-s_{j|i}=0$ dans $X_f\cap U_i\cap U_j$. Or, si $X$ est un
espace noethérien, $U_i\cap U_j$ est quasi-compact~; si $X$ est un schéma,
$U_i\cap U_j$ est un ouvert affine
(\hyperref[I.5.5.6-fr]{5.5.6}), donc encore quasi-compact. En vertu de (i),
il existe donc un entier $m$ (indépendant de $i$ et $j$) tel que
$(s_{i|j}-s_{j|i})\otimes f^{\otimes m}=0$. On en conclut aussitôt qu'il
existe une section $s'$ de
$\mathscr{F}\otimes\mathscr{L}^{\otimes(n+m)}$ au-dessus de $X$, induisant
$s_i\otimes f^{\otimes m}$ au-dessus de chaque $U_i$, et induisant par suite
$s\otimes f^{\otimes(n+m)}$ au-dessus de $X_f$.
\end{enumerate}
\end{proof}

Les corollaires qui suivent donnent une interprétation du th. (9.3.1) en un
langage plus algébrique~:

\begin{corollary}[9.3.2]
\label{I.9.3.2-fr}
Les hypothèses étant celles de (9.3.1), considérons l'anneau gradué
$A_* = \Gamma_*(\mathscr{L})$ et le $A_*$-module gradué
$M_* = \Gamma_*(\mathscr{L},\mathscr{F})$
(\hyperref[0.5.4.6-fr]{0, 5.4.6}). Si $f\in A_n$, où $n\in\mathbb{Z}$, alors
on a un isomorphisme canonique
\[
  \Gamma(X_f,\mathscr{F})\xrightarrow{\sim}((M_*)_f)_0
\]
(sous-groupe du module de fractions $(M_*)_f$ formé des éléments de degré
$0$).
\end{corollary}

\begin{corollary}[9.3.3]
\label{I.9.3.3-fr}
On suppose vérifiées les hypothèses de (9.3.1), et on suppose en outre que
$\mathscr{L}=\mathscr{O}_X$. Alors si on pose
$A=\Gamma(X,\mathscr{O}_X)$, $M=\Gamma(X,\mathscr{F})$, le $A_f$-module
$\Gamma(X_f,\mathscr{F})$ est canoniquement isomorphe à $M_f$.
\end{corollary}

\begin{proposition}[9.3.4]
\label{I.9.3.4-fr}
Soient $X$ un préschéma noethérien, $\mathscr{F}$ un
$\mathscr{O}_X$-Module cohérent, $\mathscr{J}$ un faisceau cohérent d'idéaux
dans $\mathscr{O}_X$ tels que le support de $\mathscr{F}$ soit contenu dans
celui de $\mathscr{O}_X/\mathscr{J}$. Alors il existe un entier $n>0$ tel que
$\mathscr{J}^n\mathscr{F}=0$.
\end{proposition}

\begin{proof}
Comme $X$ est réunion finie d'ouverts affines dont les anneaux sont
noethériens, on peut supposer $X$ affine d'anneau $A$ noethérien~; alors
$\mathscr{F}=\widetilde M$, où $M=\Gamma(X,\mathscr{F})$ est un $A$-module
de type fini, et $\mathscr{J}=\widetilde{\mathfrak{J}}$, où
$\mathfrak{J}=\Gamma(X,\mathscr{J})$ est un idéal de $A$
(\hyperref[I.1.4.1-fr]{1.4.1} et \hyperref[I.1.5.1-fr]{1.5.1}). Comme $A$
est noethérien, $\mathfrak{J}$ admet un système fini de générateurs $f_i$
($1\leq i\leq m$). Par hypothèse, toute section de $\mathscr{F}$ au-dessus
de $X$ est nulle dans chacun des $D(f_i)$~; si $s_j$ ($1\leq j\leq q$) sont
des sections de $\mathscr{F}$ engendrant $M$, il existe donc un entier $h$
indépendant de $i$ et $j$, tel que $f_i^h s_j=0$
(\hyperref[I.1.4.1-fr]{1.4.1}), donc $f_i^h s=0$ pour tout $s\in M$. On en
conclut que si $n=mh$, on a $\mathfrak{J}^nM=0$, et par suite le
$\mathscr{O}_X$-Module correspondant
$\mathscr{J}^n\mathscr{F}=\widetilde{\mathfrak{J}^nM}$
(\hyperref[I.1.3.13-fr]{1.3.13}) est nul.
\end{proof}

\begin{corollary}[9.3.5]
\label{I.9.3.5-fr}
Sous les hypothèses de (9.3.4), il existe un sous-préschéma fermé $Y$ de
$X$, dont l'espace sous-jacent est le support de
$\mathscr{O}_X/\mathscr{J}$, et tel que, si $j:Y\to X$ est l'injection
canonique, on ait $\mathscr{F}=j_*(j^*(\mathscr{F}))$.
\end{corollary}

\begin{proof}
Notons d'abord que les supports de $\mathscr{O}_X/\mathscr{J}$ et de
$\mathscr{O}_X/\mathscr{J}^n$ sont les mêmes car si
$\mathscr{J}_x=\mathscr{O}_x$, on a aussi
$\mathscr{J}_x^n=\mathscr{O}_x$, et on a d'autre part
$\mathscr{J}_x^n\subset\mathscr{J}_x$ pour tout $x\in X$. On peut donc en
vertu de (9.3.4) supposer que $\mathscr{J}\mathscr{F}=0$~; on prend alors
pour $Y$ le sous-préschéma fermé de $X$ défini par $\mathscr{J}$, et comme
$\mathscr{F}$ est alors un $(\mathscr{O}_X/\mathscr{J})$-Module, la
conclusion est immédiate.
\end{proof}

\oldpage[I]{174}

\subsection{Prolongement des faisceaux quasi-cohérents.}
\label{subsection:I.9.4-fr}

\begin{env}[9.4.1]
\label{I.9.4.1-fr}
Soient $X$ un espace topologique, $\mathscr{F}$ un faisceau d'ensembles
(resp. de groupes, d'anneaux) sur $X$, $U$ une partie ouverte de $X$,
$\psi:U\to X$ l'injection canonique, $\mathscr{G}$ un sous-faisceau de
$\mathscr{F}|U=\psi^*(\mathscr{F})$. Comme $\psi_*$ est exact à gauche,
$\psi_*(\mathscr{G})$ est un sous-faisceau de
$\psi_*(\psi^*(\mathscr{F}))$~; si on considère l'homomorphisme canonique
$\rho:\mathscr{F}\to\psi_*(\psi^*(\mathscr{F}))$
(\hyperref[0.3.5.3-fr]{0, 3.5.3}), nous désignerons par
$\overline{\mathscr{G}}$ le sous-faisceau
$\rho^{-1}(\psi_*(\mathscr{G}))$ de $\mathscr{F}$. Il résulte immédiatement
des définitions que pour tout ouvert $V$ de $X$,
$\Gamma(V,\overline{\mathscr{G}})$ est formé des sections
$s\in\Gamma(V,\mathscr{F})$ dont la restriction à $V\cap U$ soit une section
de $\mathscr{G}$ au-dessus de $V\cap U$. On a donc
$\overline{\mathscr{G}}|U=\psi^*(\overline{\mathscr{G}})=\mathscr{G}$, et
$\overline{\mathscr{G}}$ est le \emph{plus grand} sous-faisceau de
$\mathscr{F}$ induisant $\mathscr{G}$ sur $U$~; nous dirons que
$\overline{\mathscr{G}}$ est le \emph{prolongement canonique} du
sous-faisceau $\mathscr{G}$ de $\mathscr{F}|U$ en un sous-faisceau de
$\mathscr{F}$.
\end{env}

\begin{proposition}[9.4.2]
\label{I.9.4.2-fr}
Soient $X$ un préschéma, $U$ une partie ouverte de $X$ telle que l'injection
canonique $j:U\to X$ soit un morphisme quasi-compact
\emph{(ce qui sera vérifié pour \emph{tout} $U$ si l'espace sous-jacent $X$
est \emph{localement noethérien}
(\hyperref[I.6.6.4-fr]{6.6.4, (i)}))}. Alors~:
\begin{enumerate}
\item[(i)] Pour tout $(\mathscr{O}_X|U)$-Module quasi-cohérent
  $\mathscr{G}$, $j_*(\mathscr{G})$ est un $\mathscr{O}_X$-Module
  quasi-cohérent et on a
  $j_*(\mathscr{G})|U=j^*(j_*(\mathscr{G}))=\mathscr{G}$.
\item[(ii)] Pour tout $\mathscr{O}_X$-Module quasi-cohérent $\mathscr{F}$ et
  tout sous-$(\mathscr{O}_X|U)$-Module quasi-cohérent $\mathscr{G}$, le
  prolongement canonique $\overline{\mathscr{G}}$ de $\mathscr{G}$ (9.4.1)
  est un sous-$\mathscr{O}_X$-Module quasi-cohérent de $\mathscr{F}$.
\end{enumerate}
\end{proposition}

\begin{proof}
Si $j=(\psi,\theta)$ ($\psi$ étant l'injection $U\to X$ des espaces
sous-jacents), on a par définition $j_*(\mathscr{G})=\psi_*(\mathscr{G})$
pour tout $(\mathscr{O}_X|U)$-Module $\mathscr{G}$, et en outre
$j^*(\mathscr{H})=\psi^*(\mathscr{H})=\mathscr{H}|U$ pour tout
$\mathscr{O}_X$-Module $\mathscr{H}$ en raison de la définition d'un
préschéma induit sur un ouvert. (i) est donc un cas particulier de
(\hyperref[I.9.2.2-fr]{9.2.2, a)})~; pour la même raison,
$j_*(j^*(\mathscr{F}))$ est quasi-cohérent, et comme
$\overline{\mathscr{G}}$ est l'image réciproque de $j_*(\mathscr{G})$ par
l'homomorphisme $\rho:\mathscr{F}\to j_*(j^*(\mathscr{F}))$, (ii) résulte de
(\hyperref[I.4.1.1-fr]{4.1.1}).
\end{proof}

On notera que l'hypothèse que le morphisme $j:U\to X$ est quasi-compact est
aussi vérifiée lorsque l'ouvert $U$ est \emph{quasi-compact} et $X$ un
\emph{schéma}~: en effet, $U$ est alors réunion finie d'ouverts affines
$U_i$, et pour tout ouvert affine $V$ de $X$, $V\cap U_i$ est un ouvert
affine (\hyperref[I.5.5.6-fr]{5.5.6}), donc quasi-compact.

\begin{corollary}[9.4.3]
\label{I.9.4.3-fr}
Soient $X$ un préschéma, $U$ un ouvert quasi-compact de $X$ tel que le
morphisme d'injection $j:U\to X$ soit quasi-compact. Supposons en outre que
tout $\mathscr{O}_X$-Module quasi-cohérent soit limite inductive de ses
sous-$\mathscr{O}_X$-Modules quasi-cohérents de type fini
\emph{(ce qui a lieu lorsque $X$ est un \emph{schéma affine})}. Soient alors
$\mathscr{F}$ un $\mathscr{O}_X$-Module quasi-cohérent, $\mathscr{G}$ un
sous-$(\mathscr{O}_X|U)$-Module quasi-cohérent de type fini de
$\mathscr{F}|U$. Il existe alors un sous-$\mathscr{O}_X$-Module
quasi-cohérent de type fini $\mathscr{G}'$ de $\mathscr{F}$ tel que
$\mathscr{G}'|U=\mathscr{G}$.
\end{corollary}

\begin{proof}
En effet, on a $\mathscr{G}=\overline{\mathscr{G}}|U$, et
$\overline{\mathscr{G}}$ est quasi-cohérent d'après (9.4.2), donc limite
inductive de ses sous-$\mathscr{O}_X$-Modules quasi-cohérents de type fini
$\mathscr{H}_\lambda$. Par suite $\mathscr{G}$ est limite inductive des
$\mathscr{H}_\lambda|U$, donc égal à un des $\mathscr{H}_\lambda|U$
puisqu'il est de type fini (\hyperref[0.5.2.3-fr]{0, 5.2.3}).
\end{proof}

\begin{remark}[9.4.4]
\label{I.9.4.4-fr}
Supposons que pour \emph{tout} ouvert affine $U\subset X$ le morphisme
d'injection $U\to X$ soit quasi-compact. Alors si la conclusion de (9.4.3) a
lieu pour tout ouvert affine $U$ et tout sous-$(\mathscr{O}_X|U)$-Module
quasi-cohérent de type fini $\mathscr{G}$ de $\mathscr{F}|U$,
\oldpage[I]{175}
il en résulte que $\mathscr{F}$ est limite inductive de ses
sous-$\mathscr{O}_X$-Modules quasi-cohérents de type fini. En effet, pour
tout ouvert affine $U\subset X$, on a $\mathscr{F}|U=\widetilde{M}$, où $M$
est un $A(U)$-module, et comme ce dernier est limite inductive de ses
sous-modules de type fini, $\mathscr{F}|U$ est limite inductive de ses
sous-$(\mathscr{O}_X|U)$-Modules quasi-cohérents de type fini
(\hyperref[I.1.3.9-fr]{1.3.9}). Or, par hypothèse, chacun de ces sous-Modules
est induit sur $U$ par un sous-$\mathscr{O}_X$-Module quasi-cohérent de type
fini $\mathscr{G}_{\lambda,U}$ de $\mathscr{F}$. Les sommes finies des
$\mathscr{G}_{\lambda,U}$ sont encore des $\mathscr{O}_X$-Modules
quasi-cohérents de type fini, car la question est locale et le cas où $X$
est affine a été traité dans (\hyperref[I.1.3.10-fr]{1.3.10})~; il est clair
alors que $\mathscr{F}$ est limite inductive de ces sommes finies, d'où notre
assertion.
\end{remark}

\begin{corollary}[9.4.5]
\label{I.9.4.5-fr}
Sous les hypothèses de (9.4.3), pour tout $(\mathscr{O}_X|U)$-Module
quasi-cohérent de type fini $\mathscr{G}$, il existe un
$\mathscr{O}_X$-Module quasi-cohérent de type fini $\mathscr{G}'$ tel que
$\mathscr{G}'|U=\mathscr{G}$.
\end{corollary}

\begin{proof}
Comme $\mathscr{F}=j_*(\mathscr{G})$ est quasi-cohérent (9.4.2) et
$\mathscr{F}|U=\mathscr{G}$, il suffit d'appliquer (9.4.3) à $\mathscr{F}$.
\end{proof}

\begin{lemma}[9.4.6]
\label{I.9.4.6-fr}
Soient $X$ un préschéma, $L$ un ensemble bien ordonné,
$(V_\lambda)_{\lambda\in L}$ un recouvrement de $X$ par des ouverts affines,
$U$ un ouvert de $X$~; pour tout $\lambda\in L$, on pose
$W_\lambda=\bigcup_{\mu<\lambda}V_\mu$. On suppose que~:
\begin{enumerate}
  \item[1\textsuperscript{o}] Pour tout $\lambda\in L$,
    $V_\lambda\cap W_\lambda$ soit quasi-compact~;
  \item[2\textsuperscript{o}] Le morphisme d'immersion $U\to X$ est
    quasi-compact.
\end{enumerate}
Alors, pour tout $\mathscr{O}_X$-Module quasi-cohérent $\mathscr{F}$ et tout
sous-$(\mathscr{O}_X|U)$-Module quasi-cohérent de type fini $\mathscr{G}$ de
$\mathscr{F}|U$, il existe un sous-$\mathscr{O}_X$-Module quasi-cohérent de
type fini $\mathscr{G}'$ de $\mathscr{F}$ tel que
$\mathscr{G}'|U=\mathscr{G}$.
\end{lemma}

\begin{proof}
Posons $U_\lambda=U\cup W_\lambda$~; on va définir par récurrence transfinie
une famille $(\mathscr{G}'_\lambda)$, où $\mathscr{G}'_\lambda$ est un
sous-$(\mathscr{O}_X|U_\lambda)$-Module quasi-cohérent de type fini de
$\mathscr{F}|U_\lambda$, tel que
$\mathscr{G}'_\lambda|U_\mu=\mathscr{G}'_\mu$ pour $\mu<\lambda$ et
$\mathscr{G}'_\lambda|U=\mathscr{G}$. L'unique sous-$\mathscr{O}_X$-Module
$\mathscr{G}'$ de $\mathscr{F}$ tel que
$\mathscr{G}'|U_\lambda=\mathscr{G}'_\lambda$ pour tout $\lambda\in L$
(\hyperref[0.3.3.1-fr]{0, 3.3.1}) répondra à la question. Supposons donc les
$\mathscr{G}'_\mu$ définis et ayant les propriétés précédentes pour
$\mu<\lambda$~; si $\lambda$ n'a pas de prédécesseur on prendra pour
$\mathscr{G}'_\lambda$ l'unique sous-$(\mathscr{O}_X|U_\lambda)$-Module de
$\mathscr{F}|U_\lambda$ tel que
$\mathscr{G}'_\lambda|U_\mu=\mathscr{G}'_\mu$ pour tout $\mu<\lambda$, ce
qui est licite puisque les $U_\mu$ avec $\mu<\lambda$ forment alors un
recouvrement de $U_\lambda$. Si au contraire $\lambda=\mu+1$, on a
$U_\lambda=U_\mu\cup V_\mu$, et il suffira de définir un
sous-$(\mathscr{O}_X|V_\mu)$-Module quasi-cohérent de type fini
$\mathscr{G}''_\mu$ de $\mathscr{F}|V_\mu$ tel que
\[
  \mathscr{G}''_\mu|(U_\mu\cap V_\mu)
  =\mathscr{G}'_\mu|(U_\mu\cap V_\mu)~;
\]
on prendra ensuite pour $\mathscr{G}'_\lambda$ le
sous-$(\mathscr{O}_X|U_\lambda)$-Module de $\mathscr{F}|U_\lambda$ tel que
$\mathscr{G}'_\lambda|U_\mu=\mathscr{G}'_\mu$ et
$\mathscr{G}'_\lambda|V_\mu=\mathscr{G}''_\mu$
(\hyperref[0.3.3.1-fr]{0, 3.3.1}). Or, comme $V_\mu$ est affine, l'existence
de $\mathscr{G}''_\mu$ est assurée par (9.4.3) dès que l'on a démontré que
$U_\mu\cap V_\mu$ est quasi-compact~; mais $U_\mu\cap V_\mu$ est réunion de
$U\cap V_\mu$ et de $W_\mu\cap V_\mu$, qui sont tous deux quasi-compacts en
vertu de l'hypothèse.
\end{proof}

\begin{theorem}[9.4.7]
\label{I.9.4.7-fr}
Soient $X$ un préschéma, $U$ un ouvert de $X$. On suppose vérifiée l'une des
conditions suivantes~:
\begin{enumerate}
  \item[(a)] L'espace sous-jacent à $X$ est localement noethérien.
  \item[(b)] $X$ est un schéma quasi-compact et $U$ un ouvert quasi-compact.
\end{enumerate}
Pour tout $\mathscr{O}_X$-Module quasi-cohérent $\mathscr{F}$ et tout
sous-$(\mathscr{O}_X|U)$-Module quasi-cohérent de type fini $\mathscr{G}$ de
$\mathscr{F}|U$, il existe alors un sous-$\mathscr{O}_X$-Module
quasi-cohérent de type fini $\mathscr{G}'$ de $\mathscr{F}$ tel que
$\mathscr{G}'|U=\mathscr{G}$.
\end{theorem}

\oldpage[I]{176}

\begin{proof}
Soit $(V_\lambda)_{\lambda\in L}$ un recouvrement de $X$ par des ouverts
affines, $L$ étant supposé fini dans le cas b)~; $L$ étant muni d'une
structure d'ensemble bien ordonné, il suffit de vérifier que les conditions
du lemme (9.4.6) sont satisfaites. C'est évident dans l'hypothèse a), les
espaces $V_\lambda$ étant noethériens. Dans l'hypothèse b), les
$V_\lambda\cap V_\mu$ sont affines (\hyperref[I.5.5.6-fr]{5.5.6}), donc
quasi-compacts, et comme $L$ est fini, $V_\lambda\cap W_\lambda$ est
quasi-compact. D'où le théorème.
\end{proof}

\begin{corollary}[9.4.8]
\label{I.9.4.8-fr}
Sous les hypothèses de (9.4.7), pour tout $(\mathscr{O}_X|U)$-Module
quasi-cohérent de type fini $\mathscr{G}$, il existe un
$\mathscr{O}_X$-Module quasi-cohérent de type fini $\mathscr{G}'$ tel que
$\mathscr{G}'|U=\mathscr{G}$.
\end{corollary}

\begin{proof}
Il suffit d'appliquer (9.4.7) à $\mathscr{F}=j_*(\mathscr{G})$, qui est
quasi-cohérent (9.4.2) et tel que $\mathscr{F}|U=\mathscr{G}$.
\end{proof}

\begin{corollary}[9.4.9]
\label{I.9.4.9-fr}
Soit $X$ un préschéma dont l'espace sous-jacent est localement noethérien, ou
un schéma quasi-compact. Alors tout $\mathscr{O}_X$-Module quasi-cohérent est
limite inductive de ses sous-$\mathscr{O}_X$-Modules quasi-cohérents de type
fini.
\end{corollary}

\begin{proof}
Cela résulte de (9.4.7) et de la remarque (9.4.4).
\end{proof}

\begin{corollary}[9.4.10]
\label{I.9.4.10-fr}
Sous les hypothèses de (9.4.9), si un $\mathscr{O}_X$-Module quasi-cohérent
$\mathscr{F}$ est tel que tout sous-$\mathscr{O}_X$-Module quasi-cohérent de
type fini de $\mathscr{F}$ soit engendré par ses sections au-dessus de $X$,
alors $\mathscr{F}$ est engendré par ses sections au-dessus de $X$.
\end{corollary}

\begin{proof}
En effet, soient $U$ un voisinage ouvert affine d'un point $x\in X$, et soit
$s$ une section de $\mathscr{F}$ au-dessus de $U$~; le
sous-$\mathscr{O}_X$-Module $\mathscr{G}$ de $\mathscr{F}|U$ engendré par
$s$ est quasi-cohérent et de type fini, donc il existe un
sous-$\mathscr{O}_X$-Module quasi-cohérent de type fini $\mathscr{G}'$ de
$\mathscr{F}$ tel que $\mathscr{G}'|U=\mathscr{G}$ (9.4.7). Par hypothèse,
il y a donc un nombre fini de sections $t_i$ de $\mathscr{G}'$ au-dessus de
$X$ et des sections $a_i$ de $\mathscr{O}_X$ au-dessus d'un voisinage
$V\subset U$ de $x$ telles que
$s|V=\sum_i a_i\mathbin{\cdot}(t_i|V)$, ce qui démontre le corollaire.
\end{proof}

\subsection{Image fermée d'un préschéma. Adhérence d'un sous-préschéma.}
\label{subsection:I.9.5-fr}

\begin{proposition}[9.5.1]
\label{I.9.5.1-fr}
Soit $f:X\to Y$ un morphisme de préschémas tel que $f_*(\mathscr{O}_X)$ soit
un $\mathscr{O}_Y$-Module quasi-cohérent (ce qui a lieu si $f$ est
quasi-compact et si de plus $f$ est séparé ou $X$ localement noethérien
(\hyperref[I.9.2.2-fr]{9.2.2})). Alors il existe un plus petit sous-préschéma
$Y'$ de $Y$ tel que l'injection canonique $j:Y'\to Y$ majore $f$ (ou, ce qui
revient au même (\hyperref[I.4.4.1-fr]{4.4.1}), tel que le sous-préschéma
$f^{-1}(Y')$ de $X$ soit \emph{identique} à $X$).
\end{proposition}

De façon plus précise~:

\begin{corollary}[9.5.2]
\label{I.9.5.2-fr}
Sous les conditions de (9.5.1), soit $f=(\psi,\theta)$ et soit $\mathscr{J}$
le noyau (quasi-cohérent) de l'homomorphisme
$\theta:\mathscr{O}_Y\to f_*(\mathscr{O}_X)$. Alors le sous-préschéma fermé
$Y'$ de $Y$ défini par $\mathscr{J}$ vérifie les conditions de (9.5.1).
\end{corollary}

\begin{proof}
Comme le foncteur $\psi^*$ est exact, la factorisation canonique
$\theta:\mathscr{O}_Y\to\mathscr{O}_Y/\mathscr{J}
\xrightarrow{\theta'}\psi_*(\mathscr{O}_X)$ donne
(\hyperref[0.3.5.4.3-fr]{0, 3.5.4.3}) une factorisation
\[
  \theta^\sharp:\psi^*(\mathscr{O}_Y)
  \to\psi^*(\mathscr{O}_Y)/\psi^*(\mathscr{J})
  \xrightarrow{{\theta'}^\sharp}\mathscr{O}_X~;
\]
comme pour tout $x\in X$, $\theta_x^\sharp$ est un homomorphisme local, il
en est de même de ${\theta'_x}^\sharp$~; si on désigne par $\psi_0$
l'application continue $\psi$ considérée comme application de $X$ dans
$X'$, par $\theta_0$ la restriction
$\theta'|X':(\mathscr{O}_Y/\mathscr{J})|X'
\to\psi_*(\mathscr{O}_X)|X'=(\psi_0)_*(\mathscr{O}_X)$, on voit donc que
$f_0=(\psi_0,\theta_0)$ est un morphisme de préschémas $X\to X'$
(\hyperref[I.2.2.1-fr]{2.2.1}) tel que $f=j\circ f_0$. Si maintenant $X''$
\oldpage[I]{177}
est un second sous-préschéma fermé de $Y$, défini par un faisceau
quasi-cohérent d'idéaux $\mathscr{J}'$ de $\mathscr{O}_Y$, et tel que
l'injection $j':X''\to Y$ majore $f$, on doit tout d'abord avoir
$X''\supset\psi(X)$, donc $X'\subset X''$ puisque $X''$ est fermé. En outre,
pour tout $y\in X''$, $\theta$ doit se factoriser en
$\mathscr{O}_y\to\mathscr{O}_y/\mathscr{J}'_y
\to(\psi_*(\mathscr{O}_X))_y$, ce qui par définition entraîne
$\mathscr{J}'_y\subset\mathscr{J}_y$, et par suite $X'$ est un
sous-préschéma fermé de $X''$ (\hyperref[I.4.1.10-fr]{4.1.10}).
\end{proof}

\begin{definition}[9.5.3]
\label{I.9.5.3-fr}
Lorsqu'il existe un plus petit sous-préschéma fermé $Y'$ de $Y$ tel que
l'injection canonique $j:Y'\to Y$ majore $f$, on dit que $Y'$ est le
préschéma image fermée de $X$ par le morphisme $f$.
\end{definition}

\begin{proposition}[9.5.4]
\label{I.9.5.4-fr}
Si $f_*(\mathscr{O}_X)$ est un $\mathscr{O}_Y$-Module quasi-cohérent,
l'espace sous-jacent de l'image fermée de $X$ par $f$ est l'adhérence
$\overline{f(X)}$ dans $Y$.
\end{proposition}

\begin{proof}
Comme le support de $f_*(\mathscr{O}_X)$ est contenu dans
$\overline{f(X)}$, on a (avec les notations de
(\hyperref[I.9.5.2-fr]{9.5.2})) $\mathscr{J}_y=\mathscr{O}_y$ pour
$y\notin\overline{f(X)}$, donc le support de
$\mathscr{O}_Y/\mathscr{J}$ est contenu dans $\overline{f(X)}$. En outre, ce
support est fermé et contient $f(X)$~: en effet, si $y\in f(X)$, l'élément
unité de l'anneau $(\psi_*(\mathscr{O}_X))_y$ n'est pas nul, étant le germe
en $y$ de la section
\[
  1\in\Gamma(X,\mathscr{O}_X)=\Gamma(Y,\psi_*(\mathscr{O}_X))~;
\]
comme il est l'image par $\theta$ de l'élément unité de $\mathscr{O}_y$, ce
dernier n'appartient pas à $\mathscr{J}_y$, donc
$\mathscr{O}_y/\mathscr{J}_y\neq0$~; ceci achève la démonstration.
\end{proof}

\begin{proposition}[9.5.5]
\label{I.9.5.5-fr}
\emph{(transitivité des images fermées).} Soient $f:X\to Y$ et $g:Y\to Z$
deux morphismes de préschémas~; on suppose que l'image fermée $Y'$ de $X$
par $f$ existe, et que, si $g'$ est la restriction de $g$ à $Y'$, l'image
fermée $Z'$ de $Y'$ par $g'$ existe. Alors l'image fermée de $X$ par
$g\circ f$ existe et est égale à $Z'$.
\end{proposition}

\begin{proof}
Il suffit (\hyperref[I.9.5.1-fr]{9.5.1}) de montrer que $Z'$ est le plus
petit sous-préschéma fermé $Z_1$ de $Z$ tel que le sous-préschéma fermé
$(g\circ f)^{-1}(Z_1)$ de $X$ (égal à $f^{-1}(g^{-1}(Z_1))$ par
(\hyperref[I.4.4.2-fr]{4.4.2})) soit égal à $X$~; il revient au même de dire
que $Z'$ est le plus petit sous-préschéma fermé de $Z$ tel que $f$ soit
majoré par l'injection $g^{-1}(Z_1)\to Y$
(\hyperref[I.4.4.1-fr]{4.4.1}). Or, en vertu de l'existence de l'image
fermée $Y'$, tout $Z_1$ ayant cette propriété est tel que $g^{-1}(Z_1)$
majore $Y'$, ce qui équivaut à dire que
$j^{-1}(g^{-1}(Z_1))=g'^{-1}(Z_1)=Y'$, en désignant par $j$ l'injection
$Y'\to Y$. Par définition de $Z'$, on en conclut bien que $Z'$ est le plus
petit sous-préschéma fermé de $Z$ vérifiant la condition précédente.
\end{proof}

\begin{corollary}[9.5.6]
\label{I.9.5.6-fr}
Soit $f:X\to Y$ un $S$-morphisme tel que $Y$ soit l'image fermée de $X$ par
$f$. Soit $Z$ un $S$-schéma~; si deux $S$-morphismes $g_1$, $g_2$ de $Y$
dans $Z$ sont tels que $g_1\circ f=g_2\circ f$, alors $g_1=g_2$.
\end{corollary}

\begin{proof}
Soit $h=(g_1,g_2)_S:Y\to Z\times_S Z$~; comme la diagonale
$T=\Delta_Z(Z)$ est un sous-préschéma fermé de $Z\times_S Z$,
$Y'=h^{-1}(T)$ est un sous-préschéma fermé de $Y$
(\hyperref[I.4.4.1-fr]{4.4.1}). Posons $u=g_1\circ f=g_2\circ f$~; on a
alors par définition du produit $h'=h\circ f=(u,u)_S$, donc
$h\circ f=\Delta_Z\circ u$~; comme $\Delta_Z^{-1}(T)=Z$, on a
$h'^{-1}(T)=u^{-1}(Z)=X$, donc $f^{-1}(Y')=X$. On en conclut
(\hyperref[I.4.4.1-fr]{4.4.1}) que l'injection canonique $Y'\to Y$ majore
$f$, donc $Y'=Y$ par hypothèse~; par suite
(\hyperref[I.4.4.1-fr]{4.4.1}), $h$ se factorise en $\Delta_Z\circ v$, où
$v$ est un morphisme $Y\to Z$, ce qui entraîne $g_1=g_2=v$.
\end{proof}

\begin{remark}[9.5.7]
\label{I.9.5.7-fr}
Si $X$ et $Y$ sont des $S$-schémas, la prop.
(\hyperref[I.9.5.6-fr]{9.5.6}) signifie que
\oldpage[I]{178}
lorsque $Y$ est l'image fermée de $X$ par $f$, $f$ est un
\emph{épimorphisme} dans la catégorie des $S$-schémas (T, 1.1). Nous
démontrerons au chap. V que, réciproquement, si l'image fermée $Y'$ de $X$
par $f$ existe et si $f$ est un épimorphisme de $S$-schémas, alors on a
nécessairement $Y'=Y$.
\end{remark}

\begin{proposition}[9.5.8]
\label{I.9.5.8-fr}
Supposons vérifiées les hypothèses de (\hyperref[I.9.5.1-fr]{9.5.1}), et
soit $Y'$ l'image fermée de $X$ par $f$. Pour tout ouvert $V$ de $Y$, soit
$f_V:f^{-1}(V)\to V$ la restriction de $f$~; alors l'image fermée de
$f^{-1}(V)$ par $f_V$ dans $V$ existe et est égale au préschéma induit par
$Y'$ sur l'ouvert $V\cap Y'$ de $Y'$ (autrement dit, au sous-préschéma
$\inf(V,Y')$ de $Y$ (\hyperref[I.4.4.3-fr]{4.4.3})).
\end{proposition}

\begin{proof}
Posons $X'=f^{-1}(V)$~; comme l'image directe de $\mathscr{O}_{X'}$ par
$f_V$ n'est autre que la restriction de $f_*(\mathscr{O}_X)$ à $V$, il est
clair que le noyau $\mathscr{J}'$ de l'homomorphisme
$\mathscr{O}_V\to(f_V)_*(\mathscr{O}_{X'})$ est la restriction de
$\mathscr{J}$ à $V$, d'où aussitôt la proposition.
\end{proof}

On notera que ce résultat s'interprète en disant que la formation de l'image
fermée commute à une extension $Y_1\to Y$ du préschéma de base, qui est une
\emph{immersion ouverte}. Nous verrons au chap. IV qu'il en est de même pour
une extension $Y_1\to Y$ qui est un morphisme \emph{plat}, pourvu que $f$
soit séparé et quasi-compact.

\begin{proposition}[9.5.9]
\label{I.9.5.9-fr}
Soit $f:X\to Y$ un morphisme tel que l'image fermée $Y'$ de $X$ par $f$
existe.
\begin{enumerate}
  \item[(i)] Si $X$ est réduit, il en est de même de $Y'$.
  \item[(ii)] Si on suppose vérifiées les hypothèses de
    (\hyperref[I.9.5.1-fr]{9.5.1}) et si $X$ est irréductible (resp.
    intègre), il en est de même de $Y'$.
\end{enumerate}
\end{proposition}

\begin{proof}
Par hypothèse, le morphisme $f$ se factorise en
$X\xrightarrow{g}Y'\xrightarrow{j}Y$, où $j$ est l'injection canonique.
Comme $X$ est réduit, $g$ se factorise en
$X\xrightarrow{h}Y'_{\mathrm{red}}\xrightarrow{j'}Y'$, où $j'$ est
l'injection canonique (\hyperref[I.5.2.2-fr]{5.2.2}), et il résulte alors de
la définition de $Y'$ que $Y'_{\mathrm{red}}=Y'$. Si de plus les conditions
de (\hyperref[I.9.5.1-fr]{9.5.1}) sont remplies, il résulte de
(\hyperref[I.9.5.4-fr]{9.5.4}) que $f(X)$ est dense dans $Y'$~; si $X$ est
irréductible, il en est donc de même de $Y'$
(\hyperref[0.2.1.5-fr]{0, 2.1.5}). L'assertion relative aux préschémas
intègres résulte de la conjonction des deux autres.
\end{proof}

\begin{proposition}[9.5.10]
\label{I.9.5.10-fr}
Soit $Y$ un sous-préschéma d'un préschéma $X$, tel que l'injection canonique
$i:Y\to X$ soit un morphisme quasi-compact. Il existe alors un plus petit
sous-préschéma fermé $\overline{Y}$ de $X$ majorant $Y$~; son espace
sous-jacent est l'adhérence de celui de $Y$~; ce dernier est ouvert dans son
adhérence, et le préschéma $Y$ est induit sur cet ouvert par $\overline{Y}$.
\end{proposition}

\begin{proof}
Il suffit d'appliquer (\hyperref[I.9.5.1-fr]{9.5.1}) à l'injection $j$, qui
est séparée (\hyperref[I.5.5.1-fr]{5.5.1}) et quasi-compacte par hypothèse~;
(\hyperref[I.9.5.1-fr]{9.5.1}) prouve donc l'existence de $\overline{Y}$ et
(\hyperref[I.9.5.4-fr]{9.5.4}) montre que son espace sous-jacent est
l'adhérence de $Y$ dans $X$~; comme $Y$ est localement fermé dans $X$, il
est ouvert dans $\overline{Y}$, et la dernière assertion provient de
(\hyperref[I.9.5.8-fr]{9.5.8}) appliqué à un ouvert $V$ de $X$ tel que $Y$
soit fermé dans $V$.
\end{proof}

Avec ces notations, si l'injection $V\to X$ est quasi-compacte, et si
$\mathscr{J}$ est le faisceau quasi-cohérent d'idéaux de
$\mathscr{O}_X|V$ définissant le sous-préschéma fermé $Y$ de $V$, il résulte
de (\hyperref[I.9.5.1-fr]{9.5.1}) que le faisceau quasi-cohérent d'idéaux de
$\mathscr{O}_X$ définissant $\overline{Y}$ est le prolongement canonique
(\hyperref[I.9.4.1-fr]{9.4.1}) $\overline{\mathscr{J}}$ de $\mathscr{J}$,
car c'est évidemment le plus grand sous-faisceau d'idéaux quasi-cohérent de
$\mathscr{O}_X$ induisant $\mathscr{J}$ sur $V$.
\oldpage[I]{179}
\begin{corollary}[9.5.11]
\label{I.9.5.11-fr}
Sous les hypothèses de (\hyperref[I.9.5.10-fr]{9.5.10}), toute section de
$\mathscr{O}_{\overline{Y}}$ au-dessus d'un ouvert $V$ de $\overline{Y}$
qui est nulle dans $V\cap Y$ est nulle.
\end{corollary}

\begin{proof}
En vertu de (\hyperref[I.9.5.8-fr]{9.5.8}), on peut se ramener au cas où
$V=\overline{Y}$. Si l'on tient compte de ce que les sections de
$\mathscr{O}_{\overline{Y}}$ au-dessus de $\overline{Y}$ correspondent
canoniquement aux $\overline{Y}$-sections de
$\overline{Y}\otimes_{\mathbb{Z}}\mathbb{Z}[T]$
(\hyperref[I.3.3.15-fr]{3.3.15}) et de ce que ce dernier est séparé sur
$\overline{Y}$, le corollaire apparaît comme un cas particulier de
(\hyperref[I.9.5.6-fr]{9.5.6}).
\end{proof}

Lorsqu'il existe un plus petit sous-préschéma fermé $\overline{Y}$ de $X$
majorant un sous-préschéma $Y$ de $X$, on dit que $\overline{Y}$ est
l'\emph{adhérence} de $Y$ dans $X$, lorsqu'il n'en résulte pas de
confusion.

\subsection{Faisceaux quasi-cohérents d'algèbres~; changement du faisceau
structural.}
\label{subsection:I.9.6-fr}

\begin{proposition}[9.6.1]
\label{I.9.6.1-fr}
Soient $X$ un préschéma, $\mathscr{B}$ une $\mathscr{O}_X$-Algèbre
quasi-cohérente (\hyperref[0.5.1.3-fr]{0, 5.1.3}). Pour qu'un
$\mathscr{B}$-Module $\mathscr{F}$ soit quasi-cohérent (sur l'espace
annelé $(X,\mathscr{B})$), il faut et il suffit que $\mathscr{F}$ soit un
$\mathscr{O}_X$-Module quasi-cohérent.
\end{proposition}

\begin{proof}
Comme la question est locale, on peut supposer $X$ affine d'anneau $A$, et
alors $\mathscr{B}=\widetilde{B}$, où $B$ est une $A$-algèbre
(\hyperref[I.1.4.3-fr]{1.4.3}). Si $\mathscr{F}$ est quasi-cohérent sur
l'espace annelé $(X,\mathscr{B})$, on peut aussi supposer que
$\mathscr{F}$ est le conoyau d'un $\mathscr{B}$-homomorphisme
$\mathscr{B}^{(I)}\to\mathscr{B}^{(J)}$~; comme cet homomorphisme est
aussi un $\mathscr{O}_X$-homomorphisme de $\mathscr{O}_X$-Modules, et que
$\mathscr{B}^{(I)}$ et $\mathscr{B}^{(J)}$ sont des
$\mathscr{O}_X$-Modules quasi-cohérents
(\hyperref[I.1.3.9-fr]{1.3.9, (ii)}), $\mathscr{F}$ est aussi un
$\mathscr{O}_X$-Module quasi-cohérent
(\hyperref[I.1.3.9-fr]{1.3.9, (i)}).

Inversement, si $\mathscr{F}$ est un $\mathscr{O}_X$-Module
quasi-cohérent, on a $\mathscr{F}=\widetilde{M}$, où $M$ est un $B$-module
(\hyperref[I.1.4.3-fr]{1.4.3})~; $M$ est isomorphe à un conoyau de
$B$-homomorphisme $B^{(I)}\to B^{(J)}$, donc $\mathscr{F}$ est un
$\mathscr{B}$-Module isomorphe au conoyau de l'homomorphisme correspondant
$\mathscr{B}^{(I)}\to\mathscr{B}^{(J)}$
(\hyperref[I.1.3.13-fr]{1.3.13}), ce qui achève la démonstration.
\end{proof}

En particulier, si $\mathscr{F}$ et $\mathscr{G}$ sont deux
$\mathscr{B}$-Modules quasi-cohérents,
$\mathscr{F}\otimes_{\mathscr{B}}\mathscr{G}$ est un
$\mathscr{B}$-Module quasi-cohérent~; il en est de même de
$\mathscr{H}\!om_{\mathscr{B}}(\mathscr{F},\mathscr{G})$ lorsqu'on
suppose en outre que $\mathscr{F}$ admet une présentation finie
(\hyperref[I.1.3.13-fr]{1.3.13}).

\begin{env}[9.6.2]
\label{I.9.6.2-fr}
Étant donné un préschéma $X$, nous dirons qu'une
$\mathscr{O}_X$-Algèbre quasi-cohérente $\mathscr{B}$ est de
\emph{type fini} si pour tout $x\in X$, il existe un voisinage ouvert
\emph{affine} $U$ de $x$ tel que $\Gamma(U,\mathscr{B})=B$ soit une
algèbre de type fini sur $\Gamma(U,\mathscr{O}_X)=A$. On a alors
$\mathscr{B}|U=\widetilde{B}$ et pour tout $f\in A$, la
$(\mathscr{O}_X|D(f))$-Algèbre induite $\mathscr{B}|D(f)$ est de type
fini, car elle est isomorphe à $\widetilde{B_f}$, et
$B_f=B\otimes_A A_f$ est évidemment une algèbre de type fini sur $A_f$.
Comme les $D(f)$ forment une base de la topologie de $U$, on en conclut que
si $\mathscr{B}$ est une $\mathscr{O}_X$-Algèbre quasi-cohérente de type
fini, pour tout ouvert $V$ de $X$, $\mathscr{B}|V$ est une
$(\mathscr{O}_X|V)$-Algèbre quasi-cohérente de type fini.
\end{env}

\begin{proposition}[9.6.3]
\label{I.9.6.3-fr}
Soit $X$ un préschéma localement noethérien. Toute
$\mathscr{O}_X$-Algèbre quasi-cohérente de type fini $\mathscr{B}$ est
alors un faisceau cohérent d'anneaux
(\hyperref[0.5.3.7-fr]{0, 5.3.7}).
\end{proposition}

\begin{proof}
On peut de nouveau se limiter au cas où $X$ est un schéma affine d'anneau
noethérien $A$, et où $\mathscr{B}=\widetilde{B}$, $B$ étant une
$A$-algèbre de type fini~; $B$ est donc un anneau noethérien. Cela étant, il
faut prouver que le noyau $\mathscr{N}$ d'un
$\mathscr{B}$-homomorphisme $\mathscr{B}^m\to\mathscr{B}$ est un
$\mathscr{B}$-
\oldpage[I]{180}
Module de type fini~; or, il est isomorphe (en tant que
$\mathscr{B}$-Module) à $\widetilde{N}$, où $N$ est le noyau de
l'homomorphisme correspondant de $B$-modules $B^m\to B$
(\hyperref[I.1.3.13-fr]{1.3.13}). Comme $B$ est noethérien, le sous-module
$N$ de $B^m$ est un $B$-module de type fini, donc il existe un homomorphisme
$B^p\to B^m$ d'image $N$~; la suite $B^p\to B^m\to B$ étant exacte, il
en est de même de la suite
$\mathscr{B}^p\to\mathscr{B}^m\to\mathscr{B}$ correspondante
(\hyperref[I.1.3.5-fr]{1.3.5}) et comme $\mathscr{N}$ est l'image de
$\mathscr{B}^p\to\mathscr{B}^m$
(\hyperref[I.1.3.9-fr]{1.3.9, (i)}), la proposition est démontrée.
\end{proof}

\begin{corollary}[9.6.4]
\label{I.9.6.4-fr}
Sous les hypothèses de (\hyperref[I.9.6.3-fr]{9.6.3}), pour qu'un
$\mathscr{B}$-Module $\mathscr{F}$ soit cohérent, il faut et il suffit
qu'il soit un $\mathscr{O}_X$-Module quasi-cohérent et un
$\mathscr{B}$-Module de type fini. S'il en est ainsi, et si $\mathscr{G}$
est un sous-$\mathscr{B}$-Module ou un $\mathscr{B}$-Module quotient de
$\mathscr{F}$, pour que $\mathscr{G}$ soit un $\mathscr{B}$-Module
cohérent, il faut et il suffit que $\mathscr{G}$ soit un
$\mathscr{O}_X$-Module quasi-cohérent.
\end{corollary}

\begin{proof}
Compte tenu de (\hyperref[I.9.6.1-fr]{9.6.1}), les conditions sur
$\mathscr{F}$ sont évidemment nécessaires~; montrons qu'elles sont
suffisantes. On peut se limiter au cas où $X$ est affine d'anneau
noethérien $A$, $\mathscr{B}=\widetilde{B}$, où $B$ est une $A$-algèbre de
type fini, $\mathscr{F}=\widetilde{M}$, où $M$ est un $B$-module, et où il
existe un $\mathscr{B}$-homomorphisme surjectif
$\mathscr{B}^m\to\mathscr{F}\to 0$. On a alors une suite exacte
correspondante $B^m\to M\to 0$, donc $M$ est un $B$-module de type fini~;
en outre, le noyau $P$ de l'homomorphisme $B^m\to M$ est alors un
$B$-module de type fini, puisque $B$ est noethérien. On en conclut
(\hyperref[I.1.3.13-fr]{1.3.13}) que $\mathscr{F}$ est le conoyau d'un
$\mathscr{B}$-homomorphisme $\mathscr{B}^n\to\mathscr{B}^m$, et est donc
cohérent puisque $\mathscr{B}$ est un faisceau cohérent d'anneaux
(\hyperref[0.5.3.4-fr]{0, 5.3.4}). Le même raisonnement montre qu'un
sous-$\mathscr{B}$-Module (resp. un $\mathscr{B}$-Module quotient)
quasi-cohérent de $\mathscr{F}$ est de type fini, d'où la seconde partie du
corollaire.
\end{proof}

\begin{proposition}[9.6.5]
\label{I.9.6.5-fr}
Soit $X$ un schéma quasi-compact ou un préschéma dont l'espace sous-jacent
est noethérien. Pour toute $\mathscr{O}_X$-Algèbre quasi-cohérente de type
fini $\mathscr{B}$, il existe un sous-$\mathscr{O}_X$-Module quasi-cohérent
de type fini $\mathscr{E}$ de $\mathscr{B}$ tel que $\mathscr{E}$ engendre
(\hyperref[0.4.1.4-fr]{0, 4.1.4}) la $\mathscr{O}_X$-Algèbre
$\mathscr{B}$.
\end{proposition}

\begin{proof}
En effet, il existe par hypothèse un recouvrement fini $(U_i)$ de $X$ formé
d'ouverts affines tels que $\Gamma(U_i,\mathscr{B})=B_i$ soit une algèbre
de type fini sur $\Gamma(U_i,\mathscr{O}_X)=A_i$~; soit $E_i$ un
sous-$A_i$-module de type fini de $B_i$ engendrant la $A_i$-algèbre $B_i$~;
en vertu de (\hyperref[I.9.4.7-fr]{9.4.7}), il existe un
sous-$\mathscr{O}_X$-Module $\mathscr{E}_i$ de $\mathscr{B}$,
quasi-cohérent et de type fini, tel que
$\mathscr{E}_i|U_i=\widetilde{E_i}$. Il est clair que la somme
$\mathscr{E}$ des $\mathscr{E}_i$ répond à la question.
\end{proof}

\begin{proposition}[9.6.6]
\label{I.9.6.6-fr}
Soit $X$ un préschéma dont l'espace sous-jacent est localement noethérien,
ou un schéma quasi-compact. Alors toute $\mathscr{O}_X$-Algèbre
quasi-cohérente $\mathscr{B}$ est limite inductive de ses
sous-$\mathscr{O}_X$-Algèbres quasi-cohérentes de type fini.
\end{proposition}

\begin{proof}
En effet, il résulte de (\hyperref[I.9.4.9-fr]{9.4.9}) que $\mathscr{B}$
est limite inductive (en tant que $\mathscr{O}_X$-Module) de ses
sous-$\mathscr{O}_X$-Modules quasi-cohérents de type fini~; ces derniers
engendrent des sous-$\mathscr{O}_X$-Algèbres quasi-cohérentes de type fini
de $\mathscr{B}$ (\hyperref[I.1.3.14-fr]{1.3.14}), dont $\mathscr{B}$ est
\emph{a fortiori} limite inductive.
\end{proof}

\section{Schémas formels}
\label{section:I.10-fr}

\subsection{Schémas formels affines.}
\label{subsection:I.10.1-fr}

\begin{env}[10.1.1]
\label{I.10.1.1-fr}
Soit $A$ un anneau topologique \emph{admissible}
(\hyperref[0.7.1.2-fr]{0, 7.1.2})~; pour tout idéal de définition
$\mathfrak{J}$ de $A$, $\operatorname{Spec}(A/\mathfrak{J})$ s'identifie
au sous-espace fermé $V(\mathfrak{J})$ de $\operatorname{Spec}(A)$
(\hyperref[I.1.1.11-fr]{1.1.11}), ensemble des idéaux premiers
\emph{ouverts} de $A$~; cet espace topologique ne dépend pas de
\oldpage[I]{181}
l'idéal de définition $\mathfrak{J}$ considéré~; notons-le
$\mathfrak{X}$. Soit $(\mathfrak{J}_\lambda)$ un système fondamental de
voisinages de 0 dans $A$, formé d'idéaux de définition, et pour tout
$\lambda$, soit $\mathscr{O}_\lambda$ le faisceau structural de
$\operatorname{Spec}(A/\mathfrak{J}_\lambda)$~; ce faisceau est induit sur
$\mathfrak{X}$ par
$\widetilde{A}/\widetilde{\mathfrak{J}}_\lambda$ (lequel est nul hors de
$\mathfrak{X}$).

Pour $\mathfrak{J}_\mu\subset\mathfrak{J}_\lambda$, l'homomorphisme
canonique $A/\mathfrak{J}_\mu\to A/\mathfrak{J}_\lambda$ définit donc un
homomorphisme
$u_{\lambda\mu}:\mathscr{O}_\mu\to\mathscr{O}_\lambda$ de faisceaux
d'anneaux (\hyperref[I.1.6.1-fr]{1.6.1}), et
$(\mathscr{O}_\lambda)$ est un \emph{système projectif de faisceaux
d'anneaux} pour ces homomorphismes. Comme la topologie de $\mathfrak{X}$
admet une base formée d'ouverts quasi-compacts, on peut associer à tout
$\mathscr{O}_\lambda$ un \emph{faisceau d'anneaux topologiques
pseudo-discrets} (\hyperref[0.3.8.1-fr]{0, 3.8.1}) qui a
$\mathscr{O}_\lambda$ comme faisceau d'anneaux (sans topologie) sous-jacent,
et que nous noterons encore $\mathscr{O}_\lambda$~; et les
$\mathscr{O}_\lambda$ forment encore un \emph{système projectif de faisceaux
d'anneaux topologiques} (\hyperref[0.3.8.2-fr]{0, 3.8.2}). Nous désignerons
par $\mathscr{O}_{\mathfrak{X}}$ le \emph{faisceau d'anneaux topologiques}
sur $\mathfrak{X}$, limite projective du système
$(\mathscr{O}_\lambda)$~; pour tout ouvert \emph{quasi-compact} $U$ de
$\mathfrak{X}$, $\Gamma(U,\mathscr{O}_{\mathfrak{X}})$ est donc l'anneau
topologique limite projective du système d'anneaux \emph{discrets}
$\Gamma(U,\mathscr{O}_\lambda)$
(\hyperref[0.3.2.6-fr]{0, 3.2.6}).
\end{env}

\begin{definition}[10.1.2]
\label{I.10.1.2-fr}
Étant donné un anneau topologique admissible $A$, on appelle spectre formel
de $A$ et on note $\operatorname{Spf}(A)$ le sous-espace fermé
$\mathfrak{X}$ de $\operatorname{Spec}(A)$ formé des idéaux premiers
ouverts de $A$. On dit qu'un espace topologiquement annelé est un schéma
formel affine s'il est isomorphe à un spectre formel
$\operatorname{Spf}(A)=\mathfrak{X}$ muni du faisceau d'anneaux
topologiques $\mathscr{O}_{\mathfrak{X}}$ limite projective des faisceaux
d'anneaux pseudo-discrets
$(\widetilde{A}/\widetilde{\mathfrak{J}}_\lambda)|\mathfrak{X}$, où
$\mathfrak{J}_\lambda$ parcourt l'ensemble filtrant des idéaux de
définition de $A$.
\end{definition}

Quand nous parlerons d'un \emph{spectre formel}
$\mathfrak{X}=\operatorname{Spf}(A)$ comme d'un schéma formel affine, il
s'agira toujours de l'espace topologiquement annelé
$(\mathfrak{X},\mathscr{O}_{\mathfrak{X}})$ où
$\mathscr{O}_{\mathfrak{X}}$ est défini comme ci-dessus.

On notera que tout \emph{schéma affine} $X=\operatorname{Spec}(A)$ peut être
considéré comme un schéma formel affine d'une seule manière, en considérant
$A$ comme un anneau topologique discret~: les anneaux topologiques
$\Gamma(U,\mathscr{O}_X)$ sont alors discrets lorsque $U$ est
quasi-compact (mais non en général lorsque $U$ est un ouvert quelconque de
$X$).

\begin{proposition}[10.1.3]
\label{I.10.1.3-fr}
Si $\mathfrak{X}=\operatorname{Spf}(A)$, où $A$ est un anneau admissible,
$\Gamma(\mathfrak{X},\mathscr{O}_{\mathfrak{X}})$ est topologiquement
isomorphe à $A$.
\end{proposition}

\begin{proof}
En effet, comme $\mathfrak{X}$ est fermé dans $\operatorname{Spec}(A)$, il
est quasi-compact, et par suite
$\Gamma(\mathfrak{X},\mathscr{O}_{\mathfrak{X}})$ est topologiquement
isomorphe à la limite projective des anneaux discrets
$\Gamma(\mathfrak{X},\mathscr{O}_\lambda)$~; mais
$\Gamma(\mathfrak{X},\mathscr{O}_\lambda)$ est isomorphe à
$A/\mathfrak{J}_\lambda$ (\hyperref[I.1.3.7-fr]{1.3.7})~; comme $A$ est
séparé et complet, il est topologiquement isomorphe à
$\varprojlim A/\mathfrak{J}_\lambda$
(\hyperref[0.7.2.1-fr]{0, 7.2.1}), d'où la proposition.
\end{proof}

\begin{proposition}[10.1.4]
\label{I.10.1.4-fr}
Soient $A$ un anneau admissible, $\mathfrak{X}=\operatorname{Spf}(A)$, et
pour tout $f\in A$, soit
$\mathfrak{D}(f)=D(f)\cap\mathfrak{X}$~; l'espace topologiquement annelé
$(\mathfrak{D}(f),\mathscr{O}_{\mathfrak{X}}|\mathfrak{D}(f))$ est
isomorphe au spectre formel affine
$\operatorname{Spf}(A_{\{f\}})$
(\hyperref[0.7.6.15-fr]{0, 7.6.15}).
\end{proposition}

\begin{proof}
Pour tout idéal de définition $\mathfrak{J}$ de $A$, l'anneau discret
$S_f^{-1}A/S_f^{-1}\mathfrak{J}$ s'identifie canoniquement à
$A_{\{f\}}/\mathfrak{J}_{\{f\}}$
(\hyperref[0.7.6.9-fr]{0, 7.6.9}), donc
(\hyperref[I.1.2.5-fr]{1.2.5} et \hyperref[I.1.2.6-fr]{1.2.6}), l'espace
topologique $\operatorname{Spf}(A_{\{f\}})$ s'identifie canoniquement à
$\mathfrak{D}(f)$. En outre, pour tout ouvert quasi-compact $U$ de
$\mathfrak{X}$ contenu dans $\mathfrak{D}(f)$,
$\Gamma(U,\mathscr{O}_\lambda)$ s'identifie au module des sections du
faisceau structural de
$\operatorname{Spec}(S_f^{-1}A/S_f^{-1}\mathfrak{J}_\lambda)$ au-dessus de
$U$ (\hyperref[I.1.3.6-fr]{1.3.6}), donc, si on pose
$\mathfrak{Y}=\operatorname{Spf}(A_{\{f\}})$,
$\Gamma(U,\mathscr{O}_{\mathfrak{X}})$ s'identifie au module des sections
$\Gamma(U,\mathscr{O}_{\mathfrak{Y}})$, ce qui démontre la proposition.
\end{proof}

\oldpage[I]{182}
\begin{env}[10.1.5]
\label{I.10.1.5-fr}
En tant que faisceau d'anneaux \emph{sans topologie}, le faisceau
structural $\mathscr{O}_{\mathfrak{X}}$ de
$\operatorname{Spf}(A)$ admet pour tout $x\in\mathfrak{X}$ une fibre qui,
en vertu de (\hyperref[I.10.1.4-fr]{10.1.4}), s'identifie à la limite
inductive
$\varinjlim_{f\notin\mathfrak{j}_x}A_{\{f\}}$. Par suite
(\hyperref[0.7.6.17-fr]{0, 7.6.17} et
\hyperref[0.7.6.18-fr]{7.6.18})~:
\end{env}

\begin{proposition}[10.1.6]
\label{I.10.1.6-fr}
Pour tout $x\in\mathfrak{X}=\operatorname{Spf}(A)$, la fibre
$\mathscr{O}_x$ est un anneau local dont le corps résiduel est isomorphe à
$k(x)=A_x/\mathfrak{j}_xA_x$. Si en outre $A$ est adique et noethérien,
$\mathscr{O}_x$ est un anneau noethérien.
\end{proposition}

Comme $k(x)$ n'est pas réduit à 0, on conclut de ce résultat que le
\emph{support} du faisceau d'anneaux $\mathscr{O}_{\mathfrak{X}}$ est
\emph{égal à $\mathfrak{X}$}.

\subsection{Morphismes de schémas formels affines.}
\label{subsection:I.10.2-fr}

\begin{env}[10.2.1]
\label{I.10.2.1-fr}
Soient $A$, $B$ deux anneaux admissibles, et soit
$\varphi:B\to A$ un homomorphisme \emph{continu}. L'application continue
$ {}^a\varphi:\operatorname{Spec}(A)\to\operatorname{Spec}(B)$
(\hyperref[I.1.2.1-fr]{1.2.1}) applique alors
$\mathfrak{X}=\operatorname{Spf}(A)$ dans
$\mathfrak{Y}=\operatorname{Spf}(B)$, car l'image réciproque par
$\varphi$ d'un idéal premier ouvert de $A$ est un idéal premier ouvert de
$B$. D'autre part, pour tout $g\in B$, $\varphi$ définit un homomorphisme
continu
\[
  \Gamma(\mathfrak{D}(g),\mathscr{O}_{\mathfrak{Y}})
  \longrightarrow
  \Gamma(\mathfrak{D}(\varphi(g)),\mathscr{O}_{\mathfrak{X}})
\]
en vertu de (\hyperref[I.10.1.4-fr]{10.1.4}),
(\hyperref[I.10.1.3-fr]{10.1.3}) et
(\hyperref[0.7.6.7-fr]{0, 7.6.7})~; comme ces homomorphismes satisfont aux
conditions de compatibilité pour les restrictions correspondant au
passage de $g$ à un multiple de $g$, et que
$\mathfrak{D}(\varphi(g))={}^a\varphi^{-1}(\mathfrak{D}(g))$, ils
définissent un homomorphisme \emph{continu} de faisceaux d'anneaux
topologiques
$\mathscr{O}_{\mathfrak{Y}}\to
{}^a\varphi_*(\mathscr{O}_{\mathfrak{X}})$
(\hyperref[0.3.2.5-fr]{0, 3.2.5}), que nous noterons encore
$\widetilde{\varphi}$~; on a ainsi obtenu un morphisme
$\Phi=({}^a\varphi,\widetilde{\varphi})$ d'espaces topologiquement annelés
$\mathfrak{X}\to\mathfrak{Y}$. On notera qu'en tant qu'homomorphisme de
faisceaux d'anneaux sans topologie, $\widetilde{\varphi}$ définit un
homomorphisme
$\widetilde{\varphi}_x^\sharp:
\mathscr{O}_{{}^a\varphi(x)}\to\mathscr{O}_x$ sur les fibres, pour tout
$x\in\mathfrak{X}$.
\end{env}

\begin{proposition}[10.2.2]
\label{I.10.2.2-fr}
Soient $A$, $B$ deux anneaux topologiques admissibles, et soient
$\mathfrak{X}=\operatorname{Spf}(A)$,
$\mathfrak{Y}=\operatorname{Spf}(B)$. Pour qu'un morphisme
$u=(\psi,\theta):\mathfrak{X}\to\mathfrak{Y}$ d'espaces topologiquement
annelés soit de la forme $({}^a\varphi,\widetilde{\varphi})$, où
$\varphi$ est un homomorphisme continu d'anneaux $B\to A$, il faut et il
suffit que pour tout $x\in\mathfrak{X}$,
$\theta_x^\sharp$ soit un homomorphisme local
$\mathscr{O}_{\psi(x)}\to\mathscr{O}_x$.
\end{proposition}

\begin{proof}
La condition est nécessaire~: en effet, soit
$\mathfrak{p}=\mathfrak{j}_x\in\operatorname{Spf}(A)$, et soit
$\mathfrak{q}=\varphi^{-1}(\mathfrak{j}_x)$~; si
$g\notin\mathfrak{q}$, on a donc $\varphi(g)\notin\mathfrak{p}$, et il est
immédiat que l'homomorphisme
$B_{\{g\}}\to A_{\{\varphi(g)\}}$ déduit de $\varphi$
(\hyperref[0.7.6.7-fr]{0, 7.6.7}) transforme
$\mathfrak{q}_{\{g\}}$ en une partie de
$\mathfrak{p}_{\{\varphi(g)\}}$~; passant à la limite inductive, on voit
donc (compte tenu de (\hyperref[I.10.1.5-fr]{10.1.5}) et de
(\hyperref[0.7.6.17-fr]{0, 7.6.17})) que
$\widetilde{\varphi}_x^\sharp$ est un homomorphisme local.

Inversement, soit $(\psi,\theta)$ un morphisme vérifiant la condition de
l'énoncé~; en vertu de (\hyperref[I.10.1.3-fr]{10.1.3}), $\theta$ définit
un homomorphisme continu d'anneaux
\[
  \varphi=\Gamma(\theta):
  B=\Gamma(\mathfrak{Y},\mathscr{O}_{\mathfrak{Y}})
  \longrightarrow
  \Gamma(\mathfrak{X},\mathscr{O}_{\mathfrak{X}})=A.
\]
En vertu de l'hypothèse sur $\theta$, pour que la section $\varphi(g)$ de
$\mathscr{O}_{\mathfrak{X}}$ au-dessus de $\mathfrak{X}$ ait un germe
inversible au point $x$, il faut et il suffit que $g$ ait un germe
inversible au point $\psi(x)$. Mais en vertu de
(\hyperref[0.7.6.17-fr]{0, 7.6.17}), les sections de
$\mathscr{O}_{\mathfrak{X}}$ (resp. $\mathscr{O}_{\mathfrak{Y}}$)
au-dessus de $\mathfrak{X}$ (resp. $\mathfrak{Y}$) dont le germe n'est pas
inversible au point $x$ (resp. $\psi(x)$) sont exactement les éléments de
$\mathfrak{j}_x$
\oldpage[I]{183}
(resp. $\mathfrak{j}_{\psi(x)}$)~; la remarque précédente montre donc que
${}^a\varphi=\psi$. Enfin, pour tout $g\in B$, le diagramme
\[
  \xymatrix{
    B=\Gamma(\mathfrak{Y},\mathscr{O}_{\mathfrak{Y}})
      \ar[r]^\varphi\ar[d] &
    \Gamma(\mathfrak{X},\mathscr{O}_{\mathfrak{X}})=A\ar[d]\\
    B_{\{g\}}=\Gamma(\mathfrak{D}(g),\mathscr{O}_{\mathfrak{Y}})
      \ar[r]_{\Gamma(\theta_{\mathfrak{D}(g)})} &
    \Gamma(\mathfrak{D}(\varphi(g)),\mathscr{O}_{\mathfrak{X}})
      =A_{\{\varphi(g)\}}
  }
\]
est commutatif~; d'après la propriété universelle des anneaux complets de
fractions (\hyperref[0.7.6.6-fr]{0, 7.6.6}),
$\theta_{\mathfrak{D}(g)}$ est égale à
$\widetilde{\varphi}_{\mathfrak{D}(g)}$ pour tout $g\in B$, donc
(\hyperref[0.3.2.5-fr]{0, 3.2.5}) on a
$\theta=\widetilde{\varphi}$.
\end{proof}

Nous dirons qu'un morphisme $(\psi,\theta)$ d'espaces topologiquement
annelés vérifiant la condition de (\hyperref[I.10.2.2-fr]{10.2.2}) est un
\emph{morphisme de schémas formels affines}. On peut dire que les foncteurs
$\operatorname{Spf}(A)$ en $A$ et
$\Gamma(\mathfrak{X},\mathscr{O}_{\mathfrak{X}})$ en $\mathfrak{X}$
définissent une \emph{équivalence} de la catégorie des anneaux admissibles
et de la catégorie duale de la catégorie des schémas formels affines
(T, I, 1.2).

\begin{env}[10.2.3]
\label{I.10.2.3-fr}
Comme cas particulier de (\hyperref[I.10.2.2-fr]{10.2.2}), notons que, pour
$f\in A$, l'injection canonique du schéma formel affine induit par
$\mathfrak{X}$ sur $\mathfrak{D}(f)$ correspond à l'homomorphisme continu
canonique $A\to A_{\{f\}}$. Sous les hypothèses de
(\hyperref[I.10.2.2-fr]{10.2.2}), soient $h$ un élément de $B$, $g$ un
élément de $A$, multiple de $\varphi(h)$~; on a alors
$\psi(\mathfrak{D}(g))\subset\mathfrak{D}(h)$~; la restriction de $u$ à
$\mathfrak{D}(g)$, considérée comme morphisme de $\mathfrak{D}(g)$ dans
$\mathfrak{D}(h)$, est l'unique morphisme $v$ rendant commutatif le
diagramme
\[
  \xymatrix{
    \mathfrak{D}(g)\ar[r]^v\ar[d] &
    \mathfrak{D}(h)\ar[d]\\
    \mathfrak{X}\ar[r]_u &
    \mathfrak{Y}
  }
\]
Ce morphisme correspond à l'unique homomorphisme continu
$\varphi':B_{\{h\}}\to A_{\{g\}}$
(\hyperref[0.7.6.7-fr]{0, 7.6.7}) rendant commutatif le diagramme
\[
  \xymatrix{
    A\ar[d] &
    B\ar[l]_\varphi\ar[d]\\
    A_{\{g\}} &
    B_{\{h\}}\ar[l]_{\varphi'}
  }
\]
\end{env}

\subsection{Idéaux de définition d'un schéma formel affine.}
\label{subsection:I.10.3-fr}

\begin{env}[10.3.1]
\label{I.10.3.1-fr}
Soient $A$ un anneau admissible, $\mathfrak{J}$ un idéal ouvert de $A$,
$\mathfrak{X}$ le schéma formel affine $\operatorname{Spf}(A)$. Soit
$(\mathfrak{J}_\lambda)$ l'ensemble des idéaux de définition de $A$
contenus dans $\mathfrak{J}$~; alors
$\widetilde{\mathfrak{J}}/\widetilde{\mathfrak{J}}_\lambda$ est un
faisceau d'idéaux de
$\widetilde{A}/\widetilde{\mathfrak{J}}_\lambda$. Désignons par
$\mathfrak{J}^\Delta$ la limite projective des faisceaux induits sur
$\mathfrak{X}$ par
$\widetilde{\mathfrak{J}}/\widetilde{\mathfrak{J}}_\lambda$, qui
s'identifie à un \emph{faisceau d'idéaux} de
$\mathscr{O}_{\mathfrak{X}}$ (\hyperref[0.3.2.6-fr]{0, 3.2.6}). Pour tout
$f\in A$, $\Gamma(\mathfrak{D}(f),\mathfrak{J}^\Delta)$ est la limite
projective de
$S_f^{-1}\mathfrak{J}/S_f^{-1}\mathfrak{J}_\lambda$, autrement dit
s'identifie à l'idéal ouvert $\mathfrak{J}_{\{f\}}$ de l'anneau
$A_{\{f\}}$ (\hyperref[0.7.6.9-fr]{0, 7.6.9}), et en particulier
$\Gamma(\mathfrak{X},\mathfrak{J}^\Delta)=\mathfrak{J}$~; on en conclut
(les $\mathfrak{D}(f)$ formant une base de la topologie de $\mathfrak{X}$)
que l'on a
\[
  \mathfrak{J}^\Delta|\mathfrak{D}(f)
  =(\mathfrak{J}_{\{f\}})^\Delta.
  \tag{10.3.1.1}\label{I.10.3.1.1-fr}
\]
\end{env}

\oldpage[I]{184}
\begin{env}[10.3.2]
\label{I.10.3.2-fr}
Avec les notations de (\hyperref[I.10.3.1-fr]{10.3.1}), pour tout $f\in A$,
l'application canonique de
$A_{\{f\}}=\Gamma(\mathfrak{D}(f),\mathscr{O}_{\mathfrak{X}})$ dans
$\Gamma\bigl(\mathfrak{D}(f),
(\widetilde{A}/\widetilde{\mathfrak{J}})|\mathfrak{X}\bigr)
=S_f^{-1}A/S_f^{-1}\mathfrak{J}$
est \emph{surjective} et a pour noyau
$\Gamma(\mathfrak{D}(f),\mathfrak{J}^\Delta)=\mathfrak{J}_{\{f\}}$
(\hyperref[0.7.6.9-fr]{0, 7.6.9})~; ces applications définissent donc un
homomorphisme \emph{surjectif} continu, dit \emph{canonique}, du faisceau
d'anneaux topologiques $\mathscr{O}_{\mathfrak{X}}$ sur le faisceau
d'anneaux discrets
$(\widetilde{A}/\widetilde{\mathfrak{J}})|\mathfrak{X}$, dont le noyau est
$\mathfrak{J}^\Delta$~; cet homomorphisme n'est autre d'ailleurs que
$\widetilde{\varphi}$ (\hyperref[I.10.2.1-fr]{10.2.1}), où $\varphi$ est
l'homomorphisme continu $A\to A/\mathfrak{J}$~; le morphisme
$({}^a\varphi,\widetilde{\varphi}):
\operatorname{Spec}(A/\mathfrak{J})\to\mathfrak{X}$ de schémas formels
affines (où ${}^a\varphi$ est d'ailleurs l'homéomorphisme identique de
$\mathfrak{X}$ sur lui-même) est encore dit \emph{canonique}. On a donc,
d'après ce qui précède, un \emph{isomorphisme canonique}
\[
  \mathscr{O}_{\mathfrak{X}}/\mathfrak{J}^\Delta
  \xrightarrow{\sim}
  (\widetilde{A}/\widetilde{\mathfrak{J}})|\mathfrak{X}.
  \tag{10.3.2.1}\label{I.10.3.2.1-fr}
\]

Il est clair (en vertu de
$\Gamma(\mathfrak{X},\mathfrak{J}^\Delta)=\mathfrak{J}$) que
l'application $\mathfrak{J}\to\mathfrak{J}^\Delta$ est
\emph{strictement croissante}~; d'après ce qui précède, pour
$\mathfrak{J}\subset\mathfrak{J}'$, le faisceau
${\mathfrak{J}'}^\Delta/\mathfrak{J}^\Delta$ est canoniquement isomorphe à
$\widetilde{\mathfrak{J}'}/\widetilde{\mathfrak{J}}
=\widetilde{\mathfrak{J}'/\mathfrak{J}}$.
\end{env}

\begin{env}[10.3.3]
\label{I.10.3.3-fr}
Les hypothèses et notations étant toujours celles de
(\hyperref[I.10.3.1-fr]{10.3.1}), nous dirons qu'un faisceau d'idéaux
$\mathscr{J}$ de $\mathscr{O}_{\mathfrak{X}}$ est un
\emph{faisceau d'idéaux de définition de $\mathfrak{X}$} (ou un
\emph{Idéal de définition de $\mathfrak{X}$}) si, pour tout
$x\in\mathfrak{X}$, il existe un voisinage ouvert de $x$ de la forme
$\mathfrak{D}(f)$, où $f\in A$, tel que
$\mathscr{J}|\mathfrak{D}(f)$ soit de la forme $\mathfrak{H}^\Delta$, où
$\mathfrak{H}$ est un idéal de définition de $A_{\{f\}}$.
\end{env}

\begin{proposition}[10.3.4]
\label{I.10.3.4-fr}
Pour tout $f\in A$, tout Idéal de définition de $\mathfrak{X}$ induit un
Idéal de définition de $\mathfrak{D}(f)$.
\end{proposition}

\begin{proof}
Cela résulte de (\hyperref[I.10.3.1.1-fr]{10.3.1.1}).
\end{proof}

\begin{proposition}[10.3.5]
\label{I.10.3.5-fr}
Si $A$ est un anneau admissible, tout Idéal de définition de
$\mathfrak{X}=\operatorname{Spf}(A)$ est de la forme
$\mathfrak{J}^\Delta$, où $\mathfrak{J}$ est un idéal de définition de
$A$, uniquement déterminé.
\end{proposition}

\begin{proof}
En effet, soit $\mathscr{J}$ un Idéal de définition de $\mathfrak{X}$~;
par hypothèse, et puisque $\mathfrak{X}$ est quasi-compact, il y a un nombre
fini d'éléments $f_i\in A$ tels que les $\mathfrak{D}(f_i)$ recouvrent
$\mathfrak{X}$ et que
$\mathscr{J}|\mathfrak{D}(f_i)=\mathfrak{H}_i^\Delta$, où
$\mathfrak{H}_i$ est un idéal de définition de $A_{\{f_i\}}$. Pour tout
$i$, il existe donc un idéal ouvert $\mathfrak{K}_i$ de $A$ tel que
$(\mathfrak{K}_i)_{\{f_i\}}=\mathfrak{H}_i$
(\hyperref[0.7.6.9-fr]{0, 7.6.9})~; soit $\mathfrak{K}$ un idéal de
définition de $A$ contenu dans tous les $\mathfrak{K}_i$. L'image canonique
de $\mathscr{J}/\mathfrak{K}^\Delta$ dans le faisceau structural
$\widetilde{A/\mathfrak{K}}$ de
$\operatorname{Spec}(A/\mathfrak{K})$
(\hyperref[I.10.3.2-fr]{10.3.2}) est donc telle que sa restriction à
$\mathfrak{D}(f_i)$ soit égale à celle de
$\widetilde{\mathfrak{K}_i/\mathfrak{K}}$~; on en conclut que cette image
canonique est un faisceau quasi-cohérent sur
$\operatorname{Spec}(A/\mathfrak{K})$, donc de la forme
$\widetilde{\mathfrak{J}/\mathfrak{K}}$, où $\mathfrak{J}$ est un idéal de
$A$ contenant $\mathfrak{K}$ (\hyperref[I.1.4.1-fr]{1.4.1}) d'où
$\mathscr{J}=\mathfrak{J}^\Delta$
(\hyperref[I.10.3.2-fr]{10.3.2})~; en outre, comme pour tout $i$ il existe
un entier $n_i$ tel que
$\mathfrak{H}_i^{n_i}\subset\mathfrak{K}_{\{f_i\}}$, on aura, en désignant
par $n$ le plus grand des $n_i$,
$(\mathscr{J}/\mathfrak{K}^\Delta)^n=0$, et par suite
(\hyperref[I.10.3.2-fr]{10.3.2})
$(\widetilde{\mathfrak{J}/\mathfrak{K}})^n=0$, d'où finalement
$(\mathfrak{J}/\mathfrak{K})^n=0$
(\hyperref[I.1.3.13-fr]{1.3.13}), ce qui prouve que $\mathfrak{J}$ est un
idéal de définition de $A$ (\hyperref[0.7.1.4-fr]{0, 7.1.4}).
\end{proof}

\begin{proposition}[10.3.6]
\label{I.10.3.6-fr}
Soient $A$ un anneau adique, $\mathfrak{J}$ un idéal de définition de $A$
tel que $\mathfrak{J}/\mathfrak{J}^2$ soit un
$A/\mathfrak{J}$-module de type fini. Pour tout entier $n>0$, on a alors
$(\mathfrak{J}^\Delta)^n=(\mathfrak{J}^n)^\Delta$.
\end{proposition}

\begin{proof}
En effet, pour tout $f\in A$, on a (puisque $\mathfrak{J}^n$ est un idéal
ouvert)
\[
  \bigl(\Gamma(\mathfrak{D}(f),\mathfrak{J}^\Delta)\bigr)^n
  =(\mathfrak{J}_{\{f\}})^n
  =(\mathfrak{J}^n)_{\{f\}}
  =\Gamma\bigl(\mathfrak{D}(f^n),(\mathfrak{J}^n)^\Delta\bigr)
\]
\oldpage[I]{185}
en vertu de (\hyperref[I.10.3.1.1-fr]{10.3.1.1}) et de
(\hyperref[0.7.6.12-fr]{0, 7.6.12}). Comme
$(\mathfrak{J}^\Delta)^n$ est associé au préfaisceau
$U\mapsto(\Gamma(U,\mathfrak{J}^\Delta))^n$
(\hyperref[0.4.1.6-fr]{0, 4.1.6}), le corollaire en résulte, puisque les
$\mathfrak{D}(f)$ forment une base de la topologie de $\mathfrak{X}$.
\end{proof}

\begin{env}[10.3.7]
\label{I.10.3.7-fr}
On dit qu'une famille $(\mathscr{J}_\lambda)$ d'Idéaux de définition de
$\mathfrak{X}$ est un \emph{système fondamental d'Idéaux de définition}
si tout Idéal de définition de $\mathfrak{X}$ contient un des
$\mathscr{J}_\lambda$~; comme
$\mathscr{J}_\lambda=\mathfrak{J}_\lambda^\Delta$, il revient au même de
dire que les $\mathfrak{J}_\lambda$ forment un \emph{système fondamental
de voisinages de 0} dans $A$. Soit $(f_\alpha)$ une famille d'éléments de
$A$ tels que les $\mathfrak{D}(f_\alpha)$ recouvrent $\mathfrak{X}$. Si
$(\mathscr{J}_\lambda)$ est une famille filtrante décroissante d'idéaux de
$\mathscr{O}_{\mathfrak{X}}$ telle que pour tout $\alpha$, la famille
$(\mathscr{J}_\lambda|\mathfrak{D}(f_\alpha))$ soit un système fondamental
d'Idéaux de définition de $\mathfrak{D}(f_\alpha)$, alors
$(\mathscr{J}_\lambda)$ est un système fondamental d'Idéaux de définition
de $\mathfrak{X}$. En effet, pour tout Idéal de définition $\mathscr{J}$
de $\mathfrak{X}$, il y a un recouvrement fini de $\mathfrak{X}$ par des
$\mathfrak{D}(f_i)$ tel que, pour tout $i$,
$\mathscr{J}_{\lambda_i}|\mathfrak{D}(f_i)$ soit un Idéal de définition
de $\mathfrak{D}(f_i)$ contenu dans
$\mathscr{J}|\mathfrak{D}(f_i)$. Si $\mu$ est un indice tel que
$\mathscr{J}_\mu\subset\mathscr{J}_{\lambda_i}$ pour tout $i$, il résulte
de (\hyperref[I.10.3.3-fr]{10.3.3}) que $\mathscr{J}_\mu$ est un Idéal de
définition de $\mathfrak{X}$, évidemment contenu dans $\mathscr{J}$, d'où
notre assertion.
\end{env}

\subsection{Préschémas formels et morphismes de préschémas formels.}
\label{subsection:I.10.4-fr}

\begin{env}[10.4.1]
\label{I.10.4.1-fr}
Étant donné un espace topologiquement annelé $\mathfrak{X}$, on dit qu'un
ouvert $U\subset\mathfrak{X}$ est un \emph{ouvert formel affine} (resp. un
\emph{ouvert formel affine adique}, resp. un \emph{ouvert formel affine
noethérien}) si l'espace topologiquement annelé induit par $\mathfrak{X}$
sur $U$ est un schéma formel affine (resp. un tel schéma dont l'anneau est
adique, resp. adique et noethérien).
\end{env}

\begin{definition}[10.4.2]
\label{I.10.4.2-fr}
On appelle préschéma formel un espace topologiquement annelé $\mathfrak{X}$
dont tout point admet un voisinage ouvert formel affine. On dit que le
préschéma formel $\mathfrak{X}$ est adique (resp. localement noethérien) si
tout point de $\mathfrak{X}$ admet un voisinage ouvert formel affine adique
(resp. noethérien). On dit que $\mathfrak{X}$ est noethérien s'il est
localement noethérien et si son espace sous-jacent est quasi-compact (donc
noethérien).
\end{definition}

\begin{proposition}[10.4.3]
\label{I.10.4.3-fr}
Si $\mathfrak{X}$ est un préschéma formel (resp. localement noethérien),
les ensembles ouverts formels affines (resp. affines noethériens) forment
une base de la topologie de $\mathfrak{X}$.
\end{proposition}

\begin{proof}
Cela résulte de (\hyperref[I.10.4.2-fr]{10.4.2}) et
(\hyperref[I.10.1.4-fr]{10.1.4}) en tenant compte de ce que si $A$ est un
anneau adique noethérien, il en est de même de $A_{\{f\}}$ pour tout
$f\in A$ (\hyperref[0.7.6.11-fr]{0, 7.6.11}).
\end{proof}

\begin{corollary}[10.4.4]
\label{I.10.4.4-fr}
Si $\mathfrak{X}$ est un préschéma formel (resp. un préschéma formel
localement noethérien, resp. noethérien), l'espace topologiquement annelé
induit sur tout ouvert de $\mathfrak{X}$ est encore un préschéma formel
(resp. un préschéma formel localement noethérien, resp. noethérien).
\end{corollary}

\begin{definition}[10.4.5]
\label{I.10.4.5-fr}
Étant donnés deux préschémas formels $\mathfrak{X}$, $\mathfrak{Y}$, on
appelle morphisme (de préschémas formels) de $\mathfrak{X}$ dans
$\mathfrak{Y}$ tout morphisme $(\psi,\theta)$ d'espaces topologiquement
annelés tel que, pour tout $x\in\mathfrak{X}$, $\theta_x^\sharp$ soit un
homomorphisme local
$\mathscr{O}_{\psi(x)}\to\mathscr{O}_x$.
\end{definition}

Il est immédiat que le composé de deux morphismes de préschémas formels est
encore un tel morphisme~; les préschémas formels forment donc une
\emph{catégorie}, et on notera
$\operatorname{Hom}(\mathfrak{X},\mathfrak{Y})$ l'ensemble des morphismes
d'un préschéma formel $\mathfrak{X}$ dans un préschéma formel
$\mathfrak{Y}$.

\oldpage[I]{186}
Si $U$ est une partie ouverte de $\mathfrak{X}$, l'injection canonique
dans $\mathfrak{X}$ du préschéma formel induit par $\mathfrak{X}$ sur $U$
est un morphisme de préschémas formels (et même un
\emph{monomorphisme} d'espaces topologiquement annelés
(\hyperref[0.4.1.1-fr]{0, 4.1.1})).

\begin{proposition}[10.4.6]
\label{I.10.4.6-fr}
Soient $\mathfrak{X}$ un préschéma formel,
$\mathfrak{S}=\operatorname{Spf}(A)$ un schéma formel affine. Il existe
une correspondance biunivoque canonique entre les morphismes du préschéma
formel $\mathfrak{X}$ dans le préschéma formel $\mathfrak{S}$ et les
homomorphismes continus de l'anneau $A$ dans l'anneau topologique
$\Gamma(\mathfrak{X},\mathscr{O}_{\mathfrak{X}})$.
\end{proposition}

La démonstration est la même que celle de
(\hyperref[I.2.2.4-fr]{2.2.4}), en remplaçant «~homomorphisme~» par
«~homomorphisme continu~», «~ouvert affine~» par «~ouvert formel
affine~», et en utilisant (\hyperref[I.10.2.2-fr]{10.2.2}) au lieu de
(\hyperref[I.1.7.3-fr]{1.7.3})~; nous en laissons les détails au lecteur.

\begin{env}[10.4.7]
\label{I.10.4.7-fr}
Étant donné un préschéma formel $\mathfrak{S}$, on dit que la donnée d'un
préschéma formel $\mathfrak{X}$ et d'un morphisme
$\varphi:\mathfrak{X}\to\mathfrak{S}$ définit un préschéma formel
$\mathfrak{X}$ \emph{au-dessus de $\mathfrak{S}$} ou un
\emph{$\mathfrak{S}$-préschéma formel}, $\varphi$ étant appelé le
\emph{morphisme structural} du $\mathfrak{S}$-préschéma
$\mathfrak{X}$. Si $\mathfrak{S}=\operatorname{Spf}(A)$, où $A$ est un
anneau admissible, on dit aussi que le $\mathfrak{S}$-préschéma formel
$\mathfrak{X}$ est un \emph{$A$-préschéma formel} ou un préschéma formel
\emph{au-dessus de $A$}. Un préschéma formel quelconque peut toujours
être considéré comme un préschéma formel au-dessus de $\mathbf{Z}$ (muni
de la topologie discrète).

Si $\mathfrak{X}$, $\mathfrak{Y}$ sont deux $\mathfrak{S}$-préschémas
formels, on dit qu'un morphisme
$u:\mathfrak{X}\to\mathfrak{Y}$ est un
\emph{$\mathfrak{S}$-morphisme} si le diagramme
\[
  \xymatrix{
    \mathfrak{X}\ar[rr]^u\ar[rd] & &
    \mathfrak{Y}\ar[ld]\\
    & \mathfrak{S}
  }
\]
(où les flèches obliques sont les morphismes structuraux) est
commutatif. Avec cette définition, les $\mathfrak{S}$-préschémas formels
(pour $\mathfrak{S}$ fixé) forment une \emph{catégorie}. On désigne par
$\operatorname{Hom}_{\mathfrak{S}}(\mathfrak{X},\mathfrak{Y})$
l'ensemble des $\mathfrak{S}$-morphismes du $\mathfrak{S}$-préschéma
formel $\mathfrak{X}$ dans le $\mathfrak{S}$-préschéma formel
$\mathfrak{Y}$. Lorsque $\mathfrak{S}=\operatorname{Spf}(A)$, on dit
aussi \emph{$A$-morphisme} au lieu de
\emph{$\mathfrak{S}$-morphisme}.
\end{env}

\begin{env}[10.4.8]
\label{I.10.4.8-fr}
Comme tout schéma affine peut être considéré comme un schéma formel affine
(\hyperref[I.10.1.2-fr]{10.1.2}), tout préschéma (usuel) peut être
considéré comme un préschéma formel. Il résulte en outre de
(\hyperref[I.10.4.5-fr]{10.4.5}) que pour les préschémas \emph{usuels},
les morphismes (resp. $S$-morphismes) de préschémas \emph{formels}
coïncident avec les morphismes (resp. $S$-morphismes) définis au
\S~\hyperref[section:I.2-fr]{2}.
\end{env}

\subsection{Idéaux de définition des préschémas formels.}
\label{subsection:I.10.5-fr}

\begin{env}[10.5.1]
\label{I.10.5.1-fr}
Soit $\mathfrak{X}$ un préschéma formel~; on dit qu'un
$\mathscr{O}_{\mathfrak{X}}$-Idéal $\mathscr{J}$ est un
\emph{faisceau d'idéaux de définition} (ou un
\emph{Idéal de définition}) de $\mathfrak{X}$ si tout
$x\in\mathfrak{X}$ possède un voisinage ouvert formel affine $U$ tel que
$\mathscr{J}|U$ soit un Idéal de définition du schéma formel affine
induit par $\mathfrak{X}$ sur $U$
(\hyperref[I.10.3.3-fr]{10.3.3})~; en vertu de
(\hyperref[I.10.3.1.1-fr]{10.3.1.1}) et
(\hyperref[I.10.4.3-fr]{10.4.3}), pour tout ouvert
$V\subset\mathfrak{X}$, $\mathscr{J}|V$ est alors un Idéal de définition
du préschéma formel induit par $\mathfrak{X}$ sur $V$.

On dit qu'une famille $(\mathscr{J}_\lambda)$ d'Idéaux de définition de
$\mathfrak{X}$ est un \emph{système fondamental}
\oldpage[I]{187}
\emph{d'Idéaux de définition} s'il existe un recouvrement
$(U_\alpha)$ de $\mathfrak{X}$ par des ouverts formels affines tel que,
pour tout $\alpha$, la famille des
$\mathscr{J}_\lambda|U_\alpha$ soit un système fondamental d'Idéaux de
définition (\hyperref[I.10.3.6-fr]{10.3.6}) du schéma formel affine
induit par $\mathfrak{X}$ sur $U_\alpha$. Il résulte de la remarque finale
de (\hyperref[I.10.3.7-fr]{10.3.7}) que lorsque $\mathfrak{X}$ est un
schéma formel affine, cette définition coïncide avec la définition donnée
dans (\hyperref[I.10.3.7-fr]{10.3.7}). Pour tout ouvert $V$ de
$\mathfrak{X}$, les restrictions $\mathscr{J}_\lambda|V$ forment alors un
système fondamental d'Idéaux de définition du préschéma formel induit sur
$V$, en vertu de (\hyperref[I.10.3.1.1-fr]{10.3.1.1}). Si
$\mathfrak{X}$ est un préschéma formel \emph{localement noethérien}, et
$\mathscr{J}$ un Idéal de définition de $\mathfrak{X}$, il résulte de
(\hyperref[I.10.3.6-fr]{10.3.6}) que les puissances $\mathscr{J}^n$
forment un système fondamental d'Idéaux de définition de $\mathfrak{X}$.
\end{env}

\begin{env}[10.5.2]
\label{I.10.5.2-fr}
Soient $\mathfrak{X}$ un préschéma formel, $\mathscr{J}$ un Idéal de
définition de $\mathfrak{X}$. Alors l'espace annelé
$(\mathfrak{X},\mathscr{O}_{\mathfrak{X}}/\mathscr{J})$ est un
\emph{préschéma} (usuel), qui est affine (resp. localement noethérien,
resp. noethérien) lorsque $\mathfrak{X}$ est un schéma formel affine
(resp. un préschéma formel localement noethérien, resp. noethérien)~; on
est en effet aussitôt ramené au cas affine, et alors la proposition a déjà
été démontrée dans (\hyperref[I.10.3.2-fr]{10.3.2}). En outre, si
$\theta:\mathscr{O}_{\mathfrak{X}}\to
\mathscr{O}_{\mathfrak{X}}/\mathscr{J}$ est l'homomorphisme canonique,
$u=(1_{\mathfrak{X}},\theta)$ est un \emph{morphisme} (dit
\emph{canonique}) de préschémas formels
$(\mathfrak{X},\mathscr{O}_{\mathfrak{X}}/\mathscr{J})\to
(\mathfrak{X},\mathscr{O}_{\mathfrak{X}})$, car ici encore, cela a été vu
dans le cas affine (\hyperref[I.10.3.2-fr]{10.3.2}), auquel on se ramène
aussitôt.
\end{env}

\begin{proposition}[10.5.3]
\label{I.10.5.3-fr}
Soient $\mathfrak{X}$ un préschéma formel, $(\mathscr{J}_\lambda)$ un
système fondamental d'Idéaux de définition de $\mathfrak{X}$. Alors le
faisceau d'anneaux topologiques $\mathscr{O}_{\mathfrak{X}}$ est limite
projective des faisceaux d'anneaux pseudo-discrets
(\hyperref[0.3.8.1-fr]{0, 3.8.1})
$\mathscr{O}_{\mathfrak{X}}/\mathscr{J}_\lambda$.
\end{proposition}

\begin{proof}
Comme la topologie de $\mathfrak{X}$ admet une base d'ouverts formels
affines quasi-compacts (\hyperref[I.10.4.3-fr]{10.4.3}), on est ramené au
cas affine, où la proposition est conséquence de
(\hyperref[I.10.3.5-fr]{10.3.5}),
(\hyperref[I.10.3.2-fr]{10.3.2}) et de la définition
(\hyperref[I.10.1.1-fr]{10.1.1}).
\end{proof}

Il n'est pas certain que tout préschéma formel admette des Idéaux de
définition. Toutefois~:

\begin{proposition}[10.5.4]
\label{I.10.5.4-fr}
Soit $\mathfrak{X}$ un préschéma formel localement noethérien. Il existe
un plus grand Idéal de définition $\mathscr{T}$ de $\mathfrak{X}$~;
c'est le seul Idéal de définition $\mathscr{J}$ tel que le préschéma
$(\mathfrak{X},\mathscr{O}_{\mathfrak{X}}/\mathscr{J})$ soit réduit. Si
$\mathscr{J}$ est un Idéal de définition de $\mathfrak{X}$,
$\mathscr{T}$ est l'image réciproque par
$\mathscr{O}_{\mathfrak{X}}\to
\mathscr{O}_{\mathfrak{X}}/\mathscr{J}$ du Nilradical de
$\mathscr{O}_{\mathfrak{X}}/\mathscr{J}$.
\end{proposition}

\begin{proof}
Supposons d'abord que $\mathfrak{X}=\operatorname{Spf}(A)$, où $A$ est
un anneau adique noethérien. L'existence et les propriétés de
$\mathscr{T}$ résultent immédiatement de
(\hyperref[I.10.3.4-fr]{10.3.4}) et
(\hyperref[I.5.1.1-fr]{5.1.1}), compte tenu de l'existence et des
propriétés du plus grand idéal de définition de $A$
(\hyperref[0.7.1.6-fr]{0, 7.1.6} et
\hyperref[0.7.1.7-fr]{7.1.7}).

Pour prouver l'existence et les propriétés de $\mathscr{T}$ dans le cas
général, il suffit de montrer que si $U\supset V$ sont deux ouverts
formels affines noethériens de $\mathfrak{X}$, le plus grand Idéal de
définition $\mathscr{T}_U$ de $U$ induit le plus grand Idéal de définition
$\mathscr{T}_V$ de $V$~; mais comme
$(V,(\mathscr{O}_{\mathfrak{X}}|V)/(\mathscr{T}_U|V))$ est réduit, cela
résulte de ce qui précède.
\end{proof}

On désigne par $\mathfrak{X}_{\mathrm{red}}$ le préschéma (usuel) réduit
$(\mathfrak{X},\mathscr{O}_{\mathfrak{X}}/\mathscr{T})$.

\begin{corollary}[10.5.5]
\label{I.10.5.5-fr}
Soient $\mathfrak{X}$ un préschéma formel localement noethérien,
$\mathscr{T}$ le plus grand Idéal de définition de $\mathfrak{X}$~; pour
tout ouvert $V$ de $\mathfrak{X}$, $\mathscr{T}|V$ est le plus grand
Idéal de définition du préschéma formel induit par $\mathfrak{X}$ sur
$V$.
\end{corollary}

\oldpage[I]{188}

\begin{proposition}[10.5.6]
\label{I.10.5.6-fr}
Soient $\mathfrak{X}$, $\mathfrak{Y}$ deux préschémas formels,
$\mathscr{J}$ (resp. $\mathscr{K}$) un Idéal de définition de
$\mathfrak{X}$ (resp. $\mathfrak{Y}$),
$f:\mathfrak{X}\to\mathfrak{Y}$ un morphisme de préschémas formels.
\begin{enumerate}
  \item[\rm{(i)}] Si
    $f^*(\mathscr{K})\mathscr{O}_{\mathfrak{X}}\subset\mathscr{J}$,
    il existe un unique morphisme
    $f':(\mathfrak{X},\mathscr{O}_{\mathfrak{X}}/\mathscr{J})\to
    (\mathfrak{Y},\mathscr{O}_{\mathfrak{Y}}/\mathscr{K})$ de
    préschémas usuels rendant commutatif le diagramme
    \[
      \label{I.10.5.6.1-fr}
      \xymatrix{
        (\mathfrak{X},\mathscr{O}_{\mathfrak{X}})\ar[r]^f &
        (\mathfrak{Y},\mathscr{O}_{\mathfrak{Y}})\\
        (\mathfrak{X},\mathscr{O}_{\mathfrak{X}}/\mathscr{J})
          \ar[r]_{f'}\ar[u] &
        (\mathfrak{Y},\mathscr{O}_{\mathfrak{Y}}/\mathscr{K})\ar[u]
      }
      \tag{10.5.6.1}
    \]
    où les flèches verticales sont les morphismes canoniques.
  \item[\rm{(ii)}] Supposons que
    $\mathfrak{X}=\operatorname{Spf}(A)$,
    $\mathfrak{Y}=\operatorname{Spf}(B)$ soient des schémas formels
    affines, $\mathscr{J}=\mathfrak{J}^{\Delta}$,
    $\mathscr{K}=\mathfrak{K}^{\Delta}$, où $\mathfrak{J}$ (resp.
    $\mathfrak{K}$) est un idéal de définition de $A$ (resp. $B$), et
    $f=({}^a\varphi,\widetilde{\varphi})$, où
    $\varphi:B\to A$ est un homomorphisme continu~; pour que
    $f^*(\mathscr{K})\mathscr{O}_{\mathfrak{X}}\subset\mathscr{J}$,
    il faut et il suffit que
    $\varphi(\mathfrak{K})\subset\mathfrak{J}$, et $f'$ est alors le
    morphisme $({}^a\varphi',\widetilde{\varphi'})$, où
    $\varphi':B/\mathfrak{K}\to A/\mathfrak{J}$ est l'homomorphisme
    déduit de $\varphi$ par passage aux quotients.
\end{enumerate}
\end{proposition}

\begin{proof}
\medskip\noindent
\begin{enumerate}
  \item[\rm{(i)}] Si $f=(\psi,\theta)$, l'hypothèse entraîne que
    l'image par
    $\theta^\sharp:\psi^*(\mathscr{O}_{\mathfrak{Y}})\to
    \mathscr{O}_{\mathfrak{X}}$ du faisceau d'idéaux
    $\psi^*(\mathscr{K})$ de $\psi^*(\mathscr{O}_{\mathfrak{Y}})$ est
    contenue dans $\mathscr{J}$
    (\hyperref[0.4.3.5-fr]{0, 4.3.5}). Par passage aux quotients, on
    déduit donc de $\theta^\sharp$ un homomorphisme de faisceau d'anneaux
    \[
      \omega:\psi^*(\mathscr{O}_{\mathfrak{Y}}/\mathscr{K})
      =\psi^*(\mathscr{O}_{\mathfrak{Y}})/\psi^*(\mathscr{K})
      \to\mathscr{O}_{\mathfrak{X}}/\mathscr{J}~;
    \]
    en outre, comme pour tout $x\in\mathfrak{X}$,
    $\theta_x^\sharp$ est un homomorphisme \emph{local}, il en est de
    même de $\omega_x$. Le morphisme d'espaces annelés
    $(\psi,\omega^b)$ est donc
    (\hyperref[I.2.2.1-fr]{2.2.1}) l'unique morphisme $f'$ d'espaces
    annelés répondant à la question.
  \item[\rm{(ii)}] La correspondance canonique fonctorielle entre
    morphismes de préschémas formels affines et homomorphismes continus
    d'anneaux (\hyperref[I.10.2.2-fr]{10.2.2}) montre que dans le cas
    considéré, la relation
    $f^*(\mathscr{K})\mathscr{O}_{\mathfrak{X}}\subset\mathscr{J}$
    entraîne que l'on a
    $f'=({}^a\varphi',\widetilde{\varphi'})$, où
    $\varphi':B/\mathfrak{K}\to A/\mathfrak{J}$ est l'unique
    homomorphisme rendant commutatif le diagramme
    \[
      \label{I.10.5.6.2-fr}
      \xymatrix{
        B\ar[r]^\varphi\ar[d] & A\ar[d]\\
        B/\mathfrak{K}\ar[r]_{\varphi'} & A/\mathfrak{J}
      }
      \tag{10.5.6.2}
    \]
    L'existence de $\varphi'$ implique donc que
    $\varphi(\mathfrak{K})\subset\mathfrak{J}$. Inversement, si cette
    condition est vérifiée, en désignant par $\varphi'$ l'unique
    homomorphisme rendant commutatif le diagramme
    (\hyperref[I.10.5.6.2-fr]{10.5.6.2}) et posant
    $f'=({}^a\varphi',\widetilde{\varphi'})$, il est clair que le
    diagramme (\hyperref[I.10.5.6.1-fr]{10.5.6.1}) est commutatif~; la
    considération des homomorphismes
    ${}^a\varphi^*(\mathscr{O}_{\mathfrak{Y}})\to
    \mathscr{O}_{\mathfrak{X}}$ et
    ${}^a{\varphi'}^*(\mathscr{O}_{\mathfrak{Y}}/\mathscr{K})\to
    \mathscr{O}_{\mathfrak{X}}/\mathscr{J}$ correspondant à $f$ et
    $f'$ respectivement, montre alors que cela entraîne la relation
    $f^*(\mathscr{K})\mathscr{O}_{\mathfrak{X}}\subset\mathscr{J}$.
\end{enumerate}
\end{proof}

Il est clair que la correspondance $f\mapsto f'$ définie ci-dessus est
\emph{fonctorielle}.

\subsection{Préschémas formels comme limites inductives de préschémas.}
\label{subsection:I.10.6-fr}

\begin{env}[10.6.1]
\label{I.10.6.1-fr}
Soient $\mathfrak{X}$ un préschéma formel,
$(\mathscr{J}_\lambda)$ un système fondamental d'Idéaux de définition de
$\mathfrak{X}$~; pour tout $\lambda$, soit $f_\lambda$ le morphisme
canonique
$(\mathfrak{X},\mathscr{O}_{\mathfrak{X}}/\mathscr{J}_\lambda)
\to\mathfrak{X}$ (\hyperref[I.10.5.2-fr]{10.5.2})~; pour
$\mathscr{J}_\mu\subset\mathscr{J}_\lambda$, l'homomorphisme canonique
$\mathscr{O}_{\mathfrak{X}}/\mathscr{J}_\mu\to
\mathscr{O}_{\mathfrak{X}}/\mathscr{J}_\lambda$ définit un morphisme

\oldpage[I]{189}

canonique
$f_{\mu\lambda}:(\mathfrak{X},
\mathscr{O}_{\mathfrak{X}}/\mathscr{J}_\lambda)\to
(\mathfrak{X},\mathscr{O}_{\mathfrak{X}}/\mathscr{J}_\mu)$ de préschémas
(usuels) tel que l'on ait $f_\lambda=f_\mu\circ f_{\mu\lambda}$. Les
préschémas
$X_\lambda=(\mathfrak{X},
\mathscr{O}_{\mathfrak{X}}/\mathscr{J}_\lambda)$ et les morphismes
$f_{\mu\lambda}$ constituent donc (en vertu de
\hyperref[I.10.4.8-fr]{10.4.8}) un \emph{système inductif} dans la
catégorie des préschémas formels.
\end{env}

\begin{proposition}[10.6.2]
\label{I.10.6.2-fr}
Avec les notations de (\hyperref[I.10.6.1-fr]{10.6.1}), le préschéma
formel $\mathfrak{X}$ et les morphismes $f_\lambda$ constituent une
limite inductive (T, I, 1.8) du système $(X_\lambda,f_{\mu\lambda})$
dans la catégorie des préschémas formels.
\end{proposition}

\begin{proof}
Soit $\mathfrak{Y}$ un préschéma formel, et pour chaque indice $\lambda$,
soit
\[
  g_\lambda=(\psi_\lambda,\theta_\lambda):X_\lambda\to\mathfrak{Y}
\]
un morphisme, tel que l'on ait
$g_\lambda=g_\mu\circ f_{\mu\lambda}$ pour
$\mathscr{J}_\mu\subset\mathscr{J}_\lambda$. Cette dernière condition
et la définition des $X_\lambda$ entraînent d'abord que tous les
$\psi_\lambda$ sont identiques à une même application continue
$\psi:\mathfrak{X}\to\mathfrak{Y}$ des espaces sous-jacents~; en outre,
les homomorphismes
$\theta_\lambda^\sharp:\psi^*(\mathscr{O}_{\mathfrak{Y}})\to
\mathscr{O}_{X_\lambda}=
\mathscr{O}_{\mathfrak{X}}/\mathscr{J}_\lambda$ forment un
\emph{système projectif} d'homomorphismes de faisceaux d'anneaux. Par
passage à la limite projective, on en déduit donc un homomorphisme
\[
  \omega:\psi^*(\mathscr{O}_{\mathfrak{Y}})\to
  \varprojlim\mathscr{O}_{\mathfrak{X}}/\mathscr{J}_\lambda
  =\mathscr{O}_{\mathfrak{X}},
\]
et il est clair que le morphisme $g=(\psi,\omega^b)$ d'\emph{espaces
annelés} est le \emph{seul} rendant commutatif les diagrammes
\[
  \label{I.10.6.2.1-fr}
  \xymatrix{
    X_\lambda\ar[rr]^{g_\lambda}\ar[rd]_{f_\lambda} &&
    \mathfrak{Y}\\
    & \mathfrak{X}\ar[ru]_g
  }
  \tag{10.6.2.1}
\]
Il reste donc à prouver que $g$ est un morphisme de \emph{préschémas
formels}~; la question étant locale sur $\mathfrak{X}$ et
$\mathfrak{Y}$, on peut donc supposer
$\mathfrak{X}=\operatorname{Spf}(A)$,
$\mathfrak{Y}=\operatorname{Spf}(B)$, $A$ et $B$ étant des anneaux
admissibles, avec
$\mathscr{J}_\lambda=\mathfrak{J}_\lambda^{\Delta}$, où
$(\mathfrak{J}_\lambda)$ est un système fondamental d'idéaux de
définition de $A$ (\hyperref[I.10.3.5-fr]{10.3.5})~; comme
$A=\varprojlim A/\mathfrak{J}_\lambda$, l'existence d'un morphisme de
schémas formels affines $g$ rendant commutatifs les diagrammes
(\hyperref[I.10.6.2.1-fr]{10.6.2.1}) résulte alors de la correspondance
biunivoque (\hyperref[I.10.2.2-fr]{10.2.2}) entre morphismes de schémas
formels affines et homomorphismes continus d'anneaux, et de la définition
de la limite projective. Mais l'unicité de $g$ en tant que morphisme
d'espaces annelés montre qu'il coïncide avec le morphisme noté de même au
début de la démonstration.
\end{proof}

La proposition suivante établit, sous certaines conditions
supplémentaires, l'existence de la limite inductive d'un système inductif
donné de préschémas (usuels) dans la catégorie des préschémas formels~:

\begin{proposition}[10.6.3]
\label{I.10.6.3-fr}
Soient $\mathfrak{X}$ un espace topologique,
$(\mathscr{O}_i,u_{ji})$ un système projectif de faisceaux d'anneaux sur
$\mathfrak{X}$, ayant $\mathbf{N}$ pour ensemble d'indices. Soit
$\mathscr{J}_i$ le noyau de
$u_{0i}:\mathscr{O}_i\to\mathscr{O}_0$. On suppose que~:
\begin{enumerate}
  \item[a)] L'espace annelé $(\mathfrak{X},\mathscr{O}_i)$ est un
    préschéma $X_i$.
  \item[b)] Pour tout $x\in X$ et tout $i$, il existe un voisinage
    ouvert $U_i$ de $x$ dans $\mathfrak{X}$ tel que la restriction
    $\mathscr{J}_i|U_i$ soit nilpotente.
  \item[c)] Les homomorphismes $u_{ji}$ sont surjectifs.
\end{enumerate}

\oldpage[I]{190}

Soit $\mathscr{O}_{\mathfrak{X}}$ le faisceau d'anneaux topologiques
limite projective des faisceaux d'anneaux pseudo-discrets
$\mathscr{O}_i$, et soit
$u_i:\mathscr{O}_{\mathfrak{X}}\to\mathscr{O}_i$ l'homomorphisme
canonique. Alors l'espace topologiquement annelé
$(\mathfrak{X},\mathscr{O}_{\mathfrak{X}})$ est un préschéma formel~;
les homomorphismes $u_i$ sont surjectifs~; leurs noyaux
$\mathscr{J}^{(i)}$ forment un système fondamental d'Idéaux de définition
de $\mathfrak{X}$, et $\mathscr{J}^{(0)}$ est la limite projective des
faisceaux d'idéaux $\mathscr{J}_i$.
\end{proposition}

\begin{proof}
Notons d'abord que sur chaque fibre, $u_{ji}$ est un homomorphisme
surjectif et \emph{a fortiori} un homomorphisme local~; donc
$v_{ij}=(1_{\mathfrak{X}},u_{ji})$ est un morphisme de préschémas
$X_j\to X_i$ ($i\geq j$) (\hyperref[I.2.2.1-fr]{2.2.1}). Supposons
d'abord que chaque $X_i$ soit un schéma affine d'anneau $A_i$. Il existe
un homomorphisme \emph{d'anneaux}
$\varphi_{ji}:A_i\to A_j$ tel que
$u_{ji}=\widetilde{\varphi_{ji}}$
(\hyperref[I.1.7.3-fr]{1.7.3})~; par suite
(\hyperref[I.1.6.3-fr]{1.6.3}), le faisceau $\mathscr{O}_j$ est un
$\mathscr{O}_i$-Module quasi-cohérent sur $X_i$ (pour la loi externe
définie par $u_{ji}$), associé à $A_j$ considéré comme $A_i$-module à
l'aide de $\varphi_{ji}$. Pour tout $f\in A_i$, soit
$f'=\varphi_{ji}(f)$~; par hypothèse, les ouverts $D(f)$ et $D(f')$ sont
identiques dans $\mathfrak{X}$, et l'homomorphisme de
$\Gamma(D(f),\mathscr{O}_i)=(A_i)_f$ dans
$\Gamma(D(f),\mathscr{O}_j)=(A_j)_{f'}$ correspondant à $u_{ji}$ n'est
autre que $(\varphi_{ji})_f$ (\hyperref[I.1.6.1-fr]{1.6.1}). Mais
lorsqu'on considère $A_j$ comme $A_i$-module, $(A_j)_{f'}$ est le
$(A_i)_f$-module $(A_j)_f$, donc on a aussi
$u_{ji}=\widetilde{\varphi_{ji}}$ lorsque $\varphi_{ji}$ est cette fois
considéré comme homomorphisme de $A_i$-modules. Alors, comme $u_{ji}$ est
surjectif, on en conclut que $\varphi_{ji}$ l'est aussi
(\hyperref[I.1.3.9-fr]{1.3.9}) et si $\mathfrak{J}_{ji}$ est le noyau de
$\varphi_{ji}$, le noyau de $u_{ji}$ est un $\mathscr{O}_i$-Module
quasi-cohérent égal à $\widetilde{\mathfrak{J}_{ji}}$. En particulier,
on a $\mathscr{J}_i=\widetilde{\mathfrak{J}_i}$, où $\mathfrak{J}_i$ est
le noyau de $\varphi_{0i}:A_i\to A_0$. L'hypothèse b) entraîne que
$\mathscr{J}_i$ est \emph{nilpotent}~: en effet, comme $\mathfrak{X}$
est quasi-compact, on peut recouvrir $\mathfrak{X}$ par un nombre fini
d'ouverts $U_k$ tels que $(\mathscr{J}_i|U_k)^{n_k}=0$ et en prenant pour
$n$ le plus grand des $n_k$, on a $\mathscr{J}_i^n=0$. On en conclut que
$\mathfrak{J}_i$ est nilpotent (\hyperref[I.1.3.13-fr]{1.3.13}). Alors
l'anneau $A=\varprojlim A_i$ est admissible
(\hyperref[0.7.2.2-fr]{0, 7.2.2}), l'homomorphisme canonique
$\varphi_i:A\to A_i$ est surjectif et son noyau
$\mathfrak{J}^{(i)}$ est égal à la limite projective des
$\mathfrak{J}_{ik}$ pour $k\geq i$~; les $\mathfrak{J}^{(i)}$ forment
un système fondamental de voisinages de 0 dans $A$. Les assertions de
(\hyperref[I.10.6.3-fr]{10.6.3}) résultent dans ce cas de
(\hyperref[I.10.1.1-fr]{10.1.1}) et
(\hyperref[I.10.3.2-fr]{10.3.2}),
$(\mathfrak{X},\mathscr{O}_{\mathfrak{X}})$ n'étant autre que
$\operatorname{Spf}(A)$.

Toujours dans ce même cas particulier, notons que si $f=(f_i)$ est un
élément de la limite projective $A=\varprojlim A_i$, tous les ouverts
$D(f_i)$ (ouvert affine dans $X_i$) s'identifient à l'ouvert
$\mathfrak{D}(f)$ de $\mathfrak{X}$, le préschéma induit par $X_i$ sur
$\mathfrak{D}(f)$ s'identifiant donc au schéma affine
$\operatorname{Spec}((A_i)_{f_i})$.

Dans le cas général, remarquons d'abord que pour tout ouvert quasi-compact
$U$ de $\mathfrak{X}$, chacun des $\mathscr{J}_i|U$ est nilpotent, comme
le montre le raisonnement fait ci-dessus. Nous allons voir que pour tout
$x\in\mathfrak{X}$, il y a un voisinage ouvert $U$ de $x$ dans
$\mathfrak{X}$ qui est un \emph{ouvert affine pour tous les $X_i$}. En
effet, prenons $U$ ouvert affine pour $X_0$, et observons que
$\mathscr{O}_{X_0}=\mathscr{O}_{X_i}/\mathscr{J}_i$. Comme
$\mathscr{J}_i|U$ est nilpotent, en vertu de ce qui précède, $U$ est
ouvert affine aussi pour chaque $X_i$ en vertu de
(\hyperref[I.5.1.9-fr]{5.1.9}). Cela étant, pour tout $U$ satisfaisant
aux conditions précédentes, l'étude du cas affine faite plus haut montre
que $(U,\mathscr{O}_{\mathfrak{X}}|U)$ est un préschéma formel dont les
$\mathscr{J}^{(i)}|U$ forment un système fondamental d'idéaux de
définition et $\mathscr{J}^{(0)}|U$ est limite projective des
$\mathscr{J}_i|U$~; d'où la conclusion.
\end{proof}

\begin{corollary}[10.6.4]
\label{I.10.6.4-fr}
Supposons que pour $i\geq j$, le noyau de $u_{ji}$ soit
$\mathscr{J}_i^{j+1}$ et que $\mathscr{J}_1/\mathscr{J}_1^2$

\oldpage[I]{191}

soit de type fini sur
$\mathscr{O}_0=\mathscr{O}_1/\mathscr{J}_1$. Alors $\mathfrak{X}$ est
un préschéma formel adique, et si $\mathscr{J}^{(n)}$ est le noyau de
$\mathscr{O}_{\mathfrak{X}}\to\mathscr{O}_n$, on a
$\mathscr{J}^{(n)}=\mathscr{J}^{n+1}$ et
$\mathscr{J}/\mathscr{J}^2$ est isomorphe à $\mathscr{J}_1$. Si en
outre $X_0$ est localement noethérien (resp. noethérien),
$\mathfrak{X}$ est localement noethérien (resp. noethérien).
\end{corollary}

\begin{proof}
Comme les espaces sous-jacents à $\mathfrak{X}$ et $X_0$ sont les mêmes,
la question est locale et l'on peut supposer tous les $X_i$ affines~;
compte tenu des relations
$\mathscr{J}_{ij}=\widetilde{\mathfrak{J}_{ji}}$ (avec les notations
de (\hyperref[I.10.6.3-fr]{10.6.3})), on est aussitôt ramené aux
assertions correspondantes de
(\hyperref[0.7.2.7-fr]{0, 7.2.7} et
\hyperref[0.7.2.8-fr]{7.2.8}), en notant que
$\mathfrak{J}_1/\mathfrak{J}_1^2$ est alors un $A_0$-module de type
fini (\hyperref[I.1.3.9-fr]{1.3.9}).
\end{proof}

En particulier, \emph{tout préschéma formel localement noethérien
$\mathfrak{X}$} est limite inductive d'une suite $(X_n)$ de préschémas
(usuels) localement noethériens vérifiant les conditions de
(\hyperref[I.10.6.3-fr]{10.6.3}) et
(\hyperref[I.10.6.4-fr]{10.6.4})~: il suffit de considérer un Idéal de
définition $\mathscr{J}$ de $\mathfrak{X}$
(\hyperref[I.10.5.4-fr]{10.5.4}) et de prendre
$X_n=(\mathfrak{X},
\mathscr{O}_{\mathfrak{X}}/\mathscr{J}^{n+1})$
((\hyperref[I.10.5.1-fr]{10.5.1}) et
(\hyperref[I.10.6.2-fr]{10.6.2})).

\begin{corollary}[10.6.5]
\label{I.10.6.5-fr}
Soit $A$ un anneau admissible. Pour que le schéma formel affine
$\mathfrak{X}=\operatorname{Spf}(A)$ soit noethérien, il faut et il
suffit que $A$ soit adique et noethérien.
\end{corollary}

\begin{proof}
La condition est évidemment suffisante. Inversement, supposons que
$\mathfrak{X}$ soit noethérien, et soient $\mathfrak{J}$ un idéal de
définition de $A$, $\mathscr{J}=\mathfrak{J}^\Delta$ l'Idéal de
définition correspondant de $\mathfrak{X}$. Les préschémas (usuels)
$X_n=(\mathfrak{X},
\mathscr{O}_{\mathfrak{X}}/\mathscr{J}^{n+1})$ sont alors affines et
noethériens, donc les anneaux
$A_n=A/\mathfrak{J}^{n+1}$ sont noethériens
(\hyperref[I.6.1.3-fr]{6.1.3}), d'où on conclut que
$\mathfrak{J}/\mathfrak{J}^2$ est un $A/\mathfrak{J}$-module de type
fini. Comme les $\mathscr{J}^n$ forment un système fondamental d'Idéaux
de définition de $\mathfrak{X}$
(\hyperref[I.10.5.1-fr]{10.5.1}), on a
$\mathscr{O}_{\mathfrak{X}}
=\varprojlim(\mathscr{O}_{\mathfrak{X}}/\mathscr{J}^n)$
(\hyperref[I.10.5.3-fr]{10.5.3})~; on en conclut
(\hyperref[I.10.1.3-fr]{10.1.3}) que $A$ est topologiquement isomorphe à
$\varprojlim A/\mathfrak{J}^n$, donc est adique et noethérien
(\hyperref[0.7.2.8-fr]{0, 7.2.8}).
\end{proof}

\begin{remark}[10.6.6]
\label{I.10.6.6-fr}
Avec les notations de (\hyperref[I.10.6.3-fr]{10.6.3}), soit
$\mathscr{F}_i$ un $\mathscr{O}_i$-Module, et supposons donné, pour
$i\geq j$, un $v_{ij}$-morphisme
$\theta_{ji}:\mathscr{F}_i\to\mathscr{F}_j$, de sorte que
$\theta_{kj}\circ\theta_{ji}=\theta_{ki}$ pour $k\leq j\leq i$.
Comme l'application continue sous-jacente à $v_{ij}$ est l'identité,
$\theta_{ji}$ est un homomorphisme de faisceaux de groupes abéliens sur
l'espace $\mathfrak{X}$~; en outre, si $\mathscr{F}$ est la limite
projective du système projectif $(\mathscr{F}_i)$ de faisceaux de groupes
abéliens, le fait que les $\theta_{ji}$ sont des $v_{ij}$-morphismes
permet de définir sur $\mathscr{F}$ une structure de
$\mathscr{O}_{\mathfrak{X}}$-Module par passage à la limite projective~;
muni de cette structure, nous dirons que $\mathscr{F}$ est la
\emph{limite projective} (pour les $\theta_{ji}$) du système de
$\mathscr{O}_i$-Modules $(\mathscr{F}_i)$. Dans le cas particulier où
$v_{ij}^*(\mathscr{F}_i)=\mathscr{F}_j$ et où $\theta_{ji}$ est
l'\emph{identité}, nous dirons pour abréger, que $\mathscr{F}$ est la
limite projective d'un système $(\mathscr{F}_i)$ tel que
$v_{ij}^*(\mathscr{F}_i)=\mathscr{F}_j$ pour $j\leq i$ (sans
mentionner les $\theta_{ji}$).
\end{remark}

\begin{env}[10.6.7]
\label{I.10.6.7-fr}
Soient $\mathfrak{X}$, $\mathfrak{Y}$ deux préschémas formels,
$\mathscr{J}$ (resp. $\mathscr{K}$) un Idéal de définition de
$\mathfrak{X}$ (resp. $\mathfrak{Y}$),
$f:\mathfrak{X}\to\mathfrak{Y}$ un morphisme tel que
$f^*(\mathscr{K})\mathscr{O}_{\mathfrak{X}}\subset\mathscr{J}$. On a
alors pour tout entier $n>0$,
$f^*(\mathscr{K}^n)\mathscr{O}_{\mathfrak{X}}
=(f^*(\mathscr{K})\mathscr{O}_{\mathfrak{X}})^n
\subset\mathscr{J}^n$~; on peut donc
(\hyperref[I.10.5.6-fr]{10.5.6}) déduire de $f$ un morphisme de
préschémas (usuels) $f_n:X_n\to Y_n$, en posant
$X_n=(\mathfrak{X},
\mathscr{O}_{\mathfrak{X}}/\mathscr{J}^{n+1})$,
$Y_n=(\mathfrak{Y},
\mathscr{O}_{\mathfrak{Y}}/\mathscr{K}^{n+1})$, et il résulte
aussitôt des définitions que les diagrammes
\[
  \label{I.10.6.7.1-fr}
  \xymatrix{
    X_m\ar[r]^{f_m}\ar[d] &
    Y_m\ar[d]\\
    X_n\ar[r]_{f_n} &
    Y_n
  }
  \tag{10.6.7.1}
\]

\oldpage[I]{192}

sont commutatifs pour $m\leq n$~; autrement dit, $(f_n)$ est un
\emph{système inductif} de morphismes.
\end{env}

\begin{env}[10.6.8]
\label{I.10.6.8-fr}
Inversement, soit $(X_n)$ (resp. $(Y_n)$) un système inductif de
préschémas (usuels) satisfaisant aux conditions b) et c) de
(\hyperref[I.10.6.3-fr]{10.6.3}), et soit $\mathfrak{X}$ (resp.
$\mathfrak{Y}$) sa limite inductive. Par définition des limites
inductives, toute suite $(f_n)$ de morphismes $X_n\to Y_n$ formant un
système inductif admet une \emph{limite inductive}
$f:\mathfrak{X}\to\mathfrak{Y}$, qui est l'unique morphisme de
préschémas formels rendant commutatifs les diagrammes
\[
  \xymatrix{
    X_n\ar[r]^{f_n}\ar[d] &
    Y_n\ar[d]\\
    \mathfrak{X}\ar[r]_f &
    \mathfrak{Y}
  }
\]
\end{env}

\begin{proposition}[10.6.9]
\label{I.10.6.9-fr}
Soient $\mathfrak{X}$, $\mathfrak{Y}$ deux préschémas formels
localement noethériens, $\mathscr{J}$ (resp. $\mathscr{K}$) un Idéal
de définition de $\mathfrak{X}$ (resp. $\mathfrak{Y}$)~;
l'application $f\mapsto(f_n)$ définie dans
(\hyperref[I.10.6.7-fr]{10.6.7}) est une bijection de l'ensemble des
morphismes $f:\mathfrak{X}\to\mathfrak{Y}$ tels que
$f^*(\mathscr{K})\mathscr{O}_{\mathfrak{X}}\subset\mathscr{J}$, sur
l'ensemble des suites $(f_n)$ de morphismes rendant commutatifs les
diagrammes (\hyperref[I.10.6.7.1-fr]{10.6.7.1}).
\end{proposition}

\begin{proof}
Si $f$ est la limite inductive d'une telle suite, il faut montrer que
$f^*(\mathscr{K})\mathscr{O}_{\mathfrak{X}}\subset\mathscr{J}$. La
question étant locale sur $\mathfrak{X}$ et $\mathfrak{Y}$, on peut se
borner au cas où $\mathfrak{X}=\operatorname{Spf}(A)$,
$\mathfrak{Y}=\operatorname{Spf}(B)$ sont affines, $A$ et $B$ étant
adiques noethériens,
$\mathscr{J}=\mathfrak{J}^\Delta$,
$\mathscr{K}=\mathfrak{K}^\Delta$, où $\mathfrak{J}$ (resp.
$\mathfrak{K}$) est un idéal de définition de $A$ (resp. $B$). On a
alors $X_n=\operatorname{Spec}(A_n)$,
$Y_n=\operatorname{Spec}(B_n)$, avec
$A_n=A/\mathfrak{J}^{n+1}$ et $B_n=B/\mathfrak{K}^{n+1}$, en vertu de
(\hyperref[I.10.3.6-fr]{10.3.6}) et
(\hyperref[I.10.3.2-fr]{10.3.2})~;
$f_n=({}^a\varphi_n,\widetilde{\varphi_n})$, où les homomorphismes
$\varphi_n:B_n\to A_n$ forment un système projectif, donc
$f=({}^a\varphi,\widetilde{\varphi})$, où
$\varphi=\varprojlim\varphi_n$. La commutativité du diagramme
(\hyperref[I.10.6.7.1-fr]{10.6.7.1}) pour $m=0$ donne alors la
condition
$\varphi_n(\mathfrak{K}/\mathfrak{K}^{n+1})
\subset\mathfrak{J}/\mathfrak{J}^{n+1}$ pour tout $n$, donc, en
passant à la limite projective, $\varphi(\mathfrak{K})\subset
\mathfrak{J}$, et cela entraîne
$f^*(\mathscr{K})\mathscr{O}_{\mathfrak{X}}\subset\mathscr{J}$
(\hyperref[I.10.5.6-fr]{10.5.6, (ii)}).
\end{proof}

\begin{corollary}[10.6.10]
\label{I.10.6.10-fr}
Soient $\mathfrak{X}$, $\mathfrak{Y}$ deux préschémas formels
localement noethériens, $\mathscr{T}$ le plus grand Idéal de définition
de $\mathfrak{X}$ (\hyperref[I.10.5.4-fr]{10.5.4}).
\begin{enumerate}
  \item[\rm{(i)}] Pour tout Idéal de définition $\mathscr{K}$ de
  $\mathfrak{Y}$ et tout morphisme
  $f:\mathfrak{X}\to\mathfrak{Y}$, on a
  $f^*(\mathscr{K})\mathscr{O}_{\mathfrak{X}}\subset\mathscr{T}$.
  \item[\rm{(ii)}] Il y a correspondance biunivoque canonique entre
  $\operatorname{Hom}(\mathfrak{X},\mathfrak{Y})$ et l'ensemble des
  suites $(f_n)$ de morphismes rendant commutatifs les diagrammes
  (\hyperref[I.10.6.7.1-fr]{10.6.7.1}), où
  $X_n=(\mathfrak{X},
  \mathscr{O}_{\mathfrak{X}}/\mathscr{T}^{n+1})$,
  $Y_n=(\mathfrak{Y},
  \mathscr{O}_{\mathfrak{Y}}/\mathscr{K}^{n+1})$.
\end{enumerate}
\end{corollary}

\begin{proof}
(ii) résulte aussitôt de (i) et de
(\hyperref[I.10.6.9-fr]{10.6.9}). Pour démontrer (i), on peut se borner
au cas où $\mathfrak{X}=\operatorname{Spf}(A)$,
$\mathfrak{Y}=\operatorname{Spf}(B)$, $A$ et $B$ étant noethériens,
$\mathscr{T}=\mathfrak{T}^\Delta$,
$\mathscr{K}=\mathfrak{K}^\Delta$, où $\mathfrak{T}$ est le plus grand
idéal de définition de $A$ et $\mathfrak{K}$ un idéal de définition de
$B$. Soit $f=({}^a\varphi,\widetilde{\varphi})$, où
$\varphi:B\to A$ est un homomorphisme continu~; comme les éléments de
$\mathfrak{K}$ sont topologiquement nilpotents
(\hyperref[0.7.1.4-fr]{0, 7.1.4, (ii)}), il en est de même de ceux de
$\varphi(\mathfrak{K})$, donc
$\varphi(\mathfrak{K})\subset\mathfrak{T}$ puisque $\mathfrak{T}$ est
l'ensemble des éléments topologiquement nilpotents de $A$
(\hyperref[0.7.1.6-fr]{0, 7.1.6})~; d'où la conclusion en vertu de
(\hyperref[I.10.5.6-fr]{10.5.6, (ii)}).
\end{proof}

\begin{corollary}[10.6.11]
\label{I.10.6.11-fr}
Soient $\mathfrak{S}$, $\mathfrak{X}$, $\mathfrak{Y}$ trois
préschémas formels localement noethériens,
$f:\mathfrak{X}\to\mathfrak{S}$,
$g:\mathfrak{Y}\to\mathfrak{S}$ des morphismes faisant de
$\mathfrak{X}$ et $\mathfrak{Y}$ des $\mathfrak{S}$-préschémas
formels. Soit $\mathscr{J}$ (resp. $\mathscr{K}$, $\mathscr{L}$) un
Idéal de définition de $\mathfrak{S}$ (resp. $\mathfrak{X}$,
$\mathfrak{Y}$), et supposons que
$f^*(\mathscr{J})\mathscr{O}_{\mathfrak{X}}\subset\mathscr{K}$,
$g^*(\mathscr{J})\mathscr{O}_{\mathfrak{Y}}=\mathscr{L}$~; posons
$S_n=(\mathfrak{S},
\mathscr{O}_{\mathfrak{S}}/\mathscr{J}^{n+1})$,
$X_n=(\mathfrak{X},
\mathscr{O}_{\mathfrak{X}}/\mathscr{K}^{n+1})$,
$Y_n=(\mathfrak{Y},
\mathscr{O}_{\mathfrak{Y}}/\mathscr{L}^{n+1})$. Il y a alors
correspondance

\oldpage[I]{193}

biunivoque canonique entre
$\operatorname{Hom}_{\mathfrak{S}}(\mathfrak{X},\mathfrak{Y})$ et
l'ensemble des suites $(u_n)$ de $S_n$-morphismes
$u_n:X_n\to Y_n$ rendant commutatifs les diagrammes
(\hyperref[I.10.6.7.1-fr]{10.6.7.1}).
\end{corollary}

\begin{proof}
Pour tout $\mathfrak{S}$-morphisme
$u:\mathfrak{X}\to\mathfrak{Y}$, on a par définition $f=g\circ u$, donc
\[
  u^*(\mathscr{L})\mathscr{O}_{\mathfrak{X}}
  =u^*(g^*(\mathscr{J})\mathscr{O}_{\mathfrak{Y}})
    \mathscr{O}_{\mathfrak{X}}
  =f^*(\mathscr{J})\mathscr{O}_{\mathfrak{X}}
  \subset\mathscr{K}
\]
le corollaire découle donc de
(\hyperref[I.10.6.9-fr]{10.6.9}).
\end{proof}

On notera que, pour $m\leq n$, la donnée d'un morphisme
$f_n:X_n\to Y_n$ détermine un morphisme et un seul
$f_m:X_m\to Y_m$ rendant commutatif le diagramme
(\hyperref[I.10.6.7.1-fr]{10.6.7.1}), comme on le voit aussitôt en se
ramenant au cas affine~; on a ainsi défini une application
$\varphi_{mn}:\operatorname{Hom}_{S_n}(X_n,Y_n)
\to\operatorname{Hom}_{S_m}(X_m,Y_m)$ et les
$\operatorname{Hom}_{S_n}(X_n,Y_n)$ forment pour les $\varphi_{mn}$ un
\emph{système projectif d'ensembles}~;
(\hyperref[I.10.6.11-fr]{10.6.11}) s'énonce encore en disant qu'il existe
une bijection canonique
\[
  \operatorname{Hom}_{\mathfrak{S}}(\mathfrak{X},\mathfrak{Y})
  \xrightarrow{\sim}
  \varprojlim_n\operatorname{Hom}_{S_n}(X_n,Y_n).
\]

\subsection{Produit de préschémas formels.}
\label{subsection:I.10.7-fr}

\begin{env}[10.7.1]
\label{I.10.7.1-fr}
Soit $\mathfrak{S}$ un préschéma formel~; les $\mathfrak{S}$-préschémas
formels formant une catégorie, on peut définir la notion de
\emph{produit} de $\mathfrak{S}$-préschémas formels.
\end{env}

\begin{proposition}[10.7.2]
\label{I.10.7.2-fr}
Soient $\mathfrak{X}=\operatorname{Spf}(B)$,
$\mathfrak{Y}=\operatorname{Spf}(C)$ deux schémas formels affines
au-dessus d'un schéma formel affine
$\mathfrak{S}=\operatorname{Spf}(A)$. Soient
$\mathfrak{Z}=\operatorname{Spf}(B\widehat{\otimes}_A C)$,
$p_1$, $p_2$ les $\mathfrak{S}$-morphismes correspondant
(\hyperref[I.10.2.2-fr]{10.2.2}) aux $A$-homomorphismes canoniques
(continus) $\rho$ et $\sigma$ de $B$ et $C$ dans
$B\widehat{\otimes}_A C$~; alors $(\mathfrak{Z},p_1,p_2)$ est un produit
des $\mathfrak{S}$-schémas formels affines $\mathfrak{X}$ et
$\mathfrak{Y}$.
\end{proposition}

\begin{proof}
En vertu de (\hyperref[I.10.4.6-fr]{10.4.6}), tout revient à vérifier que
si, à tout $A$-homomorphisme continu
$\varphi:B\widehat{\otimes}_A C\to D$, où $D$ est un anneau admissible qui
est une $A$-algèbre topologique, on associe le couple
$(\varphi\circ\rho,\varphi\circ\sigma)$, on définit une bijection
\[
  \operatorname{Hom}_A(B\widehat{\otimes}_A C,D)
  \xrightarrow{\sim}
  \operatorname{Hom}_A(B,D)\times\operatorname{Hom}_A(C,D)
\]
ce qui n'est autre que la propriété universelle du produit tensoriel
complété (\hyperref[0.7.7.6-fr]{0, 7.7.6}).
\end{proof}

\begin{proposition}[10.7.3]
\label{I.10.7.3-fr}
Étant donnés deux $\mathfrak{S}$-préschémas formels
$\mathfrak{X}$, $\mathfrak{Y}$, le produit
$\mathfrak{X}\times_{\mathfrak{S}}\mathfrak{Y}$ existe.
\end{proposition}

\begin{proof}
La démonstration est identique à celle de
(\hyperref[I.3.2.6-fr]{3.2.6}), en y remplaçant les schémas affines
(resp. les ouverts affines) par les schémas formels affines (resp. les
ouverts formels affines), et la prop.
(\hyperref[I.3.2.2-fr]{3.2.2}) par
(\hyperref[I.10.7.2-fr]{10.7.2}).
\end{proof}

Toutes les propriétés formelles du produit de préschémas
(\hyperref[I.3.2.7-fr]{3.2.7} et
\hyperref[I.3.2.8-fr]{3.2.8},
\hyperref[I.3.3.1-fr]{3.3.1} à
\hyperref[I.3.3.12-fr]{3.3.12}) sont valables sans aucune modification
pour le produit de préschémas formels.

\begin{env}[10.7.4]
\label{I.10.7.4-fr}
Soient $\mathfrak{S}$, $\mathfrak{X}$, $\mathfrak{Y}$ trois préschémas
formels et soient $f:\mathfrak{X}\to\mathfrak{S}$,
$g:\mathfrak{Y}\to\mathfrak{S}$ deux morphismes. Supposons qu'il existe
dans $\mathfrak{S}$, $\mathfrak{X}$, $\mathfrak{Y}$ respectivement,
trois systèmes fondamentaux d'Idéaux de définition
$(\mathscr{J}_\lambda)$, $(\mathscr{K}_\lambda)$,
$(\mathscr{L}_\lambda)$ respectivement, ayant même ensemble d'indices
$I$, tels que
$f^*(\mathscr{J}_\lambda)\mathscr{O}_{\mathfrak{X}}
\subset\mathscr{K}_\lambda$ et
$g^*(\mathscr{J}_\lambda)\mathscr{O}_{\mathfrak{Y}}
\subset\mathscr{L}_\lambda$ pour tout $\lambda$. Posons
$S_\lambda=(\mathfrak{S},
\mathscr{O}_{\mathfrak{S}}/\mathscr{J}_\lambda)$,
$X_\lambda=(\mathfrak{X},
\mathscr{O}_{\mathfrak{X}}/\mathscr{K}_\lambda)$,
$Y_\lambda=(\mathfrak{Y},
\mathscr{O}_{\mathfrak{Y}}/\mathscr{L}_\lambda)$~; pour
$\mathscr{J}_\mu\subset\mathscr{J}_\lambda$,
$\mathscr{K}_\mu\subset\mathscr{K}_\lambda$,
$\mathscr{L}_\mu\subset\mathscr{L}_\lambda$, notons que
$S_\lambda$ (resp. $X_\lambda$, $Y_\lambda$) est un sous-préschéma fermé
de $S_\mu$ (resp. $X_\mu$, $Y_\mu$) ayant même

\oldpage[I]{194}

espace sous-jacent (\hyperref[I.10.6.1-fr]{10.6.1}). Comme
$S_\lambda\to S_\mu$ est un monomorphisme de préschémas, on voit donc
d'abord que les produits
$X_\lambda\times_{S_\lambda}Y_\lambda$ et
$X_\lambda\times_{S_\mu}Y_\lambda$ sont identiques
(\hyperref[I.3.2.4-fr]{3.2.4}), puis que
$X_\lambda\times_{S_\mu}Y_\lambda$ s'identifie à un sous-préschéma fermé
de $X_\mu\times_{S_\mu}Y_\mu$ ayant même espace sous-jacent
(\hyperref[I.4.3.1-fr]{4.3.1}). Cela étant, le produit
$\mathfrak{X}\times_{\mathfrak{S}}\mathfrak{Y}$ est la
\emph{limite inductive} des préschémas usuels
$X_\lambda\times_{S_\lambda}Y_\lambda$~: en effet, on voit comme dans
(\hyperref[I.10.6.2-fr]{10.6.2}) qu'on peut se ramener au cas où
$\mathfrak{S}$, $\mathfrak{X}$ et $\mathfrak{Y}$ sont des schémas formels
affines. Compte tenu de (\hyperref[I.10.5.6-fr]{10.5.6, (ii)}) et de
l'hypothèse sur les systèmes fondamentaux d'Idéaux de définition de
$\mathfrak{S}$, $\mathfrak{X}$ et $\mathfrak{Y}$, on voit aussitôt que
notre assertion résulte de la définition du produit tensoriel complété de
deux algèbres (\hyperref[0.7.7.1-fr]{0, 7.7.1}).

En outre, soit $\mathfrak{Z}$ un $\mathfrak{S}$-préschéma formel,
$(\mathscr{M}_\lambda)$ un système fondamental d'Idéaux de définition de
$\mathfrak{Z}$ ayant $I$ comme ensemble d'indices,
$u:\mathfrak{Z}\to\mathfrak{X}$,
$v:\mathfrak{Z}\to\mathfrak{Y}$ deux $\mathfrak{S}$-morphismes tels que
$u^*(\mathscr{K}_\lambda)\mathscr{O}_{\mathfrak{Z}}
\subset\mathscr{M}_\lambda$ et
$v^*(\mathscr{L}_\lambda)\mathscr{O}_{\mathfrak{Z}}
\subset\mathscr{M}_\lambda$. Si l'on pose
$Z_\lambda=(\mathfrak{Z},
\mathscr{O}_{\mathfrak{Z}}/\mathscr{M}_\lambda)$, et si
$u_\lambda:Z_\lambda\to X_\lambda$ et
$v_\lambda:Z_\lambda\to Y_\lambda$ sont les $S_\lambda$-morphismes
correspondant à $u$ et $v$ (\hyperref[I.10.5.6-fr]{10.5.6}), on vérifie
aussitôt que $(u,v)_{\mathfrak{S}}$ est la limite inductive des
$S_\lambda$-morphismes $(u_\lambda,v_\lambda)_{S_\lambda}$.

Les considérations de ce numéro s'appliquent en particulier lorsque
$\mathfrak{S}$, $\mathfrak{X}$ et $\mathfrak{Y}$ sont localement
noethériens, en prenant comme systèmes fondamentaux d'Idéaux de définition
les systèmes formés des puissances d'un Idéal de définition
(\hyperref[I.10.5.1-fr]{10.5.1}). Mais on notera que
$\mathfrak{X}\times_{\mathfrak{S}}\mathfrak{Y}$ n'est pas nécessairement
localement noethérien (voir toutefois
(\hyperref[I.10.13.5-fr]{10.13.5})).
\end{env}

\subsection{Complété formel d'un préschéma le long d'une partie fermée.}
\label{subsection:I.10.8-fr}

\begin{env}[10.8.1]
\label{I.10.8.1-fr}
Soient $X$ un préschéma (usuel) localement noethérien, $X'$ une partie
fermée de l'espace sous-jacent à $X$~; désignons par $\Phi$ l'ensemble des
faisceaux cohérents d'idéaux $\mathscr{J}$ dans $\mathscr{O}_X$, tels que
le support de $\mathscr{O}_X/\mathscr{J}$ soit $X'$. L'ensemble $\Phi$
n'est pas vide (\hyperref[I.5.2.1-fr]{5.2.1},
\hyperref[I.4.1.4-fr]{4.1.4} et
\hyperref[I.6.1.1-fr]{6.1.1})~; nous l'ordonnerons par la relation
$\supset$.
\end{env}

\begin{lemma}[10.8.2]
\label{I.10.8.2-fr}
L'ensemble ordonné $\Phi$ est filtrant~; si $X$ est noethérien, pour tout
$\mathscr{J}_0\in\Phi$, l'ensemble des puissances
$\mathscr{J}_0^n$ ($n>0$) est cofinal à $\Phi$.
\end{lemma}

\begin{proof}
En effet, si $\mathscr{J}_1$ et $\mathscr{J}_2$ appartiennent à $\Phi$, et
si on pose $\mathscr{J}=\mathscr{J}_1\cap\mathscr{J}_2$,
$\mathscr{J}$ est cohérent puisque $\mathscr{O}_X$ est cohérent
(\hyperref[I.6.1.1-fr]{6.1.1} et
\hyperref[0.5.3.4-fr]{0, 5.3.4}), et l'on a
$\mathscr{J}_x=(\mathscr{J}_1)_x\cap(\mathscr{J}_2)_x$ pour tout
$x\in X$, donc $\mathscr{J}_x=\mathscr{O}_x$ pour $x\notin X'$ et
$\mathscr{J}_x\neq\mathscr{O}_x$ pour $x\in X'$, ce qui prouve que
$\mathscr{J}\in\Phi$. D'autre part, si $X$ est noethérien, et si
$\mathscr{J}_0$ et $\mathscr{J}$ appartiennent à $\Phi$, il existe un
entier $n>0$ tel que
$\mathscr{J}_0^n(\mathscr{O}_X/\mathscr{J})=0$
(\hyperref[I.9.3.4-fr]{9.3.4}), ce qui signifie que
$\mathscr{J}_0^n\subset\mathscr{J}$.
\end{proof}

\begin{env}[10.8.3]
\label{I.10.8.3-fr}
Soit maintenant $\mathscr{F}$ un $\mathscr{O}_X$-Module cohérent~; pour
tout $\mathscr{J}\in\Phi$,
$\mathscr{F}\otimes_{\mathscr{O}_X}(\mathscr{O}_X/\mathscr{J})$ est un
$\mathscr{O}_X$-Module cohérent
(\hyperref[I.9.1.1-fr]{9.1.1}), de support contenu dans $X'$, et que nous
identifierons le plus souvent à sa restriction à $X'$. Lorsque
$\mathscr{J}$ parcourt $\Phi$, ces faisceaux forment un
\emph{système projectif} de faisceaux de groupes abéliens.
\end{env}

\begin{definition}[10.8.4]
\label{I.10.8.4-fr}
Étant donnés une partie fermée $X'$ d'un préschéma localement noethérien
$X$ et un $\mathscr{O}_X$-Module cohérent $\mathscr{F}$, on appelle
\emph{complété de $\mathscr{F}$ le long de $X'$} et on désigne par
$\mathscr{F}_{/X'}$ ou par $\widehat{\mathscr{F}}$ (lorsqu'aucune confusion
n'est possible) la restriction à $X'$ du faisceau

\oldpage[I]{195}

$\varprojlim_{\Phi}
(\mathscr{F}\otimes_{\mathscr{O}_X}
(\mathscr{O}_X/\mathscr{J}))$~; on dit que ses sections au-dessus de $X'$ sont
les \emph{sections formelles de $\mathscr{F}$ le long de $X'$}.
\end{definition}

Il est immédiat que pour tout ouvert $U\subset X$, on a
$(\mathscr{F}|U)_{/(U\cap X')}=(\mathscr{F}_{/X'})|(U\cap X')$.

Par passage à la limite projective, il est clair que
$(\mathscr{O}_X)_{/X'}$ est un faisceau d'anneaux, et que
$\mathscr{F}_{/X'}$ peut être considéré comme un
$(\mathscr{O}_X)_{/X'}$-Module. En outre, comme il existe une base de la
topologie de $X'$ formée d'ouverts quasi-compacts, on peut considérer
$(\mathscr{O}_X)_{/X'}$ (resp. $\mathscr{F}_{/X'}$) comme un
\emph{faisceau d'anneaux topologiques} (resp. de
\emph{groupes topologiques}) limite projective des faisceaux d'anneaux
(resp. groupes) \emph{pseudo-discrets}
$\mathscr{O}_X/\mathscr{J}$ (resp.
$\mathscr{F}\otimes_{\mathscr{O}_X}
(\mathscr{O}_X/\mathscr{J})=\mathscr{F}/\mathscr{J}\mathscr{F}$), et, par
passage à la limite projective, $\mathscr{F}_{/X'}$ devient alors un
\emph{$(\mathscr{O}_X)_{/X'}$-Module topologique}
(\hyperref[0.3.8.1-fr]{0, 3.8.1} et
\hyperref[0.3.8.2-fr]{3.8.2})~; rappelons que pour tout ouvert
\emph{quasi-compact} $U\subset X$,
$\Gamma(U\cap X',(\mathscr{O}_X)_{/X'})$ (resp.
$\Gamma(U\cap X',\mathscr{F}_{/X'})$) est alors limite projective des
anneaux (resp. groupes) discrets $\Gamma(U,\mathscr{O}_X/\mathscr{J})$
(resp. $\Gamma(U,\mathscr{F}/\mathscr{J}\mathscr{F})$).

Si maintenant $u:\mathscr{F}\to\mathscr{G}$ est un homomorphisme de
$\mathscr{O}_X$-Modules, on en déduit canoniquement des homomorphismes
$u_{\mathscr{J}}:
\mathscr{F}\otimes_{\mathscr{O}_X}(\mathscr{O}_X/\mathscr{J})
\to
\mathscr{G}\otimes_{\mathscr{O}_X}(\mathscr{O}_X/\mathscr{J})$
pour tout $\mathscr{J}\in\Phi$, et ces homomorphismes forment un système
projectif. Par passage à la limite projective et restriction à $X'$, ils
donnent donc un $(\mathscr{O}_X)_{/X'}$-homomorphisme continu
$\mathscr{F}_{/X'}\to\mathscr{G}_{/X'}$, noté $u_{/X'}$ ou $\widehat u$,
et appelé le \emph{complété de l'homomorphisme $u$ le long de $X'$}. Il est
clair que si $v:\mathscr{G}\to\mathscr{H}$ est un second homomorphisme de
$\mathscr{O}_X$-Modules, on a
$(v\circ u)_{/X'}=(v_{/X'})\circ(u_{/X'})$ donc
$\mathscr{F}_{/X'}$ est un \emph{foncteur additif covariant} en
$\mathscr{F}$, de la catégorie des $\mathscr{O}_X$-Modules cohérents, à
valeurs dans la catégorie des $(\mathscr{O}_X)_{/X'}$-Modules topologiques.

\begin{proposition}[10.8.5]
\label{I.10.8.5-fr}
Le support de $(\mathscr{O}_X)_{/X'}$ est $X'$~; l'espace topologiquement
annelé $(X',(\mathscr{O}_X)_{/X'})$ est un préschéma formel localement
noethérien, et si $\mathscr{J}\in\Phi$, $\mathscr{J}_{/X'}$ est un Idéal de
définition de ce préschéma formel. Si $X=\operatorname{Spec}(A)$ est un
schéma affine d'anneau noethérien, $\mathscr{J}=\widetilde{\mathfrak{J}}$
où $\mathfrak{J}$ est un idéal de $A$, et $X'=V(\mathfrak{J})$,
$(X',(\mathscr{O}_X)_{/X'})$ s'identifie canoniquement à
$\operatorname{Spf}(\widehat A)$, où $\widehat A$ est le séparé complété de
$A$ pour la topologie $\mathfrak{J}$-préadique.
\end{proposition}

\begin{proof}
On peut évidemment se borner à prouver la dernière assertion. On sait
(\hyperref[0.7.3.3-fr]{0, 7.3.3}) que le séparé complété
$\widehat{\mathfrak{J}}$ de $\mathfrak{J}$ pour la topologie
$\mathfrak{J}$-préadique s'identifie à l'idéal
$\mathfrak{J}\widehat A$ de $\widehat A$, et que $\widehat A$ est un anneau
$\widehat{\mathfrak{J}}$-adique noethérien tel que
$\widehat A/\widehat{\mathfrak{J}}^n=A/\mathfrak{J}^n$
(\hyperref[0.7.2.6-fr]{0, 7.2.6}). Cette dernière relation montre que les
idéaux premiers ouverts de $\widehat A$ sont les idéaux
$\widehat{\mathfrak p}=\mathfrak p\widehat A$, où $\mathfrak p$ est un
idéal premier de $A$ contenant $\mathfrak{J}$, et que l'on a
$\widehat{\mathfrak p}\cap A=\mathfrak p$, d'où
$\operatorname{Spf}(\widehat A)=X'$. Comme
$\mathscr{O}_X/\mathscr{J}^n=(A/\mathfrak{J}^n)^{\sim}$, la proposition
découle aussitôt des définitions.
\end{proof}

On dit que le préschéma formel ainsi défini est le
\emph{complété de $X$ le long de $X'$} et on le note $X_{/X'}$ ou
$\widehat X$ si aucune confusion n'est à craindre. Lorsque l'on prend
$X'=X$, on peut prendre $\mathscr{J}=0$, et on a donc $X_{/X}=X$.

Il est clair que si $U$ est un sous-préschéma induit sur un ouvert de $X$,
$U_{/(U\cap X')}$ s'identifie canoniquement au sous-préschéma formel induit
par $X_{/X'}$ sur l'ouvert $U\cap X'$ de $X'$.

\begin{corollary}[10.8.6]
\label{I.10.8.6-fr}
Le préschéma (usuel) $\widehat X_{\mathrm{red}}$ est l'unique
sous-préschéma réduit de $X$ ayant pour espace sous-jacent $X'$
(\hyperref[I.5.2.1-fr]{5.2.1}). Pour que $\widehat X$ soit noethérien, il
faut et il suffit que $\widehat X_{\mathrm{red}}$ le soit, et il suffit que
$X$ le soit.
\end{corollary}

\oldpage[I]{196}

La détermination de $\widehat X_{\mathrm{red}}$ étant locale
(\hyperref[I.10.5.4-fr]{10.5.4}), on peut encore supposer que $X$ est un
schéma affine d'anneau noethérien~; avec les notations de
(\hyperref[I.10.8.5-fr]{10.8.5}), l'idéal $\mathfrak{T}$ des éléments
topologiquement nilpotents de $\widehat A$ est l'image réciproque par
l'application canonique
$\widehat A\to\widehat A/\widehat{\mathfrak{J}}=A/\mathfrak{J}$ du
nilradical de $A/\mathfrak{J}$
(\hyperref[0.7.1.3-fr]{0, 7.1.3}), donc $\widehat A/\mathfrak{T}$ est
isomorphe au quotient de $A/\mathfrak{J}$ par son nilradical. La première
assertion résulte donc de (\hyperref[I.10.5.4-fr]{10.5.4}) et
(\hyperref[I.5.1.1-fr]{5.1.1}). Si $\widehat X_{\mathrm{red}}$ est
noethérien, son espace sous-jacent $X'$ l'est aussi, donc les
$X'_n=\operatorname{Spec}(\mathscr{O}_X/\mathscr{J}^n)$ sont noethériens
(\hyperref[I.6.1.2-fr]{6.1.2}) et il en est de même de $\widehat X$
(\hyperref[I.10.6.4-fr]{10.6.4})~; la réciproque est immédiate, en vertu
de (\hyperref[I.6.1.2-fr]{6.1.2}).

\begin{env}[10.8.7]
\label{I.10.8.7-fr}
Les homomorphismes canoniques
$\mathscr{O}_X\to\mathscr{O}_X/\mathscr{J}$ (pour
$\mathscr{J}\in\Phi$) forment un système projectif et donnent donc, par
passage à la limite projective, un homomorphisme de faisceau d'anneaux
$\theta:\mathscr{O}_X\to\psi_*((\mathscr{O}_X)_{/X'})
=\varprojlim_\Phi(\mathscr{O}_X/\mathscr{J})$, en désignant par $\psi$
l'injection canonique $X'\to X$ des espaces sous-jacents. Nous désignerons
par $i$ (ou $i_X$) le morphisme (dit canonique)
\[
  (\psi,\theta):X_{/X'}\to X
\]
d'espaces annelés.

Par tensorisation, pour tout $\mathscr{O}_X$-Module cohérent
$\mathscr{F}$, les homomorphismes canoniques
$\mathscr{O}_X\to\mathscr{O}_X/\mathscr{J}$ donnent des homomorphismes
$\mathscr{F}\to
\mathscr{F}\otimes_{\mathscr{O}_X}(\mathscr{O}_X/\mathscr{J})$ de
$\mathscr{O}_X$-Modules qui forment encore un système projectif, et
donnent donc, par passage à la limite projective, un homomorphisme
canonique fonctoriel
$\gamma:\mathscr{F}\to\psi_*(\mathscr{F}_{/X'})$ de
$\mathscr{O}_X$-Modules.
\end{env}

\begin{proposition}[10.8.8]
\label{I.10.8.8-fr}
\medskip\noindent
\begin{enumerate}
  \item[\rm{(i)}] Le foncteur $\mathscr{F}_{/X'}$ (en $\mathscr{F}$) est
    exact.
  \item[\rm{(ii)}] L'homomorphisme fonctoriel
    $\gamma^\sharp:i^*(\mathscr{F})\to\mathscr{F}_{/X'}$ de
    $(\mathscr{O}_X)_{/X'}$-Modules est un isomorphisme.
\end{enumerate}
\end{proposition}

\medskip\noindent
\begin{enumerate}
  \item[\rm{(i)}] Il suffit de prouver que si
    $0\to\mathscr{F}'\to\mathscr{F}\to\mathscr{F}''\to0$ est une suite
    exacte de $\mathscr{O}_X$-Modules cohérents, et $U$ un ouvert affine
    de $X$, d'anneau noethérien $A$, la suite
    \[
      0\to\Gamma(U\cap X',\mathscr{F}'_{/X'})
      \to\Gamma(U\cap X',\mathscr{F}_{/X'})
      \to\Gamma(U\cap X',\mathscr{F}''_{/X'})\to0
    \]
    est exacte. On a alors $\mathscr{F}|U=\widetilde M$,
    $\mathscr{F}'|U=\widetilde{M'}$, $\mathscr{F}''|U=\widetilde{M''}$,
    où $M$, $M'$, $M''$ sont trois $A$-modules de type fini tels que la
    suite $0\to M'\to M\to M''\to0$ soit exacte
    (\hyperref[I.1.5.1-fr]{1.5.1} et
    \hyperref[I.1.3.11-fr]{1.3.11})~; soit $\mathscr{J}\in\Phi$ et soit
    $\mathfrak{J}$ un idéal de $A$ tel que
    $\mathscr{J}|U=\widetilde{\mathfrak{J}}$. On a alors
    \[
      \Gamma(U\cap X',
      \mathscr{F}\otimes_{\mathscr{O}_X}(\mathscr{O}_X/\mathscr{J}^n))
      =M\otimes_A(A/\mathfrak{J}^n)
    \]
    (\hyperref[I.1.3.12-fr]{1.3.12})~; donc, par définition de la limite
    projective, on a
    \[
      \Gamma(U\cap X',\mathscr{F}_{/X'})
      =\varprojlim_n(M\otimes_A(A/\mathfrak{J}^n))=\widehat M
    \]
    séparé complété de $M$ pour la topologie $\mathfrak{J}$-préadique, et
    de même
    \[
      \Gamma(U\cap X',\mathscr{F}'_{/X'})=\widehat{M'},\qquad
      \Gamma(U\cap X',\mathscr{F}''_{/X'})=\widehat{M''}~;
    \]
    notre assertion résulte alors de ce que lorsque $A$ est noethérien,
    le foncteur $\widehat M$ en $M$ est exact sur la catégorie des
    $A$-modules de type fini (\hyperref[0.7.3.3-fr]{0, 7.3.3}).

\oldpage[I]{197}
  \item[\rm{(ii)}] La question étant locale, on peut supposer que l'on a
    une suite exacte
    $\mathscr{O}_X^m\to\mathscr{O}_X^n\to\mathscr{F}\to0$
    (\hyperref[0.5.3.2-fr]{0, 5.3.2})~; comme $\gamma^\sharp$ est
    fonctoriel, et que les foncteurs $i^*(\mathscr{F})$ et
    $\mathscr{F}_{/X'}$ sont exacts à droite (d'après (i) et
    (\hyperref[0.4.3.1-fr]{0, 4.3.1})), on a le diagramme commutatif
    \[
      \xymatrix{
        i^*(\mathscr{O}_X^m)\ar[r]\ar[d]_{\gamma^\sharp} &
        i^*(\mathscr{O}_X^n)\ar[r]\ar[d]_{\gamma^\sharp} &
        i^*(\mathscr{F})\ar[r]\ar[d]_{\gamma^\sharp} &
        0\\
        (\mathscr{O}_X^m)_{/X'}\ar[r] &
        (\mathscr{O}_X^n)_{/X'}\ar[r] &
        \mathscr{F}_{/X'}\ar[r] &
        0
      }
      \tag{10.8.8.1}\label{I.10.8.8.1-fr}
    \]
    dont les lignes sont exactes. En outre, les deux foncteurs
    $i^*(\mathscr{F})$ et $\mathscr{F}_{/X'}$ commutent aux sommes
    directes finies (\hyperref[0.3.2.6-fr]{0, 3.2.6} et
    \hyperref[0.4.3.2-fr]{4.3.2}) et on est donc ramené à démontrer notre
    assertion pour $\mathscr{F}=\mathscr{O}_X$. On a alors
    $i^*(\mathscr{O}_X)=(\mathscr{O}_X)_{/X'}=
    \mathscr{O}_{\widehat X}$
    (\hyperref[0.4.3.4-fr]{0, 4.3.4}), et $\gamma^\sharp$ est un
    homomorphisme de $\mathscr{O}_{\widehat X}$-Modules~; il suffit donc
    de vérifier que $\gamma^\sharp$ transforme la section unité de
    $\mathscr{O}_{\widehat X}$ au-dessus d'un ouvert de $X'$ en elle-même,
    ce qui est immédiat et montre donc que dans ce cas $\gamma^\sharp$ est
    l'identité.
\end{enumerate}

\begin{corollary}[10.8.9]
\label{I.10.8.9-fr}
Le morphisme d'espaces annelés $i:X_{/X'}\to X$ est plat.
\end{corollary}

Cela résulte en effet de (\hyperref[0.6.7.3-fr]{0, 6.7.3}) et de
(\hyperref[I.10.8.8-fr]{10.8.8, (i)}).

\begin{corollary}[10.8.10]
\label{I.10.8.10-fr}
Si $\mathscr{F}$ et $\mathscr{G}$ sont des $\mathscr{O}_X$-Modules
cohérents, il existe des isomorphismes canoniques fonctoriels (en
$\mathscr{F}$ et $\mathscr{G}$)
\[
  (\mathscr{F}_{/X'})\otimes_{(\mathscr{O}_X)_{/X'}}
  (\mathscr{G}_{/X'})
  \xrightarrow{\sim}
  (\mathscr{F}\otimes_{\mathscr{O}_X}\mathscr{G})_{/X'}
  \tag{10.8.10.1}\label{I.10.8.10.1-fr}
\]
\[
  (\mathscr{H}\!om_{\mathscr{O}_X}(\mathscr{F},\mathscr{G}))_{/X'}
  \xrightarrow{\sim}
  \mathscr{H}\!om_{(\mathscr{O}_X)_{/X'}}
  (\mathscr{F}_{/X'},\mathscr{G}_{/X'})
  \tag{10.8.10.2}\label{I.10.8.10.2-fr}
\]
\end{corollary}

Cela résulte de l'identification canonique de $i^*(\mathscr{F})$ et de
$\mathscr{F}_{/X'}$~; l'existence du premier isomorphisme est alors un
résultat valable pour tous les morphismes d'espaces annelés
(\hyperref[0.4.3.3.1-fr]{0, 4.3.3.1}) et celle du second est un résultat
valable pour tous les morphismes plats
(\hyperref[0.6.7.6-fr]{0, 6.7.6}), donc découle de
(\hyperref[I.10.8.9-fr]{10.8.9}).

\begin{proposition}[10.8.11]
\label{I.10.8.11-fr}
Pour tout $\mathscr{O}_X$-Module cohérent $\mathscr{F}$, le noyau de
l'homomorphisme canonique
$\Gamma(X,\mathscr{F})\to\Gamma(X',\mathscr{F}_{/X'})$ déduit de
$\mathscr{F}\to\mathscr{F}_{/X'}$ est formé des sections nulles dans un
voisinage de $X'$.
\end{proposition}

Il résulte de la définition de $\mathscr{F}_{/X'}$ que l'image canonique
d'une telle section est nulle. Réciproquement, si
$s\in\Gamma(X,\mathscr{F})$ a une image nulle dans
$\Gamma(X',\mathscr{F}_{/X'})$, il suffit de voir que tout $x\in X'$ admet
un voisinage dans $X$ dans lequel $s$ est nulle, et on peut donc se ramener
au cas où $X=\operatorname{Spec}(A)$ est affine, $A$ noethérien,
$X'=V(\mathfrak{J})$, où $\mathfrak{J}$ est un idéal de $A$, et
$\mathscr{F}=\widetilde M$, où $M$ est un $A$-module de type fini. Alors
$\Gamma(X',\mathscr{F}_{/X'})$ est le séparé complété $\widehat M$ de $M$
pour la topologie $\mathfrak{J}$-préadique, et l'homomorphisme
$\Gamma(X,\mathscr{F})\to\Gamma(X',\mathscr{F}_{/X'})$ est
l'homomorphisme canonique $M\to\widehat M$. On sait
(\hyperref[0.7.3.7-fr]{0, 7.3.7}) que le noyau de cet homomorphisme est
l'ensemble des $z\in M$ annulés par un élément de $1+\mathfrak{J}$. On a
donc $(1+f)s=0$ pour un $f\in\mathfrak{J}$~; pour tout $x\in X'$ on en
déduit $(1_x+f_x)s_x=0$, et comme $1_x+f_x$ est inversible dans
$\mathscr{O}_x$ ($\mathfrak{J}_x\mathscr{O}_x$ étant contenu dans l'idéal
maximal de $\mathscr{O}_x$), on a $s_x=0$, ce qui démontre la proposition.

\begin{corollary}[10.8.12]
\label{I.10.8.12-fr}
Le support de $\mathscr{F}_{/X'}$ est égal à
$\operatorname{Supp}(\mathscr{F})\cap X'$.
\end{corollary}

Il est clair que $\mathscr{F}_{/X'}$ est un
$(\mathscr{O}_X)_{/X'}$-Module de type fini
(\hyperref[I.10.8.8-fr]{10.8.8, (ii)}) et
(\hyperref[0.5.2.4-fr]{0, 5.2.4}),

\oldpage[I]{198}

donc son support est fermé (\hyperref[0.5.2.2-fr]{0, 5.2.2}) et
évidemment contenu dans $\operatorname{Supp}(\mathscr{F})\cap X'$. Pour
montrer qu'il est égal à ce dernier ensemble, on est aussitôt ramené à
prouver que la relation $\Gamma(X',\mathscr{F}_{/X'})=0$ entraîne
$\operatorname{Supp}(\mathscr{F})\cap X'=\varnothing$~; or cela résulte de
(\hyperref[I.10.8.11-fr]{10.8.11}) et de
(\hyperref[I.1.4.1-fr]{1.4.1}).

\begin{corollary}[10.8.13]
\label{I.10.8.13-fr}
Soit $u:\mathscr{F}\to\mathscr{G}$ un homomorphisme de
$\mathscr{O}_X$-Modules cohérents. Pour que
$u_{/X'}:\mathscr{F}_{/X'}\to\mathscr{G}_{/X'}$ soit nul, il faut et il
suffit que $u$ soit nul dans un voisinage de $X'$.
\end{corollary}

En effet, d'après (\hyperref[I.10.8.8-fr]{10.8.8, (ii)}), $u_{/X'}$
s'identifie à $i^*(u)$, donc si on considère $u$ comme une section au-dessus
de $X$ du faisceau
$\mathscr{H}=\mathscr{H}\!om_{\mathscr{O}_X}(\mathscr{F},\mathscr{G})$,
$u_{/X'}$ est la section de
$i^*(\mathscr{H})=\mathscr{H}_{/X'}$ au-dessus de $X'$ qui lui correspond
canoniquement ((\hyperref[I.10.8.10.2-fr]{10.8.10.2}) et
(\hyperref[0.4.4.6-fr]{0, 4.4.6})). Il suffit donc d'appliquer
(\hyperref[I.10.8.11-fr]{10.8.11}) au $\mathscr{O}_X$-Module cohérent
$\mathscr{H}$.

\begin{corollary}[10.8.14]
\label{I.10.8.14-fr}
Soit $u:\mathscr{F}\to\mathscr{G}$ un homomorphisme de
$\mathscr{O}_X$-Modules cohérents. Pour que $u_{/X'}$ soit un monomorphisme
(resp. un épimorphisme), il faut et il suffit que $u$ soit un monomorphisme
(resp. un épimorphisme) dans un voisinage de $X'$.
\end{corollary}

Soient $\mathscr{P}$ et $\mathscr{N}$ le conoyau et le noyau de $u$, de
sorte qu'on a la suite exacte
\[
  0\to\mathscr{N}\xrightarrow{v}\mathscr{F}\xrightarrow{u}
  \mathscr{G}\xrightarrow{w}\mathscr{P}\to0,
\]
d'où (\hyperref[I.10.8.8-fr]{10.8.8, (i)}) la suite exacte
\[
  0\to\mathscr{N}_{/X'}\xrightarrow{v_{/X'}}
  \mathscr{F}_{/X'}\xrightarrow{u_{/X'}}
  \mathscr{G}_{/X'}\xrightarrow{w_{/X'}}
  \mathscr{P}_{/X'}\to0.
\]
Si $u_{/X'}$ est un monomorphisme (resp. un épimorphisme), on a
$v_{/X'}=0$ (resp. $w_{/X'}=0$), donc il y a un voisinage de $X'$ dans
lequel $v=0$ (resp. $w=0$) en vertu de
(\hyperref[I.10.8.13-fr]{10.8.13}).

\subsection{Prolongement d'un morphisme aux complétés.}
\label{subsection:I.10.9-fr}

\begin{env}[10.9.1]
\label{I.10.9.1-fr}
Soient $X$, $Y$ deux préschémas (usuels) localement noethériens,
$f:X\to Y$ un morphisme, $X'$ (resp. $Y'$) une partie fermée de l'espace
sous-jacent $X$ (resp. $Y$), telles que $f(X')\subset Y'$. Soit
$\mathscr{J}$ (resp. $\mathscr{K}$) un faisceau d'idéaux de
$\mathscr{O}_X$ (resp. $\mathscr{O}_Y$) tel que le support de
$\mathscr{O}_X/\mathscr{J}$ (resp. $\mathscr{O}_Y/\mathscr{K}$) soit
$X'$ (resp. $Y'$) et que
$f^*(\mathscr{K})\mathscr{O}_X\subset\mathscr{J}$~; on notera qu'il
existe toujours de tels faisceaux d'idéaux, car on peut par exemple prendre
pour $\mathscr{J}$ le plus grand faisceau d'idéaux de $\mathscr{O}_X$
définissant un sous-préschéma de $X$ ayant $X'$ pour espace sous-jacent
(\hyperref[I.5.2.1-fr]{5.2.1}), et l'hypothèse $f(X')\subset Y'$ entraîne
alors $f^*(\mathscr{K})\mathscr{O}_X\subset\mathscr{J}$
(\hyperref[I.5.2.4-fr]{5.2.4}). On a donc pour tout entier $n>0$,
$f^*(\mathscr{K}^n)\mathscr{O}_X\subset\mathscr{J}^n$
(\hyperref[0.4.3.5-fr]{0, 4.3.5})~; par suite
(\hyperref[I.4.4.6-fr]{4.4.6}), si on pose
\[
  X'_n=(X',\mathscr{O}_X/\mathscr{J}^{n+1}),\qquad
  Y'_n=(Y',\mathscr{O}_Y/\mathscr{K}^{n+1}),
\]
on déduit de $f$ un morphisme $f_n:X'_n\to Y'_n$, et il est immédiat que
les $f_n$ forment un système inductif. Nous désignerons sa limite inductive
(\hyperref[I.10.6.8-fr]{10.6.8}) par
$\widehat f:X_{/X'}\to Y_{/Y'}$, et nous dirons (par abus de langage) que
$\widehat f$ est le \emph{prolongement de $f$ aux complétés de $X$ et $Y$
le long de $X'$ et $Y'$}. Il est immédiat de vérifier que ce morphisme ne
dépend pas du choix des faisceaux d'idéaux $\mathscr{J}$, $\mathscr{K}$
vérifiant les conditions ci-dessus. Il suffit de le voir en effet lorsque
$X$ et $Y$ sont des schémas affines noethériens d'anneaux $A$, $B$~; alors
$\mathscr{J}=\widetilde{\mathfrak{J}}$,
$\mathscr{K}=\widetilde{\mathfrak{K}}$, où $\mathfrak{J}$ (resp.
$\mathfrak{K}$) est un idéal de $A$ (resp. $B$), $f$ correspond à un
homomorphisme d'anneaux $\varphi:B\to A$ tel que
$\varphi(\mathfrak{K})\subset\mathfrak{J}$
(\hyperref[I.4.4.6-fr]{4.4.6} et
\hyperref[I.1.7.4-fr]{1.7.4})~; \emph{$\widehat f$ est alors le morphisme
qui correspond (\hyperref[I.10.2.2-fr]{10.2.2}) à l'homomorphisme continu
$\widehat\varphi:\widehat B\to\widehat A$}, où $\widehat A$ (resp.
$\widehat B$) est le séparé complété de $A$ (resp. $B$) pour la topologie
$\mathfrak{J}$-préadique (resp. $\mathfrak{K}$-préadique)
(\hyperref[I.10.6.8-fr]{10.6.8})~; et on sait que si on remplace
$\mathscr{J}$ par un autre

\oldpage[I]{199}

faisceau d'idéaux
$\mathscr{J}'=\widetilde{\mathfrak{J}'}$ tel que le support de
$\mathscr{O}_X/\mathscr{J}'$ soit encore $X'$, les topologies
$\mathfrak{J}$-préadique et $\mathfrak{J}'$-préadique sur $A$ sont les
mêmes (\hyperref[I.10.8.2-fr]{10.8.2}).

On notera que, d'après cette définition, l'application continue
$X'\to Y'$ des espaces sous-jacents à $X_{/X'}$ et $Y_{/Y'}$, qui
correspond à $\widehat f$, n'est autre que la restriction à $X'$ de $f$.
\end{env}

\begin{env}[10.9.2]
\label{I.10.9.2-fr}
Il résulte aussitôt de la définition précédente que le diagramme de
morphismes d'espaces annelés
\[
  \label{I.10.9.2.1-fr}
  \xymatrix{
    \widehat X\ar[r]^{\widehat f}\ar[d]_{i_X} &
    \widehat Y\ar[d]^{i_Y}\\
    X\ar[r]_f &
    Y
  }
\]
est commutatif, les flèches verticales étant les morphismes canoniques
(\hyperref[I.10.8.7-fr]{10.8.7}).
\end{env}

\begin{env}[10.9.3]
\label{I.10.9.3-fr}
Soient $Z$ un troisième préschéma, $g:Y\to Z$ un morphisme, $Z'$ une
partie fermée de $Z$ telle que $g(Y')\subset Z'$. Si $\widehat g$ désigne
le complété le long de $Y'$ et $Z'$ du morphisme $g$, il résulte aussitôt
de (\hyperref[I.10.9.1-fr]{10.9.1}) que l'on a
$\widehat{(g\circ f)}=\widehat g\circ\widehat f$.
\end{env}

\begin{proposition}[10.9.4]
\label{I.10.9.4-fr}
Soient $X$, $Y$ deux $S$-préschémas localement noethériens, $Y$ étant de
type fini sur $S$. Soient $f$, $g$, deux $S$-morphismes de $X$ dans $Y$
tels que $f(X')\subset Y'$, $g(X')\subset Y'$. Pour que
$\widehat f=\widehat g$, il faut et il suffit que $f$ et $g$ coïncident
dans un voisinage de $X'$.
\end{proposition}

La condition est évidemment suffisante (sans hypothèse de finitude sur
$Y$). Pour voir qu'elle est nécessaire, remarquons d'abord que l'hypothèse
$\widehat f=\widehat g$ implique $f(x)=g(x)$ pour tout $x\in X'$. D'autre
part, la question étant locale, on peut supposer que $X$ et $Y$ sont des
voisinages ouverts affines respectifs de $x$ et de $y=f(x)=g(x)$,
d'anneaux noethériens, que $S$ est affine et que
$\Gamma(Y,\mathscr{O}_Y)$ est une $\Gamma(S,\mathscr{O}_S)$-algèbre de
type fini (\hyperref[I.6.3.3-fr]{6.3.3}). Alors $f$ et $g$ correspondent à
deux $\Gamma(S,\mathscr{O}_S)$-homomorphismes $\rho$, $\sigma$ de
$\Gamma(Y,\mathscr{O}_Y)$ dans $\Gamma(X,\mathscr{O}_X)$
(\hyperref[I.1.7.3-fr]{1.7.3}), et par hypothèse, les prolongements par
continuité de ces homomorphismes au séparé complété de
$\Gamma(Y,\mathscr{O}_Y)$ sont les mêmes. On conclut de
(\hyperref[I.10.8.11-fr]{10.8.11}) que pour toute section
$s\in\Gamma(Y,\mathscr{O}_Y)$, les sections $\rho(s)$ et $\sigma(s)$
coïncident dans un voisinage de $X'$ (dépendant de $s$)~; comme
$\Gamma(Y,\mathscr{O}_Y)$ est une algèbre de type fini sur
$\Gamma(S,\mathscr{O}_S)$, on en déduit aussitôt qu'il existe un
voisinage $V$ de $X'$ tel que $\rho(s)$ et $\sigma(s)$ coïncident dans
$V$ pour \emph{toute} section $s\in\Gamma(Y,\mathscr{O}_Y)$. Si
$h\in\Gamma(Y,\mathscr{O}_X)$ est tel que $D(h)$ soit un voisinage de $x$
contenu dans $V$, on conclut de ce qui précède et de
(\hyperref[I.1.4.1-fr]{1.4.1, d)}) que $f$ et $g$ coïncident dans $D(h)$.

\begin{proposition}[10.9.5]
\label{I.10.9.5-fr}
Sous les hypothèses de (\hyperref[I.10.9.1-fr]{10.9.1}), pour tout
$\mathscr{O}_Y$-Module cohérent $\mathscr{G}$, il existe un isomorphisme
canonique fonctoriel de $(\mathscr{O}_X)_{/X'}$-Modules
\[
  (f^*(\mathscr{G}))_{/X'}\xrightarrow{\sim}
  \widehat f^*(\mathscr{G}_{/Y'})
\]
\end{proposition}

Si on identifie canoniquement $(f^*(\mathscr{G}))_{/X'}$ à
$i_X^*(f^*(\mathscr{G}))$ et $\widehat f^*(\mathscr{G}_{/Y'})$ à
$\widehat f^*(i_Y^*(\mathscr{G}))$
(\hyperref[I.10.8.8-fr]{10.8.8}), la proposition résulte aussitôt de la
commutativité du diagramme de
(\hyperref[I.10.9.2-fr]{10.9.2}).

\begin{env}[10.9.6]
\label{I.10.9.6-fr}
Soient maintenant $\mathscr{F}$ un $\mathscr{O}_X$-Module cohérent,
$\mathscr{G}$ un $\mathscr{O}_Y$-Module cohérent. Si
$u:\mathscr{G}\to\mathscr{F}$ est un $f$-morphisme de $\mathscr{G}$ dans
$\mathscr{F}$, il lui correspond un $\mathscr{O}_X$-homomorphisme
$u^\sharp:f^*(\mathscr{G})\to\mathscr{F}$, donc par complétion un
$(\mathscr{O}_X)_{/X'}$-homomorphisme continu
$(u^\sharp)_{/X'}:(f^*(\mathscr{G}))_{/X'}\to\mathscr{F}_{/X'}$, et en
vertu de (\hyperref[I.10.9.5-fr]{10.9.5}) il existe un
$\widehat f$-morphisme $v:\mathscr{G}_{/Y'}\to\mathscr{F}_{/X'}$,
\oldpage[I]{200}

et un seul, tel que $v^\sharp=(u^\sharp)_{/X'}$. Si on considère les
triplets $(\mathscr{F},X,X')$ ($\mathscr{F}$ étant un
$\mathscr{O}_X$-Module cohérent et $X'$ une partie fermée de $X$) comme
une \emph{catégorie}, les morphismes
$(\mathscr{F},X,X')\to(\mathscr{G},Y,Y')$ consistant en un morphisme de
préschémas $f:X\to Y$ tel que $f(X')\subset Y'$ et en un
$f$-morphisme $u:\mathscr{G}\to\mathscr{F}$, on peut donc dire que
$(X_{/X'},\mathscr{F}_{/X'})$ est un \emph{foncteur} en
$(\mathscr{F},X,X')$, prenant ses valeurs dans la catégorie des couples
$(\mathfrak{Z},\mathscr{H})$ formés d'un préschéma formel localement
noethérien $\mathfrak{Z}$ et d'un
$\mathscr{O}_{\mathfrak{Z}}$-Module $\mathscr{H}$, les morphismes de
cette dernière catégorie consistant en les couples formés d'un morphisme
$g$ de préschémas formels et d'un $g$-morphisme.
\end{env}

\begin{proposition}[10.9.7]
\label{I.10.9.7-fr}
Soient $S$, $X$, $Y$ trois préschémas localement noethériens,
$g:X\to S$, $h:Y\to S$ deux morphismes, $S'$ une partie fermée de $S$,
$X'$ (resp. $Y'$) une partie fermée de $X$ (resp. $Y$) telle que
$g(X')\subset S'$ (resp. $h(Y')\subset S'$)~; soit
$Z=X\times_S Y$~; supposons $Z$ localement noethérien, et soit
$Z'=p^{-1}(X')\cap q^{-1}(Y')$, où $p$ et $q$ sont les projections de
$X\times_S Y$. Dans ces conditions, le complété $Z_{/Z'}$ s'identifie au
produit des $S_{/S'}$-préschémas formels
$(X_{/X'})\times_{S_{/S'}}(Y_{/Y'})$, les morphismes structuraux
s'identifiant à $\widehat g$ et $\widehat h$, et les projections à
$\widehat p$ et $\widehat q$.
\end{proposition}

Il est immédiat que la question est locale pour $S$, $X$ et $Y$, et on
est donc ramené au cas où $S=\operatorname{Spec}(A)$,
$X=\operatorname{Spec}(B)$, $Y=\operatorname{Spec}(C)$,
$S'=V(\mathfrak{J})$, $X'=V(\mathfrak{K})$, $Y'=V(\mathfrak{L})$, où
$\mathfrak{J}$, $\mathfrak{K}$, $\mathfrak{L}$ sont trois idéaux tels
que $\varphi(\mathfrak{J})\subset\mathfrak{K}$ et
$\psi(\mathfrak{J})\subset\mathfrak{L}$, en désignant par $\varphi$ et
$\psi$ les homomorphismes $A\to B$ et $A\to C$ qui correspondent à $g$
et $h$. Alors on sait que $Z=\operatorname{Spec}(B\otimes_A C)$ et que
$Z'=V(\mathfrak{M})$, où $\mathfrak{M}$ est l'idéal
$\operatorname{Im}(\mathfrak{K}\otimes_A C)
+\operatorname{Im}(B\otimes_A\mathfrak{L})$. La conclusion résulte
(\hyperref[I.10.7.2-fr]{10.7.2}) de ce que le produit tensoriel complété
$(\widehat B\otimes_{\widehat A}\widehat C)^\wedge$ (où $\widehat A$,
$\widehat B$, $\widehat C$ sont respectivement les séparés complétés de
$A$, $B$, $C$ pour les topologies $\mathfrak{J}$-, $\mathfrak{K}$- et
$\mathfrak{L}$-préadiques) est le séparé complété du produit tensoriel
$B\otimes_A C$ pour la topologie $\mathfrak{M}$-préadique
(\hyperref[0.7.7.2-fr]{0, 7.7.2}).

On notera en outre que si $T$ est un $S$-préschéma localement
noethérien, $u:T\to X$, $v:T\to Y$ deux $S$-morphismes, $T'$ une partie
fermée de $T$ telle que $u(T')\subset X'$, $v(T')\subset Y'$, alors le
prolongement aux complétés $((u,v)_S)^\wedge$ s'identifie à
$(\widehat u,\widehat v)_{S_{/S'}}$.

\begin{corollary}[10.9.8]
\label{I.10.9.8-fr}
Soient $X$, $Y$ deux $S$-préschémas localement noethériens tels que
$X\times_S Y$ soit localement noethérien~; soient $S'$ une partie fermée
de $S$, $X'$ (resp. $Y'$) une partie fermée de $X$ (resp. $Y$) dont
l'image dans $S$ est contenue dans $S'$. Pour tout $S$-morphisme
$f:X\to Y$ tel que $f(X')\subset Y'$, le morphisme graphe
$\Gamma_{\widehat f}$ s'identifie au prolongement $(\Gamma_f)^\wedge$ du
morphisme graphe de $f$.
\end{corollary}

\begin{corollary}[10.9.9]
\label{I.10.9.9-fr}
Soient $X$, $Y$ deux préschémas localement noethériens, $f:X\to Y$ un
morphisme, $Y'$ une partie fermée de $Y$, $X'=f^{-1}(Y')$. Alors le
préschéma $X_{/X'}$ s'identifie par le diagramme commutatif
\[
  \label{I.10.9.9.1-fr}
  \xymatrix{
    X\ar[d]_f &
    X_{/X'}\ar[l]\ar[d]^{\widehat f}\\
    Y &
    Y_{/Y'}\ar[l]
  }
\]
au produit $X\times_Y(Y_{/Y'})$ de préschémas formels.
\end{corollary}

Il suffit d'appliquer (\hyperref[I.10.9.7-fr]{10.9.7}) en remplaçant
$S$ et $S'$ par $Y$, $X$ et $X'$ par $X$.
\oldpage[I]{201}

\begin{remark}[10.9.10]
\label{I.10.9.10-fr}
Si $X$ est la somme $X_1\amalg X_2$
(\hyperref[subsection:I.3.1-fr]{3.1}), $X'$ la réunion
$X'_1\cup X'_2$, où $X'_i$ est une partie fermée de $X_i$ ($i=1,2$),
on voit aussitôt que l'on a
$X_{/X'}={X_1}_{/X'_1}\amalg {X_2}_{/X'_2}$.
\end{remark}

\subsection{Application aux faisceaux cohérents sur les schémas formels
affines.}
\label{subsection:I.10.10-fr}

\begin{env}[10.10.1]
\label{I.10.10.1-fr}
Dans tout ce paragraphe, $A$ désignera un \emph{anneau adique
noethérien}, $\mathfrak{J}$ un idéal de définition de $A$. Soit
$X=\operatorname{Spec}(A)$, $\mathfrak{X}=\operatorname{Spf}(A)$, qui
s'identifie à la partie fermée $V(\mathfrak{J})$ de $X$
(\hyperref[I.10.1.2-fr]{10.1.2}). En outre, la définition
(\hyperref[I.10.1.2-fr]{10.1.2}) et la définition
(\hyperref[I.10.8.4-fr]{10.8.4}) montrent que le
\emph{schéma formel affine $\mathfrak{X}$} est identique au
\emph{complété $X_{/\mathfrak{X}}$} du schéma affine $X$ le long de la
partie fermée $\mathfrak{X}$ de son espace sous-jacent. A tout
$\mathscr{O}_X$-Module cohérent $\mathscr{F}$ correspond donc un
$\mathscr{O}_{\mathfrak{X}}$-Module de type fini
$\mathscr{F}_{/\mathfrak{X}}$, qui est d'ailleurs un faisceau de modules
topologiques sur le faisceau d'anneaux topologiques
$\mathscr{O}_{\mathfrak{X}}$. Mais tout $\mathscr{O}_X$-Module cohérent
$\mathscr{F}$ est de la forme $\widetilde M$, où $M$ est un $A$-module
de type fini (\hyperref[I.1.5.1-fr]{1.5.1})~; nous poserons
$(\widetilde M)_{/\mathfrak{X}}=M^\Delta$. En outre, si $u:M\to N$ est
un $A$-homomorphisme de $A$-modules de type fini, il lui correspond un
homomorphisme $\widetilde u:\widetilde M\to\widetilde N$, et par suite
aussi un homomorphisme continu
$\widetilde u_{/\mathfrak{X}}:(\widetilde M)_{/\mathfrak{X}}\to
(\widetilde N)_{/\mathfrak{X}}$, que nous noterons $u^\Delta$. Il est
immédiat que
$(v\circ u)^\Delta=v^\Delta\circ u^\Delta$~; on a ainsi défini un
\emph{foncteur additif covariant $M^\Delta$} de la catégorie des
$A$-modules de type fini dans celle des
$\mathscr{O}_{\mathfrak{X}}$-Modules de type fini. Lorsque $A$ est un
anneau \emph{discret}, on a $M^\Delta=\widetilde M$.
\end{env}

\begin{proposition}[10.10.2]
\label{I.10.10.2-fr}
\medskip\noindent
\begin{enumerate}
  \item[\rm{(i)}] $M^\Delta$ est un foncteur exact en $M$, et il existe
    un isomorphisme canonique fonctoriel de $A$-modules
    $\Gamma(\mathfrak{X},M^\Delta)\xrightarrow{\sim}M$.
  \item[\rm{(ii)}] Si $M$ et $N$ sont deux $A$-modules de type fini, il
    existe des isomorphismes canoniques fonctoriels
    \[
      (M\otimes_A N)^\Delta
      \xrightarrow{\sim}
      M^\Delta\otimes_{\mathscr{O}_{\mathfrak{X}}}N^\Delta
      \tag{10.10.2.1}\label{I.10.10.2.1-fr}
    \]
    \[
      (\operatorname{Hom}_A(M,N))^\Delta
      \xrightarrow{\sim}
      \mathscr{H}\!om_{\mathscr{O}_{\mathfrak{X}}}
      (M^\Delta,N^\Delta)
      \tag{10.10.2.2}\label{I.10.10.2.2-fr}
    \]
  \item[\rm{(iii)}] L'application $u\to u^\Delta$ est un isomorphisme
    fonctoriel
    \[
      \operatorname{Hom}_A(M,N)
      \xrightarrow{\sim}
      \operatorname{Hom}_{\mathscr{O}_{\mathfrak{X}}}
      (M^\Delta,N^\Delta)
      \tag{10.10.2.3}\label{I.10.10.2.3-fr}
    \]
\end{enumerate}
\end{proposition}

L'exactitude de $M^\Delta$ résulte de l'exactitude des foncteurs
$\widetilde M$ (\hyperref[I.1.3.5-fr]{1.3.5}) et
$\mathscr{F}_{/\mathfrak{X}}$ (\hyperref[I.10.8.8-fr]{10.8.8}). Par
définition, $\Gamma(\mathfrak{X},M^\Delta)$ est le séparé complété du
$A$-module $\Gamma(X,\widetilde M)=M$ pour la topologie
$\mathfrak{J}$-préadique~; mais comme $A$ est complet et $M$ de type
fini, on sait (\hyperref[0.7.3.6-fr]{0, 7.3.6}) que $M$ est séparé et
complet, ce qui achève de prouver (i). L'isomorphisme
(\hyperref[I.10.10.2.1-fr]{10.10.2.1}) (resp.
\hyperref[I.10.10.2.2-fr]{10.10.2.2}) provient de la composition des
isomorphismes (\hyperref[I.1.3.12-fr]{1.3.12, (i)}) et
(\hyperref[I.10.8.10.1-fr]{10.8.10.1}) (resp.
(\hyperref[I.1.3.12-fr]{1.3.12, (ii)}) et
(\hyperref[I.10.8.10.2-fr]{10.8.10.2})). Enfin, comme
$\operatorname{Hom}_A(M,N)$ est un $A$-module de type fini, on peut lui
appliquer (i), qui identifie
$\Gamma(\mathfrak{X},(\operatorname{Hom}_A(M,N))^\Delta)$ à
$\operatorname{Hom}_A(M,N)$, et utiliser
(\hyperref[I.10.10.2.2-fr]{10.10.2.2}), qui prouve que
l'homomorphisme (\hyperref[I.10.10.2.3-fr]{10.10.2.3}) est un
isomorphisme.

On déduit de (\hyperref[I.10.10.2-fr]{10.10.2}) toute une série de
conséquences analogues à celles déduites de
(\hyperref[I.1.3.7-fr]{1.3.7}) et
(\hyperref[I.1.3.12-fr]{1.3.12}), que nous laissons au lecteur le soin
de formuler.
\oldpage[I]{202}

Notons que la propriété d'exactitude de $M^\Delta$, appliqué à la suite
exacte
$0\to\mathfrak{J}\to A\to A/\mathfrak{J}\to0$ montre que le faisceau
d'idéaux de $\mathscr{O}_{\mathfrak{X}}$ désigné ici par
$\mathfrak{J}^\Delta$ coïncide avec celui qui avait été noté de la même
façon en (\hyperref[I.10.3.1-fr]{10.3.1}), en vertu de
(\hyperref[I.10.3.2-fr]{10.3.2}).

\begin{proposition}[10.10.3]
\label{I.10.10.3-fr}
Sous les hypothèses de (\hyperref[I.10.10.1-fr]{10.10.1}),
$\mathscr{O}_{\mathfrak{X}}$ est un faisceau cohérent d'anneaux.
\end{proposition}

Si $f\in A$, on sait que $A_{\{f\}}$ est un anneau adique noethérien
(\hyperref[0.7.6.11-fr]{0, 7.6.11}) et comme la question est locale, on
est ramené (\hyperref[I.10.1.4-fr]{10.1.4}) à prouver que le noyau d'un
homomorphisme
$v:\mathscr{O}_{\mathfrak{X}}^n\to\mathscr{O}_{\mathfrak{X}}$ est un
$\mathscr{O}_{\mathfrak{X}}$-Module de type fini. On a alors
$v=u^\Delta$, où $u$ est un $A$-homomorphisme $A^n\to A$
(\hyperref[I.10.10.2-fr]{10.10.2})~; comme $A$ est noethérien, le noyau
de $u$ est de type fini, autrement dit on a un homomorphisme
$A^m\xrightarrow{w}A^n$ tel que la suite
$A^m\xrightarrow{w}A^n\xrightarrow{u}A$ soit exacte. On en conclut
(\hyperref[I.10.10.2-fr]{10.10.2}) que la suite
$\mathscr{O}_{\mathfrak{X}}^m\xrightarrow{w^\Delta}
\mathscr{O}_{\mathfrak{X}}^n\xrightarrow{v}
\mathscr{O}_{\mathfrak{X}}$ est exacte, ce qui prouve que le noyau de
$v$ est de type fini.

\begin{env}[10.10.4]
\label{I.10.10.4-fr}
Avec les notations précédentes, posons
$A_n=A/\mathfrak{J}^{n+1}$, et soit $X_n$ le schéma affine
$\operatorname{Spec}(A_n)=(\mathfrak{X},
\mathscr{O}_{\mathfrak{X}}/\mathscr{J}^{n+1})$,
$\mathscr{J}=\mathfrak{J}^\Delta$ étant le faisceau d'idéaux de
définition de $\mathscr{O}_{\mathfrak{X}}$ correspondant à l'idéal
$\mathfrak{J}$. Soit $u_{mn}$ le morphisme de préschémas
$X_m\to X_n$ correspondant à l'homomorphisme canonique
$A_n\to A_m$ pour $m\leq n$~; le schéma formel $\mathfrak{X}$ est
limite inductive des $X_n$ pour les $u_{mn}$
(\hyperref[I.10.6.3-fr]{10.6.3}).
\end{env}

\begin{proposition}[10.10.5]
\label{I.10.10.5-fr}
Sous les hypothèses de (\hyperref[I.10.10.1-fr]{10.10.1}), soit
$\mathscr{F}$ un $\mathscr{O}_{\mathfrak{X}}$-Module. Les conditions
suivantes sont équivalentes~:
\begin{enumerate}
  \item[a)] $\mathscr{F}$ est un $\mathscr{O}_{\mathfrak{X}}$-Module
    cohérent.
  \item[b)] $\mathscr{F}$ est isomorphe à la limite projective
    (\hyperref[I.10.6.6-fr]{10.6.6}) d'une suite $(\mathscr{F}_n)$ de
    $\mathscr{O}_{X_n}$-Modules cohérents tels que
    $u_{mn}^*(\mathscr{F}_n)=\mathscr{F}_m$.
  \item[c)] Il existe un $A$-module de type fini $M$ (déterminé à un
    isomorphisme canonique près par
    (\hyperref[I.10.10.2-fr]{10.10.2, (i)})) tel que $\mathscr{F}$ soit
    isomorphe à $M^\Delta$.
\end{enumerate}
\end{proposition}

Montrons d'abord que b) implique c). On a
$\mathscr{F}_n=\widetilde{M_n}$, où $M_n$ est un $A_n$-module de type
fini, et l'hypothèse entraîne que
$M_m=M_n\otimes_{A_n}A_m$ pour $m\leq n$
(\hyperref[I.1.6.5-fr]{1.6.5})~; les $M_n$ forment donc un système
projectif pour les di-homomorphismes canoniques $M_n\to M_m$
($m\leq n$), et il résulte aussitôt de la définition des $A_n$ que ce
système projectif vérifie les conditions de
(\hyperref[0.7.2.9-fr]{0, 7.2.9})~; sa limite projective $M$ est par
suite un $A$-module de type fini tel que
$M_n=M\otimes_A A_n$ pour tout $n$. On en déduit que $\mathscr{F}_n$
est induit sur $X_n$ par
$\widetilde M\otimes_{\mathscr{O}_X}
(\mathscr{O}_X/\widetilde{\mathfrak{J}}^{n+1})$, donc
$\mathscr{F}=M^\Delta$ par définition
(\hyperref[I.10.8.4-fr]{10.8.4}).

Inversement, c) entraîne b)~; en effet, si $u_n$ est le morphisme
d'immersion $X_n\to X$,
$u_n^*(\widetilde M)=(M\otimes_A A_n)^{\sim}$ est induit sur $X_n$ par
$\widetilde M\otimes_{\mathscr{O}_X}
(\mathscr{O}_X/\widetilde{\mathfrak{J}}^{n+1})$, et
$M^\Delta=\varprojlim u_n^*(\widetilde M)$ par définition
(\hyperref[I.10.8.4-fr]{10.8.4})~; comme
$u_m=u_n\circ u_{mn}$ pour $m\leq n$, les
$\mathscr{F}_n=u_n^*(\widetilde M)$ vérifient les conditions de b),
d'où notre assertion.

Montrons maintenant que c) implique a)~: en effet, on a par définition
$\mathscr{O}_{\mathfrak{X}}=A^\Delta$~; $M$ étant conoyau d'un
homomorphisme $A^m\to A^n$, il résulte de
(\hyperref[I.10.10.2-fr]{10.10.2}) que $M^\Delta$ est conoyau d'un
homomorphisme
$\mathscr{O}_{\mathfrak{X}}^m\to\mathscr{O}_{\mathfrak{X}}^n$, et
comme le faisceau d'anneaux $\mathscr{O}_{\mathfrak{X}}$ est cohérent
(\hyperref[I.10.10.3-fr]{10.10.3}), il en est de même de $M^\Delta$
(\hyperref[0.5.3.4-fr]{0, 5.3.4}).
\oldpage[I]{203}

Enfin, a) entraîne b). Considéré comme
$\mathscr{O}_{\mathfrak{X}}$-Module, on a
$\mathscr{O}_{X_n}=\mathscr{O}_{\mathfrak{X}}/
\mathscr{J}^{n+1}=A_n^\Delta$~; $\mathscr{F}_n=\mathscr{F}
\otimes_{\mathscr{O}_{\mathfrak{X}}}\mathscr{O}_{X_n}$ est un
$\mathscr{O}_{\mathfrak{X}}$-Module cohérent
(\hyperref[0.5.3.5-fr]{0, 5.3.5}), et comme c'est aussi un
$\mathscr{O}_{X_n}$-Module et que $\mathscr{J}^{n+1}$ est cohérent,
on en conclut que $\mathscr{F}_n$ est un
$\mathscr{O}_{X_n}$-Module cohérent
(\hyperref[0.5.3.10-fr]{0, 5.3.10}), et il est immédiat que
$u_{mn}^*(\mathscr{F}_n)=\mathscr{F}_m$ pour $m\leq n$ (en se
souvenant que l'application continue $X_m\to X_n$ des espaces
sous-jacents est l'identité de $\mathfrak{X}$). Le faisceau
$\mathscr{G}=\varprojlim\mathscr{F}_n$ est donc un
$\mathscr{O}_{\mathfrak{X}}$-Module cohérent, puisqu'on a vu que b)
entraîne a). Les homomorphismes canoniques
$\mathscr{F}\to\mathscr{F}_n$ forment un système projectif, qui par
passage à la limite donne un homomorphisme canonique
$w:\mathscr{F}\to\mathscr{G}$, et tout revient à démontrer que $w$ est
bijectif. La question étant maintenant locale, on peut se borner au cas
où $\mathscr{F}$ est conoyau d'un homomorphisme
$\mathscr{O}_{\mathfrak{X}}^p\to\mathscr{O}_{\mathfrak{X}}^q$~; cet
homomorphisme étant de la forme $v^\Delta$, où $v$ est un homomorphisme
$A^m\to A^n$ (\hyperref[I.10.10.2-fr]{10.10.2}), $\mathscr{F}$ est
isomorphe à $M^\Delta$, où $M=\operatorname{Coker}v$
(\hyperref[I.10.10.2-fr]{10.10.2}). On a alors, en vertu de
(\hyperref[I.10.10.2-fr]{10.10.2}),
$\mathscr{F}_n=M^\Delta\otimes_{\mathscr{O}_{\mathfrak{X}}}A_n^\Delta
=(M\otimes_A A_n)^\Delta$, et comme la topologie $\mathfrak{J}$-adique
sur $M\otimes_A A_n$ est discrète, on a
$(M\otimes_A A_n)^\Delta=(M\otimes_A A_n)^{\sim}$ (en tant que
$\mathscr{O}_{X_n}$-Module)~; on a vu plus haut que
$M^\Delta=\varprojlim\mathscr{F}_n$ et $w$ est donc bien dans ce cas
l'identité.
\par\hfill C.Q.F.D.\par

\begin{corollary}[10.10.6]
\label{I.10.10.6-fr}
Si $\mathscr{F}$ vérifie la condition b) de
(\hyperref[I.10.10.5-fr]{10.10.5}), le système projectif
$(\mathscr{F}_n)$ est isomorphe au système des
$\mathscr{F}\otimes_{\mathscr{O}_{\mathfrak{X}}}\mathscr{O}_{X_n}$.
\end{corollary}

\begin{env}[10.10.7]
\label{I.10.10.7-fr}
Soient maintenant $A$, $B$ deux anneaux adiques noethériens,
$\varphi:B\to A$ un homomorphisme continu~; on désignera par
$\mathfrak{J}$ (resp. $\mathfrak{K}$) un idéal de définition de $A$
(resp. $B$), tels que $\varphi(\mathfrak{K})\subset\mathfrak{J}$, et
on posera $X=\operatorname{Spec}(A)$, $Y=\operatorname{Spec}(B)$,
$\mathfrak{X}=\operatorname{Spf}(A)$,
$\mathfrak{Y}=\operatorname{Spf}(B)$. Soient $f:X\to Y$ le morphisme
de préschémas correspondant à $\varphi$
(\hyperref[I.1.6.1-fr]{1.6.1}), $\widehat f:\mathfrak{X}\to
\mathfrak{Y}$ son prolongement aux complétés
(\hyperref[I.10.9.1-fr]{10.9.1}), qui est aussi le morphisme de
préschémas formels correspondant à $\varphi$
(\hyperref[I.10.2.2-fr]{10.2.2}).
\end{env}

\begin{proposition}[10.10.8]
\label{I.10.10.8-fr}
Pour tout $B$-module $N$ de type fini, il existe un isomorphisme
canonique fonctoriel de $\mathscr{O}_{\mathfrak{X}}$-Modules
\[
  \widehat f^*(N^\Delta)
  \xrightarrow{\sim}
  (N\otimes_B A)^\Delta.
\]
\end{proposition}

En effet, en désignant par $i_X:\mathfrak{X}\to X$ et
$i_Y:\mathfrak{Y}\to Y$ les morphismes canoniques, on a
(\hyperref[I.10.8.8-fr]{10.8.8}), à des isomorphismes canoniques
fonctoriels près, $N^\Delta=i_Y^*(\widetilde N)$ et
\[
  (N\otimes_B A)^\Delta
  =i_X^*((N\otimes_B A)^{\sim})
  =i_X^*(f^*(\widetilde N))
\]
(\hyperref[I.1.6.5-fr]{1.6.5})~; la proposition résulte donc de la
commutativité du diagramme (\hyperref[I.10.9.2-fr]{10.9.2}).

\begin{corollary}[10.10.9]
\label{I.10.10.9-fr}
Pour tout idéal $\mathfrak b$ de $B$, on a
$\widehat f^*(\mathfrak b^\Delta)\mathscr{O}_{\mathfrak{X}}
=(\mathfrak bA)^\Delta$.
\end{corollary}

En effet, soit $j$ l'injection canonique $\mathfrak b\to B$, à laquelle
il correspond l'injection canonique
$j^\Delta:\mathfrak b^\Delta\to\mathscr{O}_{\mathfrak{Y}}$ de
faisceaux de $\mathscr{O}_{\mathfrak{Y}}$-Modules~; par définition,
$\widehat f^*(\mathfrak b^\Delta)\mathscr{O}_{\mathfrak{X}}$ est
l'image de l'homomorphisme
$\widehat f^*(j^\Delta):\widehat f^*(\mathfrak b^\Delta)\to
\mathscr{O}_{\mathfrak{X}}=\widehat f^*(\mathscr{O}_{\mathfrak{Y}})$~;
mais cet homomorphisme s'identifie à
$(j\otimes1)^\Delta:(\mathfrak b\otimes_B A)^\Delta\to
\mathscr{O}_{\mathfrak{X}}=(B\otimes_B A)^\Delta$ d'après
(\hyperref[I.10.10.8-fr]{10.10.8}). Comme l'image de $j\otimes1$ est
l'idéal $\mathfrak bA$ de $A$, l'image de $(j\otimes1)^\Delta$ est
donc $(\mathfrak bA)^\Delta$ en vertu de
(\hyperref[I.10.10.2-fr]{10.10.2}), d'où la conclusion.
\oldpage[I]{204}

\subsection{Faisceaux cohérents sur les préschémas formels.}
\label{subsection:I.10.11-fr}

\begin{proposition}[10.11.1]
\label{I.10.11.1-fr}
Si $\mathfrak{X}$ est un préschéma formel localement noethérien, le
faisceau d'anneaux $\mathscr{O}_{\mathfrak{X}}$ est cohérent et tout
faisceau d'idéaux de définition de $\mathfrak{X}$ est cohérent.
\end{proposition}

La question étant locale, on est ramené au cas d'un schéma formel affine
noethérien, et la proposition résulte donc de
(\hyperref[I.10.10.3-fr]{10.10.3}) et
(\hyperref[I.10.10.5-fr]{10.10.5}).

\begin{env}[10.11.2]
\label{I.10.11.2-fr}
Soient $\mathfrak{X}$ un préschéma formel localement noethérien,
$\mathscr{J}$ un faisceau d'idéaux de définition de $\mathfrak{X}$,
$X_n$ le préschéma (usuel) localement noethérien
$(\mathfrak{X},\mathscr{O}_{\mathfrak{X}}/\mathscr{J}^{n+1})$, de
sorte que $\mathfrak{X}$ est limite inductive de la suite $(X_n)$ pour
les morphismes canoniques $u_{mn}:X_m\to X_n$
(\hyperref[I.10.6.3-fr]{10.6.3}). Avec ces notations~:
\end{env}

\begin{theorem}[10.11.3]
\label{I.10.11.3-fr}
Pour qu'un $\mathscr{O}_{\mathfrak{X}}$-Module $\mathscr{F}$ soit
cohérent, il faut et il suffit qu'il soit isomorphe à une limite
projective d'une suite $(\mathscr{F}_n)$, où $\mathscr{F}_n$ est un
$\mathscr{O}_{X_n}$-Module cohérent, tel que
$u_{mn}^*(\mathscr{F}_n)=\mathscr{F}_m$ pour $m\leq n$
(\hyperref[I.10.6.6-fr]{10.6.6}). Le système projectif
$(\mathscr{F}_n)$ est alors isomorphe au système des
$u_n^*(\mathscr{F})=\mathscr{F}
\otimes_{\mathscr{O}_{\mathfrak{X}}}\mathscr{O}_{X_n}$, $u_n$ étant
le morphisme canonique $X_n\to\mathfrak{X}$.
\end{theorem}

La question étant locale, on est ramené au cas où $\mathfrak{X}$ est
un schéma formel affine noethérien, et le théorème est alors conséquence
de (\hyperref[I.10.10.5-fr]{10.10.5}) et
(\hyperref[I.10.10.6-fr]{10.10.6}).

On peut donc dire que \emph{la donnée d'un
$\mathscr{O}_{\mathfrak{X}}$-Module cohérent équivaut à celle d'un
système projectif $(\mathscr{F}_n)$ de
$\mathscr{O}_{X_n}$-Modules cohérents tels que
$u_{mn}^*(\mathscr{F}_n)=\mathscr{F}_m$ pour $m\leq n$}.

\begin{corollary}[10.11.4]
\label{I.10.11.4-fr}
Si $\mathscr{F}$ et $\mathscr{G}$ sont deux
$\mathscr{O}_{\mathfrak{X}}$-Modules cohérents, on peut (avec les
notations de (\hyperref[I.10.11.3-fr]{10.11.3})) définir un
isomorphisme canonique fonctoriel
\[
  \label{I.10.11.4.1-fr}
  \operatorname{Hom}_{\mathscr{O}_{\mathfrak{X}}}
  (\mathscr{F},\mathscr{G})
  \xrightarrow{\sim}
  \varprojlim_n
  \operatorname{Hom}_{\mathscr{O}_{X_n}}
  (\mathscr{F}_n,\mathscr{G}_n).
  \tag{10.11.4.1}
\]
\end{corollary}

La limite projective du second membre doit s'entendre pour les
applications
$\theta_n\mapsto u_{mn}^*(\theta_n)$ ($m\leq n$) de
$\operatorname{Hom}_{\mathscr{O}_{X_n}}
(\mathscr{F}_n,\mathscr{G}_n)$ dans
$\operatorname{Hom}_{\mathscr{O}_{X_m}}
(\mathscr{F}_m,\mathscr{G}_m)$. L'homomorphisme
(\hyperref[I.10.11.4.1-fr]{10.11.4.1}) fait correspondre à un élément
$\theta\in\operatorname{Hom}_{\mathscr{O}_{\mathfrak{X}}}
(\mathscr{F},\mathscr{G})$ la suite $(u_n^*(\theta))$~; on voit
aussitôt qu'on en définit un homomorphisme réciproque du précédent en
faisant correspondre au système projectif
$(\theta_n)\in\varprojlim_n
\operatorname{Hom}_{\mathscr{O}_{X_n}}
(\mathscr{F}_n,\mathscr{G}_n)$
sa limite projective dans
$\operatorname{Hom}_{\mathscr{O}_{\mathfrak{X}}}
(\mathscr{F},\mathscr{G})$, compte tenu de
(\hyperref[I.10.11.3-fr]{10.11.3}).

\begin{corollary}[10.11.5]
\label{I.10.11.5-fr}
Pour qu'un homomorphisme $\theta:\mathscr{F}\to\mathscr{G}$ soit
surjectif, il faut et il suffit que l'homomorphisme correspondant
$\theta_0=u_0^*(\theta):\mathscr{F}_0\to\mathscr{G}_0$ le soit.
\end{corollary}

La question étant locale, on est ramené au cas où
$\mathfrak{X}=\operatorname{Spf}(A)$, $A$ étant adique noethérien,
$\mathscr{F}=M^\Delta$, $\mathscr{G}=N^\Delta$ et $\theta=u^\Delta$,
où $M$ et $N$ sont des $A$-modules de type fini et $u$ un
homomorphisme $M\to N$~; on a en outre alors
$\theta_0=\widetilde{u_0}$, où $u_0$ est l'homomorphisme
$u\otimes1:M\otimes_A A/\mathfrak{J}\to
N\otimes_A A/\mathfrak{J}$~;
la conclusion résulte de ce que $\theta$ et $u$ (resp. $\theta_0$ et
$u_0$) sont simultanément surjectifs
(\hyperref[I.1.3.9-fr]{1.3.9} et
\hyperref[I.10.10.2-fr]{10.10.2}) et de ce que $u$ et $u_0$ sont
simultanément surjectifs (\hyperref[0.7.1.14-fr]{0, 7.1.14}).

\begin{env}[10.11.6]
\label{I.10.11.6-fr}
Le th. (\hyperref[I.10.11.3-fr]{10.11.3}) montre qu'on peut considérer
tout $\mathscr{O}_{\mathfrak{X}}$-Module cohérent $\mathscr{F}$ comme
un $\mathscr{O}_{\mathfrak{X}}$-Module \emph{topologique}, en le
considérant comme limite projective des faisceaux de groupes
\emph{pseudo-discrets} $\mathscr{F}_n$
(\hyperref[0.3.8.1-fr]{0, 3.8.1}). Il résulte alors de
(\hyperref[I.10.11.4-fr]{10.11.4}) que tout homomorphisme
$u:\mathscr{F}\to\mathscr{G}$ de
$\mathscr{O}_{\mathfrak{X}}$-Modules cohérents est automatiquement
\emph{continu}

\oldpage[I]{205}

(\hyperref[0.3.8.2-fr]{0, 3.8.2}). En outre, si $\mathscr{H}$ est un
sous-$\mathscr{O}_{\mathfrak{X}}$-Module cohérent d'un
$\mathscr{O}_{\mathfrak{X}}$-Module cohérent $\mathscr{F}$, pour tout
ouvert $U\subset\mathfrak{X}$, $\Gamma(U,\mathscr{H})$ est un sous-groupe
\emph{fermé} du groupe topologique $\Gamma(U,\mathscr{F})$ car le foncteur
$\Gamma$ étant exact à gauche, $\Gamma(U,\mathscr{H})$ est le noyau de
l'homomorphisme
$\Gamma(U,\mathscr{F})\to\Gamma(U,\mathscr{F}/\mathscr{H})$, qui est
\emph{continu} d'après ce qui précède, puisque $\mathscr{F}/\mathscr{H}$
est cohérent (\hyperref[0.5.3.4-fr]{0, 5.3.4})~; notre assertion résulte
de ce que $\Gamma(U,\mathscr{F}/\mathscr{H})$ est un groupe topologique
séparé.
\end{env}

\begin{proposition}[10.11.7]
\label{I.10.11.7-fr}
Soient $\mathscr{F}$ et $\mathscr{G}$ deux
$\mathscr{O}_{\mathfrak{X}}$-Modules cohérents. On peut définir (avec les
notations de (\hyperref[I.10.11.3-fr]{10.11.3})) des isomorphismes
canoniques fonctoriels de $\mathscr{O}_{\mathfrak{X}}$-Modules
topologiques (\hyperref[I.10.11.6-fr]{10.11.6})
\[
  \label{I.10.11.7.1-fr}
  \mathscr{F}\otimes_{\mathscr{O}_{\mathfrak{X}}}\mathscr{G}
  \xrightarrow{\sim}
  \varprojlim_n
  (\mathscr{F}_n\otimes_{\mathscr{O}_{X_n}}\mathscr{G}_n)
  \tag{10.11.7.1}
\]
\[
  \label{I.10.11.7.2-fr}
  \mathscr{H}\!om_{\mathscr{O}_{\mathfrak{X}}}
  (\mathscr{F},\mathscr{G})
  \xrightarrow{\sim}
  \varprojlim_n
  \mathscr{H}\!om_{\mathscr{O}_{X_n}}(\mathscr{F}_n,\mathscr{G}_n)
  \tag{10.11.7.2}
\]
\end{proposition}

L'existence de l'isomorphisme
(\hyperref[I.10.11.7.1-fr]{10.11.7.1}) résulte de la formule
\[
  \mathscr{F}_n\otimes_{\mathscr{O}_{X_n}}\mathscr{G}_n
  =
  (\mathscr{F}\otimes_{\mathscr{O}_{\mathfrak{X}}}
  \mathscr{O}_{X_n})
  \otimes_{\mathscr{O}_{X_n}}
  (\mathscr{G}\otimes_{\mathscr{O}_{\mathfrak{X}}}
  \mathscr{O}_{X_n})
  =
  (\mathscr{F}\otimes_{\mathscr{O}_{\mathfrak{X}}}\mathscr{G})
  \otimes_{\mathscr{O}_{\mathfrak{X}}}\mathscr{O}_{X_n}
\]
et de (\hyperref[I.10.11.3-fr]{10.11.3}). L'isomorphisme
(\hyperref[I.10.11.7.2-fr]{10.11.7.2}) où les deux membres sont
considérés comme faisceaux de modules sans topologie, résulte de la
définition des sections de
$\mathscr{H}\!om_{\mathscr{O}_{\mathfrak{X}}}
(\mathscr{F},\mathscr{G})$ et
$\mathscr{H}\!om_{\mathscr{O}_{X_n}}(\mathscr{F}_n,\mathscr{G}_n)$ et de
l'existence de l'isomorphisme
(\hyperref[I.10.11.4.1-fr]{10.11.4.1}), appliqué au préschéma induit sur
un ouvert formel affine noethérien arbitraire de $\mathfrak{X}$. Reste à
prouver que l'isomorphisme
(\hyperref[I.10.11.7.2-fr]{10.11.7.2}) est bicontinu au-dessus d'un
ensemble quasi-compact, et on est donc ramené au cas où
$\mathfrak{X}=\operatorname{Spf}(A)$, $A$ étant adique noethérien, d'où
(\hyperref[I.10.10.5-fr]{10.10.5})
$\mathscr{F}=M^\Delta$, $\mathscr{G}=N^\Delta$, $M$, $N$ étant des
$A$-modules de type fini~; compte tenu de
(\hyperref[I.10.10.2.1-fr]{10.10.2.1}),
(\hyperref[I.10.10.2.3-fr]{10.10.2.3}) et
(\hyperref[I.1.3.12-fr]{1.3.12, (ii)}), on est ramené à montrer que
l'isomorphisme canonique
$\operatorname{Hom}_A(M,N)\xrightarrow{\sim}
\varprojlim_n\operatorname{Hom}_{A_n}(M_n,N_n)$ (avec
$M_n=M\otimes_A A_n$, $N_n=N\otimes_A A_n$) est continu, ce qui a été
prouvé dans (\hyperref[0.7.8.2-fr]{0, 7.8.2}).

\begin{env}[10.11.8]
\label{I.10.11.8-fr}
Comme $\operatorname{Hom}_{\mathscr{O}_{\mathfrak{X}}}
(\mathscr{F},\mathscr{G})$ est le groupe des sections du faisceau de
groupes topologiques
$\mathscr{H}\!om_{\mathscr{O}_{\mathfrak{X}}}
(\mathscr{F},\mathscr{G})$, il est muni d'une topologie de groupe. Si
$\mathfrak{X}$ est \emph{noethérien}, il résulte de
(\hyperref[I.10.11.7.2-fr]{10.11.7.2}) qu'un système fondamental de
voisinages de $0$ dans ce groupe s'obtient en prenant les sous-groupes
$\operatorname{Hom}_{\mathscr{O}_{\mathfrak{X}}}
(\mathscr{F},\mathscr{J}^n\mathscr{G})$ ($n$ arbitraire).
\end{env}

\begin{proposition}[10.11.9]
\label{I.10.11.9-fr}
Soient $\mathfrak{X}$ un préschéma formel noethérien, $\mathscr{F}$ et
$\mathscr{G}$ deux $\mathscr{O}_{\mathfrak{X}}$-Modules cohérents. Dans
le groupe topologique
$\operatorname{Hom}_{\mathscr{O}_{\mathfrak{X}}}
(\mathscr{F},\mathscr{G})$ les homomorphismes surjectifs (resp.
injectifs, bijectifs) forment une partie ouverte.
\end{proposition}

En vertu de (\hyperref[I.10.11.5-fr]{10.11.5}), l'ensemble des
homomorphismes surjectifs dans
$\operatorname{Hom}_{\mathscr{O}_{\mathfrak{X}}}
(\mathscr{F},\mathscr{G})$ est l'image réciproque par l'application
continue
$\operatorname{Hom}_{\mathscr{O}_{\mathfrak{X}}}
(\mathscr{F},\mathscr{G})\to
\operatorname{Hom}_{\mathscr{O}_{X_0}}(\mathscr{F}_0,\mathscr{G}_0)$
d'une partie du groupe discret
$\operatorname{Hom}_{\mathscr{O}_{X_0}}(\mathscr{F}_0,\mathscr{G}_0)$
d'où la première assertion. Pour démontrer la seconde, recouvrons
$\mathfrak{X}$ par un nombre fini d'ouverts formels affines noethériens
$U_i$. Pour que
$\theta\in\operatorname{Hom}_{\mathscr{O}_{\mathfrak{X}}}
(\mathscr{F},\mathscr{G})$ soit injectif, il faut et il suffit que toutes
ses images par les applications (continues) de restriction
$\operatorname{Hom}_{\mathscr{O}_{\mathfrak{X}}}
(\mathscr{F},\mathscr{G})\to
\operatorname{Hom}_{\mathscr{O}_{\mathfrak{X}}|U_i}
(\mathscr{F}|U_i,\mathscr{G}|U_i)$ le soient~; on est donc ramené au cas
affine, et alors cela a déjà été prouvé dans
(\hyperref[0.7.8.3-fr]{0, 7.8.3}).

\oldpage[I]{206}

\subsection{Morphismes adiques de préschémas formels.}
\label{subsection:I.10.12-fr}

\begin{env}[10.12.1]
\label{I.10.12.1-fr}
Soient $\mathfrak{X}$, $\mathfrak{S}$ deux préschémas formels
\emph{localement noethériens}~; nous dirons qu'un morphisme
$f:\mathfrak{X}\to\mathfrak{S}$ est \emph{adique} s'il existe un Idéal
de définition $\mathscr{J}$ de $\mathfrak{S}$ tel que
$\mathscr{K}=f^*(\mathscr{J})\mathscr{O}_{\mathfrak{X}}$ soit un Idéal de
définition de $\mathfrak{X}$~; on dit aussi alors que $\mathfrak{X}$ est
un \emph{$\mathfrak{S}$-préschéma adique} (pour $f$). Lorsqu'il en est
ainsi, pour \emph{tout} Idéal de définition $\mathscr{J}_1$ de
$\mathfrak{S}$,
$\mathscr{K}_1=f^*(\mathscr{J}_1)\mathscr{O}_{\mathfrak{X}}$ est un Idéal
de définition de $\mathfrak{X}$. En effet, la question étant locale, on
peut supposer $\mathfrak{X}$ et $\mathfrak{S}$ affines noethériens~; il
existe donc un entier $n$ tel que $\mathscr{J}^n\subset\mathscr{J}_1$ et
$\mathscr{J}_1^n\subset\mathscr{J}$
(\hyperref[I.10.3.6-fr]{10.3.6} et
\hyperref[0.7.1.4-fr]{0, 7.1.4}), d'où
$\mathscr{K}^n\subset\mathscr{K}_1$ et
$\mathscr{K}_1^n\subset\mathscr{K}$. La première de ces relations prouve
que $\mathscr{K}_1=\mathfrak{K}_1^\Delta$, où $\mathfrak{K}_1$ est un
idéal ouvert de
$A=\Gamma(\mathfrak{X},\mathscr{O}_{\mathfrak{X}})$, et la seconde prouve
que $\mathfrak{K}_1$ est un idéal de définition de $A$
(\hyperref[0.7.1.4-fr]{0, 7.1.4}), d'où notre assertion.
\end{env}

Il résulte aussitôt de ce qui précède que si $\mathfrak{X}$ et
$\mathfrak{Y}$ sont deux $\mathfrak{S}$-préschémas adiques, \emph{tout
$\mathfrak{S}$-morphisme $u:\mathfrak{X}\to\mathfrak{Y}$ est adique}~:
en effet, si $f:\mathfrak{X}\to\mathfrak{S}$,
$g:\mathfrak{Y}\to\mathfrak{S}$ sont les morphismes structuraux, et
$\mathscr{J}$ un Idéal de définition de $\mathfrak{S}$, on a
$f=g\circ u$, donc
$u^*(g^*(\mathscr{J})\mathscr{O}_{\mathfrak{Y}})
\mathscr{O}_{\mathfrak{X}}=f^*(\mathscr{J})
\mathscr{O}_{\mathfrak{X}}$ est un Idéal de définition de
$\mathfrak{X}$, et par hypothèse
$g^*(\mathscr{J})\mathscr{O}_{\mathfrak{Y}}$ est un Idéal de définition
de $\mathfrak{Y}$.

\begin{env}[10.12.2]
\label{I.10.12.2-fr}
Dans ce qui suit, nous supposerons fixé un préschéma formel localement
noethérien $\mathfrak{S}$ et un Idéal de définition $\mathscr{J}$ de
$\mathfrak{S}$~; nous poserons
$S_n=(\mathfrak{S},\mathscr{O}_{\mathfrak{S}}/\mathscr{J}^{n+1})$. Les
$\mathfrak{S}$-préschémas adiques (localement noethériens) forment
évidemment une \emph{catégorie}. Nous dirons qu'un système inductif
$(X_n)$ de $S_n$-préschémas (usuels) localement noethériens est un
\emph{$(S_n)$-système inductif adique} si les morphismes structuraux
$f_n:X_n\to S_n$ sont tels que, pour $m\leq n$, les diagrammes
\[
  \label{I.10.12.2.1-fr}
  \xymatrix{
    X_n \ar[d]_{f_n}
    & X_m \ar[l] \ar[d]^{f_m}\\
    S_n
    & S_m \ar[l]
  }
  \tag{10.12.2.1}
\]
soient commutatifs et \emph{identifient $X_m$ au produit
$X_n\times_{S_n}S_m=(X_n)_{(S_m)}$}. Les systèmes inductifs adiques
forment une \emph{catégorie}~: il suffit en effet de définir un morphisme
$(X_n)\to(Y_n)$ de tels systèmes comme un \emph{système inductif de
$S_n$-morphismes $u_n:X_n\to Y_n$} tel que $u_m$ s'identifie à
$(u_n)_{(S_m)}$ pour $m\leq n$. Cela étant~:
\end{env}

\begin{theorem}[10.12.3]
\label{I.10.12.3-fr}
Il y a équivalence canonique entre la catégorie des
$\mathfrak{S}$-préschémas adiques et la catégorie des
$(S_n)$-systèmes inductifs adiques.
\end{theorem}

L'équivalence en question s'obtient de la façon suivante~: si
$\mathfrak{X}$ est un $\mathfrak{S}$-préschéma adique,
$f:\mathfrak{X}\to\mathfrak{S}$ le morphisme structural,
$\mathscr{K}=f^*(\mathscr{J})\mathscr{O}_{\mathfrak{X}}$ est un Idéal de
définition de $\mathfrak{X}$ et on fait correspondre à $\mathfrak{X}$ le
système inductif des
$X_n=(\mathfrak{X},\mathscr{O}_{\mathfrak{X}}/\mathscr{K}^{n+1})$, le
morphisme structural $f_n:X_n\to S_n$ correspondant à $f$
(\hyperref[I.10.5.6-fr]{10.5.6}). Montrons d'abord que $(X_n)$ est un
\emph{système inductif adique}~: si $f=(\psi,\theta)$, on a
$\psi^*(\mathscr{J})\mathscr{O}_{\mathfrak{X}}=\mathscr{K}$, donc
$\psi^*(\mathscr{J}^n)\mathscr{O}_{\mathfrak{X}}=\mathscr{K}^n$ pour tout
$n$, et (par l'exactitude du foncteur $\psi^*$)
$\mathscr{K}^{m+1}/\mathscr{K}^{n+1}
=\psi^*(\mathscr{J}^{m+1}/\mathscr{J}^{n+1})
(\mathscr{O}_{\mathfrak{X}}/\mathscr{K}^{n+1})$ pour $m\leq n$~; notre
conclusion résulte donc de (\hyperref[I.4.4.5-fr]{4.4.5}). Il est immédiat
en outre de vérifier qu'à un $\mathfrak{S}$-morphisme
$u:\mathfrak{X}\to\mathfrak{Y}$ de $\mathfrak{S}$-préschémas adiques
correspond (avec des notations évidentes)

\oldpage[I]{207}

un système inductif de $S_n$-morphismes $u_n:X_n\to Y_n$ tel que $u_m$
s'identifie à $(u_n)_{(S_m)}$ pour $m\leq n$.

Le fait qu'on a bien défini ainsi une équivalence résultera de la proposition
plus précise suivante~:

\begin{proposition}[10.12.3.1]
\label{I.10.12.3.1-fr}
Soit $(X_n)$ un système inductif de $S_n$-préschémas~; on suppose que les
morphismes structuraux $f_n:X_n\to S_n$ sont tels que les diagrammes
(\hyperref[I.10.12.2.1-fr]{10.12.2.1}) soient commutatifs et identifient
$X_m$ à $X_n\times_{S_n}S_m$ pour $m\leq n$. Alors le système inductif
$(X_n)$ vérifie les conditions b) et c) de
(\hyperref[I.10.6.3-fr]{10.6.3})~; soient $\mathfrak{X}$ sa limite inductive,
$f:\mathfrak{X}\to\mathfrak{S}$ le morphisme limite inductive du système
inductif $(f_n)$. Alors, si $X_0$ est localement noethérien,
$\mathfrak{X}$ est localement noethérien et $f$ est un morphisme adique.
\end{proposition}

Comme le faisceau d'idéaux de $\mathscr{O}_{S_n}$ qui définit le
sous-préschéma $S_m$ de $S_n$ est nilpotent, il en est de même, en vertu de
(\hyperref[I.4.4.5-fr]{4.4.5}) du faisceau d'idéaux de
$\mathscr{O}_{X_n}$ définissant le sous-préschéma $X_m$ de $X_n$, donc les
conditions de (\hyperref[I.10.6.3-fr]{10.6.3}) sont bien vérifiées. La
question est par suite locale sur $\mathfrak{X}$ et $\mathfrak{S}$ et on
peut supposer que $\mathfrak{S}=\operatorname{Spf}(A)$,
$\mathscr{J}=\mathfrak{J}^\Delta$, $A$ étant un anneau $\mathfrak{J}$-adique
noethérien, et $X_n=\operatorname{Spec}(B_n)$~; si
$A_n=A/\mathfrak{J}^{n+1}$, l'hypothèse entraîne que $B_0$ est noethérien et
que si on pose $\mathfrak{J}_n=\mathfrak{J}/\mathfrak{J}^{n+1}$,
$B_m=B_n/\mathfrak{J}_n^{m+1}B_n$. Le noyau de $B_n\to B_0$ est donc
$\mathfrak{K}_n=\mathfrak{J}_nB_n$ et le noyau de $B_n\to B_m$ est
$\mathfrak{K}_n^{m+1}$ pour $m\leq n$~; en outre, comme $A_1$ est
noethérien, $\mathfrak{J}_1$ est de type fini sur $A_1$, donc
$\overline{\mathfrak{K}}_1=\mathfrak{K}_1/\mathfrak{K}_1^2$ est de type fini
sur $B_1$, et \emph{a fortiori} sur
$B_0=B_1/\mathfrak{K}_1$~; le fait que $\mathfrak{X}$ soit noethérien
résulte alors de (\hyperref[I.10.6.4-fr]{10.6.4})~; si
$B=\varprojlim B_n$, on a $\mathfrak{X}=\operatorname{Spf}(B)$, et si
$\mathfrak{K}$ est le noyau de $B\to B_0$,
$B_n=B/\mathfrak{K}^{n+1}$. Si
$\rho_n:A/\mathfrak{J}^{n+1}\to B/\mathfrak{K}^{n+1}$ est l'homomorphisme
correspondant à $f_n$, on a donc
\[
  \mathfrak{K}/\mathfrak{K}^{n+1}
  =(B/\mathfrak{K}^{n+1})\rho_n(\mathfrak{J}/\mathfrak{J}^{n+1})
\]
comme l'homomorphisme $\rho:A\to B$ correspondant à $f$ est égal à
$\varprojlim\rho_n$, l'idéal $\mathfrak{J}B$ de $B$ est dense dans
$\mathfrak{K}$, et comme tout idéal de $B$ est fermé
(\hyperref[0.7.3.5-fr]{0, 7.3.5}), on a $\mathfrak{K}=\mathfrak{J}B$. Si
$\mathscr{K}=\mathfrak{K}^\Delta$, la relation
$f^*(\mathscr{J})\mathscr{O}_{\mathfrak{X}}=\mathscr{K}$ résulte alors de
(\hyperref[I.10.10.9-fr]{10.10.9}) et achève la démonstration.

\begin{env}[10.12.3.2]
\label{I.10.12.3.2-fr}
L'équivalence précédente fournit, pour deux $\mathfrak{S}$-préschémas
adiques $\mathfrak{X}$, $\mathfrak{Y}$, une \emph{bijection canonique}
\[
  \operatorname{Hom}_{\mathfrak{S}}(\mathfrak{X},\mathfrak{Y})
  \xrightarrow{\sim}
  \varprojlim_n\operatorname{Hom}_{S_n}(X_n,Y_n)
\]
la limite projective étant relative aux applications
$u_n\to(u_n)_{(S_m)}$ pour $m\leq n$.
\end{env}

\subsection{Morphismes de type fini.}
\label{subsection:I.10.13-fr}

\begin{proposition}[10.13.1]
\label{I.10.13.1-fr}
Soient $\mathfrak{Y}$ un préschéma formel localement noethérien,
$\mathscr{K}$ un Idéal de définition de $\mathfrak{Y}$,
$f:\mathfrak{X}\to\mathfrak{Y}$ un morphisme de préschémas formels. Les
conditions suivantes sont équivalentes~:
\begin{enumerate}
  \item[a)] $\mathfrak{X}$ est localement noethérien, $f$ est un morphisme
    adique (\hyperref[I.10.12.1-fr]{10.12.1}) et si l'on pose
    $\mathscr{J}=f^*(\mathscr{K})\mathscr{O}_{\mathfrak{X}}$, le morphisme
    $f_0:(\mathfrak{X},\mathscr{O}_{\mathfrak{X}}/\mathscr{J})
    \to(\mathfrak{Y},\mathscr{O}_{\mathfrak{Y}}/\mathscr{K})$ déduit de
    $f$ est de type fini.
  \item[b)] $\mathfrak{X}$ est localement noethérien, et est limite
    inductive d'un $(Y_n)$-système inductif adique $(X_n)$ tel que le
    morphisme $X_0\to Y_0$ soit de type fini.

\oldpage[I]{208}

  \item[c)] Tout point de $\mathfrak{Y}$ possède un voisinage ouvert formel
    affine noethérien $V$ ayant la propriété suivante~:
    \begin{enumerate}
      \item[(\textbf{Q})] $f^{-1}(V)$ est réunion d'une famille finie
        d'ouverts formels affines noethériens $U_i$ tels que l'anneau adique
        noethérien $\Gamma(U_i,\mathscr{O}_{\mathfrak{X}})$ soit
        topologiquement isomorphe au quotient d'une algèbre de séries
        formelles restreintes (\hyperref[0.7.5.1-fr]{0, 7.5.1}) sur
        $\Gamma(V,\mathscr{O}_{\mathfrak{Y}})$, par un idéal (nécessairement
        fermé).
    \end{enumerate}
\end{enumerate}
\end{proposition}

Il est immédiat que a) entraîne b) en vertu de
(\hyperref[I.10.12.3-fr]{10.12.3}). Pour montrer que b) entraîne c), on
peut, puisque la question est locale sur $\mathfrak{Y}$, supposer que
$\mathfrak{Y}=\operatorname{Spf}(B)$, où $B$ est adique noethérien~; soit
$\mathscr{K}=\mathfrak{K}^\Delta$, $\mathfrak{K}$ étant un idéal de
définition de $B$. Comme par hypothèse, $X_0$ est de type fini sur $Y_0$,
$X_0$ est réunion finie d'ouverts affines $U_i$ tels que l'anneau $A_{i0}$
du schéma affine induit par $X_0$ sur $U_i$ soit une algèbre de type fini sur
l'anneau $B/\mathfrak{K}$ de $Y_0$
(\hyperref[I.6.3.2-fr]{6.3.2}). En vertu de
(\hyperref[I.5.1.9-fr]{5.1.9}), $U_i$ est aussi un ouvert affine dans chacun
des préschémas noethériens $X_n$, et si $A_{in}$ est l'anneau du schéma
affine induit par $X_n$ sur $U_i$, l'hypothèse b) entraîne que pour
$m\leq n$, $A_{im}$ est isomorphe à
$A_{in}/\mathfrak{K}^{m+1}A_{in}$. Par suite, le préschéma formel induit sur
$U_i$ par $\mathfrak{X}$ est isomorphe à $\operatorname{Spf}(A_i)$, où
$A_i=\varprojlim_n A_{in}$ (\hyperref[I.10.6.4-fr]{10.6.4})~; $A_i$ est un
anneau $\mathfrak{K}A_i$-adique, et $A_i/\mathfrak{K}A_i$, isomorphe à
$A_{i0}$, est une algèbre de type fini sur $B/\mathfrak{K}$. On en conclut
(\hyperref[0.7.5.5-fr]{0, 7.5.5}) que $A_i$ est topologiquement isomorphe à
un quotient d'une algèbre de séries formelles restreintes sur $B$ (par un
idéal nécessairement fermé, puisqu'une telle algèbre est noethérienne
(\hyperref[0.7.5.4-fr]{0, 7.5.4})).

Pour démontrer que c) entraîne a), on peut se limiter au cas où
$\mathfrak{X}=\operatorname{Spf}(A)$ est aussi affine, $A$ étant un anneau
adique noethérien, isomorphe au quotient d'une algèbre de séries formelles
restreintes sur $B$ par un idéal fermé. Alors
(\hyperref[0.7.5.5-fr]{0, 7.5.5}), $A/\mathfrak{K}A$ est une algèbre de type
fini sur $B/\mathfrak{K}$, et $\mathfrak{K}A=\mathfrak{J}$ est un idéal de
définition de $A$, donc, en vertu de
(\hyperref[I.10.10.9-fr]{10.10.9}), les conditions de a) sont satisfaites.

On notera que si les conditions de la prop.
(\hyperref[I.10.13.1-fr]{10.13.1}) sont remplies, la propriété a) est
valable pour \emph{tout Idéal de définition $\mathscr{K}$} de
$\mathfrak{Y}$ (en vertu de c)), et par suite, dans la propriété b),
\emph{tous les $f_n$ sont des morphismes de type fini}.

\begin{corollary}[10.13.2]
\label{I.10.13.2-fr}
Si les conditions de (\hyperref[I.10.13.1-fr]{10.13.1}) sont vérifiées,
tout ouvert formel affine noethérien $V$ de $\mathfrak{Y}$ possède la
propriété (\textbf{Q}) et si $\mathfrak{Y}$ est noethérien, il en est de
même de $\mathfrak{X}$.
\end{corollary}

Cela résulte aussitôt de (\hyperref[I.10.13.1-fr]{10.13.1}) et de
(\hyperref[I.6.3.2-fr]{6.3.2}).

\begin{definition}[10.13.3]
\label{I.10.13.3-fr}
Lorsque les propriétés équivalentes a), b), c) de
(\hyperref[I.10.13.1-fr]{10.13.1}) sont vérifiées, on dit que le morphisme
$f$ est \emph{de type fini}, ou que $\mathfrak{X}$ est un
\emph{$\mathfrak{Y}$-préschéma formel de type fini}, ou un
\emph{préschéma formel de type fini au-dessus de $\mathfrak{Y}$}.
\end{definition}

\begin{corollary}[10.13.4]
\label{I.10.13.4-fr}
Soient $\mathfrak{X}=\operatorname{Spf}(A)$,
$\mathfrak{Y}=\operatorname{Spf}(B)$ deux schémas formels affines
noethériens~; pour que $\mathfrak{X}$ soit de type fini sur
$\mathfrak{Y}$, il faut et il suffit que l'anneau adique noethérien $A$ soit
isomorphe au quotient d'une algèbre de séries formelles restreintes sur $B$
par un idéal fermé.
\end{corollary}

En effet, avec les notations de (\hyperref[I.10.13.1-fr]{10.13.1}), si
$\mathfrak{X}$ est de type fini sur $\mathfrak{Y}$,
$A/\mathfrak{K}A$ est alors une $(B/\mathfrak{K})$-algèbre de type fini en
vertu de (\hyperref[I.6.3.3-fr]{6.3.3}) et $\mathfrak{K}A$ est un idéal de
définition de $A$ (\hyperref[I.10.10.9-fr]{10.10.9}). On conclut donc par
(\hyperref[0.7.5.5-fr]{0, 7.5.5}).

\begin{proposition}[10.13.5]
\label{I.10.13.5-fr}
\medskip\noindent
\begin{enumerate}
  \item[\rm{(i)}] Le composé de deux morphismes de préschémas formels qui
    sont de type fini est de type fini.

\oldpage[I]{209}

  \item[\rm{(ii)}] Soient $\mathfrak{X}$, $\mathfrak{S}$,
    $\mathfrak{S}'$ trois préschémas formels localement noethériens (resp.
    noethériens), $f:\mathfrak{X}\to\mathfrak{S}$,
    $g:\mathfrak{S}'\to\mathfrak{S}$ deux morphismes. Si $f$ est de type
    fini, $\mathfrak{X}\times_{\mathfrak{S}}\mathfrak{S}'$ est localement
    noethérien (resp. noethérien) et est de type fini sur $\mathfrak{S}'$.
  \item[\rm{(iii)}] Soient $\mathfrak{S}$ un préschéma formel localement
    noethérien, $\mathfrak{X}'$, $\mathfrak{Y}'$ deux
    $\mathfrak{S}$-préschémas formels localement noethériens tels que
    $\mathfrak{X}'\times_{\mathfrak{S}}\mathfrak{Y}'$ soit localement
    noethérien. Si $\mathfrak{X}$, $\mathfrak{Y}$ sont des
    $\mathfrak{S}$-préschémas formels localement noethériens,
    $f:\mathfrak{X}\to\mathfrak{X}'$,
    $g:\mathfrak{Y}\to\mathfrak{Y}'$ deux $\mathfrak{S}$-morphismes de
    type fini, $\mathfrak{X}\times_{\mathfrak{S}}\mathfrak{Y}$ est
    localement noethérien et $f\times_{\mathfrak{S}}g$ est un
    $\mathfrak{S}$-morphisme de type fini.
\end{enumerate}
\end{proposition}

(iii) se déduit de (i) et (ii) par le raisonnement formel de
(\hyperref[I.3.5.1-fr]{3.5.1}) et il suffit donc de prouver (i) et (ii).

Soient $\mathfrak{X}$, $\mathfrak{Y}$, $\mathfrak{Z}$ trois préschémas
formels localement noethériens, $f:\mathfrak{X}\to\mathfrak{Y}$,
$g:\mathfrak{Y}\to\mathfrak{Z}$ deux morphismes de type fini. Si
$\mathscr{L}$ est un Idéal de définition de $\mathfrak{Z}$,
$\mathscr{K}=g^*(\mathscr{L})\mathscr{O}_{\mathfrak{Y}}$ en est un pour
$\mathfrak{Y}$ et
$\mathscr{J}=f^*(g^*(\mathscr{L}))\mathscr{O}_{\mathfrak{X}}$ en est un
pour $\mathfrak{X}$. Posons
$X_0=(\mathfrak{X},\mathscr{O}_{\mathfrak{X}}/\mathscr{J})$,
$Y_0=(\mathfrak{Y},\mathscr{O}_{\mathfrak{Y}}/\mathscr{K})$,
$Z_0=(\mathfrak{Z},\mathscr{O}_{\mathfrak{Z}}/\mathscr{L})$ et soient
$f_0:X_0\to Y_0$, $g_0:Y_0\to Z_0$ les morphismes correspondant à $f$ et
$g$. Comme par hypothèse $f_0$ et $g_0$ sont de type fini, il en est de même
de $g_0\circ f_0$ (\hyperref[I.6.3.4-fr]{6.3.4}) qui correspond à
$g\circ f$~; donc $g\circ f$ est de type fini par
(\hyperref[I.10.13.1-fr]{10.13.1}).

Sous les conditions de (ii), $\mathfrak{S}$ (resp. $\mathfrak{X}$,
$\mathfrak{S}'$) est limite inductive d'une suite $(S_n)$ (resp. $(X_n)$,
$(S'_n)$) de préschémas localement noethériens et on peut supposer
(\hyperref[I.10.13.1-fr]{10.13.1}) que
$X_m=X_n\times_{S_n}S_m$ pour $m\leq n$. Le préschéma formel
$\mathfrak{X}\times_{\mathfrak{S}}\mathfrak{S}'$ est alors limite inductive
des préschémas $X_n\times_{S_n}S'_n$
(\hyperref[I.10.7.4-fr]{10.7.4}), et on a
\[
  X_m\times_{S_m}S'_m
  =(X_n\times_{S_n}S_m)\times_{S_m}S'_m
  =(X_n\times_{S_n}S'_n)\times_{S'_n}S'_m.
\]
En outre, $X_0\times_{S_0}S'_0$ est localement noethérien puisque $X_0$
est de type fini sur $S_0$ (\hyperref[I.6.3.8-fr]{6.3.8}). On en conclut
d'abord (\hyperref[I.10.12.3.1-fr]{10.12.3.1}) que
$\mathfrak{X}\times_{\mathfrak{S}}\mathfrak{S}'$ est localement
noethérien~; en outre, comme $X_0\times_{S_0}S'_0$ est de type fini sur
$S'_0$ (\hyperref[I.6.3.8-fr]{6.3.8}), il résulte de
(\hyperref[I.10.12.3.1-fr]{10.12.3.1}) et de
(\hyperref[I.10.13.1-fr]{10.13.1}) que
$\mathfrak{X}\times_{\mathfrak{S}}\mathfrak{S}'$ est de type fini sur
$\mathfrak{S}'$, ce qui achève de prouver (ii) (l'assertion relative aux
préschémas noethériens étant conséquence immédiate de
(\hyperref[I.6.3.8-fr]{6.3.8})).

\begin{corollary}[10.13.6]
\label{I.10.13.6-fr}
Sous les hypothèses de (\hyperref[I.10.9.9-fr]{10.9.9}), si $f$ est un
morphisme de type fini, il en est de même de son prolongement $\widehat f$
aux complétés.
\end{corollary}

\subsection{Sous-préschémas fermés des préschémas formels.}
\label{subsection:I.10.14-fr}

\begin{proposition}[10.14.1]
\label{I.10.14.1-fr}
Soient $\mathfrak{X}$ un préschéma formel localement noethérien,
$\mathscr{A}$ un faisceau d'idéaux cohérent de
$\mathscr{O}_{\mathfrak{X}}$. Si $\mathfrak{Y}$ est le support (fermé) de
$\mathscr{O}_{\mathfrak{X}}/\mathscr{A}$, l'espace topologiquement annelé
$(\mathfrak{Y},(\mathscr{O}_{\mathfrak{X}}/\mathscr{A})|\mathfrak{Y})$ est
un préschéma formel localement noethérien, qui est noethérien si
$\mathfrak{X}$ l'est.
\end{proposition}

Notons que $\mathscr{O}_{\mathfrak{X}}/\mathscr{A}$ est cohérent en vertu de
(\hyperref[I.10.10.3-fr]{10.10.3}) et
(\hyperref[0.5.3.4-fr]{0, 5.3.4}), donc son support $\mathfrak{Y}$ est fermé
(\hyperref[0.5.2.2-fr]{0, 5.2.2}). Soit $\mathscr{J}$ un Idéal de définition
de $\mathfrak{X}$, et soit
$X_n=(\mathfrak{X},\mathscr{O}_{\mathfrak{X}}/\mathscr{J}^{n+1})$~; le
faisceau d'anneaux $\mathscr{O}_{\mathfrak{X}}/\mathscr{A}$ est limite
projective des faisceaux
$\mathscr{O}_{\mathfrak{X}}/(\mathscr{A}+\mathscr{J}^{n+1})
=(\mathscr{O}_{\mathfrak{X}}/\mathscr{A})
\otimes_{\mathscr{O}_{\mathfrak{X}}}
(\mathscr{O}_{\mathfrak{X}}/\mathscr{J}^{n+1})$
(\hyperref[I.10.11.3-fr]{10.11.3}), qui ont tous pour support
$\mathfrak{Y}$. Le faisceau
$(\mathscr{A}+\mathscr{J}^{n+1})/\mathscr{J}^{n+1}$ est un
$\mathscr{O}_{\mathfrak{X}}$-Module cohérent, puisque
$\mathscr{J}^{n+1}$ est cohérent, donc
$(\mathscr{A}+\mathscr{J}^{n+1})/\mathscr{J}^{n+1}$ est aussi un
$(\mathscr{O}_{\mathfrak{X}}/\mathscr{J}^{n+1})$-Module cohérent
(\hyperref[0.5.3.10-fr]{0, 5.3.10})~; si $Y_n$ est le sous-préschéma fermé
de $X_n$ défini par ce faisceau d'idéaux, il est immédiat que
$(\mathfrak{Y},(\mathscr{O}_{\mathfrak{X}}/\mathscr{A})|\mathfrak{Y})$

\oldpage[I]{210}
est le préschéma formel limite inductive des $Y_n$, et comme les conditions de
(\hyperref[I.10.6.4-fr]{10.6.4}) sont satisfaites, cela prouve que ce préschéma
formel est localement noethérien, et noethérien si $\mathfrak{X}$ l'est (puisque
alors $Y_0$ l'est en vertu de (\hyperref[I.6.1.4-fr]{6.1.4})).

\begin{definition}[10.14.2]
\label{I.10.14.2-fr}
On appelle \emph{sous-préschéma fermé} d'un préschéma formel $\mathfrak{X}$ tout
préschéma formel
$(\mathfrak{Y},(\mathscr{O}_{\mathfrak{X}}/\mathscr{A})|\mathfrak{Y})$ où
$\mathscr{A}$ est un $\mathscr{O}_{\mathfrak{X}}$-Module cohérent~; on dit que ce
préschéma est le sous-préschéma fermé défini par $\mathscr{A}$.
\end{definition}

Il est clair que la correspondance ainsi définie entre
$\mathscr{O}_{\mathfrak{X}}$-Modules cohérents et sous-préschémas fermés de
$\mathfrak{X}$ est biunivoque.

Le morphisme d'espaces topologiquement annelés
$j=(\psi,\theta):\mathfrak{Y}\to\mathfrak{X}$, où $\psi$ est l'injection
$\mathfrak{Y}\to\mathfrak{X}$ et $\theta^\sharp$ l'homomorphisme canonique
$\mathscr{O}_{\mathfrak{X}}\to\mathscr{O}_{\mathfrak{X}}/\mathscr{A}$, est
évidemment (\hyperref[I.10.4.5-fr]{10.4.5}) un morphisme de préschémas formels,
qu'on appelle l'\emph{injection canonique} de $\mathfrak{Y}$ dans $\mathfrak{X}$.
On notera que si $\mathfrak{X}=\operatorname{Spf}(A)$, où $A$ est adique
noethérien, on a $\mathscr{A}=\mathfrak{a}^\Delta$, où $\mathfrak{a}$ est un
idéal de $A$ (\hyperref[I.10.10.5-fr]{10.10.5}), et il résulte aussitôt de ce
qui précède que l'on a alors
$\mathfrak{Y}=\operatorname{Spf}(A/\mathfrak{a})$ à un isomorphisme près, et que
$j$ correspond (\hyperref[I.10.2.2-fr]{10.2.2}) à l'homomorphisme canonique
$A\to A/\mathfrak{a}$.

On dit qu'un morphisme $f:\mathfrak{Z}\to\mathfrak{X}$ de préschémas formels
localement noethériens est une \emph{immersion fermée} si elle se factorise en
$\mathfrak{Z}\xrightarrow{g}\mathfrak{Y}\xrightarrow{j}\mathfrak{X}$, où $g$
est un isomorphisme de $\mathfrak{Z}$ sur un sous-préschéma fermé
$\mathfrak{Y}$ de $\mathfrak{X}$ et $j$ l'injection canonique. Comme $j$ est un
monomorphisme d'espaces annelés, $g$ et $\mathfrak{Y}$ sont nécessairement
\emph{uniques}.

\begin{proposition}[10.14.3]
\label{I.10.14.3-fr}
Une immersion fermée est un morphisme de type fini.
\end{proposition}

On se ramène aussitôt au cas où $\mathfrak{X}$ est un schéma formel affine
$\operatorname{Spf}(A)$ et $\mathfrak{Y}=\operatorname{Spf}(A/\mathfrak{a})$~;
la proposition résulte de (\hyperref[I.10.13.1-fr]{10.13.1, c)}).

\begin{lemma}[10.14.4]
\label{I.10.14.4-fr}
Soit $f:\mathfrak{Y}\to\mathfrak{X}$ un morphisme de préschémas formels
localement noethériens, et soit $(U_\alpha)$ un recouvrement de
$f(\mathfrak{Y})$ par des ouverts formels affines noethériens de
$\mathfrak{X}$, tels que les $f^{-1}(U_\alpha)$ soient des ouverts formels
affines noethériens de $\mathfrak{Y}$. Pour que $f$ soit une immersion fermée,
il faut et il suffit que $f(\mathfrak{Y})$ soit une partie fermée de
$\mathfrak{X}$ et que, pour tout $\alpha$, la restriction de $f$ à
$f^{-1}(U_\alpha)$ corresponde (\hyperref[I.10.4.6-fr]{10.4.6}) à un
homomorphisme surjectif
$\Gamma(U_\alpha,\mathscr{O}_{\mathfrak{X}})\to
\Gamma(f^{-1}(U_\alpha),\mathscr{O}_{\mathfrak{Y}})$.
\end{lemma}

Les conditions sont évidemment nécessaires. Inversement, si elles sont
remplies, et si on désigne par $\mathfrak{a}_\alpha$ le noyau de
$\Gamma(U_\alpha,\mathscr{O}_{\mathfrak{X}})\to
\Gamma(f^{-1}(U_\alpha),\mathscr{O}_{\mathfrak{Y}})$, on définit un faisceau
cohérent d'idéaux $\mathscr{A}$ de $\mathscr{O}_{\mathfrak{X}}$ en prenant
$\mathscr{A}|U_\alpha=\mathfrak{a}_\alpha^\Delta$, et en prenant $\mathscr{A}$
nul dans le complémentaire de la réunion des $U_\alpha$. En effet, puisque
$f(\mathfrak{Y})$ est fermé et que le support de
$\mathfrak{a}_\alpha^\Delta$ est $U_\alpha\cap f(\mathfrak{Y})$, tout revient
à vérifier que $\mathfrak{a}_\alpha^\Delta$ et
$\mathfrak{a}_\beta^\Delta$ induisent le même faisceau sur un ouvert formel
affine noethérien $V\subset U_\alpha\cap U_\beta$. Or, la restriction de $f$
à $f^{-1}(U_\alpha)$ étant une immersion fermée de ce préschéma formel dans
$U_\alpha$, $f^{-1}(V)$ est un ouvert formel affine noethérien dans
$f^{-1}(U_\alpha)$ et la restriction de $f$ à $f^{-1}(V)$ est une immersion
fermée~; si $\mathfrak{b}$ est le noyau de l'homomorphisme surjectif
$\Gamma(V,\mathscr{O}_{\mathfrak{X}})\to
\Gamma(f^{-1}(V),\mathscr{O}_{\mathfrak{Y}})$ correspondant à cette
restriction, il est immédiat (\hyperref[I.10.10.2-fr]{10.10.2}) que
$\mathfrak{a}_\alpha^\Delta$ induit $\mathfrak{b}^\Delta$ sur $V$. Le faisceau
d'idéaux $\mathscr{A}$ étant ainsi défini, il est alors clair que
$f=g\circ j$, où $j:\mathfrak{Z}\to\mathfrak{X}$ est l'injection canonique du
sous-préschéma fermé $\mathfrak{Z}$ de $\mathfrak{X}$ défini par
$\mathscr{A}$, et $g$ un isomorphisme de $\mathfrak{Y}$ sur $\mathfrak{Z}$.

\begin{proposition}[10.14.5]
\label{I.10.14.5-fr}
\begin{enumerate}
  \item[\rm{(i)}] Si $f:\mathfrak{Z}\to\mathfrak{Y}$,
  $g:\mathfrak{Y}\to\mathfrak{X}$ sont des immersions fermées de préschémas
  formels localement noethériens, $g\circ f$ est une immersion fermée.

\oldpage[I]{211}
  \item[\rm{(ii)}] Soient $\mathfrak{X}$, $\mathfrak{Y}$, $\mathfrak{S}$ trois
  préschémas formels localement noethériens,
  $f:\mathfrak{X}\to\mathfrak{S}$ une immersion fermée,
  $g:\mathfrak{Y}\to\mathfrak{S}$ un morphisme. Alors le morphisme
  $\mathfrak{X}\times_{\mathfrak{S}}\mathfrak{Y}\to\mathfrak{Y}$ est une
  immersion fermée.
  \item[\rm{(iii)}] Soient $\mathfrak{S}$ un préschéma formel localement
  noethérien, $\mathfrak{X}'$, $\mathfrak{Y}'$ deux
  $\mathfrak{S}$-préschémas formels localement noethériens tels que
  $\mathfrak{X}'\times_{\mathfrak{S}}\mathfrak{Y}'$ soit localement
  noethérien. Si $\mathfrak{X}$, $\mathfrak{Y}$ sont des
  $\mathfrak{S}$-préschémas formels localement noethériens,
  $f:\mathfrak{X}\to\mathfrak{X}'$, $g:\mathfrak{Y}\to\mathfrak{Y}'$ deux
  $\mathfrak{S}$-morphismes qui sont des immersions fermées, alors
  $f\times_{\mathfrak{S}}g$ est une immersion fermée.
\end{enumerate}
\end{proposition}

En vertu de (\hyperref[I.3.5.1-fr]{3.5.1}), il suffit encore de prouver (i)
et (ii).

Pour démontrer (i), on peut supposer que $\mathfrak{Y}$ (resp.
$\mathfrak{Z}$) est un sous-préschéma fermé de $\mathfrak{X}$ (resp.
$\mathfrak{Y}$) défini par un faisceau cohérent $\mathscr{J}$ (resp.
$\mathscr{K}$) d'idéaux de $\mathscr{O}_{\mathfrak{X}}$ (resp.
$\mathscr{O}_{\mathfrak{Y}}$)~; si $\psi$ est l'injection
$\mathfrak{Y}\to\mathfrak{X}$ des espaces sous-jacents,
$\psi_*(\mathscr{K})$ est un faisceau cohérent d'idéaux de
$\psi_*(\mathscr{O}_{\mathfrak{Y}})
=\mathscr{O}_{\mathfrak{X}}/\mathscr{J}$
(\hyperref[0.5.3.12-fr]{0, 5.3.12}), donc aussi un
$\mathscr{O}_{\mathfrak{X}}$-Module cohérent
(\hyperref[0.5.3.10-fr]{0, 5.3.10})~; le noyau $\mathscr{K}_1$ de
$\mathscr{O}_{\mathfrak{X}}\to
(\mathscr{O}_{\mathfrak{X}}/\mathscr{J})/\psi_*(\mathscr{K})$ est donc un
faisceau cohérent d'idéaux de $\mathscr{O}_{\mathfrak{X}}$
(\hyperref[0.5.3.4-fr]{0, 5.3.4}), et
$\mathscr{O}_{\mathfrak{X}}/\mathscr{K}_1$ est isomorphe à
$\psi_*(\mathscr{O}_{\mathfrak{Y}}/\mathscr{K})$, ce qui prouve que
$\mathfrak{Z}$ est isomorphe à un sous-préschéma fermé de $\mathfrak{X}$.

Pour démontrer (ii), il est immédiat qu'on peut se limiter au cas où
$\mathfrak{S}=\operatorname{Spf}(A)$,
$\mathfrak{X}=\operatorname{Spf}(B)$,
$\mathfrak{Y}=\operatorname{Spf}(C)$, $A$ étant un anneau
$\mathfrak{J}$-adique noethérien, $B=A/\mathfrak{a}$, où $\mathfrak{a}$ est
un idéal de $A$, $C$ une $A$-algèbre topologique adique et noethérienne. Tout
revient à prouver que l'homomorphisme
$C\to C\widehat{\otimes}_A(A/\mathfrak{a})$ est \emph{surjectif}~: or,
$A/\mathfrak{a}$ est un $A$-module de type fini, et sa topologie est la
topologie $\mathfrak{J}$-adique~; il résulte alors de
(\hyperref[0.7.7.8-fr]{0, 7.7.8}) que
$C\widehat{\otimes}_A(A/\mathfrak{a})$ s'identifie à
$C\otimes_A(A/\mathfrak{a})=C/\mathfrak{a}C$, d'où notre assertion.

\begin{corollary}[10.14.6]
\label{I.10.14.6-fr}
Sous les hypothèses de (\hyperref[I.10.14.5-fr]{10.14.5, (ii)}), soient
$p:\mathfrak{X}\times_{\mathfrak{S}}\mathfrak{Y}\to\mathfrak{X}$,
$q:\mathfrak{X}\times_{\mathfrak{S}}\mathfrak{Y}\to\mathfrak{Y}$ les
projections, de sorte que le diagramme
\[
  \xymatrix{
    \mathfrak{X} \ar[d]_{f}
    & \mathfrak{X}\times_{\mathfrak{S}}\mathfrak{Y}
      \ar[l]_{p} \ar[d]^{q}\\
    \mathfrak{S}
    & \mathfrak{Y} \ar[l]^{g}
  }
\]
est commutatif. Pour tout $\mathscr{O}_{\mathfrak{X}}$-Module cohérent
$\mathscr{F}$, on a alors un isomorphisme canonique de
$\mathscr{O}_{\mathfrak{Y}}$-Modules
\[
  \label{I.10.14.6.1-fr}
  u:g^*(f_*(\mathscr{F}))
  \xrightarrow{\sim}
  q_*(p^*(\mathscr{F})).
  \tag{10.14.6.1}
\]
\end{corollary}

Pour définir un homomorphisme
$g^*(f_*(\mathscr{F}))\to q_*(p^*(\mathscr{F}))$, on sait qu'il revient au
même de définir un homomorphisme
$f_*(\mathscr{F})\to g_*(q_*(p^*(\mathscr{F})))
=f_*(p_*(p^*(\mathscr{F})))$
(\hyperref[0.4.4.3-fr]{0, 4.4.3})~: nous prendrons $u=f_*(\rho)$, où $\rho$
est l'homomorphisme canonique
$\mathscr{F}\to p_*(p^*(\mathscr{F}))$
(\hyperref[0.4.4.3-fr]{0, 4.4.3}). Pour voir que $u$ est un isomorphisme, on
se ramène aussitôt au cas où $\mathfrak{S}$, $\mathfrak{X}$,
$\mathfrak{Y}$ sont des spectres formels d'anneaux adiques noethériens $A$,
$B$, $C$, avec les conditions vues ci-dessus dans
(\hyperref[I.10.14.5-fr]{10.14.5, (ii)})~; on a alors
$\mathscr{F}=M^\Delta$, où $M$ est un $(A/\mathfrak{a})$-module de type fini
(\hyperref[I.10.10.5-fr]{10.10.5}), et les deux membres de
(\hyperref[I.10.14.6.1-fr]{10.14.6.1}) s'identifient respectivement, en vertu
de (\hyperref[I.10.10.8-fr]{10.10.8}), à
$(C\otimes_A M)^\Delta$ et à
$((C/\mathfrak{a}C)\otimes_{A/\mathfrak{a}}M)^\Delta$, d'où le corollaire,
puisque
$(C/\mathfrak{a}C)\otimes_{A/\mathfrak{a}}M
=(C\otimes_A(A/\mathfrak{a}))\otimes_{A/\mathfrak{a}}M$ s'identifie
canoniquement à $C\otimes_A M$.

\oldpage[I]{212}
\begin{corollary}[10.14.7]
\label{I.10.14.7-fr}
Soient $X$ un préschéma usuel localement noethérien, $Y$ un sous-préschéma
fermé de $X$, $j$ l'injection canonique $Y\to X$, $X'$ une partie fermée de
$X$ et $Y'=Y\cap X'$~; alors
$\widehat{j}:Y_{/Y'}\to X_{/X'}$ est une immersion fermée, et pour tout
$\mathscr{O}_Y$-Module cohérent $\mathscr{F}$, on a
\[
  \widehat{j}_*(\mathscr{F}_{/Y'})=(j_*(\mathscr{F}))_{/X'}.
\]
\end{corollary}

Comme $Y'=j^{-1}(X')$, il suffit d'utiliser
(\hyperref[I.10.9.9-fr]{10.9.9}) et d'appliquer
(\hyperref[I.10.14.5-fr]{10.14.5}) et
(\hyperref[I.10.14.6-fr]{10.14.6}).

\subsection{Préschémas formels séparés.}
\label{subsection:I.10.15-fr}

\begin{definition}[10.15.1]
\label{I.10.15.1-fr}
Soient $\mathfrak{S}$ un préschéma formel, $\mathfrak{X}$ un
$\mathfrak{S}$-préschéma formel,
$f:\mathfrak{X}\to\mathfrak{S}$ le morphisme structural. On appelle
morphisme diagonal
$\Delta_{\mathfrak{X}|\mathfrak{S}}:\mathfrak{X}\to
\mathfrak{X}\times_{\mathfrak{S}}\mathfrak{X}$ (noté aussi
$\Delta_{\mathfrak{X}}$) le morphisme
$(1_{\mathfrak{X}},1_{\mathfrak{X}})_{\mathfrak{S}}$. On dit que
$\mathfrak{X}$ est séparé sur $\mathfrak{S}$, ou un
$\mathfrak{S}$-schéma formel, ou que $f$ est un morphisme séparé, si l'image
par $\Delta_{\mathfrak{X}}$ de l'espace sous-jacent $\mathfrak{X}$ est une
partie fermée de l'espace sous-jacent à
$\mathfrak{X}\times_{\mathfrak{S}}\mathfrak{X}$. On dit qu'un préschéma
formel $\mathfrak{X}$ est séparé, ou est un schéma formel, s'il est séparé
sur $\mathbf{Z}$.
\end{definition}

\begin{proposition}[10.15.2]
\label{I.10.15.2-fr}
Supposons que les préschémas formels $\mathfrak{S}$, $\mathfrak{X}$ soient
respectivement limites inductives des suites $(S_n)$, $(X_n)$ de
préschémas usuels, et que le morphisme
$f:\mathfrak{X}\to\mathfrak{S}$ soit limite inductive d'une suite de
morphismes $f_n:X_n\to S_n$. Pour que $f$ soit séparé, il faut et il suffit
que le morphisme $f_0:X_0\to S_0$ le soit.
\end{proposition}

En effet, $\Delta_{\mathfrak{X}|\mathfrak{S}}$ est alors limite inductive de
la suite de morphismes $\Delta_{X_n|S_n}$
(\hyperref[I.10.7.4-fr]{10.7.4}), et l'image par
$\Delta_{\mathfrak{X}|\mathfrak{S}}$ de l'espace sous-jacent
$\mathfrak{X}$ (resp. l'espace sous-jacent
$\mathfrak{X}\times_{\mathfrak{S}}\mathfrak{X}$) est identique à l'image
par $\Delta_{X_0|S_0}$ de l'espace sous-jacent $X_0$ (resp. à l'espace
sous-jacent $X_0\times_{S_0}X_0$)~; d'où la conclusion.

\begin{proposition}[10.15.3]
\label{I.10.15.3-fr}
On suppose dans ce qui suit que tous les préschémas formels (resp.
morphismes de préschémas formels) considérés soient limites inductives de
suites de préschémas usuels (resp. de morphismes de préschémas usuels).
\begin{enumerate}
  \item[\rm{(i)}] Le composé de deux morphismes séparés est séparé.
  \item[\rm{(ii)}] Si $f:\mathfrak{X}\to\mathfrak{X}'$,
  $g:\mathfrak{Y}\to\mathfrak{Y}'$ sont deux $\mathfrak{S}$-morphismes
  séparés, $f\times_{\mathfrak{S}}g$ est séparé.
  \item[\rm{(iii)}] Si $f:\mathfrak{X}\to\mathfrak{Y}$ est un
  $\mathfrak{S}$-morphisme séparé, le $\mathfrak{S}'$-morphisme
  $f_{(\mathfrak{S}')}$ est séparé pour toute extension
  $\mathfrak{S}'\to\mathfrak{S}$ du préschéma formel de base.
  \item[\rm{(iv)}] Si le composé $g\circ f$ de deux morphismes est séparé,
  $f$ est séparé.
\end{enumerate}
\emph{(On sous-entend dans cet énoncé que si un même préschéma formel
$\mathfrak{Z}$ intervient plusieurs fois dans une même proposition, on le
considère comme limite inductive de la \emph{même} suite $(Z_n)$ de
préschémas usuels partout où il figure, et les morphismes de $\mathfrak{Z}$
dans un préschéma formel (resp. d'un préschéma formel dans $\mathfrak{Z}$)
comme limites inductives de morphismes des $Z_n$ dans des préschémas usuels
(resp. de préschémas usuels dans les $Z_n$)).}
\end{proposition}

Avec les notations de (\hyperref[I.10.15.2-fr]{10.15.2}), on a en effet
$(g\circ f)_0=g_0\circ f_0$, et
$(f\times_{\mathfrak{S}}g)_0=f_0\times_{S_0}g_0$~; les assertions de
(\hyperref[I.10.15.3-fr]{10.15.3}) sont alors conséquences immédiates de
(\hyperref[I.10.15.2-fr]{10.15.2}) et des assertions correspondantes de
(\hyperref[I.5.5.1-fr]{5.5.1}) pour les préschémas usuels.

Nous laissons au lecteur le soin d'énoncer pour la même espèce de
préschémas formels et de morphismes que dans
(\hyperref[I.10.15.3-fr]{10.15.3}) les propositions correspondant à
(\hyperref[I.5.5.5-fr]{5.5.5}),

\oldpage[I]{213}
(\hyperref[I.5.5.9-fr]{5.5.9}) et
(\hyperref[I.5.5.10-fr]{5.5.10}) (en y remplaçant «~ouvert affine~» par
«~ouvert formel affine vérifiant la condition b) de
(\hyperref[I.10.6.3-fr]{10.6.3})~»).

Un raisonnement analogue montre aussi que tout schéma formel affine
\emph{noethérien} est séparé, ce qui justifie la terminologie.

\begin{proposition}[10.15.4]
\label{I.10.15.4-fr}
Soient $\mathfrak{S}$ un préschéma formel localement noethérien,
$\mathfrak{X}$, $\mathfrak{Y}$ deux $\mathfrak{S}$-préschémas formels
localement noethériens, tels que $\mathfrak{X}$ ou $\mathfrak{Y}$ soit de
type fini sur $\mathfrak{S}$ (de sorte que
$\mathfrak{X}\times_{\mathfrak{S}}\mathfrak{Y}$ est localement noethérien)
et que $\mathfrak{Y}$ soit séparé sur $\mathfrak{S}$. Soit
$f:\mathfrak{X}\to\mathfrak{Y}$ un $\mathfrak{S}$-morphisme~; alors le
morphisme graphe
$\Gamma_f=(1_{\mathfrak{X}},f)_{\mathfrak{S}}:\mathfrak{X}\to
\mathfrak{X}\times_{\mathfrak{S}}\mathfrak{Y}$ est une immersion fermée.
\end{proposition}

On peut supposer que $\mathfrak{S}$ est limite inductive d'une suite
$(S_n)$ de préschémas localement noethériens, $\mathfrak{X}$ (resp.
$\mathfrak{Y}$) limite inductive d'une suite $(X_n)$ (resp. $(Y_n)$) de
$S_n$-préschémas, $f$ limite inductive d'une suite de $S_n$-morphismes
$f_n:X_n\to Y_n$~; alors
$\mathfrak{X}\times_{\mathfrak{S}}\mathfrak{Y}$ est limite inductive de la
suite $(X_n\times_{S_n}Y_n)$ et $\Gamma_f$ de la suite
$(\Gamma_{f_n})$ (\hyperref[I.10.7.4-fr]{10.7.4})~; par hypothèse, $Y_0$ est
séparé sur $S_0$ (\hyperref[I.10.15.2-fr]{10.15.2}), donc l'espace
$\Gamma_{f_0}(X_0)$ est un sous-espace fermé de
$X_0\times_{S_0}Y_0$~; comme les espaces sous-jacents de
$\mathfrak{X}\times_{\mathfrak{S}}\mathfrak{Y}$ (resp.
$\Gamma_f(\mathfrak{X})$) et $X_0\times_{S_0}Y_0$ (resp.
$\Gamma_{f_0}(X_0)$) sont les mêmes, on voit déjà que
$\Gamma_f(\mathfrak{X})$ est un sous-espace \emph{fermé} de
$\mathfrak{X}\times_{\mathfrak{S}}\mathfrak{Y}$. Remarquons maintenant que
lorsque $(U,V)$ parcourt l'ensemble des couples formés d'un ouvert formel
affine noethérien $U$ (resp. $V$) de $\mathfrak{X}$ (resp.
$\mathfrak{Y}$) tels que $f(U)\subset V$, les ouverts
$U\times_{\mathfrak{S}}V$ forment un recouvrement de
$\Gamma_f(\mathfrak{X})$ dans
$\mathfrak{X}\times_{\mathfrak{S}}\mathfrak{Y}$, et si
$f_U:U\to V$ est la restriction de $f$ à $U$,
$\Gamma_{f_U}:U\to U\times_{\mathfrak{S}}V$ est la restriction de
$\Gamma_f$ à $U$. Si nous montrons que $\Gamma_{f_U}$ est une immersion
fermée, il en sera donc de même de $\Gamma_f$
(\hyperref[I.10.14.4-fr]{10.14.4}), autrement dit, on est ramené au cas où
$\mathfrak{S}=\operatorname{Spf}(A)$,
$\mathfrak{X}=\operatorname{Spf}(B)$,
$\mathfrak{Y}=\operatorname{Spf}(C)$ sont affines ($A$, $B$, $C$ adiques
noethériens), $f$ correspondant à un $A$-homomorphisme continu
$\varphi:C\to B$~; $\Gamma_f$ correspond alors à l'unique homomorphisme
continu $\omega:B\widehat{\otimes}_A C\to B$ qui, composé avec les
homomorphismes canoniques $B\to B\widehat{\otimes}_A C$ et
$C\to B\widehat{\otimes}_A C$, donne respectivement l'identité et
$\varphi$. Or, il est clair que $\omega$ est \emph{surjectif}, d'où notre
assertion.

\begin{corollary}[10.15.5]
\label{I.10.15.5-fr}
Soient $\mathfrak{S}$ un préschéma formel localement noethérien,
$\mathfrak{X}$ un $\mathfrak{S}$-préschéma de type fini~; pour que
$\mathfrak{X}$ soit séparé sur $\mathfrak{S}$, il faut et il suffit que le
morphisme diagonal
$\mathfrak{X}\to\mathfrak{X}\times_{\mathfrak{S}}\mathfrak{X}$ soit une
immersion fermée.
\end{corollary}

\begin{proposition}[10.15.6]
\label{I.10.15.6-fr}
Une immersion fermée $j:\mathfrak{Y}\to\mathfrak{X}$ de préschémas formels
localement noethériens est un morphisme séparé.
\end{proposition}

Avec les notations de (\hyperref[I.10.14.2-fr]{10.14.2}),
$j_0:Y_0\to X_0$ est une immersion fermée, donc un morphisme séparé, et il
suffit d'appliquer (\hyperref[I.10.15.2-fr]{10.15.2}).

\begin{proposition}[10.15.7]
\label{I.10.15.7-fr}
Soient $X$ un préschéma (usuel) localement noethérien, $X'$ une partie
fermée de $X$ et $\widehat{X}=X_{/X'}$. Pour que $\widehat{X}$ soit séparé,
il faut et il suffit que $\widehat{X}_{\mathrm{red}}$ le soit, et il suffit
que $X$ le soit.
\end{proposition}

En effet, avec les notations de (\hyperref[I.10.8.5-fr]{10.8.5}), pour que
$\widehat{X}$ soit séparé, il faut et il suffit que $X'_0$ le soit
(\hyperref[I.10.15.2-fr]{10.15.2}), et comme
$\widehat{X}_{\mathrm{red}}=(X'_0)_{\mathrm{red}}$, il est équivalent de
dire que $\widehat{X}_{\mathrm{red}}$ l'est
(\hyperref[I.5.5.1-fr]{5.5.1, (vi)}).

\input{ega1/bib-fr.tex}
\clearpage
\thispagestyle{plain}
\markboth{CHAPITRE II}{CHAPITRE II}
\oldpage[II]{5}
\phantomsection
\label{chapter:II-fr}

\begin{center}
{\large CHAPITRE II\par}
\vspace{1.5em}
{\Large\bfseries ÉTUDE GLOBALE ÉLÉMENTAIRE\par}
{\Large\bfseries DE QUELQUES CLASSES DE MORPHISMES\par}
\vspace{1.8em}
{\bfseries Sommaire\par}
\end{center}

\medskip
\begin{tabular}{@{}r@{\ }p{0.82\textwidth}@{}}
\S~1. & Morphismes affines.\\
\S~2. & Spectres premiers homogènes.\\
\S~3. & Spectre homogène d'un faisceau d'algèbres graduées.\\
\S~4. & Fibrés projectifs. Faisceaux amples.\\
\S~5. & Morphismes quasi-affines~; morphismes quasi-projectifs~; morphismes
         propres~; morphismes projectifs.\\
\S~6. & Morphismes entiers et morphismes finis.\\
\S~7. & Critères valuatifs.\\
\S~8. & Schémas éclatés~; cônes projetants~; fermeture projective.
\end{tabular}

\medskip
Les diverses classes de morphismes étudiées dans ce chapitre le sont sans
faire grand usage des méthodes cohomologiques~; une étude plus poussée,
utilisant ces dernières méthodes, en sera faite au chapitre~III, où on
utilisera surtout les \S\S~2, 4 et~5 du chapitre~II. Le \S~8 peut être omis en
première lecture~: il donne quelques compléments au formalisme développé dans
les \S\S~1 à~3, se réduisant à des applications faciles de ce formalisme, et
nous en ferons un usage moins constant que des autres résultats de ce
chapitre.

\setcounter{section}{0}
\section{Morphismes affines}
\label{section:II.1-fr}

La plupart des résultats de ce paragraphe sont les contreparties
<<~globales~>> de ceux du chapitre~I\textsuperscript{er}, \S~1~; ils ne sont
donc pas essentiellement nouveaux et fournissent simplement un langage
commode pour la suite.

\setcounter{subsection}{0}
\subsection{$S$-préschémas et $\mathcal{O}_S$-Algèbres}
\label{subsection:II.1.1-fr}

\begin{env}[1.1.1]
\label{II.1.1.1-fr}
Soient $S$ un préschéma, $X$ un $S$-préschéma, $f:X\to S$ son morphisme
structural. On sait (\hyperref[0.4.2.4-fr]{0, 4.2.4}) que l'image directe
$f_*(\mathcal{O}_X)$ est une $\mathcal{O}_S$-Algèbre, que nous
\oldpage[II]{6}
noterons $\mathcal{A}(X)$ lorsqu'il n'en résultera pas de confusion~; si $U$
est un ouvert de $S$, on a
\[
  \mathcal{A}(f^{-1}(U))=\mathcal{A}(X)|U.
\]
De même, pour tout $\mathcal{O}_X$-Module $\mathcal{F}$ (resp. toute
$\mathcal{O}_X$-Algèbre $\mathcal{B}$), nous écrirons
$\mathcal{A}(\mathcal{F})$ (resp. $\mathcal{A}(\mathcal{B})$) l'image directe
$f_*(\mathcal{F})$ (resp. $f_*(\mathcal{B})$) qui est un
$\mathcal{A}(X)$-Module (resp. une $\mathcal{A}(X)$-Algèbre), et non seulement
un $\mathcal{O}_S$-Module (resp. $\mathcal{O}_S$-Algèbre).
\end{env}

\begin{env}[1.1.2]
\label{II.1.1.2-fr}
Soient $Y$ un second $S$-préschéma, $g:Y\to S$ son morphisme structural,
$h:X\to Y$ un $S$-morphisme~; on a donc le diagramme commutatif
\[
  \xymatrix{
    X\ar[rr]^h\ar[rd]_f && Y\ar[ld]^g\\
    & S
  }
  \tag{1.1.2.1}\label{II.1.1.2.1-fr}
\]

On a par définition $h=(\psi,\theta)$, où
$\theta:\mathcal{O}_Y\to h_*(\mathcal{O}_X)=\psi_*(\mathcal{O}_X)$ est un
homomorphisme de faisceaux d'anneaux~; on en déduit
(\hyperref[0.4.2.2-fr]{0, 4.2.2}) un homomorphisme de
$\mathcal{O}_S$-Algèbres
\[
  g_*(\theta):g_*(\mathcal{O}_Y)\to
  g_*(h_*(\mathcal{O}_X))=f_*(\mathcal{O}_X),
\]
autrement dit, un homomorphisme de $\mathcal{O}_S$-Algèbres
$\mathcal{A}(Y)\to\mathcal{A}(X)$, que nous noterons aussi
$\mathcal{A}(h)$. Si $h':Y\to Z$ est un second $S$-morphisme, il est immédiat
que $\mathcal{A}(h'\circ h)=\mathcal{A}(h)\circ\mathcal{A}(h')$. Nous avons
donc défini un \emph{foncteur contravariant} $\mathcal{A}(X)$ sur la catégorie
des $S$-préschémas, à valeurs dans la catégorie des
$\mathcal{O}_S$-Algèbres.

Soient maintenant $\mathcal{F}$ un $\mathcal{O}_X$-Module, $\mathcal{G}$ un
$\mathcal{O}_Y$-Module, et $u:\mathcal{G}\to\mathcal{F}$ un $h$-morphisme,
c'est-à-dire (\hyperref[0.4.4.1-fr]{0, 4.4.1}) un homomorphisme de
$\mathcal{O}_Y$-Modules $\mathcal{G}\to h_*(\mathcal{F})$. Alors
\[
  g_*(u):g_*(\mathcal{G})\to g_*(h_*(\mathcal{F}))=f_*(\mathcal{F})
\]
est un homomorphisme $\mathcal{A}(\mathcal{G})\to
\mathcal{A}(\mathcal{F})$ de $\mathcal{O}_S$-Modules, que nous noterons
$\mathcal{A}(u)$~; en outre, le couple
$(\mathcal{A}(h),\mathcal{A}(u))$ constitue un \emph{di-homomorphisme} du
$\mathcal{A}(Y)$-Module $\mathcal{A}(\mathcal{G})$ dans le
$\mathcal{A}(X)$-Module $\mathcal{A}(\mathcal{F})$.
\end{env}

\begin{env}[1.1.3]
\label{II.1.1.3-fr}
Le préschéma $S$ étant fixé, on peut considérer les couples
$(X,\mathcal{F})$, où $X$ est un $S$-préschéma et $\mathcal{F}$ un
$\mathcal{O}_X$-Module, comme formant une \emph{catégorie}, en définissant un
\emph{morphisme} $(X,\mathcal{F})\to(Y,\mathcal{G})$ comme un couple $(h,u)$,
où $h:X\to Y$ est un $S$-morphisme et $u:\mathcal{G}\to\mathcal{F}$ un
$h$-morphisme. On peut alors dire que
$(\mathcal{A}(X),\mathcal{A}(\mathcal{F}))$ est un \emph{foncteur
contravariant} à valeurs dans la catégorie dont les objets sont les couples
formés d'une $\mathcal{O}_S$-Algèbre et d'un Module sur cette Algèbre, et les
morphismes sont les di-homomorphismes.
\end{env}

\subsection{Préschémas affines sur un préschéma}
\label{subsection:II.1.2-fr}

\begin{definition}[1.2.1]
\label{II.1.2.1-fr}
Soient $X$ un $S$-préschéma, $f:X\to S$ son morphisme structural. On dit que
$X$ est \emph{affine au-dessus de $S$} s'il existe un recouvrement
$(S_\alpha)$ de $S$ par des ouverts affines tels que, pour tout $\alpha$, le
préschéma induit par $X$ sur l'ensemble ouvert $f^{-1}(S_\alpha)$ soit affine.
\end{definition}

\begin{example}[1.2.2]
\label{II.1.2.2-fr}
Tout sous-préschéma fermé de $S$ est un $S$-préschéma affine au-dessus de $S$
(\hyperref[I.4.2.3-fr]{I, 4.2.3} et
\hyperref[I.4.2.4-fr]{4.2.4}).
\end{example}

\begin{remark}[1.2.3]
\label{II.1.2.3-fr}
Un préschéma $X$ affine au-dessus de $S$ n'est pas nécessairement un schéma
affine, comme le montre l'exemple $X=S$ (\hyperref[II.1.2.2-fr]{1.2.2}).
D'autre part, si un schéma affine $X$ est un $S$-préschéma, $X$ n'est pas
nécessairement affine au-dessus
\oldpage[II]{7}
de $S$ (voir (\hyperref[II.1.3.3-fr]{1.3.3})). Toutefois, rappelons que si
$S$ est un \emph{schéma}, tout $S$-préschéma qui est un schéma affine est
affine au-dessus de $S$ (\hyperref[I.5.5.10-fr]{I, 5.5.10}).
\end{remark}

\begin{proposition}[1.2.4]
\label{II.1.2.4-fr}
Tout $S$-préschéma qui est affine au-dessus de $S$ est séparé au-dessus de
$S$ (autrement dit, est un $S$-schéma).
\end{proposition}

Cela résulte aussitôt de (\hyperref[I.5.5.5-fr]{I, 5.5.5}) et
(\hyperref[I.5.5.8-fr]{I, 5.5.8}).

\begin{proposition}[1.2.5]
\label{II.1.2.5-fr}
Soient $X$ un $S$-schéma affine au-dessus de $S$, $f:X\to S$ le morphisme
structural. Pour tout ouvert $U\subset S$, $f^{-1}(U)$ est affine au-dessus
de $U$.
\end{proposition}

En vertu de la déf. (\hyperref[II.1.2.1-fr]{1.2.1}), on est ramené au cas où
$S=\operatorname{Spec}(A)$ et $X=\operatorname{Spec}(B)$ sont affines et
alors $f=({}^a\varphi,\widetilde{\varphi})$, où $\varphi:A\to B$ est un
homomorphisme. Comme les $D(g)$, où $g\in A$, forment une base de $S$, on est
ramené au cas où $U=D(g)$~; mais on sait alors que
$f^{-1}(U)=D(\varphi(g))$ (\hyperref[I.1.2.2.2-fr]{I, 1.2.2.2}), d'où la
proposition.

\begin{proposition}[1.2.6]
\label{II.1.2.6-fr}
Soient $X$ un $S$-schéma affine au-dessus de $S$, $f:X\to S$ le morphisme
structural. Pour tout $\mathcal{O}_X$-Module quasi-cohérent $\mathcal{F}$,
$f_*(\mathcal{F})$ est un $\mathcal{O}_S$-Module quasi-cohérent.
\end{proposition}

Compte tenu de (\hyperref[II.1.2.4-fr]{1.2.4}), cela résulte de
(\hyperref[I.9.2.2-fr]{I, 9.2.2, a)}).

En particulier, la $\mathcal{O}_S$-Algèbre
$\mathcal{A}(X)=f_*(\mathcal{O}_X)$ est \emph{quasi-cohérente}.

\begin{proposition}[1.2.7]
\label{II.1.2.7-fr}
Soient $X$ un $S$-schéma affine au-dessus de $S$. Pour tout $S$-préschéma
$Y$, l'application $h\mapsto\mathcal{A}(h)$ de l'ensemble
$\operatorname{Hom}_S(Y,X)$ dans l'ensemble
$\operatorname{Hom}(\mathcal{A}(X),\mathcal{A}(Y))$
(\hyperref[II.1.1.2-fr]{1.1.2}) est bijective.
\end{proposition}

Soient $f:X\to S$, $g:Y\to S$ les morphismes structuraux. Supposons d'abord
$S=\operatorname{Spec}(A)$ et $X=\operatorname{Spec}(B)$ affines~; il faut
prouver que pour tout homomorphisme
$\omega:f_*(\mathcal{O}_X)\to g_*(\mathcal{O}_Y)$ de
$\mathcal{O}_S$-Algèbres, il existe un $S$-morphisme $h:Y\to X$ et un seul
tel que $\mathcal{A}(h)=\omega$. Par définition, pour tout ouvert $U\subset
S$, $\omega$ définit un homomorphisme
\[
  \omega_U=\Gamma(U,\omega):
  \Gamma(f^{-1}(U),\mathcal{O}_X)\to
  \Gamma(g^{-1}(U),\mathcal{O}_Y)
\]
de $\Gamma(U,\mathcal{O}_S)$-algèbres. En particulier, pour $U=S$, cela
donne un homomorphisme
$\varphi:\Gamma(X,\mathcal{O}_X)\to\Gamma(Y,\mathcal{O}_Y)$ de
$\Gamma(S,\mathcal{O}_S)$-algèbres, auquel correspond un $S$-morphisme
$h:Y\to X$ bien déterminé, puisque $X$ est affine
(\hyperref[I.2.2.4-fr]{I, 2.2.4}). Reste à prouver que
$\mathcal{A}(h)=\omega$, autrement dit, que pour tout ouvert $U$ d'une base
de $S$, $\omega_U$ coïncide avec l'homomorphisme d'algèbres $\varphi_U$
correspondant au $S$-morphisme $g^{-1}(U)\to f^{-1}(U)$ restriction de $h$.
On peut se borner au cas où $U=D(\lambda)$, avec $\lambda\in S$~; alors, si
$f=({}^a\rho,\widetilde{\rho})$, où $\rho:A\to B$ est un homomorphisme
d'anneaux, on a $f^{-1}(U)=D(\mu)$, où $\mu=\rho(\lambda)$, et
$\Gamma(f^{-1}(U),\mathcal{O}_X)$ est l'anneau de fractions $B_\mu$~; or,
le diagramme
\[
  \xymatrix{
    B\ar[r]^\varphi\ar[d] & \Gamma(Y,\mathcal{O}_Y)\ar[d]\\
    B_\mu\ar[r]_{\varphi_U} & \Gamma(g^{-1}(U),\mathcal{O}_Y)
  }
\]
est commutatif, et il en est de même du diagramme analogue où $\varphi_U$
est remplacé par $\omega_U$~; l'égalité $\varphi_U=\omega_U$ résulte alors
de la propriété universelle des anneaux de fractions
(\hyperref[0.1.2.4-fr]{0, 1.2.4}).

Passons au cas général, et soit $(S_\alpha)$ un recouvrement de $S$ par des
ouverts affines
\oldpage[II]{8}
tels que les $f^{-1}(S_\alpha)$ soient affines. Alors tout homomorphisme
$\omega:\mathcal{A}(X)\to\mathcal{A}(Y)$ de $\mathcal{O}_S$-Algèbres donne
par restriction une famille d'homomorphismes
\[
  \omega_\alpha:
  \mathcal{A}(f^{-1}(S_\alpha))\to
  \mathcal{A}(g^{-1}(S_\alpha))
\]
de $\mathcal{O}_{S_\alpha}$-Algèbres, donc une famille de
$S_\alpha$-morphismes
$h_\alpha:g^{-1}(S_\alpha)\to f^{-1}(S_\alpha)$ d'après ce qui précède.
Tout revient à voir que pour tout ouvert affine $U$ d'une base de
$S_\alpha\cap S_\beta$, les restrictions de $h_\alpha$ et de $h_\beta$ à
$g^{-1}(U)$ coïncident, ce qui est évident, puisque, en vertu de ce qui
précède, ces restrictions correspondent toutes deux à l'homomorphisme
$\mathcal{A}(X)|U\to\mathcal{A}(Y)|U$ restriction de $\omega$.

\begin{corollary}[1.2.8]
\label{II.1.2.8-fr}
Soient $X$, $Y$ deux $S$-schémas qui sont affines au-dessus de $S$. Pour
qu'un $S$-morphisme $h:Y\to X$ soit un isomorphisme, il faut et il suffit que
$\mathcal{A}(h):\mathcal{A}(X)\to\mathcal{A}(Y)$ soit un isomorphisme.
\end{corollary}

Cela résulte aussitôt de (\hyperref[II.1.2.7-fr]{1.2.7}) et du caractère
fonctoriel de $\mathcal{A}(X)$.

\subsection{Préschéma affine au-dessus de $S$ associé à une
$\mathcal{O}_S$-Algèbre}
\label{subsection:II.1.3-fr}

\begin{proposition}[1.3.1]
\label{II.1.3.1-fr}
Soit $S$ un préschéma. Pour toute $\mathcal{O}_S$-Algèbre quasi-cohérente
$\mathcal{B}$, il existe un préschéma $X$ affine au-dessus de $S$, défini à
un $S$-isomorphisme unique près, tel que $\mathcal{A}(X)=\mathcal{B}$.
\end{proposition}

L'unicité résulte de (\hyperref[II.1.2.8-fr]{1.2.8})~; démontrons l'existence
de $X$. Pour tout ouvert affine $U\subset S$, soit $X_U$ le préschéma
$\operatorname{Spec}(\Gamma(U,\mathcal{B}))$~; comme
$\Gamma(U,\mathcal{B})$ est une $\Gamma(U,\mathcal{O}_S)$-algèbre, $X_U$
est un $S$-préschéma (\hyperref[I.1.6.1-fr]{I, 1.6.1}). En outre, comme
$\mathcal{B}$ est quasi-cohérente, la $\mathcal{O}_S$-Algèbre
$\mathcal{A}(X_U)$ s'identifie canoniquement à $\mathcal{B}|U$
(\hyperref[I.1.3.7-fr]{I, 1.3.7},
\hyperref[I.1.3.13-fr]{1.3.13} et
\hyperref[I.1.6.3-fr]{1.6.3}). Soit $V$ un second ouvert affine de $S$, et
soit $X_{U,V}$ le préschéma induit par $X_U$ sur $f_U^{-1}(U\cap V)$, en
désignant par $f_U$ le morphisme structural $X_U\to S$~; $X_{U,V}$ et
$X_{V,U}$ sont affines sur $U\cap V$
(\hyperref[II.1.2.5-fr]{1.2.5}), et par définition $\mathcal{A}(X_{U,V})$
et $\mathcal{A}(X_{V,U})$ s'identifient canoniquement à
$\mathcal{B}|(U\cap V)$. Il y a donc
(\hyperref[II.1.2.8-fr]{1.2.8}) un $S$-isomorphisme canonique
$\theta_{U,V}:X_{V,U}\to X_{U,V}$~; en outre, si $W$ est un troisième
ouvert affine de $S$, et si $\theta'_{U,V}$, $\theta'_{V,W}$,
$\theta'_{U,W}$ sont les restrictions de $\theta_{U,V}$, $\theta_{V,W}$,
$\theta_{U,W}$ aux images réciproques de $U\cap V\cap W$ dans $X_V$,
$X_W$ et $X_W$ respectivement par les morphismes structuraux, on a
$\theta'_{U,V}\circ\theta'_{V,W}=\theta'_{U,W}$. Il existe par suite un
préschéma $X$, un recouvrement $(T_U)$ de $X$ par des ouverts affines, et
pour chaque $U$ un isomorphisme $\varphi_U:X_U\to T_U$, de sorte que
$\varphi_U$ applique $f_U^{-1}(U\cap V)$ sur $T_U\cap T_V$ et que l'on ait
$\theta_{U,V}=\varphi_U^{-1}\circ\varphi_V$
(\hyperref[I.2.3.1-fr]{I, 2.3.1}). Le morphisme
$g_U=f_U\circ\varphi_U^{-1}$ fait de $T_U$ un $S$-préschéma, et les
morphismes $g_U$ et $g_V$ coïncident dans $T_U\cap T_V$, donc $X$ est un
$S$-préschéma. En outre, il est clair par définition que $X$ est affine
au-dessus de $S$ et que $\mathcal{A}(T_U)=\mathcal{B}|U$, donc
$\mathcal{A}(X)=\mathcal{B}$.

On dit que le $S$-schéma $X$ ainsi défini est \emph{associé à la
$\mathcal{O}_S$-Algèbre $\mathcal{B}$} ou est le \emph{spectre de
$\mathcal{B}$}, et on le note $\operatorname{Spec}(\mathcal{B})$.

\begin{corollary}[1.3.2]
\label{II.1.3.2-fr}
Soient $X$ un préschéma affine au-dessus de $S$, $f:X\to S$ le morphisme
structural. Pour tout ouvert affine $U\subset S$, le préschéma induit sur
$f^{-1}(U)$ est un schéma affine d'anneau
$\Gamma(U,\mathcal{A}(X))$.
\end{corollary}

\oldpage[II]{9}
Comme on peut supposer $X$ associé à une $\mathcal{O}_S$-Algèbre en vertu de
(\hyperref[II.1.2.6-fr]{1.2.6}) et (\hyperref[II.1.3.1-fr]{1.3.1}), le
corollaire résulte de la construction de $X$ décrite dans
(\hyperref[II.1.3.1-fr]{1.3.1}).

\begin{example}[1.3.3]
\label{II.1.3.3-fr}
Soit $S$ le plan affine sur un corps $K$, où le point $0$ a été dédoublé
(\hyperref[I.5.5.11-fr]{I, 5.5.11})~; avec les notations de
(\hyperref[I.5.5.11-fr]{I, 5.5.11}), $S$ est réunion de deux ouverts affines
$Y_1$, $Y_2$~; si $f$ est l'immersion ouverte $Y_1\to S$,
$f^{-1}(Y_2)=Y_1\cap Y_2$ n'est pas un ouvert affine dans $Y_1$
(\emph{loc. cit.}), donc on a un exemple d'un schéma affine qui n'est pas
affine au-dessus de $S$.
\end{example}

\begin{corollary}[1.3.4]
\label{II.1.3.4-fr}
Soit $S$ un schéma affine~; pour qu'un $S$-préschéma $X$ soit affine au-dessus
de $S$, il faut et il suffit que $X$ soit un schéma affine.
\end{corollary}

\begin{corollary}[1.3.5]
\label{II.1.3.5-fr}
Soient $X$ un préschéma affine au-dessus d'un préschéma $S$, $Y$ un
$X$-préschéma. Pour que $Y$ soit affine au-dessus de $S$, il faut et il
suffit que $Y$ soit affine au-dessus de $X$.
\end{corollary}

On est aussitôt ramené au cas où $S$ est un schéma affine, et il en est alors
de même de $X$ (\hyperref[II.1.3.4-fr]{1.3.4})~; les deux conditions de
l'énoncé expriment alors que $Y$ est un schéma affine
(\hyperref[II.1.3.4-fr]{1.3.4}).

\begin{env}[1.3.6]
\label{II.1.3.6-fr}
Soit $X$ un préschéma affine au-dessus de $S$. Pour définir un préschéma $Y$
affine \emph{au-dessus de $X$}, il revient au même, en vertu de
(\hyperref[II.1.3.5-fr]{1.3.5}), de se donner un préschéma $Y$ affine
\emph{au-dessus de $S$}, et un $S$-morphisme $g:Y\to X$~; en d'autres termes
(\hyperref[II.1.3.1-fr]{1.3.1} et
\hyperref[II.1.2.7-fr]{1.2.7}), cela revient à se donner une
$\mathcal{O}_S$-Algèbre quasi-cohérente $\mathcal{B}$ et un homomorphisme
$\mathcal{A}(X)\to\mathcal{B}$ de $\mathcal{O}_S$-Algèbres (que l'on peut
envisager comme définissant sur $\mathcal{B}$ une structure de
$\mathcal{A}(X)$-Algèbre). Si $f:X\to S$ est le morphisme structural, on a
alors $\mathcal{B}=f_*(g_*(\mathcal{O}_Y))$.
\end{env}

\begin{corollary}[1.3.7]
\label{II.1.3.7-fr}
Soit $X$ un préschéma affine au-dessus de $S$~; pour que $X$ soit de type fini
sur $S$, il faut et il suffit que la $\mathcal{O}_S$-Algèbre quasi-cohérente
$\mathcal{A}(X)$ soit de type fini
(\hyperref[I.9.6.2-fr]{I, 9.6.2}).
\end{corollary}

Par définition (\hyperref[I.9.6.2-fr]{I, 9.6.2}), on est ramené au cas où
$S$ est affine~; alors $X$ est un schéma affine
(\hyperref[II.1.3.4-fr]{1.3.4}), et si $S=\operatorname{Spec}(A)$,
$X=\operatorname{Spec}(B)$, $\mathcal{A}(X)$ est la
$\mathcal{O}_S$-Algèbre $\widetilde{B}$~; comme
$\Gamma(U,\widetilde{B})=B$, le corollaire résulte de
(\hyperref[I.9.6.2-fr]{I, 9.6.2}) et
(\hyperref[I.6.3.3-fr]{I, 6.3.3}).

\begin{corollary}[1.3.8]
\label{II.1.3.8-fr}
Soit $X$ un préschéma affine au-dessus de $S$~; pour que $X$ soit réduit il
faut et il suffit que la $\mathcal{O}_S$-Algèbre quasi-cohérente
$\mathcal{A}(X)$ soit réduite
(\hyperref[0.4.1.4-fr]{0, 4.1.4}).
\end{corollary}

En effet, la question est évidemment locale sur $S$~; en vertu de
(\hyperref[II.1.3.2-fr]{1.3.2}), le corollaire résulte de
(\hyperref[I.5.1.1-fr]{I, 5.1.1}) et
(\hyperref[I.5.1.4-fr]{I, 5.1.4}).

\subsection{Faisceaux quasi-cohérents sur un préschéma affine au-dessus de
$S$}
\label{subsection:II.1.4-fr}

\begin{proposition}[1.4.1]
\label{II.1.4.1-fr}
Soient $X$ un préschéma affine au-dessus de $S$, $Y$ un $S$-préschéma,
$\mathcal{F}$ (resp. $\mathcal{G}$) un $\mathcal{O}_X$-Module (resp. un
$\mathcal{O}_Y$-Module) quasi-cohérent. Alors l'application
$(h,u)\mapsto(\mathcal{A}(h),\mathcal{A}(u))$ de l'ensemble des morphismes
$(Y,\mathcal{G})\to(X,\mathcal{F})$ dans l'ensemble des di-homomorphismes
$(\mathcal{A}(X),\mathcal{A}(\mathcal{F}))\to
(\mathcal{A}(Y),\mathcal{A}(\mathcal{G}))$
(\hyperref[II.1.1.2-fr]{1.1.2} et
\hyperref[II.1.1.3-fr]{1.1.3}) est bijective.
\end{proposition}

La démonstration suit exactement la même marche que celle de
(\hyperref[II.1.2.7-fr]{1.2.7}) en utilisant
(\hyperref[I.2.2.5-fr]{I, 2.2.5}) et
(\hyperref[I.2.2.4-fr]{I, 2.2.4}), et les détails en sont laissés au
lecteur.

\begin{corollary}[1.4.2]
\label{II.1.4.2-fr}
Si, en outre des hypothèses de (\hyperref[II.1.4.1-fr]{1.4.1}), on suppose
$Y$ affine au-dessus de $S$, alors, pour que $(h,u)$ soit un isomorphisme,
il faut et il suffit que $(\mathcal{A}(h),\mathcal{A}(u))$ soit un
di-isomorphisme.
\end{corollary}

\oldpage[II]{10}
\begin{proposition}[1.4.3]
\label{II.1.4.3-fr}
Pour tout couple $(\mathcal{B},\mathcal{M})$ formé d'une
$\mathcal{O}_S$-Algèbre quasi-cohérente $\mathcal{B}$ et d'un
$\mathcal{B}$-Module $\mathcal{M}$ quasi-cohérent (en tant que
$\mathcal{O}_S$-Module ou en tant que $\mathcal{B}$-Module, ce qui revient au
même (\hyperref[I.9.6.1-fr]{I, 9.6.1})), il existe un couple
$(X,\mathcal{F})$ formé d'un préschéma $X$ affine au-dessus de $S$ et d'un
$\mathcal{O}_X$-Module quasi-cohérent $\mathcal{F}$, tels que
$\mathcal{A}(X)=\mathcal{B}$ et $\mathcal{A}(\mathcal{F})=\mathcal{M}$~; en
outre ce couple est déterminé à un isomorphisme unique près.
\end{proposition}

L'unicité résulte de (\hyperref[II.1.4.1-fr]{1.4.1}) et
(\hyperref[II.1.4.2-fr]{1.4.2})~; l'existence se démontre comme dans
(\hyperref[II.1.3.1-fr]{1.3.1}), et nous laissons encore les détails au
lecteur. On note $\widetilde{\mathcal{M}}$ le $\mathcal{O}_X$-Module
$\mathcal{F}$, et on dit qu'il est \emph{associé} au $\mathcal{B}$-Module
quasi-cohérent $\mathcal{M}$~; pour tout ouvert affine $U\subset S$,
$\widetilde{\mathcal{M}}|p^{-1}(U)$ (où $p$ est le morphisme structural
$X\to S$) est canoniquement isomorphe à
$(\Gamma(U,\mathcal{M}))^{\sim}$.

\begin{corollary}[1.4.4]
\label{II.1.4.4-fr}
Dans la catégorie des $\mathcal{B}$-Modules quasi-cohérents,
$\widetilde{\mathcal{M}}$ est un foncteur covariant additif exact en
$\mathcal{M}$, qui commute aux limites inductives et aux sommes directes.
\end{corollary}

On est en effet aussitôt ramené au cas où $S$ est affine, et le corollaire
résulte alors de (\hyperref[I.1.3.5-fr]{I, 1.3.5}),
(\hyperref[I.1.3.9-fr]{I, 1.3.9}) et
(\hyperref[I.1.3.11-fr]{I, 1.3.11}).

\begin{corollary}[1.4.5]
\label{II.1.4.5-fr}
Sous les hypothèses de (\hyperref[II.1.4.3-fr]{1.4.3}), pour que
$\widetilde{\mathcal{M}}$ soit un $\mathcal{O}_X$-Module de type fini, il faut
et il suffit que $\mathcal{M}$ soit un $\mathcal{B}$-Module de type fini.
\end{corollary}

On est aussitôt ramené au cas où $S=\operatorname{Spec}(A)$ est un schéma
affine. Alors $\mathcal{B}=\widetilde{B}$, où $B$ est une $A$-algèbre de type
fini (\hyperref[I.9.6.2-fr]{I, 9.6.2}), et
$\mathcal{M}=\widetilde{M}$, où $M$ est un $B$-module
(\hyperref[I.1.3.13-fr]{I, 1.3.13})~; \emph{sur le préschéma $X$},
$\mathcal{O}_X$ est associé à l'anneau $B$ et
$\widetilde{\mathcal{M}}$ au $B$-module $M$~; pour que
$\widetilde{\mathcal{M}}$ soit de type fini, il faut et il suffit donc que
$M$ soit de type fini (\hyperref[I.1.3.13-fr]{I, 1.3.13}), d'où notre
assertion.

\begin{proposition}[1.4.6]
\label{II.1.4.6-fr}
Soient $Y$ un préschéma affine au-dessus de $S$, $X$, $X'$ deux préschémas
affines au-dessus de $Y$ (donc aussi au-dessus de $S$
(\hyperref[II.1.3.5-fr]{1.3.5})). Soient
$\mathcal{B}=\mathcal{A}(Y)$, $\mathcal{A}=\mathcal{A}(X)$,
$\mathcal{A}'=\mathcal{A}(X')$. Alors $X\times_Y X'$ est affine au-dessus de
$Y$ (donc aussi au-dessus de $S$) et $\mathcal{A}(X\times_Y X')$
s'identifie canoniquement à $\mathcal{A}\otimes_{\mathcal{B}}\mathcal{A}'$.
\end{proposition}

En effet (\hyperref[I.9.6.1-fr]{I, 9.6.1})
$\mathcal{A}\otimes_{\mathcal{B}}\mathcal{A}'$ est une
$\mathcal{B}$-Algèbre quasi-cohérente, donc aussi une
$\mathcal{O}_S$-Algèbre quasi-cohérente
(\hyperref[I.9.6.1-fr]{I, 9.6.1})~; soit $Z$ le spectre de
$\mathcal{A}\otimes_{\mathcal{B}}\mathcal{A}'$
(\hyperref[II.1.3.1-fr]{1.3.1}). Les $\mathcal{B}$-homomorphismes canoniques
$\mathcal{A}\to\mathcal{A}\otimes_{\mathcal{B}}\mathcal{A}'$ et
$\mathcal{A}'\to\mathcal{A}\otimes_{\mathcal{B}}\mathcal{A}'$ correspondent
(\hyperref[II.1.2.7-fr]{1.2.7}) à des $Y$-morphismes $p:Z\to X$ et
$p':Z\to X'$. Pour voir que le triplet $(Z,p,p')$ est un produit
$X\times_Y X'$, on peut se borner au cas où $S$ est un schéma affine
d'anneau $C$ (\hyperref[I.3.2.6.4-fr]{I, 3.2.6.4}). Mais alors $Y$, $X$,
$X'$ sont des schémas affines (\hyperref[II.1.3.4-fr]{1.3.4}) dont les
anneaux $B$, $A$, $A'$ sont des $C$-algèbres telles que
$\mathcal{B}=\widetilde{B}$, $\mathcal{A}=\widetilde{A}$,
$\mathcal{A}'=\widetilde{A'}$. On sait alors
(\hyperref[I.1.3.13-fr]{I, 1.3.13}) que
$\mathcal{A}\otimes_{\mathcal{B}}\mathcal{A}'$ s'identifie canoniquement à
la $\mathcal{O}_S$-Algèbre $\widetilde{A\otimes_B A'}$, donc l'anneau
$A(Z)$ s'identifie à $A\otimes_B A'$ et les morphismes $p$, $p'$ correspondent
aux homomorphismes canoniques $A\to A\otimes_B A'$,
$A'\to A\otimes_B A'$. La proposition résulte alors de
(\hyperref[I.3.2.2-fr]{I, 3.2.2}).

\begin{corollary}[1.4.7]
\label{II.1.4.7-fr}
Soit $\mathcal{F}$ (resp. $\mathcal{F}'$) un $\mathcal{O}_X$-Module (resp. un
$\mathcal{O}_{X'}$-Module) quasi-cohérent~; alors
$\mathcal{A}(\mathcal{F}\otimes_Y\mathcal{F}')$ s'identifie canoniquement à
$\mathcal{A}(\mathcal{F})\otimes_{\mathcal{A}(Y)}
\mathcal{A}(\mathcal{F}')$.
\end{corollary}

On sait que $\mathcal{F}\otimes_Y\mathcal{F}'$ est quasi-cohérent sur
$X\times_Y X'$ (\hyperref[I.9.1.2-fr]{I, 9.1.2}). Soient $g:Y\to S$,
$f:X\to Y$, $f':X'\to Y$ les morphismes structuraux, de sorte que le
morphisme structural
\oldpage[II]{11}
$h:Z\to S$ est égal à $g\circ f\circ p$ et à $g\circ f'\circ p'$. On
définit un homomorphisme canonique
\[
  \mathcal{A}(\mathcal{F})\otimes_{\mathcal{A}(Y)}\mathcal{A}(\mathcal{F}')
  \longrightarrow \mathcal{A}(\mathcal{F}\otimes_Y\mathcal{F}')
\]
de la façon suivante : pour tout ouvert $U\subset S$, on a des homomorphismes
canoniques
\[
  \begin{aligned}
  \Gamma(f^{-1}(g^{-1}(U)),\mathcal{F})
  &\longrightarrow \Gamma(h^{-1}(U),p^*(\mathcal{F}))\quad\text{et}\quad
  \Gamma(f^{\prime-1}(g^{-1}(U)),\mathcal{F}')
  \longrightarrow \Gamma(h^{-1}(U),p^{\prime*}(\mathcal{F}')).
  \end{aligned}
\]
(\hyperref[0.4.4.3-fr]{0, 4.4.3}), d'où on déduit un homomorphisme canonique
\[
  \begin{aligned}
  &\Gamma(f^{-1}(g^{-1}(U)),\mathcal{F})
    \otimes_{\Gamma(g^{-1}(U),\mathcal{O}_Y)}
    \Gamma(f^{\prime-1}(g^{-1}(U)),\mathcal{F}')\longrightarrow\\
  &\hspace{17em}
    \Gamma(h^{-1}(U),p^*(\mathcal{F}))
    \otimes_{\Gamma(h^{-1}(U),\mathcal{O}_Z)}
    \Gamma(h^{-1}(U),p^{\prime*}(\mathcal{F}')).
  \end{aligned}
\]

Pour voir qu'on a bien là un isomorphisme de $\mathcal{A}(Z)$-Modules, on
peut se borner au cas où $S$ (et par suite $X$, $X'$, $Y$,
$X\times_Y X'$) sont des schémas affines, et (avec les notations de
(\hyperref[II.1.4.6-fr]{1.4.6})), $\mathcal{F}=\widetilde{M}$,
$\mathcal{F}'=\widetilde{M'}$, $M$ (resp. $M'$) étant un $A$-module (resp. un
$A'$-module). Alors $\mathcal{F}\otimes_Y\mathcal{F}'$ s'identifie au
faisceau sur $X\times_Y X'$ associé au $(A\otimes_B A')$-module
$M\otimes_B M'$ (\hyperref[I.9.1.3-fr]{I, 9.1.3}) et le corollaire résulte de
l'identification canonique des $\mathcal{O}_S$-Modules
$\widetilde{M\otimes_B M'}$ et
$\widetilde{M}\otimes_{\widetilde{B}}\widetilde{M'}$ (où $M$, $M'$ et $B$
sont considérés comme des $C$-modules)
(\hyperref[I.1.3.13-fr]{I, 1.3.13} et
\hyperref[I.1.6.3-fr]{I, 1.6.3}).

Si on applique en particulier (\hyperref[II.1.4.7-fr]{1.4.7}) au cas où
$X=Y$ et $\mathcal{F}'=\mathcal{O}_{X'}$, on voit que le
$\mathcal{A}'$-Module $\mathcal{A}(f^{\prime*}(\mathcal{F}))$ s'identifie à
$\mathcal{A}(\mathcal{F})\otimes_{\mathcal{B}}\mathcal{A}'$.

\begin{env}[1.4.8]
\label{II.1.4.8-fr}
En particulier, lorsque $X=X'=Y$ ($X$ étant affine au-dessus de $S$), on voit
que si $\mathcal{F}$, $\mathcal{G}$ sont deux $\mathcal{O}_X$-Modules
quasi-cohérents, on a
\[
  \mathcal{A}(\mathcal{F}\otimes_{\mathcal{O}_X}\mathcal{G})
  =\mathcal{A}(\mathcal{F})\otimes_{\mathcal{A}(X)}\mathcal{A}(\mathcal{G})
  \tag{1.4.8.1}\label{II.1.4.8.1-fr}
\]
à un isomorphisme canonique fonctoriel près. Si de plus $\mathcal{F}$ admet
une présentation finie, il résulte de
(\hyperref[I.1.6.3-fr]{I, 1.6.3} et
\hyperref[I.1.3.12-fr]{I, 1.3.12}) que
\[
  \mathcal{A}(\mathcal{H}om_X(\mathcal{F},\mathcal{G}))
  =\mathcal{H}om_{\mathcal{A}(X)}
    (\mathcal{A}(\mathcal{F}),\mathcal{A}(\mathcal{G}))
  \tag{1.4.8.2}\label{II.1.4.8.2-fr}
\]
à un isomorphisme canonique près.
\end{env}

\begin{remark}[1.4.9]
\label{II.1.4.9-fr}
Si $X$, $X'$ sont deux préschémas affines au-dessus de $S$, la somme
$X\sqcup X'$ est aussi affine au-dessus de $S$, la somme de deux schémas
affines étant un schéma affine.
\end{remark}

\begin{proposition}[1.4.10]
\label{II.1.4.10-fr}
Soient $S$ un préschéma, $\mathcal{B}$ une $\mathcal{O}_S$-Algèbre
quasi-cohérente, $X=\operatorname{Spec}(\mathcal{B})$. Pour tout faisceau
quasi-cohérent d'idéaux $\mathcal{J}$ de $\mathcal{B}$,
$\widetilde{\mathcal{J}}$ est un faisceau quasi-cohérent d'idéaux de
$\mathcal{O}_X$, et le sous-préschéma fermé $Y$ de $X$ défini par
$\widetilde{\mathcal{J}}$ est canoniquement isomorphe à
$\operatorname{Spec}(\mathcal{B}/\mathcal{J})$.
\end{proposition}

En effet, il résulte aussitôt de (\hyperref[I.4.2.3-fr]{I, 4.2.3}) que $Y$
est affine au-dessus de $S$~; en vertu de
(\hyperref[II.1.3.1-fr]{1.3.1}), on est donc ramené au cas où $S$ est affine,
et la proposition résulte alors aussitôt de
(\hyperref[I.4.1.2-fr]{I, 4.1.2}).

On peut encore exprimer le résultat de
(\hyperref[II.1.4.10-fr]{1.4.10}) en disant que si
$h:\mathcal{B}\to\mathcal{B}'$ est un homomorphisme \emph{surjectif} de
$\mathcal{O}_S$-Algèbres quasi-cohérentes,
${}^a h:\operatorname{Spec}(\mathcal{B}')\to
\operatorname{Spec}(\mathcal{B})$ est une \emph{immersion fermée}.

\oldpage[II]{12}
\begin{proposition}[1.4.11]
\label{II.1.4.11-fr}
Soient $S$ un préschéma, $\mathcal{B}$ une $\mathcal{O}_S$-Algèbre
quasi-cohérente, $X=\operatorname{Spec}(\mathcal{B})$. Pour tout Idéal
quasi-cohérent $\mathcal{K}$ de $\mathcal{O}_S$, on a (en désignant par $f$
le morphisme structural $X\to S$),
$f^*(\mathcal{K})\mathcal{O}_X=\widetilde{\mathcal{K}\mathcal{B}}$ à un
isomorphisme canonique près.
\end{proposition}

La question étant locale sur $S$, on est ramené au cas où
$S=\operatorname{Spec}(A)$ est affine, et dans ce cas la proposition n'est
autre que (\hyperref[I.1.6.9-fr]{I, 1.6.9}).

\subsection{Changement du préschéma de base}
\label{II.1.5-fr}

\begin{proposition}[1.5.1]
\label{II.1.5.1-fr}
Soit $X$ un préschéma affine au-dessus de $S$. Pour toute extension
$g:S'\to S$ du préschéma de base, $X'=X_{(S')}=X\times_S S'$ est affine
au-dessus de $S'$.
\end{proposition}

Si $f'$ est la projection $X'\to S'$, il suffit de prouver que
$f^{\prime-1}(U')$ est un ouvert affine pour tout ouvert affine $U'$ de $S'$
tel que $g(U')$ soit contenu dans un ouvert affine $U$ de $S$
(\hyperref[II.1.2.1-fr]{1.2.1})~; on peut donc se limiter au cas où $S$ et
$S'$ sont affines, et il suffit de prouver que $X'$ est alors un schéma
affine (\hyperref[II.1.3.4-fr]{1.3.4}). Mais alors
(\hyperref[II.1.3.4-fr]{1.3.4}) $X$ est un schéma affine, et si $A$, $A'$ et
$B$ sont les anneaux de $S$, $S'$ et $X$ respectivement, on sait que $X'$ est
un schéma affine d'anneau $A'\otimes_A B$
(\hyperref[I.3.2.2-fr]{I, 3.2.2}).

\begin{corollary}[1.5.2]
\label{II.1.5.2-fr}
Sous les hypothèses de (\hyperref[II.1.5.1-fr]{1.5.1}), soient $f:X\to S$
le morphisme structural, $f':X'\to S'$, $g':X'\to X$ les projections, de
sorte que le diagramme
\[
  \xymatrix{
    X\ar[d] & X'\ar[l]_{g'}\ar[d]^{f'}\\
    S & S'\ar[l]_g
  }
\]
est commutatif. Pour tout $\mathcal{O}_X$-Module quasi-cohérent $\mathcal{F}$,
il existe un isomorphisme canonique de $\mathcal{O}_{S'}$-Modules
\[
  u:g^*(f_*(\mathcal{F}))\xrightarrow{\sim}
  f'_*(g^{\prime*}(\mathcal{F})).
  \tag{1.5.2.1}\label{II.1.5.2.1-fr}
\]
En particulier, il existe un isomorphisme canonique de $\mathcal{A}(X')$ sur
$g^*(\mathcal{A}(X))$.
\end{corollary}

Pour définir $u$, il suffit de définir un homomorphisme
\[
  v:f_*(\mathcal{F})\longrightarrow
  g_*(f'_*(g^{\prime*}(\mathcal{F})))
  =f_*(g'_*(g^{\prime*}(\mathcal{F})))
\]
et de prendre $u=v^\sharp$ (\hyperref[0.4.4.3-fr]{0, 4.4.3}). Nous prendrons
$v=f_*(\rho)$, où $\rho$ est l'homomorphisme canonique
$\mathcal{F}\to g'_*(g^{\prime*}(\mathcal{F}))$
(\hyperref[0.4.4.3-fr]{0, 4.4.3}). Pour prouver que $u$ est un isomorphisme,
on peut se borner au cas où $S$ et $S'$, donc $X$ et $X'$, sont affines~;
avec les notations de (\hyperref[II.1.5.1-fr]{1.5.1}), on a alors
$\mathcal{F}=\widetilde{M}$, où $M$ est un $B$-module. On constate alors
aussitôt que $g^*(f_*(\mathcal{F}))$ et
$f'_*(g^{\prime*}(\mathcal{F}))$ sont tous deux égaux au
$\mathcal{O}_{S'}$-Module associé au $A'$-module $A'\otimes_A M$ (où $M$ est
considéré comme $A$-module), et que $u$ est l'homomorphisme associé à
l'identité (\hyperref[I.1.6.3-fr]{I, 1.6.3},
\hyperref[I.1.6.5-fr]{1.6.5} et \hyperref[I.1.6.7-fr]{1.6.7}).

\begin{remark}[1.5.3]
\label{II.1.5.3-fr}
On se gardera de croire que (\hyperref[II.1.5.2-fr]{1.5.2}) reste valable
lorsque $X$ n'est plus supposé affine au-dessus de $S$, même lorsque
$S'=\operatorname{Spec}(k(s))$ ($s\in S$) et que $S'\to S$ est le morphisme
canonique (\hyperref[I.2.4.5-fr]{I, 2.4.5}) --- auquel cas $X'$ n'est autre
que la \emph{fibre} $f^{-1}(s)$ (\hyperref[I.3.6.2-fr]{I, 3.6.2}). En
d'autres termes, lorsque $X$ n'est pas affine au-dessus de $S$, l'opération
\oldpage[II]{13}
<<~image directe de faisceaux quasi-cohérents~>> ne permute pas à l'opération
de <<~passage aux fibres~>>. Nous verrons cependant au chapitre~III
(\hyperref[III.4.2.4-fr]{III, 4.2.4}) un résultat dans ce sens, de nature
<<~asymptotique~>>, valable pour les faisceaux \emph{cohérents} sur $X$ lorsque
$f$ est propre (\hyperref[II.5.4-fr]{5.4}) et $S$ noethérien.
\end{remark}

\begin{corollary}[1.5.4]
\label{II.1.5.4-fr}
Pour tout préschéma $X$ affine au-dessus de $S$ et tout $s\in S$, la fibre
$f^{-1}(s)$ (où $f$ désigne le morphisme structural $X\to S$) est un schéma
affine.
\end{corollary}

Il suffit d'appliquer (\hyperref[II.1.5.1-fr]{1.5.1}) à
$S'=\operatorname{Spec}(k(s))$ et d'utiliser
(\hyperref[II.1.3.4-fr]{1.3.4}).

\begin{corollary}[1.5.5]
\label{II.1.5.5-fr}
Soient $X$ un $S$-préschéma, $S'$ un préschéma affine au-dessus de $S$~; alors
$X'=X_{(S')}$ est un préschéma affine au-dessus de $X$. En outre, si
$f:X\to S$ est le morphisme structural, il y a un isomorphisme canonique de
$\mathcal{O}_X$-Algèbres
$\mathcal{A}(X')\xrightarrow{\sim}f^*(\mathcal{A}(S'))$, et pour tout
$\mathcal{A}(S')$-Module quasi-cohérent $\mathcal{M}$, un di-isomorphisme
canonique
$f^*(\mathcal{M})\xrightarrow{\sim}
\mathcal{A}(f^{\prime*}(\widetilde{\mathcal{M}}))$, en désignant par
$f'=f_{(S')}$ le morphisme structural $X'\to S'$.
\end{corollary}

Il suffit d'intervertir les rôles de $X$ et de $S'$ dans
(\hyperref[II.1.5.1-fr]{1.5.1}) et
(\hyperref[II.1.5.2-fr]{1.5.2}).

\begin{env}[1.5.6]
\label{II.1.5.6-fr}
Soient maintenant $S$, $S'$ deux préschémas, $q:S'\to S$ un morphisme,
$\mathcal{B}$ (resp. $\mathcal{B}'$) une $\mathcal{O}_S$-Algèbre (resp. une
$\mathcal{O}_{S'}$-Algèbre) quasi-cohérente, $u:\mathcal{B}\to\mathcal{B}'$
un $q$-morphisme (c'est-à-dire un homomorphisme
$\mathcal{B}\to q_*(\mathcal{B}')$ de $\mathcal{O}_S$-Algèbres). Si
$X=\operatorname{Spec}(\mathcal{B})$,
$X'=\operatorname{Spec}(\mathcal{B}')$, on en déduit canoniquement un
morphisme
\[
  v=\operatorname{Spec}(u):X'\to X
\]
tel que le diagramme
\[
  \xymatrix{
    X'\ar[r]^v\ar[d] & X\ar[d]\\
    S'\ar[r]_q & S
  }
  \tag{1.5.6.1}\label{II.1.5.6.1-fr}
\]
soit commutatif (les flèches verticales étant les morphismes structuraux). En
effet, la donnée de $u$ équivaut à celle d'un homomorphisme de
$\mathcal{O}_{S'}$-Algèbres quasi-cohérentes
$u^\sharp:q^*(\mathcal{B})\to\mathcal{B}'$
(\hyperref[0.4.4.3-fr]{0, 4.4.3})~; elle définit donc canoniquement un
$S'$-morphisme
\[
  w:\operatorname{Spec}(\mathcal{B}')\to
  \operatorname{Spec}(q^*(\mathcal{B}))
\]
tel que $\mathcal{A}(w)=u^\sharp$
(\hyperref[II.1.2.7-fr]{1.2.7}). D'autre part, il résulte de
(\hyperref[II.1.5.2-fr]{1.5.2}) que
$\operatorname{Spec}(q^*(\mathcal{B}))$ s'identifie canoniquement à
$X\times_S S'$~; le morphisme $v$ est le composé
$X'\xrightarrow{w}X\times_S S'\xrightarrow{p_1}X$ de $w$ avec la première
projection, et la commutativité de
(\hyperref[II.1.5.6.1-fr]{1.5.6.1}) résulte des définitions. Soient $U$ (resp.
$U'$) un ouvert affine de $S$ (resp. $S'$) tels que $q(U')\subset U$,
$A=\Gamma(U,\mathcal{O}_S)$, $A'=\Gamma(U',\mathcal{O}_{S'})$ leurs anneaux,
$B=\Gamma(U,\mathcal{B})$, $B'=\Gamma(U',\mathcal{B}')$~; la restriction de
$u$ en un $(q|U')$-morphisme
$\mathcal{B}|U\to\mathcal{B}'|U'$ correspond à un di-homomorphisme d'algèbres
$B\to B'$~; si $V$, $V'$ sont les images réciproques de $U$, $U'$ dans $X$,
$X'$ respectivement, par les morphismes structuraux, le morphisme $V'\to V$,
restriction de $v$, correspond (\hyperref[I.1.7.3-fr]{I, 1.7.3}) au
di-homomorphisme précédent.
\end{env}

\begin{env}[1.5.7]
\label{II.1.5.7-fr}
Sous les mêmes hypothèses que dans (\hyperref[II.1.5.6-fr]{1.5.6}), soit
$\mathcal{M}$ un $\mathcal{B}$-Module quasi-cohérent~; il y a alors un
isomorphisme canonique de $\mathcal{O}_{X'}$-Modules
\[
  v^*(\widetilde{\mathcal{M}})\xrightarrow{\sim}
  \bigl(q^*(\mathcal{M})\otimes_{q^*(\mathcal{B})}\mathcal{B}'\bigr)^{\sim}
  \tag{1.5.7.1}\label{II.1.5.7.1-fr}
\]

\oldpage[II]{14}
En effet, l'isomorphisme canonique
(\hyperref[II.1.5.2.1-fr]{1.5.2.1}) fournit un isomorphisme canonique de
$p_1^*(\widetilde{\mathcal{M}})$ sur le faisceau sur
$\operatorname{Spec}(q^*(\mathcal{B}))$ associé au
$q^*(\mathcal{B})$-Module $q^*(\mathcal{M})$, et il suffit ensuite d'appliquer
(\hyperref[II.1.4.7-fr]{1.4.7}).
\end{env}

\subsection{Morphismes affines}
\label{II.1.6-fr}

\begin{env}[1.6.1]
\label{II.1.6.1-fr}
Nous dirons qu'un morphisme $f:X\to Y$ de préschémas est \emph{affine} s'il
définit $X$ comme un préschéma affine au-dessus de $Y$. Les propriétés des
préschémas affines au-dessus d'un autre se traduisent comme suit dans ce
langage~:
\end{env}

\begin{proposition}[1.6.2]
\label{II.1.6.2-fr}
\begin{enumerate}
  \item[\rm(i)] Une immersion fermée est affine.
  \item[\rm(ii)] Le composé de deux morphismes affines est affine.
  \item[\rm(iii)] Si $f:X\to Y$ est un $S$-morphisme affine,
    $f_{(S')}:X_{(S')}\to Y_{(S')}$ est affine pour toute extension $S'\to S$
    de la base.
  \item[\rm(iv)] Si $f:X\to Y$ et $f':X'\to Y'$ sont deux $S$-morphismes
    affines,
    \[
      f\times_S f':X\times_S X'\to Y\times_S Y'
    \]
    est affine.
  \item[\rm(v)] Si $f:X\to Y$ et $g:Y\to Z$ sont deux morphismes tels que
    $g\circ f$ soit affine et $g$ séparé, alors $f$ est affine.
  \item[\rm(vi)] Si $f$ est affine, il en est de même de
    $f_{\mathrm{red}}$.
\end{enumerate}
\end{proposition}

En vertu de (\hyperref[I.5.5.12-fr]{I, 5.5.12}), il suffit de démontrer (i),
(ii) et (iii). Or (i) n'est autre que l'exemple
(\hyperref[II.1.2.2-fr]{1.2.2}), et (ii) n'est autre que
(\hyperref[II.1.3.5-fr]{1.3.5})~; enfin, (iii) découle de
(\hyperref[II.1.5.1-fr]{1.5.1}) puisque $X_{(S')}$ s'identifie au produit
$X\times_Y Y_{(S')}$ (\hyperref[I.3.3.11-fr]{I, 3.3.11}).

\begin{corollary}[1.6.3]
\label{II.1.6.3-fr}
Si $X$ est un schéma affine et $Y$ un schéma, tout morphisme $X\to Y$ est
affine.
\end{corollary}

\begin{proposition}[1.6.4]
\label{II.1.6.4-fr}
Soient $Y$ un préschéma localement noethérien, $f:X\to Y$ un morphisme de type
fini. Pour que $f$ soit affine, il faut et il suffit que $f_{\mathrm{red}}$ le
soit.
\end{proposition}

Vu (\hyperref[II.1.6.2-fr]{1.6.2, (vi)}), nous n'avons à démontrer que la
suffisance de la condition. Il suffit de prouver que si $Y$ est affine et
noethérien, $X$ est affine~; or $Y_{\mathrm{red}}$ est alors affine, donc il
en est de même de $X_{\mathrm{red}}$ par hypothèse. Or, $X$ est noethérien,
donc la conclusion résulte de (\hyperref[I.6.1.7-fr]{I, 6.1.7}).

\subsection{Fibré vectoriel associé à un Faisceau de modules}
\label{II.1.7-fr}

\begin{env}[1.7.1]
\label{II.1.7.1-fr}
Soient $A$ un anneau, $E$ un $A$-module. Rappelons qu'on appelle \emph{algèbre
symétrique} sur $E$ et qu'on note $\mathbf{S}(E)$ (ou $\mathbf{S}_A(E)$)
l'algèbre quotient de l'algèbre tensorielle $\mathbf{T}(E)$ par l'idéal
bilatère engendré par les éléments $x\otimes y-y\otimes x$, où $x,y$
parcourent $E$. L'algèbre $\mathbf{S}(E)$ est caractérisée par la propriété
universelle suivante~: si $\sigma$ est l'application canonique
$E\to\mathbf{S}(E)$ (obtenue par composition de $E\to\mathbf{T}(E)$ et de
l'application canonique $\mathbf{T}(E)\to\mathbf{S}(E)$), toute application
$A$-linéaire $E\to B$, où $B$ est une $A$-algèbre \emph{commutative}, se
factorise de façon unique en
$E\xrightarrow{\sigma}\mathbf{S}(E)\xrightarrow{g}B$, où $g$ est un
$A$-homomorphisme \emph{d'algèbres}. On déduit aussitôt de cette
caractérisation que pour deux $A$-modules $E$, $F$, on a
\[
  \mathbf{S}(E\oplus F)=\mathbf{S}(E)\otimes\mathbf{S}(F)
\]

\oldpage[II]{15}
à un isomorphisme canonique près~; en outre, $\mathbf{S}(E)$ est un foncteur
covariant en $E$, de la catégorie des $A$-modules dans celle des $A$-algèbres
commutatives~; enfin, la caractérisation précédente montre aussi que si
$E=\varinjlim E_\lambda$, on a
$\mathbf{S}(E)=\varinjlim\mathbf{S}(E_\lambda)$ à un isomorphisme canonique
près. Par abus de langage, un produit
$\sigma(x_1)\sigma(x_2)\cdots\sigma(x_n)$, où les $x_i\in E$, se note souvent
$x_1x_2\cdots x_n$ si aucune confusion n'en résulte. L'algèbre
$\mathbf{S}(E)$ est \emph{graduée}, $\mathbf{S}_n(E)$ étant l'ensemble des
combinaisons linéaires des produits de $n$ éléments de $E$ ($n\geq 0$)~;
l'algèbre $\mathbf{S}(A)$ est canoniquement isomorphe à l'algèbre des
polynômes $A[T]$ à une indéterminée et l'algèbre $\mathbf{S}(A^n)$ à
l'algèbre des polynômes à $n$ indéterminées $A[T_1,\ldots,T_n]$.
\end{env}

\begin{env}[1.7.2]
\label{II.1.7.2-fr}
Soit $\varphi$ un homomorphisme d'anneaux $A\to B$. Si $F$ est un $B$-module,
l'application canonique $F\to\mathbf{S}(F)$ donne une application canonique
$F_{[\varphi]}\to\mathbf{S}(F)_{[\varphi]}$, qui se factorise donc en
$F_{[\varphi]}\to\mathbf{S}(F_{[\varphi]})\to
\mathbf{S}(F)_{[\varphi]}$~; l'homomorphisme canonique
$\mathbf{S}(F_{[\varphi]})\to\mathbf{S}(F)_{[\varphi]}$ est surjectif, mais
non nécessairement bijectif. Si $E$ est un $A$-module, tout di-homomorphisme
$E\to F$ (c'est-à-dire tout $A$-homomorphisme $E\to F_{[\varphi]}$) donne donc
canoniquement un $A$-homomorphisme d'algèbres
$\mathbf{S}(E)\to\mathbf{S}(F_{[\varphi]})\to
\mathbf{S}(F)_{[\varphi]}$, c'est-à-dire un di-homomorphisme d'algèbres
$\mathbf{S}(E)\to\mathbf{S}(F)$.

Avec les mêmes notations, pour tout $A$-module $E$,
$\mathbf{S}(E\otimes_A B)$ s'identifie canoniquement à l'algèbre
$\mathbf{S}(E)\otimes_A B$~; cela résulte aussitôt de la propriété universelle
de $\mathbf{S}(E)$ (\hyperref[II.1.7.1-fr]{1.7.1}).
\end{env}

\begin{env}[1.7.3]
\label{II.1.7.3-fr}
Soit $R$ une partie multiplicative de l'anneau $A$~; appliquant
(\hyperref[II.1.7.2-fr]{1.7.2}) à l'anneau $B=R^{-1}A$, et se rappelant que
$R^{-1}E=E\otimes_A R^{-1}A$, on voit que l'on a
$\mathbf{S}(R^{-1}E)=R^{-1}\mathbf{S}(E)$ à un isomorphisme canonique près.
En outre, si $R'\supset R$ est une seconde partie multiplicative de $A$, le
diagramme
\[
  \xymatrix{
    R^{-1}E\ar[r]\ar[d] & {R'}^{-1}E\ar[d]\\
    \mathbf{S}(R^{-1}E)\ar[r] & \mathbf{S}({R'}^{-1}E)
  }
\]
est commutatif.
\end{env}

\begin{env}[1.7.4]
\label{II.1.7.4-fr}
Soit maintenant $(S,\mathcal{A})$ un espace annelé, et soit $\mathcal{E}$ un
$\mathcal{A}$-Module sur $S$. Si, à tout ouvert $U\subset S$ on associe le
$\Gamma(U,\mathcal{A})$-module
$\mathbf{S}(\Gamma(U,\mathcal{E}))$, on définit (vu le caractère fonctoriel
de $\mathbf{S}(E)$ (\hyperref[II.1.7.2-fr]{1.7.2})) un préfaisceau d'algèbres~;
on dit que le faisceau associé, que l'on note $\mathbf{S}(\mathcal{E})$ ou
$\mathbf{S}_{\mathcal{A}}(\mathcal{E})$ est la $\mathcal{A}$-Algèbre
\emph{symétrique} du $\mathcal{A}$-Module $\mathcal{E}$. Il résulte aussitôt
de (\hyperref[II.1.7.1-fr]{1.7.1}) que $\mathbf{S}(\mathcal{E})$ est solution
d'un problème universel~: tout homomorphisme de $\mathcal{A}$-Modules
$\mathcal{E}\to\mathcal{B}$, où $\mathcal{B}$ est une $\mathcal{A}$-Algèbre,
se factorise d'une seule manière en
$\mathcal{E}\to\mathbf{S}(\mathcal{E})\to\mathcal{B}$, la seconde flèche étant
un homomorphisme de $\mathcal{A}$-Algèbres. Il y a donc correspondance
biunivoque entre homomorphismes $\mathcal{E}\to\mathcal{B}$ de
$\mathcal{A}$-Modules et homomorphismes
$\mathbf{S}(\mathcal{E})\to\mathcal{B}$ de $\mathcal{A}$-Algèbres. En
particulier, tout homomorphisme $u:\mathcal{E}\to\mathcal{F}$ de
$\mathcal{A}$-Modules définit un homomorphisme
$\mathbf{S}(u):\mathbf{S}(\mathcal{E})\to\mathbf{S}(\mathcal{F})$ de
$\mathcal{A}$-Algèbres et $\mathbf{S}(\mathcal{E})$ est donc un foncteur
covariant en $\mathcal{E}$.

\oldpage[II]{16}
En vertu de (\hyperref[II.1.7.2-fr]{1.7.2}) et de la permutabilité de
$\mathbf{S}$ avec les limites inductives, on a
$(\mathbf{S}(\mathcal{E}))_s=\mathbf{S}(\mathcal{E}_s)$ pour tout point
$s\in S$. Si $\mathcal{E}$, $\mathcal{F}$ sont deux $\mathcal{A}$-Modules,
$\mathbf{S}(\mathcal{E}\oplus\mathcal{F})$ s'identifie canoniquement à
$\mathbf{S}(\mathcal{E})\otimes_{\mathcal{A}}\mathbf{S}(\mathcal{F})$,
comme on le voit sur les préfaisceaux correspondants.

Notons aussi que $\mathbf{S}(\mathcal{E})$ est une $\mathcal{A}$-Algèbre
graduée, somme directe infinie des $\mathbf{S}_n(\mathcal{E})$, où le
$\mathcal{A}$-Module $\mathbf{S}_n(\mathcal{E})$ est le faisceau associé au
préfaisceau $U\mapsto\mathbf{S}_n(\Gamma(U,\mathcal{E}))$. Si on prend en
particulier $\mathcal{E}=\mathcal{A}$, on voit que
$\mathbf{S}_{\mathcal{A}}(\mathcal{A})$ s'identifie à
$\mathcal{A}[T]=\mathcal{A}\otimes_{\mathbf{Z}}\mathbf{Z}[T]$
($T$ indéterminée, $\mathbf{Z}$ étant considéré comme faisceau simple).
\end{env}

\begin{env}[1.7.5]
\label{II.1.7.5-fr}
Soient $(T,\mathcal{B})$ un second espace annelé, $f$ un morphisme
$(S,\mathcal{A})\to(T,\mathcal{B})$. Si $\mathcal{F}$ est un
$\mathcal{B}$-Module, $\mathbf{S}(f^*(\mathcal{F}))$ s'identifie
canoniquement à $f^*(\mathbf{S}(\mathcal{F}))$~; en effet, si
$f=(\psi,\theta)$, on a par définition
(\hyperref[0.4.3.1-fr]{0, 4.3.1})
\[
  \mathbf{S}(f^*(\mathcal{F}))
  =\mathbf{S}(\psi^*(\mathcal{F})
    \otimes_{\psi^*(\mathcal{B})}\mathcal{A})
  =\mathbf{S}(\psi^*(\mathcal{F}))
    \otimes_{\psi^*(\mathcal{B})}\mathcal{A}
\]
(\hyperref[II.1.7.2-fr]{1.7.2})~; pour tout ouvert $U$ de $S$ et toute
section $h$ de $\mathbf{S}(\psi^*(\mathcal{F}))$ au-dessus de $U$, $h$
coïncide, dans un voisinage $V$ de chaque point $s\in U$, avec un élément de
$\mathbf{S}(\Gamma(V,\psi^*(\mathcal{F})))$~; si on se reporte à la
définition de $\psi^*(\mathcal{F})$
(\hyperref[0.3.7.1-fr]{0, 3.7.1}) et qu'on tient compte de ce que tout élément
de $\mathbf{S}(E)$ pour un module $E$ est combinaison linéaire d'un nombre
fini de produits d'éléments de $E$, on voit qu'il y a un voisinage $W$ de
$\psi(s)$ dans $T$, une section $h'$ de $\mathbf{S}(\mathcal{F})$ au-dessus
de $W$ et un voisinage $V'\subset V\cap\psi^{-1}(W)$ de $s$ tels que $h$
coïncide avec $t\mapsto h'(\psi(t))$ au-dessus de $V'$~; d'où notre
assertion.
\end{env}

\begin{proposition}[1.7.6]
\label{II.1.7.6-fr}
Soient $A$ un anneau, $S=\operatorname{Spec}(A)$ son spectre premier,
$\mathcal{E}=\widetilde{M}$ le $\mathcal{O}_S$-Module associé à un $A$-module
$M$~; alors la $\mathcal{O}_S$-Algèbre $\mathbf{S}(\mathcal{E})$ est associée
à la $A$-algèbre $\mathbf{S}(M)$.
\end{proposition}

\begin{proof}
En effet, pour tout $f\in A$,
$\mathbf{S}(M_f)=(\mathbf{S}(M))_f$
(\hyperref[II.1.7.3-fr]{1.7.3}) et la proposition résulte donc des définitions
(\hyperref[I.1.3.4-fr]{I, 1.3.4}).
\end{proof}

\begin{corollary}[1.7.7]
\label{II.1.7.7-fr}
Si $S$ est un préschéma, $\mathcal{E}$ un $\mathcal{O}_S$-Module
quasi-cohérent, la $\mathcal{O}_S$-Algèbre $\mathbf{S}(\mathcal{E})$ est
quasi-cohérente. Si en outre $\mathcal{E}$ est de type fini, chacun des
$\mathcal{O}_S$-Modules $\mathbf{S}_n(\mathcal{E})$ est de type fini.
\end{corollary}

\begin{proof}
La première assertion est conséquence immédiate de
(\hyperref[II.1.7.6-fr]{1.7.6}) et de
(\hyperref[I.1.4.1-fr]{I, 1.4.1})~; la seconde résulte de ce que, si $E$ est
un $A$-module de type fini, $\mathbf{S}_n(E)$ est un $A$-module de type
fini~; on applique alors (\hyperref[I.1.3.13-fr]{I, 1.3.13}).
\end{proof}

\begin{definition}[1.7.8]
\label{II.1.7.8-fr}
Soit $\mathcal{E}$ un $\mathcal{O}_S$-Module quasi-cohérent. On appelle
\emph{fibré vectoriel sur $S$ défini par $\mathcal{E}$} et on note
$\mathbf{V}(\mathcal{E})$ le spectre
(\hyperref[II.1.3.1-fr]{1.3.1}) de la $\mathcal{O}_S$-Algèbre quasi-cohérente
$\mathbf{S}(\mathcal{E})$.
\end{definition}

En vertu de (\hyperref[II.1.2.7-fr]{1.2.7}), pour tout $S$-préschéma $X$, il
y a correspondance biunivoque canonique entre les $S$-morphismes
$X\to\mathbf{V}(\mathcal{E})$ et les homomorphismes de
$\mathcal{O}_S$-Algèbres $\mathbf{S}(\mathcal{E})\to\mathcal{A}(X)$, donc
aussi entre ces $S$-morphismes et les homomorphismes de
$\mathcal{O}_S$-Modules
$\mathcal{E}\to\mathcal{A}(X)=f_*(\mathcal{O}_X)$ (où $f$ est le morphisme
structural $X\to S$). En particulier~:

\begin{env}[1.7.9]
\label{II.1.7.9-fr}
Prenons pour $X$ un sous-préschéma induit par $S$ sur un ouvert $U\subset S$.
Alors, les $S$-morphismes $U\to\mathbf{V}(\mathcal{E})$ ne sont autres que
les $U$-sections (\hyperref[I.2.5.5-fr]{I, 2.5.5}) du $U$-préschéma induit
par $\mathbf{V}(\mathcal{E})$ sur l'ouvert $p^{-1}(U)$ (où $p$ est le
morphisme structural $\mathbf{V}(\mathcal{E})\to S$). D'après ce qu'on vient
de voir, ces $U$-sections correspondent biunivoquement aux homomorphismes de
$\mathcal{O}_S$-Modules
$\mathcal{E}\to j_*(\mathcal{O}_S|U)$ (où $j$ est l'injection canonique
$U\to S$), ou,

\oldpage[II]{17}
ce qui revient au même (\hyperref[0.4.4.3-fr]{0, 4.4.3}) aux
$(\mathcal{O}_S|U)$-homomorphismes
$j^*(\mathcal{E})=\mathcal{E}|U\to\mathcal{O}_S|U$. En outre, il est
immédiat qu'à la restriction à un ouvert $U'\subset U$ d'un $S$-morphisme
$U\to\mathbf{V}(\mathcal{E})$ correspond la restriction à $U'$ de
l'homomorphisme correspondant
$\mathcal{E}|U\to\mathcal{O}_S|U$. On en conclut que \emph{le faisceau des
germes de $S$-sections} de $\mathbf{V}(\mathcal{E})$ s'identifie
canoniquement au \emph{dual $\check{\mathcal{E}}$} de $\mathcal{E}$.

En particulier, si on prend $X=U=S$, à l'homomorphisme \emph{nul}
$\mathcal{E}\to\mathcal{O}_S$ correspond une $S$-section canonique de
$\mathbf{V}(\mathcal{E})$, dite \emph{$S$-section nulle}
(cf. (\hyperref[II.8.3.3-fr]{8.3.3})).
\end{env}

\begin{env}[1.7.10]
\label{II.1.7.10-fr}
Prenons maintenant pour $X$ le spectre $\{\xi\}$ d'un corps $K$~; le
morphisme structural $f:X\to S$ correspond donc à un monomorphisme
$k(s)\to K$, où $s=f(\xi)$
(\hyperref[I.2.4.6-fr]{I, 2.4.6})~; les $S$-morphismes
$\{\xi\}\to\mathbf{V}(\mathcal{E})$ ne sont donc autres que les
\emph{points géométriques de $\mathbf{V}(\mathcal{E})$ à valeurs dans
l'extension $K$ de $k(s)$} (\hyperref[I.3.4.5-fr]{I, 3.4.5}), points qui
sont localisés aux points de $p^{-1}(s)$. L'ensemble de ces points, qu'on
peut encore appeler \emph{la fibre géométrique rationnelle sur $K$ de
$\mathbf{V}(\mathcal{E})$ au-dessus du point $s$}, s'identifie d'après
(\hyperref[II.1.7.8-fr]{1.7.8}) à l'ensemble des homomorphismes de
$\mathcal{O}_S$-Modules
$\mathcal{E}\to f_*(\mathcal{O}_X)$, ou, ce qui revient au même
(\hyperref[0.4.4.3-fr]{0, 4.4.3}) à l'ensemble des homomorphismes de
$\mathcal{O}_X$-Modules
$f^*(\mathcal{E})\to\mathcal{O}_X=K$. Mais on a par définition
(\hyperref[0.4.3.1-fr]{0, 4.3.1})
$f^*(\mathcal{E})=\mathcal{E}_s\otimes_{\mathcal{O}_s}K
=\mathcal{E}^s\otimes_{k(s)}K$, en posant
$\mathcal{E}^s=\mathcal{E}_s/\mathfrak{m}_s\mathcal{E}_s$~; la fibre
géométrique de $\mathbf{V}(\mathcal{E})$ rationnelle sur $K$ au-dessus de
$s$ s'identifie donc au \emph{dual} du \emph{$K$-espace vectoriel
$\mathcal{E}^s\otimes_{k(s)}K$}~; si $\mathcal{E}^s$ ou $K$ est de
dimension finie sur $k(s)$, ce dual s'identifie aussi à
$(\mathcal{E}^s)^\vee\otimes_{k(s)}K$, en désignant par
$(\mathcal{E}^s)^\vee$ le dual du $k(s)$-espace vectoriel
$\mathcal{E}^s$.
\end{env}

\begin{proposition}[1.7.11]
\label{II.1.7.11-fr}
\begin{enumerate}
  \item[\rm(i)] $\mathbf{V}(\mathcal{E})$ est un foncteur contravariant en
    $\mathcal{E}$ de la catégorie des $\mathcal{O}_S$-Modules
    quasi-cohérents dans la catégorie des $S$-schémas affines.
  \item[\rm(ii)] Si $\mathcal{E}$ est un $\mathcal{O}_S$-Module de type
    fini, $\mathbf{V}(\mathcal{E})$ est de type fini sur $S$.
  \item[\rm(iii)] Si $\mathcal{E}$ et $\mathcal{F}$ sont deux
    $\mathcal{O}_S$-Modules quasi-cohérents,
    $\mathbf{V}(\mathcal{E}\oplus\mathcal{F})$ s'identifie canoniquement à
    $\mathbf{V}(\mathcal{E})\times_S\mathbf{V}(\mathcal{F})$.
  \item[\rm(iv)] Soit $g:S'\to S$ un morphisme~; pour tout
    $\mathcal{O}_S$-Module quasi-cohérent $\mathcal{E}$,
    $\mathbf{V}(g^*(\mathcal{E}))$ s'identifie canoniquement à
    $\mathbf{V}(\mathcal{E})_{(S')}
    =\mathbf{V}(\mathcal{E})\times_S S'$.
  \item[\rm(v)] À un homomorphisme surjectif
    $\mathcal{E}\to\mathcal{F}$ de $\mathcal{O}_S$-Modules quasi-cohérents
    correspond une immersion fermée
    $\mathbf{V}(\mathcal{F})\to\mathbf{V}(\mathcal{E})$.
\end{enumerate}
\end{proposition}

\begin{proof}
(i) est conséquence immédiate de
(\hyperref[II.1.2.7-fr]{1.2.7}), compte tenu de ce que tout homomorphisme de
$\mathcal{O}_S$-Modules $\mathcal{E}\to\mathcal{F}$ définit canoniquement un
homomorphisme de $\mathcal{O}_S$-Algèbres
$\mathbf{S}(\mathcal{E})\to\mathbf{S}(\mathcal{F})$. (ii) découle
immédiatement des définitions (\hyperref[I.6.3.1-fr]{I, 6.3.1}) et du fait
que si $E$ est un $A$-module de type fini, $\mathbf{S}(E)$ est une
$A$-algèbre de type fini. Pour démontrer (iii), il suffit de partir de
l'isomorphisme canonique
$\mathbf{S}(\mathcal{E}\oplus\mathcal{F})
\simeq\mathbf{S}(\mathcal{E})
\otimes_{\mathcal{O}_S}\mathbf{S}(\mathcal{F})$
(\hyperref[II.1.7.4-fr]{1.7.4}) et d'appliquer
(\hyperref[II.1.4.6-fr]{1.4.6}). De même, pour démontrer (iv), il suffit de
partir de l'isomorphisme canonique
$\mathbf{S}(g^*(\mathcal{E}))
\simeq g^*(\mathbf{S}(\mathcal{E}))$
(\hyperref[II.1.7.5-fr]{1.7.5}) et d'appliquer
(\hyperref[II.1.5.2-fr]{1.5.2}). Enfin, pour établir (v), il suffit de
remarquer que si l'homomorphisme $\mathcal{E}\to\mathcal{F}$ est surjectif,
il en est de même de l'homomorphisme correspondant
$\mathbf{S}(\mathcal{E})\to\mathbf{S}(\mathcal{F})$ de
$\mathcal{O}_S$-Algèbres, et la conclusion résulte de
(\hyperref[II.1.4.10-fr]{1.4.10}).
\end{proof}

\begin{env}[1.7.12]
\label{II.1.7.12-fr}
Prenons en particulier $\mathcal{E}=\mathcal{O}_S$~; le préschéma
$\mathbf{V}(\mathcal{O}_S)$ est le $S$-schéma affine, spectre de la
$\mathcal{O}_S$-Algèbre $\mathbf{S}(\mathcal{O}_S)$ qui s'identifie à la
$\mathcal{O}_S$-Algèbre
$\mathcal{O}_S[T]=\mathcal{O}_S\otimes_{\mathbf{Z}}\mathbf{Z}[T]$
\oldpage[II]{18}
($T$ indéterminée)~; c'est évident lorsque
$S=\operatorname{Spec}(\mathbf{Z})$, en vertu de
(\hyperref[II.1.7.6-fr]{1.7.6}), et on passe de là au cas général en
considérant le morphisme structural
$S\to\operatorname{Spec}(\mathbf{Z})$ et utilisant
(\hyperref[II.1.7.11-fr]{1.7.11, (iv)}). En raison de ce résultat, on pose
encore $\mathbf{V}(\mathcal{O}_S)=S[T]$, et on a donc la formule
\[
  S[T]=S\otimes_{\mathbf{Z}}\mathbf{Z}[T].
  \tag{1.7.12.1}\label{II.1.7.12.1-fr}
\]

L'identification du faisceau des germes de $S$-sections de $S[T]$ avec
$\mathcal{O}_S$, déjà vue dans
(\hyperref[I.3.3.15-fr]{I, 3.3.15}), se retrouve ici dans un contexte plus
général, comme cas particulier de (\hyperref[II.1.7.9-fr]{1.7.9}).
\end{env}

\begin{env}[1.7.13]
\label{II.1.7.13-fr}
Pour tout $S$-préschéma $X$, on a vu
(\hyperref[II.1.7.8-fr]{1.7.8}) que
$\operatorname{Hom}_S(X,S[T])$ s'identifie canoniquement à
$\operatorname{Hom}_{\mathcal{O}_S}(\mathcal{O}_S,\mathcal{A}(X))$, lequel
est canoniquement isomorphe à $\Gamma(S,\mathcal{A}(X))$, et est par suite
muni d'une structure d'anneau~; en outre, à tout $S$-morphisme $h:X\to Y$
correspond un morphisme
$\Gamma(\mathcal{A}(h)):\Gamma(S,\mathcal{A}(Y))
\to\Gamma(S,\mathcal{A}(X))$ pour les structures d'anneau
(\hyperref[II.1.1.2-fr]{1.1.2}). Quand on munit
$\operatorname{Hom}_S(X,S[T])$ de la structure d'anneau ainsi définie, on
voit donc qu'on peut considérer $\operatorname{Hom}_S(X,S[T])$ comme un
\emph{foncteur contravariant en $X$}, de la catégorie des $S$-préschémas
dans celle des anneaux. D'autre part,
$\operatorname{Hom}_S(X,\mathbf{V}(\mathcal{E}))$ s'identifie de même à
$\operatorname{Hom}_{\mathcal{O}_S}(\mathcal{E},\mathcal{A}(X))$ (où
$\mathcal{A}(X)$ est considéré comme un
\emph{$\mathcal{O}_S$-Module})~; on peut par suite le munir canoniquement
d'une structure de \emph{module} sur l'anneau
$\operatorname{Hom}_S(X,S[T])$, et on voit comme ci-dessus que le couple
\[
  \bigl(\operatorname{Hom}_S(X,S[T]),
  \operatorname{Hom}_S(X,\mathbf{V}(\mathcal{E}))\bigr)
\]
est un foncteur contravariant en $X$, à valeurs dans la catégorie dont les
éléments sont les couples $(A,M)$ formés d'un anneau $A$ et d'un $A$-module
$M$, les morphismes étant les di-homomorphismes.

Nous interpréterons ces faits en disant que $S[T]$ est un
\emph{$S$-schéma d'anneaux} et que $\mathbf{V}(\mathcal{E})$ est un
\emph{$S$-schéma de modules} sur le $S$-schéma d'anneaux $S[T]$
(cf. chap.~0, \textsection8).
\end{env}

\begin{env}[1.7.14]
\label{II.1.7.14-fr}
Nous allons voir que la structure de $S$-schéma de modules définie sur le
$S$-schéma $\mathbf{V}(\mathcal{E})$ permet de reconstituer le
$\mathcal{O}_S$-Module $\mathcal{E}$ à un isomorphisme unique près~: pour
cela, nous allons montrer que $\mathcal{E}$ est canoniquement isomorphe à un
sous-$\mathcal{O}_S$-Module de
$\mathbf{S}(\mathcal{E})=\mathcal{A}(\mathbf{V}(\mathcal{E}))$, défini au
moyen de cette structure. En effet
(\hyperref[II.1.7.4-fr]{1.7.4}) l'ensemble
$\operatorname{Hom}_{\mathcal{O}_S}(\mathbf{S}(\mathcal{E}),
\mathcal{A}(X))$ des homomorphismes de
\emph{$\mathcal{O}_S$-Algèbres} s'identifie canoniquement à
$\operatorname{Hom}_{\mathcal{O}_S}(\mathcal{E},\mathcal{A}(X))$, ensemble
d'homomorphismes de \emph{$\mathcal{O}_S$-Modules}~: si $h$, $h'$ sont deux
éléments de ce dernier ensemble, $s_i$ ($1\leq i\leq n$) des sections de
$\mathcal{E}$ au-dessus d'un ouvert $U\subset S$, $t$ une section de
$\mathcal{A}(X)$ au-dessus de $U$, on a par définition
\[
  (h+h')(s_1s_2\cdots s_n)
  =\prod_{i=1}^n\bigl(h(s_i)+h'(s_i)\bigr)
\]
et
\[
  (t.h)(s_1s_2\cdots s_n)=t^n\prod_{i=1}^n h(s_i).
\]

Cela étant, si $z$ est une section de $\mathbf{S}(\mathcal{E})$ au-dessus
de $U$, $h\mapsto h(z)$ est une application de
$\operatorname{Hom}_S(X,\mathbf{V}(\mathcal{E}))
=\operatorname{Hom}_{\mathcal{O}_S}(\mathbf{S}(\mathcal{E}),
\mathcal{A}(X))$ dans $\Gamma(U,\mathcal{A}(X))$. Nous allons
\oldpage[II]{19}
montrer que $\mathcal{E}$ s'identifie au sous-Module de
$\mathbf{S}(\mathcal{E})$ tel que, \emph{pour tout ouvert $U\subset S$,
toute section $z$ de ce sous-$\mathcal{O}_S$-Module au-dessus de $U$ et tout
$S$-préschéma $X$, l'application $h\mapsto h(z)$ de
$\operatorname{Hom}_{\mathcal{O}_S|U}(\mathbf{S}(\mathcal{E})|U,
\mathcal{A}(X)|U)$ dans $\Gamma(U,\mathcal{A}(X))$ soit un homomorphisme de
$\Gamma(U,\mathcal{A}(X))$-modules}.

Il est immédiat en effet que $\mathcal{E}$ possède cette propriété~; pour
démontrer la réciproque, on peut se borner à prouver que lorsque
$S=\operatorname{Spec}(A)$, $\mathcal{E}=\widetilde{M}$, une section $z$ de
$\mathbf{S}(\mathcal{E})$ au-dessus de $S$ qui (pour $U=S$) possède la
propriété énoncée ci-dessus, est nécessairement une section de
$\mathcal{E}$~; on a alors $z=\sum_{n=0}^{\infty}z_n$, où
$z_n\in\mathbf{S}_n(M)$, et il s'agit de prouver que $z_n=0$ pour $n\ne1$.
Posons $B=\mathbf{S}(M)$ et prenons pour $X$ le préschéma
$\operatorname{Spec}(B[T])$, où $T$ est une indéterminée. L'ensemble
$\operatorname{Hom}_{\mathcal{O}_S}(\mathbf{S}(\mathcal{E}),
\mathcal{A}(X))$ s'identifie à l'ensemble des homomorphismes d'anneaux
$h:B\to B[T]$ (\hyperref[I.1.3.13-fr]{I, 1.3.13}), et d'après ce qu'on a vu
plus haut, on a
$(T.h)(z)=\sum_{n=0}^{\infty}T^nh(z_n)$~: l'hypothèse sur $z$ implique que
l'on a
$\sum_{n=0}^{\infty}T^nh(z_n)=T.\sum_{n=0}^{\infty}h(z_n)$ pour tout
homomorphisme $h$. Si on prend en particulier pour $h$ l'injection
canonique, il vient
$\sum_{n=0}^{\infty}T^nz_n=T.\sum_{n=0}^{\infty}z_n$, ce qui entraîne bien
la conclusion $z_n=0$ pour $n\ne1$.
\end{env}

\begin{proposition}[1.7.15]
\label{II.1.7.15-fr}
Soit $Y$ un préschéma dont l'espace sous-jacent est noethérien, ou un schéma
quasi-compact. Tout $Y$-schéma affine $X$ de type fini sur $Y$ est
$Y$-isomorphe à un sous-$Y$-schéma fermé d'un $Y$-schéma de la forme
$\mathbf{V}(\mathcal{E})$, où $\mathcal{E}$ est un $\mathcal{O}_Y$-Module
quasi-cohérent de type fini.
\end{proposition}

\begin{proof}
En effet, la $\mathcal{O}_Y$-Algèbre quasi-cohérente $\mathcal{A}(X)$ est de
type fini (\hyperref[II.1.3.7-fr]{1.3.7}). Les hypothèses faites entraînent
que $\mathcal{A}(X)$ est engendrée par un sous-$\mathcal{O}_Y$-Module
quasi-cohérent de type fini $\mathcal{E}$
(\hyperref[I.9.6.5-fr]{I, 9.6.5})~; par définition, cela entraîne que
l'homomorphisme canonique
$\mathbf{S}(\mathcal{E})\to\mathcal{A}(X)$ prolongeant canoniquement
l'injection $\mathcal{E}\to\mathcal{A}(X)$ est \emph{surjectif}~; la
conclusion résulte alors de (\hyperref[II.1.4.10-fr]{1.4.10}).
\end{proof}

\section{Spectres premiers homogènes}
\label{section:II.2-fr}

\subsection{Généralités sur les anneaux et modules gradués.}
\label{subsection:II.2.1-fr}

\begin{notation}[2.1.1]
\label{II.2.1.1-fr}
Étant donné un anneau $S$ \emph{gradué} en degrés positifs, nous désignerons
par $S_n$ la partie de $S$ formée des éléments homogènes de degré $n$
($n\geq0$), par $S_+$ la somme (directe) des $S_n$ tels que $n>0$~; on a
$1\in S_0$, $S_0$ est un sous-anneau de $S$, $S_+$ un idéal gradué de $S$ et
$S$ est somme directe de $S_0$ et de $S_+$. Si $M$ est un \emph{module
gradué} sur $S$ (à degrés positifs ou négatifs), on désignera de même par
$M_n$ le $S_0$-module formé des éléments homogènes de degré $n$ de $M$ (avec
$n\in\mathbf{Z}$).

Pour tout entier $d>0$, on désigne par $S^{(d)}$ la somme directe des
$S_{nd}$~; en considérant les éléments de $S_{nd}$ comme homogènes de degré
$n$, les $S_{nd}$ définissent sur $S^{(d)}$ une structure d'anneau gradué.

Pour tout entier $k$ tel que $0\leq k\leq d-1$, on désignera par
$M^{(d,k)}$ la somme directe
\oldpage[II]{20}
des $M_{nd+k}$ ($n\in\mathbf{Z}$)~; c'est un $S^{(d)}$-module gradué, quand
on considère les éléments de $M_{nd+k}$ comme homogènes de degré $n$. On
écrira $M^{(d)}$ au lieu de $M^{(d,0)}$.

Avec les notations précédentes, pour tout entier $n$ (positif ou négatif),
nous désignerons par $M(n)$ le $S$-module gradué défini par
$(M(n))_k=M_{n+k}$ pour tout $k\in\mathbf{Z}$. En particulier, $S(n)$ sera
un $S$-module gradué tel que $(S(n))_k=S_{n+k}$, en convenant de poser
$S_n=0$ pour $n<0$. On dit qu'un $S$-module gradué $M$ est \emph{libre}
s'il est isomorphe, en tant que \emph{module gradué}, à une somme directe de
modules de la forme $S(n)$~; comme $S(n)$ est un $S$-module monogène,
engendré par l'élément $1$ de $S$ considéré comme élément de degré $-n$, il
revient au même de dire que $M$ admet une \emph{base} sur $S$ formée
d'éléments \emph{homogènes}.

On dit qu'un $S$-module gradué $M$ \emph{admet une présentation finie} s'il
existe une suite exacte $P\to Q\to M\to0$, où $P$ et $Q$ sont des sommes
directes finies de modules de la forme $S(n)$ et les homomorphismes sont de
degré $0$ (cf. (\hyperref[II.2.1.2-fr]{2.1.2})).
\end{notation}

\begin{env}[2.1.2]
\label{II.2.1.2-fr}
Soient $M$, $N$ deux $S$-modules gradués~; on définit sur $M\otimes_S N$ une
structure de $S$-module \emph{gradué} de la façon suivante. Sur le produit
tensoriel $M\otimes_{\mathbf{Z}}N$, on peut définir une structure de
$\mathbf{Z}$-module gradué (où $\mathbf{Z}$ est gradué par
$\mathbf{Z}_0=\mathbf{Z}$, $\mathbf{Z}_n=0$ pour $n\ne0$) en posant
\[
  (M\otimes_{\mathbf{Z}}N)_q
  =\bigoplus_{m+n=q}M_m\otimes_{\mathbf{Z}}N_n
\]
(comme $M$ et $N$ sont respectivement sommes directes des $M_m$ et des
$N_n$, on sait que l'on peut identifier canoniquement
$M\otimes_{\mathbf{Z}}N$ à la somme directe de tous les
$M_m\otimes_{\mathbf{Z}}N_n$). Cela étant, on a
$M\otimes_S N=(M\otimes_{\mathbf{Z}}N)/P$, où $P$ est le
sous-$\mathbf{Z}$-module de $M\otimes_{\mathbf{Z}}N$ engendré par les
éléments $(xs)\otimes y-x\otimes(sy)$ pour $x\in M$, $y\in N$, $s\in S$~;
il est clair que $P$ est un sous-$\mathbf{Z}$-module \emph{gradué} de
$M\otimes_{\mathbf{Z}}N$, et on voit aussitôt qu'on obtient bien en passant
au quotient une structure de $S$-module gradué sur $M\otimes_S N$.

Pour deux $S$-modules gradués $M$, $N$, rappelons qu'un homomorphisme
$u:M\to N$ de $S$-modules est dit \emph{de degré $k$} si
$u(M_j)\subset N_{j+k}$ pour tout $j\in\mathbf{Z}$. Si $H_n$ désigne
l'ensemble de tous les homomorphismes de degré $n$ de $M$ dans $N$, on
désigne par $\operatorname{Hom}_S(M,N)$ la \emph{somme} (directe) des $H_n$
($n\in\mathbf{Z}$) dans le $S$-module $H$ de tous les homomorphismes (de
$S$-module) de $M$ dans $N$~; en général $\operatorname{Hom}_S(M,N)$ n'est
pas égal à ce dernier. Toutefois, on a $H=\operatorname{Hom}_S(M,N)$ lorsque
$M$ est \emph{de type fini}~; on peut en effet supposer alors $M$ engendré
par un nombre fini d'éléments homogènes $x_i$ ($1\leq i\leq n$), et tout
homomorphisme $u\in H$ s'écrit d'une seule manière
$\sum_{k\in\mathbf{Z}}u_k$, où pour chaque $k$, $u_k(x_i)$ est égal au
composant homogène de degré $k+\deg(x_i)$ de $u(x_i)$ ($1\leq i\leq n$), ce
qui entraîne $u_k=0$ sauf pour un nombre fini d'indices~; on a par définition
$u_k\in H_k$, d'où la conclusion.

On dit que les éléments de degré $0$ de $\operatorname{Hom}_S(M,N)$ sont des
\emph{homomorphismes de $S$-modules gradués}. Il est clair que
$S_mH_n\subset H_{m+n}$, donc les $H_n$ définissent sur
$\operatorname{Hom}_S(M,N)$ une structure de $S$-module gradué.

Il résulte aussitôt de ces définitions que l'on a
\[
  \label{II.2.1.2.1-fr}
  M(m)\otimes_S N(n)=(M\otimes_S N)(m+n),
  \tag{2.1.2.1}
\]
\[
  \label{II.2.1.2.2-fr}
  \operatorname{Hom}_S(M(m),N(n))
  =(\operatorname{Hom}_S(M,N))(n-m)
  \tag{2.1.2.2}
\]
pour deux $S$-modules gradués $M$, $N$.
\oldpage[II]{21}
Soient $S$, $S'$ deux anneaux gradués~; un homomorphisme d'\emph{anneaux
gradués} $\varphi:S\to S'$ est un homomorphisme d'anneaux tel que
$\varphi(S_n)\subset S'_n$, pour tout $n\in\mathbf{Z}$ (autrement dit,
$\varphi$ doit être un homomorphisme \emph{de degré $0$} de
$\mathbf{Z}$-modules gradués). La donnée d'un tel homomorphisme définit sur
$S'$ une structure de $S$-module \emph{gradué}~; muni de cette structure et
de sa structure d'anneau gradué, on dit que $S'$ est une \emph{$S$-algèbre
graduée}.

Si alors $M$ est un $S$-module gradué, le produit tensoriel $M\otimes_S S'$
de $S$-modules gradués est muni de façon naturelle d'une structure de
$S'$-module \emph{gradué}, la graduation étant celle définie plus haut.
\end{env}

\begin{lemma}[2.1.3]
\label{II.2.1.3-fr}
Soit $S$ un anneau gradué à degrés positifs. Pour qu'une partie $E$ de $S_+$
formée d'éléments homogènes engendre $S_+$ en tant que $S$-module, il faut et
il suffit que $E$ engendre $S$ en tant que $S_0$-algèbre.
\end{lemma}

\begin{proof}
La condition est évidemment suffisante~; montrons qu'elle est nécessaire.
Soit $E_n$ (resp. $E^n$) l'ensemble des éléments de $E$ de degré égal à $n$
(resp. $\leq n$)~; il suffira de montrer, par récurrence sur $n>0$, que $S_n$
est le $S_0$-module engendré par les éléments de degré $n$ qui sont des
produits d'éléments de $E^n$. Or, cela est évident pour $n=1$ en vertu de
l'hypothèse~; cette dernière montre en outre que
$S_n=\sum_{p=0}^{n-1}S_pE_{n-p}$, et le raisonnement par récurrence est alors
immédiat.
\end{proof}

\begin{corollary}[2.1.4]
\label{II.2.1.4-fr}
Pour que $S_+$ soit un idéal de type fini, il faut et il suffit que $S$ soit
une $S_0$-algèbre de type fini.
\end{corollary}

\begin{proof}
On peut en effet toujours supposer qu'un système fini de générateurs de la
$S_0$-algèbre $S$ (resp. du $S$-idéal $S_+$) est formé d'éléments homogènes,
en remplaçant chacun des générateurs considérés par ses composants homogènes.
\end{proof}

\begin{corollary}[2.1.5]
\label{II.2.1.5-fr}
Pour que $S$ soit noethérien, il faut et il suffit que $S_0$ soit noethérien
et que $S$ soit une $S_0$-algèbre de type fini.
\end{corollary}

\begin{proof}
La condition est évidemment suffisante~; elle est nécessaire, puisque $S_0$
est isomorphe à $S/S_+$ et que $S_+$ doit être un idéal de type fini
(\hyperref[II.2.1.4-fr]{2.1.4}).
\end{proof}

\begin{lemma}[2.1.6]
\label{II.2.1.6-fr}
Soit $S$ un anneau gradué à degrés positifs, qui soit une $S_0$-algèbre de
type fini. Soit $M$ un $S$-module gradué de type fini. Alors~:
\begin{enumerate}
  \item[\rm(i)] Les $M_n$ sont des $S_0$-modules de type fini, et il existe
    un entier $n_0$ tel que $M_n=0$ pour $n\leq n_0$.
  \item[\rm(ii)] Il existe un entier $n_1$ et un entier $h>0$ tel que, pour
    tout entier $n\geq n_1$, on ait $M_{n+h}=S_hM_n$.
  \item[\rm(iii)] Pour tout couple d'entiers $(d,k)$ tels que $d>0$,
    $0\leq k\leq d-1$, $M^{(d,k)}$ est un $S^{(d)}$-module de type fini.
  \item[\rm(iv)] Pour tout entier $d>0$, $S^{(d)}$ est une $S_0$-algèbre de
    type fini.
  \item[\rm(v)] Il existe un entier $h>0$ tel que $S_{mh}=(S_h)^m$ pour tout
    $m>0$.
  \item[\rm(vi)] Pour tout entier $n>0$, il existe un entier $m_0$ tel que
    $S_m\subset S_+^n$ pour tout $m\geq m_0$.
\end{enumerate}
\end{lemma}

\begin{proof}
On peut supposer $S$ engendré (en tant que $S_0$-algèbre) par des éléments
homogènes $f_i$, de degré $h_i$ ($1\leq i\leq r$), et $M$ engendré (en tant
que $S$-module) par des éléments homogènes $x_j$ de degré $k_j$
($1\leq j\leq s$). Il est clair que $M_n$ est formé des combinaisons
\oldpage[II]{22}
linéaires, à coefficients dans $S_0$, des éléments
$f_1^{\alpha_1}\cdots f_r^{\alpha_r}x_j$ tels que les $\alpha_i$ soient des
entiers $\geq 0$ vérifiant $k_j+\sum_i\alpha_i h_i=n$~; pour chaque $j$, il
n'y a qu'un nombre fini de systèmes $(\alpha_i)$ vérifiant cette équation,
puisque les $h_i$ sont $>0$, d'où la première assertion de (i)~; la seconde
est évidente. Soit d'autre part $h$ le p.p.c.m. des $h_i$ et posons
$g_i=f_i^{h/h_i}$ ($1\leq i\leq r$) de sorte que tous les $g_i$ sont de
degré $h$~; soient $z_\mu$ les éléments de $M$ de la forme
$f_1^{\alpha_1}\cdots f_r^{\alpha_r}x_j$ avec
$0\leq\alpha_i<h/h_i$ pour $1\leq i\leq r$~; ces éléments sont en nombre
fini, soit $n_1$ le plus grand de leurs degrés. Il est clair que pour
$n\geq n_1$, tout élément de $M_{n+h}$ est combinaison linéaire des $z_\mu$
dont les coefficients sont des monômes de degré $>0$ par rapport aux $g_i$,
donc on a $M_{n+h}=S_hM_n$, ce qui établit (ii). De la même manière, on voit
(pour tout $d>0$) qu'un élément de $M^{(d,k)}$ est combinaison linéaire, à
coefficients dans $S_0$, d'éléments de la forme
$g^df_1^{\alpha_1}\cdots f_r^{\alpha_r}x_j$ avec
$0\leq\alpha_i<d$, $g$ étant un élément homogène de $S$~; d'où (iii)~; (iv)
résulte alors de (iii) et du lemme (\hyperref[II.2.1.3-fr]{2.1.3}), en
prenant $M=S_+$, puisque $(S_+)^{(d)}=(S^{(d)})_+$. L'assertion de (v) se
déduit de (ii) en prenant $M=S$. Enfin, pour un $n$ donné, il n'y a qu'un
nombre fini de systèmes $(\alpha_i)$ tels que $\alpha_i\geq 0$ et
$\sum_i\alpha_i<n$, donc si $m_0$ est la plus grande valeur de la somme
$\sum_i\alpha_i h_i$ pour ces systèmes, on a $S_m\subset S_+^n$ pour
$m>m_0$, ce qui prouve (vi).
\end{proof}

\begin{corollary}[2.1.7]
\label{II.2.1.7-fr}
Si $S$ est noethérien, il en est de même de $S^{(d)}$ pour tout entier
$d>0$.
\end{corollary}

\begin{proof}
Cela résulte de (\hyperref[II.2.1.5-fr]{2.1.5}) et de
(\hyperref[II.2.1.6-fr]{2.1.6, (iv)}).
\end{proof}

\begin{env}[2.1.8]
\label{II.2.1.8-fr}
Soit $\mathfrak{p}$ un idéal premier gradué de l'anneau gradué $S$~;
$\mathfrak{p}$ est donc somme directe des sous-groupes
$\mathfrak{p}_n=\mathfrak{p}\cap S_n$. Supposons que $\mathfrak{p}$ ne
contienne pas $S_+$. Alors, si $f\in S_+$, n'appartient pas à
$\mathfrak{p}$, la relation $f^nx\in\mathfrak{p}$ est équivalente à
$x\in\mathfrak{p}$~; en particulier, si $f\in S_d$ ($d>0$), pour tout
$x\in S_{m-nd}$, la relation $f^nx\in\mathfrak{p}_m$ équivaut à
$x\in\mathfrak{p}_{m-nd}$.
\end{env}

\begin{proposition}[2.1.9]
\label{II.2.1.9-fr}
Soit $n_0$ un entier $>0$~; pour tout $n\geq n_0$ soit $\mathfrak{p}_n$ un
sous-groupe de $S_n$. Pour qu'il existe un idéal premier gradué
$\mathfrak{p}$ de $S$ ne contenant pas $S_+$ et tel que
$\mathfrak{p}\cap S_n=\mathfrak{p}_n$ pour tout $n\geq n_0$, il faut et il
suffit que les conditions suivantes soient vérifiées~:
\begin{enumerate}
  \item[$1^\circ$] $S_m\mathfrak{p}_n\subset\mathfrak{p}_{m+n}$ pour tout
    $m\geq 0$ et tout $n\geq n_0$.
  \item[$2^\circ$] Pour $m\geq n_0$, $n\geq n_0$, $f\in S_m$, $g\in S_n$,
    la relation $fg\in\mathfrak{p}_{m+n}$ entraîne
    $f\in\mathfrak{p}_m$ ou $g\in\mathfrak{p}_n$.
  \item[$3^\circ$] $\mathfrak{p}_n\neq S_n$ pour un $n\geq n_0$ au moins.
\end{enumerate}
En outre l'idéal premier gradué $\mathfrak{p}$ est alors unique.
\end{proposition}

\begin{proof}
Il est évident que les conditions $1^\circ$ et $2^\circ$ sont nécessaires.
En outre, si $\mathfrak{p}\not\supset S_+$, il existe au moins un $k>0$ tel que
$\mathfrak{p}\cap S_k\neq S_k$~; si $f\in S_k$ n'appartient pas à
$\mathfrak{p}$, la relation $\mathfrak{p}\cap S_n=S_n$ entraîne
$\mathfrak{p}\cap S_{n-mk}=S_{n-mk}$ d'après
(\hyperref[II.2.1.8-fr]{2.1.8})~; donc, si
$\mathfrak{p}\cap S_n=S_n$ à partir d'une certaine valeur de $n$, on aurait
$\mathfrak{p}\supset S_+$ contrairement à l'hypothèse, ce qui prouve que
$3^\circ$ est nécessaire. Inversement, supposons satisfaites les conditions
$1^\circ$, $2^\circ$ et $3^\circ$. Notons que, si pour un entier
$d\geq n_0$, $f\in S_d$ n'appartient pas à $\mathfrak{p}_d$, alors, si
$\mathfrak{p}$ existe, $\mathfrak{p}_m$, pour $m<n_0$, est nécessairement
égal à l'ensemble des $x\in S_m$ tels que
$f^rx\in\mathfrak{p}_{m+rd}$, sauf pour un nombre fini de valeurs de $r$.
Cela prouve déjà que si $\mathfrak{p}$ existe, il est unique. Reste à montrer
que si on définit les $\mathfrak{p}_m$ pour $m<n_0$ par la condition
précédente, $\mathfrak{p}=\sum_{n=0}^{\infty}\mathfrak{p}_n$ est un idéal
premier. Notons d'abord qu'en vertu de $2^\circ$, pour $m\geq n_0$,
$\mathfrak{p}_m$ est aussi défini comme l'ensemble des $x\in S_m$ tels que
$f^rx\in\mathfrak{p}_{m+rd}$, sauf pour un nombre fini de valeurs de $r$.
Cela
\oldpage[II]{23}
étant, si $g\in S_m$, $x\in\mathfrak{p}_n$, on a
$f^rgx\in\mathfrak{p}_{m+n+rd}$ sauf pour un nombre fini de valeurs de $r$,
donc $gx\in\mathfrak{p}_{m+n}$, ce qui prouve que $\mathfrak{p}$ est un
idéal de $S$. Pour établir que cet idéal est premier, autrement dit que
l'anneau $S/\mathfrak{p}$, gradué par les sous-groupes
$S_n/\mathfrak{p}_n$, est un anneau intègre, il suffit (en considérant les
composants de plus haut degré de deux éléments de $S/\mathfrak{p}$) de
prouver que si $x\in S_m$, $y\in S_n$ sont tels que
$x\notin\mathfrak{p}_m$, $y\notin\mathfrak{p}_n$, alors
$xy\notin\mathfrak{p}_{m+n}$. Sinon, pour $r$ assez grand, on aurait
$f^{2r}xy\in\mathfrak{p}_{m+n+2rd}$~; mais on a
$f^ry\notin\mathfrak{p}_{n+rd}$ pour tout $r>0$~; il résulte alors de
$2^\circ$ que, sauf pour un nombre fini de valeurs de $r$, on a
$f^rx\in\mathfrak{p}_{m+rd}$, et on en conclurait que
$x\in\mathfrak{p}_m$ contrairement à l'hypothèse.
\end{proof}

\begin{env}[2.1.10]
\label{II.2.1.10-fr}
Nous dirons qu'une partie $\mathfrak{J}$ de $S_+$ est un
\emph{idéal de $S_+$} si c'est un idéal de $S$, et que $\mathfrak{J}$ est un
\emph{idéal premier gradué de $S_+$} si elle est l'intersection de $S_+$ et
d'un idéal premier gradué de $S$ \emph{ne contenant pas $S_+$} (cet idéal
premier est d'ailleurs unique d'après
(\hyperref[II.2.1.9-fr]{2.1.9})). Si $\mathfrak{J}$ est un idéal de $S_+$,
la \emph{racine de $\mathfrak{J}$ dans $S_+$} est l'ensemble des éléments de
$S_+$ dont une puissance appartient à $\mathfrak{J}$, autrement dit c'est
l'ensemble $r_+(\mathfrak{J})=r(\mathfrak{J})\cap S_+$~; en particulier, la
racine de $0$ dans $S_+$ est encore appelée le \emph{nilradical} de $S_+$ et
notée $\mathfrak{N}_+$~: c'est donc l'ensemble des éléments nilpotents de
$S_+$. Si $\mathfrak{J}$ est un idéal \emph{gradué} de $S_+$, sa racine
$r_+(\mathfrak{J})$ est un idéal gradué~: en passant à l'anneau quotient
$S/\mathfrak{J}$, on peut se ramener au cas $\mathfrak{J}=0$, et tout
revient à voir que si $x=x_h+x_{h+1}+\cdots+x_k$ est nilpotent, il en est de
même des $x_i\in S_i$ ($1\leq h\leq i\leq k$)~; on peut supposer
$x_k\neq 0$ et le composant de plus grand degré de $x^n$ est alors $x_k^n$,
donc $x_k$ est nilpotent, et on raisonne alors par récurrence sur $k$. On dit
que l'anneau gradué $S$ est \emph{essentiellement réduit} si
$\mathfrak{N}_+=0$, autrement dit si $S_+$ ne contient pas d'éléments
nilpotents $\neq 0$.
\end{env}

\begin{env}[2.1.11]
\label{II.2.1.11-fr}
On notera que si, dans l'anneau gradué $S$, un élément $x$ est diviseur de
$0$, il en est de même de son composant de plus haut degré. On dit qu'un
anneau $S$ est \emph{essentiellement intègre} si l'anneau $S_+$
(\emph{sans élément unité}) ne contient pas de diviseur de $0$ et est
$\neq 0$~; il suffit pour cela qu'un élément homogène $\neq 0$ dans $S_+$ ne
soit pas diviseur de $0$ dans cet anneau. Il est clair que si $\mathfrak{p}$
est un idéal premier gradué de $S_+$, $S/\mathfrak{p}$ est essentiellement
intègre.

Soit $S$ un anneau gradué essentiellement intègre, et soit $x_0\in S_0$~;
s'il existe alors \emph{un} élément homogène $f\neq 0$ de $S_+$ tel que
$x_0f=0$, on a $x_0S_+=0$, car on a $(x_0g)f=(x_0f)g=0$ pour tout
$g\in S_+$, et l'hypothèse entraîne donc $x_0g=0$. Pour que $S$ soit
intègre, il faut et il suffit donc que $S_0$ soit intègre et que l'annulateur
de $S_+$ dans $S_0$ soit réduit à $0$.
\end{env}

\subsection{Anneaux de fractions d'un anneau gradué.}
\label{subsection:II.2.2-fr}

\begin{env}[2.2.1]
\label{II.2.2.1-fr}
Soient $S$ un anneau gradué, à degrés positifs, $f$ un élément
\emph{homogène} de $S$, de degré $d>0$~; alors l'anneau de fractions
$S'=S_f$ est gradué, en prenant pour $S'_n$ l'ensemble des $x/f^k$, où
$x\in S_{n+kd}$ avec $k\geq 0$ (on observera ici que $n$ peut prendre des
valeurs négatives arbitraires)~; nous désignerons le sous-anneau
$S'_0=(S_f)_0$ de $S'$ formé des éléments \emph{de degré $0$} par la
notation $S_{(f)}$.

Si $f\in S_d$, les monômes $(f/1)^h$ dans $S_f$ ($h$ entier positif ou
négatif) forment un \emph{système libre} sur l'anneau $S_{(f)}$, et
l'ensemble de leurs combinaisons linéaires n'est autre que
\oldpage[II]{24}
l'anneau $(S^{(d)})_f$, qui est donc \emph{isomorphe à
$S_{(f)}[T,T^{-1}]=S_{(f)}\otimes_{\mathbf{Z}}\mathbf{Z}[T,T^{-1}]$}
(où $T$ est une indéterminée). En effet, si on a une relation
$\sum_{h=-a}^{b}z_h(f/1)^h=0$ avec $z_h=x_h/f^m$, où tous les $x_h$
appartiennent à $S_{md}$, cette relation équivaut par définition à
l'existence d'un $k\geq a$ tel que
$\sum_{h=-a}^{b}f^{h+k}x_h=0$, et comme les degrés des termes de cette somme
sont distincts, on a $f^{h+k}x_h=0$ pour tout $h$, d'où $z_h=0$ pour tout
$h$.

Si $M$ est un $S$-module gradué, $M'=M_f$ est un $S_f$-module gradué,
$M'_n$ étant l'ensemble des $z/f^k$ avec $z\in M_{n+kd}$ ($k\geq 0$)~; on
désigne par $M_{(f)}$ l'ensemble des éléments homogènes de degré $0$ de
$M'$~; il est immédiat que $M_{(f)}$ est un $S_{(f)}$-module et que l'on a
$(M^{(d)})_f=M_{(f)}\otimes_{S_{(f)}}(S^{(d)})_f$.
\end{env}

\begin{lemma}[2.2.2]
\label{II.2.2.2-fr}
Soient $d$, $e$ deux entiers $>0$, $f\in S_d$, $g\in S_e$. Il existe un
isomorphisme canonique d'anneaux
\[
  S_{(fg)}\xrightarrow{\sim}(S_{(f)})_{g^d/f^e}~;
\]
si on identifie canoniquement ces deux anneaux, il existe un isomorphisme
canonique de modules
\[
  M_{(fg)}\xrightarrow{\sim}(M_{(f)})_{g^d/f^e}.
\]
\end{lemma}

\begin{proof}
En effet, $fg$ divise $f^e g^d$, et ce dernier élément divise $(fg)^{de}$,
donc les anneaux gradués $S_{fg}$ et $S_{f^e g^d}$ s'identifient
canoniquement~; d'autre part, $S_{f^e g^d}$ s'identifie aussi à
$(S_{f^e})_{g^d/1}$ (\hyperref[0.1.4.6-fr]{0, I.4.6}), et comme $f^e/1$ est
inversible dans $S_{f^e}$, $S_{f^e g^d}$ s'identifie aussi à
$(S_{f^e})_{g^d/f^e}$. Or l'élément $g^d/f^e$ est de degré $0$ dans
$S_{f^e}$~; on en conclut aussitôt que le sous-anneau de
$(S_{f^e})_{g^d/f^e}$ formé des éléments de degré $0$ est
$(S_{(f^e)})_{g^d/f^e}$, et comme on a évidemment
$S_{(f^e)}=S_{(f)}$, cela démontre la première partie de la proposition~;
la seconde s'établit de même.
\end{proof}

\begin{env}[2.2.3]
\label{II.2.2.3-fr}
Avec les hypothèses de (\hyperref[II.2.2.2-fr]{2.2.2}), il est clair que
l'homomorphisme canonique $S_f\to S_{fg}$
(\hyperref[0.1.4.1-fr]{0, I.4.1}), qui transforme $x/f^k$ en
$g^k x/(fg)^k$ est de degré $0$, donc donne par restriction un
\emph{homomorphisme canonique $S_{(f)}\to S_{(fg)}$}, tel que le diagramme
\[
  \xymatrix{
    & S_{(f)}\ar[dl]\ar[dr]\\
    S_{(fg)}\ar[rr]^-{\sim} && (S_{(f)})_{g^d/f^e}
  }
\]
soit commutatif. On définit de même un homomorphisme canonique
$M_{(f)}\to M_{(fg)}$.
\end{env}

\begin{lemma}[2.2.4]
\label{II.2.2.4-fr}
Si $f$, $g$ sont deux éléments homogènes de $S_+$, l'anneau $S_{(fg)}$ est
engendré par la réunion des images canoniques de $S_{(f)}$ et $S_{(g)}$.
\end{lemma}

\begin{proof}
En vertu de (\hyperref[II.2.2.2-fr]{2.2.2}), il suffit de voir que
$1/(g^d/f^e)=f^{d+e}/(fg)^d$ appartient à l'image canonique de $S_{(g)}$
dans $S_{(fg)}$, ce qui est évident par définition.
\end{proof}

\begin{proposition}[2.2.5]
\label{II.2.2.5-fr}
Soit $d$ un entier $>0$ et soit $f\in S_d$. Alors il existe un isomorphisme
canonique d'anneaux
$S_{(f)}\xrightarrow{\sim}S^{(d)}/(f-1)S^{(d)}$~; si on identifie ces deux
anneaux par cet isomorphisme, il existe un isomorphisme canonique de modules
$M_{(f)}\xrightarrow{\sim}M^{(d)}/(f-1)M^{(d)}$.
\end{proposition}

\begin{proof}
Le premier de ces isomorphismes se définit en faisant correspondre à
$x/f^n$, où $x\in S_{nd}$, l'élément $\overline{x}$, classe de $x$ mod.
$(f-1)S^{(d)}$~; cette application est bien définie, car on a la congruence
$f^h x\equiv x\pmod{(f-1)S^{(d)}}$ pour tout $x\in S^{(d)}$, donc si
$f^h x=0$ pour un $h>0$,
\oldpage[II]{25}
on a $\overline{x}=0$. D'autre part, si $x\in S_{nd}$ est tel que
$x=(f-1)y$ avec
$y=y_{hd}+y_{(h+1)d}+\cdots+y_{kd}$ avec $y_{jd}\in S_{jd}$ et
$y_{hd}\neq 0$, on a nécessairement $h=n$ et $x=-y_{hd}$, ainsi que les
relations $y_{(j+1)d}=fy_{jd}$ pour $h\leq j\leq k-1$, $fy_{kd}=0$, ce qui
donne finalement $f^{k-n+1}x=0$~; à toute classe $\overline{x}$ mod.
$(f-1)S^{(d)}$ d'un élément $x\in S_{nd}$, on peut donc faire correspondre
l'élément $x/f^n$ de $S_{(f)}$, car la remarque qui précède montre que cette
application est bien définie. Il est immédiat que les deux applications ainsi
définies sont des homomorphismes d'anneaux, réciproques l'un de l'autre. On
procède exactement de même pour $M$.
\end{proof}

\begin{corollary}[2.2.6]
\label{II.2.2.6-fr}
Si $S$ est noethérien, il en est de même de $S_{(f)}$ pour tout $f$ homogène
de degré $>0$.
\end{corollary}

\begin{proof}
Cela résulte aussitôt de (\hyperref[II.2.1.7-fr]{2.1.7}) et
(\hyperref[II.2.2.5-fr]{2.2.5}).
\end{proof}

\begin{env}[2.2.7]
\label{II.2.2.7-fr}
Soit $T$ une partie multiplicative de $S_+$ formée d'éléments
\emph{homogènes}~; $T_0=T\cup\{1\}$ est alors une partie multiplicative de
$S$~; comme les éléments de $T_0$ sont homogènes, l'anneau $T_0^{-1}S$ est
encore gradué de façon évidente~; on désignera par $S_{(T)}$ le sous-anneau
de $T_0^{-1}S$ formé des éléments d'ordre $0$, c'est-à-dire des éléments de
la forme $x/h$, où $h\in T$ et $x$ est homogène de degré égal à celui de
$h$. On sait (\hyperref[0.1.4.5-fr]{0, 1.4.5}) que $T_0^{-1}S$ s'identifie
canoniquement à la limite inductive des anneaux $S_f$, où $f$ parcourt $T$
(pour les homomorphismes canoniques $S_f\to S_{fg}$)~; comme cette
identification respecte les degrés, elle identifie $S_{(T)}$ à la
\emph{limite inductive} des $S_{(f)}$ pour $f\in T$. Pour tout $S$-module
gradué $M$, on définit de même le module $M_{(T)}$ (sur l'anneau $S_{(T)}$)
formé des éléments de degré $0$ de $T_0^{-1}M$, et on voit que ce module est
limite inductive des $M_{(f)}$ pour $f\in T$.

Si $\mathfrak{p}$ est un idéal premier gradué de $S_+$, on désignera par
$S_{(\mathfrak{p})}$ et $M_{(\mathfrak{p})}$ l'anneau $S_{(T)}$ et le
module $M_{(T)}$ respectivement, où $T$ est l'ensemble des éléments
\emph{homogènes} de $S_+$ qui n'appartiennent pas à $\mathfrak{p}$.
\end{env}

\subsection{Spectre premier homogène d'un anneau gradué.}
\label{subsection:II.2.3-fr}

\begin{env}[2.3.1]
\label{II.2.3.1-fr}
Étant donné un anneau gradué $S$, à degrés positifs, on appelle
\emph{spectre premier homogène} de $S$ et on désigne par
$\operatorname{Proj}(S)$ l'ensemble des idéaux premiers gradués de $S_+$
(\hyperref[II.2.1.10-fr]{2.1.10}), ou, ce qui revient au même, l'ensemble
des idéaux premiers gradués de $S$ \emph{ne contenant pas} $S_+$~; nous
allons définir une structure de \emph{schéma} ayant
$\operatorname{Proj}(S)$ comme ensemble de base.
\end{env}

\begin{env}[2.3.2]
\label{II.2.3.2-fr}
Pour toute partie $E$ de $S$, soit $V_+(E)$ l'ensemble des idéaux premiers
gradués de $S$ contenant $E$ et ne contenant pas $S_+$~; c'est donc la
partie $V(E)\cap\operatorname{Proj}(S)$ de
$\operatorname{Spec}(S)$. De (\hyperref[I.1.1.2-fr]{I, 1.1.2}) on déduit
donc~:
\[
  \label{II.2.3.2.1-fr}
  V_+(0)=\operatorname{Proj}(S),\quad
  V_+(S)=V_+(S_+)=\varnothing
  \tag{2.3.2.1}
\]
\[
  \label{II.2.3.2.2-fr}
  V_+\left(\bigcup_\lambda E_\lambda\right)
  =\bigcap_\lambda V_+(E_\lambda)
  \tag{2.3.2.2}
\]
\[
  \label{II.2.3.2.3-fr}
  V_+(EE')=V_+(E)\cup V_+(E')
  \tag{2.3.2.3}
\]
On ne change pas $V_+(E)$ en remplaçant $E$ par l'idéal gradué engendré par
$E$~; en outre, si $\mathfrak{J}$ est un idéal gradué de $S$, on a
\[
  \label{II.2.3.2.4-fr}
  V_+(\mathfrak{J})
  =V_+\left(\bigcup_{q\geq n}(\mathfrak{J}\cap S_q)\right)
  \tag{2.3.2.4}
\]
\oldpage[II]{26}
pour tout $n>0$~: en effet, si
$\mathfrak{p}\in\operatorname{Proj}(S)$ contient les éléments homogènes de
$\mathfrak{J}$ de degré $\geq n$, comme par hypothèse il existe un élément
homogène $f\in S_d$ non contenu dans $\mathfrak{p}$, pour tout $m\geq 0$ et
tout $x\in S_m\cap\mathfrak{J}$, on a
$f^r x\in\mathfrak{J}\cap S_{m+rd}$ sauf pour un nombre fini de valeurs de
$r$, donc $f^r x\in\mathfrak{p}\cap S_{m+rd}$, ce qui entraîne
$x\in\mathfrak{p}\cap S_m$
(\hyperref[II.2.1.9-fr]{2.1.9}).

Enfin, on a, pour tout idéal gradué $\mathfrak{J}$ de $S$
\[
  \label{II.2.3.2.5-fr}
  V_+(\mathfrak{J})=V_+(r_+(\mathfrak{J})).
  \tag{2.3.2.5}
\]
\end{env}

\begin{env}[2.3.3]
\label{II.2.3.3-fr}
Par définition, les $V_+(E)$ sont les parties fermées de
$X=\operatorname{Proj}(S)$ pour la topologie induite par la topologie
spectrale de $\operatorname{Spec}(S)$, et qu'on appellera encore
\emph{topologie spectrale} sur $X$. On pose, pour tout $f\in S$
\[
  \label{II.2.3.3.1-fr}
  D_+(f)=D(f)\cap\operatorname{Proj}(S)
  =\operatorname{Proj}(S)\setminus V_+(f)
  \tag{2.3.3.1}
\]
et on a par suite, pour deux éléments $f$, $g$ de $S$
(\hyperref[I.1.1.9.1-fr]{I, 1.1.9.1})
\[
  \label{II.2.3.3.2-fr}
  D_+(fg)=D_+(f)\cap D_+(g).
  \tag{2.3.3.2}
\]
\end{env}

\begin{proposition}[2.3.4]
\label{II.2.3.4-fr}
Lorsque $f$ parcourt l'ensemble des éléments homogènes de $S_+$, les
$D_+(f)$ forment une base de la topologie de
$X=\operatorname{Proj}(S)$.
\end{proposition}

\begin{proof}
Il résulte en effet de (\hyperref[II.2.3.2.2-fr]{2.3.2.2}) et
(\hyperref[II.2.3.2.4-fr]{2.3.2.4}) que toute partie fermée de $X$ est
intersection d'ensembles de la forme $V_+(f)$, où $f$ est homogène de degré
$>0$.
\end{proof}

\begin{env}[2.3.5]
\label{II.2.3.5-fr}
Soit $f$ un élément \emph{homogène} de $S_+$, de degré $d>0$~; pour tout
idéal premier gradué $\mathfrak{p}$ de $S$, ne contenant pas $f$, on sait
que l'ensemble des $x/f^n$, où $x\in\mathfrak{p}$ et $n\geq 0$, est un
idéal premier de l'anneau de fractions $S_f$
(\hyperref[0.1.2.6-fr]{0, 1.2.6})~; sa trace sur $S_{(f)}$ est donc un idéal
premier de cet anneau, que nous désignerons par
$\psi_f(\mathfrak{p})$~: c'est l'ensemble des $x/f^n$ pour $n\geq 0$,
$x\in\mathfrak{p}\cap S_{nd}$. On a ainsi défini une application
\[
  \psi_f:D_+(f)\longrightarrow\operatorname{Spec}(S_{(f)})~;
\]
en outre, si $g\in S_e$ est un second élément homogène de $S_+$, on a un
diagramme commutatif
\[
  \label{II.2.3.5.1-fr}
  \xymatrix{
    D_+(f)\ar[r]^{\psi_f}
    & \operatorname{Spec}(S_{(f)})\\
    D_+(fg)\ar[u]\ar[r]_{\psi_{fg}}
    & \operatorname{Spec}(S_{(fg)})\ar[u]
  }
  \tag{2.3.5.1}
\]
où la flèche verticale de gauche est l'inclusion, et celle de droite est
l'application ${}^a\omega_{fg,f}$ déduite de l'homomorphisme canonique
$\omega=\omega_{fg,f}:S_{(f)}\to S_{(fg)}$
(\hyperref[I.1.2.1-fr]{I, 1.2.1}). En effet, si
$x/f^n\in\omega^{-1}(\psi_{fg}(\mathfrak{p}))$, où
$fg\notin\mathfrak{p}$, on a par définition
$g^n x/(fg)^n\in\psi_{fg}(\mathfrak{p})$, donc
$g^n x\in\mathfrak{p}$ et par suite $x\in\mathfrak{p}$, et la réciproque
est évidente.
\end{env}

\begin{proposition}[2.3.6]
\label{II.2.3.6-fr}
L'application $\psi_f$ est un homéomorphisme de $D_+(f)$ sur
$\operatorname{Spec}(S_{(f)})$.
\end{proposition}

\begin{proof}
En premier lieu, $\psi_f$ est continue~; car si $h\in S_{nd}$ est tel que
$h/f^n\in\psi_f(\mathfrak{p})$, on a par définition
$h\in\mathfrak{p}$ et réciproquement, donc
\[
  \psi_f^{-1}(D(h/f^n))=D_+(hf)
\]
et notre assertion résulte de la formule
(\hyperref[II.2.3.3.2-fr]{2.3.3.2}). En outre, les $D_+(hf)$, où $h$
parcourt les ensembles $S_{nd}$, forment une base de la topologie de
$D_+(f)$, en vertu de (\hyperref[II.2.3.4-fr]{2.3.4}) et de la formule
(\hyperref[II.2.3.3.2-fr]{2.3.3.2})~; ce
\oldpage[II]{27}
qui précède prouve donc, compte tenu de l'axiome $(T_0)$ valable dans
$D_+(f)$ et dans $\operatorname{Spec}(S_{(f)})$, que $\psi_f$ est injective
et que l'application réciproque
$\psi_f(D_+(f))\longrightarrow D_+(f)$ est continue. Enfin, pour voir que
$\psi_f$ est surjective, on remarque que, si $\mathfrak{q}_0$ est un idéal
premier de $S_{(f)}$, et si, pour tout $n>0$, on désigne par
$\mathfrak{p}_n$ l'ensemble des $x\in S_n$ tels que
$x^d/f^n\in\mathfrak{q}_0$, les $\mathfrak{p}_n$ vérifient les conditions
de (\hyperref[II.2.1.9-fr]{2.1.9})~: en effet, si $x\in S_n$, $y\in S_n$
sont tels que $x^d/f^n\in\mathfrak{q}_0$ et
$y^d/f^n\in\mathfrak{q}_0$, on a
$(x+y)^{2d}/f^{2n}\in\mathfrak{q}_0$, d'où
$(x+y)^d/f^n\in\mathfrak{q}_0$ puisque $\mathfrak{q}_0$ est premier~;
cela prouve que les $\mathfrak{p}_n$ sont des sous-groupes des $S_n$, et la
vérification des autres conditions de
(\hyperref[II.2.1.9-fr]{2.1.9}) est immédiate, compte tenu de ce que
$\mathfrak{q}_0$ est premier. Si $\mathfrak{p}$ est l'idéal premier gradué
de $S$ ainsi défini, on a bien
$\psi_f(\mathfrak{p})=\mathfrak{q}_0$, car si $x\in S_{nd}$, les relations
$x/f^n\in\mathfrak{q}_0$ et $x^d/f^{nd}\in\mathfrak{q}_0$ sont
équivalentes, $\mathfrak{q}_0$ étant premier.
\end{proof}

\begin{corollary}[2.3.7]
\label{II.2.3.7-fr}
Pour que $D_+(f)=\varnothing$, il faut et il suffit que $f$ soit nilpotent.
\end{corollary}

\begin{proof}
En effet, pour que $\operatorname{Spec}(S_{(f)})=\varnothing$, il faut et
il suffit que $S_{(f)}=0$, ou encore que $1=0$ dans $S_f$, ce qui signifie
par définition que $f$ est nilpotent.
\end{proof}

\begin{corollary}[2.3.8]
\label{II.2.3.8-fr}
Soit $E$ une partie de $S_+$. Les conditions suivantes sont équivalentes~:
\begin{enumerate}
  \item[\rm a)] $V_+(E)=X=\operatorname{Proj}(S)$.
  \item[\rm b)] Tout élément de $E$ est nilpotent.
  \item[\rm c)] Les composants homogènes de tout élément de $E$ sont
    nilpotents.
\end{enumerate}
\end{corollary}

\begin{proof}
Il est clair que c) entraîne b) et que b) entraîne a). Si $\mathfrak{J}$ est
l'idéal gradué de $S$ engendré par $E$, la condition a) équivaut à
$V_+(\mathfrak{J})=X$~; \emph{a fortiori}, a) entraîne que tout élément
homogène $f\in\mathfrak{J}$ est tel que $V_+(f)=X$, donc $f$ est nilpotent
en vertu de (\hyperref[II.2.3.7-fr]{2.3.7}).
\end{proof}

\begin{corollary}[2.3.9]
\label{II.2.3.9-fr}
Si $\mathfrak{J}$ est un idéal gradué de $S_+$,
$r_+(\mathfrak{J})$ est l'intersection des idéaux premiers gradués de
$S_+$ qui contiennent $\mathfrak{J}$.
\end{corollary}

\begin{proof}
En considérant l'anneau gradué $S/\mathfrak{J}$, on peut se ramener au cas
où $\mathfrak{J}=0$. Il faut prouver que si $f\in S_+$ n'est pas
nilpotent, il y a un idéal premier gradué de $S$ ne contenant pas $f$~; or,
un au moins des composants homogènes de $f$ n'est pas nilpotent, et on peut
donc supposer $f$ homogène~; l'assertion résulte alors de
(\hyperref[II.2.3.7-fr]{2.3.7}).
\end{proof}

\begin{env}[2.3.10]
\label{II.2.3.10-fr}
Pour toute partie $Y$ de $X=\operatorname{Proj}(S)$, soit
$\mathfrak{j}_+(Y)$ l'ensemble des $f\in S_+$ tels que
$Y\subset V_+(f)$~; il revient au même de dire que
$\mathfrak{j}_+(Y)=\mathfrak{j}(Y)\cap S_+$~; $\mathfrak{j}_+(Y)$ est
donc un idéal de $S_+$, égal à sa racine dans $S_+$.
\end{env}

\begin{proposition}[2.3.11]
\label{II.2.3.11-fr}
\begin{enumerate}
  \item[\rm(i)] Pour toute partie $E$ de $S_+$,
    $\mathfrak{j}_+(V_+(E))$ est la racine dans $S_+$ de l'idéal gradué de
    $S_+$ engendré par $E$.
  \item[\rm(ii)] Pour toute partie $Y$ de $X$,
    $V_+(\mathfrak{j}_+(Y))=\overline{Y}$, adhérence de $Y$ dans $X$.
\end{enumerate}
\end{proposition}

\begin{proof}
\begin{enumerate}
  \item[\rm(i)] Si $\mathfrak{J}$ est l'idéal gradué de $S_+$ engendré
    par $E$, on a $V_+(E)=V_+(\mathfrak{J})$ et l'assertion résulte alors de
    (\hyperref[II.2.3.9-fr]{2.3.9}).
  \item[\rm(ii)] Comme
    $V_+(\mathfrak{J})=\bigcap_{f\in\mathfrak{J}}V_+(f)$, la relation
    $Y\subset V_+(\mathfrak{J})$ entraîne $Y\subset V_+(f)$ pour tout
    $f\in\mathfrak{J}$, et par suite
    $\mathfrak{j}_+(Y)\supset\mathfrak{J}$, d'où
    $V_+(\mathfrak{j}_+(Y))\subset V_+(\mathfrak{J})$, ce qui prouve (ii)
    par définition des ensembles fermés.
\end{enumerate}
\end{proof}

\begin{corollary}[2.3.12]
\label{II.2.3.12-fr}
Les parties fermées $Y$ de $X=\operatorname{Proj}(S)$ et les idéaux gradués
de $S_+$ égaux à leur racine dans $S_+$ se correspondent biunivoquement par
les applications décroissantes
$Y\longmapsto\mathfrak{j}_+(Y)$, $\mathfrak{J}\longmapsto
V_+(\mathfrak{J})$~; à la réunion $Y_1\cup Y_2$ de deux parties fermées de
$X$ correspond
\oldpage[II]{28}
$\mathfrak{j}_+(Y_1)\cap\mathfrak{j}_+(Y_2)$, et à l'intersection d'une
famille quelconque $(Y_\lambda)$ de parties fermées correspond la racine
dans $S_+$ de la somme des $\mathfrak{j}_+(Y_\lambda)$.
\end{corollary}

\begin{corollary}[2.3.13]
\label{II.2.3.13-fr}
Soit $\mathfrak{J}$ un idéal gradué dans $S_+$~; pour que
$V_+(\mathfrak{J})=\varnothing$, il faut et il suffit que tout élément de
$S_+$ ait une puissance dans $\mathfrak{J}$.
\end{corollary}

Ce dernier corollaire s'exprime aussi sous l'une des formes équivalentes~:

\begin{corollary}[2.3.14]
\label{II.2.3.14-fr}
Soit $(f_\alpha)$ une famille d'éléments homogènes de $S_+$. Pour que les
$D_+(f_\alpha)$ forment un recouvrement de
$X=\operatorname{Proj}(S)$, il faut et il suffit que tout élément de $S_+$
ait une puissance dans l'idéal engendré par les $f_\alpha$.
\end{corollary}

\begin{corollary}[2.3.15]
\label{II.2.3.15-fr}
Soient $(f_\alpha)$ une famille d'éléments homogènes de $S_+$, $f$ un
élément de $S_+$. Les relations suivantes sont équivalentes~:
\begin{enumerate}
  \item[\rm a)] $D_+(f)\subset\bigcup_\alpha D_+(f_\alpha)$~;
  \item[\rm b)] $V_+(f)\supset\bigcap_\alpha V_+(f_\alpha)$~;
  \item[\rm c)] une puissance de $f$ appartient à l'idéal engendré par les
    $f_\alpha$.
\end{enumerate}
\end{corollary}

\begin{corollary}[2.3.16]
\label{II.2.3.16-fr}
Pour que $X=\operatorname{Proj}(S)$ soit vide, il faut et il suffit que tout
élément de $S_+$ soit nilpotent.
\end{corollary}

\begin{corollary}[2.3.17]
\label{II.2.3.17-fr}
Dans la correspondance biunivoque décrite dans
(\hyperref[II.2.3.12-fr]{2.3.12}), aux parties fermées irréductibles de $X$
correspondent les idéaux gradués premiers dans $S_+$.
\end{corollary}

\begin{proof}
En effet, si $Y=Y_1\cup Y_2$, où $Y_1$ et $Y_2$ sont fermés et distincts de
$Y$, on a
\[
  \mathfrak{j}_+(Y)
  =\mathfrak{j}_+(Y_1)\cap\mathfrak{j}_+(Y_2)
\]
les idéaux $\mathfrak{j}_+(Y_1)$ et $\mathfrak{j}_+(Y_2)$ étant distincts
de $\mathfrak{j}_+(Y)$, donc $\mathfrak{j}_+(Y)$ n'est pas premier.
Inversement, si $\mathfrak{J}$ est un idéal gradué non premier de $S_+$, il
existe deux éléments $f$, $g$ de $S_+$ tels que
$f\notin\mathfrak{J}$, $g\notin\mathfrak{J}$, $fg\in\mathfrak{J}$~; alors
$V_+(f)\not\supset V_+(\mathfrak{J})$,
$V_+(g)\not\supset V_+(\mathfrak{J})$ mais
$V_+(\mathfrak{J})\subset V_+(f)\cup V_+(g)$ par
(\hyperref[II.2.3.2.3-fr]{2.3.2.3})~; on en conclut que
$V_+(\mathfrak{J})$ est réunion des ensembles fermés
$V_+(f)\cap V_+(\mathfrak{J})$ et $V_+(g)\cap V_+(\mathfrak{J})$, qui sont
distincts de $V_+(\mathfrak{J})$.
\end{proof}

\subsection{La structure de schéma sur $\operatorname{Proj}(S)$.}
\label{subsection:II.2.4-fr}

\begin{env}[2.4.1]
\label{II.2.4.1-fr}
Soient $f$, $g$ deux éléments homogènes de $S_+$~; considérons les schémas
affines $Y_f=\operatorname{Spec}(S_{(f)})$,
$Y_g=\operatorname{Spec}(S_{(g)})$ et
$Y_{fg}=\operatorname{Spec}(S_{(fg)})$. En vertu de
(\hyperref[II.2.2.2-fr]{2.2.2}), le morphisme
$w_{fg,f}=({}^a\omega_{fg,f},\widetilde{\omega}_{fg,f})$ de $Y_{fg}$ dans
$Y_f$, correspondant à l'homomorphisme canonique
$\omega_{fg,f}:S_{(f)}\to S_{(fg)}$, est une \emph{immersion ouverte}
(\hyperref[I.1.3.6-fr]{I, 1.3.6}). Au moyen de l'homéomorphisme réciproque
de $\psi_f:D_+(f)\to Y_f$ (\hyperref[II.2.3.6-fr]{2.3.6}), on peut
transporter à $D_+(f)$ la structure de schéma affine de $Y_f$~; en vertu de
la commutativité du diagramme (\hyperref[II.2.3.5.1-fr]{2.3.5.1}), le
schéma affine $D_+(fg)$ s'identifie ainsi au schéma induit sur l'ensemble
ouvert $D_+(fg)$ de l'espace sous-jacent au schéma affine $D_+(f)$. Il est
clair alors (compte tenu de (\hyperref[II.2.3.4-fr]{2.3.4})) que
$X=\operatorname{Proj}(S)$ est muni d'une unique structure de
\emph{préschéma}, dont la restriction à chaque $D_+(f)$ est le schéma
affine qui vient d'être défini. En outre~:
\end{env}

\begin{proposition}[2.4.2]
\label{II.2.4.2-fr}
Le préschéma $\operatorname{Proj}(S)$ est un schéma.
\end{proposition}

\begin{proof}
Il suffit (\hyperref[I.5.5.6-fr]{I, 5.5.6}) de montrer que, quels que soient
$f$, $g$ homogènes dans $S_+$,
$D_+(f)\cap D_+(g)=D_+(fg)$ est affine et que son anneau est engendré par
les images canoniques des anneaux de $D_+(f)$ et $D_+(g)$~; le premier
point est évident par définition, et le second résulte de
(\hyperref[II.2.2.4-fr]{2.2.4}).
\end{proof}
\oldpage[II]{29}
Lorsqu'on parlera du spectre premier homogène $\operatorname{Proj}(S)$ comme
d'un \emph{schéma}, il s'agira toujours de la structure qui vient d'être
définie.

\begin{example}[2.4.3]
\label{II.2.4.3-fr}
Prenons $S=K[T_1,T_2]$, $K$ étant un corps, $T_1$, $T_2$ deux
indéterminées, et $S$ étant gradué par le degré total. Il résulte de
(\hyperref[II.2.3.14-fr]{2.3.14}) que $\operatorname{Proj}(S)$ est réunion
de $D_+(T_1)$ et $D_+(T_2)$~; on voit aussitôt que ces schémas affines sont
canoniquement isomorphes à $K[T]$, et que $\operatorname{Proj}(S)$
s'obtient par le recollement de ces deux schémas affines décrit dans
(\hyperref[I.2.3.2-fr]{I, 2.3.2}) (cf.
(\hyperref[II.7.4.14-fr]{7.4.14})).
\end{example}

\begin{proposition}[2.4.4]
\label{II.2.4.4-fr}
Soit $S$ un anneau gradué à degrés positifs, et soit $X$ le schéma
$\operatorname{Proj}(S)$.
\begin{enumerate}
  \item[\rm(i)] Si $\mathfrak{N}_+$ est le nilradical de $S_+$
    (\hyperref[II.2.1.10-fr]{2.1.10}), le schéma $X_{\mathrm{red}}$ est
    canoniquement isomorphe à $\operatorname{Proj}(S/\mathfrak{N}_+)$~; en
    particulier, si $S$ est essentiellement réduite,
    $\operatorname{Proj}(S)$ est réduit.
  \item[\rm(ii)] Supposons $S$ essentiellement réduit. Alors, pour que $X$
    soit intègre, il faut et il suffit que $S$ soit essentiellement intègre.
\end{enumerate}
\end{proposition}

\begin{proof}
\begin{enumerate}
  \item[\rm(i)] Soit $\overline{S}$ l'anneau gradué
    $S/\mathfrak{N}_+$, et désignons par $x\mapsto\overline{x}$
    l'homomorphisme canonique $S\to\overline{S}$, de degré $0$. Pour tout
    $f\in S_d$ ($d>0$), l'homomorphisme canonique
    $S_f\to\overline{S}_{\overline{f}}$
    (\hyperref[0.1.5.1-fr]{0, 1.5.1}) est surjectif et de degré $0$, donc
    par restriction donne un homomorphisme surjectif
    $S_{(f)}\to\overline{S}_{(\overline{f})}$~; si on suppose
    $f\notin\mathfrak{N}_+$, on vérifie immédiatement que
    $\overline{S}_{(\overline{f})}$ est réduit, et que le noyau de
    l'homomorphisme précédent est le nilradical de $S_{(f)}$, autrement dit
    $\overline{S}_{(\overline{f})}=(S_{(f)})_{\mathrm{red}}$. À cet
    homomorphisme correspond donc une immersion fermée
    $D_+(\overline{f})\to D_+(f)$ qui identifie $D_+(\overline{f})$ à
    $(D_+(f))_{\mathrm{red}}$ (\hyperref[I.5.1.2-fr]{I, 5.1.2}), et est en
    particulier un homéomorphisme des espaces sous-jacents à ces deux
    schémas affines. En outre, si $g\notin\mathfrak{N}_+$ est un second
    élément homogène de $S_+$, le diagramme
    \[
      \xymatrix{
        S_{(f)}\ar[r]\ar[d]
        & \overline{S}_{(\overline{f})}\ar[d]\\
        S_{(fg)}\ar[r]
        & \overline{S}_{(\overline{fg})}
      }
    \]
    est commutatif~; comme en outre les $D_+(f)$, pour $f$ homogène de
    degré $>0$ et $f\notin\mathfrak{N}_+$, forment un recouvrement de
    $X=\operatorname{Proj}(S)$ (\hyperref[II.2.3.7-fr]{2.3.7}), on voit
    que les morphismes $D_+(\overline{f})\to D_+(f)$ sont les restrictions
    d'une immersion fermée
    $\operatorname{Proj}(\overline{S})\to\operatorname{Proj}(S)$ qui est
    un homéomorphisme des espaces sous-jacents~; d'où la conclusion
    (\hyperref[I.5.1.2-fr]{I, 5.1.2}).
  \item[\rm(ii)] Supposons que $S$ soit essentiellement intègre, autrement
    dit que $(0)$ soit un idéal gradué premier de $S_+$ distinct de $S_+$~;
    alors $X$ est réduit en vertu de (i) et irréductible en vertu de
    (\hyperref[II.2.3.17-fr]{2.3.17}). Inversement, supposons que $S$ soit
    essentiellement réduit et $X$ intègre~; alors, pour $f\neq0$ homogène
    dans $S_+$, on a $D_+(f)\neq\varnothing$
    (\hyperref[II.2.3.7-fr]{2.3.7})~; l'hypothèse que $X$ est irréductible
    entraîne que $D_+(f)\cap D_+(g)\neq\varnothing$ pour $f$, $g$
    homogènes et $\neq0$ dans $S_+$~; donc $fg\neq0$ en vertu de
    (\hyperref[II.2.3.3.2-fr]{2.3.3.2}), et on en conclut que $S_+$ n'a pas
    de diviseurs de zéro, d'où la première assertion.
\end{enumerate}
\end{proof}

\begin{env}[2.4.5]
\label{II.2.4.5-fr}
Étant donné un anneau commutatif $A$, rappelons qu'on dit qu'un anneau
gradué $S$ est une \emph{$A$-algèbre graduée} s'il est muni d'une structure
de $A$-algèbre telle que chacun des sous-groupes $S_n$ soit un $A$-module~;
il suffit d'ailleurs pour cela que $S_0$ soit
\oldpage[II]{30}
une $A$-algèbre, autrement dit on définit sur $S$ une structure de
$A$-algèbre graduée en définissant une structure de $A$-algèbre sur $S_0$,
et en posant, pour $\alpha\in A$, $x\in S_n$,
$\alpha\cdot x=(\alpha\cdot1)x$.
\end{env}

\begin{proposition}[2.4.6]
\label{II.2.4.6-fr}
Supposons que $S$ soit une $A$-algèbre graduée. Alors sur
$X=\operatorname{Proj}(S)$, le faisceau structural $\mathcal{O}_X$ est une
$A$-Algèbre ($A$ étant considéré comme faisceau simple sur $X$)~; en
d'autres termes, $X$ est un schéma au-dessus de $\operatorname{Spec}(A)$.
\end{proposition}

\begin{proof}
Il suffit de noter que pour tout $f$ homogène dans $S_+$, $S_{(f)}$ est une
algèbre sur $A$, et que le diagramme
\[
  \xymatrix{
    S_{(f)}\ar[rr] && S_{(fg)}\\
    & A\ar[ul]\ar[ur] &
  }
\]
est commutatif, pour $f$, $g$ homogènes dans $S_+$.
\end{proof}

\begin{proposition}[2.4.7]
\label{II.2.4.7-fr}
Soit $S$ un anneau gradué à degrés positifs.
\begin{enumerate}
  \item[\rm(i)] Pour tout entier $d>0$, il existe un isomorphisme canonique
    du schéma $\operatorname{Proj}(S)$ sur le schéma
    $\operatorname{Proj}(S^{(d)})$.
  \item[\rm(ii)] Soit $S'$ l'anneau gradué tel que $S'_0=\mathbf{Z}$,
    $S'_n=S_n$ (considéré comme $\mathbf{Z}$-module) pour $n>0$. Il existe un
    isomorphisme canonique du schéma $\operatorname{Proj}(S)$ sur le schéma
    $\operatorname{Proj}(S')$.
\end{enumerate}
\end{proposition}

\begin{proof}
\begin{enumerate}
  \item[\rm(i)] Montrons d'abord que l'application
    $\mathfrak{p}\mapsto\mathfrak{p}\cap S^{(d)}$ est une bijection de
    l'ensemble $\operatorname{Proj}(S)$ sur
    $\operatorname{Proj}(S^{(d)})$. En effet, supposons donné un idéal
    premier gradué $\mathfrak{p}'\in\operatorname{Proj}(S^{(d)})$, et posons
    $\mathfrak{p}_{nd}=\mathfrak{p}'\cap S_{nd}$ ($n\geq0$). Pour tout
    $n>0$, non multiple de $d$, définissons $\mathfrak{p}_n$ comme
    l'ensemble des $x\in S_n$ tels que
    $x^d\in\mathfrak{p}_{nd}$~; si $x\in\mathfrak{p}_n$,
    $y\in\mathfrak{p}_n$, on a
    $(x+y)^{2d}\in\mathfrak{p}_{2nd}$, donc
    $(x+y)^d\in\mathfrak{p}_{nd}$ puisque $\mathfrak{p}'$ est premier~; il
    est immédiat que les $\mathfrak{p}_n$ ainsi définis pour tout $n\geq0$
    vérifient les conditions de
    (\hyperref[II.2.1.9-fr]{2.1.9}), donc il existe un idéal premier unique
    $\mathfrak{p}\in\operatorname{Proj}(S)$ tel que
    $\mathfrak{p}\cap S^{(d)}=\mathfrak{p}'$. Comme pour tout $f$ homogène
    dans $S_+$, on a
    $V_+(f)=V_+(f^d)$ (\hyperref[II.2.3.2.3-fr]{2.3.2.3}), on voit que la
    bijection précédente est un homéomorphisme d'espaces topologiques.
    Enfin, avec les mêmes notations, $S_{(f)}$ et
    $S^{(d)}_{(f^d)}$ s'identifient canoniquement
    (\hyperref[II.2.2.2-fr]{2.2.2}), donc $\operatorname{Proj}(S)$ et
    $\operatorname{Proj}(S^{(d)})$ s'identifient canoniquement en tant que
    schémas.
  \item[\rm(ii)] Si, à tout
    $\mathfrak{p}\in\operatorname{Proj}(S)$ on fait correspondre l'unique
    idéal premier $\mathfrak{p}'\in\operatorname{Proj}(S')$ tel que
    $\mathfrak{p}'\cap S_n=\mathfrak{p}\cap S_n$ pour tout $n>0$
    (\hyperref[II.2.1.9-fr]{2.1.9}), il est clair qu'on définit un
    homéomorphisme canonique
    $\operatorname{Proj}(S)\xrightarrow{\sim}\operatorname{Proj}(S')$ des
    espaces sous-jacents, car $V_+(f)$ est le même ensemble pour $S$ et
    $S'$ lorsque $f$ est un élément homogène de $S_+$. Comme en outre
    $S_{(f)}=S'_{(f)}$, $\operatorname{Proj}(S)$ et
    $\operatorname{Proj}(S')$ s'identifient en tant que schémas.
\end{enumerate}
\end{proof}

\begin{corollary}[2.4.8]
\label{II.2.4.8-fr}
Si $S$ est une $A$-algèbre graduée, $S'_A$ la $A$-algèbre graduée telle que
$(S'_A)_0=A$, $(S'_A)_n=S_n$ pour $n>0$, il existe un isomorphisme canonique
de $\operatorname{Proj}(S)$ sur $\operatorname{Proj}(S'_A)$.
\end{corollary}

\begin{proof}
En effet, ces deux schémas sont isomorphes canoniquement à
$\operatorname{Proj}(S')$, avec les notations de
(\hyperref[II.2.4.7-fr]{2.4.7, (ii)}).
\end{proof}

\subsection{Faisceau associé à un module gradué.}
\label{subsection:II.2.5-fr}

\begin{env}[2.5.1]
\label{II.2.5.1-fr}
Soit $M$ un $S$-module gradué. Pour tout $f$ homogène dans $S_+$,
$M_{(f)}$ est un $S_{(f)}$-module, et il lui correspond donc un faisceau
quasi-cohérent associé $\widetilde{M_{(f)}}$ sur le schéma affine
$\operatorname{Spec}(S_{(f)})$, identifié à $D_+(f)$
(\hyperref[I.1.3.4-fr]{I, 1.3.4}).
\end{env}
\oldpage[II]{31}

\begin{proposition}[2.5.2]
\label{II.2.5.2-fr}
Il existe sur $X=\operatorname{Proj}(S)$ un $\mathcal{O}_X$-Module
quasi-cohérent et un seul $\widetilde{M}$ tel que pour tout $f$ homogène
dans $S_+$, on ait
$\Gamma(D_+(f),\widetilde{M})=M_{(f)}$, l'homomorphisme de restriction
$\Gamma(D_+(f),\widetilde{M})\to
\Gamma(D_+(fg),\widetilde{M})$ pour $f$, $g$ homogènes dans $S_+$, étant
l'homomorphisme canonique $M_{(f)}\to M_{(fg)}$
(\hyperref[II.2.2.3-fr]{2.2.3}).
\end{proposition}

\begin{proof}
Supposons $f\in S_d$, $g\in S_e$. Comme $D_+(fg)$ s'identifie au spectre
premier de $(S_{(f)})_{g^d/f^e}$ en vertu de
(\hyperref[II.2.2.2-fr]{2.2.2}), la restriction à $D_+(fg)$ du faisceau
$\widetilde{M_{(f)}}$ sur $D_+(f)$ s'identifie canoniquement au faisceau
associé au module $(M_{(f)})_{g^d/f^e}$
(\hyperref[I.1.3.6-fr]{I, 1.3.6}), donc aussi à
$\widetilde{M_{(fg)}}$ (\hyperref[II.2.2.2-fr]{2.2.2})~; on en conclut
qu'il existe un isomorphisme canonique
\[
  \theta_{g,f}:\widetilde{M_{(f)}}|D_+(fg)
  \xrightarrow{\sim}\widetilde{M_{(g)}}|D_+(fg)
\]
tel que, si $h$ est un troisième élément homogène de $S_+$, on ait
$\theta_{f,h}=\theta_{f,g}\circ\theta_{g,h}$ dans $D_+(fgh)$. Par suite
(\hyperref[0.3.3.1-fr]{0, 3.3.1}) il existe sur $X$ un
$\mathcal{O}_X$-Module quasi-cohérent $\mathcal{F}$, et pour chaque $f$
homogène dans $S_+$, un isomorphisme $\eta_f$ de
$\mathcal{F}|D_+(f)$ sur $\widetilde{M_{(f)}}$ tels que
$\theta_{g,f}=\eta_g\circ\eta_f^{-1}$. Si alors on considère le faisceau
$\mathcal{G}$ associé au préfaisceau (sur la base de la topologie de $X$
formée des $D_+(f)$) défini par $D_+(f)\to M_{(f)}$, avec les
homomorphismes canoniques $M_{(f)}\to M_{(fg)}$ comme homomorphismes de
restriction, ce qui précède prouve que $\mathcal{F}$ et $\mathcal{G}$ sont
isomorphes (compte tenu de (\hyperref[I.1.3.7-fr]{I, 1.3.7}))~; le
faisceau $\mathcal{G}$ sera désigné par $\widetilde{M}$ et vérifie bien les
conditions de l'énoncé. On a en particulier
$\widetilde{S}=\mathcal{O}_X$.
\end{proof}

\begin{definition}[2.5.3]
\label{II.2.5.3-fr}
On dit que le $\mathcal{O}_X$-Module quasi-cohérent $\widetilde{M}$ défini
dans (\hyperref[II.2.5.2-fr]{2.5.2}) est \emph{associé} au $S$-module
gradué $M$.
\end{definition}

Rappelons que les $S$-modules gradués forment une catégorie lorsqu'on
restreint les homomorphismes de modules gradués aux homomorphismes
\emph{de degré $0$}. Avec cette convention~:

\begin{proposition}[2.5.4]
\label{II.2.5.4-fr}
Le foncteur $M\mapsto\widetilde{M}$ est un foncteur covariant additif exact,
de la catégorie des $S$-modules gradués dans celle des
$\mathcal{O}_X$-Modules quasi-cohérents, qui commute aux limites
inductives et aux sommes directes.
\end{proposition}

\begin{proof}
En effet, ces propriétés étant locales, il suffit de les vérifier sur les
faisceaux
$\widetilde{M}|D_+(f)=\widetilde{M_{(f)}}$~; or, le foncteur
$M\mapsto M_f$, le foncteur $N\mapsto N_0$ (dans la catégorie des
$S_f$-modules gradués) et le foncteur $P\mapsto\widetilde{P}$ (dans la
catégorie des $S_{(f)}$-modules) ont tous trois les propriétés d'exactitude
et de permutabilité aux limites inductives et aux sommes directes
(\hyperref[I.1.3.5-fr]{I, 1.3.5} et
\hyperref[I.1.3.9-fr]{I, 1.3.9})~; d'où la proposition.
\end{proof}

On notera $\widetilde{u}$ l'homomorphisme
$\widetilde{M}\to\widetilde{N}$ correspondant à un homomorphisme de degré
$0$, $u:M\to N$. On déduit aussitôt de
(\hyperref[II.2.5.4-fr]{2.5.4}) que les résultats de
(\hyperref[I.1.3.9-fr]{I, 1.3.9} et
\hyperref[I.1.3.10-fr]{I, 1.3.10}) sont encore valables pour les
$S$-modules gradués et les homomorphismes de degré $0$ (avec le sens donné
ici à $\widetilde{M}$), les démonstrations étant purement formelles.

\begin{proposition}[2.5.5]
\label{II.2.5.5-fr}
Pour tout $\mathfrak{p}\in X=\operatorname{Proj}(S)$, on a
$\widetilde{M}_{\mathfrak{p}}=M_{(\mathfrak{p})}$.
\end{proposition}

\begin{proof}
En effet, par définition
$\widetilde{M}_{\mathfrak{p}}=
\varinjlim\Gamma(D_+(f),\widetilde{M})$, où $f$ parcourt l'ensemble des
éléments homogènes de $S_+$ tels que $f\notin\mathfrak{p}$~; comme
$\Gamma(D_+(f),\widetilde{M})=M_{(f)}$, la proposition résulte de la
définition de $M_{(\mathfrak{p})}$
(\hyperref[II.2.2.7-fr]{2.2.7}).
\end{proof}
\oldpage[II]{32}

En particulier, \emph{l'anneau local $(\mathcal{O}_X)_\mathfrak{p}$} n'est
autre que l'anneau $S_{(\mathfrak{p})}$, ensemble des fractions $x/f$, où
$f$ est homogène dans $S_+$ et n'appartient pas à $\mathfrak{p}$, et $x$
homogène est de même degré que $f$.

Plus particulièrement, si $S$ est \emph{essentiellement intègre}, de sorte
que $\operatorname{Proj}(S)=X$ est \emph{intègre}
(\hyperref[II.2.4.4-fr]{2.4.4}), et si $\xi=(0)$ est le point générique de
$X$, le \emph{corps des fonctions rationnelles}
$R(X)=\mathcal{O}_\xi$, est le corps formé des éléments $fg^{-1}$, où $f$
et $g$ sont homogènes de même degré dans $S_+$ et $g\neq0$.

\begin{proposition}[2.5.6]
\label{II.2.5.6-fr}
Si, pour tout $z\in M$ et tout $f$ homogène dans $S_+$, il existe une
puissance de $f$ annulant $z$, on a $\widetilde{M}=0$. Cette condition
suffisante est aussi nécessaire lorsque la $S_0$-algèbre $S$ est engendrée
par l'ensemble $S_1$ des éléments homogènes de degré $1$.
\end{proposition}

\begin{proof}
En effet, la condition $\widetilde{M}=0$ équivaut à $M_{(f)}=0$ pour tout
$f$ homogène dans $S_+$. D'autre part, si $f\in S_d$, dire que $M_{(f)}=0$
signifie que pour tout $z\in M$ homogène et de degré multiple de $d$, il y
a une puissance $f^n$ telle que $f^nz=0$. Si $M_{(f)}=0$ pour tout
$f\in S_1$, la condition de l'énoncé est donc vérifiée pour tout
$f\in S_1$~; elle l'est \emph{a fortiori} pour tout $f$ homogène dans
$S_+$ lorsque $S_1$ engendre $S$, puisque tout élément homogène de $S_+$
est alors combinaison linéaire de produits d'éléments de $S_1$.
\end{proof}

\begin{proposition}[2.5.7]
\label{II.2.5.7-fr}
Soient $d$ un entier $>0$, $f\in S_d$. Alors, pour tout $n\in\mathbf{Z}$,
le $(\mathcal{O}_X|D_+(f))$-Module
$\widetilde{S(nd)}|D_+(f)$ est canoniquement isomorphe à
$\mathcal{O}_X|D_+(f)$.
\end{proposition}

\begin{proof}
En effet, la multiplication par l'élément inversible $(f/1)^n$ de $S_f$
est une bijection de $S_{(f)}=(S_f)_0$ sur
$(S_f)_{nd}=(S_f(nd))_0=(S(nd)_f)_0=S(nd)_{(f)}$~; autrement dit, les
$S_{(f)}$-modules $S_{(f)}$ et $S(nd)_{(f)}$ sont canoniquement
isomorphes, d'où la proposition.
\end{proof}

\begin{corollary}[2.5.8]
\label{II.2.5.8-fr}
Sur l'ensemble ouvert
\[
  U=\bigcup_{f\in S_d}D_+(f),
\]
la restriction du $\mathcal{O}_X$-Module $\widetilde{S(nd)}$ est un
$(\mathcal{O}_X|U)$-Module inversible
(\hyperref[0.5.4.1-fr]{0, 5.4.1}).
\end{corollary}

\begin{corollary}[2.5.9]
\label{II.2.5.9-fr}
Si l'idéal $S_+$ de $S$ est engendré par l'ensemble $S_1$ des éléments
homogènes de degré $1$, les $\mathcal{O}_X$-Modules
$\widetilde{S(n)}$ sont inversibles pour tout $n\in\mathbf{Z}$.
\end{corollary}

\begin{proof}
Il suffit de remarquer que
$X=\bigcup_{f\in S_1}D_+(f)$ en vertu de l'hypothèse
(\hyperref[II.2.3.14-fr]{2.3.14}) et d'appliquer
(\hyperref[II.2.5.8-fr]{2.5.8}) avec $U=X$.
\end{proof}

\begin{env}[2.5.10]
\label{II.2.5.10-fr}
Nous poserons, dans toute la suite de ce paragraphe
\[
\label{II.2.5.10.1-fr}
  \mathcal{O}_X(n)=\widetilde{S(n)}
\tag{2.5.10.1}
\]
pour tout $n\in\mathbf{Z}$, et pour toute partie ouverte $U$ de $X$ et tout
$(\mathcal{O}_X|U)$-Module $\mathcal{F}$
\[
\label{II.2.5.10.2-fr}
  \mathcal{F}(n)=
  \mathcal{F}\otimes_{\mathcal{O}_X|U}(\mathcal{O}_X(n)|U)
\tag{2.5.10.2}
\]
pour tout $n\in\mathbf{Z}$. Si l'idéal $S_+$ est engendré par $S_1$, le
foncteur $\mathcal{F}(n)$ est \emph{exact} en
$\mathcal{F}$ pour tout $n\in\mathbf{Z}$, car $\mathcal{O}_X(n)$ est alors
un $\mathcal{O}_X$-Module inversible.
\end{env}

\begin{env}[2.5.11]
\label{II.2.5.11-fr}
Soient $M$ et $N$ deux $S$-modules gradués. Pour tout $f\in S_d$ ($d>0$)
on définit un homomorphisme canonique fonctoriel de $S_{(f)}$-modules
\[
\label{II.2.5.11.1-fr}
  \lambda_f:M_{(f)}\otimes_{S_{(f)}}N_{(f)}
  \longrightarrow(M\otimes_S N)_{(f)}
\tag{2.5.11.1}
\]
\oldpage[II]{33}

en composant l'homomorphisme
$M_{(f)}\otimes_{S_{(f)}}N_{(f)}\to M_f\otimes_{S_f}N_f$
(provenant des injections $M_{(f)}\to M_f$, $N_{(f)}\to N_f$ et
$S_{(f)}\to S_f$) et l'isomorphisme canonique
$M_f\otimes_{S_f}N_f\xrightarrow{\sim}(M\otimes_S N)_f$
(\hyperref[0.1.3.4-fr]{0, 1.3.4}) et en notant que, par définition du
produit tensoriel de deux modules gradués, ce dernier isomorphisme conserve
les degrés~; pour $x\in M_{md}$, $y\in N_{nd}$ ($m\geq0$, $n\geq0$), on a
donc
\[
  \lambda_f((x/f^m)\otimes(y/f^n))=(x\otimes y)/f^{m+n}.
\]

Il résulte aussitôt de cette définition que si $g\in S_e$ ($e>0$), le
diagramme
\[
  \xymatrix{
    M_{(f)}\otimes_{S_{(f)}}N_{(f)}\ar[r]^{\lambda_f}\ar[d]
    & (M\otimes_S N)_{(f)}\ar[d]\\
    M_{(fg)}\otimes_{S_{(fg)}}N_{(fg)}\ar[r]_{\lambda_{fg}}
    & (M\otimes_S N)_{(fg)}
  }
\]
(où la flèche verticale de droite est l'homomorphisme canonique et celle de
gauche provient des homomorphismes canoniques) est commutatif. On déduit
donc des $\lambda$ un homomorphisme canonique fonctoriel de
$\mathcal{O}_X$-Modules
\[
\label{II.2.5.11.2-fr}
  \lambda:\widetilde{M}\otimes_{\mathcal{O}_X}\widetilde{N}
  \longrightarrow\widetilde{M\otimes_S N}.
\tag{2.5.11.2}
\]

Considérons en particulier deux idéaux gradués $\mathfrak{J}$,
$\mathfrak{K}$ de $S$~; comme $\widetilde{\mathfrak{J}}$ et
$\widetilde{\mathfrak{K}}$ sont des faisceaux d'idéaux de
$\mathcal{O}_X$, on a un homomorphisme canonique
$\widetilde{\mathfrak{J}}\otimes_{\mathcal{O}_X}
\widetilde{\mathfrak{K}}\to\mathcal{O}_X$~; le diagramme
\[
\label{II.2.5.11.3-fr}
  \xymatrix{
    \widetilde{\mathfrak{J}}\otimes_{\mathcal{O}_X}
    \widetilde{\mathfrak{K}}\ar[rr]^\lambda\ar[dr]
    && \widetilde{\mathfrak{J}\otimes_S\mathfrak{K}}\ar[dl]\\
    & \mathcal{O}_X &
  }
\tag{2.5.11.3}
\]
est alors commutatif. On se ramène en effet à le vérifier dans chaque
ouvert $D_+(f)$ ($f$ homogène dans $S_+$) et cela résulte aussitôt de la
définition (\hyperref[II.2.5.11.1-fr]{2.5.11.1}) de $\lambda_f$ et de
(\hyperref[I.1.3.13-fr]{I, 1.3.13}).

Notons enfin que si $M$, $N$, $P$ sont trois $S$-modules gradués, le
diagramme
\[
\label{II.2.5.11.4-fr}
  \xymatrix{
    \widetilde{M}\otimes_{\mathcal{O}_X}\widetilde{N}
    \otimes_{\mathcal{O}_X}\widetilde{P}
    \ar[r]^{\lambda\otimes1}\ar[d]_{1\otimes\lambda}
    & \widetilde{M\otimes_S N}\otimes_{\mathcal{O}_X}\widetilde{P}
    \ar[d]^\lambda\\
    \widetilde{M}\otimes_{\mathcal{O}_X}\widetilde{N\otimes_S P}
    \ar[r]_\lambda
    & \widetilde{M\otimes_S N\otimes_S P}
  }
\tag{2.5.11.4}
\]
\oldpage[II]{34}

est commutatif. Il suffit encore de le vérifier dans chaque $D_+(f)$ et
cela découle aussitôt des définitions et de
(\hyperref[I.1.3.13-fr]{I, 1.3.13}).
\end{env}

\begin{env}[2.5.12]
\label{II.2.5.12-fr}
Sous les hypothèses de (\hyperref[II.2.5.11-fr]{2.5.11}), on définit un
homomorphisme canonique fonctoriel de $S_{(f)}$-modules
\[
\label{II.2.5.12.1-fr}
  \mu_f:\bigl(\operatorname{Hom}_S(M,N)\bigr)_{(f)}
  \longrightarrow
  \operatorname{Hom}_{S_{(f)}}(M_{(f)},N_{(f)})
\tag{2.5.12.1}
\]
en faisant correspondre à $u/f^n$, où $u$ est un homomorphisme de degré
$nd$, l'homomorphisme $M_{(f)}\to N_{(f)}$ qui transforme $x/f^m$
($x\in M_{md}$) en $u(x)/f^{m+n}$. Pour $g\in S_e$ ($e>0$), on a encore
un diagramme commutatif
\[
  \xymatrix{
    \bigl(\operatorname{Hom}_S(M,N)\bigr)_{(f)}
    \ar[r]^{\mu_f}\ar[d]
    & \operatorname{Hom}_{S_{(f)}}(M_{(f)},N_{(f)})\ar[d]\\
    \bigl(\operatorname{Hom}_S(M,N)\bigr)_{(fg)}
    \ar[r]_{\mu_{fg}}
    & \operatorname{Hom}_{S_{(fg)}}(M_{(fg)},N_{(fg)})
  }
\]
(la flèche de gauche étant l'homomorphisme canonique, celle de droite
provenant des homomorphismes canoniques). On en conclut encore (compte
tenu de (\hyperref[I.1.3.8-fr]{I, 1.3.8})) que les $\mu_f$ définissent un
homomorphisme canonique fonctoriel de $\mathcal{O}_X$-Modules
\[
\label{II.2.5.12.2-fr}
  \mu:\widetilde{\operatorname{Hom}_S(M,N)}
  \longrightarrow
  \mathcal{H}om_{\mathcal{O}_X}(\widetilde{M},\widetilde{N}).
\tag{2.5.12.2}
\]
\end{env}

\begin{proposition}[2.5.13]
\label{II.2.5.13-fr}
Supposons l'idéal $S_+$ engendré par $S_1$. Alors l'homomorphisme
$\lambda$ (\hyperref[II.2.5.11.2-fr]{2.5.11.2}) est un isomorphisme~; il
en est de même de l'homomorphisme $\mu$
(\hyperref[II.2.5.12.2-fr]{2.5.12.2}) lorsque le $S$-module gradué $M$
admet une présentation finie (\hyperref[II.2.1.1-fr]{2.1.1}).
\end{proposition}

\begin{proof}
Comme $X$ est réunion des $D_+(f)$ pour $f\in S_1$
(\hyperref[II.2.3.14-fr]{2.3.14}), on est ramené à prouver que $\lambda_f$
et $\mu_f$ sont des isomorphismes, sous les hypothèses envisagées, lorsque
$f$ est homogène et \emph{de degré $1$}. Or, on définit alors une
application $\mathbf{Z}$-bilinéaire
\[
  M_m\times N_n\longrightarrow
  M_{(f)}\otimes_{S_{(f)}}N_{(f)}
\]
en faisant correspondre à $(x,y)$ l'élément
$(x/f^m)\otimes(y/f^n)$ (lorsque $m<0$, nous écrivons $x/f^m$ pour
$f^{-m}x/1$)~; ces applications définissent une application
$\mathbf{Z}$-linéaire
\[
  M\otimes_{\mathbf{Z}}N\longrightarrow
  M_{(f)}\otimes_{S_{(f)}}N_{(f)},
\]
et si $s\in S_q$, cette application transforme $(sx)\otimes y$ en
\[
  (s/f^q)\bigl((x/f^m)\otimes(y/f^n)\bigr)
\]
(pour $x\in M_m$, $y\in N_n$). On en déduit donc un di-homomorphisme de
modules
\[
  \gamma_f:M\otimes_S N\longrightarrow
  M_{(f)}\otimes_{S_{(f)}}N_{(f)},
\]
relatif à l'homomorphisme canonique $S\to S_{(f)}$ (transformant
$s\in S_q$ en $s/f^q$). Supposons en outre que pour un élément
$\sum_i(x_i\otimes y_i)$ de $M\otimes_S N$ ($x_i$, $y_i$ homogènes de
degrés respectifs $m_i$, $n_i$) on ait
\[
  f^r\sum_i(x_i\otimes y_i)=0,
\]
autrement dit $\sum_i(f^rx_i\otimes y_i)=0$. On en déduit en vertu de
(\hyperref[0.1.3.4-fr]{0, 1.3.4}) que l'on a
\[
  \sum_i(f^rx_i/f^{m_i+r})\otimes(y_i/f^{n_i})=0,
\]
c'est-à-dire $\gamma_f(\sum_i(x_i\otimes y_i))=0$. Par suite $\gamma_f$
se factorise en
\[
  M\otimes_S N\longrightarrow(M\otimes_S N)_f
  \xrightarrow{\gamma'_f}
  M_{(f)}\otimes_{S_{(f)}}N_{(f)}~;
\]
si $\lambda'_f$ est la restriction de $\gamma'_f$
\oldpage[II]{35}
à $(M\otimes_S N)_{(f)}$, on vérifie aussitôt que $\lambda_f$ et
$\lambda'_f$ sont des $S_{(f)}$-homomorphismes réciproques, d'où la
première partie de la proposition.

Pour démontrer la seconde partie, supposons que $M$ soit le conoyau d'un
homomorphisme $P\to Q$, de $S$-modules gradués, $P$ et $Q$ étant sommes
directes d'un nombre fini de modules de la forme $S(n)$~; utilisant
l'exactitude à gauche de $\operatorname{Hom}_S(L,N)$ en $L$ et l'exactitude
de $M_{(f)}$ en $M$, on est aussitôt ramené à prouver que $\mu_f$ est un
isomorphisme lorsque $M=S(n)$. Or, pour tout $z$ homogène dans $N$, soit
$u_z$ l'homomorphisme de $S(n)$ dans $N$ tel que $u_z(1)=z$~; on voit
aussitôt que $\eta:z\mapsto u_z$ est un isomorphisme de degré $0$ de
$N(-n)$ sur $\operatorname{Hom}_S(S(n),N)$. Il lui correspond un
isomorphisme
\[
  \eta_f:\bigl(N(-n)\bigr)_{(f)}
  \longrightarrow
  \bigl(\operatorname{Hom}_S(S(n),N)\bigr)_{(f)}.
\]
D'autre part, soit $\eta'_f$ l'isomorphisme
\[
  N_{(f)}\longrightarrow
  \operatorname{Hom}_{S_{(f)}}\bigl(S(n)_{(f)},N_{(f)}\bigr)
\]
qui, à tout $z'\in N_{(f)}$, fait correspondre l'homomorphisme $v_{z'}$ tel
que $v_{z'}(s/f^k)=sz'/f^{n+k}$ (pour
$s\in S_{n+k}=(S(n))_k$). On constate aisément que l'application composée
\[
  \bigl(N(-n)\bigr)_{(f)}
  \xrightarrow{\eta_f}
  \bigl(\operatorname{Hom}_S(S(n),N)\bigr)_{(f)}
  \xrightarrow{\mu_f}
  \operatorname{Hom}_{S_{(f)}}\bigl(S(n)_{(f)},N_{(f)}\bigr)
  \xrightarrow{{\eta'_f}^{-1}}N_{(f)}
\]
est l'isomorphisme $z/f^h\mapsto z/f^{h-n}$ de
$\bigl(N(-n)\bigr)_{(f)}$ sur $N_{(f)}$, donc $\mu_f$ est un isomorphisme.
\end{proof}

Si l'idéal $S_+$ est engendré par $S_1$, on déduit de
(\hyperref[II.2.5.13-fr]{2.5.13}) que pour tout idéal gradué $\mathfrak{J}$
de $S$ et tout $S$-module gradué $M$, on a
\[
\label{II.2.5.13.1-fr}
  \widetilde{\mathfrak{J}}\cdot\widetilde{M}
  =\widetilde{\mathfrak{J}\cdot M}
\tag{2.5.13.1}
\]
à un isomorphisme canonique près~; cela résulte en effet de la commutativité
du diagramme
\[
  \xymatrix{
    \widetilde{\mathfrak{J}}\otimes_{\mathcal{O}_X}\widetilde{M}
    \ar[rr]^\lambda\ar[dr]
    && \widetilde{\mathfrak{J}\otimes_S M}\ar[dl]\\
    & \widetilde{M} &
  }
\]
que l'on vérifie comme celle de
(\hyperref[II.2.5.11.3-fr]{2.5.11.3}).

\begin{corollary}[2.5.14]
\label{II.2.5.14-fr}
Supposons $S$ engendrée par $S_1$. Quels que soient $m$, $n$ dans
$\mathbf{Z}$, on a alors~:
\[
\label{II.2.5.14.1-fr}
  \mathcal{O}_X(m)\otimes_{\mathcal{O}_X}\mathcal{O}_X(n)
  =\mathcal{O}_X(m+n)
\tag{2.5.14.1}
\]
\[
\label{II.2.5.14.2-fr}
  \mathcal{O}_X(n)=\bigl(\mathcal{O}_X(1)\bigr)^{\otimes n}
\tag{2.5.14.2}
\]
à des isomorphismes canoniques près.
\end{corollary}

\begin{proof}
La première formule résulte de (\hyperref[II.2.5.13-fr]{2.5.13}) et de
l'existence de l'isomorphisme canonique de degré $0$,
$S(m)\otimes_S S(n)\xrightarrow{\sim}S(m+n)$, qui, à l'élément
$1\otimes 1$ (où le premier facteur $1$ est dans $(S(m))_{-m}$ et le second
dans $(S(n))_{-n}$) fait correspondre l'élément
$1\in(S(m+n))_{-m-n}$. Il suffit ensuite de démontrer la seconde formule
pour $n=-1$, et en vertu de (\hyperref[II.2.5.13-fr]{2.5.13}), cela revient
à voir que $\operatorname{Hom}_S(S(1),S)$ est canoniquement isomorphe à
$S(-1)$, vérification immédiate en remontant aux définitions
(\hyperref[II.2.1.2-fr]{2.1.2}) et en se souvenant que $S(1)$ est un
$S$-module monogène.
\end{proof}

\oldpage[II]{36}
\begin{corollary}[2.5.15]
\label{II.2.5.15-fr}
Supposons $S$ engendrée par $S_1$. Pour tout $S$-module gradué $M$ et
tout $n\in\mathbf{Z}$, on a
\[
\label{II.2.5.15.1-fr}
  \widetilde{M(n)}=\widetilde{M}(n)
\tag{2.5.15.1}
\]
à un isomorphisme canonique près.
\end{corollary}

\begin{proof}
Cela résulte des définitions (\hyperref[II.2.5.10.2-fr]{2.5.10.2}) et
(\hyperref[II.2.5.10.1-fr]{2.5.10.1}), de la prop.
(\hyperref[II.2.5.13-fr]{2.5.13}) et de l'existence d'un isomorphisme
canonique $M(n)\xrightarrow{\sim}M\otimes_S S(n)$ de degré $0$ qui, à tout
$z\in(M(n))_h=M_{n+h}$, fait correspondre
$z\otimes1\in M_{n+h}\otimes(S(n))_{-n}\subset(M\otimes_S S(n))_h$.
\end{proof}

\begin{env}[2.5.16]
\label{II.2.5.16-fr}
Désignons par $S'$ l'anneau gradué tel que $S'_0=\mathbf{Z}$,
$S'_n=S_n$ pour $n>0$. Alors, si $f\in S_d$ ($d>0$), on a
$(S(n))_{(f)}=(S'(n))_{(f)}$ pour tout $n\in\mathbf{Z}$, car un élément de
$(S'(n))_{(f)}$ est de la forme $x/f^k$ avec $x\in S'_{n+kd}$ ($k>0$), et
on peut toujours prendre $k$ tel que $n+kd\neq0$. Comme
$X=\operatorname{Proj}(S)$ et $X'=\operatorname{Proj}(S')$ s'identifient
canoniquement (\hyperref[II.2.4.7-fr]{2.4.7, (ii)}), on voit que pour tout
$n\in\mathbf{Z}$, $\mathcal{O}_X(n)$ et $\mathcal{O}_{X'}(n)$ sont images
l'un de l'autre par l'identification précédente.

Notons d'autre part que pour tout $d>0$ et tout $n\in\mathbf{Z}$, on a
\[
  (S^{(d)}(n))_h=S_{(n+h)d}=(S(nd))_{hd}
\]
pour $f\in S_d$, on a donc $(S^{(d)}(n))_{(f)}=(S(nd))_{(f)}$. On sait que
les schémas $X=\operatorname{Proj}(S)$ et
$X^{(d)}=\operatorname{Proj}(S^{(d)})$ s'identifient canoniquement
(\hyperref[II.2.4.7-fr]{2.4.7, (i)})~; ce qui précède montre que si la
$S_0$-algèbre $S^{(d)}$ est engendrée par $S_d$, $\mathcal{O}_X(nd)$ et
$\mathcal{O}_{X^{(d)}}(n)$ sont images l'un de l'autre par cette
identification, pour tout $n\in\mathbf{Z}$.
\end{env}

\begin{proposition}[2.5.17]
\label{II.2.5.17-fr}
Soit $d$ un entier $>0$, et soit
$U=\bigcup_{f\in S_d}D_+(f)$. La restriction à $U$ de l'homomorphisme
canonique
$\mathcal{O}_X(nd)\otimes_{\mathcal{O}_X}\mathcal{O}_X(-nd)
\to\mathcal{O}_X$ est un isomorphisme pour tout entier $n$.
\end{proposition}

\begin{proof}
En vertu de (\hyperref[II.2.5.16-fr]{2.5.16}), on peut se borner au cas où
$d=1$, et la conclusion résulte alors de la démonstration de
(\hyperref[II.2.5.13-fr]{2.5.13}).
\end{proof}

\subsection{$S$-module gradué associé à un faisceau sur
$\operatorname{Proj}(S)$.}
\label{subsection:II.2.6-fr}

\emph{On suppose dans tout ce numéro que l'idéal $S_+$ de $S$ est engendré
par l'ensemble $S_1$ des éléments homogènes de degré $1$.}

\begin{env}[2.6.1]
\label{II.2.6.1-fr}
Le $\mathcal{O}_X$-Module $\mathcal{O}_X(1)$ est alors
\emph{inversible} (\hyperref[II.2.5.9-fr]{2.5.9})~; on pose alors, pour tout
$\mathcal{O}_X$-Module $\mathcal{F}$
(\hyperref[0.5.4.6-fr]{0, 5.4.6}),
\[
\label{II.2.6.1.1-fr}
  \Gamma_*(\mathcal{F})
  =\Gamma_*(\mathcal{O}_X(1),\mathcal{F})
  =\bigoplus_{n\in\mathbf{Z}}\Gamma(X,\mathcal{F}(n))
\tag{2.6.1.1}
\]
compte tenu de (\hyperref[II.2.5.14.2-fr]{2.5.14.2}). Rappelons
(\hyperref[0.5.4.6-fr]{0, 5.4.6}) que $\Gamma_*(\mathcal{O}_X)$ est muni
d'une structure d'\emph{anneau gradué}, et $\Gamma_*(\mathcal{F})$ d'une
structure de
\emph{module gradué} sur $\Gamma_*(\mathcal{O}_X)$.

Puisque $\mathcal{O}_X(n)$ est localement libre, $\Gamma_*(\mathcal{F})$
est un foncteur covariant additif
\emph{exact à gauche} en $\mathcal{F}$~; en particulier, si $\mathcal{J}$
est un faisceau d'idéaux de $\mathcal{O}_X$, $\Gamma_*(\mathcal{J})$
s'identifie canoniquement à un
\emph{idéal gradué} de $\Gamma_*(\mathcal{O}_X)$.
\end{env}

\begin{env}[2.6.2]
\label{II.2.6.2-fr}
Soit $M$ un $S$-module gradué~; pour tout $f\in S_d$ ($d>0$),
$x\mapsto x/1$ est un homomorphisme de groupes abéliens
$M_0\to M_{(f)}$, et comme $M_{(f)}$ s'identifie canoni-

\oldpage[II]{37}

quement à $\Gamma(D_+(f),\widetilde{M})$, on obtient ainsi un
homomorphisme de groupes abéliens
\[
  \alpha_0^f:M_0\longrightarrow\Gamma(D_+(f),\widetilde{M}).
\]
Il est clair que, pour tout $g\in S_e$ ($e>0$), le diagramme
\[
  \xymatrix{
    & \Gamma(D_+(f),\widetilde{M})\ar[dd]\\
    M_0\ar[ur]^{\alpha_0^f}\ar[dr]_{\alpha_0^{fg}}\\
    & \Gamma(D_+(fg),\widetilde{M})
  }
\]
est commutatif~; cela signifie que pour tout $x\in M_0$, les sections
$\alpha_0^f(x)$ et $\alpha_0^g(x)$ de $\widetilde{M}$ coïncident dans
$D_+(f)\cap D_+(g)$, et par suite il existe une section unique
$\alpha_0(x)\in\Gamma(X,\widetilde{M})$ dont la restriction à chaque
$D_+(f)$ est $\alpha_0^f(x)$. On a ainsi défini (sans utiliser l'hypothèse
que $S$ est engendrée par $S_1$) un homomorphisme de groupes abéliens
\[
\label{II.2.6.2.1-fr}
  \alpha_0:M_0\longrightarrow\Gamma(X,\widetilde{M}).
\tag{2.6.2.1}
\]

Appliquant ce résultat au $S$-module gradué $M(n)$ (pour tout
$n\in\mathbf{Z}$), on obtient pour chaque $n\in\mathbf{Z}$ un homomorphisme
de groupes abéliens
\[
\label{II.2.6.2.2-fr}
  \alpha_n:M_n=(M(n))_0\longrightarrow\Gamma(X,\widetilde{M}(n))
\tag{2.6.2.2}
\]
(compte tenu de (\hyperref[II.2.5.15-fr]{2.5.15}))~; d'où un homomorphisme
fonctoriel (de degré $0$) de groupes abéliens gradués
\[
\label{II.2.6.2.3-fr}
  \alpha:M\longrightarrow\Gamma_*(\widetilde{M})
\tag{2.6.2.3}
\]
(aussi noté $\alpha_M$) qui, dans chaque $M_n$, coïncide avec $\alpha_n$.

Si on prend en particulier $M=S$, on vérifie aussitôt (compte tenu de la
définition de la multiplication dans $\Gamma_*(\mathcal{O}_X)$
(\hyperref[0.5.4.6-fr]{0, 5.4.6})) que
$\alpha:S\to\Gamma_*(\mathcal{O}_X)$ est un homomorphisme d'anneaux
gradués, et que pour tout $S$-module gradué $M$,
(\hyperref[II.2.6.2.3-fr]{2.6.2.3}) est un di-homomorphisme de modules
gradués.
\end{env}

\begin{proposition}[2.6.3]
\label{II.2.6.3-fr}
Pour tout $f\in S_d$ ($d>0$), $D_+(f)$ est identique à l'ensemble des
$\mathfrak{p}\in X$ où la section $\alpha_d(f)$ de $\mathcal{O}_X(d)$ ne
s'annule pas (\hyperref[0.5.5.2-fr]{0, 5.5.2}).
\end{proposition}

\begin{proof}
Comme $X=\bigcup_{g\in S_1}D_+(g)$ par hypothèse, il suffit de vérifier que
pour tout $g\in S_1$, l'ensemble des $\mathfrak{p}\in D_+(g)$ où
$\alpha_d(f)$ ne s'annule pas est identique à $D_+(fg)$. Or, la restriction
de $\alpha_d(f)$ à $D_+(g)$ est par définition la section correspondant à
l'élément $f/1$ de $(S(d))_{(g)}$~; par l'isomorphisme canonique
$(S(d))_{(g)}\xrightarrow{\sim}S_{(g)}$
(\hyperref[II.2.5.7-fr]{2.5.7}), cette section de $\mathcal{O}_X(d)$
au-dessus de $D_+(g)$ s'identifie à la section de $\mathcal{O}_X$ au-dessus
de $D_+(g)$ qui correspond à l'élément $f/g^d$ de $S_{(g)}$~; dire que cette
section s'annule en $\mathfrak{p}\in D_+(g)$ signifie que
$f/g^d\in\mathfrak{q}$, où $\mathfrak{q}$ est l'idéal premier de $S_{(g)}$
correspondant à $\mathfrak{p}$ (\hyperref[II.2.3.6-fr]{2.3.6})~; par
définition cela veut dire que $f\in\mathfrak{p}$, d'où la proposition.
\end{proof}

\begin{env}[2.6.4]
\label{II.2.6.4-fr}
Soit maintenant $\mathcal{F}$ un $\mathcal{O}_X$ module, et posons
$M=\Gamma_*(\mathcal{F})$~; en vertu de l'existence de l'homomorphisme
d'anneaux gradués $\alpha:S\to\Gamma_*(\mathcal{O}_X)$, $M$ peut être
considéré comme un $S$-module gradué. Pour tout $f\in S_d$ ($d>0$), il
résulte de (\hyperref[II.2.6.3-fr]{2.6.3}) que la restriction à $D_+(f)$ de
la section $\alpha_d(f)$ de $\mathcal{O}_X(d)$ est inversible~; il en est
donc de même de la restriction à $D_+(f)$ de la section $\alpha_d(f^n)$ de
$\mathcal{O}_X(nd)$, pour tout $n>0$. Soit alors
$z\in M_{nd}=\Gamma(X,\mathcal{F}(nd))$ ($n>0$)~; s'il existe un entier
$k\geq0$ tel que la restriction à $D_+(f)$ de $f^kz$, c'est-à-dire la
\oldpage[II]{38}
section
$(z|D_+(f))(\alpha_d(f^k)|D_+(f))$ de $\mathcal{F}((n+k)d)$, soit nulle,
alors, en raison de la remarque précédente, on a aussi $z|D_+(f)=0$. Cela
montre que l'on définit un $S_{(f)}$-homomorphisme
$\beta_f:M_{(f)}\to\Gamma(D_+(f),\mathcal{F})$ en faisant correspondre à
l'élément $z/f^n$ la section
$(z|D_+(f))(\alpha_d(f^n)|D_+(f))^{-1}$ de $\mathcal{F}$ au-dessus de
$D_+(f)$. On vérifie en outre aussitôt que le diagramme
\[
\label{II.2.6.4.1-fr}
  \xymatrix{
    M_{(f)}\ar[r]^{\beta_f}\ar[d]
    & \Gamma(D_+(f),\mathcal{F})\ar[d]\\
    M_{(fg)}\ar[r]_{\beta_{fg}}
    & \Gamma(D_+(fg),\mathcal{F})
  }
\tag{2.6.4.1}
\]
est commutatif pour $g\in S_e$ ($e>0$). Si on se rappelle que $M_{(f)}$
s'identifie canoniquement à $\Gamma(D_+(f),\widetilde{M})$ et que les
$D_+(f)$ forment une base de la topologie de $X$
(\hyperref[II.2.3.4-fr]{2.3.4}), on voit que les $\beta_f$ proviennent d'un
unique homomorphisme canonique de $\mathcal{O}_X$-Modules
\[
\label{II.2.6.4.2-fr}
  \beta:\widetilde{\Gamma_*(\mathcal{F})}
  \longrightarrow\mathcal{F}
\tag{2.6.4.2}
\]
(aussi noté $\beta_{\mathcal{F}}$) qui est évidemment fonctoriel.
\end{env}

\begin{proposition}[2.6.5]
\label{II.2.6.5-fr}
Soient $M$ un $S$-module gradué, $\mathcal{F}$ un $\mathcal{O}_X$-Module~;
alors les homomorphismes composés
\[
\label{II.2.6.5.1-fr}
  \widetilde{M}\xrightarrow{\widetilde{\alpha}}
  \widetilde{\Gamma_*(\widetilde{M})}
  \xrightarrow{\beta}\widetilde{M}
\tag{2.6.5.1}
\]
\[
\label{II.2.6.5.2-fr}
  \Gamma_*(\mathcal{F})\xrightarrow{\alpha}
  \Gamma_*(\widetilde{\Gamma_*(\mathcal{F})})
  \xrightarrow{\Gamma_*(\beta)}\Gamma_*(\mathcal{F})
\tag{2.6.5.2}
\]
sont les isomorphismes identiques.
\end{proposition}

\begin{proof}
La vérification de (\hyperref[II.2.6.5.1-fr]{2.6.5.1}) est locale~: dans un
ouvert $D_+(f)$, elle résulte aussitôt des définitions et de ce que $\beta$,
appliqué à des faisceaux quasi-cohérents, est déterminé par son action sur les
sections au-dessus de $D_+(f)$ (\hyperref[I.1.3.8-fr]{I, 1.3.8}). La
vérification de (\hyperref[II.2.6.5.2-fr]{2.6.5.2}) se fait pour chaque degré
séparément~: si on pose $M=\Gamma_*(\mathcal{F})$, on a
$M_n=\Gamma(X,\mathcal{F}(n))$ et
$(\Gamma_*(\widetilde{M}))_n=\Gamma(X,\widetilde{M}(n))
=\Gamma(X,\widetilde{M(n)})$. Or, si $f\in S_1$ et $z\in M_n$,
$\alpha_n^f(z)$ est l'élément $z/1$ de $(M(n))_{(f)}$, égal à
$(f/1)^n(z/f^n)$~; il lui correspond par $\beta_f$ la section
\[
  ((\alpha_1(f))^n|D_+(f))
  ((z|D_+(f))((\alpha_1(f))^n|D_+(f))^{-1})
\]
au-dessus de $D_+(f)$, c'est-à-dire la restriction de $z$ à $D_+(f)$, ce qui
vérifie (\hyperref[II.2.6.5.2-fr]{2.6.5.2}).
\end{proof}

\subsection{Conditions de finitude.}
\label{subsection:II.2.7-fr}

\begin{proposition}[2.7.1]
\label{II.2.7.1-fr}
\begin{enumerate}
  \item[\rm(i)] Si $S$ est un anneau gradué noethérien,
    $X=\operatorname{Proj}(S)$ est un schéma noethérien.
  \item[\rm(ii)] Si $S$ est une $A$-algèbre graduée de type fini,
    $X=\operatorname{Proj}(S)$ est un schéma de type fini au-dessus de
    $Y=\operatorname{Spec}(A)$.
\end{enumerate}
\end{proposition}

\oldpage[II]{39}
\begin{proof}
\begin{enumerate}
  \item[\rm(i)] Si $S$ est noethérien, l'idéal $S_+$ admet un système fini de
    générateurs homogènes $f_i$ ($1\leq i\leq p$), donc
    (\hyperref[II.2.3.14-fr]{2.3.14}) l'espace sous-jacent $X$ est réunion des
    $D_+(f_i)=\operatorname{Spec}(S_{(f_i)})$, et tout revient à voir que
    chacun des $S_{(f_i)}$ est noethérien, ce qui résulte de
    (\hyperref[II.2.2.6-fr]{2.2.6}).
  \item[\rm(ii)] L'hypothèse entraîne que $S_0$ est une $A$-algèbre de type
    fini et $S$ une $S_0$-algèbre de type fini, donc $S_+$ est un idéal de
    type fini (\hyperref[II.2.1.4-fr]{2.1.4}). On est alors ramené, comme dans
    (i), à prouver que pour $f\in S_d$, $S_{(f)}$ est une $A$-algèbre de type
    fini. En vertu de (\hyperref[II.2.2.5-fr]{2.2.5}), il suffit de montrer
    que $S^{(d)}$ est une $A$-algèbre de type fini, ce qui résulte de
    (\hyperref[II.2.1.6-fr]{2.1.6}).
\end{enumerate}
\end{proof}

\begin{env}[2.7.2]
\label{II.2.7.2-fr}
Nous considérerons dans ce qui suit les conditions de finitude suivantes pour
un $S$-module gradué $M$~:
\begin{enumerate}
  \item[\rm(TF)] Il existe un entier $n$ tel que le sous-module
    $\bigoplus_{k\geq n}M_k$ soit un $S$-module de type fini.
  \item[\rm(TN)] Il existe un entier $n$ tel que $M_k=0$ pour $k\geq n$.
\end{enumerate}

Si $M$ vérifie (TN), on a $M_{(f)}=0$ pour tout $f$ homogène dans $S_+$,
donc $\widetilde{M}=0$.

Soient $M$, $N$ deux $S$-modules gradués~; on dit qu'un homomorphisme
$u:M\to N$ de degré $0$ est (TN)-\emph{injectif} (resp.
(TN)-\emph{surjectif}, (TN)-\emph{bijectif}) s'il existe un entier $n$ tel
que $u_k:M_k\to N_k$ soit injectif (resp. surjectif, bijectif) pour
$k\geq n$. Dire que $u$ est (TN)-injectif (resp. (TN)-surjectif) revient donc
à dire que $\operatorname{Ker}u$ (resp. $\operatorname{Coker}u$) vérifie
(TN). En vertu de (\hyperref[II.2.5.4-fr]{2.5.4}), si $u$ est (TN)-injectif
(resp. (TN)-surjectif, (TN)-bijectif), $\widetilde{u}$ est injectif (resp.
surjectif, bijectif)~; lorsque $u$ est (TN)-bijectif, on dit encore que $u$
est un (TN)-\emph{isomorphisme}.
\end{env}

\begin{proposition}[2.7.3]
\label{II.2.7.3-fr}
Soient $S$ un anneau gradué tel que l'idéal $S_+$ soit de type fini, $M$ un
$S$-module gradué.
\begin{enumerate}
  \item[\rm(i)] Si $M$ vérifie la condition (TF), le $\mathcal{O}_X$-Module
    $\widetilde{M}$ est de type fini.
  \item[\rm(ii)] Supposons que $M$ vérifie (TF)~; pour que $\widetilde{M}=0$,
    il faut et il suffit que $M$ vérifie (TN).
\end{enumerate}
\end{proposition}

\begin{proof}
On vient de voir que la condition (TN) implique $\widetilde{M}=0$. Si $M$
vérifie (TF), le sous-module gradué $M'=\bigoplus_{k\geq n}M_k$ qui est par
hypothèse de type fini, est tel que $M/M'$ satisfasse à (TN)~; on a par suite
$\widetilde{M/M'}=0$, et l'exactitude du foncteur
$M\mapsto\widetilde{M}$ (\hyperref[II.2.5.4-fr]{2.5.4}) entraîne
$\widetilde{M}=\widetilde{M'}$~; pour prouver que $\widetilde{M}$ est de type
fini, on est donc ramené au cas où $M$ est de type fini. La question étant
locale, il suffit de prouver que $M_{(f)}$ est un $S_{(f)}$-module de type
fini (\hyperref[I.1.3.9-fr]{I, 1.3.9})~; mais $M^{(d)}$ est un
$S^{(d)}$-module de type fini (\hyperref[II.2.1.6-fr]{2.1.6, (iii)}) et notre
assertion résulte alors de (\hyperref[II.2.2.5-fr]{2.2.5}).

Supposons maintenant que $M$ vérifie (TF) et que $\widetilde{M}=0$~; alors,
avec les mêmes notations que ci-dessus, on a $\widetilde{M'}=0$, et la
condition (TN) pour $M'$ est équivalente à la condition (TN) pour $M$, donc
pour prouver que $\widetilde{M}=0$ implique que $M$ vérifie (TN), on peut
encore se borner au cas où $M$ est engendré par un nombre fini d'éléments
homogènes $x_i$ ($1\leq i\leq p$)~; soit d'autre part
$(f_j)_{1\leq j\leq q}$ un système de générateurs homogènes de l'idéal $S_+$.
On a par hypothèse $M_{(f_j)}=0$ pour tout $j$, donc il existe un entier $n$
tel que $f_j^nx_i=0$ quels que soient $i$ et $j$. Soit $n_j$ le degré de
$f_j$, et soit $m$ la plus grande valeur de $\sum_jr_jn_j$ pour les systèmes
d'entiers $(r_j)$ en nombre fini tels que $\sum_jr_j\leq nq$~; il est clair
alors

\oldpage[II]{40}
que si $k>m$, on a $S_kx_i=0$ pour tout $i$~; si $h$ est le plus grand des
degrés des $x_i$, on en conclut que $M_k=0$ pour $k>h+m$, ce qui termine la
démonstration.
\end{proof}

\begin{corollary}[2.7.4]
\label{II.2.7.4-fr}
Soit $S$ un anneau gradué tel que l'idéal $S_+$ soit de type fini~; pour que
$X=\operatorname{Proj}(S)=\varnothing$, il faut et il suffit qu'il existe $n$
tel que $S_k=0$ pour $k\geq n$.
\end{corollary}

\begin{proof}
La condition $X=\varnothing$ est en effet équivalente à
$\mathcal{O}_X=\widetilde{S}=0$, et $S$ est un $S$-module monogène.
\end{proof}

\begin{theorem}[2.7.5]
\label{II.2.7.5-fr}
On suppose que l'idéal $S_+$ est engendré par un nombre fini d'éléments
homogènes de degré $1$~; soit $X=\operatorname{Proj}(S)$. Alors, pour tout
$\mathcal{O}_X$-Module quasi-cohérent $\mathcal{F}$, l'homomorphisme
canonique $\beta:\widetilde{\Gamma_*(\mathcal{F})}
\to\mathcal{F}$ (\hyperref[II.2.6.4-fr]{2.6.4}) est un isomorphisme.
\end{theorem}

\begin{proof}
En effet, si $S_+$ est engendré par un nombre fini d'éléments $f_i\in S_1$,
$X$ est réunion des sous-espaces $\operatorname{Spec}(S_{(f_i)})$
(\hyperref[II.2.3.6-fr]{2.3.6}) qui sont quasi-compacts, donc est
quasi-compact~; en outre $X$ est un schéma
(\hyperref[II.2.4.2-fr]{2.4.2})~; d'après
(\hyperref[I.9.3.2-fr]{I, 9.3.2}),
(\hyperref[II.2.5.14.2-fr]{2.5.14.2}) et
(\hyperref[II.2.6.3-fr]{2.6.3}) on a, pour tout $f\in S_d$ ($d>0$), un
isomorphisme canonique
$(\Gamma_*(\mathcal{F}))_{(\alpha_d(f))}\xrightarrow{\sim}
\Gamma(D_+(f),\mathcal{F})$~; d'ailleurs, par définition,
$(\Gamma_*(\mathcal{F}))_{(\alpha_d(f))}$ (où $\Gamma_*(\mathcal{F})$ est
considéré comme $\Gamma_*(\mathcal{O}_X)$-module) n'est autre que
$(\Gamma_*(\mathcal{F}))_{(f)}$ (où $\Gamma_*(\mathcal{F})$ est considéré
comme $S$-module)~; si on se reporte à la définition
(\hyperref[I.9.3.1-fr]{I, 9.3.1}) de l'isomorphisme canonique précédent, on
voit qu'il coïncide avec l'homomorphisme $\beta_f$, d'où le théorème.
\end{proof}

\begin{remark}[2.7.6]
\label{II.2.7.6-fr}
Si on suppose l'anneau gradué $S$ \emph{noethérien}, la condition de
(\hyperref[II.2.7.5-fr]{2.7.5}) est vérifiée \emph{ipso facto} dès que l'on
suppose l'idéal $S_+$ engendré par l'ensemble $S_1$ des éléments homogènes de
degré $1$.
\end{remark}

\begin{corollary}[2.7.7]
\label{II.2.7.7-fr}
Sous les hypothèses de (\hyperref[II.2.7.5-fr]{2.7.5}), tout
$\mathcal{O}_X$-Module quasi-cohérent $\mathcal{F}$ est isomorphe à un
$\mathcal{O}_X$-Module de la forme $\widetilde{M}$, où $M$ est un $S$-module
gradué.
\end{corollary}

\begin{corollary}[2.7.8]
\label{II.2.7.8-fr}
Sous les hypothèses de (\hyperref[II.2.7.5-fr]{2.7.5}), tout
$\mathcal{O}_X$-Module quasi-cohérent de type fini $\mathcal{F}$ est
isomorphe à un $\mathcal{O}_X$-Module de la forme $\widetilde{N}$, où $N$ est
un $S$-module gradué de type fini.
\end{corollary}

\begin{proof}
On peut supposer $\mathcal{F}=\widetilde{M}$, où $M$ est un $S$-module gradué
(\hyperref[II.2.7.7-fr]{2.7.7}). Soit $(f_\lambda)_{\lambda\in L}$ un
système de générateurs homogènes de $M$~; pour toute partie finie $H$ de $L$,
soit $M_H$ le sous-module gradué de $M$ engendré par les $f_\lambda$ tels que
$\lambda\in H$~; il est clair que $M$ est limite inductive de ses sous-modules
$M_H$, donc $\mathcal{F}$ est limite inductive de ses sous-$\mathcal{O}_X$-Modules
$\widetilde{M_H}$ (\hyperref[II.2.5.4-fr]{2.5.4}). Mais comme $\mathcal{F}$
est de type fini et l'espace sous-jacent $X$ quasi-compact, il résulte de
(\hyperref[0.5.2.3-fr]{0, 5.2.3}) que $\mathcal{F}=\widetilde{M_H}$ pour une
partie finie $H$ de $L$.
\end{proof}

\begin{corollary}[2.7.9]
\label{II.2.7.9-fr}
Sous les hypothèses de (\hyperref[II.2.7.5-fr]{2.7.5}), soit $\mathcal{F}$ un
$\mathcal{O}_X$-Module quasi-cohérent de type fini. Il existe alors un entier
$n_0$ tel que, pour tout $n\geq n_0$, $\mathcal{F}(n)$ soit isomorphe à un
quotient d'un $\mathcal{O}_X$-Module de la forme $\mathcal{O}_X^k$ ($k>0$
dépendant de $n$), et par suite est engendré par un nombre fini de ses sections
au-dessus de $X$ (\hyperref[0.5.1.1-fr]{0, 5.1.1}).
\end{corollary}

\begin{proof}
En vertu de (\hyperref[II.2.7.8-fr]{2.7.8}), on peut supposer que
$\mathcal{F}=\widetilde{M}$, où $M$ est quotient d'une somme directe finie de
$S$-modules de la forme $S(m_i)$~; en vertu de
(\hyperref[II.2.5.4-fr]{2.5.4}) on peut donc se borner au cas où $M=S(m)$,
donc $\mathcal{F}(n)=\widetilde{S(m+n)}
=\mathcal{O}_X(m+n)$ (\hyperref[II.2.5.15-fr]{2.5.15}). Il suffira donc de
démontrer le

\oldpage[II]{41}
\begin{lemma}[2.7.9.1]
\label{II.2.7.9.1-fr}
Sous les hypothèses de (\hyperref[II.2.7.5-fr]{2.7.5}), pour tout $n\geq0$,
il existe un entier $k$ (dépendant de $n$) et un homomorphisme surjectif
$\mathcal{O}_X^k\to\mathcal{O}_X(n)$.
\end{lemma}

Il suffit en effet (\hyperref[II.2.7.2-fr]{2.7.2}) de montrer que, pour un $k$
convenable, il y a un homomorphisme (TN)-surjectif $u$ \emph{de degré $0$},
du $S$-module gradué produit $S^k$ dans $S(n)$. Or, on a $(S(n))_0=S_n$, et
par hypothèse $S_h=S_1^h$ pour tout $h>0$, donc
$SS_n=S_n+S_{n+1}+\cdots$. Comme $S_n$ est un $S_0$-module de type fini
(\hyperref[II.2.1.5-fr]{2.1.5} et
\hyperref[II.2.1.6-fr]{2.1.6, (i)}), considérons un système de générateurs
$(a_i)_{1\leq i\leq k}$ de ce module~; l'homomorphisme $u$ fera correspondre
$a_i$ au $i$-ème élément de la base canonique de $S^k$ ($1\leq i\leq k$)~;
comme $\operatorname{Coker}u$ s'identifie alors à
$(S(n))_{-n}+\cdots+(S(n))_{-1}$, $u$ répond bien à la question.
\end{proof}

\begin{corollary}[2.7.10]
\label{II.2.7.10-fr}
Sous les hypothèses de (\hyperref[II.2.7.5-fr]{2.7.5}), soit $\mathcal{F}$ un
$\mathcal{O}_X$-Module quasi-cohérent de type fini. Il existe un entier $n_0$
tel que, pour tout $n\geq n_0$, $\mathcal{F}$ soit isomorphe à un quotient
d'un $\mathcal{O}_X$-Module de la forme $(\mathcal{O}_X(-n))^k$ ($k$
dépendant de $n$).
\end{corollary}

\begin{proposition}[2.7.11]
\label{II.2.7.11-fr}
Supposons vérifiées les hypothèses de
(\hyperref[II.2.7.5-fr]{2.7.5}) et soit $M$ un $S$-module gradué. Alors~:
\begin{enumerate}
  \item[\rm(i)] L'homomorphisme canonique
    $\widetilde{\alpha}:\widetilde{M}\to
    \widetilde{\Gamma_*(\widetilde{M})}$ est un isomorphisme.
  \item[\rm(ii)] Soit $\mathcal{G}$ un sous-$\mathcal{O}_X$-Module
    quasi-cohérent de $\widetilde{M}$, et soit $N$ le sous-$S$-module gradué
    de $M$, image réciproque de $\Gamma_*(\mathcal{G})$ par $\alpha$. Alors
    on a $\widetilde{N}=\mathcal{G}$ ($\widetilde{N}$ étant identifié, en
    vertu de (\hyperref[II.2.5.4-fr]{2.5.4}), à un
    sous-$\mathcal{O}_X$-Module de $\widetilde{M}$).
\end{enumerate}
\end{proposition}

\begin{proof}
Comme $\beta:\widetilde{\Gamma_*(\widetilde{M})}
\to\widetilde{M}$ est un isomorphisme en vertu de
(\hyperref[II.2.7.5-fr]{2.7.5}), $\widetilde{\alpha}$ est l'isomorphisme
réciproque en vertu de (\hyperref[II.2.6.5.1-fr]{2.6.5.1}), d'où (i). Soit
$P$ le sous-module gradué $\alpha(M)$ de $\Gamma_*(\widetilde{M})$~; comme
$M\mapsto\widetilde{M}$ est un foncteur exact
(\hyperref[II.2.5.4-fr]{2.5.4}), l'image de $\widetilde{M}$ par
$\widetilde{\alpha}$ est égale à $\widetilde{P}$, donc, en vertu de (i),
$\widetilde{P}=\widetilde{\Gamma_*(\widetilde{M})}$. Posons
$Q=\Gamma_*(\mathcal{G})\cap P$~; en vertu de ce qui précède et de
(\hyperref[II.2.5.4-fr]{2.5.4}), on a
$\widetilde{Q}=\widetilde{\Gamma_*(\mathcal{G})}$, donc la
restriction de $\beta$ à $\widetilde{Q}$ est un \emph{isomorphisme} de ce
$\mathcal{O}_X$-Module sur $\mathcal{G}$ par
(\hyperref[II.2.7.5-fr]{2.7.5}). Mais par définition de $N$ et
(\hyperref[II.2.5.4-fr]{2.5.4}), la restriction de l'isomorphisme
$\widetilde{\alpha}$ à $\widetilde{N}$ est un isomorphisme de $\widetilde{N}$
sur $\widetilde{Q}$, d'où la conclusion par
(\hyperref[II.2.6.5.1-fr]{2.6.5.1}).
\end{proof}

\subsection{Comportements fonctoriels.}
\label{subsection:II.2.8-fr}

\begin{env}[2.8.1]
\label{II.2.8.1-fr}
Soient $S$, $S'$ deux anneaux gradués à degrés positifs,
$\varphi:S'\to S$ un homomorphisme d'anneaux gradués. Nous désignerons par
$G(\varphi)$ la partie ouverte de $X=\operatorname{Proj}(S)$ complémentaire
de $V_+(\varphi(S'_+))$, ou, ce qui revient au même, la réunion des
$D_+(\varphi(f'))$, où $f'$ parcourt l'ensemble des éléments homogènes de
$S'_+$. La restriction à $G(\varphi)$ de l'application continue
${}^a\varphi$ de $\operatorname{Spec}(S)$ dans $\operatorname{Spec}(S')$
(\hyperref[I.1.2.1-fr]{I, 1.2.1}) est donc une application continue de
$G(\varphi)$ dans $\operatorname{Proj}(S')$, que nous noterons encore
${}^a\varphi$ par abus de langage. Si $f'\in S'_+$ est homogène, on a
\[
\label{II.2.8.1.1-fr}
  {}^a\varphi^{-1}(D_+(f'))=D_+(\varphi(f'))
\tag{2.8.1.1}
\]
compte tenu du fait que ${}^a\varphi$ applique $G(\varphi)$ dans
$\operatorname{Proj}(S')$ et de
(\hyperref[I.1.2.2.2-fr]{I, 1.2.2.2}). D'autre part, l'homomorphisme
$\varphi$ définit canoniquement (avec les mêmes notations) un homomorphisme
d'anneaux gradués $S'_{f'}\to S_f$, d'où, par restriction aux éléments de
degré $0$,

\oldpage[II]{42}
un homomorphisme $S'_{(f')}\to S_{(f)}$ que nous désignerons par
$\varphi_{(f)}$~; il lui correspond
(\hyperref[I.1.6.1-fr]{I, 1.6.1}) un morphisme
$({}^a\varphi_{(f)},\widetilde{\varphi}_{(f)}):
\operatorname{Spec}(S_{(f)})\to\operatorname{Spec}(S'_{(f')})$ de schémas
affines. Si on identifie canoniquement $\operatorname{Spec}(S_{(f)})$ au
schéma induit par $\operatorname{Proj}(S)$ sur $D_+(f)$
(\hyperref[II.2.3.6-fr]{2.3.6}), on a défini un morphisme
$\Phi_f:D_+(f)\to D_+(f')$ et ${}^a\varphi_{(f)}$ s'identifie à la
restriction de ${}^a\varphi$ à $D_+(f)$. D'autre part, il est immédiat que si
$g'$ est un second élément homogène de $S'_+$ et $g=\varphi(g')$, le
diagramme
\[
  \xymatrix{
    D_+(f)\ar[r]^{\Phi_f} & D_+(f')\\
    D_+(fg)\ar[u]\ar[r]_{\Phi_{fg}} & D_+(f'g')\ar[u]
  }
\]
est commutatif, en raison de la commutativité du diagramme
\[
  \xymatrix{
    S'_{(f')}\ar[r]^{\varphi_{(f)}}\ar[d]_{\omega_{f'g',f'}}
    & S_{(f)}\ar[d]^{\omega_{fg,f}}\\
    S'_{(f'g')}\ar[r]_{\varphi_{(fg)}} & S_{(fg)}
  }
\]
Compte tenu de la définition de $G(\varphi)$ et de
(\hyperref[II.2.3.3.2-fr]{2.3.3.2}), on voit donc que~:
\end{env}

\begin{proposition}[2.8.2]
\label{II.2.8.2-fr}
Étant donné un homomorphisme d'anneaux gradués $\varphi:S'\to S$, il existe
un morphisme et un seul $({}^a\varphi,\widetilde{\varphi})$ du préschéma
induit $G(\varphi)$ dans $\operatorname{Proj}(S')$ (dit \emph{associé à
$\varphi$} et noté $\operatorname{Proj}(\varphi)$), tel que pour tout élément
homogène $f'\in S'_+$, la restriction de ce morphisme à
$D_+(\varphi(f'))$ coïncide avec le morphisme associé à l'homomorphisme
$S'_{(f')}\to S_{(\varphi(f'))}$ correspondant à $\varphi$.
\end{proposition}

\begin{proof}
On notera qu'avec les notations précédentes, si $f'\in S'_d$, le diagramme
\[
\label{II.2.8.2.1-fr}
  \xymatrix{
    S'_{(f')}\ar[r]^{\varphi_{(f)}}\ar[d]_{\sim}
    & S_{(f)}\ar[d]^{\sim}\\
    {S'}^{(d)}/(f'-1){S'}^{(d)}\ar[r]
    & S^{(d)}/(f-1)S^{(d)}
  }
\tag{2.8.2.1}
\]
est commutatif (les flèches verticales étant les isomorphismes
(\hyperref[II.2.2.5-fr]{2.2.5})).
\end{proof}

\begin{corollary}[2.8.3]
\label{II.2.8.3-fr}
\begin{enumerate}
  \item[\rm(i)] Le morphisme $\operatorname{Proj}(\varphi)$ est affine.
  \item[\rm(ii)] Si $\operatorname{Ker}(\varphi)$ est nilpotent (et en
    particulier si $\varphi$ est injectif), le morphisme
    $\operatorname{Proj}(\varphi)$ est dominant.
\end{enumerate}
\end{corollary}

\begin{proof}
(i) est conséquence immédiate de (\hyperref[II.2.8.2-fr]{2.8.2}) et de
(\hyperref[II.2.8.1.1-fr]{2.8.1.1}). D'autre part, si
$\operatorname{Ker}(\varphi)$ est nilpotent, pour tout $f'$ homogène dans
$S'_+$, on vérifie aussitôt que $\operatorname{Ker}(\varphi_f)$ (avec
$f=\varphi(f')$) est nilpotent, donc aussi
$\operatorname{Ker}(\varphi_{(f)})$~; la conclusion résulte donc de
(\hyperref[II.2.8.2-fr]{2.8.2}) et de
(\hyperref[I.1.2.7-fr]{I, 1.2.7}).
\end{proof}

\oldpage[II]{43}
On notera qu'il y a en général des morphismes de $\operatorname{Proj}(S)$
dans $\operatorname{Proj}(S')$ qui ne sont pas affines, et ne proviennent
donc pas d'homomorphismes d'anneaux gradués $S'\to S$~; un exemple est le
morphisme structural $\operatorname{Proj}(S)\to\operatorname{Spec}(A)$ quand
$A$ est un corps ($\operatorname{Spec}(A)$ étant identifié à
$\operatorname{Proj}(A[T])$ (cf. (\hyperref[II.3.1.7-fr]{3.1.7})))~; cela
résulte en effet de (\hyperref[I.2.3.2-fr]{I, 2.3.2}).

\begin{env}[2.8.4]
\label{II.2.8.4-fr}
Soient $S''$ un troisième anneau gradué à degrés positifs,
$\varphi':S''\to S'$ un homomorphisme d'anneaux gradués, et posons
$\varphi''=\varphi\circ\varphi'$. En raison de
(\hyperref[II.2.8.1.1-fr]{2.8.1.1}) et de la formule
${}^a\varphi''=({}^a\varphi')\circ({}^a\varphi)$ on vérifie aussitôt que l'on
a $G(\varphi'')\subset G(\varphi)$ et que, si $\Phi$, $\Phi'$ et $\Phi''$
sont les morphismes associés à $\varphi$, $\varphi'$ et $\varphi''$, on a
$\Phi''=\Phi'\circ(\Phi|G(\varphi''))$.
\end{env}

\begin{env}[2.8.5]
\label{II.2.8.5-fr}
Supposons que $S$ (resp. $S'$) soit une $A$-algèbre graduée (resp. une
$A'$-algèbre graduée), et soit $\psi:A'\to A$ un homomorphisme d'anneaux tel
que le diagramme
\[
\label{II.2.8.5.1-fr}
  \xymatrix{
    A'\ar[r]^{\psi}\ar[d] & A\ar[d]\\
    S'\ar[r]_{\varphi} & S
  }
\tag{2.8.5.1}
\]
soit commutatif. On peut alors considérer $G(\varphi)$ et
$\operatorname{Proj}(S')$ comme des schémas au-dessus de
$\operatorname{Spec}(A)$ et $\operatorname{Spec}(A')$ respectivement~; si
$\Phi$ et $\Psi$ sont les morphismes associés respectivement à $\varphi$ et
$\psi$, le diagramme
\[
  \xymatrix{
    G(\varphi)\ar[r]^{\Phi}\ar[d] & \operatorname{Proj}(S')\ar[d]\\
    \operatorname{Spec}(A)\ar[r]_{\Psi} & \operatorname{Spec}(A')
  }
\]
est alors commutatif~: il suffit en effet de le démontrer pour la restriction
de $\Phi$ à $D_+(f)$, où $f=\varphi(f')$, $f'$ homogène dans $S'_+$~; et
alors cela résulte de la commutativité du diagramme
\[
  \xymatrix{
    A'\ar[r]^{\psi}\ar[d] & A\ar[d]\\
    S'_{(f')}\ar[r]_{\varphi_{(f)}} & S_{(f)}
  }
\]
\end{env}

\begin{env}[2.8.6]
\label{II.2.8.6-fr}
Soit maintenant $M$ un $S$-module gradué, et considérons le $S'$-module
$M_{[\varphi]}$, qui est évidemment gradué. Soit $f'$ homogène dans $S'_+$,
et soit $f=\varphi(f')$~; on sait
(\hyperref[0.1.5.2-fr]{0, 1.5.2}) qu'il y a un isomorphisme canonique
$(M_{[\varphi]})_{f'}\xrightarrow{\sim}(M_f)_{[\varphi_f]}$, et il est
immédiat que cet isomorphisme conserve les degrés, donc il y a un isomorphisme
canonique
$(M_{[\varphi]})_{(f')}\xrightarrow{\sim}
(M_{(f)})_{[\varphi_{(f)}]}$. Il correspond canoniquement à cet isomorphisme
un isomorphisme de faisceaux
$\widetilde{M_{[\varphi]}}|D_+(f')
\xrightarrow{\sim}(\Phi_f)_*(\widetilde{M}|D_+(f))$
(\hyperref[II.2.5.2-fr]{2.5.2} et
\hyperref[I.1.6.3-fr]{I, 1.6.3}). En outre,

\oldpage[II]{44}
si $g'$ est un second élément homogène de $S'_+$ et $g=\varphi(g')$, le
diagramme
\[
  \xymatrix{
    (M_{[\varphi]})_{(f')}\ar[r]^{\sim}\ar[d]
    & (M_{(f)})_{[\varphi_{(f)}]}\ar[d]\\
    (M_{[\varphi]})_{(f'g')}\ar[r]^{\sim}
    & (M_{(fg)})_{[\varphi_{(fg)}]}
  }
\]
est commutatif, d'où on conclut aussitôt que l'isomorphisme
\[
  \widetilde{M_{[\varphi]}}|D_+(f'g')
  \xrightarrow{\sim}(\Phi_{fg})_*(\widetilde{M}|D_+(fg))
\]
est la restriction à $D_+(f'g')$ de l'isomorphisme
$\widetilde{M_{[\varphi]}}|D_+(f')
\xrightarrow{\sim}(\Phi_f)_*(\widetilde{M}|D_+(f))$. Comme $\Phi_f$ est la
restriction à $D_+(f)$ du morphisme $\Phi$, on voit donc, compte tenu de
(\hyperref[II.2.8.1.1-fr]{2.8.1.1}), et en posant
$X'=\operatorname{Proj}(S')$~:
\end{env}

\begin{proposition}[2.8.7]
\label{II.2.8.7-fr}
Il existe un isomorphisme canonique fonctoriel du $\mathcal{O}_{X'}$-Module
$\widetilde{M_{[\varphi]}}$ sur le
$\mathcal{O}_{X'}$-Module $\Phi_*(\widetilde{M}|G(\varphi))$.
\end{proposition}

On en déduit aussitôt une application canonique fonctorielle de l'ensemble
des $\varphi$-morphismes $M'\to M$ d'un $S'$-module gradué dans le $S$-module
gradué $M$, dans l'ensemble des $\Phi$-morphismes
$\widetilde{M'}\to\widetilde{M}|G(\varphi)$. Avec les notations de
(\hyperref[II.2.8.4-fr]{2.8.4}), si $M''$ est un $S''$-module gradué, au
composé d'un $\varphi$-morphisme $M'\to M$ et d'un $\varphi'$-morphisme
$M''\to M'$ correspond canoniquement le composé de
$\widetilde{M'}|G(\varphi')\to\widetilde{M}|G(\varphi'')$ et de
$\widetilde{M''}\to\widetilde{M'}|G(\varphi')$.

\begin{proposition}[2.8.8]
\label{II.2.8.8-fr}
Sous les hypothèses de (\hyperref[II.2.8.1-fr]{2.8.1}), soit $M'$ un
$S'$-module gradué. Il existe un homomorphisme canonique fonctoriel $\nu$ du
$(\mathcal{O}_X|G(\varphi))$-Module $\Phi^*(\widetilde{M'})$ dans le
$(\mathcal{O}_X|G(\varphi))$-Module
$\widetilde{M'\otimes_{S'}S}|G(\varphi)$. Si l'idéal $S'_+$
est engendré par $S'_1$, $\nu$ est un isomorphisme.
\end{proposition}

\begin{proof}
En effet, pour $f'\in S'_d$ ($d>0$), on définit un homomorphisme canonique
fonctoriel de $S_{(f)}$-modules (où $f=\varphi(f')$)
\[
\label{II.2.8.8.1-fr}
  \nu_f:M'_{(f')}\otimes_{S'_{(f')}}S_{(f)}
  \longrightarrow(M'\otimes_{S'}S)_{(f)}
\tag{2.8.8.1}
\]
en composant l'homomorphisme
$M'_{(f')}\otimes_{S'_{(f')}}S_{(f)}\to
M'_{f'}\otimes_{S'_{f'}}S_f$ et l'isomorphisme canonique
$M'_{f'}\otimes_{S'_{f'}}S_f\xrightarrow{\sim}(M'\otimes_{S'}S)_f$
(\hyperref[0.1.5.4-fr]{0, 1.5.4}) et notant que ce dernier conserve les
degrés. On vérifie aussitôt la compatibilité des $\nu_f$ avec les opérateurs
de restriction de $D_+(f)$ à $D_+(fg)$ (pour un second $g'\in S'_+$ et
$g=\varphi(g')$), d'où la définition de l'homomorphisme
\[
  \nu:\Phi^*(\widetilde{M'})
  \longrightarrow\widetilde{M'\otimes_{S'}S}|G(\varphi)
\]
compte tenu de (\hyperref[I.1.6.5-fr]{I, 1.6.5}). Pour prouver la seconde
assertion, il suffit de montrer que $\nu_f$ est un isomorphisme pour tout
$f'\in S'_1$, puisque $G(\varphi)$ est alors réunion des
$D_+(\varphi(f'))$. On définit d'abord une application $\mathbf{Z}$-bilinéaire
$M'_m\times S_n\to M'_{(f')}\otimes_{S'_{(f')}}S_{(f)}$ en faisant
correspondre à $(x',s)$ l'élément
$(x'/{f'}^m)\otimes(s/f^n)$ (avec la convention que $x'/{f'}^m$ est
${f'}^{-m}x'/1$

\oldpage[II]{45}
lorsque $m<0$). On constate comme dans la démonstration de
(\hyperref[II.2.5.13-fr]{2.5.13}) que cette application donne naissance à un
di-homomorphisme de modules
\[
  \eta_f:M'\otimes_{S'}S
  \longrightarrow M'_{(f')}\otimes_{S'_{(f')}}S_{(f)}.
\]
En outre, si, pour $r>0$, on a
$f^r\sum_i(x'_i\otimes s_i)=0$, cela s'écrit aussi
$\sum_i({f'}^rx'_i\otimes s_i)=0$, d'où, par
(\hyperref[0.1.5.4-fr]{0, 1.5.4}),
$\sum_i({f'}^rx'_i/{f'}^{m_i+r})\otimes(s_i/f^{n_i})=0$, c'est-à-dire
$\eta_f(\sum_i x'_i\otimes s_i)=0$, ce qui prouve que $\eta_f$ se factorise
en
\[
  M'\otimes_{S'}S\longrightarrow(M'\otimes_{S'}S)_f
  \xrightarrow{\eta'_f}M'_{(f')}\otimes_{S'_{(f')}}S_{(f)}~;
\]
on vérifie enfin que $\eta'_f$ et $\nu_f$ sont des isomorphismes réciproques
l'un de l'autre.

En particulier, il résulte de (\hyperref[II.2.1.2.1-fr]{2.1.2.1}) que l'on a
un homomorphisme canonique
\[
\label{II.2.8.8.2-fr}
  \Phi^*(\mathcal{O}_{X'}(n))\xrightarrow{\sim}
  \mathcal{O}_X(n)|G(\varphi)
\tag{2.8.8.2}
\]
pour tout $n\in\mathbf{Z}$.
\end{proof}

\begin{env}[2.8.9]
\label{II.2.8.9-fr}
Soient $A$, $A'$ deux anneaux, $\psi:A'\to A$ un homomorphisme d'anneaux,
définissant un morphisme
$\Psi:\operatorname{Spec}(A)\to\operatorname{Spec}(A')$. Soit $S'$ une
$A'$-algèbre graduée à degrés positifs, et posons $S=S'\otimes_{A'}A$, qui
est de façon évidente une $A$-algèbre graduée par les $S'_n\otimes_{A'}A$~;
l'application $\varphi:s'\mapsto s'\otimes1$ est alors un homomorphisme
d'anneaux gradués rendant commutatif le diagramme
(\hyperref[II.2.8.5.1-fr]{2.8.5.1}). Comme ici $S_+$ est le $A$-module
engendré par $\varphi(S'_+)$, on a $G(\varphi)=\operatorname{Proj}(S)=X$~;
d'où, en posant $X'=\operatorname{Proj}(S')$, un diagramme commutatif
\[
\label{II.2.8.9.1-fr}
  \xymatrix{
    X\ar[r]^{\Phi}\ar[d]_p & X'\ar[d]\\
    Y\ar[r]_{\Psi} & Y'
  }
\tag{2.8.9.1}
\]

Soit maintenant $M'$ un $S'$-module gradué, et posons
$M=M'\otimes_{A'}A=M'\otimes_{S'}S$. Dans ces conditions~:
\end{env}

\begin{proposition}[2.8.10]
\label{II.2.8.10-fr}
Le diagramme (\hyperref[II.2.8.9.1-fr]{2.8.9.1}) identifie le schéma $X$ au
produit $X'\times_{Y'}Y$~; en outre, l'homomorphisme canonique
$\nu:\Phi^*(\widetilde{M'})\to\widetilde{M}$
(\hyperref[II.2.8.8-fr]{2.8.8}) est un isomorphisme.
\end{proposition}

\begin{proof}
La première assertion sera démontrée si l'on prouve que pour tout $f'$
homogène dans $S'_+$, en posant $f=\varphi(f')$, les restrictions de $\Phi$
et de $p$ à $D_+(f)$ identifient ce schéma au produit
$D_+(f')\times_{Y'}Y$ (\hyperref[I.3.2.6.2-fr]{I, 3.2.6.2})~; autrement
dit, il suffit (\hyperref[I.3.2.2-fr]{I, 3.2.2}) de prouver que $S_{(f)}$
s'identifie canoniquement à $S'_{(f')}\otimes_{A'}A$, ce qui est immédiat
en vertu de l'existence de l'isomorphisme canonique
$S_f\xrightarrow{\sim}S'_{f'}\otimes_{A'}A$ conservant les degrés
(\hyperref[0.1.5.4-fr]{0, 1.5.4}). La seconde assertion résulte de ce que
$M'_{(f')}\otimes_{S'_{(f')}}S_{(f)}$ s'identifie d'après ce qui précède à
$M'_{(f')}\otimes_{A'}A$, et ce dernier à $M_{(f)}$, puisque $M_f$
s'identifie canoniquement à $M'_{f'}\otimes_{A'}A$ par un isomorphisme
conservant les degrés.
\end{proof}

\begin{corollary}[2.8.11]
\label{II.2.8.11-fr}
Pour tout entier $n\in\mathbf{Z}$, $\widetilde{M(n)}$ s'identifie à
$\Phi^*(\widetilde{M'(n)})=\widetilde{M'(n)}\otimes_{Y'}\mathcal{O}_Y$~;
en particulier
$\mathcal{O}_X(n)=\Phi^*(\mathcal{O}_{X'}(n))
=\mathcal{O}_{X'}(n)\otimes_{Y'}\mathcal{O}_Y$.
\end{corollary}

\begin{proof}
Cela résulte de (\hyperref[II.2.8.10-fr]{2.8.10}) et de
(\hyperref[II.2.5.15-fr]{2.5.15}).
\end{proof}

\oldpage[II]{46}

\begin{env}[2.8.12]
\label{II.2.8.12-fr}
Sous les hypothèses de (\hyperref[II.2.8.9-fr]{2.8.9}), pour
$f'\in S'_d$ ($d>0$) et $f=\varphi(f')$, le diagramme
\[
\label{II.2.8.12.1-fr}
  \xymatrix{
    M'_{(f')}\ar[r]^{\sim}\ar[d]
    & M'^{(d)}/(f'-1)M'^{(d)}\ar[d]\\
    M_{(f)}\ar[r]^{\sim}
    & M^{(d)}/(f-1)M^{(d)}
  }
\tag{2.8.12.1}
\]
(cf. (\hyperref[II.2.2.5-fr]{2.2.5})) est commutatif.
\end{env}

\begin{env}[2.8.13]
\label{II.2.8.13-fr}
Gardons les notations et hypothèses de (\hyperref[II.2.8.9-fr]{2.8.9}), et
soit $\mathcal{F}'$ un $\mathcal{O}_{X'}$-Module~; si on pose
$\mathcal{F}=\Phi^*(\mathcal{F}')$, on a, pour tout $n\in\mathbf{Z}$,
$\mathcal{F}(n)=\Phi^*(\mathcal{F}'(n))$ en vertu de
(\hyperref[II.2.8.11-fr]{2.8.11}) et
(\hyperref[0.4.3.3-fr]{0, 4.3.3}). Par suite
(\hyperref[0.3.7.1-fr]{0, 3.7.1}), on a un homomorphisme canonique
\[
  \Gamma(\rho):\Gamma(X',\mathcal{F}'(n))
  \longrightarrow\Gamma(X,\mathcal{F}(n))
\]
ce qui donne un di-homomorphisme canonique de modules gradués
\[
  \Gamma_\bullet(\mathcal{F}')\longrightarrow
  \Gamma_\bullet(\mathcal{F}).
\]

Supposons l'idéal $S_+$ engendré par $S_1$, et
$\mathcal{F}'=\widetilde{M'}$, donc $\mathcal{F}=\widetilde{M}$ avec
$M=M'\otimes_{A'}A$. Si $f'$ est homogène dans $S'_+$ et
$f=\varphi(f')$, on a vu que
$M_{(f)}=M'_{(f')}\otimes_{A'}A$, et le diagramme
\[
  \xymatrix{
    M'_0\ar[r]\ar[d]
    & M'_{(f')}\ar[d]
    & =\Gamma(D_+(f'),\widetilde{M'})\\
    M_0\ar[r]
    & M_{(f)}
    & =\Gamma(D_+(f),\widetilde{M})
  }
\]
est donc commutatif~; on conclut aussitôt de cette remarque et de la
définition de l'homomorphisme $\alpha$
(\hyperref[II.2.6.2-fr]{2.6.2}) que le diagramme
\[
\label{II.2.8.13.1-fr}
  \xymatrix{
    M'\ar[r]^{\alpha_{M'}}\ar[d]
    & \Gamma_\bullet(\widetilde{M'})\ar[d]\\
    M\ar[r]_{\alpha_M}
    & \Gamma_\bullet(\widetilde{M})
  }
\tag{2.8.13.1}
\]
est commutatif. De même, le diagramme
\[
\label{II.2.8.13.2-fr}
  \xymatrix{
    \widetilde{\Gamma_\bullet(\mathcal{F}')}\ar[r]^{\beta_{\mathcal{F}'}}
      \ar[d]
    & \mathcal{F}'\ar[d]\\
    \widetilde{\Gamma_\bullet(\mathcal{F})}\ar[r]_{\beta_{\mathcal{F}}}
    & \mathcal{F}
  }
\tag{2.8.13.2}
\]
est commutatif (la flèche verticale de droite étant le $\Phi$-morphisme
canonique
$\mathcal{F}'\to\Phi^*(\mathcal{F}')=\mathcal{F}$).
\end{env}

\oldpage[II]{47}

\begin{env}[2.8.14]
\label{II.2.8.14-fr}
Conservant toujours les hypothèses et notations de
(\hyperref[II.2.8.9-fr]{2.8.9}), soit $N'$ un second $S'$-module gradué, et
soit $N=N'\otimes_{A'}A$. Il est immédiat que les di-homomorphismes
canoniques $M'\to M$, $N'\to N$ donnent un di-homomorphisme
$M'\otimes_{S'}N'\to M\otimes_S N$ (relatif à l'homomorphisme canonique
d'anneaux $S'\to S$), et par suite aussi un $S$-homomorphisme de degré~$0$
\[
  (M'\otimes_{S'}N')\otimes_{A'}A\longrightarrow M\otimes_S N,
\]
auquel correspond (compte tenu de
(\hyperref[II.2.8.10-fr]{2.8.10})) un homomorphisme de
$\mathcal{O}_X$-Modules
\[
\label{II.2.8.14.1-fr}
  \Phi^*(\widetilde{M'\otimes_{S'}N'})
  \longrightarrow\widetilde{M\otimes_S N}.
\tag{2.8.14.1}
\]

En outre, on vérifie aussitôt que le diagramme
\[
\label{II.2.8.14.2-fr}
  \xymatrix{
    \Phi^*(\widetilde{M'}\otimes_{\mathcal{O}_{X'}}\widetilde{N'})
      \ar[r]^{\sim}\ar[d]_{\Phi^*(\lambda)}
    & \widetilde{M}\otimes_{\mathcal{O}_X}\widetilde{N}
      \ar[d]^{\lambda}
    & =\Phi^*(\widetilde{M'})\otimes_{\mathcal{O}_X}
      \Phi^*(\widetilde{N'})\\
    \Phi^*(\widetilde{M'\otimes_{S'}N'})\ar[r]
    & \widetilde{M\otimes_S N}
  }
\tag{2.8.14.2}
\]
est commutatif, la première ligne étant l'isomorphisme canonique
(\hyperref[0.4.3.3-fr]{0, 4.3.3}). Si l'idéal $S'_+$ est engendré par
$S'_1$, il est clair que $S_+$ est engendré par $S_1$, et les deux flèches
verticales de (\hyperref[II.2.8.14.2-fr]{2.8.14.2}) sont alors des
isomorphismes (\hyperref[II.2.5.13-fr]{2.5.13})~; il en est donc de même de
(\hyperref[II.2.8.14.1-fr]{2.8.14.1}).

On a de même un di-homomorphisme canonique
$\operatorname{Hom}_{S'}(M',N')\to\operatorname{Hom}_S(M,N)$ en faisant
correspondre à un homomorphisme $u'$ de degré $k$ l'homomorphisme
$u'\otimes1$, qui est aussi de degré $k$~; on en déduit encore un
$S$-homomorphisme de degré~$0$
\[
  (\operatorname{Hom}_{S'}(M',N'))\otimes_{A'}A
  \longrightarrow\operatorname{Hom}_S(M,N)
\]
d'où un homomorphisme de $\mathcal{O}_X$-Modules
\[
\label{II.2.8.14.3-fr}
  \Phi^*(\widetilde{\operatorname{Hom}_{S'}(M',N')})
  \longrightarrow\widetilde{\operatorname{Hom}_S(M,N)}.
\tag{2.8.14.3}
\]

En outre, le diagramme
\[
  \xymatrix{
    \Phi^*(\widetilde{\operatorname{Hom}_{S'}(M',N')})\ar[r]
      \ar[d]_{\Phi^*(\mu)}
    & \widetilde{\operatorname{Hom}_S(M,N)}\ar[d]^{\mu}\\
    \Phi^*(\mathcal{H}om_{\mathcal{O}_{X'}}
      (\widetilde{M'},\widetilde{N'}))\ar[r]
    & \mathcal{H}om_{\mathcal{O}_X}(\widetilde{M},\widetilde{N})
  }
\]
est commutatif (la seconde ligne horizontale étant l'homomorphisme canonique
(\hyperref[0.4.4.6-fr]{0, 4.4.6})).
\end{env}

\oldpage[II]{48}

\begin{env}[2.8.15]
\label{II.2.8.15-fr}
Les notations et hypothèses étant celles de
(\hyperref[II.2.8.1-fr]{2.8.1}), il résulte de
(\hyperref[II.2.4.7-fr]{2.4.7}) que l'on ne change pas le morphisme $\Phi$,
à des isomorphismes près, en remplaçant $S_0$ et $S'_0$ par
$\mathbf{Z}$ et $\varphi_0$ par l'application identique, puis en remplaçant
$S$ et $S'$ par $S^{(d)}$ et ${S'}^{(d)}$ respectivement ($d>0$) et
$\varphi$ par sa restriction $\varphi^{(d)}$ à $S^{(d)}$.
\end{env}

\subsection{Sous-schémas fermés d'un schéma $\operatorname{Proj}(S)$.}
\label{subsection:II.2.9-fr}

\begin{env}[2.9.1]
\label{II.2.9.1-fr}
Si $\varphi:S\to S'$ est un homomorphisme d'anneaux gradués, on dit que
$\varphi$ est (TN)-\emph{surjectif} (resp. (TN)-\emph{injectif},
(TN)-\emph{bijectif}) s'il existe un entier $n$ tel que, pour $k\geq n$,
$\varphi_k:S_k\to S'_k$ soit surjectif (resp. injectif, bijectif). Alors il
résulte de (\hyperref[II.2.8.15-fr]{2.8.15}) que l'étude de $\Phi$ peut se
ramener au cas où $\varphi$ est \emph{surjectif} (resp. \emph{injectif},
\emph{bijectif}). Au lieu de dire que $\varphi$ est (TN)-bijectif, on dit aussi
que c'est alors un (TN)-\emph{isomorphisme}.
\end{env}

\begin{proposition}[2.9.2]
\label{II.2.9.2-fr}
Soit $S$ un anneau gradué à degrés positifs et soit
$X=\operatorname{Proj}(S)$.
\begin{enumerate}
  \item[\rm(i)] Si $\varphi:S\to S'$ est un homomorphisme (TN)-surjectif
    d'anneaux gradués, le morphisme correspondant $\Phi$
    (\hyperref[II.2.8.1-fr]{2.8.1}) est défini dans
    $\operatorname{Proj}(S')$ tout entier et est une immersion fermée de
    $\operatorname{Proj}(S')$ dans $X$. Si $\mathfrak{J}$ est le noyau de
    $\varphi$, le sous-schéma fermé de $X$ associé à $\Phi$ est défini par
    le faisceau d'idéaux quasi-cohérent $\widetilde{\mathfrak{J}}$ de
    $\mathcal{O}_X$.
  \item[\rm(ii)] Supposons en outre l'idéal $S_+$ engendré par un nombre
    fini d'éléments homogènes de degré~$1$. Soit $X'$ un sous-schéma fermé
    de $X$, défini par un faisceau quasi-cohérent d'idéaux $\mathcal{J}$ de
    $\mathcal{O}_X$. Soit $\mathfrak{J}$ l'idéal gradué de $S$, image
    réciproque de $\Gamma_\bullet(\mathcal{J})$ par l'homomorphisme canonique
    $\alpha:S\to\Gamma_\bullet(\mathcal{O}_X)$
    (\hyperref[II.2.6.2-fr]{2.6.2}), et posons
    $S'=S/\mathfrak{J}$. Alors $X'$ est le sous-schéma associé à
    l'immersion fermée $\operatorname{Proj}(S')\to X$ correspondant à
    l'homomorphisme canonique d'anneaux gradués $S\to S'$.
\end{enumerate}
\end{proposition}

\begin{proof}
\begin{enumerate}
  \item[\rm(i)] On peut supposer $\varphi$ surjectif
    (\hyperref[II.2.9.1-fr]{2.9.1}). Comme par hypothèse
    $\varphi(S_+)$ engendre $S'_+$, on a
    $G(\varphi)=\operatorname{Proj}(S')$. D'autre part, la seconde assertion
    se vérifie localement sur $X$~; soit donc $f$ un élément homogène de
    $S_+$ et posons $f'=\varphi(f)$. Comme $\varphi$ est un homomorphisme
    surjectif d'anneaux gradués, on vérifie aussitôt que
    $\varphi_{(f')}:S_{(f)}\to S'_{(f')}$ est surjectif et que son noyau est
    $\mathfrak{J}_{(f)}$, ce qui achève de prouver (i)
    (\hyperref[I.4.2.3-fr]{I, 4.2.3}).
  \item[\rm(ii)] En vertu de (i), on est ramené à vérifier que
    l'homomorphisme $\widetilde{j}:\widetilde{\mathfrak{J}}\to\mathcal{O}_X$
    déduit de l'injection canonique $j:\mathfrak{J}\to S$ est un
    isomorphisme de $\widetilde{\mathfrak{J}}$ sur $\mathcal{J}$, ce qui
    résulte de (\hyperref[II.2.7.11-fr]{2.7.11}).
\end{enumerate}
\end{proof}

On notera que $\mathfrak{J}$ est le \emph{plus grand} des idéaux gradués
$\mathfrak{J}'$ de $S$ tels que
$\widetilde{j}(\widetilde{\mathfrak{J}'})=\mathcal{J}$, car on vérifie
immédiatement en remontant aux définitions
(\hyperref[II.2.6.2-fr]{2.6.2}) que cette relation implique
$\alpha(\mathfrak{J}')\subset\Gamma_\bullet(\mathcal{J})$.

\begin{corollary}[2.9.3]
\label{II.2.9.3-fr}
On suppose vérifiées les hypothèses de
(\hyperref[II.2.9.2-fr]{2.9.2}, (i)) et en outre que l'idéal $S_+$ est
engendré par $S_1$~; alors $\Phi^*(\widetilde{S(n)})$ est canoniquement
isomorphe à $\widetilde{S'(n)}$ pour tout $n\in\mathbf{Z}$, et par suite
$\Phi^*(\mathcal{F}(n))$ à $\Phi^*(\mathcal{F})(n)$ pour tout
$\mathcal{O}_X$-Module $\mathcal{F}$.
\end{corollary}

\begin{proof}
C'est un cas particulier de (\hyperref[II.2.8.8-fr]{2.8.8}), compte tenu de
(\hyperref[II.2.5.10.2-fr]{2.5.10.2}).
\end{proof}

\begin{corollary}[2.9.4]
\label{II.2.9.4-fr}
On suppose vérifiées les hypothèses de
(\hyperref[II.2.9.2-fr]{2.9.2}, (ii)). Pour que le sous-préschéma fermé $X'$
de $X$ soit intègre, il faut et il suffit que l'idéal gradué
$\mathfrak{J}$ soit premier dans $S$.
\end{corollary}

\oldpage[II]{49}

\begin{proof}
Comme $X'$ est isomorphe à $\operatorname{Proj}(S/\mathfrak{J})$, la
condition est suffisante en vertu de (\hyperref[II.2.4.4-fr]{2.4.4}). Pour
voir qu'elle est nécessaire, considérons la suite exacte
$0\to\mathcal{J}\to\mathcal{O}_X\to\mathcal{O}_X/\mathcal{J}$, qui donne une
suite exacte
\[
  0\longrightarrow\Gamma_\bullet(\mathcal{J})
  \longrightarrow\Gamma_\bullet(\mathcal{O}_X)
  \longrightarrow\Gamma_\bullet(\mathcal{O}_X/\mathcal{J}).
\]
Il suffit de prouver que si $f\in S_m$, $g\in S_n$ sont tels que l'image dans
$\Gamma_\bullet(\mathcal{O}_X/\mathcal{J})$ de $\alpha_{n+m}(fg)$ est nulle,
alors l'une des images de $\alpha_m(f)$, $\alpha_n(g)$ est nulle. Or, par
définition, ces images sont des sections des
$(\mathcal{O}_X/\mathcal{J})$-Modules inversibles
$\mathcal{L}=(\mathcal{O}_X/\mathcal{J})(m)$ et
$\mathcal{L}'=(\mathcal{O}_X/\mathcal{J})(n)$ au-dessus du schéma intègre
$X'$~; l'hypothèse entraîne que le produit de ces deux sections est nul dans
$\mathcal{L}\otimes\mathcal{L}'$
((\hyperref[II.2.9.3-fr]{2.9.3}) et
(\hyperref[II.2.5.14.1-fr]{2.5.14.1})), donc l'une d'elles est nulle en vertu
de (\hyperref[I.7.4.4-fr]{I, 7.4.4}).
\end{proof}

\begin{corollary}[2.9.5]
\label{II.2.9.5-fr}
Soient $A$ un anneau, $M$ un $A$-module, $S$ une $A$-algèbre graduée
engendrée par l'ensemble $S_1$ des éléments homogènes de degré~$1$,
$u:M\to S_1$ un homomorphisme surjectif de $A$-modules,
$\overline{u}:\mathbf{S}(M)\to S$ l'homomorphisme (de $A$-algèbres) de
l'algèbre symétrique $\mathbf{S}(M)$ de $M$ dans $S$ qui prolonge $u$. Le
morphisme correspondant à $\overline{u}$ est alors une immersion fermée de
$\operatorname{Proj}(S)$ dans $\operatorname{Proj}(\mathbf{S}(M))$.
\end{corollary}

\begin{proof}
En effet, $\overline{u}$ est surjectif par hypothèse et il suffit d'appliquer
(\hyperref[II.2.9.2-fr]{2.9.2}).
\end{proof}

\section{Spectre homogène d'un faisceau d'algèbres graduées}
\label{section:II.3-fr}

\subsection{Spectre homogène d'une $\mathcal{O}_Y$-Algèbre graduée
quasi-cohérente.}
\label{subsection:II.3.1-fr}
\label{II.3.1-fr}

\begin{env}[3.1.1]
\label{II.3.1.1-fr}
Soient $Y$ un préschéma, $\mathcal{S}$ une $\mathcal{O}_Y$-Algèbre graduée,
$\mathcal{M}$ un $\mathcal{S}$-Module gradué. Si $\mathcal{S}$ est
\emph{quasi-cohérente}, chacun de ses composants homogènes $\mathcal{S}_n$ est
un $\mathcal{O}_Y$-Module \emph{quasi-cohérent}, étant l'image de
$\mathcal{S}$ par un homomorphisme de $\mathcal{S}$ dans lui-même
((\hyperref[I.1.3.8-fr]{I, 1.3.8}) et
(\hyperref[I.1.3.9-fr]{I, 1.3.9}))~; de même, si $\mathcal{M}$ est
quasi-cohérent en tant que $\mathcal{O}_Y$-Module, ses composants homogènes
$\mathcal{M}_n$ le sont aussi, et réciproquement. Si $d$ est un entier $>0$,
on désignera par $\mathcal{S}^{(d)}$ la somme directe des
$\mathcal{O}_Y$-Modules $\mathcal{S}_{nd}$ ($n\in\mathbf{Z}$), qui est
quasi-cohérente si $\mathcal{S}$ l'est
(\hyperref[I.1.3.9-fr]{I, 1.3.9})~; pour tout entier $k$ tel que
$0\leq k\leq d-1$, on désignera par $\mathcal{M}^{(d,k)}$ (ou
$\mathcal{M}^{(d)}$ pour $k=0$) la somme directe des
$\mathcal{M}_{nd+k}$ ($n\in\mathbf{Z}$), qui est un
$\mathcal{S}^{(d)}$-Module gradué, quasi-cohérent lorsque $\mathcal{S}$ et
$\mathcal{M}$ sont quasi-cohérents (\hyperref[I.9.6.1-fr]{I, 9.6.1}). On
notera $\mathcal{M}(n)$ le $\mathcal{S}$-Module gradué tel que
$(\mathcal{M}(n))_k=\mathcal{M}_{n+k}$ pour tout $k\in\mathbf{Z}$~; si
$\mathcal{S}$ et $\mathcal{M}$ sont quasi-cohérents, $\mathcal{M}(n)$ est un
$\mathcal{S}$-Module gradué quasi-cohérent
(\hyperref[I.9.6.1-fr]{I, 9.6.1}).

Nous dirons que $\mathcal{M}$ est un $\mathcal{S}$-Module gradué
\emph{de type fini} (resp. qu'il admet une \emph{présentation finie}) si pour
tout $y\in Y$, il existe un voisinage ouvert $U$ de $y$ et des entiers $n_i$
(resp. des entiers $m_i$ et $n_j$) tels qu'il y ait un homomorphisme surjectif
de degré~$0$
\[
  \bigoplus_{i=1}^r(\mathcal{S}(n_i)|U)\longrightarrow\mathcal{M}|U
\]
(resp. que $\mathcal{M}|U$ soit isomorphe au conoyau d'un homomorphisme de
degré~$0$
\[
  \bigoplus_{i=1}^r(\mathcal{S}(m_i)|U)
  \longrightarrow\bigoplus_{j=1}^s(\mathcal{S}(n_j)|U)).
\]

Soient $U$ un ouvert affine de $Y$, $A=\Gamma(U,\mathcal{O}_Y)$ son anneau~;
par hypothèse, la $(\mathcal{O}_Y|U)$-Algèbre graduée $\mathcal{S}|U$ est
isomorphe à $\widetilde{S}$, où $S=\Gamma(U,\mathcal{S})$ est une
$A$-algèbre

\oldpage[II]{50}

graduée (\hyperref[I.1.4.3-fr]{I, 1.4.3})~; posons
$X_U=\operatorname{Proj}(\Gamma(U,\mathcal{S}))$. Soient $U'\subset U$ un
second ouvert affine de $Y$, $A'$ son anneau, $j:U'\to U$ l'injection
canonique, qui correspond à l'homomorphisme de restriction $A\to A'$~; on a
$\mathcal{S}|U'=j^*(\mathcal{S}|U)$, et par suite
$S'=\Gamma(U',\mathcal{S})$ s'identifie canoniquement à $S\otimes_A A'$
(\hyperref[I.1.6.4-fr]{I, 1.6.4}). On en conclut
(\hyperref[II.2.8.10-fr]{2.8.10}) que $X_{U'}$ s'identifie canoniquement à
$X_U\times_U U'$, et par suite aussi à $f_U^{-1}(U')$, en désignant par
$f_U$ le morphisme structural $X_U\to U$
(\hyperref[I.4.4.1-fr]{I, 4.4.1}). Nous désignerons par $\sigma_{U',U}$
l'isomorphisme canonique $f_U^{-1}(U')\xrightarrow{\sim}X_{U'}$ ainsi défini,
par $\rho_{U',U}$ l'immersion ouverte $X_{U'}\to X_U$ obtenue en composant
$\sigma_{U',U}^{-1}$ et l'injection canonique $f_U^{-1}(U')\to X_U$. Il est
immédiat que si $U''\subset U'$ est un troisième ouvert affine de $Y$, on a
$\rho_{U'',U}=\rho_{U',U}\circ\rho_{U'',U'}$.
\end{env}

\begin{proposition}[3.1.2]
\label{II.3.1.2-fr}
Soit $Y$ un préschéma. Pour toute $\mathcal{O}_Y$-Algèbre graduée
quasi-cohérente $\mathcal{S}$ à degrés positifs, il existe un préschéma $X$
au-dessus de $Y$, et un seul à un $Y$-isomorphisme près, ayant la propriété
suivante~: si $f$ est le morphisme structural $X\to Y$, pour tout ouvert
affine $U$ de $Y$, il existe un \emph{isomorphisme} $\eta_U$ du préschéma
induit $f^{-1}(U)$ sur $X_U=\operatorname{Proj}(\Gamma(U,\mathcal{S}))$ tel
que, si $V$ est un second ouvert affine de $Y$ contenu dans $U$, le diagramme
\[
\label{II.3.1.2.1-fr}
  \xymatrix{
    f^{-1}(V)\ar[r]^{\eta_V}\ar[d]
    & X_V\ar[d]^{\rho_{V,U}}\\
    f^{-1}(U)\ar[r]_{\eta_U}
    & X_U
  }
\tag{3.1.2.1}
\]
soit commutatif.
\end{proposition}

\begin{proof}
Pour deux ouverts affines $U$, $V$ de $Y$, soit $X_{U,V}$ le préschéma
induit sur $f_U^{-1}(U\cap V)$ par $X_U$~; nous allons définir un
$Y$-isomorphisme $\theta_{U,V}:X_{V,U}\xrightarrow{\sim}X_{U,V}$. Pour cela,
considérons un ouvert affine $W\subset U\cap V$~: en composant les
isomorphismes
\[
  f_U^{-1}(W)\xrightarrow{\sigma_{W,U}}X_W
  \xrightarrow{\sigma_{W,V}^{-1}}f_V^{-1}(W),
\]
on obtient un isomorphisme $\tau_W$, et on vérifie aussitôt que si
$W'\subset W$ est un ouvert affine, $\tau_{W'}$ est la restriction à
$f_U^{-1}(W')$ de $\tau_W$~; les $\tau_W$ sont donc bien les restrictions
d'un $Y$-isomorphisme $\theta_{V,U}$. En outre, si $U$, $V$, $W$ sont trois
ouverts affines de $Y$, $\theta'_{U,V}$, $\theta'_{V,W}$ et
$\theta'_{U,W}$ les restrictions de $\theta_{U,V}$, $\theta_{V,W}$,
$\theta_{U,W}$ aux images réciproques de $U\cap V\cap W$ dans $X_V$, $X_W$,
$X_W$ respectivement, il résulte des définitions précédentes que l'on a
$\theta'_{U,V}\circ\theta'_{V,W}=\theta'_{U,W}$. L'existence de $X$
vérifiant les propriétés énoncées résulte donc de
(\hyperref[I.2.3.1-fr]{I, 2.3.1})~; son unicité à un $Y$-isomorphisme près
est triviale, compte tenu de (\hyperref[II.3.1.2.1-fr]{3.1.2.1}).
\end{proof}

\begin{env}[3.1.3]
\label{II.3.1.3-fr}
Nous dirons que le préschéma $X$ défini dans
(\hyperref[II.3.1.2-fr]{3.1.2}) est le \emph{spectre homogène} de la
$\mathcal{O}_Y$-Algèbre graduée quasi-cohérente $\mathcal{S}$ et nous le
noterons $\operatorname{Proj}(\mathcal{S})$. Il est immédiat que
$\operatorname{Proj}(\mathcal{S})$ est \emph{séparé au-dessus de $Y$}
((\hyperref[II.2.4.2-fr]{2.4.2}) et
(\hyperref[I.5.5.5-fr]{I, 5.5.5}))~; si $\mathcal{S}$ est une
$\mathcal{O}_Y$-Algèbre \emph{de type fini}
(\hyperref[I.9.6.2-fr]{I, 9.6.2}), $\operatorname{Proj}(\mathcal{S})$ est
\emph{de type fini} sur $Y$ ((\hyperref[II.2.7.1-fr]{2.7.1, (ii)}) et
(\hyperref[I.6.3.1-fr]{I, 6.3.1})).

Si $f$ est le morphisme structural $X\to Y$, il est immédiat que pour tout
préschéma induit par $Y$ sur un ouvert $U$ de $Y$, $f^{-1}(U)$ s'identifie au
spectre homogène $\operatorname{Proj}(\mathcal{S}|U)$.
\end{env}

\begin{proposition}[3.1.4]
\label{II.3.1.4-fr}
Soit $f\in\Gamma(Y,\mathcal{S}_d)$ pour $d>0$. Il existe une partie ouverte
$X_f$ de l'espace sous-jacent à $X=\operatorname{Proj}(\mathcal{S})$ ayant la
propriété suivante~: pour tout ouvert affine $U$ de $Y$,

\oldpage[II]{51}

$X_f\cap\varphi^{-1}(U)=D_+(f|U)$ dans $\varphi^{-1}(U)$ identifié à
$X_U=\operatorname{Proj}(\Gamma(U,\mathcal{S}))$ ($\varphi$ désignant le
morphisme structural $X\to Y$). En outre, le préschéma induit sur $X_f$ est
affine au-dessus de $Y$ et est canoniquement isomorphe à
$\operatorname{Spec}(\mathcal{S}^{(d)}/(f-1)\mathcal{S}^{(d)})$
(\hyperref[II.1.3.1-fr]{1.3.1}).
\end{proposition}

\begin{proof}
On a $f|U\in\Gamma(U,\mathcal{S}_d)=(\Gamma(U,\mathcal{S}))_d$. Si $U$,
$U'$ sont deux ouverts affines de $Y$ tels que $U'\subset U$, $f|U'$ est
l'image de $f|U$ par l'homomorphisme de restriction
\[
  \Gamma(U,\mathcal{S})\longrightarrow\Gamma(U',\mathcal{S}),
\]
donc $D_+(f|U')$ est égal (avec les notations de
(\hyperref[II.3.1.1-fr]{3.1.1})) au préschéma induit sur l'image réciproque
$\rho_{U',U}^{-1}(D_+(f|U))$ dans $X_{U'}$
(\hyperref[II.2.8.1-fr]{2.8.1})~; d'où la première assertion. En outre, le
préschéma induit sur $D_+(f|U)$ par $X_U$ s'identifie canoniquement à
$\operatorname{Spec}((\Gamma(U,\mathcal{S}))_{(f|U)})$, ces identifications
étant compatibles avec les homomorphismes de restriction
(\hyperref[II.2.8.1-fr]{2.8.1})~; la seconde assertion résulte alors de
(\hyperref[II.2.2.5-fr]{2.2.5}) et de la commutativité du diagramme
(\hyperref[II.2.8.2.1-fr]{2.8.2.1}).
\end{proof}

On dira encore que $X_f$ (en tant qu'ouvert dans l'espace sous-jacent $X$)
est l'ensemble des $x\in X$ où $f$ \emph{ne s'annule pas}.

\begin{corollary}[3.1.5]
\label{II.3.1.5-fr}
Si $f\in\Gamma(Y,\mathcal{S}_d)$, $g\in\Gamma(Y,\mathcal{S}_e)$, on a
\[
\label{II.3.1.5.1-fr}
  X_{fg}=X_f\cap X_g.
\tag{3.1.5.1}
\]
\end{corollary}

\begin{proof}
Il suffit en effet de considérer l'intersection des deux membres avec un
ensemble $\varphi^{-1}(U)$, où $U$ est un ouvert affine dans $Y$, et
d'appliquer la formule (\hyperref[II.2.3.3.2-fr]{2.3.3.2}).
\end{proof}

\begin{corollary}[3.1.6]
\label{II.3.1.6-fr}
Soit $(f_\alpha)$ une famille de sections de $\mathcal{S}$ au-dessus de $Y$
telle que $f_\alpha\in\Gamma(Y,\mathcal{S}_{d_\alpha})$~; si le faisceau
d'idéaux de $\mathcal{S}$ engendré par cette famille
(\hyperref[0.5.1.1-fr]{0, 5.1.1}) contient tous les $\mathcal{S}_n$ à partir
d'un certain rang, l'espace sous-jacent $X$ est réunion des $X_{f_\alpha}$.
\end{corollary}

\begin{proof}
En effet, pour tout ouvert affine $U$ de $Y$, $\varphi^{-1}(U)$ est réunion
des $X_{f_\alpha}\cap\varphi^{-1}(U)$
(\hyperref[II.2.3.14-fr]{2.3.14}).
\end{proof}

\begin{corollary}[3.1.7]
\label{II.3.1.7-fr}
Soit $\mathcal{A}$ une $\mathcal{O}_Y$-Algèbre quasi-cohérente~; posons
\[
  \mathcal{S}=\mathcal{A}[T]
  =\mathcal{A}\otimes_{\mathbf{Z}}\mathbf{Z}[T]
\]
où $T$ est une indéterminée ($\mathbf{Z}$ et $\mathbf{Z}[T]$ étant
considérés comme faisceaux simples sur $Y$). Alors
$X=\operatorname{Proj}(\mathcal{S})$ s'identifie canoniquement à
$\operatorname{Spec}(\mathcal{A})$. En particulier,
$\operatorname{Proj}(\mathcal{O}_Y[T])$ s'identifie à $Y$.
\end{corollary}

\begin{proof}
En appliquant (\hyperref[II.3.1.6-fr]{3.1.6}) à l'unique section
$f\in\Gamma(Y,\mathcal{S})$ égale à $T$ en chaque point de $Y$, on voit que
$X_f=X$. En outre, on a ici $d=1$, et
$\mathcal{S}^{(1)}/(f-1)\mathcal{S}^{(1)}
=\mathcal{S}/(f-1)\mathcal{S}$ est canoniquement isomorphe à $\mathcal{A}$,
d'où le corollaire (\hyperref[II.1.2.2-fr]{1.2.2}).
\end{proof}

Soit $g\in\Gamma(Y,\mathcal{O}_Y)$~; si on prend
$\mathcal{S}=\mathcal{O}_Y[T]$, on a donc $g\in\Gamma(Y,\mathcal{S}_0)$~; soit
\[
  h=gT\in\Gamma(Y,\mathcal{S}_1).
\]
Si $X=\operatorname{Proj}(\mathcal{S})$, l'identification canonique définie
dans (\hyperref[II.3.1.7-fr]{3.1.7}) identifie $X_h$ à l'ensemble ouvert
$Y_g$ de $Y$ (au sens de (\hyperref[0.5.5.2-fr]{0, 5.5.2}))~: en effet, on
peut se borner au cas où $Y=\operatorname{Spec}(A)$ est affine, et tout se
ramène alors (compte tenu de (\hyperref[II.2.2.5-fr]{2.2.5})) au fait que
l'anneau de fractions $A_g$ s'identifie canoniquement à
$A[T]/(gT-1)A[T]$ (\hyperref[0.1.2.3-fr]{0, 1.2.3}).

\begin{proposition}[3.1.8]
\label{II.3.1.8-fr}
Soit $\mathcal{S}$ une $\mathcal{O}_Y$-Algèbre graduée quasi-cohérente à
degrés positifs.
\begin{enumerate}
  \item[\rm(i)] Pour tout $d>0$, il existe un $Y$-isomorphisme canonique de
    $\operatorname{Proj}(\mathcal{S})$ sur
    $\operatorname{Proj}(\mathcal{S}^{(d)})$.

\oldpage[II]{52}

  \item[\rm(ii)] Soit $\mathcal{S}'$ la $\mathcal{O}_Y$-Algèbre graduée
    somme directe de $\mathcal{O}_Y$ et des $\mathcal{S}_n$ ($n\geq0$)~;
    alors $\operatorname{Proj}(\mathcal{S}')$ et
    $\operatorname{Proj}(\mathcal{S})$ sont canoniquement $Y$-isomorphes.
  \item[\rm(iii)] Soit $\mathcal{L}$ un $\mathcal{O}_Y$-Module inversible
    (\hyperref[0.5.4.1-fr]{0, 5.4.1}) et soit $\mathcal{S}_{(\mathcal{L})}$ la
    $\mathcal{O}_Y$-Algèbre graduée somme directe des
    $\mathcal{S}_d\otimes\mathcal{L}^{\otimes d}$ ($d\geq0$)~; alors
    $\operatorname{Proj}(\mathcal{S})$ et
    $\operatorname{Proj}(\mathcal{S}_{(\mathcal{L})})$ sont canoniquement
    $Y$-isomorphes.
\end{enumerate}
\end{proposition}

\begin{proof}
Dans chacun des trois cas, il suffit de définir l'isomorphisme localement sur
$Y$, la vérification des compatibilités avec les opérations de restriction
d'un ouvert à un ouvert plus petit étant triviale. On peut donc supposer $Y$
affine, et alors (i) résulte de (\hyperref[II.2.4.7-fr]{2.4.7, (i)}) et (ii)
de (\hyperref[II.2.4.8-fr]{2.4.8}). En ce qui concerne (iii), lorsqu'on
suppose en outre $\mathcal{L}$ isomorphe à $\mathcal{O}_Y$ (ce qui est permis,
la question étant locale sur $Y$), l'isomorphie de
$\operatorname{Proj}(\mathcal{S})$ et
$\operatorname{Proj}(\mathcal{S}_{(\mathcal{L})})$ est évidente~; pour
définir un isomorphisme \emph{canonique}, soient
$Y=\operatorname{Spec}(A)$, $\mathcal{S}=\widetilde{S}$, où $S$ est une
$A$-algèbre graduée, et soit $c$ un générateur du $A$-module libre $L$ tel
que $\mathcal{L}=\widetilde{L}$~; alors pour tout $n>0$,
$x_n\mapsto x_n\otimes c^{\otimes n}$ est un $A$-isomorphisme de $S_n$ sur
$S_n\otimes L^{\otimes n}$, et ces $A$-isomorphismes définissent un
$A$-isomorphisme d'algèbres graduées
\[
  \varphi_c:S\longrightarrow S_{(L)}
  =\bigoplus_{n\geq0}S_n\otimes L^{\otimes n}.
\]
Soit alors $f\in S_+$, homogène de degré $d$~; pour tout $x\in S_{nd}$, on a
\[
  \frac{x\otimes c^{\otimes nd}}{(f\otimes c^{\otimes d})^n}
  =\frac{x\otimes(\varepsilon c)^{\otimes nd}}
  {(f\otimes(\varepsilon c)^{\otimes d})^n}
\]
pour tout élément inversible $\varepsilon\in A$, ce qui montre que
l'isomorphisme $S_{(f)}\to(S_{(L)})_{(f\otimes c^d)}$ déduit de $\varphi_c$
est \emph{indépendant} du générateur $c$ de $L$ considéré et achève la
démonstration.
\end{proof}

\begin{env}[3.1.9]
\label{II.3.1.9-fr}
Rappelons ((\hyperref[0.4.1.3-fr]{0, 4.1.3}) et
(\hyperref[I.1.3.14-fr]{I, 1.3.14})) que pour que la
$\mathcal{O}_Y$-Algèbre graduée quasi-cohérente $\mathcal{S}$ soit
\emph{engendrée par le $\mathcal{O}_Y$-Module $\mathcal{S}_1$}, il faut et il
suffit qu'il existe un recouvrement $(U_\alpha)$ de $Y$ par des ouverts
affines tels que l'algèbre graduée $\Gamma(U_\alpha,\mathcal{S})$ sur
$\Gamma(U_\alpha,\mathcal{S}_0)$ soit engendrée par l'ensemble
$\Gamma(U_\alpha,\mathcal{S}_1)$ de ses éléments homogènes de degré~$1$. Pour
tout ouvert $V$ de $Y$, $\mathcal{S}|V$ est alors engendrée par le
$(\mathcal{O}_Y|V)$-Module $\mathcal{S}_1|V$.
\end{env}

\begin{proposition}[3.1.10]
\label{II.3.1.10-fr}
Supposons qu'il existe un recouvrement ouvert affine fini $(U_i)$ de $Y$ tel
que chacune des algèbres graduées $\Gamma(U_i,\mathcal{S})$ soit de type fini
sur $\Gamma(U_i,\mathcal{O}_Y)$. Alors il existe $d>0$ tel que
$\mathcal{S}^{(d)}$ soit engendrée par $\mathcal{S}_d$, $\mathcal{S}_d$ étant
un $\mathcal{O}_Y$-Module de type fini.
\end{proposition}

\begin{proof}
En effet, il résulte de (\hyperref[II.2.1.6-fr]{2.1.6, (v)}) que pour chaque
$i$, il existe un entier $m_i$ tel que
$\Gamma(U_i,\mathcal{S}_{nm_i})=(\Gamma(U_i,\mathcal{S}_{m_i}))^n$ pour tout
$n>0$~; il suffit de prendre pour $d$ un multiple commun des $m_i$, compte
tenu de (\hyperref[II.2.1.6-fr]{2.1.6, (i)}).
\end{proof}

\begin{corollary}[3.1.11]
\label{II.3.1.11-fr}
Sous les hypothèses de (\hyperref[II.3.1.10-fr]{3.1.10}),
$\operatorname{Proj}(\mathcal{S})$ est $Y$-isomorphe à un spectre homogène
$\operatorname{Proj}(\mathcal{S}')$, où $\mathcal{S}'$ est une
$\mathcal{O}_Y$-Algèbre graduée engendrée par $\mathcal{S}'_1$,
$\mathcal{S}'_1$ étant supposé être un $\mathcal{O}_Y$-Module de type fini.
\end{corollary}

\begin{proof}
Il suffit en effet de prendre $\mathcal{S}'=\mathcal{S}^{(d)}$, $d$ étant
déterminé par la propriété de (\hyperref[II.3.1.10-fr]{3.1.10}), et
d'appliquer (\hyperref[II.3.1.8-fr]{3.1.8, (i)}).
\end{proof}

\begin{env}[3.1.12]
\label{II.3.1.12-fr}
Si $\mathcal{S}$ est une $\mathcal{O}_Y$-Algèbre graduée quasi-cohérente à
degrés positifs, on sait (\hyperref[I.5.1.1-fr]{I, 5.1.1}) que son
\emph{nilradical} $\mathcal{N}$ est un $\mathcal{O}_Y$-Module
quasi-cohérent~; nous dirons que $\mathcal{N}_+=\mathcal{N}\cap\mathcal{S}_+$
est le \emph{nilradical de $\mathcal{S}_+$}~; c'est un
$\mathcal{S}_0$-Module quasi-cohérent gradué, car on se ramène immédiatement
au cas où $Y$ est affine, et la proposition résulte alors de
(\hyperref[II.2.1.10-fr]{2.1.10}). Pour tout $y\in Y$,
$(\mathcal{N}_+)_y$ est alors le nilradical de
$(\mathcal{S}_+)_y=(\mathcal{S}_y)_+$
(\hyperref[I.5.1.1-fr]{I, 5.1.1}). On dit que la
$\mathcal{O}_Y$-Algèbre graduée $\mathcal{S}$ est \emph{essentiellement
réduite} si $\mathcal{N}_+=0$, ce qui équivaut

\oldpage[II]{53}

à dire que $\mathcal{S}_y$
est une $\mathcal{O}_y$-algèbre graduée essentiellement réduite pour tout
$y\in Y$. Pour toute $\mathcal{O}_Y$-Algèbre graduée $\mathcal{S}$,
$\mathcal{S}/\mathcal{N}_+$ est essentiellement réduite.

Nous dirons que $\mathcal{S}$ est \emph{intègre} si $\mathcal{S}_y$ est un
anneau intègre pour tout $y\in Y$ et si en outre
$(\mathcal{S}_y)_+=(\mathcal{S}_+)_y\neq0$ pour tout $y\in Y$.
\end{env}

\begin{proposition}[3.1.13]
\label{II.3.1.13-fr}
Soit $\mathcal{S}$ une $\mathcal{O}_Y$-Algèbre graduée à degrés positifs.
Si $X=\operatorname{Proj}(\mathcal{S})$, le $Y$-schéma $X_{\mathrm{red}}$ est
canoniquement isomorphe à $\operatorname{Proj}(\mathcal{S}/\mathcal{N}_+)$~;
en particulier, si $\mathcal{S}$ est essentiellement réduite, $X$ est réduit.
\end{proposition}

\begin{proof}
Le fait que $X'=\operatorname{Proj}(\mathcal{S}/\mathcal{N}_+)$ est réduit
découle immédiatement de (\hyperref[II.2.4.4-fr]{2.4.4, (i)}), la propriété
étant locale~; en outre, pour tout ouvert affine $U\subset Y$,
${\varphi'}^{-1}(U)$ est égal à $(\varphi^{-1}(U))_{\mathrm{red}}$ (en
désignant par $\varphi$ et $\varphi'$ les morphismes structuraux $X\to Y$,
$X'\to Y$)~; on vérifie aussitôt que les $U$-morphismes canoniques
${\varphi'}^{-1}(U)\to\varphi^{-1}(U)$ sont compatibles avec les opérations
de restriction et définissent par suite une immersion fermée $X'\to X$ qui
est un homéomorphisme des espaces sous-jacents~; d'où la conclusion
(\hyperref[I.5.1.2-fr]{I, 5.1.2}).
\end{proof}

\begin{proposition}[3.1.14]
\label{II.3.1.14-fr}
Soient $Y$ un préschéma intègre, $\mathcal{S}$ une
$\mathcal{O}_Y$-Algèbre graduée quasi-cohérente telle que
$\mathcal{S}_0=\mathcal{O}_Y$.
\begin{enumerate}
  \item[\rm(i)] Si $\mathcal{S}$ est intègre
    (\hyperref[II.3.1.12-fr]{3.1.12}), alors
    $X=\operatorname{Proj}(\mathcal{S})$ est intègre et le morphisme
    structural $\varphi:X\to Y$ est dominant.
  \item[\rm(ii)] Supposons en outre que $\mathcal{S}$ soit essentiellement
    réduite. Alors, inversement, si $X$ est intègre et si $\varphi$ est
    dominant, $\mathcal{S}$ est intègre.
\end{enumerate}
\end{proposition}

\begin{proof}
\begin{enumerate}
  \item[\rm(i)] Si $(U_\alpha)$ est une base de $Y$ formée d'ouverts affines
    non vides, il suffit de démontrer la proposition lorsque $Y$ est remplacé
    par un des $U_\alpha$ et $\mathcal{S}$ par $\mathcal{S}|U_\alpha$~: en
    effet, il en résultera d'une part que les espaces sous-jacents
    $\varphi^{-1}(U_\alpha)$ sont des ouverts irréductibles (donc non vides)
    de $X$ tels que
    $\varphi^{-1}(U_\alpha)\cap\varphi^{-1}(U_\beta)\neq\varnothing$ pour tout
    couple d'indices (puisque $U_\alpha\cap U_\beta$ contient un $U_\gamma$),
    donc que $X$ est irréductible (\hyperref[0.2.1.4-fr]{0, 2.1.4})~; d'autre
    part $X$ sera réduit, puisque c'est là une propriété locale, donc $X$
    sera bien intègre et $\varphi(X)$ dense dans $Y$.

    Supposons donc que $Y=\operatorname{Spec}(A)$ où $A$ est intègre
    (\hyperref[I.5.1.4-fr]{I, 5.1.4}) et $\mathcal{S}=\widetilde{S}$, où $S$
    est une $A$-algèbre graduée~; l'hypothèse est que pour tout $y\in Y$,
    $\widetilde{S}_y=S_y$ est un anneau gradué intègre tel que
    $(S_y)_+\neq0$. Il suffit de prouver que $S$ est un anneau \emph{intègre},
    car alors on aura $S_+\neq0$ et on pourra appliquer
    (\hyperref[II.2.4.4-fr]{2.4.4, (ii)}). Or, soient $f$, $g$ deux éléments
    $\neq0$ de $S$ et supposons que $fg=0$~; pour tout $y\in Y$ on aura donc
    $(f/1)(g/1)=0$ dans $S_y$, donc $f/1=0$ ou $g/1=0$ par hypothèse.
    Supposons par exemple $f/1=0$ dans $S_y$~; cela signifie qu'il existe
    $a\in A$ tel que $a\notin\mathfrak{j}_y$ et $af=0$~; on a alors pour tout
    $z\in Y$, $(a/1)(f/1)=0$ dans l'anneau \emph{intègre} $S_z$, et comme
    $a/1\neq0$ (puisque $A$ est intègre), $f/1=0$, ce qui implique $f=0$.
  \item[\rm(ii)] La question étant locale sur $Y$, on peut encore supposer
    que $Y=\operatorname{Spec}(A)$, $A$ étant intègre, et
    $\mathcal{S}=\widetilde{S}$. Par hypothèse, pour tout $y\in Y$, $(S_y)_+$
    ne contient pas d'élément nilpotent $\neq0$, et il en est de même de
    $(S_0)_y=A_y$ par hypothèse~; donc $S_y$ est un anneau réduit pour tout
    $y\in Y$, et on en conclut d'abord que $S$ lui-même est réduit
    (\hyperref[I.5.1.1-fr]{I, 5.1.1}). L'hypothèse que $X$ est intègre
    entraîne d'autre part que $S$ est essentiellement intègre
    (\hyperref[II.2.4.4-fr]{2.4.4, (ii)}), et tout revient alors à voir que
    l'annulateur $\mathfrak{J}$ de $S_+$ dans $A=S_0$ est réduit

\oldpage[II]{54}

    à $0$
    (\hyperref[II.2.1.11-fr]{2.1.11}). Dans le cas contraire, on aurait
    $(S_h)_+=0$ pour un $h\neq0$ dans $\mathfrak{J}$, donc
    (\hyperref[II.3.1.1-fr]{3.1.1}) $\varphi^{-1}(D(h))=\varnothing$, et
    $\varphi(X)$ ne serait pas dense dans $Y$ contrairement à l'hypothèse
    (puisque $D(h)\neq\varnothing$, $h$ n'étant pas nilpotent).
\end{enumerate}
\end{proof}

\subsection{Faisceau sur $\operatorname{Proj}(\mathcal{S})$ associé à un
$\mathcal{S}$-Module gradué.}
\label{subsection:II.3.2-fr}

\begin{env}[3.2.1]
\label{II.3.2.1-fr}
Soient $Y$ un préschéma, $\mathcal{S}$ une $\mathcal{O}_Y$-Algèbre graduée
quasi-cohérente à degrés positifs, $\mathcal{M}$ un $\mathcal{S}$-module
gradué quasi-cohérent (sur $(Y,\mathcal{O}_Y)$ ou sur l'espace annelé
$(Y,\mathcal{S})$, ce qui revient au même
(\hyperref[I.9.6.1-fr]{I, 9.6.1})). Avec les notations de
(\hyperref[II.3.1.1-fr]{3.1.1}), désignons par $\widetilde{\mathcal{M}}_U$ le
$\mathcal{O}_{X_U}$-Module quasi-cohérent
$\widetilde{\Gamma(U,\mathcal{M})}$~; pour $U'\subset U$,
$\Gamma(U',\mathcal{M})$ s'identifie canoniquement à
$\Gamma(U,\mathcal{M})\otimes_A A'$ (\hyperref[I.1.6.4-fr]{I, 1.6.4})~; donc
on a $\widetilde{\mathcal{M}}_{U'}
=\rho_{U',U}^*(\widetilde{\mathcal{M}}_U)$
(\hyperref[II.2.8.11-fr]{2.8.11}).
\end{env}

\begin{proposition}[3.2.2]
\label{II.3.2.2-fr}
Il existe sur $\operatorname{Proj}(\mathcal{S})=X$ un
$\mathcal{O}_X$-Module quasi-cohérent et un seul $\widetilde{\mathcal{M}}$ tel
que, pour tout ouvert affine $U$ de $Y$, on ait
\[
  \eta_U^*(\widetilde{\Gamma(U,\mathcal{M})})
  =\widetilde{\mathcal{M}}|f^{-1}(U)
\]
(en désignant par $\eta_U$ l'isomorphisme
$f^{-1}(U)\xrightarrow{\sim}\operatorname{Proj}(\Gamma(U,\mathcal{S}))$, où
$f$ est le morphisme structural $X\to Y$).
\end{proposition}

\begin{proof}
Comme $\rho_{U',U}$ s'identifie au morphisme d'injection
$f^{-1}(U')\to f^{-1}(U)$ (\hyperref[II.3.1.2.1-fr]{3.1.2.1}), la proposition
résulte aussitôt de la relation
$\widetilde{\mathcal{M}}_{U'}=\rho_{U',U}^*(\widetilde{\mathcal{M}}_U)$ et du
principe de recollement des faisceaux (\hyperref[0.3.3.1-fr]{0, 3.3.1}).
\end{proof}

On dit que $\widetilde{\mathcal{M}}$ est le $\mathcal{O}_X$-Module associé au
$\mathcal{S}$-Module gradué quasi-cohérent $\mathcal{M}$.

\begin{proposition}[3.2.3]
\label{II.3.2.3-fr}
Soit $\mathcal{M}$ un $\mathcal{S}$-Module gradué quasi-cohérent, et soit
$f\in\Gamma(Y,\mathcal{S}_d)$ ($d>0$). Si $\xi_f$ est l'isomorphisme canonique
de $X_f$ sur le $Y$-préschéma
$Z_f=\operatorname{Spec}(\mathcal{S}^{(d)}/(f-1)\mathcal{S}^{(d)})$
(\hyperref[II.3.1.4-fr]{3.1.4}), $(\xi_f)_*(\widetilde{\mathcal{M}}|X_f)$ est
le $\mathcal{O}_{Z_f}$-Module
$\widetilde{\mathcal{M}^{(d)}/(f-1)\mathcal{M}^{(d)}}$
(\hyperref[II.1.4.3-fr]{1.4.3}).
\end{proposition}

\begin{proof}
La question étant locale sur $Y$, on est aussitôt ramené à
(\hyperref[II.2.2.5-fr]{2.2.5}), compte tenu de la commutativité du diagramme
(\hyperref[II.2.8.12.1-fr]{2.8.12.1}).
\end{proof}

\begin{proposition}[3.2.4]
\label{II.3.2.4-fr}
Le $\mathcal{O}_X$-Module $\widetilde{\mathcal{M}}$ est un foncteur covariant
additif exact en $\mathcal{M}$, de la catégorie des $\mathcal{S}$-Modules
gradués quasi-cohérents dans celle des $\mathcal{O}_X$-Modules
quasi-cohérents, qui commute aux limites inductives et aux sommes directes.
\end{proposition}

\begin{proof}
La question étant locale sur $Y$, se ramène à
(\hyperref[I.1.3.11-fr]{I, 1.3.11} et
\hyperref[I.1.3.9-fr]{I, 1.3.9}) et à
(\hyperref[II.2.5.4-fr]{2.5.4}).
\end{proof}

En particulier, si $\mathcal{N}$ est un sous-$\mathcal{S}$-Module gradué
quasi-cohérent de $\mathcal{M}$, $\widetilde{\mathcal{N}}$ s'identifie
canoniquement à un sous-$\mathcal{O}_X$-Module quasi-cohérent de
$\widetilde{\mathcal{M}}$~; plus particulièrement, pour tout faisceau gradué
quasi-cohérent $\mathcal{J}$ d'idéaux de $\mathcal{S}$,
$\widetilde{\mathcal{J}}$ est un faisceau quasi-cohérent d'idéaux de
$\mathcal{O}_X$.

Si $\mathcal{M}$ est un $\mathcal{S}$-Module gradué quasi-cohérent et
$\mathcal{I}$ un faisceau quasi-cohérent d'idéaux de $\mathcal{O}_Y$,
$\mathcal{I}\mathcal{M}$ est un sous-$\mathcal{S}$-Module gradué
quasi-cohérent de $\mathcal{M}$ et on a
\[
\label{II.3.2.4.1-fr}
  \widetilde{\mathcal{I}\mathcal{M}}
  =\mathcal{I}\cdot\widetilde{\mathcal{M}}
\tag{3.2.4.1}
\]
(le second membre ayant le sens défini en
(\hyperref[0.4.3.5-fr]{0, 4.3.5})). Il suffit en effet de vérifier cette
formule lorsque $Y=\operatorname{Spec}(A)$ est affine,
$\mathcal{S}=\widetilde{S}$, où $S$ est une $A$-algèbre graduée,
$\mathcal{M}=\widetilde{M}$,

\oldpage[II]{55}

où $M$ est un $S$-module gradué, et
$\mathcal{I}=\widetilde{\mathfrak{I}}$, où $\mathfrak{I}$ est un idéal de
$A$. Pour tout $f$ homogène de $S_+$, la restriction à
$D_+(f)=\operatorname{Spec}(S_{(f)})$ du premier membre de
(\hyperref[II.3.2.4.1-fr]{3.2.4.1}) est associée à
$(\mathfrak{I}M)_{(f)}=\mathfrak{I}\cdot M_{(f)}$, et il en est de même de la
restriction du second membre, vu
(\hyperref[I.1.3.13-fr]{I, 1.3.13} et
\hyperref[I.1.6.9-fr]{I, 1.6.9}).

\begin{proposition}[3.2.5]
\label{II.3.2.5-fr}
Soit $f\in\Gamma(Y,\mathcal{S}_d)$ ($d>0$). Sur l'ensemble ouvert $X_f$, le
$(\mathcal{O}_X|X_f)$-Module
$\widetilde{\mathcal{S}(nd)}|X_f$ est canoniquement isomorphe à
$\mathcal{O}_X|X_f$ pour tout $n\in\mathbf{Z}$. En particulier, si la
$\mathcal{O}_Y$-Algèbre $\mathcal{S}$ est engendrée par $\mathcal{S}_1$
(\hyperref[II.3.1.9-fr]{3.1.9}), les $\mathcal{O}_X$-Modules
$\widetilde{\mathcal{S}(n)}$ sont inversibles pour tout $n\in\mathbf{Z}$.
\end{proposition}

\begin{proof}
En effet, pour tout ouvert affine $U$ de $Y$, on a défini dans
(\hyperref[II.2.5.7-fr]{2.5.7}) un isomorphisme canonique de
$\widetilde{\mathcal{S}(nd)}|(X_f\cap\varphi^{-1}(U))$ sur
$\mathcal{O}_X|(X_f\cap\varphi^{-1}(U))$, compte tenu de
(\hyperref[II.3.1.4-fr]{3.1.4}) (où $\varphi$ est le morphisme structural
$X\to Y$)~; il est immédiat de vérifier que ces isomorphismes sont compatibles
avec la restriction de $U$ à un ouvert affine $U'\subset U$, d'où la première
assertion. Pour établir la seconde, il suffit de remarquer que si
$\mathcal{S}$ est engendrée par $\mathcal{S}_1$, il y a un recouvrement
$(U_\alpha)$ de $Y$ par des ouverts affines tels que
$\Gamma(U_\alpha,\mathcal{S})$ soit engendrée par
$\Gamma(U_\alpha,\mathcal{S})_1=\Gamma(U_\alpha,\mathcal{S}_1)$~; on est alors
ramené à appliquer le résultat de (\hyperref[II.2.5.9-fr]{2.5.9}), la
propriété d'être inversible étant locale.
\end{proof}

On posera encore, pour tout $n\in\mathbf{Z}$,
\[
\label{II.3.2.5.1-fr}
  \mathcal{O}_X(n)=\widetilde{\mathcal{S}(n)}
\tag{3.2.5.1}
\]
et pour tout $\mathcal{O}_X$-Module $\mathcal{F}$
\[
\label{II.3.2.5.2-fr}
  \mathcal{F}(n)=\mathcal{F}\otimes_{\mathcal{O}_X}\mathcal{O}_X(n).
\tag{3.2.5.2}
\]
Il résulte aussitôt de ces définitions que pour tout ouvert $U$ de $Y$, on a
\[
  \widetilde{(\mathcal{S}|U)(n)}
  =\mathcal{O}_X(n)|f^{-1}(U)
\]
$f$ étant le morphisme structural $X\to Y$.

\begin{proposition}[3.2.6]
\label{II.3.2.6-fr}
Soient $\mathcal{M}$ et $\mathcal{N}$ deux $\mathcal{S}$-Modules gradués
quasi-cohérents, il existe un homomorphisme canonique fonctoriel (en
$\mathcal{M}$ et $\mathcal{N}$)
\[
\label{II.3.2.6.1-fr}
  \lambda:\widetilde{\mathcal{M}}
  \otimes_{\mathcal{O}_X}\widetilde{\mathcal{N}}
  \longrightarrow\widetilde{\mathcal{M}\otimes_{\mathcal{S}}\mathcal{N}}
\tag{3.2.6.1}
\]
et un homomorphisme canonique fonctoriel (en $\mathcal{M}$ et
$\mathcal{N}$)
\[
\label{II.3.2.6.2-fr}
  \mu:\widetilde{\mathcal{H}om_{\mathcal{S}}(\mathcal{M},\mathcal{N})}
  \longrightarrow
  \mathcal{H}om_{\mathcal{O}_X}
  (\widetilde{\mathcal{M}},\widetilde{\mathcal{N}}).
\tag{3.2.6.2}
\]
En outre, si $\mathcal{S}$ est engendrée par $\mathcal{S}_1$
(\hyperref[II.3.1.9-fr]{3.1.9}), $\lambda$ est un isomorphisme~; si de plus
$\mathcal{M}$ admet une présentation finie
(\hyperref[II.3.1.1-fr]{3.1.1}), $\mu$ est un isomorphisme.
\end{proposition}

\begin{proof}
Les isomorphismes $\lambda$ et $\mu$ ont été définis dans
(\hyperref[II.2.5.11.2-fr]{2.5.11.2}) et
(\hyperref[II.2.5.12.2-fr]{2.5.12.2}) lorsque $Y$ est affine~; ces définitions
étant locales se transportent aussitôt au cas général considéré ici, compte
tenu de (\hyperref[II.2.8.14-fr]{2.8.14}).
\end{proof}

\begin{corollary}[3.2.7]
\label{II.3.2.7-fr}
Si $\mathcal{S}$ est engendrée par $\mathcal{S}_1$, on a, quels que soient
$m$, $n$ dans $\mathbf{Z}$,
\[
\label{II.3.2.7.1-fr}
  \mathcal{O}_X(m)\otimes_{\mathcal{O}_X}\mathcal{O}_X(n)
  =\mathcal{O}_X(m+n)
\tag{3.2.7.1}
\]
\[
\label{II.3.2.7.2-fr}
  \mathcal{O}_X(n)=(\mathcal{O}_X(1))^{\otimes n}
\tag{3.2.7.2}
\]
à des isomorphismes canoniques près.
\end{corollary}

\oldpage[II]{56}

\begin{corollary}[3.2.8]
\label{II.3.2.8-fr}
Supposons $\mathcal{S}$ engendrée par $\mathcal{S}_1$. Pour tout
$\mathcal{S}$-Module gradué $\mathcal{M}$ et tout $n\in\mathbf{Z}$, on a
\[
\label{II.3.2.8.1-fr}
  (\mathcal{M}(n))^{\sim}=\widetilde{\mathcal{M}}(n)
\tag{3.2.8.1}
\]
à un isomorphisme canonique près.
\end{corollary}

\begin{proof}
Cela résulte des propriétés correspondantes pour $Y$ affine
(\hyperref[II.2.5.14-fr]{2.5.14} et
\hyperref[II.2.5.15-fr]{2.5.15}) ainsi que de
(\hyperref[II.2.8.11-fr]{2.8.11}).
\end{proof}

\par\medskip\noindent\textit{Remarques (3.2.9). ---\ }
\label{II.3.2.9-fr}
\begin{enumerate}
  \item[\rm(i)] Si $\mathcal{S}=\mathcal{A}[T]$, $\mathcal{A}$ étant une
    $\mathcal{O}_Y$-Algèbre quasi-cohérente
    (\hyperref[II.3.1.7-fr]{3.1.7}), on vérifie immédiatement que tous les
    $\mathcal{O}_X$-Modules inversibles $\mathcal{O}_X(n)$ sont canoniquement
    isomorphes à $\mathcal{O}_X$.

    En outre, soit $\mathcal{N}$ un $\mathcal{A}$-Module quasi-cohérent, et
    posons $\mathcal{M}=\mathcal{N}\otimes_{\mathcal{A}}\mathcal{A}[T]$. Il
    résulte alors de (\hyperref[II.3.2.3-fr]{3.2.3}) et de
    (\hyperref[II.3.1.7-fr]{3.1.7}) que dans l'identification canonique de
    $X=\operatorname{Proj}(\mathcal{A}[T])$ et de
    $X'=\operatorname{Spec}(\mathcal{A})$, le $\mathcal{O}_X$-Module
    $\widetilde{\mathcal{M}}$ s'identifie au $\mathcal{O}_{X'}$-Module
    $\widetilde{\mathcal{N}}$ associé à $\mathcal{N}$ (au sens de
    (\hyperref[II.1.4.3-fr]{1.4.3})).
  \item[\rm(ii)] Soit $\mathcal{S}$ une $\mathcal{O}_Y$-Algèbre graduée
    quelconque, $\mathcal{S}'$ la $\mathcal{O}_Y$-Algèbre graduée telle que
    $\mathcal{S}'_0=\mathcal{O}_Y$, $\mathcal{S}'_n=\mathcal{S}_n$ pour tout
    $n>0$~; l'isomorphisme canonique de
    $X=\operatorname{Proj}(\mathcal{S})$ sur
    $X'=\operatorname{Proj}(\mathcal{S}')$
    (\hyperref[II.3.1.8-fr]{3.1.8, (ii)}) identifie $\mathcal{O}_X(n)$ et
    $\mathcal{O}_{X'}(n)$ pour tout $n\in\mathbf{Z}$~: cela résulte de la même
    proposition pour le cas affine (\hyperref[II.2.5.16-fr]{2.5.16}) et du
    fait que ces identifications, pour les ouverts affines de $Y$, commutent
    avec les opérations de restriction. De même, soit
    $X^{(d)}=\operatorname{Proj}(\mathcal{S}^{(d)})$~; l'isomorphisme
    canonique de $X$ sur $X^{(d)}$
    (\hyperref[II.3.1.8-fr]{3.1.8, (i)}) identifie $\mathcal{O}_X(nd)$ et
    $\mathcal{O}_{X^{(d)}}(n)$ pour tout $n\in\mathbf{Z}$.
\end{enumerate}

\begin{proposition}[3.2.10]
\label{II.3.2.10-fr}
Soient $\mathcal{L}$ un $\mathcal{O}_Y$-Module inversible, $g$ l'isomorphisme
canonique
$X_{(\mathcal{L})}=\operatorname{Proj}(\mathcal{S}_{(\mathcal{L})})
\to X=\operatorname{Proj}(\mathcal{S})$
(\hyperref[II.3.1.8-fr]{3.1.8, (iii)}). Pour tout entier $n\in\mathbf{Z}$,
$g_*(\mathcal{O}_{X_{(\mathcal{L})}}(n))$ est canoniquement isomorphe à
$\mathcal{O}_X(n)\otimes_Y\mathcal{L}^{\otimes n}$.
\end{proposition}

\begin{proof}
Supposons d'abord $Y$ affine, d'anneau $A$, et
$\mathcal{L}=\widetilde{L}$, où $L$ est un $A$-module libre monogène. Avec les
notations de la démonstration de
(\hyperref[II.3.1.8-fr]{3.1.8, (iii)}) on définit, pour $f\in S_d$, un
isomorphisme de $(S(n))_{(f)}\otimes_A L^{\otimes n}$ sur
$(S_{(L)}(n))_{(f\otimes c^{\otimes d})}$ en faisant correspondre à
$(x/f^k)\otimes c^{\otimes n}$, où $x\in S_{kd+n}$, l'élément
$(x\otimes c^{\otimes(n+kd)})/(f\otimes c^{\otimes d})^k$~; il est immédiat
que cet isomorphisme est indépendant du générateur $c$ choisi pour $L$~; en
outre, les isomorphismes ainsi définis pour chaque $f\in S_+$ sont compatibles
avec les opérations de restriction $D_+(f)\to D_+(fg)$. Enfin, dans le cas
général, on voit sans peine en revenant aux définitions
(\hyperref[II.3.1.1-fr]{3.1.1}) que les isomorphismes ainsi définis pour
chaque ouvert affine $U$ de $Y$ sont compatibles avec le passage de $U$ à un
ouvert affine $U'\subset U$.
\end{proof}

\subsection{$\mathcal{S}$-Module gradué associé à un faisceau sur
$\operatorname{Proj}(\mathcal{S})$.}
\label{subsection:II.3.3-fr}

\emph{On suppose dans tout ce numéro que la $\mathcal{O}_Y$-Algèbre graduée
$\mathcal{S}$ est engendrée par $\mathcal{S}_1$
(\hyperref[II.3.1.9-fr]{3.1.9}).} Rappelons qu'en raison de
(\hyperref[II.3.1.8-fr]{3.1.8, (i)}), cette restriction n'est pas essentielle,
moyennant les conditions de finitude (\hyperref[II.3.1.10-fr]{3.1.10}).

\begin{env}[3.3.1]
\label{II.3.3.1-fr}
Soit $p$ le morphisme structural $X=\operatorname{Proj}(\mathcal{S})\to Y$.
Pour tout $\mathcal{O}_X$-Module $\mathcal{F}$, nous poserons
\[
\label{II.3.3.1.1-fr}
  \Gamma_*(\mathcal{F})
  =\bigoplus_{n\in\mathbf{Z}}p_*(\mathcal{F}(n))
\tag{3.3.1.1}
\]

\oldpage[II]{57}

et en particulier
\[
\label{II.3.3.1.2-fr}
  \Gamma_*(\mathcal{O}_X)
  =\bigoplus_{n\in\mathbf{Z}}p_*(\mathcal{O}_X(n)).
\tag{3.3.1.2}
\]
On sait (\hyperref[0.4.2.2-fr]{0, 4.2.2}) qu'il existe un homomorphisme
canonique
\[
  p_*(\mathcal{F})\otimes_{\mathcal{O}_Y}p_*(\mathcal{G})
  \longrightarrow
  p_*(\mathcal{F}\otimes_{\mathcal{O}_X}\mathcal{G})
\]
pour deux $\mathcal{O}_X$-Modules $\mathcal{F}$ et $\mathcal{G}$~; on déduit
donc de (\hyperref[II.3.2.7.1-fr]{3.2.7.1}) que
$\Gamma_*(\mathcal{O}_X)$ est muni d'une structure de
$\mathcal{O}_Y$-\emph{Algèbre graduée}, et
(\hyperref[II.3.2.5.2-fr]{3.2.5.2}) définit de même sur
$\Gamma_*(\mathcal{F})$ une structure de \emph{Module gradué} sur
$\Gamma_*(\mathcal{O}_X)$.

En vertu de (\hyperref[II.3.2.5-fr]{3.2.5}), et de l'exactitude à gauche du
foncteur $p_*$ (\hyperref[0.4.2.1-fr]{0, 4.2.1}),
$\Gamma_*(\mathcal{F})$ est un foncteur covariant additif et
\emph{exact à gauche} en $\mathcal{F}$ dans la catégorie des
$\mathcal{O}_X$-Modules, à valeurs dans la catégorie des
$\mathcal{O}_Y$-Modules gradués (où les morphismes sont les homomorphismes de
degré $0$). En particulier, si $\mathcal{J}$ est un faisceau d'idéaux dans
$\mathcal{O}_X$, $\Gamma_*(\mathcal{J})$ s'identifie à un
\emph{faisceau gradué d'idéaux} de $\Gamma_*(\mathcal{O}_X)$.
\end{env}

\begin{env}[3.3.2]
\label{II.3.3.2-fr}
Soit $\mathcal{M}$ un $\mathcal{S}$-Module gradué quasi-cohérent. Pour tout
ouvert affine $U$ de $Y$, on a défini en
(\hyperref[II.2.6.2-fr]{2.6.2}) un homomorphisme de groupes abéliens
\[
  \alpha_{0,U}:\Gamma(U,\mathcal{M}_0)
  \longrightarrow\Gamma(p^{-1}(U),\widetilde{\mathcal{M}}).
\]
Il est immédiat que ces homomorphismes commutent aux opérations de restriction
(\hyperref[II.2.8.13.1-fr]{2.8.13.1}) et définissent donc (sans utiliser
l'hypothèse que $\mathcal{S}$ est engendrée par $\mathcal{S}_1$) un
homomorphisme de faisceaux de groupes abéliens
\[
\label{II.3.3.2.1-fr}
  \alpha_0:\mathcal{M}_0\longrightarrow p_*(\widetilde{\mathcal{M}}).
\tag{3.3.2.1}
\]
Appliquant ce résultat à chacun des
$\mathcal{M}_n=(\mathcal{M}(n))_0$, et tenant compte de
(\hyperref[II.3.2.8.1-fr]{3.2.8.1}), on définit un homomorphisme de faisceaux
de groupes abéliens
\[
\label{II.3.3.2.2-fr}
  \alpha_n:\mathcal{M}_n\longrightarrow
  p_*(\widetilde{\mathcal{M}}(n))
  \qquad\text{pour tout }n\in\mathbf{Z},
\tag{3.3.2.2}
\]
d'où un homomorphisme fonctoriel (de degré $0$) de faisceaux gradués de groupes
abéliens
\[
\label{II.3.3.2.3-fr}
  \alpha:\mathcal{M}\longrightarrow\Gamma_*(\widetilde{\mathcal{M}})
\tag{3.3.2.3}
\]
(aussi noté $\alpha_{\mathcal{M}}$).

En prenant en particulier $\mathcal{M}=\mathcal{S}$, on vérifie que
$\alpha:\mathcal{S}\to\Gamma_*(\mathcal{O}_X)$ est un homomorphisme de
$\mathcal{O}_Y$-Algèbres graduées et que
(\hyperref[II.3.3.2.3-fr]{3.3.2.3}) est un di-homomorphisme de Modules
gradués, relatif à cet homomorphisme d'Algèbres graduées.

Remarquons encore qu'à chacun des $\alpha_n$ il correspond
(\hyperref[0.4.4.3-fr]{0, 4.4.3}) un homomorphisme canonique de
$\mathcal{O}_X$-Modules
\[
\label{II.3.3.2.4-fr}
  \alpha_n^\sharp:p^*(\mathcal{M}_n)
  \longrightarrow\widetilde{\mathcal{M}}(n).
\tag{3.3.2.4}
\]
On vérifie sans peine que cet homomorphisme n'est autre que celui qui
correspond fonctoriellement (\hyperref[II.3.2.4-fr]{3.2.4}) à
l'homomorphisme canonique (de degré $0$) de $\mathcal{O}_Y$-Modules gradués
\[
\label{II.3.3.2.5-fr}
  \mathcal{M}_n\otimes_{\mathcal{O}_Y}\mathcal{S}
  \longrightarrow\mathcal{M}(n)
\tag{3.3.2.5}
\]

\oldpage[II]{58}

où la graduation du second membre provient naturellement de celle de
$\mathcal{S}$. On peut en effet se borner au cas où
$Y=\operatorname{Spec}(A)$ est affine, $\mathcal{M}=\widetilde{M}$ et
$\mathcal{S}=\widetilde{S}$, la $A$-algèbre graduée $S$ étant engendrée par
$S_1$, de sorte que lorsque $f$ parcourt $S_1$, les $D_+(f)$ forment un
recouvrement de $X$. Revenant aux définitions
(\hyperref[II.2.6.2-fr]{2.6.2}) on voit alors, compte tenu de
(\hyperref[I.1.6.7-fr]{I, 1.6.7}), que la restriction à $D_+(f)$ de
l'homomorphisme (\hyperref[II.3.3.2.4-fr]{3.3.2.4}) correspond
(\hyperref[I.1.3.8-fr]{I, 1.3.8}) à l'homomorphisme de $S_{(f)}$-modules
$M_n\otimes_A S_{(f)}\to(S(n))_{(f)}$ qui, à $x\otimes1$ ($x\in M_n$), fait
correspondre $x/1$~; cela démontre notre assertion.
\end{env}

\begin{proposition}[3.3.3]
\label{II.3.3.3-fr}
Pour toute section $f\in\Gamma(Y,\mathcal{S}_d)$ ($d>0$), $X_f$ est identique
à l'ensemble des points de $X$ où $\alpha_d(f)$ (considérée comme section de
$\mathcal{O}_X(d)$) ne s'annule pas
(\hyperref[0.5.5.2-fr]{0, 5.5.2}).
\end{proposition}

\begin{proof}
($\alpha_d(f)$ est une section de $p_*(\mathcal{O}_X(d))$ au-dessus de $Y$,
mais par définition une telle section est aussi une section de
$\mathcal{O}_X(d)$ au-dessus de $X$
(\hyperref[0.4.2.1-fr]{0, 4.2.1})). La définition de $X_f$
(\hyperref[II.3.1.4-fr]{3.1.4}) ramène au cas où $Y$ est affine, qui a été
traité dans (\hyperref[II.2.6.3-fr]{2.6.3}).
\end{proof}

\begin{env}[3.3.4]
\label{II.3.3.4-fr}
Nous supposerons désormais, outre l'hypothèse du début de ce numéro, que pour
tout $\mathcal{O}_X$-Module quasi-cohérent $\mathcal{F}$, les
$p_*(\mathcal{F}(n))$ sont \emph{quasi-cohérents} sur $Y$, et par suite
$\Gamma_*(\mathcal{F})=\bigoplus_{n\in\mathbf{Z}}p_*(\mathcal{F}(n))$ est
aussi un $\mathcal{O}_Y$-Module quasi-cohérent
(\hyperref[I.1.4.1-fr]{I, 1.4.1} et
\hyperref[I.1.3.9-fr]{I, 1.3.9})~; cette circonstance se produira toujours en
particulier si $X$ est \emph{de type fini} sur $Y$
(\hyperref[I.9.2.2-fr]{I, 9.2.2}). On en conclut que
$(\Gamma_*(\mathcal{F}))^{\sim}$ est défini et est un
$\mathcal{O}_X$-Module quasi-cohérent. Pour tout ouvert affine $U$ de $Y$, on
a (\hyperref[I.1.3.9-fr]{I, 1.3.9} et
\hyperref[II.2.5.4-fr]{2.5.4})
\begin{align*}
  \left(\Gamma\left(U,\bigoplus_{n\in\mathbf{Z}}
  p_*(\mathcal{F}(n))\right)\right)^{\sim}
  &=\bigoplus_{n\in\mathbf{Z}}
  \bigl(\Gamma(U,p_*(\mathcal{F}(n)))\bigr)^{\sim}\\
  &=\bigoplus_{n\in\mathbf{Z}}
  \bigl(\Gamma(p^{-1}(U),\mathcal{F}(n))\bigr)^{\sim}\\
  &=\left(\bigoplus_{n\in\mathbf{Z}}
  \Gamma(p^{-1}(U),\mathcal{F}(n))\right)^{\sim}\\
  &=\bigl(\Gamma_*(\mathcal{F}|p^{-1}(U))\bigr)^{\sim}
\end{align*}
et par suite (\hyperref[II.2.6.4-fr]{2.6.4}) on a un homomorphisme canonique
\[
  \beta_U:\left(\Gamma\left(U,\bigoplus_{n\in\mathbf{Z}}
  p_*(\mathcal{F}(n))\right)\right)^{\sim}
  \longrightarrow\mathcal{F}|p^{-1}(U).
\]

En outre, la commutativité de
(\hyperref[II.2.8.13.2-fr]{2.8.13.2}) montre que ces homomorphismes commutent
avec les opérations de restriction sur $Y$~; on en déduit donc un
homomorphisme canonique fonctoriel
\[
\label{II.3.3.4.1-fr}
  \beta:(\Gamma_*(\mathcal{F}))^{\sim}\longrightarrow\mathcal{F}
\tag{3.3.4.1}
\]
(aussi noté $\beta_{\mathcal{F}}$) pour les $\mathcal{O}_X$-Modules
quasi-cohérents.
\end{env}

\begin{proposition}[3.3.5]
\label{II.3.3.5-fr}
Soient $\mathcal{M}$ un $\mathcal{S}$-Module gradué quasi-cohérent,
$\mathcal{F}$ un $\mathcal{O}_X$-Module quasi-cohérent~; les homomorphismes
composés
\[
\label{II.3.3.5.1-fr}
  \widetilde{\mathcal{M}}
  \xrightarrow{\widetilde{\alpha}}
  (\Gamma_*(\widetilde{\mathcal{M}}))^{\sim}
  \xrightarrow{\beta}\widetilde{\mathcal{M}}
\tag{3.3.5.1}
\]
\[
\label{II.3.3.5.2-fr}
  \Gamma_*(\mathcal{F})\xrightarrow{\alpha}
  \Gamma_*((\Gamma_*(\mathcal{F}))^{\sim})
  \xrightarrow{\Gamma_*(\beta)}\Gamma_*(\mathcal{F})
\tag{3.3.5.2}
\]
sont les isomorphismes identiques.
\end{proposition}

\begin{proof}
La question est en effet locale sur $Y$, et on est ramené à
(\hyperref[II.2.6.5-fr]{2.6.5}).
\end{proof}

\subsection{Conditions de finitude.}
\label{subsection:II.3.4-fr}

\oldpage[II]{59}

\begin{proposition}[3.4.1]
\label{II.3.4.1-fr}
Soient $Y$ un préschéma, $\mathcal{S}$ une $\mathcal{O}_Y$-Algèbre
quasi-cohérente engendrée par $\mathcal{S}_1$
(\hyperref[II.3.1.9-fr]{3.1.9})~; on suppose en outre $\mathcal{S}_1$ de type
fini. Alors $X=\operatorname{Proj}(\mathcal{S})$ est de type fini au-dessus de
$Y$.
\end{proposition}

\begin{proof}
En effet, la question étant locale sur $Y$, on peut supposer $Y$ affine
d'anneau $A$~; on a alors $\mathcal{S}=\widetilde{S}$, où
$S=\Gamma(Y,\mathcal{S})$, et par hypothèse $S$ est une $A$-algèbre engendrée
par $S_1=\Gamma(Y,\mathcal{S}_1)$, où on peut en outre supposer que $S_1$ est
un $A$-module de type fini
(\hyperref[I.1.3.9-fr]{I, 1.3.9} et
\hyperref[I.1.3.12-fr]{I, 1.3.12}). Par suite $S$ est une $A$-algèbre graduée
de type fini, et on est ramené à
(\hyperref[II.2.7.1-fr]{2.7.1, (ii)}).
\end{proof}

\begin{env}[3.4.2]
\label{II.3.4.2-fr}
Soit $\mathcal{S}$ une $\mathcal{O}_Y$-Algèbre graduée quasi-cohérente~; pour
un $\mathcal{S}$-Module gradué quasi-cohérent $\mathcal{M}$, nous considérerons
les conditions de finitude suivantes~:
\begin{enumerate}
  \item[(\textbf{TF})] \emph{Il existe un entier $n$ tel que le
    $\mathcal{S}$-Module $\bigoplus_{k\geq n}\mathcal{M}_k$ soit de type fini.}
  \item[(\textbf{TN})] \emph{Il existe un entier $n$ tel que
    $\mathcal{M}_k=0$ pour $k\geq n$.}
\end{enumerate}
Si $\mathcal{M}$ vérifie (\textbf{TN}), on a
$\widetilde{\mathcal{M}}=0$, la question étant locale sur $Y$
(\hyperref[II.2.7.2-fr]{2.7.2}).

Soient $\mathcal{M}$, $\mathcal{N}$ deux $\mathcal{S}$-Modules gradués
quasi-cohérents~; on dit qu'un homomorphisme
$u:\mathcal{M}\to\mathcal{N}$ de degré $0$ est
(\textbf{TN})-\emph{injectif} (resp. (\textbf{TN})-\emph{surjectif},
(\textbf{TN})-\emph{bijectif}) s'il existe un entier $n$ tel que
$u_k:\mathcal{M}_k\to\mathcal{N}_k$ soit injectif (resp. surjectif, bijectif)
pour $k\geq n$~; alors
$\widetilde{u}:\widetilde{\mathcal{M}}\to\widetilde{\mathcal{N}}$ est injectif
(resp. surjectif, bijectif) en vertu de
(\hyperref[II.2.7.2-fr]{2.7.2}), la question étant locale sur $Y$, et compte
tenu de (\hyperref[I.1.3.9-fr]{I, 1.3.9})~; lorsque $u$ est
(\textbf{TN})-bijectif, on dit aussi que $u$ est un
(\textbf{TN})-\emph{isomorphisme}.
\end{env}

\begin{proposition}[3.4.3]
\label{II.3.4.3-fr}
Soient $Y$ un préschéma, $\mathcal{S}$ une $\mathcal{O}_Y$-Algèbre graduée
quasi-cohérente engendrée par $\mathcal{S}_1$, $\mathcal{S}_1$ étant supposé
de type fini. Soit $\mathcal{M}$ un $\mathcal{S}$-Module gradué
quasi-cohérent.
\begin{enumerate}
  \item[\rm(i)] Si $\mathcal{M}$ vérifie la condition (\textbf{TF}),
    $\widetilde{\mathcal{M}}$ est de type fini.
  \item[\rm(ii)] Supposons que $\mathcal{M}$ vérifie (\textbf{TF})~; pour que
    $\widetilde{\mathcal{M}}=0$, il faut et il suffit que $\mathcal{M}$
    vérifie (\textbf{TN}).
\end{enumerate}
\end{proposition}

\begin{proof}
Les questions étant locales sur $Y$, on est ramené au cas où $Y$ est affine
d'anneau $A$, $\mathcal{S}=\widetilde{S}$, où $S$ est une $A$-algèbre graduée
telle que l'idéal $S_+$ soit de type fini,
$\mathcal{M}=\widetilde{M}$ où $M$ est un $S$-module gradué~; la proposition
résulte alors de (\hyperref[II.2.7.3-fr]{2.7.3}).
\end{proof}

\begin{theorem}[3.4.4]
\label{II.3.4.4-fr}
Soient $Y$ un préschéma, $\mathcal{S}$ une $\mathcal{O}_Y$-Algèbre graduée
quasi-cohérente engendrée par $\mathcal{S}_1$, $\mathcal{S}_1$ étant supposé
de type fini~; soit $X=\operatorname{Proj}(\mathcal{S})$. Pour tout
$\mathcal{O}_X$-Module quasi-cohérent $\mathcal{F}$, l'homomorphisme canonique
$\beta$ (\hyperref[II.3.3.4-fr]{3.3.4}) est un isomorphisme.
\end{theorem}

\begin{proof}
Notons d'abord que $\beta$ est défini en vertu de
(\hyperref[II.3.4.1-fr]{3.4.1}). Pour voir que $\beta$ est un isomorphisme,
on est ramené au cas où $Y$ est affine d'anneau $A$,
$\mathcal{S}=\widetilde{S}$, où $S$ est une $A$-algèbre graduée engendrée par
$S_1$ et $S_1$ est un $A$-module de type fini. Il suffit alors d'appliquer
(\hyperref[II.2.7.5-fr]{2.7.5}).
\end{proof}

\begin{corollary}[3.4.5]
\label{II.3.4.5-fr}
Sous les hypothèses de (\hyperref[II.3.4.4-fr]{3.4.4}), tout
$\mathcal{O}_X$-Module quasi-cohérent $\mathcal{F}$ est isomorphe à un
$\mathcal{O}_X$-Module de la forme $\widetilde{\mathcal{M}}$, où
$\mathcal{M}$ est un $\mathcal{S}$-Module gradué quasi-cohérent. Si en outre
$\mathcal{F}$ est de type fini, et si on suppose que $Y$ est un schéma
quasi-compact, ou que l'espace sous-jacent à $Y$ est noethérien, on peut
supposer $\mathcal{M}$ de type fini.
\end{corollary}

\oldpage[II]{60}

\begin{proof}
La première assertion résulte aussitôt de
(\hyperref[II.3.4.4-fr]{3.4.4}) en prenant
$\mathcal{M}=\Gamma_*(\mathcal{F})$. Pour établir la seconde, il suffira
d'établir que $\mathcal{M}$ est limite inductive de ses sous-$\mathcal{S}$-Modules
gradués de type fini $\mathcal{N}_\lambda$~: en effet, il en résultera que
$\widetilde{\mathcal{M}}$ est limite inductive des
$\widetilde{\mathcal{N}}_\lambda$
(\hyperref[II.3.2.4-fr]{3.2.4}), donc $\mathcal{F}$ est limite inductive des
$\beta(\widetilde{\mathcal{N}}_\lambda)$~; comme $X$ est quasi-compact
((\hyperref[II.3.4.1-fr]{3.4.1}) et
\hyperref[I.6.3.1-fr]{I, 6.3.1}) et que $\mathcal{F}$ est de type fini,
$\mathcal{F}$ sera nécessairement égal à un des
$\beta(\widetilde{\mathcal{N}}_\lambda)$
(\hyperref[0.5.2.3-fr]{0, 5.2.3}).

Pour définir les $\mathcal{N}_\lambda$ ayant $\mathcal{M}$ pour limite
inductive, il suffit de considérer pour chaque $n\in\mathbf{Z}$ le
$\mathcal{O}_Y$-Module quasi-cohérent $\mathcal{M}_n$, qui est limite
inductive de ses sous-$\mathcal{O}_Y$-Modules
$\mathcal{M}_n^{(\mu_n)}$ de type fini, en raison des hypothèses sur $Y$
(\hyperref[I.9.4.9-fr]{I, 9.4.9})~; il est immédiat que
$\mathcal{P}_{\mu_n}=\mathcal{S}\cdot\mathcal{M}_n^{(\mu_n)}$ est un
$\mathcal{S}$-Module gradué de type fini, et on vérifie aussitôt qu'en prenant
pour $\mathcal{N}_\lambda$ les sommes finies de $\mathcal{S}$-Modules de la
forme $\mathcal{P}_{\mu_n}$, on répond bien à la question.
\end{proof}

\begin{corollary}[3.4.6]
\label{II.3.4.6-fr}
Supposons vérifiées les hypothèses de (\hyperref[II.3.4.4-fr]{3.4.4}), et en
outre que l'espace sous-jacent $Y$ soit quasi-compact~; soit $\mathcal{F}$ un
$\mathcal{O}_X$-Module quasi-cohérent de type fini. Il existe $n_0$ tel que
pour $n\geq n_0$, l'homomorphisme canonique
$\sigma:p^*(p_*(\mathcal{F}(n)))\to\mathcal{F}(n)$
(\hyperref[0.4.4.3-fr]{0, 4.4.3}) soit surjectif.
\end{corollary}

\begin{proof}
En effet, pour tout $y\in Y$, soit $U$ un voisinage ouvert affine de $y$ dans
$Y$. Il existe un entier $n_0(U)$ tel que, pour $n\geq n_0(U)$,
$\mathcal{F}(n)|p^{-1}(U)$ soit engendré par un nombre fini de ses sections
au-dessus de $p^{-1}(U)$ (\hyperref[II.2.7.9-fr]{2.7.9})~; mais ces dernières
sont images canoniques de sections de $p^*(p_*(\mathcal{F}(n)))$ au-dessus de
$p^{-1}(U)$ (\hyperref[0.3.7.1-fr]{0, 3.7.1} et
\hyperref[0.4.4.3-fr]{0, 4.4.3}), donc
$\mathcal{F}(n)|p^{-1}(U)$ est égal à l'image canonique de
$p^*(p_*(\mathcal{F}(n)))|p^{-1}(U)$. Enfin, comme $Y$ est quasi-compact, il
existe un recouvrement fini de $Y$ par des ouverts affines $U_i$, et en
prenant pour $n_0$ le plus grand des $n_0(U_i)$, on achève la démonstration.
\end{proof}

\par\medskip\noindent\textit{Remarques (3.4.7). ---\ }
\label{II.3.4.7-fr}
Si $p=(\psi,\theta):X\to Y$ est un morphisme d'espaces annelés et
$\mathcal{F}$ un $\mathcal{O}_X$-Module, le fait que l'homomorphisme canonique
$\sigma:p^*(p_*(\mathcal{F}))\to\mathcal{F}$ soit surjectif s'explicite de la
façon suivante (\hyperref[0.4.4.1-fr]{0, 4.4.1})~: pour tout $x\in X$ et toute
section $s$ de $\mathcal{F}$ au-dessus d'un voisinage ouvert $V$ de $x$, il
existe un voisinage ouvert $U$ de $p(x)$ dans $Y$, un nombre fini de sections
$t_i$ ($1\leq i\leq m$) de $\mathcal{F}$ au-dessus de $p^{-1}(U)$, un
voisinage $W\subset V\cap p^{-1}(U)$ de $x$ et des sections $a_i$
($1\leq i\leq m$) de $\mathcal{O}_X$ au-dessus de $W$ telles que
\[
  s|W=\sum_i a_i\cdot(t_i|W).
\]
Lorsque $Y$ est un \emph{schéma affine} et $p_*(\mathcal{F})$
\emph{quasi-cohérent}, cette condition équivaut au fait que $\mathcal{F}$ est
\emph{engendré par ses sections au-dessus de $X$}
(\hyperref[0.5.5.1-fr]{0, 5.5.1})~: en effet, si
$Y=\operatorname{Spec}(A)$, on peut supposer que $U=D(f)$ avec $f\in A$~; il
existe un entier $n>0$ et des sections $s_i$ de $\mathcal{F}$ au-dessus de
$X$ telles que $t_i$ soit la restriction à $p^{-1}(U)$ de $s_ig^n$, avec
$g=\theta(f)$ (en appliquant (\hyperref[I.1.4.1-fr]{I, 1.4.1}) à
$p_*(\mathcal{F})$)~; comme $g$ est inversible au-dessus de $p^{-1}(U)$, on a
donc
\[
  s|W=\sum_i b_i\cdot(s_i|W)
\]
avec $b_i=a_i(g|W)^{-n}$, d'où notre assertion. Lorsque $Y$ est affine, le
corollaire (\hyperref[II.3.4.6-fr]{3.4.6}) redonne donc
(\hyperref[II.2.7.9-fr]{2.7.9}).

\oldpage[II]{61}

On en conclut que lorsque $Y$ est un préschéma quelconque, les trois
conditions suivantes, pour un $\mathcal{O}_X$-Module quasi-cohérent
$\mathcal{F}$ tel que $p_*(\mathcal{F})$ soit \emph{quasi-cohérent}, sont
équivalentes~:
\begin{enumerate}
  \item[a)] \emph{L'homomorphisme canonique
    $\sigma:p^*(p_*(\mathcal{F}))\to\mathcal{F}$ est surjectif.}
  \item[b)] \emph{Il existe un $\mathcal{O}_Y$-Module quasi-cohérent
    $\mathcal{G}$ et un homomorphisme surjectif
    $p^*(\mathcal{G})\to\mathcal{F}$.}
  \item[c)] \emph{Pour tout ouvert affine $U$ de $Y$,
    $\mathcal{F}|p^{-1}(U)$ est engendré par ses sections au-dessus de
    $p^{-1}(U)$.}
\end{enumerate}
On vient en effet de prouver l'équivalence de \emph{a)} et \emph{c)}. D'autre
part, il est clair que \emph{a)} entraîne \emph{b)}, $p_*(\mathcal{F})$ étant
par hypothèse quasi-cohérent. Inversement, tout homomorphisme
$u:p^*(\mathcal{G})\to\mathcal{F}$ se factorise en
$p^*(\mathcal{G})\to p^*(p_*(\mathcal{F}))\xrightarrow{\sigma}\mathcal{F}$
(\hyperref[0.3.5.4.4-fr]{0, 3.5.4.4}), donc si $u$ est surjectif il en est de
même de $\sigma$, ce qui prouve que \emph{b)} entraîne \emph{a)}.

\begin{corollary}[3.4.8]
\label{II.3.4.8-fr}
Supposons vérifiées les hypothèses de (\hyperref[II.3.4.4-fr]{3.4.4}) et
supposons en outre que $Y$ soit un schéma quasi-compact ou que l'espace
sous-jacent à $Y$ soit noethérien. Soit $\mathcal{F}$ un
$\mathcal{O}_X$-Module quasi-cohérent de type fini~; il existe alors un entier
$n_0$ tel que pour $n\geq n_0$, $\mathcal{F}$ soit isomorphe à un quotient
d'un $\mathcal{O}_X$-Module de la forme $(p^*(\mathcal{G}))(-n)$ où
$\mathcal{G}$ est un $\mathcal{O}_Y$-Module quasi-cohérent de type fini
(dépendant de $n$).
\end{corollary}

\begin{proof}
Comme le morphisme structural $X\to Y$ est séparé et de type fini,
$p_*(\mathcal{F}(n))$ est quasi-cohérent
(\hyperref[I.9.2.2-fr]{I, 9.2.2, b)}), donc limite inductive de ses
sous-$\mathcal{O}_Y$-Modules quasi-cohérents de type fini, en vertu de
l'hypothèse sur $Y$ (\hyperref[I.9.4.9-fr]{I, 9.4.9}). On en déduit par
(\hyperref[II.3.4.6-fr]{3.4.6}),
(\hyperref[0.4.3.2-fr]{0, 4.3.2}) et
(\hyperref[0.5.2.3-fr]{0, 5.2.3}) que $\mathcal{F}(n)$ est l'image canonique
d'un $\mathcal{O}_X$-Module de la forme $p^*(\mathcal{G})$, où $\mathcal{G}$
est un sous-$\mathcal{O}_Y$-Module quasi-cohérent de type fini de
$p_*(\mathcal{F}(n))$~; le corollaire résulte alors de
(\hyperref[II.3.2.5.2-fr]{3.2.5.2}) et
(\hyperref[II.3.2.7.1-fr]{3.2.7.1}).
\end{proof}

\subsection{Comportements fonctoriels.}
\label{subsection:II.3.5-fr}

\begin{env}[3.5.1]
\label{II.3.5.1-fr}
Soient $Y$ un préschéma, $\mathcal{S}$, $\mathcal{S}'$ deux
$\mathcal{O}_Y$-Algèbres graduées quasi-cohérentes à degrés positifs~; posons
$X=\operatorname{Proj}(\mathcal{S})$, $X'=\operatorname{Proj}(\mathcal{S}')$,
et soient $p$, $p'$ les morphismes structuraux de $X$ et $X'$ dans $Y$. Soit
$\varphi:\mathcal{S}'\to\mathcal{S}$ un $\mathcal{O}_Y$-homomorphisme
d'Algèbres graduées. Pour tout ouvert affine $U$ de $Y$, posons
$S_U=\Gamma(U,\mathcal{S})$, $S'_U=\Gamma(U,\mathcal{S}')$~;
l'homomorphisme $\varphi$ définit un homomorphisme
$\varphi_U:S'_U\to S_U$ de $A_U$-algèbres graduées, en posant
$A_U=\Gamma(U,\mathcal{O}_Y)$. Il lui correspond dans $p^{-1}(U)$ un ensemble
ouvert $G(\varphi_U)$ et un morphisme
$\Phi_U:G(\varphi_U)\to p'^{-1}(U)$
(\hyperref[II.2.8.1-fr]{2.8.1}). En outre, si $V\subset U$ est un ouvert
affine, le diagramme
\[
\label{II.3.5.1.1-fr}
  \xymatrix{
    S'_U\ar[r]^{\varphi_U}\ar[d]
    & S_U\ar[d]\\
    S'_V\ar[r]_{\varphi_V}
    & S_V
  }
\tag{3.5.1.1}
\]
est commutatif, et on vérifie aussitôt, à partir des définitions
(\hyperref[II.2.8.1-fr]{2.8.1}), que l'on a
\[
  G(\varphi_V)=G(\varphi_U)\cap p^{-1}(V)
\]
et que $\Phi_V$ est la restriction de $\Phi_U$ à $G(\varphi_V)$. On a ainsi
défini une partie ouverte $G(\varphi)$ de $X$ telle que
$G(\varphi)\cap p^{-1}(U)=G(\varphi_U)$ pour tout ouvert affine $U\subset Y$,
et un $Y$-morphisme

\oldpage[II]{62}

affine $\Phi:G(\varphi)\to X'$, qui est dit \emph{associé} à $\varphi$ et que
nous noterons $\operatorname{Proj}(\varphi)$. Lorsque, pour tout $y\in Y$, il
y a un voisinage affine $U$ de $y$ tel que le
$\Gamma(U,\mathcal{O}_Y)$-module $\Gamma(U,\mathcal{S}_+)$ soit engendré par
$\varphi(\Gamma(U,\mathcal{S}'_+))$, on a
$G(\varphi_U)=p^{-1}(U)$, et par suite $G(\varphi)=X$.
\end{env}

\begin{proposition}[3.5.2]
\label{II.3.5.2-fr}
\begin{enumerate}
  \item[\rm(i)] Si $\mathcal{M}$ est un $\mathcal{S}$-Module gradué
    quasi-cohérent, il existe un isomorphisme canonique fonctoriel du
    $\mathcal{O}_{X'}$-Module $(\mathcal{M}_{[\varphi]})^{\sim}$ sur le
    $\mathcal{O}_{X'}$-Module
    $\Phi_*(\widetilde{\mathcal{M}}|G(\varphi))$.
  \item[\rm(ii)] Si $\mathcal{M}'$ est un $\mathcal{S}'$-Module gradué
    quasi-cohérent, il existe un homomorphisme canonique fonctoriel $\nu$ du
    $(\mathcal{O}_X|G(\varphi))$-Module
    $\Phi^*(\widetilde{\mathcal{M}'})$ dans le
    $(\mathcal{O}_X|G(\varphi))$-Module
    $(\mathcal{M}'\otimes_{\mathcal{S}'}\mathcal{S})^{\sim}|G(\varphi)$. Si
    $\mathcal{S}'$ est engendrée par $\mathcal{S}'_1$, $\nu$ est un
    isomorphisme.
\end{enumerate}
\end{proposition}

\begin{proof}
Les homomorphismes considérés ont en effet déjà été définis lorsque $Y$ est
affine (\hyperref[II.2.8.7-fr]{2.8.7} et
\hyperref[II.2.8.8-fr]{2.8.8}), et dans le cas général il suffit de vérifier
qu'ils sont compatibles avec les opérations de restriction d'un ouvert
affine de $Y$ à un ouvert plus petit, ce qui résulte aussitôt de la
commutativité de (\hyperref[II.3.5.1.1-fr]{3.5.1.1}).
\end{proof}

En particulier, pour tout $n\in\mathbf{Z}$, on a un homomorphisme canonique
\[
\label{II.3.5.2.1-fr}
  \Phi^*(\mathcal{O}_{X'}(n))\longrightarrow
  \mathcal{O}_X(n)|G(\varphi).
\tag{3.5.2.1}
\]

\begin{proposition}[3.5.3]
\label{II.3.5.3-fr}
Soient $Y$, $Y'$ deux préschémas, $\psi:Y'\to Y$ un morphisme,
$\mathcal{S}$ une $\mathcal{O}_Y$-Algèbre graduée quasi-cohérente, et posons
$\mathcal{S}'=\psi^*(\mathcal{S})$. Alors le $Y'$-schéma
$X'=\operatorname{Proj}(\mathcal{S}')$ s'identifie canoniquement à
$\operatorname{Proj}(\mathcal{S})\times_Y Y'$. En outre, si $\mathcal{M}$ est
un $\mathcal{S}$-Module gradué quasi-cohérent, le $\mathcal{O}_{X'}$-Module
$(\psi^*(\mathcal{M}))^{\sim}$ s'identifie à
$\widetilde{\mathcal{M}}\otimes_Y\mathcal{O}_{Y'}$.
\end{proposition}

\begin{proof}
Notons en premier lieu que $\psi^*(\mathcal{S})$ et
$\psi^*(\mathcal{M})$ sont des $\mathcal{O}_{Y'}$-Modules quasi-cohérents
ainsi que leurs composants homogènes
(\hyperref[0.5.1.4-fr]{0, 5.1.4}). Soient $U$ un ouvert affine de $Y$,
$U'\subset\psi^{-1}(U)$ un ouvert affine de $Y'$, $A$, $A'$ les anneaux de
$U$ et $U'$ respectivement~; on a alors
$\mathcal{S}|U=\widetilde{S}$, où $S$ est une $A$-algèbre graduée, et
$\mathcal{S}'|U'$ s'identifie à $(S\otimes_A A')^{\sim}$
(\hyperref[I.1.6.5-fr]{I, 1.6.5})~; la première assertion résulte alors de
(\hyperref[II.2.8.10-fr]{2.8.10}) et de
(\hyperref[I.3.2.6.2-fr]{I, 3.2.6.2}), car on vérifie aussitôt que la
projection
$\operatorname{Proj}(\mathcal{S}'|U')\to
\operatorname{Proj}(\mathcal{S}|U)$ définie par l'identification précédente
est compatible avec les opérations de restrictions sur $U$ et $U'$ et
définit donc bien un morphisme
$\operatorname{Proj}(\mathcal{S}')\to\operatorname{Proj}(\mathcal{S})$.
Soient maintenant
\[
  q:\operatorname{Proj}(\mathcal{S})\to Y,
  \qquad q':\operatorname{Proj}(\mathcal{S}')\to Y'
\]
les morphismes structuraux~; $q'^{-1}(U')$ s'identifie alors à
$q^{-1}(U)\times_U U'$, et les deux faisceaux
$(\psi^*(\mathcal{M}))^{\sim}|q'^{-1}(U')$ et
$(\widetilde{\mathcal{M}}\otimes_Y\mathcal{O}_{Y'})|q'^{-1}(U')$
s'identifient alors canoniquement tous deux à $(M\otimes_A A')^{\sim}$, où
on a posé $M=\Gamma(U,\mathcal{M})$, en vertu de
(\hyperref[II.2.8.10-fr]{2.8.10}) et de
(\hyperref[I.1.6.5-fr]{I, 1.6.5})~; d'où la seconde assertion, car on vérifie
encore de façon immédiate la compatibilité des identifications précédentes
avec les opérations de restriction.
\end{proof}

\begin{corollary}[3.5.4]
\label{II.3.5.4-fr}
Avec les notations de (\hyperref[II.3.5.3-fr]{3.5.3}),
$\mathcal{O}_{X'}(n)$ s'identifie canoniquement à
$\mathcal{O}_X(n)\otimes_Y\mathcal{O}_{Y'}$ pour tout $n\in\mathbf{Z}$ (avec
$X=\operatorname{Proj}(\mathcal{S})$).
\end{corollary}

\begin{proof}
En effet, avec les notations de (\hyperref[II.3.5.3-fr]{3.5.3}), il est clair
que $\psi^*(\mathcal{S}(n))=\mathcal{S}'(n)$ pour tout $n\in\mathbf{Z}$.
\end{proof}

\begin{env}[3.5.5]
\label{II.3.5.5-fr}
Gardant les notations précédentes, désignons par $\Psi$ la projection
canonique $X'\to X$, et posons $\mathcal{M}'=\psi^*(\mathcal{M})$~; nous
supposerons en outre que $\mathcal{S}$ est engendrée par $\mathcal{S}_1$

\oldpage[II]{63}

et que $X$ est de type fini sur $Y$~; il en résulte alors que $\mathcal{S}'$
est engendrée par $\mathcal{S}'_1$ (comme on le voit en se ramenant au cas où
$Y$ et $Y'$ sont affines) et que $X'$ est de type fini sur $Y'$
(\hyperref[I.6.3.4-fr]{I, 6.3.4}). Soit $\mathcal{F}$ un
$\mathcal{O}_X$-Module et posons $\mathcal{F}'=\Psi^*(\mathcal{F})$~; il
résulte alors de (\hyperref[II.3.5.4-fr]{3.5.4}) et de
(\hyperref[0.4.3.3-fr]{0, 4.3.3}) que l'on a
$\mathcal{F}'(n)=\Psi^*(\mathcal{F}(n))$ pour tout $n\in\mathbf{Z}$. On
définit en outre un $\psi$-homomorphisme canonique
$\theta_n:q_*(\mathcal{F}(n))\to q'_*(\mathcal{F}'(n))$ de la façon
suivante~: vu la commutativité du diagramme
\[
  \xymatrix{
    X\ar[d]_{q}
    & X'\ar[l]_{\Psi}\ar[d]^{q'}\\
    Y
    & Y'\ar[l]^{\psi}
  }
\]
il s'agit de définir un homomorphisme
$q_*(\mathcal{F}(n))\to
\psi_*(q'_*(\Psi^*(\mathcal{F}(n))))
=q_*(\Psi_*(\Psi^*(\mathcal{F}(n))))$, et il suffit de prendre
l'homomorphisme $\theta_n=q_*(\rho_n)$, où $\rho_n$ est l'homomorphisme
canonique
$\mathcal{F}(n)\to\Psi_*(\Psi^*(\mathcal{F}(n)))$
(\hyperref[0.4.4.3-fr]{0, 4.4.3}). Il est immédiat que pour tout ouvert affine
$U$ de $Y$ et tout ouvert affine $U'$ de $Y'$ tel que
$U'\subset\psi^{-1}(U)$, l'homomorphisme $\theta_n$ donne pour les sections
l'homomorphisme canonique (\hyperref[0.3.7.2-fr]{0, 3.7.2})
$\Gamma(q^{-1}(U),\mathcal{F}(n))\to
\Gamma(q'^{-1}(U'),\mathcal{F}'(n))$.

La commutativité de (\hyperref[II.2.8.13.2-fr]{2.8.13.2}) montre alors que si
$\mathcal{F}$ est quasi-cohérent, le diagramme
\[
  \xymatrix{
    \mathcal{F}\ar[r]^{\rho}
    & \mathcal{F}'\\
    (\Gamma_*(\mathcal{F}))^{\sim}
      \ar[u]^{\beta_{\mathcal{F}}}\ar[r]_{\widetilde{\theta}}
    & (\Gamma_*(\mathcal{F}'))^{\sim}
      \ar[u]_{\beta_{\mathcal{F}'}}
  }
\]
est commutatif (la flèche horizontale du haut étant le $\Psi$-morphisme
canonique $\mathcal{F}\to\Psi^*(\mathcal{F})$).

De même, la commutativité de
(\hyperref[II.2.8.13.1-fr]{2.8.13.1}) montre que le diagramme
\[
  \xymatrix{
    \Gamma_*(\widetilde{\mathcal{M}})\ar[r]^{\theta}
    & \Gamma_*(\widetilde{\mathcal{M}'})\\
    \mathcal{M}\ar[r]_{\rho}\ar[u]^{\alpha_{\mathcal{M}}}
    & \mathcal{M}'\ar[u]_{\alpha_{\mathcal{M}'}}
  }
\]
est commutatif (la flèche horizontale du bas étant le $\psi$-morphisme
canonique $\mathcal{M}\to\psi^*(\mathcal{M})$).
\end{env}

\begin{env}[3.5.6]
\label{II.3.5.6-fr}
Considérons maintenant deux préschémas $Y$, $Y'$, un morphisme
$g:Y'\to Y$, une $\mathcal{O}_Y$-Algèbre (resp.
$\mathcal{O}_{Y'}$-Algèbre) graduée quasi-cohérente $\mathcal{S}$ (resp.
$\mathcal{S}'$) et un $g$-morphisme d'Algèbres graduées
$u:\mathcal{S}\to\mathcal{S}'$, c'est-à-dire un
$\mathcal{O}_Y$-homomorphisme d'Algèbres graduées
$\mathcal{S}\to g_*(\mathcal{S}')$~; on sait d'ailleurs qu'il revient au même
de se donner un $\mathcal{O}_{Y'}$-homomorphisme d'Algèbres graduées
$u^\sharp:g^*(\mathcal{S})\to\mathcal{S}'$. On déduit alors canoniquement de
$u^\sharp$ un $Y'$-morphisme
$W=\operatorname{Proj}(u^\sharp):G(u^\sharp)\to
\operatorname{Proj}(g^*(\mathcal{S}))$ où $G(u^\sharp)$ est un ouvert de
$X'=\operatorname{Proj}(\mathcal{S}')$
(\hyperref[II.3.5.1-fr]{3.5.1}). D'autre part,
$X''=\operatorname{Proj}(g^*(\mathcal{S}))$ s'identifie

\oldpage[II]{64}

canoniquement à $X\times_Y Y'$, en posant
$X=\operatorname{Proj}(\mathcal{S})$
(\hyperref[II.3.5.3-fr]{3.5.3})~; composant avec
$\operatorname{Proj}(u^\sharp)$ la première projection
$p:X\times_Y Y'\to X$, on obtient donc un morphisme
$v:G(u^\sharp)\to X$, que nous désignerons par $\operatorname{Proj}(u)$, et
qui est tel que le diagramme
\[
  \xymatrix{
    G(u^\sharp)\ar[r]^{v}\ar[d]
    & X\ar[d]\\
    Y'\ar[r]_{g}
    & Y
  }
\]
soit commutatif.

En outre, pour tout $\mathcal{S}$-Module gradué quasi-cohérent $\mathcal{M}$,
on a un $v$-morphisme canonique
\[
\label{II.3.5.6.1-fr}
  \nu:\widetilde{\mathcal{M}}\longrightarrow
  (g^*(\mathcal{M})\otimes_{g^*(\mathcal{S})}\mathcal{S}')^{\sim}
  |G(u^\sharp).
\tag{3.5.6.1}
\]

En effet, $\nu^\sharp$ s'obtient en composant les homomorphismes
\[
  v^*(\widetilde{\mathcal{M}})
  =W^*(p^*(\widetilde{\mathcal{M}}))
  \longrightarrow W^*((g^*(\mathcal{M}))^{\sim})
  \longrightarrow
  (g^*(\mathcal{M})\otimes_{g^*(\mathcal{S})}\mathcal{S}')^{\sim}
  |G(u^\sharp)
\]
où la première flèche provient de l'isomorphisme
(\hyperref[II.3.5.3-fr]{3.5.3}) et la seconde est l'homomorphisme
(\hyperref[II.3.5.2-fr]{3.5.2, (ii)})~; lorsque $\mathcal{S}$ est engendrée
par $\mathcal{S}_1$, il résulte de (\hyperref[II.3.5.2-fr]{3.5.2}) que
$\nu^\sharp$ est un isomorphisme.

Comme cas particulier de (\hyperref[II.3.5.6.1-fr]{3.5.6.1}), on a, pour
tout $n\in\mathbf{Z}$, un $v$-morphisme canonique
\[
\label{II.3.5.6.2-fr}
  \nu:\mathcal{O}_X(n)\longrightarrow
  \mathcal{O}_{X'}(n)|G(u^\sharp).
\tag{3.5.6.2}
\]
\end{env}

\subsection{Sous-préschémas fermés d'un préschéma
$\operatorname{Proj}(\mathcal{S})$.}
\label{subsection:II.3.6-fr}

\begin{env}[3.6.1]
\label{II.3.6.1-fr}
Soient $Y$ un préschéma,
$\varphi:\mathcal{S}\to\mathcal{S}'$ un homomorphisme de
$\mathcal{O}_Y$-Algèbres graduées quasi-cohérentes, de degré $0$. On dit que
$\varphi$ est (\textbf{TN})-\emph{surjectif} (resp.
(\textbf{TN})-\emph{injectif}, (\textbf{TN})-\emph{bijectif}) s'il existe $n$
tel que, pour $k\geq n$, $\varphi_k:\mathcal{S}_k\to\mathcal{S}'_k$ soit
surjectif (resp. injectif, bijectif). Lorsqu'il en est ainsi on ramène l'étude
du morphisme
$\Phi:\operatorname{Proj}(\mathcal{S}')\to\operatorname{Proj}(\mathcal{S})$
correspondant au cas où $\varphi$ est surjectif (resp. injectif, bijectif).
Cela se démontre comme dans (\hyperref[II.2.9.1-fr]{2.9.1}) (qui en est le cas
particulier correspondant à $Y$ affine) en utilisant
(\hyperref[II.3.1.8-fr]{3.1.8}). Au lieu de dire que $\varphi$ est
(\textbf{TN})-bijectif, on dit aussi que c'est alors un
(\textbf{TN})-isomorphisme.
\end{env}

\begin{proposition}[3.6.2]
\label{II.3.6.2-fr}
Soient $Y$ un préschéma, $\mathcal{S}$ une $\mathcal{O}_Y$-Algèbre graduée
quasi-cohérente, et soit $X=\operatorname{Proj}(\mathcal{S})$.
\begin{enumerate}
  \item[\rm(i)] Si $\varphi:\mathcal{S}\to\mathcal{S}'$ est un
    homomorphisme (\textbf{TN})-surjectif de $\mathcal{O}_Y$-Algèbres
    graduées, le morphisme correspondant
    $\Phi=\operatorname{Proj}(\varphi)$
    (\hyperref[II.3.5.1-fr]{3.5.1}) est défini dans
    $\operatorname{Proj}(\mathcal{S}')$ tout entier et est une immersion
    fermée de $\operatorname{Proj}(\mathcal{S}')$ dans $X$. Si $\mathcal{J}$
    est le noyau de $\varphi$, le sous-préschéma fermé de $X$ associé à
    $\Phi$ est défini par le faisceau d'idéaux quasi-cohérent
    $\widetilde{\mathcal{J}}$ de $\mathcal{O}_X$.
  \item[\rm(ii)] Supposons en outre $\mathcal{S}_0=\mathcal{O}_Y$,
    $\mathcal{S}$ engendrée par $\mathcal{S}_1$ et $\mathcal{S}_1$ de type
    fini. Soit $X'$ un sous-préschéma fermé de
    $X=\operatorname{Proj}(\mathcal{S})$, défini par un faisceau
    quasi-cohérent $\mathcal{I}$ d'idéaux de $\mathcal{O}_X$. Soit
    $\mathcal{J}$

\oldpage[II]{65}

    le faisceau gradué quasi-cohérent d'idéaux de $\mathcal{S}$, image
    réciproque de $\Gamma_*(\mathcal{I})$ par l'homomorphisme canonique
    $\alpha:\mathcal{S}\to\Gamma_*(\mathcal{O}_X)$
    (\hyperref[II.3.3.2-fr]{3.3.2}), et posons
    $\mathcal{S}'=\mathcal{S}/\mathcal{J}$. Alors $X'$ est le sous-préschéma
    associé (\hyperref[I.4.2.1-fr]{I, 4.2.1}) à l'immersion fermée
    $\operatorname{Proj}(\mathcal{S}')\to X$ correspondant à l'homomorphisme
    canonique $\mathcal{S}\to\mathcal{S}'$ de $\mathcal{O}_Y$-Algèbres
    graduées.
\end{enumerate}
\end{proposition}

\begin{proof}
\begin{enumerate}
  \item[\rm(i)] On peut supposer $\varphi$ surjectif
    (\hyperref[II.3.6.1-fr]{3.6.1}). Alors, pour tout ouvert affine $U$ de
    $Y$, $\Gamma(U,\mathcal{S})\to\Gamma(U,\mathcal{S}')$ est surjectif
    (\hyperref[I.1.3.9-fr]{I, 1.3.9}), donc
    (\hyperref[II.3.5.1-fr]{3.5.1}) on a $G(\varphi)=X$. On est
    immédiatement ramené à démontrer la proposition lorsque $Y$ est affine,
    et alors elle découle de (\hyperref[II.2.9.2-fr]{2.9.2, (i)}).
  \item[\rm(ii)] On est ramené à prouver que l'homomorphisme
    $\widetilde{\mathcal{J}}\to\mathcal{O}_X$ déduit de l'injection canonique
    $\mathcal{J}\to\mathcal{S}$ est un isomorphisme de
    $\widetilde{\mathcal{J}}$ sur $\mathcal{I}$~; comme la question est locale
    sur $Y$, on peut supposer $Y$ affine d'anneau $A$, ce qui entraîne
    $\mathcal{S}=\widetilde{S}$, où $S$ est une $A$-algèbre graduée engendrée
    par $S_1$, $S_1$ étant de type fini sur $A$. Il suffit alors d'appliquer
    (\hyperref[II.2.9.2-fr]{2.9.2, (ii)}).
\end{enumerate}
\end{proof}

\begin{corollary}[3.6.3]
\label{II.3.6.3-fr}
Sous les conditions de (\hyperref[II.3.6.2-fr]{3.6.2, (i)}), et en supposant
de plus $\mathcal{S}$ engendrée par $\mathcal{S}_1$,
$\Phi^*(\mathcal{O}_X(n))$ s'identifie canoniquement à
$\mathcal{O}_{X'}(n)$ pour tout $n\in\mathbf{Z}$.
\end{corollary}

\begin{proof}
On a défini un tel isomorphisme canonique lorsque $Y$ est affine
(\hyperref[II.2.9.3-fr]{2.9.3})~; dans le cas général, il suffit de vérifier
que les isomorphismes ainsi définis pour chaque ouvert affine $U$ de $Y$ sont
compatibles avec le passage de $U$ à un ouvert affine $U'\subset U$, ce qui
est immédiat.
\end{proof}

\begin{corollary}[3.6.4]
\label{II.3.6.4-fr}
Soient $Y$ un préschéma, $\mathcal{S}$ une $\mathcal{O}_Y$-Algèbre graduée
quasi-cohérente engendrée par $\mathcal{S}_1$, $\mathcal{M}$ un
$\mathcal{O}_Y$-Module quasi-cohérent, $u$ un $\mathcal{O}_Y$-homomorphisme
surjectif $\mathcal{M}\to\mathcal{S}_1$,
$\overline{u}:\mathbf{S}_{\mathcal{O}_Y}(\mathcal{M})\to\mathcal{S}$
l'homomorphisme de $\mathcal{O}_Y$-Algèbres graduées prolongeant $u$
(\hyperref[II.1.7.4-fr]{1.7.4}). Le morphisme correspondant à
$\overline{u}$ est alors une immersion fermée de
$\operatorname{Proj}(\mathcal{S})$ dans
$\operatorname{Proj}(\mathbf{S}_{\mathcal{O}_Y}(\mathcal{M}))$.
\end{corollary}

\begin{proof}
En effet, $\overline{u}$ est surjectif par hypothèse et on applique
(\hyperref[II.3.6.1-fr]{3.6.1, (i)}).
\end{proof}

\subsection{Morphismes d'un préschéma dans un spectre homogène.}
\label{subsection:II.3.7-fr}

\begin{env}[3.7.1]
\label{II.3.7.1-fr}
Soient $q:X\to Y$ un morphisme de préschémas, $\mathcal{L}$ un
$\mathcal{O}_X$-Module inversible, $\mathcal{S}$ une $\mathcal{O}_Y$-Algèbre
graduée quasi-cohérente à degrés positifs~; $q^*(\mathcal{S})$ est alors une
$\mathcal{O}_X$-Algèbre graduée quasi-cohérente à degrés positifs.
Considérons la $\mathcal{O}_X$-Algèbre graduée quasi-cohérente
$\mathcal{S}'=\bigoplus_{n\geq0}\mathcal{L}^{\otimes n}$ et supposons donné
un $\mathcal{O}_X$-homomorphisme d'Algèbres graduées
\[
  \psi:q^*(\mathcal{S})\longrightarrow
  \mathcal{S}'=\bigoplus_{n\geq0}\mathcal{L}^{\otimes n}
\]
ce qui revient d'ailleurs à se donner un $q$-morphisme d'Algèbres graduées
\[
  \psi^\flat:\mathcal{S}\longrightarrow q_*(\mathcal{S}').
\]
On sait que $\operatorname{Proj}(\mathcal{S}')$ s'identifie canoniquement à
$X$ (\hyperref[II.3.1.7-fr]{3.1.7} et
\hyperref[II.3.1.8-fr]{3.1.8, (iii)})~; on déduit canoniquement de $\psi$ un
ensemble ouvert $G(\psi)$ de $X$ et un $Y$-morphisme
\[
\label{II.3.7.1.1-fr}
  r_{\mathcal{L},\psi}:G(\psi)\longrightarrow
  \operatorname{Proj}(\mathcal{S})=P
\tag{3.7.1.1}
\]

\oldpage[II]{66}

que nous appellerons le morphisme \emph{associé à $\mathcal{L}$ et à $\psi$}~;
rappelons (\hyperref[II.3.5.6-fr]{3.5.6}) que ce morphisme s'obtient en
composant avec la première projection
$\pi:\operatorname{Proj}(q^*(\mathcal{S}))=P\times_YX\to P$ le
$Y$-morphisme
\[
  \tau=\operatorname{Proj}(\psi):G(\psi)\longrightarrow
  \operatorname{Proj}(q^*(\mathcal{S})).
\]
\end{env}

\begin{env}[3.7.2]
\label{II.3.7.2-fr}
Explicitons $r=r_{\mathcal{L},\psi}$ lorsque $Y=\operatorname{Spec}(A)$ est
affine, et par suite $\mathcal{S}=\widetilde{S}$, où $S$ est une $A$-algèbre
graduée à degrés positifs. Supposons d'abord que $X=\operatorname{Spec}(B)$
soit aussi affine et que l'on ait $\mathcal{L}=\widetilde{L}$, où $L$ est un
$B$-module libre de rang $1$. On a alors
$q^*(\mathcal{S})=(S\otimes_AB)^{\sim}$
(\hyperref[I.1.6.5-fr]{I, 1.6.5})~; si $c$ est un générateur de $L$,
$\psi_n:q^*(\mathcal{S}_n)\to\mathcal{L}^{\otimes n}$ correspond à un
homomorphisme
$w_n:s\otimes b\mapsto bv_n(s)c^{\otimes n}$ de $S_n\otimes_AB$ dans
$L^{\otimes n}$, où $v_n:S_n\to B$ est un homomorphisme de $A$-modules, les
$v_n$ constituant un homomorphisme d'algèbres $S\to B$. Soit $f\in S_d$
($d>0$) et posons $g=v_d(f)$~; on a
$\pi^{-1}(D_+(f))=D_+(f\otimes1)$ en vertu de
(\hyperref[II.2.8.10-fr]{2.8.10}) et de l'identification de $D_+(f)$ à
$\operatorname{Spec}(S_{(f)})$ (\hyperref[II.2.3.6-fr]{2.3.6})~; d'autre
part, la formule (\hyperref[II.2.8.1.1-fr]{2.8.1.1}) montre (compte tenu de
l'identification canonique de $X$ et de $\operatorname{Proj}(\mathcal{S}')$)
que
\[
  \tau^{-1}(D_+(f\otimes1))=D(g)
\]
d'où
\[
\label{II.3.7.2.1-fr}
  r^{-1}(D_+(f))=D(g).
\tag{3.7.2.1}
\]

En outre, le morphisme $\tau=\operatorname{Proj}(\psi)$, restreint à $D(g)$,
correspond à l'homomorphisme qui applique
$(s\otimes1)/(f\otimes1)^n$ (pour $s\in S_{nd}$) sur
$v_{nd}(s)/g^n$ (\hyperref[II.2.8.1-fr]{2.8.1}), et la projection $\pi$,
restreinte à $D_+(f\otimes1)$, correspond à l'homomorphisme
$s/f^n\mapsto(s\otimes1)/(f\otimes1)^n$~; on en conclut que $r$, restreint à
$D(g)$, correspond à l'homomorphisme de $A$-algèbres
$\omega:S_{(f)}\to B_g$ tel que $\omega(s/f^n)=v_{nd}(s)/g^n$ pour
$s\in S_{nd}$ ($n>0$). Passant au cas où $X$ est quelconque ($Y$ étant
toujours affine), on obtient donc, compte tenu de
(\hyperref[II.2.8.1-fr]{2.8.1})~:
\end{env}

\begin{proposition}[3.7.3]
\label{II.3.7.3-fr}
Si $Y=\operatorname{Spec}(A)$ est affine et $\mathcal{S}=\widetilde{S}$, où
$S$ est une $A$-algèbre graduée, pour tout
$f\in S_d=\Gamma(Y,\mathcal{S}_d)$, on a
\[
\label{II.3.7.3.1-fr}
  r_{\mathcal{L},\psi}^{-1}(D_+(f))=X_{\psi^\flat(f)}
  \qquad\bigl(\text{où }\psi^\flat(f)\in
  \Gamma(X,\mathcal{L}^{\otimes d})\bigr)
\tag{3.7.3.1}
\]
et la restriction $X_{\psi^\flat(f)}\to D_+(f)=\operatorname{Spec}(S_{(f)})$
de $r_{\mathcal{L},\psi}$ correspond
(\hyperref[I.2.2.4-fr]{I, 2.2.4}) à l'homomorphisme d'algèbres
\[
\label{II.3.7.3.2-fr}
  \psi_{(f)}^\flat:S_{(f)}\longrightarrow
  \Gamma(X_{\psi^\flat(f)},\mathcal{O}_X)
\tag{3.7.3.2}
\]
tel que, pour $s\in S_{nd}=\Gamma(Y,\mathcal{S}_{nd})$,
\[
\label{II.3.7.3.3-fr}
  \psi_{(f)}^\flat(s/f^n)
  =(\psi^\flat(s)|X_{\psi^\flat(f)})
   (\psi^\flat(f)|X_{\psi^\flat(f)})^{-n}.
\tag{3.7.3.3}
\]
\end{proposition}

Nous dirons que $r_{\mathcal{L},\psi}$ est \emph{partout défini} si l'on a
$G(\psi)=X$. Pour qu'il en soit ainsi, il faut et il suffit évidemment que
$G(\psi)\cap q^{-1}(U)=q^{-1}(U)$ pour tout ouvert affine $U\subset Y$~;
autrement dit, la question est \emph{locale sur $Y$}. Si $Y$ est affine,
$G(\psi)$ est réunion des $r^{-1}(D_+(f))$ pour $f$ homogène dans $S_+$
(\hyperref[II.2.8.1-fr]{2.8.1})~; en vertu de
(\hyperref[II.3.7.3.1-fr]{3.7.3.1}), les $X_{\psi^\flat(f)}$ doivent donc
former un \emph{recouvrement} de $X$, autrement dit~:

\begin{corollary}[3.7.4]
\label{II.3.7.4-fr}
Sous les hypothèses de (\hyperref[II.3.7.3-fr]{3.7.3}), pour que
$r_{\mathcal{L},\psi}$ soit partout défini, il faut et il suffit que pour tout
$x\in X$, il existe un entier $n>0$ et une section $s$ de $\mathcal{S}_n$
au-dessus de $Y$ telle que si on pose
$t=\psi^\flat(s)\in\Gamma(X,\mathcal{L}^{\otimes n})$, on ait $t(x)\neq0$.
\end{corollary}

\oldpage[II]{67}

On notera que cette condition est toujours vérifiée si $\psi$ est
(\textbf{TN})-surjectif.

De même, la question de savoir si $r_{\mathcal{L},\psi}$ est \emph{dominant}
est locale sur $Y$, et l'on a~:

\begin{corollary}[3.7.5]
\label{II.3.7.5-fr}
Sous les hypothèses de (\hyperref[II.3.7.3-fr]{3.7.3}), pour que
$r_{\mathcal{L},\psi}$ soit dominant, il faut et il suffit que, pour tout
entier $n>0$, toute section $s\in S_n$ telle que
$\psi^\flat(s)\in\Gamma(X,\mathcal{L}^{\otimes n})$ soit localement
nilpotente, soit elle-même nilpotente.
\end{corollary}

\begin{proof}
On doit en effet exprimer que
$r_{\mathcal{L},\psi}^{-1}(D_+(s))$ n'est vide que si $D_+(s)$ est vide,
et le corollaire résulte donc de
(\hyperref[II.3.7.3.1-fr]{3.7.3.1}) et de
(\hyperref[II.2.3.7-fr]{2.3.7}).
\end{proof}

\begin{proposition}[3.7.6]
\label{II.3.7.6-fr}
Soient $q:X\to Y$ un morphisme, $\mathcal{L}$ un $\mathcal{O}_X$-Module
inversible, $\mathcal{S}$, $\mathcal{S}'$ deux $\mathcal{O}_Y$-Algèbres
graduées quasi-cohérentes,
$u:\mathcal{S}'\to\mathcal{S}$ un homomorphisme d'Algèbres graduées,
$\psi:q^*(\mathcal{S})\to\bigoplus_{n\geq0}\mathcal{L}^{\otimes n}$ un
homomorphisme d'Algèbres graduées,
$\psi'=\psi\circ q^*(u)$ l'homomorphisme composé. Si
$r_{\mathcal{L},\psi'}$ est partout défini, il en est de même de
$r_{\mathcal{L},\psi}$~; si $u$ est (\textbf{TN})-surjectif et si
$r_{\mathcal{L},\psi'}$ est dominant, il en est de même de
$r_{\mathcal{L},\psi}$~; inversement, si $u$ est (\textbf{TN})-injectif et si
$r_{\mathcal{L},\psi}$ est dominant, $r_{\mathcal{L},\psi'}$ est dominant.
\end{proposition}

\begin{proof}
On a en effet $G(\psi')\subset G(\psi)$
(\hyperref[II.2.8.4-fr]{2.8.4}), d'où la première assertion~; si $u$ est
(\textbf{TN})-surjectif,
$\operatorname{Proj}(u):\operatorname{Proj}(\mathcal{S})\to
\operatorname{Proj}(\mathcal{S}')$ est partout défini et est une immersion
fermée~; comme $r_{\mathcal{L},\psi'}$ est composé de
$\operatorname{Proj}(u)$ et de la restriction de $r_{\mathcal{L},\psi}$ à
$G(\psi')$, on en conclut que si $r_{\mathcal{L},\psi'}$ est dominant, il en
est de même de $r_{\mathcal{L},\psi}$. Enfin, si $u$ est
(\textbf{TN})-injectif, on sait que $\operatorname{Proj}(u)$ est un morphisme
dominant de $G(u)$ dans $\operatorname{Proj}(\mathcal{S}')$
(\hyperref[II.2.8.3-fr]{2.8.3})~; comme $G(\psi')$ est l'image réciproque de
$G(u)$ par $r_{\mathcal{L},\psi}$, on voit que si $r_{\mathcal{L},\psi}$ est
dominant, il en est de même de $r_{\mathcal{L},\psi'}$.
\end{proof}

\begin{proposition}[3.7.7]
\label{II.3.7.7-fr}
Soient $Y$ un préschéma quasi-compact, $q:X\to Y$ un morphisme quasi-compact,
$\mathcal{L}$ un $\mathcal{O}_X$-Module inversible, $\mathcal{S}$ une
$\mathcal{O}_Y$-Algèbre graduée quasi-cohérente, limite inductive filtrante
d'un système inductif $(\mathcal{S}^\lambda)$ de $\mathcal{O}_Y$-Algèbres
quasi-cohérentes. Soient
$\varphi_\lambda:\mathcal{S}^\lambda\to\mathcal{S}$ l'homomorphisme
canonique,
$\psi:q^*(\mathcal{S})\to\bigoplus_{n\geq0}\mathcal{L}^{\otimes n}$ un
homomorphisme d'Algèbres graduées, et posons
$\psi_\lambda=\psi\circ q^*(\varphi_\lambda)$. Pour que
$r_{\mathcal{L},\psi}$ soit partout défini, il faut et il suffit qu'il existe
un $\lambda$ tel que $r_{\mathcal{L},\psi_\lambda}$ soit partout défini~;
$r_{\mathcal{L},\psi_\mu}$ est alors partout défini pour $\mu\geq\lambda$.
\end{proposition}

\begin{proof}
La condition est suffisante en vertu de
(\hyperref[II.3.7.6-fr]{3.7.6}). Inversement, supposons
$r_{\mathcal{L},\psi}$ partout défini~; on peut se ramener au cas où $Y$ est
affine, car si pour tout ouvert affine $U\subset Y$, il existe $\lambda(U)$
tel que la restriction de $r_{\mathcal{L},\psi_{\lambda(U)}}$ à $q^{-1}(U)$
soit partout définie, il suffira ($Y$ étant quasi-compact) de recouvrir $Y$
par un nombre fini d'ouverts affines $U_i$, et de prendre
$\lambda\geq\lambda(U_i)$ pour tous les indices $i$, en vertu de
(\hyperref[II.3.7.6-fr]{3.7.6}). Si $Y$ est affine, l'hypothèse entraîne que
pour tout $x\in X$, il y a une section $s^{(x)}$ d'un $S_n$ telle que, si on
pose $t^{(x)}=\psi^\flat(s^{(x)})$, on ait $t^{(x)}(x)\neq0$
($t^{(x)}$ étant considérée comme section de $\mathcal{L}^{\otimes n}$
au-dessus de $X$), ce qui entraîne $t^{(x)}(z)\neq0$ pour tout $z$ dans un
voisinage $V(x)$ de $x$. Couvrons $X$ par un nombre fini de $V(x_i)$ et
soient $s^{(i)}$ les sections correspondantes de $S$~; il existe alors un
$\lambda$ tel que les $s^{(i)}$ soient toutes de la forme
$\varphi_\lambda(s_\lambda^{(i)})$, avec $s_\lambda^{(i)}\in S^\lambda$ pour
tous les $i$~; il résulte donc de (\hyperref[II.3.7.4-fr]{3.7.4}) que
$r_{\mathcal{L},\psi_\lambda}$ est partout défini. La dernière assertion est
conséquence triviale de (\hyperref[II.3.7.6-fr]{3.7.6}).
\end{proof}

\begin{corollary}[3.7.8]
\label{II.3.7.8-fr}
Sous les hypothèses de (\hyperref[II.3.7.7-fr]{3.7.7}), si les
$r_{\mathcal{L},\psi_\lambda}$ sont dominants, il en est de même de
$r_{\mathcal{L},\psi}$~; la réciproque est vraie lorsque les
$\varphi_\lambda$ sont injectifs.
\end{corollary}

\oldpage[II]{68}

\begin{proof}
La seconde assertion est un cas particulier de
(\hyperref[II.3.7.6-fr]{3.7.6})~; d'autre part, pour montrer que
$r_{\mathcal{L},\psi}$ est dominant, on peut se borner au cas où $Y$ est
affine~; si $s\in S$ est telle que $\psi^\flat(s)$ soit localement
nilpotente, comme on peut écrire $s=\varphi_\lambda(s_\lambda)$ pour un
$\lambda$ au moins, on conclut de l'hypothèse et de
(\hyperref[II.3.7.5-fr]{3.7.5}) que $s_\lambda$ est nilpotente, donc il en est
de même de $s$, et le critère de (\hyperref[II.3.7.5-fr]{3.7.5}) s'applique.
\end{proof}

\begin{namedtheorem}[3.7.9]{Remarques}
\label{II.3.7.9-fr}
\begin{enumerate}
  \item[(i)] Avec les notations de (\hyperref[II.3.7.1-fr]{3.7.1}) et en
    tenant compte de (\hyperref[II.3.2.10-fr]{3.2.10}), on a, pour tout
    $n\in\mathbf{Z}$, un homomorphisme canonique
    \[
    \label{II.3.7.9.1-fr}
      \theta:r_{\mathcal{L},\psi}^*(\mathcal{O}_P(n))\longrightarrow
      \mathcal{L}^{\otimes n}
    \tag{3.7.9.1}
    \]
    défini de façon générale dans
    (\hyperref[II.3.5.6.2-fr]{3.5.6.2}). On voit aussitôt que sous les
    conditions de (\hyperref[II.3.7.3-fr]{3.7.3}), la restriction de cet
    homomorphisme à $X_{\psi^\flat(f)}$ s'explicite comme faisant
    correspondre à $s/f^k$ ($s\in S_{n+kd}$) l'élément
    $(\psi^\flat(s)|X_{\psi^\flat(f)})
    (\psi^\flat(f)|X_{\psi^\flat(f)})^{-k}$.
  \item[(ii)] Soit $\mathcal{F}$ un $\mathcal{O}_X$-Module quasi-cohérent, et
    supposons que $q$ soit quasi-compact et séparé, de sorte que pour tout
    $n\geq0$, $q_*(\mathcal{F}\otimes\mathcal{L}^{\otimes n})$ soit un
    $\mathcal{O}_Y$-Module quasi-cohérent
    (\hyperref[I.9.2.2-fr]{I, 9.2.2}). Soit
    $\mathcal{M}'=\bigoplus_{n\geq0}\mathcal{F}\otimes
    \mathcal{L}^{\otimes n}$, qui est un $\mathcal{S}'$-Module gradué
    quasi-cohérent, et considérons son image
    $\mathcal{M}=q_*(\mathcal{M}')=
    \bigoplus_{n\geq0}q_*(\mathcal{F}\otimes\mathcal{L}^{\otimes n})$
    (qui est un $\mathcal{S}$-Module quasi-cohérent, au moyen de
    l'homomorphisme $\psi^\flat$). Nous allons voir qu'il y a un
    homomorphisme canonique de $\mathcal{O}_X$-Modules
    \[
    \label{II.3.7.9.2-fr}
      \xi:r_{\mathcal{L},\psi}^*(\widetilde{\mathcal{M}})
      \longrightarrow\mathcal{F}|G(\psi).
    \tag{3.7.9.2}
    \]

    En effet, on a défini (\hyperref[II.3.5.6.1-fr]{3.5.6.1}) un
    homomorphisme canonique
    \[
    \label{II.3.7.9.3-fr}
      r_{\mathcal{L},\psi}^*(\widetilde{\mathcal{M}})
      \longrightarrow
      (q^*(\mathcal{M})\otimes_{q^*(\mathcal{S})}\mathcal{S}')^{\sim}
      |G(\psi),
    \tag{3.7.9.3}
    \]
    où le second membre est considéré comme un faisceau quasi-cohérent sur
    $\operatorname{Proj}(\mathcal{S}')$. D'autre part, on a un
    homomorphisme canonique
    \[
    \label{II.3.7.9.4-fr}
      q^*(q_*(\mathcal{M}'))\otimes_{q^*(\mathcal{S})}\mathcal{S}'
      \longrightarrow\mathcal{M}'
    \tag{3.7.9.4}
    \]
    pour tout $\mathcal{S}'$-Module gradué quasi-cohérent $\mathcal{M}'$~:
    pour tout ouvert $U$ de $X$, toute section $t'$ de
    $q^*(q_*(\mathcal{M}'_h))$ au-dessus de $U$ et toute section $b'$ de
    $\mathcal{S}'_k$ au-dessus de $U$, on fait correspondre à $t'\otimes b'$ la
    section $b'\sigma(t')$ de $\mathcal{M}'_{h+k}$, où $\sigma(t')$ est la
    section de $\mathcal{M}'_h$ au-dessus de $U$ qui correspond canoniquement
    (\hyperref[0.4.4.3-fr]{0, 4.4.3}) à $t'$. On en conclut un homomorphisme
    canonique
    \[
    \label{II.3.7.9.5-fr}
      (q^*(q_*(\mathcal{M}'))\otimes_{q^*(\mathcal{S})}\mathcal{S}')^{\sim}
      |G(\psi)\longrightarrow\widetilde{\mathcal{M}'}|G(\psi)
    \tag{3.7.9.5}
    \]
    et comme finalement $\widetilde{\mathcal{M}'}$ s'identifie
    canoniquement à $\mathcal{F}$
    (\hyperref[II.3.2.9-fr]{3.2.9, (i)}), on obtient bien l'homomorphisme
    canonique cherché.

    Sous les conditions de (\hyperref[II.3.7.3-fr]{3.7.3}), la restriction
    de cet homomorphisme à $X_{\psi^\flat(f)}$ s'explicite de la façon
    suivante~: étant donnée une section $t_{nd}$ de
    $\mathcal{F}\otimes\mathcal{L}^{\otimes nd}$ au-dessus de $X$, si
    $t'_{nd}$ est la section $t_{nd}$ considérée comme section de
    $q_*(\mathcal{F}\otimes\mathcal{L}^{\otimes nd})$ au-dessus de $Y$, à
    l'élément $t'_{nd}/f^n$ il correspond la section
    $(t_{nd}|X_{\psi^\flat(f)})
    (\psi^\flat(f)|X_{\psi^\flat(f)})^{-n}$ de $\mathcal{F}$ au-dessus de
    $X_{\psi^\flat(f)}$.
\end{enumerate}
\end{namedtheorem}

\oldpage[II]{69}

\subsection{Critères d'immersion dans un spectre homogène.}
\label{subsection:II.3.8-fr}

\begin{env}[3.8.1]
\label{II.3.8.1-fr}
Avec les notations de (\hyperref[II.3.7.1-fr]{3.7.1}), la question de savoir
si $r_{\mathcal{L},\psi}$ est une immersion (resp. une immersion ouverte, une
immersion fermée) est évidemment \emph{locale sur $Y$}.
\end{env}

\begin{proposition}[3.8.2]
\label{II.3.8.2-fr}
Sous les hypothèses de (\hyperref[II.3.7.3-fr]{3.7.3}), pour que
$r_{\mathcal{L},\psi}$ soit partout définie et soit une immersion, il faut et
il suffit qu'il existe une famille de sections $s_\alpha\in S_{n_\alpha}$
($n_\alpha>0$) telle que, si on pose
$f_\alpha=\psi^\flat(s_\alpha)$, les conditions suivantes soient vérifiées~:
\begin{enumerate}
  \item[(i)] Les $X_{f_\alpha}$ forment un recouvrement de $X$.
  \item[(ii)] Les $X_{f_\alpha}$ sont des ouverts affines.
  \item[(iii)] Pour tout $\alpha$ et tout
    $t\in\Gamma(X_{f_\alpha},\mathcal{O}_X)$, il existe un entier $m>0$ et
    un $s\in S_{mn_\alpha}$ tels que
    $t=(\psi^\flat(s)|X_{f_\alpha})(f_\alpha|X_{f_\alpha})^{-m}$.
\end{enumerate}

Pour que $r_{\mathcal{L},\psi}$ soit partout définie et soit une immersion
ouverte, il faut et il suffit qu'il existe une famille $(s_\alpha)$ vérifiant
les conditions (i), (ii), (iii) et~:
\begin{enumerate}
  \item[(iv)] Pour tout $n>0$ et tout $s\in S_{nn_\alpha}$ tel que
    $\psi^\flat(s)|X_{f_\alpha}=0$, il existe un entier $k>0$ tel que
    $s_\alpha^ks=0$.
\end{enumerate}

Pour que $r_{\mathcal{L},\psi}$ soit partout définie et soit une immersion
fermée, il faut et il suffit qu'il existe une famille $(s_\alpha)$ vérifiant
(i), (ii), (iii) et~:
\begin{enumerate}
  \item[(v)] Les $D_+(s_\alpha)$ forment un recouvrement de
    $P=\operatorname{Proj}(S)$.
\end{enumerate}
\end{proposition}

\begin{proof}
En effet, pour que $r$ soit une immersion (resp. une immersion fermée), il
faut et il suffit qu'il existe un recouvrement de $r(G(\psi))$ (resp. de
$P$) par des $D_+(s_\alpha)$ tel que si on pose
$V_\alpha=r^{-1}(D_+(s_\alpha))$, la restriction de $r$ à $U_\alpha$ soit
une immersion fermée de $V_\alpha$ dans $D_+(s_\alpha)$
(\hyperref[I.4.2.4-fr]{I, 4.2.4}). La condition (i) exprime à la fois que
$r$ est partout définie et que les $D_+(s_\alpha)$ recouvrent $r(X)$,
d'après (\hyperref[II.3.7.3.1-fr]{3.7.3.1})~; comme $D_+(s_\alpha)$ est
affine, les conditions (ii) et (iii) expriment que la restriction de $r$ à
$X_{f_\alpha}$ est une immersion fermée dans $D_+(s_\alpha)$, d'après
(\hyperref[I.4.2.3-fr]{I, 4.2.3})~; enfin, comme (iii) et (iv) expriment que
$\psi_{(s_\alpha)}^\flat$ est bijective (notation de
(\hyperref[II.3.7.3.2-fr]{3.7.3.2})), (ii), (iii) et (iv) expriment que la
restriction de $r$ à $X_{f_\alpha}$ est un isomorphisme sur
$D_+(s_\alpha)$ pour tout $\alpha$, donc (i), (ii), (iii) et (iv) expriment
que $r$ est une immersion ouverte.
\end{proof}

\begin{corollary}[3.8.3]
\label{II.3.8.3-fr}
Sous les hypothèses de (\hyperref[II.3.7.6-fr]{3.7.6}), si
$r_{\mathcal{L},\psi'}$ est partout définie et est une immersion, il en est
de même de $r_{\mathcal{L},\psi}$. Si on suppose en outre que $u$ est
(\textbf{TN})-surjectif et si $r_{\mathcal{L},\psi'}$ est une immersion
ouverte (resp. fermée), il en est de même de $r_{\mathcal{L},\psi}$.
\end{corollary}

\begin{proof}
En effet, d'après (\hyperref[II.3.8.2-fr]{3.8.2}), il y a une famille de
$s'_\alpha\in S'_{n_\alpha}$ telle que si on pose
$f_\alpha=\psi'^\flat(s'_\alpha)$, les conditions (i), (ii), (iii) soient
vérifiées. Or, si on pose $s_\alpha=u(s'_\alpha)$, on a aussi
$f_\alpha=\psi^\flat(s_\alpha)$, et si on a
$t=(\psi'^\flat(s')|X_{f_\alpha})(f_\alpha|X_{f_\alpha})^{-m}$, on a aussi
$t=(\psi^\flat(s)|X_{f_\alpha})(f_\alpha|X_{f_\alpha})^{-m}$ en posant
$s=u(s')$, d'où la première assertion. La seconde résulte aussitôt de ce que
$\operatorname{Proj}(u)$ est une immersion fermée.
\end{proof}

\begin{proposition}[3.8.4]
\label{II.3.8.4-fr}
Supposons vérifiées les hypothèses de
(\hyperref[II.3.7.7-fr]{3.7.7}) et en outre que $q:X\to Y$ soit un morphisme
de type fini. Alors, pour que $r_{\mathcal{L},\psi}$ soit partout définie et
soit une immersion, il faut et il suffit qu'il existe $\lambda$ tel que
$r_{\mathcal{L},\psi_\lambda}$ soit partout définie et soit une immersion et
dans ce cas $r_{\mathcal{L},\psi_\mu}$ est partout définie et est une
immersion pour $\mu\geq\lambda$.
\end{proposition}

\oldpage[II]{70}

\begin{proof}
Compte tenu de (\hyperref[II.3.8.3-fr]{3.8.3}), il suffit de démontrer que si
$r_{\mathcal{L},\psi}$ est partout définie et est une immersion, il en est de
même de $r_{\mathcal{L},\psi_\lambda}$ pour un $\lambda$ au moins. Par le
même raisonnement que dans (\hyperref[II.3.7.7-fr]{3.7.7}), utilisant la
quasi-compacité de $Y$, on se ramène d'abord au cas où $Y$ est affine. Comme
$X$ est quasi-compact, (\hyperref[II.3.8.2-fr]{3.8.2}) montre l'existence
d'une famille \emph{finie} $(s_i)$ d'éléments de $S$ ($s_i\in S_{n_i}$)
ayant les propriétés (i), (ii) et (iii). Le morphisme $X_{f_i}\to Y$ (où
$f_i=\psi^\flat(s_i)$) est de type fini~: en effet, c'est un morphisme de
schémas affines, donc il est quasi-compact
(\hyperref[I.6.6.1-fr]{I, 6.6.1}), et localement de type fini puisque $q$ est
de type fini (\hyperref[I.6.3.2-fr]{I, 6.3.2}), et la conclusion résulte de
(\hyperref[I.6.6.3-fr]{I, 6.6.3}). L'anneau $B_i$ de $X_{f_i}$ est par suite
une $A$-algèbre de type fini (\hyperref[I.6.3.3-fr]{I, 6.3.3})~; soit
$(t_{ij})$ une famille de générateurs de cette algèbre. Il y a par hypothèse
des éléments $s'_{ij}\in S_{m_{ij}n_i}$ tels que
\[
  t_{ij}=(\psi^\flat(s'_{ij})|X_{f_i})
  (\psi^\flat(s_i)|X_{f_i})^{-m_{ij}}.
\]
Il y a par hypothèse un $\lambda$ et des éléments
$s_{i\lambda}\in S_{n_i}^\lambda$,
$s'_{ij\lambda}\in S_{m_{ij}n_i}^\lambda$ dont les images par
$\varphi_\lambda$ sont respectivement $s_i$ et $s'_{ij}$~; il est clair que
$r_{\mathcal{L},\psi_\lambda}$ vérifie les conditions (i), (ii) et (iii) de
(\hyperref[II.3.8.2-fr]{3.8.2}).
\end{proof}

\begin{proposition}[3.8.5]
\label{II.3.8.5-fr}
Soient $Y$ un schéma quasi-compacte, ou un préschéma dont l'espace sous-jacent
est noethérien, $q:X\to Y$ un morphisme de type fini, $\mathcal{L}$ un
$\mathcal{O}_X$-Module inversible, $\mathcal{S}$ une $\mathcal{O}_Y$-Algèbre
graduée quasi-cohérente,
$\psi:q^*(\mathcal{S})\to\bigoplus_{n\geq0}\mathcal{L}^{\otimes n}$ un
homomorphisme d'Algèbres graduées. Pour que $r_{\mathcal{L},\psi}$ soit
partout défini et soit une immersion, il faut et il suffit qu'il existe un
entier $n>0$ et un sous-$\mathcal{O}_Y$-Module quasi-cohérent de type fini
$\mathcal{E}$ de $\mathcal{S}_n$, tels que~:
\begin{enumerate}
  \item[a)] l'homomorphisme
    $\psi_n\circ q^*(j_n):q^*(\mathcal{E})\to\mathcal{L}^{\otimes n}$
    (où $j_n:\mathcal{E}\to\mathcal{S}$ est l'injection canonique) soit
    surjectif~;
  \item[b)] si on désigne par $\mathcal{S}'$ la sous-$\mathcal{O}_Y$-Algèbre
    (graduée) de $\mathcal{S}$ engendrée par $\mathcal{E}$, par $\psi'$
    l'homomorphisme $\psi\circ q^*(j')$ ($j'$ injection de $\mathcal{S}'$
    dans $\mathcal{S}$), $r_{\mathcal{L},\psi'}$ soit partout défini et soit
    une immersion.
\end{enumerate}

Lorsqu'il en est ainsi, tout sous-$\mathcal{O}_Y$-Module quasi-cohérent
$\mathcal{E}'$ de $\mathcal{S}_n$, contenant $\mathcal{E}$, possède la même
propriété, et il en est de même de l'image de $\bigotimes^k\mathcal{E}$ dans
$\mathcal{S}_{kn}$ pour tout $k>0$.
\end{proposition}

\begin{proof}
La suffisance de la condition et les deux dernières assertions sont des cas
particuliers de (\hyperref[II.3.8.3-fr]{3.8.3}), compte tenu de l'isomorphie
canonique entre $\operatorname{Proj}(\mathcal{S})$ et
$\operatorname{Proj}(\mathcal{S}^{(k)})$
(\hyperref[II.3.1.8-fr]{3.1.8}).

Soit $(U_i)$ un recouvrement ouvert fini de $Y$ par des ouverts affines et
posons $A_i=A(U_i)$. Comme $q^{-1}(U_i)$ est compact, l'hypothèse que
$r_{\mathcal{L},\psi}$ est une immersion partout définie et
(\hyperref[II.3.8.2-fr]{3.8.2}) entraînent l'existence d'une famille finie
$(s_{ij})$ d'éléments de
$S^{(i)}=\Gamma(U_i,\mathcal{S})$ ($s_{ij}\in S_{n_{ij}}^{(i)}$) ayant les
propriétés (i), (ii), (iii). On voit comme dans la démonstration de
(\hyperref[II.3.8.4-fr]{3.8.4}) que le morphisme $X_{f_{ij}}\to U_i$ (où
$f_{ij}=\psi^\flat(s_{ij})$) est de type fini, et l'anneau $B_{ij}$ de
$X_{f_{ij}}$ est par suite une $A_i$-algèbre de type fini, ayant un système de
générateurs de la forme
$(\psi^\flat(t_{ijk})|X_{f_{ij}})(f_{ij}|X_{f_{ij}})^{-m_{ijk}}$, où
$t_{ijk}\in S_{m_{ijk}n_{ij}}^{(i)}$. Soit $n$ un multiple commun des
$m_{ijk}n_{ij}$~; remplaçant (pour chaque $(i,j,k)$) $s_{ij}$ par une
puissance $s_{ij}^\rho$ telle que $\rho m_{ijk}n_{ij}=n$, et multipliant
$t_{ijk}$ par $s_{ij}^{\rho-m_{ijk}}$, on peut supposer que pour chaque $i$
tous les $s_{ij}$ et $t_{ijk}$ appartiennent à $S_n^{(i)}$ et que
$m_{ijk}=1$. Soit $E_i$ le sous-$A_i$-Module de $S_n^{(i)}$ engendré par ces
éléments (pour $i$ fixé). Il existe un sous-$\mathcal{O}_Y$-Module cohérent de
type

\oldpage[II]{71}

fini $\mathcal{E}_i$ de $\mathcal{S}_n$ tel que
$\mathcal{E}_i|U_i=(E_i)^{\sim}$
(\hyperref[I.9.4.7-fr]{I, 9.4.7}). Il est clair que le
sous-$\mathcal{O}_Y$-Module $\mathcal{E}$ de $\mathcal{S}_n$, somme des
$\mathcal{E}_i$, répond à la question (chaque section $f_{ij}$ étant telle
que pour tout $x\in X_{f_{ij}}$, il y a un voisinage affine
$V\subset X_{f_{ij}}$ de $x$ tel que $f|V$ soit une base de
$\Gamma(V,\mathcal{L}^{\otimes n})$).
\end{proof}

\section{Fibres projectifs. Faisceaux amples}
\label{section:II.4-fr}

\subsection{Définition des fibres projectifs.}
\label{subsection:II.4.1-fr}

\begin{definition}[4.1.1]
\label{II.4.1.1-fr}
Soient $Y$ un préschéma, $\mathcal{E}$ un $\mathcal{O}_Y$-Module
quasi-cohérent, $\mathbf{S}_{\mathcal{O}_Y}(\mathcal{E})$ la
$\mathcal{O}_Y$-Algèbre symétrique de $\mathcal{E}$
(\hyperref[II.1.7.4-fr]{1.7.4}) qui est quasi-cohérente
(\hyperref[II.1.7.7-fr]{1.7.7}). On appelle \emph{fibré projectif sur $Y$
défini par $\mathcal{E}$} et on note $\mathbf{P}(\mathcal{E})$ le $Y$-schéma
$P=\operatorname{Proj}(\mathbf{S}_{\mathcal{O}_Y}(\mathcal{E}))$. Le
$\mathcal{O}_P$-Module $\mathcal{O}_P(1)$ est appelé le \emph{faisceau
fondamental sur $P$}.
\end{definition}

Lorsque $Y$ est affine d'anneau $A$, on a $\mathcal{E}=\widetilde{E}$, où
$E$ est un $A$-module, et on écrit alors $\mathbf{P}(E)$ au lieu de
$\mathbf{P}(\widetilde{E})$.

Lorsque l'on prend $\mathcal{E}=\mathcal{O}_Y^n$, on écrit
$\mathbf{P}_Y^{n-1}$ au lieu de $\mathbf{P}(\mathcal{E})$~; si de plus $Y$
est affine d'anneau $A$, on écrit aussi $\mathbf{P}_A^{n-1}$ au lieu de
$\mathbf{P}_Y^{n-1}$. Comme
$\mathbf{S}_{\mathcal{O}_Y}(\mathcal{O}_Y)$ s'identifie canoniquement à
$\mathcal{O}_Y[T]$ (\hyperref[II.1.7.4-fr]{1.7.4}),
$\mathbf{P}_Y^0$ s'identifie canoniquement à $Y$
(\hyperref[II.3.1.7-fr]{3.1.7})~; l'exemple
(\hyperref[II.2.4.3-fr]{2.4.3}) n'est autre que $\mathbf{P}_K^1$.

\begin{env}[4.1.2]
\label{II.4.1.2-fr}
Soient $\mathcal{E}$, $\mathcal{F}$ deux $\mathcal{O}_Y$-Modules
quasi-cohérents~; $u:\mathcal{E}\to\mathcal{F}$ un
$\mathcal{O}_Y$-homomorphisme~; il lui correspond canoniquement un
homomorphisme
$\mathbf{S}(u):\mathbf{S}_{\mathcal{O}_Y}(\mathcal{E})\to
\mathbf{S}_{\mathcal{O}_Y}(\mathcal{F})$ de $\mathcal{O}_Y$-Algèbres
graduées (\hyperref[II.1.7.4-fr]{1.7.4}). Si $u$ est \emph{surjectif}, il en
est de même de $\mathbf{S}(u)$, et par suite
(\hyperref[II.3.6.2-fr]{3.6.2, (i)}),
$\operatorname{Proj}(\mathbf{S}(u))$ est une immersion fermée
$\mathbf{P}(\mathcal{F})\to\mathbf{P}(\mathcal{E})$, que nous noterons
$\mathbf{P}(u)$. On peut donc dire que $\mathbf{P}(\mathcal{E})$ est un
\emph{foncteur contravariant en $\mathcal{E}$}, à condition de ne prendre
pour morphismes des $\mathcal{O}_Y$-Modules quasi-cohérents que les
homomorphismes surjectifs.

En supposant toujours $u$ surjectif et posant
$P=\mathbf{P}(\mathcal{E})$, $Q=\mathbf{P}(\mathcal{F})$ et
$j=\mathbf{P}(u)$, on a, à un isomorphisme près,
\[
\label{II.4.1.2.1-fr}
  j^*(\mathcal{O}_P(n))=\mathcal{O}_Q(n)
  \qquad\text{pour tout }n\in\mathbf{Z}
\tag{4.1.2.1}
\]
comme il résulte de (\hyperref[II.3.6.3-fr]{3.6.3}).
\end{env}

\begin{env}[4.1.3]
\label{II.4.1.3-fr}
Soit maintenant $\psi:Y'\to Y$ un morphisme, et soit
$\mathcal{E}'=\psi^*(\mathcal{E})$~; on a alors
$\mathbf{S}_{\mathcal{O}_{Y'}}(\mathcal{E}')=
\psi^*(\mathbf{S}_{\mathcal{O}_Y}(\mathcal{E}))$
(\hyperref[II.1.7.5-fr]{1.7.5})~; par suite
(\hyperref[II.3.5.3-fr]{3.5.3}), on a
\[
\label{II.4.1.3.1-fr}
  \mathbf{P}(\psi^*(\mathcal{E}))=
  \mathbf{P}(\mathcal{E})\times_YY'
\tag{4.1.3.1}
\]
à un isomorphisme canonique près~; en outre, on a évidemment
\[
  \psi^*((\mathbf{S}_{\mathcal{O}_Y}(\mathcal{E}))(n))
  =(\mathbf{S}_{\mathcal{O}_{Y'}}(\mathcal{E}'))(n)
\]
pour tout $n\in\mathbf{Z}$, donc, en posant $P=\mathbf{P}(\mathcal{E})$,
$P'=\mathbf{P}(\mathcal{E}')$, on a
(\hyperref[II.3.5.4-fr]{3.5.4}), à un isomorphisme près,
\[
\label{II.4.1.3.2-fr}
  \mathcal{O}_{P'}(n)=\mathcal{O}_P(n)\otimes_Y\mathcal{O}_{Y'}
  \qquad\text{pour tout }n\in\mathbf{Z}.
\tag{4.1.3.2}
\]
\end{env}

\oldpage[II]{72}

\begin{proposition}[4.1.4]
\label{II.4.1.4-fr}
Soit $\mathcal{L}$ un $\mathcal{O}_Y$-Module inversible. Pour tout
$\mathcal{O}_Y$-Module quasi-cohérent $\mathcal{E}$, il existe un
$Y$-isomorphisme canonique
$i_{\mathcal{L}}:\mathbf{P}(\mathcal{E})\xrightarrow{\sim}
\mathbf{P}(\mathcal{E}\otimes\mathcal{L})$~; en outre, si l'on pose
$P=\mathbf{P}(\mathcal{E})$, $Q=\mathbf{P}(\mathcal{E}\otimes\mathcal{L})$,
$i_{\mathcal{L}}^*(\mathcal{O}_Q(n))$ est canoniquement isomorphe à
$\mathcal{O}_P(n)\otimes_Y\mathcal{L}^{\otimes n}$ pour tout
$n\in\mathbf{Z}$.
\end{proposition}

\begin{proof}
Remarquons d'abord que si $A$ est un anneau, $E$ un $A$-module, $L$ un
$A$-module \emph{monogène libre}, on définit canoniquement un homomorphisme
de $A$-modules
\[
  \mathbf{S}_n(E\otimes L)\longrightarrow
  \mathbf{S}_n(E)\otimes L^{\otimes n}
\]
en faisant correspondre à
$(x_1\otimes y_1)\cdots(x_n\otimes y_n)$ l'élément
\[
  (x_1x_2\cdots x_n)\otimes(y_1\otimes y_2\otimes\cdots\otimes y_n)
  \qquad(x_i\in E,\ y_i\in L\text{ pour }1\leq i\leq n)~;
\]
on vérifie immédiatement (en se ramenant au cas où $L=A$) que cet
homomorphisme est en fait un isomorphisme. On en conclut un isomorphisme
canonique de $A$-algèbres graduées
$\mathbf{S}_A(E\otimes L)\xrightarrow{\sim}
\bigoplus_{n\geq0}\mathbf{S}_n(E)\otimes L^{\otimes n}$. En revenant aux
conditions de (4.1.4), les remarques précédentes permettent de définir un
isomorphisme canonique de $\mathcal{O}_Y$-Algèbres graduées
\[
\label{II.4.1.4.1-fr}
  \mathbf{S}_{\mathcal{O}_Y}(\mathcal{E}\otimes_{\mathcal{O}_Y}\mathcal{L})
  \xrightarrow{\sim}
  \bigoplus_{n\geq0}\mathbf{S}_n(\mathcal{E})
  \otimes_{\mathcal{O}_Y}\mathcal{L}^{\otimes n}
\tag{4.1.4.1}
\]
en définissant cet isomorphisme comme un isomorphisme de préfaisceaux et
tenant compte de (\hyperref[II.1.7.4-fr]{1.7.4}),
(\hyperref[I.1.3.9-fr]{I, 1.3.9}) et
(\hyperref[I.1.3.12-fr]{I, 1.3.12}). La proposition résulte alors de
(\hyperref[II.3.1.8-fr]{3.1.8, (iii)}) et
(\hyperref[II.3.2.10-fr]{3.2.10}).
\end{proof}

\begin{env}[4.1.5]
\label{II.4.1.5-fr}
Avec les hypothèses de (\hyperref[II.4.1.1-fr]{4.1.1}), posons
$P=\mathbf{P}(\mathcal{E})$, et désignons par $p$ le morphisme structural
$P\to Y$. Comme on a par définition
$\mathcal{E}=(\mathbf{S}_{\mathcal{O}_Y}(\mathcal{E}))_1$, on a un
homomorphisme canonique
$\alpha_1:\mathcal{E}\to p_*(\mathcal{O}_P(1))$
(\hyperref[II.3.3.2.2-fr]{3.3.2.2}) et par suite aussi
(\hyperref[0.4.4.3-fr]{0, 4.4.3}) un homomorphisme canonique
\[
\label{II.4.1.5.1-fr}
  \alpha_1^\sharp:p^*(\mathcal{E})\longrightarrow\mathcal{O}_P(1).
\tag{4.1.5.1}
\]
\end{env}

\begin{proposition}[4.1.6]
\label{II.4.1.6-fr}
L'homomorphisme canonique (\hyperref[II.4.1.5.1-fr]{4.1.5.1}) est surjectif.
\end{proposition}

\begin{proof}
En effet, on a vu dans (\hyperref[II.3.3.2-fr]{3.3.2}) que
$\alpha_1^\sharp$ correspond fonctoriellement à l'homomorphisme canonique
\[
  \mathcal{E}\otimes_{\mathcal{O}_Y}
  \mathbf{S}_{\mathcal{O}_Y}(\mathcal{E})
  \longrightarrow
  (\mathbf{S}_{\mathcal{O}_Y}(\mathcal{E}))(1)~;
\]
comme par définition $\mathcal{E}$ engendre
$\mathbf{S}_{\mathcal{O}_Y}(\mathcal{E})$, cet homomorphisme est surjectif,
d'où la conclusion en vertu de (\hyperref[II.3.2.4-fr]{3.2.4}).
\end{proof}

\subsection{Morphismes d'un préschéma dans un fibré projectif.}
\label{subsection:II.4.2-fr}

\begin{env}[4.2.1]
\label{II.4.2.1-fr}
Gardant les notations de (\hyperref[II.4.1.5-fr]{4.1.5}), soient $X$ un
$Y$-préschéma, $q:X\to Y$ le morphisme structural, et soit $r:X\to P$ un
$Y$-\emph{morphisme}, de sorte que l'on a le diagramme commutatif
\[
  \xymatrix{
    P \ar[d]_p & X \ar[l]_r \ar[dl]^q
    \\Y
  }
\]

\oldpage[II]{73}

Comme le foncteur $r^*$ est exact à droite
(\hyperref[0.4.3.1-fr]{0, 4.3.1}), on déduit de l'homomorphisme
(\hyperref[II.4.1.5.1-fr]{4.1.5.1}), qui est surjectif
(\hyperref[II.4.1.6-fr]{4.1.6}), un homomorphisme surjectif
\[
  r^*(\alpha_1^\sharp):r^*(p^*(\mathcal{E}))\longrightarrow
  r^*(\mathcal{O}_P(1)).
\]
Mais $r^*(p^*(\mathcal{E}))=q^*(\mathcal{E})$, et
$r^*(\mathcal{O}_P(1))$ est localement isomorphe à
$r^*(\mathcal{O}_P)=\mathcal{O}_X$, autrement dit c'est un faisceau
inversible $\mathcal{L}_r$ sur $\mathcal{O}_X$ et on a ainsi défini à partir
de $r$ un $\mathcal{O}_X$-homomorphisme canonique surjectif
\[
\label{II.4.2.1.1-fr}
  \varphi_r:q^*(\mathcal{E})\longrightarrow\mathcal{L}_r.
\tag{4.2.1.1}
\]

Lorsque $Y=\operatorname{Spec}(A)$ est affine, et
$\mathcal{E}=\widetilde{E}$ on explicite encore cet homomorphisme de la façon
suivante~: étant donné $f\in E$, il résulte d'abord de
(\hyperref[II.2.6.3-fr]{2.6.3}) que l'on a
\[
\label{II.4.2.1.2-fr}
  r^{-1}(D_+(f))=X_{\varphi_r^\flat(f)}.
\tag{4.2.1.2}
\]
D'autre part, soit $V$ un ouvert affine de $X$ contenu dans
$r^{-1}(D_+(f))$, et soit $B$ son anneau, qui est une $A$-algèbre~; posons
$S=\mathbf{S}_A(E)$~; la restriction de $r$ à $V$ correspond à un
$A$-homomorphisme $\omega:S_{(f)}\to B$, on a
$q^*(\mathcal{E})|V=\widetilde{E\otimes_A B}$ et
$\mathcal{L}_r|V=\widetilde{L_r}$, où
$L_r=(S(1))_{(f)}\otimes_{S_{(f)}}B_{[\omega]}$
(\hyperref[I.1.6.5-fr]{I, 1.6.5}). La restriction de $\varphi_r$ à
$q^*(\mathcal{E})|V$ correspond au $B$-homomorphisme
$u:E\otimes_A B\to L_r$ qui transforme $x\otimes1$ en
$(x/1)\otimes f=(f/1)\otimes\omega(x/f)$. L'extension canonique de
$\varphi_r$ en un homomorphisme de $\mathcal{O}_X$-Algèbres
\[
  \psi_r:q^*(\mathbf{S}(\mathcal{E}))=
  \mathbf{S}(q^*(\mathcal{E}))\longrightarrow
  \mathbf{S}(\mathcal{L}_r)=\bigoplus_{n\geq0}\mathcal{L}_r^{\otimes n}
\]
est donc telle que la restriction de $\psi_r$ à
$q^*(\mathbf{S}_n(\mathcal{E}))|V$ correspond à l'homomorphisme
$\mathbf{S}_n(E\otimes_A B)=\mathbf{S}_n(E)\otimes_A B\to L_r^{\otimes n}$
qui transforme $s\otimes1$ en
$(f/1)^{\otimes n}\otimes\omega(s/f^n)$.
\end{env}

\begin{env}[4.2.2]
\label{II.4.2.2-fr}
Inversement, supposons donnés un morphisme $q:X\to Y$, un
$\mathcal{O}_X$-Module inversible $\mathcal{L}$ et un
$\mathcal{O}_Y$-Module quasi-cohérent $\mathcal{E}$~; à un homomorphisme
$\varphi:q^*(\mathcal{E})\to\mathcal{L}$ correspond canoniquement un
homomorphisme de $\mathcal{O}_X$-Algèbres quasi-cohérentes
\[
  \psi:\mathbf{S}(q^*(\mathcal{E}))=q^*(\mathbf{S}(\mathcal{E}))
  \longrightarrow\bigoplus_{n\geq0}\mathcal{L}^{\otimes n},
\]
et par suite (\hyperref[II.3.7.1-fr]{3.7.1}) un $Y$-morphisme
$r_{\mathcal{L},\psi}:G(\psi)\to
\operatorname{Proj}(\mathbf{S}(\mathcal{E}))=\mathbf{P}(\mathcal{E})$, que
nous noterons $r_{\mathcal{L},\varphi}$. Si $\varphi$ est surjectif, il en est
de même de $\psi$, donc (\hyperref[II.3.7.4-fr]{3.7.4})
$r_{\mathcal{L},\varphi}$ est partout défini. En outre, avec les notations de
(\hyperref[II.4.2.1-fr]{4.2.1}) et de (4.2.2)~:
\end{env}

\begin{proposition}[4.2.3]
\label{II.4.2.3-fr}
Étant donnés un morphisme $q:X\to Y$ et un $\mathcal{O}_Y$-Module
quasi-cohérent $\mathcal{E}$, les applications
$r\mapsto(\mathcal{L}_r,\varphi_r)$ et
$(\mathcal{L},\varphi)\mapsto r_{\mathcal{L},\varphi}$ mettent en
correspondance biunivoque l'ensemble des $Y$-morphismes
$r:X\to\mathbf{P}(\mathcal{E})$ et l'ensemble des classes d'équivalence des
couples $(\mathcal{L},\varphi)$ formés d'un $\mathcal{O}_X$-Module inversible
$\mathcal{L}$ et d'un homomorphisme surjectif
$\varphi:q^*(\mathcal{E})\to\mathcal{L}$, deux couples
$(\mathcal{L},\varphi)$, $(\mathcal{L}',\varphi')$ étant équivalents s'il
existe un $\mathcal{O}_X$-isomorphisme
$\tau:\mathcal{L}\xrightarrow{\sim}\mathcal{L}'$ tel que
$\varphi'=\tau\circ\varphi$.
\end{proposition}

\begin{proof}
Partons d'abord d'un $Y$-morphisme $r:X\to\mathbf{P}(\mathcal{E})$, formons
$\mathcal{L}_r$ et $\varphi_r$ (\hyperref[II.4.2.1-fr]{4.2.1}), puis
$r'=r_{\mathcal{L}_r,\varphi_r}$~; il résulte aussitôt de
(\hyperref[II.4.2.1-fr]{4.2.1}) et de
(\hyperref[II.3.7.2-fr]{3.7.2}) que les morphismes $r$ et $r'$ sont
identiques (en prenant pour générateur de $\mathcal{L}_r$ l'élément
$(f/1)\otimes1$ pour définir les homomorphismes $v_n$ de
(\hyperref[II.3.7.2-fr]{3.7.2})). Inversement, partons d'un couple
$(\mathcal{L},\varphi)$ et formons

\oldpage[II]{74}

$r''=r_{\mathcal{L},\varphi}$, puis $\mathcal{L}_{r''}$ et
$\varphi_{r''}$~; montrons qu'il y a un isomorphisme canonique
$\tau:\mathcal{L}_{r''}\xrightarrow{\sim}\mathcal{L}$ tel que
$\varphi=\tau\circ\varphi_{r''}$~; pour le définir on peut se placer dans le
cas où $Y=\operatorname{Spec}(A)$ et $X=\operatorname{Spec}(B)$ sont affines,
et (avec les notations de (\hyperref[II.4.2.1-fr]{4.2.1}) et de
(\hyperref[II.3.7.2-fr]{3.7.2})) faire correspondre à tout élément
$(x/1)\otimes1$ de $L_{r''}$ (avec $x\in E$) l'élément $v_1(x)c$ de $L$. On
vérifie aussitôt que $\tau$ ne dépend pas du générateur choisi $c$ de $L$~;
comme $v_1$ est par hypothèse surjectif, pour prouver que $\tau$ est un
isomorphisme, il suffit de voir que si $x/1=0$ dans $(S(1))_{(f)}$, on a
$v_1(x)/1=0$ dans $B_g$~; mais la première relation signifie que $f^nx=0$
dans $\mathbf{S}_{n+1}(E)$ pour un certain $n$, et on en déduit que
$v_{n+1}(f^nx)=g^nv_1(x)=0$ dans $B$, d'où la conclusion. Enfin, il est
immédiat que pour deux couples équivalents $(\mathcal{L},\varphi)$,
$(\mathcal{L}',\varphi')$, on a
$r_{\mathcal{L},\varphi}=r_{\mathcal{L}',\varphi'}$.
\end{proof}

En particulier, pour $X=Y$~:

\begin{theorem}[4.2.4]
\label{II.4.2.4-fr}
L'ensemble des $Y$-sections de $\mathbf{P}(\mathcal{E})$ est en
correspondance biunivoque canonique avec l'ensemble des
sous-$\mathcal{O}_Y$-Modules quasi-cohérents $\mathcal{F}$ de $\mathcal{E}$
tels que $\mathcal{E}/\mathcal{F}$ soit inversible.
\end{theorem}

On notera que cette propriété correspond à la définition classique de
l'\og espace projectif\fg{} comme ensemble des hyperplans d'un espace
vectoriel (le cas classique correspondant à $Y=\operatorname{Spec}(K)$, où
$K$ est un corps, et $\mathcal{E}=\widetilde{E}$, $E$ étant un $K$-espace
vectoriel de dimension finie~; les $\mathcal{F}$ ayant la propriété énoncée
dans (\hyperref[II.4.2.4-fr]{4.2.4}) correspondent alors aux hyperplans de
$E$, et on sait d'autre part que les $Y$-sections de
$\mathbf{P}(\mathcal{E})$ sont alors les points de
$\mathbf{P}(\mathcal{E})$ rationnels sur $K$
(\hyperref[I.3.4.5-fr]{I, 3.4.5})).

\begin{remark}[4.2.5]
\label{II.4.2.5-fr}
Comme il y a correspondance biunivoque canonique entre $Y$-morphismes de $X$
dans $P$ et leurs morphismes graphes, $X$-sections de $P\times_YX$
(\hyperref[I.3.3.14-fr]{I, 3.3.14}), on voit qu'inversement
(\hyperref[II.4.2.3-fr]{4.2.3}) peut se déduire de
(\hyperref[II.4.2.4-fr]{4.2.4}). Désignons par
$\operatorname{Hyp}_Y(X,\mathcal{E})$ l'ensemble des
sous-$\mathcal{O}_X$-Modules quasi-cohérents $\mathcal{F}$ de
$\mathcal{E}\otimes_Y\mathcal{O}_X=q^*(\mathcal{E})$ tels que
$q^*(\mathcal{E})/\mathcal{F}$ soit un $\mathcal{O}_X$-Module inversible. Si
$g:X'\to X$ est un $Y$-morphisme, il résulte du fait que $g^*$ est exact à
droite que l'on a
\[
  g^*(q^*(\mathcal{E})/\mathcal{F})=
  g^*(q^*(\mathcal{E}))/g^*(\mathcal{F}),
\]
donc que ce dernier faisceau est inversible, et par suite
$\operatorname{Hyp}_Y(X,\mathcal{E})$ est un foncteur contravariant dans la
catégorie des $Y$-préschémas. On peut alors interpréter le th.
(\hyperref[II.4.2.4-fr]{4.2.4}) comme définissant un isomorphisme canonique
des foncteurs $\operatorname{Hom}_Y(X,\mathbf{P}(\mathcal{E}))$ et
$\operatorname{Hyp}_Y(X,\mathcal{E})$ contravariants en $X$ variable dans la
catégorie des $Y$-préschémas. Ceci fournit aussi une caractérisation du fibré
projectif $P=\mathbf{P}(\mathcal{E})$ par la propriété universelle suivante,
plus proche de l'intuition géométrique que les constructions des \S\S~2 et
3~: pour tout morphisme $q:X\to Y$ et tout $\mathcal{O}_X$-Module inversible
$\mathcal{L}$, quotient de $\mathcal{E}\otimes_{\mathcal{O}_Y}\mathcal{O}_X$,
il existe un $Y$-morphisme unique $r:X\to P$ tel que
$\mathcal{L}=r^*(\mathcal{O}_P(1))$.

On verra plus tard comment on peut de la même manière définir, entre autres,
les schémas \og grassmanniennes\fg{}.
\end{remark}

\begin{corollary}[4.2.6]
\label{II.4.2.6-fr}
Supposons que tout $\mathcal{O}_Y$-Module inversible soit trivial
(\hyperref[I.2.4.8-fr]{I, 2.4.8}). Soit $V$ le groupe
$\operatorname{Hom}_{\mathcal{O}_Y}(\mathcal{E},\mathcal{O}_Y)$, considéré
comme module sur l'anneau $A=\Gamma(Y,\mathcal{O}_Y)$, et soit $V^*$ le
sous-ensemble de $V$ formé des homomorphismes surjectifs. Alors l'ensemble
des $Y$-sections de $\mathbf{P}(\mathcal{E})$ s'identifie canoniquement à
$V^*/A^*$, où $A^*$ est le groupe des unités de $A$.
\end{corollary}

\oldpage[II]{75}

En particulier~:
\begin{enumerate}
  \item[1\textsuperscript{o}] Le corollaire
  (\hyperref[II.4.2.6-fr]{4.2.6}) s'applique lorsque $Y$ est un
  \emph{schéma local} (\hyperref[I.2.4.8-fr]{I, 2.4.8}). Soient $Y$ un
  préschéma quelconque, $y$ un point de $Y$,
  $Y'=\operatorname{Spec}(k(y))$~; la fibre $p^{-1}(y)$ de
  $\mathbf{P}(\mathcal{E})$ est, en vertu de
  (\hyperref[II.4.1.3.1-fr]{4.1.3.1}), identifiée à
  $\mathbf{P}(\mathcal{E}^y)$, où
  $\mathcal{E}^y=\mathcal{E}_y\otimes_{\mathcal{O}_y}k(y)
  =\mathcal{E}_y/\mathfrak{m}_y\mathcal{E}_y$ est considéré comme espace
  vectoriel sur $k(y)$. Plus généralement, si $K$ est une extension de
  $k(y)$, $p^{-1}(y)\otimes_{k(y)}K$ s'identifie à
  $\mathbf{P}(\mathcal{E}^y\otimes_{k(y)}K)$. Le cor.
  (\hyperref[II.4.2.6-fr]{4.2.6}) montre donc que l'ensemble des
  \emph{points géométriques de $\mathbf{P}(\mathcal{E})$ à valeurs dans
  l'extension $K$ de $k(y)$} (\hyperref[I.3.4.5-fr]{I, 3.4.5}), que l'on
  peut encore appeler \emph{fibre géométrique rationnelle sur $K$ de
  $\mathbf{P}(\mathcal{E})$ au-dessus du point $y$}, s'identifie à
  l'\emph{espace projectif} associé au \emph{dual} du $K$-espace vectoriel
  $\mathcal{E}^y\otimes_{k(y)}K$.
  \item[2\textsuperscript{o}] Supposons que $Y$ soit affine d'anneau $A$, et
  en outre que tout $\mathcal{O}_Y$-Module inversible soit trivial~; prenons
  en outre $\mathcal{E}=\mathcal{O}_Y^n$~; alors, dans
  (\hyperref[II.4.2.6-fr]{4.2.6}), $V$ s'identifie à $A^n$
  (\hyperref[I.1.3.8-fr]{I, 1.3.8}), $V^*$ à l'ensemble des systèmes
  $(f_i)_{1\leq i\leq n}$ d'éléments de $A$ engendrant l'idéal $A$~; deux
  tels systèmes définissent la même $Y$-section de
  $\mathbf{P}_Y^{n-1}=\mathbf{P}_A^{n-1}$, autrement dit le même point de
  $\mathbf{P}_A^{n-1}$ à valeurs dans $A$ si et seulement si l'un se déduit
  de l'autre par multiplication par un élément inversible de $A$.
\end{enumerate}

Ces propriétés justifient la terminologie de \og fibré projectif\fg{} pour
$\mathbf{P}(\mathcal{E})$. On notera que la définition que nous obtenons
ainsi de l'\og espace projectif\fg{} est en fait duale de la définition
classique~; cela nous est imposé par la nécessité de pouvoir définir
$\mathbf{P}(\mathcal{E})$ pour un $\mathcal{O}_Y$-Module quasi-cohérent
quelconque $\mathcal{E}$, non nécessairement localement libre.

\begin{remark}[4.2.7]
\label{II.4.2.7-fr}
Nous verrons au chapitre~V que si $Y$ est localement noethérien et connexe,
et si $\mathcal{E}$ est localement libre, alors, en posant
$P=\mathbf{P}(\mathcal{E})$, tout $\mathcal{O}_P$-Module inversible
$\mathcal{L}$ est isomorphe à un $\mathcal{O}_P$-Module de la forme
$\mathcal{L}'\otimes_{\mathcal{O}_Y}\mathcal{O}_P(m)$, où $\mathcal{L}'$ est
un $\mathcal{O}_Y$-Module inversible, bien déterminé à un isomorphisme près,
et $m$ un entier bien déterminé. En d'autres termes,
$\mathrm{H}^1(P,\mathcal{O}_P^*)$ est isomorphe à
$\mathbf{Z}\times\mathrm{H}^1(Y,\mathcal{O}_Y^*)$
(\hyperref[0.5.4.7-fr]{0, 5.4.7}). Nous verrons aussi
(\hyperref[III.2.1.14-fr]{III, 2.1.14}, compte tenu de
(\hyperref[0.5.4.10-fr]{0, 5.4.10})) que
$p_*(\mathcal{L}^{\otimes m})=0$ si $m<0$ et que
$p_*(\mathcal{L}^{\otimes m})$ est isomorphe à
$\mathcal{L}'\otimes_{\mathcal{O}_Y}
(\mathbf{S}_{\mathcal{O}_Y}(\mathcal{E}))_m$ si $m\geq0$. Si $\mathcal{F}$
est un $\mathcal{O}_Y$-Module quasi-cohérent, tout $Y$-morphisme
$\mathbf{P}(\mathcal{E})\to\mathbf{P}(\mathcal{F})$ est donc déterminé par
la donnée d'un $\mathcal{O}_Y$-Module inversible $\mathcal{L}'$, d'un entier
$m\geq0$ et d'un $\mathcal{O}_Y$-homomorphisme
$\psi:\mathcal{F}\to\mathcal{L}'\otimes_{\mathcal{O}_Y}
(\mathbf{S}_{\mathcal{O}_Y}(\mathcal{E}))_m$ tel que l'homomorphisme
correspondant $\psi^\sharp$ de $\mathcal{O}_{\mathbf{P}(\mathcal{F})}$-Modules
soit surjectif. Nous verrons aussi que si le $Y$-morphisme considéré est un
isomorphisme, alors $m=1$ et $\mathcal{F}$ est isomorphe à
$\mathcal{E}\otimes_{\mathcal{O}_Y}\mathcal{L}'$ (réciproque de
(\hyperref[II.4.1.4-fr]{4.1.4})). Ceci permettra de déterminer le faisceau
des germes d'automorphismes de $\mathbf{P}(\mathcal{E})$ comme quotient du
faisceau de groupes $\mathcal{A}ut(\mathcal{E})$ (localement isomorphe à
$\operatorname{GL}(n,\mathcal{O}_Y)$ si $\mathcal{E}$ est de rang $n$) par
$\mathcal{O}_Y^*$.
\end{remark}

\begin{env}[4.2.8]
\label{II.4.2.8-fr}
Conservant les notations de (\hyperref[II.4.2.1-fr]{4.2.1}), soit
$u:X'\to X$ un morphisme~; si le $Y$-morphisme $r:X\to P$ correspond à
l'homomorphisme $\varphi:q^*(\mathcal{E})\to\mathcal{L}$, le $Y$-morphisme
$r\circ u$ correspond à
$u^*(\varphi):u^*(q^*(\mathcal{E}))\to u^*(\mathcal{L})$, comme il résulte
aussitôt des définitions.
\end{env}

\begin{env}[4.2.9]
\label{II.4.2.9-fr}
Soient $\mathcal{E}$, $\mathcal{F}$ deux $\mathcal{O}_Y$-Modules
quasi-cohérents, $v:\mathcal{E}\to\mathcal{F}$ un homomorphisme surjectif,
$j=\mathbf{P}(v)$ l'immersion fermée correspondante
$\mathbf{P}(\mathcal{F})\to\mathbf{P}(\mathcal{E})$
(\hyperref[II.4.1.2-fr]{4.1.2}). Si le $Y$-morphisme
$r:X\to\mathbf{P}(\mathcal{F})$ correspond à l'homomorphisme
$\varphi:q^*(\mathcal{F})\to\mathcal{L}$, le

\oldpage[II]{76}

$Y$-morphisme $j\circ r$ correspond à
\[
  q^*(\mathcal{E})\xrightarrow{q^*(v)}q^*(\mathcal{F})
  \xrightarrow{\varphi}\mathcal{L}~;
\]
cela résulte encore de la définition donnée en
(\hyperref[II.4.2.1-fr]{4.2.1}).
\end{env}

\begin{env}[4.2.10]
\label{II.4.2.10-fr}
Soit $\psi:Y'\to Y$ un morphisme, et posons
$\mathcal{E}'=\psi^*(\mathcal{E})$. Si le $Y$-morphisme $r:X\to P$ correspond
à l'homomorphisme $\varphi:q^*(\mathcal{E})\to\mathcal{L}$, le
$Y'$-morphisme
\[
  r_{(Y')}:X_{(Y')}\longrightarrow P'=\mathbf{P}(\mathcal{E}')
\]
correspond à
$\varphi_{(Y')}:q_{(Y')}^*(\mathcal{E}')=
  q^*(\mathcal{E})\otimes_Y\mathcal{O}_{Y'}\to
\mathcal{L}\otimes_Y\mathcal{O}_{Y'}$. En effet, en vertu de
(\hyperref[II.4.1.3.1-fr]{4.1.3.1}), on a le diagramme commutatif
\[
  \xymatrix{
    Y' \ar[d]
    & P'=P_{(Y')} \ar[l]_{p_{(Y')}} \ar[d]^{u}
    & X_{(Y')} \ar[l]_{r_{(Y')}} \ar[d]^{v}
    \\Y
    & P \ar[l]_{p}
    & X \ar[l]_{r}
  }
\]
D'après (\hyperref[II.4.1.3.2-fr]{4.1.3.2}), on a
\[
  (r_{(Y')})^*(\mathcal{O}_{P'}(1))
  =(r_{(Y')})^*(u^*(\mathcal{O}_P(1)))
  =v^*(r^*(\mathcal{O}_P(1)))
  =v^*(\mathcal{L})=\mathcal{L}\otimes_Y\mathcal{O}_{Y'}~;
\]
d'autre part, $u^*(\alpha_1^\sharp)$ n'est autre que l'homomorphisme
canonique
$\alpha_1^\sharp:(p_{(Y')})^*(\mathcal{E}')\to\mathcal{O}_{P'}(1)$~; cela
se voit en explicitant les homomorphismes canoniques $\alpha_1^\sharp$ dans
$P$ et $P'$ comme dans (\hyperref[II.4.1.6-fr]{4.1.6}). D'où notre assertion.
\end{env}

\subsection{Le morphisme de Segre.}
\label{subsection:II.4.3-fr}

\begin{env}[4.3.1]
\label{II.4.3.1-fr}
Soient $Y$ un préschéma, $\mathcal{E}$, $\mathcal{F}$ deux
$\mathcal{O}_Y$-Modules quasi-cohérents~; posons
$P_1=\mathbf{P}(\mathcal{E})$, $P_2=\mathbf{P}(\mathcal{F})$, et désignons
par $p_1:P_1\to Y$, $p_2:P_2\to Y$ les morphismes structuraux. Soit
$Q=P_1\times_YP_2$, et soient $q_1:Q\to P_1$, $q_2:Q\to P_2$ les
projections canoniques~; le $\mathcal{O}_Q$-Module
\[
  \mathcal{L}=\mathcal{O}_{P_1}(1)\otimes_Y\mathcal{O}_{P_2}(1)
  :=q_1^*(\mathcal{O}_{P_1}(1))\otimes_{\mathcal{O}_Q}
  q_2^*(\mathcal{O}_{P_2}(1))
\]
est inversible comme produit tensoriel de deux $\mathcal{O}_Q$-Modules
inversibles (\hyperref[0.5.4.4-fr]{0, 5.4.4}). D'autre part, si
$r=p_1\circ q_1=p_2\circ q_2$ est le morphisme structural $Q\to Y$, on a
\[
  r^*(\mathcal{E}\otimes_{\mathcal{O}_Y}\mathcal{F})
  =q_1^*(p_1^*(\mathcal{E}))\otimes_{\mathcal{O}_Q}
  q_2^*(p_2^*(\mathcal{F}))
\]
(\hyperref[0.4.3.3-fr]{0, 4.3.3})~; les homomorphismes canoniques surjectifs
(\hyperref[II.4.1.5.1-fr]{4.1.5.1})
$p_1^*(\mathcal{E})\to\mathcal{O}_{P_1}(1)$,
$p_2^*(\mathcal{F})\to\mathcal{O}_{P_2}(1)$ donnent donc, par produit
tensoriel, un homomorphisme canonique
\[
\label{II.4.3.1.1-fr}
  s:r^*(\mathcal{E}\otimes_{\mathcal{O}_Y}\mathcal{F})\longrightarrow
  \mathcal{L}
\tag{4.3.1.1}
\]
évidemment surjectif~; on en déduit donc
(\hyperref[II.4.2.2-fr]{4.2.2}) un morphisme canonique, dit
\emph{morphisme de Segre}~:
\[
\label{II.4.3.1.2-fr}
  \varsigma:\mathbf{P}(\mathcal{E})\times_Y\mathbf{P}(\mathcal{F})
  \longrightarrow
  \mathbf{P}(\mathcal{E}\otimes_{\mathcal{O}_Y}\mathcal{F}).
\tag{4.3.1.2}
\]
Explicitons le morphisme $\varsigma$, lorsque $Y=\operatorname{Spec}(A)$ est
affine, $\mathcal{E}=\widetilde{E}$, $\mathcal{F}=\widetilde{F}$, $E$ et $F$
étant deux $A$-modules, d'où
$\mathcal{E}\otimes_{\mathcal{O}_Y}\mathcal{F}=\widetilde{E\otimes_A F}$
(\hyperref[I.1.3.12-fr]{I, 1.3.12})~; posons
$R=\mathbf{S}_A(E)$, $S=\mathbf{S}_A(F)$,
$T=\mathbf{S}_A(E\otimes_A F)$~; soient $f\in E$, $g\in F$, et considérons
l'ouvert affine
\[
  D_+(f)\times_YD_+(g)=\operatorname{Spec}(B)
\]

\oldpage[II]{77}

de $Q$, où $B=R_{(f)}\otimes_AS_{(g)}$~; la restriction de $\mathcal{L}$ à
cet ouvert affine est $\widetilde{L}$, où
\[
  L=(R(1))_{(f)}\otimes_A(S(1))_{(g)}
\]
et l'élément $c=(f/1)\otimes(g/1)$ est un générateur de $L$ considéré comme
$B$-module libre (\hyperref[II.2.5.7-fr]{2.5.7}). L'homomorphisme
(\hyperref[II.4.3.1.1-fr]{4.3.1.1}) correspond à l'homomorphisme
\[
  (x\otimes y)\otimes b\longmapsto b((x/1)\otimes(y/1))
\]
de $(E\otimes_A F)\otimes_AB$ dans $L$. Avec les notations de
(\hyperref[II.3.7.2-fr]{3.7.2}), on a donc
$v_1(x\otimes y)=(x/f)\otimes(y/g)$~; la restriction de $\varsigma$ à
$D_+(f)\times_YD_+(g)$ est un morphisme de ce schéma affine dans
$D_+(f\otimes g)$, correspondant à l'homomorphisme d'anneaux
$\omega:T_{(f\otimes g)}\to R_{(f)}\otimes_AS_{(g)}$ tel que
\[
\label{II.4.3.1.3-fr}
  \omega((x\otimes y)/(f\otimes g))=(x/f)\otimes(y/g)
\tag{4.3.1.3}
\]
pour $x\in E$ et $y\in F$.
\end{env}

\begin{env}[4.3.2]
\label{II.4.3.2-fr}
Il résulte de (\hyperref[II.4.2.3-fr]{4.2.3}) que l'on a un isomorphisme
canonique
\[
\label{II.4.3.2.1-fr}
  \tau:\varsigma^*(\mathcal{O}_P(1))\xrightarrow{\sim}
  \mathcal{O}_{P_1}(1)\otimes_Y\mathcal{O}_{P_2}(1)
\tag{4.3.2.1}
\]
où on a posé
$P=\mathbf{P}(\mathcal{E}\otimes_{\mathcal{O}_Y}\mathcal{F})$. Montrons
que, pour $x\in\Gamma(Y,\mathcal{E})$ et
$y\in\Gamma(Y,\mathcal{F})$, on a
\[
\label{II.4.3.2.2-fr}
  \tau(\alpha_1(x\otimes y))=\alpha_1(x)\otimes\alpha_1(y).
\tag{4.3.2.2}
\]
En effet, on est ramené au cas où $Y$ est affine, et on a alors, avec les
notations de (\hyperref[II.4.3.1-fr]{4.3.1}) et de
(\hyperref[II.2.6.2-fr]{2.6.2}),
$\alpha_1^{f\otimes g}(x\otimes y)=(x\otimes y)/1$ dans
$(T(1))_{(f\otimes g)}$, $\alpha_1^f(x)=x/1$ dans $(R(1))_{(f)}$ et
$\alpha_1^g(y)=y/1$ dans $(S(1))_{(g)}$. La définition de $\tau$ donnée dans
(\hyperref[II.4.2.3-fr]{4.2.3}) et le calcul de $v_1$ fait dans
(\hyperref[II.4.3.1-fr]{4.3.1}) prouvent alors aussitôt
(\hyperref[II.4.3.2.2-fr]{4.3.2.2}). On en déduit la formule
\[
\label{II.4.3.2.3-fr}
  \varsigma^{-1}(P_{x\otimes y})=(P_1)_x\times_Y(P_2)_y
\tag{4.3.2.3}
\]
avec les notations de (\hyperref[II.3.1.4-fr]{3.1.4}). En effet, compte tenu
de (\hyperref[II.3.3.3-fr]{3.3.3}), la formule
(\hyperref[II.4.3.2.2-fr]{4.3.2.2}) ramène (en revenant au cas affine, à
l'aide de (\hyperref[I.3.2.7-fr]{I, 3.2.7} et
(\hyperref[I.3.2.3-fr]{3.2.3})) à prouver le lemme suivant~:

\begin{lemma}[4.3.2.4]
\label{II.4.3.2.4-fr}
Soient $B$, $B'$ deux $A$-algèbres, et soient
$Y=\operatorname{Spec}(A)$, $Z=\operatorname{Spec}(B)$,
$Z'=\operatorname{Spec}(B')$~; quels que soient $t\in B$, $t'\in B'$, on a
$D(t\otimes t')=D(t)\times_YD(t')$.
\end{lemma}

\begin{proof}
En effet, si $p:Z\times_YZ'\to Z$ et $p':Z\times_YZ'\to Z'$ sont les
projections canoniques, il résulte de
(\hyperref[I.1.2.2.2-fr]{I, 1.2.2.2}) que l'on a
$p^{-1}(D(t))=D(t\otimes1)$ et
$p'^{-1}(D(t'))=D(1\otimes t')$~; la conclusion résulte de
(\hyperref[I.3.2.7-fr]{I, 3.2.7}) et
(\hyperref[I.1.1.9.1-fr]{I, 1.1.9.1}), puisque
$(t\otimes1)(1\otimes t')=t\otimes t'$.
\end{proof}
\end{env}

\begin{proposition}[4.3.3]
\label{II.4.3.3-fr}
Le morphisme de Segre est une immersion fermée.
\end{proposition}

\begin{proof}
En effet, la question étant locale sur $Y$, on est ramené au cas où $Y$ est
affine. Avec les notations de (\hyperref[II.4.3.1-fr]{4.3.1}) et
(\hyperref[II.4.3.2-fr]{4.3.2}), les $D_+(f\otimes g)$ forment alors une base
de la topologie de $P$, puisque les $f\otimes g$ engendrent $T$ lorsque $f$
parcourt $E$ et $g$ parcourt $F$. D'autre part, on a en vertu de
(\hyperref[II.4.3.2.3-fr]{4.3.2.3})
$\varsigma^{-1}(D_+(f\otimes g))=D_+(f)\times_YD_+(g)$. Il suffit donc
(\hyperref[I.4.2.4-fr]{I, 4.2.4}) de prouver que la restriction de
$\varsigma$ à $D_+(f)\times_YD_+(g)$ est une immersion fermée dans
$D_+(f\otimes g)$. Or, avec les mêmes notations, la formule
(\hyperref[II.4.3.1.3-fr]{4.3.1.3}) montre que $\omega$ est surjectif, ce qui
termine la démonstration.
\end{proof}

\begin{env}[4.3.4]
\label{II.4.3.4-fr}
Le morphisme de Segre est fonctoriel en $\mathcal{E}$ et $\mathcal{F}$,
lorsqu'on restreint les

\oldpage[II]{78}

homomorphismes de $\mathcal{O}_Y$-Modules quasi-cohérents aux homomorphismes
surjectifs. Il faut en effet montrer que si
$\mathcal{E}\to\mathcal{E}'$ est un $\mathcal{O}_Y$-homomorphisme surjectif,
le diagramme
\[
  \xymatrix{
    \mathbf{P}(\mathcal{E}')\times_Y\mathbf{P}(\mathcal{F})
      \ar[r]^{j\times1} \ar[d]_{\varsigma}
    & \mathbf{P}(\mathcal{E})\times_Y\mathbf{P}(\mathcal{F})
      \ar[d]^{\varsigma}
    \\\mathbf{P}(\mathcal{E}'\otimes\mathcal{F}) \ar[r]
    & \mathbf{P}(\mathcal{E}\otimes\mathcal{F})
  }
\]
est commutatif, $j$ désignant l'immersion fermée canonique
$\mathbf{P}(\mathcal{E}')\to\mathbf{P}(\mathcal{E})$. Posons
$P_1'=\mathbf{P}(\mathcal{E}')$ et gardons par ailleurs les notations de
(\hyperref[II.4.3.1-fr]{4.3.1})~; $j\times1$ est une immersion fermée
(\hyperref[I.4.3.1-fr]{I, 4.3.1}) et on a, à des isomorphismes près,
\[
  (j\times1)^*(\mathcal{O}_{P_1}(1)\otimes\mathcal{O}_{P_2}(1))
  =j^*(\mathcal{O}_{P_1}(1))\otimes\mathcal{O}_{P_2}(1)
  =\mathcal{O}_{P_1'}(1)\otimes\mathcal{O}_{P_2}(1)
\]
en vertu de (\hyperref[II.4.1.2.1-fr]{4.1.2.1}) et de
(\hyperref[I.9.1.5-fr]{I, 9.1.5})~; notre assertion résulte donc de
(\hyperref[II.4.2.8-fr]{4.2.8}) et
(\hyperref[II.4.2.9-fr]{4.2.9}).
\end{env}

\begin{env}[4.3.5]
\label{II.4.3.5-fr}
Avec les notations de (\hyperref[II.4.3.1-fr]{4.3.1}), soit
$\psi:Y'\to Y$ un morphisme, et posons
$\mathcal{E}'=\psi^*(\mathcal{E})$, $\mathcal{F}'=\psi^*(\mathcal{F})$~;
alors le morphisme de Segre
$\mathbf{P}(\mathcal{E}')\times\mathbf{P}(\mathcal{F}')
\to\mathbf{P}(\mathcal{E}'\otimes\mathcal{F}')$ s'identifie à
$\varsigma_{(Y')}$. En effet, gardant les notations de
(\hyperref[II.4.3.1-fr]{4.3.1}), posons par ailleurs
$P_1'=\mathbf{P}(\mathcal{E}')$, $P_2'=\mathbf{P}(\mathcal{F}')$~; on sait
(\hyperref[II.4.1.3.1-fr]{4.1.3.1}) que $P_i'$ s'identifie à $(P_i)_{(Y')}$
($i=1,2$), donc le morphisme structural $P_1'\times P_2'\to Y'$ s'identifie
à $r_{(Y')}$. D'autre part,
$\mathcal{E}'\otimes\mathcal{F}'$ s'identifie à
$\psi^*(\mathcal{E}\otimes\mathcal{F})$, donc
$\mathbf{P}(\mathcal{E}'\otimes\mathcal{F}')$ s'identifie à
$(\mathbf{P}(\mathcal{E}\otimes\mathcal{F}))_{(Y')}$
(\hyperref[II.4.1.3.1-fr]{4.1.3.1}). Enfin,
$\mathcal{O}_{P_1'}(1)\otimes_{Y'}\mathcal{O}_{P_2'}(1)=\mathcal{L}'$
s'identifie à $\mathcal{L}\otimes_Y\mathcal{O}_{Y'}$, en vertu de
(\hyperref[II.4.1.3.2-fr]{4.1.3.2}) et de
(\hyperref[I.9.1.11-fr]{I, 9.1.11}). L'homomorphisme canonique
$(r_{(Y')})^*(\mathcal{E}'\otimes\mathcal{F}')\to\mathcal{L}'$ s'identifie
alors à $s_{(Y')}$, et notre assertion résulte de
(\hyperref[II.4.2.10-fr]{4.2.10}).
\end{env}

\begin{remark}[4.3.6]
\label{II.4.3.6-fr}
Le préschéma somme de $\mathbf{P}(\mathcal{E})$ et
$\mathbf{P}(\mathcal{F})$ est de même canoniquement isomorphe à un
sous-préschéma fermé de $\mathbf{P}(\mathcal{E}\oplus\mathcal{F})$. En
effet, les homomorphismes surjectifs
$\mathcal{E}\oplus\mathcal{F}\to\mathcal{E}$,
$\mathcal{E}\oplus\mathcal{F}\to\mathcal{F}$ correspondent à des immersions
fermées
$\mathbf{P}(\mathcal{E})\to\mathbf{P}(\mathcal{E}\oplus\mathcal{F})$,
$\mathbf{P}(\mathcal{F})\to\mathbf{P}(\mathcal{E}\oplus\mathcal{F})$~;
tout revient à voir que les espaces sous-jacents des sous-préschémas fermés
de $\mathbf{P}(\mathcal{E}\oplus\mathcal{F})$ ainsi obtenus sont sans point
commun. La question étant locale sur $Y$, on peut adopter les notations de
(\hyperref[II.4.3.1-fr]{4.3.1})~; or $\mathbf{S}_n(E)$ et
$\mathbf{S}_n(F)$ s'identifient à des sous-modules de
$\mathbf{S}_n(E\oplus F)$ d'intersection réduite à $0$~; et si
$\mathfrak{p}$ est un idéal premier gradué de $\mathbf{S}(E)$ tel que
$\mathfrak{p}\cap\mathbf{S}_n(E)\neq\mathbf{S}_n(E)$ pour tout $n\geq0$,
il lui correspond dans $\mathbf{S}(E\oplus F)$ un idéal premier gradué dont
la trace sur $\mathbf{S}_n(E)$ est
$\mathfrak{p}\cap\mathbf{S}_n(E)$, mais qui contient $\mathbf{S}_+(F)$,
comme on le voit aussitôt~; deux points de
$\operatorname{Proj}(\mathbf{S}(E))$ et
$\operatorname{Proj}(\mathbf{S}(F))$ respectivement ne peuvent donc avoir
même image dans $\operatorname{Proj}(\mathbf{S}(E\oplus F))$.
\end{remark}

\subsection{Immersions dans les fibrés projectifs. Faisceaux très amples}
\label{subsection:II.4.4-fr}

\begin{proposition}[4.4.1]
\label{II.4.4.1-fr}
Soient $Y$ un schéma quasi-compact, ou un préschéma dont l'espace sous-jacent
est noethérien, $q:X\to Y$ un morphisme de type fini, $\mathcal{L}$ un
$\mathcal{O}_X$-Module inversible.
\begin{enumerate}
  \item[\rm{(i)}] Soient $\mathcal{S}$ une $\mathcal{O}_Y$-Algèbre graduée
  quasi-cohérente à degrés positifs,
  $\psi:q^*(\mathcal{S})\to\bigoplus_{n\geq0}\mathcal{L}^{\otimes n}$ un
  homomorphisme d'Algèbres graduées. Pour que $r_{\mathcal{L},\psi}$ soit
  partout défini et soit une immersion, il faut et

\oldpage[II]{79}

  il suffit qu'il existe un entier $n\geq0$, un sous-$\mathcal{O}_Y$-Module
  quasi-cohérent de type fini $\mathcal{E}$ de $\mathcal{S}_n$, tels que
  l'homomorphisme
  \[
    \psi'=\psi_n\circ q^*(j):q^*(\mathcal{E})
      \longrightarrow\mathcal{L}^{\otimes n}=\mathcal{L}'
  \]
  ($j$ étant l'injection $\mathcal{E}\to\mathcal{S}_n$) soit surjectif et que
  le morphisme
  $r_{\mathcal{L}',\psi'}:X\to\mathbf{P}(\mathcal{E})$ soit une immersion.
  \item[\rm{(ii)}] Soient $\mathcal{F}$ un $\mathcal{O}_Y$-Module
  quasi-cohérent, $\varphi:q^*(\mathcal{F})\to\mathcal{L}$ un homomorphisme
  surjectif. Pour que le morphisme $r_{\mathcal{L},\varphi}$ soit une
  immersion $X\to\mathbf{P}(\mathcal{F})$, il faut et il suffit qu'il existe
  un sous-$\mathcal{O}_Y$-Module quasi-cohérent de type fini $\mathcal{E}$ de
  $\mathcal{F}$ tel que l'homomorphisme
  $\varphi'=\varphi\circ q^*(j):q^*(\mathcal{E})\to\mathcal{L}$ (où
  $j:\mathcal{E}\to\mathcal{F}$ est l'injection canonique) soit surjectif et
  que $r_{\mathcal{L},\varphi'}:X\to\mathbf{P}(\mathcal{E})$ soit une
  immersion.
\end{enumerate}
\end{proposition}

\begin{proof}
\begin{enumerate}
  \item[\rm{(i)}] Le fait que $r_{\mathcal{L},\psi}$ soit partout défini et
  soit une immersion équivaut d'après
  (\hyperref[II.3.8.5-fr]{3.8.5}) à l'existence de $n\geq0$ et de
  $\mathcal{E}$ tels que, si $\mathcal{S}'$ est la sous-algèbre de
  $\mathcal{S}$ engendrée par $\mathcal{E}$, l'homomorphisme
  $q^*(\mathcal{E})\to\mathcal{L}^{\otimes n}$ soit surjectif et que
  $r_{\mathcal{L}',\psi'}:X\to\operatorname{Proj}(\mathcal{S}')$ soit partout
  défini et soit une immersion. On a d'autre part un homomorphisme canonique
  surjectif $\mathbf{S}(\mathcal{E})\to\mathcal{S}'$ auquel correspond une
  immersion fermée
  $\operatorname{Proj}(\mathcal{S}')\to\mathbf{P}(\mathcal{E})$
  (\hyperref[II.3.6.2-fr]{3.6.2})~; d'où la conclusion.
  \item[\rm{(ii)}] Comme $\mathcal{F}$ est limite inductive de ses
  sous-Modules quasi-cohérents de type fini $\mathcal{E}_\lambda$
  (\hyperref[I.9.4.9-fr]{I, 9.4.9}), $\mathbf{S}(\mathcal{F})$ est limite
  inductive des $\mathbf{S}(\mathcal{E}_\lambda)$~; la conclusion résulte
  alors de (\hyperref[II.3.8.4-fr]{3.8.4}), en observant que, dans la
  démonstration de (\hyperref[II.3.8.4-fr]{3.8.4}), on peut prendre tous les
  $n_i$ égaux à $1$~: en effet ($Y$ étant supposé affine), si
  $r=r_{\mathcal{L},\varphi}$ est une immersion, $r(X)$ est un sous-espace
  quasi-compact de $\mathbf{P}(\mathcal{F})$ que l'on peut recouvrir par un
  nombre fini d'ouverts de $\mathbf{P}(\mathcal{F})$ de la forme $D_+(f)$ avec
  $f\in F$, tels que $D_+(f)\cap r(X)$ soit fermé.
\end{enumerate}
\end{proof}

\begin{definition}[4.4.2]
\label{II.4.4.2-fr}
Soient $Y$ un préschéma, $q:X\to Y$ un morphisme. On dit qu'un
$\mathcal{O}_X$-Module inversible $\mathcal{L}$ est \emph{très ample pour
$q$} (ou \emph{très ample pour $Y$}, ou simplement \emph{très ample} si
aucune confusion n'en résulte) s'il existe un $\mathcal{O}_Y$-Module
quasi-cohérent $\mathcal{E}$ et une $Y$-immersion $i$ de $X$ dans
$P=\mathbf{P}(\mathcal{E})$ tels que $\mathcal{L}$ soit isomorphe à
$i^*(\mathcal{O}_P(1))$.
\end{definition}

Il revient au même (\hyperref[II.4.2.3-fr]{4.2.3}) de dire qu'il existe un
$\mathcal{O}_Y$-Module quasi-cohérent $\mathcal{E}$ et un homomorphisme
surjectif $\varphi:q^*(\mathcal{E})\to\mathcal{L}$ tels que
$r_{\mathcal{L},\varphi}:X\to\mathbf{P}(\mathcal{E})$ soit une immersion.

On notera que l'existence d'un $\mathcal{O}_X$-Module très ample relativement
à $Y$ implique que $q$ est séparé
((\hyperref[II.3.1.3-fr]{3.1.3}) et
(\hyperref[I.5.5.1-fr]{I, 5.5.1, (i) et (ii)})).

\begin{corollary}[4.4.3]
\label{II.4.4.3-fr}
Supposons qu'il existe une $\mathcal{O}_Y$-Algèbre graduée quasi-cohérente
$\mathcal{S}$ engendrée par $\mathcal{S}_1$, et une $Y$-immersion
$i:X\to P=\operatorname{Proj}(\mathcal{S})$ telle que $\mathcal{L}$ soit
isomorphe à $i^*(\mathcal{O}_P(1))$~; alors $\mathcal{L}$ est très ample
relativement à $q$.
\end{corollary}

\begin{proof}
En effet, si $\mathcal{F}=\mathcal{S}_1$, l'homomorphisme canonique
$\mathbf{S}(\mathcal{F})\to\mathcal{S}$ est surjectif, donc, en composant
l'immersion fermée correspondante
$\operatorname{Proj}(\mathcal{S})\to\mathbf{P}(\mathcal{F})$
(\hyperref[II.3.6.2-fr]{3.6.2}) et l'immersion $i$, on obtient une immersion
$j:X\to\mathbf{P}(\mathcal{F})=P'$ telle que $\mathcal{L}$ soit isomorphe à
$j^*(\mathcal{O}_{P'}(1))$ (\hyperref[II.3.6.3-fr]{3.6.3}).
\end{proof}

\begin{proposition}[4.4.4]
\label{II.4.4.4-fr}
Soient $q:X\to Y$ un morphisme quasi-compact, $\mathcal{L}$ un
$\mathcal{O}_X$-Module inversible. Les propriétés suivantes sont
équivalentes~:
\begin{enumerate}
  \item[\rm{a)}] $\mathcal{L}$ est très ample relativement à $q$.
  \item[\rm{b)}] $q_*(\mathcal{L})$ est quasi-cohérent, l'homomorphisme
  canonique $\sigma:q^*(q_*(\mathcal{L}))\to\mathcal{L}$ est surjectif, et le
  morphisme
  $r_{\mathcal{L},\sigma}:X\to\mathbf{P}(q_*(\mathcal{L}))$ est une immersion.
\end{enumerate}
\end{proposition}

\begin{proof}
Comme $q$ est quasi-compact, on sait que $q_*(\mathcal{L})$ est quasi-cohérent
si $q$ est séparé (\hyperref[I.9.2.2-fr]{I, 9.2.2}).

\oldpage[II]{80}

On sait (\hyperref[II.3.4.7-fr]{3.4.7}) que l'existence d'un homomorphisme
surjectif $\varphi:q^*(\mathcal{E})\to\mathcal{L}$ ($\mathcal{E}$ étant un
$\mathcal{O}_Y$-Module quasi-cohérent) entraîne que $\sigma$ est surjectif~;
en outre, à la factorisation
$q^*(\mathcal{E})\to q^*(q_*(\mathcal{L}))\xrightarrow{\sigma}\mathcal{L}$
de $\varphi$ correspond canoniquement une factorisation
\[
  q^*(\mathbf{S}(\mathcal{E}))
    \longrightarrow q^*(\mathbf{S}(q_*(\mathcal{L})))
    \longrightarrow\bigoplus_{n\geq0}\mathcal{L}^{\otimes n}
\]
donc (\hyperref[II.3.8.3-fr]{3.8.3}) l'hypothèse que
$r_{\mathcal{L},\varphi}$ est une immersion entraîne qu'il en est de même de
$j=r_{\mathcal{L},\sigma}$~; en outre
(\hyperref[II.4.2.4-fr]{4.2.4}), $\mathcal{L}$ est isomorphe à
$j^*(\mathcal{O}_{P'}(1))$ en posant
$P'=\mathbf{P}(q_*(\mathcal{L}))$. On voit donc que a) et b) sont
équivalents.
\end{proof}

\begin{corollary}[4.4.5]
\label{II.4.4.5-fr}
Supposons $q$ quasi-compact. Pour que $\mathcal{L}$ soit très ample
relativement à $Y$, il faut et il suffit qu'il existe un recouvrement ouvert
$(U_\alpha)$ de $Y$ tel que
$\mathcal{L}|q^{-1}(U_\alpha)$ soit très ample relativement à $U_\alpha$ pour
tout $\alpha$.
\end{corollary}

\begin{proof}
En effet, la condition b) de (\hyperref[II.4.4.4-fr]{4.4.4}) est locale sur
$Y$.
\end{proof}

\begin{proposition}[4.4.6]
\label{II.4.4.6-fr}
Soient $Y$ un schéma quasi-compact, ou un préschéma dont l'espace sous-jacent
est noethérien, $q:X\to Y$ un morphisme de type fini, $\mathcal{L}$ un
$\mathcal{O}_X$-Module inversible. Les conditions a) et b) de
(\hyperref[II.4.4.4-fr]{4.4.4}) sont alors équivalentes aux suivantes~:
\begin{enumerate}
  \item[\rm{a')}] Il existe un $\mathcal{O}_Y$-Module quasi-cohérent de type
  fini $\mathcal{E}$ et un homomorphisme surjectif
  $\varphi:q^*(\mathcal{E})\to\mathcal{L}$ tel que
  $r_{\mathcal{L},\varphi}$ soit une immersion.
  \item[\rm{b')}] Il existe un sous-$\mathcal{O}_Y$-Module cohérent de type
  fini $\mathcal{E}$ de $q_*(\mathcal{L})$ ayant les propriétés énoncées dans
  a').
\end{enumerate}
\end{proposition}

\begin{proof}
Il est clair que a') ou b') implique a)~; d'autre part a) implique a') en
vertu de (\hyperref[II.4.4.1-fr]{4.4.1}) et de même b) implique b').
\end{proof}

\begin{corollary}[4.4.7]
\label{II.4.4.7-fr}
Supposons que $Y$ soit un schéma quasi-compact ou un préschéma noethérien. Si
$\mathcal{L}$ est très ample pour $q$, il existe une $\mathcal{O}_Y$-Algèbre
graduée quasi-cohérente $\mathcal{S}$, telle que $\mathcal{S}_1$ soit de type
fini et engendre $\mathcal{S}$, et une $Y$-immersion ouverte dominante
$i:X\to P=\operatorname{Proj}(\mathcal{S})$ telle que $\mathcal{L}$ soit
isomorphe à $i^*(\mathcal{O}_P(1))$.
\end{corollary}

\begin{proof}
En effet, la condition b') de (\hyperref[II.4.4.6-fr]{4.4.6}) est vérifiée~;
le morphisme structural $p:\mathbf{P}(\mathcal{E})=P'\to Y$ est alors séparé
et de type fini (\hyperref[II.3.1.3-fr]{3.1.3}), donc $P'$ est un schéma
quasi-compact (resp. un préschéma noethérien) si $Y$ est un schéma
quasi-compact (resp. un préschéma noethérien). Soit $Z$ l'adhérence
(\hyperref[I.9.5.11-fr]{I, 9.5.11}) du sous-préschéma $X'$ de $P'$ associé à
l'immersion $j=r_{\mathcal{L},\varphi}$ de $X$ dans $P'$~; il est clair que
$j$ se factorise en l'injection canonique $Z\to P'$ et en une immersion
ouverte dominante $i:X\to Z$. Mais $Z$ s'identifie à un préschéma
$\operatorname{Proj}(\mathcal{S})$, où $\mathcal{S}$ est une
$\mathcal{O}_Y$-Algèbre graduée quotient de $\mathbf{S}(\mathcal{E})$ par un
faisceau quasi-cohérent gradué d'idéaux
(\hyperref[II.3.6.2-fr]{3.6.2}), et il est clair que $\mathcal{S}_1$ est de
type fini et engendre $\mathcal{S}$~; en outre $\mathcal{O}_Z(1)$ est l'image
réciproque de $\mathcal{O}_{P'}(1)$ par l'injection canonique
(\hyperref[II.3.6.3-fr]{3.6.3}), donc
$\mathcal{L}=i^*(\mathcal{O}_Z(1))$.
\end{proof}

\begin{proposition}[4.4.8]
\label{II.4.4.8-fr}
Soient $q:X\to Y$ un morphisme, $\mathcal{L}$ un $\mathcal{O}_X$-Module très
ample relativement à $q$, $\mathcal{L}'$ un $\mathcal{O}_X$-Module inversible
tel qu'il existe un $\mathcal{O}_Y$-Module quasi-cohérent $\mathcal{E}'$ et un
homomorphisme surjectif $q^*(\mathcal{E}')\to\mathcal{L}'$. Alors
$\mathcal{L}\otimes_{\mathcal{O}_X}\mathcal{L}'$ est très ample relativement
à $q$.
\end{proposition}

\begin{proof}
En effet, l'hypothèse entraîne l'existence d'un $Y$-morphisme
$r':X\to P'=\mathbf{P}(\mathcal{E}')$ tel que
$\mathcal{L}'=r'^*(\mathcal{O}_{P'}(1))$
(\hyperref[II.4.2.1-fr]{4.2.1}). Il y a par hypothèse un
$\mathcal{O}_Y$-Module quasi-cohérent $\mathcal{E}$ et une

\oldpage[II]{81}

$Y$-immersion $r:X\to P=\mathbf{P}(\mathcal{E})$ telle que
$\mathcal{L}=r^*(\mathcal{O}_P(1))$. Posons
$Q=\mathbf{P}(\mathcal{E}\otimes\mathcal{E}')$ et considérons le morphisme de
Segre $\varsigma:P\times_YP'\to Q$ (\hyperref[II.4.3.1-fr]{4.3.1}). Comme
$r$ est une immersion, il en est de même de
$(r,r')_Y:X\to P\times_YP'$ (\hyperref[I.5.3.14-fr]{I, 5.3.14})~; mais
$\varsigma$ étant une immersion (\hyperref[II.4.3.3-fr]{4.3.3}), il en est de
même de
$r'':X\xrightarrow{(r,r')_Y}P\times_YP'\xrightarrow{\varsigma}Q$. D'autre
part (\hyperref[II.4.3.2.1-fr]{4.3.2.1}),
$\varsigma^*(\mathcal{O}_Q(1))$ est isomorphe à
$\mathcal{O}_P(1)\otimes_Y\mathcal{O}_{P'}(1)$, donc
(\hyperref[I.9.1.4-fr]{I, 9.1.4}), $r''^*(\mathcal{O}_Q(1))$ est isomorphe à
$\mathcal{L}\otimes\mathcal{L}'$, ce qui démontre la proposition.
\end{proof}

\begin{corollary}[4.4.9]
\label{II.4.4.9-fr}
Soit $q:X\to Y$ un morphisme.
\begin{enumerate}
  \item[\rm{(i)}] Soient $\mathcal{L}$ un $\mathcal{O}_X$-Module inversible,
  $\mathcal{K}$ un $\mathcal{O}_Y$-Module inversible. Pour que $\mathcal{L}$
  soit très ample relativement à $q$, il faut et il suffit que
  $\mathcal{L}\otimes q^*(\mathcal{K})$ le soit.
  \item[\rm{(ii)}] Si $\mathcal{L}$ et $\mathcal{L}'$ sont deux
  $\mathcal{O}_X$-Modules très amples relativement à $q$, il en est de même
  de $\mathcal{L}\otimes\mathcal{L}'$~; en particulier
  $\mathcal{L}^{\otimes n}$ est très ample relativement à $q$ pour tout
  $n>0$.
\end{enumerate}
\end{corollary}

\begin{proof}
(ii) est conséquence immédiate de (\hyperref[II.4.4.8-fr]{4.4.8}), ainsi que
la nécessité de la condition (i)~; d'autre part, si
$\mathcal{L}\otimes q^*(\mathcal{K})$ est très ample, il en est de même de
$(\mathcal{L}\otimes q^*(\mathcal{K}))\otimes q^*(\mathcal{K}^{-1})$
d'après ce qui précède, et ce dernier $\mathcal{O}_X$-Module est isomorphe à
$\mathcal{L}$ (\hyperref[0.5.4.3-fr]{0, 5.4.3} et
\hyperref[0.5.4.5-fr]{5.4.5}).
\end{proof}

\begin{proposition}[4.4.10]
\label{II.4.4.10-fr}
\begin{enumerate}
  \item[\rm{(i)}] Pour tout préschéma $Y$, tout $\mathcal{O}_Y$-Module
  inversible $\mathcal{L}$ est très ample relativement au morphisme identique
  $1_Y$.
  \item[\rm{(i bis)}] Soient $f:X\to Y$ un morphisme, $j:X'\to X$ une
  immersion. Si $\mathcal{L}$ est un $\mathcal{O}_X$-Module très ample
  relativement à $f$, $j^*(\mathcal{L})$ est très ample relativement à
  $f\circ j$.
  \item[\rm{(ii)}] Soient $Z$ un préschéma quasi-compact, $f:X\to Y$ un
  morphisme de type fini, $g:Y\to Z$ un morphisme quasi-compact,
  $\mathcal{L}$ un $\mathcal{O}_X$-Module très ample relativement à $f$,
  $\mathcal{K}$ un $\mathcal{O}_Y$-Module très ample relativement à $g$.
  Alors il existe un entier $n_0>0$ tel que
  $\mathcal{L}\otimes f^*(\mathcal{K}^{\otimes n})$ soit très ample
  relativement à $g\circ f$ pour tout $n\geq n_0$.
  \item[\rm{(iii)}] Soient $f:X\to Y$, $g:Y'\to Y$ deux morphismes, et posons
  $X'=X_{(Y')}$. Si $\mathcal{L}$ est un $\mathcal{O}_X$-Module très ample
  relativement à $f$, $\mathcal{L}'=\mathcal{L}\otimes_Y\mathcal{O}_{Y'}$
  est un $\mathcal{O}_{X'}$-Module très ample relativement à $f_{(Y')}$.
  \item[\rm{(iv)}] Soient $f_i:X_i\to Y_i$ ($i=1,2$) deux $S$-morphismes.
  Si $\mathcal{L}_i$ est un $\mathcal{O}_{X_i}$-Module très ample relativement
  à $f_i$ ($i=1,2$), $\mathcal{L}_1\otimes_S\mathcal{L}_2$ est très ample
  relativement à $f_1\times_Sf_2$.
  \item[\rm{(v)}] Soient $f:X\to Y$, $g:Y\to Z$ deux morphismes. Si un
  $\mathcal{O}_X$-Module $\mathcal{L}$ est très ample relativement à
  $g\circ f$, alors $\mathcal{L}$ est très ample relativement à $f$.
  \item[\rm{(vi)}] Soient $f:X\to Y$ un morphisme, $j$ l'injection canonique
  $X_{\mathrm{red}}\to X$. Si $\mathcal{L}$ est un $\mathcal{O}_X$-Module
  très ample relativement à $f$, alors $j^*(\mathcal{L})$ est très ample
  relativement à $f_{\mathrm{red}}$.
\end{enumerate}
\end{proposition}

\begin{proof}
La propriété (i bis) découle aussitôt de la définition
(\hyperref[II.4.4.2-fr]{4.4.2}) et il est immédiat que (vi) se déduit
formellement de (i bis) et de (v) par un raisonnement calqué sur celui de
(\hyperref[I.5.5.12-fr]{I, 5.5.12}), que nous laissons au lecteur. Pour
démontrer (v), considérons, comme dans
(\hyperref[I.5.5.12-fr]{I, 5.5.12}), la factorisation
$X\xrightarrow{\Gamma_f}X\times_ZY\xrightarrow{p_2}Y$, en notant que
$p_2=(g\circ f)\times1_Y$. Il résulte de l'hypothèse et de (i) et (iv) que
$\mathcal{L}\otimes_{\mathcal{O}_Z}\mathcal{O}_Y$ est très ample pour $p_2$~;
d'autre part, on a
$\mathcal{L}=\Gamma_f^*(\mathcal{L}\otimes_{\mathcal{O}_Z}\mathcal{O}_Y)$
(\hyperref[I.9.1.4-fr]{I, 9.1.4}), et $\Gamma_f$ est une immersion
(\hyperref[I.5.3.11-fr]{I, 5.3.11})~; on peut donc appliquer (i bis).

\oldpage[II]{82}

Pour prouver (i), on applique la définition
(\hyperref[II.4.4.2-fr]{4.4.2}) avec $\mathcal{E}=\mathcal{L}$, en notant
qu'alors $\mathbf{P}(\mathcal{E})$ s'identifie à $Y$
(\hyperref[II.4.1.4-fr]{4.1.4}).

Démontrons (iii). Il existe un $\mathcal{O}_Y$-Module quasi-cohérent
$\mathcal{E}$ et une $Y$-immersion $i:X\to\mathbf{P}(\mathcal{E})=P$ telle
que $\mathcal{L}=i^*(\mathcal{O}_P(1))$~; alors, si on pose
$\mathcal{E}'=g^*(\mathcal{E})$, $\mathcal{E}'$ est un
$\mathcal{O}_{Y'}$-Module quasi-cohérent et on a
$P'=\mathbf{P}(\mathcal{E}')=P_{(Y')}$
(\hyperref[II.4.1.3.1-fr]{4.1.3.1}), $i_{(Y')}$ est une immersion de
$X_{(Y')}$ dans $P'$ (\hyperref[I.4.3.2-fr]{I, 4.3.2}) et $\mathcal{L}'$ est
isomorphe à $(i_{(Y')})^*(\mathcal{O}_{P'}(1))$
(\hyperref[II.4.2.10-fr]{4.2.10}).

Pour prouver (iv), remarquons qu'il y a par hypothèse une $Y_i$-immersion
$r_i:X_i\to P_i=\mathbf{P}(\mathcal{E}_i)$, où $\mathcal{E}_i$ est un
$\mathcal{O}_{Y_i}$-Module quasi-cohérent, et
$\mathcal{L}_i=r_i^*(\mathcal{O}_{P_i}(1))$ ($i=1,2$)~;
$r_1\times_Sr_2$ est une $S$-immersion de $X_1\times_SX_2$ dans
$P_1\times_SP_2$ (\hyperref[I.4.3.1-fr]{I, 4.3.1}) et l'image réciproque de
$\mathcal{O}_{P_1}(1)\otimes_S\mathcal{O}_{P_2}(1)$ par cette immersion est
$\mathcal{L}_1\otimes_S\mathcal{L}_2$. D'autre part, posons
$T=Y_1\times_SY_2$, et soient $p_1,p_2$ les projections de $T$ dans $Y_1,Y_2$
respectivement. Si on pose
$P_i'=\mathbf{P}(p_i^*(\mathcal{E}_i))$ ($i=1,2$), on a par
(\hyperref[II.4.1.3.1-fr]{4.1.3.1}), $P_i'=P_i\times_{Y_i}T$ d'où
\[
  P_1'\times_TP_2'
  =(P_1\times_{Y_1}T)\times_T(P_2\times_{Y_2}T)
  =P_1\times_{Y_1}(T\times_{Y_2}P_2)
  =P_1\times_{Y_1}(Y_1\times_SP_2)
  =P_1\times_SP_2
\]
à un isomorphisme canonique près. De même, on a
$\mathcal{O}_{P_i'}(1)=\mathcal{O}_{P_i}(1)\otimes_{Y_i}\mathcal{O}_T$
(\hyperref[II.4.1.3.2-fr]{4.1.3.2}), et un calcul analogue (basé notamment
sur (\hyperref[I.9.1.9.1-fr]{I, 9.1.9.1} et
\hyperref[I.9.1.2-fr]{9.1.2})) montre que dans l'identification précédente,
$\mathcal{O}_{P_1'}(1)\otimes_T\mathcal{O}_{P_2'}(1)$ s'identifie à
$\mathcal{O}_{P_1}(1)\otimes_S\mathcal{O}_{P_2}(1)$. On peut donc considérer
$r_1\times_Sr_2$ comme une $T$-immersion de $X_1\times_SX_2$ dans
$P_1'\times_TP_2'$, l'image réciproque de
$\mathcal{O}_{P_1'}(1)\otimes_T\mathcal{O}_{P_2'}(1)$ par cette immersion
étant $\mathcal{L}_1\otimes_S\mathcal{L}_2$. On termine alors le raisonnement
comme dans (\hyperref[II.4.4.8-fr]{4.4.8}) en utilisant le morphisme de Segre.

Il reste à démontrer (ii). On peut d'abord se restreindre au cas où $Z$ est un
schéma affine, car en général il y a un recouvrement fini $(U_i)$ de $Z$ par
des ouverts affines~; si la proposition est démontrée pour
$\mathcal{K}|g^{-1}(U_i)$, $\mathcal{L}|f^{-1}(g^{-1}(U_i))$ et un entier
$n_i$, il suffira de prendre pour $n_0$ le plus grand des $n_i$ pour démontrer
la proposition pour $\mathcal{K}$ et $\mathcal{L}$
(\hyperref[II.4.4.5-fr]{4.4.5}). L'hypothèse implique que $f$ et $g$ sont des
morphismes séparés, donc $X$ et $Y$ sont des schémas quasi-compacts.

Il y a une immersion $r:X\to P=\mathbf{P}(\mathcal{E})$, où $\mathcal{E}$ est
un $\mathcal{O}_Y$-Module quasi-cohérent de type fini et
$\mathcal{L}=r^*(\mathcal{O}_P(1))$, en vertu de
(\hyperref[II.4.4.6-fr]{4.4.6}). Nous allons voir qu'il existe un
$\mathcal{O}_P$-Module $\mathcal{M}$ très ample relativement au morphisme
composé $P\xrightarrow{h}Y\xrightarrow{g}Z$, tel que $\mathcal{O}_P(1)$ soit
isomorphe à $\mathcal{M}\otimes_Y\mathcal{K}^{\otimes(-m)}$ pour un entier
$m$. Pour $n\geq m+1$,
$\mathcal{O}_P(1)\otimes_Y\mathcal{K}^{\otimes n}$ sera donc très ample pour
$Z$ en vertu de l'hypothèse et de (iv) appliqué aux morphismes $h:P\to Y$ et
$1_Y$~; comme $r$ est une immersion et que
$\mathcal{L}\otimes f^*(\mathcal{K}^{\otimes n})
=r^*(\mathcal{O}_P(1)\otimes_Y\mathcal{K}^{\otimes n})$, la conclusion
résultera de (i bis). Pour établir notre assertion sur $\mathcal{O}_P(1)$,
nous utiliserons le lemme suivant~:

\begin{lemma}[4.4.10.1]
\label{II.4.4.10.1-fr}
Soient $Z$ un schéma quasi-compact ou un préschéma dont l'espace sous-jacent
est noethérien, $g:Y\to Z$ un morphisme quasi-compact, $\mathcal{K}$ un
$\mathcal{O}_Y$-Module inversible très ample pour $g$, $\mathcal{E}$ un
$\mathcal{O}_Y$-Module quasi-cohérent de type fini. Il existe alors un entier
$m_0$ tel que, pour tout $m\geq m_0$, $\mathcal{E}$ soit isomorphe à un
quotient d'un $\mathcal{O}_Y$-Module de la forme
$g^*(\mathcal{F})\otimes\mathcal{K}^{\otimes(-m)}$, où $\mathcal{F}$ est un
$\mathcal{O}_Z$-Module quasi-cohérent de type fini (dépendant de $m$).
\end{lemma}

Ce lemme sera démontré en (\hyperref[II.4.5.10.1-fr]{4.5.10.1})~; le lecteur
pourra vérifier que (\hyperref[II.4.4.10-fr]{4.4.10}) n'est utilisé nulle part
dans le n\textsuperscript{o}~4.5.

\oldpage[II]{83}

Ce lemme étant admis, il existe une immersion fermée $j_1$ de $P$ dans
\[
  P_1=\mathbf{P}(g^*(\mathcal{F})\otimes\mathcal{K}^{\otimes(-m)})
\]
telle que $\mathcal{O}_P(1)$ soit isomorphe à
$j_1^*(\mathcal{O}_{P_1}(1))$ (\hyperref[II.4.1.2-fr]{4.1.2}). D'autre part,
il existe un isomorphisme de $P_1$ sur
$P_2=\mathbf{P}(g^*(\mathcal{F}))$, identifiant
$\mathcal{O}_{P_1}(1)$ à
$\mathcal{O}_{P_2}(1)\otimes_Y\mathcal{K}^{\otimes(-m)}$
(\hyperref[II.4.1.4-fr]{4.1.4})~; on a donc une immersion fermée
$j_2:P\to P_2$ telle que $\mathcal{O}_P(1)$ soit isomorphe à
$j_2^*(\mathcal{O}_{P_2}(1))\otimes_Y\mathcal{K}^{\otimes(-m)}$. Enfin,
$P_2$ s'identifie à $P_3\times_ZY$, où $P_3=\mathbf{P}(\mathcal{F})$, et
$\mathcal{O}_{P_2}(1)$ à
$\mathcal{O}_{P_3}(1)\otimes_Z\mathcal{O}_Y$
(\hyperref[II.4.1.3-fr]{4.1.3}). Par définition,
$\mathcal{O}_{P_3}(1)$ est très ample pour $Z$~; comme il en est de même de
$\mathcal{K}$, on conclut de (iv) que
$\mathcal{O}_{P_2}(1)\otimes_Y\mathcal{K}$ est très ample pour $Z$~; il en
est de même de
$\mathcal{M}=j_2^*(\mathcal{O}_{P_2}(1)\otimes_Y\mathcal{K})$ en vertu de
(i bis), et $\mathcal{O}_P(1)$ est isomorphe à
$\mathcal{M}\otimes_Y\mathcal{K}^{\otimes(-m-1)}$, ce qui achève la
démonstration.
\end{proof}

\begin{proposition}[4.4.11]
\label{II.4.4.11-fr}
Soient $f:X\to Y$, $f':X'\to Y$ deux morphismes, $X''$ le préschéma somme
$X\sqcup X'$, $f''$ le morphisme $X''\to Y$ qui coïncide avec $f$ (resp.
$f'$) dans $X$ (resp. $X'$). Soit $\mathcal{L}$ (resp. $\mathcal{L}'$) un
$\mathcal{O}_X$-Module (resp. un $\mathcal{O}_{X'}$-Module) inversible, et
soit $\mathcal{L}''$ le $\mathcal{O}_{X''}$-Module inversible qui coïncide
avec $\mathcal{L}$ dans $X$ et avec $\mathcal{L}'$ dans $X'$. Pour que
$\mathcal{L}''$ soit très ample relativement à $f''$, il faut et il suffit
que $\mathcal{L}$ soit très ample relativement à $f$ et $\mathcal{L}'$ très
ample relativement à $f'$.
\end{proposition}

\begin{proof}
On est aussitôt ramené au cas où $Y$ est affine. Si $\mathcal{L}''$ est très
ample il en est de même de $\mathcal{L}$ et $\mathcal{L}'$ en vertu de
(\hyperref[II.4.4.10-fr]{4.4.10, (i bis)}). Inversement, si $\mathcal{L}$ et
$\mathcal{L}'$ sont très amples, il résulte aussitôt de la définition
(\hyperref[II.4.4.2-fr]{4.4.2}) et de
(\hyperref[II.4.3.6-fr]{4.3.6}) que $\mathcal{L}''$ est très ample.
\end{proof}

\subsection{Faisceaux amples}
\label{subsection:II.4.5-fr}

\begin{env}[4.5.1]
\label{II.4.5.1-fr}
Étant donnés un préschéma $X$ et un $\mathcal{O}_X$-Module inversible
$\mathcal{L}$, nous poserons, pour tout $\mathcal{O}_X$-Module $\mathcal{F}$
(lorsque aucune confusion ne sera possible sur $\mathcal{L}$),
$\mathcal{F}(n)=\mathcal{F}\otimes_{\mathcal{O}_X}\mathcal{L}^{\otimes n}$
($n\in\mathbf{Z}$)~; nous poserons d'autre part
\[
  S=\bigoplus_{n\geq0}\Gamma(X,\mathcal{L}^{\otimes n})
\]
(sous-anneau gradué de l'anneau $\Gamma_\bullet(\mathcal{L})$ défini dans
(\hyperref[0.5.4.6-fr]{0, 5.4.6})). Si on considère $X$ comme un
$\mathbf{Z}$-préschéma et qu'on désigne par $p$ le morphisme structural
$X\to\operatorname{Spec}(\mathbf{Z})$, il y a correspondance biunivoque
entre les homomorphismes
\[
  p^*(\widetilde{S})\longrightarrow
  \bigoplus_{n\geq0}\mathcal{L}^{\otimes n}
\]
de $\mathcal{O}_X$-Algèbres graduées et les endomorphismes de l'anneau gradué
$S$ (\hyperref[I.2.2.5-fr]{I, 2.2.5})~; l'homomorphisme
\[
  \varepsilon:p^*(\widetilde{S})\longrightarrow
  \bigoplus_{n\geq0}\mathcal{L}^{\otimes n}
\]
qui correspond ainsi à l'automorphisme identique de $S$ sera dit canonique.
Il lui correspond (\hyperref[II.3.7.1-fr]{3.7.1}) un morphisme
$G(\varepsilon)\to\operatorname{Proj}(S)$ qui sera aussi dit canonique.
\end{env}

\begin{theorem}[4.5.2]
\label{II.4.5.2-fr}
Soient $X$ un schéma quasi-compact ou un préschéma dont l'espace sous-jacent
est noethérien, $\mathcal{L}$ un $\mathcal{O}_X$-Module inversible, $S$
l'anneau gradué
$\bigoplus_{n\geq0}\Gamma(X,\mathcal{L}^{\otimes n})$. Les conditions
suivantes sont équivalentes~:
\begin{enumerate}
  \item[\rm{a)}] Lorsque $f$ parcourt l'ensemble des éléments homogènes de
  $S_+$, les $X_f$ forment une base de la topologie de $X$.
  \item[\rm{a')}] Lorsque $f$ parcourt l'ensemble des éléments homogènes de
  $S_+$, ceux des $X_f$ qui sont affines forment un recouvrement de $X$.
  \item[\rm{b)}] Le morphisme canonique
  $G(\varepsilon)\to\operatorname{Proj}(S)$
  (\hyperref[II.4.5.1-fr]{4.5.1}) est partout défini et est une immersion
  ouverte dominante.

\oldpage[II]{84}

  \item[\rm{b')}] Le morphisme canonique
  $G(\varepsilon)\to\operatorname{Proj}(S)$ est partout défini et est un
  homéomorphisme de l'espace sous-jacent à $X$ sur un sous-espace de
  $\operatorname{Proj}(S)$.
  \item[\rm{c)}] Pour tout $\mathcal{O}_X$-Module quasi-cohérent
  $\mathcal{F}$, si on désigne par $\mathcal{F}_n$ le
  sous-$\mathcal{O}_X$-Module de $\mathcal{F}(n)$ engendré par les sections de
  $\mathcal{F}(n)$ au-dessus de $X$, alors $\mathcal{F}$ est la somme des
  sous-$\mathcal{O}_X$-Modules $\mathcal{F}_n(-n)$, pour les entiers $n>0$.
  \item[\rm{c')}] La propriété c) a lieu pour tout faisceau quasi-cohérent
  d'idéaux de $\mathcal{O}_X$.
\end{enumerate}
En outre, si $(f_\alpha)$ est une famille d'éléments homogènes de $S_+$,
telle que les $X_{f_\alpha}$ soient affines, alors la restriction à
$\bigcup_\alpha X_{f_\alpha}$ du morphisme canonique
$X\to\operatorname{Proj}(S)$ est un isomorphisme de
$\bigcup_\alpha X_{f_\alpha}$ sur
$\bigcup_\alpha(\operatorname{Proj}(S))_{f_\alpha}$.
\end{theorem}

\begin{proof}
Il est clair que b) implique b'), et b') implique a) en vertu de la formule
(\hyperref[II.3.7.3.1-fr]{3.7.3.1}) (tenant compte de ce que
$\varepsilon^\flat$ est l'identité). La condition a) implique a'), car tout
$x\in X$ a un voisinage affine $U$ tel que
$\mathcal{L}|U$ soit isomorphe à $\mathcal{O}_X|U$~; si
$f\in\Gamma(X,\mathcal{L}^{\otimes n})$ est telle que $x\in X_f\subset U$,
$X_f$ est aussi l'ensemble des $x'\in U$ tels que $(f|U)(x')\neq0$, et c'est
par suite un ouvert affine (\hyperref[I.1.3.6-fr]{I, 1.3.6}). Pour prouver
que a') entraîne b), il suffit de prouver la dernière assertion de l'énoncé,
et de vérifier de plus que si $X=\bigcup_\alpha X_{f_\alpha}$, la condition
(iv) de (\hyperref[II.3.8.2-fr]{3.8.2}) est satisfaite. Ce dernier point
résulte aussitôt de (\hyperref[I.9.3.1-fr]{I, 9.3.1, (i)}). Quant à la
dernière assertion de (\hyperref[II.4.5.2-fr]{4.5.2}), comme
$X_{f_\alpha}$ est l'image réciproque de
$(\operatorname{Proj}(S))_{f_\alpha}$ par
$G(\varepsilon)\to\operatorname{Proj}(S)$, il suffit d'appliquer
(\hyperref[I.9.3.2-fr]{I, 9.3.2}). Donc a), a'), b), b') sont équivalentes.

Pour montrer que a') entraîne c), notons que si $X_f$ est affine (avec
$f\in S_k$), $\mathcal{F}|X_f$ est engendré par ses sections au-dessus de
$X_f$ (\hyperref[I.1.3.9-fr]{I, 1.3.9})~; d'autre part
(\hyperref[I.9.3.1-fr]{I, 9.3.1, (ii)}) une telle section $s$ est de la forme
$(t|X_f)\otimes(f|X_f)^{-m}$ où
$t\in\Gamma(X,\mathcal{F}(km))$~; par définition, $t$ est aussi une section
de $\mathcal{F}_{km}$, donc $s$ est bien une section de
$\mathcal{F}_{km}(-km)$ au-dessus de $X_f$, ce qui prouve c). Il est clair
que c) implique c'), et il reste à montrer que c') entraîne a). Or, soit $U$
un voisinage ouvert de $x\in X$, et soit $\mathcal{J}$ un faisceau d'idéaux
quasi-cohérent de $\mathcal{O}_X$ définissant un sous-préschéma fermé de $X$
ayant $X-U$ comme espace sous-jacent
(\hyperref[I.5.2.1-fr]{I, 5.2.1}). L'hypothèse c') entraîne qu'il existe un
entier $n>0$ et une section $f$ de $\mathcal{J}(n)$ au-dessus de $X$ telle
que $f(x)\neq0$. Mais on a évidemment $f\in S_n$, et $x\in X_f\subset U$, ce
qui prouve a).
\end{proof}

Lorsque $X$ est un préschéma dont l'espace sous-jacent est noethérien, les
conditions équivalentes de (\hyperref[II.4.5.2-fr]{4.5.2}) impliquent que
$X$ est un schéma, puisqu'il est isomorphe à un sous-préschéma du schéma
$S=\operatorname{Proj}(A)$ en vertu de
(\hyperref[II.4.5.2-fr]{4.5.2, b)}).

\begin{definition}[4.5.3]
\label{II.4.5.3-fr}
On dit qu'un $\mathcal{O}_X$-Module inversible $\mathcal{L}$ est ample si
$X$ est un schéma quasi-compact et si les conditions équivalentes de
(\hyperref[II.4.5.2-fr]{4.5.2}) sont vérifiées.
\end{definition}

Il résulte évidemment du critère (\hyperref[II.4.5.2-fr]{4.5.2, a)}) que si
$\mathcal{L}$ est un $\mathcal{O}_X$-Module ample, alors, pour tout ouvert
$U$ de $X$, $\mathcal{L}|U$ est un $(\mathcal{O}_X|U)$-Module ample.

Il résulte de la démonstration de (\hyperref[II.4.5.2-fr]{4.5.2}) que les
$X_f$ affines forment déjà une base de la topologie de $X$. De plus~:

\begin{corollary}[4.5.4]
\label{II.4.5.4-fr}
Soit $\mathcal{L}$ un $\mathcal{O}_X$-Module ample. Pour tout sous-espace fini
$Z$ de $X$ et tout voisinage $U$ de $Z$, il existe un entier $n$ et un
$f\in\Gamma(X,\mathcal{L}^{\otimes n})$ tels que $X_f$ soit un voisinage
affine de $Z$ contenu dans $U$.
\end{corollary}

\begin{proof}
\oldpage[II]{85}

En vertu de (\hyperref[II.4.5.2-fr]{4.5.2, b)}), on peut se borner à prouver
que pour toute partie finie $Z'$ de $\operatorname{Proj}(S)$ et tout
voisinage ouvert $U$ de $Z'$, il existe un élément homogène $f\in S_+$ tel
que $Z'\subset(\operatorname{Proj}(S))_f\subset U$
(\hyperref[II.2.4.1-fr]{2.4.1}). Or, par définition, l'ensemble fermé $Y$,
complémentaire de $U$ dans $\operatorname{Proj}(S)$, est de la forme
$V_+(\mathfrak{I})$, où $\mathfrak{I}$ est un idéal gradué de $S$, ne
contenant pas $S_+$ (\hyperref[II.2.3.2-fr]{2.3.2})~; d'autre part, les
points de $Z'$ sont par définition des idéaux premiers gradués
$\mathfrak{p}_i$ de $S_+$ ne contenant pas $\mathfrak{I}$
(\hyperref[II.2.3.1-fr]{2.3.1}). Il existe donc un élément
$f\in\mathfrak{I}$ n'appartenant à aucun des $\mathfrak{p}_i$ (Bourbaki,
\emph{Alg. comm.}, chap.~II, \S~1, n\textsuperscript{o}~1, prop.~2), et comme
les $\mathfrak{p}_i$ sont gradués, le raisonnement fait \emph{loc. cit.}
montre qu'on peut même supposer $f$ homogène~; cet élément répond alors à la
question.
\end{proof}

\begin{proposition}[4.5.5]
\label{II.4.5.5-fr}
Supposons que $X$ soit un schéma quasi-compact, ou un préschéma dont l'espace
sous-jacent est noethérien. Les conditions a) à c') de
(\hyperref[II.4.5.2-fr]{4.5.2}) sont aussi équivalentes aux suivantes~:
\begin{enumerate}
  \item[\rm{d)}] Pour tout $\mathcal{O}_X$-Module quasi-cohérent
  $\mathcal{F}$ de type fini, il existe un entier $n_0$ tel que, pour tout
  $n\geq n_0$, $\mathcal{F}(n)$ soit engendré par ses sections au-dessus de
  $X$.
  \item[\rm{d')}] Pour tout $\mathcal{O}_X$-Module quasi-cohérent
  $\mathcal{F}$ de type fini, il existe deux entiers $n>0$, $k>0$ tels que
  $\mathcal{F}$ soit isomorphe à un quotient du $\mathcal{O}_X$-Module
  $\mathcal{L}^{\otimes(-n)}\otimes\mathcal{O}_X^k$.
  \item[\rm{d'')}] La propriété d') a lieu pour tout faisceau d'idéaux
  quasi-cohérent de type fini de $\mathcal{O}_X$.
\end{enumerate}
\end{proposition}

\begin{proof}
Comme $X$ est quasi-compact, si un $\mathcal{O}_X$-Module quasi-cohérent de
type fini $\mathcal{F}$ est tel que $\mathcal{F}(n)$ (qui est de type fini)
est engendré par ses sections au-dessus de $X$, $\mathcal{F}(n)$ est alors
engendré par un nombre fini de ces sections
(\hyperref[0.5.2.3-fr]{0, 5.2.3}), donc d) implique d') et il est clair que
d') implique d''). Comme tout $\mathcal{O}_X$-Module quasi-cohérent
$\mathcal{G}$ est limite inductive de ses sous-$\mathcal{O}_X$-Modules de
type fini (\hyperref[I.9.4.9-fr]{I, 9.4.9}), pour vérifier la condition c')
de (\hyperref[II.4.5.2-fr]{4.5.2}), il suffit de le faire pour un faisceau
quasi-cohérent d'idéaux de $\mathcal{O}_X$ qui est de type fini, et d'')
entraîne donc c'). Reste à voir que si $\mathcal{L}$ est ample la propriété
d) est vérifiée. Considérons un recouvrement fini de $X$ par des
$X_{f_i}$ ($f_i\in S_{n_i}$) qu'on peut supposer affines~; en remplaçant les
$f_i$ par des puissances convenables (ce qui ne change pas les $X_{f_i}$),
on peut supposer tous les $n_i$ égaux à un même entier $m$. Le faisceau
$\mathcal{F}|X_{f_i}$, étant de type fini par hypothèse, est engendré par un
nombre fini de ses sections $h_{ij}$ au-dessus de $X_{f_i}$
(\hyperref[I.1.3.13-fr]{I, 1.3.13})~; il existe donc un entier $k_0$ tel que
la section $h_{ij}\otimes f_i^{\otimes k_0}$ se prolonge en une section de
$\mathcal{F}(k_0m)$ au-dessus de $X$ pour tout couple $(i,j)$
(\hyperref[I.9.3.1-fr]{I, 9.3.1}). \emph{A fortiori} les
$h_{ij}\otimes f_i^{\otimes k}$ se prolongent en sections de
$\mathcal{F}(km)$ au-dessus de $X$ pour tout $k\geq k_0$, et pour ces valeurs
de $k$, $\mathcal{F}(km)$ est donc engendré par ses sections au-dessus de
$X$. Pour tout $p$ tel que $0<p<m$, $\mathcal{F}(p)$ est aussi de type fini,
donc il y a un entier $k_p$ tel que
$\mathcal{F}(p)(km)=\mathcal{F}(p+km)$ soit engendré par ses sections
au-dessus de $X$ pour $k\geq k_p$. Prenant $n_0$ plus grand que tous les
$k_pm$, on en conclut que $\mathcal{F}(n)$ est engendré par ses sections
au-dessus de $X$ pour tout $n\geq n_0$, car un tel $n$ s'écrit $n=km+p$ avec
$k\geq k_p$ et $0\leq p<m$.
\end{proof}

\begin{proposition}[4.5.6]
\label{II.4.5.6-fr}
Soient $X$ un schéma quasi-compact, $\mathcal{L}$ un $\mathcal{O}_X$-Module
inversible.
\begin{enumerate}
  \item[\rm{(i)}] Soit $n$ un entier $>0$. Pour que $\mathcal{L}$ soit
  ample, il faut et il suffit que $\mathcal{L}^{\otimes n}$ soit ample.
  \item[\rm{(ii)}] Soit $\mathcal{L}'$ un $\mathcal{O}_X$-Module inversible
  tel que, pour tout $x\in X$, il existe un entier $n>0$

\oldpage[II]{86}

  et une section $s'$ de $\mathcal{L}'^{\otimes n}$ au-dessus de $X$ tels
  que $s'(x)\neq0$. Alors, si $\mathcal{L}$ est ample, il en est de même de
  $\mathcal{L}\otimes\mathcal{L}'$.
\end{enumerate}
\end{proposition}

\begin{proof}
La propriété (i) est conséquence évidente du critère a) de
(\hyperref[II.4.5.2-fr]{4.5.2}) puisque $X_{f^{\otimes n}}=X_f$. D'autre
part, si $\mathcal{L}$ est ample, pour tout $x\in X$ et tout voisinage $U$ de
$x$, il existe $m>0$ et $f\in\Gamma(X,\mathcal{L}^{\otimes m})$ tel que
$x\in X_f\subset U$ (\hyperref[II.4.5.2-fr]{4.5.2, a)})~; si
$f'\in\Gamma(X,\mathcal{L}'^{\otimes n})$ est tel que $f'(x)\neq0$, on aura
$s(x)\neq0$ pour
$s=f^{\otimes n}\otimes f'^{\otimes m}\in
\Gamma(X,(\mathcal{L}\otimes\mathcal{L}')^{\otimes mn})$, donc
$x\in X_s\subset X_f\subset U$, ce qui prouve que
$\mathcal{L}\otimes\mathcal{L}'$ est ample
(\hyperref[II.4.5.2-fr]{4.5.2, a)}).
\end{proof}

\begin{corollary}[4.5.7]
\label{II.4.5.7-fr}
Le produit tensoriel de deux $\mathcal{O}_X$-Modules amples est ample.
\end{corollary}

\begin{corollary}[4.5.8]
\label{II.4.5.8-fr}
Soient $\mathcal{L}$ un $\mathcal{O}_X$-Module ample, $\mathcal{L}'$ un
$\mathcal{O}_X$-Module inversible~; il existe alors un entier $n_0>0$ tel
que $\mathcal{L}^{\otimes n}\otimes\mathcal{L}'$ soit ample et engendré par
ses sections au-dessus de $X$ pour $n\geq n_0$.
\end{corollary}

\begin{proof}
En effet, il résulte de (\hyperref[II.4.5.5-fr]{4.5.5}) qu'il existe un entier
$m_0$ tel que $\mathcal{L}^{\otimes m}\otimes\mathcal{L}'$ soit engendré par
ses sections au-dessus de $X$ pour $m\geq m_0$~; d'après
(\hyperref[II.4.5.6-fr]{4.5.6}) on peut donc prendre $n_0=m_0+1$.
\end{proof}

\begin{remark}[4.5.9]
\label{II.4.5.9-fr}
Soit $P=\mathrm{H}^1(X,\mathcal{O}_X^*)$ le groupe des classes de
$\mathcal{O}_X$-Modules inversibles (\hyperref[0.5.4.7-fr]{0, 5.4.7}), et
soit $P^+$ la partie de $P$ formée des classes de faisceaux amples. Supposons
que $P^+$ soit non vide. Alors il résulte de
(\hyperref[II.4.5.7-fr]{4.5.7}) et
(\hyperref[II.4.5.8-fr]{4.5.8}) que l'on a
\[
  P^++P^+\subset P^+
  \qquad\text{et}\qquad
  P^+-P^+=P
\]
autrement dit $P^+\cup\{0\}$ est l'ensemble des éléments positifs dans $P$
pour une structure de préordre sur $P$ compatible avec sa structure de
groupe, et qui est même archimédienne en vertu de
(\hyperref[II.4.5.8-fr]{4.5.8}). C'est pourquoi on dit parfois «~faisceau
positif~» au lieu de faisceau ample et «~faisceau négatif~» pour l'inverse
d'un tel faisceau (terminologie que nous ne suivrons pas).
\end{remark}

\begin{proposition}[4.5.10]
\label{II.4.5.10-fr}
Soient $Y$ un schéma affine, $q:X\to Y$ un morphisme séparé quasi-compact,
$\mathcal{L}$ un $\mathcal{O}_X$-Module inversible.
\begin{enumerate}
  \item[\rm{(i)}] Si $\mathcal{L}$ est très ample relativement à $q$, alors
  $\mathcal{L}$ est ample.
  \item[\rm{(ii)}] Supposons en outre que le morphisme $q$ soit de type fini.
  Alors, pour que $\mathcal{L}$ soit ample, il faut et il suffit qu'il possède
  l'une des propriétés équivalentes suivantes~:
  \begin{enumerate}
    \item[\rm{e)}] Il existe $n_0>0$ tel que pour tout entier $n\geq n_0$,
    $\mathcal{L}^{\otimes n}$ soit très ample relativement à $q$.
    \item[\rm{e')}] Il existe $n>0$ tel que $\mathcal{L}^{\otimes n}$ soit
    très ample relativement à $q$.
  \end{enumerate}
\end{enumerate}
\end{proposition}

\begin{proof}
La première assertion résulte de la définition
(\hyperref[II.4.4.2-fr]{4.4.2}) d'un $\mathcal{O}_X$-Module très ample~: si
$A$ est l'anneau de $Y$, il y a un $A$-module $E$ et un homomorphisme
surjectif
\[
  \psi:q^*((\mathbf{S}(E))^{\sim})\longrightarrow
  \bigoplus_{n\geq0}\mathcal{L}^{\otimes n}
\]
tel que $i=r_{\mathcal{L},\psi}$ soit une immersion partout définie
$X\to P=\mathbf{P}(\widetilde{E})$ et que l'on ait
$\mathcal{L}=i^*(\mathcal{O}_P(1))$~; comme les $D_+(f)$ pour $f$ homogène
dans $(\mathbf{S}(E))_+$ forment une base de la topologie de $P$, et que
$i^{-1}(D_+(f))=X_{\psi^\flat(f)}$ par
(\hyperref[II.3.7.3.1-fr]{3.7.3.1}), on voit que la condition a) de
(\hyperref[II.4.5.2-fr]{4.5.2}) est vérifiée, donc $\mathcal{L}$ est ample.

Prouvons maintenant que si $q$ est de type fini et si $\mathcal{L}$ est
ample, il vérifie la condition e). En premier lieu, il résulte du critère b)
de (\hyperref[II.4.5.2-fr]{4.5.2}) et de
(\hyperref[II.4.4.1-fr]{4.4.1, (i)}) qu'il existe

\oldpage[II]{87}

un entier $k_0>0$ tel que $\mathcal{L}^{\otimes k_0}$ soit très ample
relativement à $q$. D'autre part, en vertu de
(\hyperref[II.4.5.5-fr]{4.5.5}), il existe un entier $m_0$ tel que, pour
$m\geq m_0$, $\mathcal{L}^{\otimes m}$ soit engendré par ses sections
au-dessus de $X$. Posons $n_0=k_0+m_0$~; si $n\geq n_0$, on peut écrire
$n=k_0+m$ avec $m\geq m_0$, d'où
$\mathcal{L}^{\otimes n}
=\mathcal{L}^{\otimes k_0}\otimes\mathcal{L}^{\otimes m}$. Comme
$\mathcal{L}^{\otimes m}$ est engendré par ses sections au-dessus de $X$, il
résulte du critère (\hyperref[II.4.4.8-fr]{4.4.8}) et de
(\hyperref[II.3.4.7-fr]{3.4.7}) que $\mathcal{L}^{\otimes n}$ est très ample
relativement à $q$. Enfin, il est clair que e) entraîne e'), et e') implique
que $\mathcal{L}$ est ample en vertu de (i) et de
(\hyperref[II.4.5.6-fr]{4.5.6, (i)}).

\begin{env}[4.5.10.1]
\label{II.4.5.10.1-fr}
\emph{Démonstration du lemme (\hyperref[II.4.4.10.1-fr]{4.4.10.1}).} Posons
$\mathcal{E}(n)=\mathcal{E}\otimes\mathcal{K}^{\otimes n}$~; comme $g$ est
séparé (\hyperref[II.4.4.2-fr]{4.4.2}), le raisonnement de
(\hyperref[II.3.4.8-fr]{3.4.8}) s'applique et ramène à voir que
l'homomorphisme canonique
$g^*(g_*(\mathcal{E}(n)))\to\mathcal{E}(n)$ est surjectif pour $n$ assez
grand. En outre, comme $Z$ est quasi-compact, le raisonnement de
(\hyperref[II.3.4.6-fr]{3.4.6}) ramène à prouver la proposition dans le cas
où $Z$ est affine. Or, $\mathcal{K}$ est alors ample en vertu de
(\hyperref[II.4.5.10-fr]{4.5.10, (i)}), et la conclusion résulte de
(\hyperref[II.4.5.5-fr]{4.5.5, d')}).
\end{env}
\end{proof}

\begin{corollary}[4.5.11]
\label{II.4.5.11-fr}
Soient $Y$ un schéma affine, $q:X\to Y$ un morphisme séparé de type fini,
$\mathcal{L}$ un $\mathcal{O}_X$-Module ample, $\mathcal{L}'$ un
$\mathcal{O}_X$-Module inversible. Il existe un entier $n_0$ tel que, pour
$n\geq n_0$, $\mathcal{L}^{\otimes n}\otimes\mathcal{L}'$ soit très ample
relativement à $q$.
\end{corollary}

\begin{proof}
En effet, il existe $m_0$ tel que pour $m\geq m_0$,
$\mathcal{L}^{\otimes m}\otimes\mathcal{L}'$ soit engendré par ses sections
au-dessus de $X$ (\hyperref[II.4.5.8-fr]{4.5.8})~; d'autre part, il existe
$k_0$ tel que $\mathcal{L}^{\otimes k}$ soit très ample relativement à $q$
pour $k\geq k_0$. Par suite
$\mathcal{L}^{\otimes(k+m_0)}\otimes\mathcal{L}'$ est très ample pour
$k\geq k_0$ ((\hyperref[II.4.4.8-fr]{4.4.8}) et
(\hyperref[II.3.4.7-fr]{3.4.7})).
\end{proof}

\begin{remark}[4.5.12]
\label{II.4.5.12-fr}
On ignore si l'hypothèse qu'un $\mathcal{O}_X$-Module $\mathcal{L}$ est tel
que $\mathcal{L}^{\otimes n}$ soit très ample (relativement à $q$) entraîne
la même conclusion pour $\mathcal{L}^{\otimes(n+1)}$.
\end{remark}

\begin{proposition}[4.5.13]
\label{II.4.5.13-fr}
Soient $X$ un préschéma quasi-compact, $Z$ un sous-préschéma fermé de $X$
défini par un faisceau quasi-cohérent nilpotent $\mathcal{J}$ d'idéaux de
$\mathcal{O}_X$, $j$ l'injection canonique $Z\to X$. Pour qu'un
$\mathcal{O}_X$-Module inversible $\mathcal{L}$ soit ample, il faut et il
suffit que $\mathcal{L}'=j^*(\mathcal{L})$ soit un
$\mathcal{O}_Z$-Module ample.
\end{proposition}

\begin{proof}
La condition est nécessaire. En effet, pour toute section $f$ de
$\mathcal{L}^{\otimes n}$ au-dessus de $X$, soit $f'$ son image canonique
$f\otimes1$, section de
$\mathcal{L}'^{\otimes n}=\mathcal{L}^{\otimes n}
\otimes_{\mathcal{O}_X}(\mathcal{O}_X/\mathcal{J})$ au-dessus de l'espace
$Z$ (qui est identique à $X$)~; il est clair que $X_f=Z_{f'}$, donc le
critère a) de (\hyperref[II.4.5.2-fr]{4.5.2}) montre que $\mathcal{L}'$ est
ample.

Pour voir que la condition est suffisante, notons d'abord qu'on peut se
ramener au cas où $\mathcal{J}^2=0$, en considérant la suite (finie) des
préschémas
$X_k=(X,\mathcal{O}_X/\mathcal{J}^{k+1})$ dont chacun est sous-préschéma
fermé du suivant, défini par un faisceau d'idéaux de carré nul. D'ailleurs
$X$ est un schéma puisque $X_{\mathrm{red}}$ l'est par hypothèse
((\hyperref[II.4.5.3-fr]{4.5.3}) et
(\hyperref[I.5.5.1-fr]{I, 5.5.1})). Le critère a) de
(\hyperref[II.4.5.2-fr]{4.5.2}) montre qu'il suffira de prouver le

\begin{lemma}[4.5.13.1]
\label{II.4.5.13.1-fr}
Sous les hypothèses de (\hyperref[II.4.5.13-fr]{4.5.13}), supposons de plus
$\mathcal{J}$ de carré nul~; $\mathcal{L}$ étant un
$\mathcal{O}_X$-Module inversible, soit $g$ une section de
$\mathcal{L}'^{\otimes n}$ au-dessus de $Z$ telle que $Z_g$ soit affine.
Alors il existe un entier $m>0$ tel que $g^{\otimes m}$ soit l'image
canonique d'une section $f$ de $\mathcal{L}^{\otimes nm}$ au-dessus de $X$.
\end{lemma}

On a la suite exacte de $\mathcal{O}_X$-Modules
\[
  0\longrightarrow\mathcal{J}(n)\longrightarrow
  \mathcal{O}_X(n)=\mathcal{L}^{\otimes n}\longrightarrow
  \mathcal{O}_Z(n)=\mathcal{L}'^{\otimes n}\longrightarrow0
\]

\oldpage[II]{88}

puisque $\mathcal{F}(n)$ est un foncteur exact en $\mathcal{F}$~; d'où la
suite exacte de cohomologie
\[
  0\longrightarrow\Gamma(X,\mathcal{J}(n))\longrightarrow
  \Gamma(X,\mathcal{L}^{\otimes n})\longrightarrow
  \Gamma(X,\mathcal{L}'^{\otimes n})\xrightarrow{\partial}
  \mathrm{H}^1(X,\mathcal{J}(n))
\]
qui associe en particulier à $g$ un élément
$\partial g\in\mathrm{H}^1(X,\mathcal{J}(n))$.

Notons que puisque $\mathcal{J}^2=0$, $\mathcal{J}$ peut être considéré
comme un $\mathcal{O}_Z$-Module quasi-cohérent et l'on a, pour tout $k$,
$\mathcal{L}'^{\otimes k}\otimes_{\mathcal{O}_Z}\mathcal{J}(n)
=\mathcal{J}(n+k)$~; pour toute section
$s\in\Gamma(X,\mathcal{L}'^{\otimes k})$ la multiplication tensorielle par
$s$ est donc un homomorphisme $\mathcal{J}(n)\to\mathcal{J}(n+k)$ de
$\mathcal{O}_Z$-Modules, qui donne par suite un homomorphisme
$\mathrm{H}^i(X,\mathcal{J}(n))\to
\mathrm{H}^i(X,\mathcal{J}(n+k))$ de groupes de cohomologie.

Cela étant, nous allons voir que l'on a
\begin{equation}
\label{II.4.5.13.2-fr}
  g^{\otimes m}\otimes\partial g=0
\tag{4.5.13.2}
\end{equation}
pour un $m>0$ assez grand. En effet, $Z_g$ est un ouvert affine de $Z$ et
on a donc $\mathrm{H}^1(Z_g,\mathcal{J}(n))=0$ quand $\mathcal{J}(n)$ est
considéré comme un $\mathcal{O}_Z$-Module
(\hyperref[I.5.1.9.2-fr]{I, 5.1.9.2}). En particulier, si on pose
$g'=g|Z_g$, et si on considère son image par l'application
$\partial:\Gamma(Z_g,\mathcal{L}'^{\otimes n})\to
\mathrm{H}^1(Z_g,\mathcal{J}(n))$, on a $\partial g'=0$. Pour expliciter
cette relation, observons qu'en dimension $1$ la cohomologie d'un faisceau
de groupes abéliens est la même que sa cohomologie de Čech (G, II, 5.9)~;
pour former $\partial g$, il faut donc considérer un recouvrement ouvert
$(U_\alpha)$ assez fin de $X$, que l'on peut supposer fini et formé
d'ouverts affines, prendre pour chaque $\alpha$ une section
$g_\alpha\in\Gamma(U_\alpha,\mathcal{L}^{\otimes n})$ dont l'image
canonique dans $\Gamma(U_\alpha,\mathcal{L}'^{\otimes n})$ est
$g|U_\alpha$, et considérer la classe du cocycle
$(g_{\alpha|\beta}-g_{\beta|\alpha})$, $g_{\alpha|\beta}$ étant la
restriction de $g_\alpha$ à $U_\alpha\cap U_\beta$ (ce cocycle étant à
valeurs dans $\mathcal{J}(n)$). On peut en outre supposer que $\partial g'$
est calculé de la même manière à l'aide du recouvrement formé des
$U_\alpha\cap Z_g$ et des restrictions $g_\alpha|(U_\alpha\cap Z_g)$ (en
remplaçant au besoin $(U_\alpha)$ par un recouvrement plus fin)~; la
relation $\partial g'=0$ signifie alors qu'il existe pour chaque $\alpha$
une section $h_\alpha\in\Gamma(U_\alpha\cap Z_g,\mathcal{J}(n))$ telle que
\[
  (g_{\alpha|\beta}-g_{\beta|\alpha})
  |(U_\alpha\cap U_\beta\cap Z_g)
  =h_{\alpha|\beta}-h_{\beta|\alpha}
\]
en désignant par $h_{\alpha|\beta}$ la restriction de $h_\alpha$ à
$U_\alpha\cap U_\beta\cap Z_g$ (G, II, 5.11). Cela étant, il y a un entier
$m>0$ tel que $g^{\otimes m}\otimes h_\alpha$ soit la restriction à
$U_\alpha\cap Z_g$ d'une section
$t_\alpha\in\Gamma(U_\alpha,\mathcal{J}(n+nm))$ pour tout $\alpha$
(\hyperref[I.9.3.1-fr]{I, 9.3.1})~; on a donc
\[
  g^{\otimes m}\otimes(g_{\alpha|\beta}-g_{\beta|\alpha})
  =t_{\alpha|\beta}-t_{\beta|\alpha}
\]
pour tout couple d'indices, ce qui prouve
(\hyperref[II.4.5.13.2-fr]{4.5.13.2}).

Remarquons d'autre part que si
$s\in\Gamma(X,\mathcal{O}_Z(p))$,
$t\in\Gamma(X,\mathcal{O}_Z(q))$, on a, dans le groupe
$\mathrm{H}^1(X,\mathcal{J}(p+q))$
\begin{equation}
\label{II.4.5.13.3-fr}
  \partial(s\otimes t)=(\partial s)\otimes t+s\otimes(\partial t).
\tag{4.5.13.3}
\end{equation}

En effet, on peut encore, pour calculer les deux membres, considérer un
recouvrement ouvert $(U_\alpha)$ de $X$, pour chaque $\alpha$ une section
$s_\alpha\in\Gamma(U_\alpha,\mathcal{O}_X(p))$ (resp.
$t_\alpha\in\Gamma(U_\alpha,\mathcal{O}_X(q))$) dont l'image canonique dans
$\Gamma(U_\alpha,\mathcal{O}_Z(p))$ (resp.
$\Gamma(U_\alpha,\mathcal{O}_Z(q))$ soit $s|U_\alpha$ (resp.
$t|U_\alpha$)~; la relation
(\hyperref[II.4.5.13.3-fr]{4.5.13.3}) résulte alors des relations
\[
  (s_{\alpha|\beta}\otimes t_{\alpha|\beta})
  -(s_{\beta|\alpha}\otimes t_{\beta|\alpha})
  =(s_{\alpha|\beta}-s_{\beta|\alpha})\otimes t_{\alpha|\beta}
  +s_{\beta|\alpha}\otimes(t_{\alpha|\beta}-t_{\beta|\alpha})
\]
avec les mêmes notations que ci-dessus. Par récurrence sur $k$, on a donc
\begin{equation}
\label{II.4.5.13.4-fr}
  \partial(g^{\otimes k})=(kg^{\otimes(k-1)})\otimes(\partial g)
\tag{4.5.13.4}
\end{equation}

\oldpage[II]{89}

et on conclut de (\hyperref[II.4.5.13.2-fr]{4.5.13.2}) et
(\hyperref[II.4.5.13.4-fr]{4.5.13.4}) que l'on a
$\partial(g^{\otimes(m+1)})=0$~; $g^{\otimes(m+1)}$ est donc l'image
canonique d'une section $f$ de $\mathcal{L}^{\otimes n(m+1)}$ au-dessus de
$X$, ce qui achève de démontrer (\hyperref[II.4.5.13-fr]{4.5.13}).
\end{proof}

\begin{corollary}[4.5.14]
\label{II.4.5.14-fr}
Soient $X$ un schéma noethérien, $j$ l'injection canonique
$X_{\mathrm{red}}\to X$. Pour qu'un $\mathcal{O}_X$-Module inversible
$\mathcal{L}$ soit ample, il faut et il suffit que $j^*(\mathcal{L})$ soit
un $\mathcal{O}_{X_{\mathrm{red}}}$-Module ample.
\end{corollary}

Cela résulte de (\hyperref[I.6.1.6-fr]{I, 6.1.6}).

\subsection{Faisceaux relativement amples}
\label{II.4.6-fr}

\begin{definition}[4.6.1]
\label{II.4.6.1-fr}
Soient $f:X\to Y$ un morphisme quasi-compact, $\mathcal{L}$ un
$\mathcal{O}_X$-Module inversible. On dit que $\mathcal{L}$ est ample
relativement à $f$, ou relativement à $Y$, ou $f$-ample, ou $Y$-ample (ou
même ample si aucune confusion n'en résulte avec la notion définie dans
(\hyperref[II.4.5.3-fr]{4.5.3})) s'il existe un recouvrement $(U_\alpha)$
de $Y$ formé d'ouverts affines tel que si on pose
$X_\alpha=f^{-1}(U_\alpha)$, $\mathcal{L}|X_\alpha$ soit un
$\mathcal{O}_{X_\alpha}$-Module ample pour tout $\alpha$.
\end{definition}

L'existence d'un $\mathcal{O}_X$-Module $f$-ample entraîne que $f$ est
nécessairement séparé ((\hyperref[II.4.5.3-fr]{4.5.3}) et
(\hyperref[I.5.5.5-fr]{I, 5.5.5})).

\begin{proposition}[4.6.2]
\label{II.4.6.2-fr}
Soit $f:X\to Y$ un morphisme quasi-compact, $\mathcal{L}$ un
$\mathcal{O}_X$-Module inversible. Si $\mathcal{L}$ est très ample
relativement à $f$, il est ample relativement à $f$.
\end{proposition}

Cela résulte du caractère local (sur $Y$) de la notion de faisceau très
ample (\hyperref[II.4.4.5-fr]{4.4.5}), de la définition
(\hyperref[II.4.6.1-fr]{4.6.1}) et du critère
(\hyperref[II.4.5.10-fr]{4.5.10, (i)}).

\begin{proposition}[4.6.3]
\label{II.4.6.3-fr}
Soient $f:X\to Y$ un morphisme quasi-compact, $\mathcal{L}$ un
$\mathcal{O}_X$-Module inversible, et soit $\mathcal{S}$ la
$\mathcal{O}_Y$-Algèbre graduée
$\bigoplus_{n\geq0}f_*(\mathcal{L}^{\otimes n})$. Les conditions suivantes
sont équivalentes~:
\begin{enumerate}
  \item[\rm{a)}] $\mathcal{L}$ est $f$-ample.
  \item[\rm{b)}] $\mathcal{S}$ est quasi-cohérente et l'homomorphisme
  canonique $\sigma:f^*(\mathcal{S})\to
  \bigoplus_{n\geq0}\mathcal{L}^{\otimes n}$
  (\hyperref[0.4.4.3-fr]{0, 4.4.3}) est tel que le $Y$-morphisme
  $r_{\mathcal{L},\sigma}:G(\sigma)\to\operatorname{Proj}(\mathcal{S})=P$
  soit partout défini et soit une immersion ouverte dominante.
  \item[\rm{b')}] Le morphisme $f$ est séparé, le $Y$-morphisme
  $r_{\mathcal{L},\sigma}$ est partout défini et est un homéomorphisme de
  l'espace sous-jacent à $X$ sur un sous-espace de
  $\operatorname{Proj}(\mathcal{S})$.
\end{enumerate}
En outre, lorsqu'il en est ainsi, pour tout $n\in\mathbf{Z}$,
l'homomorphisme canonique
\begin{equation}
\label{II.4.6.3.1-fr}
  r_{\mathcal{L},\sigma}^*(\mathcal{O}_P(n))
  \longrightarrow\mathcal{L}^{\otimes n}
\tag{4.6.3.1}
\end{equation}
défini dans (\hyperref[II.3.7.9.1-fr]{3.7.9.1}), est un isomorphisme.

Enfin, pour tout $\mathcal{O}_X$-Module quasi-cohérent $\mathcal{F}$, si
l'on pose
$\mathcal{M}=\bigoplus_{n\geq0}f_*(\mathcal{F}\otimes
\mathcal{L}^{\otimes n})$, l'homomorphisme canonique
\begin{equation}
\label{II.4.6.3.2-fr}
  r_{\mathcal{L},\sigma}^*(\widetilde{\mathcal{M}})
  \longrightarrow\mathcal{F}
\tag{4.6.3.2}
\end{equation}
défini dans (\hyperref[II.3.7.9.2-fr]{3.7.9.2}), est un isomorphisme.
\end{proposition}

On a remarqué que a) implique que $f$ est séparé, donc $\mathcal{S}$ est
quasi-cohérente (\hyperref[I.9.2.2-fr]{I, 9.2.2, a)}). Comme le fait que
$r_{\mathcal{L},\sigma}$ soit une immersion partout définie est de
caractère local sur $Y$, pour prouver que a) entraîne b), on peut supposer
$Y$ affine et $\mathcal{L}$ ample~;

\oldpage[II]{90}

l'assertion résulte alors de (\hyperref[II.4.5.2-fr]{4.5.2, b)}). Il est
clair que b) entraîne b')~; enfin, pour prouver que b') entraîne a), il
suffit de considérer un recouvrement de $Y$ par des ouverts affines
$U_\alpha$ et d'appliquer le critère
(\hyperref[II.4.5.2-fr]{4.5.2, b')}) à chaque faisceau
$\mathcal{L}|f^{-1}(U_\alpha)$.

Pour les deux dernières assertions, on utilise le fait que
$\sigma^\flat$ est ici l'identité, et l'explicitation des homomorphismes
(\hyperref[II.3.7.9.1-fr]{3.7.9.1}) et
(\hyperref[II.3.7.9.2-fr]{3.7.9.2})~; il en résulte aussitôt que
(\hyperref[II.4.6.3.1-fr]{4.6.3.1}) est un isomorphisme. Quant à
(\hyperref[II.4.6.3.2-fr]{4.6.3.2}), on peut se ramener au cas où $Y$ est
affine, donc $\mathcal{L}$ ample~; il est clair que l'homomorphisme
(\hyperref[II.4.6.3.2-fr]{4.6.3.2}) est injectif, et le critère
(\hyperref[II.4.5.2-fr]{4.5.2, c)}) montre qu'il est surjectif, d'où la
conclusion.

\begin{corollary}[4.6.4]
\label{II.4.6.4-fr}
Soit $(U_\alpha)$ un recouvrement ouvert de $Y$. Pour que $\mathcal{L}$
soit ample relativement à $Y$, il faut et il suffit que
$\mathcal{L}|f^{-1}(U_\alpha)$ soit ample relativement à $U_\alpha$ pour
tout $\alpha$.
\end{corollary}

La condition b) est en effet locale sur $Y$.

\begin{corollary}[4.6.5]
\label{II.4.6.5-fr}
Soit $\mathcal{K}$ un $\mathcal{O}_Y$-Module inversible. Pour que
$\mathcal{L}$ soit $Y$-ample, il faut et il suffit que
$\mathcal{L}\otimes f^*(\mathcal{K})$ le soit.
\end{corollary}

C'est une conséquence évidente de (\hyperref[II.4.6.4-fr]{4.6.4}) en
prenant les $U_\alpha$ tels que $\mathcal{K}|U_\alpha$ soit isomorphe à
$\mathcal{O}_Y|U_\alpha$ pour tout $\alpha$.

\begin{corollary}[4.6.6]
\label{II.4.6.6-fr}
Supposons $Y$ affine~; pour que $\mathcal{L}$ soit $Y$-ample, il faut et il
suffit que $\mathcal{L}$ soit ample.
\end{corollary}

C'est une conséquence immédiate de la définition
(\hyperref[II.4.6.1-fr]{4.6.1}), et des critères
(\hyperref[II.4.6.3-fr]{4.6.3, b)}) et
(\hyperref[II.4.5.2-fr]{4.5.2, b)}), car ici
$\operatorname{Proj}(\mathcal{S})
=\operatorname{Proj}(\Gamma(Y,\mathcal{S}))$ par définition.

\begin{corollary}[4.6.7]
\label{II.4.6.7-fr}
Soit $f:X\to Y$ un morphisme quasi-compact. Supposons qu'il existe un
$\mathcal{O}_Y$-Module quasi-cohérent $\mathcal{E}$ et un $Y$-morphisme
$g:X\to P=\mathbf{P}(\mathcal{E})$ qui soit un homéomorphisme de l'espace
sous-jacent à $X$ sur un sous-espace de $P$~; alors
$\mathcal{L}=g^*(\mathcal{O}_P(1))$ est $Y$-ample.
\end{corollary}

On peut en effet supposer $Y$ affine~; le corollaire résulte alors du
critère (\hyperref[II.4.5.2-fr]{4.5.2, a)}), de la formule
(\hyperref[II.3.7.3.1-fr]{3.7.3.1}) et de
(\hyperref[II.4.2.3-fr]{4.2.3}).

\begin{proposition}[4.6.8]
\label{II.4.6.8-fr}
Soient $X$ un schéma quasi-compact ou un préschéma dont l'espace
sous-jacent est noethérien, $f:X\to Y$ un morphisme séparé quasi-compact.
Pour qu'un $\mathcal{O}_X$-Module inversible $\mathcal{L}$ soit $f$-ample,
il faut et il suffit que l'une des conditions équivalentes suivantes soit
vérifiée~:
\begin{enumerate}
  \item[\rm{c)}] Pour tout $\mathcal{O}_X$-Module quasi-cohérent
  $\mathcal{F}$ de type fini, il existe un entier $n_0>0$ tel que, pour
  tout $n\geq n_0$, l'homomorphisme canonique
  $\sigma:f^*(f_*(\mathcal{F}\otimes\mathcal{L}^{\otimes n}))\to
  \mathcal{F}\otimes\mathcal{L}^{\otimes n}$ soit surjectif.
  \item[\rm{c')}] Pour tout faisceau d'idéaux quasi-cohérent de type fini
  $\mathcal{J}$ de $\mathcal{O}_X$, il existe un entier $n>0$ tel que
  l'homomorphisme canonique
  $\sigma:f^*(f_*(\mathcal{J}\otimes\mathcal{L}^{\otimes n}))\to
  \mathcal{J}\otimes\mathcal{L}^{\otimes n}$ soit surjectif.
\end{enumerate}
\end{proposition}

Comme $X$ est quasi-compact, il en est de même de $f(X)$ et il existe donc
un recouvrement fini $(U_i)$ de $f(X)$ par des ouverts affines de $Y$. Pour
démontrer la condition c) lorsque $\mathcal{L}$ est $f$-ample, on peut
remplacer $Y$ par les $U_i$ et $X$ par les $f^{-1}(U_i)$, car si on obtient
pour chaque $i$ un entier $n_i$ tel que c) ait lieu (pour $U_i$,
$f^{-1}(U_i)$ et $\mathcal{L}|f^{-1}(U_i)$) dès que $n\geq n_i$, il
suffira de prendre pour $n_0$ le plus grand des $n_i$ pour obtenir c) pour
$Y$, $X$ et $\mathcal{L}$. Or lorsque $Y$ est affine, la condition c)
découle de (\hyperref[II.4.5.5-fr]{4.5.5, d)}), compte tenu de
(\hyperref[II.4.6.6-fr]{4.6.6}). Il est trivial que c) entraîne c'). Enfin,
pour prouver que c') implique que $\mathcal{L}$ est $f$-ample, on peut
encore se borner au cas où $Y$ est affine~: en effet, tout faisceau
d'idéaux quasi-cohérent de type fini $\mathcal{J}_i$ de
$\mathcal{O}_X|f^{-1}(U_i)$ est la restriction d'un faisceau cohérent
d'idéaux de type fini de $\mathcal{O}_X$
(\hyperref[I.9.4.7-fr]{I, 9.4.7}) et l'hypothèse c') entraîne que

\oldpage[II]{91}

$\mathcal{J}_i\otimes(\mathcal{L}^{\otimes n}|f^{-1}(U_i))$ est engendré
par ses sections (compte tenu de (\hyperref[I.9.2.2-fr]{I, 9.2.2}) et de
(\hyperref[II.3.4.7-fr]{3.4.7}))~; il suffit donc d'appliquer le critère
(\hyperref[II.4.5.5-fr]{4.5.5, d'')}).

\begin{proposition}[4.6.9]
\label{II.4.6.9-fr}
Soient $f:X\to Y$ un morphisme quasi-compact, $\mathcal{L}$ un
$\mathcal{O}_X$-Module inversible.
\begin{enumerate}
  \item[\rm{(i)}] Soit $n$ un entier $>0$. Pour que $\mathcal{L}$ soit
  $f$-ample, il faut et il suffit que $\mathcal{L}^{\otimes n}$ soit
  $f$-ample.
  \item[\rm{(ii)}] Soit $\mathcal{L}'$ un $\mathcal{O}_X$-Module
  inversible, et supposons qu'il existe un entier $n>0$ tel que
  l'homomorphisme canonique
  $\sigma:f^*(f_*(\mathcal{L}'^{\otimes n}))\to
  \mathcal{L}'^{\otimes n}$ soit surjectif. Alors, si $\mathcal{L}$ est
  $f$-ample, il en est de même de $\mathcal{L}\otimes\mathcal{L}'$.
\end{enumerate}
\end{proposition}

On se ramène en effet aussitôt au cas où $Y$ est affine, et la proposition
est alors conséquence immédiate de (\hyperref[II.4.5.6-fr]{4.5.6}).

\begin{corollary}[4.6.10]
\label{II.4.6.10-fr}
Le produit tensoriel de deux $\mathcal{O}_X$-Modules $f$-amples est
$f$-ample.
\end{corollary}

\begin{proposition}[4.6.11]
\label{II.4.6.11-fr}
Soient $Y$ un préschéma quasi-compact, $f:X\to Y$ un morphisme de type
fini, $\mathcal{L}$ un $\mathcal{O}_X$-Module inversible. Pour que
$\mathcal{L}$ soit $f$-ample, il faut et il suffit qu'il possède une des
propriétés équivalentes suivantes~:
\begin{enumerate}
  \item[\rm{d)}] Il existe $n_0>0$ tel que, pour tout entier $n\geq n_0$,
  $\mathcal{L}^{\otimes n}$ soit très ample relativement à $f$.
  \item[\rm{d')}] Il existe $n>0$ tel que $\mathcal{L}^{\otimes n}$ soit
  très ample relativement à $f$.
\end{enumerate}
\end{proposition}

Si $\mathcal{L}$ est ample relativement à $f$, il y a un recouvrement fini
$(U_i)$ de $Y$ par des ouverts affines tel que les
$\mathcal{L}|f^{-1}(U_i)$ soient amples. On en conclut
(\hyperref[II.4.5.10-fr]{4.5.10}) qu'il existe un entier $n_0$ tel que
$\mathcal{L}^{\otimes n}|f^{-1}(U_i)$ soit très ample relativement à
$f^{-1}(U_i)\to U_i$ pour tout $n\geq n_0$ et tout $i$, donc
$\mathcal{L}^{\otimes n}$ est très ample relativement à $f$
(\hyperref[II.4.4.5-fr]{4.4.5}). Inversement, d') entraîne déjà que
$\mathcal{L}^{\otimes n}$ est $f$-ample
(\hyperref[II.4.6.2-fr]{4.6.2}), donc il en est de même de $\mathcal{L}$
(\hyperref[II.4.6.9-fr]{4.6.9, (i)}).

\begin{corollary}[4.6.12]
\label{II.4.6.12-fr}
Soient $Y$ un préschéma quasi-compact, $f:X\to Y$ un morphisme de type
fini, $\mathcal{L}$, $\mathcal{L}'$ deux $\mathcal{O}_X$-Modules
inversibles. Si $\mathcal{L}$ est $f$-ample, il existe $n_0$ tel que
$\mathcal{L}^{\otimes n}\otimes\mathcal{L}'$ soit très ample relativement
à $f$ pour tout $n\geq n_0$.
\end{corollary}

On raisonne comme dans (\hyperref[II.4.6.11-fr]{4.6.11}) en utilisant un
recouvrement ouvert affine fini de $Y$ et
(\hyperref[II.4.5.11-fr]{4.5.11}).

\begin{proposition}[4.6.13]
\label{II.4.6.13-fr}
\begin{enumerate}
  \item[\rm{(i)}] Pour tout préschéma $Y$, tout
  $\mathcal{O}_Y$-Module inversible $\mathcal{L}$ est ample relativement
  au morphisme identique $1_Y$.
  \item[\rm{(i bis)}] Soient $f:X\to Y$ un morphisme quasi-compact,
  $j:X'\to X$ un morphisme quasi-compact qui est un homéomorphisme de
  l'espace sous-jacent à $X'$ sur un sous-espace de $X$. Si $\mathcal{L}$
  est un $\mathcal{O}_X$-Module ample relativement à $f$,
  $j^*(\mathcal{L})$ est ample relativement à $f\circ j$.
  \item[\rm{(ii)}] Soient $Z$ un préschéma quasi-compact, $f:X\to Y$,
  $g:Y\to Z$ deux morphismes quasi-compacts, $\mathcal{L}$ un
  $\mathcal{O}_X$-Module ample relativement à $f$, $\mathcal{K}$ un
  $\mathcal{O}_Y$-Module ample relativement à $g$. Alors il existe un
  entier $n_0>0$ tel que
  $\mathcal{L}\otimes f^*(\mathcal{K}^{\otimes n})$ soit ample relativement
  à $g\circ f$ pour tout $n\geq n_0$.
  \item[\rm{(iii)}] Soient $f:X\to Y$ un morphisme quasi-compact,
  $g:Y'\to Y$ un morphisme, et posons $X'=X_{(Y')}$. Si $\mathcal{L}$ est
  un $\mathcal{O}_X$-Module ample relativement à $f$,
  $\mathcal{L}'=\mathcal{L}\otimes_Y\mathcal{O}_{Y'}$ est un
  $\mathcal{O}_{X'}$-Module ample relativement à $f_{(Y')}$.
  \item[\rm{(iv)}] Soient $f_i:X_i\to Y_i$ ($i=1,2$) deux
  $S$-morphismes quasi-compacts. Si $\mathcal{L}_i$ est un
  $\mathcal{O}_{X_i}$-Module ample relativement à $f_i$ ($i=1,2$),
  $\mathcal{L}_1\otimes_S\mathcal{L}_2$ est ample relativement à
  $f_1\times_S f_2$.
  \item[\rm{(v)}] Soient $f:X\to Y$, $g:Y\to Z$ deux morphismes tels que
  $g\circ f$ soit quasi-compact. Si un $\mathcal{O}_X$-Module
  $\mathcal{L}$ est ample relativement à $g\circ f$, et si $g$ est séparé
  ou l'espace sous-jacent à $X$ localement noethérien, alors $\mathcal{L}$
  est ample relativement à $f$.

  \oldpage[II]{92}

  \item[\rm{(vi)}] Soient $f:X\to Y$ un morphisme quasi-compact, $j$
  l'injection canonique $X_{\mathrm{red}}\to X$. Si $\mathcal{L}$ est un
  $\mathcal{O}_X$-Module ample relativement à $f$, alors
  $j^*(\mathcal{L})$ est ample relativement à $f_{\mathrm{red}}$.
\end{enumerate}
\end{proposition}

Notons d'abord que (v) et (vi) se dérivent de (i), (i bis) et (iv) par le
même raisonnement que dans (\hyperref[II.4.4.10-fr]{4.4.10}), en utilisant
(\hyperref[II.4.6.4-fr]{4.6.4}) au lieu de
(\hyperref[II.4.4.5-fr]{4.4.5})~; nous laissons le détail du raisonnement au
lecteur. L'assertion (i) est trivialement conséquence de
(\hyperref[II.4.4.10-fr]{4.4.10, (i)}) et de
(\hyperref[II.4.6.2-fr]{4.6.2}). Pour démontrer (i bis), (iii) et (iv), nous
utiliserons le lemme suivant~:

\begin{lemma}[4.6.13.1]
\label{II.4.6.13.1-fr}
\begin{enumerate}
  \item[\rm{(i)}] Soient $u:Z\to S$ un morphisme, $\mathcal{L}$ un
  $\mathcal{O}_S$-Module inversible, $s$ une section de $\mathcal{L}$
  au-dessus de $S$, $s'$ la section de $u^*(\mathcal{L})=\mathcal{L}'$
  au-dessus de $Z$ qui lui correspond canoniquement. On a alors
  $Z_{s'}=u^{-1}(S_s)$.
  \item[\rm{(ii)}] Soient $Z$, $Z'$ deux $S$-préschémas, $p$, $p'$ les
  projections de $T=Z\times_S Z'$, $\mathcal{L}$ (resp.
  $\mathcal{L}'$) un $\mathcal{O}_Z$-Module (resp. un
  $\mathcal{O}_{Z'}$-Module) inversible, $t$ (resp. $t'$) une section de
  $\mathcal{L}$ (resp. $\mathcal{L}'$) au-dessus de $Z$ (resp. $Z'$), $s$
  (resp. $s'$) la section de $p^*(\mathcal{L})$ (resp.
  $p'^*(\mathcal{L}')$) au-dessus de $Z\times_S Z'$ qui lui correspond.
  Alors on a $T_{s\otimes s'}=Z_t\times_S Z'_{t'}$.
\end{enumerate}
\end{lemma}

Il résulte des définitions que l'on peut se ramener au cas où tous les
préschémas considérés sont affines. En outre, dans (i), on peut supposer
$\mathcal{L}=\mathcal{O}_S$~; l'assertion (i) résulte alors aussitôt de
(\hyperref[I.1.2.2.2-fr]{I, 1.2.2.2}). De même, dans (ii), on peut se
limiter au cas où $\mathcal{L}=\mathcal{O}_Z$,
$\mathcal{L}'=\mathcal{O}_{Z'}$, et alors l'assertion se réduit au lemme
(\hyperref[II.4.3.2.4-fr]{4.3.2.4}).

Démontrons alors (i bis). On peut supposer $Y$ affine
(\hyperref[II.4.6.4-fr]{4.6.4}), donc $\mathcal{L}$ ample
(\hyperref[II.4.6.6-fr]{4.6.6})~; lorsque $s$ parcourt l'ensemble réunion
des $\Gamma(X,\mathcal{L}^{\otimes n})$ ($n>0$), les $X_s$ forment une
base de la topologie de $X$ (\hyperref[II.4.5.2-fr]{4.5.2, a)}), donc par
hypothèse les $j^{-1}(X_s)$ forment une base de la topologie de $X'$~; on
conclut donc par le lemme
(\hyperref[II.4.6.13.1-fr]{4.6.13.1, (i)}) et par
(\hyperref[II.4.5.2-fr]{4.5.2, a)}) que $j^*(\mathcal{L})$ est ample.

Démontrons ensuite (iii). On peut de nouveau supposer $Y$ et $Y'$ affines
(\hyperref[II.4.6.4-fr]{4.6.4}), d'où il résulte que la projection
$h:X'\to X$ est affine (\hyperref[II.1.5.5-fr]{1.5.5}). Comme
$\mathcal{L}$ est ample (\hyperref[II.4.6.6-fr]{4.6.6}), lorsque $s$
parcourt l'ensemble des sections des $\mathcal{L}^{\otimes n}$ ($n>0$)
au-dessus de $X$ telles que $X_s$ soit affine, les $X_s$ recouvrent $X$
(\hyperref[II.4.5.2-fr]{4.5.2, a')})~; donc les $h^{-1}(X_s)$ sont affines
(\hyperref[II.1.2.5-fr]{1.2.5}) et recouvrent $X'$~; il résulte donc encore
du lemme (\hyperref[II.4.6.13.1-fr]{4.6.13.1, (i)}) et de
(\hyperref[II.4.5.2-fr]{4.5.2, a')}) que $\mathcal{L}'$ est ample, le
morphisme $f_{(Y')}$ étant quasi-compact
(\hyperref[I.6.6.4-fr]{I, 6.6.4, (iii)}).

Pour prouver (iv), notons d'abord que $f_1\times_S f_2$ est quasi-compact
(\hyperref[I.6.6.4-fr]{I, 6.6.4, (iv)}). On peut en outre supposer $S$,
$Y_1$ et $Y_2$ affines ((\hyperref[II.4.6.4-fr]{4.6.4}) et
(\hyperref[I.3.2.7-fr]{I, 3.2.7})), donc $\mathcal{L}_i$ ample ($i=1,2$)
(\hyperref[II.4.6.6-fr]{4.6.6}). Les ouverts
$(X_1)_{s_1}\times_S(X_2)_{s_2}$ forment un recouvrement de
$X_1\times_S X_2$ lorsque $s_i$ parcourt les sections des
$\mathcal{L}_i^{\otimes n_i}$ telles que $(X_i)_{s_i}$ soit affine
(\hyperref[II.4.5.2-fr]{4.5.2, a')}). D'ailleurs, en remplaçant $s_1$ et
$s_2$ par des puissances convenables, ce qui ne change pas les
$(X_i)_{s_i}$, on peut toujours supposer $n_1=n_2$. On déduit alors de
(\hyperref[II.4.6.13.1-fr]{4.6.13.1, (ii)}) et de
(\hyperref[II.4.5.2-fr]{4.5.2, a')}) que
$\mathcal{L}_1\otimes_S\mathcal{L}_2$ est ample, d'où l'assertion, puisque
$Y_1\times_S Y_2$ est affine (\hyperref[II.4.6.6-fr]{4.6.6}).

Reste à prouver (ii). Par le même raisonnement que dans
(\hyperref[II.4.4.10-fr]{4.4.10}), mais utilisant ici
(\hyperref[II.4.6.4-fr]{4.6.4}), on peut se borner au cas où $Z$ est
affine. Comme $\mathcal{K}$ est alors ample, et $Y$ quasi-compact, il existe
un nombre fini de sections $s_i\in\Gamma(Y,\mathcal{K}^{\otimes k_i})$
telles que les $Y_{s_i}$

\oldpage[II]{93}

soient affines et recouvrent $Y$
(\hyperref[II.4.5.2-fr]{4.5.2, a')})~;
remplaçant les $s_i$ par des puissances convenables, on peut en outre
supposer tous les $k_i$ égaux à un même entier $k$. Soient $s'_i$ les
sections de $f^*(\mathcal{K}^{\otimes k})$ au-dessus de $X$ correspondant
canoniquement aux $s_i$, de sorte que les
$X_{s'_i}=f^{-1}(Y_{s_i})$ (4.6.1.13, (i))
recouvrent $X$. Comme $\mathcal{L}|X_{s'_i}$ est ample
((\hyperref[II.4.6.4-fr]{4.6.4}) et
(\hyperref[II.4.6.6-fr]{4.6.6})), il existe pour chaque $i$ des sections en
nombre fini $t_{ij}\in\Gamma(X,\mathcal{L}^{\otimes n_{ij}})$ tels que les
$X_{t_{ij}}$ soient affines, contenus dans $X_{s'_i}$ et recouvrent
$X_{s'_i}$ (\hyperref[II.4.5.2-fr]{4.5.2, a')})~; on peut d'ailleurs
supposer tous les $n_{ij}$ égaux à un même nombre $n$. Cela étant, $X$ est
séparé et quasi-compact, donc il existe un entier $m>0$, et pour tout
$(i,j)$ une section
\[
  u_{ij}\in\Gamma(X,\mathcal{L}^{\otimes n}\otimes
  f^*(\mathcal{K}^{\otimes mk}))
\]
tels que $t_{ij}\otimes s_i'^{\otimes m}$ soit la restriction à $X_{s'_i}$
de $u_{ij}$ (\hyperref[I.9.3.1-fr]{I, 9.3.1})~; en outre on a
$X_{u_{ij}}=X_{t_{ij}}$, donc les $X_{u_{ij}}$ sont affines et recouvrent
$X$. On peut d'ailleurs supposer que $m$ est de la forme $nr$~; si on pose
$n_0=rk$, on voit donc (\hyperref[II.4.5.2-fr]{4.5.2, a')}) que
$\mathcal{L}\otimes_{\mathcal{O}_X}f^*(\mathcal{K}^{\otimes n_0})$ est
ample. En outre,
il existe $h_0>0$ tel que $\mathcal{K}^{\otimes h}$ soit engendré par ses
sections au-dessus de $Y$ dès que $h\geq h_0$
(\hyperref[II.4.5.5-fr]{4.5.5})~; a fortiori,
$f^*(\mathcal{K}^{\otimes h})$ est engendré par ses sections au-dessus de
$X$ pour $h\geq h_0$, par définition des images réciproques
((\hyperref[0.3.7.1-fr]{0, 3.7.1}) et
(\hyperref[II.4.4.1-fr]{4.4.1})). On en conclut que
$\mathcal{L}\otimes f^*(\mathcal{K}^{\otimes n_0+h})$ est ample dès que
$h\geq h_0$ (\hyperref[II.4.5.6-fr]{4.5.6}), ce qui achève la
démonstration.

\begin{remark}[4.6.14]
\label{II.4.6.14-fr}
Sous les conditions de (ii), on se gardera de croire que
$\mathcal{L}\otimes f^*(\mathcal{K})$ soit ample pour $g\circ f$~; en effet,
comme $\mathcal{L}\otimes f^*(\mathcal{K}^{-1})$ est aussi ample pour $f$
(\hyperref[II.4.6.5-fr]{4.6.5}), on en conclurait que $\mathcal{L}$ serait
ample pour $g\circ f$~; prenant en particulier pour $g$ le morphisme
identique, tout $\mathcal{O}_X$-Module inversible serait ample pour $f$, ce
qui n'est pas le cas en général (voir (\hyperref[II.5.1.6-fr]{5.1.6}),
(\hyperref[II.5.3.4-fr]{5.3.4, (i)}) et
(\hyperref[II.5.3.1-fr]{5.3.1})).
\end{remark}

\begin{proposition}[4.6.15]
\label{II.4.6.15-fr}
Soient $f:X\to Y$ un morphisme quasi-compact, $\mathcal{J}$ un faisceau
quasi-cohérent d'idéaux localement nilpotent de $\mathcal{O}_X$, $Z$ le
sous-préschéma fermé de $X$ défini par $\mathcal{J}$, $j:Z\to X$
l'injection canonique. Pour qu'un $\mathcal{O}_X$-Module inversible
$\mathcal{L}$ soit ample pour $f$, il faut et il suffit que
$j^*(\mathcal{L})$ soit ample pour $f\circ j$.
\end{proposition}

En effet, la question étant locale sur $Y$
(\hyperref[II.4.6.4-fr]{4.6.4}), on peut supposer $Y$ affine~; $X$ étant
alors quasi-compact, on peut supposer $\mathcal{J}$ nilpotent. Compte tenu
de (\hyperref[II.4.6.6-fr]{4.6.6}), la proposition n'est autre alors que
(\hyperref[II.4.5.13-fr]{4.5.13}).

\begin{corollary}[4.6.16]
\label{II.4.6.16-fr}
Soient $X$ un préschéma localement noethérien, $f:X\to Y$ un morphisme
quasi-compact. Pour qu'un $\mathcal{O}_X$-Module inversible $\mathcal{L}$
soit ample pour $f$, il faut et il suffit que son image réciproque
$\mathcal{L}'$ par l'injection canonique $X_{\mathrm{red}}\to X$ soit ample
pour $f_{\mathrm{red}}$.
\end{corollary}

On a déjà vu que la condition est nécessaire
(\hyperref[II.4.6.13-fr]{4.6.13, (vi)})~; inversement, si elle est remplie,
on peut se borner, pour prouver que $\mathcal{L}$ est ample pour $f$, au
cas où $Y$ est affine (\hyperref[II.4.6.4-fr]{4.6.4})~; alors
$Y_{\mathrm{red}}$ est aussi affine, donc $\mathcal{L}'$ est ample
(\hyperref[II.4.6.6-fr]{4.6.6}), et il en est de même de $\mathcal{L}$ en
vertu de (\hyperref[II.4.5.13-fr]{4.5.13}), puisque alors $X$ est
noethérien et $X_{\mathrm{red}}$ un sous-préschéma fermé de $X$ défini par
un faisceau quasi-cohérent nilpotent d'idéaux
(\hyperref[I.6.1.6-fr]{I, 6.1.6}).

\begin{proposition}[4.6.17]
\label{II.4.6.17-fr}
Les notations et hypothèses étant celles de
(\hyperref[II.4.4.11-fr]{4.4.11}), pour que $\mathcal{L}''$ soit ample
relativement à $f''$, il faut et il suffit que $\mathcal{L}$ soit ample
relativement à $f$ et $\mathcal{L}'$ ample relativement à $f'$.
\end{proposition}

\oldpage[II]{94}

La nécessité de la condition résulte de
(\hyperref[II.4.6.13-fr]{4.6.13, (i bis)}). Pour voir que la condition est
suffisante, on peut se borner au cas où $Y$ est affine, et alors le fait
que $\mathcal{L}''$ est ample résulte du critère
(\hyperref[II.4.5.2-fr]{4.5.2, a)}) appliqué à $\mathcal{L}$,
$\mathcal{L}'$ et $\mathcal{L}''$, en observant qu'une section de
$\mathcal{L}$ au-dessus de $X$ se prolonge (par $0$) en une section de
$\mathcal{L}''$ au-dessus de $X''$.

\begin{proposition}[4.6.18]
\label{II.4.6.18-fr}
Soient $Y$ un préschéma quasi-compact, $\mathcal{S}$ une
$\mathcal{O}_Y$-Algèbre graduée quasi-cohérente de type fini,
$X=\operatorname{Proj}(\mathcal{S})$, $f:X\to Y$ le morphisme structural.
Alors $f$ est de type fini, et il existe un entier $d>0$ tel que
$\mathcal{O}_X(d)$ soit inversible et $f$-ample.
\end{proposition}

En vertu de (\hyperref[II.3.1.10-fr]{3.1.10}), il existe un entier $d>0$
tel que $\mathcal{S}^{(d)}$ soit engendrée par $\mathcal{S}_d$. On sait que,
dans l'isomorphisme canonique entre $X$ et
$X^{(d)}=\operatorname{Proj}(\mathcal{S}^{(d)})$, $\mathcal{O}_X(d)$
s'identifie à $\mathcal{O}_{X^{(d)}}(1)$
(\hyperref[II.3.2.9-fr]{3.2.9, (ii)}). On voit donc qu'on est ramené au cas
où $\mathcal{S}$ est engendrée par $\mathcal{S}_1$~; la proposition résulte
alors de (\hyperref[II.4.4.3-fr]{4.4.3}) et de
(\hyperref[II.4.6.2-fr]{4.6.2}) (compte tenu de ce que $f$ est un morphisme
de type fini (\hyperref[II.3.4.1-fr]{3.4.1})).

\section{Morphismes quasi-affines~; morphismes quasi-projectifs. Morphismes
propres~; morphismes projectifs}

\subsection{Morphismes quasi-affines.}

\begin{definition}[5.1.1]
\label{II.5.1.1-fr}
On appelle \emph{schéma quasi-affine} un schéma isomorphe à un sous-schéma
induit sur un ouvert quasi-compact d'un schéma affine. On dit qu'un morphisme
$f:X\to Y$ est \emph{quasi-affine}, ou encore que $X$ (considéré comme
$Y$-préschéma au moyen de $f$) est un \emph{$Y$-schéma quasi-affine}, s'il
existe un recouvrement $(U_\alpha)$ de $Y$ par des ouverts affines tels que
les $f^{-1}(U_\alpha)$ soient des schémas quasi-affines.
\end{definition}

Il est clair qu'un morphisme quasi-affine est \emph{séparé}
(\hyperref[I.5.5.5-fr]{I, 5.5.5} et \hyperref[I.5.5.8-fr]{5.5.8}) et
\emph{quasi-compact} (\hyperref[I.6.6.1-fr]{I, 6.6.1})~; tout morphisme
affine est évidemment quasi-affine.

Rappelons que pour tout préschéma $X$, si on pose
$A=\Gamma(X,\mathcal{O}_X)$, l'homomorphisme identique
$A\to A=\Gamma(X,\mathcal{O}_X)$ définit un morphisme
$X\to\operatorname{Spec}(A)$ dit \emph{canonique}
(\hyperref[I.2.2.4-fr]{I, 2.2.4})~; ce n'est autre d'ailleurs que le morphisme
canonique défini dans (\hyperref[II.4.5.1-fr]{4.5.1}) pour le cas particulier
où $\mathcal{L}=\mathcal{O}_X$, si on se souvient que
$\operatorname{Proj}(A[T])$ s'identifie canoniquement à
$\operatorname{Spec}(A)$ (\hyperref[II.3.1.7-fr]{3.1.7}).

\begin{proposition}[5.1.2]
\label{II.5.1.2-fr}
Soient $X$ un schéma quasi-compact ou un préschéma dont l'espace sous-jacent
est noethérien, $A$ l'anneau $\Gamma(X,\mathcal{O}_X)$. Les conditions
suivantes sont équivalentes~:
\begin{enumerate}
  \item[\rm{a)}] $X$ est un schéma quasi-affine.
  \item[\rm{b)}] Le morphisme canonique
  $u:X\to\operatorname{Spec}(A)$ est une immersion ouverte.
  \item[\rm{b')}] Le morphisme canonique
  $u:X\to\operatorname{Spec}(A)$ est un homéomorphisme de $X$ sur un
  sous-espace de l'espace sous-jacent à $\operatorname{Spec}(A)$.
  \item[\rm{c)}] Le $\mathcal{O}_X$-Module $\mathcal{O}_X$ est très ample
  relativement à $u$ (\hyperref[II.4.4.2-fr]{4.4.2}).
  \item[\rm{c')}] Le $\mathcal{O}_X$-Module $\mathcal{O}_X$ est ample
  (\hyperref[II.4.5.1-fr]{4.5.1}).
  \item[\rm{d)}] Lorsque $f$ parcourt $A$, les $X_f$ forment une base de la
  topologie de $X$.
  \item[\rm{d')}] Lorsque $f$ parcourt $A$, ceux des $X_f$ qui sont affines
  forment un recouvrement de $X$.

\oldpage[II]{95}

  \item[\rm{e)}] Tout $\mathcal{O}_X$-Module quasi-cohérent est engendré par
  ses sections au-dessus de $X$.
  \item[\rm{e')}] Tout faisceau quasi-cohérent d'idéaux de type fini de
  $\mathcal{O}_X$ est engendré par ses sections au-dessus de $X$.
\end{enumerate}
\end{proposition}

Il est clair que b) entraîne a), et a) entraîne c) en vertu du critère b) de
(\hyperref[II.4.4.4-fr]{4.4.4}) appliqué au morphisme identique (compte tenu
de la remarque précédant l'énoncé)~; d'autre part c) entraîne c')
(\hyperref[II.4.5.10-fr]{4.5.10, (i)}), et c'), b) et b') sont équivalents par
(\hyperref[II.4.5.2-fr]{4.5.2, critères b) et b')}). Enfin c') équivaut à
chacun des critères d), d'), e), e') en vertu de
(\hyperref[II.4.5.2-fr]{4.5.2, critères a), a'), c)}) et
(\hyperref[II.4.5.5-fr]{4.5.5, critère d'')}).

On observera en outre qu'avec les notations précédentes, ceux des $X_f$ qui
sont affines forment une \emph{base} de la topologie de $X$, et que le
morphisme canonique $u$ est \emph{dominant}
(\hyperref[II.4.5.2-fr]{4.5.2}).

\begin{corollary}[5.1.3]
\label{II.5.1.3-fr}
Soit $X$ un préschéma quasi-compact. S'il existe un morphisme $v:X\to Y$ de
$X$ dans un schéma affine $Y$, qui soit un homéomorphisme de $X$ sur un
sous-espace ouvert de $Y$, alors $X$ est quasi-affine.
\end{corollary}

En effet, il existe une famille $(g_\alpha)$ de sections de $\mathcal{O}_Y$
au-dessus de $Y$ telles que les $D(g_\alpha)$ forment une base de la
topologie de $v(X)$~; si $v=(\psi,\theta)$ et si l'on pose
$f_\alpha=\theta(g_\alpha)$, on a
$X_{f_\alpha}=\psi^{-1}(D(g_\alpha))$
(\hyperref[I.2.2.4.1-fr]{I, 2.2.4.1}), donc les $X_{f_\alpha}$ forment une
base de la topologie de $X$, et le critère d) de
(\hyperref[II.5.1.2-fr]{5.1.2}) est vérifié.

\begin{corollary}[5.1.4]
\label{II.5.1.4-fr}
Si $X$ est un schéma quasi-affine, tout $\mathcal{O}_X$-Module inversible est
très ample (relativement au morphisme canonique) et a fortiori ample.
\end{corollary}

En effet, un tel Module $\mathcal{L}$ est engendré par ses sections au-dessus
de $X$ (\hyperref[II.5.1.2-fr]{5.1.2, e)}), donc
$\mathcal{L}\otimes\mathcal{O}_X=\mathcal{L}$ est très ample
(\hyperref[II.4.4.8-fr]{4.4.8}).

\begin{corollary}[5.1.5]
\label{II.5.1.5-fr}
Soit $X$ un préschéma quasi-compact. S'il existe un $\mathcal{O}_X$-Module
inversible $\mathcal{L}$ tel que $\mathcal{L}$ et $\mathcal{L}^{-1}$ soient
amples, $X$ est un schéma quasi-affine.
\end{corollary}

En effet, $\mathcal{O}_X=\mathcal{L}\otimes\mathcal{L}^{-1}$ est alors ample
(\hyperref[II.4.5.7-fr]{4.5.7}).

\begin{proposition}[5.1.6]
\label{II.5.1.6-fr}
Soit $f:X\to Y$ un morphisme quasi-compact. Les conditions suivantes sont
équivalentes~:
\begin{enumerate}
  \item[\rm{a)}] $f$ est quasi-affine.
  \item[\rm{b)}] La $\mathcal{O}_Y$-Algèbre
  $f_*(\mathcal{O}_X)=\mathcal{A}(X)$ est quasi-cohérente et le morphisme
  canonique $X\to\operatorname{Spec}(\mathcal{A}(X))$ correspondant à
  l'homomorphisme identique $\mathcal{A}(X)\to\mathcal{A}(X)$
  (\hyperref[II.1.2.7-fr]{1.2.7}) est une immersion ouverte.
  \item[\rm{b')}] La $\mathcal{O}_Y$-Algèbre $\mathcal{A}(X)$ est
  quasi-cohérente et le morphisme canonique
  $X\to\operatorname{Spec}(\mathcal{A}(X))$ est un homéomorphisme de $X$
  sur un sous-espace de $\operatorname{Spec}(\mathcal{A}(X))$.
  \item[\rm{c)}] Le $\mathcal{O}_X$-Module $\mathcal{O}_X$ est très ample
  pour $f$.
  \item[\rm{c')}] Le $\mathcal{O}_X$-Module $\mathcal{O}_X$ est ample pour
  $f$.
  \item[\rm{d)}] Le morphisme $f$ est séparé et, pour tout
  $\mathcal{O}_X$-Module quasi-cohérent $\mathcal{F}$, l'homomorphisme
  canonique $\sigma:f^*(f_*(\mathcal{F}))\to\mathcal{F}$
  (\hyperref[0.4.4.3-fr]{0, 4.4.3}) est surjectif.
\end{enumerate}
En outre, lorsque $f$ est quasi-affine, tout $\mathcal{O}_X$-Module
inversible $\mathcal{L}$ est très ample relativement à $f$.
\end{proposition}

L'équivalence de a) et c') résulte du caractère local sur $Y$ de la
$f$-amplitude (\hyperref[II.4.6.4-fr]{4.6.4}), de la définition
(\hyperref[II.5.1.1-fr]{5.1.1}) et du critère
(\hyperref[II.5.1.2-fr]{5.1.2, c')}). Les autres propriétés sont locales sur
$Y$

\oldpage[II]{96}

et résultent donc aussitôt de (\hyperref[II.5.1.2-fr]{5.1.2}) et
(\hyperref[II.5.1.4-fr]{5.1.4}), compte tenu de ce que
$f_*(\mathcal{F})$ est quasi-cohérent lorsque $f$ est séparé
(\hyperref[I.9.2.2-fr]{I, 9.2.2, a)}).

\begin{corollary}[5.1.7]
\label{II.5.1.7-fr}
Soit $f:X\to Y$ un morphisme quasi-affine. Pour tout ouvert $U$ de $Y$, la
restriction $f^{-1}(U)\to U$ de $f$ est quasi-affine.
\end{corollary}

\begin{corollary}[5.1.8]
\label{II.5.1.8-fr}
Soient $Y$ un schéma affine, $f:X\to Y$ un morphisme quasi-compact. Pour que
$f$ soit quasi-affine, il faut et il suffit que $X$ soit un schéma
quasi-affine.
\end{corollary}

C'est une conséquence immédiate de (\hyperref[II.5.1.6-fr]{5.1.6}) et
(\hyperref[II.4.6.6-fr]{4.6.6}).

\begin{corollary}[5.1.9]
\label{II.5.1.9-fr}
Soient $Y$ un schéma quasi-compact ou un préschéma dont l'espace sous-jacent
est noethérien, $f:X\to Y$ un morphisme de type fini. Si $f$ est
quasi-affine, il existe une sous-$\mathcal{O}_Y$-Algèbre quasi-cohérente
$\mathcal{B}$ de $\mathcal{A}(X)=f_*(\mathcal{O}_X)$, de type fini
(\hyperref[I.9.6.2-fr]{I, 9.6.2}), telle que le morphisme
$X\to\operatorname{Spec}(\mathcal{B})$ correspondant à l'injection canonique
$\mathcal{B}\to\mathcal{A}(X)$ (\hyperref[II.1.2.7-fr]{1.2.7}) soit une
immersion. En outre, toute sous-$\mathcal{O}_Y$-Algèbre quasi-cohérente de
type fini $\mathcal{B}'$ de $\mathcal{A}(X)$, contenant $\mathcal{B}$,
possède la même propriété.
\end{corollary}

En effet, $\mathcal{A}(X)$ est limite inductive de ses
sous-$\mathcal{O}_Y$-Algèbres quasi-cohérentes de type fini
(\hyperref[I.9.6.5-fr]{I, 9.6.5})~; le résultat est alors un cas particulier
de (\hyperref[II.3.8.4-fr]{3.8.4}), compte tenu de l'identification de
$\operatorname{Spec}(\mathcal{A}(X))$ et de
$\operatorname{Proj}(\mathcal{A}(X)[T])$
(\hyperref[II.3.1.7-fr]{3.1.7}).

\begin{proposition}[5.1.10]
\label{II.5.1.10-fr}
\begin{enumerate}
  \item[\rm{(i)}] Un morphisme quasi-compact $X\to Y$ qui est un
  homéomorphisme de l'espace sous-jacent à $X$ sur un sous-espace de
  l'espace sous-jacent à $Y$ (en particulier une immersion fermée) est
  quasi-affine.
  \item[\rm{(ii)}] Le composé de deux morphismes quasi-affines est
  quasi-affine.
  \item[\rm{(iii)}] Si $f:X\to Y$ est un $S$-morphisme quasi-affine,
  $f_{(S')}:X_{(S')}\to Y_{(S')}$ est un morphisme quasi-affine pour toute
  extension $S'\to S$ du préschéma de base.
  \item[\rm{(iv)}] Si $f:X\to Y$ et $g:X'\to Y'$ sont deux $S$-morphismes
  quasi-affines, $f\times_S g$ est quasi-affine.
  \item[\rm{(v)}] Si $f:X\to Y$, $g:Y\to Z$ sont deux morphismes tels que
  $g\circ f$ soit quasi-affine, et si $g$ est séparé ou l'espace sous-jacent
  à $X$ localement noethérien, $f$ est quasi-affine.
  \item[\rm{(vi)}] Si $f$ est un morphisme quasi-affine, il en est de même
  de $f_{\mathrm{red}}$.
\end{enumerate}
\end{proposition}

Compte tenu du critère (\hyperref[II.5.1.6-fr]{5.1.6, c')}), (i), (iii),
(iv), (v) et (vi) découlent aussitôt de
(\hyperref[II.4.6.13-fr]{4.6.13, (i bis), (iii), (iv), (v) et (vi)})
respectivement. Pour démontrer (ii), on peut se borner au cas où $Z$ est
affine, et alors l'assertion résulte directement de
(\hyperref[II.4.6.13-fr]{4.6.13, (ii)}), appliqué à
$\mathcal{L}=\mathcal{O}_X$ et $\mathcal{K}=\mathcal{O}_Y$.

\begin{remark}[5.1.11]
\label{II.5.1.11-fr}
Soient $f:X\to Y$, $g:Y\to Z$ deux morphismes tels que $X\times_ZY$ soit
localement noethérien. Alors l'immersion graphe
$\Gamma_f:X\to X\times_ZY$ est quasi-affine, étant quasi-compacte
(\hyperref[I.6.3.5-fr]{I, 6.3.5}), et
(\hyperref[I.5.5.12-fr]{I, 5.5.12}) montre que, dans (v), la conclusion reste
valable si on supprime l'hypothèse que $g$ est séparé.
\end{remark}

\begin{proposition}[5.1.12]
\label{II.5.1.12-fr}
Soient $f:X\to Y$ un morphisme quasi-compact, $g:X'\to X$ un morphisme
quasi-affine. Si $\mathcal{L}$ est un $\mathcal{O}_X$-Module ample pour
$f$, $g^*(\mathcal{L})$ est un $\mathcal{O}_{X'}$-Module ample pour
$f\circ g$.
\end{proposition}

En effet, comme $\mathcal{O}_{X'}$ est très ample pour $g$, et que la
question est locale sur $Y$ (\hyperref[II.4.6.4-fr]{4.6.4}), il résulte de
(\hyperref[II.4.6.13-fr]{4.6.13, (ii)}) qu'il existe (pour $Y$ affine) un
entier $n$ tel que $g^*(\mathcal{L})^{\otimes n}$ soit ample pour
$f\circ g$, donc $g^*(\mathcal{L})$ est ample pour $f\circ g$
(\hyperref[II.4.6.9-fr]{4.6.9}).

\oldpage[II]{97}

\subsection{Le critère de Serre.}

\begin{theorem}[5.2.1]
\label{II.5.2.1-fr}
\emph{(Critère de Serre.)} Soit $X$ un schéma quasi-compact ou un préschéma
dont l'espace sous-jacent est noethérien. Les conditions suivantes sont
équivalentes~:
\begin{enumerate}
  \item[\rm{a)}] $X$ est un schéma affine.
  \item[\rm{b)}] Il existe une famille d'éléments
  $f_\alpha\in A=\Gamma(X,\mathcal{O}_X)$ tels que les $X_{f_\alpha}$ soient
  affines et que l'idéal engendré par les $f_\alpha$ dans $A$ soit égal à
  $A$.
  \item[\rm{c)}] Le foncteur $\Gamma(X,\mathcal{F})$ est exact en
  $\mathcal{F}$ dans la catégorie des $\mathcal{O}_X$-Modules
  quasi-cohérents~; autrement dit, si
  \[
    0\to\mathcal{F}'\to\mathcal{F}\to\mathcal{F}''\to0 \tag{*}
  \]
  est une suite exacte de $\mathcal{O}_X$-Modules quasi-cohérents, la suite
  \[
    0\to\Gamma(X,\mathcal{F}')\to\Gamma(X,\mathcal{F})
      \to\Gamma(X,\mathcal{F}'')\to0
  \]
  est exacte.
  \item[\rm{c')}] La propriété c) a lieu pour toute suite exacte (*) de
  $\mathcal{O}_X$-Modules quasi-cohérents telle que $\mathcal{F}$ soit
  isomorphe à un sous-$\mathcal{O}_X$-Module d'un produit fini
  $\mathcal{O}_X^n$.
  \item[\rm{d)}] $\mathrm{H}^1(X,\mathcal{F})=0$ pour tout
  $\mathcal{O}_X$-Module quasi-cohérent $\mathcal{F}$.
  \item[\rm{d')}] $\mathrm{H}^1(X,\mathcal{J})=0$ pour tout faisceau
  quasi-cohérent d'idéaux $\mathcal{J}$ de $\mathcal{O}_X$.
\end{enumerate}
\end{theorem}

Il est évident que a) implique b)~; b) implique d'autre part que les
$X_{f_\alpha}$ recouvrent $X$, puisque par hypothèse la section $1$ est
combinaison linéaire des $f_\alpha$ et que les $D(f_\alpha)$ recouvrent
alors $\operatorname{Spec}(A)$. La dernière assertion de
(\hyperref[II.4.5.2-fr]{4.5.2}) implique donc que
$X\to\operatorname{Spec}(A)$ est un isomorphisme.

On sait que a) implique c) (\hyperref[I.1.3.11-fr]{I, 1.3.11}) et c)
entraîne trivialement c'). Prouvons que c') implique b). Tout d'abord, c')
entraîne que, pour tout point fermé $x\in X$ et tout voisinage ouvert $U$ de
$x$, il existe $f\in A$ tel que $x\in X_f\subset X-U$. Soit en effet
$\mathcal{J}$ (resp. $\mathcal{J}'$) le faisceau quasi-cohérent d'idéaux de
$\mathcal{O}_X$ définissant le sous-préschéma fermé réduit de $X$ ayant pour
espace sous-jacent $X-U$ (resp. $(X-U)\cup\{x\}$)
(\hyperref[I.5.2.1-fr]{I, 5.2.1})~; il est clair que l'on a
$\mathcal{J}'\subset\mathcal{J}$ et que
$\mathcal{J}''=\mathcal{J}/\mathcal{J}'$ est un $\mathcal{O}_X$-Module
quasi-cohérent ayant pour support $\{x\}$ et tel que
$\mathcal{J}''_x=k(x)$. L'hypothèse c') appliquée à la suite exacte
$0\to\mathcal{J}'\to\mathcal{J}\to\mathcal{J}''\to0$ montre que
$\Gamma(X,\mathcal{J})\to\Gamma(X,\mathcal{J}'')$ est surjectif. La section
de $\mathcal{J}''$ dont le germe en $x$ est $1_x$ est donc l'image d'une
section $f\in\Gamma(X,\mathcal{J})\subset\Gamma(X,\mathcal{O}_X)$, et on a
par définition $f(x)=1_x$ et $f(y)=0$ dans $X-U$, ce qui établit notre
assertion. D'ailleurs, si $U$ est affine, il en est de même de $X_f$
(\hyperref[I.1.3.6-fr]{I, 1.3.6}), donc la réunion des $X_f$ qui sont affines
($f\in A$) est un ensemble ouvert $Z$ contenant tous les points fermés de
$X$~; comme $X$ est un espace de Kolmogoroff quasi-compact, on a
nécessairement $Z=X$ (\hyperref[0.2.1.3-fr]{0, 2.1.3}). Puisque $X$ est
quasi-compact, il y a un nombre fini d'éléments $f_i\in A$
($1\leq i\leq n$) tels que les $X_{f_i}$ soient affines et recouvrent $X$.
Considérons alors l'homomorphisme $\mathcal{O}_X^n\to\mathcal{O}_X$ défini
par les sections $f_i$ (\hyperref[0.5.1.1-fr]{0, 5.1.1})~; comme, pour tout
$x\in X$, un au moins des $(f_i)_x$ est inversible, cet homomorphisme est
surjectif, et on a donc une suite exacte
$0\to\mathcal{R}\to\mathcal{O}_X^n\to\mathcal{O}_X\to0$, où $\mathcal{R}$
est un sous-$\mathcal{O}_X$-Module quasi-cohérent de $\mathcal{O}_X^n$. Il
résulte alors

\oldpage[II]{98}

de c') que l'homomorphisme correspondant
$\Gamma(X,\mathcal{O}_X^n)\to\Gamma(X,\mathcal{O}_X)$ est surjectif, ce qui
démontre b).

Enfin, a) implique d) (\hyperref[I.5.1.9.2-fr]{I, 5.1.9.2}) et d) implique
trivialement d'). Reste à montrer que d') entraîne c'). Or, si
$\mathcal{F}'$ est un sous-$\mathcal{O}_X$-Module quasi-cohérent de
$\mathcal{O}_X^n$, la filtration
$0\subset\mathcal{O}_X\subset\mathcal{O}_X^2\subset\cdots\subset
\mathcal{O}_X^n$ définit sur $\mathcal{F}'$ une filtration formée des
$\mathcal{F}'_k=\mathcal{O}_X^k\cap\mathcal{F}'$ ($0\leq k\leq n$), qui
sont des $\mathcal{O}_X$-Modules quasi-cohérents
(\hyperref[I.4.1.1-fr]{I, 4.1.1}), et
$\mathcal{F}'_{k+1}/\mathcal{F}'_k$ est isomorphe à un
sous-$\mathcal{O}_X$-Module quasi-cohérent de
$\mathcal{O}_X^{k+1}/\mathcal{O}_X^k=\mathcal{O}_X$, c'est-à-dire à un
faisceau quasi-cohérent d'idéaux de $\mathcal{O}_X$. L'hypothèse d')
entraîne donc
$\mathrm{H}^1(X,\mathcal{F}'_{k+1}/\mathcal{F}'_k)=0$~; la suite exacte de
cohomologie
$\mathrm{H}^1(X,\mathcal{F}'_k)\to\mathrm{H}^1(X,\mathcal{F}'_{k+1})
\to\mathrm{H}^1(X,\mathcal{F}'_{k+1}/\mathcal{F}'_k)=0$ permet alors de
prouver par récurrence sur $k$ que
$\mathrm{H}^1(X,\mathcal{F}'_k)=0$ pour tout $k$. C.Q.F.D.

\begin{remark}[5.2.1.1]
\label{II.5.2.1.1-fr}
Lorsque $X$ est un préschéma noethérien, on peut dans l'énoncé de c') et de
d'), remplacer «~quasi-cohérent~» par «~cohérent~». En effet, dans la
démonstration du fait que c') entraîne b), $\mathcal{J}$ et $\mathcal{J}'$
sont alors des faisceaux cohérents d'idéaux, et, d'autre part, tout
sous-Module quasi-cohérent d'un Module cohérent est cohérent
(\hyperref[I.6.1.1-fr]{I, 6.1.1})~; d'où la conclusion.
\end{remark}

\begin{corollary}[5.2.2]
\label{II.5.2.2-fr}
Soit $f:X\to Y$ un morphisme séparé quasi-compact. Les conditions suivantes
sont équivalentes~:
\begin{enumerate}
  \item[\rm{a)}] $f$ est un morphisme affine.
  \item[\rm{b)}] Le foncteur $f_*$ est exact dans la catégorie des
  $\mathcal{O}_X$-Modules quasi-cohérents.
  \item[\rm{c)}] Pour tout $\mathcal{O}_X$-Module quasi-cohérent
  $\mathcal{F}$, on a $\mathrm{R}^1f_*(\mathcal{F})=0$.
  \item[\rm{c')}] Pour tout faisceau d'idéaux quasi-cohérent $\mathcal{J}$
  de $\mathcal{O}_X$, on a $\mathrm{R}^1f_*(\mathcal{J})=0$.
\end{enumerate}
\end{corollary}

Toutes ces conditions étant locales sur $Y$, par définition du foncteur
$\mathrm{R}^1f_*$ (T, 3.7.3), on peut supposer que $Y$ est affine. Si $f$
est affine, $X$ est alors affine et la propriété b) n'est autre que
(\hyperref[I.1.6.4-fr]{I, 1.6.4}). Inversement, montrons que b) implique
a)~: pour tout $\mathcal{O}_X$-Module quasi-cohérent $\mathcal{F}$,
$f_*(\mathcal{F})$ est un $\mathcal{O}_Y$-Module quasi-cohérent
(\hyperref[I.9.2.2-fr]{I, 9.2.2, a)}). Par hypothèse le foncteur
$f_*(\mathcal{F})$ est exact en $\mathcal{F}$, et le foncteur
$\Gamma(Y,\mathcal{G})$ est exact en $\mathcal{G}$ (dans la catégorie des
$\mathcal{O}_Y$-Modules quasi-cohérents) puisque $Y$ est affine
(\hyperref[I.1.3.11-fr]{I, 1.3.11})~; donc
$\Gamma(Y,f_*(\mathcal{F}))=\Gamma(X,\mathcal{F})$ est exact en
$\mathcal{F}$, ce qui prouve notre assertion en vertu de
(\hyperref[II.5.2.1-fr]{5.2.1, c)}).

Si $f$ est affine, $f^{-1}(U)$ est affine pour tout ouvert affine $U$ de
$Y$ (\hyperref[II.1.3.2-fr]{1.3.2}), donc
$\mathrm{H}^1(f^{-1}(U),\mathcal{F})=0$
(\hyperref[II.5.2.1-fr]{5.2.1, d)}), ce qui par définition entraîne
$\mathrm{R}^1f_*(\mathcal{F})=0$. Supposons enfin vérifiée la condition
c')~; la suite exacte des termes de bas degré dans la suite spectrale de
Leray (G, II, 4.17.1 et I, 4.5.1) donne en particulier la suite exacte
\[
  0\to\mathrm{H}^1(Y,f_*(\mathcal{J}))\to
  \mathrm{H}^1(X,\mathcal{J})\to
  \mathrm{H}^0(Y,\mathrm{R}^1f_*(\mathcal{J})).
\]
Comme $Y$ est affine et $f_*(\mathcal{J})$ quasi-cohérent
(\hyperref[I.9.2.2-fr]{I, 9.2.2, a)}), on a
$\mathrm{H}^1(Y,f_*(\mathcal{J}))=0$
(\hyperref[II.5.2.1-fr]{5.2.1})~; l'hypothèse c') entraîne donc
$\mathrm{H}^1(X,\mathcal{J})=0$, et on conclut par
(\hyperref[II.5.2.1-fr]{5.2.1}) que $X$ est un schéma affine.

\begin{corollary}[5.2.3]
\label{II.5.2.3-fr}
Si $f:X\to Y$ est un morphisme affine, alors, pour tout
$\mathcal{O}_X$-Module quasi-cohérent $\mathcal{F}$, l'homomorphisme
canonique
$\mathrm{H}^1(Y,f_*(\mathcal{F}))\to\mathrm{H}^1(X,\mathcal{F})$ est
bijectif.
\end{corollary}

\oldpage[II]{99}

En effet, on a la suite exacte
\[
  0\to\mathrm{H}^1(Y,f_*(\mathcal{F}))\to
  \mathrm{H}^1(X,\mathcal{F})\to
  \mathrm{H}^0(Y,\mathrm{R}^1f_*(\mathcal{F}))
\]
des termes de bas degré de la suite spectrale de Leray, et la conclusion
résulte de (\hyperref[II.5.2.2-fr]{5.2.2}).

\begin{remark}[5.2.4]
\label{II.5.2.4-fr}
Au chapitre III, § 1, nous prouverons que si $X$ est affine, on a
$\mathrm{H}^i(X,\mathcal{F})=0$ pour tout $i>0$ et tout
$\mathcal{O}_X$-Module quasi-cohérent $\mathcal{F}$.
\end{remark}

\subsection{Morphismes quasi-projectifs.}

\begin{definition}[5.3.1]
\label{II.5.3.1-fr}
On dit qu'un morphisme $f:X\to Y$ est \emph{quasi-projectif}, ou que $X$
(considéré comme $Y$-préschéma au moyen de $f$) est \emph{quasi-projectif
sur $Y$}, ou que $X$ est un \emph{$Y$-schéma quasi-projectif}, si $f$ est de
type fini et s'il existe un $\mathcal{O}_X$-Module inversible qui est
$f$-ample.
\end{definition}

On notera que cette notion n'est pas locale sur $Y$~: les contre-exemples de
Nagata [26] et de Hironaka montrent que, même si $X$ et $Y$ sont des schémas
algébriques non singuliers sur un corps algébriquement clos, tout point de
$Y$ peut avoir un voisinage affine $U$ tel que $f^{-1}(U)$ soit
quasi-projectif sur $U$, sans que $f$ soit quasi-projectif.

On notera qu'un morphisme quasi-projectif est nécessairement séparé
(\hyperref[II.4.6.1-fr]{4.6.1}). Lorsque $Y$ est quasi-compact, il revient au
même de dire que $f$ est quasi-projectif ou que $f$ est de type fini et qu'il
existe un $\mathcal{O}_X$-Module très ample relativement à $f$
(\hyperref[II.4.6.2-fr]{4.6.2} et \hyperref[II.4.6.11-fr]{4.6.11}). De plus~:

\begin{proposition}[5.3.2]
\label{II.5.3.2-fr}
Soit $Y$ un schéma quasi-compact ou un préschéma dont l'espace sous-jacent
est noethérien, et soit $X$ un $Y$-préschéma. Les conditions suivantes sont
équivalentes~:
\begin{enumerate}
  \item[\rm{a)}] $X$ est un $Y$-schéma quasi-projectif.
  \item[\rm{b)}] $X$ est de type fini sur $Y$, et il existe un
  $\mathcal{O}_Y$-Module quasi-cohérent de type fini $\mathcal{E}$ tel que
  $X$ soit $Y$-isomorphe à un sous-préschéma de $\mathbf{P}(\mathcal{E})$.
  \item[\rm{c)}] $X$ est de type fini sur $Y$, et il existe une
  $\mathcal{O}_Y$-Algèbre graduée quasi-cohérente $\mathcal{S}$ telle que
  $\mathcal{S}_1$ soit de type fini et engendre $\mathcal{S}$ et que $X$
  soit $Y$-isomorphe à un sous-préschéma induit sur un ouvert partout dense
  de $\operatorname{Proj}(\mathcal{S})$.
\end{enumerate}
\end{proposition}

Cela résulte aussitôt de la remarque précédente et de
(\hyperref[II.4.4.3-fr]{4.4.3}), (\hyperref[II.4.4.6-fr]{4.4.6}) et
(\hyperref[II.4.4.7-fr]{4.4.7}).

On notera que lorsque $Y$ est un préschéma noethérien, on peut, dans les
conditions b) et c) de (\hyperref[II.5.3.2-fr]{5.3.2}), supprimer
l'hypothèse que $X$ est de type fini sur $Y$, qui est alors automatiquement
vérifiée (\hyperref[I.6.3.5-fr]{I, 6.3.5}).

\begin{corollary}[5.3.3]
\label{II.5.3.3-fr}
Soit $Y$ un schéma quasi-compact tel qu'il existe un $\mathcal{O}_Y$-Module
ample $\mathcal{L}$ (\hyperref[II.4.5.3-fr]{4.5.3}). Pour qu'un $Y$-schéma
$X$ soit quasi-projectif, il faut et il suffit qu'il soit de type fini sur
$Y$ et isomorphe à un sous-$Y$-schéma d'un fibré projectif de la forme
$\mathbf{P}_Y^r$.
\end{corollary}

En effet, si $\mathcal{E}$ est un $\mathcal{O}_Y$-Module quasi-cohérent de
type fini, $\mathcal{E}$ est isomorphe à un quotient d'un
$\mathcal{O}_Y$-Module
$\mathcal{L}^{\otimes(-n)}\otimes_{\mathcal{O}_Y}\mathcal{O}_Y^k$
(\hyperref[II.4.5.5-fr]{4.5.5}), donc $\mathbf{P}(\mathcal{E})$ est
isomorphe à un sous-schéma fermé de $\mathbf{P}_Y^{k-1}$
(\hyperref[II.4.1.2-fr]{4.1.2} et \hyperref[II.4.1.4-fr]{4.1.4}).

\begin{proposition}[5.3.4]
\label{II.5.3.4-fr}
\begin{enumerate}
  \item[\rm{(i)}] Un morphisme quasi-affine de type fini (et en particulier
  une immersion quasi-compacte, ou un morphisme affine de type fini) est
  quasi-projectif.
  \item[\rm{(ii)}] Si $f:X\to Y$ et $g:Y\to Z$ sont quasi-projectifs et si
  $Z$ est quasi-compact, $g\circ f$ est quasi-projectif.

\oldpage[II]{100}

  \item[\rm{(iii)}] Si $f:X\to Y$ est un $S$-morphisme quasi-projectif,
  $f_{(S')}:X_{(S')}\to Y_{(S')}$ est quasi-projectif pour toute extension
  $S'\to S$ du préschéma de base.
  \item[\rm{(iv)}] Si $f:X\to Y$ et $g:X'\to Y'$ sont deux $S$-morphismes
  quasi-projectifs, $f\times_Sg$ est quasi-projectif.
  \item[\rm{(v)}] Si $f:X\to Y$, $g:Y\to Z$ sont deux morphismes tels que
  $g\circ f$ soit quasi-projectif, et si $g$ est séparé ou $X$ localement
  noethérien, alors $f$ est quasi-projectif.
  \item[\rm{(vi)}] Si $f$ est un morphisme quasi-projectif, il en est de
  même de $f_{\mathrm{red}}$.
\end{enumerate}
\end{proposition}

(i) résulte de (\hyperref[II.5.1.6-fr]{5.1.6}) et
(\hyperref[II.5.1.10-fr]{5.1.10, (i)}). Les autres assertions sont des
conséquences immédiates de la définition (\hyperref[II.5.3.1-fr]{5.3.1}),
des propriétés des morphismes de type fini (\hyperref[I.6.3.4-fr]{I, 6.3.4})
et de (\hyperref[II.4.6.13-fr]{4.6.13}).

\begin{remark}[5.3.5]
\label{II.5.3.5-fr}
On notera qu'il peut se faire que $f_{\mathrm{red}}$ soit quasi-projectif
sans que $f$ le soit, même en supposant que $Y$ est le spectre d'une algèbre
de rang fini sur $\mathbf{C}$ et que $f$ est propre.
\end{remark}

\begin{corollary}[5.3.6]
\label{II.5.3.6-fr}
Si $X$ et $X'$ sont deux $Y$-schémas quasi-projectifs, $X\amalg X'$ est un
$Y$-schéma quasi-projectif.
\end{corollary}

Cela résulte de (\hyperref[II.4.6.18-fr]{4.6.18}).

\subsection{Morphismes propres et morphismes universellement fermés.}

\begin{definition}[5.4.1]
\label{II.5.4.1-fr}
On dit qu'un morphisme de préschémas $f:X\to Y$ est \emph{propre} s'il
vérifie les deux conditions suivantes~:
\begin{enumerate}
  \item[\rm{a)}] $f$ est séparé et de type fini.
  \item[\rm{b)}] Pour tout préschéma $Y'$ et tout morphisme $Y'\to Y$, la
  projection $f_{(Y')}:X\times_YY'\to Y'$ est un morphisme fermé
  (\hyperref[I.2.2.6-fr]{I, 2.2.6}).
\end{enumerate}
Lorsqu'il en est ainsi, on dit aussi que $X$ (considéré comme $Y$-préschéma
de morphisme structural $f$) est propre au-dessus de $Y$.
\end{definition}

Il est immédiat que les conditions a) et b) sont locales sur $Y$. Pour
vérifier que l'image d'une partie fermée $Z$ de $X\times_YY'$ par la
projection $q:X\times_YY'\to Y'$ est fermée dans $Y'$, il suffit de voir que
$q(Z)\cap U'$ est fermé dans $U'$ pour tout ouvert affine $U'$ de $Y'$~;
comme $q(Z)\cap U'=q(Z\cap q^{-1}(U'))$ et que $q^{-1}(U')$ s'identifie à
$X\times_YU'$ (\hyperref[I.4.4.1-fr]{I, 4.4.1}), on voit que pour vérifier
la condition b) de la déf. (\hyperref[II.5.4.1-fr]{5.4.1}), on peut se
limiter au cas où $Y'$ est un schéma affine. On verra plus loin
(\hyperref[II.5.6.3-fr]{5.6.3}) que si $Y$ est localement noethérien, on
peut même se borner à vérifier b) lorsque $Y'$ est de type fini sur $Y$.

Il est clair que tout morphisme propre est \emph{fermé}.

\begin{proposition}[5.4.2]
\label{II.5.4.2-fr}
\begin{enumerate}
  \item[\rm{(i)}] Une immersion fermée est un morphisme propre.
  \item[\rm{(ii)}] Le composé de deux morphismes propres est propre.
  \item[\rm{(iii)}] Si $X$, $Y$ sont des $S$-préschémas,
  $f:X\to Y$ un $S$-morphisme propre, alors
  $f_{(S')}:X_{(S')}\to Y_{(S')}$ est propre pour toute extension
  $S'\to S$ du préschéma de base.
  \item[\rm{(iv)}] Si $f:X\to Y$ et $g:X'\to Y'$ sont deux $S$-morphismes
  propres, le $S$-morphisme
  $f\times_Sg:X\times_SX'\to Y\times_SY'$ est propre.
\end{enumerate}
\end{proposition}

\oldpage[II]{101}

Il suffit de démontrer (i), (ii) et (iii) (\hyperref[I.3.5.1-fr]{I, 3.5.1}).
Dans chacun de ces trois cas, la vérification de la condition a) de
(\hyperref[II.5.4.1-fr]{5.4.1}) découle de résultats antérieurs
(\hyperref[I.5.5.1-fr]{I, 5.5.1} et \hyperref[I.6.3.4-fr]{6.3.4})~; reste à
vérifier la condition b). C'est immédiat dans le cas (i), car si $X\to Y$
est une immersion fermée, il en est de même de
$X\times_YY'\to Y\times_YY'=Y'$ (\hyperref[I.4.3.2-fr]{I, 4.3.2} et
\hyperref[II.3.3.3-fr]{3.3.3}). Pour démontrer (ii), considérons deux
morphismes propres $X\to Y$, $Y\to Z$, et un morphisme $Z'\to Z$. On peut
écrire $X\times_ZZ'=X\times_Y(Y\times_ZZ')$
(\hyperref[I.3.3.9.1-fr]{I, 3.3.9.1}), et par suite la projection
$X\times_ZZ'\to Z'$ se factorise en
$X\times_Y(Y\times_ZZ')\to Y\times_ZZ'\to Z'$. Compte tenu de la remarque
du début, (ii) résulte donc de ce que le composé de deux morphismes fermés
est fermé. Enfin, pour tout morphisme $S'\to S$, $X_{(S')}$ s'identifie à
$X\times_YY_{(S')}$ (\hyperref[I.3.3.11-fr]{I, 3.3.11})~; pour tout
morphisme $Z\to Y_{(S')}$, on peut écrire
\[
  X_{(S')}\times_{Y_{(S')}}Z
  =(X\times_YY_{(S')})\times_{Y_{(S')}}Z=X\times_YZ~;
\]
$X\times_YZ\to Z$ est fermé par hypothèse, donc cela prouve (iii).

\begin{corollary}[5.4.3]
\label{II.5.4.3-fr}
Soient $f:X\to Y$, $g:Y\to Z$ deux morphismes tels que $g\circ f$ soit
propre.
\begin{enumerate}
  \item[\rm{(i)}] Si $g$ est séparé, $f$ est propre.
  \item[\rm{(ii)}] Si $g$ est séparé et de type fini, et si $f$ est
  surjectif, $g$ est propre.
\end{enumerate}
\end{corollary}

(i) découle de (\hyperref[II.5.4.2-fr]{5.4.2}) par le procédé général
(\hyperref[I.5.5.12-fr]{I, 5.5.12}). Pour démontrer (ii), il n'y a à
vérifier que la condition b) de la déf. (\hyperref[II.5.4.1-fr]{5.4.1}).
Pour tout morphisme $Z'\to Z$, le diagramme
\[
  \xymatrix{
    X\times_ZZ'\ar[r]^{f\times1_{Z'}}\ar[dr]_p &
    Y\times_ZZ'\ar[d]^{p'}\\
    & Z'
  }
\]
(où $p$ et $p'$ sont les projections) est commutatif
(\hyperref[I.3.2.1-fr]{I, 3.2.1})~; en outre, $f\times1_{Z'}$ est surjectif
puisqu'il en est ainsi de $f$ (\hyperref[I.3.5.2-fr]{I, 3.5.2}), et $p$ est
un morphisme fermé par hypothèse. Toute partie fermée $F$ de
$Y\times_ZZ'$ est alors l'image par $f\times1_{Z'}$ d'une partie fermée $E$
de $X\times_ZZ'$, donc $p'(F)=p(E)$ est fermé dans $Z'$ par hypothèse, d'où
le corollaire.

\begin{corollary}[5.4.4]
\label{II.5.4.4-fr}
Si $X$ est un préschéma propre au-dessus de $Y$, et $\mathcal{S}$ une
$\mathcal{O}_Y$-Algèbre quasi-cohérente, tout $Y$-morphisme
$f:X\to\operatorname{Proj}(\mathcal{S})$ est propre (et a fortiori fermé).
\end{corollary}

En effet, le morphisme structural
$p:\operatorname{Proj}(\mathcal{S})\to Y$ est séparé, et $p\circ f$ est
propre par hypothèse.

\begin{corollary}[5.4.5]
\label{II.5.4.5-fr}
Soit $f:X\to Y$ un morphisme séparé de type fini. Soient
$(X_i)_{1\leq i\leq n}$ (resp. $(Y_i)_{1\leq i\leq n}$) une famille finie
de sous-préschémas fermés de $X$ (resp. $Y$), $j_i$ (resp. $h_i$)
l'injection canonique $X_i\to X$ (resp. $Y_i\to Y$). Supposons que l'espace
sous-jacent $X$ soit réunion des $X_i$, et que pour tout $i$, il y ait un
morphisme $f_i:X_i\to Y_i$ tel que le diagramme
\[
  \xymatrix{
    X_i\ar[r]^{f_i}\ar[d]_{j_i} & Y_i\ar[d]^{h_i}\\
    X\ar[r]_f & Y
  }
\]
soit commutatif. Alors, pour que $f$ soit propre, il faut et il suffit que
chacun des $f_i$ le soit.
\end{corollary}

\oldpage[II]{102}

Si $f$ est propre, il en est de même de $f\circ j_i$, puisque $j_i$ est une
immersion fermée (\hyperref[II.5.4.2-fr]{5.4.2})~; comme $h_i$ est une
immersion fermée, donc un morphisme séparé, $f_i$ est propre par
(\hyperref[II.5.4.3-fr]{5.4.3}). Supposons inversement que chacun des $f_i$
soit propre, et considérons le préschéma $Z$ somme des $X_i$, soit $u$ le
morphisme $Z\to X$ qui se réduit à $j_i$ dans chacun des $X_i$. La
restriction de $f\circ u$ à chaque $X_i$ est égale à
$f\circ j_i=h_i\circ f_i$, donc est propre puisqu'il en est ainsi de $h_i$
et de $f_i$ (\hyperref[II.5.4.2-fr]{5.4.2})~; il résulte alors aussitôt de
la déf. (\hyperref[II.5.4.1-fr]{5.4.1}) que $f\circ u$ est propre. Mais
comme $u$ est surjectif par hypothèse, on conclut que $f$ est propre par
(\hyperref[II.5.4.3-fr]{5.4.3}).

\begin{corollary}[5.4.6]
\label{II.5.4.6-fr}
Soit $f:X\to Y$ un morphisme séparé et de type fini~; pour que $f$ soit
propre, il faut et il suffit que
$f_{\mathrm{red}}:X_{\mathrm{red}}\to Y_{\mathrm{red}}$ le soit.
\end{corollary}

C'est un cas particulier de (\hyperref[II.5.4.5-fr]{5.4.5}) avec $n=1$,
$X_1=X_{\mathrm{red}}$ et $Y_1=Y_{\mathrm{red}}$
(\hyperref[I.5.1.5-fr]{I, 5.1.5}).

\begin{env}[5.4.7]
\label{II.5.4.7-fr}
Si $X$ et $Y$ sont des préschémas noethériens et $f:X\to Y$ un morphisme
séparé de type fini, on peut, pour vérifier que $f$ est propre, se ramener à
le faire pour des morphismes dominants de préschémas intègres. En effet,
soient $X_i$ ($1\leq i\leq n$) les composantes irréductibles en nombre fini
de $X$ et considérons pour chaque $i$ l'unique sous-préschéma fermé réduit de
$X$ ayant $X_i$ pour espace sous-jacent, que nous désignerons encore par
$X_i$ (\hyperref[I.5.2.1-fr]{I, 5.2.1}). Soit d'autre part $Y_i$ l'unique
sous-préschéma fermé réduit de $Y$ ayant $\overline{f(X_i)}$ pour espace
sous-jacent. Si $g_i$ (resp. $h_i$) est le morphisme d'injection $X_i\to X$
(resp. $Y_i\to Y$), on en conclut que $f\circ g_i=h_i\circ f_i$, où $f_i$
est un morphisme dominant $X_i\to Y_i$
(\hyperref[I.5.2.2-fr]{I, 5.2.2})~; on est alors dans les conditions
d'application de (\hyperref[II.5.4.5-fr]{5.4.5}), et pour que $f$ soit
propre, il faut et il suffit que les $f_i$ le soient.
\end{env}

\begin{corollary}[5.4.8]
\label{II.5.4.8-fr}
Soient $X$, $Y$ des $S$-préschémas séparés de type fini sur $S$, et soit
$f:X\to Y$ un $S$-morphisme. Pour que $f$ soit propre, il faut et il suffit
que, pour tout $S$-préschéma $S'$, le morphisme
$f\times_S1_{S'}:X\times_SS'\to Y\times_SS'$ soit fermé.
\end{corollary}

Notons d'abord que si $g:X\to S$ et $h:Y\to S$ sont les morphismes
structuraux, on a par définition $g=h\circ f$, donc $f$ est séparé et de
type fini (\hyperref[I.5.5.1-fr]{I, 5.5.1} et
\hyperref[I.6.3.4-fr]{6.3.4}). Si $f$ est propre, il en est de même de
$f\times_S1_{S'}$ (\hyperref[II.5.4.2-fr]{5.4.2})~; a fortiori,
$f\times_S1_{S'}$ est fermé. Inversement, supposons vérifiée la condition de
l'énoncé et soit $Y'$ un $Y$-préschéma~; $Y'$ peut aussi être considéré
comme un $S$-préschéma, et comme $Y\to S$ est séparé, $X\times_YY'$
s'identifie à un sous-préschéma fermé de $X\times_SY'$
(\hyperref[I.5.4.2-fr]{I, 5.4.2}). Dans le diagramme commutatif
\[
  \xymatrix{
    X\times_YY'\ar[r]^{f\times1_{Y'}}\ar[d] &
    Y\times_YY'=Y'\ar[d]\\
    X\times_SY'\ar[r]_{f\times_S1_{Y'}} & Y\times_SY'
  }
\]
les flèches verticales sont des immersions fermées~; il en résulte aussitôt
que si $f\times_S1_{Y'}$ est un morphisme fermé, il en est de même de
$f\times_Y1_{Y'}$.

\begin{remark}[5.4.9]
\label{II.5.4.9-fr}
Nous dirons qu'un morphisme $f:X\to Y$ est \emph{universellement fermé} s'il
satisfait à la condition b) de la définition
(\hyperref[II.5.4.1-fr]{5.4.1}). Le lecteur observera que

\oldpage[II]{103}

dans les propositions (\hyperref[II.5.4.2-fr]{5.4.2}) à
(\hyperref[II.5.4.8-fr]{5.4.8}), on peut partout remplacer «~propre~» par
«~universellement fermé~» sans modifier la validité des résultats (et dans
les hypothèses de (\hyperref[II.5.4.3-fr]{5.4.3}),
(\hyperref[II.5.4.5-fr]{5.4.5}), (\hyperref[II.5.4.6-fr]{5.4.6}) et
(\hyperref[II.5.4.8-fr]{5.4.8}), on peut omettre les conditions de
finitude).
\end{remark}

\begin{env}[5.4.10]
\label{II.5.4.10-fr}
Soit $f:X\to Y$ un morphisme de type fini. Nous dirons qu'une partie fermée
$Z$ de $X$ est \emph{propre sur $Y$} (ou \emph{$Y$-propre}, ou \emph{propre
pour $f$}) si la restriction de $f$ à un sous-préschéma fermé de $X$,
d'espace sous-jacent $Z$ (\hyperref[I.5.2.1-fr]{I, 5.2.1}), est propre.
Comme cette restriction est alors séparée, il résulte de
(\hyperref[II.5.4.6-fr]{5.4.6}) et (\hyperref[I.5.5.1-fr]{I, 5.5.1, (vi)})
que la propriété précédente ne dépend pas du sous-préschéma fermé de $X$
ayant $Z$ pour espace sous-jacent. Si $g:X'\to X$ est un morphisme propre,
$g^{-1}(Z)$ est une partie propre de $X'$~: si $T$ est un sous-préschéma de
$X$ ayant pour espace sous-jacent $Z$, il suffit en effet de remarquer que
la restriction de $g$ au sous-préschéma fermé $g^{-1}(T)$ de $X'$ est un
morphisme propre $g^{-1}(T)\to T$ par
(\hyperref[II.5.4.2-fr]{5.4.2, (iii)}), et d'appliquer ensuite
(\hyperref[II.5.4.2-fr]{5.4.2, (ii)}). D'autre part, si $X''$ est un
$Y$-schéma de type fini et $u:X\to X''$ un $Y$-morphisme, $u(Z)$ est une
partie propre de $X''$~; en effet, prenons pour $T$ le sous-préschéma fermé
réduit de $X$ ayant $Z$ pour espace sous-jacent~; la restriction de $f$ à
$T$ étant propre, il en est de même de la restriction de $u$ à $T$
(\hyperref[II.5.4.3-fr]{5.4.3, (i)}), donc $u(Z)$ est fermé dans $X''$~;
soit $T''$ un sous-préschéma fermé de $X''$ ayant $u(Z)$ pour espace
sous-jacent (\hyperref[I.5.2.1-fr]{I, 5.2.1}), de sorte que $u|T$ se
factorise en
\[
  T\xrightarrow{v}T''\xrightarrow{j}X'',
\]
où $j$ est l'injection canonique (\hyperref[I.5.2.2-fr]{I, 5.2.2}) et $v$
est donc propre et surjectif (\hyperref[II.5.4.5-fr]{5.4.5})~; si $g$ est la
restriction à $T''$ du morphisme structural $X''\to Y$, $g$ est séparé et
de type fini, et on a $f|T=g\circ v$~; il résulte donc de
(\hyperref[II.5.4.3-fr]{5.4.3, (ii)}) que $g$ est propre, d'où notre
assertion.
\end{env}

Il résulte en particulier de ces remarques que si $Z$ est une partie
$Y$-propre de $X$, alors~:
\begin{enumerate}
  \item[\rm{1\textsuperscript{o}}] Pour tout sous-préschéma fermé $X'$ de
  $X$, $Z\cap X'$ est une partie $Y$-propre de $X'$.
  \item[\rm{2\textsuperscript{o}}] Si $X$ est un sous-préschéma d'un
  $Y$-schéma de type fini $X''$, $Z$ est aussi une partie $Y$-propre de
  $X''$ (et en particulier est fermée dans $X''$).
\end{enumerate}

\subsection{Morphismes projectifs.}

\begin{proposition}[5.5.1]
\label{II.5.5.1-fr}
Soit $X$ un $Y$-préschéma. Les conditions suivantes sont équivalentes~:
\begin{enumerate}
  \item[\rm{a)}] $X$ est $Y$-isomorphe à un sous-préschéma fermé d'un fibré
  projectif $\mathbf{P}(\mathcal{E})$, où $\mathcal{E}$ est un
  $\mathcal{O}_Y$-Module quasi-cohérent de type fini.
  \item[\rm{b)}] Il existe une $\mathcal{O}_Y$-Algèbre graduée
  quasi-cohérente $\mathcal{S}$, telle que $\mathcal{S}_1$ soit de type fini
  et engendre $\mathcal{S}$, et que $X$ soit $Y$-isomorphe à
  $\operatorname{Proj}(\mathcal{S})$.
\end{enumerate}
\end{proposition}

La condition a) entraîne b) en vertu de
(\hyperref[II.3.6.2-fr]{3.6.2, (ii)})~: si $\mathcal{J}$ est un faisceau
gradué quasi-cohérent d'idéaux de $\mathbf{S}(\mathcal{E})$, la
$\mathcal{O}_Y$-Algèbre graduée quasi-cohérente
$\mathcal{S}=\mathbf{S}(\mathcal{E})/\mathcal{J}$ est engendrée par
$\mathcal{S}_1$ et ce dernier, image canonique de $\mathcal{E}$, est un
$\mathcal{O}_Y$-Module de type fini. La condition b) implique a) en vertu de
(\hyperref[II.3.6.2-fr]{3.6.2}) appliqué au cas où
$\mathcal{M}\to\mathcal{S}_1$ est l'application identique.

\oldpage[II]{104}

\begin{definition}[5.5.2]
\label{II.5.5.2-fr}
On dit qu'un $Y$-préschéma $X$ est \emph{projectif sur $Y$} ou est un
\emph{$Y$-schéma projectif} s'il vérifie les conditions équivalentes a) et b)
de (\hyperref[II.5.5.1-fr]{5.5.1}). On dit qu'un morphisme $f:X\to Y$ est
\emph{projectif} s'il fait de $X$ un $Y$-schéma projectif.
\end{definition}

Il est clair que si $f:X\to Y$ est projectif, il existe un
$\mathcal{O}_X$-Module \emph{très ample relativement à $f$}
(\hyperref[II.4.4.2-fr]{4.4.2}).

\begin{theorem}[5.5.3]
\label{II.5.5.3-fr}
\begin{enumerate}
  \item[\rm{(i)}] Tout morphisme projectif est quasi-projectif et propre.
  \item[\rm{(ii)}] Inversement, soit $Y$ un schéma quasi-compact ou un
  préschéma dont l'espace sous-jacent est noethérien~; alors tout morphisme
  $f:X\to Y$ qui est quasi-projectif et propre est projectif.
\end{enumerate}
\end{theorem}

(i) Il est clair que si $f:X\to Y$ est projectif, il est de type fini et
quasi-projectif (donc en particulier séparé)~; d'autre part, il résulte
aussitôt de (\hyperref[II.5.5.1-fr]{5.5.1, b)}) et de
(\hyperref[II.3.5.3-fr]{3.5.3}) que si $f$ est projectif, il en est de même de
$f\times_Y1_{Y'}:X\times_YY'\to Y'$ pour tout morphisme $Y'\to Y$. Pour
montrer que $f$ est universellement fermé, tout revient donc à prouver qu'un
morphisme projectif $f$ est \emph{fermé}. La question étant locale sur $Y$, on
peut supposer que $Y=\operatorname{Spec}(A)$, donc
(\hyperref[II.5.5.1-fr]{5.5.1}) $X=\operatorname{Proj}(S)$, où $S$ est une
$A$-algèbre graduée engendrée par un nombre fini d'éléments de $S_1$. Pour
tout $y\in Y$, la fibre $f^{-1}(y)$ s'identifie à
$\operatorname{Proj}(S)\times_Y\operatorname{Spec}(k(y))$
(\hyperref[I.3.6.1-fr]{I, 3.6.1}), donc à
$\operatorname{Proj}(S\otimes_Ak(y))$
(\hyperref[II.2.8.10-fr]{2.8.10})~; par suite, $f^{-1}(y)$ est vide si et
seulement si $S\otimes_Ak(y)$ vérifie la condition (TN)
(\hyperref[II.2.7.4-fr]{2.7.4}), autrement dit $S_n\otimes_Ak(y)=0$ pour $n$
assez grand. Mais comme $(S_n)_y$ est un $\mathcal{O}_y$-module de type fini,
la condition précédente signifie que $(S_n)_y=0$ pour $n$ assez grand, en
vertu du lemme de Nakayama. Si $\mathfrak a_n$ est l'annulateur dans $A$ du
$A$-module $S_n$, la condition précédente signifie encore que
$\mathfrak a_n\not\subset\mathfrak j_y$ pour $n$ assez grand
(\hyperref[0.1.7.4-fr]{0, 1.7.4}). Or, comme $S_nS_1=S_{n+1}$ par hypothèse,
on a $\mathfrak a_n\subset\mathfrak a_{n+1}$, et si $\mathfrak a$ est la
réunion des $\mathfrak a_n$, on voit donc que $f(X)=V(\mathfrak a)$, ce qui
prouve que $f(X)$ est fermé dans $Y$. Si maintenant $X'$ est une partie
fermée quelconque de $X$, il existe un sous-préschéma fermé de $X$ ayant $X'$
pour espace sous-jacent (\hyperref[I.5.2.1-fr]{I, 5.2.1}) et il est clair
(\hyperref[II.5.5.1-fr]{5.5.1, a)}) que le morphisme
$X'\to X\xrightarrow{f}Y$ est projectif, donc $f(X')$ est fermé dans $Y$.

(ii) L'hypothèse sur $Y$ et le fait que $f$ est quasi-projectif entraînent
l'existence d'un $\mathcal{O}_Y$-Module quasi-cohérent de type fini
$\mathcal{E}$ et d'une $Y$-immersion $j:X\to\mathbf{P}(\mathcal{E})$
(\hyperref[II.5.3.2-fr]{5.3.2}). Mais comme $f$ est propre, $j$ est fermée par
(\hyperref[II.5.4.4-fr]{5.4.4}), donc $f$ est projectif.

\par\medskip\noindent\textit{Remarques (5.5.4). ---\ }
\label{II.5.5.4-fr}
\begin{enumerate}
  \item[\rm{(i)}] Soit $f:X\to Y$ un morphisme tel que~:
  1\textsuperscript{o} $f$ soit propre~; 2\textsuperscript{o} il existe un
  $\mathcal{O}_X$-Module $\mathcal{L}$ très ample relativement à $f$~;
  3\textsuperscript{o} le $\mathcal{O}_Y$-Module quasi-cohérent
  $\mathcal{E}=f_*(\mathcal{L})$ soit de type fini. Alors $f$ est un morphisme
  projectif~: en effet (\hyperref[II.4.4.4-fr]{4.4.4}), il y a alors une
  $Y$-immersion $r:X\to\mathbf{P}(\mathcal{E})$, et comme $f$ est propre, $r$
  est une immersion fermée (\hyperref[II.5.4.4-fr]{5.4.4}). On verra au
  chapitre III, \S~3, que lorsque $Y$ est localement noethérien, la condition
  3\textsuperscript{o} ci-dessus est une conséquence des deux autres, donc les
  conditions 1\textsuperscript{o} et 2\textsuperscript{o} caractérisent dans
  ce cas les morphismes projectifs et si $Y$ est quasi-compact, on peut
  remplacer la condition 2\textsuperscript{o} par l'hypothèse de l'existence
  d'un $\mathcal{O}_X$-Module ample pour $f$
  (\hyperref[II.4.6.11-fr]{4.6.11}).
  \item[\rm{(ii)}] Soit $Y$ un schéma quasi-compact tel qu'il existe un
  $\mathcal{O}_Y$-Module ample. Pour qu'un $Y$-schéma $X$ soit projectif, il
  faut et il suffit qu'il soit $Y$-isomorphe à un sous-$Y$-schéma fermé d'un
  fibré projectif de la forme $\mathbf{P}_Y^r$. La condition est évidemment
  suffisante.

\oldpage[II]{105}

  Inversement, si $X$ est projectif sur $Y$, il est quasi-projectif, donc il
  existe une $Y$-immersion $j$ de $X$ dans un $\mathbf{P}_Y^r$
  (\hyperref[II.5.3.3-fr]{5.3.3}) qui est fermée par
  (\hyperref[II.5.4.4-fr]{5.4.4}) et
  (\hyperref[II.5.5.3-fr]{5.5.3}).
  \item[\rm{(iii)}] Le raisonnement de
  (\hyperref[II.5.5.3-fr]{5.5.3}) montre que, pour tout préschéma $Y$, et tout
  entier $r\geq0$, le morphisme structural $\mathbf{P}_Y^r\to Y$ est
  surjectif, car si on pose
  $\mathcal{S}=\mathbf{S}_{\mathcal{O}_Y}(\mathcal{O}_Y^{r+1})$, on a
  évidemment $\mathcal{S}_y=\mathbf{S}_{k(y)}((k(y))^{r+1})$
  (\hyperref[II.1.7.3-fr]{1.7.3}), donc $(\mathcal{S}_n)_y\neq0$ pour tout
  $y\in Y$ et tout $n\geq0$.
  \item[\rm{(iv)}] Il résulte des exemples de Nagata \cite{II-26} qu'il y a
  des morphismes propres qui ne sont pas quasi-projectifs.
\end{enumerate}

\begin{proposition}[5.5.5]
\label{II.5.5.5-fr}
\begin{enumerate}
  \item[\rm{(i)}] Une immersion fermée est un morphisme projectif.
  \item[\rm{(ii)}] Si $f:X\to Y$ et $g:Y\to Z$ sont des morphismes
  projectifs, et si $Z$ est un schéma quasi-compact ou un préschéma dont
  l'espace sous-jacent est noethérien, $g\circ f$ est projectif.
  \item[\rm{(iii)}] Si $f:X\to Y$ est un $S$-morphisme projectif,
  $f_{(S')}:X_{(S')}\to Y_{(S')}$ est projectif pour toute extension
  $S'\to S$ du préschéma de base.
  \item[\rm{(iv)}] Si $f:X\to Y$ et $g:X'\to Y'$ sont des $S$-morphismes
  projectifs, il en est de même de $f\times_Sg$.
  \item[\rm{(v)}] Si $g\circ f$ est un morphisme projectif et si $g$ est
  séparé, $f$ est projectif.
  \item[\rm{(vi)}] Si $f$ est projectif, il en est de même de
  $f_{\mathrm{red}}$.
\end{enumerate}
\end{proposition}

(i) résulte aussitôt de (\hyperref[II.3.1.7-fr]{3.1.7}). Il faut ici
démontrer (iii) et (iv) séparément, à cause de la restriction introduite sur
$Z$ dans (ii) (cf. (\hyperref[I.3.5.1-fr]{I, 3.5.1})). Pour prouver (iii), on
se ramène au cas où $S=Y$ (\hyperref[I.3.3.11-fr]{I, 3.3.11}) et l'assertion
résulte alors aussitôt de (\hyperref[II.5.5.1-fr]{5.5.1, b)}) et de
(\hyperref[II.3.5.3-fr]{3.5.3}). Pour démontrer (iv), on est aussitôt ramené
au cas où $X=\mathbf{P}(\mathcal{E})$, $X'=\mathbf{P}(\mathcal{E}')$,
$\mathcal{E}$ (resp. $\mathcal{E}'$) étant un $\mathcal{O}_Y$-Module (resp.
un $\mathcal{O}_{Y'}$-Module) quasi-cohérent de type fini. Soient $p$, $p'$ les
projections canoniques de $T=Y\times_SY'$ dans $Y$ et $Y'$ respectivement~;
d'après (\hyperref[II.4.1.3.1-fr]{4.1.3.1}), on a
$\mathbf{P}(p^*(\mathcal{E}))=\mathbf{P}(\mathcal{E})\times_YT$,
$\mathbf{P}(p'^*(\mathcal{E}'))=\mathbf{P}(\mathcal{E}')\times_{Y'}T$, d'où
\begin{align*}
\mathbf{P}(p^*(\mathcal{E}))\times_T\mathbf{P}(p'^*(\mathcal{E}'))
&=(\mathbf{P}(\mathcal{E})\times_YT)\times_T
(T\times_{Y'}\mathbf{P}(\mathcal{E}'))\\
&=\mathbf{P}(\mathcal{E})\times_Y
(T\times_{Y'}\mathbf{P}(\mathcal{E}'))
=\mathbf{P}(\mathcal{E})\times_S\mathbf{P}(\mathcal{E}')
\end{align*}
en remplaçant $T$ par $Y\times_SY'$ et utilisant
(\hyperref[I.3.3.9.1-fr]{I, 3.3.9.1}). Or, $p^*(\mathcal{E})$ et
$p'^*(\mathcal{E}')$ sont de type fini sur $T$
(\hyperref[0.5.2.4-fr]{0, 5.2.4}), donc il en est de même de
$p^*(\mathcal{E})\otimes_{\mathcal{O}_T}p'^*(\mathcal{E}')$~; comme
$\mathbf{P}(p^*(\mathcal{E}))\times_T\mathbf{P}(p'^*(\mathcal{E}'))$
s'identifie à un sous-préschéma fermé de
$\mathbf{P}(p^*(\mathcal{E})\otimes_{\mathcal{O}_T}p'^*(\mathcal{E}'))$
(\hyperref[II.4.3.3-fr]{4.3.3}), cela achève de prouver (iv). Pour obtenir (v)
et (vi), on peut appliquer (\hyperref[I.5.5.13-fr]{I, 5.5.13}), car tout
sous-préschéma fermé d'un $Y$-schéma projectif est un $Y$-schéma projectif par
(\hyperref[II.5.5.1-fr]{5.5.1, a)}).

Il reste à prouver (ii)~; en vertu de l'hypothèse sur $Z$, cela résulte de
(\hyperref[II.5.5.3-fr]{5.5.3}), de
(\hyperref[II.5.3.4-fr]{5.3.4, (ii)}) et
(\hyperref[II.5.4.2-fr]{5.4.2, (ii)}).

\begin{proposition}[5.5.6]
\label{II.5.5.6-fr}
Si $X$ et $X'$ sont deux $Y$-schémas projectifs, $X\amalg X'$ est un
$Y$-schéma projectif.
\end{proposition}

C'est une conséquence évidente de (\hyperref[II.5.5.2-fr]{5.5.2}) et
(\hyperref[II.4.3.6-fr]{4.3.6}).

\begin{proposition}[5.5.7]
\label{II.5.5.7-fr}
Soit $X$ un $Y$-schéma projectif, $\mathcal{L}$ un $\mathcal{O}_X$-Module
$Y$-ample~; pour toute section $f$ de $\mathcal{L}$ au-dessus de $X$, $X_f$
est affine sur $Y$.
\end{proposition}

\oldpage[II]{106}

La question étant locale sur $Y$, on peut supposer
$Y=\operatorname{Spec}(A)$~; en outre
$X_{f^{\otimes n}}=X_f$, donc en remplaçant $\mathcal{L}$ par un
$\mathcal{L}^{\otimes n}$ convenable, on peut supposer $\mathcal{L}$ très
ample pour le morphisme structural $q:X\to Y$
(\hyperref[II.4.6.11-fr]{4.6.11}). L'homomorphisme canonique
$\sigma:q^*(q_*(\mathcal{L}))\to\mathcal{L}$ est donc surjectif et le
morphisme correspondant
\[
r=r_{\mathcal{L},\sigma}:X\to P=\mathbf{P}(q_*(\mathcal{L}))
\]
une immersion telle que $\mathcal{L}=r^*(\mathcal{O}_P(1))$
(\hyperref[II.4.4.4-fr]{4.4.4})~; en outre, comme $X$ est propre sur $Y$,
l'immersion $r$ est fermée (\hyperref[II.5.4.4-fr]{5.4.4}). Par ailleurs, on
a par définition $f\in\Gamma(Y,q_*(\mathcal{L}))$ et $\sigma^\flat$ est
l'identité de $q_*(\mathcal{L})$~; il résulte alors de la formule
(\hyperref[II.3.7.3.1-fr]{3.7.3.1}) que l'on a $X_f=r^{-1}(D_+(f))$~; $X_f$
est par suite un sous-préschéma fermé du schéma affine $D_+(f)$ et est donc
bien un schéma affine.

Lorsqu'on prend en particulier $Y=X$, on obtient (compte tenu de
(\hyperref[II.4.6.13-fr]{4.6.13, (i)})) le corollaire suivant, dont la
démonstration directe est d'ailleurs immédiate~:

\begin{corollary}[5.5.8]
\label{II.5.5.8-fr}
Soit $X$ un préschéma, $\mathcal{L}$ un $\mathcal{O}_X$-Module inversible.
Pour toute section $f$ de $\mathcal{L}$ au-dessus de $X$, $X_f$ est affine
sur $X$ (et par suite un schéma affine lorsque $X$ est un schéma affine).
\end{corollary}

\subsection{Le lemme de Chow.}

\begin{theorem}[5.6.1]
\label{II.5.6.1-fr}
(Lemme de Chow). Soit $S$ un préschéma, $X$ un $S$-schéma de type fini. On
suppose vérifiée l'une des hypothèses suivantes~:
\begin{enumerate}
  \item[\rm{a)}] $S$ est noethérien.
  \item[\rm{b)}] $S$ est un schéma quasi-compact, et $X$ a un nombre fini de
  composantes irréductibles.
\end{enumerate}
Dans ces conditions~:
\begin{enumerate}
  \item[\rm{(i)}] Il existe un $S$-schéma $X'$ quasi-projectif sur $S$ et un
  $S$-morphisme $f:X'\to X$, projectif et surjectif.
  \item[\rm{(ii)}] On peut prendre $X'$ et $f$ tels qu'il existe un ouvert
  $U\subset X$ pour lequel $U'=f^{-1}(U)$ soit dense dans $X'$ et la
  restriction de $f$ à $U'$ un isomorphisme $U'\simeq U$.
  \item[\rm{(iii)}] Si $X$ est réduit (resp. irréductible, intègre), on peut
  supposer $X'$ réduit (resp. irréductible, intègre).
\end{enumerate}
\end{theorem}

La démonstration se fait en plusieurs étapes.

A) On peut d'abord se ramener au cas où $X$ est irréductible. En effet, dans
l'hypothèse a), $X$ est noethérien, donc, dans les deux hypothèses, les
composantes irréductibles $X_i$ de $X$ sont en nombre fini. Si le théorème est
démontré pour chacun des sous-préschémas fermés réduits de $X$ ayant les $X_i$
pour espaces sous-jacents, et si $X'_i$ et $f_i:X'_i\to X_i$ sont le préschéma
et le morphisme correspondant à $X_i$, le préschéma $X'$ somme des $X'_i$, et
le morphisme $f:X'\to X$ dont la restriction à chaque $X'_i$ est $j_i\circ f_i$
($j_i$ injection canonique $X_i\to X$), répondent à la question. Il est
immédiat en effet que $X'$ est réduit si les $X'_i$ le sont~; d'autre part on
satisfait à (ii) en prenant pour $U$ la réunion des ensembles
$U_i\cap(\bigcup_{j\neq i}X_j)^{\complement}$. Enfin, comme les $X'_i$ sont
quasi-projectifs sur $S$, il en est de même de $X'$

\oldpage[II]{107}

(\hyperref[II.5.3.6-fr]{5.3.6})~; de même, les morphismes $X'_i\to X$ sont
projectifs par (\hyperref[II.5.5.5-fr]{5.5.5, (i) et (ii)}), donc $f$ est
projectif (\hyperref[II.5.5.6-fr]{5.5.6}), et est évidemment surjectif par
définition.

B) Supposons maintenant $X$ irréductible. Comme le morphisme structural
$r:X\to S$ est de type fini, il y a un recouvrement fini $(S_i)$ de $S$ par
des ouverts affines, et pour chaque $i$ un recouvrement fini $(T_{ij})$ de
$r^{-1}(S_i)$ par des ouverts affines, le morphisme $T_{ij}\to S_i$ étant
affine et de type fini, donc quasi-projectif
(\hyperref[II.5.3.4-fr]{5.3.4, (i)})~; comme dans chacune des hypothèses a),
b), l'immersion $S_i\to S$ est quasi-compacte, c'est un morphisme
quasi-projectif (\hyperref[II.5.3.4-fr]{5.3.4, (ii)}), donc la restriction de
$r$ à $T_{ij}$ est un morphisme quasi-projectif
(\hyperref[II.5.3.4-fr]{5.3.4, (iii)}). Désignons les $T_{ij}$ par $U_k$
$(1\leq k\leq n)$. Il existe, pour chaque indice $k$ une immersion ouverte
$\varphi_k:U_k\to P_k$, où $P_k$ est projectif sur $S$
(\hyperref[II.5.3.2-fr]{5.3.2} et \hyperref[II.5.5.2-fr]{5.5.2}). Soit
$U=\bigcap_kU_k$~; comme $X$ est irréductible et les $U_k$ non vides, $U$ est
non vide, et par suite partout dense dans $X$~; les restrictions à $U$ des
$\varphi_k$ définissent un morphisme
\[
\varphi:U\to P=P_1\times_SP_2\times_S\cdots\times_SP_n
\]
de sorte que les diagrammes
\[
\tag{5.6.1.1}
\label{II.5.6.1.1-fr}
\xymatrix{
U\ar[r]^\varphi\ar[d]_{j_k} & P\ar[d]^{p_k}\\
U_k\ar[r]_{\varphi_k} & P_k
}
\]
soient commutatifs, $j_k$ étant l'injection canonique $U\to U_k$ et $p_k$ la
projection canonique $P\to P_k$. Si $j$ est l'injection canonique $U\to X$,
le morphisme $\psi=(j,\varphi)_S:U\to X\times_SP$ est une immersion
(\hyperref[I.5.3.14-fr]{I, 5.3.14}). Dans l'hypothèse a), $X\times_SP$ est
localement noethérien (\hyperref[II.3.4.1-fr]{3.4.1} et
\hyperref[I.6.3.7-fr]{I, 6.3.7 et 6.3.8})~; dans l'hypothèse b),
$X\times_SP$ est un schéma quasi-compact
(\hyperref[I.5.5.1-fr]{I, 5.5.1 et 6.6.4})~; dans les deux cas l'adhérence
$X'$ dans $X\times_SP$ du sous-préschéma $Z$ associé à $\psi$ (et dont
$\psi(U)$ est l'espace sous-jacent) existe, et $\psi$ se factorise en
\[
\tag{5.6.1.2}
\label{II.5.6.1.2-fr}
\psi:U\xrightarrow{\psi'}X'\xrightarrow{h}X\times_SP
\]
où $\psi'$ est une immersion ouverte et $h$ une immersion fermée
(\hyperref[I.9.5.10-fr]{I, 9.5.10}). Soient
$q_1:X\times_SP\to X$ et $q_2:X\times_SP\to P$ les projections canoniques~;
nous poserons
\[
\tag{5.6.1.3}
\label{II.5.6.1.3-fr}
f:X'\xrightarrow{h}X\times_SP\xrightarrow{q_1}X
\]
et
\[
\tag{5.6.1.4}
\label{II.5.6.1.4-fr}
g:X'\xrightarrow{h}X\times_SP\xrightarrow{q_2}P.
\]
Nous allons voir que $X'$ et $f$ répondent à la question.

C) Montrons d'abord que $f$ est projectif et surjectif, et que la restriction
de $f$ à $U'=f^{-1}(U)$ est un isomorphisme de $U'$ sur $U$. Comme les $P_k$
sont projectifs sur $S$, il en est de même de $P$
(\hyperref[II.5.5.5-fr]{5.5.5, (iv)}), donc $X\times_SP$ est projectif sur
$X$ (\hyperref[II.5.5.5-fr]{5.5.5, (iii)}), et il en est de même de $X'$ qui
est un sous-préschéma fermé de $X\times_SP$. D'autre part, on a
$f\circ\psi'=q_1\circ(h\circ\psi')=q_1\circ\psi=j$, donc $f(X')$ contient
l'ensemble ouvert partout dense $U$ de $X$~; mais $f$ est un morphisme fermé
(\hyperref[II.5.5.3-fr]{5.5.3}), donc $f(X')=X$. Notons maintenant que
$q_1^{-1}(U)=U\times_SP$ est induit sur un ouvert de $X\times_SP$, et par
définition le préschéma $U'=h^{-1}(U\times_SP)$ est induit par $X'$ sur
l'ouvert $U'$~; c'est par suite l'adhérence par

\oldpage[II]{108}

rapport à $U\times_SP$ du préschéma $Z$
(\hyperref[I.9.5.8-fr]{I, 9.5.8}). Or l'immersion $\psi$ se factorise en
\[
U\xrightarrow{\Gamma_\varphi}U\times_SP\xrightarrow{j\times1}X\times_SP,
\]
et comme $P$ est séparé sur $S$, le morphisme graphe $\Gamma_\varphi$ est une
immersion fermée (\hyperref[I.5.4.3-fr]{I, 5.4.3}), donc $Z$ est un
sous-préschéma fermé de $U\times_SP$, d'où on conclut que $U'=Z$. Comme
d'ailleurs $\psi$ est une immersion, la restriction de $f$ à $U'$ est un
isomorphisme sur $U$, réciproque de $\psi'$~; enfin, par définition de $X'$,
$U'$ est dense dans $X'$.

D) Prouvons maintenant que $g$ est une immersion, ce qui établira que $X'$ est
quasi-projectif sur $S$, puisque $P$ est projectif sur $S$. Posons
\begin{align*}
V_k&=\varphi_k(U_k) &&\text{(ouvert de $P_k$)},\\
W_k&=p_k^{-1}(V_k) &&\text{(ouvert de $P$)},\\
U'_k&=f^{-1}(U_k),\qquad U''_k=g^{-1}(W_k)
&&\text{(ouverts de $X'$)}.
\end{align*}
Il est clair que les $U'_k$ forment un recouvrement ouvert de $X'$~; nous
allons d'abord voir qu'il en est de même des $U''_k$, en prouvant que
$U'_k\subset U''_k$. Il suffira pour cela de montrer que le diagramme
\[
\tag{5.6.1.5}
\label{II.5.6.1.5-fr}
\xymatrix{
U'_k\ar[r]^{g|U'_k}\ar[d]_{f|U'_k} & P\ar[d]^{p_k}\\
U_k\ar[r]_{\varphi_k} & P_k
}
\]
est commutatif. Or le préschéma $U'_k=h^{-1}(U_k\times_SP)$ est induit par
$X'$ sur l'ouvert $U'_k$, et est donc l'adhérence de $Z=U'\subset U'_k$ par
rapport à $U'_k$ (\hyperref[I.9.5.8-fr]{I, 9.5.8}). Pour montrer la
commutativité du diagramme (\hyperref[II.5.6.1.5-fr]{5.6.1.5}), il suffit
donc, puisque $P_k$ est un $S$-schéma, de montrer qu'en composant ce diagramme
avec l'injection canonique $U'\to U'_k$ (ou, ce qui revient au même en vertu de
l'isomorphisme de $U'$ et $U$, avec $\psi$), on obtient un diagramme commutatif
(\hyperref[I.9.5.6-fr]{I, 9.5.6}). Mais par définition le diagramme que l'on
obtient alors n'est autre que (\hyperref[II.5.6.1.1-fr]{5.6.1.1}), d'où notre
assertion.

Les $W_k$ forment donc un recouvrement ouvert de $g(X')$~; pour prouver que
$g$ est une immersion, il suffit de voir que chacune des restrictions
$g|U''_k$ est une immersion dans $W_k$
(\hyperref[I.4.2.4-fr]{I, 4.2.4}). Pour cela, considérons le morphisme
$u_k:W_k\xrightarrow{p_k}V_k\xrightarrow{\varphi_k^{-1}}U_k\to X$~; comme
$X$ est séparé sur $S$, le morphisme graphe
$\Gamma_{u_k}:W_k\to X\times_SW_k$ est une immersion fermée
(\hyperref[I.5.4.3-fr]{I, 5.4.3}), donc le graphe
$T_k=\Gamma_{u_k}(W_k)$ est un sous-préschéma fermé de $X\times_SW_k$~; si
nous montrons que ce sous-préschéma majore $U'$, il majorera aussi le
sous-préschéma induit par $X'$ sur l'ouvert $X''_k$ de $X'$, par
(\hyperref[I.9.5.8-fr]{I, 9.5.8}). Comme la restriction de $q_2$ à $T_k$ est
un isomorphisme sur $W_k$, la restriction de $g$ à $X''_k$ sera une immersion
dans $W_k$, et notre assertion sera démontrée. Soit $v_k$ l'injection
canonique $U'\to X\times_SW_k$~; nous devons montrer qu'il existe un morphisme
$w_k:U'\to W_k$ tel que $v_k=\Gamma_{u_k}\circ w_k$. Par définition du
produit, il suffit de prouver que
$q_1\circ v_k=u_k\circ q_2\circ v_k$
(\hyperref[I.3.2.1-fr]{I, 3.2.1}), ou, en composant à droite

\oldpage[II]{109}

avec l'isomorphisme $\psi':U\to U'$, que
$q_1\circ\psi=u_k\circ q_2\circ\psi$. Mais comme $q_1\circ\psi=j$,
$q_2\circ\psi=\varphi$, notre assertion résulte de la commutativité du
diagramme (\hyperref[II.5.6.1.1-fr]{5.6.1.1}), compte tenu de la définition de
$u_k$.

E) Il est clair que comme $U$, et par suite $U'$, est irréductible, il en est
de même de $X'$ dans la construction précédente, et le morphisme $f$ est donc
birationnel (\hyperref[I.2.2.9-fr]{I, 2.2.9}). Si en outre $X$ est réduit, il
en est de même de $U'$, donc $X'$ est réduit
(\hyperref[I.9.5.9-fr]{I, 9.5.9}). Cela achève la démonstration.

\begin{corollary}[5.6.2]
\label{II.5.6.2-fr}
On suppose vérifiée l'une des hypothèses a), b) de
(\hyperref[II.5.6.1-fr]{5.6.1}). Pour que $X$ soit propre sur $S$, il faut et
il suffit qu'il existe un schéma $X'$ projectif sur $S$ et un $S$-morphisme
surjectif $f:X'\to X$ (qui est alors projectif par
(\hyperref[II.5.5.5-fr]{5.5.5, (v)})). Lorsqu'il en est ainsi, on peut en
outre choisir $f$ de sorte qu'il existe un ouvert dense $U$ de $X$ tel que la
restriction de $f$ à $f^{-1}(U)$ soit un isomorphisme
$f^{-1}(U)\simeq U$ et que $f^{-1}(U)$ soit dense dans $X'$. Si en outre $X$
est irréductible (resp. réduit) on peut supposer qu'il en est de même de
$X'$~; lorsque $X$ et $X'$ sont irréductibles, $f$ est un morphisme
birationnel.
\end{corollary}

La condition est suffisante en vertu de (\hyperref[II.5.5.3-fr]{5.5.3}) et
(\hyperref[II.5.4.3-fr]{5.4.3, (ii)}). Elle est nécessaire, car avec les
notations de (\hyperref[II.5.6.1-fr]{5.6.1}), si $X$ est propre sur $S$, $X'$
est propre sur $S$, car il est projectif sur $X$, donc propre sur $X$
(\hyperref[II.5.5.3-fr]{5.5.3}), et notre assertion résulte de
(\hyperref[II.5.4.2-fr]{5.4.2, (ii)})~; en outre, comme $X'$ est
quasi-projectif sur $S$, il est projectif sur $S$ en vertu de
(\hyperref[II.5.5.3-fr]{5.5.3}).

\begin{corollary}[5.6.3]
\label{II.5.6.3-fr}
Soit $S$ un préschéma localement noethérien, $X$ un $S$-schéma de type fini
sur $S$, $f_0:X\to S$ son morphisme structural. Pour que $X$ soit propre sur
$S$, il faut et il suffit que pour tout morphisme de type fini $S'\to S$,
$(f_0)_{(S')}:X_{(S')}\to S'$ soit un morphisme fermé. Il suffit même que
cette condition soit vérifiée pour tout $S$-préschéma de la forme
$S'=S\otimes_{\mathbf Z}\mathbf Z[T_1,\ldots,T_n]$ ($T_i$ indéterminées).
\end{corollary}

La condition étant évidemment nécessaire, montrons qu'elle est suffisante.
La question étant locale sur $S$ et $S'$ (\hyperref[II.5.4.1-fr]{5.4.1}), on
peut supposer $S$ et $S'$ affines et noethériens. En vertu du lemme de Chow,
il existe un $S$-schéma projectif $P$, une immersion $j:X'\to P$ et un
morphisme projectif et surjectif $f:X'\to X$, tels que le diagramme
\[
\xymatrix{
X\ar[d]_{f_0} & X'\ar[l]_f\ar[d]^j\\
S & P\ar[l]_r
}
\]
soit commutatif. Comme $P$ est de type fini sur $S$, la première hypothèse
entraîne que la projection $q_2:X\times_SP\to P$ est un morphisme fermé. Mais
l'immersion $j$ est composée de $q_2$ et du morphisme $f\times1$ de
$X'\times_SP$ dans $X\times_SP$~; or $f$, étant projectif, est propre
(\hyperref[II.5.5.3-fr]{5.5.3}), donc $f\times1$ est fermé. On en conclut que
$j$ est une immersion fermée, et par suite est propre
(\hyperref[II.5.4.2-fr]{5.4.2, (ii)}). En outre le morphisme structural
$r:P\to S$ est projectif, donc propre (\hyperref[II.5.5.3-fr]{5.5.3}), donc
$f_0\circ f=r\circ j$ est propre
(\hyperref[II.5.4.2-fr]{5.4.2, (ii)})~; finalement, comme $f$ est surjectif,
$f_0$ est propre en vertu de (\hyperref[II.5.4.3-fr]{5.4.3}).

Pour prouver la proposition en utilisant seulement la seconde hypothèse, il
suffit de montrer qu'elle entraîne la première. Or, si $S'$ est affine et de
type fini sur $S=\operatorname{Spec}(A)$,

\oldpage[II]{110}

on a $S'=\operatorname{Spec}(A[c_1,\ldots,c_n])$
(\hyperref[I.6.3.3-fr]{I, 6.3.3}), et $S'$ est donc isomorphe à un
sous-préschéma fermé de $S''=\operatorname{Spec}(A[T_1,\ldots,T_n])$ ($T_i$
indéterminées). Dans le diagramme commutatif
\[
\xymatrix{
X\times_SS'\ar[r]^{1_X\times j}\ar[d]_{(f_0)_{(S')}} &
X\times_SS''\ar[d]^{(f_0)_{(S'')}}\\
S'\ar[r]_j & S''
}
\]
$j$ et $1_X\times j$ sont des immersions fermées
(\hyperref[I.4.3.1-fr]{I, 4.3.1}) et $(f_0)_{(S'')}$ est fermé par hypothèse~;
il en est donc de même de $(f_0)_{(S')}$. 

\section{Morphismes entiers et morphismes finis}

\subsection{Préschémas entiers sur un autre.}

\begin{definition}[6.1.1]
\label{II.6.1.1-fr}
Soit $X$ un $S$-préschéma, $f:X\to S$ son morphisme structural. On dit que
$X$ est \emph{entier au-dessus de $S$}, ou que $f$ est un \emph{morphisme
entier}, s'il existe un recouvrement $(S_\alpha)$ de $S$ par des ouverts
affines tels que, pour tout $\alpha$, le préschéma induit $f^{-1}(S_\alpha)$
soit un schéma affine, dont l'anneau $B_\alpha$ est une algèbre entière sur
l'anneau $A_\alpha$ de $S_\alpha$. On dit que $X$ est \emph{fini au-dessus de
$S$}, ou que $f$ est un \emph{morphisme fini} si $X$ est entier sur $S$ et de
type fini sur $S$.
\end{definition}

Lorsque $S$ est affine d'anneau $A$, on dira aussi «~entier (resp. fini) sur
$A$~» au lieu de «~entier (resp. fini) sur $S$~».

\begin{env}[6.1.2]
\label{II.6.1.2-fr}
Il est clair que si $X$ est entier au-dessus de $S$, il est affine au-dessus
de $S$. Pour qu'un préschéma $X$ affine au-dessus de $S$ soit entier (resp.
fini) au-dessus de $S$, il faut et il suffit que la $\mathcal{O}_S$-Algèbre
quasi-cohérente associée $\mathcal{A}(X)$ soit telle qu'il existe un
recouvrement $(S_\alpha)$ de $S$ par des ouverts affines ayant la propriété
que pour tout $\alpha$, $\Gamma(S_\alpha,\mathcal{A}(X))$ est une algèbre
entière (resp. entière et de type fini) sur
$\Gamma(S_\alpha,\mathcal{O}_S)$. Une $\mathcal{O}_S$-Algèbre quasi-cohérente
ayant cette propriété est dite \emph{entière} (resp. \emph{finie}) sur
$\mathcal{O}_S$. Se donner un préschéma entier (resp. fini) au-dessus de $S$
revient donc (\hyperref[II.1.3.1-fr]{1.3.1}) à se donner une
$\mathcal{O}_S$-Algèbre quasi-cohérente entière (resp. finie) sur
$\mathcal{O}_S$. On notera qu'une $\mathcal{O}_S$-Algèbre quasi-cohérente
$\mathcal{B}$ est finie si et seulement si c'est un $\mathcal{O}_S$-Module de
type fini (\hyperref[I.1.3.9-fr]{I, 1.3.9})~; il revient au même de dire que
$\mathcal{B}$ est une $\mathcal{O}_S$-Algèbre entière et de type fini, car une
algèbre entière et de type fini sur un anneau $A$ est un $A$-module de type
fini.
\end{env}

\begin{proposition}[6.1.3]
\label{II.6.1.3-fr}
Soit $S$ un préschéma localement noethérien. Pour qu'un préschéma $X$ affine
au-dessus de $S$ soit fini au-dessus de $S$, il faut et il suffit que la
$\mathcal{O}_S$-Algèbre $\mathcal{A}(X)$ soit cohérente.
\end{proposition}

Compte tenu de la remarque précédente, cela revient à remarquer que si $S$ est
localement noethérien, les $\mathcal{O}_S$-Modules quasi-cohérents de type fini
ne sont autres que les $\mathcal{O}_S$-Modules cohérents
(\hyperref[I.1.5.1-fr]{I, 1.5.1}).

\begin{proposition}[6.1.4]
\label{II.6.1.4-fr}
Soit $X$ un préschéma entier (resp. fini) au-dessus de $S$, $f:X\to S$ le
morphisme structural. Alors, pour tout ouvert affine $U\subset S$, d'anneau
$A$, $f^{-1}(U)$ est un schéma affine dont l'anneau $B$ est une algèbre
entière (resp. finie) sur $A$.
\end{proposition}

\oldpage[II]{111}

Nous démontrerons d'abord le lemme suivant~:

\begin{lemma}[6.1.4.1]
\label{II.6.1.4.1-fr}
Soient $A$ un anneau, $M$ un $A$-module, $(g_i)_{1\leq i\leq m}$ un système
fini d'éléments de $A$ tels que les $D(g_i)$ ($1\leq i\leq m$) recouvrent
$\operatorname{Spec}(A)$. Si, pour tout $i$, $M_{g_i}$ est un
$A_{g_i}$-module de type fini, alors $M$ est un $A$-module de type fini.
\end{lemma}

On peut en effet supposer que $M_{g_i}$ admet un système de générateurs fini
$(m_{ij}/g_i^n)$ avec $m_{ij}\in M$, $n$ étant le même pour tous les indices
$i$. Montrons que les $m_{ij}$ forment un système de générateurs de $M$. Soit
$M'$ le sous-$A$-module de $M$ engendré par ces éléments, et soit $m$ un
élément de $M$. Par hypothèse, pour chaque $i$, il existe des $a_{ij}\in A$ et
un entier $p$ (indépendant de $i$) tels que, dans $M_{g_i}$,
$m/1=(\sum_j a_{ij}m_{ij})/g_i^p$~; cela entraîne qu'il existe un entier
$r\geq p$ tel que, pour tout $i$, on ait $g_i^r m\in M'$. Or, comme les
$D(g_i^r)=D(g_i)$ recouvrent $\operatorname{Spec}(A)$, l'idéal de $A$
engendré par les $g_i^r$ est égal à $A$, autrement dit il existe des éléments
$a_i\in A$ tels que $\sum_i a_i g_i^r=1$~; par suite
$m=(\sum_i a_i g_i^r)m\in M'$, d'où le lemme.

Cela étant, on sait déjà (\hyperref[II.1.3.2-fr]{1.3.2}) que $f^{-1}(U)$ est
affine. Si $\varphi$ est l'homomorphisme $A\to B$ correspondant à $f$, il
existe un recouvrement fini de $U$ par des ouverts $D(g_i)$ ($g_i\in A$) tels
que, si $h_i=\varphi(g_i)$, $B_{h_i}$ soit une algèbre entière (resp. entière
finie) sur $A_{g_i}$. En effet, il y a un recouvrement de $U$ par des ouverts
affines $V_\alpha\subset U$ tels que si
$A_\alpha=A(V_\alpha)$, $B_\alpha=A(f^{-1}(V_\alpha))$ soit une algèbre
entière (resp. finie) sur $A_\alpha$. Tout $x\in U$ appartient à un
$V_\alpha$, donc il existe $g\in A$ tel que $x\in D(g)\subset V_\alpha$~; si
$g_\alpha$ est l'image de $g$ dans $A_\alpha$, on a
$A(D(g))=A_g=(A_\alpha)_{g_\alpha}$~; soit $h=\varphi(g)$, et soit $h_\alpha$
l'image de $g_\alpha$ dans $B_\alpha$~; on a
\[
A(D(h))=B_h=(B_\alpha)_{h_\alpha}
\]
et comme $B_\alpha$ est entière (resp. finie) sur $A_\alpha$,
$(B_\alpha)_{h_\alpha}$ est entière (resp. finie) sur
$(A_\alpha)_{g_\alpha}$. Il suffit maintenant d'utiliser le fait que $U$ est
quasi-compact pour obtenir le recouvrement cherché.

Si on suppose d'abord que les $B_{h_i}$ sont entières finies sur les
$A_{g_i}$, comme, en tant que $A_{g_i}$-module, $B_{h_i}$ s'écrit aussi
$B_{g_i}$, le lemme (\hyperref[II.6.1.4.1-fr]{6.1.4.1}) montre que dans ce cas
$B$ est un $A$-module de type fini.

Supposons seulement maintenant chaque $B_{h_i}$ entière sur $A_{g_i}$~; soit
$b\in B$, et soit $C$ la sous-$A$-algèbre de $B$ engendrée par $b$. Pour tout
$i$, $C_{h_i}$ est l'algèbre sur $A_{g_i}$ engendrée par $b/1$ dans
$B_{h_i}$~; il résulte de l'hypothèse que chaque $C_{h_i}$ est un
$A_{g_i}$-module de type fini, donc
(\hyperref[II.6.1.4.1-fr]{6.1.4.1}) $C$ est un $A$-module de type fini, ce qui
prouve que $B$ est entière sur $A$.

\begin{proposition}[6.1.5]
\label{II.6.1.5-fr}
\begin{enumerate}
\item[(i)] Une immersion fermée est finie (et a fortiori entière).
\item[(ii)] Le composé de deux morphismes finis (resp. entiers) est fini
(resp. entier).
\item[(iii)] Si $f:X\to Y$ est un $S$-morphisme fini (resp. entier),
$f_{(S')}:X_{(S')}\to Y_{(S')}$ est fini (resp. entier) pour toute extension
$S'\to S$ de la base.
\item[(iv)] Si $f:X\to Y$ et $g:X'\to Y'$ sont deux $S$-morphismes finis
(resp. entiers), il en est de même de
$f\times_S g:X\times_S X'\to Y\times_S Y'$.
\item[(v)] Si $f:X\to Y$ et $g:Y\to Z$ sont deux morphismes tels que $g\circ
f$ soit fini (resp. entier) et $g$ séparé, alors $f$ est fini (resp. entier).
\item[(vi)] Si $f:X\to Y$ est un morphisme fini (resp. entier), il en est de
même de $f_{\mathrm{red}}$.
\end{enumerate}
\end{proposition}

\oldpage[II]{112}

En vertu de (\hyperref[I.5.5.12-fr]{I, 5.5.12}), il suffit de démontrer (i),
(ii) et (iii). Pour prouver qu'une immersion fermée $X\to S$ est finie, on
peut se borner au cas où $S=\operatorname{Spec}(A)$, et tout revient alors à
remarquer qu'un anneau quotient $A/\mathfrak{J}$ est un $A$-module monogène.
Pour prouver que le composé de deux morphismes finis (resp. entiers) $X\to Y$,
$Y\to Z$ est fini (resp. entier), on peut encore supposer $Z$ (et par suite
$X$, $Y$ (\hyperref[II.1.3.4-fr]{1.3.4})) affines, et alors l'assertion
équivaut à dire que si $B$ est une $A$-algèbre finie (resp. entière) et $C$
une $B$-algèbre finie (resp. entière), alors $C$ est une $A$-algèbre finie
(resp. entière), ce qui est immédiat. Enfin, pour démontrer (iii), on peut se
borner au cas où $S=Y$, puisque $X_{(S')}$ s'identifie à
$X\times_Y Y_{(S')}$ (\hyperref[I.3.3.11-fr]{I, 3.3.11})~; on peut en outre
supposer que $S=\operatorname{Spec}(A)$, $S'=\operatorname{Spec}(A')$~; alors
$X$ est affine d'anneau $B$ (\hyperref[II.1.3.4-fr]{1.3.4}), $X_{(S')}$
affine d'anneau $A'\otimes_A B$, et il suffit de remarquer que si $B$ est une
$A$-algèbre finie (resp. entière), $A'\otimes_A B$ est une $A'$-algèbre finie
(resp. entière).

Notons aussi que si $X$ et $Y$ sont deux $S$-préschémas qui sont finis (resp.
entiers) au-dessus de $S$, leur \emph{somme} $X\amalg Y$ est un préschéma fini
(resp. entier) au-dessus de $S$, car cela revient à voir que si $B$ et $C$
sont deux $A$-algèbres finies (resp. entières) sur $A$, il en est de même de
$B\times C$.

\begin{corollary}[6.1.6]
\label{II.6.1.6-fr}
Si $X$ est un préschéma entier (resp. fini) au-dessus de $S$, alors, pour tout
ouvert $U\subset S$, $f^{-1}(U)$ est entier (resp. fini) au-dessus de $U$.
\end{corollary}

C'est un cas particulier de (\hyperref[II.6.1.5-fr]{6.1.5, (iii)}).

\begin{corollary}[6.1.7]
\label{II.6.1.7-fr}
Soit $f:X\to Y$ un morphisme fini. Alors, pour tout $y\in Y$, la fibre
$f^{-1}(y)$ est un schéma algébrique fini sur $k(y)$, et a fortiori son
espace sous-jacent est discret et fini.
\end{corollary}

En effet, en tant que $k(y)$-préschéma, $f^{-1}(y)$ est identifié à
$X\times_Y\operatorname{Spec}(k(y))$ (\hyperref[I.3.6.1-fr]{I, 3.6.1}), qui
est fini au-dessus de $\operatorname{Spec}(k(y))$
(\hyperref[II.6.1.5-fr]{6.1.5, (iii)})~; c'est donc un schéma affine dont
l'anneau est une algèbre de rang fini sur $k(y)$
(\hyperref[II.6.1.4-fr]{6.1.4}). La proposition résulte alors de
(\hyperref[I.6.4.4-fr]{I, 6.4.4}).

\begin{corollary}[6.1.8]
\label{II.6.1.8-fr}
Soient $X$, $S$ deux préschémas intègres, $f:X\to S$ un morphisme dominant. Si
$f$ est entier (resp. fini), le corps $R(X)$ des fonctions rationnelles sur
$X$ est algébrique (resp. algébrique de degré fini) sur le corps $R(S)$ des
fonctions rationnelles sur $S$.
\end{corollary}

En effet, soit $s$ le point générique de $S$~; le $k(s)$-préschéma
$f^{-1}(s)$ est entier (resp. fini) au-dessus de
$\operatorname{Spec}(k(s))$ (\hyperref[II.6.1.5-fr]{6.1.5, (iii)}), et
contient par hypothèse le point générique $x$ de $X$~; l'anneau local de $x$
dans $f^{-1}(s)$ étant égal à $k(x)$ (\hyperref[I.3.6.5-fr]{I, 3.6.5}) est un
anneau local d'une algèbre entière (resp. finie) sur $k(s)$
(\hyperref[II.6.1.4-fr]{6.1.4}), d'où le corollaire.

\begin{remark}[6.1.9]
\label{II.6.1.9-fr}
L'hypothèse que $g$ est \emph{séparé} est essentielle pour la validité de
(\hyperref[II.6.1.5-fr]{6.1.5, (v)})~: en effet, si $Y$ n'est pas séparé sur
$Z$, l'identité $1_Y$ est le morphisme composé
$Y\xrightarrow{\Delta_Y}Y\times_ZY\xrightarrow{p_1}Y$, mais $\Delta_Y$ n'est
pas un morphisme entier, comme il résulte de
(\hyperref[II.6.1.10-fr]{6.1.10})~:
\end{remark}

\begin{proposition}[6.1.10]
\label{II.6.1.10-fr}
Tout morphisme entier est universellement fermé.
\end{proposition}

Soit $f:X\to Y$ un morphisme entier~; en vertu de
(\hyperref[II.6.1.5-fr]{6.1.5, (iii)}), il suffit de montrer que $f$ est
fermé. Soit $Z$ une partie fermée de $X$~; il existe un sous-préschéma de $X$
ayant

\oldpage[II]{113}

pour espace sous-jacent $Z$ (\hyperref[I.5.2.1-fr]{I, 5.2.1}), et il résulte
donc de (\hyperref[II.6.1.5-fr]{6.1.5, (i) et (ii)}) qu'on peut se borner à
prouver que $f(X)$ est fermé dans $Y$. En vertu de
(\hyperref[II.6.1.5-fr]{6.1.5, (vii)}), on peut supposer $X$ et $Y$ réduits~;
en outre, si $T$ est le sous-préschéma fermé réduit de $Y$ ayant pour espace
sous-jacent $\overline{f(X)}$ (\hyperref[I.5.2.1-fr]{I, 5.2.1}), on sait que
$f$ se factorise en $X\xrightarrow{g}T\xrightarrow{j}Y$, où $j$ est le
morphisme d'injection (\hyperref[I.5.2.2-fr]{I, 5.2.2}), et comme $j$ est
séparé (\hyperref[I.5.5.1-fr]{I, 5.5.1, (i)}), il résulte de
(\hyperref[II.6.1.5-fr]{6.1.5, (v)}) que $g$ est un morphisme entier. On peut
donc supposer que $f(X)$ est dense dans $Y$. Enfin, la question étant locale
sur $Y$, on peut se borner au cas où $Y=\operatorname{Spec}(A)$. Alors
$X=\operatorname{Spec}(B)$, où $B$ est une $A$-algèbre entière sur $A$
(\hyperref[II.6.1.4-fr]{6.1.4})~; en outre $A$ est réduit
(\hyperref[I.5.1.4-fr]{I, 5.1.4}), et l'hypothèse que $f(X)$ est dense dans
$Y$ entraîne que l'homomorphisme $\varphi:A\to B$ correspondant à $f$ est
injectif (\hyperref[I.1.2.7-fr]{I, 1.2.7}). Dire que dans ces conditions
$f(X)=Y$ signifie que tout idéal premier de $A$ est la trace sur $A$ d'un
idéal premier de $B$, ce qui n'est autre que le premier théorème de
Cohen-Seidenberg (\cite[t.~I, p.~257, th.~3]{I-13}).

\begin{corollary}[6.1.11]
\label{II.6.1.11-fr}
Tout morphisme fini $f:X\to Y$ est projectif.
\end{corollary}

En effet, comme $f$ est affine, $\mathcal{O}_X$ est un
$\mathcal{O}_X$-Module très ample relativement à $f$
(\hyperref[II.5.1.2-fr]{5.1.2})~; en outre $f_*(\mathcal{O}_X)$ est un
$\mathcal{O}_Y$-Module quasi-cohérent de type fini
(\hyperref[II.6.1.2-fr]{6.1.2})~; enfin $f$ est séparé, de type fini et
universellement fermé (\hyperref[II.6.1.10-fr]{6.1.10}), et on est donc dans
les conditions d'applications du critère
(\hyperref[II.5.5.4-fr]{5.5.4, (ii)}).

\begin{proposition}[6.1.12]
\label{II.6.1.12-fr}
Soient $f:X'\to X$ un morphisme fini, et soit
$\mathcal{B}=f_*(\mathcal{O}_{X'})$ (qui est une $\mathcal{O}_X$-Algèbre
quasi-cohérente et un $\mathcal{O}_X$-Module de type fini). Soit
$\mathcal{F}'$ un $\mathcal{O}_{X'}$-Module quasi-cohérent~; pour que
$\mathcal{F}'$ soit localement libre, de rang $r$, il faut et il suffit que
$f_*(\mathcal{F}')$ soit un $\mathcal{B}$-Module localement libre de rang
$r$.
\end{proposition}

Il est clair que si $f_*(\mathcal{F}')|U$ est isomorphe à
$\mathcal{B}^r|U$ ($U$ ouvert de $X$), $\mathcal{F}'|f^{-1}(U)$ est isomorphe
à $\mathcal{O}_{X'}^r|f^{-1}(U)$
(\hyperref[II.1.4.2-fr]{1.4.2}). Inversement, supposons que $\mathcal{F}'$
soit localement libre de rang $r$ et montrons que $f_*(\mathcal{F}')$ est
localement isomorphe à $\mathcal{B}^r$ en tant que $\mathcal{B}$-Module. Soit
$x$ un point de $X$~; lorsque $U$ parcourt un système fondamental de
voisinages affines de $x$, $f^{-1}(U)$ parcourt un système fondamental de
voisinages affines (\hyperref[II.1.2.5-fr]{1.2.5}) de l'ensemble fini
$f^{-1}(x)$, puisque $f$ est fermé
(\hyperref[II.6.1.10-fr]{6.1.10}). La proposition résulte alors du lemme
suivant~:

\begin{lemma}[6.1.12.1]
\label{II.6.1.12.1-fr}
Soient $Y$ un préschéma, $\mathcal{E}$ un $\mathcal{O}_Y$-Module localement
libre de rang $r$, $Z$ une partie finie de $Y$ contenue dans un ouvert affine
$V$. Il existe alors un voisinage $U\subset V$ de $Z$ tel que
$\mathcal{E}|U$ soit isomorphe à $\mathcal{O}_Y^r|U$.
\end{lemma}

On peut évidemment supposer $Y$ affine~; pour tout $z_i\in Z$, il existe dans
l'adhérence $\overline{\{z_i\}}$ au moins un point fermé $z'_i$
(\hyperref[0.2.1.3-fr]{0, 2.1.3})~; si $Z'$ est l'ensemble des $z'_i$, tout
voisinage de $Z'$ est un voisinage de $Z$, et on peut donc supposer $Z$ fermé
dans $Y$ et discret. Considérons le sous-préschéma fermé réduit de $Y$ ayant
$Z$ pour espace sous-jacent (\hyperref[I.5.2.1-fr]{I, 5.2.1}) et soit
$j:Z\to Y$ l'injection canonique~; $j^*(\mathcal{E})=\mathcal{E}\otimes_Y
\mathcal{O}_Z$ est localement libre de rang $r$ sur le schéma discret $Z$,
et par suite est isomorphe à $\mathcal{O}_Z^r$~; autrement dit il existe $r$
sections $s_i$ ($1\leq i\leq r$) de $\mathcal{E}\otimes_Y\mathcal{O}_Z$
au-dessus de $Z$ telles que l'homomorphisme
$\mathcal{O}_Z^r\to\mathcal{E}\otimes_Y\mathcal{O}_Z$ défini par ces sections
soit bijectif. Mais $Y=\operatorname{Spec}(A)$ est affine, $Z$ est défini par
un idéal $\mathfrak{J}$ de $A$ et on a $\mathcal{E}=\widetilde{M}$, où $M$
est un $A$-module~; les $s_i$ sont des éléments

\oldpage[II]{114}

de $M\otimes_A(A/\mathfrak{J})$ et sont donc images de $r$ éléments
$t_i\in M=\Gamma(Y,\mathcal{E})$. Pour tout $z_j\in Z$ il y a alors un
voisinage $V_j$ de $z_j$ tel que les restrictions des $t_i$ à $V_j$
définissent un isomorphisme
$\mathcal{O}_Y^r|V_j\to\mathcal{E}|V_j$
(\hyperref[0.5.5.4-fr]{0, 5.5.4})~; le voisinage $U$ réunion des $V_j$ répond
donc à la question.

\begin{proposition}[6.1.13]
\label{II.6.1.13-fr}
Soient $g:X'\to X$ un morphisme entier de préschémas, $Y$ un préschéma normal
localement intègre, $f$ une application rationnelle de $Y$ dans $X'$ telle
que $g\circ f$ soit une application rationnelle partout définie
(\hyperref[I.7.2.1-fr]{I, 7.2.1})~; alors $f$ est partout définie.
\end{proposition}

Si $f_1$ et $f_2$ sont deux morphismes (d'ouverts denses de $Y$ dans $X'$) de
la classe $f$, il est clair que $g\circ f_1$ et $g\circ f_2$ sont des
morphismes équivalents, ce qui justifie la notation $g\circ f$ pour leur
classe d'équivalence. On rappelle aussi que lorsqu'on suppose en outre que
$Y$ est localement noethérien, l'hypothèse que $Y$ est normal entraîne déjà
que $Y$ est localement intègre (\hyperref[I.6.1.13-fr]{I, 6.1.13}).

Pour démontrer (\hyperref[II.6.1.13-fr]{6.1.13}), notons d'abord que la
question est locale sur $Y$ et on peut par suite supposer qu'il existe dans
la classe de $g\circ f$ un morphisme $h:Y\to X$. Considérons l'image
réciproque $Y'=X'_{(h)}=X'_{(Y)}$, et notons que le morphisme
$g'=g_{(Y)}:Y'\to Y$ est entier
(\hyperref[II.6.1.5-fr]{6.1.5, (iii)}). Vu la correspondance entre
applications rationnelles de $Y$ dans $X'$ et $Y$-sections rationnelles de
$Y'$ (\hyperref[I.7.1.2-fr]{I, 7.1.2}), on voit qu'on est ramené à prouver le
cas particulier de (\hyperref[II.6.1.13-fr]{6.1.13}) où $X=Y$, autrement dit
le

\begin{corollary}[6.1.14]
\label{II.6.1.14-fr}
Soient $X$ un préschéma normal localement intègre, $g:X'\to X$ un morphisme
entier, $f$ une $X$-section rationnelle de $X'$. Alors $f$ est partout
définie.
\end{corollary}

La question étant locale sur $X$, on peut supposer $X$ intègre, et $f$
s'identifie alors à un morphisme d'un ouvert $U$ de $X$ dans $X'$
(\hyperref[I.7.2.2-fr]{I, 7.2.2}) qui est une $U$-section de $g^{-1}(U)$.
Comme $g$ est séparé, $f$ est une immersion fermée de $U$ dans $g^{-1}(U)$
(\hyperref[I.5.4.6-fr]{I, 5.4.6})~; soit $Z$ le sous-préschéma fermé de
$g^{-1}(U)$ associé à $f$ (\hyperref[I.4.2.1-fr]{I, 4.2.1}), qui est
isomorphe à $U$, donc intègre~; soit $X_1$ le sous-préschéma réduit de $X'$
ayant pour espace sous-jacent l'adhérence $\overline{Z}$ de $Z$ dans $X'$
(\hyperref[I.5.2.1-fr]{I, 5.2.1})~; $Z$ est alors un sous-préschéma induit
sur un ouvert de $X_1$ (\hyperref[I.5.2.3-fr]{I, 5.2.3}) et comme il est
irréductible, il en est de même de $X_1$, qui est par suite intègre. Le
morphisme $f$ peut alors être considéré comme une $X$-section rationnelle de
$X_1$~; comme la restriction de $g$ à $X_1$ est un morphisme entier
(\hyperref[II.6.1.5-fr]{6.1.5, (i) et (ii)}), on est finalement ramené à
prouver (\hyperref[II.6.1.14-fr]{6.1.14}) dans le cas particulier où
$X'=X_1$, autrement dit, le

\begin{corollary}[6.1.15]
\label{II.6.1.15-fr}
Soient $X$ un préschéma intègre normal, $X'$ un préschéma intègre,
$g:X'\to X$ un morphisme entier. S'il existe une $X$-section rationnelle $f$
de $X'$, $g$ est un isomorphisme.
\end{corollary}

La question étant locale sur $X$, on peut supposer $X$ affine d'anneau
intègre $A$, et alors $X'$ est affine d'anneau $A'$ entier sur $A$
(\hyperref[II.6.1.4-fr]{6.1.4}) et intègre~; en outre, le raisonnement de
(\hyperref[II.6.1.14-fr]{6.1.14}) montre qu'il y a un ouvert dense dans $X$
isomorphe à un ouvert dense de $X'$, donc $A$ et $A'$ ont même corps de
fractions. D'autre part, en vertu de
(\hyperref[I.8.2.1.1-fr]{I, 8.2.1.1}) et de l'hypothèse que les
$\mathcal{O}_x$ sont intégralement clos, l'anneau $A$ est intégralement clos,
donc $A'=A$, ce qui achève la démonstration de
(\hyperref[II.6.1.13-fr]{6.1.13}).

\subsection{Morphismes quasi-finis.}

\begin{proposition}[6.2.1]
\label{II.6.2.1-fr}
Soient $f:X\to Y$ un morphisme localement de type fini, $x$ un point de $X$.
Les conditions suivantes sont équivalentes~:

\oldpage[II]{115}

\begin{enumerate}
\item[a)] Le point $x$ est isolé dans sa fibre $f^{-1}(f(x))$.
\item[b)] L'anneau $\mathcal{O}_x$ est un $\mathcal{O}_{f(x)}$-module
quasi-fini (\hyperref[0.7.4.1-fr]{0, 7.4.1}).
\end{enumerate}
\end{proposition}

La question étant évidemment locale sur $X$ et sur $Y$, on peut supposer
$X=\operatorname{Spec}(A)$ et $Y=\operatorname{Spec}(B)$ affines, $A$ étant
une $B$-algèbre de type fini (\hyperref[I.6.3.3-fr]{I, 6.3.3}). En outre, on
peut remplacer $X$ par $X\times_Y\operatorname{Spec}(\mathcal{O}_{f(x)})$
sans changer la fibre $f^{-1}(f(x))$ ni l'anneau local $\mathcal{O}_x$
(\hyperref[I.3.6.5-fr]{I, 3.6.5})~; on peut donc supposer que $B$ est un
anneau local (égal à $\mathcal{O}_{f(x)}$)~; si $\mathfrak{n}$ est l'idéal
maximal de $B$, $f^{-1}(f(x))$ est alors un schéma affine d'anneau
$A/\mathfrak{n}A$, de type fini sur $k(f(x))=B/\mathfrak{n}$
(\hyperref[I.6.4.11-fr]{I, 6.4.11}). Cela étant, si a) est vérifiée, on peut
en outre supposer que $f^{-1}(f(x))$ est réduite au point $x$~; donc
$A/\mathfrak{n}A$ est de rang fini sur $B/\mathfrak{n}$
(\hyperref[I.6.4.4-fr]{I, 6.4.4}), autrement dit $A$ est un $B$-module
quasi-fini. Inversement, s'il en est ainsi, $f^{-1}(f(x))$ est un schéma
affine artinien, donc discret (\hyperref[I.6.4.4-fr]{I, 6.4.4})~; par suite
$x$ est isolé dans sa fibre, et cela montre que b) entraîne a).

\begin{corollary}[6.2.2]
\label{II.6.2.2-fr}
Soit $f:X\to Y$ un morphisme de type fini. Les conditions suivantes sont
équivalentes~:
\begin{enumerate}
\item[a)] Tout point $x\in X$ est isolé dans sa fibre $f^{-1}(f(x))$
(autrement dit le sous-espace $f^{-1}(f(x))$ est discret).
\item[b)] Pour tout $x\in X$, le préschéma $f^{-1}(f(x))$ est un
$k(f(x))$-préschéma fini.
\item[c)] Pour tout $x\in X$, l'anneau $\mathcal{O}_x$ est un
$\mathcal{O}_{f(x)}$-module quasi-fini.
\end{enumerate}
\end{corollary}

L'équivalence de a) et c) résulte de
(\hyperref[II.6.2.1-fr]{6.2.1}). D'autre part, comme $f^{-1}(f(x))$ est un
$k(f(x))$-préschéma algébrique (\hyperref[I.6.4.11-fr]{I, 6.4.11}),
l'équivalence de a) et b) résulte de (\hyperref[I.6.4.4-fr]{I, 6.4.4}).

\begin{definition}[6.2.3]
\label{II.6.2.3-fr}
Lorsque $f:X\to Y$ est un morphisme de type fini vérifiant les conditions
équivalentes de (\hyperref[II.6.2.2-fr]{6.2.2}), on dit que $f$ est
\emph{quasi-fini}, ou que $X$ est quasi-fini sur $Y$.
\end{definition}

Il est clair que tout morphisme fini est quasi-fini
(\hyperref[II.6.1.8-fr]{6.1.8}).

\begin{proposition}[6.2.4]
\label{II.6.2.4-fr}
\begin{enumerate}
\item[(i)] Une immersion $X\to Y$, qui est fermée, ou telle que $X$ soit
noethérien, est un morphisme quasi-fini.
\item[(ii)] Si $f:X\to Y$ et $g:Y\to Z$ sont des morphismes quasi-finis,
$g\circ f$ est quasi-fini.
\item[(iii)] Si $X$ et $Y$ sont des $S$-préschémas, $f:X\to Y$ un
$S$-morphisme quasi-fini, alors $f_{(S')}:X_{(S')}\to Y_{(S')}$ est quasi-fini
pour toute extension $g:S'\to S$ de la base.
\item[(iv)] Si $f:X\to Y$ et $g:X'\to Y'$ sont deux $S$-morphismes
quasi-finis,
\[
f\times_S g:X\times_S Y\to X'\times_S Y'
\]
est quasi-fini.
\item[(v)] Soient $f:X\to Y$ et $g:Y\to Z$ deux morphismes tels que $g\circ
f$ soit quasi-fini~; alors, si en outre $g$ est séparé, ou $X$ noethérien, ou
$X\times_ZY$ localement noethérien, $f$ est quasi-fini.
\item[(vi)] Si $f$ est quasi-fini, il en est de même de $f_{\mathrm{red}}$.
\end{enumerate}
\end{proposition}

Si $f:X\to Y$ est une immersion, toute fibre est réduite à un point, et
l'assertion (i) résulte donc de (\hyperref[I.6.3.4-fr]{I, 6.3.4 (i) et
6.3.5}). Pour prouver (ii), on remarque d'abord que $h=g\circ f$ est de type
fini (\hyperref[I.6.3.4-fr]{I, 6.3.4 (ii)})~; en outre, si $z=h(x)$ et
$y=f(x)$, $y$ est isolé dans $g^{-1}(z)$, donc il existe un voisinage ouvert
$V$ de $y$ dans $Y$ ne rencontrant aucun point de $g^{-1}(z)$ distinct de
$y$~; $f^{-1}(V)$ est donc un voisinage ouvert de $x$ ne rencontrant aucun

\oldpage[II]{116}

des $f^{-1}(y')$, où $y'\neq y$ est dans $g^{-1}(z)$~; comme $x$ est isolé
dans $f^{-1}(y)$, il est isolé dans
$h^{-1}(z)=f^{-1}(g^{-1}(z))$. Pour prouver (iii), on peut se limiter au cas
où $Y=S$ (\hyperref[I.3.3.11-fr]{I, 3.3.11})~; on remarque encore que tout
d'abord $f'=f_{(S')}$ est de type fini
(\hyperref[I.6.3.4-fr]{I, 6.3.4, (iii)})~; d'autre part, si
$x'\in X'=X_{(S')}$ et si on pose $y'=f'(x')$ et $y=g(y')$,
$f'^{-1}(y')$ s'identifie à
$f^{-1}(y)\otimes_{k(y)}k(y')$ (\hyperref[I.3.6.5-fr]{I, 3.6.5})~; comme
$f^{-1}(y)$ est de rang fini sur $k(y)$ par hypothèse, $f'^{-1}(y')$ est de
rang fini sur $k(y')$, donc discret. Les assertions (iv), (v), (vi) se
déduisent de (i), (ii), (iii) par la méthode générale
(\hyperref[I.5.5.12-fr]{I, 5.5.12}), sauf lorsque les hypothèses dans (v) sont
autres que l'hypothèse «~$g$ séparé~»~; pour traiter ce cas, on remarque
d'abord que si $x$ est isolé dans
$f^{-1}(g^{-1}(g(f(x))))$, il l'est a fortiori dans $f^{-1}(f(x))$~; le fait
que $f$ est de type fini résulte alors de
(\hyperref[I.6.3.6-fr]{I, 6.3.6}).

\begin{proposition}[6.2.5]
\label{II.6.2.5-fr}
Soient $A$ un anneau local noethérien complet, $X$ un $A$-schéma localement
de type fini, $x$ un point de $X$ au-dessus du point fermé $y$ de
$Y=\operatorname{Spec}(A)$, et supposons $x$ isolé dans sa fibre
$f^{-1}(y)$ ($f$ étant le morphisme structural $X\to Y$). Alors
$\mathcal{O}_x$ est un $A$-module de type fini et $X$ est $Y$-isomorphe à la
somme (I, 3.1) de
$X'=\operatorname{Spec}(\mathcal{O}_x)$ (qui est un $Y$-schéma fini) et d'un
$A$-schéma $X''$.
\end{proposition}

Il résulte de (\hyperref[II.6.2.1-fr]{6.2.1}) que $\mathcal{O}_x$ est un
$A$-module quasi-fini. Comme $\mathcal{O}_x$ est noethérien
(\hyperref[I.6.3.7-fr]{I, 6.3.7}) et l'homomorphisme
$A\to\mathcal{O}_x$ local, l'hypothèse que $A$ est complet entraîne que
$\mathcal{O}_x$ est un $A$-module de type fini
(\hyperref[0.7.4.3-fr]{0, 7.4.3}). Soit
$X'=\operatorname{Spec}(\mathcal{O}_x)$ le schéma local de $X$ au point $x$
(\hyperref[I.2.4.1-fr]{I, 2.4.1}), et soit $g:X'\to X$ le morphisme
canonique. Comme le composé $X'\xrightarrow{g}X\xrightarrow{f}Y$ est fini
(\hyperref[II.6.1.1-fr]{6.1.1}) et que $f$ est séparé, $g$ est fini
(\hyperref[II.6.1.5-fr]{6.1.5, (v)}), donc $g(X')$ est fermé dans $X$
(\hyperref[II.6.1.10-fr]{6.1.10})~; d'autre part, comme $g$ est de type fini,
$g$ est un isomorphisme local au point fermé $x'$ de $X'$ en vertu de la
définition de $g$ et de (\hyperref[I.6.5.4-fr]{I, 6.5.4})~; mais comme $X'$
est le seul voisinage ouvert de $x'$, cela signifie que $g$ est une immersion
ouverte, donc $g(X')$ est aussi ouvert dans $X$, ce qui achève la
démonstration.

\begin{corollary}[6.2.6]
\label{II.6.2.6-fr}
Soient $A$ un anneau local noethérien complet, $Y=\operatorname{Spec}(A)$,
$f:X\to Y$ un morphisme séparé et quasi-fini. Alors $X$ est $Y$-isomorphe à
une somme $X'\amalg X''$, où $X'$ est un $Y$-schéma fini et $X''$ un
$Y$-schéma quasi-fini tel que, si $y$ est le point fermé de $Y$,
$X''\cap f^{-1}(y)=\varnothing$.
\end{corollary}

En effet, la fibre $f^{-1}(y)$ est finie et discrète par hypothèse, et le
corollaire résulte donc, par récurrence sur le nombre de points de cette
fibre, de (\hyperref[II.6.2.5-fr]{6.2.5}).

\begin{remark}[6.2.7]
\label{II.6.2.7-fr}
Au chapitre V, nous verrons que si $Y$ est localement noethérien, un morphisme
quasi-fini séparé $X\to Y$ est nécessairement quasi-affine.
\end{remark}

\subsection{Fermeture intégrale d'un préschéma.}

\begin{proposition}[6.3.1]
\label{II.6.3.1-fr}
Soient $(X,\mathcal{A})$ un espace annelé, $\mathcal{B}$ une
$\mathcal{A}$-Algèbre (commutative), $f$ une section de $\mathcal{B}$
au-dessus de $X$. Les propriétés suivantes sont équivalentes~:
\begin{enumerate}
\item[a)] Le sous-$\mathcal{A}$-Module de $\mathcal{B}$ engendré par les
$f^n$ pour $n\geq 0$ (\hyperref[0.5.1.1-fr]{0, 5.1.1}) est de type fini.
\item[b)] Il existe une sous-$\mathcal{A}$-Algèbre $\mathcal{C}$ de
$\mathcal{B}$, qui est un $\mathcal{A}$-Module de type fini, telle que
$f\in\Gamma(X,\mathcal{C})$.
\item[c)] Pour tout $x\in X$, $f_x$ est entier sur la fibre
$\mathcal{A}_x$.
\end{enumerate}
\end{proposition}

\oldpage[II]{117}

Comme le sous-$\mathcal{A}$-Module de $\mathcal{B}$ engendré par les $f^n$ est
une $\mathcal{A}$-Algèbre, il est clair que a) entraîne b). D'autre part, b)
implique que pour tout $x\in X$, le $\mathcal{A}_x$-module $\mathcal{C}_x$ est
de type fini, ce qui entraîne que tout élément de l'algèbre $\mathcal{C}_x$ et
en particulier $f_x$, est entier sur $\mathcal{A}_x$. Enfin, si pour un
$x\in X$, on a une relation de la forme
\[
f_x^n+(a_1)_xf_x^{n-1}+\ldots+(a_n)_x=0
\]
les $a_i$ ($1\leq i\leq n$) étant des sections de $\mathcal{A}$ au-dessus
d'un voisinage ouvert $U$ de $x$, la section
$f^n|U+a_1\cdot f^{n-1}|U+\ldots+a_n$ est nulle au-dessus d'un voisinage
$V\subset U$ de $x$, d'où résulte aussitôt que tous les $f^k|V$ ($k\geq0$)
sont combinaisons linéaires à coefficients dans $\Gamma(V,\mathcal{A})$ des
$f^j|V$ pour $0\leq j\leq n-1$~; on en conclut que c) entraîne a).

Lorsque les conditions équivalentes de (\hyperref[II.6.3.1-fr]{6.3.1}) sont
remplies on dit que la section $f$ est \emph{entière sur $\mathcal{A}$}.

\begin{corollary}[6.3.2]
\label{II.6.3.2-fr}
Sous les hypothèses de (\hyperref[II.6.3.1-fr]{6.3.1}), il existe un
sous-$\mathcal{A}$-Module (unique) $\mathcal{A}'$ de $\mathcal{B}$ tel que pour
tout $x\in X$, $\mathcal{A}'_x$ soit l'ensemble des germes
$f_x\in\mathcal{B}_x$ qui sont entiers sur $\mathcal{A}_x$. Pour tout ouvert
$U\subset X$, les sections de $\mathcal{A}'$ au-dessus de $U$ sont les
sections $f\in\Gamma(U,\mathcal{B})$ qui sont entières sur
$\mathcal{A}|U$.
\end{corollary}

L'existence de $\mathcal{A}'$ est immédiate, en prenant pour
$\Gamma(U,\mathcal{A}')$ l'ensemble des $f\in\Gamma(U,\mathcal{B})$ tels que
$f_x$ soit entier sur $\mathcal{A}_x$ pour tout $x\in U$. La seconde
assertion résulte aussitôt de (\hyperref[II.6.3.1-fr]{6.3.1}).

Il est clair que $\mathcal{A}'$ est une sous-$\mathcal{A}$-Algèbre de
$\mathcal{B}$~; on dit que c'est la \emph{fermeture intégrale de
$\mathcal{A}$ dans $\mathcal{B}$}.

\begin{env}[6.3.3]
\label{II.6.3.3-fr}
Soient $(X,\mathcal{A})$, $(Y,\mathcal{B})$ deux espaces annelés,
\[
g=(\psi,\theta):X\to Y
\]
un morphisme. Soit $\mathcal{C}$ (resp. $\mathcal{D}$) une
$\mathcal{A}$-Algèbre (resp. une $\mathcal{B}$-Algèbre) et soit
\[
u:\mathcal{D}\to\mathcal{C}
\]
un $g$-morphisme (\hyperref[0.4.4.1-fr]{0, 4.4.1}). Alors, si
$\mathcal{A}'$ (resp. $\mathcal{B}'$) est la fermeture intégrale de
$\mathcal{A}$ (resp. $\mathcal{B}$) dans $\mathcal{C}$ (resp.
$\mathcal{D}$), la \emph{restriction} de $u$ à $\mathcal{B}'$ est un
$g$-morphisme
\[
u':\mathcal{B}'\to\mathcal{A}'.
\]

En effet, si $j$ est l'injection canonique $\mathcal{B}'\to\mathcal{D}$, il
suffit de montrer que
\[
v=u^\sharp\circ g^*(j):g^*(\mathcal{B}')\to\mathcal{C}
\]
applique $g^*(\mathcal{B}')$ dans $\mathcal{A}'$. Or, un élément de
$(g^*(\mathcal{B}'))_x=\mathcal{B}'_{\psi(x)}\otimes_{\mathcal{B}_{\psi(x)}}
\mathcal{A}_x$ est entier sur $\mathcal{A}_x$ par définition de
$\mathcal{B}'$, donc il en est de même de son image par $v_x$, ce qui établit
notre assertion.
\end{env}

\begin{proposition}[6.3.4]
\label{II.6.3.4-fr}
Soient $X$ un préschéma, $\mathcal{A}$ une $\mathcal{O}_X$-Algèbre
quasi-cohérente. La fermeture intégrale $\mathcal{O}'_X$ de
$\mathcal{O}_X$ dans $\mathcal{A}$ est alors une $\mathcal{O}_X$-Algèbre
quasi-cohérente, et pour tout ouvert affine $U$ de $X$,
$\Gamma(U,\mathcal{O}'_X)$ est la fermeture intégrale de
$\Gamma(U,\mathcal{O}_X)$ dans $\Gamma(U,\mathcal{A})$.
\end{proposition}

On peut se borner au cas où $X=\operatorname{Spec}(B)$ est affine et
$\mathcal{A}=\widetilde{A}$, $A$ étant une

\oldpage[II]{118}

$B$-algèbre~; soit $B'$ la fermeture intégrale de $B$ dans $A$. Tout revient
à voir que pour tout $x\in X$, un élément de $A_x$, entier sur $B_x$,
appartient nécessairement à $B'_x$, ce qui résulte du fait que, pour un anneau
commutatif $C$, les opérations de fermeture intégrale dans une $C$-algèbre et
de passage à un anneau de fractions (pour une partie multiplicative de $C$)
commutent (\cite[t.~I, p.~261 et 257]{I-13}).

Le $X$-schéma $X'=\operatorname{Spec}(\mathcal{O}'_X)$ est alors appelé la
\emph{fermeture intégrale de $X$ relativement à $\mathcal{A}$} (ou
\emph{relativement à $\operatorname{Spec}(\mathcal{A})$})~; il est clair que
$X'$ est \emph{entier} sur $X$ (\hyperref[II.6.1.2-fr]{6.1.2}).

On déduit aussitôt de (\hyperref[II.6.3.4-fr]{6.3.4}) que si $f:X'\to X$ est
le morphisme structural, alors, pour tout ouvert $U$ de $X$, $f^{-1}(U)$ est
la \emph{fermeture intégrale du préschéma induit par $X$ sur $U$, relativement
à $\mathcal{A}|U$}.

\begin{env}[6.3.5]
\label{II.6.3.5-fr}
Soient $X$, $Y$ deux préschémas, $f:X\to Y$ un morphisme, $\mathcal{A}$
(resp. $\mathcal{B}$) une $\mathcal{O}_X$-Algèbre (resp. une
$\mathcal{O}_Y$-Algèbre) quasi-cohérente, et soit
$u:\mathcal{B}\to\mathcal{A}$ un $f$-morphisme. On a vu
(\hyperref[II.6.3.3-fr]{6.3.3}) qu'on en déduit un $f$-morphisme
$u':\mathcal{O}'_Y\to\mathcal{O}'_X$, où $\mathcal{O}'_X$ (resp.
$\mathcal{O}'_Y$) est la fermeture intégrale de $\mathcal{O}_X$ (resp.
$\mathcal{O}_Y$) dans $\mathcal{A}$ (resp. $\mathcal{B}$). Par suite, si
$X'$ (resp. $Y'$) est la fermeture intégrale de $X$ (resp. $Y$) relativement
à $\mathcal{A}$ (resp. $\mathcal{B}$), on déduit canoniquement de $u$ un
morphisme $f'=\operatorname{Spec}(u'):X'\to Y'$
(\hyperref[II.1.5.6-fr]{1.5.6}) rendant commutatif le diagramme
\[
\label{II.6.3.5.1-fr}
\xymatrix{
X'\ar[r]^{f'}\ar[d]&Y'\ar[d]\\
X\ar[r]_{f}&Y
}
\tag{6.3.5.1}
\]
\end{env}

\begin{env}[6.3.6]
\label{II.6.3.6-fr}
Supposons que $X$ n'ait qu'un nombre \emph{fini} de composantes irréductibles
$X_i$ ($1\leq i\leq r$), de points génériques $\xi_i$, et considérons en
particulier la fermeture intégrale de $X$ relativement à une
$\mathcal{R}(X)$-Algèbre $\mathcal{A}$ quasi-cohérente (en tant que
$\mathcal{O}_X$-Algèbre ou en tant que $\mathcal{R}(X)$-Algèbre, ce qui
revient au même). On sait (\hyperref[I.7.3.5-fr]{I, 7.3.5}) que
$\mathcal{A}$ est composée directe de $r$ $\mathcal{O}_X$-Algèbres
quasi-cohérentes $\mathcal{A}_i$, le support de $\mathcal{A}_i$ étant contenu
dans $X_i$, et le faisceau induit par $\mathcal{A}_i$ sur $X_i$ étant un
faisceau simple dont la fibre $A_i$ est une algèbre sur
$\mathcal{O}_{\xi_i}$. Il est clair alors (\hyperref[II.6.3.4-fr]{6.3.4}) que
la fermeture intégrale $\mathcal{O}'_X$ de $\mathcal{O}_X$ dans
$\mathcal{A}$ est \emph{composée directe} des fermetures intégrales
$\mathcal{O}^{(i)}_X$ de $\mathcal{O}_X$ dans chacune des $\mathcal{A}_i$, et
par suite la fermeture intégrale $X'=\operatorname{Spec}(\mathcal{O}'_X)$ de
$X$ relativement à $\mathcal{A}$ est un $X$-schéma \emph{somme} des
$\operatorname{Spec}(\mathcal{O}^{(i)}_X)=X'_i$ ($1\leq i\leq r$).

Supposons en outre que la $\mathcal{O}_X$-Algèbre $\mathcal{A}$ soit
\emph{réduite}, ou, ce qui revient au même, que chacune des algèbres $A_i$
soit réduite, et puisse par suite être considérée comme une algèbre sur le
\emph{corps} $k(\xi_i)$ (égal au corps des fonctions rationnelles du
sous-préschéma réduit de $X$ ayant $X_i$ pour espace sous-jacent)~; alors
(\hyperref[II.1.3.8-fr]{1.3.8}) chacun des $X'_i$ est un $X$-schéma
\emph{réduit} et $X'$ est aussi la fermeture intégrale de $X_{\mathrm{red}}$.
Supposons de plus que chacune des algèbres $A_i$ soit \emph{composée directe
d'un nombre fini de corps} $K_{ij}$ ($1\leq j\leq s_i$)~; si
$\mathcal{K}_{ij}$ est la sous-Algèbre de $\mathcal{A}_i$ correspondant à
$K_{ij}$, il est clair que $\mathcal{O}^{(i)}_X$ est \emph{composée directe}
des fermetures intégrales $\mathcal{O}^{(ij)}_X$ de $\mathcal{O}_X$ dans
chacune des $\mathcal{K}_{ij}$. Par suite, $X'_i$ est alors \emph{somme} des
$X$-schémas $X'_{ij}=\operatorname{Spec}(\mathcal{O}^{(ij)}_X)$
($1\leq j\leq s_i$). En outre, sous ces hypothèses et avec ces notations~:
\end{env}

\oldpage[II]{119}

\begin{proposition}[6.3.7]
\label{II.6.3.7-fr}
Chacun des $X'_{ij}$ est un $X$-préschéma intègre et normal, et son corps des
fonctions rationnelles $R(X'_{ij})$ s'identifie canoniquement à la fermeture
algébrique $K'_{ij}$ de $k(\xi_i)$ dans $K_{ij}$.
\end{proposition}

En vertu de ce qui précède, on peut supposer $X$ intègre, donc $r=1$, et
$s_1=1$, de sorte que l'unique algèbre $A_1$ est un corps $K$~; soit $\xi$ le
point générique de $X$, et soit $f:X'\to X$ le morphisme structural. Pour tout
ouvert affine non vide $U$ de $X$, $f^{-1}(U)$ s'identifie au spectre de la
fermeture intégrale $B'_U$ dans le corps $K$ de l'anneau intègre
$B_U=\Gamma(U,\mathcal{O}_X)$ (\hyperref[II.6.3.4-fr]{6.3.4})~; comme l'anneau
$B'_U$ est intègre et intégralement clos, il en est de même des anneaux locaux
des points de son spectre, donc $f^{-1}(U)$ est par définition un schéma
\emph{intègre et normal} ((\hyperref[0.4.1.4-fr]{0, 4.1.4}) et
(\hyperref[I.5.1.4-fr]{I, 5.1.4})). En outre, comme $(0)$ est le seul idéal
premier de $B'_U$ au-dessus de l'idéal premier $(0)$ de $B_U$
(\cite[t.~I, p.~259]{I-13}), $f^{-1}(\xi)$ se réduit à un seul point $\xi'$,
et $k(\xi')$ est le corps des fractions $K'$ de $B'_U$, qui n'est autre que la
fermeture algébrique de $k(\xi)$ dans $K$. Enfin, $X'$ est irréductible, car
lorsque $U$ parcourt l'ensemble des ouverts affines non vides de $X$, les
$f^{-1}(U)$ constituent un recouvrement ouvert de $X'$ formé d'ouverts
irréductibles~; en outre l'intersection $f^{-1}(U\cap V)$ de deux de ces
ouverts contient $\xi'$, donc n'est pas vide, et on conclut par
(\hyperref[0.2.1.4-fr]{0, 2.1.4}).

\begin{corollary}[6.3.8]
\label{II.6.3.8-fr}
Soit $X$ un préschéma réduit n'ayant qu'un nombre fini de composantes
irréductibles $X_i$ ($1\leq i\leq r$), et soit $\xi_i$ le point générique de
$X_i$. La fermeture intégrale $X'$ de $X$ relativement à $\mathcal{R}(X)$ est
somme de $r$ $X$-schémas $X'_i$ qui sont intègres et normaux. Si $f:X'\to X$
est le morphisme structural, $f^{-1}(\xi_i)$ se réduit au point générique
$\xi'_i$ de $X'_i$ et on a $k(\xi'_i)=k(\xi_i)$, autrement dit $f$ est
birationnel.
\end{corollary}

On dit dans ce cas que $X'$ est le \emph{normalisé} du préschéma réduit $X$~;
on notera que $f$, étant birationnel et entier, est \emph{surjectif}
(\hyperref[II.6.1.10-fr]{6.1.10}). Pour que l'on ait alors $X'=X$, il faut et
il suffit que $X$ soit normal. Lorsque $X$ est un préschéma intègre, il résulte
de (\hyperref[II.6.3.8-fr]{6.3.8}) que son normalisé $X'$ est intègre.

\begin{env}[6.3.9]
\label{II.6.3.9-fr}
Soient $X$, $Y$ deux préschémas intègres, $f:X\to Y$ un morphisme dominant,
$L=R(X)$, $K=R(Y)$ les corps des fonctions rationnelles de $X$ et $Y$~; il
correspond canoniquement à $f$ une injection $K\to L$, et lorsqu'on identifie
$K$ (resp. $L$) au faisceau simple $\mathcal{R}(Y)$ (resp.
$\mathcal{R}(X)$), cette injection est un $f$-morphisme. Soit $K_1$ (resp.
$L_1$) une extension de $K$ (resp. $L$) et supposons donné un monomorphisme
$K_1\to L_1$ tel que le diagramme
\[
\xymatrix{
K_1\ar[r]&L_1\\
K\ar[u]\ar[r]&L\ar[u]
}
\]
soit commutatif~; si $K_1$ (resp. $L_1$) est considéré comme faisceau simple
sur $Y$ (resp. $X$), donc comme une $\mathcal{R}(Y)$-Algèbre (resp. une
$\mathcal{R}(X)$-Algèbre), cela signifie que $K_1\to L_1$ est un
$f$-morphisme. Cela étant, si $X'$ (resp. $Y'$) est la fermeture intégrale de
$X$ (resp. $Y$) relativement à $L_1$ (resp. $K_1$), $X'$ (resp. $Y'$) est un
préschéma intègre et normal (\hyperref[II.6.3.6-fr]{6.3.6}) dont le corps des
fonctions rationnelles s'identifie canoniquement à la fermeture algébrique

\oldpage[II]{120}

$L'$ (resp. $K'$) de $L$ (resp. $K$) dans $L_1$ (resp. $K_1$), et il existe
un morphisme canonique (nécessairement dominant) $f':X'\to Y'$ rendant
commutatif le diagramme (\hyperref[II.6.3.5.1-fr]{6.3.5.1}).

Le cas le plus important est celui où on prend $L_1=L$, $K_1$ étant donc une
extension de $K$ contenue dans $L$, et où on suppose $X$ intègre et
\emph{normal}, donc $X'=X$. Ce qui précède montre alors que lorsque $X$ est
normal, et que $Y'$ est la fermeture intégrale de $Y$ relative à un corps
$K_1\subset L=R(X)$, tout morphisme dominant $f:X\to Y$ se factorise en
\[
f:X\xrightarrow{f'}Y'\to Y
\]
où $f'$ est dominant~; d'ailleurs, lorsque le monomorphisme $K_1\to L$ est
donné, $f'$ est nécessairement unique, comme on le voit en se ramenant au cas
où $X$ et $Y$ sont affines. On voit donc que pour la donnée de $Y$, $L$ et
d'un $K$-monomorphisme $K_1\to L$, la fermeture intégrale $Y'$ de $Y$
relativement à $K_1$ est solution d'un \emph{problème universel}.
\end{env}

\begin{remark}[6.3.10]
\label{II.6.3.10-fr}
Reprenons les hypothèses de (\hyperref[II.6.3.6-fr]{6.3.6}), en supposant de
plus que chacune des algèbres $A_i$ soit de \emph{rang fini} sur $k(\xi_i)$
(ce qui entraîne que $A_i$ est composée directe d'un nombre fini de corps)~;
on peut affirmer dans certains cas que le morphisme structural $X'\to X$ est
non seulement entier, mais même \emph{fini}. Bornons-nous au cas où $X$ est
\emph{réduit}~; la question étant locale sur $X$, on peut en outre supposer
que $X$ est affine d'anneau $C$, et que $C$ n'a qu'un nombre fini d'idéaux
minimaux $\mathfrak{p}_i$ ($1\leq i\leq r$) tels que les
$C_i=C/\mathfrak{p}_i$ soient intègres~; $X'$ sera alors fini sur $X$ si la
fermeture intégrale de chaque $C_i$ dans toute extension de degré fini de son
corps des fractions est un $C$-module de type fini
(\hyperref[II.6.3.4-fr]{6.3.4}). On sait que cette condition est toujours
vérifiée lorsque $C$ est une \emph{algèbre de type fini sur un corps}
(\cite[t.~I, p.~267, th.~9]{I-13}), ou \emph{sur $\mathbf{Z}$}
(\cite[t.~I, p.~93, th.~3]{I-9}), ou \emph{sur un anneau local noethérien
complet} ([25], p.~298). On en conclut que $X'\to X$ sera un morphisme
fini lorsque $X$ est un schéma de \emph{type fini sur un corps}, ou sur
$\mathbf{Z}$, ou sur un anneau local noethérien complet.
\end{remark}

\subsection{Déterminant d'un endomorphisme de $\mathcal{O}_X$-Module.}

\begin{env}[6.4.1]
\label{II.6.4.1-fr}
Soient $A$ un anneau (commutatif), $E$ un $A$-module libre de rang $n$, $u$
un endomorphisme de $E$~; rappelons que pour définir le \emph{polynôme
caractéristique} de $u$, on considère l'endomorphisme $u\otimes1$ du
$A[T]$-module libre de rang $n$, $E\otimes_AA[T]$ ($T$ indéterminée), et on
pose
\[
\label{II.6.4.1.1-fr}
P(u,T)=\det(T\cdot I-(u\otimes1))
\tag{6.4.1.1}
\]
($I$ automorphisme identique de $E\otimes_AA[T]$). On a
\[
\label{II.6.4.1.2-fr}
P(u,T)=T^n-\sigma_1(u)T^{n-1}+\ldots+(-1)^n\sigma_n(u)
\tag{6.4.1.2}
\]
où $\sigma_i(u)$ est un élément de $A$, égal à un polynôme homogène de degré
$i$ (à coefficients entiers) par rapport aux éléments de la matrice de $u$
relative à une base quelconque de $E$~; on dit que les $\sigma_i(u)$ sont les
\emph{fonctions symétriques élémentaires de $u$}, et on a en particulier
$\sigma_1(u)=\operatorname{Tr}u$ et $\sigma_n(u)=\det u$. Rappelons qu'en
vertu du th. de Hamilton-Cayley, on a
\[
\label{II.6.4.1.3-fr}
P(u,u)=u^n-\sigma_1(u)u^{n-1}+\ldots+(-1)^n\sigma_n(u)=0
\tag{6.4.1.3}
\]

\oldpage[II]{121}

ce qui s'écrit aussi
\[
\label{II.6.4.1.4-fr}
(\det u)\cdot I_E=uQ(u)
\tag{6.4.1.4}
\]
($I_E$ automorphisme identique de $E$), avec
\[
\label{II.6.4.1.5-fr}
Q(u)=(-1)^{n+1}(u^{n-1}-\sigma_1(u)u^{n-2}+\ldots+(-1)^{n-1}\sigma_{n-1}(u)).
\tag{6.4.1.5}
\]

Soit $\varphi:A\to B$ un homomorphisme d'anneaux, faisant de $B$ une
$A$-algèbre~; considérons le $B$-module
$E_{(B)}=E\otimes_AB$ qui est libre de rang $n$, et l'extension $u\otimes1$
de $u$ à un endomorphisme de $E_{(B)}$~; il est immédiat que l'on a
$\sigma_i(u\otimes1)=\varphi(\sigma_i(u))$ pour tout indice $i$.
\end{env}

\begin{env}[6.4.2]
\label{II.6.4.2-fr}
Supposons maintenant que $A$ soit un anneau \emph{intègre}, de corps des
fractions $K$, et $E$ un $A$-module de \emph{type fini} (mais non
nécessairement libre). Soit $n$ le \emph{rang} de $E$, c'est-à-dire le rang du
$K$-module libre $E\otimes_AK$~; à tout endomorphisme $u$ de $E$ correspond
canoniquement l'endomorphisme $u\otimes1$ de $E\otimes_AK$~; par abus de
langage, on appelle encore \emph{polynôme caractéristique} de $u$ et on note
$P(u,T)$ le polynôme $P(u\otimes1,T)$, dont les coefficients
$\sigma_i(u\otimes1)$ (qui appartiennent à $K$) seront encore notés
$\sigma_i(u)$ et appelés les \emph{fonctions symétriques élémentaires de
$u$}~; en particulier $\det u=\det(u\otimes1)$ par définition. Avec ces
notations, les formules (\hyperref[II.6.4.1.3-fr]{6.4.1.3}) à
(\hyperref[II.6.4.1.5-fr]{6.4.1.5}) ont un sens et sont encore valables,
lorsque l'on interprète $u^j$ comme l'homomorphisme $E\to E\otimes_AK$
composé de l'endomorphisme $u^j\otimes1=(u\otimes1)^j$ de $E\otimes_AK$ et
de l'homomorphisme canonique $x\mapsto x\otimes1$.

Si $F$ est le sous-module de torsion de $E$, et si $E_0=E/F$, on a
$u(F)\subset F$, donc, par passage aux quotients, $u$ donne un endomorphisme
$u_0$ de $E_0$~; en outre $E\otimes_AK$ s'identifie à $E_0\otimes_AK$ et
$u\otimes1$ à $u_0\otimes1$, donc $\sigma_i(u)=\sigma_i(u_0)$ pour
$1\leq i\leq n$.

Lorsque $E$ est \emph{sans torsion}, $E$ s'identifie canoniquement à un
sous-$A$-module de $E\otimes_AK$, et la relation $u\otimes1=0$ est
équivalente à $u=0$. Lorsque $E$ est un $A$-module libre, les deux définitions
des $\sigma_i(u)$ données dans (\hyperref[II.6.4.1-fr]{6.4.1}) et
(\hyperref[II.6.4.2-fr]{6.4.2}) coïncident d'après ce qui précède, ce qui
justifie les notations adoptées.

On notera aussi que lorsque $E$ est un \emph{module de torsion}, autrement dit
$E_0=\{0\}$, l'algèbre extérieure de $E_0$ est réduite à $K$ et le déterminant
de l'unique endomorphisme $u_0$ de $E_0$ est \emph{égal à $1$}.
\end{env}

\begin{proposition}[6.4.3]
\label{II.6.4.3-fr}
Soient $A$ un anneau intègre, $E$ un $A$-module de type fini, $u$ un
endomorphisme de $E$~; alors les fonctions symétriques élémentaires
$\sigma_i(u)$ de $u$ (et en particulier $\det u$) sont des éléments de $K$
entiers sur $A$.
\end{proposition}

Soit $A'$ la clôture intégrale de $A$~; comme $A'[T]$ est un anneau
intégralement clos ([24], p.~99), c'est la clôture intégrale de
$A[T]$ dans son corps des fractions $K(T)$. Remplaçant $u$ par
$T\cdot I-u\otimes1$ et $A$ par $A[T]$, on voit qu'on est ramené à prouver que
$\det u$ est entier sur $A$. Si $n$ est le rang de $E$, on a
$\det u=\det(\bigwedge^n u)$ et
$(\bigwedge^n u)\otimes1=\bigwedge^n(u\otimes1)$, donc on peut supposer
$n=1$. Mais alors l'application $u\mapsto\det u$ est un homomorphisme du
$A$-module $\operatorname{Hom}_A(E,E)$ dans $K$~; comme $E$ est de type fini,
$\operatorname{Hom}_A(E,E)$ est isomorphe à un sous-module du $A$-module de
type fini $E^n$ (si $E$ admet un système de $n$ générateurs)~; donc les
éléments $\det u$ appartiennent à un sous-$A$-module de $K$ de type fini, et
sont par suite entiers sur $A$.

\oldpage[II]{122}

\begin{corollary}[6.4.4]
\label{II.6.4.4-fr}
Sous les hypothèses de (\hyperref[II.6.4.3-fr]{6.4.3}), si on suppose en outre
$A$ normal, les $\sigma_i(u)$ (en particulier $\operatorname{Tr}u$ et
$\det u$) appartiennent à $A$.
\end{corollary}

\begin{proposition}[6.4.5]
\label{II.6.4.5-fr}
Soient $A$ un anneau intègre, $E$ un $A$-module de type fini, de rang $n$, $u$
un endomorphisme de $E$ tel que les $\sigma_i(u)$ appartiennent à $A$. Pour
que $u$ soit un automorphisme de $E$, il est nécessaire que $\det u$ soit
inversible dans $A$~; cette condition est suffisante lorsque $E$ est sans
torsion.
\end{proposition}

La condition est \emph{suffisante}, car l'hypothèse et les formules
(\hyperref[II.6.4.1.4-fr]{6.4.1.4}) et
(\hyperref[II.6.4.1.5-fr]{6.4.1.5}) (valables dans $E$, et non seulement dans
$E\otimes_AK$, puisque $E$ est sans torsion) prouvent que
$(\det u)^{-1}Q(u)$ est l'inverse de $u$.

La condition est \emph{nécessaire}, car si $u$ est inversible, il résulte de
(\hyperref[II.6.4.3-fr]{6.4.3}) que $\det(u^{-1})$ appartient à la clôture
intégrale $A'$ de $A$ dans son corps des fractions $K$ et est évidemment
inverse de $\det u$ dans $A'$. Notre assertion résulte alors du~:

\begin{lemma}[6.4.5.1]
\label{II.6.4.5.1-fr}
Soit $A$ un sous-anneau d'un anneau $A'$ tel que $A'$ soit entier sur $A$. Si
un élément $x\in A$ est inversible dans $A'$, il est aussi inversible dans
$A$.
\end{lemma}

Dans le cas contraire, $x$ appartiendrait à un idéal maximal $\mathfrak{m}$
de $A$ et il résulte du premier th. de Cohen-Seidenberg
(\cite[t.~I, p.~257, th.~3]{I-13}) qu'il y aurait un idéal maximal
$\mathfrak{m}'$ de $A'$ tel que $\mathfrak{m}=\mathfrak{m}'\cap A$~; on aurait
donc $x\in\mathfrak{m}'$, ce qui est absurde.

\begin{corollary}[6.4.6]
\label{II.6.4.6-fr}
Soient $A$ un anneau intègre et intégralement clos, $E$ un $A$-module de type
fini sans torsion, $u$ un endomorphisme de $E$. Pour que $u$ soit un
automorphisme de $E$, il faut et il suffit que $\det u$ soit inversible dans
$A$.
\end{corollary}

Cela résulte de (\hyperref[II.6.4.4-fr]{6.4.4}) et
(\hyperref[II.6.4.5-fr]{6.4.5}).

\begin{remark}[6.4.7]
\label{II.6.4.7-fr}
Nous aurons besoin plus loin d'une généralisation des résultats précédents.
Considérons un anneau noethérien \emph{réduit} $A$~; soient
$\mathfrak{p}_\alpha$ ($1\leq\alpha\leq r$) ses idéaux minimaux, $K_\alpha$
le corps des fractions de l'anneau intègre $A/\mathfrak{p}_\alpha$, $K$
l'anneau total des fractions de $A$, \emph{composé direct} des corps
$K_\alpha$. Soit $E$ un $A$-module de type fini, et \emph{supposons que}
$E\otimes_AK$ soit un $K$-module libre de rang $n$ (ce qui ici n'est plus
conséquence des autres hypothèses)~; il revient au même de dire que tous les
$K_\alpha$-espaces vectoriels $E\otimes_AK_\alpha=E_\alpha$ ont \emph{même
dimension $n$}~; si alors $u$ est un endomorphisme de $E$, on pose encore
$P(u,T)=P(u\otimes1,T)$ et $\sigma_j(u)=\sigma_j(u\otimes1)$, et en
particulier $\det u=\det(u\otimes1)$~; les $\sigma_j(u)$ sont donc des
éléments de $K$. Il est immédiat que $E\otimes_AK$ est somme directe des
$E_\alpha$ et que ces derniers sont stables par $u\otimes1$~; la restriction
de $u\otimes1$ à $E_\alpha$ n'est autre d'ailleurs que l'extension $u_\alpha$
de $u$ à ce $K_\alpha$-espace vectoriel~; on en conclut que $\sigma_j(u)$ est
l'élément de $K$ dont les composantes dans les $K_\alpha$ sont les
$\sigma_j(u_\alpha)$. Comme la fermeture intégrale de $A$ dans $K$ est
\emph{composée directe} des fermetures intégrales de $A$ dans les
$K_\alpha$, les $\sigma_j(u)$ sont \emph{entiers} sur $A$.

\begin{lemma}[6.4.7.1]
\label{II.6.4.7.1-fr}
La sous-$A$-algèbre de $K$ engendrée par tous les éléments $\sigma_j(u)$
($1\leq j\leq n$) lorsque $u$ parcourt $\operatorname{Hom}_A(E,E)$, est un
$A$-module de type fini.
\end{lemma}

Il suffit de démontrer que la sous-$A[T]$-algèbre de $K[T]$ engendrée par les
$P(u,T)$ est un $A[T]$-module de type fini, car si les $F_i(T)$
($1\leq i\leq m$) forment un système de générateurs de ce $A[T]$-module, les
coefficients des $F_i(T)$ sont entiers sur $A$, donc engendrent une
$A$-algèbre qui est un $A$-module de type fini
(\cite[t.~I, p.~255, th.~1]{I-13}). On peut donc

\oldpage[II]{123}

remplacer $A$ par $A[T]$ (qui est noethérien) et $E$ par
$E\otimes_AA[T]=E'$ qui est tel que
$E'\otimes_{A[T]}K[T]=E\otimes_AK[T]$ soit un $K[T]$-module libre de rang $n$.
Revenant aux notations initiales, on voit donc qu'il suffit de prouver que le
$A$-module engendré par les éléments $\det u$, où $u$ parcourt
$\operatorname{Hom}_A(E,E)$, est de type fini~; \emph{a fortiori} (puisque
tout sous-module d'un $A$-module de type fini est de type fini) il suffit de
prouver que lorsque $v$ parcourt l'ensemble des endomorphismes de
$\bigwedge^nE$, le $A$-module engendré par les $\det v$ est de type fini~;
autrement dit, on est encore ramené au cas où $n=1$. Mais alors la proposition
résulte de ce que $\operatorname{Hom}_A(E,E)$ est un $A$-module de type fini
et de ce que $v\mapsto\det v$ est un homomorphisme de ce $A$-module dans $K$.

Soit $F$ le noyau de l'homomorphisme canonique $E\to E\otimes_AK$ et soit
$E_0=E/F$~; on voit comme ci-dessus que $E\otimes_AK$ s'identifie à
$E_0\otimes_AK$, que $u(F)\subset F$, et que si $u_0$ est l'endomorphisme de
$E_0$ déduit de $u$ par passage aux quotients, $u\otimes1$ s'identifie à
$u_0\otimes1$, et $\sigma_j(u)=\sigma_j(u_0)$ pour tout $j$. \emph{Si l'on a
$F=0$}, les formules (\hyperref[II.6.4.1.3-fr]{6.4.1.3}) et
(\hyperref[II.6.4.1.5-fr]{6.4.1.5}) ont un sens et sont valables lorsque $E$
est identifié à un sous-module de $E\otimes_AK$ et les $u^j$ à des
homomorphismes $E\to E\otimes_AK$~; par suite, la prop.
(\hyperref[II.6.4.5-fr]{6.4.5}) s'étend aussi à ce cas, avec sa démonstration.
\end{remark}

\begin{env}[6.4.8]
\label{II.6.4.8-fr}
Soient $(X,\mathcal{A})$ un espace annelé, $\mathcal{E}$ un
$\mathcal{A}$-Module \emph{localement libre} (de rang fini), $u$ un
endomorphisme de $\mathcal{E}$. Il y a par hypothèse une base $\mathfrak{B}$
de la topologie de $X$ telle que pour tout $V\in\mathfrak{B}$,
$\mathcal{E}|V$ soit isomorphe à $\mathcal{A}^n|V$ (pour un $n$ pouvant varier
avec $V$). Soit $u$ un endomorphisme de $\mathcal{E}$~; pour tout
$V\in\mathfrak{B}$, $u_V$ est donc un endomorphisme du
$\Gamma(V,\mathcal{A})$-module $\Gamma(V,\mathcal{E})$, qui est libre par
hypothèse~; le déterminant $\det u_V$ est donc défini et appartient à
$\Gamma(V,\mathcal{A})$. En outre, si $e_1,\ldots,e_n$ forment une base de
$\Gamma(V,\mathcal{E})$, leurs restrictions à tout ouvert $W\subset V$
forment une base de $\Gamma(W,\mathcal{E})$ sur $\Gamma(W,\mathcal{A})$,
donc $\det u_W$ est la restriction de $\det u_V$ à $W$. Il existe donc une
section et une seule de $\mathcal{A}$ au-dessus de $X$, que nous noterons
$\det u$ et appellerons le \emph{déterminant de $u$}, telle que la restriction
de $\det u$ à tout $V\in\mathfrak{B}$ soit $\det u_V$. Il est clair que pour
tout $x\in X$, on a $(\det u)_x=\det u_x$~; pour deux endomorphismes $u$, $v$
de $\mathcal{E}$, on a
\[
\label{II.6.4.8.1-fr}
\det(u\circ v)=(\det u)(\det v)
\tag{6.4.8.1}
\]
ainsi que
\[
\label{II.6.4.8.2-fr}
\det(1_{\mathcal{E}})=1_{\mathcal{A}}
\tag{6.4.8.2}
\]
et, si le rang de $\mathcal{E}$ est constant (ce qui sera le cas
(\hyperref[0.5.4.1-fr]{0, 5.4.1}) pour $X$ connexe) et égal à $n$
\[
\label{II.6.4.8.3-fr}
\det(s\cdot u)=s^n\det u
\tag{6.4.8.3}
\]
pour tout $s\in\Gamma(X,\mathcal{A})$ (on notera que
$\det(0)=0_{\mathcal{A}}$ si $n\geq1$, mais $\det(0)=1_{\mathcal{A}}$ pour
$n=0$). En outre, pour que $u$ soit un \emph{automorphisme} de
$\mathcal{E}$, il faut et il suffit que $\det u$ soit \emph{inversible} dans
$\Gamma(X,\mathcal{A})$.

Si le rang de $\mathcal{E}$ est constant, on peut définir de la même manière
les fonctions symétriques élémentaires $\sigma_i(u)$, qui sont des éléments
de $\Gamma(X,\mathcal{A})$~; on a encore les relations
(\hyperref[II.6.4.1.3-fr]{6.4.1.3}) à
(\hyperref[II.6.4.1.5-fr]{6.4.1.5}).

\oldpage[II]{124}

On a ainsi défini un homomorphisme
$\det:\operatorname{Hom}_{\mathcal{A}}(\mathcal{E},\mathcal{E})
\to\Gamma(X,\mathcal{A})$ de monoïdes multiplicatifs~; si on note que
$\operatorname{Hom}_{\mathcal{A}}(\mathcal{E},\mathcal{E})=
\Gamma(X,\mathcal{H}om_{\mathcal{A}}(\mathcal{E},\mathcal{E}))$ par
définition, on voit qu'on peut dans cette définition remplacer $X$ par un
ouvert quelconque $U$ de $X$, ce qui définit immédiatement un homomorphisme
$\det:\mathcal{H}om_{\mathcal{A}}(\mathcal{E},\mathcal{E})\to\mathcal{A}$ de
faisceaux de monoïdes multiplicatifs. Lorsque $\mathcal{E}$ a un rang
constant, on définit de même des homomorphismes
$\sigma_i:\mathcal{H}om_{\mathcal{A}}(\mathcal{E},\mathcal{E})
\to\mathcal{A}$ de faisceaux d'ensembles~; pour $i=1$, l'homomorphisme
$\sigma_1=\operatorname{Tr}$ est un homomorphisme de $\mathcal{A}$-Modules.

Soit $(Y,\mathcal{B})$ un second espace annelé, et soit
$f:(X,\mathcal{A})\to(Y,\mathcal{B})$ un morphisme d'espaces annelés~; si
$\mathcal{F}$ est un $\mathcal{B}$-Module localement libre,
$f^*(\mathcal{F})$ est un $\mathcal{A}$-Module localement libre (qui a même
rang que $\mathcal{F}$ si ce dernier est de rang constant)
(\hyperref[0.5.4.5-fr]{0, 5.4.5}). Pour tout endomorphisme $v$ de
$\mathcal{F}$, $f^*(v)$ est un endomorphisme de $f^*(\mathcal{F})$, et il
résulte aussitôt des définitions que $\det f^*(v)$ est la section de
$\mathcal{A}=f^*(\mathcal{B})$ au-dessus de $X$ qui correspond canoniquement
à $\det v\in\Gamma(Y,\mathcal{B})$. On peut encore dire que l'homomorphisme
$f^*(\det):f^*(\mathcal{H}om_{\mathcal{B}}(\mathcal{F},\mathcal{F}))
\to f^*(\mathcal{B})=\mathcal{A}$ est le composé
\[
f^*(\mathcal{H}om_{\mathcal{B}}(\mathcal{F},\mathcal{F}))
\xrightarrow{\gamma^\sharp}
\mathcal{H}om_{\mathcal{A}}(f^*(\mathcal{F}),f^*(\mathcal{F}))
\xrightarrow{\det}\mathcal{A}
\]
(\hyperref[0.4.4.6-fr]{0, 4.4.6}). On a des résultats analogues pour les
$\sigma_i$.
\end{env}

\begin{env}[6.4.9]
\label{II.6.4.9-fr}
Supposons maintenant que $X$ soit un \emph{préschéma localement intègre}, de
sorte que le faisceau $\mathcal{R}(X)$ des fonctions rationnelles sur $X$ est
un faisceau de corps localement simple (\hyperref[I.7.3.4-fr]{I, 7.3.4}),
quasi-cohérent en tant que $\mathcal{O}_X$-Module. Si $\mathcal{E}$ est un
$\mathcal{O}_X$-Module quasi-cohérent \emph{de type fini},
$\mathcal{E}'=\mathcal{E}\otimes_{\mathcal{O}_X}\mathcal{R}(X)$ est alors un
$\mathcal{R}(X)$-Module localement libre
(\hyperref[I.7.3.6-fr]{I, 7.3.6})~; pour tout endomorphisme $u$ de
$\mathcal{E}$, $u\otimes1_{\mathcal{R}(X)}$ est alors un endomorphisme de
$\mathcal{E}'$, et $\det(u\otimes1)$ une section de $\mathcal{R}(X)$ au-dessus
de $X$, que l'on appelle encore \emph{déterminant de $u$} et que l'on note
encore $\det u$. Il résulte de (\hyperref[II.6.4.3-fr]{6.4.3}) que $\det u$ est
une section de la \emph{fermeture intégrale} de $\mathcal{O}_X$ dans
$\mathcal{R}(X)$ (\hyperref[II.6.3.2-fr]{6.3.2})~; si en outre $X$ est
\emph{normal}, $\det u$ est donc une \emph{section de $\mathcal{O}_X$}
au-dessus de $X$, et si on suppose de plus que $\mathcal{E}$ est \emph{sans
torsion}, pour que $u$ soit un automorphisme de $\mathcal{E}$, il faut et il
suffit que $\det u$ soit \emph{inversible} en vertu de
(\hyperref[II.6.4.6-fr]{6.4.6}). Les formules
(\hyperref[II.6.4.8.1-fr]{6.4.8.1}) à
(\hyperref[II.6.4.8.3-fr]{6.4.8.3}) sont encore valables~; de l'homomorphisme
$u\mapsto\det u$, appliqué aux modules de sections de
$\mathcal{H}om_{\mathcal{O}_X}(\mathcal{E},\mathcal{E})$, on déduit un
homomorphisme de faisceaux
$\det:\mathcal{H}om_{\mathcal{O}_X}(\mathcal{E},\mathcal{E})
\to\mathcal{R}(X)$, qui prend ses valeurs dans $\mathcal{O}_X$ quand $X$ est
normal. On a des définitions et résultats analogues pour les autres fonctions
symétriques élémentaires $\sigma_j(u)$, lorsque $\mathcal{E}'$ est de
\emph{rang constant}~; si de plus $X$ est \emph{normal}, les $\sigma_j(u)$
sont des \emph{sections de $\mathcal{O}_X$} au-dessus de $X$.

Enfin, soient $X$, $Y$ deux préschémas intègres, et soit $f:X\to Y$ un
morphisme \emph{dominant}. On sait qu'il existe alors un homomorphisme
canonique $f^*(\mathcal{R}(Y))\to\mathcal{R}(X)$
(\hyperref[I.7.3.8-fr]{I, 7.3.8}\footnote{Voir \emph{Errata et Addenda},
liste 1.}), d'où résulte, pour tout
$\mathcal{O}_Y$-Module quasi-cohérent de type fini $\mathcal{F}$, un
homomorphisme canonique
$\theta:f^*(\mathcal{F}\otimes_{\mathcal{O}_Y}\mathcal{R}(Y))
\to f^*(\mathcal{F})\otimes_{\mathcal{O}_X}\mathcal{R}(X)$. Si $v$ est un
endo-

\oldpage[II]{125}

morphisme de $\mathcal{F}$,
$f^*(v\otimes1_{\mathcal{R}(Y)})$ est un endomorphisme de
$f^*(\mathcal{F}\otimes_{\mathcal{O}_Y}\mathcal{R}(Y))$, et on a un diagramme
commutatif
\[
\xymatrix{
f^*(\mathcal{F}\otimes_{\mathcal{O}_Y}\mathcal{R}(Y))
  \ar[r]^{f^*(v\otimes1)}\ar[d]_{\theta}
& f^*(\mathcal{F}\otimes_{\mathcal{O}_Y}\mathcal{R}(Y))\ar[d]^{\theta}\\
f^*(\mathcal{F})\otimes_{\mathcal{O}_X}\mathcal{R}(X)
  \ar[r]_{f^*(v)\otimes1}
& f^*(\mathcal{F})\otimes_{\mathcal{O}_X}\mathcal{R}(X).
}
\]
On en conclut aisément que $\det f^*(v)$ est l'image canonique, par
l'homomorphisme $f^*(\mathcal{R}(Y))\to\mathcal{R}(X)$, de la section
$\det v$ de $\mathcal{R}(Y)$~; on est en effet aussitôt ramené au cas où
$X=\operatorname{Spec}(A)$ et $Y=\operatorname{Spec}(B)$ sont affines, $A$ et
$B$ étant des anneaux intègres de corps de fractions respectifs $K$ et $L$,
l'homomorphisme $B\to A$ étant injectif et s'étendant donc en un monomorphisme
$L\to K$~; si $\mathcal{F}=\widetilde{M}$, où $M$ est un $B$-module de type
fini, le rang sur $L$ de $M\otimes_B L$ est \emph{égal} à celui de
$(M\otimes_B A)\otimes_A K$ sur $K$, et $\det((u\otimes1)\otimes1)$ est
l'image dans $K$ de $\det(u\otimes1)$ pour tout endomorphisme $u$ de $M$,
d'où la conclusion.
\end{env}

\begin{env}[6.4.10]
\label{II.6.4.10-fr}
Supposons enfin que $X$ soit un \emph{préschéma localement noethérien réduit},
de sorte que le faisceau $\mathcal{R}(X)$ des fonctions rationnelles sur $X$
est encore un $\mathcal{O}_X$-Module quasi-cohérent
(\hyperref[I.7.3.4-fr]{I, 7.3.4})~; soit d'autre part $\mathcal{E}$ un
$\mathcal{O}_X$-Module \emph{cohérent} tel que
$\mathcal{E}'=\mathcal{E}\otimes_{\mathcal{O}_X}\mathcal{R}(X)$ soit
\emph{localement libre} (de rang fini). En vertu de
(\hyperref[II.6.4.7-fr]{6.4.7}), si $\mathcal{E}'$ est de rang constant, on
peut, pour tout endomorphisme $u$ de $\mathcal{E}$, définir les $\sigma_j(u)$,
qui sont des sections de $\mathcal{R}(X)$ au-dessus de $X$. Lorsqu'on ne
suppose plus $\mathcal{E}'$ de rang constant, on peut encore définir
l'homomorphisme
$\det:\mathcal{H}om_{\mathcal{O}_X}(\mathcal{E},\mathcal{E})
\to\mathcal{R}(X)$.
\end{env}

\subsection{Norme d'un faisceau inversible.}
\label{subsection:II.6.5-fr}

\begin{env}[6.5.1]
\label{II.6.5.1-fr}
Soit $(X,\mathcal{A})$ un espace annelé et soit $\mathcal{B}$ une
$\mathcal{A}$-Algèbre (commutative). Le $\mathcal{A}$-Module $\mathcal{B}$
s'identifie canoniquement à un sous-$\mathcal{A}$-Module de
$\mathcal{H}om_{\mathcal{A}}(\mathcal{B},\mathcal{B})$, une section $f$ de
$\mathcal{B}$ au-dessus d'un ouvert $U$ de $X$ s'identifiant à la
multiplication par cette section. Si $(X,\mathcal{A})$ et $\mathcal{B}$
vérifient l'une des conditions énumérées dans
(\hyperref[II.6.4.8-fr]{6.4.8}), (\hyperref[II.6.4.9-fr]{6.4.9}) ou
(\hyperref[II.6.4.10-fr]{6.4.10}), on peut donc définir $\det(f)$ (et dans
certains cas les $\sigma_j(f)$) qui sont des sections de $\mathcal{A}$ ou de
$\mathcal{R}(X)$ au-dessus de $U$, que nous appellerons la \emph{norme de
$f$} (resp. les \emph{fonctions symétriques élémentaires de $f$}) et noterons
$N_{\mathcal{B}/\mathcal{A}}(f)$. Nous allons supposer vérifiée \emph{l'une}
des conditions suivantes~:
\begin{enumerate}
\item[(I)] \emph{$\mathcal{B}$ est un $\mathcal{A}$-Module localement libre
(de rang fini).}
\item[(II)] \emph{$(X,\mathcal{A})$ est un préschéma localement noethérien
réduit, $\mathcal{B}$ est un $\mathcal{A}$-Module cohérent tel que
$\mathcal{B}\otimes_{\mathcal{A}}\mathcal{R}(X)$ soit un
$\mathcal{R}(X)$-Module localement libre, et pour toute section
$f\in\Gamma(U,\mathcal{B})$ au-dessus d'un ouvert $U\subset X$,
$N_{\mathcal{B}/\mathcal{A}}(f)$ est une section de $\mathcal{A}$ au-dessus
de $U$.}
\end{enumerate}

\oldpage[II]{126}

On notera que cette dernière condition est automatiquement vérifiée lorsque
le préschéma localement noethérien $X$ est \emph{normal}
(\hyperref[II.6.4.9-fr]{6.4.9}).

D'autre part, l'hypothèse que
$\mathcal{B}\otimes_{\mathcal{A}}\mathcal{R}(X)$ soit localement libre peut
encore s'exprimer de la façon suivante~: désignons par $X_\alpha$ les
sous-préschémas fermés réduits de $X$ ayant pour espaces sous-jacents les
composantes irréductibles de $X$ (\hyperref[I.5.2.1-fr]{I, 5.2.1}), qui sont
donc des préschémas intègres localement noethériens. Tout $x\in X$ appartient
à un nombre fini des sous-espaces $X_\alpha$~; d'autre part,
$\mathcal{B}\otimes_{\mathcal{A}}\mathcal{R}(X_\alpha)$ est un
$\mathcal{R}(X_\alpha)$-Module localement libre de rang constant $k_\alpha$
(\hyperref[I.7.3.6-fr]{I, 7.3.6})~; dire que
$\mathcal{B}\otimes_{\mathcal{A}}\mathcal{R}(X)$ est un
$\mathcal{R}(X)$-Module localement libre signifie que, pour tout $x\in X$,
\emph{les rangs $k_\alpha$ correspondant aux indices tels que $x\in X_\alpha$
sont égaux}. La question est en effet locale, et on est ramené au cas
$X=\operatorname{Spec}(C)$, où $C$ est un anneau noethérien réduit, et
$\mathcal{B}=\widetilde{D}$, où $D$ est une $C$-algèbre qui est un $C$-module
de type fini~; si $\mathfrak{p}_i$ ($1\leq i\leq m$) sont les idéaux premiers
minimaux de $C$, l'anneau total des fractions $L$ de $C$ est composé direct
des corps des fractions $K_i$ des anneaux intègres $C/\mathfrak{p}_i$, et
$D\otimes_C L$ est somme directe des $D\otimes_C K_i$, d'où la conclusion.

Il est clair que sous les hypothèses (I) ou (II), on a défini ainsi un
\emph{homomorphisme de faisceaux de monoïdes multiplicatifs}
$N_{\mathcal{B}/\mathcal{A}}:\mathcal{B}\to\mathcal{A}$, aussi noté $N$ si
aucune confusion n'en résulte, et appelé l'homomorphisme \emph{norme}. Pour
deux sections $f$, $g$ de $\mathcal{B}$ au-dessus d'un même ouvert $U$, on a
donc
\[
\label{II.6.5.1.1-fr}
N_{\mathcal{B}/\mathcal{A}}(fg)
=N_{\mathcal{B}/\mathcal{A}}(f)N_{\mathcal{B}/\mathcal{A}}(g)
\tag{6.5.1.1}
\]
pour les sections correspondantes de $\mathcal{A}$ au-dessus de $U$~;
\[
\label{II.6.5.1.2-fr}
N_{\mathcal{B}/\mathcal{A}}(1_{\mathcal{B}})=1_{\mathcal{A}}~;
\tag{6.5.1.2}
\]
enfin, pour toute section $s$ de $\mathcal{A}$ au-dessus de $U$
\[
\label{II.6.5.1.3-fr}
N_{\mathcal{B}/\mathcal{A}}(s\cdot1_{\mathcal{B}})=s^n
\tag{6.5.1.3}
\]
si le rang de $\mathcal{B}$ est constant et égal à $n$ (pour
$s=0_{\mathcal{A}}$, cette formule donne
$N(0_{\mathcal{B}})=0_{\mathcal{A}}$ si $n\geq1$,
$N(0_{\mathcal{B}})=N(1_{\mathcal{B}})=1_{\mathcal{A}}$ si $n=0$).

Dans l'hypothèse (I), pour que $f\in\Gamma(U,\mathcal{B})$ soit inversible,
il faut et il suffit que $N(f)\in\Gamma(U,\mathcal{A})$ le soit. Dans
l'hypothèse (II), cette condition est nécessaire~; elle est suffisante (par
(\hyperref[II.6.4.7-fr]{6.4.7})) lorsqu'on suppose que
$\mathcal{B}\to\mathcal{B}\otimes_{\mathcal{A}}\mathcal{R}(X)$ est
\emph{injectif} et que l'hypothèse plus restrictive suivante est vérifiée~:
\begin{enumerate}
\item[(II \emph{bis})] \emph{$(X,\mathcal{A})$ est un préschéma localement
noethérien réduit, $\mathcal{B}$ est un $\mathcal{A}$-Module cohérent tel que
$\mathcal{B}\otimes_{\mathcal{A}}\mathcal{R}(X)$ soit un
$\mathcal{R}(X)$-Module localement libre, et pour toute section
$f\in\Gamma(U,\mathcal{B})$ au-dessus d'un ouvert tel que
$\mathcal{B}\otimes_{\mathcal{A}}\mathcal{R}(X)|U$ soit de rang constant $n$
sur $\mathcal{R}(X)|U$, les $\sigma_j(f)$ ($1\leq j\leq n$) sont des sections
de $\mathcal{A}$ au-dessus de $U$.}
\end{enumerate}
(On notera encore que cette condition est vérifiée lorsque $X$ est
\emph{normal}.)
\end{env}

\begin{env}[6.5.2]
\label{II.6.5.2-fr}
Supposons vérifiée une des hypothèses (I), (II) de
(\hyperref[II.6.5.1-fr]{6.5.1}), et soit $\mathcal{L}'$ un
$\mathcal{B}$-Module \emph{inversible}. Nous allons lui associer canoniquement
(à un isomorphisme unique près) un $\mathcal{A}$-Module \emph{inversible}
$\mathcal{L}$ de la façon suivante. Désignons par $\mathcal{A}^*$ (resp.
$\mathcal{B}^*$) le sous-faisceau de $\mathcal{A}$ (resp. $\mathcal{B}$) tel
que $\Gamma(U,\mathcal{A}^*)$ (resp. $\Gamma(U,\mathcal{B}^*)$) soit
l'ensemble des

\oldpage[II]{127}

éléments inversibles de $\Gamma(U,\mathcal{A})$ (resp.
$\Gamma(U,\mathcal{B})$) pour tout ouvert $U\subset X$~; ce sont des faisceaux
de groupes multiplicatifs, et $N_{\mathcal{B}/\mathcal{A}}$, restreint à
$\mathcal{B}^*$, est un \emph{homomorphisme}
$\mathcal{B}^*\to\mathcal{A}^*$ de faisceaux de groupes
(\hyperref[II.6.5.1-fr]{6.5.1}). Soit $\mathfrak{L}$ l'ensemble des couples
$(U_\lambda,\eta_\lambda)$, ayant la propriété suivante~: $U_\lambda$ est un
ouvert de $X$, et $\eta_\lambda$ est un isomorphisme
$\mathcal{L}'|U_\lambda\simeq\mathcal{B}|U_\lambda$ de
$(\mathcal{B}|U_\lambda)$-Modules. Par hypothèse, les $U_\lambda$ forment un
recouvrement de $X$~; pour deux indices quelconques $\lambda$, $\mu$, on pose
$\omega_{\lambda\mu}=(\eta_\lambda|U_\lambda\cap U_\mu)
\circ(\eta_\mu|U_\lambda\cap U_\mu)^{-1}$, automorphisme de
$\mathcal{B}|U_\lambda\cap U_\mu$, qui s'identifie canoniquement à une section
de $\mathcal{B}^*$ au-dessus de $U_\lambda\cap U_\mu$, et
$(\omega_{\lambda\mu})$ est un $1$-cocycle du recouvrement
$\mathfrak{U}=(U_\lambda)$ à valeurs dans $\mathcal{B}^*$
(\hyperref[0.5.4.7-fr]{0, 5.4.7}). Le fait que
$N_{\mathcal{B}/\mathcal{A}}:\mathcal{B}^*\to\mathcal{A}^*$ est un
homomorphisme entraîne alors que
$(N_{\mathcal{B}/\mathcal{A}}\omega_{\lambda\mu})$ est un $1$-cocycle de
$\mathfrak{U}$ à valeurs dans $\mathcal{A}^*$, et il lui correspond donc (à un
isomorphisme unique près) un $\mathcal{A}$-Module inversible $\mathcal{L}$
(\hyperref[0.5.4.7-fr]{0, 5.4.7}) que nous désignerons par
$N_{\mathcal{B}/\mathcal{A}}(\mathcal{L}')$ et appellerons \emph{la norme du
$\mathcal{B}$-Module inversible $\mathcal{L}'$}.

Soit $\mathfrak{M}$ une partie de $\mathfrak{L}$ telle que les $U_\lambda$
correspondants forment encore un recouvrement de $X$, et soit $\mathfrak{B}$
ce recouvrement~; la restriction du cocycle $(\omega_{\lambda\mu})$ à
$\mathfrak{B}$ définit un $1$-cocycle
$(N_{\mathcal{B}/\mathcal{A}}\omega_{\lambda\mu})$ de $\mathfrak{B}$ à
valeurs dans $\mathcal{A}^*$, restriction du $1$-cocycle
$(N_{\mathcal{B}/\mathcal{A}}\omega_{\lambda\mu})$ de $\mathfrak{U}$~; il est
clair qu'il y a un isomorphisme canonique du $\mathcal{A}$-Module inversible
défini par ce $1$-cocycle de $\mathfrak{B}$ sur
$N_{\mathcal{B}/\mathcal{A}}(\mathcal{L}')$, ce qui permet de définir ce
$\mathcal{A}$-Module inversible par un sous-recouvrement quelconque de
$\mathfrak{U}$. Cette possibilité montre aussitôt que, si
$\mathcal{L}'_1$, $\mathcal{L}'_2$ sont deux $\mathcal{B}$-Modules inversibles,
on a, en vertu de (\hyperref[II.6.5.1.1-fr]{6.5.1.1}) et
(\hyperref[II.6.5.1.2-fr]{6.5.1.2}),
\[
\label{II.6.5.2.1-fr}
N(\mathcal{L}'_1\otimes_{\mathcal{B}}\mathcal{L}'_2)
=N(\mathcal{L}'_1)\otimes_{\mathcal{A}}N(\mathcal{L}'_2)
\tag{6.5.2.1}
\]
et
\[
\label{II.6.5.2.2-fr}
N_{\mathcal{B}/\mathcal{A}}(\mathcal{B})=\mathcal{A}
\tag{6.5.2.2}
\]
ainsi que
\[
\label{II.6.5.2.3-fr}
N((\mathcal{L}')^{-1})=(N(\mathcal{L}'))^{-1}
\tag{6.5.2.3}
\]
à des isomorphismes canoniques près. En outre il résulte de
(\hyperref[II.6.5.1.3-fr]{6.5.1.3}) que si $\mathcal{L}$ est un
$\mathcal{A}$-Module inversible et si $\mathcal{B}$ est de rang constant $n$
sur $\mathcal{A}$ dans le cas (I) (resp.
$\mathcal{B}\otimes_{\mathcal{A}}\mathcal{R}(X)$ de rang constant $n$ sur
$\mathcal{R}(X)$ dans le cas (II)), on a, à un isomorphisme canonique près
\[
\label{II.6.5.2.4-fr}
N_{\mathcal{B}/\mathcal{A}}(\mathcal{L}\otimes_{\mathcal{A}}\mathcal{B})
=\mathcal{L}^{\otimes n}.
\tag{6.5.2.4}
\]
\end{env}

\begin{env}[6.5.3]
\label{II.6.5.3-fr}
Montrons maintenant que $N_{\mathcal{B}/\mathcal{A}}(\mathcal{L}')$ est un
\emph{foncteur covariant} dans la catégorie des $\mathcal{B}$-Modules
inversibles. Soit en effet $h':\mathcal{L}'_1\to\mathcal{L}'_2$ un
homomorphisme de $\mathcal{B}$-Modules inversibles, et soit
$\mathfrak{B}=(U_\lambda)$ un recouvrement ouvert de $X$ tel que pour tout
$\lambda$, on ait des $(\mathcal{B}|U_\lambda)$-isomorphismes
$\eta^{(1)}_\lambda:\mathcal{L}'_1|U_\lambda\simeq\mathcal{B}|U_\lambda$ et
$\eta^{(2)}_\lambda:\mathcal{L}'_2|U_\lambda\simeq\mathcal{B}|U_\lambda$~;
il y a donc pour chaque $\lambda$ un endomorphisme $h'_\lambda$ de
$\mathcal{B}|U_\lambda$ tel que
$h'_\lambda\circ\eta^{(1)}_\lambda
=\eta^{(2)}_\lambda\circ(h'|U_\lambda)$, et on peut évidemment identifier
$h'_\lambda$ à une section de $\mathcal{B}$ au-dessus de $U_\lambda$
(\hyperref[0.5.1.1-fr]{0, 5.1.1}). Donc, pour tout couple $(\lambda,\mu)$
d'indices, les restrictions à $U_\lambda\cap U_\mu$ de
$(\eta^{(2)}_\lambda)^{-1}\circ h'_\lambda\circ\eta^{(1)}_\lambda$ et de
$(\eta^{(2)}_\mu)^{-1}\circ h'_\mu\circ\eta^{(1)}_\mu$ coïncident. On en
déduit pour les $1$-cocycles $(\omega^{(1)}_{\lambda\mu})$ et
$(\omega^{(2)}_{\lambda\mu})$ à valeurs dans $\mathcal{B}^*$ correspondant à
$\mathcal{L}'_1$ et $\mathcal{L}'_2$, la relation
\[
\omega^{(2)}_{\lambda\mu}h'_\mu
=h'_\lambda\omega^{(1)}_{\lambda\mu}.
\]

\oldpage[II]{128}

Si on pose $h_\lambda=N(h'_\lambda)$, on aura donc les relations analogues
\[
N(\omega^{(2)}_{\lambda\mu})h_\mu
=h_\lambda N(\omega^{(1)}_{\lambda\mu})
\]
et par suite les $h_\lambda$ définissent un homomorphisme
$N(\mathcal{L}'_1)\to N(\mathcal{L}'_2)$ que nous désignerons par
$N_{\mathcal{B}/\mathcal{A}}(h')$ ou $N(h')$. Dans l'hypothèse (I), pour que
$h'$ soit un \emph{isomorphisme}, il faut et il suffit (puisqu'il s'agit d'une
question locale) que $N_{\mathcal{B}/\mathcal{A}}(h')$ le soit. Dans
l'hypothèse (II), cette condition est encore nécessaire~; elle est suffisante
si l'hypothèse (II \emph{bis}) est vérifiée et si
$\mathcal{B}\to\mathcal{B}\otimes_{\mathcal{A}}\mathcal{R}(X)$ est injectif.

Prenons en particulier $\mathcal{L}'_1=\mathcal{B}$~; les homomorphismes
$\mathcal{B}\to\mathcal{L}'$ s'identifient alors
(\hyperref[0.5.1.1-fr]{0, 5.1.1}) aux sections de $\mathcal{L}'$ au-dessus de
$X$, d'où une application canonique
\[
N_{\mathcal{B}/\mathcal{A}}:\Gamma(X,\mathcal{L}')
\to\Gamma(X,N_{\mathcal{B}/\mathcal{A}}(\mathcal{L}')).
\]
Il résulte encore de (\hyperref[II.6.5.1.1-fr]{6.5.1.1}) que si
$f'_1\in\Gamma(X,\mathcal{L}'_1)$, $f'_2\in\Gamma(X,\mathcal{L}'_2)$, on a
\[
\label{II.6.5.3.1-fr}
N(f'_1\otimes f'_2)=N(f'_1)\otimes N(f'_2).
\tag{6.5.3.1}
\]
Pour tout $\mathcal{A}$-Module inversible $\mathcal{L}$ et toute section
$f\in\Gamma(X,\mathcal{L})$, on a, compte tenu de
(\hyperref[II.6.5.2.4-fr]{6.5.2.4})
\[
\label{II.6.5.3.2-fr}
N_{\mathcal{B}/\mathcal{A}}(f\otimes1_{\mathcal{B}})=f^{\otimes n}
\tag{6.5.3.2}
\]
lorsque $\mathcal{B}$ est de rang constant $n$ dans l'hypothèse (I) (resp.
lorsque $\mathcal{B}\otimes_{\mathcal{A}}\mathcal{R}(X)$ est de rang constant
$n$ dans l'hypothèse (II)). Enfin, pour que l'homomorphisme
$\mathcal{B}\to\mathcal{L}'$ correspondant à une section $f'$ de
$\mathcal{L}'$ au-dessus de $X$ soit un isomorphisme, il faut et il suffit que
$f'_x$ soit une base de $\mathcal{L}'_x$ pour tout $x\in X$~; dans l'hypothèse
(I), cette condition équivaut donc à dire que $(N(f'))_x$ est une base de
$(N(\mathcal{L}'))_x$ pour tout $x$~; dans l'hypothèse (II), cette condition
est encore nécessaire, et elle est suffisante lorsque $\mathcal{B}$ vérifie
(II \emph{bis}) et que
$\mathcal{B}\to\mathcal{B}\otimes_{\mathcal{A}}\mathcal{R}(X)$ est injectif.
\end{env}

\begin{env}[6.5.4]
\label{II.6.5.4-fr}
Soient $(X,\mathcal{A})$, $(X',\mathcal{A}')$ deux espaces annelés,
$f:X'\to X$ un morphisme, $\mathcal{B}$ une $\mathcal{A}$-Algèbre,
$\mathcal{B}'=f^*(\mathcal{B})$ la $\mathcal{A}'$-Algèbre image réciproque. On
suppose vérifiée l'une des hypothèses suivantes~:
\begin{enumerate}
\item[$1^\circ$] $\mathcal{B}$ vérifie l'hypothèse (I) de
(\hyperref[II.6.5.1-fr]{6.5.1}).
\item[$2^\circ$] $(X,\mathcal{A})$ et $\mathcal{B}$ vérifient l'hypothèse
(II) de (\hyperref[II.6.5.1-fr]{6.5.1}), $(X',\mathcal{A}')$ est un préschéma
localement noethérien réduit, et si l'on désigne par $X_\alpha$ et
$X'_\beta$ les sous-préschémas fermés réduits respectifs de $X$ et $X'$ ayant
pour espaces sous-jacents les composantes irréductibles de ces espaces, la
restriction de $f$ à chaque $X'_\beta$ est un morphisme \emph{dominant} de
$X'_\beta$ dans un $X_\alpha$.
\end{enumerate}
Sous ces conditions, $\mathcal{B}'$ vérifie respectivement l'hypothèse (I) ou
l'hypothèse (II) de (\hyperref[II.6.5.1-fr]{6.5.1})~; la première assertion
est immédiate. Pour établir la seconde, il suffit de voir que pour tout
$x'\in X'$, les rangs des
$\mathcal{B}'\otimes_{\mathcal{O}_{X'}}\mathcal{R}(X'_\beta)$ pour les $\beta$
tels que $x'\in X'_\beta$ sont \emph{les mêmes}. Or, si la restriction de $f$
à $X'_\beta$ est un morphisme dominant dans $X_\alpha$, le rang de
$\mathcal{B}'\otimes_{\mathcal{O}_{X'}}\mathcal{R}(X'_\beta)$ est égal à celui
de $\mathcal{B}\otimes_{\mathcal{O}_X}\mathcal{R}(X_\alpha)$ (on le voit
aussitôt en se ramenant au cas affine comme dans
(\hyperref[II.6.4.9-fr]{6.4.9})), d'où notre assertion en vertu de l'hypothèse
(II) et de (\hyperref[II.6.5.1-fr]{6.5.1}).

\oldpage[II]{129}

Cela étant, il résulte de (\hyperref[II.6.4.8-fr]{6.4.8}),
(\hyperref[II.6.4.9-fr]{6.4.9}) et (\hyperref[II.6.4.10-fr]{6.4.10}) que si
$s$ est une section de $\mathcal{B}$ au-dessus d'un ouvert $U\subset X$,
$s'$ la section correspondante de $\mathcal{B}'$ au-dessus de $f^{-1}(U)$,
$N_{\mathcal{B}'/\mathcal{A}'}(s')$ est la section de $\mathcal{A}'$ au-dessus
de $f^{-1}(U)$ qui correspond à la section
$N_{\mathcal{B}/\mathcal{A}}(s)$ de $\mathcal{A}$ au-dessus de $U$.

Si $\mathcal{M}$ est un $\mathcal{B}$-Module inversible, on déduit de ce qui
précède que si $\mathcal{M}'=f^*(\mathcal{M})$ (qui est un
$\mathcal{B}'$-Module inversible), on a
$N_{\mathcal{B}'/\mathcal{A}'}(\mathcal{M}')
=f^*(N_{\mathcal{B}/\mathcal{A}}(\mathcal{M}))$ à un isomorphisme canonique
près.
\end{env}

\begin{env}[6.5.5]
\label{II.6.5.5-fr}
Supposons désormais que $(X,\mathcal{A})$ soit un \emph{préschéma}. La donnée
d'une $\mathcal{A}$-Algèbre quasi-cohérente $\mathcal{B}$, qui est un
$\mathcal{A}$-Module \emph{de type fini} équivaut alors, comme on sait, à
celle d'un morphisme \emph{fini} $g:X'\to X$ tel que
$g_*(\mathcal{O}_{X'})=\mathcal{B}$, défini à un $X$-isomorphisme près
(\hyperref[II.6.1.2-fr]{6.1.2} et \hyperref[II.1.3.1-fr]{1.3.1}). En outre,
se donner un $\mathcal{O}_{X'}$-Module quasi-cohérent $\mathcal{F}'$ équivaut
à se donner un $\mathcal{B}$-Module quasi-cohérent $\mathcal{F}$ tel que
$g_*(\mathcal{F}')=\mathcal{F}$ (\hyperref[II.1.4.3-fr]{1.4.3}), et pour que
$\mathcal{F}'$ soit inversible, il faut et il suffit que $\mathcal{F}$ le soit
(\hyperref[II.6.1.12-fr]{6.1.12}). Pour traduire les résultats précédents en
termes du morphisme fini $g$, il faudra donc supposer, soit que
$g_*(\mathcal{O}_{X'})$ est un $\mathcal{O}_X$-Module \emph{localement libre}
(de type fini), soit que $(X,\mathcal{O}_X)$ et $g_*(\mathcal{O}_{X'})$
vérifient l'hypothèse (II). Pour tout $\mathcal{O}_{X'}$-Module inversible
$\mathcal{L}'$, on posera alors
\[
\label{II.6.5.5.1-fr}
N_{X'/X}(\mathcal{L}')
=N_{g_*(\mathcal{O}_{X'})/\mathcal{O}_X}(g_*(\mathcal{L}'))
\tag{6.5.5.1}
\]
et on dira que c'est la \emph{norme} (relative à $g$) de $\mathcal{L}'$. De
même, si $h':\mathcal{L}'_1\to\mathcal{L}'_2$ est un homomorphisme de
$\mathcal{O}_{X'}$-Modules inversibles, on posera
\[
\label{II.6.5.5.2-fr}
N_{X'/X}(h')
=N_{g_*(\mathcal{O}_{X'})/\mathcal{O}_X}(g_*(h')):
N_{X'/X}(\mathcal{L}'_1)\to N_{X'/X}(\mathcal{L}'_2).
\tag{6.5.5.2}
\]
En particulier, pour $\mathcal{L}'_1=\mathcal{O}_{X'}$, on obtient ainsi une
application canonique
\[
\label{II.6.5.5.3-fr}
N_{X'/X}:\Gamma(X',\mathcal{L}')
\to\Gamma(X,N_{X'/X}(\mathcal{L}')).
\tag{6.5.5.3}
\]
Nous laisserons au lecteur la plupart de ces traductions, et nous bornerons à
expliciter les suivantes~:
\end{env}

\begin{proposition}[6.5.6]
\label{II.6.5.6-fr}
Soit $g:X'\to X$ un morphisme fini, et supposons, soit que
$g_*(\mathcal{O}_{X'})$ est un $\mathcal{O}_X$-Module localement libre, soit
que $(X,\mathcal{O}_X)$ et $g_*(\mathcal{O}_{X'})$ vérifient (II \emph{bis})
(ce qui sera le cas en particulier lorsque $X$ est localement noethérien et
normal). Pour qu'un homomorphisme $h':\mathcal{L}'_1\to\mathcal{L}'_2$ de
$\mathcal{O}_{X'}$-Modules inversibles soit un isomorphisme, il faut et il
suffit, dans la première hypothèse, que $N_{X'/X}(h')$ soit un isomorphisme~;
dans la seconde hypothèse, cette condition est encore nécessaire, et elle est
suffisante lorsque l'homomorphisme
$g_*(\mathcal{O}_{X'})\to
g_*(\mathcal{O}_{X'})\otimes_{\mathcal{O}_X}\mathcal{R}(X)$ est injectif.
\end{proposition}

On notera qu'on utilise ici le fait que, pour que $g_*(h')$ soit un
isomorphisme, il faut et il suffit que $h'$ en soit un
(\hyperref[II.1.4.2-fr]{1.4.2}).

\begin{corollary}[6.5.7]
\label{II.6.5.7-fr}
Soit $g:X'\to X$ un morphisme fini, et supposons, soit que
$g_*(\mathcal{O}_{X'})$ est un $\mathcal{O}_X$-Module localement libre, soit
que $(X,\mathcal{O}_X)$ et $g_*(\mathcal{O}_{X'})$ vérifient (II \emph{bis})
et que
$g_*(\mathcal{O}_{X'})\to
g_*(\mathcal{O}_{X'})\otimes_{\mathcal{O}_X}\mathcal{R}(X)$ est injectif.
Soient $\mathcal{L}'$ un $\mathcal{O}_{X'}$-Module inversible, $f'$ une
section de $\mathcal{L}'$ au-dessus de $X'$, $f=N_{X'/X}(f')$ la section de
$\mathcal{L}=N_{X'/X}(\mathcal{L}')$ au-dessus de $X$ qui lui correspond
(\hyperref[II.6.5.5.3-fr]{6.5.5.3}).

\oldpage[II]{130}

Alors on a $g(X'-X'_{f'})=X-X_f$ et $X_f$ est le plus grand ouvert $U$ de
$X$ tel que $g^{-1}(U)\subset X'_{f'}$.
\end{corollary}

En effet, $g(X'-X'_{f'})$ est fermé dans $X$
(\hyperref[II.6.1.10-fr]{6.1.10})~; il suffit donc de prouver la dernière
assertion. Or la relation $U\subset X_f$ équivaut au fait que l'homomorphisme
$\mathcal{O}_X|U\to\mathcal{L}|U$ défini par $f|U$ est un isomorphisme. En
vertu de (\hyperref[II.6.5.6-fr]{6.5.6}), cela équivaut à dire que
l'homomorphisme
$\mathcal{O}_{X'}|g^{-1}(U)\to\mathcal{L}'|g^{-1}(U)$ défini par
$f'|g^{-1}(U)$ est un isomorphisme, c'est-à-dire à la relation
$g^{-1}(U)\subset X'_{f'}$.

\begin{proposition}[6.5.8]
\label{II.6.5.8-fr}
Soient $g:X'\to X$ un morphisme fini, $f:Y\to X$ un morphisme~; soient
$Y'=X'_{(Y)}$, $g'=g_{(Y)}$, $f'=f_{(X')}$ de sorte que l'on a le diagramme
commutatif
\[
\xymatrix{
X'\ar[d]_g & Y'\ar[l]_{f'}\ar[d]^{g'}\\
X & Y\ar[l]^f
}
\]
On suppose, soit que $g_*(\mathcal{O}_{X'})$ est localement libre, soit que
$(X,\mathcal{O}_X)$ et $g_*(\mathcal{O}_{X'})$ vérifient (II), que $Y$ est un
préschéma localement noethérien réduit, et que la restriction de $f$ à toute
composante irréductible de $Y$ est un morphisme dominant dans une composante
irréductible de $X$. Alors, pour tout $\mathcal{O}_{X'}$-Module inversible
$\mathcal{L}'$, on a
\[
N_{Y'/Y}({f'}^*(\mathcal{L}'))=f^*(N_{X'/X}(\mathcal{L}'))
\]
à un isomorphisme canonique près.
\end{proposition}

Notons que l'on a $f^*(g_*(\mathcal{L}'))=g'_*({f'}^*(\mathcal{L}'))$ en
vertu de (\hyperref[II.1.5.2-fr]{1.5.2}), et en particulier
$g'_*(\mathcal{O}_{Y'})=f^*(g_*(\mathcal{O}_{X'}))$~; si
$g_*(\mathcal{O}_{X'})$ est localement libre, il en est donc de même de
$g'_*(\mathcal{O}_{Y'})$. La conclusion résulte alors des définitions et de
(\hyperref[II.6.5.4-fr]{6.5.4}).

\begin{remark}[6.5.9]
\label{II.6.5.9-fr}
Nous généraliserons ultérieurement la notion de norme développée ci-dessus,
et la mettrons en relation avec la notion d'image directe d'un diviseur.
\end{remark}

\subsection{Application : critères d'amplitude.}
\label{subsection:II.6.6-fr}

\begin{proposition}[6.6.1]
\label{II.6.6.1-fr}
Soient $Y$ un préschéma, $f:X\to Y$ un morphisme quasi-compact,
$g:X'\to X$ un morphisme fini et surjectif. On suppose, soit que
$g_*(\mathcal{O}_{X'})$ est un $\mathcal{O}_X$-Module localement libre, soit
que $(X,\mathcal{O}_X)$ et $g_*(\mathcal{O}_{X'})$ vérifient (II
\emph{bis}). Alors, pour tout $\mathcal{O}_{X'}$-Module inversible
$\mathcal{L}'$, ample pour $f\circ g$, $N_{X'/X}(\mathcal{L}')=\mathcal{L}$
est ample pour $f$.
\end{proposition}

On peut supposer $Y$ affine (\hyperref[II.4.6.4-fr]{4.6.4}), et alors, en
vertu de (\hyperref[II.4.6.6-fr]{4.6.6}), l'énoncé est équivalent au

\begin{corollary}[6.6.2]
\label{II.6.6.2-fr}
Soient $X$ un préschéma quasi-compact, $g:X'\to X$ un morphisme fini
surjectif, tel que $g_*(\mathcal{O}_{X'})$ soit un $\mathcal{O}_X$-Module
localement libre, ou que $(X,\mathcal{O}_X)$ et $g_*(\mathcal{O}_{X'})$
vérifient (II \emph{bis}). Alors, pour tout $\mathcal{O}_{X'}$-Module ample
$\mathcal{L}'$, $\mathcal{L}=N_{X'/X}(\mathcal{L}')$ est ample.
\end{corollary}

Dans la seconde hypothèse, on peut supposer en outre que l'homomorphisme
canonique
$g_*(\mathcal{O}_{X'})\to
g_*(\mathcal{O}_{X'})\otimes_{\mathcal{O}_X}\mathcal{R}(X)$ est
\emph{injectif}. En effet, dans le cas contraire, soit $\mathcal{T}$

\oldpage[II]{131}

le noyau de cet homomorphisme, qui est un Idéal cohérent de
$g_*(\mathcal{O}_{X'})=\mathcal{B}$
(\hyperref[I.6.1.1-fr]{I, 6.1.1}), et posons
$X''=\operatorname{Spec}(\mathcal{B}/\mathcal{T})$~; on a donc un diagramme
commutatif
\[
\xymatrix{
X''\ar[rr]^h\ar[dr]_{g'} && X'\ar[dl]^g\\
& X
}
\]
où $h$ est une immersion fermée (\hyperref[II.1.4.10-fr]{1.4.10}). En outre,
on sait que le support de $\mathcal{T}$ est un ensemble fermé
(\hyperref[0.5.2.2-fr]{0, 5.2.2}) rare dans $X$
(\hyperref[I.7.4.6-fr]{I, 7.4.6}), d'où on conclut que pour le point générique
$x$ d'une composante irréductible de $X$, il y a un voisinage ouvert affine
$U$ de $x$ tel que
$\mathcal{B}|U=(\mathcal{B}/\mathcal{T})|U$. Comme $g$ est par hypothèse
surjectif, on en conclut que $x\in g'(X'')$~; $g'$ est donc dominant, et étant
un morphisme fini, il est \emph{surjectif}
(\hyperref[II.6.1.10-fr]{6.1.10})~; on a par définition
\[
g'_*(\mathcal{O}_{X''})\otimes_{\mathcal{O}_X}\mathcal{R}(X)
=(\mathcal{B}/\mathcal{T})\otimes_{\mathcal{O}_X}\mathcal{R}(X)
=g_*(\mathcal{O}_{X'})\otimes_{\mathcal{O}_X}\mathcal{R}(X),
\]
donc $(X,\mathcal{O}_X)$ et $g'_*(\mathcal{O}_{X''})$ vérifient
(II \emph{bis}), et en outre
\[
g'_*(\mathcal{O}_{X''})\longrightarrow
g'_*(\mathcal{O}_{X''})\otimes_{\mathcal{O}_X}\mathcal{R}(X)
\]
est injectif. Enfin, $h^*(\mathcal{L}')=\mathcal{L}''$ est un
$\mathcal{O}_{X''}$-Module ample
(\hyperref[II.4.6.13-fr]{4.6.13, (i \emph{bis})}), et on a
$N_{X''/X}(\mathcal{L}'')=N_{X'/X}(\mathcal{L}')$. En effet, pour définir
ces deux $\mathcal{O}_X$-Modules inversibles on peut utiliser un même
recouvrement ouvert affine $(U_\lambda)$ de $X$ tel que les restrictions de
$g_*(\mathcal{L}')$ et $g'_*(\mathcal{L}'')$ à $U_\lambda$ soient
respectivement isomorphes à $\mathcal{B}|U_\lambda$ et
$(\mathcal{B}/\mathcal{T})|U_\lambda$. On voit aussitôt qu'à tout
isomorphisme
\[
\eta_\lambda:g_*(\mathcal{L}')|U_\lambda\longrightarrow
\mathcal{B}|U_\lambda
\]
correspond canoniquement un isomorphisme
\[
\eta'_\lambda:g'_*(\mathcal{L}'')|U_\lambda\longrightarrow
(\mathcal{B}/\mathcal{T})|U_\lambda,
\]
de sorte que, si $(\omega_{\lambda\mu})$ et
$(\omega'_{\lambda\mu})$ sont les $1$-cocycles correspondant aux systèmes
d'isomorphismes $(\eta_\lambda)$ et $(\eta'_\lambda)$
(\hyperref[II.6.5.2-fr]{6.5.2}), $\omega'_{\lambda\mu}$ soit l'image
canonique dans
$\Gamma(U_\lambda\cap U_\mu,\mathcal{B}/\mathcal{T})$ de
$\omega_{\lambda\mu}\in\Gamma(U_\lambda\cap U_\mu,\mathcal{B})$. En vertu
de la définition de $\mathcal{T}$, on en conclut que
\[
N_{\mathcal{B}/\mathcal{A}}\omega_{\lambda\mu}
=N_{(\mathcal{B}/\mathcal{T})/\mathcal{A}}\omega'_{\lambda\mu}
\]
(avec $\mathcal{A}=\mathcal{O}_X$), d'où l'égalité énoncée.

Supposons donc l'homomorphisme
$g_*(\mathcal{O}_{X'})\to
g_*(\mathcal{O}_{X'})\otimes_{\mathcal{O}_X}\mathcal{R}(X)$ injectif
lorsqu'on est dans l'hypothèse (II \emph{bis}). Il suffit de prouver que
lorsque $f$ parcourt les sections de $\mathcal{L}^{\otimes n}$ ($n>0$)
au-dessus de $X$, les $X_f$ forment une base de la topologie de $X$
(\hyperref[II.4.5.2-fr]{4.5.2}). Or, soit $x\in X$, et soit $U$ un voisinage
quelconque de $x$~; comme $g^{-1}(x)$ est fini
(\hyperref[II.6.1.7-fr]{6.1.7}) et que $\mathcal{L}'$ est ample, il existe un
entier $n>0$ et une section $f'$ de ${\mathcal{L}'}^{\otimes n}$ au-dessus de
$X'$, tels que $X'_{f'}$ soit un voisinage de $g^{-1}(x)$ contenu dans
$g^{-1}(U)$ (\hyperref[II.4.5.4-fr]{4.5.4}). Comme
\[
\mathcal{L}^{\otimes n}=N_{X'/X}({\mathcal{L}'}^{\otimes n}),
\]
il suffit de prendre $f=N_{X'/X}(f')$~; en effet, alors
$X-X_f=g(X'-X'_{f'})$ (\hyperref[II.6.5.7-fr]{6.5.7}), donc
$x\in X_f\subset U$.

\begin{corollary}[6.6.3]
\label{II.6.6.3-fr}
Sous les hypothèses de (\hyperref[II.6.6.1-fr]{6.6.1}), pour qu'un
$\mathcal{O}_X$-Module inversible $\mathcal{L}$ soit ample pour $f$, il faut
et il suffit que $\mathcal{L}'=g^*(\mathcal{L})$ soit ample pour $f\circ g$.
\end{corollary}

La condition est nécessaire, puisque $g$ est affine
(\hyperref[II.5.1.12-fr]{5.1.12}). Pour démontrer que la condition est
suffisante, on peut supposer $Y$ affine (\hyperref[II.4.6.4-fr]{4.6.4}), donc
$X$ et $X'$ quasi-compacts et $\mathcal{L}'$ ample
(\hyperref[II.4.6.6-fr]{4.6.6}), et il faut montrer que $\mathcal{L}$ est
ample. Or, l'ensemble des points

\oldpage[II]{132}

$x\in X$ au voisinage desquels $g_*(\mathcal{O}_{X'})$ (resp.
$g_*(\mathcal{O}_{X'})\otimes_{\mathcal{O}_X}\mathcal{R}(X)$) a un rang
donné $n$ dans la première (resp. la seconde) hypothèse, est à la fois ouvert
et fermé dans $X$, donc $X$ est le préschéma somme d'un nombre fini de ces
ouverts, et on peut par suite supposer qu'il est égal à l'un d'eux
(\hyperref[II.4.6.17-fr]{4.6.17}). Mais on a alors
$N_{X'/X}(\mathcal{L}')=\mathcal{L}^{\otimes n}$, donc
$\mathcal{L}^{\otimes n}$ est ample en vertu de
(\hyperref[II.6.6.2-fr]{6.6.2}), et il en est donc de même de $\mathcal{L}$
(\hyperref[II.4.5.6-fr]{4.5.6}).

\begin{corollary}[6.6.4]
\label{II.6.6.4-fr}
Supposons vérifiées les hypothèses de
(\hyperref[II.6.6.1-fr]{6.6.1}) et supposons en outre que
$f:X\to Y$ soit de type fini. Alors, pour que $f$ soit quasi-projectif, il
faut et il suffit que $f\circ g$ le soit. Si on suppose de plus que $Y$ est
un schéma quasi-compact ou un préschéma dont l'espace sous-jacent est
noethérien, alors, pour que $f$ soit projectif, il faut et il suffit que
$f\circ g$ le soit.
\end{corollary}

L'hypothèse entraîne que $f\circ g$ est de type fini. Compte tenu de la
définition des morphismes quasi-projectifs
(\hyperref[II.5.3.1-fr]{5.3.1}), la première assertion résulte de
(\hyperref[II.6.6.1-fr]{6.6.1}) et
(\hyperref[II.6.6.3-fr]{6.6.3}). Compte tenu de ce résultat et de
(\hyperref[II.5.5.3-fr]{5.5.3, (ii)}), il reste à prouver que lorsque $f$ est
supposé quasi-projectif, pour que $f$ soit propre, il faut et il suffit que
$f\circ g$ le soit. Or $f$ est alors séparé
(\hyperref[II.5.3.1-fr]{5.3.1}) et de type fini~; comme $g$ est surjectif,
notre assertion résulte de (\hyperref[II.5.4.2-fr]{5.4.2, (ii)}) et
(\hyperref[II.5.4.3-fr]{5.4.3, (ii)}).

En particulier~:

\begin{corollary}[6.6.5]
\label{II.6.6.5-fr}
Soient $X$ un préschéma de type fini sur un corps $K$, $K'$ une extension de
degré fini de $K$. Pour que $X$ soit projectif (resp. quasi-projectif) sur
$K$, il faut et il suffit que $X'=X\otimes_K K'$ soit projectif (resp.
quasi-projectif) sur $K'$.
\end{corollary}

La condition est en effet nécessaire
(\hyperref[II.5.3.4-fr]{5.3.4, (iii)} et
\hyperref[II.5.5.5-fr]{5.5.5, (iii)}). Inversement, supposons-la vérifiée,
et soit $g:X'\to X$ la projection canonique. Il est clair que $g$ est un
morphisme fini (\hyperref[II.6.1.5-fr]{6.1.5, (iii)}) et surjectif
(\hyperref[I.3.5.2-fr]{I, 3.5.2, (ii)}). En outre,
$g_*(\mathcal{O}_{X'})$ est un $\mathcal{O}_X$-Module localement libre, étant
isomorphe à $\mathcal{O}_X\otimes_K K'$
(\hyperref[II.1.5.2-fr]{1.5.2}). Il résulte alors de l'hypothèse et de
(\hyperref[II.6.1.11-fr]{6.1.11}) et
(\hyperref[II.5.5.5-fr]{5.5.5, (ii)}) que $X'$ est projectif (resp.
quasi-projectif) \emph{sur $K$}~; on déduit alors de
(\hyperref[II.6.6.4-fr]{6.6.4}) que $X$ est projectif (resp.
quasi-projectif) sur $K$.

Au chapitre V, nous montrerons que l'énoncé
(\hyperref[II.6.6.5-fr]{6.6.5}) reste valable lorsque $K'$ est une extension
\emph{quelconque} de $K$.

La fin de ce numéro est consacrée à la démonstration du critère
(\hyperref[II.6.6.11-fr]{6.6.11}), qui est un raffinement assez technique de
(\hyperref[II.6.6.1-fr]{6.6.1})~; elle peut être omise en première lecture.

\begin{lemma}[6.6.6]
\label{II.6.6.6-fr}
Soient $X$ un préschéma noethérien réduit, $\mathcal{E}$ un
$\mathcal{O}_X$-Module cohérent tel que
$\mathcal{E}\otimes_{\mathcal{O}_X}\mathcal{R}(X)$ soit un
$\mathcal{R}(X)$-Module localement libre de rang constant $n$. Il existe
alors un préschéma noethérien réduit $Z$ et un morphisme fini birationnel
$h:Z\to X$ ayant la propriété suivante~: les morphismes de faisceaux
d'ensembles
$\sigma_i:\mathcal{H}om_{\mathcal{O}_X}(\mathcal{E},\mathcal{E})
\to\mathcal{R}(X)$ ($1\leq i\leq n$)
(cf. \hyperref[II.6.4.10-fr]{6.4.10}) appliquent
$\mathcal{H}om_{\mathcal{O}_X}(\mathcal{E},\mathcal{E})$ dans la
$\mathcal{O}_X$-Algèbre cohérente $h_*(\mathcal{O}_Z)$.
\end{lemma}

Considérons en effet un ouvert affine $U$ de $X$, d'anneau $A(U)=A$~; soit
$E=\Gamma(U,\mathcal{E})$, et soit $C_U$ la sous-algèbre de $R(U)$ engendrée
par les $\sigma_i(u)$ lorsque $u$ parcourt $\operatorname{Hom}_A(E,E)$~; on
a vu (\hyperref[II.6.4.7.1-fr]{6.4.7.1}) que cette $A$-algèbre est de
\emph{rang fini}. En outre, il est clair que la formation des algèbres $C_U$
commute avec les opérations de restriction d'un ouvert affine $U$ à un ouvert
affine $U'\subset U$. On a donc défini ainsi une sous-$\mathcal{O}_X$-Algèbre
\emph{finie} $\mathcal{C}$ de $\mathcal{R}(X)$ telle que
$\Gamma(U,\mathcal{C})=C_U$ pour tout ouvert affine $U$ de $X$. On

\oldpage[II]{133}

prendra $Z=\operatorname{Spec}(\mathcal{C})$, et pour $h$ le morphisme
structural, qui est donc fini (\hyperref[II.6.1.2-fr]{6.1.2})~; comme
$\mathcal{C}$ est réduite, $Z$ est un préschéma noethérien réduit
(\hyperref[II.1.3.8-fr]{1.3.8}). Enfin, l'anneau total des fractions de $C_U$
est $R(U)$ par définition, et comme $C_U$ est contenu dans la fermeture
intégrale de $A(U)$ dans $R(U)$, il y a correspondance biunivoque entre idéaux
premiers minimaux de $A(U)$ et idéaux premiers minimaux de $C_U$
(\cite[t.~I, p.~259]{I-13}), ce qui prouve que $h$ est birationnel et achève
la démonstration.

\begin{corollary}[6.6.7]
\label{II.6.6.7-fr}
Sous les hypothèses de (\hyperref[II.6.6.6-fr]{6.6.6}), soit $W$ un ouvert
de $X$ tel que, pour tout $x\in W$, ou bien $X$ est normal au point $x$, ou
bien $\mathcal{E}_x$ est un $\mathcal{O}_x$-Module libre. Alors on peut
supposer $h$ défini de sorte que la restriction de $h$ à $h^{-1}(W)$ soit un
isomorphisme de $h^{-1}(W)$ sur $W$.
\end{corollary}

En effet, l'une ou l'autre hypothèse entraîne que si $U\subset W$ est un
ouvert affine, on a, avec les notations de
(\hyperref[II.6.6.6-fr]{6.6.6}), $(\sigma_i(u))_x\in A_x$ pour tout
$x\in U$ (\hyperref[II.6.4.3-fr]{6.4.3}), donc $\sigma_i(u)\in A$, et la
conclusion résulte de la définition de $h$ donnée dans
(\hyperref[II.6.6.6-fr]{6.6.6}).

\begin{env}[6.6.8]
\label{II.6.6.8-fr}
Soient $X$ un préschéma noethérien réduit, $g:X'\to X$ un morphisme fini
surjectif, de sorte que $\mathcal{B}=g_*(\mathcal{O}_{X'})$ est une
$\mathcal{O}_X$-Algèbre \emph{cohérente}~; supposons en outre que
$\mathcal{B}\otimes_{\mathcal{O}_X}\mathcal{R}(X)$ soit un
$\mathcal{R}(X)$-Module \emph{localement libre de rang constant $n$}. On peut
alors appliquer le lemme (\hyperref[II.6.6.6-fr]{6.6.6}) en prenant
$\mathcal{E}=\mathcal{B}$, d'où, avec les notations de
(\hyperref[II.6.6.6-fr]{6.6.6}), un homomorphisme de faisceaux de monoïdes
multiplicatifs
$\sigma_n:\mathcal{H}om_{\mathcal{O}_X}(\mathcal{B},\mathcal{B})
\to h_*(\mathcal{O}_Z)$, et en composant cet homomorphisme avec
l'homomorphisme canonique
$\mathcal{B}\to
\mathcal{H}om_{\mathcal{O}_X}(\mathcal{B},\mathcal{B})$
(\hyperref[II.6.5.1-fr]{6.5.1}), on obtient donc un homomorphisme de
faisceaux de monoïdes multiplicatifs~:
\[
\label{II.6.6.8.1-fr}
N':\mathcal{B}=g_*(\mathcal{O}_{X'})\longrightarrow
h_*(\mathcal{O}_Z)=\mathcal{C}.
\tag{6.6.8.1}
\]

Cela étant, pour tout $\mathcal{O}_{X'}$-Module inversible $\mathcal{L}'$,
$g_*(\mathcal{L}')$ est un $\mathcal{B}$-Module inversible
(\hyperref[II.6.1.12-fr]{6.1.12}), et la méthode de
(\hyperref[II.6.5.2-fr]{6.5.2}) permet d'associer fonctoriellement à
$\mathcal{L}'$ un $\mathcal{C}$-Module inversible que nous noterons
$N'(g_*(\mathcal{L}'))$.
\end{env}

\begin{lemma}[6.6.9]
\label{II.6.6.9-fr}
Soient $X$ un préschéma noethérien réduit, $g:X'\to X$ un morphisme fini
surjectif tel que
$g_*(\mathcal{O}_{X'})\otimes_{\mathcal{O}_X}\mathcal{R}(X)$ soit un
$\mathcal{R}(X)$-Module localement libre de rang constant. Il existe alors
un préschéma noethérien réduit $Z$ et un morphisme fini birationnel
$h:Z\to X$ ayant la propriété suivante~: pour tout
$\mathcal{O}_{X'}$-Module ample $\mathcal{L}'$, le $\mathcal{O}_Z$-Module
inversible $\mathcal{M}$ tel que
$h_*(\mathcal{M})=N'(g_*(\mathcal{L}'))$ (notations de
\hyperref[II.6.6.8-fr]{6.6.8}) est ample.
\end{lemma}

Supposons d'abord que l'homomorphisme
$\mathcal{B}\to
\mathcal{B}\otimes_{\mathcal{O}_X}\mathcal{R}(X)$ soit \emph{injectif}.
Définissons $Z$ et $h$ comme dans (\hyperref[II.6.6.6-fr]{6.6.6}) (avec
$\mathcal{E}=g_*(\mathcal{O}_{X'})$). Soit $z\in Z$~; il faut montrer qu'il
existe un entier $m>0$ et une section $t$ de $\mathcal{M}^{\otimes m}$
au-dessus de $Z$ telle que $Z_t$ soit un ouvert affine contenant $z$
(\hyperref[II.4.5.2-fr]{4.5.2}). Soient $x=h(z)$, $U$ un ouvert affine de
$X$ contenant $x$~; $h^{-1}(U)$ est alors un voisinage ouvert affine de $z$,
et il suffira de trouver $t$ tel que $z\in Z_t\subset h^{-1}(U)$, car $Z_t$
sera alors nécessairement affine (\hyperref[II.5.5.8-fr]{5.5.8}). Il existe
par hypothèse un entier $n>0$ et une section $s'$ de
${\mathcal{L}'}^{\otimes n}$ au-dessus de $X'$ telle que l'on ait
\[
\label{II.6.6.9.1-fr}
g^{-1}(x)\subset X'_{s'}\subset g^{-1}(U)
\tag{6.6.9.1}
\]
en vertu de (\hyperref[II.4.5.4-fr]{4.5.4}). Par définition, $s'$ est aussi
une section de $g_*({\mathcal{L}'}^{\otimes n})$ au-dessus de $X$, et il lui
correspond comme dans (\hyperref[II.6.5.2-fr]{6.5.2}) une section
$s=N'(s')$ de $N'(g_*({\mathcal{L}'}^{\otimes n}))$ au-dessus

\oldpage[II]{134}

de $X$. Nous allons voir que si $t$ est la section $s$ considérée comme
section de $\mathcal{M}^{\otimes n}$ au-dessus de $Z$, $t$ répond à la
question. Posons en effet
\[
\label{II.6.6.9.2-fr}
V=X-g(X'-X'_{s'})
\tag{6.6.9.2}
\]
qui est un ouvert de $X$ contenant $x$ et contenu dans $U$, en vertu de
(\hyperref[II.6.6.9.1-fr]{6.6.9.1}) et
(\hyperref[II.6.1.10-fr]{6.1.10}). Nous allons montrer que l'on a
\[
\label{II.6.6.9.3-fr}
h^{-1}(V)\subset Z_t\subset h^{-1}(U),
\tag{6.6.9.3}
\]
ce qui achèvera la démonstration. Il revient au même de dire que l'ensemble
$T$ des $y\in X$ tels que $s_y$ soit inversible, contient $V$ et est contenu
dans $U$. Pour cela, considérons d'abord un ouvert affine $W$ contenu dans
$V$~; $g^{-1}(W)$ est ouvert affine dans $X'$ et en vertu de
(\hyperref[II.6.6.9.2-fr]{6.6.9.2}), $s'_{y'}$ est inversible pour tout
$y'\in g^{-1}(W)$~; en vertu des hypothèses sur $X$ et $\mathcal{B}$, on
peut appliquer les résultats de (\hyperref[II.6.4.7-fr]{6.4.7}) et on voit
que si $y=g(y')$, $s_y$ est inversible, autrement dit $V\subset T$. D'autre
part, il résulte aussi de (\hyperref[II.6.4.7-fr]{6.4.7}) que, inversement,
si $s_y$ est inversible, il en est de même de $s'_{y'}$, ce qui implique
$y'\in g^{-1}(U)$ par (\hyperref[II.6.6.9.1-fr]{6.6.9.1}) et par suite
$y\in U$, d'où $T\subset U$ dans ce cas.

On passe de là au cas général par le même raisonnement que dans
(\hyperref[II.6.6.2-fr]{6.6.2}), remplaçant $X'$ par $X''$ tel que
\[
g'_*(\mathcal{O}_{X''})\longrightarrow
g'_*(\mathcal{O}_{X''})\otimes_{\mathcal{O}_X}\mathcal{R}(X)
\]
soit injectif, et $\mathcal{L}'$ par un $\mathcal{O}_{X''}$-Module ample
$\mathcal{L}''$ tel que
$N'(g_*(\mathcal{L}'))=N'(g'_*(\mathcal{L}''))$. Le lemme
(\hyperref[II.6.6.9-fr]{6.6.9}) est donc démontré dans tous les cas (avec un
choix convenable de $h$).

\begin{corollary}[6.6.10]
\label{II.6.6.10-fr}
Supposons vérifiées les hypothèses de
(\hyperref[II.6.6.9-fr]{6.6.9})~; pour tout $\mathcal{O}_X$-Module inversible
$\mathcal{L}$ tel que $g^*(\mathcal{L})$ soit ample, $h^*(\mathcal{L})$ est
ample.
\end{corollary}

En effet, si on pose $\mathcal{L}'=g^*(\mathcal{L})$, on a alors
$g_*(\mathcal{L}')=\mathcal{L}\otimes_{\mathcal{O}_X}\mathcal{B}$
(\hyperref[0.5.4.10-fr]{0, 5.4.10}), donc
\[
N'(g_*(\mathcal{L}'))=
(\mathcal{L}\otimes_{\mathcal{O}_X}\mathcal{C})^{\otimes n}
\]
(par le même raisonnement que pour
\hyperref[II.6.5.2.4-fr]{6.5.2.4}). On en conclut que l'on a
$\mathcal{M}=(h^*(\mathcal{L}))^{\otimes n}$, et comme $\mathcal{M}$ est
ample, il en est de même de $h^*(\mathcal{L})$
(\hyperref[II.4.5.6-fr]{4.5.6}).

\begin{proposition}[6.6.11]
\label{II.6.6.11-fr}
Soient $Y$ un schéma affine, $X$ un préschéma noethérien réduit,
$f:X\to Y$ un morphisme quasi-compact, $g:X'\to X$ un morphisme fini
surjectif. Soit $W$ une partie ouverte de $X$ telle que, pour tout $x\in W$,
ou bien $X$ est normal au point $x$, ou bien il existe un voisinage ouvert
$T\subset W$ de $x$ tel que $(g_*(\mathcal{O}_{X'}))|T$ soit un
$(\mathcal{O}_X|T)$-Module localement libre. Il existe alors un
$Y$-préschéma réduit $Z$ et un $Y$-morphisme fini et birationnel
$h:Z\to X$ tel que la restriction de $h$ à $h^{-1}(W)$ soit un isomorphisme
$h^{-1}(W)\xrightarrow{\sim}W$ et qui possède la propriété suivante~: pour
tout $\mathcal{O}_X$-Module inversible $\mathcal{L}$ tel que
$g^*(\mathcal{L})$ soit ample relativement à $f\circ g$, $h^*(\mathcal{L})$
est ample relativement à $f\circ h$.
\end{proposition}

Comme $Y$ est affine, $g^*(\mathcal{L})$ est ample, et il s'agit alors de
prouver que pour un choix convenable de $h$, $h^*(\mathcal{L})$ est ample
(\hyperref[II.4.6.6-fr]{4.6.6}). Nous allons voir qu'on peut remplacer $g$
par un morphisme fini surjectif $g':X''\to X$ tel que
${g'}^*(\mathcal{L})$ soit ample et que
$g'_*(\mathcal{O}_{X''})\otimes_{\mathcal{O}_X}\mathcal{R}(X)$ soit un
$\mathcal{R}(X)$-Module \emph{localement libre de rang constant}~; on sera
alors ramené aux conditions de (\hyperref[II.6.6.10-fr]{6.6.10}) et la
proposition sera démontrée.

Pour cela, soit $\mathcal{B}=g_*(\mathcal{O}_{X'})$~; désignons par $X_i$
($1\leq i\leq n$) les sous-préschémas fermés réduits de $X$ qui ont pour
espaces sous-jacents les composantes irréductibles de $X$
(\hyperref[I.5.2.1-fr]{I, 5.2.1})~; ils sont intègres par hypothèse. Soient
$X'_i$ le sous-préschéma fermé

\oldpage[II]{135}

$g^{-1}(X_i)$ de $X'$, $g_i:X'_i\to X_i$ le morphisme $g$ restreint à
$X'_i$, qui est fini (\hyperref[II.6.1.5-fr]{6.1.5, (iii)}) et surjectif~;
soit $k_i$ le \emph{rang} de la $\mathcal{O}_{X_i}$-Algèbre
$\mathcal{B}_i=(g_i)_*(\mathcal{O}_{X'_i})$. Comme
$\mathcal{B}_i\otimes_{\mathcal{O}_{X_i}}\mathcal{R}(X_i)$ est un
préfaisceau constant (\hyperref[I.7.3.5-fr]{I, 7.3.5}), le rang $k_i$ est
aussi le rang de la $(\mathcal{O}_X|U)$-Algèbre $\mathcal{B}|U$ pour tout
ouvert $U$ de $X$ ne rencontrant que la \emph{seule} composante irréductible
$X_i$. Si $T$ est un ensemble ouvert dans $X$ tel que $\mathcal{B}|T$ soit
isomorphe à $\mathcal{O}_X^m|T$, il résulte de la remarque précédente que les
nombres $k_i$ sont égaux à $m$ pour tous les indices $i$ tels que
$T\cap X_i\neq\varnothing$. Soit alors $U$ l'ouvert de $X$ formé des points
au voisinage desquels $\mathcal{B}$ est un $\mathcal{O}_X$-Module localement
libre, et soient $U_j$ ($1\leq j\leq s$) ses composantes connexes, qui sont
ouvertes dans $X$ et en nombre fini (puisque $U$ est noethérien)~; désignons
par $V_j$ le sous-préschéma fermé de $X'$, \emph{adhérence} du sous-préschéma
induit sur l'ouvert $g^{-1}(U_j)$
(\hyperref[I.9.5.11-fr]{I, 9.5.11}). D'après ce qui précède, pour tous les
indices $i$ tels que $X_i\cap U_j\neq\varnothing$, les rangs $k_i$ sont tous
égaux à un même entier $m_j$~; on notera d'ailleurs qu'un même $X_i$ ne peut
rencontrer deux $U_j$ d'indices distincts. Soient $i_\lambda$ les indices
$i$ tels que $X_i\cap U=\varnothing$. Considérons le produit $k$ de tous les
$k_i$, posons $n_i=k/k_i$, et soit $X''$ le préschéma défini comme suit.
Pour chaque $j$ ($1\leq j\leq s$), on considère $k/m_j$ préschémas
isomorphes à $V_j$, et pour chaque $\lambda$, $k/k_{i_\lambda}$ préschémas
isomorphes à $X'_{i_\lambda}$~; $X''$ est la \emph{somme} de tous ces
préschémas. On définit un morphisme $g'':X''\to X'$, se réduisant à
l'injection canonique dans chacun des summandes de $X''$~; il est clair que
$g''$ est un morphisme fini dominant, donc surjectif (puisqu'un morphisme
fini est fermé (\hyperref[II.6.1.10-fr]{6.1.10}))~; posons
$g'=g\circ g''$, qui est un morphisme fini surjectif $X''\to X$~; on a
${g'}^*(\mathcal{L})={g''}^*(g^*(\mathcal{L}))$, donc
${g'}^*(\mathcal{L})$ est un $\mathcal{O}_{X''}$-Module ample
(\hyperref[II.5.1.12-fr]{5.1.12}). Il est clair alors que pour ce nouveau
préschéma $X''$, les rangs définis comme les $k_i$ pour $X'$ sont tous
\emph{égaux à $k$}~; compte tenu de
(\hyperref[I.7.3.3-fr]{I, 7.3.3}), on en conclut aussitôt que pour tout ouvert
affine $T$ de $X$,
\[
\bigl(g'_*(\mathcal{O}_{X''})
\otimes_{\mathcal{O}_X}\mathcal{R}(X)\bigr)|T
\]
est un $(\mathcal{R}(X)|T)$-Module isomorphe à
$(\mathcal{R}(X)|T)^k$.

\begin{corollary}[6.6.12]
\label{II.6.6.12-fr}
Si, dans l'énoncé de (\hyperref[II.6.6.11-fr]{6.6.11}), on a $W=X$, alors
pour qu'un $\mathcal{O}_X$-Module inversible $\mathcal{L}$ soit ample
relativement à $f$, il faut et il suffit que $g^*(\mathcal{L})$ soit ample
relativement à $f\circ g$.
\end{corollary}

\begin{remark}[6.6.13]
\label{II.6.6.13-fr}
Au chapitre III, nous verrons que si $Y$ est noethérien, $f$ de type fini, et
si la restriction de $f$ au sous-préschéma fermé réduit de $X$ ayant $X-W$
pour espace sous-jacent est \emph{propre}, alors la conclusion de
(\hyperref[II.6.6.12-fr]{6.6.12}) est encore valable. Mais nous donnerons au
chapitre V des exemples de schémas algébriques $X$ sur un corps $K$ (le
morphisme structural $X\to\operatorname{Spec}(K)$ n'étant pas propre) dont
le normalisé $X'$ est quasi-affine, mais qui n'est pas quasi-affine (de sorte
que $\mathcal{O}_X$ n'est pas ample, bien que $\mathcal{O}_{X'}$ le soit
(\hyperref[II.5.1.12-fr]{5.1.12}) et que le morphisme $X'\to X$ soit fini
surjectif (\hyperref[II.6.3.10-fr]{6.3.10})). Nous allons voir au numéro
suivant que cette circonstance ne peut se produire lorsqu'on remplace
«~quasi-affine~» par «~affine~».
\end{remark}

\subsection{Le théorème de Chevalley.}
\label{subsection:II.6.7-fr}

Nous allons établir (à l'aide du critère de Serre
(\hyperref[II.5.2.1-fr]{5.2.1})) le théorème suivant, qui a été démontré par
C.~Chevalley par d'autres méthodes, dans le cas des schémas
\emph{algébriques}.

\oldpage[II]{136}

\begin{theorem}[6.7.1]
\label{II.6.7.1-fr}
Soient $X$ un schéma affine, $Y$ un préschéma noethérien, $f:X\to Y$ un
morphisme fini surjectif. Alors $Y$ est un schéma affine.
\end{theorem}

Il est clair que
$f_{\mathrm{red}}:X_{\mathrm{red}}\to Y_{\mathrm{red}}$ est fini
(\hyperref[II.6.1.5-fr]{6.1.5, (vi)})~; comme $X_{\mathrm{red}}$ est un
schéma affine, et qu'il revient au même de dire que $Y$ est affine ou que
$Y_{\mathrm{red}}$ est affine, puisque $Y$ est noethérien
(\hyperref[I.6.1.7-fr]{I, 6.1.7}), on voit qu'on peut supposer $Y$
\emph{réduit}. Pour toute partie fermée $Y'$ de $Y$, il y a alors un seul
sous-préschéma réduit de $Y$ ayant $Y'$ comme espace sous-jacent
(\hyperref[I.5.1.2-fr]{I, 5.1.2})~; son image réciproque $f^{-1}(Y')$,
canoniquement isomorphe à $X\times_Y Y'$
(\hyperref[I.4.4.1-fr]{I, 4.4.1}), est affine comme sous-préschéma fermé de
$X$, et la restriction de $f$ à $f^{-1}(Y')$, qui s'identifie à
$f\times_Y 1_{Y'}$, est un morphisme fini surjectif
(\hyperref[II.6.1.5-fr]{6.1.5, (iii)}). En vertu du principe de récurrence
noethérienne (\hyperref[0.2.2.2-fr]{0, 2.2.2}), on peut donc (compte tenu de
\hyperref[I.6.1.7-fr]{I, 6.1.7}) se ramener à démontrer le théorème sous
l'hypothèse que pour \emph{toute partie fermée $Y'\neq Y$}, tout
sous-préschéma fermé de $Y$ ayant $Y'$ pour espace sous-jacent est affine.
On en conclut que, \emph{pour tout $\mathcal{O}_Y$-Module cohérent
$\mathcal{F}$ dont le support (fermé) $Z$ est distinct de $Y$, on a
$\mathrm{H}^1(Y,\mathcal{F})=0$}. En effet, il existe un sous-préschéma fermé
$Y'$ de $Y$ ayant $Z$ pour espace sous-jacent et tel que, si $j:Y'\to Y$ est
l'injection canonique, on ait
$\mathcal{F}=j_*(j^*(\mathcal{F}))$
(\hyperref[I.9.3.5-fr]{I, 9.3.5})~; par suite
(\hyperref[II.5.2.3-fr]{5.2.3}),
$\mathrm{H}^1(Y,\mathcal{F})
=\mathrm{H}^1(Y',j^*(\mathcal{F}))=0$ d'après
(\hyperref[I.5.1.9.2-fr]{I, 5.1.9.2}).

Supposons d'abord que $Y$ ne soit pas irréductible, et soit $Y'$ une
composante irréductible de $Y$, et $Y''=Y-Y'$~; on désignera encore par
$Y'$ le sous-préschéma fermé réduit de $Y$ ayant $Y'$ pour espace
sous-jacent, et par $j$ l'injection canonique $Y'\to Y$. Soit $\mathcal{F}$
un $\mathcal{O}_Y$-Module cohérent, et considérons l'homomorphisme canonique
\[
\rho:\mathcal{F}\longrightarrow
\mathcal{F}'=j_*(j^*(\mathcal{F}))
\]
(\hyperref[0.4.4.3-fr]{0, 4.4.3})~; $\mathcal{F}'$ est un
$\mathcal{O}_Y$-Module cohérent en vertu de
(\hyperref[0.5.3.10-fr]{0, 5.3.10}) et
(\hyperref[0.5.3.12-fr]{0, 5.3.12}), car on a
$j_*(\mathcal{O}_{Y'})=\mathcal{O}_Y/\mathcal{J}$, en désignant par
$\mathcal{J}$ le faisceau d'idéaux de $\mathcal{O}_Y$ définissant le
sous-préschéma $Y'$. Par suite $\mathcal{G}=\operatorname{Ker}\rho$ et
$\mathcal{K}=\operatorname{Im}\rho$ sont aussi des $\mathcal{O}_Y$-Modules
cohérents (\hyperref[0.5.3.4-fr]{0, 5.3.4})~; or par définition la fibre
$\mathcal{F}'_y$ de $\mathcal{F}'$ au point générique $y$ de $Y'$ est égale
à $\mathcal{F}_y$, car $y$ est intérieur à $Y'$ et on a donc
$\mathcal{J}_y=0$, puisque $Y$ est réduit. On en conclut que $y$ n'est pas
contenu dans le support (fermé) de $\mathcal{G}$~; par ailleurs, le support de
$\mathcal{F}'$ (et \emph{a fortiori} celui de $\mathcal{K}$) est contenu
dans $Y'$~; autrement dit les supports de $\mathcal{G}$ et de $\mathcal{K}$
sont \emph{distincts de $Y$}. On en déduit que
$\mathrm{H}^1(Y,\mathcal{G})=\mathrm{H}^1(Y,\mathcal{K})=0$, et la suite
exacte de cohomologie appliquée à la suite exacte
$0\to\mathcal{G}\to\mathcal{F}\to\mathcal{K}\to0$ montre que
$\mathrm{H}^1(Y,\mathcal{F})=0$. On conclut alors par le critère de Serre
(\hyperref[II.5.2.1-fr]{5.2.1}).

Supposons donc $Y$ irréductible, et par suite \emph{intègre}. On peut aussi
supposer que $X$ est \emph{intègre}~: en effet, si on désigne par $X_i$ les
sous-préschémas fermés réduits de $X$ ayant pour espaces sous-jacents les
composantes irréductibles de $X$ (\hyperref[I.5.2.1-fr]{I, 5.2.1}) et par
$f_i$ la restriction de $f$ à $X_i$, un des $f_i$ au moins est dominant, et
comme c'est un morphisme fini (\hyperref[II.6.1.5-fr]{6.1.5}), il est
surjectif (\hyperref[II.6.1.10-fr]{6.1.10})~; comme d'autre part $X_i$ est un
schéma affine, on voit qu'on peut remplacer $X$ par $X_i$ dans l'énoncé.

\begin{lemma}[6.7.1.1]
\label{II.6.7.1.1-fr}
Soient $X$, $Y$ deux préschémas intègres noethériens, $x$ (resp. $y$) le
point générique de $X$ (resp. $Y$), $f:X\to Y$ un morphisme fini surjectif.
Soit $\mathcal{L}$ un $\mathcal{O}_X$-Module inversible tel qu'il existe un
voisinage ouvert affine $U$ de $y$ et une section
$g\in\Gamma(X,\mathcal{L})$ pour lesquels

\oldpage[II]{137}

$x\in X_g\subset f^{-1}(U)$. Alors il existe deux entiers $m>0$, $n>0$, un
homomorphisme
$u:\mathcal{O}_Y^m\to f_*(\mathcal{L}^{\otimes n})$ et un voisinage ouvert
$V$ de $y$ tels que la restriction $u|V$ soit un isomorphisme de
$\mathcal{O}_Y^m|V$ sur $f_*(\mathcal{L}^{\otimes n})|V$.
\end{lemma}

Soient $C$ l'anneau (intègre) de $U$, $k=\mathcal{O}_y$ son corps des
fractions~: comme $f$ est fini, $U'=f^{-1}(U)$ est affine
(\hyperref[II.1.3.2-fr]{1.3.2})~; soit $D$ son anneau (intègre) de corps des
fractions $K=\mathcal{O}_x$~; par hypothèse $D$ est un $C$-module de type
fini (\hyperref[II.6.1.4-fr]{6.1.4}), donc $K$ est une extension de rang fini
de $k$. La fibre
\[
f^{-1}(y)=X\times_Y\operatorname{Spec}(k(y))
=X\times_Y\operatorname{Spec}(\mathcal{O}_y)
\]
s'identifie à $\operatorname{Spec}(K)$
(\hyperref[I.3.6.6-fr]{I, 3.6.6})~; soient $s_i$ ($1\leq i\leq m$) des
éléments de $D$ formant une base de $K$ sur $k$. Il existe $n>0$ tel que les
sections $(s_i|X_g)g^{\otimes n}$ de $\mathcal{L}^{\otimes n}$ au-dessus de
$X_g$ se prolongent en sections $b_i$ ($1\leq i\leq m$) de
$\mathcal{L}^{\otimes n}$ au-dessus de $X$
(\hyperref[I.9.3.1-fr]{I, 9.3.1}). Les $b_i$ sont aussi par définition des
sections de $f_*(\mathcal{L}^{\otimes n})$ au-dessus de $Y$ et définissent
donc un homomorphisme
$u:\mathcal{O}_Y^m\to f_*(\mathcal{L}^{\otimes n})$
(\hyperref[0.5.1.1-fr]{0, 5.1.1})~; nous allons voir que $u$ répond à la
question. On a $\mathcal{L}^{\otimes n}|U'=\widetilde{M}$, où $M$ est un
$D$-module de type fini, donc si $\varphi$ est l'injection $C\to D$
correspondant au morphisme $f^{-1}(U)\to U$ restriction de $f$,
$M_{[\varphi]}$ est un $C$-module de type fini~; par suite
\[
f_*(\mathcal{L}^{\otimes n})|U=(M_{[\varphi]})^\sim
\]
(\hyperref[I.1.6.3-fr]{I, 1.6.3}) est cohérent et comme $U$ est un ouvert
affine quelconque de $Y$, $f_*(\mathcal{L}^{\otimes n})$ est cohérent~; en
outre $u|U=\widetilde{\theta}$, où $\theta$ est un $C$-homomorphisme
$C^m\to M_{[\varphi]}$, et $u_y=\theta_y$ est l'homomorphisme
\[
\theta\otimes1:K^m=C^m\otimes K\longrightarrow
M_{[\varphi]}\otimes K,
\]
or ce dernier est par définition un \emph{isomorphisme}, car les $(b_i)_x$
forment une base de $M_{[\varphi]}\otimes K$ sur $k$, $g_x$ étant par
hypothèse un générateur de $(\mathcal{L}^{\otimes n})_x$. On en conclut que
les supports des $\mathcal{O}_Y$-Modules $\operatorname{Ker}u$ et
$\operatorname{Coker}u$ ne contiennent pas $y$~; comme ces
$\mathcal{O}_Y$-Modules sont cohérents
(\hyperref[0.5.3.3-fr]{0, 5.3.3}), leurs supports sont fermés
(\hyperref[0.5.2.2-fr]{0, 5.2.2}), d'où le lemme.

Cela étant, les hypothèses du lemme
(\hyperref[II.6.7.1.1-fr]{6.7.1.1}) sont remplies dans le cas que nous
considérons en prenant $\mathcal{L}=\mathcal{O}_X$, puisque $X$ est affine
(\hyperref[I.1.1.10-fr]{I, 1.1.10})~; nous poserons
$\mathcal{A}=\mathcal{O}_Y$, $\mathcal{B}=f_*(\mathcal{O}_X)$. En vertu du
critère de Serre (\hyperref[II.5.2.1.1-fr]{5.2.1.1}), il suffit de prouver
que pour tout $\mathcal{O}_Y$-Module cohérent $\mathcal{F}$, on a
$\mathrm{H}^1(Y,\mathcal{F})=0$~; il suffit même de le prouver lorsque
$\mathcal{F}\subset\mathcal{O}_Y$, ce qui entraîne que $\mathcal{F}$ est
sans torsion puisque $Y$ est intègre~; en fait nous allons montrer que
$\mathrm{H}^1(Y,\mathcal{F})=0$ pour \emph{tout} $\mathcal{O}_Y$-Module
cohérent \emph{sans torsion} $\mathcal{F}$. Or, l'homomorphisme
$u:\mathcal{A}^m\to\mathcal{B}$ définit un homomorphisme
\[
v:\mathcal{G}
=\mathcal{H}om_{\mathcal{A}}(\mathcal{B},\mathcal{F})
\longrightarrow
\mathcal{H}om_{\mathcal{A}}(\mathcal{A}^m,\mathcal{F})
=\mathcal{F}^m.
\]
Montrons d'abord que $v$ est \emph{injectif}~: par hypothèse,
$\mathcal{T}=\operatorname{Coker}u$ a un support ne rencontrant pas $V$, donc
est un $\mathcal{O}_Y$-Module de torsion
(\hyperref[I.7.4.6-fr]{I, 7.4.6})~; la suite exacte
\[
\mathcal{A}^m\longrightarrow\mathcal{B}\longrightarrow\mathcal{T}
\longrightarrow0
\]
donne, par exactitude à gauche du foncteur
$\mathcal{H}om_{\mathcal{A}}$, la suite exacte
\[
0\longrightarrow
\mathcal{H}om_{\mathcal{A}}(\mathcal{T},\mathcal{F})
\longrightarrow\mathcal{G}\xrightarrow{v}\mathcal{F}^m.
\]
Mais comme $\mathcal{F}$ est sans torsion, on a
$\mathcal{H}om_{\mathcal{A}}(\mathcal{T},\mathcal{F})=0$
(\hyperref[0.5.2.6-fr]{0, 5.2.6}), d'où notre assertion. On a donc la suite
exacte
\[
0\longrightarrow\mathcal{G}\longrightarrow\mathcal{F}^m
\longrightarrow\operatorname{Coker}v\longrightarrow0
\]

\oldpage[II]{138}

où $\mathcal{G}$ et $\operatorname{Coker}v$ sont des
$\mathcal{O}_Y$-Modules cohérents
(\hyperref[0.5.3.4-fr]{0, 5.3.4} et
\hyperref[0.5.3.5-fr]{5.3.5})~; en vertu de la suite exacte de cohomologie,
il suffira de montrer que
$\mathrm{H}^1(Y,\mathcal{G})
=\mathrm{H}^1(Y,\operatorname{Coker}v)=0$ pour en déduire
$\mathrm{H}^1(Y,\mathcal{F}^m)
=(\mathrm{H}^1(Y,\mathcal{F}))^m=0$, et par suite
$\mathrm{H}^1(Y,\mathcal{F})=0$. Or la restriction $v|V$ est un isomorphisme,
donc le support de $\operatorname{Coker}v$ est distinct de $Y$, d'où
$\mathrm{H}^1(Y,\operatorname{Coker}v)=0$ par l'hypothèse du début. D'autre
part, $\mathcal{G}$ est un $\mathcal{B}$-Module cohérent
(\hyperref[I.9.6.4-fr]{I, 9.6.4})~; comme $X$ est affine au-dessus de $Y$, il
existe un $\mathcal{O}_X$-Module quasi-cohérent $\mathcal{K}$ tel que
$\mathcal{G}$ soit isomorphe à $f_*(\mathcal{K})$
(\hyperref[II.1.4.3-fr]{1.4.3})~; comme $X$ est affine, on a
$\mathrm{H}^1(X,\mathcal{K})=0$
(\hyperref[I.5.1.9.2-fr]{I, 5.1.9.2}), donc aussi
$\mathrm{H}^1(Y,\mathcal{G})=0$ par
(\hyperref[II.5.2.3-fr]{5.2.3}), ce qui achève la démonstration du théorème
(\hyperref[II.6.7.1-fr]{6.7.1}).

\begin{corollary}[6.7.2]
\label{II.6.7.2-fr}
Soient $X$ un préschéma noethérien, $(X_i)_{1\leq i\leq n}$ un recouvrement
fini de l'espace $X$ formé d'ensembles fermés. Pour que $X$ soit affine, il
faut et il suffit que pour chaque $i$, il existe un sous-préschéma fermé de
$X$, ayant $X_i$ pour espace sous-jacent et qui soit affine.
\end{corollary}

En effet, s'il en est ainsi, soit $X'$ le schéma \emph{somme} des $X_i$~; il
est clair que $X'$ est affine, et on définit un morphisme surjectif
$f:X'\to X$ en prenant pour restriction de $f$ à $X_i$ l'injection
canonique. Tout revient à vérifier que $f$ est \emph{fini}, en vertu de
(\hyperref[II.6.7.1-fr]{6.7.1}), et cela a été vu en
(\hyperref[II.6.1.5-fr]{6.1.5}).

\begin{corollary}[6.7.3]
\label{II.6.7.3-fr}
Pour qu'un préschéma noethérien $X$ soit affine, il faut et il suffit que les
sous-préschémas fermés réduits ayant pour espaces sous-jacents les composantes
irréductibles de $X$ soient affines.
\end{corollary}

\section{Critères valuatifs}
\label{section:II.7-fr}

Nous donnons dans ce paragraphe des critères valuatifs de séparation et de
propreté pour un morphisme, c'est-à-dire des critères qui font intervenir un
schéma auxiliaire variable de la forme $\operatorname{Spec}(A)$, où $A$ est
un anneau de valuation. Moyennant des hypothèses «~noethériennes~»
convenables, on peut dans ces critères se restreindre au cas où $A$ est un
anneau de valuation \emph{discrète}. Ce sera sans doute le seul cas que nous
aurons à appliquer par la suite, et nous n'avons introduit les anneaux de
valuation quelconques, dans le cas général, que pour faire le lien avec des
développements classiques.

\subsection{Rappels sur les anneaux de valuation.}
\label{subsection:II.7.1-fr}

\begin{env}[7.1.1]
\label{II.7.1.1-fr}
Parmi les diverses propriétés équivalentes qui caractérisent les anneaux de
valuation, nous utiliserons la suivante~: un anneau $A$ est dit
\emph{anneau de valuation} si c'est un anneau intègre qui n'est pas un
corps, et si, dans l'ensemble des anneaux locaux contenus dans le corps des
fractions $K$ de $A$ et distincts de $K$, $A$ est \emph{maximal} pour la
relation de domination (\hyperref[I.8.1.1-fr]{I, 8.1.1}). Rappelons qu'un
anneau de valuation est \emph{intégralement clos}. Si $A$ est un anneau de
valuation, $A_{\mathfrak{p}}$ est aussi un anneau de valuation pour tout
idéal premier $\mathfrak{p}\neq0$ de $A$.
\end{env}

\begin{env}[7.1.2]
\label{II.7.1.2-fr}
Soient $K$ un corps, $A$ un sous-anneau local de $K$ qui n'est pas un corps~;

\oldpage[II]{139}

il existe alors un anneau de valuation ayant $K$ pour corps des fractions et
qui domine $A$ (\cite[p.~1-07, lemme~2]{I-1}).

D'autre part, soient $B$ un anneau de valuation, $k$ son corps des restes,
$K$ son corps des fractions, $L$ une extension de $k$. Alors il existe un
anneau de valuation \emph{complet} $C$ qui domine $B$ et dont le corps des
restes est $L$. En effet, $L$ est extension algébrique d'une extension
transcendante pure $L'=k(T_\mu)_{\mu\in M}$~; on sait qu'on peut prolonger la
valuation de $K$ correspondant à $B$ à une valuation de
$K'=K(T_\mu)_{\mu\in M}$ de sorte que $L'$ soit le corps des restes de cette
valuation (\cite[p.~98]{II-24})~; remplaçant $B$ par le complété de l'anneau
de cette valuation prolongée, on voit qu'on peut se limiter au cas où $B$ est
complet et $L$ est une clôture algébrique de $k$. Si $\overline{K}$ est une
clôture algébrique de $K$, on peut alors prolonger à $\overline{K}$ la
valuation qui définit $B$, et le corps des restes correspondant est une
clôture algébrique de $k$, comme on le voit en relevant dans $\overline{K}$
les coefficients d'un polynôme unitaire de $k[T]$. On est donc finalement
ramené au cas où $L=k$ et il suffit alors de prendre pour $C$ le complété de
$B$ pour répondre à la question.
\end{env}

\begin{env}[7.1.3]
\label{II.7.1.3-fr}
Soient $K$ un corps, $A$ un sous-anneau de $K$~; la fermeture intégrale $A'$
de $A$ dans $K$ est l'intersection des anneaux de valuation contenant $A$ et
ayant $K$ pour corps des fractions
(\cite[p.~51, th.~2]{I-11}). La proposition
(\hyperref[II.7.1.2-fr]{7.1.2}) s'exprime sous la forme géométrique
équivalente~:
\end{env}

\begin{proposition}[7.1.4]
\label{II.7.1.4-fr}
Soient $Y$ un préschéma, $p:X\to Y$ un morphisme, $x$ un point de $X$,
$y=p(x)$, $y'\neq y$ une spécialisation
(\hyperref[0.2.1.2-fr]{0, 2.1.2}) de $y$. Il existe alors un schéma local
$Y'$, spectre d'un anneau de valuation, et un morphisme séparé $f:Y'\to Y$
tel que, si $a$ désigne l'unique point fermé de $Y'$ et $b$ le point
générique de $Y'$, on ait $f(a)=y'$ et $f(b)=y$. On peut en outre supposer
vérifiée l'une des deux propriétés additionnelles suivantes~:
\begin{enumerate}
\item[\rm{(i)}] $Y'$ est le spectre d'un anneau de valuation \emph{complet}
dont le corps des restes est algébriquement clos, et il existe un
$k(y)$-homomorphisme $k(x)\to k(b)$.
\item[\rm{(ii)}] Il existe un $k(y)$-isomorphisme
$k(x)\xrightarrow{\sim}k(b)$.
\end{enumerate}
\end{proposition}

Soit $Y_1$ le sous-préschéma fermé réduit de $Y$ ayant
$\overline{\{y\}}$ comme espace sous-jacent
(\hyperref[I.5.2.1-fr]{I, 5.2.1}), et soit $X_1$ le sous-préschéma fermé
image réciproque $p^{-1}(Y_1)$~; comme $y'\in\overline{\{y\}}$ par hypothèse
et que $k(x)$ est le même dans $X$ et dans $X_1$, on peut supposer $Y$
\emph{intègre}, de point générique $y$~; $\mathcal{O}_{y'}$ est alors un
anneau local intègre qui n'est pas un corps, et dont le corps des fractions
est $\mathcal{O}_y=k(y)$, et $k(x)$ est une extension de $k(y)$. Pour
réaliser les conditions $f(a)=y'$ et $f(b)=y$ ainsi que la condition
additionnelle (i) (resp. (ii)), on prendra
$Y'=\operatorname{Spec}(A')$, où $A'$ est un anneau de valuation dominant
$\mathcal{O}_{y'}$ et qui est complet et a un corps des restes
algébriquement clos extension de $k(x)$ (resp. un anneau de valuation $A'$
dominant $\mathcal{O}_{y'}$ et dont $k(x)$ est le corps des fractions)~;
l'existence de $A'$ est garantie par
(\hyperref[II.7.1.2-fr]{7.1.2}).

\begin{env}[7.1.5]
\label{II.7.1.5-fr}
Rappelons qu'un anneau local $A$ est dit \emph{de dimension $1$} s'il existe
un idéal premier distinct de l'idéal maximal $\mathfrak{m}$, et si tout idéal
premier de $A$ distinct de $\mathfrak{m}$ est un idéal premier
\emph{minimal}~; lorsque $A$ est \emph{intègre}, il revient au même de dire
que $\mathfrak{m}$ et $(0)$ sont les seuls idéaux premiers et
$\mathfrak{m}\neq(0)$~; autrement dit, $Y=\operatorname{Spec}(A)$ est
réduit à deux

\oldpage[II]{140}

points $a$, $b$~: $a$ est l'unique point \emph{fermé}, on a
$\mathfrak{j}_a=\mathfrak{m}$ et $k(a)=k$ est le \emph{corps résiduel}
$k=A/\mathfrak{m}$~; $b$ est le \emph{point générique} de $Y$,
$\mathfrak{j}_b=(0)$, l'ensemble $\{b\}$ étant l'unique ensemble ouvert de
$Y$ distinct de $\varnothing$ et de $Y$ (ensemble ouvert qui est donc
\emph{partout dense}), et $k(b)=K$ est le \emph{corps des fractions} de $A$.
\end{env}

\begin{env}[7.1.6]
\label{II.7.1.6-fr}
Pour un anneau local $A$, noethérien et de dimension $1$, on sait que les
conditions suivantes sont équivalentes
(\cite[p.~2-08 et 17-01]{I-1})~:
\begin{enumerate}
\item[\rm{a)}] $A$ est normal~;
\item[\rm{b)}] $A$ est régulier~;
\item[\rm{c)}] $A$ est un anneau de valuation~;
\end{enumerate}
en outre, $A$ est alors un \emph{anneau de valuation discrète}. Les prop.
(\hyperref[II.7.1.2-fr]{7.1.2}) et
(\hyperref[II.7.1.3-fr]{7.1.3}) ont alors les analogues suivants pour les
anneaux de valuation discrète~:
\end{env}

\begin{proposition}[7.1.7]
\label{II.7.1.7-fr}
Soient $A$ un anneau local intègre noethérien qui n'est pas un corps, $K$ son
corps des fractions, $L$ une extension de type fini de $K$~; il existe alors
un anneau de valuation discrète ayant $L$ pour corps des fractions et
dominant $A$.
\end{proposition}

Supposons d'abord $L=K$. Soient $\mathfrak{m}$ l'idéal maximal de $A$,
$(x_1,\ldots,x_n)$ un système de générateurs non nuls de $\mathfrak{m}$,
$B$ le sous-anneau
$A[x_2/x_1,\ldots,x_n/x_1]$ de $K$, qui est noethérien. Il est immédiat que
l'idéal $\mathfrak{m}B$ de $B$ est identique à l'idéal principal $x_1B$~;
si $\mathfrak{p}$ est un idéal premier minimal de $x_1B$, on sait que
$\mathfrak{p}$ est de rang $1$ (\cite[t.~I, p.~277]{I-13})~; autrement dit
$B_{\mathfrak{p}}$ est un anneau local noethérien \emph{de dimension $1$}~;
il est clair que
$\mathfrak{p}B_{\mathfrak{p}}\cap A$ est un idéal de $A$ contenant
$\mathfrak{m}$ et ne contenant pas $1$, donc égal à $\mathfrak{m}$, et par
suite $B_{\mathfrak{p}}$ \emph{domine} $A$
(\hyperref[I.8.1.1-fr]{I, 8.1.1}). Il résulte du th. de Krull-Akizuki
(\cite[p.~293]{II-25}) que la clôture intégrale $C$ de $B_{\mathfrak{p}}$
est un anneau noethérien (bien que $C$ ne soit pas nécessairement un
$B_{\mathfrak{p}}$-module de type fini)~; si $\mathfrak{n}$ est un idéal
maximal de $C$, $C_{\mathfrak{n}}$ est un anneau local noethérien normal de
dimension $1$ (\cite[p.~295]{II-25}), donc un anneau de valuation discrète,
qui domine $B_{\mathfrak{p}}$ et \emph{a fortiori} $A$.

Si maintenant $L$ est une extension de type fini de $K$, on peut, d'après ce
qui précède, se limiter au cas où $A$ est déjà un anneau de valuation
discrète. Soit $w$ une valuation de $K$ associée à $A$~; il existe une
valuation discrète $w'$ de $L$ qui \emph{prolonge} $w$~: on se ramène en
effet par récurrence sur le nombre de générateurs de $L$ au cas où
$L=K(\alpha)$, et alors la proposition est classique
(\cite[p.~106]{II-24}).

\begin{corollary}[7.1.8]
\label{II.7.1.8-fr}
Soient $A$ un anneau intègre noethérien, $K$ son corps des fractions, $L$ une
extension de type fini de $K$. Alors la fermeture intégrale de $A$ dans $L$
est l'intersection des anneaux de valuation discrète ayant $L$ pour corps des
fractions et contenant $A$.
\end{corollary}

En effet, un tel anneau de valuation discrète, étant normal, contient
\emph{a fortiori} tout élément de $L$ entier sur $A$. Il suffit donc de
prouver que si $x\in L$ n'est pas entier sur $A$, il existe un anneau de
valuation discrète $C$ ayant $L$ pour corps des fractions, contenant $A$ et
ne contenant pas $x$. L'hypothèse sur $x$ signifie en effet que l'on a
$x\notin B=A[1/x]$, autrement dit $1/x$ n'est pas inversible dans l'anneau
noethérien $B$. Il y a donc un idéal premier $\mathfrak{p}$ de $B$ contenant
$1/x$. L'anneau local intègre $B_{\mathfrak{p}}$ est noethérien et contenu
dans $L$, qui est une extension de type fini du corps des fractions de
$B_{\mathfrak{p}}$ (ce dernier contenant $K$). En vertu de
(\hyperref[II.7.1.7-fr]{7.1.7}), il y a donc un anneau de valuation discrète
$C$, dominant $B_{\mathfrak{p}}$ et ayant $L$ pour corps des fractions~;
comme $1/x\in\mathfrak{p}B_{\mathfrak{p}}$ appartient à l'idéal maximal de
$C$, on a $x\notin C$, ce qui conclut la démonstration.

La forme géométrique de (\hyperref[II.7.1.7-fr]{7.1.7}) est la suivante~:

\oldpage[II]{141}
\begin{proposition}[7.1.9]
\label{II.7.1.9-fr}
Soient $Y$ un préschéma localement noethérien, $p:X\to Y$ un morphisme
localement de type fini, $x$ un point de $X$, $y=p(x)$, $y'\neq y$ une
spécialisation de $y$. Il existe alors un schéma local $Y'$, spectre d'un
anneau de valuation discrète, un morphisme séparé $f:Y'\to Y$ et une
$Y$-application rationnelle $g$ de $Y'$ dans $X$, tels que, si $a$ désigne le
point fermé de $Y'$ et $b$ son point générique, on ait $f(a)=y'$,
$f(b)=y$, $g(b)=x$, et que dans le diagramme commutatif
\[
\tag{7.1.9.1}
\label{II.7.1.9.1-fr}
\begin{tikzcd}[column sep=large,row sep=large]
& k(x) \arrow[dl,"\gamma"'] \\
k(b) & k(y) \arrow[u,"\pi"] \arrow[l,"\varphi"']
\end{tikzcd}
\]
(où $\pi$, $\varphi$, $\gamma$ sont les homomorphismes correspondant
respectivement à $p$, $f$ et $g$), $\gamma$ soit une bijection.
\end{proposition}

Comme dans (\hyperref[II.7.1.4-fr]{7.1.4}), on peut se ramener au cas où $Y$
est intègre de point générique $y$ (compte tenu de
(\hyperref[I.6.3.4-fr]{I, 6.3.4, (iv)})) et comme la question est locale sur
$X$ et $Y$, on peut supposer $p$ de type fini~; on est alors dans la situation
de (\hyperref[II.7.1.4-fr]{7.1.4}), avec la propriété supplémentaire que
$k(x)$ est une extension \emph{de type fini} de $k(y)$
(\hyperref[I.6.4.11-fr]{I, 6.4.11}) et que $\mathcal{O}_{y'}$ est noethérien~;
cela permet d'appliquer (\hyperref[II.7.1.7-fr]{7.1.7}) et de prendre
$Y'=\operatorname{Spec}(A')$, où $A'$ est un anneau de valuation discrète
dominant $\mathcal{O}_{y'}$ et dont $k(x)$ est le corps des fractions. On a
donc ainsi défini un diagramme commutatif
(\hyperref[II.7.1.9.1-fr]{7.1.9.1}), où $\gamma$ est une bijection, et $\pi$
et $\varphi$ correspondent aux morphismes $p$ et $f$. En outre, comme $X$ et
$Y$ sont localement noethériens (\hyperref[I.6.6.2-fr]{I, 6.6.2}) et que $Y'$
est intègre, il y a une $Y$-application rationnelle $g$ et une seule de $Y'$
dans $X$ à laquelle correspond l'isomorphisme $\gamma$
(\hyperref[I.7.1.15-fr]{I, 7.1.15}), ce qui achève la démonstration.

\subsection{Critère valuatif de séparation}
\label{subsection:II.7.2-fr}

\begin{proposition}[7.2.1]
\label{II.7.2.1-fr}
Soient $X$, $Y$ deux préschémas, $f:X\to Y$ un morphisme quasi-compact. Les
deux conditions suivantes sont équivalentes~:
\begin{enumerate}
\item[\rm a)] Le morphisme $f$ est fermé.
\item[\rm b)] Pour tout $x\in X$ et toute spécialisation $y'$ de $y=f(x)$,
distincte de $y$, il existe une spécialisation $x'$ de $x$ telle que
$f(x')=x$.
\end{enumerate}
\end{proposition}

La condition b) exprime que $f(\overline{\{x\}})=\overline{\{y\}}$ et est donc
conséquence de a). Pour montrer que b) entraîne a), considérons une partie
fermée $X'$ de l'espace sous-jacent $X$~; soit $Y'=\overline{f(X')}$ et
montrons que $Y'=f(X')$. Considérons les sous-préschémas fermés réduits de
$X$ et $Y$ respectivement ayant pour espaces sous-jacents $X'$ et $Y'$
(\hyperref[I.5.2.1-fr]{I, 5.2.1})~; il y a alors un morphisme $f':X'\to Y'$ tel
que le diagramme
\[
\begin{tikzcd}[column sep=large,row sep=large]
X' \arrow[r,"f'"] \arrow[d] & Y' \arrow[d] \\
X \arrow[r,"f"'] & Y
\end{tikzcd}
\]
soit commutatif (\hyperref[I.5.2.2-fr]{I, 5.2.2}), et comme $f$ est
quasi-compact, il en est de même de $f'$. On est donc ramené à prouver que si
$f$ est un morphisme quasi-compact et \emph{dominant}, alors la
\oldpage[II]{142}
condition b) implique que $f(X)=Y$. Or, soit $y'$ un point de $Y$ et soit $y$
le point générique d'une composante irréductible de $Y$ contenant $y'$~; en
vertu de b), il suffit de montrer que $f^{-1}(y)$ n'est pas vide. Mais on sait
que cette propriété est conséquence du fait que $f$ est dominant et
quasi-compact (\hyperref[I.6.6.5-fr]{I, 6.6.5}).

\begin{corollary}[7.2.2]
\label{II.7.2.2-fr}
Soit $f:X\to Y$ un morphisme d'immersion quasi-compact. Pour que l'espace
sous-jacent $X$ soit fermé dans $Y$, il faut et il suffit qu'il contienne toute
spécialisation (dans $Y$) de chacun de ses points.
\end{corollary}

\begin{proposition}[7.2.3]
\label{II.7.2.3-fr}
Soient $Y$ un préschéma (resp. un préschéma localement noethérien), $f:X\to Y$
un morphisme (resp. un morphisme localement de type fini). Les conditions
suivantes sont équivalentes~:
\begin{enumerate}
\item[\rm a)] $f$ est séparé.
\item[\rm b)] Le morphisme diagonal $X\to X\times_YX$ est quasi-compact, et
pour tout $Y$-préschéma de la forme $Y'=\operatorname{Spec}(A)$, où $A$ est un
anneau de valuation (resp. un anneau de valuation discrète), deux
$Y$-morphismes de $Y'$ dans $X$ qui coïncident au point générique de $Y'$ sont
égaux.
\item[\rm c)] Le morphisme diagonal $X\to X\times_YX$ est quasi-compact, et
pour tout $Y$-préschéma de la forme $Y'=\operatorname{Spec}(A)$, où $A$ est un
anneau de valuation (resp. un anneau de valuation discrète), deux
$Y'$-sections de $X'=X_{(Y')}$ qui coïncident au point générique de $Y'$ sont
égales.
\end{enumerate}
\end{proposition}

L'équivalence de b) et de c) résulte de la correspondance biunivoque entre
$Y$-morphismes de $Y'$ dans $X$ et $Y'$-sections de $X'$
(\hyperref[I.3.3.14-fr]{I, 3.3.14}). Si $X$ est séparé sur $Y$, la condition b)
est vérifiée en vertu de (\hyperref[I.7.2.2.1-fr]{I, 7.2.2.1}), puisque $Y'$
est intègre. Il reste à prouver que b) entraîne que le morphisme diagonal
$\Delta:X\to X\times_YX$ est fermé, et il revient au même de montrer qu'il
vérifie le critère de (\hyperref[II.7.2.2-fr]{7.2.2}). Or, soit $z$ un point de
la diagonale $\Delta(X)$, $z'\neq z$ une spécialisation de $z$ dans
$X\times_YX$. Il existe alors (\hyperref[II.7.1.4-fr]{7.1.4}) un anneau de
valuation $A$ et un morphisme $f$ de $Y'=\operatorname{Spec}(A)$ dans
$X\times_YX$, qui applique le point fermé $a$ de $Y'$ sur $z'$ et le point
générique $b$ de $Y'$ sur $z$~; ce morphisme fait de $Y'$ un
$(X\times_YX)$-préschéma, et \emph{a fortiori} un $Y$-préschéma. Si on compose
avec $f$ les deux projections de $X\times_YX$, on obtient deux $Y$-morphismes
$g_1$, $g_2$ de $Y'$ dans $X$, qui par hypothèse coïncident au point $b$~; ils
sont donc égaux à un même morphisme $g$, ce qui signifie
(\hyperref[I.5.3.1-fr]{I, 5.3.1}) que $f$ se factorise en $f=\Delta\circ g$, et
par suite $z'\in\Delta(X)$. Lorsqu'on suppose $Y$ localement noethérien et $f$
localement de type fini, $X\times_YX$ est localement noethérien
(\hyperref[I.6.6.7-fr]{I, 6.6.7})~; on peut alors reprendre le raisonnement
précédent en y supposant que $A$ est un anneau de valuation discrète, en vertu
de (\hyperref[II.7.1.9-fr]{7.1.9}).

\begin{remark}[7.2.4]
\label{II.7.2.4-fr}
\begin{enumerate}
\item[\rm (i)] L'hypothèse que le morphisme $\Delta$ est quasi-compact est
toujours vérifiée lorsque $Y$ est localement noethérien et $f$ localement de
type fini, car $X\times_YX$ est alors localement noethérien
(\hyperref[I.6.6.4-fr]{I, 6.6.4, (i)}). Dans le cas général, elle signifie aussi
que pour tout recouvrement $(U_\alpha)$ de $X$ par des ouverts affines, les
ensembles $U_\alpha\cap U_\beta$ sont \emph{quasi-compacts}.
\item[\rm (ii)] Pour que $f$ soit séparé, il \emph{suffit} que la condition b)
ou la condition c) soit vérifiée pour un anneau de valuation $A$ qui est
\emph{complet} et dont le corps des restes est \emph{algébriquement clos}~;
cela résulte en effet de la démonstration de
(\hyperref[II.7.2.3-fr]{7.2.3}) et de
(\hyperref[II.7.1.4-fr]{7.1.4}).
\end{enumerate}
\end{remark}

\oldpage[II]{143}
\subsection{Critère valuatif de propreté}
\label{subsection:II.7.3-fr}

\begin{proposition}[7.3.1]
\label{II.7.3.1-fr}
Soient $A$ un anneau de valuation, $Y=\operatorname{Spec}(A)$, $b$ le point
générique de $Y$, $X$ un schéma intègre, $f:X\to Y$ un morphisme fermé tel
que $f^{-1}(b)$ se réduise à un point $x$ et que l'homomorphisme correspondant
$k(b)\to k(x)$ soit bijectif. Alors $f$ est un isomorphisme.
\end{proposition}

Comme $f$ est fermé et dominant, on a $f(X)=Y$~; il suffit
(\hyperref[I.4.2.2-fr]{I, 4.2.2}) de prouver que pour tout $y'\neq b$ dans $Y$,
il existe un \emph{seul} point $x'$ tel que $f(x')=y'$ et que l'homomorphisme
correspondant $\mathcal{O}_{y'}\to\mathcal{O}_{x'}$ soit bijectif, car $f$ sera
alors un homéomorphisme. Or, si $f(x')=y'$, $\mathcal{O}_{x'}$ est un anneau
local contenu dans $K=k(x)=k(b)$ et dominant $\mathcal{O}_{y'}$~; ce dernier
est l'anneau local $A_{y'}$, donc est un anneau de valuation
(\hyperref[II.7.1.1-fr]{7.1.1}) ayant $K$ pour corps des fractions. Par
ailleurs on a $\mathcal{O}_{x'}\neq K$ puisque $x'$ n'est pas le point
générique de $X$ (\hyperref[0.2.1.3-fr]{0, 2.1.3})~; on en conclut que
$\mathcal{O}_{x'}=\mathcal{O}_{y'}$. Comme $X$ est un schéma intègre, la
relation $\mathcal{O}_{x'}=\mathcal{O}_{x''}$ entraîne $x'=x''$
(\hyperref[I.8.2.2-fr]{I, 8.2.2}), ce qui achève la démonstration.

\begin{env}[7.3.2]
\label{II.7.3.2-fr}
Soient $A$ un anneau de valuation, $Y=\operatorname{Spec}(A)$, $b$ le point
générique de $Y$, de sorte que $\mathcal{O}_b=k(b)$ est égal à $K$, corps des
fractions de $A$~; soit $f:X\to Y$ un morphisme. On sait
(\hyperref[I.7.1.4-fr]{I, 7.1.4}) que les $Y$-sections \emph{rationnelles} de
$X$ sont en correspondance biunivoque avec les \emph{germes} de $Y$-sections
(définies dans des voisinages de $b$) au point $b$, d'où une application
canonique
\[
\tag{7.3.2.1}
\label{II.7.3.2.1-fr}
\Gamma_{\mathrm{rat}}(X/Y)\longrightarrow
\Gamma(f^{-1}(b)/\operatorname{Spec}(K))
\]
les éléments de $\Gamma(f^{-1}(b)/\operatorname{Spec}(K))$ s'identifiant
d'ailleurs, par définition (\hyperref[I.3.4.5-fr]{I, 3.4.5}) aux \emph{points
de $f^{-1}(b)=X\otimes_AK$ rationnels sur $K$}. Lorsque $f$ est \emph{séparé},
il résulte de (\hyperref[I.5.4.7-fr]{I, 5.4.7}) que l'application
(\hyperref[II.7.3.2.1-fr]{7.3.2.1}) est \emph{injective}, puisque $Y$ est un
schéma intègre.

Composant (\hyperref[II.7.3.2.1-fr]{7.3.2.1}) avec l'application canonique
$\Gamma(X/Y)\to\Gamma_{\mathrm{rat}}(X/Y)$
(\hyperref[I.7.1.2-fr]{I, 7.1.2}), on obtient une application canonique
\[
\tag{7.3.2.2}
\label{II.7.3.2.2-fr}
\Gamma(X/Y)\longrightarrow \Gamma(f^{-1}(b)/\operatorname{Spec}(K)).
\]
Lorsque $f$ est \emph{séparé}, cette application est encore \emph{injective}
(\hyperref[I.5.4.7-fr]{I, 5.4.7}).
\end{env}

\begin{proposition}[7.3.3]
\label{II.7.3.3-fr}
Soient $A$ un anneau de valuation de corps des fractions $K$,
$Y=\operatorname{Spec}(A)$, $b$ le point générique de $Y$, $f:X\to Y$ un
morphisme \emph{séparé} et \emph{fermé}. Alors l'application canonique
(\hyperref[II.7.3.2.2-fr]{7.3.2.2}) est bijective (ce qui revient à dire
qu'elle est \emph{surjective} et entraîne que les $Y$-sections rationnelles de
$X$ sont \emph{partout définies}).
\end{proposition}

Soit donc $x$ un point de $f^{-1}(b)$, \emph{rationnel sur $K$}. Comme $f$ est
séparé, il en est de même du morphisme
$f^{-1}(b)\to\operatorname{Spec}(K)$ correspondant à $f$
(\hyperref[I.5.5.1-fr]{I, 5.5.1, (iv)}), et toute section de $f^{-1}(b)$ étant
une immersion fermée (\hyperref[I.5.4.6-fr]{I, 5.4.6}), $\{x\}$ est fermé
\emph{dans $f^{-1}(b)$}. Considérons le sous-préschéma fermé réduit $X'$ de
$X$ ayant pour espace sous-jacent l'adhérence $\overline{\{x\}}$ de $\{x\}$
dans $X$. Il est clair que la restriction de $f$ à $X'$ vérifie les hypothèses
de (\hyperref[II.7.3.1-fr]{7.3.1}), donc est un \emph{isomorphisme} de $X'$ sur
$Y$, dont l'isomorphisme réciproque est la $Y$-section de $X$ cherchée.

\begin{env}[7.3.4]
\label{II.7.3.4-fr}
Pour énoncer les deux résultats suivants, nous nous servirons d'une
terminologie qui sera justifiée et discutée dans le chapitre IV~: si $F$ est
une partie
\oldpage[II]{144}
d'un préschéma $Y$, on appelle \emph{codimension} de $F$ dans $Y$ et on note
$\operatorname{codim}_YF$ la borne inférieure des entiers
$\dim(\mathcal{O}_z)$, où $z$ parcourt $F$.
\end{env}

\begin{corollary}[7.3.5]
\label{II.7.3.5-fr}
Soient $Y$ un préschéma localement noethérien réduit, $N$ l'ensemble des
points $y\in Y$ où $Y$ n'est pas régulier
(\hyperref[0.4.1.4-fr]{0, 4.1.4})~; on suppose que
$\operatorname{codim}_YN\geq2$. Soit $f:X\to Y$ un morphisme de type fini,
séparé et fermé et soit $g$ une $Y$-section rationnelle de $X$~; si $Y'$ est
l'ensemble des points de $Y$ où $g$ n'est pas définie (ensemble qui est
\emph{fermé} (\hyperref[I.7.2.1-fr]{I, 7.2.1})) on a
$\operatorname{codim}_YY'\geq2$.
\end{corollary}

Il suffit de prouver que $g$ est définie en tout point $z\in Y$ tel que
$\dim\mathcal{O}_z\leq1$. Si $\dim\mathcal{O}_z=0$, $z$ est point générique
d'une composante irréductible de $Y$ (\hyperref[I.1.1.14-fr]{I, 1.1.14}), donc
appartient à tout ouvert partout dense de $Y$, et en particulier au domaine de
définition de $g$. Supposons donc que $\dim\mathcal{O}_z=1$~;
$\mathcal{O}_z$ est alors par hypothèse un anneau local noethérien régulier,
donc (\hyperref[II.7.1.6-fr]{7.1.6}) un \emph{anneau de valuation discrète}.
Soit $Z=\operatorname{Spec}(\mathcal{O}_z)$~; comme $U=Y-Y'$ est partout
dense, $U\cap Z$ n'est pas vide (\hyperref[I.2.4.2-fr]{I, 2.4.2})~; soit $g'$
l'application rationnelle de $Z$ dans $X$ induite par $g$
(\hyperref[I.7.2.8-fr]{I, 7.2.8})~; il suffit de montrer que $g'$ est un
\emph{morphisme} (\hyperref[I.7.2.9-fr]{I, 7.2.9}). Or, $g'$ peut être
considérée comme une $Z$-section rationnelle du $Z$-préschéma
$f^{-1}(Z)=X\times_YZ$~; il est clair que le morphisme $f^{-1}(Z)\to Z$
correspondant à $f$ est fermé, et il résulte de
(\hyperref[I.5.5.1-fr]{I, 5.5.1, (i)}) qu'il est séparé~; on conclut alors de
(\hyperref[II.7.3.3-fr]{7.3.3}) que $g'$ est partout définie~; comme $Z$ est
réduit et $X$ séparé sur $Y$, $g'$ est un morphisme
(\hyperref[I.7.2.2-fr]{I, 7.2.2}).

\begin{corollary}[7.3.6]
\label{II.7.3.6-fr}
Soient $S$ un préschéma localement noethérien, $X$ et $Y$ deux
$S$-préschémas~; on suppose $Y$ réduit, et en outre que l'ensemble $N$ des
points $y\in Y$ où $Y$ n'est pas régulier soit tel que
$\operatorname{codim}_YN\geq2$~; on suppose enfin que le morphisme structural
$X\to S$ soit propre. Soit $f$ une $S$-application rationnelle de $Y$ dans
$X$ et soit $Y'$ l'ensemble des points de $Y$ où $f$ n'est pas définie~;
alors on a $\operatorname{codim}_YY'\geq2$.
\end{corollary}

On sait (\hyperref[I.7.1.2-fr]{I, 7.1.2}) qu'on peut identifier les
$S$-applications rationnelles de $Y$ dans $X$ aux $Y$-sections rationnelles
de $X\times_SY$~; comme le morphisme structural $X\times_SY\to Y$ est fermé
(\hyperref[II.5.4.1-fr]{5.4.1}), on peut appliquer
(\hyperref[II.7.3.5-fr]{7.3.5}), d'où le corollaire.

\begin{remark}[7.3.7]
\label{II.7.3.7-fr}
Les hypothèses faites sur $Y$ dans (\hyperref[II.7.3.5-fr]{7.3.5}) et
(\hyperref[II.7.3.6-fr]{7.3.6}) seront en particulier réalisées lorsque $Y$
est \emph{normal} (\hyperref[0.4.1.4-fr]{0, 4.1.4}), en vertu de
(\hyperref[II.7.1.6-fr]{7.1.6}).
\end{remark}

Nous pouvons caractériser les morphismes universellement fermés (resp.
propres) par une réciproque de (\hyperref[II.7.3.3-fr]{7.3.3})~:

\begin{theorem}[7.3.8]
\label{II.7.3.8-fr}
Soient $Y$ un préschéma (resp. un préschéma localement noethérien), $f:X\to Y$
un morphisme quasi-compact séparé (resp. de type fini). Les conditions
suivantes sont équivalentes~:
\begin{enumerate}
\item[\rm a)] $f$ est universellement fermé (resp. propre).
\item[\rm b)] Pour tout $Y$-schéma de la forme
$Y'=\operatorname{Spec}(A)$, où $A$ est un anneau de valuation (resp. un anneau
de valuation discrète) de corps des fractions $K$, l'application canonique
\[
\operatorname{Hom}_Y(Y',X)\longrightarrow
\operatorname{Hom}_Y(\operatorname{Spec}(K),X)
\]
correspondant à l'injection canonique $A\to K$, est surjective (resp.
bijective).
\oldpage[II]{145}
\item[\rm c)] Pour tout $Y$-schéma de la forme $Y'=\operatorname{Spec}(A)$,
où $A$ est un anneau de valuation (resp. un anneau de valuation discrète)
l'application canonique (\hyperref[II.7.3.2.2-fr]{7.3.2.2}) relative au
$Y'$-préschéma $X_{(Y')}$ est surjective (resp. bijective).
\end{enumerate}
\end{theorem}

L'équivalence de b) et c) résulte aussitôt de
(\hyperref[I.3.3.14-fr]{I, 3.3.14})~; a) implique c), car a) entraîne, dans
l'une ou l'autre hypothèse, que $f_{(Y')}$ est séparé
(\hyperref[I.5.5.1-fr]{I, 5.5.1, (iv)}) et fermé, et il suffit d'appliquer
(\hyperref[II.7.3.3-fr]{7.3.3}). Reste à prouver que b) entraîne a).
Plaçons-nous d'abord dans le cas où $Y$ est quelconque, $f$ séparé et
quasi-compact. Si la condition b) est vérifiée par $f$, elle l'est aussi par
$f_{(Y'')}:X_{(Y'')}\to Y''$, où $Y''$ est un $Y$-préschéma quelconque, en
vertu de l'équivalence de b) et c), et du fait que
$X_{(Y'')}\times_{Y''}Y'=X\times_YY'$ pour tout morphisme $Y'\to Y''$
(\hyperref[I.3.3.9.1-fr]{I, 3.3.9.1})~; comme en outre $f_{(Y'')}$ est séparé
et quasi-compact quand $f$ l'est
(\hyperref[I.5.5.1-fr]{I, 5.5.1, (iv)} et
\hyperref[I.6.6.4-fr]{6.6.4, (iii)}), on est ramené à prouver que b) entraîne
que $f$ est \emph{fermé}. Pour cela, il suffit de vérifier la condition b) de
(\hyperref[II.7.2.1-fr]{7.2.1}). Soient donc $x\in X$, $y'$ une spécialisation
de $y=f(x)$, distincte de $y$~; en vertu de
(\hyperref[II.7.1.4-fr]{7.1.4}), il y a un schéma $Y'$, spectre d'un anneau de
valuation, et un morphisme séparé $g:Y'\to Y$ tels que, si $a$ désigne le point
fermé et $b$ le point générique de $Y'$, on ait $g(a)=y'$, $g(b)=y$, et qu'il
existe un $k(y)$-homomorphisme $k(x)\to k(b)$. Ce dernier correspond
canoniquement à un $Y$-morphisme $\operatorname{Spec}(k(b))\to X$
(\hyperref[I.2.4.6-fr]{I, 2.4.6}), et il résulte donc de b) qu'il existe un
$Y$-morphisme $h:Y'\to X$ auquel correspond le morphisme précédent. On a
alors $h(b)=x$~; si l'on pose $h(a)=x'$, $x'$ est spécialisation de $x$, et
l'on a $f(x')=f(h(a))=g(a)=y'$.

Si maintenant $Y$ est localement noethérien et $f$ de type fini, l'hypothèse
b) entraîne d'abord que $f$ est \emph{séparé}, en vertu de
(\hyperref[II.7.2.3-fr]{7.2.3}), le morphisme diagonal $X\to X\times_YX$
étant quasi-compact (\hyperref[II.7.2.4-fr]{7.2.4}). En outre, pour vérifier
que $f$ est propre, il suffit de montrer que $f_{(Y'')}:X_{(Y'')}\to Y''$ est
\emph{fermé} pour tout $Y$-préschéma $Y''$ \emph{de type fini}, compte tenu de
(\hyperref[II.5.6.3-fr]{5.6.3}). Comme alors $Y''$ est localement noethérien,
on peut reprendre le raisonnement fait dans le premier cas en prenant pour
$Y'$ un spectre d'anneau de valuation discrète, et en appliquant
(\hyperref[II.7.1.9-fr]{7.1.9}) au lieu de
(\hyperref[II.7.1.4-fr]{7.1.4}).

\begin{remark}[7.3.9]
\label{II.7.3.9-fr}
\begin{enumerate}
\item[\rm (i)] Lorsque $Y$ est un préschéma quelconque et $f$ un morphisme
séparé, pour que $f$ soit universellement fermé, il \emph{suffit} que la
condition b) ou la condition c) soit vérifiée pour les anneaux de valuation
$A$ \emph{complets} et dont le corps des restes est \emph{algébriquement
clos}~; cela résulte en effet de la démonstration et de
(\hyperref[II.7.1.4-fr]{7.1.4}).
\item[\rm (ii)] On déduit du critère c) de
(\hyperref[II.7.3.8-fr]{7.3.8}) une nouvelle démonstration du fait qu'un
morphisme projectif $X\to Y$ est fermé (\hyperref[II.5.5.3-fr]{5.5.3}), plus
proche des méthodes classiques. On peut en effet supposer $Y$ affine, et par
suite $X$ identifié à un sous-préschéma fermé d'un fibré projectif
$\mathbf{P}_Y^n$ (\hyperref[II.5.3.3-fr]{5.3.3})~; pour prouver que $X\to Y$
est fermé, il suffit de vérifier qu'il en est ainsi du morphisme structural
$\mathbf{P}_Y^n\to Y$, et le critère c) de
(\hyperref[II.7.3.8-fr]{7.3.8}), joint à
(\hyperref[II.4.1.3.1-fr]{4.1.3.1}), prouve qu'on est ramené à démontrer le
fait suivant~: \emph{si $Y$ est le spectre d'un anneau de valuation $A$, de
corps des fractions $K$, tout point de $\mathbf{P}_Y^n$ à valeurs dans $K$
provient (par restriction au point générique de $Y$) d'un point de
$\mathbf{P}_Y^n$ à valeurs dans $A$}. Or, tout $\mathcal{O}_Y$-Module
inversible est trivial (\hyperref[I.2.4.8-fr]{I, 2.4.8})~; donc, il résulte de
(\hyperref[II.4.2.6-fr]{4.2.6}) qu'un point de $\mathbf{P}_Y^n$ à valeurs dans
$K$ s'identifie à une classe d'éléments $(\zeta c_0,\zeta c_1,\ldots,\zeta
c_n)$ de $K$, où $\zeta\neq0$ et les $c_i$ sont des éléments non tous nuls de
$K$. Or, en multipliant les $c_i$ par un élément de $A$ de
\oldpage[II]{146}
valuation convenable, on peut supposer que les $c_i$ appartiennent tous à
$A$, et que l'un au moins est inversible. Mais alors
(\hyperref[II.4.2.6-fr]{4.2.6}), le système $(c_0,\ldots,c_n)$ définit aussi
un point de $\mathbf{P}_Y^n$ à valeurs dans $A$, ce qui démontre notre
assertion.
\item[\rm (iii)] Les critères (\hyperref[II.7.2.3-fr]{7.2.3}) et
(\hyperref[II.7.3.8-fr]{7.3.8}) sont surtout commodes quand on considère la
donnée d'un $Y$-préschéma $X$ comme équivalente à la donnée du foncteur
\[
X(Y')=\operatorname{Hom}_Y(Y',X)
\]
en un $Y$-préschéma $Y'$~; ces critères nous permettront par exemple de
prouver que sous certaines conditions les «~schémas de Picard~» sont propres.
\end{enumerate}
\end{remark}

\begin{corollary}[7.3.10]
\label{II.7.3.10-fr}
Soient $Y$ un schéma intègre (resp. un schéma intègre localement noethérien),
$X$ un schéma intègre, $f:X\to Y$ un morphisme dominant.
\begin{enumerate}
\item[\rm (i)] Si $f$ est quasi-compact et universellement fermé, tout anneau
de valuation dont le corps des fractions est le corps $R(X)$ des fonctions
rationnelles sur $X$, et qui domine un anneau local de $Y$, domine aussi un
anneau local de $X$.
\item[\rm (ii)] Inversement, supposons que $f$ soit de type fini, et que la
propriété énoncée dans (i) soit vérifiée par tout anneau de valuation (resp.
tout anneau de valuation discrète) ayant $R(X)$ pour corps des fractions.
Alors $f$ est propre.
\end{enumerate}
\end{corollary}

Notons d'abord que les hypothèses impliquent dans tous les cas que $f$ est
séparé (\hyperref[I.5.5.9-fr]{I, 5.5.9}).
\begin{enumerate}
\item[\rm (i)] Soient $K=R(Y)$, $L=R(X)$, $y$ un point de $Y$, $A$ un anneau
de valuation ayant $L$ pour corps des fractions et dominant
$\mathcal{O}_y$~; l'injection $\mathcal{O}_y\to A$ définit donc un morphisme
$h$ de $Y'=\operatorname{Spec}(A)$ dans $Y$
(\hyperref[I.2.4.4-fr]{I, 2.4.4}) tel que $h(a)=y$, en désignant par $a$ le
point fermé de $Y'$~; en outre, si $\eta$ est le point générique de $Y$, qui
est aussi celui de $\operatorname{Spec}(\mathcal{O}_y)$, on a $h(b)=\eta$, en
désignant par $b$ le point générique de $Y'$ (puisque $K\subset L$ par
hypothèse). Si $\xi$ est le point générique de $X$, on a $k(\xi)=k(b)=L$ par
hypothèse, d'où un $Y$-morphisme $g:\operatorname{Spec}(L)\to X$ tel que
$g(b)=\xi$~; en vertu de (\hyperref[II.7.3.8-fr]{7.3.8}), $g$ provient d'un
$Y$-morphisme $g':Y'\to X$. Si $x=g'(a)$, il est clair que $A$ domine
$\mathcal{O}_x$.
\item[\rm (ii)] La question étant locale sur $Y$, on peut toujours supposer
que $Y$ est affine (resp. affine et noethérien). Comme $f$ est de type fini,
on peut appliquer dans les deux cas le lemme de Chow
(\hyperref[II.5.6.1-fr]{5.6.1}). Il y a donc un morphisme projectif
$p:P\to Y$, un morphisme d'immersion $j:X'\to P$ et un morphisme projectif,
surjectif et birationnel $g:X'\to X$ ($X'$ étant intègre) tels que le diagramme
\[
\begin{tikzcd}[column sep=large,row sep=large]
P \arrow[d,"p"'] & X' \arrow[l,"j"'] \arrow[d,"g"] \\
Y & X \arrow[l,"f"]
\end{tikzcd}
\]
soit commutatif. Il suffit de prouver que $j$ est une immersion \emph{fermée},
car alors $f\circ g=p\circ j$ sera un morphisme projectif, donc propre, et
comme $g$ est surjectif, $f$ sera aussi propre
(\hyperref[II.5.4.3-fr]{5.4.3}). Soit $Z$ le sous-préschéma fermé réduit de
$P$ ayant $\overline{j(X')}$ comme espace sous-jacent
(\hyperref[I.5.2.1-fr]{I, 5.2.1})~; comme $X'$ est intègre, $j$ se factorise
en $i\circ h$, où $i:Z\to P$ est l'injection canonique, $h:X'\to Z$ une
immersion ouverte dominante (\hyperref[I.5.2.3-fr]{I, 5.2.3}), et $Z$ est
intègre~;

\oldpage[II]{147}
en outre $Z$ est projectif sur $Y$, et on voit qu'on peut se borner au cas où
$P$ est \emph{intègre} et $j$ \emph{dominant et birationnel}, et tout revient à
voir que $j$ est \emph{surjectif}. Or, soit $z\in P$~; $\mathcal{O}_z$ est un
anneau local intègre (resp. intègre et noethérien) dont le corps des fractions
est
\[
L=R(P)=R(X')=R(X).
\]
On peut se borner au cas où $z$ n'est pas le point générique de $P$. Il y a par
suite (\hyperref[II.7.1.2-fr]{7.1.2} et
\hyperref[II.7.1.7-fr]{7.1.7}) un anneau de valuation (resp. un anneau de
valuation discrète) $A$, ayant $L$ pour corps des fractions et dominant
$\mathcal{O}_z$. \emph{A fortiori}, $A$ domine $\mathcal{O}_y$, en posant
$y=p(z)$, et par hypothèse il y a donc un $x\in X$ tel que $A$ domine
$\mathcal{O}_x$. Comme $g$ est propre, la première partie de la démonstration
prouve que $A$ domine aussi un $\mathcal{O}_{x'}$, pour un $x'\in X'$~; il
s'ensuit que $\mathcal{O}_z$ et $\mathcal{O}_{j(x')}=\mathcal{O}_{x'}$ sont
apparentés (\hyperref[I.8.1.4-fr]{I, 8.1.4}), et, comme $P$ est un schéma, cela
entraîne $z=j(x')$ (\hyperref[I.8.2.2-fr]{I, 8.2.2}) et achève la démonstration.
\end{enumerate}

\begin{corollary}[7.3.11]
\label{II.7.3.11-fr}
Soient $X$, $Y$ deux schémas intègres, $f:X\to Y$ un morphisme dominant,
quasi-compact et universellement fermé. Supposons en outre $Y$ affine d'anneau
(intègre) $B$. Alors $\Gamma(X,\mathcal{O}_X)$ est canoniquement isomorphe à
un sous-anneau de la fermeture intégrale de $B$ dans $R(X)$.
\end{corollary}

En effet (\hyperref[I.8.2.1.1-fr]{I, 8.2.1.1}),
$\Gamma(X,\mathcal{O}_X)$ s'identifie à l'intersection des $\mathcal{O}_x$ pour
$x\in X$~; en vertu de (\hyperref[II.7.3.10-fr]{7.3.10}), de
(\hyperref[II.7.1.2-fr]{7.1.2}) et (\hyperref[II.7.1.3-fr]{7.1.3}),
$\Gamma(X,\mathcal{O}_X)$ est par suite contenu dans l'intersection des
anneaux de valuation contenant $B$ et ayant $R(X)$ pour corps des fractions~;
la conclusion résulte alors de (\hyperref[II.7.1.3-fr]{7.1.3}).

\begin{remark}[7.3.12]
\label{II.7.3.12-fr}
Sous les hypothèses de (\hyperref[II.7.3.11-fr]{7.3.11}), et lorsqu'on suppose
que $R(X)$ est une extension de type fini de $R(Y)$, on pourra dans de nombreux
cas conclure que $\Gamma(X,\mathcal{O}_X)$ est un \emph{module de type fini}
sur l'anneau $B=\Gamma(Y,\mathcal{O}_Y)$. Ce sera le cas par exemple lorsque
$B$ est une \emph{algèbre de type fini sur un corps}, car on sait alors que la
fermeture intégrale de $B$ dans une extension de type fini de son corps des
fractions est un $B$-module de type fini ([13], t.~I, p.~267, th.~9)~; la
conclusion résulte alors de (\hyperref[II.7.3.11-fr]{7.3.11}) et du fait que
$B$ est noethérien.

En particulier, \emph{un schéma $X$ propre et affine sur un corps $K$ est
fini}. En effet, en vertu de (\hyperref[II.1.6.4-fr]{1.6.4}),
(\hyperref[II.5.4.6-fr]{5.4.6}) et
(\hyperref[I.6.4.4-fr]{I, 6.4.4, c)}), on peut se borner au cas où $X$ est
\emph{réduit}. En outre, il suffira de prouver que chacun des sous-préschémas
fermés de $X$ ayant pour espace sous-jacent une composante irréductible de $X$
(en nombre fini) est fini sur $K$, si bien que (tenant compte de
(\hyperref[II.5.4.5-fr]{5.4.5})) on est finalement ramené au cas où $X$ est
\emph{intègre}. Mais alors le résultat découle des remarques faites ci-dessus.

Au chapitre III, nous retrouverons cette dernière proposition par d'autres
méthodes et comme conséquence de résultats plus généraux, montrant que si
$f:X\to Y$ est propre et $Y$ localement noethérien, $f_*(\mathcal{F})$ est
cohérent pour tout $\mathcal{O}_X$-Module cohérent $\mathcal{F}$
(\hyperref[III.4.4.2-fr]{III, 4.4.2}).

Notons enfin que le critère (\hyperref[II.7.3.10-fr]{7.3.10}) est pris comme
\emph{définition} des morphismes propres en Géométrie algébrique classique.
Nous ne l'avons mentionné que pour cette raison, le critère
(\hyperref[II.7.3.8-fr]{7.3.8}) semblant plus maniable dans toutes les
applications à notre connaissance.
\end{remark}

\oldpage[II]{148}
\subsection{Courbes algébriques et corps de fonctions de dimension 1}
\label{subsection:II.7.4-fr}

Le but de ce numéro est de montrer comment se formule en langage des schémas
la notion classique de courbe algébrique (telle qu'elle est exposée par
exemple dans le livre de C.~Chevalley [23]). Dans tout le numéro, \emph{on
désigne par $k$ un corps, tous les schémas considérés sont des $k$-schémas de
type fini, et tous les morphismes des $k$-morphismes}.

\begin{proposition}[7.4.1]
\label{II.7.4.1-fr}
Soit $X$ un préschéma de type fini sur $k$ (donc noethérien)~; soient $x_i$
$(1\leq i\leq n)$ les points génériques des composantes irréductibles $X_i$ de
$X$, et soit $K_i=k(x_i)$ $(1\leq i\leq n)$. Les conditions suivantes sont
équivalentes~:
\begin{enumerate}
\item[\rm a)] Chacun des $K_i$ est une extension de $k$ dont le degré de
transcendance est égal à 1.
\item[\rm b)] Pour tout point fermé $x$ de $X$, l'anneau local
$\mathcal{O}_x$ est de dimension 1 (\hyperref[II.7.1.5-fr]{7.1.5}).
\item[\rm c)] Les parties fermées irréductibles de $X$ distinctes des $X_i$
sont les points fermés de $X$.
\end{enumerate}
\end{proposition}

Comme $X$ est quasi-compact, toute partie fermée irréductible $F$ de $X$
contient un point fermé (\hyperref[0.2.1.3-fr]{0, 2.1.3}). En vertu de
(\hyperref[I.2.4.2-fr]{I, 2.4.2}), il y a correspondance biunivoque entre les
idéaux premiers de $\mathcal{O}_x$ et les parties fermées irréductibles de $X$
contenant $x$ (\hyperref[I.1.1.14-fr]{I, 1.1.14})~; l'équivalence de b) et c)
en découle aussitôt. D'autre part, si $\mathfrak{p}_\alpha$
$(1\leq\alpha\leq r)$ sont les idéaux premiers minimaux de l'anneau local
noethérien $\mathcal{O}_x$, les anneaux locaux
$\mathcal{O}_x/\mathfrak{p}_\alpha$ sont intègres, et ont pour corps des
fractions ceux des $K_i$ tels que $x\in X_i$. En outre, on sait ([1],
p.~4-06, th.~2) que la dimension d'une $k$-algèbre locale intègre de type fini
est égale au degré de transcendance sur $k$ de son corps des fractions. Enfin,
la dimension de $\mathcal{O}_x$ est borne supérieure des dimensions des
$\mathcal{O}_x/\mathfrak{p}_\alpha$~; or, la condition a) implique que ces
dimensions sont égales à 1, donc a) implique b)~; inversement, si
$\mathcal{O}_x$ est de dimension 1, aucun des $\mathfrak{p}_\alpha$ ne peut
être égal à l'idéal maximal de $\mathcal{O}_x$, sans quoi $\mathcal{O}_x$
serait de dimension 0~; donc chacun des
$\mathcal{O}_x/\mathfrak{p}_\alpha$ est de dimension 1, ce qui montre que b)
entraîne a).

On notera que sous les conditions de (\hyperref[II.7.4.1-fr]{7.4.1}),
l'ensemble $X$ est \emph{vide ou infini}, comme il résulte aussitôt de
(\hyperref[I.6.4.4-fr]{I, 6.4.4}).

\begin{definition}[7.4.2]
\label{II.7.4.2-fr}
On appelle courbe algébrique sur $k$ un schéma algébrique non vide sur $k$
vérifiant les conditions de (\hyperref[II.7.4.1-fr]{7.4.1}).
\end{definition}

Dans le langage de la dimension, qui sera introduit au chapitre IV, cela
s'exprimera en disant qu'une courbe algébrique sur $k$ est un $k$-schéma
algébrique non vide dont \emph{toutes les composantes irréductibles sont de
dimension 1}.

On notera que si $X$ est une courbe algébrique sur $k$, les sous-préschémas
fermés réduits $X_i$ $(1\leq i\leq n)$ de $X$ ayant pour espaces sous-jacents
les composantes irréductibles de $X$ sont des courbes algébriques sur $k$.

\begin{corollary}[7.4.3]
\label{II.7.4.3-fr}
Soit $X$ une courbe algébrique irréductible. Le seul point non fermé de $X$ est
son point générique. Les parties fermées de $X$ distinctes de $X$ sont les
ensembles finis de points fermés~; ce sont aussi les seules parties de $X$ qui
ne soient pas partout denses.
\end{corollary}

Si un point $x\in X$ n'est pas fermé, son adhérence dans $X$ est une partie
fermée irréductible de $X$, donc nécessairement $X$ tout entier en vertu de
(\hyperref[II.7.4.1-fr]{7.4.1}), et par suite $x$ est le point générique de
$X$. Une partie fermée $F$ de $X$, distincte de $X$, ne peut contenir

\oldpage[II]{149}
le point générique de $X$, donc tous ses points sont fermés (dans $X$ et
\emph{a fortiori} dans $F$)~; en considérant le sous-préschéma fermé réduit de
$X$ ayant $F$ pour espace sous-jacent (\hyperref[I.5.2.1-fr]{I, 5.2.1}), il
résulte donc de (\hyperref[I.6.2.2-fr]{I, 6.2.2}) que $F$ est fini et discret.
L'adhérence dans $X$ d'une partie infinie de $X$ est donc nécessairement égale
à $X$.

Lorsque $X$ est une courbe algébrique quelconque, en appliquant
(\hyperref[II.7.4.3-fr]{7.4.3}) aux composantes irréductibles de $X$, on voit
que les seuls points non fermés de $X$ sont les points génériques de ces
composantes.

\begin{corollary}[7.4.4]
\label{II.7.4.4-fr}
Soient $X$ et $Y$ deux courbes algébriques irréductibles sur $k$, $f:X\to Y$
un $k$-morphisme. Pour que $f$ soit dominant, il faut et il suffit que
$f^{-1}(y)$ soit fini pour tout $y\in Y$.
\end{corollary}

En effet, si $f$ n'est pas dominant, $f(X)$ est nécessairement une partie
\emph{finie} de $Y$ en vertu de (\hyperref[II.7.4.3-fr]{7.4.3}), donc il n'est
pas possible que $f^{-1}(y)$ soit fini pour tout point de $Y$, sinon $X$ serait
fini, ce qui est absurde (\hyperref[II.7.4.1-fr]{7.4.1}). Inversement, si $f$
est dominant, pour tout $y\in Y$ distinct du point générique $\eta$ de $Y$,
$f^{-1}(y)$ est fermé dans $X$ puisque $\{y\}$ est fermé dans $Y$
(\hyperref[II.7.4.3-fr]{7.4.3})~; d'autre part, par hypothèse, $f^{-1}(y)$ ne
contient pas le point générique $\xi$ de $X$, donc est fini en vertu de
(\hyperref[II.7.4.3-fr]{7.4.3}). Enfin, pour voir que lorsque $f$ est dominant,
$f^{-1}(\eta)$ est fini, on note que la fibre $f^{-1}(\eta)$ est un schéma
irréductible de type fini sur $k(\eta)$, et de point générique $\xi$
(\hyperref[I.6.3.9-fr]{I, 6.3.9} et
\hyperref[I.6.4.11-fr]{6.4.11}). Comme $k(\xi)$ et $k(\eta)$ sont des
extensions de type fini de $k$, de même degré de transcendance 1, $k(\xi)$ est
nécessairement une extension de degré fini de $k(\eta)$, donc $\xi$ est fermé
dans $f^{-1}(\eta)$ (\hyperref[I.6.4.2-fr]{I, 6.4.2}), et $f^{-1}(\eta)$ est
par suite réduit au point $\xi$.

Nous verrons au chapitre III qu'un morphisme \emph{propre} $f:X\to Y$ de
préschémas noethériens, tel que $f^{-1}(y)$ soit fini pour tout $y\in Y$, est
nécessairement \emph{fini}~; il s'ensuivra donc de
(\hyperref[II.7.4.4-fr]{7.4.4}) qu'un morphisme dominant propre d'une courbe
algébrique irréductible dans une courbe algébrique est \emph{fini}.

\begin{corollary}[7.4.5]
\label{II.7.4.5-fr}
Soit $X$ une courbe algébrique sur $k$. Pour que $X$ soit régulière, il faut et
il suffit que $X$ soit normale, ou encore que les anneaux locaux de ses points
fermés soient des anneaux de valuation discrète.
\end{corollary}

Cela résulte aussitôt de la condition b) de
(\hyperref[II.7.4.1-fr]{7.4.1}) et de
(\hyperref[II.7.1.6-fr]{7.1.6}).

\begin{corollary}[7.4.6]
\label{II.7.4.6-fr}
Soient $X$ une courbe algébrique réduite, $\mathcal{A}$ une
$\mathcal{R}(X)$-Algèbre cohérente réduite~; alors la fermeture intégrale $X'$
de $X$ relativement à $\mathcal{A}$ (\hyperref[II.6.3.4-fr]{6.3.4}) est une
courbe algébrique normale, et le morphisme canonique $X'\to X$ est fini.
\end{corollary}

Le fait que $X'\to X$ soit fini résulte de
(\hyperref[II.6.3.10-fr]{6.3.10})~; $X'$ est donc un $k$-schéma algébrique~;
en outre, si $x_i$ $(1\leq i\leq n)$ sont les points génériques des composantes
irréductibles de $X$, $x'_j$ $(1\leq j\leq m)$ ceux des composantes
irréductibles de $X'$, chacun des $k(x'_j)$ est extension algébrique finie de
l'un des $k(x_i)$ (\hyperref[II.6.3.6-fr]{6.3.6}), donc a un degré de
transcendance 1 sur $k$. $X'$ est donc bien une courbe algébrique sur $k$, et
en outre on sait que $X'$ est somme d'un nombre fini de schémas intègres et
normaux (\hyperref[II.6.3.6-fr]{6.3.6} et
\hyperref[II.6.3.7-fr]{6.3.7}).

\begin{env}[7.4.7]
\label{II.7.4.7-fr}
On dit qu'une courbe algébrique $X$ sur $k$ est \emph{complète} si elle est
\emph{propre} sur $k$.
\end{env}

\begin{corollary}[7.4.8]
\label{II.7.4.8-fr}
Pour qu'une courbe algébrique réduite $X$ sur $k$ soit complète, il faut et il
suffit que sa normalisée $X'$ le soit.
\end{corollary}

\oldpage[II]{150}
En effet, le morphisme canonique $f:X'\to X$ est alors fini
(\hyperref[II.7.4.6-fr]{7.4.6}), donc propre
(\hyperref[II.6.1.11-fr]{6.1.11}) et surjectif
(\hyperref[II.6.3.8-fr]{6.3.8})~; si $g:X\to\operatorname{Spec}(k)$ est le
morphisme structural, $g$ et $g\circ f$ sont donc propres simultanément,
comme il résulte de (\hyperref[II.5.4.2-fr]{5.4.2, (iii)}) et
(\hyperref[II.5.4.3-fr]{5.4.3, (ii)}), $g$ étant séparé par hypothèse.

\begin{proposition}[7.4.9]
\label{II.7.4.9-fr}
Soient $X$ une courbe algébrique normale sur $k$, $Y$ un $k$-schéma algébrique
propre sur $k$. Toute $k$-application rationnelle de $X$ dans $Y$ est alors
partout définie, autrement dit est un morphisme.
\end{proposition}

En effet, il résulte de (\hyperref[II.7.3.7-fr]{7.3.7}) qu'aux points $x\in X$
où une telle application n'est pas définie, la dimension de $\mathcal{O}_x$
devrait être $\geq2$, donc l'ensemble de ces points est vide~; la dernière
assertion provient de (\hyperref[I.7.2.3-fr]{I, 7.2.3}).

\begin{corollary}[7.4.10]
\label{II.7.4.10-fr}
Une courbe algébrique normale sur $k$ est quasi-projective sur $k$.
\end{corollary}

Comme $X$ est somme d'un nombre fini de courbes algébriques intègres et
normales (\hyperref[II.6.3.8-fr]{6.3.8}), on peut se borner au cas où $X$ est
intègre (\hyperref[II.5.3.6-fr]{5.3.6}). Comme $X$ est quasi-compact, il est
recouvert par un nombre fini d'ouverts affines $U_i$ $(1\leq i\leq n)$, et
comme chacun de ces derniers est de type fini sur $k$, il existe un entier
$n_i$ et une $k$-immersion $f_i:U_i\to\mathbf{P}_k^{n_i}$
(\hyperref[II.5.3.3-fr]{5.3.3} et
\hyperref[II.5.3.4-fr]{5.3.4, (i)}). Comme $U_i$ est dense dans $X$, il résulte
de (\hyperref[II.7.4.9-fr]{7.4.9}) que $f_i$ se prolonge en un $k$-morphisme
$g_i:X\to\mathbf{P}_k^{n_i}$, d'où un $k$-morphisme
$g=(g_1,\ldots,g_n)_k$ de $X$ dans le produit $P$ des
$\mathbf{P}_k^{n_i}$ sur $k$. En outre, pour chaque indice $i$, comme la
restriction de $g_i$ à $U_i$ est une immersion, il en est de même de la
restriction de $g$ à $U_i$ (\hyperref[I.5.3.14-fr]{I, 5.3.14}). Comme les
$U_i$ recouvrent $X$ et que $g$ est séparé
(\hyperref[I.5.5.1-fr]{I, 5.5.1, (v)}), $g$ est une immersion de $X$ dans $P$
(\hyperref[I.8.2.8-fr]{I, 8.2.8}). Comme le morphisme de Segre
(\hyperref[II.4.3.3-fr]{4.3.3}) fournit une immersion de $P$ dans un
$\mathbf{P}_k^N$, cela achève de prouver que $X$ est quasi-projectif.

\begin{corollary}[7.4.11]
\label{II.7.4.11-fr}
Une courbe algébrique normale $X$ est isomorphe au schéma induit sur un ouvert
partout dense d'une courbe algébrique normale et complète $\widehat{X}$,
déterminée à un isomorphisme unique près.
\end{corollary}

Si $X_1$, $X_2$ sont deux courbes normales et complètes, il résulte aussitôt
de (\hyperref[II.7.4.9-fr]{7.4.9}) que tout isomorphisme d'un ouvert $U_1$
dense dans $X_1$, sur un ouvert $U_2$ dense dans $X_2$, se prolonge de façon
unique en un isomorphisme de $X_1$ sur $X_2$~; d'où l'assertion d'unicité. Pour
prouver l'existence de $\widehat{X}$, il suffit de remarquer que l'on peut
considérer $X$ comme un sous-schéma d'un fibré projectif $\mathbf{P}_k^n$
(\hyperref[II.7.4.10-fr]{7.4.10}). Soit $\overline{X}$ l'adhérence de $X$ dans
$\mathbf{P}_k^n$ (\hyperref[I.9.5.11-fr]{I, 9.5.11})~; comme $X$ est induit
par $\overline{X}$ sur un ouvert dense dans $\overline{X}$
(\hyperref[I.9.5.10-fr]{I, 9.5.10}), les points génériques $x_i$ des
composantes irréductibles de $X$ sont aussi ceux des composantes irréductibles
de $\overline{X}$, et les $k(x_i)$ sont les mêmes pour ces deux schémas, donc
(\hyperref[II.7.4.1-fr]{7.4.1}) $\overline{X}$ est une courbe algébrique sur
$k$, qui est réduite (\hyperref[I.9.5.9-fr]{I, 9.5.9}) et projective sur $k$
(\hyperref[II.5.5.1-fr]{5.5.1}), donc complète
(\hyperref[II.5.5.3-fr]{5.5.3}). Prenons alors pour $\widehat{X}$ la
\emph{normalisée} de $\overline{X}$, qui est encore complète
(\hyperref[II.7.4.8-fr]{7.4.8})~; en outre, si
$h:\widehat{X}\to\overline{X}$ est le morphisme canonique, la restriction de
$h$ à $h^{-1}(X)$ est un isomorphisme sur $X$ puisque $X$ est normale
(\hyperref[II.6.3.4-fr]{6.3.4}), et comme $h^{-1}(X)$ contient les points
génériques des composantes irréductibles de $\widehat{X}$
(\hyperref[II.6.3.8-fr]{6.3.8}), il est dense dans $\widehat{X}$, ce qui achève
la démonstration.

\oldpage[II]{151}

\begin{remark}[7.4.12]
\label{II.7.4.12-fr}
Nous montrerons au chapitre V que la conclusion de
(\hyperref[II.7.4.10-fr]{7.4.10}) est encore valable sans supposer la courbe
normale (ni même réduite)~; nous montrerons aussi que pour qu'une courbe
algébrique (réduite ou non) soit affine, il faut et il suffit que ses
composantes irréductibles (réduites) ne soient pas complètes.
\end{remark}

\begin{corollary}[7.4.13]
\label{II.7.4.13-fr}
Soient $X$ une courbe normale irréductible de corps $R(X)=K$, $Y$ une courbe
intègre complète de corps $L=R(Y)$. Il y a une correspondance biunivoque
canonique entre les $k$-morphismes dominants $X\to Y$ et les
$k$-monomorphismes $L\to K$.
\end{corollary}

En vertu de (\hyperref[II.7.4.9-fr]{7.4.9}), les $k$-applications rationnelles
de $X$ dans $Y$ s'identifient aux $k$-morphismes $u:X\to Y$. Les morphismes
$u$ dominants étant caractérisés par le fait que $u(x)=y$ (en désignant par
$x$ et $y$ les points génériques respectifs de $X$ et de $Y$), le corollaire
résulte de ces remarques et de (\hyperref[I.7.1.13-fr]{I, 7.1.13}).

\begin{env}[7.4.14]
\label{II.7.4.14-fr}
On peut préciser le résultat de (\hyperref[II.7.4.13-fr]{7.4.13}) lorsqu'on
prend pour $Y$ la \emph{droite projective}
$\mathbf{P}_k^1=\operatorname{Proj}(k[T_0,T])$, $T_0$ et $T$ étant deux
indéterminées. C'est bien là un schéma intègre
(\hyperref[II.2.4.4-fr]{2.4.4}), et le schéma induit sur l'ouvert $D_+(T_0)$
de $Y$ est isomorphe à $\operatorname{Spec}(k[T])$
(\hyperref[II.2.3.6-fr]{2.3.6}), donc le point générique de $Y$ est l'idéal
$(0)$ de $k[T]$ et son corps de fonctions rationnelles $k(T)$, ce qui montre
que $Y$ est une courbe algébrique complète sur $k$. En outre, le seul idéal
premier gradué de $S=k[T_0,T]$ contenant $T_0$ et distinct de $S_+$ est
l'idéal principal $(T_0)$, donc le complémentaire de $D_+(T_0)$ dans
$Y=\mathbf{P}_k^1$ est réduit à \emph{un point fermé}, dit «~point à
l'infini~» que nous noterons $\infty$ (pour une étude générale des rapports
entre fibrés vectoriels et fibrés projectifs, voir
\hyperref[subsection:II.8.4-fr]{8.4}). Avec ces notations~:
\end{env}

\begin{corollary}[7.4.15]
\label{II.7.4.15-fr}
Soit $X$ une courbe normale irréductible de corps $R(X)=K$. Il existe une
application bijective canonique de $K$ sur l'ensemble des morphismes $u$ de
$X$ dans $\mathbf{P}_k^1$, distincts du morphisme constant de valeur
$\infty$. Pour que $u$ soit dominant, il faut et il suffit que l'élément
correspondant de $K$ soit transcendant sur $k$.
\end{corollary}

Cet énoncé résulte immédiatement de (\hyperref[II.7.4.9-fr]{7.4.9}) et du

\begin{lemma}[7.4.15.1]
\label{II.7.4.15.1-fr}
Soit $X$ un préschéma intègre sur $k$, et soit $K=R(X)$ son corps des fonctions
rationnelles. Il existe une application bijective canonique de l'ensemble
$K$ sur l'ensemble des applications rationnelles $u$ de $X$ dans
$\mathbf{P}_k^1$, distinctes du morphisme constant de valeur $\infty$. Pour
qu'une telle application rationnelle soit dominante, il faut et il suffit que
l'élément correspondant de $K$ soit transcendant sur $k$.
\end{lemma}

Tout d'abord, les applications rationnelles de $X$ dans $\mathbf{P}_k^1$
correspondent biunivoquement aux points de $\mathbf{P}_k^1$ à valeurs dans
l'extension $K$ de $k$ (\hyperref[I.7.1.12-fr]{I, 7.1.12}). Si un tel point
est localisé (\hyperref[I.3.4.5-fr]{I, 3.4.5}) au point générique de
$\mathbf{P}_k^1$, l'application rationnelle correspondante est évidemment
dominante. Dans le cas contraire, comme tout point de $\mathbf{P}_k^1$
distinct du point générique est fermé (\hyperref[II.7.4.3-fr]{7.4.3}),
l'image du domaine de définition $U$ de $u$ par l'unique morphisme
$U\to\mathbf{P}_k^1$ de la classe $u$ (\hyperref[I.7.2.2-fr]{I, 7.2.2}),
est réduite à un point fermé $y$ de $\mathbf{P}_k^1$, et ce morphisme (qui
n'est pas nécessairement partout défini dans $X$) n'est donc pas dominant~;
par abus de langage, on dit alors que l'application rationnelle $u$ est
«~constante, de valeur $y$~». Il reste à mettre en correspondance biunivoque
les points de $\mathbf{P}_k^1$ à valeurs dans $K$, de localité
(\hyperref[I.3.4.5-fr]{I, 3.4.5}) distincte de $\infty$, et les éléments de
$K$, et à vérifier que la localité d'un tel point est le point générique de
$\mathbf{P}_k^1$ si et seulement s'il correspond à un

\oldpage[II]{152}
élément transcendant sur $k$. Or, cette vérification est immédiate sur
(\hyperref[II.4.2.6-fr]{4.2.6, exemple 1\textsuperscript{o}}).

\begin{corollary}[7.4.16]
\label{II.7.4.16-fr}
Soient $X$, $Y$ deux courbes algébriques sur $k$, normales, complètes et
irréductibles~; soient $K=R(X)$, $L=R(Y)$ leurs corps. Il existe une
correspondance biunivoque canonique entre l'ensemble des $k$-isomorphismes
$X\xrightarrow{\sim}Y$ et l'ensemble des $k$-isomorphismes
$L\xrightarrow{\sim}K$.
\end{corollary}

C'est une conséquence évidente de (\hyperref[II.7.4.13-fr]{7.4.13}).

\begin{env}[7.4.17]
\label{II.7.4.17-fr}
Le corollaire (\hyperref[II.7.4.16-fr]{7.4.16}) montre qu'une courbe algébrique
sur $k$, normale, complète et irréductible, est \emph{déterminée par son corps
de fonctions rationnelles $K$ à un isomorphisme unique près}~; par définition,
$K$ est une extension de type fini de $k$, de degré de transcendance $1$ (ce
qu'on appelle classiquement un \emph{corps de fonctions algébriques d'une
variable}). De plus~:
\end{env}

\begin{proposition}[7.4.18]
\label{II.7.4.18-fr}
Pour toute extension $K$ de $k$, de type fini et de degré de transcendance
$1$, il existe une courbe algébrique normale, complète et irréductible $X$
telle que $R(X)=K$ (déterminée à isomorphisme unique près). L'ensemble des
anneaux locaux de $X$ s'identifie (\hyperref[I.8.2.1-fr]{I, 8.2.1}) à
l'ensemble formé de $K$ et des anneaux de valuation contenant $k$ et ayant
$K$ comme corps des fractions.
\end{proposition}

En effet, $K$ est une extension de degré fini d'une extension transcendante
pure $k(T)$ de $k$, qui s'identifie, comme on l'a vu, au corps des fonctions
rationnelles de la droite projective $Y=\mathbf{P}_k^1$. Soit $X$ la
\emph{fermeture intégrale} de $Y$ relativement à $K$
(\hyperref[II.6.3.4-fr]{6.3.4})~; $X$ est une courbe algébrique normale de
corps $K$ (\hyperref[II.6.3.7-fr]{6.3.7}), et elle est complète puisque le
morphisme $X\to Y$ est fini (\hyperref[II.7.4.6-fr]{7.4.6}). Les anneaux
locaux $\mathcal{O}_x$ de $X$ sont~: le corps $K$ lorsque $x$ est le point
générique~; si $x$ est distinct du point générique, $\mathcal{O}_x$ est un
anneau de valuation discrète contenant $k$ et ayant $K$ comme corps des
fractions (\hyperref[II.7.4.5-fr]{7.4.5}). Inversement, soit $A$ un tel
anneau~; comme le morphisme $X\to\operatorname{Spec}(k)$ est propre et que
$A$ domine $k$, il domine un anneau local $\mathcal{O}_x$ de $X$
(\hyperref[II.7.3.10-fr]{7.3.10})~; ce dernier étant un anneau de valuation
ayant $K$ comme corps des fractions, est donc nécessairement égal à $A$.

\begin{namedtheorem}[7.4.19]{Remarques}
\label{II.7.4.19-fr}
Il résulte de (\hyperref[II.7.4.16-fr]{7.4.16}) et
(\hyperref[II.7.4.18-fr]{7.4.18}) que \emph{la donnée d'une courbe algébrique
sur $k$, normale, complète et irréductible, est essentiellement équivalente à
la donnée d'une extension $K$ de $k$, de type fini et de degré de
transcendance $1$}. On notera que si $k'$ est une extension du corps de base
$k$, $X\otimes_k k'$ sera encore une courbe algébrique complète sur $k'$
(\hyperref[II.5.4.2-fr]{5.4.2, (iii)}), mais en général elle ne sera ni réduite
ni irréductible. Elle le sera cependant si $K$ est une extension séparable de
$k$ et si $k$ est algébriquement fermé dans $K$ (ce qui s'exprime, dans une
terminologie classique que nous ne suivrons pas, en disant que $K$ est une
«~extension régulière~» de $k$). Mais même dans ce cas, il peut se faire que
$X\otimes_k k'$ ne soit pas normale. Le lecteur trouvera des détails sur ces
questions dans le chapitre IV.
\end{namedtheorem}

\section{Schémas éclatés~; cônes projetants~; fermeture projective}
\label{section:II.8-fr}

\subsection{Préschémas éclatés}
\label{subsection:II.8.1-fr}

\begin{env}[8.1.1]
\label{II.8.1.1-fr}
Soient $Y$ un préschéma, et pour tout entier $n\geq0$, soit $\mathcal{I}_n$ un
Idéal quasi-cohérent de $\mathcal{O}_Y$~; on suppose vérifiées les conditions
suivantes~:
\begin{equation}
\label{II.8.1.1.1-fr}
  \mathcal{I}_0=\mathcal{O}_Y,\qquad
  \mathcal{I}_n\subset\mathcal{I}_m\quad\text{pour }m\leq n
\tag{8.1.1.1}
\end{equation}
\begin{equation}
\label{II.8.1.1.2-fr}
  \mathcal{I}_m\mathcal{I}_n\subset\mathcal{I}_{m+n}
  \quad\text{quels que soient }m,n.
\tag{8.1.1.2}
\end{equation}

\oldpage[II]{153}
On notera que ces hypothèses entraînent
\begin{equation}
\label{II.8.1.1.3-fr}
  \mathcal{I}_1^n\subset\mathcal{I}_n.
\tag{8.1.1.3}
\end{equation}

Posons
\begin{equation}
\label{II.8.1.1.4-fr}
  \mathcal{S}=\bigoplus_{n\geq0}\mathcal{I}_n.
\tag{8.1.1.4}
\end{equation}

Il résulte de (\hyperref[II.8.1.1.1-fr]{8.1.1.1}) et
(\hyperref[II.8.1.1.2-fr]{8.1.1.2}) que $\mathcal{S}$ est une
$\mathcal{O}_Y$-Algèbre graduée quasi-cohérente, et définit donc un $Y$-schéma
$X=\operatorname{Proj}(\mathcal{S})$. Si $\mathcal{J}$ est un Idéal
\emph{inversible} de $\mathcal{O}_Y$,
$\mathcal{I}_n\otimes_{\mathcal{O}_Y}\mathcal{J}^{\otimes n}$ s'identifie
canoniquement à $\mathcal{I}_n\mathcal{J}^n$. Si on remplace donc les
$\mathcal{I}_n$ par les $\mathcal{I}_n\mathcal{J}^n$, ce qui remplace
$\mathcal{S}$ par une $\mathcal{O}_Y$-Algèbre quasi-cohérente
$\mathcal{S}_{(\mathcal{J})}$,
$X_{(\mathcal{J})}=\operatorname{Proj}(\mathcal{S}_{(\mathcal{J})})$ est
canoniquement isomorphe à $X$ (\hyperref[II.3.1.8-fr]{3.1.8}).
\end{env}

\begin{env}[8.1.2]
\label{II.8.1.2-fr}
Supposons $Y$ \emph{localement intègre}, de sorte que le faisceau
$\mathcal{R}(Y)$ des fonctions rationnelles est une $\mathcal{O}_Y$-Algèbre
quasi-cohérente (\hyperref[I.7.3.7-fr]{I, 7.3.7}). Nous dirons qu'un
sous-$\mathcal{O}_Y$-Module $\mathcal{I}$ de $\mathcal{R}(Y)$ est un
\emph{Idéal fractionnaire} de $\mathcal{R}(Y)$ s'il est de \emph{type fini}
(\hyperref[0.5.2.1-fr]{0, 5.2.1}). Supposons donné, pour tout $n\geq0$, un
Idéal fractionnaire quasi-cohérent $\mathcal{I}_n$ de $\mathcal{R}(Y)$, tel
que $\mathcal{I}_0=\mathcal{O}_Y$, et que la condition
(\hyperref[II.8.1.1.2-fr]{8.1.1.2}) soit vérifiée (mais non nécessairement la
seconde condition (\hyperref[II.8.1.1.1-fr]{8.1.1.1}))~; on peut encore alors
définir par la formule (\hyperref[II.8.1.1.4-fr]{8.1.1.4}) une
$\mathcal{O}_Y$-Algèbre graduée quasi-cohérente, et le $Y$-schéma correspondant
$X=\operatorname{Proj}(\mathcal{S})$~; on aura encore un isomorphisme
canonique de $X$ sur $X_{(\mathcal{J})}$ pour tout Idéal fractionnaire
inversible $\mathcal{J}$ de $\mathcal{R}(Y)$.
\end{env}

\begin{definition}[8.1.3]
\label{II.8.1.3-fr}
Soit $Y$ un préschéma (resp. un préschéma localement intègre), et soit
$\mathcal{I}$ un Idéal quasi-cohérent de $\mathcal{O}_Y$ (resp. un Idéal
fractionnaire quasi-cohérent de $\mathcal{R}(Y)$). On dit que le $Y$-schéma
$X=\operatorname{Proj}(\bigoplus_{n\geq0}\mathcal{I}^n)$ est obtenu en
\emph{faisant éclater l'Idéal $\mathcal{I}$}, ou est le \emph{préschéma éclaté
de $Y$ relativement à $\mathcal{I}$}. Lorsque $\mathcal{I}$ est un Idéal
quasi-cohérent de $\mathcal{O}_Y$ et $Y'$ le sous-préschéma fermé de $Y$ défini
par $\mathcal{I}$, on dit aussi que $X$ est le $Y$-schéma obtenu en faisant
éclater $Y'$.
\end{definition}

Par définition, $\mathcal{S}=\bigoplus_{n\geq0}\mathcal{I}^n$ est alors
engendrée par $\mathcal{S}_1=\mathcal{I}$~; si $\mathcal{I}$ est un
$\mathcal{O}_Y$-Module de \emph{type fini}, $X$ est donc \emph{projectif} sur
$Y$ (\hyperref[II.5.5.2-fr]{5.5.2}). Sans hypothèse sur $\mathcal{I}$, le
$\mathcal{O}_X$-Module $\mathcal{O}_X(1)$ est \emph{inversible}
(\hyperref[II.3.2.5-fr]{3.2.5}) et \emph{très ample} en vertu de
(\hyperref[II.4.4.3-fr]{4.4.3}) pour le morphisme structural $X\to Y$.

On notera que si $f:X\to Y$ est le morphisme structural, la restriction de
$f$ à $f^{-1}(Y-Y')$ est un \emph{isomorphisme} sur $Y-Y'$ lorsque
$\mathcal{I}$ est un Idéal de $\mathcal{O}_Y$ et $Y'$ le sous-préschéma fermé
qu'il définit~: en effet, la question étant locale sur $Y$, il suffit de
supposer $\mathcal{I}=\mathcal{O}_Y$, et notre assertion résulte de
(\hyperref[II.3.1.7-fr]{3.1.7}).

Lorsqu'on remplace $\mathcal{I}$ par $\mathcal{I}^d$ ($d>0$), le $Y$-schéma
éclaté $X$ est remplacé par un $Y$-schéma canoniquement isomorphe $X'$
(\hyperref[II.8.1.1-fr]{8.1.1})~; de même, pour tout Idéal (resp. Idéal
fractionnaire) inversible $\mathcal{J}$, le préschéma éclaté
$X_{(\mathcal{J})}$ relativement à l'Idéal $\mathcal{I}\mathcal{J}$ est
canoniquement isomorphe à $X$ (\hyperref[II.8.1.1-fr]{8.1.1}).

En particulier, lorsque $\mathcal{I}$ est un Idéal (resp. un Idéal
fractionnaire) \emph{inversible}, le $Y$-schéma obtenu en faisant éclater
$\mathcal{I}$ est \emph{isomorphe} à $Y$ (\hyperref[II.3.1.7-fr]{3.1.7}).

\begin{proposition}[8.1.4]
\label{II.8.1.4-fr}
Soit $Y$ un préschéma intègre.
\begin{enumerate}
\item[\rm(i)] Pour toute suite $(\mathcal{I}_n)$ d'Idéaux fractionnaires
quasi-cohérents de $\mathcal{R}(Y)$ vérifiant
(\hyperref[II.8.1.1.2-fr]{8.1.1.2})

\oldpage[II]{154}
et tels que $\mathcal{I}_0=\mathcal{O}_Y$, le $Y$-schéma
$X=\operatorname{Proj}(\bigoplus_{n\geq0}\mathcal{I}_n)$ est intègre et le
morphisme structural $f:X\to Y$ dominant.
\item[\rm(ii)] Soit $\mathcal{I}$ un Idéal fractionnaire quasi-cohérent de
$\mathcal{R}(Y)$, et soit $X$ le $Y$-schéma éclaté de $Y$ relativement à
$\mathcal{I}$. Si $\mathcal{I}\neq0$, le morphisme structural $f:X\to Y$ est
alors birationnel et surjectif.
\end{enumerate}
\end{proposition}

\begin{enumerate}
\item[\rm(i)] résulte de ce que $\mathcal{S}=\bigoplus_{n\geq0}\mathcal{I}_n$
est une $\mathcal{O}_Y$-Algèbre \emph{intègre}
(\hyperref[II.3.1.12-fr]{3.1.12} et \hyperref[II.3.1.14-fr]{3.1.14}) puisque
pour tout $y\in Y$, $\mathcal{O}_y$ est un anneau intègre
(\hyperref[I.5.1.4-fr]{I, 5.1.4}).
\item[\rm(ii)] D'après (i), $X$ est intègre~; si en outre, $x$ et $y$ sont les
points génériques de $X$ et $Y$, on a $f(x)=y$, et il faut prouver que $k(x)$
est de rang $1$ sur $k(y)$. Or $x$ est aussi le point générique de la fibre
$f^{-1}(y)$~; si $\psi$ est le morphisme canonique $Z\to Y$, où
$Z=\operatorname{Spec}(k(y))$, le préschéma $f^{-1}(y)$ s'identifie à
$\operatorname{Proj}(\mathcal{S}')$, où
$\mathcal{S}'=\psi^*(\mathcal{S})$ (\hyperref[II.3.5.3-fr]{3.5.3}). Mais il est
clair que $\mathcal{S}'=\bigoplus_{n\geq0}(\mathcal{I}_y)^n$, et comme
$\mathcal{I}$ est un Idéal fractionnaire quasi-cohérent non nul de
$\mathcal{R}(Y)$, $\mathcal{I}_y\neq0$
(\hyperref[I.7.3.6-fr]{I, 7.3.6}), d'où $\mathcal{I}_y=k(y)$~;
$\operatorname{Proj}(\mathcal{S}')$ s'identifie donc à
$\operatorname{Spec}(k(y))$ (\hyperref[II.3.1.7-fr]{3.1.7}), d'où la
conclusion.
\end{enumerate}

Nous démontrerons une \emph{réciproque} de
(\hyperref[II.8.1.4-fr]{8.1.4}) dans
(\hyperref[III.2.3.8-fr]{III, 2.3.8}).

\begin{env}[8.1.5]
\label{II.8.1.5-fr}
Revenons à la situation et aux notations de
(\hyperref[II.8.1.1-fr]{8.1.1}). Par définition, les homomorphismes
d'injection $\mathcal{I}_{n+1}\to\mathcal{I}_n$
(\hyperref[II.8.1.1.1-fr]{8.1.1.1}) définissent pour chaque $k\in\mathbf{Z}$ un
homomorphisme injectif de degré $0$ de $\mathcal{S}$-Modules gradués
\begin{equation}
\label{II.8.1.5.1-fr}
  u_k:\mathcal{S}_+(k+1)\to\mathcal{S}(k)~;
\tag{8.1.5.1}
\end{equation}
comme $\mathcal{S}_+(k+1)$ et $\mathcal{S}(k+1)$ sont canoniquement
(TN)-isomorphes, il correspond canoniquement à $u_k$ un homomorphisme injectif
de $\mathcal{O}_X$-Modules (\hyperref[II.3.4.2-fr]{3.4.2})~:
\begin{equation}
\label{II.8.1.5.2-fr}
  \widetilde{u}_k:\mathcal{O}_X(k+1)\to\mathcal{O}_X(k).
\tag{8.1.5.2}
\end{equation}

Rappelons d'autre part (\hyperref[II.3.2.6-fr]{3.2.6}) que l'on a défini des
homomorphismes canoniques
\begin{equation}
\label{II.8.1.5.3-fr}
  \lambda:\mathcal{O}_X(h)\otimes_{\mathcal{O}_X}\mathcal{O}_X(k)
  \to\mathcal{O}_X(h+k)
\tag{8.1.5.3}
\end{equation}
et comme le diagramme
\[
\xymatrix{
  \mathcal{S}(h)\otimes_{\mathcal{S}}\mathcal{S}(k)
    \otimes_{\mathcal{S}}\mathcal{S}(l)\ar[r]\ar[d]
  & \mathcal{S}(h+k)\otimes_{\mathcal{S}}\mathcal{S}(l)\ar[d]\\
  \mathcal{S}(h)\otimes_{\mathcal{S}}\mathcal{S}(k+l)\ar[r]
  & \mathcal{S}(h+k+l)
}
\]
est commutatif, il résulte de la fonctorialité des $\lambda$
(\hyperref[II.3.2.6-fr]{3.2.6}) que les homomorphismes
(\hyperref[II.8.1.5.3-fr]{8.1.5.3}) définissent sur
\begin{equation}
\label{II.8.1.5.4-fr}
  \mathcal{S}_X=\bigoplus_{n\in\mathbf{Z}}\mathcal{O}_X(n)
\tag{8.1.5.4}
\end{equation}

\oldpage[II]{155}
une structure de $\mathcal{O}_X$-Algèbre graduée quasi-cohérente. En outre, le
diagramme
\[
\xymatrix{
  \mathcal{S}(h)\otimes_{\mathcal{S}}\mathcal{S}_+(k+1)\ar[r]\ar[d]_{1\otimes u_k}
  & \mathcal{S}_+(h+k+1)\ar[d]^{u_{k+h}}\\
  \mathcal{S}(h)\otimes_{\mathcal{S}}\mathcal{S}(k)\ar[r]
  & \mathcal{S}(h+k)
}
\]
est commutatif~; la fonctorialité des $\lambda$ montre alors que l'on a un
diagramme commutatif
\begin{equation}
\label{II.8.1.5.5-fr}
\xymatrix{
  \mathcal{O}_X(h)\otimes_{\mathcal{O}_X}\mathcal{O}_X(k+1)
    \ar[r]^{\lambda}\ar[d]_{1\otimes\widetilde{u}_k}
  & \mathcal{O}_X(h+k+1)\ar[d]^{\widetilde{u}_{k+h}}\\
  \mathcal{O}_X(h)\otimes_{\mathcal{O}_X}\mathcal{O}_X(k)
    \ar[r]^{\lambda}
  & \mathcal{O}_X(h+k)
}
\tag{8.1.5.5}
\end{equation}
les flèches horizontales étant les homomorphismes canoniques. On peut donc
dire que les $\widetilde{u}_k$ définissent un \emph{homomorphisme injectif} (de
degré $0$) \emph{de $\mathcal{S}_X$-Modules gradués}
\begin{equation}
\label{II.8.1.5.6-fr}
  \widetilde{u}:\mathcal{S}_X(1)\to\mathcal{S}_X.
\tag{8.1.5.6}
\end{equation}
\end{env}

\begin{env}[8.1.6]
\label{II.8.1.6-fr}
Gardant les notations de (\hyperref[II.8.1.5-fr]{8.1.5}), remarquons maintenant
que, pour $n\geq0$, l'homomorphisme composé
$\widetilde{v}_n=\widetilde{u}_{n-1}\widetilde{u}_{n-2}\cdots\widetilde{u}_0$
est un homomorphisme \emph{injectif} $\mathcal{O}_X(n)\to\mathcal{O}_X$~;
nous désignerons par $\mathcal{I}_{n,X}$ son image, qui est donc un Idéal
quasi-cohérent de $\mathcal{O}_X$, \emph{isomorphe} à $\mathcal{O}_X(n)$. En
outre, le diagramme
\[
\xymatrix{
  \mathcal{O}_X(m)\otimes_{\mathcal{O}_X}\mathcal{O}_X(n)
    \ar[r]^{\lambda}\ar[d]_{\widetilde{v}_m\otimes\widetilde{v}_n}
  & \mathcal{O}_X(m+n)\ar[d]^{\widetilde{v}_{m+n}}\\
  \mathcal{O}_X\ar[r]_{\mathrm{id}}
  & \mathcal{O}_X
}
\]
est commutatif pour $m\geq0$, $n\geq0$. On en conclut les inclusions suivantes
\begin{equation}
\label{II.8.1.6.1-fr}
  \mathcal{I}_{0,X}=\mathcal{O}_X,\qquad
  \mathcal{I}_{n,X}\subset\mathcal{I}_{m,X}
  \quad\text{pour }0\leq m\leq n
\tag{8.1.6.1}
\end{equation}
\begin{equation}
\label{II.8.1.6.2-fr}
  \mathcal{I}_{m,X}\mathcal{I}_{n,X}\subset\mathcal{I}_{m+n,X}
  \quad\text{pour }m\geq0,\ n\geq0.
\tag{8.1.6.2}
\end{equation}
\end{env}

\oldpage[II]{156}

\begin{proposition}[8.1.7]
\label{II.8.1.7-fr}
Soient $Y$ un préschéma, $\mathcal{I}$ un Idéal quasi-cohérent de
$\mathcal{O}_Y$, $X=\operatorname{Proj}(\bigoplus_{n\geq0}\mathcal{I}^n)$ le
$Y$-schéma obtenu en faisant éclater $\mathcal{I}$. On a alors, pour tout
$n>0$, un isomorphisme canonique
\begin{equation}
\label{II.8.1.7.1-fr}
  \mathcal{O}_X(n)\xrightarrow{\sim}\mathcal{I}^n\mathcal{O}_X
  =\mathcal{I}_{n,X}
\tag{8.1.7.1}
\end{equation}
(cf. (\hyperref[0.4.3.5-fr]{0, 4.3.5})), et par suite
$\mathcal{I}^n\mathcal{O}_X$ est un $\mathcal{O}_X$-Module inversible très
ample si $n>0$.
\end{proposition}

La dernière assertion est immédiate, puisque $\mathcal{O}_X(1)$ est inversible
(\hyperref[II.3.2.5-fr]{3.2.5}) et très ample pour $Y$ par définition
(\hyperref[II.4.4.3-fr]{4.4.3} et \hyperref[II.4.4.9-fr]{4.4.9}). D'autre
part, par définition, l'image de $v_n$ n'est autre que
$\mathcal{I}^n\mathcal{S}$, et (\hyperref[II.8.1.7.1-fr]{8.1.7.1}) résulte donc
de l'exactitude du foncteur $\widetilde{\mathcal{M}}$
(\hyperref[II.3.2.4-fr]{3.2.4}) et de la formule
(\hyperref[II.3.2.4.1-fr]{3.2.4.1}).

\begin{corollary}[8.1.8]
\label{II.8.1.8-fr}
Sous les hypothèses de (\hyperref[II.8.1.7-fr]{8.1.7}), si $f:X\to Y$ est le
morphisme structural et $Y'$ le sous-préschéma fermé de $Y$ défini par
$\mathcal{I}$, le sous-préschéma fermé $X'=f^{-1}(Y')$ de $X$ est défini par
$\mathcal{I}\mathcal{O}_X$ (isomorphe canoniquement à $\mathcal{O}_X(1)$),
d'où une suite exacte canonique
\begin{equation}
\label{II.8.1.8.1-fr}
  0\to\mathcal{O}_X(1)\to\mathcal{O}_X\to\mathcal{O}_{X'}\to0.
\tag{8.1.8.1}
\end{equation}
\end{corollary}

Cela résulte de (\hyperref[II.8.1.7.1-fr]{8.1.7.1}) et de
(\hyperref[I.4.4.5-fr]{I, 4.4.5}).

\begin{env}[8.1.9]
\label{II.8.1.9-fr}
Sous les hypothèses de (\hyperref[II.8.1.7-fr]{8.1.7}), on peut préciser la
structure des $\mathcal{I}_{n,X}$. Remarquons en effet que l'homomorphisme
\[
  \widetilde{u}_{-1}:\mathcal{O}_X\to\mathcal{O}_X(-1)
\]
correspond canoniquement à une section $s$ de $\mathcal{O}_X(-1)$ au-dessus de
$X$, que nous appellerons la \emph{section canonique} (relative à
$\mathcal{I}$) (\hyperref[0.5.1.1-fr]{0, 5.1.1}). D'autre part, dans le
diagramme (\hyperref[II.8.1.5.5-fr]{8.1.5.5}), les flèches horizontales sont
des isomorphismes (\hyperref[II.3.2.7-fr]{3.2.7})~; en remplaçant dans ce
diagramme $h$ par $k$ et $k$ par $-1$, on obtient
$\widetilde{u}_k=1_k\otimes\widetilde{u}_{-1}$ ($1_h$ désignant l'identité de
$\mathcal{O}_X(h)$), autrement dit l'homomorphisme $\widetilde{u}_k$ n'est
autre que la \emph{multiplication tensorielle par la section canonique $s$}
(pour tout $k\in\mathbf{Z}$). L'homomorphisme $\widetilde{u}$
(\hyperref[II.8.1.5.6-fr]{8.1.5.6}) s'interprète donc aussi de la même manière.

Par suite, pour tout $n\geq0$, l'homomorphisme
$\widetilde{v}_n:\mathcal{O}_X(n)\to\mathcal{O}_X$ n'est autre que la
multiplication tensorielle par $s^{\otimes n}$~; on en déduit~:
\end{env}

\begin{corollary}[8.1.10]
\label{II.8.1.10-fr}
Avec les notations de (\hyperref[II.8.1.8-fr]{8.1.8}), l'espace sous-jacent à
$X'$ est l'ensemble des $x\in X$ tels que $s(x)=0$, $s$ désignant la section
canonique de $\mathcal{O}_X(-1)$.
\end{corollary}

En effet, si $c_x$ est un générateur de la fibre $(\mathcal{O}_X(1))_x$ en un
point $x$, $s_x\otimes c_x$ s'identifie canoniquement à un générateur de la
fibre de $\mathcal{I}_{1,X}$ au point $x$, et est donc inversible si et
seulement si l'on a
$s_x\notin\mathfrak{m}_x(\mathcal{O}_X(-1))_x$, autrement dit si
$s(x)\neq0$.

\begin{proposition}[8.1.11]
\label{II.8.1.11-fr}
Soient $Y$ un préschéma intègre, $\mathcal{I}$ un Idéal fractionnaire
quasi-cohérent de $\mathcal{R}(Y)$, $X$ le $Y$-schéma obtenu en faisant éclater
$\mathcal{I}$. Alors $\mathcal{I}\mathcal{O}_X$ est un
$\mathcal{O}_X$-Module inversible très ample pour $Y$.
\end{proposition}

La question étant locale sur $Y$ (\hyperref[II.4.4.5-fr]{4.4.5}), on peut se
borner au cas où $Y=\operatorname{Spec}(A)$, où $A$ est un anneau intègre de
corps des fractions $K$ et $\mathcal{I}=\widetilde{\mathfrak{I}}$,
$\mathfrak{I}$ étant un idéal fractionnaire de $K$~; il existe par suite un
élément $a\neq0$ de $A$ tel que $a\mathfrak{I}\subset A$. Posons
$S=\bigoplus_{n\geq0}\mathfrak{I}^n$~; l'application $x\mapsto ax$ est un
$A$-isomorphisme de $\mathfrak{I}^{n+1}=(S(1))_n$ sur
$a\mathfrak{I}^{n+1}=a\mathfrak{I}S_n\subset\mathfrak{I}^n=S_n$,

\oldpage[II]{157}
donc définit un (TN)-isomorphisme de degré $0$ de $S$-modules gradués
$S_+(1)\to a\mathfrak{I}S$. Par ailleurs $x\mapsto a^{-1}x$ est un
isomorphisme de degré $0$ de $S$-modules gradués
$a\mathfrak{I}S\xrightarrow{\sim}\mathfrak{I}S$. On obtient donc ainsi par
composition (\hyperref[II.3.2.4-fr]{3.2.4}) un isomorphisme de
$\mathcal{O}_X$-Modules
$\mathcal{O}_X(1)\xrightarrow{\sim}\mathcal{I}\mathcal{O}_X$, et comme $S$
est engendrée par $S_1=\mathfrak{I}$, $\mathcal{O}_X(1)$ est inversible
(\hyperref[II.3.2.5-fr]{3.2.5}) et très ample
(\hyperref[II.4.4.3-fr]{4.4.3} et \hyperref[II.4.4.9-fr]{4.4.9}), d'où notre
assertion.

\subsection{Résultats préliminaires sur la localisation dans les anneaux gradués}
\label{subsection:II.8.2-fr}

\begin{env}[8.2.1]
\label{II.8.2.1-fr}
Soit $S$ un anneau gradué, où nous ne supposons pas pour le moment les degrés
positifs. On posera
\begin{equation}
\label{II.8.2.1.1-fr}
  S^{\geq}=\bigoplus_{n\geq0}S_n,
  \qquad
  S^{\leq}=\bigoplus_{n\leq0}S_n,
\tag{8.2.1.1}
\end{equation}
qui sont des sous-anneaux gradués de $S$, à degrés respectivement tous
positifs et tous négatifs. Si $f$ est un élément homogène de degré $d$
(positif ou négatif) de $S$, l'anneau des fractions $S_f=S'$ est encore muni
d'une structure d'anneau gradué, en prenant pour $S'_n$ ($n\in\mathbf{Z}$)
l'ensemble des $x/f^k$, où $x\in S_{n+kd}$ ($k\geq0$)~; on posera encore
$S_{(f)}=S'_0$, et on écrira $S_f^{\geq}$ et $S_f^{\leq}$ pour $S'^{\geq}$ et
$S'^{\leq}$ respectivement. Si $d>0$, on a
\begin{equation}
\label{II.8.2.1.2-fr}
  (S^{\geq})_f=S_f
\tag{8.2.1.2}
\end{equation}
puisque, si $x\in S_{n+kd}$, avec $n+kd<0$, on peut écrire
$x/f^k=xf^h/f^{h+k}$ et que $n+(h+k)d>0$ pour $h$ assez grand et $>0$. On en
conclut par définition que
\begin{equation}
\label{II.8.2.1.3-fr}
  (S^{\geq})_{(f)}=(S_f^{\geq})_0=S_{(f)}.
\tag{8.2.1.3}
\end{equation}

Si $M$ est un $S$-module gradué, on posera de même
\begin{equation}
\label{II.8.2.1.4-fr}
  M^{\geq}=\bigoplus_{n\geq0}M_n,
  \qquad
  M^{\leq}=\bigoplus_{n\leq0}M_n
\tag{8.2.1.4}
\end{equation}
qui sont respectivement un $S^{\geq}$-module gradué et un $S^{\leq}$-module
gradué et ont pour intersection le $S_0$-module $M_0$. Si $f\in S_d$, on
définit encore $M_f$ comme $S_f$-module gradué en prenant comme éléments de
degré $n$ les $z/f^k$, où $z\in M_{n+kd}$ ($k\geq0$)~; on désignera par
$M_{(f)}$ l'ensemble des éléments de degré $0$ de $M_f$, qui est un
$S_{(f)}$-module, et on écrira $M_f^{\geq}$ et $M_f^{\leq}$ au lieu de
$(M_f)^{\geq}$ et $(M_f)^{\leq}$ respectivement. Si $d>0$, on voit comme
ci-dessus que
\begin{equation}
\label{II.8.2.1.5-fr}
  (M^{\geq})_f=M_f
\tag{8.2.1.5}
\end{equation}
et
\begin{equation}
\label{II.8.2.1.6-fr}
  (M^{\geq})_{(f)}=(M_f^{\geq})_0=M_{(f)}.
\tag{8.2.1.6}
\end{equation}
\end{env}

\begin{env}[8.2.2]
\label{II.8.2.2-fr}
Soit $z$ une indéterminée, que nous nommerons \emph{variable
d'homogénéisation}. Si $S$ est un anneau gradué (à degrés positifs ou
négatifs), l'algèbre de polynômes\footnote{Il ne saurait ici y avoir de
confusion avec l'usage de la notation $\widehat{S}$ pour désigner le séparé
complété d'un anneau.}
\begin{equation}
\label{II.8.2.2.1-fr}
  \widehat{S}=S[z]
\tag{8.2.2.1}
\end{equation}

\oldpage[II]{158}
est une $S$-algèbre graduée, quand on prend pour degré de $fz^n$ ($n\geq0$) où
$f$ est homogène,
\begin{equation}
\label{II.8.2.2.2-fr}
  \deg(fz^n)=n+\deg f.
\tag{8.2.2.2}
\end{equation}
\end{env}

\begin{lemma}[8.2.3]
\label{II.8.2.3-fr}
\begin{enumerate}
\item[\rm(i)] On a des isomorphismes canoniques d'anneaux (non gradués)
\begin{equation}
\label{II.8.2.3.1-fr}
  \widehat{S}_{(z)}\xrightarrow{\sim}
  \widehat{S}/(z-1)\widehat{S}\xrightarrow{\sim}S.
\tag{8.2.3.1}
\end{equation}
\item[\rm(ii)] On a un isomorphisme canonique d'anneaux (non gradués)
\begin{equation}
\label{II.8.2.3.2-fr}
  \widehat{S}_{(f)}\xrightarrow{\sim}S_f^{\leq}
\tag{8.2.3.2}
\end{equation}
pour tout $f\in S_d$, avec $d>0$.
\end{enumerate}
\end{lemma}

Le premier des isomorphismes de (\hyperref[II.8.2.3.1-fr]{8.2.3.1}) a été
défini dans (\hyperref[II.2.2.5-fr]{2.2.5}) et le second est trivial~;
l'isomorphisme $\widehat{S}_{(z)}\xrightarrow{\sim}S$ ainsi défini fait donc
correspondre à $xz^n/z^{n+k}$, où $\deg(x)=k$ ($k\geq-n$) l'élément $x$.
L'homomorphisme (\hyperref[II.8.2.3.2-fr]{8.2.3.2}) fait correspondre à
$xz^n/f^k$, où $\deg(x)=kd-n$, l'élément $x/f^k$, de degré $-n$ dans
$S_f^{\leq}$, et il est encore évident qu'on a bien ici un isomorphisme.

\begin{env}[8.2.4]
\label{II.8.2.4-fr}
Soit $M$ un $S$-module gradué. Il est clair que le $S$-module
\begin{equation}
\label{II.8.2.4.1-fr}
  \widehat{M}=M\otimes_S\widehat{S}=M\otimes_S S[z]
\tag{8.2.4.1}
\end{equation}
est somme directe des $S$-modules $M\otimes_S Sz^n$, donc des groupes abéliens
$M_k\otimes_S Sz^n$ ($k\in\mathbf{Z}$, $n\geq0$)~; on définit sur
$\widehat{M}$ une structure de $\widehat{S}$-module gradué en prenant
\begin{equation}
\label{II.8.2.4.2-fr}
  \deg(x\otimes z^n)=n+\deg x
\tag{8.2.4.2}
\end{equation}
pour tout $x$ homogène dans $M$. Nous laissons au lecteur le soin de démontrer
l'analogue de (\hyperref[II.8.2.3-fr]{8.2.3})~:
\end{env}

\begin{lemma}[8.2.5]
\label{II.8.2.5-fr}
\begin{enumerate}
\item[\rm(i)] On a un di-isomorphisme canonique de modules (non gradués)
\begin{equation}
\label{II.8.2.5.1-fr}
  \widehat{M}_{(z)}\xrightarrow{\sim}M.
\tag{8.2.5.1}
\end{equation}
\item[\rm(ii)] Pour tout $f\in S_d$ ($d>0$), on a un di-isomorphisme de
modules (non gradués)
\begin{equation}
\label{II.8.2.5.2-fr}
  \widehat{M}_{(f)}\xrightarrow{\sim}M_f^{\leq}.
\tag{8.2.5.2}
\end{equation}
\end{enumerate}
\end{lemma}

\begin{env}[8.2.6]
\label{II.8.2.6-fr}
Soit $S$ un anneau gradué à degrés \emph{positifs}, et considérons dans $S$
la suite décroissante d'idéaux gradués
\begin{equation}
\label{II.8.2.6.1-fr}
  S_{[n]}=\bigoplus_{m\geq n}S_m\qquad(n\geq0)
\tag{8.2.6.1}
\end{equation}
(en particulier on a $S_{[0]}=S$, $S_{[1]}=S_+$). Comme il est clair que
$S_{[m]}S_{[n]}\subset S_{[m+n]}$, on peut définir un \emph{anneau gradué}
$S^\natural$ en prenant
\begin{equation}
\label{II.8.2.6.2-fr}
  S^\natural=\bigoplus_{n\geq0}S_n^\natural
  \quad\text{avec}\quad S_n^\natural=S_{[n]}.
\tag{8.2.6.2}
\end{equation}
$S_0^\natural$ est donc l'anneau $S$ considéré comme anneau \emph{non gradué},
et $S^\natural$ est par suite une $S_0^\natural$-algèbre. Pour tout élément
homogène $f\in S_d$ ($d>0$), nous désignerons par $f^\natural$ l'élément
$f$ considéré comme appartenant à $S_{[d]}=S_d^\natural$. Avec ces
notations~:
\end{env}

\oldpage[II]{159}

\begin{lemma}[8.2.7]
\label{II.8.2.7-fr}
Soient $S$ un anneau gradué à degrés positifs, $f$ un élément homogène de
$S_d$ ($d>0$). On a des isomorphismes canoniques d'anneaux
\begin{equation}
\label{II.8.2.7.1-fr}
  S_f\xrightarrow{\sim}\bigoplus_{n\in\mathbf{Z}}S(n)_{(f)}
\tag{8.2.7.1}
\end{equation}
\begin{equation}
\label{II.8.2.7.2-fr}
  (S_f^{\geq})_{f/1}\xrightarrow{\sim}S_f
\tag{8.2.7.2}
\end{equation}
\begin{equation}
\label{II.8.2.7.3-fr}
  S^\natural_{(f^\natural)}\xrightarrow{\sim}S_f^{\geq}
\tag{8.2.7.3}
\end{equation}
dont les deux premiers sont des isomorphismes d'anneaux gradués.
\end{lemma}

Il est immédiat par définition que l'on a
$(S_f)_n=(S(n)_f)_0$, d'où l'isomorphisme
(\hyperref[II.8.2.7.1-fr]{8.2.7.1}), qui n'est autre que l'identité. D'autre
part, comme $f/1$ est inversible dans $S_f$, il y a un isomorphisme canonique
$S_f\xrightarrow{\sim}(S_f^{\geq})_{f/1}=(S_f)_{f/1}$ en vertu de
(\hyperref[II.8.2.1.2-fr]{8.2.1.2}) appliqué à $S_f$~; l'isomorphisme
réciproque est par définition l'isomorphisme
(\hyperref[II.8.2.7.2-fr]{8.2.7.2}). Enfin, si
$x=\sum_{m\geq n}y_m$ est un élément de $S_{[n]}$, avec $n=kd$, on fait
correspondre à l'élément $x/(f^\natural)^k$ l'élément
$\sum_m y_m/f^k$ de $S_f^{\geq}$, et on vérifie aussitôt que cela définit un
isomorphisme (\hyperref[II.8.2.7.3-fr]{8.2.7.3}).

\begin{env}[8.2.8]
\label{II.8.2.8-fr}
Si $M$ est un $S$-module gradué, on pose de même pour tout $n\in\mathbf{Z}$
\begin{equation}
\label{II.8.2.8.1-fr}
  M_{[n]}=\bigoplus_{m\geq n}M_m
\tag{8.2.8.1}
\end{equation}
et comme $S_{[m]}M_{[n]}\subset M_{[m+n]}$ ($m\geq0$), on peut définir un
$S^\natural$-module gradué $M^\natural$ en prenant
\begin{equation}
\label{II.8.2.8.2-fr}
  M^\natural=\bigoplus_{n\in\mathbf{Z}}M_n^\natural,
  \quad\text{avec}\quad M_n^\natural=M_{[n]}.
\tag{8.2.8.2}
\end{equation}
\end{env}

Nous laissons encore au lecteur la démonstration du~:

\begin{lemma}[8.2.9]
\label{II.8.2.9-fr}
Avec les notations de (\hyperref[II.8.2.7-fr]{8.2.7}) et
(\hyperref[II.8.2.8-fr]{8.2.8}), on a les di-isomorphismes canoniques de
modules
\begin{equation}
\label{II.8.2.9.1-fr}
  M_f\xrightarrow{\sim}\bigoplus_{n\in\mathbf{Z}}M(n)_{(f)}
\tag{8.2.9.1}
\end{equation}
\begin{equation}
\label{II.8.2.9.2-fr}
  (M_f^{\geq})_{f/1}\xrightarrow{\sim}M_f
\tag{8.2.9.2}
\end{equation}
\begin{equation}
\label{II.8.2.9.3-fr}
  M^\natural_{(f^\natural)}\xrightarrow{\sim}M_f^{\geq}
\tag{8.2.9.3}
\end{equation}
dont les deux premiers sont des di-isomorphismes de modules gradués.
\end{lemma}

\begin{lemma}[8.2.10]
\label{II.8.2.10-fr}
Soit $S$ un anneau gradué à degrés positifs.
\begin{enumerate}
\item[\rm(i)] Pour que $S^\natural$ soit une $S_0^\natural$-algèbre de type
fini (resp. noethérienne) il faut et il suffit que $S$ soit une $S_0$-algèbre
de type fini (resp. noethérienne).
\item[\rm(ii)] Pour que $S_{n+1}^\natural=S_1^\natural S_n^\natural$ pour
$n\geq n_0$, il faut et il suffit que $S_{n+1}=S_1S_n$ pour $n\geq n_0$.
\item[\rm(iii)] Pour que $S_n^\natural=S_1^{\natural n}$ pour $n\geq n_0$, il
faut et il suffit que $S_n=S_1^n$ pour $n\geq n_0$.
\item[\rm(iv)] Si $(f_\alpha)$ est un ensemble d'éléments homogènes de $S_+$
telle que $S_+$ soit la racine dans $S_+$ de l'idéal de $S_+$ engendré par
les $f_\alpha$, alors $S_+^\natural$ est la racine dans $S_+^\natural$ de
l'idéal de $S_+^\natural$ engendré par les $f_\alpha^\natural$.
\end{enumerate}
\end{lemma}

\begin{enumerate}
\item[\rm(i)] Si $S^\natural$ est une $S_0^\natural$-algèbre de type fini,
$S_+=S_1^\natural$ est un module de type fini sur $S=S_0^\natural$ par
(\hyperref[II.2.1.6-fr]{2.1.6}, (i)), donc $S$ est une $S_0$-algèbre de type
fini (\hyperref[II.2.1.4-fr]{2.1.4})~; si $S^\natural$ est un anneau
noethérien, il en est de même de $S_0^\natural=S$
(\hyperref[II.2.1.5-fr]{2.1.5}). Inversement, si $S$ est une $S_0$-algèbre

\oldpage[II]{160}
de type fini, on sait (\hyperref[II.2.1.6-fr]{2.1.6}, (ii)) qu'il existe
$h>0$ et $m_0>0$ tels que $S_{n+h}=S_hS_n$ pour $n\geq m_0$~; on peut
évidemment supposer $m_0\geq h$. En outre, les $S_m$ sont des $S_0$-modules
de type fini (\hyperref[II.2.1.6-fr]{2.1.6}, (i)). Cela étant, si
$n\geq m_0+h$, on a $S_n^\natural=S_hS_{n-h}^\natural
=S_h^\natural S_{n-h}^\natural$~; et si $m<m_0+h$, on a, en posant
$E=S_{m_0}+\cdots+S_{m_0+h-1}$,
\[
  S_m^\natural=S_m+\cdots+S_{m_0+h-1}+S_hE+S_h^2E+\cdots.
\]
Pour $1\leq m\leq m_0$, soit $G_m$ la réunion de systèmes finis de
générateurs des $S_0$-modules $S_i$, pour $m\leq i\leq m_0+h-1$,
considérée comme partie de $S_{[m]}$. Pour $m_0+1\leq m\leq m_0+h-1$, soit
de même $G_m$ la réunion de systèmes finis de générateurs des $S_0$-modules
$S_i$, pour $m\leq i\leq m_0+h-1$, et de $S_hE$, considérée comme partie de
$S_{[m]}$. Il est clair que l'on a $S_m^\natural=S_0^\natural G_m$ pour
$1\leq m\leq m_0+h-1$, et par suite la réunion $G$ des $G_m$ pour
$1\leq m\leq m_0+h-1$ est un système de générateurs de la
$S_0^\natural$-algèbre $S^\natural$. On en conclut que si $S=S_0^\natural$
est un anneau noethérien, il en est de même de $S^\natural$.

\item[\rm(ii)] Il est clair que si $S_{n+1}=S_1S_n$ pour $n\geq n_0$, on a
$S_{n+1}^\natural=S_1S_n^\natural$, et \emph{a fortiori}
$S_{n+1}^\natural=S_1^\natural S_n^\natural$ pour $n\geq n_0$.
Inversement, cette dernière relation s'écrit
\[
  S_{n+1}+S_{n+2}+\cdots=(S_1+S_2+\cdots)(S_n+S_{n+1}+\cdots)
\]
et la comparaison des termes de degré $n+1$ (dans $S$) aux deux membres donne
$S_{n+1}=S_1S_n$.

\item[\rm(iii)] Si $S_n=S_1^n$ pour $n\geq n_0$, on a
$S_n^\natural=S_1^n+S_1^{n+1}+\cdots$~; comme $S_1^\natural$ contient
$S_1+S_1^2+\cdots$, on a $S_n^\natural\subset S_1^{\natural n}$, et par suite
$S_n^\natural=S_1^{\natural n}$ pour $n\geq n_0$. Inversement, les seuls
termes de $S_1^{\natural n}=(S_1+S_2+\cdots)^n$ qui soient de degré $n$ dans
$S$ sont ceux de $S_1^n$~; la relation $S_n^\natural=S_1^{\natural n}$
implique donc $S_n=S_1^n$.

\item[\rm(iv)] Il suffit de prouver que si un élément $g\in S_{k+h}$ est
considéré comme un élément de $S_k^\natural$ ($k>0$, $h\geq0$), alors il
existe un entier $n>0$ tel que dans $S_{kn}^\natural$, $g^n$ soit combinaison
linéaire des $f_\alpha^\natural$ à coefficients dans $S^\natural$. Par
hypothèse, il y a un entier $m_0$ tel que pour $m\geq m_0$, on ait, \emph{dans
$S$}, $g^m=\sum_\alpha c_{\alpha m}f_\alpha$, les indices $\alpha$ figurant
dans cette formule étant \emph{indépendants de $m$}~; en outre, on peut
évidemment supposer les $c_{\alpha m}$ homogènes, avec
\[
  \deg(c_{\alpha m})=m(k+h)-\deg f_\alpha
\]
dans $S$. Prenons alors $m_0$ assez grand pour que l'on ait
$km_0>\deg f_\alpha$ pour tous les $f_\alpha$ qui figurent dans $g^{m_0}$~;
pour tout $\alpha$, soit $c'_{\alpha m}$ l'élément $c_{\alpha m}$ considéré
comme de degré $km-\deg(f_\alpha)$ dans $S^\natural$~; on a alors, dans
$S^\natural$, $g^m=\sum_\alpha c'_{\alpha m}f_\alpha^\natural$, ce qui
termine la démonstration.
\end{enumerate}

\begin{env}[8.2.11]
\label{II.8.2.11-fr}
Considérons la $S_0$-algèbre graduée
\begin{equation}
\label{II.8.2.11.1-fr}
  S^\natural\otimes_S S_0=S^\natural/S_+S^\natural
  =\bigoplus_{n\geq0}S_{[n]}/S_+S_{[n]}.
\tag{8.2.11.1}
\end{equation}
Comme $S_n$ est un $S_0$-module quotient de $S_{[n]}/S_+S_{[n]}$, on a un
homomorphisme canonique de $S_0$-algèbres graduées
\begin{equation}
\label{II.8.2.11.2-fr}
  S^\natural\otimes_S S_0\to S
\tag{8.2.11.2}
\end{equation}
qui est évidemment \emph{surjectif}, et correspond par suite
(\hyperref[II.2.9.2-fr]{2.9.2}) à une \emph{immersion fermée} canonique
\begin{equation}
\label{II.8.2.11.3-fr}
  \operatorname{Proj}(S)\to
  \operatorname{Proj}(S^\natural\otimes_S S_0).
\tag{8.2.11.3}
\end{equation}
\end{env}

\oldpage[II]{161}

\begin{proposition}[8.2.12]
\label{II.8.2.12-fr}
Le morphisme canonique (\hyperref[II.8.2.11.3-fr]{8.2.11.3}) est bijectif.
Pour que l'homomorphisme (\hyperref[II.8.2.11.2-fr]{8.2.11.2}) soit
(TN)-bijectif, il faut et il suffit qu'il existe $n_0$ tel que
$S_{n+1}=S_1S_n$ pour $n\geq n_0$. Si cette dernière condition est vérifiée,
(\hyperref[II.8.2.11.3-fr]{8.2.11.3}) est un isomorphisme~; la réciproque est
vraie lorsque $S$ est noethérien.
\end{proposition}

Pour démontrer la première assertion, il suffit
(\hyperref[II.2.8.3-fr]{2.8.3}) de prouver que le noyau $\mathfrak{I}$ de
l'homomorphisme (\hyperref[II.8.2.11.2-fr]{8.2.11.2}) est formé d'éléments
\emph{nilpotents}. Or si $f\in S_{[n]}$ est un élément dont la classe mod.
$S_+S_{[n]}$ appartient à ce noyau, cela signifie que $f\in S_{[n+1]}$~;
l'élément $f^{n+1}$, considéré comme appartenant à $S_{[n(n+1)]}$, appartient
alors à $S_+S_{[n(n+1)]}$, puisqu'il s'écrit $f\cdot f^n$~; donc la classe de
$f^{n+1}$ mod. $S_+S_{[n(n+1)]}$ est nulle, ce qui prouve notre assertion.
Comme l'hypothèse $S_{n+1}=S_1S_n$ pour $n\geq n_0$ équivaut à
$S_{n+1}^{\natural}=S_1^{\natural}S_n^{\natural}$ pour $n\geq n_0$
(\hyperref[II.8.2.10-fr]{8.2.10}, (ii)), cette hypothèse équivaut par définition
au fait que (\hyperref[II.8.2.11.2-fr]{8.2.11.2}) est (TN)-injectif, donc
(TN)-bijectif, et alors (\hyperref[II.8.2.11.3-fr]{8.2.11.3}) est un
isomorphisme en vertu de (\hyperref[II.2.9.1-fr]{2.9.1}). Inversement, si
(\hyperref[II.8.2.11.3-fr]{8.2.11.3}) est un isomorphisme, le faisceau
$\widetilde{\mathfrak{I}}$ sur
$\operatorname{Proj}(S^{\natural}\otimes_S S_0)$ est nul
(\hyperref[II.2.9.2-fr]{2.9.2}, (i))~; comme
$S^{\natural}\otimes_S S_0$ est noethérien comme quotient de $S^{\natural}$
(\hyperref[II.8.2.10-fr]{8.2.10}, (i)), on conclut de
(\hyperref[II.2.7.3-fr]{2.7.3}) que $\mathfrak{I}$ vérifie la condition (TN),
donc $S_{n+1}^{\natural}=S_1^{\natural}S_n^{\natural}$ pour $n\geq n_0$, et
cela achève la démonstration en vertu de
(\hyperref[II.8.2.10-fr]{8.2.10}, (ii)).

\begin{env}[8.2.13]
\label{II.8.2.13-fr}
Considérons maintenant les injections canoniques $(S_+)^n\to S_{[n]}$, qui
définissent un homomorphisme injectif de degré $0$ d'anneaux gradués
\begin{equation}
\label{II.8.2.13.1-fr}
  \bigoplus_{n\geq0}(S_+)^n\to S^{\natural}.
\tag{8.2.13.1}
\end{equation}
\end{env}

\begin{proposition}[8.2.14]
\label{II.8.2.14-fr}
Pour que l'homomorphisme (\hyperref[II.8.2.13.1-fr]{8.2.13.1}) soit un
(TN)-isomorphisme, il faut et il suffit qu'il existe $n_0$ tel que
$S_n=S_1^n$ pour tout $n\geq n_0$. Lorsqu'il en est ainsi, le morphisme
correspondant à (\hyperref[II.8.2.13.1-fr]{8.2.13.1}) est partout défini et
est un isomorphisme
\[
  \operatorname{Proj}(S^{\natural})\xrightarrow{\sim}
  \operatorname{Proj}\left(\bigoplus_{n\geq0}(S_+)^n\right)~;
\]
la réciproque est vraie lorsque $S$ est noethérien.
\end{proposition}

Les deux premières assertions sont évidentes, vu
(\hyperref[II.8.2.10-fr]{8.2.10}, (iii)) et
(\hyperref[II.2.9.1-fr]{2.9.1}). La troisième résultera de
(\hyperref[II.8.2.10-fr]{8.2.10}, (i) et (iii)) et du lemme suivant~:

\begin{lemma}[8.2.14.1]
\label{II.8.2.14.1-fr}
Soit $T$ un anneau gradué en degrés positifs qui est une $T_0$-algèbre de type
fini. Si le morphisme correspondant à l'homomorphisme injectif
$\bigoplus_{n\geq0}T_1^n\to T$ est partout défini et est un isomorphisme
\[
  \operatorname{Proj}(T)\to
  \operatorname{Proj}\left(\bigoplus_{n\geq0}T_1^n\right),
\]
il existe $n_0$ tel que $T_n=T_1^n$ pour $n\geq n_0$.
\end{lemma}

En effet, soient $g_i$ ($1\leq i\leq r$) des générateurs du $T_0$-module
$T_1$. L'hypothèse entraîne d'abord que les $D_+(g_i)$ recouvrent
$\operatorname{Proj}(T)$ (\hyperref[II.2.8.1-fr]{2.8.1}). Soit
$(h_j)_{1\leq j\leq s}$ un système d'éléments homogènes de $T_+$, avec
$\deg(h_j)=n_j$, qui forment avec les $g_i$ un système de générateurs de
l'idéal $T_+$, ou (ce qui revient au même
(\hyperref[II.2.1.3-fr]{2.1.3})) un système de générateurs de $T$ en tant que
$T_0$-algèbre~; si on pose $T'=\bigoplus_{n\geq0}T_1^n$, l'élément
$h_j/g_i^{n_j}$ de l'anneau $T_{(g_i)}$ doit par hypothèse appartenir au
sous-anneau $T'_{(g_i)}$, donc il existe un entier $k$ tel que
$T_1^k h_j\subset T_1^{k+n_j}$ pour tout $j$. On en conclut par récurrence
sur $r$ que $T_1^k h_j^r\subset T'$ pour tout $r\geq1$, et par définition des
$h_j$, on a donc $T_1^kT\subset T'$. Il y a d'autre part pour tout $j$ un
entier $m_j$ tel que $h_j^{m_j}$ appartienne à l'idéal de $T$ engendré par
les $g_i$ (\hyperref[II.2.3.14-fr]{2.3.14}), donc $h_j^{m_j}\in T_1T$, et

\oldpage[II]{162}
$h_j^{m_jk}\in T_1^kT\subset T'$. Il y a par suite un entier $m_0\geq k$ tel
que $h_j^m\in T_1^{mn_j}$ pour $m\geq m_0$. Cela étant, si $q$ est le plus
grand des entiers $n_j$, le nombre $n_0=qsm_0+k$ répond à la question. En
effet, un élément de $T_n$, pour $n\geq n_0$, est somme de monômes appartenant
à $T_1^\alpha u$, où $u$ est un produit de puissances des $h_j$~; si
$\alpha\geq k$, il résulte de ce qui précède que $T_1^\alpha u\subset T_1^n$~;
dans le cas contraire, un des exposants des $h_j$ est $\geq m_0$, donc
$u\in T_1^\beta v$, où $\beta\geq k$ et $v$ est encore un produit de
puissances des $h_j$~; on est alors ramené au cas précédent, et on voit
finalement que $T_1^\alpha u\subset T_1^n$ dans tous les cas.

\begin{remark}[8.2.15]
\label{II.8.2.15-fr}
La condition $S_n=S_1^n$ pour $n\geq n_0$ entraîne évidemment
$S_{n+1}=S_1S_n$ pour $n\geq n_0$, mais la réciproque est inexacte, même
quand on suppose $S$ noethérien. Par exemple, soient $K$ un corps,
$A=K[x]$, $B=K[y]/y^2K[y]$, où $x$ et $y$ sont deux indéterminées, $x$ étant
prise de degré $1$ et $y$ de degré $2$, et soit $S=A\otimes_KB$, de sorte que
$S$ est une algèbre graduée sur $K$ ayant une base formée des éléments $1$,
$x^n$ ($n\geq1$) et $x^ny$ ($n\geq0$). Il est immédiat que
$S_{n+1}=S_1S_n$ pour $n\geq2$, mais $S_1^n=Kx^n$ tandis que
$S_n=Kx^n+Kx^{n-2}y$ pour $n\geq2$.
\end{remark}

\subsection{Cônes projetants}
\label{subsection:II.8.3-fr}

\begin{env}[8.3.1]
\label{II.8.3.1-fr}
Soit $Y$ un préschéma~; dans tout ce numéro, il ne sera question que de
\emph{$Y$-préschémas} et de \emph{$Y$-morphismes}. Soit $\mathcal{S}$ une
$\mathcal{O}_Y$-Algèbre quasi-cohérente graduée à \emph{degrés positifs}~;
\emph{nous supposerons en outre que l'on a $\mathcal{S}_0=\mathcal{O}_Y$}.
Conformément aux notations introduites en
(\hyperref[II.8.2.2-fr]{8.2.2}), nous poserons
\begin{equation}
\label{II.8.3.1.1-fr}
  \widehat{\mathcal{S}}=\mathcal{S}[z]
  =\mathcal{S}\otimes_{\mathcal{O}_Y}\mathcal{O}_Y[z]
\tag{8.3.1.1}
\end{equation}
que l'on considère comme $\mathcal{O}_Y$-Algèbre graduée à degrés positifs en
définissant les degrés par la formule
(\hyperref[II.8.2.2.2-fr]{8.2.2.2}), de sorte que pour tout ouvert affine $U$
de $Y$, on a
\[
  \Gamma(U,\widehat{\mathcal{S}})=(\Gamma(U,\mathcal{S}))[z].
\]
Dans ce qui suit, nous poserons
\begin{equation}
\label{II.8.3.1.2-fr}
  X=\operatorname{Proj}(\mathcal{S}),\qquad
  C=\operatorname{Spec}(\mathcal{S}),\qquad
  \widehat{C}=\operatorname{Proj}(\widehat{\mathcal{S}})
\tag{8.3.1.2}
\end{equation}
(où dans la définition de $C$, $\mathcal{S}$ est considérée comme
$\mathcal{O}_Y$-Algèbre non graduée), et nous dirons que $C$ (resp.
$\widehat{C}$) est le \emph{cône affine} (resp. \emph{cône projectif}) défini
par $\mathcal{S}$~; on dira aussi «~cône~» au lieu de «~cône affine~». Par
abus de langage, nous dirons encore que $C$ (resp. $\widehat{C}$) est le
\emph{cône projetant affine de $X$} (resp. le \emph{cône projetant projectif
de $X$}), étant sous-entendu que le préschéma $X$ est donné sous la forme
$\operatorname{Proj}(\mathcal{S})$~; enfin, nous dirons que $\widehat{C}$ est
la \emph{fermeture projective} de $C$ (étant entendu que la donnée de
$\mathcal{S}$ figure dans la structure de $C$).
\end{env}

\begin{proposition}[8.3.2]
\label{II.8.3.2-fr}
Il existe des $Y$-morphismes canoniques
\begin{equation}
\label{II.8.3.2.1-fr}
  Y\xrightarrow{\varepsilon}C\xrightarrow{i}\widehat{C}
\tag{8.3.2.1}
\end{equation}
\begin{equation}
\label{II.8.3.2.2-fr}
  X\xrightarrow{j}\widehat{C}
\tag{8.3.2.2}
\end{equation}

\oldpage[II]{163}
tels que $\varepsilon$ et $j$ soient des immersions fermées, et $i$ un
morphisme affine, qui est une immersion ouverte dominante, pour laquelle
\begin{equation}
\label{II.8.3.2.3-fr}
  i(C)=\widehat{C}-j(X)~;
\tag{8.3.2.3}
\end{equation}
en outre $\widehat{C}$ est le plus petit sous-préschéma fermé de $\widehat{C}$
majorant $i(C)$.
\end{proposition}

Pour définir $i$, on considère l'ouvert de $\widehat{C}$,
\begin{equation}
\label{II.8.3.2.4-fr}
  \widehat{C}_z=
  \operatorname{Spec}(\widehat{\mathcal{S}}/(z-1)\widehat{\mathcal{S}})
\tag{8.3.2.4}
\end{equation}
(\hyperref[II.3.1.4-fr]{3.1.4}), $z$ s'identifiant canoniquement à une
section de $\widehat{\mathcal{S}}$ au-dessus de $Y$. L'isomorphisme
$i:C\xrightarrow{\sim}\widehat{C}_z$ correspond alors à l'isomorphisme
canonique (\hyperref[II.8.2.3.1-fr]{8.2.3.1})
\[
  \widehat{\mathcal{S}}/(z-1)\widehat{\mathcal{S}}
  \xrightarrow{\sim}\mathcal{S}.
\]
Le morphisme $\varepsilon$ correspond à l'homomorphisme d'augmentation
$\mathcal{S}\to\mathcal{S}_0=\mathcal{O}_Y$ ayant pour noyau
$\mathcal{S}_+$ (\hyperref[II.1.2.7-fr]{1.2.7}) et comme ce dernier est
surjectif, $\varepsilon$ est une immersion fermée
(\hyperref[II.1.4.10-fr]{1.4.10}). Enfin, $j$ correspond de même
(\hyperref[II.3.5.1-fr]{3.5.1}) à l'homomorphisme surjectif de degré $0$,
$\widehat{\mathcal{S}}\to\mathcal{S}$, qui se réduit à l'identité dans
$\mathcal{S}$ et est nul dans $z\widehat{\mathcal{S}}$, qui est son noyau~;
$j$ est partout défini et est une immersion fermée en vertu de
(\hyperref[II.3.6.2-fr]{3.6.2}).

Pour démontrer les autres assertions de
(\hyperref[II.8.3.2-fr]{8.3.2}), on peut évidemment se borner au cas où
$Y=\operatorname{Spec}(A)$ est affine, $\mathcal{S}=\widetilde{S}$, où $S$
est une $A$-algèbre graduée, d'où
$\widehat{\mathcal{S}}=(\widehat{S})^{\sim}$~; les éléments homogènes $f$ de
$S_+$ s'identifient alors à des sections de $\widehat{\mathcal{S}}$ au-dessus
de $Y$, et l'ouvert de $\widehat{C}$ noté $D_+(f)$ dans
(\hyperref[II.2.3.3-fr]{2.3.3}) s'écrit aussi $\widehat{C}_f$
(\hyperref[II.3.1.4-fr]{3.1.4})~; de même l'ouvert de $C$ noté $D(f)$ dans
(\hyperref[I.1.1.1-fr]{I, 1.1.1}) s'écrit aussi $C_f$
(\hyperref[0.5.5.2-fr]{0, 5.5.2}). Cela étant, il résulte de
(\hyperref[II.2.3.14-fr]{2.3.14}) et de la définition de $\widehat{S}$ que
dans le cas présent, les ouverts $\widehat{C}_z=i(C)$ et $\widehat{C}_f$
($f$ homogène dans $S_+$) constituent un \emph{recouvrement} de
$\widehat{C}$. En outre, on a, avec ces notations,
\begin{equation}
\label{II.8.3.2.5-fr}
  i^{-1}(\widehat{C}_f)=C_f~;
\tag{8.3.2.5}
\end{equation}
en effet,
$\widehat{C}_f\cap i(C)=\widehat{C}_f\cap\widehat{C}_z
=\widehat{C}_{fz}=\operatorname{Spec}(\widehat{S}_{(fz)})$. Or, si
$d=\deg(f)$, $\widehat{S}_{(fz)}$ est canoniquement isomorphe à
$(\widehat{S}_{(z)})_{f/z^d}$ (\hyperref[II.2.2.2-fr]{2.2.2}), et il résulte
de la définition de l'isomorphisme
(\hyperref[II.8.2.3.1-fr]{8.2.3.1}) que l'image de
$(\widehat{S}_{(z)})_{f/z^d}$ par l'isomorphisme correspondant des anneaux de
fractions est exactement $S_f$. Comme $C_f=\operatorname{Spec}(S_f)$, cela
prouve (\hyperref[II.8.3.2.5-fr]{8.3.2.5}) et montre en même temps que le
morphisme $i$ est affine~; en outre, la restriction de $i$ à $C_f$,
considérée comme morphisme dans $\widehat{C}_f$, correspond
(\hyperref[I.1.7.3-fr]{I, 1.7.3}) à l'homomorphisme canonique
$\widehat{S}_{(f)}\to\widehat{S}_{(fz)}$ et en vertu de ce qui précède et de
(\hyperref[II.8.2.3.2-fr]{8.2.3.2}), on peut énoncer le résultat suivant~:

\begin{env}[8.3.2.6]
\label{II.8.3.2.6-fr}
Si $Y=\operatorname{Spec}(A)$ est affine, et $\mathcal{S}=\widetilde{S}$,
alors, pour tout $f$ homogène dans $S_+$, $\widehat{C}_f$ s'identifie
canoniquement à $\operatorname{Spec}(S_f^{\leq})$, et le morphisme
$C_f\to\widehat{C}_f$ restriction de $i$ correspond alors à l'injection
canonique $S_f^{\leq}\to S_f$.
\end{env}

Notons maintenant que (pour $Y$ affine) le complémentaire de
$\widehat{C}_z$ dans $\widehat{C}=\operatorname{Proj}(\widehat{S})$

\oldpage[II]{164}
est, par définition, l'ensemble des idéaux premiers gradués de $\widehat{S}$
contenant $z$, c'est-à-dire $j(X)$ par définition de $j$, ce qui démontre
(\hyperref[II.8.3.2.3-fr]{8.3.2.3}).

Pour prouver enfin la dernière assertion de
(\hyperref[II.8.3.2-fr]{8.3.2}), on peut toujours supposer $Y$ affine. Avec
les notations précédentes, remarquons que dans l'anneau $\widehat{S}$, $z$
n'est pas diviseur de $0$~; comme $\overline{i(C)}=\widehat{C}$, il suffira
de prouver le

\begin{lemma}[8.3.2.7]
\label{II.8.3.2.7-fr}
Soient $T$ un anneau gradué à degrés positifs,
$Z=\operatorname{Proj}(T)$, $g$ un élément homogène de $T$ de degré $d>0$.
Si $g$ n'est pas diviseur de $0$ dans $T$, $Z$ est le plus petit
sous-préschéma fermé de $Z$ qui majore $Z_g=D_+(g)$.
\end{lemma}

En vertu de (\hyperref[I.4.1.9-fr]{I, 4.1.9}) la question est locale sur $Z$~;
pour tout élément homogène $h\in T_e$ ($e>0$), il suffit donc de prouver que
$Z_h$ est le plus petit sous-préschéma fermé de $Z_h$ qui majore $Z_{gh}$~;
il résulte des définitions et de (\hyperref[I.4.3.2-fr]{I, 4.3.2}) que cette
condition équivaut au fait que l'homomorphisme canonique
$T_{(h)}\to T_{(gh)}$ est \emph{injectif}. Or cet homomorphisme s'identifie à
l'homomorphisme canonique
$T_{(h)}\to(T_{(h)})_{g^e/h^d}$ (\hyperref[II.2.2.3-fr]{2.2.3}). Mais comme
$g^e$ n'est pas diviseur de $0$ dans $T$, $g^e/h^d$ ne l'est pas dans $T_h$
(ni \emph{a fortiori} dans $T_{(h)}$), car la relation
$(g^e/h^d)(t/h^m)=0$ avec $t\in T$ et $m>0$ entraîne l'existence d'un $n>0$
tel que $h^ng^et=0$, d'où $h^nt=0$, et par suite $t/h^m=0$ dans $T_h$. Cela
achève donc la démonstration (\hyperref[0.1.2.2-fr]{0, 1.2.2}).

\begin{env}[8.3.3]
\label{II.8.3.3-fr}
On identifiera souvent le cône affine $C$ au sous-préschéma induit par le cône
projectif $\widehat{C}$ sur l'ouvert $i(C)$, au moyen de l'immersion ouverte
$i$. Le sous-préschéma fermé de $C$, associé à l'immersion fermée
$\varepsilon$, est appelé le \emph{préschéma des sommets de $C$}~; on dit
aussi que $\varepsilon$, qui est une $Y$-section de $C$, est la \emph{section
sommet} ou la \emph{section nulle de $C$}~; on peut identifier $Y$ au
préschéma des sommets de $C$ au moyen de $\varepsilon$. D'ailleurs
$i\circ\varepsilon$ est une $Y$-section de $\widehat{C}$, donc aussi une
immersion fermée (\hyperref[I.5.4.6-fr]{I, 5.4.6}), correspondant à
l'homomorphisme canonique surjectif de degré $0$,
$\widehat{\mathcal{S}}=\mathcal{S}[z]\to\mathcal{O}_Y[z]$
(cf. (\hyperref[II.3.1.7-fr]{3.1.7})), de noyau
$\mathcal{S}_+[z]=\mathcal{S}_+\widehat{\mathcal{S}}$~; le sous-préschéma de
$\widehat{C}$ associé à cette immersion fermée est encore appelé le
\emph{préschéma des sommets de $\widehat{C}$}, et $i\circ\varepsilon$ la
\emph{section sommet de $\widehat{C}$}~; il peut s'identifier par
$i\circ\varepsilon$ à $Y$. Enfin, le sous-préschéma fermé de $\widehat{C}$
associé à $j$ est appelé le \emph{lieu à l'infini de $\widehat{C}$}, et peut
être identifié à $X$ au moyen de $j$.
\end{env}

\begin{env}[8.3.4]
\label{II.8.3.4-fr}
Les sous-préschémas de $C$ (resp. $\widehat{C}$) induits respectivement sur
les \emph{ouverts}
\begin{equation}
\label{II.8.3.4.1-fr}
  E=C-\varepsilon(Y),\qquad
  \widehat{E}=\widehat{C}-i(\varepsilon(Y))
\tag{8.3.4.1}
\end{equation}
sont appelés respectivement (par abus de langage) le \emph{cône affine
épointé} et le \emph{cône projectif épointé} définis par $\mathcal{S}$~; on
notera qu'en dépit de cette terminologie, $E$ \emph{n'est pas nécessairement
affine} sur $Y$, ni $\widehat{E}$ projectif sur $Y$
(cf. (\hyperref[II.8.4.3-fr]{8.4.3})). Lorsqu'on identifie $C$ à $i(C)$, on a
donc pour les espaces sous-jacents
\begin{equation}
\label{II.8.3.4.2-fr}
  C\cup\widehat{E}=\widehat{C},\qquad C\cap\widehat{E}=E
\tag{8.3.4.2}
\end{equation}
de sorte que $\widehat{C}$ peut être considéré comme obtenu par
\emph{recollement} des sous-préschémas ouverts $C$ et $\widehat{E}$~; en
outre, en vertu de (\hyperref[II.8.3.2.3-fr]{8.3.2.3}),
\begin{equation}
\label{II.8.3.4.3-fr}
  E=\widehat{E}-j(X).
\tag{8.3.4.3}
\end{equation}

\oldpage[II]{165}
Lorsque $Y=\operatorname{Spec}(A)$ est affine, on a, avec les notations de
(\hyperref[II.8.3.2-fr]{8.3.2})
\begin{equation}
\label{II.8.3.4.4-fr}
  E=\bigcup C_f,\qquad
  \widehat{E}=\bigcup\widehat{C}_f,\qquad
  C_f=C\cap\widehat{C}_f
\tag{8.3.4.4}
\end{equation}
$f$ parcourant l'ensemble des éléments homogènes de $S_+$ (ou seulement une
partie $M$ de cet ensemble engendrant un idéal de $S_+$ dont la racine dans
$S_+$ est $S_+$ lui-même, autrement dit, telle que les $X_f$ pour $f\in M$
recouvrent $X$ (\hyperref[II.2.3.14-fr]{2.3.14})). Le recollement de $C$ et de
$\widehat{C}_f$ le long de $C_f$ est alors déterminé par des morphismes
d'injection $C_f\to C$, $C_f\to\widehat{C}_f$, qui, ainsi qu'on l'a vu
(\hyperref[II.8.3.2.6-fr]{8.3.2.6}) correspondent respectivement aux
homomorphismes canoniques $S\to S_f$, $S_f^{\leq}\to S_f$.
\end{env}

\begin{proposition}[8.3.5]
\label{II.8.3.5-fr}
Avec les notations de (\hyperref[II.8.3.1-fr]{8.3.1}) et
(\hyperref[II.8.3.4-fr]{8.3.4}), le morphisme associé
(\hyperref[II.3.5.1-fr]{3.5.1}) à l'injection canonique
$\varphi:\mathcal{S}\to\widehat{\mathcal{S}}=\mathcal{S}[z]$ est un morphisme
affine surjectif (dit \emph{rétraction canonique})
\begin{equation}
\label{II.8.3.5.1-fr}
  p:\widehat{E}\to X
\tag{8.3.5.1}
\end{equation}
tel que l'on ait
\begin{equation}
\label{II.8.3.5.2-fr}
  p\circ j=1_X.
\tag{8.3.5.2}
\end{equation}
\end{proposition}

Pour démontrer la proposition, on peut se borner au cas où $Y$ est affine.
Compte tenu de l'expression (\hyperref[II.8.3.4.4-fr]{8.3.4.4}) de
$\widehat{E}$, le fait que le domaine de définition $G(\varphi)$ de $p$ est
égal à $\widehat{E}$ résultera de la première des assertions suivantes~:

\begin{env}[8.3.5.3]
\label{II.8.3.5.3-fr}
Si $Y=\operatorname{Spec}(A)$ est affine, et $\mathcal{S}=\widetilde{S}$, on
a, pour tout $f\in S_+$ homogène
\begin{equation}
\label{II.8.3.5.4-fr}
  p^{-1}(X_f)=\widehat{C}_f
\tag{8.3.5.4}
\end{equation}
et la restriction de $p$ à
$\widehat{C}_f=\operatorname{Spec}(S_f^{\leq})$, considérée comme morphisme de
$\widehat{C}_f$ dans $X_f$, correspond à l'injection canonique
$S_{(f)}\to S_f^{\leq}$. Si de plus $f\in S_1$, $\widehat{C}_f$ est isomorphe
à $X_f\otimes_{\mathbf{Z}}\mathbf{Z}[T]$ ($T$ indéterminée).
\end{env}

En effet, la formule (\hyperref[II.8.3.5.4-fr]{8.3.5.4}) n'est qu'un cas
particulier de (\hyperref[II.2.8.1.1-fr]{2.8.1.1}) et la seconde assertion
n'est autre que la définition de $\operatorname{Proj}(\varphi)$ lorsque $Y$
est affine (\hyperref[II.2.8.1-fr]{2.8.1}). Alors la formule
(\hyperref[II.8.3.5.2-fr]{8.3.5.2}) et le fait que $p$ est surjectif résultent
de ce que le composé $\mathcal{S}\to\widehat{\mathcal{S}}\to\mathcal{S}$ des
homomorphismes canoniques est l'identité dans $\mathcal{S}$. Enfin, la
dernière assertion de (\hyperref[II.8.3.5.3-fr]{8.3.5.3}) provient de ce que
$S_f^{\leq}$ est isomorphe à $S_{(f)}[T]$ lorsque $f\in S_1$
(\hyperref[II.2.2.1-fr]{2.2.1}).

\begin{corollary}[8.3.6]
\label{II.8.3.6-fr}
La restriction
\begin{equation}
\label{II.8.3.6.1-fr}
  \pi:E\to X
\tag{8.3.6.1}
\end{equation}
de $p$ à $E$ est un morphisme affine surjectif. Si $Y$ est affine et $f$
homogène dans $S_+$, on a
\begin{equation}
\label{II.8.3.6.2-fr}
  \pi^{-1}(X_f)=C_f
\tag{8.3.6.2}
\end{equation}
et la restriction de $\pi$ à $C_f$ correspond à l'injection canonique
$S_{(f)}\to S_f$. Si en outre $f\in S_1$, $C_f$ est isomorphe à
$X_f\otimes_{\mathbf{Z}}\mathbf{Z}[T,T^{-1}]$ ($T$ indéterminée).
\end{corollary}

La formule (\hyperref[II.8.3.6.2-fr]{8.3.6.2}) résulte aussitôt de
(\hyperref[II.8.3.5.3-fr]{8.3.5.3}) et
(\hyperref[II.8.3.2.5-fr]{8.3.2.5}), et démontre la surjectivité de $\pi$~; on
a vu par ailleurs que l'immersion $i$, restreinte à $C_f$, correspondait

\oldpage[II]{166}
à l'injection $S_f^{\leq}\to S_f$ (\hyperref[II.8.3.2-fr]{8.3.2}). Enfin, la
dernière assertion est conséquence de ce que pour $f\in S_1$, $S_f$ est
isomorphe à $S_{(f)}[T,T^{-1}]$ (\hyperref[II.2.2.1-fr]{2.2.1}).

\begin{remark}[8.3.7]
\label{II.8.3.7-fr}
Lorsque $Y$ est affine, les éléments de l'espace sous-jacent à $E$ sont les
idéaux premiers $\mathfrak{p}$ (non nécessairement gradués) de $S$ ne
contenant pas $S_+$, en vertu de la définition de l'immersion $\varepsilon$
(\hyperref[II.8.3.2-fr]{8.3.2}). Pour un tel idéal $\mathfrak{p}$, les
$\mathfrak{p}\cap S_n$ vérifient de façon évidente les conditions de
(\hyperref[II.2.1.9-fr]{2.1.9}), donc il existe un seul idéal premier
\emph{gradué} $\mathfrak{q}$ de $S$ tel que
$\mathfrak{q}\cap S_n=\mathfrak{p}\cap S_n$ pour tout $n$~; l'application
$\pi:E\to X$ des espaces sous-jacents s'interprète alors grâce à la relation
\begin{equation}
\label{II.8.3.7.1-fr}
  \pi(\mathfrak{p})=\mathfrak{q}.
\tag{8.3.7.1}
\end{equation}
En effet, pour vérifier cette relation, il suffit de considérer un $f$
homogène dans $S_+$ tel que $\mathfrak{p}\in D(f)$, et d'observer que
$\mathfrak{q}_{(f)}$ est l'image réciproque de $\mathfrak{p}_f$ par
l'injection $S_{(f)}\to S_f$.
\end{remark}

\begin{corollary}[8.3.8]
\label{II.8.3.8-fr}
Si $\mathcal{S}$ est engendrée par $\mathcal{S}_1$, les morphismes $p$ et
$\pi$ sont de type fini~; pour tout $x\in X$, la fibre $p^{-1}(x)$ est
isomorphe à $\operatorname{Spec}(k(x)[T])$ et la fibre $\pi^{-1}(x)$ est
isomorphe à $\operatorname{Spec}(k(x)[T,T^{-1}])$.
\end{corollary}

Cela résulte aussitôt de (\hyperref[II.8.3.5-fr]{8.3.5}) et
(\hyperref[II.8.3.6-fr]{8.3.6}) en observant que, lorsque $Y$ est affine et
$S$ engendrée par $S_1$, les $X_f$, pour $f\in S_1$, forment un recouvrement
de $X$ (\hyperref[II.2.3.14-fr]{2.3.14}).

\begin{remark}[8.3.9]
\label{II.8.3.9-fr}
Le cône épointé affine correspondant à la $\mathcal{O}_Y$-Algèbre graduée
$\mathcal{O}_Y[T]$ ($T$ indéterminée) s'identifie à
$G_m=\operatorname{Spec}(\mathcal{O}_Y[T,T^{-1}])$, puisqu'il n'est autre que
$C_T$, comme on l'a vu dans (\hyperref[II.8.3.2-fr]{8.3.2}) (voir
(\hyperref[II.8.4.4-fr]{8.4.4}) pour un résultat plus général). Ce préschéma
est canoniquement muni d'une structure de «~\emph{$Y$-schéma en groupes
commutatif}~». Il s'agit là d'une notion qui sera développée en détail plus
tard, et que l'on peut ici résumer rapidement comme suit. Un $Y$-schéma en
groupes est un $Y$-schéma $G$, muni de deux $Y$-morphismes
$p:G\times_YG\to G$ et $s:G\to G$ qui satisfont aux conditions formellement
analogues aux axiomes de la loi de composition et de l'application de
symétrie dans un groupe~: le diagramme
\[
  \xymatrix{
    G\times G\times G \ar[r]^{p\times1}\ar[d]_{1\times p}
      & G\times G \ar[d]^{p} \\
    G\times G \ar[r]_{p} & G
  }
\]
doit être commutatif («~associativité~») et l'on doit avoir une condition qui
correspond dans les groupes au fait que les applications
\[
  (x,y)\mapsto(x,x^{-1},y)\mapsto(x,x^{-1}y)\mapsto x(x^{-1}y)
\]
et
\[
  (x,y)\mapsto(x,x^{-1},y)\mapsto(x,yx^{-1})\mapsto(yx^{-1})x
\]
doivent toutes deux se réduire à $(x,y)\mapsto y$~; la suite de morphismes
correspondant par exemple à la première application composée est
\[
  G\times G\xrightarrow{(1,s)\times1}G\times G\times G
  \xrightarrow{1\times p}G\times G\xrightarrow{p}G
\]
et le lecteur écrira de même la seconde suite.

\oldpage[II]{167}
Il est immédiat (\hyperref[I.3.4.3-fr]{I, 3.4.3}) que la donnée d'une structure
de $Y$-schéma en groupes sur un $Y$-schéma $G$ équivaut à la donnée, pour
\emph{tout} $Y$-préschéma $Z$, d'une structure de \emph{groupe} sur
l'ensemble $\operatorname{Hom}_Y(Z,G)$, ces structures devant être telles que
pour tout $Y$-morphisme $Z\to Z'$, l'application correspondante
$\operatorname{Hom}_Y(Z',G)\to\operatorname{Hom}_Y(Z,G)$ soit un
homomorphisme de groupes. Dans le cas particulier de $G_m$ que nous
considérons ici, $\operatorname{Hom}_Y(Z,G)$ s'identifie à l'ensemble des
$Z$-sections de $Z\times_YG_m$ (\hyperref[I.3.3.14-fr]{I, 3.3.14}), donc à
l'ensemble des $Z$-sections de
$\operatorname{Spec}(\mathcal{O}_Z[T,T^{-1}])$~; finalement, le même
raisonnement que dans (\hyperref[I.3.3.15-fr]{I, 3.3.15}) montre que cet
ensemble s'identifie canoniquement à l'ensemble des éléments \emph{inversibles}
de l'anneau $\Gamma(Z,\mathcal{O}_Z)$, et la structure de groupe sur cet
ensemble est la structure provenant de la multiplication dans l'anneau
$\Gamma(Z,\mathcal{O}_Z)$. Le lecteur pourra vérifier que les morphismes $p$
et $s$ considérés plus haut sont obtenus de la façon suivante~: ils
correspondent d'après (\hyperref[II.1.2.7-fr]{1.2.7}) et
(\hyperref[II.1.4.6-fr]{1.4.6}) à des homomorphismes de $\mathcal{O}_Y$-Algèbres
\begin{align*}
  \pi &: \mathcal{O}_Y[T,T^{-1}]
  \to\mathcal{O}_Y[T,T^{-1},T',T'^{-1}],\\
  \sigma &: \mathcal{O}_Y[T,T^{-1}]\to\mathcal{O}_Y[T,T^{-1}]
\end{align*}
et sont entièrement définis par les données de $\pi(T)=TT'$ et
$\sigma(T)=T^{-1}$.

Cela étant, $G_m$ peut être considéré comme un «~\emph{domaine d'opérateurs
universel}~» pour tout \emph{cône affine}
$C=\operatorname{Spec}(\mathcal{S})$, où $\mathcal{S}$ est une
$\mathcal{O}_Y$-Algèbre graduée quasi-cohérente à degrés positifs. Cela
signifie qu'on peut définir canoniquement un $Y$-morphisme
$G_m\times_YC\to C$ qui a les propriétés formelles d'une loi externe d'un
ensemble muni d'un groupe d'opérateurs~; ou encore, comme ci-dessus pour les
schémas en groupes, on se donne pour tout $Y$-préschéma $Z$ une loi externe
sur $\operatorname{Hom}_Y(Z,C)$, ayant le groupe
$\operatorname{Hom}_Y(Z,G_m)$ comme ensemble d'opérateurs, avec les axiomes
usuels des ensembles munis d'un groupe d'opérateurs, et une condition de
compatibilité avec les $Y$-morphismes $Z\to Z'$. Dans le cas présent, le
morphisme $G_m\times_YC\to C$ est défini par la donnée de l'homomorphisme de
$\mathcal{O}_Y$-Algèbres
\[
  \mathcal{S}\to
  \mathcal{S}\otimes_{\mathcal{O}_Y}\mathcal{O}_Y[T,T^{-1}]
  =\mathcal{S}[T,T^{-1}]
\]
qui, à toute section $s_n\in\Gamma(U,\mathcal{S}_n)$ ($U$ ouvert de $Y$) fait
correspondre la section
$s_nT^n\in\Gamma(U,\mathcal{S}\otimes_{\mathcal{O}_Y}
\mathcal{O}_Y[T,T^{-1}])$.

Inversement, supposons donnée une $\mathcal{O}_Y$-Algèbre quasi-cohérente
$\mathcal{S}$ \emph{non graduée a priori}, et, sur
$C=\operatorname{Spec}(\mathcal{S})$, une structure de «~$Y$-schéma en
ensembles \emph{munis de groupes d'opérateurs}~» ayant pour domaine
d'opérateurs le $Y$-schéma en groupes $G_m$~; alors on en déduit
canoniquement une \emph{graduation} de $\mathcal{O}_Y$-Algèbre sur
$\mathcal{S}$. En effet, la donnée d'un $Y$-morphisme $G_m\times_YC\to C$
équivaut à celle d'un homomorphisme de $\mathcal{O}_Y$-Algèbres
$\psi:\mathcal{S}\to\mathcal{S}[T,T^{-1}]$, qui s'écrira donc
$\psi=\sum_{n\in\mathbf{Z}}\psi_nT^n$, où les
$\psi_n:\mathcal{S}\to\mathcal{S}$ sont des homomorphismes de
$\mathcal{O}_Y$-Modules (avec $\psi_n(s)=0$ sauf pour un nombre fini d'indices
pour toute section $s\in\Gamma(U,\mathcal{S})$, $U$ ouvert dans $Y$). On peut
alors vérifier que les axiomes des ensembles munis d'un groupe d'opérateurs
entraînent que les $\psi_n(\mathcal{S})=\mathcal{S}_n$ définissent une
graduation (à degrés positifs ou négatifs) de $\mathcal{O}_Y$-Algèbre sur
$\mathcal{S}$, les $\psi_n$ étant les projecteurs correspondants. On a ainsi
une notion de structure de «~\emph{cône affine}~» sur tout $Y$-schéma affine,
définie de façon «~géométrique~» sans appel à une graduation préalable.

\oldpage[II]{168}
Nous ne développerons pas davantage ce point de vue ici et laissons au
lecteur le soin de formuler de façon précise les définitions et résultats qui
correspondent aux indications précédentes.
\end{remark}

\subsection{Fermeture projective d'un fibré vectoriel}
\label{subsection:II.8.4-fr}

\begin{env}[8.4.1]
\label{II.8.4.1-fr}
Soient $Y$ un préschéma, $\mathcal{E}$ un $\mathcal{O}_Y$-Module
quasi-cohérent. Si on prend pour $\mathcal{S}$ la $\mathcal{O}_Y$-Algèbre
graduée $\mathbf{S}_{\mathcal{O}_Y}(\mathcal{E})$, la définition
(\hyperref[II.8.3.1.1-fr]{8.3.1.1}) montre que
$\widehat{\mathcal{S}}$ s'identifie à
$\mathbf{S}_{\mathcal{O}_Y}(\mathcal{E}\oplus\mathcal{O}_Y)$. Le cône affine
$\operatorname{Spec}(\mathcal{S})$ défini par $\mathcal{S}$ étant alors par
définition $\mathbf{V}(\mathcal{E})$, et $\operatorname{Proj}(\mathcal{S})$
étant par définition $\mathbf{P}(\mathcal{E})$, on voit donc que~:
\end{env}

\begin{proposition}[8.4.2]
\label{II.8.4.2-fr}
La fermeture projective d'un fibré vectoriel $\mathbf{V}(\mathcal{E})$ sur
$Y$ est canoniquement isomorphe à
$\mathbf{P}(\mathcal{E}\oplus\mathcal{O}_Y)$, et le lieu à l'infini de cette
dernière est canoniquement isomorphe à $\mathbf{P}(\mathcal{E})$.
\end{proposition}

\begin{remark}[8.4.3]
\label{II.8.4.3-fr}
Prenons par exemple $\mathcal{E}=\mathcal{O}_Y^r$ avec $r\geq2$~; alors les
cônes épointés $E$, $\widehat{E}$ définis par $\mathcal{S}$ ne sont ni affines
ni projectifs sur $Y$ si $Y\neq\varnothing$. La seconde assertion est
immédiate, car $\widehat{C}=\mathbf{P}(\mathcal{O}_Y^{r+1})$ est projectif sur
$Y$ et les espaces sous-jacents à $E$ et $\widehat{E}$ sont des ouverts non
fermés dans $\widehat{C}$, donc les immersions canoniques
$E\to\widehat{E}$ et $\widehat{E}\to\widehat{C}$ ne sont pas projectives
(\hyperref[II.5.5.3-fr]{5.5.3}), et on conclut par
(\hyperref[II.5.5.5-fr]{5.5.5, (v)}). D'autre part, supposant
$Y=\operatorname{Spec}(A)$ affine et par exemple $r=2$, on a
$C=\operatorname{Spec}(A[T_1,T_2])$ et $E$ est alors le préschéma induit par
$C$ sur l'ouvert $D(T_1)\cup D(T_2)$~; or nous avons vu que ce dernier n'est
pas affine (\hyperref[I.5.5.11-fr]{I, 5.5.11})~; 
\emph{a fortiori} $\widehat{E}$ ne peut être affine, puisque $E$ est l'ouvert
où ne s'annule pas la section $z$ au-dessus de $\widehat{E}$
(\hyperref[II.8.3.2-fr]{8.3.2}).

Toutefois~:
\end{remark}

\begin{proposition}[8.4.4]
\label{II.8.4.4-fr}
Si $\mathcal{L}$ est un $\mathcal{O}_Y$-Module inversible, on a des
isomorphismes canoniques pour les cônes épointés $E$ et $\widehat{E}$
correspondants à $C=\mathbf{V}(\mathcal{L})$,
\begin{equation}
\label{II.8.4.4.1-fr}
  \operatorname{Spec}\left(
    \bigoplus_{n\in\mathbf{Z}}\mathcal{L}^{\otimes n}
  \right)\xrightarrow{\sim}E
\tag{8.4.4.1}
\end{equation}
\begin{equation}
\label{II.8.4.4.2-fr}
  \mathbf{V}(\mathcal{L}^{-1})\xrightarrow{\sim}\widehat{E}.
\tag{8.4.4.2}
\end{equation}

En outre, il existe un isomorphisme canonique de la fermeture projective de
$\mathbf{V}(\mathcal{L})$ sur celle de $\mathbf{V}(\mathcal{L}^{-1})$,
transformant la section nulle (resp. le lieu à l'infini) de la première en le
lieu à l'infini (resp. la section nulle) de la seconde.
\end{proposition}

\begin{proof}
On a ici
$\mathcal{S}=\bigoplus_{n\geq0}\mathcal{L}^{\otimes n}$~; l'injection
canonique
\[
  \mathcal{S}\to\bigoplus_{n\in\mathbf{Z}}\mathcal{L}^{\otimes n}
\]
définit un morphisme dominant canonique
\begin{equation}
\label{II.8.4.4.3-fr}
  \operatorname{Spec}\left(
    \bigoplus_{n\in\mathbf{Z}}\mathcal{L}^{\otimes n}
  \right)
  \to\mathbf{V}(\mathcal{L})
  =\operatorname{Spec}\left(
    \bigoplus_{n\geq0}\mathcal{L}^{\otimes n}
  \right)
\tag{8.4.4.3}
\end{equation}
et il suffit de prouver que ce morphisme est un isomorphisme du schéma
$\operatorname{Spec}(\bigoplus_{n\in\mathbf{Z}}\mathcal{L}^{\otimes n})$
sur $E$. La question étant locale sur $Y$, on peut supposer que
$Y=\operatorname{Spec}(A)$ est affine

\oldpage[II]{169}
et $\mathcal{L}=\mathcal{O}_Y$, donc $\mathcal{S}=(A[T])^{\sim}$ et
$\bigoplus_{n\in\mathbf{Z}}\mathcal{L}^{\otimes n}
=(A[T,T^{-1}])^{\sim}$. Or $A[T,T^{-1}]$ est l'anneau de fractions
$A[T]_T$ de $A[T]$, donc (\hyperref[II.8.4.4.3-fr]{8.4.4.3}) identifie le
premier membre au préschéma induit par $C=\mathbf{V}(\mathcal{L})$ sur
l'ouvert $D(T)$~; le complémentaire $V(T)$ de cet ouvert dans $C$ est
l'espace sous-jacent au sous-préschéma fermé de $C$ défini par l'idéal
$TA[T]$, c'est-à-dire la section nulle de $C$, donc $E=D(T)$.

L'isomorphisme (\hyperref[II.8.4.4.2-fr]{8.4.4.2}) sera d'autre part
conséquence de la dernière assertion, $\mathbf{V}(\mathcal{L}^{-1})$ étant le
complémentaire du lieu à l'infini de sa fermeture projective et $\widehat{E}$
le complémentaire de la section nulle de la fermeture projective de
$C=\mathbf{V}(\mathcal{L})$. Or, ces fermetures projectives sont
respectivement $\mathbf{P}(\mathcal{L}^{-1}\oplus\mathcal{O}_Y)$ et
$\mathbf{P}(\mathcal{L}\oplus\mathcal{O}_Y)$~; mais on peut écrire
$\mathcal{L}\oplus\mathcal{O}_Y
=\mathcal{L}\otimes(\mathcal{L}^{-1}\oplus\mathcal{O}_Y)$. L'existence de
l'isomorphisme canonique cherché résulte alors de
(\hyperref[II.4.1.4-fr]{4.1.4}), et tout revient à voir que cet isomorphisme
échange sections nulles et lieux à l'infini. On peut pour cela se borner au
cas où $Y=\operatorname{Spec}(A)$ est affine,
$L=Ac$, $L^{-1}=Ac'$, l'isomorphisme canonique
$L\otimes L^{-1}\to A$ faisant correspondre à $c\otimes c'$ l'élément $1$
de $A$. Alors $\mathbf{S}(L\oplus A)$ est produit tensoriel de $A[z]$ et de
$\bigoplus_{n\geq0}Ac^{\otimes n}$, $\mathbf{S}(L^{-1}\oplus A)$ produit
tensoriel de $A[z]$ et de $\bigoplus_{n\geq0}Ac'^{\otimes n}$, et
l'isomorphisme défini dans (\hyperref[II.4.1.4-fr]{4.1.4}) fait correspondre à
$z^h\otimes c'^{\otimes(n-h)}$ l'élément
$z^{n-h}\otimes c^{\otimes h}$. Or, dans
$\mathbf{P}(\mathcal{L}^{-1}\oplus\mathcal{O}_Y)$, le lieu à l'infini est
l'ensemble des points où s'annule la section $z$ et la section nulle est
l'ensemble des points où s'annule la section $c'$~; comme on a des définitions
analogues pour $\mathbf{P}(\mathcal{L}\oplus\mathcal{O}_Y)$, notre conclusion
résulte aussitôt de l'explicitation qui précède.
\end{proof}

\subsection{Comportements fonctoriels}
\label{subsection:II.8.5-fr}

\begin{env}[8.5.1]
\label{II.8.5.1-fr}
Soient $Y$, $Y'$ deux préschémas, $q:Y'\to Y$ un morphisme,
$\mathcal{S}$ (resp. $\mathcal{S}'$) une $\mathcal{O}_Y$-Algèbre
(resp. une $\mathcal{O}_{Y'}$-Algèbre) graduée quasi-cohérente à degrés
\emph{positifs}. Considérons un $q$-morphisme d'Algèbres graduées
\begin{equation}
\label{II.8.5.1.1-fr}
  \varphi:\mathcal{S}\to\mathcal{S}'.
\tag{8.5.1.1}
\end{equation}
On sait (\hyperref[II.1.5.6-fr]{1.5.6}) qu'il lui correspond canoniquement un
morphisme
\[
  \Phi=\operatorname{Spec}(\varphi):
  \operatorname{Spec}(\mathcal{S}')\to\operatorname{Spec}(\mathcal{S})
\]
tel que le diagramme
\begin{equation}
\label{II.8.5.1.2-fr}
  \xymatrix{
    C'\ar[r]^{\Phi}\ar[d] & C\ar[d]\\
    Y'\ar[r]_{q} & Y
  }
\tag{8.5.1.2}
\end{equation}
où on a posé $C=\operatorname{Spec}(\mathcal{S})$,
$C'=\operatorname{Spec}(\mathcal{S}')$, soit commutatif. 
\emph{Supposons en outre que
$\mathcal{S}_0=\mathcal{O}_Y$, $\mathcal{S}'_0=\mathcal{O}_{Y'}$}~; soient
$\varepsilon:Y\to C$ et $\varepsilon':Y'\to C'$ les immersions canoniques
(\hyperref[II.8.3.2-fr]{8.3.2})~; on a alors un diagramme commutatif
\begin{equation}
\label{II.8.5.1.3-fr}
  \xymatrix{
    Y'\ar[r]^{q}\ar[d]_{\varepsilon'} & Y\ar[d]^{\varepsilon}\\
    C'\ar[r]_{\Phi} & C
  }
\tag{8.5.1.3}
\end{equation}

\oldpage[II]{170}
qui correspond au diagramme
\[
  \xymatrix{
    \mathcal{S}\ar[r]^{\varphi}\ar[d] & \mathcal{S}'\ar[d]\\
    \mathcal{O}_Y\ar[r] & \mathcal{O}_{Y'}
  }
\]
où les flèches verticales sont les homomorphismes d'augmentation, et dont la
commutativité résulte de l'hypothèse que $\varphi$ est supposé être un
homomorphisme d'Algèbres \emph{graduées}.
\end{env}

\begin{proposition}[8.5.2]
\label{II.8.5.2-fr}
Si $E$ (resp. $E'$) est le cône épointé affine défini par $\mathcal{S}$
(resp. $\mathcal{S}'$), on a $\Phi^{-1}(E)\subset E'$~; si en outre
$\operatorname{Proj}(\varphi):G(\varphi)\to\operatorname{Proj}(\mathcal{S})$
est partout défini (autrement dit si
$G(\varphi)=\operatorname{Proj}(\mathcal{S}')$), alors on a
$\Phi^{-1}(E)=E'$, et réciproquement.
\end{proposition}

\begin{proof}
La première assertion résulte de la commutativité de
(\hyperref[II.8.5.1.3-fr]{8.5.1.3}). Pour démontrer la seconde, on peut se
borner au cas où $Y=\operatorname{Spec}(A)$ et
$Y'=\operatorname{Spec}(A')$ sont affines,
$\mathcal{S}=\widetilde{S}$, $\mathcal{S}'=\widetilde{S'}$. Pour tout $f$
homogène dans $S_+$, si on pose $f'=\varphi(f)$, on a
$\Phi^{-1}(C_f)=C'_{f'}$ (\hyperref[I.2.2.4.1-fr]{I, 2.2.4.1})~; dire que
$G(\varphi)=\operatorname{Proj}(S')$ signifie que dans $S'_+$ la racine de
l'idéal engendré par les $f'=\varphi(f)$ est $S'_+$ lui-même
(\hyperref[II.2.8.1-fr]{2.8.1}) et (\hyperref[II.2.3.14-fr]{2.3.14}), et cela
équivaut encore à dire que les $C'_{f'}$ recouvrent $E'$
(\hyperref[II.8.3.4.4-fr]{8.3.4.4}).
\end{proof}

\begin{env}[8.5.3]
\label{II.8.5.3-fr}
Le $q$-morphisme $\varphi$ se prolonge canoniquement en un $q$-morphisme
d'Algèbres graduées
\begin{equation}
\label{II.8.5.3.1-fr}
  \widehat{\varphi}:\widehat{\mathcal{S}}
  \to\widehat{\mathcal{S}'}
\tag{8.5.3.1}
\end{equation}
en posant $\widehat{\varphi}(z)=z$. On en déduit par suite un morphisme
\[
  \widehat{\Phi}=\operatorname{Proj}(\widehat{\varphi}):
  G(\widehat{\varphi})\to\widehat{C}
  =\operatorname{Proj}(\widehat{\mathcal{S}})
\]
tel que le diagramme
\[
  \xymatrix{
    G(\widehat{\varphi})\ar[r]^{\widehat{\Phi}}\ar[d] &
    \widehat{C}\ar[d]\\
    Y'\ar[r]_{q} & Y
  }
\]
soit commutatif (\hyperref[II.3.5.6-fr]{3.5.6}). Il résulte aussitôt des
définitions que si on désigne par $i:C\to\widehat{C}$ et
$i':C'\to\widehat{C'}$ les immersions ouvertes canoniques
(\hyperref[II.8.3.2-fr]{8.3.2}), on a
$i'(C')\subset G(\widehat{\varphi})$ et le diagramme
\begin{equation}
\label{II.8.5.3.2-fr}
  \xymatrix{
    C'\ar[r]^{\Phi}\ar[d]_{i'} & C\ar[d]^{i}\\
    G(\widehat{\varphi})\ar[r]_{\widehat{\Phi}} & \widehat{C}
  }
\tag{8.5.3.2}
\end{equation}
est commutatif. Enfin, si on pose
$X=\operatorname{Proj}(\mathcal{S})$, $X'=\operatorname{Proj}(\mathcal{S}')$,
et si $j:X\to\widehat{C}$, $j':X'\to\widehat{C'}$ sont les immersions
fermées canoniques (\hyperref[II.8.3.2-fr]{8.3.2}), il résulte encore des
définitions de ces immersions que l'on a
$j'(G(\varphi))\subset G(\widehat{\varphi})$ et que le diagramme

\oldpage[II]{171}
\begin{equation}
\label{II.8.5.3.3-fr}
  \xymatrix{
    G(\varphi)\ar[r]^{\operatorname{Proj}(\varphi)}\ar[d]_{j'} &
    X\ar[d]^{j}\\
    G(\widehat{\varphi})\ar[r]_{\widehat{\Phi}} & \widehat{C}
  }
\tag{8.5.3.3}
\end{equation}
est commutatif.
\end{env}

\begin{proposition}[8.5.4]
\label{II.8.5.4-fr}
Si $\widehat{E}$ (resp. $\widehat{E'}$) est le cône épointé projectif défini
par $\mathcal{S}$ (resp. $\mathcal{S}'$), on a
$\widehat{\Phi}^{-1}(\widehat{E})\subset\widehat{E'}$~; en outre, si
$p:\widehat{E}\to X$ et $p':\widehat{E'}\to X'$ sont les rétractions
canoniques, on a
$p'(\widehat{\Phi}^{-1}(\widehat{E}))\subset G(\varphi)$, et le diagramme
\begin{equation}
\label{II.8.5.4.1-fr}
  \xymatrix{
    \widehat{\Phi}^{-1}(\widehat{E})
      \ar[r]^{\widehat{\Phi}}\ar[d]_{p'} &
    \widehat{E}\ar[d]^{p}\\
    G(\varphi)\ar[r]_{\operatorname{Proj}(\varphi)} & X
  }
\tag{8.5.4.1}
\end{equation}
est commutatif. Si $\operatorname{Proj}(\varphi)$ est partout défini, il en
est de même de $\widehat{\Phi}$, et on a
$\widehat{\Phi}^{-1}(\widehat{E})=\widehat{E'}$.
\end{proposition}

\begin{proof}
La première assertion résulte de la commutativité des diagrammes
(\hyperref[II.8.5.1.3-fr]{8.5.1.3}) et
(\hyperref[II.8.5.3.2-fr]{8.5.3.2}), et les deux suivantes de la définition
des rétractions canoniques (\hyperref[II.8.3.5-fr]{8.3.5}) et de la définition
de $\widehat{\varphi}$. D'autre part, pour voir que $\widehat{\Phi}$ est
partout défini lorsqu'il en est ainsi de $\operatorname{Proj}(\varphi)$, on
peut se borner à considérer le cas où $Y=\operatorname{Spec}(A)$ et
$Y'=\operatorname{Spec}(A')$ sont affines,
$\mathcal{S}=\widetilde{S}$, $\mathcal{S}'=\widetilde{S'}$~; l'hypothèse est
que lorsque $f$ parcourt l'ensemble des éléments homogènes de $S_+$, l'idéal
engendré dans $S'_+$ par les $\varphi(f)$ a pour racine dans $S'_+$ l'idéal
$S'_+$ lui-même~; on en conclut aussitôt que la racine dans $(S'[z])_+$ de
l'idéal engendré par $z$ et les $\varphi(f)$ est $(S'[z])_+$ lui-même, d'où
notre assertion~; cela prouve en même temps que $\widehat{E'}$ est réunion des
$\widehat{C'}_{\varphi(f)}$, donc égal à
$\widehat{\Phi}^{-1}(\widehat{E})$.
\end{proof}

\begin{corollary}[8.5.5]
\label{II.8.5.5-fr}
Lorsque $\operatorname{Proj}(\varphi)$ est partout défini, l'image réciproque
par $\widehat{\Phi}$ de l'espace sous-jacent au lieu à l'infini (resp. au
préschéma des sommets) de $\widehat{C}$ est l'espace sous-jacent au lieu à
l'infini (resp. au préschéma des sommets) de $\widehat{C'}$.
\end{corollary}

\begin{proof}
Cela résulte aussitôt de (\hyperref[II.8.5.4-fr]{8.5.4}) et
(\hyperref[II.8.5.2-fr]{8.5.2}), compte tenu des relations
(\hyperref[II.8.3.4.1-fr]{8.3.4.1}) et
(\hyperref[II.8.3.4.2-fr]{8.3.4.2}).
\end{proof}

\subsection{Un isomorphisme canonique pour les cônes épointés}
\label{subsection:II.8.6-fr}

\begin{env}[8.6.1]
\label{II.8.6.1-fr}
Soient $Y$ un préschéma, $\mathcal{S}$ une $\mathcal{O}_Y$-Algèbre graduée
quasi-cohérente à degrés positifs \emph{telle que
$\mathcal{S}_0=\mathcal{O}_Y$}, $X$ le $Y$-schéma
$\operatorname{Proj}(\mathcal{S})$. Nous allons appliquer les résultats de
(\hyperref[subsection:II.8.5-fr]{8.5}) au cas où on prend $Y'=X$,
$q:X\to Y$ étant le morphisme structural~; soit
\begin{equation}
\label{II.8.6.1.1-fr}
  \mathcal{S}_X=\bigoplus_{n\in\mathbf{Z}}\mathcal{O}_X(n)
\tag{8.6.1.1}
\end{equation}

\oldpage[II]{172}
qui est une $\mathcal{O}_X$-Algèbre graduée quasi-cohérente, la multiplication
y étant définie au moyen des homomorphismes canoniques
(\hyperref[II.3.2.6.1-fr]{3.2.6.1})
\[
  \mathcal{O}_X(m)\otimes_{\mathcal{O}_X}\mathcal{O}_X(n)
  \to\mathcal{O}_X(m+n)
\]
dont l'associativité est garantie par le diagramme commutatif
(\hyperref[II.2.5.11.4-fr]{2.5.11.4}). Prenons pour $\mathcal{S}'$ la
sous-$\mathcal{O}_X$-Algèbre graduée quasi-cohérente
$\mathcal{S}_X^{\geq}=\bigoplus_{n\geq0}\mathcal{O}_X(n)$ de
$\mathcal{S}_X$, à degrés positifs.

Enfin, nous considérerons le $q$-morphisme canonique
\begin{equation}
\label{II.8.6.1.2-fr}
  \alpha:\mathcal{S}\to\mathcal{S}_X^{\geq}
\tag{8.6.1.2}
\end{equation}
défini dans (\hyperref[II.3.3.2.3-fr]{3.3.2.3}) comme un homomorphisme
$\mathcal{S}\to q_*(\mathcal{S}_X)$, mais qui applique évidemment
$\mathcal{S}$ dans $q_*(\mathcal{S}_X^{\geq})$. Nous poserons
\begin{equation}
\label{II.8.6.1.3-fr}
  C_X=\operatorname{Spec}(\mathcal{S}_X^{\geq}),\qquad
  \widehat{C}_X=\operatorname{Proj}(\mathcal{S}_X^{\geq}[z]),\qquad
  X'=\operatorname{Proj}(\mathcal{S}_X^{\geq})
\tag{8.6.1.3}
\end{equation}
et nous désignerons par $E_X$ et $\widehat{E}_X$ le cône épointé affine et le
cône épointé projectif correspondants~; on notera
$\varepsilon_X:X\to C_X$, $i_X:C_X\to\widehat{C}_X$,
$j_X:X'\to\widehat{C}_X$, $p_X:\widehat{E}_X\to X'$,
$\pi_X:E_X\to X'$ les morphismes canoniques définis dans
(\hyperref[subsection:II.8.3-fr]{8.3}).
\end{env}

\begin{proposition}[8.6.2]
\label{II.8.6.2-fr}
Le morphisme structural $u:X'\to X$ est un isomorphisme, et le morphisme
$\operatorname{Proj}(\alpha)$ est partout défini et identique à $u$. Le
morphisme
$\operatorname{Proj}(\widehat{\alpha}):\widehat{C}_X\to\widehat{C}$ est
partout défini, et ses restrictions à $\widehat{E}_X$ et $E_X$ sont des
isomorphismes sur $\widehat{E}$ et $E$ respectivement. Enfin, lorsqu'on
identifie $X'$ et $X$ au moyen de $u$, les morphismes $p_X$ et $\pi_X$
s'identifient aux morphismes structuraux des $X$-préschémas $\widehat{E}_X$ et
$E_X$.
\end{proposition}

\begin{proof}
On peut évidemment se borner au cas où $Y=\operatorname{Spec}(A)$ est affine
et $\mathcal{S}=\widetilde{S}$~; alors $X$ est réunion des ouverts affines
$X_f$, où $f$ parcourt l'ensemble des éléments homogènes de $S_+$, l'anneau
de $X_f$ étant $S_{(f)}$. Il résulte en outre de
(\hyperref[II.8.2.7.1-fr]{8.2.7.1}) que l'on a
\begin{equation}
\label{II.8.6.2.1-fr}
  \Gamma(X_f,\mathcal{S}_X^{\geq})=S_f^{\geq}.
\tag{8.6.2.1}
\end{equation}

On a donc $u^{-1}(X_f)=\operatorname{Proj}(S_f^{\geq})$. Mais si
$f\in S_d$ ($d>0$), $\operatorname{Proj}(S_f^{\geq})$ est canoniquement
isomorphe à $\operatorname{Proj}((S_f^{\geq})^{(d)})$
(\hyperref[II.2.4.7-fr]{2.4.7}), et d'autre part
$(S_f^{\geq})^{(d)}=(S^{(d)})_f^{\geq}$ s'identifie à $S_{(f)}[T]$
(\hyperref[II.2.2.1-fr]{2.2.1}) par l'application $T\to f/1$~; on en conclut
(\hyperref[II.3.1.7-fr]{3.1.7}) que le morphisme structural
$u^{-1}(X_f)\to X_f$ est un isomorphisme, d'où la première assertion. Pour
démontrer la seconde, remarquons que la restriction
$u^{-1}(X_f)\cap G(\alpha)\to X=\operatorname{Proj}(S)$ de
$\operatorname{Proj}(\alpha)$ correspond à l'application canonique
$x\to x/1$ de $S$ dans $S_f^{\geq}$
(\hyperref[II.2.6.2-fr]{2.6.2})~; on en déduit tout d'abord que
$G(\alpha)=X'$, puis, en tenant compte de ce que
$u^{-1}(X_f)=(u^{-1}(X_f))_{f/1}$, il résulte de
(\hyperref[II.2.8.1.1-fr]{2.8.1.1}) que l'image de $u^{-1}(X_f)$ par
$\operatorname{Proj}(\alpha)$ est contenue dans $X_f$, et la restriction de
$\operatorname{Proj}(\alpha)$ à $u^{-1}(X_f)$, considérée comme morphisme
dans $X_f=\operatorname{Spec}(S_{(f)})$, est bien identique à celle de $u$.
Enfin, la formule (\hyperref[II.8.3.5.4-fr]{8.3.5.4}), appliquée à $p_X$ au
lieu de $p$, montre que
$p_X^{-1}(u^{-1}(X_f))=
\operatorname{Spec}((S_f^{\geq})_{f/1}^{\leq})$, et cet ouvert est, d'après
(\hyperref[II.8.5.4.1-fr]{8.5.4.1}), l'image réciproque par
$\operatorname{Proj}(\widehat{\alpha})$ de
$p^{-1}(X_f)=\operatorname{Spec}(S_f^{\leq})$
(\hyperref[II.8.3.5.3-fr]{8.3.5.3}). Compte tenu de
(\hyperref[II.8.2.3.2-fr]{8.2.3.2}), la restriction de
$\operatorname{Proj}(\widehat{\alpha})$ à
$p_X^{-1}(u^{-1}(X_f))$ correspond à l'isomorphisme réciproque de
(\hyperref[II.8.2.7.2-fr]{8.2.7.2}), restreint à $S_f^{\leq}$, d'où le
troisième point~; la dernière assertion est évidente par définition.

\oldpage[II]{173}
On notera aussi qu'il résulte du diagramme commutatif
(\hyperref[II.8.5.3.2-fr]{8.5.3.2}) que \emph{la restriction à $C_X$ de
$\operatorname{Proj}(\widehat{\alpha})$ n'est autre que le morphisme
$\operatorname{Spec}(\alpha)$}.
\end{proof}

\begin{corollary}[8.6.3]
\label{II.8.6.3-fr}
Considérés comme $X$-schémas, $\widehat{E}_X$ est canoniquement isomorphe à
$\operatorname{Spec}(\mathcal{S}_X^{\leq})$, et $E_X$ à
$\operatorname{Spec}(\mathcal{S}_X)$.
\end{corollary}

\begin{proof}
Comme on sait que les morphismes $p_X$ et $\pi_X$ sont affines
(\hyperref[II.8.3.5-fr]{8.3.5} et
\hyperref[II.8.3.6-fr]{8.3.6}), il suffit (vu
(\hyperref[II.1.3.1-fr]{1.3.1})) de vérifier le corollaire lorsque
$Y=\operatorname{Spec}(A)$ est affine et $\mathcal{S}=\widetilde{S}$. La
première assertion résulte de l'existence de l'isomorphisme canonique
(\hyperref[II.8.2.7.2-fr]{8.2.7.2})
$(S_f^{\geq})_{f/1}^{\leq}\xrightarrow{\sim}S_f^{\leq}$ et du fait que ces
isomorphismes sont compatibles avec le passage de $f$ à $fg$ ($f$ et $g$
homogènes dans $S_+$). De même, la formule
(\hyperref[II.8.3.6.2-fr]{8.3.6.2}), appliquée à $\pi_X$ au lieu de $\pi$,
montre que
$\pi_X^{-1}(u^{-1}(X_f))=
\operatorname{Spec}((S_f^{\geq})_{f/1})$ pour $f$ homogène dans $S_+$, et la
seconde assertion découle de l'existence de l'isomorphisme canonique
(\hyperref[II.8.2.7.2-fr]{8.2.7.2})
$(S_f^{\geq})_{f/1}\xrightarrow{\sim}S_f$.

On peut ainsi dire que $\widehat{C}_X$, considéré comme $X$-schéma, s'obtient
par \emph{recollement} des deux $X$-schémas affines sur $X$,
$C_X=\operatorname{Spec}(\mathcal{S}_X^{\geq})$ et
$\widehat{E}_X=\operatorname{Spec}(\mathcal{S}_X^{\leq})$, dont
l'intersection est l'ouvert $E_X=\operatorname{Spec}(\mathcal{S}_X)$.
\end{proof}

\begin{corollary}[8.6.4]
\label{II.8.6.4-fr}
Supposons que $\mathcal{O}_X(1)$ soit un $\mathcal{O}_X$-Module inversible et
que $\mathcal{S}_X$ soit isomorphe à
$\bigoplus_{n\in\mathbf{Z}}(\mathcal{O}_X(1))^{\otimes n}$ (ce qui sera en
particulier le cas si $\mathcal{S}$ est engendrée par $\mathcal{S}_1$
(\hyperref[II.3.2.5-fr]{3.2.5} et
\hyperref[II.3.2.7-fr]{3.2.7})). Alors le cône projectif épointé
$\widehat{E}$ s'identifie au fibré vectoriel de rang un
$\mathbf{V}(\mathcal{O}_X(-1))$ sur $X$, et le cône affine épointé $E$ au
sous-préschéma de ce fibré vectoriel induit sur le complémentaire de la
section nulle. Avec cette identification, la rétraction canonique
$\widehat{E}\to X$ s'identifie au morphisme structural du $X$-schéma
$\mathbf{V}(\mathcal{O}_X(-1))$. Enfin, il existe un $Y$-morphisme canonique
$\mathbf{V}(\mathcal{O}_X(1))\to C$, dont la restriction au complémentaire de
la section nulle de $\mathbf{V}(\mathcal{O}_X(1))$ est un isomorphisme de ce
complémentaire sur le cône affine épointé $E$.
\end{corollary}

\begin{proof}
En effet, si on pose $\mathcal{L}=\mathcal{O}_X(1)$,
$\mathcal{S}_X^{\geq}$ est alors identique à
$\mathbf{S}_{\mathcal{O}_X}(\mathcal{L})$, donc $\widehat{E}_X$ s'identifie
canoniquement à $\mathbf{V}(\mathcal{L}^{-1})$ en vertu de
(\hyperref[II.8.6.3-fr]{8.6.3}) et $C_X$ à $\mathbf{V}(\mathcal{L})$. Le
morphisme $\mathbf{V}(\mathcal{L})\to C$ est la restriction de
$\operatorname{Proj}(\widehat{\alpha})$, et les assertions du corollaire sont
donc des cas particuliers de (\hyperref[II.8.6.2-fr]{8.6.2}).
\end{proof}

On notera que l'image réciproque par le morphisme
$\mathbf{V}(\mathcal{O}_X(1))\to C$ de l'espace sous-jacent au préschéma des
sommets de $C$ est l'espace sous-jacent à la section nulle de
$\mathbf{V}(\mathcal{O}_X(1))$ (\hyperref[II.8.5.5-fr]{8.5.5})~; mais en
général les sous-préschémas correspondants de $C$ et de
$\mathbf{V}(\mathcal{O}_X(1))$ ne sont pas isomorphes. Cette question va être
étudiée ci-dessous.

\subsection{Éclatement de cônes projetants}
\label{subsection:II.8.7-fr}

\begin{env}[8.7.1]
\label{II.8.7.1-fr}
Sous les conditions de (\hyperref[II.8.6.1-fr]{8.6.1}), on a, en posant
$r=\operatorname{Proj}(\widehat{\alpha})$, un diagramme commutatif
\begin{equation}
\label{II.8.7.1.1-fr}
  \xymatrix{
    X\ar[r]^{i_X\circ\varepsilon_X}\ar[d]_{q} &
    \widehat{C}_X\ar[d]^{r}\\
    Y\ar[r]_{i\circ\varepsilon} & \widehat{C}
  }
\tag{8.7.1.1}
\end{equation}

\oldpage[II]{174}
en vertu de (\hyperref[II.8.5.1.3-fr]{8.5.1.3}) et
(\hyperref[II.8.5.3.2-fr]{8.5.3.2})~; en outre, la restriction de $r$ au
complémentaire $\widehat{C}_X-i_X(\varepsilon_X(X))$ de la section nulle est
un \emph{isomorphisme} sur le complémentaire
$\widehat{C}-i(\varepsilon(Y))$ de la section nulle en vertu de
(\hyperref[II.8.6.2-fr]{8.6.2}). Si on suppose pour simplifier $Y$ affine,
$\mathcal{S}$ de type fini et engendré par $\mathcal{S}_1$, $X$ est projectif
sur $Y$ et $\widehat{C}_X$ projectif sur $X$
(\hyperref[II.5.5.1-fr]{5.5.1}), donc $\widehat{C}_X$ est projectif sur $Y$
(\hyperref[II.5.5.5-fr]{5.5.5, (ii)}), et \emph{a fortiori} sur
$\widehat{C}$ (\hyperref[II.5.5.5-fr]{5.5.5, (v)}). On a ainsi un
$Y$-morphisme projectif $r:\widehat{C}_X\to\widehat{C}$ (dont la restriction
à $C_X$ est un $Y$-morphisme projectif $C_X\to C$) qui \emph{contracte $X$
dans $Y$} et induit un \emph{isomorphisme} quand on le restreint aux
\emph{complémentaires de $X$ et de $Y$}. On a donc une relation entre $C_X$
et $C$, analogue à celle qui a lieu entre un préschéma éclaté et son préschéma
de départ (\hyperref[II.8.1.3-fr]{8.1.3}). Nous allons effectivement montrer
qu'on peut identifier $C_X$ au spectre homogène d'une
$\mathcal{O}_C$-Algèbre graduée.
\end{env}

\begin{env}[8.7.2]
\label{II.8.7.2-fr}
Gardant les notations de (\hyperref[II.8.6.1-fr]{8.6.1}), considérons pour
chaque $n\geqslant 0$, l'Idéal quasi-cohérent
\begin{equation}
\label{II.8.7.2.1-fr}
  \mathcal{S}_{[n]}=\bigoplus_{m\geqslant n}\mathcal{S}_m
\tag{8.7.2.1}
\end{equation}
de la $\mathcal{O}_Y$-Algèbre graduée $\mathcal{S}$. Il est clair que l'on a
\begin{equation}
\label{II.8.7.2.2-fr}
  \mathcal{S}_{[0]}=\mathcal{S},\qquad
  \mathcal{S}_{[n]}\subset\mathcal{S}_{[m]}
  \qquad\text{pour }m\leqslant n
\tag{8.7.2.2}
\end{equation}
\begin{equation}
\label{II.8.7.2.3-fr}
  \mathcal{S}_{[n]}\mathcal{S}_{[m]}\subset\mathcal{S}_{[m+n]}.
\tag{8.7.2.3}
\end{equation}

Considérons le $\mathcal{O}_C$-Module associé à $\mathcal{S}_{[n]}$, qui est un
Idéal quasi-cohérent de $\mathcal{O}_C=\widetilde{\mathcal{S}}$
(\hyperref[II.1.4.4-fr]{1.4.4})
\begin{equation}
\label{II.8.7.2.4-fr}
  \mathcal{I}_n=(\mathcal{S}_{[n]})^{\sim}.
\tag{8.7.2.4}
\end{equation}

On déduit alors de (\hyperref[II.8.7.2.2-fr]{8.7.2.2}) et
(\hyperref[II.8.7.2.3-fr]{8.7.2.3}), utilisant
(\hyperref[II.1.4.4-fr]{1.4.4}) et (\hyperref[II.1.4.8.1-fr]{1.4.8.1}),
les formules analogues
\begin{equation}
\label{II.8.7.2.5-fr}
  \mathcal{I}_0=\mathcal{O}_C,\qquad
  \mathcal{I}_n\subset\mathcal{I}_m
  \qquad\text{pour }m\leqslant n
\tag{8.7.2.5}
\end{equation}
\begin{equation}
\label{II.8.7.2.6-fr}
  \mathcal{I}_n\mathcal{I}_m\subset\mathcal{I}_{n+m}.
\tag{8.7.2.6}
\end{equation}

On est donc dans les conditions de (\hyperref[II.8.1.1-fr]{8.1.1}), ce qui
conduit à introduire la $\mathcal{O}_C$-Algèbre graduée quasi-cohérente
\begin{equation}
\label{II.8.7.2.7-fr}
  \mathcal{S}^{\natural}
  =\bigoplus_{n\geqslant 0}\mathcal{I}_n
  =\left(\bigoplus_{n\geqslant 0}\mathcal{S}_{[n]}\right)^{\sim}.
\tag{8.7.2.7}
\end{equation}
\end{env}

\begin{proposition}[8.7.3]
\label{II.8.7.3-fr}
Il existe un $C$-isomorphisme canonique
\begin{equation}
\label{II.8.7.3.1-fr}
  h:C_X\xrightarrow{\sim}\operatorname{Proj}(\mathcal{S}^{\natural}).
\tag{8.7.3.1}
\end{equation}
\end{proposition}

\begin{proof}
Supposons d'abord $Y=\operatorname{Spec}(A)$ affine, d'où
$\mathcal{S}=\widetilde{S}$, où $S$ est une $A$-algèbre graduée à degrés
positifs et $C=\operatorname{Spec}(S)$. La définition
(\hyperref[II.8.2.7.4-fr]{8.2.7.4}) montre que l'on a alors, avec les
notations de (\hyperref[II.8.2.6-fr]{8.2.6}),
$\mathcal{S}^{\natural}=(S^{\natural})^{\sim}$. Pour définir
(\hyperref[II.8.7.3.1-fr]{8.7.3.1}), considérons un élément homogène
$f\in S_d$ ($d>0$) et l'élément correspondant $f^{\natural}\in S^{\natural}$
(\hyperref[II.8.2.6-fr]{8.2.6})~; le $S$-isomorphisme
(\hyperref[II.8.2.7.3-fr]{8.2.7.3}) définit donc un $C$-isomorphisme
\begin{equation}
\label{II.8.7.3.2-fr}
  \operatorname{Spec}(S_f^{\geqslant})
  \xrightarrow{\sim}
  \operatorname{Spec}(S_{(f^{\natural})}^{\natural}).
\tag{8.7.3.2}
\end{equation}

\oldpage[II]{175}
Mais avec les notations de (\hyperref[II.8.6.2-fr]{8.6.2}), si
$v:C_X\to X$ est le morphisme structural, il résulte de
(\hyperref[II.8.6.2.1-fr]{8.6.2.1}) que l'on a
$v^{-1}(X_f)=\operatorname{Spec}(S_f^{\geqslant})$. On a d'autre part
$\operatorname{Spec}(S_{(f^{\natural})}^{\natural})=D_+(f^{\natural})$, si
bien que (\hyperref[II.8.7.3.2-fr]{8.7.3.2}) définit un isomorphisme
$v^{-1}(X_f)\to D_+(f^{\natural})$. En outre, si $g\in S_e$ ($e>0$), le
diagramme
\[
  \xymatrix{
    v^{-1}(X_{fg})\ar[r]^{\sim}\ar[d] &
    D_+(f^{\natural}g^{\natural})\ar[d]\\
    v^{-1}(X_f)\ar[r]^{\sim} & D_+(f^{\natural})
  }
\]
est commutatif, comme il résulte aussitôt de la définition de l'isomorphisme
(\hyperref[II.8.2.7.3-fr]{8.2.7.3}). Enfin par définition $S_+$ est engendré
par les $f$ homogènes, donc il résulte de
(\hyperref[II.8.2.10-fr]{8.2.10, (iv)}) et de
(\hyperref[II.2.3.14-fr]{2.3.14}) que les $D_+(f^{\natural})$ forment un
recouvrement de $\operatorname{Proj}(S^{\natural})$ et les $v^{-1}(X_f)$ un
recouvrement de $C_X$, puisque les $X_f$ forment un recouvrement de $X$~;
cela achève donc de définir dans ce cas l'isomorphisme
(\hyperref[II.8.7.3.1-fr]{8.7.3.1}).

Pour démontrer (\hyperref[II.8.7.3-fr]{8.7.3}) dans le cas général, il suffit
de voir que si $U$, $U'$ sont deux ouverts affines de $Y$ tels que
$U'\subset U$, d'anneaux $A$ et $A'$, et si on pose
$\mathcal{S}|U=\widetilde{S}$, $\mathcal{S}|U'=\widetilde{S'}$, le diagramme
\begin{equation}
\label{II.8.7.3.3-fr}
  \xymatrix{
    C_{U'}\ar[r]\ar[d] & \operatorname{Proj}(S'^{\natural})\ar[d]\\
    C_U\ar[r] & \operatorname{Proj}(S^{\natural})
  }
\tag{8.7.3.3}
\end{equation}
est commutatif. Mais $S'$ s'identifie canoniquement à $S\otimes_A A'$, donc
$S'^{\natural}$ à
\[
  S^{\natural}\otimes_S S'=S^{\natural}\otimes_A A'~;
\]
on a donc
$\operatorname{Proj}(S'^{\natural})=\operatorname{Proj}(S^{\natural})
\times_U U'$ (\hyperref[II.2.8.10-fr]{2.8.10})~; de même, si
$X=\operatorname{Proj}(S)$, $X'=\operatorname{Proj}(S')$, on a
$X'=X\times_U U'$ et
$\mathcal{S}_{X'}=\mathcal{S}_X\otimes_{\mathcal{O}_U}\mathcal{O}_{U'}$
(\hyperref[II.3.5.4-fr]{3.5.4}), autrement dit
$\mathcal{S}_{X'}=j^*(\mathcal{S}_X)$, où $j$ est la projection $X'\to X$.
On a par suite (\hyperref[II.1.5.2-fr]{1.5.2})
$C_{U'}=C_U\times_X X'=C_U\times_U U'$, et la commutativité de
(\hyperref[II.8.7.3.3-fr]{8.7.3.3}) est alors immédiate.
\end{proof}

\begin{remark}[8.7.4]
\label{II.8.7.4-fr}
\begin{enumerate}
  \item[\rm{(i)}] La fin du raisonnement de
  (\hyperref[II.8.7.3-fr]{8.7.3}) se généralise aussitôt de la façon
  suivante. Soient $g:Y'\to Y$ un morphisme,
  $\mathcal{S}'=g^*(\mathcal{S})$, $X'=\operatorname{Proj}(\mathcal{S}')$~;
  on a alors un diagramme commutatif
  \begin{equation}
  \label{II.8.7.4.1-fr}
    \xymatrix{
      C_{X'}\ar[r]\ar[d] &
      \operatorname{Proj}(\mathcal{S}'^{\natural})\ar[d]\\
      C_X\ar[r] & \operatorname{Proj}(\mathcal{S}^{\natural})
    }
  \tag{8.7.4.1}
  \end{equation}

  D'autre part, soit $\varphi:\mathcal{S}''\to\mathcal{S}$ un homomorphisme
  de $\mathcal{O}_Y$-Algèbres graduées, tel que si on pose
  $X''=\operatorname{Proj}(\mathcal{S}'')$,
  $u=\operatorname{Proj}(\varphi):X\to X''$ soit partout défini~; on a
  d'autre part

  \oldpage[II]{176}
  un $Y$-morphisme $v:C\to C''$ (avec
  $C''=\operatorname{Spec}(\mathcal{S}'')$) tel que
  $\mathcal{A}(v)=\varphi$, et comme $\varphi$ est un homomorphisme
  d'Algèbres graduées, on déduit de $\varphi$ un $v$-morphisme d'Algèbres
  graduées
  $\psi:\mathcal{S}''^{\natural}\to\mathcal{S}^{\natural}$
  (\hyperref[II.1.4.1-fr]{1.4.1}). En outre, il résulte aisément de
  (\hyperref[II.8.2.10-fr]{8.2.10, (iv)}) et de l'hypothèse sur $\varphi$,
  que $\operatorname{Proj}(\psi)$ est partout défini. Enfin, compte tenu de
  (\hyperref[II.3.5.6.1-fr]{3.5.6.1}), on a un $u$-morphisme canonique
  $\mathcal{S}_{X''}\to\mathcal{S}_X$, d'où
  (\hyperref[II.1.5.6-fr]{1.5.6}) un morphisme $w:C_{X''}\to C_X$. Cela
  étant, le diagramme
  \begin{equation}
  \label{II.8.7.4.2-fr}
    \xymatrix{
      C_{X''}\ar[r]^{\sim}\ar[d]_{w} &
      \operatorname{Proj}(\mathcal{S}''^{\natural})
        \ar[d]^{\operatorname{Proj}(\psi)}\\
      C_X\ar[r]^{\sim} & \operatorname{Proj}(\mathcal{S}^{\natural})
    }
  \tag{8.7.4.2}
  \end{equation}
  est commutatif, comme on le vérifie aussitôt en se ramenant au cas où $Y$
  est affine.

  \item[\rm{(ii)}] Notons qu'en vertu de
  (\hyperref[II.8.7.2.5-fr]{8.7.2.5}) et
  (\hyperref[II.8.7.2.6-fr]{8.7.2.6}), on a
  $\mathcal{I}_1^m\subset\mathcal{I}_m\subset\mathcal{I}_1$ pour tout $m>0$.
  Or, par définition, $\mathcal{I}_1=(\mathcal{S}_+)^{\sim}$, donc
  $\mathcal{I}_1$ définit dans $C$ le sous-préschéma fermé
  $\varepsilon(Y)$ ((\hyperref[II.1.4.10-fr]{1.4.10}) et
  (\hyperref[II.8.3.2-fr]{8.3.2}))~; on en conclut que pour tout $m>0$,
  \emph{le support de $\mathcal{O}_C/\mathcal{I}_m$ est contenu dans l'espace
  sous-jacent au préschéma des sommets $\varepsilon(Y)$}~; dans l'image
  réciproque du cône épointé affine $E$, le morphisme structural
  $\operatorname{Proj}(\mathcal{S}^{\natural})\to C$ se réduit donc à un
  \emph{isomorphisme} (comme il résulte d'ailleurs de
  (\hyperref[II.8.7.3-fr]{8.7.3}) et
  (\hyperref[II.8.7.1-fr]{8.7.1})). En outre, identifiant canoniquement $C$ à
  un ouvert de $\widehat{C}$ (\hyperref[II.8.3.3-fr]{8.3.3}), on peut
  évidemment prolonger les Idéaux $\mathcal{I}_m$ de $\mathcal{O}_C$ en des
  Idéaux $\mathcal{J}_m$ de $\mathcal{O}_{\widehat{C}}$, par la condition de
  coïncider avec $\mathcal{O}_{\widehat{C}}$ dans l'ouvert $\widehat{E}$ de
  $\widehat{C}$. Si on pose
  $\mathcal{T}=\bigoplus_{n\geqslant 0}\mathcal{J}_m$, qui est une
  $\mathcal{O}_{\widehat{C}}$-Algèbre graduée quasi-cohérente, on peut
  prolonger l'isomorphisme (\hyperref[II.8.7.3.1-fr]{8.7.3.1}) en un
  $\widehat{C}$-isomorphisme
  \begin{equation}
  \label{II.8.7.4.3-fr}
    \widehat{C}_X\xrightarrow{\sim}\operatorname{Proj}(\mathcal{T}).
  \tag{8.7.4.3}
  \end{equation}

  En effet, au-dessus de $\widehat{E}$, il résulte de ce qui précède que
  $\operatorname{Proj}(\mathcal{T})$ s'identifie canoniquement à
  $\widehat{E}$, et on définit donc l'isomorphisme
  (\hyperref[II.8.7.4.3-fr]{8.7.4.3}) au-dessus de $\widehat{E}$ en lui
  imposant de coïncider avec l'isomorphisme canonique
  $\widehat{E}_X\to\widehat{E}$ (\hyperref[II.8.6.2-fr]{8.6.2})~; il est clair
  qu'alors cet isomorphisme et (\hyperref[II.8.7.3.1-fr]{8.7.3.1}) coïncident
  au-dessus de $E$.
\end{enumerate}
\end{remark}

\begin{corollary}[8.7.5]
\label{II.8.7.5-fr}
Supposons qu'il existe $n_0>0$ tel que
\begin{equation}
\label{II.8.7.5.1-fr}
  \mathcal{S}_{n+1}=\mathcal{S}_1\mathcal{S}_n
  \qquad\text{pour }n\geqslant n_0.
\tag{8.7.5.1}
\end{equation}
Alors le sous-préschéma des sommets de $C_X$ (isomorphe à $X$) est l'image
réciproque par le morphisme canonique $r:C_X\to C$ du sous-préschéma des
sommets de $C$ (isomorphe à $Y$). Inversement, si cette propriété a lieu et
si on suppose de plus $Y$ noethérien et $\mathcal{S}$ de type fini, il existe
$n_0>0$ tel que l'on ait (\hyperref[II.8.7.5.1-fr]{8.7.5.1}).
\end{corollary}

\begin{proof}
La première assertion étant locale sur $Y$, on peut pour l'établir supposer
$Y=\operatorname{Spec}(A)$ affine, et $\mathcal{S}=\widetilde{S}$, où $S$ est
une $A$-algèbre graduée à degrés positifs. Elle résulte alors de
(\hyperref[II.8.2.12-fr]{8.2.12}), car on a
$\operatorname{Proj}(S^{\natural}\otimes_S S_0)=
C_X\times_C\varepsilon(Y)$ (en vertu de l'identification
(\hyperref[II.8.7.3.1-fr]{8.7.3.1})), c'est-à-dire que ce préschéma est
l'image réciproque de $\varepsilon(Y)$ dans $C_X$
(\hyperref[I.4.4.1-fr]{I, 4.4.1}). La réciproque résulte aussi de
(\hyperref[II.8.2.12-fr]{8.2.12}) lorsque $Y$ est affine noethérien et $S$ de
type fini.

\oldpage[II]{177}
Si $Y$ est noethérien (non nécessairement affine) et $\mathcal{S}$ de type
fini, il y a un recouvrement fini de $Y$ par des ouverts affines noethériens
$U_i$, et on déduit alors de ce qui précède que pour tout $i$, il y a un
entier $n_i$ tel que
$\mathcal{S}_{n+1}|U_i=(\mathcal{S}_1|U_i)(\mathcal{S}_n|U_i)$ dès que
$n\geqslant n_i$~; le plus grand des $n_i$ vérifie donc
(\hyperref[II.8.7.5.1-fr]{8.7.5.1}).
\end{proof}

\begin{env}[8.7.6]
\label{II.8.7.6-fr}
Considérons maintenant le $C$-préschéma $Z$ obtenu en faisant \emph{éclater}
dans le cône affine $C$ le \emph{sous-préschéma des sommets}
$\varepsilon(Y)$~; par définition (\hyperref[II.8.1.3-fr]{8.1.3}) c'est donc
le préschéma
$\operatorname{Proj}(\bigoplus_{n\geqslant 0}\mathcal{S}_+^n)$~; l'injection
canonique
\begin{equation}
\label{II.8.7.6.1-fr}
  \iota:\bigoplus_{n\geqslant 0}\mathcal{S}_+^n
  \longrightarrow\mathcal{S}^{\natural}
\tag{8.7.6.1}
\end{equation}
définit (grâce à l'identification (\hyperref[II.8.7.3-fr]{8.7.3})) un
$C$-morphisme dominant canonique
\begin{equation}
\label{II.8.7.6.2-fr}
  G(\iota)\longrightarrow Z
\tag{8.7.6.2}
\end{equation}
où $G(\iota)$ est un ouvert de $C_X$ (\hyperref[II.3.5.1-fr]{3.5.1})~; on
notera qu'on peut avoir $G(\iota)\ne C_X$, comme le montre l'exemple où
$Y=\operatorname{Spec}(K)$, $K$ étant un corps, $\mathcal{S}=\widetilde{S}$
avec $S=K[\mathbf{y}]$, où $\mathbf{y}$ est une indéterminée \emph{de degré
2}~; si $R_n$ désigne l'ensemble $(S_+)^n$, considéré comme partie de
$S_{[n]}=S_n^{\natural}$, $S_+^{\natural}$ n'est pas la racine dans
$S^{\natural}$ de l'idéal engendré par la réunion des $R_n$ (cf.
(\hyperref[II.2.3.14-fr]{2.3.14})).
\end{env}

\begin{corollary}[8.7.7]
\label{II.8.7.7-fr}
Supposons qu'il existe $n_0>0$ tel que
\begin{equation}
\label{II.8.7.7.1-fr}
  \mathcal{S}_n=\mathcal{S}_1^n
  \qquad\text{pour }n\geqslant n_0.
\tag{8.7.7.1}
\end{equation}
Alors le morphisme canonique (\hyperref[II.8.7.6.2-fr]{8.7.6.2}) est partout
défini et est un isomorphisme $C_X\xrightarrow{\sim}Z$. Inversement, si cette
propriété a lieu et si on suppose de plus $Y$ noethérien et $\mathcal{S}$ de
type fini, il existe $n_0$ tel que l'on ait
(\hyperref[II.8.7.7.1-fr]{8.7.7.1}).
\end{corollary}

\begin{proof}
En effet, la première assertion est locale sur $Y$, et découle donc de
(\hyperref[II.8.2.14-fr]{8.2.14})~; il en est de même de la réciproque, en
raisonnant comme dans (\hyperref[II.8.7.5-fr]{8.7.5}).
\end{proof}

\begin{remark}[8.7.8]
\label{II.8.7.8-fr}
Comme la condition (\hyperref[II.8.7.7.1-fr]{8.7.7.1}) implique
(\hyperref[II.8.7.5.1-fr]{8.7.5.1}), on voit que lorsqu'elle est vérifiée,
non seulement $C_X$ s'identifie au préschéma obtenu en faisant éclater le
sommet (identifié à $Y$) du cône affine $C$, mais encore le sommet (identifié
à $X$) de $C_X$ s'identifie au sous-préschéma fermé image réciproque du
sommet $Y$ de $C$. En outre, l'hypothèse
(\hyperref[II.8.7.7.1-fr]{8.7.7.1}) implique que sur
$X=\operatorname{Proj}(\mathcal{S})$, les $\mathcal{O}_X$-Modules
$\mathcal{O}_X(n)$ sont inversibles (\hyperref[II.3.2.5-fr]{3.2.5} et
\hyperref[II.3.2.9-fr]{3.2.9}) et que l'on a
$\mathcal{O}_X(n)=\mathcal{L}^{\otimes n}$ avec
$\mathcal{L}=\mathcal{O}_X(1)$ (\hyperref[II.3.2.7-fr]{3.2.7} et
\hyperref[II.3.2.9-fr]{3.2.9})~; par définition
(\hyperref[II.8.6.1.1-fr]{8.6.1.1}), $C_X$ est donc le \emph{fibré vectoriel}
$\mathbf{V}(\mathcal{L})$ sur $X$, et son sommet la \emph{section nulle} de ce
fibré vectoriel.
\end{remark}

\subsection{Faisceaux amples et contractions}
\label{subsection:II.8.8-fr}

\begin{env}[8.8.1]
\label{II.8.8.1-fr}
Soient $Y$ un préschéma, $f:X\to Y$ un morphisme \emph{séparé et
quasi-compact}, $\mathcal{L}$ un $\mathcal{O}_X$-Module inversible \emph{ample
relativement à $f$}. Considérons la $\mathcal{O}_Y$-Algèbre graduée à degrés
positifs
\begin{equation}
\label{II.8.8.1.1-fr}
  \mathcal{S}=\mathcal{O}_Y\oplus
  \bigoplus_{n\geqslant 1}f_*(\mathcal{L}^{\otimes n})
\tag{8.8.1.1}
\end{equation}

\oldpage[II]{178}
qui est quasi-cohérente (\hyperref[I.9.2.2-fr]{I, 9.2.2, a)}). On a un
homomorphisme canonique de $\mathcal{O}_X$-Algèbres graduées
\begin{equation}
\label{II.8.8.1.2-fr}
  \tau:f^*(\mathcal{S})\longrightarrow
  \bigoplus_{n\geqslant 0}\mathcal{L}^{\otimes n}
\tag{8.8.1.2}
\end{equation}
qui, pour les degrés $\geqslant 1$, coïncide avec l'homomorphisme canonique
$\sigma:f^*(f_*(\mathcal{L}^{\otimes n}))\to\mathcal{L}^{\otimes n}$
(\hyperref[0.4.4.3-fr]{0, 4.4.3}), et pour le degré $0$ est l'identité.
L'hypothèse que $\mathcal{L}$ est $f$-ample entraîne alors
(\hyperref[II.4.6.3-fr]{4.6.3} et \hyperref[II.3.6.1-fr]{3.6.1}) que le
$Y$-morphisme correspondant
\begin{equation}
\label{II.8.8.1.3-fr}
  r=r_{\mathcal{L},\tau}:X\longrightarrow
  P=\operatorname{Proj}(\mathcal{S})
\tag{8.8.1.3}
\end{equation}
est partout défini et est une \emph{immersion ouverte dominante}, et que l'on
a
\begin{equation}
\label{II.8.8.1.4-fr}
  r^*(\mathcal{O}_P(n))=\mathcal{L}^{\otimes n}
  \qquad\text{pour tout }n\in\mathbf{Z}.
\tag{8.8.1.4}
\end{equation}
\end{env}

\begin{proposition}[8.8.2]
\label{II.8.8.2-fr}
Soit $C=\operatorname{Spec}(\mathcal{S})$ \emph{le cône affine défini par
$\mathcal{S}$}~; si $\mathcal{L}$ est $f$-ample, il existe un $Y$-morphisme
canonique
\begin{equation}
\label{II.8.8.2.1-fr}
  g:V=\mathbf{V}(\mathcal{L})\longrightarrow C
\tag{8.8.2.1}
\end{equation}
tel que le diagramme
\begin{equation}
\label{II.8.8.2.2-fr}
  \xymatrix{
    X\ar[r]^{j}\ar[d]_{f} &
    \mathbf{V}(\mathcal{L})\ar[r]^{\pi}\ar[d]^{g} & X\ar[d]^{f}\\
    Y\ar[r]_{\varepsilon} & C\ar[r]_{\psi} & Y
  }
\tag{8.8.2.2}
\end{equation}
soit commutatif, $\psi$ et $\pi$ étant les morphismes structuraux, $j$ et
$\varepsilon$ les immersions canoniques appliquant respectivement $X$ et $Y$
sur la section nulle de $\mathbf{V}(\mathcal{L})$ et le préschéma des sommets
de $C$. En outre, la restriction de $g$ à
$\mathbf{V}(\mathcal{L})-j(X)$ est une immersion ouverte
\begin{equation}
\label{II.8.8.2.3-fr}
  \mathbf{V}(\mathcal{L})-j(X)\longrightarrow
  E=C-\varepsilon(Y)
\tag{8.8.2.3}
\end{equation}
dans le cône épointé affine $E$ correspondant à $\mathcal{S}$.
\end{proposition}

\begin{proof}
Avec les notations de (\hyperref[II.8.8.1-fr]{8.8.1}), soient
$\mathcal{S}_P^{\geqslant}=\bigoplus_{n\geqslant 0}\mathcal{O}_P(n)$, et
$C_P=\operatorname{Spec}(\mathcal{S}_P^{\geqslant})$. On sait
(\hyperref[II.8.6.2-fr]{8.6.2}) qu'on a un morphisme canonique
$h=\operatorname{Spec}(\alpha):C_P\to C$ tel que le diagramme
\begin{equation}
\label{II.8.8.2.4-fr}
  \xymatrix{
    C_P\ar[r]\ar[d]_{h} & P\ar[d]^{p}\\
    C\ar[r]_{\psi} & Y
  }
\tag{8.8.2.4}
\end{equation}
soit commutatif~; en outre, si $\varepsilon_P:P\to C_P$ est l'immersion
canonique, le diagramme
\begin{equation}
\label{II.8.8.2.5-fr}
  \xymatrix{
    P\ar[r]^{\varepsilon_P}\ar[d]_{p} & C_P\ar[d]^{h}\\
    Y\ar[r]_{\varepsilon} & C
  }
\tag{8.8.2.5}
\end{equation}
est commutatif (\hyperref[II.8.7.1.1-fr]{8.7.1.1}), et enfin, la restriction
de $h$ au cône épointé affine $E_P$ est un \emph{isomorphisme}
$E_P\xrightarrow{\sim}E$ (\hyperref[II.8.6.2-fr]{8.6.2}). D'autre part, il
résulte de (\hyperref[II.8.8.1.4-fr]{8.8.1.4}) que l'on a
\[
  r^*(\mathcal{S}_P^{\geqslant})=
  \mathbf{S}_{\mathcal{O}_X}(\mathcal{L})
\]

\oldpage[II]{179}
et par suite on a un $P$-morphisme canonique
$q:\mathbf{V}(\mathcal{L})\to C_P$, le diagramme commutatif
\begin{equation}
\label{II.8.8.2.6-fr}
  \xymatrix{
    \mathbf{V}(\mathcal{L})\ar[r]^{\pi}\ar[d]_{q} & X\ar[d]^{r}\\
    C_P\ar[r] & P
  }
\tag{8.8.2.6}
\end{equation}
identifiant $\mathbf{V}(\mathcal{L})$ au produit $C_P\times_P X$
(\hyperref[II.1.5.2-fr]{1.5.2})~; comme $r$ est une immersion ouverte, il en
est donc de même de $q$ (\hyperref[I.4.3.2-fr]{I, 4.3.2}). En outre, la
restriction de $q$ à $\mathbf{V}(\mathcal{L})-j(X)$ applique ce préschéma dans
$E_P$ par (\hyperref[II.8.5.2-fr]{8.5.2}), et le diagramme
\begin{equation}
\label{II.8.8.2.7-fr}
  \xymatrix{
    X\ar[r]^{j}\ar[d]_{r} & \mathbf{V}(\mathcal{L})\ar[d]^{q}\\
    P\ar[r]_{\varepsilon_P} & C_P
  }
\tag{8.8.2.7}
\end{equation}
est commutatif (cas particulier de
(\hyperref[II.8.5.1.3-fr]{8.5.1.3})). Les assertions de
(\hyperref[II.8.8.2-fr]{8.8.2}) résultent aussitôt de ces faits, en prenant
pour $g$ le morphisme composé $h\circ q$.
\end{proof}

\begin{remark}[8.8.3]
\label{II.8.8.3-fr}
Supposons en outre que $Y$ soit un préschéma \emph{noethérien} et $f$ un
morphisme \emph{propre}. Comme $r$ est alors \emph{propre}
(\hyperref[II.5.4.4-fr]{5.4.4}), donc fermé, et que c'est une immersion ouverte
dominante, $r$ est nécessairement un \emph{isomorphisme}
$X\xrightarrow{\sim}P$. On verra en outre au chapitre III
(\hyperref[III.2.3.5.1-fr]{III, 2.3.5.1}), que $\mathcal{S}$ est alors
nécessairement une $\mathcal{O}_Y$-Algèbre \emph{de type fini}. Alors il en
résulte que $\mathcal{S}^{\natural}$ est une
$\mathcal{S}_0^{\natural}$-Algèbre \emph{de type fini}
(\hyperref[II.8.2.10-fr]{8.2.10, (i)} et
\hyperref[II.8.7.2.7-fr]{8.7.2.7})~; comme $C_P$ est $C$-isomorphe à
$\operatorname{Proj}(\mathcal{S}^{\natural})$
(\hyperref[II.8.7.3-fr]{8.7.3}), on voit que le morphisme $h:C_P\to C$ est
\emph{projectif}~; comme le morphisme $r$ est un isomorphisme, il en est de
même de $q:\mathbf{V}(\mathcal{L})\to C_P$, et on en conclut que le morphisme
$g:\mathbf{V}(\mathcal{L})\to C$ est \emph{projectif}. En outre, comme la
restriction de $h$ à $E_P$ est un isomorphisme sur $E$ et que $q$ est un
isomorphisme, la restriction (\hyperref[II.8.8.2.3-fr]{8.8.2.3}) de $g$ est un
\emph{isomorphisme}
$\mathbf{V}(\mathcal{L})-j(X)\xrightarrow{\sim}E$.

Si on suppose de plus que $\mathcal{L}$ est \emph{très ample} pour $f$, on
verra aussi au chapitre III (\hyperref[III.2.3.5.1-fr]{III, 2.3.5.1}), qu'il
existe un entier $n_0>0$ tel que
$\mathcal{S}_n=\mathcal{S}_1^n$ pour $n\geqslant n_0$. On en conclura par
(\hyperref[II.8.7.7-fr]{8.7.7}) que $\mathbf{V}(\mathcal{L})$ s'identifie au
préschéma $Z$ obtenu en \emph{faisant éclater dans le cône affine $C$ le
sous-préschéma des sommets} (identifié à $Y$), et que la \emph{section nulle}
de $\mathbf{V}(\mathcal{L})$ (identifiée à $Y$) est \emph{l'image réciproque}
du sous-préschéma des sommets $Y$ de $C$. Une partie des résultats précédents
peut en fait s'établir sans hypothèse noethérienne~:
\end{remark}

\begin{corollary}[8.8.4]
\label{II.8.8.4-fr}
Soient $Y$ un préschéma (resp. un schéma quasi-compact), $f:X\to Y$ un
morphisme propre, $\mathcal{L}$ un $\mathcal{O}_X$-Module inversible ample pour
$f$. Alors le morphisme (\hyperref[II.8.8.2.1-fr]{8.8.2.1}) est propre (resp.
projectif) et sa restriction (\hyperref[II.8.8.2.3-fr]{8.8.2.3}) est un
isomorphisme.
\end{corollary}

\begin{proof}
Pour prouver que $g$ est propre, on peut se ramener au cas où $Y$ est affine,
et il suffit donc de considérer le cas où $Y$ est un schéma quasi-compact. Le
même raisonnement que dans (\hyperref[II.8.8.3-fr]{8.8.3}) montre d'abord que
$r$ est un \emph{isomorphisme} $X\xrightarrow{\sim}P$~; par suite $q$ est aussi
un isomorphisme, et comme la restriction de $h$ à $E_P$ est un isomorphisme
$E_P\xrightarrow{\sim}E$, on voit déjà que
(\hyperref[II.8.8.2.3-fr]{8.8.2.3}) est un isomorphisme. Reste à prouver que
$g$ est \emph{projectif}.

Comme $f$ est de type fini par hypothèse, on peut appliquer
(\hyperref[II.3.8.5-fr]{3.8.5}) à l'homomorphisme

\oldpage[II]{180}
$\tau$ de (\hyperref[II.8.8.1.2-fr]{8.8.1.2})~: il y a par suite un entier
$d>0$, un sous-$\mathcal{O}_Y$-Module quasi-cohérent de type fini
$\mathcal{E}$ de $\mathcal{S}_d$ tel que si $\mathcal{S}'$ est la
sous-$\mathcal{O}_Y$-Algèbre de $\mathcal{S}$ engendrée par $\mathcal{E}$ et
$\tau'=\tau\circ f^*(\varphi)$ ($\varphi$ injection canonique
$\mathcal{S}'\to\mathcal{S}$), $r'=r_{\mathcal{L},\tau'}$ soit une immersion
\[
  X\longrightarrow P'=\operatorname{Proj}(\mathcal{S}').
\]
En outre, comme $\varphi$ est injectif, $r'$ est aussi une \emph{immersion
dominante} (\hyperref[II.3.7.6-fr]{3.7.6})~; le même raisonnement que pour $r$
montre donc que $r'$ est une \emph{immersion fermée surjective}~; comme $r'$ se
factorise en
$X\xrightarrow{r}\operatorname{Proj}(\mathcal{S})
\xrightarrow{\Phi}\operatorname{Proj}(\mathcal{S}')$, où
$\Phi=\operatorname{Proj}(\varphi)$, on en conclut que $\Phi$ est aussi une
\emph{immersion fermée surjective}. Mais cela entraîne que $\Phi$ est un
\emph{isomorphisme}~; en effet, on peut se ramener au cas où
$Y=\operatorname{Spec}(A)$ est affine, $\mathcal{S}=\widetilde{S}$,
$\mathcal{S}'=\widetilde{S'}$, où $S$ est une $A$-algèbre graduée, $S'$ une
sous-algèbre graduée de $S$. Pour tout élément homogène $t\in S'$,
$S'_{(t)}$ est alors un sous-anneau de $S_{(t)}$~; si on revient à la
définition de $\operatorname{Proj}(\varphi)$
(\hyperref[II.2.8.1-fr]{2.8.1}), on voit qu'on est ramené à prouver que si
$B'$ est un sous-anneau d'un anneau $B$, et si le morphisme
$\operatorname{Spec}(B)\to\operatorname{Spec}(B')$ correspondant à l'injection
canonique $B'\to B$ est une immersion fermée, ce morphisme est nécessairement
un \emph{isomorphisme}~; mais cela résulte de
(\hyperref[I.4.2.3-fr]{I, 4.2.3}).

On a en outre
$\Phi^*(\mathcal{O}_{P'}(n))=\mathcal{O}_P(n)$
(\hyperref[II.3.5.2-fr]{3.5.2, (ii)} et
\hyperref[II.3.5.4-fr]{3.5.4}), donc
$r'^*(\mathcal{O}_{P'}(n))$ est isomorphe à $\mathcal{L}^{\otimes n}$
(\hyperref[II.4.6.3-fr]{4.6.3}). Posons
$\mathcal{S}''=\mathcal{S}'^{(d)}$, de sorte
(\hyperref[II.3.1.8-fr]{3.1.8, (i)}) que $X$ s'identifie canoniquement à
$P''=\operatorname{Proj}(\mathcal{S}'')$ et
$\mathcal{L}''=\mathcal{L}^{\otimes d}$ à $\mathcal{O}_{P''}(1)$
(\hyperref[II.3.2.9-fr]{3.2.9, (ii)}).

Cela étant, si $C''=\operatorname{Spec}(\mathcal{S}'')$,
$\mathcal{S}_{P''}^{\geqslant}=
\bigoplus_{n\geqslant 0}\mathcal{O}_{P''}(n)$ s'identifie à
$\bigoplus_{n\geqslant 0}\mathcal{L}''^{\otimes n}$, donc
$C_{P''}=\operatorname{Spec}(\mathcal{S}_{P''}^{\geqslant})$ à
$\mathbf{V}(\mathcal{L}'')$~; on sait d'autre part
(\hyperref[II.8.7.3-fr]{8.7.3}) que $C_{P''}$ est $C''$-isomorphe à
$\operatorname{Proj}(\mathcal{S}''^{\natural})$~; par définition de
$\mathcal{S}''$, $\mathcal{S}''^{\natural}$ est engendré par
$\mathcal{S}_1''^{\natural}$ et $\mathcal{S}_1''^{\natural}$ est de type fini
sur $\mathcal{S}_0''^{\natural}=\mathcal{S}''$
(\hyperref[II.8.2.10-fr]{8.2.10, (i) et (iii)}), donc
$\operatorname{Proj}(\mathcal{S}''^{\natural})$ est \emph{projectif} sur
$C''$ (\hyperref[II.5.5.1-fr]{5.5.1}). Considérons alors le diagramme
\begin{equation}
\label{II.8.8.4.1-fr}
  \xymatrix{
    \mathbf{V}(\mathcal{L})\ar[r]^{g}\ar[d]_{u} &
    \operatorname{Spec}(\mathcal{S})=C\ar[d]^{v}\\
    \mathbf{V}(\mathcal{L}'')\ar[r]_{g''} &
    \operatorname{Spec}(\mathcal{S}'')=C''
  }
\tag{8.8.4.1}
\end{equation}
où $g$ et $g''$ correspondent par (\hyperref[II.1.5.6-fr]{1.5.6}) aux
$f$-morphismes canoniques
\[
  \mathcal{S}\longrightarrow
  \bigoplus_{n\geqslant 0}\mathcal{L}^{\otimes n}
  \quad\text{et}\quad
  \mathcal{S}''\longrightarrow
  \bigoplus_{n\geqslant 0}\mathcal{L}''^{\otimes n}
\]
(\hyperref[II.3.3.2.3-fr]{3.3.2.3}) (voir
(\hyperref[II.8.8.5-fr]{8.8.5}) ci-dessous), $v$ au morphisme d'inclusion
$\mathcal{S}''\to\mathcal{S}$, et $u$ au morphisme d'inclusion
$\bigoplus_{n\geqslant 0}\mathcal{L}^{\otimes nd}\to
\bigoplus_{n\geqslant 0}\mathcal{L}^{\otimes n}$~; il est immédiat
(\hyperref[II.3.3.2-fr]{3.3.2}) que ce diagramme est commutatif. Nous venons
de voir que $g''$ est un morphisme projectif~; d'autre part, $u$ est un
morphisme \emph{fini}. En effet, la question étant locale sur $X$, on peut
supposer que $X$ est affine d'anneau $A$ et que $\mathcal{L}=\mathcal{O}_X$~;
tout revient alors à remarquer que l'anneau $A[T]$ est un module de type fini
sur son sous-anneau $A[T^d]$ ($T$ indéterminée). Comme $Y$ est un schéma
quasi-compact, et que $C''$ est affine sur $Y$, $C''$ est aussi un schéma
quasi-compact,

\oldpage[II]{181}
et par suite $g''\circ u$ est un morphisme projectif
(\hyperref[II.5.5.5-fr]{5.5.5, (ii)})~; par la commutativité de
(\hyperref[II.8.8.4.1-fr]{8.8.4.1}), il en est de même de $v\circ g$, et comme
$v$ est affine, donc séparé, on en conclut finalement que $g$ est projectif
(\hyperref[II.5.5.5-fr]{5.5.5, (v)}).
\end{proof}

\begin{env}[8.8.5]
\label{II.8.8.5-fr}
Reprenons la situation de (\hyperref[II.8.8.1-fr]{8.8.1}). Nous allons voir
qu'on peut donner du morphisme $g:\mathbf{V}(\mathcal{L})\to C$ une autre
définition valable pour tout $\mathcal{O}_X$-Module inversible $\mathcal{L}$
(non nécessairement ample). Pour cela, considérons le $f$-morphisme
\begin{equation}
\label{II.8.8.5.1-fr}
  \tau^\flat:\mathcal{S}\longrightarrow
  \bigoplus_{n\geqslant 0}\mathcal{L}^{\otimes n}
\tag{8.8.5.1}
\end{equation}
correspondant au morphisme $\tau$ de
(\hyperref[II.8.8.1.2-fr]{8.8.1.2}). On en déduit
(\hyperref[II.1.5.6-fr]{1.5.6}) un morphisme $g':V\to C$ tel que, si
$\pi:V\to X$ et $\psi:C\to Y$ sont les morphismes structuraux, les diagrammes
\begin{equation}
\label{II.8.8.5.2-fr}
  \xymatrix{
    X\ar[d]_{f} & V\ar[l]^{\pi}\ar[d]^{g'}\\
    Y & C\ar[l]_{\psi}
  }
  \qquad
  \xymatrix{
    X\ar[r]^{j}\ar[d]_{f} & V\ar[d]^{g'}\\
    Y\ar[r]_{\varepsilon} & C
  }
\tag{8.8.5.2}
\end{equation}
soient commutatifs
(\hyperref[II.8.5.1.2-fr]{8.5.1.2} et
\hyperref[II.8.5.1.3-fr]{8.5.1.3}). Nous allons montrer que (si l'on suppose
$\mathcal{L}$ ample pour $f$) les morphismes $g$ et $g'$ sont identiques.

En effet, la question est locale sur $Y$, et on peut donc supposer que
$Y=\operatorname{Spec}(A)$ est affine, et (en raison de
(\hyperref[II.8.8.1.3-fr]{8.8.1.3})) identifier $X$ à un ouvert de
$P=\operatorname{Proj}(S)$, où
\[
  S=A\oplus\bigoplus_{n\geqslant 1}
  \Gamma(X,\mathcal{L}^{\otimes n})~;
\]
on déduit alors de (\hyperref[II.8.8.1.4-fr]{8.8.1.4}) que
$\Gamma(X,\mathcal{O}_P(n))=\Gamma(X,\mathcal{L}^{\otimes n})$ pour tout
$n\in\mathbf{Z}$. Tenant compte de la définition de
$h=\operatorname{Spec}(\alpha)$, où $\alpha$ est le $p$-morphisme canonique
$\widetilde{S}\to\mathcal{S}_P^{\geqslant}$
(\hyperref[II.8.6.1.2-fr]{8.6.1.2}), il faut vérifier que la restriction à
$X$ de
$\alpha^\sharp:p^*(\widetilde{S})\to\mathcal{S}_P^{\geqslant}$ est
identique à $\tau$. Compte tenu de
(\hyperref[0.4.4.3-fr]{0, 4.4.3}), on est ramené à voir que si on compose
l'homomorphisme canonique
$\alpha_n:S_n\to\Gamma(P,\mathcal{O}_P(n))$ avec l'homomorphisme de
restriction
\[
  \Gamma(P,\mathcal{O}_P(n))\longrightarrow
  \Gamma(X,\mathcal{O}_P(n))=
  \Gamma(X,\mathcal{L}^{\otimes n}),
\]
on obtient l'identité pour tout $n>0$~; or, cela résulte aussitôt de la
définition de l'algèbre $S$ et de celle de $\alpha_n$
(\hyperref[II.2.6.2-fr]{2.6.2}).
\end{env}

\begin{proposition}[8.8.6]
\label{II.8.8.6-fr}
Supposons (avec les notations de (\hyperref[II.8.8.5-fr]{8.8.5})) que si
l'on pose $f=(f_0,\lambda)$, l'homomorphisme
$\lambda:\mathcal{O}_Y\to f_*(\mathcal{O}_X)$ soit bijectif~; alors~:
\begin{enumerate}
  \item[\rm{(i)}] Si on pose $g=(g_0,\mu)$,
  $\mu:\mathcal{O}_C\to g_*(\mathcal{O}_V)$ est un isomorphisme.
  \item[\rm{(ii)}] Si $X$ est intègre (resp. localement intègre et normal),
  $C$ est intègre (resp. normal).
\end{enumerate}
\end{proposition}

\begin{proof}
En effet, le $f$-morphisme $\tau^\flat$ est alors un \emph{isomorphisme}
\[
  \tau^\flat:\mathcal{S}=\psi_*(\mathcal{O}_C)
  \longrightarrow f_*(\pi_*(\mathcal{O}_V))
  =\psi_*(g_*(\mathcal{O}_V))
\]
et le $Y$-morphisme $g$ peut être considéré comme celui pour lequel
l'homomorphisme $\mathcal{A}(g)$
(\hyperref[II.1.1.2-fr]{1.1.2}) est égal à $\tau^\flat$. Pour voir que
$\mu$ est un isomorphisme de $\mathcal{O}_C$-Modules, il suffit
(\hyperref[II.1.4.2-fr]{1.4.2}) de voir que
\[
  \mathcal{A}(\mu):\psi_*(\mathcal{O}_C)
  \longrightarrow\psi_*(g_*(\mathcal{O}_V))
\]
est un isomorphisme. Mais par définition
(\hyperref[II.1.1.2-fr]{1.1.2}), on a
$\mathcal{A}(\mu)=\mathcal{A}(g)$, d'où la conclusion de (i).

Pour prouver (ii), on peut se ramener au cas où $Y$ est affine, donc
$\mathcal{S}=\widetilde{S}$, avec

\oldpage[II]{182}
$S=\bigoplus_{n\geqslant 0}\Gamma(X,\mathcal{L}^{\otimes n})$~;
l'hypothèse que $X$ est intègre entraîne que l'anneau $S$ est intègre
(\hyperref[I.7.4.4-fr]{I, 7.4.4}), donc il en est de même de $C$
(\hyperref[I.5.1.4-fr]{I, 5.1.4}). Pour démontrer que $C$ est normal, nous
utiliserons le lemme suivant~:

\begin{lemma}[8.8.6.1]
\label{II.8.8.6.1-fr}
Soit $Z$ un préschéma intègre normal. Alors l'anneau
$\Gamma(Z,\mathcal{O}_Z)$ est intègre et intégralement clos.
\end{lemma}

\begin{proof}
En effet, il résulte de
(\hyperref[I.8.2.1.1-fr]{I, 8.2.1.1}) que
$\Gamma(Z,\mathcal{O}_Z)$ est l'intersection, dans le corps des fonctions
rationnelles $R(Z)$, des anneaux intégralement clos $\mathcal{O}_z$ pour
$z\in Z$.
\end{proof}

Cela étant, montrons d'abord que $V$ est \emph{localement intègre et normal}~;
on peut pour cela se borner au cas où $X=\operatorname{Spec}(A)$ est affine,
d'anneau $A$ intègre et intégralement clos
(\hyperref[II.6.3.8-fr]{6.3.8}), et où
$\mathcal{L}=\mathcal{O}_X$. Comme alors
$V=\operatorname{Spec}(A[T])$ et que $A[T]$ est intègre et intégralement clos
([24], p.~99), cela établit notre assertion. D'autre part, pour tout ouvert
affine $U$ de $C$, $g^{-1}(U)$ est quasi-compact puisque le morphisme $g$ est
quasi-compact~; puisque $V$ est localement intègre, les composantes connexes
de $g^{-1}(U)$ sont des préschémas intègres ouverts dans $g^{-1}(U)$, donc en
nombre fini, et comme $V$ est normal, ces préschémas sont normaux
(\hyperref[II.6.3.8-fr]{6.3.8}). Alors $\Gamma(U,\mathcal{O}_C)$, qui est égal
à $\Gamma(g^{-1}(U),\mathcal{O}_V)$ en vertu de (i), est composé direct d'un
nombre fini d'anneaux intègres et intégralement clos
(\hyperref[II.8.8.6.1-fr]{8.8.6.1}), ce qui montre que $C$ est normal
(\hyperref[II.6.3.4-fr]{6.3.4}).
\end{proof}

\subsection{Le critère d'amplitude de Grauert : énoncé}
\label{subsection:II.8.9-fr}

Nous nous proposons de montrer que les propriétés démontrées dans
(\hyperref[II.8.8.2-fr]{8.8.2}) \emph{caractérisent} les
$\mathcal{O}_X$-Modules amples pour $f$, et de façon précise, de démontrer le
critère suivant~:

\begin{theorem}[8.9.1]
\label{II.8.9.1-fr}
\emph{(critère de Grauert).} Soient $Y$ un préschéma,
$p:X\to Y$ un morphisme séparé et quasi-compact, $\mathcal{L}$ un
$\mathcal{O}_X$-Module inversible. Pour que $\mathcal{L}$ soit ample
relativement à $p$, il faut et il suffit qu'il existe un $Y$-préschéma $C$,
une $Y$-section $\varepsilon:Y\to C$ de $C$, et un $Y$-morphisme
$q:\mathbf{V}(\mathcal{L})\to C$, ayant les propriétés suivantes~:
\begin{enumerate}
  \item[\rm{(i)}] Le diagramme
  \begin{equation}
  \label{II.8.9.1.1-fr}
    \xymatrix{
      X\ar[r]^{j}\ar[d]_{p} & \mathbf{V}(\mathcal{L})\ar[d]^{q}\\
      Y\ar[r]_{\varepsilon} & C
    }
  \tag{8.9.1.1}
  \end{equation}
  où $j$ est la section nulle du fibré vectoriel $\mathbf{V}(\mathcal{L})$,
  est commutatif.
  \item[\rm{(ii)}] La restriction de $q$ à
  $\mathbf{V}(\mathcal{L})-j(X)$ est une immersion ouverte quasi-compacte
  \[
    \mathbf{V}(\mathcal{L})-j(X)\longrightarrow C
  \]
  dont l'image ne rencontre pas $\varepsilon(Y)$.
\end{enumerate}
\end{theorem}

On notera que si $C$ est \emph{séparé} sur $Y$, on peut, dans la condition
(ii), supprimer l'hypothèse que l'immersion ouverte est quasi-compacte~; pour
voir en effet que cette dernière propriété est conséquence des autres
conditions, on peut se borner au cas où $Y$ est affine, et l'assertion découle
alors de (\hyperref[I.5.5.1-fr]{I, 5.5.1, (i)} et
\hyperref[I.5.5.10-fr]{5.5.10}). On peut aussi supprimer

\oldpage[II]{183}
la même hypothèse lorsque l'on suppose $X$ noethérien, car alors $V$ est aussi
noethérien et l'assertion découle cette fois de
(\hyperref[I.6.3.5-fr]{I, 6.3.5}).

\begin{corollary}[8.9.2]
\label{II.8.9.2-fr}
Si le morphisme $p:X\to Y$ est propre, on peut dans l'énoncé de
(\hyperref[II.8.9.1-fr]{8.9.1}), supposer $q$ propre, et remplacer
«~immersion ouverte~» par «~isomorphisme~».
\end{corollary}

De façon imagée, on peut dire (lorsque $p:X\to Y$ est propre) que
$\mathcal{L}$ est ample relativement à $p$ si et seulement si on peut
«~contracter~» la section nulle du fibré vectoriel $\mathbf{V}(\mathcal{L})$
en le préschéma de base $Y$. Un cas particulier important est celui où $Y$
est le spectre d'un corps, et où l'opération de «~contraction~» consiste donc
à contracter la section nulle de $\mathbf{V}(\mathcal{L})$ \emph{en un seul
point}.

\begin{env}[8.9.3]
\label{II.8.9.3-fr}
La nécessité des conditions de (\hyperref[II.8.9.1-fr]{8.9.1}) et le
corollaire (\hyperref[II.8.9.2-fr]{8.9.2}) découlent aussitôt de
(\hyperref[II.8.8.2-fr]{8.8.2}) et
(\hyperref[II.8.8.4-fr]{8.8.4}).

Pour démontrer la suffisance de (\hyperref[II.8.9.1-fr]{8.9.1}), nous
considérerons une situation un peu plus générale. Pour cela, posons (avec les
notations de (\hyperref[II.8.8.2-fr]{8.8.2}))
\[
  \mathcal{S}'=\bigoplus_{n\geqslant 0}\mathcal{L}^{\otimes n}
\]
et
\[
  V=\mathbf{V}(\mathcal{L})=\operatorname{Spec}(\mathcal{S}').
\]
Le sous-préschéma fermé $j(X)$, section nulle de $\mathbf{V}(\mathcal{L})$,
est défini par le faisceau quasi-cohérent d'idéaux
$\mathcal{J}=(\mathcal{S}'_+)^{\sim}$ de $\mathcal{O}_V$
(\hyperref[II.1.4.10-fr]{1.4.10}). Ce $\mathcal{O}_V$-Module est
\emph{inversible}, car la question est locale sur $X$ et cela revient à
remarquer que l'idéal $TA[T]$ dans un anneau de polynômes $A[T]$ est un
$A[T]$-module libre monogène. En outre, il est immédiat (toujours parce que la
question est locale sur $X$) que
\[
  \mathcal{L}=j^*(\mathcal{J})
\]
et
\[
  j_*(\mathcal{L})=\mathcal{J}/\mathcal{J}^2.
\]
D'autre part, si
\[
  \pi:\mathbf{V}(\mathcal{L})\longrightarrow X
\]
est le morphisme structural, on a
$\pi_*(\mathcal{J})=\mathcal{S}'_+$ et
$\pi_*(\mathcal{J}/\mathcal{J}^2)=\mathcal{L}$~; on a donc deux
homomorphismes canoniques
$\mathcal{L}\to\pi_*(\mathcal{J})\to\mathcal{L}$, le premier étant
l'injection canonique $\mathcal{L}\to\mathcal{S}'_+$, le second la projection
canonique de $\mathcal{S}'_+$ sur $\mathcal{S}'_1=\mathcal{L}$, et leur
composé est l'identité. On peut d'ailleurs plonger canoniquement
\[
  \pi_*(\mathcal{J})=\mathcal{S}'_+
  =\bigoplus_{n\geqslant 1}\mathcal{L}^{\otimes n}
\]
dans le \emph{produit}
\[
  \prod_{n\geqslant 1}\mathcal{L}^{\otimes n}
  =\varprojlim_n\pi_*(\mathcal{J}/\mathcal{J}^{n+1})
\]
(puisque
$\pi_*(\mathcal{J}/\mathcal{J}^{n+1})=
\mathcal{L}\oplus\mathcal{L}^{\otimes2}\oplus\cdots\oplus
\mathcal{L}^{\otimes n}$), et on a donc encore deux homomorphismes canoniques
\begin{equation}
\label{II.8.9.3.1-fr}
  \mathcal{L}\longrightarrow
  \varprojlim_n\pi_*(\mathcal{J}/\mathcal{J}^{n+1})
  \longrightarrow\mathcal{L}
\tag{8.9.3.1}
\end{equation}
dont le composé est l'identité.

Cela étant, la généralisation de (\hyperref[II.8.9.1-fr]{8.9.1}) que nous
allons prouver est la suivante~:
\end{env}

\begin{proposition}[8.9.4]
\label{II.8.9.4-fr}
Soient $Y$ un préschéma, $V$ un $Y$-préschéma, $X$ un sous-préschéma fermé de
$V$, défini par un Idéal $\mathcal{J}$ de $\mathcal{O}_V$, qui est un
$\mathcal{O}_V$-Module inversible~; si $j:X\to V$ est

\oldpage[II]{184}
l'injection canonique, on pose
$\mathcal{L}=j^*(\mathcal{J})=
\mathcal{J}\otimes_{\mathcal{O}_V}\mathcal{O}_X$, de sorte que
$j_*(\mathcal{L})=\mathcal{J}/\mathcal{J}^2$. On suppose que le morphisme
structural $p:X\to Y$ est séparé et quasi-compact, et que les conditions
suivantes sont remplies~:
\begin{enumerate}
  \item[\rm{(i)}] Il existe un $Y$-morphisme de type fini $\pi:V\to X$ tel
  que $\pi\circ j=1_X$, et par suite
  $\pi_*(\mathcal{J}/\mathcal{J}^2)=\mathcal{L}$.
  \item[\rm{(ii)}] Il existe un homomorphisme de $\mathcal{O}_X$-Modules
  \[
    \varphi:\mathcal{L}\longrightarrow
    \varprojlim_n\pi_*(\mathcal{J}/\mathcal{J}^{n+1})
  \]
  tel que le composé
  \[
    \mathcal{L}\xrightarrow{\varphi}
    \varprojlim_n\pi_*(\mathcal{J}/\mathcal{J}^{n+1})
    \xrightarrow{\alpha}\pi_*(\mathcal{J}/\mathcal{J}^2)=\mathcal{L}
  \]
  (où $\alpha$ est l'homomorphisme canonique) soit l'identité.
  \item[\rm{(iii)}] Il existe un $Y$-préschéma $C$, une $Y$-section
  $\varepsilon$ de $C$ et un $Y$-morphisme $q:V\to C$ tels que le diagramme
  \begin{equation}
  \label{II.8.9.4.1-fr}
    \xymatrix{
      X\ar[r]^{j}\ar[d]_{p} & V\ar[d]^{q}\\
      Y\ar[r]_{\varepsilon} & C
    }
  \tag{8.9.4.1}
  \end{equation}
  soit commutatif.
  \item[\rm{(iv)}] La restriction de $q$ à $W=V-j(X)$ est une immersion
  ouverte quasi-compacte dans $C$, dont l'image ne rencontre pas
  $\varepsilon(Y)$.
\end{enumerate}
Dans ces conditions, $\mathcal{L}$ est ample relativement à $p$.
\end{proposition}

\subsection{Le critère d'amplitude de Grauert : démonstration}
\label{subsection:II.8.10-fr}

\begin{lemma}[8.10.1]
\label{II.8.10.1-fr}
Soient $\pi:V\to X$ un morphisme, $j:X\to V$ une $X$-section de $V$ qui soit
une immersion fermée, $\mathcal{J}$ l'Idéal quasi-cohérent de
$\mathcal{O}_V$ définissant le sous-préschéma fermé de $V$ associé à $j$.
\begin{enumerate}
  \item[\rm{(i)}] Pour tout $n\geqslant 0$,
  $\pi_*(\mathcal{O}_V/\mathcal{J}^{n+1})$ et
  $\pi_*(\mathcal{J}/\mathcal{J}^{n+1})$ sont des
  $\mathcal{O}_X$-Modules quasi-cohérents, et on a
  $\pi_*(\mathcal{O}_V/\mathcal{J})=\mathcal{O}_X$,
  $\pi_*(\mathcal{J}/\mathcal{J}^2)=j^*(\mathcal{J})$.
  \item[\rm{(ii)}] Si $X=\{\xi\}=\operatorname{Spec}(k)$, où $k$ est un
  corps, $\varprojlim_n\pi_*(\mathcal{O}_V/\mathcal{J}^{n+1})$ est isomorphe
  au séparé complété de l'anneau local $\mathcal{O}_{j(\xi)}$ pour la
  topologie $\mathfrak{m}_{j(\xi)}$-préadique.
  \item[\rm{(iii)}] Supposons que $\mathcal{J}$ soit un
  $\mathcal{O}_V$-Module inversible (ce qui entraîne que
  \[
    \mathcal{L}=j^*(\mathcal{J})=
    \pi_*(\mathcal{J}/\mathcal{J}^2)
  \]
  est un $\mathcal{O}_X$-Module inversible) et qu'il existe un homomorphisme
  \[
    \varphi:\mathcal{L}\longrightarrow
    \varprojlim_n\pi_*(\mathcal{J}/\mathcal{J}^{n+1})
  \]
  tel que le composé
  \[
    \mathcal{L}\xrightarrow{\varphi}
    \varprojlim_n\pi_*(\mathcal{J}/\mathcal{J}^{n+1})
    \xrightarrow{\alpha}\pi_*(\mathcal{J}/\mathcal{J}^2)
  \]
  ($\alpha$ homomorphisme canonique) soit l'identité. Si on pose
  $\mathcal{S}=\bigoplus_{n\geqslant0}\mathcal{L}^{\otimes n}$, on déduit
  alors canoniquement de $\varphi$ un isomorphisme de
  $\mathcal{O}_X$-Algèbres du complété $\widehat{\mathcal{S}}$ de
  $\mathcal{S}$ relatif à sa filtration canonique (complété qui est isomorphe
  au produit $\prod_{n\geqslant0}\mathcal{L}^{\otimes n}$) sur
  $\varprojlim_n\pi_*(\mathcal{O}_V/\mathcal{J}^{n+1})$.
\end{enumerate}
\end{lemma}

\begin{proof}
Notons en premier lieu que le support du $\mathcal{O}_V$-Module
$\mathcal{O}_V/\mathcal{J}^{n+1}$ est $j(X)$ et celui de
$\mathcal{J}/\mathcal{J}^{n+1}$ est contenu dans $j(X)$. Dans le cas (ii),
$j(X)$ est un point fermé $j(\xi)$ de $V$,

\oldpage[II]{185}
et par définition $\pi_*(\mathcal{O}_V/\mathcal{J}^{n+1})$ est la fibre de
$\mathcal{O}_V/\mathcal{J}^{n+1}$ au point $j(\xi)$, c'est-à-dire, en posant
$C=\mathcal{O}_{j(\xi)}$ et en désignant par $\mathfrak{m}$ l'idéal maximal
de $C$, le $C$-module $C/\mathfrak{m}^{n+1}$~; l'assertion (ii) est alors
évidente.

Pour démontrer (i), notons que la question est locale sur $X$~; on peut donc
se restreindre au cas où $X$ est affine. Soit $U$ un ouvert affine dans $V$~;
$j(X)\cap U$ est un ouvert affine dans $j(X)$, donc
$U_0=\pi(j(X)\cap U)$, qui lui est isomorphe, un ouvert affine dans $X$~;
pour tout ouvert affine $W_0\subset U_0$ dans $X$,
$W=\pi^{-1}(W_0)\cap U$ est un ouvert affine dans $V$, puisque $X$ est un
schéma (\hyperref[I.5.5.10-fr]{I, 5.5.10})~; en particulier
$U'=U\cap\pi^{-1}(U_0)$ est un ouvert affine dans $V$ et on a évidemment
$\pi(U')=U_0$ et $j(U_0)=j(X)\cap U$. Cela étant, par définition
\[
  \Gamma(W_0,\pi_*(\mathcal{O}_V/\mathcal{J}^{n+1}))
  =\Gamma(\pi^{-1}(W_0),\mathcal{O}_V/\mathcal{J}^{n+1})~;
\]
mais comme tout point de $\pi^{-1}(W_0)$ n'appartenant pas à $j(W_0)$ a un
voisinage ouvert dans $\pi^{-1}(W_0)$ ne rencontrant pas $j(X)$ et dans
lequel $\mathcal{O}_V/\mathcal{J}^{n+1}$ est donc nul, il est clair que les
sections de $\mathcal{O}_V/\mathcal{J}^{n+1}$ au-dessus de $\pi^{-1}(W_0)$ et
au-dessus de $W$ se correspondent biunivoquement. Autrement dit, si $\pi'$ est
la restriction de $\pi$ à $U'$, les $(\mathcal{O}_X|U_0)$-Modules
\[
  \pi_*(\mathcal{O}_V/\mathcal{J}^{n+1})|U_0
  \quad\hbox{et}\quad
  \pi'_*\bigl((\mathcal{O}_V/\mathcal{J}^{n+1})|U'\bigr)
\]
sont identiques. Comme $U'$ et $U_0$ sont affines et que les $U_0$ recouvrent
$X$, on en conclut (\hyperref[I.1.6.3-fr]{I, 1.6.3}) que
$\pi_*(\mathcal{O}_V/\mathcal{J}^{n+1})$ est quasi-cohérent, et la
démonstration est identique pour
$\pi_*(\mathcal{J}/\mathcal{J}^{n+1})$.

Pour démontrer enfin (iii), notons que $\mathcal{S}$ n'est autre que
$\mathbf{S}_{\mathcal{O}_X}(\mathcal{L})$~; on déduit donc canoniquement de
$\varphi$ un homomorphisme de $\mathcal{O}_X$-Algèbres
\[
  \psi:\mathcal{S}\longrightarrow
  \varprojlim_n\pi_*(\mathcal{O}_V/\mathcal{J}^{n+1})
\]
(\hyperref[II.1.7.4-fr]{1.7.4})~; en outre, cet homomorphisme applique
$\mathcal{L}^{\otimes n}$ dans
$\varprojlim_m\pi_*(\mathcal{J}^n/\mathcal{J}^{m+1})$, donc est continu pour
les topologies considérées, et se prolonge bien par suite en un homomorphisme
\[
  \widehat{\psi}:\widehat{\mathcal{S}}\longrightarrow
  \varprojlim_n\pi_*(\mathcal{O}_V/\mathcal{J}^{n+1}).
\]
Pour voir qu'il s'agit d'un isomorphisme, on peut, comme dans la démonstration
de (i), se réduire au cas où $X=\operatorname{Spec}(A)$ et
$V=\operatorname{Spec}(B)$ sont affines, avec
$\mathcal{J}=\widetilde{\mathfrak{J}}$, où $\mathfrak{J}$ est un idéal de
$B$~; il correspond à $\pi$ une injection $A\to B$, identifiant $A$ à un
sous-anneau de $B$ supplémentaire de $\mathfrak{J}$, et $\mathcal{L}$ (resp.
$\pi_*(\mathcal{O}_V/\mathcal{J}^{n+1})$) est le $\mathcal{O}_X$-Module
quasi-cohérent associé au $A$-module $L=\mathfrak{J}/\mathfrak{J}^2$ (resp.
$B/\mathfrak{J}^{n+1}$). Comme $\mathcal{J}$ est un
$\mathcal{O}_V$-Module \emph{inversible}, on peut en outre supposer que
$\mathfrak{J}=Bt$, où $t$ est non diviseur de 0 dans $B$. De la relation
$B=A\oplus Bt$ on déduit alors, pour tout $n>0$,
\[
  B=A\oplus At\oplus At^2\oplus\cdots\oplus At^n\oplus Bt^{n+1}
\]
donc on a un $A$-isomorphisme canonique de l'anneau de séries formelles
$A[[T]]$ sur $C=\varprojlim B/\mathfrak{J}^{n+1}$ faisant correspondre $t$ à
$T$. D'autre part, on a $L=A\bar t$, où $\bar t$ est la classe de $t$
mod. $Bt^2$, et l'homomorphisme $\varphi$ applique par hypothèse $\bar t$ sur
un élément $t'\in C$ congru à $t$ mod. $Ct^2$. On en déduit par récurrence
sur $n$ que
\[
  A\oplus At'\oplus\cdots\oplus At'{}^n\oplus Ct^{n+1}
  =A\oplus At\oplus\cdots\oplus At^n\oplus Ct^{n+1}
\]
ce qui prouve que l'homomorphisme $\widehat{\psi}$ correspond bien à un
isomorphisme de $\prod_{n\geqslant0}L^{\otimes n}$ sur $C$.
\end{proof}

\begin{lemma}[8.10.2]
\label{II.8.10.2-fr}
Sous les hypothèses de (\hyperref[II.8.10.1-fr]{8.10.1}), soit
$g:X'\to X$ un morphisme,

\oldpage[II]{186}
soit $V'=V\times_X X'$, et soient $\pi':V'\to X'$,
$g':V'\to V$ les projections canoniques, de sorte qu'on a le diagramme
commutatif
\[
  \xymatrix{
    V\ar[d]_{\pi} & V'\ar[l]^{g'}\ar[d]^{\pi'}\\
    X & X'\ar[l]_{g}
  }
\]
Alors $j'=j\times1_{X'}$ est une $X'$-section de $V'$ qui est une immersion
fermée, et $\mathcal{J}'=(g')^*(\mathcal{J})\mathcal{O}_{V'}$ est l'Idéal
quasi-cohérent de $\mathcal{O}_{V'}$ définissant le sous-préschéma fermé de
$V'$ associé à $j'$. En outre, on a
\[
  \pi'_*(\mathcal{O}_{V'}/\mathcal{J}'^{n+1})
  =g^*(\pi_*(\mathcal{O}_V/\mathcal{J}^{n+1})).
\]
Enfin, $\mathcal{J}'$ est un $\mathcal{O}_{V'}$-Module canoniquement isomorphe
à $(g')^*(\mathcal{J})$ et est en particulier inversible si $\mathcal{J}$ est
un $\mathcal{O}_V$-Module inversible.
\end{lemma}

\begin{proof}
Le fait que $j'$ est une immersion fermée découle de
(\hyperref[I.4.3.1-fr]{I, 4.3.1}), et c'est une $X'$-section de $V'$ par
fonctorialité de l'extension du préschéma de base. En outre, si $Z$ (resp.
$Z'$) est le sous-préschéma fermé de $V$ (resp. $V'$) associé à $j$ (resp.
$j'$), on a $Z'=g'{}^{-1}(Z)$
(\hyperref[I.4.3.1-fr]{I, 4.3.1}), et la seconde assertion résulte alors de
(\hyperref[I.4.4.5-fr]{I, 4.4.5}). Pour démontrer les autres assertions, on
voit comme dans (\hyperref[II.8.10.1-fr]{8.10.1}) qu'on peut se ramener au cas
où $X$, $V$ et $X'$ (donc aussi $V'$) sont affines~; gardons les notations de
la démonstration de (\hyperref[II.8.10.1-fr]{8.10.1}) et soit
$X'=\operatorname{Spec}(A')$. On a alors $V'=\operatorname{Spec}(B')$ où
$B'=B\otimes_A A'$, et $\mathcal{J}'=\widetilde{\mathfrak{J}'}$, avec
$\mathfrak{J}'=\operatorname{Im}(\mathfrak{J}\otimes_A A')$. On a donc
\[
  B'/\mathfrak{J}'^{n+1}=(B/\mathfrak{J}^{n+1})\otimes_A A'~;
\]
en outre, comme $\mathfrak{J}$ est facteur direct (en tant que $A$-module) de
$B$, $\mathfrak{J}\otimes_A A'$ est facteur direct (en tant que $A'$-module)
de $B'$, donc s'identifie canoniquement à $\mathfrak{J}'$.
\end{proof}

\begin{corollary}[8.10.3]
\label{II.8.10.3-fr}
Supposons vérifiées les hypothèses de (\hyperref[II.8.10.1-fr]{8.10.1}) et
supposons en outre que $\pi$ soit de type fini et que $\mathcal{J}$ soit un
$\mathcal{O}_V$-Module inversible. Alors, pour tout $x\in X$, l'anneau local
au point $j(x)$ de la fibre $\pi^{-1}(x)$ est un anneau régulier (donc
intègre) de dimension 1, dont le complété est isomorphe à l'anneau de séries
formelles $k(x)[[T]]$ ($T$ indéterminée)~; en outre il n'existe qu'une seule
composante irréductible de $\pi^{-1}(x)$ contenant $j(x)$.
\end{corollary}

\begin{proof}
Comme $\pi^{-1}(x)=V\times_X\operatorname{Spec}(k(x))$, on est ramené par
(\hyperref[II.8.10.2-fr]{8.10.2}) au cas où $X$ est le spectre d'un corps
$K$. Comme $\pi$ est de type fini
(\hyperref[I.6.3.4-fr]{I, 6.3.4, (iv)}), $\mathcal{O}_{j(x)}$ est alors un
anneau local noethérien, donc séparé pour la topologie
$\mathfrak{m}_{j(x)}$-préadique
(\hyperref[0.7.3.5-fr]{0, 7.3.5})~; il résulte de
(\hyperref[II.8.10.1-fr]{8.10.1, (ii) et (iii)}) que le complété de cet anneau
est isomorphe à $K[[T]]$, et par suite $\mathcal{O}_{j(x)}$ est régulier et de
dimension 1 ([1], p.~17-01, th.~1)~; enfin, puisque $\mathcal{O}_{j(x)}$ est
intègre, $j(x)$ n'appartient qu'à une seule des composantes irréductibles (en
nombre fini) de $V$ (\hyperref[I.5.1.4-fr]{I, 5.1.4}).
\end{proof}

\begin{corollary}[8.10.4]
\label{II.8.10.4-fr}
Supposons vérifiées les hypothèses de (\hyperref[II.8.10.1-fr]{8.10.1}) et en
outre supposons que $\mathcal{J}$ soit un $\mathcal{O}_V$-Module inversible.
Soit $W=V-j(X)$~; pour tout Idéal quasi-cohérent $\mathcal{K}$ de
$\mathcal{O}_X$, posons
$\mathcal{K}_V=\pi^*(\mathcal{K})\mathcal{O}_V$ et
$\mathcal{K}_W=\mathcal{K}_V|W$. Alors $\mathcal{K}_V$ est le plus grand
Idéal quasi-cohérent de $\mathcal{O}_V$ dont la restriction à $W$ soit
$\mathcal{K}_W$.
\end{corollary}

\begin{proof}
En effet, on voit comme dans (\hyperref[II.8.10.1-fr]{8.10.1}) que la question
est locale sur $X$ et sur $V$~; on peut donc garder les notations de la
démonstration de (\hyperref[II.8.10.1-fr]{8.10.1}), avec
$\mathfrak{J}=Bt$, où $t$ n'est pas diviseur de 0 dans $B$. En outre, on a
$W=\operatorname{Spec}(B_t)$ et $\mathcal{K}=\widetilde{\mathfrak{K}}$, où
$\mathfrak{K}$ est un idéal de $A$~; d'où
$\pi^*(\mathcal{K})\mathcal{O}_V=(\mathfrak{K}.B)^{\sim}$
(\hyperref[I.1.6.9-fr]{I, 1.6.9}),
$\mathcal{K}_W=(\mathfrak{K}.B_t)^{\sim}$, et le plus grand idéal

\oldpage[II]{187}
de $B$ dont l'image canonique dans $B_t$ soit $\mathfrak{K}.B_t$ est l'image
réciproque de $\mathfrak{K}.B_t$, c'est-à-dire l'ensemble des $s\in B$ tels
que pour un entier $n>0$, on ait $t^ns\in\mathfrak{K}.B$. Il faut montrer que
cette dernière relation entraîne $s\in\mathfrak{K}.B$, ou encore que l'image
canonique de $t$ n'est pas diviseur de 0 dans
$B/\mathfrak{K}B=(A/\mathfrak{K})\otimes_A B$, ce qui résulte de
(\hyperref[II.8.10.2-fr]{8.10.2}) appliqué à
$X'=\operatorname{Spec}(A/\mathfrak{K})$.
\end{proof}

\begin{corollary}[8.10.5]
\label{II.8.10.5-fr}
Supposons vérifiées les hypothèses de (\hyperref[II.8.10.3-fr]{8.10.3})~;
soient $W=V-j(X)$, $x$ un point de $X$, $\mathcal{K}$ un Idéal
quasi-cohérent de $\mathcal{O}_X$, $z$ le point générique de la composante
irréductible de $\pi^{-1}(x)$ contenant $j(x)$
(\hyperref[II.8.10.3-fr]{8.10.3}).
\begin{enumerate}
  \item[\rm{(i)}] Soit $g$ une section de $\mathcal{O}_V$ au-dessus de $V$
  telle que $g|W$ soit une section de $\mathcal{K}_W$ au-dessus de $W$
  (notations de (\hyperref[II.8.10.4-fr]{8.10.4})). Alors $g$ est une section
  de $\mathcal{K}_V$~; si de plus $g(z)\neq0$, et si, pour tout entier $m>0$,
  on désigne par $g_m^x$ le germe au point $x$ de l'image canonique $g_m$ de
  $g$ dans
  $\Gamma(X,\pi_*(\mathcal{O}_V/\mathcal{J}^{m+1}))$, alors il existe un
  entier $m>0$ tel que l'image de $g_m^x$ dans
  \[
    (\pi_*(\mathcal{O}_V/\mathcal{J}^{m+1}))_x
    \otimes_{\mathcal{O}_x}k(x)
  \]
  soit $\neq0$.
  \item[\rm{(ii)}] Supposons en outre que les conditions de
  (\hyperref[II.8.10.1-fr]{8.10.1, (iii)}) soient remplies. Alors, s'il existe
  une section $g$ de $\mathcal{K}_V$ au-dessus de $V$ telle que $g(z)\neq0$,
  il existe un entier $n\geqslant0$ et une section $f$ de
  $\mathcal{K}.\mathcal{L}^{\otimes n}=
  \mathcal{K}\otimes\mathcal{L}^{\otimes n}\subset\mathcal{L}^{\otimes n}$
  telle que $f(x)\neq0$. Si $g$ est une section de $\mathcal{J}$, on peut
  prendre $n>0$.
\end{enumerate}
\end{corollary}

\begin{proof}
\begin{enumerate}
  \item[\rm{(i)}] Comme l'Idéal de $\mathcal{O}_W$ engendré par $g|W$ est
  contenu dans $\mathcal{K}_W$ par hypothèse, l'Idéal de $\mathcal{O}_V$
  engendré par $g$ est contenu dans $\mathcal{K}_V$ par
  (\hyperref[II.8.10.4-fr]{8.10.4}), autrement dit $g$ est une section de
  $\mathcal{K}_V$. Pour démontrer la seconde assertion de (i), on peut encore
  supposer $X$ et $V$ affines et conserver les notations de
  (\hyperref[II.8.10.1-fr]{8.10.1})~; la fibre $\pi^{-1}(x)$ est alors affine
  d'anneau $B'=B\otimes_A k(x)$, et il existe dans $B'$ un élément $t'$ non
  diviseur de 0 tel que $B'=k(x)\oplus B't'$. Comme $j(x)$ est une
  spécialisation de $z$ et que $g(z)\neq0$, on a nécessairement
  $g_{j(x)}\neq0$. Mais $\mathcal{O}_{j(x)}$ est un anneau local séparé
  (\hyperref[II.8.10.3-fr]{8.10.3}), qui se plonge donc dans son complété, et
  l'image de $g$ dans ce complété n'est par suite pas nulle. Or, ce complété
  est isomorphe à $\varprojlim_n(B'/B't'^{n+1})$
  (\hyperref[II.8.10.3-fr]{8.10.3})~; si $g'=g\otimes1\in B'$, il existe donc
  un entier $m$ tel que $g'\notin B't'^{m+1}$, ou encore l'image $g'_m$ de
  $g'$ dans $B'/B't'^{m+1}$ n'est pas nulle. Mais comme $g'_m$ n'est autre que
  l'image de $g_m^x$, notre assertion est démontrée.
  \item[\rm{(ii)}] En vertu de
  (\hyperref[II.8.10.1-fr]{8.10.1, (iii)})
  $\pi_*(\mathcal{O}_V/\mathcal{J}^{m+1})$ est isomorphe à la somme directe
  des $\mathcal{L}^{\otimes k}$ pour $0\leqslant k\leqslant m$~; désignons
  par $f_k$ la section de $\mathcal{L}^{\otimes k}$ au-dessus de $X$,
  composante de l'élément de
  $\bigoplus_{k=0}^m\Gamma(X,\mathcal{L}^{\otimes k})$ qui correspond à $g_m$
  par cet isomorphisme. Choisissant $m$ comme dans (i), il y a donc un indice
  $k$ tel que $f_k(x)\neq0$ d'après (i). Pour voir que $f_k$ est une section
  de $\mathcal{K}\mathcal{L}^{\otimes k}$, il suffit de considérer comme
  ci-dessus le cas où $X$ et $V$ sont affines, et cela résulte aussitôt alors
  de ce que $g\in\mathfrak{K}.B$ (notations de
  (\hyperref[II.8.10.4-fr]{8.10.4})). La dernière assertion résulte de ce que
  l'hypothèse $g\in\Gamma(V,\mathcal{J})$ entraîne $f_0=0$.
\end{enumerate}
\end{proof}

\begin{env}[8.10.6]
\label{II.8.10.6-fr}
\emph{Démonstration de (\hyperref[II.8.9.4-fr]{8.9.4}).} La question est
locale sur $Y$ (\hyperref[II.4.6.4-fr]{4.6.4})~; comme $\varepsilon$ est une
$Y$-section, on peut donc remplacer $C$ par un voisinage ouvert affine $U$
d'un point de $\varepsilon(Y)$ tel que $\varepsilon(Y)\cap U$ soit fermé dans
$U$. Autrement dit, on peut supposer $C$ affine, et $Y$ sous-préschéma fermé
de $C$ (donc affine) défini par un faisceau quasi-

\oldpage[II]{188}
cohérent $\mathcal{I}$ d'idéaux de $\mathcal{O}_C$. Comme $p$ est séparé et
quasi-compact, $X$ est alors un schéma quasi-compact, et on est ramené à
prouver que $\mathcal{L}$ est ample
(\hyperref[II.4.6.4-fr]{4.6.4}). En vertu du critère
(\hyperref[II.4.5.2-fr]{4.5.2, a)}), on doit donc prouver qui ce suit~: pour
tout Idéal quasi-cohérent $\mathcal{K}$ de $\mathcal{O}_X$ et tout point
$x\in X$ n'appartenant pas au support de $\mathcal{O}_X/\mathcal{K}$, il
existe un entier $n>0$ et une section $f$ de
$\mathcal{K}\otimes\mathcal{L}^{\otimes n}$ au-dessus de $X$ telle que
$f(x)\neq0$.

Pour cela, posons
\[
  \mathcal{K}_V=\pi^*(\mathcal{K})\mathcal{O}_V
\]
et
\[
  \mathcal{K}_W=\mathcal{K}_V|W,\qquad\hbox{où}\qquad W=V-j(X)~;
\]
comme la restriction de $q$ à $W$ est une immersion quasi-compacte dans $C$,
il résulte de (\hyperref[I.9.4.2-fr]{I, 9.4.2}) que $\mathcal{K}_W$ est la
restriction à $W$ d'un Idéal quasi-cohérent $\mathcal{K}'_V$ de
$\mathcal{O}_V$, de la forme
\[
  \mathcal{K}'_V=q^*(\mathcal{K}_C)\mathcal{O}_V
\]
où $\mathcal{K}_C$ est un Idéal quasi-cohérent de $\mathcal{O}_C$. En outre,
comme par hypothèse $q^{-1}(Y)\subset j(X)$ et que $Y$ est défini par l'Idéal
$\mathcal{I}$, la restriction à $W$ de $q^*(\mathcal{I})\mathcal{O}_V$ est
identique à celle de $\mathcal{O}_V$, donc $\mathcal{K}_W$ est aussi la
restriction à $W$ de $q^*(\mathcal{I}\mathcal{K}_C)\mathcal{O}_V$, et l'on
peut donc supposer que l'on a $\mathcal{K}_C\subset\mathcal{I}$, d'où
\begin{equation}
\label{II.8.10.6.1-fr}
  \mathcal{K}'_V\subset q^*(\mathcal{I})\mathcal{O}_V\subset\mathcal{J}
\tag{8.10.6.1}
\end{equation}
compte tenu de (\hyperref[I.4.4.6-fr]{I, 4.4.6}) et de la commutativité de
(\hyperref[II.8.9.4.1-fr]{8.9.4.1}). En outre, on déduit de
(\hyperref[II.8.10.4-fr]{8.10.4}) que l'on a
\begin{equation}
\label{II.8.10.6.2-fr}
  \mathcal{K}'_V\subset\mathcal{K}_V.
\tag{8.10.6.2}
\end{equation}

Cela étant, il résulte de (\hyperref[II.8.10.3-fr]{8.10.3}) que $j(x)$
appartient à une seule composante irréductible de $\pi^{-1}(x)$~; soit $z$ le
point générique de cette composante, et posons $z'=q(z)$. En vertu de
(\hyperref[II.8.10.5-fr]{8.10.5}), la démonstration sera achevée (compte tenu
de (\hyperref[II.8.10.6.1-fr]{8.10.6.1}) et
(\hyperref[II.8.10.6.2-fr]{8.10.6.2})) si nous prouvons l'existence d'une
section $g$ de $\mathcal{K}'_V$ au-dessus de $V$ telle que $g(z)\neq0$. Or,
par hypothèse, $\mathcal{K}$ a une restriction égale à celle de
$\mathcal{O}_X$ dans un voisinage ouvert de $x$~; d'autre part, il résulte de
(\hyperref[II.8.10.3-fr]{8.10.3}) que l'on a $z\neq j(x)$, donc $z\in W$, et
par suite $(\mathcal{K}_W)_z=\mathcal{O}_{V,z}$, d'où par définition
$(\mathcal{K}_C)_{z'}=\mathcal{O}_{C,z'}$. Puisque $C$ est affine, il y a donc
une section $g'$ de $\mathcal{K}_C$ au-dessus de $C$ telle que
$g'(z')\neq0$, et en prenant pour $g$ la section de $\mathcal{K}'_V$
correspondant canoniquement à $g'$, on a bien $g(z)\neq0$, ce qui achève la
démonstration.
\end{env}

\begin{remark}[8.10.7]
\label{II.8.10.7-fr}
Nous ignorons si, dans l'énoncé de (\hyperref[II.8.9.4-fr]{8.9.4}), la
condition (ii) est ou non superflue. En tout cas, la conclusion ne subsiste
plus si l'on ne suppose pas l'existence d'un $Y$-morphisme $\pi:V\to X$ tel
que $\pi\circ j=1_X$~; indiquons brièvement comment on peut en effet former un
contre-exemple, dont les détails ne pourront être développés que plus tard.
On prend $Y=\operatorname{Spec}(k)$ où $k$ est un corps,
$C=\operatorname{Spec}(A)$, où $A=k[T_1,T_2]$, la $Y$-section $\varepsilon$
correspondant à l'homomorphisme d'augmentation $A\to k$. On désigne par $C'$
le schéma déduit de $C$ par éclatement du point fermé $a=\varepsilon(Y)$ de
$C$~; si $D$ est l'image réciproque de $a$ dans $C'$, on considère dans $D$
un point fermé $b$,

\oldpage[II]{189}
et on désigne par $V$ le schéma déduit de $C'$ par éclatement de $b$~; $X$ est
le sous-préschéma fermé de $V$, image réciproque de $a$ par le morphisme
structural $q:V\to C$. On montre que $X$ est réunion de deux composantes
irréductibles $X_1$, $X_2$, où $X_1$ est l'image réciproque de $b$ dans $V$.
Il est immédiat que l'Idéal $\mathcal{J}$ de $\mathcal{O}_V$ qui définit $X$
est encore inversible, mais on prouve que
$j^*(\mathcal{J})=\mathcal{L}$ ($j$ injection canonique $X\to V$) n'est pas
ample, par la considération du «~degré~» de l'image réciproque de
$\mathcal{L}$ sur $X_1$, qui devrait être $>0$ si $\mathcal{L}$ était ample,
et qu'on montre (par un calcul élémentaire d'intersections) être égal à 0.
\end{remark}

\subsection{Unicité des contractions}
\label{subsection:II.8.11-fr}

\begin{lemma}[8.11.1]
\label{II.8.11.1-fr}
Soient $U$, $V$ deux préschémas,
$h=(h_0,\lambda):U\to V$ un morphisme surjectif. On suppose que~:
\begin{enumerate}
  \item[$1^\circ$]
  $\lambda:\mathcal{O}_V\to h_*(\mathcal{O}_U)
  =(h_0)_*(\mathcal{O}_U)$ est un isomorphisme.
  \item[$2^\circ$] L'espace sous-jacent à $V$ s'identifie à l'espace quotient
  de l'espace sous-jacent à $U$ par la relation
  $R:h_0(x)=h_0(y)$ (condition toujours vérifiée lorsque le morphisme $h$ est
  ouvert, ou fermé, ou a fortiori lorsque $h$ est propre).
\end{enumerate}
Alors, pour tout préschéma $W$, l'application
\begin{equation}
\label{II.8.11.1.1-fr}
  \operatorname{Hom}(V,W)\longrightarrow\operatorname{Hom}(U,W)
\tag{8.11.1.1}
\end{equation}
qui, à tout morphisme $v=(v_0,\nu)$ de $V$ dans $W$, fait correspondre le
morphisme $u=v\circ h=(u_0,\mu)$, est une bijection de
$\operatorname{Hom}(V,W)$ sur l'ensemble des $u$ tels que $u_0$ soit
constante sur toute fibre $h_0^{-1}(x)$.
\end{lemma}

\begin{proof}
Il est clair que si $u=v\circ h$, donc $u_0=v_0\circ h_0$, $u_0$ est
constante sur tout ensemble $h_0^{-1}(x)$. Inversement, si $u$ a cette
propriété, montrons qu'il existe un $v\in\operatorname{Hom}(V,W)$ et un seul
tel que $u=v\circ h$. L'existence et l'unicité de l'application continue
$v_0:V\to W$ telle que $u_0=v_0\circ h_0$ résultent de l'hypothèse, puisque
$h_0$ s'identifie à l'application canonique de $U$ sur $U/R$. On peut d'autre
part, en remplaçant au besoin $V$ par un préschéma isomorphe, supposer que
$\lambda$ est l'identité~; par hypothèse, $\mu$ est alors un homomorphisme
\[
  \mu:\mathcal{O}_W\longrightarrow(u_0)_*(\mathcal{O}_U)
  =(v_0)_*((h_0)_*(\mathcal{O}_U))
\]
tel que l'homomorphisme correspondant
$\mu^\sharp:u_0^*(\mathcal{O}_W)\to\mathcal{O}_U$ soit local sur chaque
fibre. Comme
$(v_0)_*((h_0)_*(\mathcal{O}_U))=(v_0)_*(\mathcal{O}_V)$, on doit
nécessairement avoir $\nu=\mu$, et tout revient à voir que l'homomorphisme
correspondant $\nu^\sharp:v_0^*(\mathcal{O}_W)\to\mathcal{O}_V$ est local sur
chaque fibre. Or, tout $y\in V$ est de la forme $h_0(x)$ pour un $x\in U$~;
posons $z=v_0(y)=u_0(x)$. Alors
(\hyperref[0.3.5.5-fr]{0, 3.5.5}) l'homomorphisme $\mu_x^\sharp$ se factorise
en
\[
  \mu_x^\sharp:\mathcal{O}_z
  \xrightarrow{\nu_y^\sharp}\mathcal{O}_y
  \xrightarrow{\lambda_x^\sharp}\mathcal{O}_x.
\]
Par hypothèse $\lambda_x^\sharp$ et $\mu_x^\sharp$ sont des homomorphismes
locaux~; donc $\lambda_x^\sharp$ transforme tout élément inversible de
$\mathcal{O}_y$ en un élément inversible de $\mathcal{O}_x$~; si
$\nu_y^\sharp$ transformait un élément non inversible de $\mathcal{O}_z$ en
un élément inversible de $\mathcal{O}_y$, $\mu_x^\sharp$ transformerait cet
élément de $\mathcal{O}_z$ en un élément inversible de $\mathcal{O}_x$,
contrairement à l'hypothèse, d'où le lemme.
\end{proof}

\begin{corollary}[8.11.2]
\label{II.8.11.2-fr}
Soient $U$ un préschéma intègre, $V$ un préschéma normal~; alors tout morphisme
$h:U\to V$ qui est universellement fermé, birationnel et radiciel est un
isomorphisme.
\end{corollary}

\begin{proof}
Si $h=(h_0,\lambda)$, il résulte des hypothèses que $h_0$ est injectif et
fermé, et que $h_0(U)$

\oldpage[II]{190}
est dense dans $V$, donc $h_0$ est un \emph{homéomorphisme} de $U$ sur $V$.
Pour démontrer le corollaire, il suffira de voir que
$\lambda:\mathcal{O}_V\to(h_0)_*(\mathcal{O}_U)$ est un isomorphisme~: on
pourra appliquer alors (\hyperref[II.8.11.1-fr]{8.11.1}) qui prouvera que
l'application (\hyperref[II.8.11.1.1-fr]{8.11.1.1}) est bijective (les fibres
$h_0^{-1}(x)$ étant chacune réduite à un seul point)~; donc $h$ sera un
isomorphisme. La question étant évidemment locale sur $V$, on peut supposer
que $V=\operatorname{Spec}(A)$ est affine d'anneau intègre et intégralement
clos (\hyperref[II.8.8.6.1-fr]{8.8.6.1})~; $h$ correspond alors
(\hyperref[I.2.2.4-fr]{I, 2.2.4}) à un homomorphisme
$\varphi:A\to\Gamma(U,\mathcal{O}_U)$ et tout revient à voir que $\varphi$
est un isomorphisme. Or, si $K$ est le corps des fractions de $A$,
$\Gamma(U,\mathcal{O}_U)$ a par hypothèse $K$ pour corps des fractions et $A$
est un sous-anneau de $\Gamma(U,\mathcal{O}_U)$, $\varphi$ étant l'injection
canonique (\hyperref[I.8.2.7-fr]{I, 8.2.7}). Comme le morphisme $h$ vérifie
les hypothèses de (\hyperref[II.7.3.11-fr]{7.3.11}),
$\Gamma(U,\mathcal{O}_U)$ est un sous-anneau de la fermeture intégrale de $A$
dans $K$, donc est identique à $A$ par hypothèse.
\end{proof}

\begin{remark}[8.11.3]
\label{II.8.11.3-fr}
On verra au chapitre III (\hyperref[III.4.4.11-fr]{III, 4.4.11}) que lorsque
$V$ est un préschéma \emph{localement noethérien}, tout morphisme $h:U\to V$
qui est propre et quasi-fini (en particulier tout morphisme vérifiant les
hypothèses de (\hyperref[II.8.11.2-fr]{8.11.2})) est nécessairement
\emph{fini}. La conclusion de (\hyperref[II.8.11.2-fr]{8.11.2}) résulte donc
dans ce cas de (\hyperref[II.6.1.15-fr]{6.1.15}).
\end{remark}

\begin{env}[8.11.4]
\label{II.8.11.4-fr}
Nous allons maintenant voir que, dans le critère de Grauert
(\hyperref[II.8.9.1-fr]{8.9.1}), on peut souvent affirmer que le préschéma
$C$ et la «~contraction~» $q$ sont déterminés de façon essentiellement
unique.
\end{env}

\begin{lemma}[8.11.5]
\label{II.8.11.5-fr}
Soient $Y$ un préschéma, $p:X\to Y$ un morphisme propre, $\mathcal{L}$ un
$\mathcal{O}_X$-Module inversible $p$-ample, $C$ un $Y$-préschéma,
$\varepsilon:Y\to C$ une $Y$-section,
$q:V=\mathbf{V}(\mathcal{L})\to C$ un $Y$-morphisme, tels que le diagramme
(\hyperref[II.8.9.1.1-fr]{8.9.1.1}) soit commutatif. On suppose en outre que,
si $p=(p_0,\theta)$,
$\theta:\mathcal{O}_Y\to p_*(\mathcal{O}_X)$ est un isomorphisme. Soient alors
$\mathcal{S}'=\bigoplus_{n\geqslant0}p_*(\mathcal{L}^{\otimes n})$,
$C'=\operatorname{Spec}(\mathcal{S}')$, et
$q':\mathbf{V}(\mathcal{L})\to C'$ le $Y$-morphisme canonique
(\hyperref[II.8.8.5-fr]{8.8.5}). Il existe un $Y$-morphisme $u:C'\to C$ et
un seul tel que $q=u\circ q'$.
\end{lemma}

\begin{proof}
L'hypothèse sur $\theta$ entraîne en particulier que $p$ est surjectif~;
comme en vertu de (\hyperref[II.8.8.4-fr]{8.8.4}) la restriction de $q'$ à
$\mathbf{V}(\mathcal{L})-j(X)$ est un \emph{isomorphisme} sur
$C'-\varepsilon'(Y)$ ($\varepsilon'$ étant la section sommet de $C'$), il
résulte de (\hyperref[II.8.8.4-fr]{8.8.4}) que $q'$ est \emph{propre et
surjectif}~; en outre, en vertu de (\hyperref[II.8.8.6-fr]{8.8.6}), si on pose
$q'=(q'_0,\tau)$,
$\tau:\mathcal{O}_{C'}\to q'_*(\mathcal{O}_V)$ est un isomorphisme. On est
donc dans les conditions d'application du lemme
(\hyperref[II.8.11.1-fr]{8.11.1}), et le lemme sera démontré si l'on prouve
que $q$ est constant sur toute fibre $q'{}^{-1}(z')$, où $z'\in C'$. Or, la
condition est trivialement vérifiée pour $z'\notin\varepsilon'(Y)$. D'autre
part, si $z'\in\varepsilon'(Y)$, il existe un $y\in Y$ et un seul tel que
$z'=\varepsilon'(y)$~; et en vertu de la commutativité de
(\hyperref[II.8.8.5.2-fr]{8.8.5.2}) et le fait que $q'$ applique
$\mathbf{V}(\mathcal{L})-j(X)$ dans $C'-\varepsilon'(Y)$,
$q'{}^{-1}(z')=j(p^{-1}(y))$~; la commutativité du diagramme
(\hyperref[II.8.9.1.1-fr]{8.9.1.1}) établit donc notre assertion.
\end{proof}

\begin{corollary}[8.11.6]
\label{II.8.11.6-fr}
Sous les hypothèses de (\hyperref[II.8.11.5-fr]{8.11.5}), supposons outre que
$q$ soit propre et que la restriction de $q$ à
$\mathbf{V}(\mathcal{L})-j(X)$ soit un isomorphisme sur
$C-\varepsilon(Y)$. Alors le morphisme $u$ est universellement fermé,
surjectif et radiciel, et sa restriction à $C'-\varepsilon'(Y)$ est un
isomorphisme sur $C-\varepsilon(Y)$.
\end{corollary}

\begin{proof}
Comme $q'$ est un isomorphisme de $\mathbf{V}(\mathcal{L})-j(X)$ sur
$C'-\varepsilon'(Y)$ (\hyperref[II.8.8.4-fr]{8.8.4}), la dernière assertion
résulte aussitôt de $q=u\circ q'$. En outre, la commutativité des diagrammes

\oldpage[II]{191}
(\hyperref[II.8.8.5.2-fr]{8.8.5.2}) et
(\hyperref[II.8.9.1.1-fr]{8.9.1.1}) montre que la restriction de $u$ au
sous-préschéma fermé $\varepsilon'(Y)$ de $C'$ est un isomorphisme sur le
sous-préschéma fermé $\varepsilon(Y)$ de $C$, d'où résulte aussitôt que pour
tout $z'\in\varepsilon'(Y)$, si $z=u(z')$, $u$ définit un isomorphisme de
$k(z)$ sur $k(z')$. Ces remarques prouvent que $u$ est bijectif et radiciel~;
en outre, si $\psi$ et $\psi'$ sont les morphismes structuraux $C\to Y$,
$C'\to Y$, on a $\psi'=\psi\circ u$ et comme $\psi'$ est séparé
(\hyperref[II.1.2.4-fr]{1.2.4}), il en est de même de $u$
(\hyperref[I.5.5.1-fr]{I, 5.5.1, (v)}). On a vu d'autre part dans la
démonstration du lemme (\hyperref[II.8.11.5-fr]{8.11.5}) que $q'$ est
surjectif~; comme $q=u\circ q'$ est propre, on conclut finalement de
(\hyperref[II.5.4.3-fr]{5.4.3}) et (\hyperref[II.5.4.9-fr]{5.4.9}) que $u$
est universellement fermé.
\end{proof}

\begin{proposition}[8.11.7]
\label{II.8.11.7-fr}
Soient $Y$ un préschéma, $X$ un préschéma intègre, $p:X\to Y$ un morphisme
propre, $\mathcal{L}$ un $\mathcal{O}_X$-Module inversible $p$-ample, $C$ un
$Y$-préschéma normal, $\varepsilon:Y\to C$ une $Y$-section,
$q:V=\mathbf{V}(\mathcal{L})\to C$ un $Y$-morphisme, tels que le diagramme
(\hyperref[II.8.9.1.1-fr]{8.9.1.1}) soit commutatif. On suppose en outre que,
si $p=(p_0,\theta)$, $\theta:\mathcal{O}_Y\to p_*(\mathcal{O}_X)$ soit un
isomorphisme. Soient alors
$\mathcal{S}'=\bigoplus_{n\geqslant0}p_*(\mathcal{L}^{\otimes n})$,
$C'=\operatorname{Spec}(\mathcal{S}')$, et
$q':\mathbf{V}(\mathcal{L})\to C'$ le $Y$-morphisme canonique
(\hyperref[II.8.8.5-fr]{8.8.5}). Alors l'unique $Y$-morphisme $u:C'\to C$ tel
que $q=u\circ q'$ est un isomorphisme.
\end{proposition}

\begin{proof}
Il résulte de (\hyperref[II.8.8.6-fr]{8.8.6}) que $C'$ est intègre~; comme
$u$ est un homéomorphisme des espaces sous-jacents $C'\to C$ ($u$ étant
bijectif et fermé d'après (\hyperref[II.8.11.6-fr]{8.11.6})), $C$ est
irréductible, donc intègre, et puisque la restriction de $u$ à un ouvert non
vide de $C'$ est un isomorphisme sur un ouvert de $C$, $u$ est birationnel.
Puisque $C$ est supposé normal, il suffit d'appliquer
(\hyperref[II.8.11.2-fr]{8.11.2}) pour obtenir la conclusion.
\end{proof}

\begin{remark}[8.11.8]
\label{II.8.11.8-fr}
\begin{enumerate}
  \item[\rm{(i)}] Notons que l'hypothèse que $C$ est normal implique qu'il en
  est de même de $X$. En effet, $C'=\operatorname{Spec}(\mathcal{S}')$ est
  alors normal, étant isomorphe à $C$, et intègre en vertu de
  (\hyperref[II.8.8.6-fr]{8.8.6})~; on en conclut que
  $\operatorname{Proj}(\mathcal{S}')$ est \emph{normal}. En effet, la question
  est locale sur $Y$~; si $Y$ est affine, $\mathcal{S}'=\widetilde{S'}$,
  l'anneau $S'=\Gamma(C',\mathcal{O}_{C'})$ est intègre et intégralement clos
  (\hyperref[II.8.8.6.1-fr]{8.8.6.1}), donc, pour tout élément homogène
  $f\in S'_+$, l'anneau gradué $S'_f$ est intègre et intégralement clos
  ([13], t.~I, p.~257 et 261), et par suite aussi l'anneau $S'_{(f)}$ de ses
  termes de degré 0, car l'intersection de $S'_f$ et du corps des fractions de
  $S'_{(f)}$ est égale à $S'_{(f)}$~; ce qui prouve notre assertion
  (\hyperref[II.6.3.4-fr]{6.3.4}). Enfin, comme $X$ est isomorphe à un
  sous-préschéma ouvert de $\operatorname{Proj}(\mathcal{S}')$
  (\hyperref[II.8.8.1-fr]{8.8.1}), $X$ est bien normal. On peut donc exprimer
  la prop. (\hyperref[II.8.11.7-fr]{8.11.7}) sous la forme suivante~:
  \emph{Si $X$ est intègre et normal,
  $p=(p_0,\theta):X\to Y$ un morphisme propre tel que
  $\theta:\mathcal{O}_Y\to p_*(\mathcal{O}_X)$ soit un isomorphisme, alors,
  pour tout $\mathcal{O}_X$-Module $p$-ample $\mathcal{L}$, il existe une
  façon et une seule de contracter la section nulle de
  $V=\mathbf{V}(\mathcal{L})$ de façon à obtenir un $Y$-schéma normal $C$ et
  un $Y$-morphisme propre $q:V\to C$.}
  \item[\rm{(ii)}] Lorsque $p$ est propre, l'hypothèse
  $p_*(\mathcal{O}_X)=\mathcal{O}_Y$ peut être considérée comme une hypothèse
  auxiliaire, ne constituant pas une vraie restriction de la généralité du
  résultat. En effet, si elle n'est pas vérifiée, il suffit de remplacer $Y$
  par le $Y$-schéma
  $Y'=\operatorname{Spec}(p_*(\mathcal{O}_X))$, et de considérer $X$ comme un
  $Y'$-schéma. Nous reviendrons sur cette méthode générale au chapitre III,
  \S~4.
\end{enumerate}
\end{remark}
\subsection{Faisceaux quasi-cohérents sur les cônes projetants}
\label{subsection:II.8.12-fr}

\begin{env}[8.12.1]
\label{II.8.12.1-fr}
Reprenons les hypothèses et notations de
(\hyperref[II.8.3.1-fr]{8.3.1}). Soit $\mathcal{M}$ un
$\mathcal{S}$-Module gradué quasi-cohérent~; pour éviter toute confusion,
nous désignerons par $\widetilde{\mathcal{M}}$ le $\mathcal{O}_C$-Module

\oldpage[II]{192}
quasi-cohérent associé à $\mathcal{M}$
(\hyperref[II.1.4.3-fr]{1.4.3}) lorsque $\mathcal{M}$ est considéré comme
$\mathcal{S}$-Module \emph{non gradué}, et par
$\operatorname{Proj}_0(\mathcal{M})$ le $\mathcal{O}_X$-Module
quasi-cohérent associé à $\mathcal{M}$, $\mathcal{M}$ étant cette fois
considéré comme $\mathcal{S}$-Module gradué (autrement dit le
$\mathcal{O}_X$-Module noté $\widetilde{\mathcal{M}}$ dans
(\hyperref[II.3.2.2-fr]{3.2.2})). Nous poserons de plus
\begin{equation}
\label{II.8.12.1.1-fr}
  \mathcal{M}_X=\operatorname{Proj}(\mathcal{M})
  =\bigoplus_{n\in\mathbf{Z}}\operatorname{Proj}_0(\mathcal{M}(n))~;
\tag{8.12.1.1}
\end{equation}
la $\mathcal{O}_X$-Algèbre graduée quasi-cohérente $\mathcal{S}_X$ étant
définie par (\hyperref[II.8.6.1.1-fr]{8.6.1.1}),
$\operatorname{Proj}(\mathcal{M})$ est muni d'une structure de
$\mathcal{S}_X$-Module gradué (quasi-cohérent), au moyen des homomorphismes
canoniques (\hyperref[II.3.2.6.1-fr]{3.2.6.1})
\begin{equation}
\label{II.8.12.1.2-fr}
  \mathcal{O}_X(m)\otimes_{\mathcal{O}_X}
  \operatorname{Proj}_0(\mathcal{M}(n))
  \to \operatorname{Proj}_0(\mathcal{S}(m)\otimes_{\mathcal{S}}\mathcal{M}(n))
  \to \operatorname{Proj}_0(\mathcal{M}(m+n))
\tag{8.12.1.2}
\end{equation}
la vérification des axiomes des Modules se faisant à l'aide du diagramme
commutatif (\hyperref[II.2.5.11.4-fr]{2.5.11.4}).

Si $Y=\operatorname{Spec}(A)$ est affine, $\mathcal{S}=\widetilde{S}$ et
$\mathcal{M}=\widetilde{M}$, où $S$ est une $A$-algèbre graduée et $M$ un
$S$-module gradué, alors, pour tout élément homogène $f\in S_+$, on a
\begin{equation}
\label{II.8.12.1.3-fr}
  \Gamma(X_f,\operatorname{Proj}(\widetilde{M}))=M_f
\tag{8.12.1.3}
\end{equation}
en vertu des définitions et de (\hyperref[II.8.2.9.1-fr]{8.2.9.1}).

Considérons maintenant le $\widehat{\mathcal{S}}$-Module gradué
quasi-cohérent
\begin{equation}
\label{II.8.12.1.4-fr}
  \widehat{\mathcal{M}}
  =\mathcal{M}\otimes_{\mathcal{S}}\widehat{\mathcal{S}}
\tag{8.12.1.4}
\end{equation}
($\widehat{\mathcal{S}}$ étant défini par
(\hyperref[II.8.3.1.1-fr]{8.3.1.1}))~; on en déduit un
$\mathcal{O}_{\widehat{C}}$-Module gradué quasi-cohérent
$\operatorname{Proj}_0(\widehat{\mathcal{M}})$, que nous noterons aussi
\begin{equation}
\label{II.8.12.1.5-fr}
  \mathcal{M}^{\square}=\operatorname{Proj}_0(\widehat{\mathcal{M}}).
\tag{8.12.1.5}
\end{equation}
Il est clair (\hyperref[II.3.2.4-fr]{3.2.4}) que
$\mathcal{M}^{\square}$ est un foncteur additif exact en $\mathcal{M}$,
commutant aux sommes directes et aux limites inductives.
\end{env}

\begin{proposition}[8.12.2]
\label{II.8.12.2-fr}
Avec les notations de (\hyperref[II.8.3.2-fr]{8.3.2}), on a des
isomorphismes canoniques fonctoriels
\begin{equation}
\label{II.8.12.2.1-fr}
  i^*(\mathcal{M}^{\square})\xrightarrow{\sim}\widetilde{\mathcal{M}},
  \qquad
  j^*(\mathcal{M}^{\square})\xrightarrow{\sim}
  \operatorname{Proj}_0(\mathcal{M}).
\tag{8.12.2.1}
\end{equation}
\end{proposition}

\begin{proof}
En effet, $i^*(\mathcal{M}^{\square})$ s'identifie canoniquement à
$\left(\widehat{\mathcal{M}}/(z-1)\widehat{\mathcal{M}}\right)^{\sim}$
sur
$\operatorname{Spec}(\widehat{\mathcal{S}}/(z-1)\widehat{\mathcal{S}})$
en vertu de (\hyperref[II.3.2.3-fr]{3.2.3})~; le premier des isomorphismes
canoniques (\hyperref[II.8.12.2.1-fr]{8.12.2.1}) se déduit alors aussitôt
(\hyperref[II.1.4.1-fr]{1.4.1}) de l'isomorphisme canonique
$\widehat{\mathcal{M}}/(z-1)\widehat{\mathcal{M}}
\xrightarrow{\sim}\mathcal{M}$. D'autre part, l'immersion canonique
$j:X\to\widehat{C}$ correspond à l'homomorphisme canonique
$\widehat{\mathcal{S}}\to\mathcal{S}$ de noyau
$z\widehat{\mathcal{S}}$ (\hyperref[II.8.3.2-fr]{8.3.2})~; le second
homomorphisme (\hyperref[II.8.12.2.1-fr]{8.12.2.1}) est le cas particulier
de l'homomorphisme canonique (\hyperref[II.3.5.2-fr]{3.5.2, (ii)}), du fait
que l'on a ici
$\widehat{\mathcal{M}}\otimes_{\widehat{\mathcal{S}}}\mathcal{S}
=\mathcal{M}$~; pour vérifier que c'est un isomorphisme, on peut se borner au
cas où $Y=\operatorname{Spec}(A)$ est affine,
$\mathcal{S}=\widetilde{S}$, $\mathcal{M}=\widetilde{M}$~; en se reportant à
(\hyperref[II.2.8.8-fr]{2.8.8}), la vérification que pour tout $f$ homogène
dans $S_+$, l'homomorphisme précédent, restreint à $X_f$, se réduit à un
isomorphisme, est alors immédiate.
\end{proof}

\oldpage[II]{193}
Par abus de langage, on dira encore, en raison de l'existence du premier
isomorphisme (\hyperref[II.8.12.2.1-fr]{8.12.2.1}), que
$\mathcal{M}^{\square}$ est la \emph{fermeture projective} du
$\mathcal{O}_C$-Module $\widetilde{\mathcal{M}}$ (étant sous-entendu que dans
la donnée du $\mathcal{O}_C$-Module $\widetilde{\mathcal{M}}$ figure la
graduation du $\mathcal{S}$-Module $\mathcal{M}$).

\begin{env}[8.12.3]
\label{II.8.12.3-fr}
Avec les notations de (\hyperref[II.8.3.5-fr]{8.3.5}), on a un
homomorphisme canonique fonctoriel
\begin{equation}
\label{II.8.12.3.1-fr}
  p^*(\operatorname{Proj}_0(\mathcal{M}))
  \to\mathcal{M}^{\square}|\widehat{E}.
\tag{8.12.3.1}
\end{equation}
En effet, c'est un cas particulier de l'homomorphisme $u^{\sharp}$ défini de
façon générale dans (\hyperref[II.3.5.6-fr]{3.5.6}). Si
$Y=\operatorname{Spec}(A)$ est affine, $\mathcal{S}=\widetilde{S}$,
$\mathcal{M}=\widetilde{M}$, on voit, en se reportant à
(\hyperref[II.2.8.8-fr]{2.8.8}), que la restriction de
(\hyperref[II.8.12.3.1-fr]{8.12.3.1}) à
$p^{-1}(X_f)=\widehat{C}_f$ (pour un $f$ homogène dans $S_+$) correspond à
l'homomorphisme canonique
\begin{equation}
\label{II.8.12.3.2-fr}
  M_{(f)}\otimes_{S_{(f)}}S_f^{\leq}\to M_f^{\leq}
\tag{8.12.3.2}
\end{equation}
compte tenu de (\hyperref[II.8.2.3.2-fr]{8.2.3.2}) et
(\hyperref[II.8.2.5.2-fr]{8.2.5.2}).
\end{env}

\begin{env}[8.12.4]
\label{II.8.12.4-fr}
Plaçons-nous dans les hypothèses de (\hyperref[II.8.5.1-fr]{8.5.1}), dont
nous conservons les notations. Il résulte de
(\hyperref[II.1.5.6-fr]{1.5.6}) que pour tout $\mathcal{S}$-Module gradué
quasi-cohérent $\mathcal{M}$, on a d'une part un isomorphisme canonique
\begin{equation}
\label{II.8.12.4.1-fr}
  \Phi^*(\widetilde{\mathcal{M}})
  \xrightarrow{\sim}
  \left(q^*(\mathcal{M})\otimes_{q^*(\mathcal{S})}\mathcal{S}'\right)^{\sim}
\tag{8.12.4.1}
\end{equation}
de $\mathcal{O}_{C'}$-Modules~; et d'autre part
(\hyperref[II.3.5.6-fr]{3.5.6}) entraîne l'existence d'un
$\operatorname{Proj}(\varphi)$-morphisme canonique
\begin{equation}
\label{II.8.12.4.2-fr}
  \operatorname{Proj}_0(\mathcal{M})\to
  \left(\operatorname{Proj}_0
  (q^*(\mathcal{M})\otimes_{q^*(\mathcal{S})}\mathcal{S}')\right)
  |G(\varphi)
\tag{8.12.4.2}
\end{equation}
et aussi d'un $\widehat{\Phi}$-morphisme canonique
\begin{equation}
\label{II.8.12.4.3-fr}
  \operatorname{Proj}_0(\widehat{\mathcal{M}})\to
  \left(\operatorname{Proj}_0
  (q^*(\widehat{\mathcal{M}})
  \otimes_{q^*(\widehat{\mathcal{S}})}\widehat{\mathcal{S}'})\right)
  |G(\widehat{\varphi}).
\tag{8.12.4.3}
\end{equation}
\end{env}

\begin{env}[8.12.5]
\label{II.8.12.5-fr}
Considérons maintenant la situation de (\hyperref[II.8.6.1-fr]{8.6.1}),
avec les mêmes notations~; on prend donc dans ce qui précède $Y'=X$, le
morphisme $q:X\to Y$ étant le morphisme structural, et $\varphi$ est le
$q$-morphisme canonique (\hyperref[II.8.6.1.2-fr]{8.6.1.2}). On a alors un
isomorphisme canonique
\begin{equation}
\label{II.8.12.5.1-fr}
  q^*(\mathcal{M})\otimes_{q^*(\mathcal{S})}\mathcal{S}_X^{\geq}
  \xrightarrow{\sim}\mathcal{M}_X^{\geq}
\tag{8.12.5.1}
\end{equation}
en posant
$\mathcal{M}_X^{\geq}=\bigoplus_{n\geq0}\operatorname{Proj}_0(\mathcal{M}(n))$.
On peut en effet se borner au cas où $Y=\operatorname{Spec}(A)$ est affine,
$\mathcal{S}=\widetilde{S}$ et $\mathcal{M}=\widetilde{M}$, et définir
l'isomorphisme (\hyperref[II.8.12.5.1-fr]{8.12.5.1}) dans chacun des ouverts
affines $X_f$ ($f$ homogène dans $S_+$), en vérifiant la compatibilité avec le
passage à un multiple homogène de $f$. Or, la restriction à $X_f$ du premier
membre de (\hyperref[II.8.12.5.1-fr]{8.12.5.1}) est
$\widetilde{M}'=
\left((M\otimes_A S_{(f)})
\otimes_{S\otimes_A S_{(f)}}S_f^{\geq}\right)^{\sim}$ par
(\hyperref[II.8.6.2.1-fr]{8.6.2.1})~; comme on a un isomorphisme canonique de
$M\otimes_A S_{(f)}$ sur $M\otimes_S(S\otimes_A S_{(f)})$, on en déduit un
isomorphisme de $\widetilde{M}'$ sur
$\left(M\otimes_S S_f^{\geq}\right)^{\sim}$, et ce dernier est
canoniquement isomorphe par (\hyperref[II.8.2.9.1-fr]{8.2.9.1}) à la
restriction à $X_f$ du second membre de
(\hyperref[II.8.12.5.1-fr]{8.12.5.1}), et satisfait aux conditions de
compatibilité requises.

\oldpage[II]{194}
Remplaçant $\mathcal{M}$ par $\widehat{\mathcal{M}}$, $\mathcal{S}$ par
$\widehat{\mathcal{S}}$ et $\mathcal{S}_X$ par
$(\mathcal{S}_X^{\geq})^{\wedge}$ dans le raisonnement précédent, on a de
même un isomorphisme canonique
\begin{equation}
\label{II.8.12.5.2-fr}
  q^*(\widehat{\mathcal{M}})
  \otimes_{q^*(\widehat{\mathcal{S}})}(\mathcal{S}_X^{\geq})^{\wedge}
  \xrightarrow{\sim}(\mathcal{M}_X^{\geq})^{\wedge}.
\tag{8.12.5.2}
\end{equation}
Si l'on se souvient (\hyperref[II.8.6.2-fr]{8.6.2}) que le morphisme
structural $u:\operatorname{Proj}(\mathcal{S}_X^{\geq})\to X$ est un
isomorphisme, on déduit d'abord de ce qui précède que l'on a un
$u$-isomorphisme canonique
\begin{equation}
\label{II.8.12.5.3-fr}
  \operatorname{Proj}_0\mathcal{M}
  \xrightarrow{\sim}\operatorname{Proj}_0(\mathcal{M}_X^{\geq})
\tag{8.12.5.3}
\end{equation}
comme cas particulier de (\hyperref[II.8.12.4.2-fr]{8.12.4.2}). On constate
en effet, avec les notations de la démonstration de
(\hyperref[II.8.6.2-fr]{8.6.2}), que cela revient à voir que
l'homomorphisme canonique
$M_{(f)}\otimes_{S_{(f)}}(S_f^{\geq})^{(d)}\to(M_f^{\geq})^{(d)}$ est un
isomorphisme lorsque $f\in S_d$, ce qui est immédiat.

En second lieu, l'isomorphisme (\hyperref[II.8.12.5.2-fr]{8.12.5.2}) permet,
en appliquant cette fois (\hyperref[II.8.12.4.3-fr]{8.12.4.3}) au morphisme
canonique $r=\operatorname{Proj}(\widehat{\alpha}):\widehat{C}_X\to
\widehat{C}$, d'obtenir un $r$-morphisme canonique
\begin{equation}
\label{II.8.12.5.4-fr}
  \mathcal{M}^{\square}\to(\mathcal{M}_X^{\geq})^{\square}.
\tag{8.12.5.4}
\end{equation}
Rappelons maintenant (\hyperref[II.8.6.2-fr]{8.6.2}) que les restrictions de
$r$ aux cônes épointés $\widehat{E}_X$ et $E_X$ sont des \emph{isomorphismes}
sur $\widehat{E}$ et $E$ respectivement. En outre~:
\end{env}

\begin{proposition}[8.12.6]
\label{II.8.12.6-fr}
Les restrictions à $\widehat{E}_X$ et à $E_X$ du $r$-morphisme canonique
(\hyperref[II.8.12.5.4-fr]{8.12.5.4}) sont des isomorphismes
\begin{equation}
\label{II.8.12.6.1-fr}
  \mathcal{M}^{\square}|\widehat{E}
  \xrightarrow{\sim}(\mathcal{M}_X^{\geq})^{\square}|\widehat{E}_X
\tag{8.12.6.1}
\end{equation}
\begin{equation}
\label{II.8.12.6.2-fr}
  \widetilde{\mathcal{M}}|E
  \xrightarrow{\sim}(\mathcal{M}_X^{\geq})^{\sim}|E_X.
\tag{8.12.6.2}
\end{equation}
\end{proposition}

\begin{proof}
On se ramène au cas où $Y$ est affine comme dans la démonstration de
(\hyperref[II.8.6.2-fr]{8.6.2})~; avec les notations de cette dernière, et
en se ramenant aux définitions (\hyperref[II.2.8.8-fr]{2.8.8}), on doit
montrer que l'homomorphisme canonique
\[
  \widehat{M}_{(f)}\otimes_{\widehat{S}_{(f)}}
  (S_f^{\geq})_{(f/1)}^{\wedge}
  \to(M\otimes_S S_f^{\geq})_{(f/1)}^{\wedge}
\]
est un isomorphisme~; mais en vertu de
(\hyperref[II.8.2.3.2-fr]{8.2.3.2}) et
(\hyperref[II.8.2.5.2-fr]{8.2.5.2}), le premier membre s'identifie
canoniquement à
$M_f^{\leq}\otimes_{S_f^{\leq}}(S_f^{\geq})_{f/1}^{\leq}$, donc à
$M_f^{\leq}$ en vertu de (\hyperref[II.8.2.7.2-fr]{8.2.7.2}), et le second à
$(M_f^{\geq})_{f/1}^{\leq}$, donc aussi à $M_f^{\leq}$ en vertu de
(\hyperref[II.8.2.9.2-fr]{8.2.9.2}), d'où la conclusion en ce qui concerne
(\hyperref[II.8.12.6.1-fr]{8.12.6.1})~; la formule
(\hyperref[II.8.12.6.2-fr]{8.12.6.2}) résulte alors de
(\hyperref[II.8.12.6.1-fr]{8.12.6.1}) et de
(\hyperref[II.8.12.2.1-fr]{8.12.2.1}).
\end{proof}

\begin{corollary}[8.12.7]
\label{II.8.12.7-fr}
Avec les identifications de (\hyperref[II.8.6.3-fr]{8.6.3}), la restriction
de $(\mathcal{M}_X^{\geq})^{\square}$ à $\widehat{E}_X$ s'identifie à
$(\mathcal{M}_X^{\leq})^{\sim}$ et la restriction de
$(\mathcal{M}_X^{\geq})^{\square}$ à $E_X$ s'identifie à
$\widetilde{\mathcal{M}}_X$.
\end{corollary}

\begin{proof}
On est en effet ramené au cas affine, et cela résulte de l'identification de
$(M_f^{\geq})_{f/1}^{\leq}$ à $M_f^{\leq}$ et de
$(M_f^{\geq})_{f/1}$ à $M_f$ (\hyperref[II.8.2.9.2-fr]{8.2.9.2}).
\end{proof}

\begin{proposition}[8.12.8]
\label{II.8.12.8-fr}
Sous les hypothèses de (\hyperref[II.8.6.4-fr]{8.6.4}), l'homomorphisme
canonique (\hyperref[II.8.12.3.1-fr]{8.12.3.1}) est un isomorphisme.
\end{proposition}

\begin{proof}
Compte tenu du fait que $\operatorname{Proj}(\mathcal{S}_X^{\geq})\to X$
est un isomorphisme (\hyperref[II.8.6.2-fr]{8.6.2}) et des

\oldpage[II]{195}
isomorphismes (\hyperref[II.8.12.5.4-fr]{8.12.5.4}) et
(\hyperref[II.8.12.6.1-fr]{8.12.6.1}), on est ramené à démontrer la
proposition correspondante pour l'homomorphisme canonique
$p_X^*(\operatorname{Proj}_0(\mathcal{M}_X^{\geq}))
\to(\mathcal{M}_X^{\geq})^{\square}|\widehat{E}_X$, autrement dit, on est
ramené au cas où $\mathcal{S}_1$ est un $\mathcal{O}_Y$-Module inversible et
où $\mathcal{S}$ est engendrée par $\mathcal{S}_1$. Avec les notations de
(\hyperref[II.8.12.3-fr]{8.12.3}), on a alors, pour un $f\in S_1$,
$S_f^{\leq}=S_{(f)}[1/f]$ et l'homomorphisme canonique
$M_{(f)}\otimes_{S_{(f)}}S_f^{\leq}\to M_f^{\leq}$ est un isomorphisme par
définition de $M_f^{\leq}$.
\end{proof}

\begin{env}[8.12.9]
\label{II.8.12.9-fr}
Considérons maintenant les $\mathcal{S}$-Modules quasi-cohérents
\[
  \mathcal{M}_{[n]}=\bigoplus_{m\geqslant n}\mathcal{M}_m
\]
et (avec les notations de (\hyperref[II.8.7.2-fr]{8.7.2})) le
$\mathcal{S}^{\natural}$-Module gradué quasi-cohérent
\begin{equation}
\label{II.8.12.9.1-fr}
  \mathcal{M}^{\natural}
  =\left(\bigoplus_{n\geqslant0}\mathcal{M}_{[n]}\right)^{\sim}.
\tag{8.12.9.1}
\end{equation}
On a vu (\hyperref[II.8.7.3-fr]{8.7.3}) qu'il existe un $C$-isomorphisme
canonique $h:C_X\xrightarrow{\sim}\operatorname{Proj}(\mathcal{S}^{\natural})$.
En outre~:
\end{env}

\begin{proposition}[8.12.10]
\label{II.8.12.10-fr}
Il existe un $h$-isomorphisme canonique
\begin{equation}
\label{II.8.12.10.1-fr}
  \operatorname{Proj}_0(\mathcal{M}^{\natural})
  \xrightarrow{\sim}\widetilde{\mathcal{M}}_X.
\tag{8.12.10.1}
\end{equation}
\end{proposition}

\begin{proof}
On raisonne comme dans (\hyperref[II.8.7.3-fr]{8.7.3}), en utilisant cette
fois l'existence du di-isomorphisme (\hyperref[II.8.2.9.3-fr]{8.2.9.3}) au
lieu de (\hyperref[II.8.2.7.3-fr]{8.2.7.3}). Nous laissons les détails au
lecteur.
\end{proof}

\subsection{Fermetures projectives de sous-faisceaux et de sous-schémas
fermés}
\label{subsection:II.8.13-fr}

\begin{env}[8.13.1]
\label{II.8.13.1-fr}
Les hypothèses et notations étant celles de
(\hyperref[II.8.12.1-fr]{8.12.1}), considérons un sous-$\mathcal{S}$-Module
quasi-cohérent $\mathcal{N}$ de $\mathcal{M}$, \emph{non nécessairement
gradué}. On peut donc considérer le $\mathcal{O}_C$-Module quasi-cohérent
$\widetilde{\mathcal{N}}$ associé à $\mathcal{N}$, qui est un
sous-$\mathcal{O}_C$-Module de $\widetilde{\mathcal{M}}$. On a vu d'autre
part (\hyperref[II.8.12.2.1-fr]{8.12.2.1}) que
$\widetilde{\mathcal{M}}$ s'identifie à la restriction de
$\mathcal{M}^{\square}$ à $C$. Comme l'injection canonique
$i:C\to\widehat{C}$ est un morphisme affine
(\hyperref[II.8.3.2-fr]{8.3.2}), et \emph{a fortiori} quasi-compact, le
\emph{prolongement canonique} $(\widetilde{\mathcal{N}})^{-}$, plus grand
sous-$\mathcal{O}_{\widehat{C}}$-Module contenu dans
$\mathcal{M}^{\square}$ et induisant $\widetilde{\mathcal{N}}$ sur $C$, est
un $\mathcal{O}_{\widehat{C}}$-Module quasi-cohérent
(\hyperref[I.9.4.2-fr]{I, 9.4.2}). Nous allons en donner une description
explicite à l'aide d'un $\widehat{\mathcal{S}}$-Module gradué.
\end{env}

\begin{env}[8.13.2]
\label{II.8.13.2-fr}
Pour cela, considérons, pour tout entier $n\geqslant0$, l'homomorphisme
$\bigoplus_{i\leqslant n}\mathcal{M}_i\to\mathcal{M}$ qui, pour tout ouvert
$U$ de $Y$, fait correspondre à la famille
\[
  (s_i)\in\bigoplus_{i\leqslant n}\Gamma(U,\mathcal{M}_i)
\]
la section $\sum_i s_i\in\Gamma(U,\mathcal{M})$. Désignons par
$\mathcal{N}'_n$ l'image réciproque de $\mathcal{N}$ par cet homomorphisme,
qui est un sous-$\mathcal{O}_Y$-Module quasi-cohérent de
$\bigoplus_{i\leqslant n}\mathcal{M}_i$. Considérons maintenant
l'homomorphisme
$\bigoplus_{i\leqslant n}\mathcal{M}_i\to
\widehat{\mathcal{M}}=\mathcal{M}[z]$ qui fait correspondre à $(s_i)$ la
section
$\sum_{i\leqslant n}s_i z^{n-i}\in\Gamma(U,\widehat{\mathcal{M}}_n)$, et
soit $\mathcal{N}_n$ l'image de $\mathcal{N}'_n$ par cet homomorphisme~; on
vérifie aussitôt que
$\overline{\mathcal{N}}=\bigoplus_{n\geqslant0}\mathcal{N}_n$ est un
sous-$\widehat{\mathcal{S}}$-Module gradué (quasi-cohérent) de
$\widehat{\mathcal{M}}$~; on dit que $\overline{\mathcal{N}}$ est déduit de
$\mathcal{N}$ par \emph{homogénéisation}, à l'aide de la «~variable
homogénéisante~» $z$. On notera

\oldpage[II]{196}
que si $\mathcal{N}$ est déjà un sous-$\mathcal{S}$-Module \emph{gradué} de
$\mathcal{M}$, alors $\overline{\mathcal{N}}$ s'identifie à la somme directe
des composantes $\widehat{\mathcal{N}}_n$ de degré $n\geqslant0$ dans
$\widehat{\mathcal{N}}=\mathcal{N}[z]$.
\end{env}

\begin{proposition}[8.13.3]
\label{II.8.13.3-fr}
Le $\mathcal{O}_{\widehat{C}}$-Module
$\operatorname{Proj}_0(\overline{\mathcal{N}})$ est le prolongement canonique
$(\widetilde{\mathcal{N}})^{-}$ de $\widetilde{\mathcal{N}}$ à
$\widehat{C}$.
\end{proposition}

\begin{proof}
La question est locale sur $Y$ et $\widehat{C}$ en vertu de la définition du
prolongement canonique (\hyperref[I.9.4.1-fr]{I, 9.4.1}). On peut donc déjà
supposer que $Y=\operatorname{Spec}(A)$ est affine, avec
$\mathcal{S}=\widetilde{S}$, $\mathcal{M}=\widetilde{M}$ et
$\mathcal{N}=\widetilde{N}$, où $N$ est un sous-$S$-module non
nécessairement gradué de $M$. En outre
(\hyperref[II.8.3.2.6-fr]{8.3.2.6}) $\widehat{C}$ est réunion des ouverts
affines $\widehat{C}_z=C$ et
$\widehat{C}_f=\operatorname{Spec}(S_f^{\leq})$ ($f$ homogène dans $S_+$).
Il suffira donc de montrer que~: $1^\circ$ la restriction de
$\operatorname{Proj}_0(\overline{\mathcal{N}})$ à $C$ est
$\widetilde{\mathcal{N}}$~; $2^\circ$ la restriction de
$\operatorname{Proj}_0(\overline{\mathcal{N}})$ à chaque $\widehat{C}_f$ est
le prolongement canonique de la restriction de $\mathcal{N}$ à
$C\cap\widehat{C}_f=\operatorname{Spec}(S_f)$
(\hyperref[II.8.3.2.6-fr]{8.3.2.6}). Pour le premier point, notons que
$\operatorname{Proj}_0(\overline{\mathcal{N}})|C$ s'identifie à
$(\overline{N}_{(z)})^{\sim}$
(\hyperref[II.8.3.2.4-fr]{8.3.2.4})~; or $\overline{N}_{(z)}$ s'identifie
canoniquement (\hyperref[II.2.2.5-fr]{2.2.5}) à l'image de
$\overline{N}$ dans $\widehat{M}/(z-1)\widehat{M}$, et par l'isomorphisme
canonique de ce dernier sur $M$ (\hyperref[II.8.2.5-fr]{8.2.5}), cette image
s'identifie à $N$, en vertu de la définition de $\overline{N}$ donnée dans
(\hyperref[II.8.13.2-fr]{8.13.2}).

Pour prouver le second point, notons que l'injection
$i:C\cap\widehat{C}_f\to\widehat{C}_f$ correspond à l'injection canonique
$S_f^{\leq}\to S_f$ (\hyperref[II.8.3.2.6-fr]{8.3.2.6})~; d'autre part, on
a $\Gamma(\widehat{C}_f,\mathcal{M}^{\square})=M_f^{\leq}$,
$\Gamma(\widehat{C}_f,i_*(\widetilde{\mathcal{N}}))=N_f$ et (par
(\hyperref[II.8.12.2.1-fr]{8.12.2.1}))
$\Gamma(\widehat{C}_f,i_*(i^*(\mathcal{M}^{\square})))=M_f$. Compte tenu de
(\hyperref[I.9.4.2-fr]{I, 9.4.2}), on est donc ramené à montrer que
$\overline{N}_{(f)}\subset\widehat{M}_{(f)}=M_f^{\leq}$ s'identifie
canoniquement à l'image réciproque de $N_f$, par l'injection canonique
$M_f^{\leq}\to M_f$. En effet, soit $d=\deg(f)>0$, et supposons qu'un
élément $(\sum_{k\leqslant md}x_k)/f^m$ de $M_f$ (avec $x_k\in M_k$) soit de
la forme $y/f^m$ avec $y\in N$. En multipliant $y$ et les $x_k$ par un même
$f^h$ convenable, on peut déjà supposer que $\sum_{k\leqslant md}x_k=y$.
Mais dans l'identification (\hyperref[II.8.2.5.2-fr]{8.2.5.2}),
$(\sum_{k\leqslant md}x_k)/f^m$ correspond à
$\sum_{k\leqslant md}x_kz^{md-k}/f^m$, et cet élément appartient bien à
$\overline{N}_{(f)}$, puisque $\sum_{k\leqslant md}x_k\in N$~; la réciproque
est évidente.
\end{proof}

\begin{remark}[8.13.4]
\label{II.8.13.4-fr}
\begin{enumerate}
  \item[\rm{(i)}] Le cas d'application le plus important de
  (\hyperref[II.8.13.3-fr]{8.13.3}) est celui où
  $\mathcal{M}=\mathcal{S}$, $\widetilde{\mathcal{N}}$ étant donc un faisceau
  quasi-cohérent \emph{arbitraire} d'idéaux $\mathcal{J}$ de
  $\mathcal{O}_C$ (\hyperref[II.1.4.3-fr]{1.4.3}), correspondant
  biunivoquement à un \emph{sous-préschéma fermé} $Z$ de $C$. Alors le
  prolongement canonique $\overline{\mathcal{J}}$ de $\mathcal{J}$ est le
  faisceau quasi-cohérent d'idéaux de $\mathcal{O}_{\widehat{C}}$ qui définit
  l'\emph{adhérence} $\overline{Z}$ de $Z$ dans $\widehat{C}$
  (\hyperref[I.9.5.10-fr]{I, 9.5.10})~; la prop.
  (\hyperref[II.8.13.3-fr]{8.13.3}) donne un moyen canonique de définir
  $\overline{Z}$ à l'aide d'un Idéal gradué dans
  $\widehat{\mathcal{S}}=\mathcal{S}[z]$.
  \item[\rm{(ii)}] Supposons pour simplifier que $Y$ soit affine, et gardons
  les notations de la démonstration de
  (\hyperref[II.8.13.3-fr]{8.13.3}). Pour tout $x\in N$ non nul, soit $d(x)$
  le plus grand des degrés des composantes homogènes $x_i$ de $x$ dans $M$~;
  $\overline{N}$ est par définition le sous-module de $\widehat{M}$ formé de
  $0$ et des éléments de la forme
  $h(x,k)=z^k\sum_{i\leqslant d(x)}x_i z^{d(x)-i}$ ($k$ entier
  $\geqslant0$)~; il est donc engendré, en tant que module sur
  $\widehat{S}=S[z]$, par les éléments de la forme
  \[
    h(x,0)=\sum_{i\leqslant d(x)}x_i z^{d(x)-i}.
  \]

  \oldpage[II]{197}
  On dit que $h(x,0)$ est déduit de $x$ par \emph{homogénéisation} à l'aide
  de la «~variable homogénéisante~» $z$. Mais comme $h(x,0)$ ne dépend pas
  additivement de $x$ (ni \emph{a fortiori} $S$-linéairement), \emph{on se
  gardera de croire} (même lorsque $M=S$) que les $h(x,0)$ parcourent un
  \emph{système de générateurs} du $\widehat{S}$-module gradué
  $\overline{N}$ lorsqu'on fait parcourir à $x$ un \emph{système de
  générateurs} du $S$-module $N$. Il en est cependant bien ainsi dans le cas
  (seul considéré dans la géométrie algébrique élémentaire) où $N$ est un
  $S$-module \emph{libre monogène}, car si $t$ est une base de $N$, $h(t,0)$
  engendre le $\widehat{S}$-module $\overline{N}$.
\end{enumerate}
\end{remark}

\subsection{Compléments sur les faisceaux associés aux
$\mathcal{S}$-Modules gradués}
\label{subsection:II.8.14-fr}

\begin{env}[8.14.1]
\label{II.8.14.1-fr}
Soient $Y$ un préschéma, $\mathcal{S}$ une $\mathcal{O}_Y$-Algèbre graduée
quasi-cohérente à \emph{degrés positifs},
$X=\operatorname{Proj}(\mathcal{S})$, $q:X\to Y$ le morphisme structural
(séparé en vertu de (\hyperref[II.3.1.3-fr]{3.1.3})). Conservons les
notations de (\hyperref[II.8.12.1-fr]{8.12.1}), de sorte que nous avons
défini un foncteur
$\mathcal{M}_X=\operatorname{Proj}(\mathcal{M})$ en $\mathcal{M}$, de la
catégorie des $\mathcal{S}$-Modules gradués quasi-cohérents dans celle des
$\mathcal{S}_X$-Modules gradués quasi-cohérents~; il est clair en outre
(\hyperref[II.3.2.4-fr]{3.2.4}) que c'est un foncteur \emph{additif exact},
commutant aux limites inductives.

Notons en outre qu'il résulte aussitôt de la définition
(\hyperref[II.8.12.1.1-fr]{8.12.1.1}) que l'on a
\begin{equation}
\label{II.8.14.1.1-fr}
  \operatorname{Proj}(\mathcal{M}(n))
  =(\operatorname{Proj}(\mathcal{M}))(n)
  \qquad\text{pour tout }n\in\mathbf{Z}.
\tag{8.14.1.1}
\end{equation}
\end{env}

\begin{env}[8.14.2]
\label{II.8.14.2-fr}
Nous allons d'abord étendre aux $\mathcal{S}_X$-Modules de la forme
$\operatorname{Proj}(\mathcal{M})$ les homomorphismes canoniques $\lambda$
et $\mu$ définis dans (\hyperref[II.3.2.6-fr]{3.2.6}). Pour cela, notons
que, quels que soient $m\in\mathbf{Z}$ et $n\in\mathbf{Z}$, on a, en vertu
de (\hyperref[II.2.1.2.1-fr]{2.1.2.1}), un homomorphisme canonique de
$\mathcal{O}_X$-Modules
\begin{equation}
\label{II.8.14.2.1-fr}
  \lambda_{mn}:\operatorname{Proj}_0(\mathcal{M}(m))
  \otimes_{\mathcal{O}_X}\operatorname{Proj}_0(\mathcal{N}(n))
  \to\operatorname{Proj}_0
  ((\mathcal{M}\otimes_{\mathcal{S}}\mathcal{N})(m+n))
\tag{8.14.2.1}
\end{equation}
et de même, en vertu de (\hyperref[II.2.1.2.2-fr]{2.1.2.2}), un
homomorphisme canonique de $\mathcal{O}_X$-Modules
\begin{equation}
\label{II.8.14.2.2-fr}
  \mu_{mn}:\operatorname{Proj}_0
  ((\mathcal{H}om_{\mathcal{S}}(\mathcal{M},\mathcal{N}))(n-m))
  \to\mathcal{H}om_{\mathcal{O}_X}
  (\operatorname{Proj}_0(\mathcal{M}(m)),
   \operatorname{Proj}_0(\mathcal{N}(n)))
\tag{8.14.2.2}
\end{equation}
quels que soient les $\mathcal{S}$-Modules gradués quasi-cohérents
$\mathcal{M}$, $\mathcal{N}$. On en déduit un homomorphisme
\[
  \mu_k:\operatorname{Proj}_0
  ((\mathcal{H}om_{\mathcal{S}}(\mathcal{M},\mathcal{N}))(k))
  \to
  (\mathcal{H}om_{\mathcal{S}_X}
  (\operatorname{Proj}(\mathcal{M}),\operatorname{Proj}(\mathcal{N})))_k
\]
obtenu en faisant correspondre à tout
$u\in\Gamma(U,\operatorname{Proj}_0
((\mathcal{H}om_{\mathcal{S}}(\mathcal{M},\mathcal{N}))(k)))$
l'homomorphisme $\mu_k(u)$, de degré $k$, de $\mathbf{Z}$-modules gradués
$\Gamma(U,\operatorname{Proj}(\mathcal{M}))\to
\Gamma(U,\operatorname{Proj}(\mathcal{N}))$ ($U$ ouvert de $X$) qui dans
chaque $\Gamma(U,\operatorname{Proj}_0(\mathcal{M}(m)))$ coïncide avec
$\mu_{m,m+k}(u)$~; en outre, en revenant à la définition des $\mu_{mn}$
(\hyperref[II.2.5.12.1-fr]{2.5.12.1}), on voit aussitôt que $\mu_k(u)$ est
en fait un homomorphisme de degré $k$ de
$\Gamma(U,\mathcal{S}_X)$-modules gradués, et en outre que les $\mu_k$
définissent un homomorphisme de $\mathcal{S}_X$-Modules gradués
\begin{equation}
\label{II.8.14.2.3-fr}
  \operatorname{Proj}
  (\mathcal{H}om_{\mathcal{S}}(\mathcal{M},\mathcal{N}))
  \to\mathcal{H}om_{\mathcal{S}_X}
  (\operatorname{Proj}(\mathcal{M}),\operatorname{Proj}(\mathcal{N})).
\tag{8.14.2.3}
\end{equation}

De même, compte tenu du diagramme d'associativité
(\hyperref[II.2.5.11.4-fr]{2.5.11.4}), les homomorphismes
(\hyperref[II.8.14.2.1-fr]{8.14.2.1}) donnent un homomorphisme de
$\mathcal{S}_X$-Modules gradués
\begin{equation}
\label{II.8.14.2.4-fr}
  \lambda:\operatorname{Proj}(\mathcal{M})
  \otimes_{\mathcal{S}_X}\operatorname{Proj}(\mathcal{N})
  \to\operatorname{Proj}(\mathcal{M}\otimes_{\mathcal{S}}\mathcal{N}).
\tag{8.14.2.4}
\end{equation}
\end{env}

\oldpage[II]{198}
\begin{proposition}[8.14.3]
\label{II.8.14.3-fr}
L'homomorphisme (\hyperref[II.8.14.2.4-fr]{8.14.2.4}) est bijectif~; il en
est de même de (\hyperref[II.8.14.2.3-fr]{8.14.2.3}) lorsque le
$\mathcal{S}$-Module gradué $\mathcal{M}$ admet une présentation finie
(\hyperref[II.3.1.1-fr]{3.1.1}).
\end{proposition}

\begin{proof}
La question est évidemment locale sur $X$ et sur $Y$~; on peut donc supposer
$Y=\operatorname{Spec}(A)$ affine, $\mathcal{S}=\widetilde{S}$,
$\mathcal{M}=\widetilde{M}$, $\mathcal{N}=\widetilde{N}$, où $S$ est une
$A$-algèbre graduée à degrés positifs, $M$ et $N$ deux $S$-modules gradués.
Si $f$ est un élément homogène de $S_+$, les homomorphismes
(\hyperref[II.8.14.2.1-fr]{8.14.2.1}) et
(\hyperref[II.8.14.2.2-fr]{8.14.2.2}), restreints à l'ouvert affine
$D_+(f)$, correspondent aux homomorphismes canoniques
(\hyperref[II.2.5.11.1-fr]{2.5.11.1}) et
(\hyperref[II.2.5.12.1-fr]{2.5.12.1})~:
\begin{align*}
  M(m)_{(f)}\otimes_{S_{(f)}}N(n)_{(f)}
  &\to(M\otimes_S N)(m+n)_{(f)},\\
  (\operatorname{Hom}_S(M,N))(n-m)_{(f)}
  &\to\operatorname{Hom}_{S_{(f)}}(M(m)_{(f)},N(n)_{(f)}).
\end{align*}
Si on se reporte aux définitions de ces homomorphismes, on voit donc (compte
tenu de (\hyperref[II.8.2.9.1-fr]{8.2.9.1})) que la restriction de
(\hyperref[II.8.14.2.4-fr]{8.14.2.4}) à $D_+(f)$ correspond à
l'homomorphisme canonique
\[
  M_f\otimes_{S_f}N_f\to(M\otimes_S N)_f
\]
défini dans (\hyperref[0.1.3.4-fr]{0, 1.3.4}), et on sait que ce dernier est
un isomorphisme. De même, la restriction de
(\hyperref[II.8.14.2.3-fr]{8.14.2.3}) à $D_+(f)$ correspond à
l'homomorphisme canonique (\hyperref[0.1.3.5-fr]{0, 1.3.5})
\[
  (\operatorname{Hom}_S(M,N))_f\to
  \operatorname{Hom}_{S_f}(M_f,N_f)
\]
compte tenu de ce que, $M$ étant de type fini, le module
$\operatorname{Hom}_S(M,N)$, somme directe des sous-groupes formés des
homomorphismes \emph{homogènes} de $S$-modules
(\hyperref[II.2.1.2-fr]{2.1.2}), coïncide avec l'ensemble de \emph{tous} les
homomorphismes $M\to N$ de $S$-modules. L'hypothèse que $M$ admet une
présentation finie entraîne alors (\hyperref[0.1.3.5-fr]{0, 1.3.5}) que
l'homomorphisme canonique considéré est bien un isomorphisme.
\end{proof}

\begin{proposition}[8.14.4]
\label{II.8.14.4-fr}
Si $U$ est un ouvert quasi-compact de $X$, il existe un entier $d$ tel que
pour tout entier $n$ multiple de $d$, $\mathcal{O}_X(n)|U$ soit inversible,
son inverse étant $\mathcal{O}_X(-n)|U$.
\end{proposition}

\begin{proof}
Comme $q(U)$ est quasi-compact, il est recouvert par un nombre fini d'ouverts
affines $V_i$, donc tout $x\in U$ est contenu dans un ouvert affine de la
forme $D_+(f)$, où $f$ est un élément homogène de degré $>0$ de l'un des
anneaux $\Gamma(V_i,\mathcal{S})$. Comme $U$ est quasi-compact, on peut le
recouvrir par un nombre fini de tels ouverts $D_+(f_j)$~; soit $d$ un
multiple commun des degrés des $f_j$. Le nombre $d$ répond à la question en
vertu de (\hyperref[II.2.5.17-fr]{2.5.17}).
\end{proof}

\begin{env}[8.14.5]
\label{II.8.14.5-fr}
Avec les hypothèses et notations de (\hyperref[II.8.14.1-fr]{8.14.1}), on a
défini dans (\hyperref[II.3.3.2-fr]{3.3.2}) des homomorphismes canoniques de
$\mathcal{O}_Y$-Modules
\begin{equation}
\label{II.8.14.5.1-fr}
  \alpha_n:\mathcal{M}_n\to
  q_*(\operatorname{Proj}_0(\mathcal{M}(n)))
  \qquad(n\in\mathbf{Z})
\tag{8.14.5.1}
\end{equation}
Généralisant les notations de (\hyperref[II.3.3.1-fr]{3.3.1}), nous
poserons, pour tout $\mathcal{S}_X$-Module gradué $\mathcal{F}$,
\begin{equation}
\label{II.8.14.5.2-fr}
  \Gamma_*(\mathcal{F})=
  \bigoplus_{n\in\mathbf{Z}}q_*(\mathcal{F}_n).
\tag{8.14.5.2}
\end{equation}
En particulier,
$\Gamma_*(\mathcal{S}_X)=\bigoplus_{n\in\mathbf{Z}}q_*(\mathcal{O}_X(n))$
est la $\mathcal{O}_Y$-Algèbre graduée notée
$\Gamma_*(\mathcal{O}_X)$ dans
(\hyperref[II.3.3.1.2-fr]{3.3.1.2})~; il est clair que
$\Gamma_*(\mathcal{F})$ est un $\Gamma_*(\mathcal{S}_X)$-Module gradué
(\hyperref[0.4.2.2-fr]{0, 4.2.2}). Lorsque

\oldpage[II]{199}
dans les homomorphismes (\hyperref[II.8.14.5.1-fr]{8.14.5.1}), on prend
$\mathcal{M}=\mathcal{S}$, on obtient l'homomorphisme de
$\mathcal{O}_Y$-Algèbres graduées
\begin{equation}
\label{II.8.14.5.3-fr}
  \alpha:\mathcal{S}\to\Gamma_*(\mathcal{S}_X)
\tag{8.14.5.3}
\end{equation}
déjà défini dans (\hyperref[II.3.3.2-fr]{3.3.2}), et qui fait donc de
$\Gamma_*(\mathcal{F})$ un $\mathcal{S}$-Module gradué~; les homomorphismes
(\hyperref[II.8.14.5.1-fr]{8.14.5.1}) définissent alors un homomorphisme (de
degré $0$) de $\mathcal{S}$-Modules gradués
\begin{equation}
\label{II.8.14.5.4-fr}
  \alpha:\mathcal{M}\to\Gamma_*(\operatorname{Proj}(\mathcal{M})).
\tag{8.14.5.4}
\end{equation}
\end{env}

\begin{env}[8.14.6]
\label{II.8.14.6-fr}
En général, pour un $\mathcal{S}_X$-Module gradué quasi-cohérent
$\mathcal{F}$, il n'est pas certain que le $\mathcal{S}$-Module gradué
$\Gamma_*(\mathcal{F})$ soit nécessairement quasi-cohérent. Considérons un
ouvert $X'$ de $X$ tel que la restriction $q':X'\to Y$ de $q$ à $X'$ soit
un morphisme \emph{quasi-compact}. Comme $q'$ est en outre séparé,
$q'_*(\mathcal{F}')$ est alors un $\mathcal{O}_Y$-Module quasi-cohérent pour
tout $\mathcal{O}_{X'}$-Module quasi-cohérent $\mathcal{F}'$
(\hyperref[I.9.2.2-fr]{I, 9.2.2, b)}). Nous poserons
\begin{equation}
\label{II.8.14.6.1-fr}
  \mathcal{S}_{X'}=\mathcal{S}_X|X'
  =\bigoplus_{n\in\mathbf{Z}}\mathcal{O}_X(n)|X'
\tag{8.14.6.1}
\end{equation}
et, pour tout $\mathcal{S}_{X'}$-Module gradué $\mathcal{F}'$,
\begin{equation}
\label{II.8.14.6.2-fr}
  \Gamma'_*(\mathcal{F}')=
  \bigoplus_{n\in\mathbf{Z}}q'_*(\mathcal{F}'_n).
\tag{8.14.6.2}
\end{equation}
La remarque précédente montre alors que si $\mathcal{F}'$ est un
$\mathcal{S}_{X'}$-Module quasi-cohérent, $\Gamma'_*(\mathcal{F}')$ est un
$\mathcal{S}$-Module gradué quasi-cohérent
(\hyperref[I.9.6.1-fr]{I, 9.6.1}).

On notera d'ailleurs que l'injection canonique $j:X'\to X$ est
\emph{quasi-compacte} puisque $q'=q\circ j$ l'est et que $q$ est séparé
(\hyperref[I.6.6.4-fr]{I, 6.6.4, (v)}). Par suite
$\mathcal{F}=j_*(\mathcal{F}')$ est un $\mathcal{S}_X$-Module gradué
quasi-cohérent pour tout $\mathcal{S}_{X'}$-Module gradué quasi-cohérent
$\mathcal{F}'$, et il résulte des définitions précédentes que l'on a
\begin{equation}
\label{II.8.14.6.3-fr}
  \Gamma'_*(\mathcal{F}')=\Gamma_*(\mathcal{F}).
\tag{8.14.6.3}
\end{equation}
Avec les mêmes hypothèses sur $X'$, pour tout $\mathcal{S}$-Module gradué
quasi-cohérent $\mathcal{M}$, nous poserons
\begin{equation}
\label{II.8.14.6.4-fr}
  \operatorname{Proj}'(\mathcal{M})
  =\operatorname{Proj}(\mathcal{M})|X'
\tag{8.14.6.4}
\end{equation}
qui est un $\mathcal{S}_{X'}$-Module gradué quasi-cohérent. L'homomorphisme
canonique
\[
  \operatorname{Proj}(\mathcal{M})\to
  j_*(\operatorname{Proj}'(\mathcal{M}))
\]
(\hyperref[0.4.4.3-fr]{0, 4.4.3}) donne donc un homomorphisme canonique
$\Gamma_*(\operatorname{Proj}(\mathcal{M}))\to
\Gamma'_*(\operatorname{Proj}'(\mathcal{M}))$ de $\mathcal{S}$-Modules
gradués, et par suite, en le composant avec
(\hyperref[II.8.14.5.4-fr]{8.14.5.4}), on obtient un homomorphisme canonique
fonctoriel (de degré $0$) de $\mathcal{S}$-Modules gradués quasi-cohérents
\begin{equation}
\label{II.8.14.6.5-fr}
  \alpha':\mathcal{M}\to
  \Gamma'_*(\operatorname{Proj}'(\mathcal{M})).
\tag{8.14.6.5}
\end{equation}
\end{env}

\begin{env}[8.14.7]
\label{II.8.14.7-fr}
Conservons sur $X'$ les hypothèses de
(\hyperref[II.8.14.6-fr]{8.14.6}), et soit $\mathcal{F}'$ un
$\mathcal{S}_{X'}$-Module gradué quasi-cohérent, de sorte que
$\operatorname{Proj}'(\Gamma'_*(\mathcal{F}'))$ est aussi un
$\mathcal{S}_{X'}$-Module gradué quasi-cohérent. Nous allons définir un
homomorphisme canonique fonctoriel (de degré $0$) de
$\mathcal{S}_{X'}$-Modules gradués
\begin{equation}
\label{II.8.14.7.1-fr}
  \beta':\operatorname{Proj}'(\Gamma'_*(\mathcal{F}'))\to\mathcal{F}'.
\tag{8.14.7.1}
\end{equation}

\oldpage[II]{200}
Supposons d'abord $Y=\operatorname{Spec}(A)$ affine,
$\mathcal{S}=\widetilde{S}$, où $S$ est une $A$-algèbre graduée à degrés
positifs~; alors $\Gamma'_*(\mathcal{F}')=\widetilde{M}$, où
$M=\bigoplus_{n\in\mathbf{Z}}\Gamma(X',\mathcal{F}'_n)$ est un $S$-module
gradué. Soit $f\in S_d$ tel que $D_+(f)\subset X'$~; par définition
(\hyperref[II.2.6.2-fr]{2.6.2}), $\alpha_d(f)$ restreinte à $D_+(f)$ est la
section de $\mathcal{O}_X(d)$ au-dessus de $D_+(f)$ correspondant à
l'élément $f/1$ de $(S(d))_{(f)}$ et est par suite inversible~; il en est donc
de même de $\alpha_d(f^n)$ pour tout $n>0$. On en conclut aussitôt que
l'on définit un $S_f$-homomorphisme (de degré $0$) de modules gradués
$\beta_f:M_f\to\Gamma(D_+(f),\mathcal{F}')$ en faisant correspondre à tout
élément $z/f^n\in M_f$ (où $z\in M$), la section
$(z|D_+(f))(\alpha_d(f^n)|D_+(f))^{-1}$ de $\mathcal{F}'$ au-dessus de
$D_+(f)$. En outre, on a un diagramme commutatif correspondant à
(\hyperref[II.2.6.4.1-fr]{2.6.4.1}), d'où la définition de $\beta'$ dans ce
cas. Pour passer au cas général, il faut considérer une $A$-algèbre $A'$,
la $A'$-algèbre graduée $S'=S\otimes_A A'$, et utiliser un diagramme
commutatif analogue à (\hyperref[II.2.8.13.2-fr]{2.8.13.2})~; nous laissons
les détails aux lecteurs.
\end{env}

\begin{proposition}[8.14.8]
\label{II.8.14.8-fr}
Si $X'$ est un ouvert de $X=\operatorname{Proj}(\mathcal{S})$ tel que
$q':X'\to Y$ soit quasi-compact, l'homomorphisme $\beta'$ défini dans
(\hyperref[II.8.14.7-fr]{8.14.7}) est bijectif.
\end{proposition}

\begin{proof}
On peut évidemment se restreindre au cas où $Y$ est affine et tout revient à
prouver (avec les notations de (\hyperref[II.8.14.7-fr]{8.14.7})) que
l'homomorphisme $\beta_f:M_f\to\Gamma(D_+(f),\mathcal{F}')$ est un
isomorphisme. Or, lorsqu'on remplace $f$ par une de ses puissances, on ne
change pas $D_+(f)$ ni $\beta_f$~; comme $X'$ est \emph{quasi-compact} par
hypothèse, on peut donc toujours supposer, en vertu de
(\hyperref[II.8.14.4-fr]{8.14.4}), que le faisceau $\mathcal{O}_X(d)$ est
\emph{inversible}. Comme $X'$ est un schéma (puisque $q'$ est séparé), la
proposition n'est autre alors que (\hyperref[I.9.3.1-fr]{I, 9.3.1}).
\end{proof}

\begin{corollary}[8.14.9]
\label{II.8.14.9-fr}
Sous les hypothèses de (\hyperref[II.8.14.8-fr]{8.14.8}), tout
$\mathcal{S}_{X'}$-Module gradué quasi-cohérent $\mathcal{F}'$ est isomorphe
à un $\mathcal{S}_{X'}$-Module gradué de la forme
$\operatorname{Proj}'(\mathcal{M})$, où $\mathcal{M}$ est un
$\mathcal{S}$-Module gradué quasi-cohérent. Si en outre $\mathcal{F}'$ est de
type fini, et si on suppose que $Y$ est un schéma quasi-compact ou un
préschéma dont l'espace sous-jacent est noethérien, on peut supposer
$\mathcal{M}$ de type fini.
\end{corollary}

\begin{proof}
La démonstration à partir de (\hyperref[II.8.14.8-fr]{8.14.8}) suit
exactement la marche de celle de (\hyperref[II.3.4.5-fr]{3.4.5}) à partir de
(\hyperref[II.3.4.4-fr]{3.4.4}), et nous en laissons les détails au lecteur.
\end{proof}

\begin{proposition}[8.14.10]
\label{II.8.14.10-fr}
Sous les hypothèses de (\hyperref[II.8.14.7-fr]{8.14.7}), soient
$\mathcal{M}$ un $\mathcal{S}$-Module gradué quasi-cohérent,
$\mathcal{F}'$ un $\mathcal{S}_{X'}$-Module gradué quasi-cohérent~; les
homomorphismes composés
\begin{equation}
\label{II.8.14.10.1-fr}
  \operatorname{Proj}'(\mathcal{M})
  \xrightarrow{\operatorname{Proj}'(\alpha')}
  \operatorname{Proj}'(\Gamma'_*(\operatorname{Proj}'(\mathcal{M})))
  \xrightarrow{\beta'}\operatorname{Proj}'(\mathcal{M})
\tag{8.14.10.1}
\end{equation}
\begin{equation}
\label{II.8.14.10.2-fr}
  \Gamma'_*(\mathcal{F}')
  \xrightarrow{\alpha'}
  \Gamma'_*(\operatorname{Proj}'(\Gamma'_*(\mathcal{F}')))
  \xrightarrow{\Gamma'_*(\beta')}\Gamma'_*(\mathcal{F}')
\tag{8.14.10.2}
\end{equation}
sont les isomorphismes identiques.
\end{proposition}

\begin{proof}
La question est locale sur $Y$ et la vérification se fait comme dans
(\hyperref[II.2.6.5-fr]{2.6.5})~; nous en laissons les détails au lecteur.
\end{proof}

\begin{remark}[8.14.11]
\label{II.8.14.11-fr}
Au chapitre III (\hyperref[III.2.3.1-fr]{III, 2.3.1}), nous verrons que
lorsque $Y$ est \emph{localement noethérien} et $\mathcal{S}$ une
$\mathcal{O}_Y$-Algèbre graduée quasi-cohérente \emph{de type fini} (auquel
cas

\oldpage[II]{201}
on peut prendre $X'=X$), alors l'homomorphisme $\alpha$
(\hyperref[II.8.14.5.4-fr]{8.14.5.4}) est \textbf{(TN)}-\emph{bijectif} pour
tout $\mathcal{S}$-Module gradué quasi-cohérent $\mathcal{M}$ vérifiant la
condition \textbf{(TF)}.
\end{remark}

\begin{remark}[8.14.12]
\label{II.8.14.12-fr}
La situation décrite dans (\hyperref[II.8.14.4-fr]{8.14.4}) est un cas
particulier de la suivante. Soient $X$ un espace annelé, $\mathcal{S}$ une
$\mathcal{O}_X$-Algèbre graduée (à degrés positifs et négatifs)~; supposons
qu'il existe un entier $d>0$ tel que $\mathcal{S}_d$ et
$\mathcal{S}_{-d}$ soient \emph{inversibles}, l'homomorphisme canonique
\begin{equation}
\label{II.8.14.12.1-fr}
  \mathcal{S}_d\otimes_{\mathcal{O}_X}\mathcal{S}_{-d}
  \to\mathcal{O}_X
\tag{8.14.12.1}
\end{equation}
étant un \emph{isomorphisme} (de sorte que $\mathcal{S}_{-d}$ s'identifie à
$\mathcal{S}_d^{-1}$). On dit alors que la $\mathcal{O}_X$-Algèbre graduée
$\mathcal{S}$ est \emph{périodique de période $d$}. Cette dénomination
provient de la propriété suivante~: \emph{sous les hypothèses précédentes,
pour tout $\mathcal{S}$-Module gradué $\mathcal{F}$, l'homomorphisme
canonique}
\begin{equation}
\label{II.8.14.12.2-fr}
  \mathcal{S}_d\otimes\mathcal{F}_n\to\mathcal{F}_{n+d}
\tag{8.14.12.2}
\end{equation}
\emph{est un isomorphisme pour tout $n\in\mathbf{Z}$.} En effet, la question
est locale sur $X$ et on peut donc supposer que $\mathcal{S}_d$ possède une
section \emph{inversible} $s$ au-dessus de $X$, son inverse $s'$ étant une
section de $\mathcal{S}_{-d}$. L'homomorphisme
$\mathcal{F}_{n+d}\to\mathcal{S}_d\otimes\mathcal{F}_n$ qui, à toute section
$z\in\Gamma(U,\mathcal{F}_{n+d})$ fait correspondre la section
$(s|U)\otimes(s'|U)z$ de $\mathcal{S}_d\otimes\mathcal{F}_n$ sur $U$, est
alors réciproque de (\hyperref[II.8.14.12.2-fr]{8.14.12.2}), d'où notre
assertion. On en déduit pour tout $k\in\mathbf{Z}$ un isomorphisme canonique
\[
  (\mathcal{S}_d)^{\otimes k}\otimes\mathcal{F}_n
  \xrightarrow{\sim}\mathcal{F}_{n+kd}.
\]
Par suite, \emph{un $\mathcal{S}$-Module gradué $\mathcal{F}$ est connu
lorsqu'on connaît les $\mathcal{S}_0$-Modules $\mathcal{F}_i$
$(0\leqslant i\leqslant d-1)$ et les homomorphismes canoniques}
\[
  \mathcal{S}_i\otimes\mathcal{F}_j\to\mathcal{F}_{i+j}
  \qquad\text{pour }0\leqslant i,j\leqslant d-1
\]
(en posant
$\mathcal{F}_{i+j}=\mathcal{S}_d\otimes_{\mathcal{S}_0}
\mathcal{F}_{i+j-d}$ lorsque $i+j\geqslant d$). Bien entendu, pour que ces
homomorphismes définissent bien une structure de $\mathcal{S}$-Module sur la
somme directe des $(\mathcal{S}_d)^{\otimes k}\otimes\mathcal{F}_i$
$(k\in\mathbf{Z},\ 0\leqslant i\leqslant d-1)$, ils doivent satisfaire à des
conditions d'associativité que nous n'expliciterons pas.

Dans le cas où $d=1$ (qui est celui considéré dans
(\hyperref[subsection:II.3.3-fr]{3.3})) on peut donc dire que la catégorie des
$\mathcal{S}$-Modules gradués (resp. quasi-cohérents si $X$ est un préschéma
et $\mathcal{S}$ est quasi-cohérente) est \emph{équivalente} à la catégorie
des $\mathcal{S}_0$-Modules quelconques (resp. quasi-cohérents)~; c'est dans
ce sens qu'on peut considérer les développements de ce numéro comme
généralisant ceux du \S~\hyperref[section:II.3-fr]{3}. En outre, on voit que,
sous des conditions de finitude convenables, les résultats de ce numéro
(joints à (\hyperref[II.8.14.11-fr]{8.14.11})) ramènent dans une certaine
mesure l'étude des Algèbres graduées quasi-cohérentes sur un préschéma, et des
Modules gradués «~mod. \textbf{(TN)}~» sur de telles Algèbres, à l'étude du
cas particulier où les Algèbres considérées sont \emph{périodiques} (et où
la condition \textbf{(TN)} pour $\mathcal{M}$
(\hyperref[II.3.4.2-fr]{3.4.2}) implique donc $\mathcal{M}=0$).
\end{remark}

\begin{remark}[8.14.13]
\label{II.8.14.13-fr}
Sous les hypothèses de (\hyperref[II.8.14.1-fr]{8.14.1}), soit $d$ un entier
$>0$~; on a défini un $Y$-isomorphisme canonique $h$ de $X$ sur
$X^{(d)}=\operatorname{Proj}(\mathcal{S}^{(d)})$
(\hyperref[II.3.1.8-fr]{3.1.8}). Pour tout

\oldpage[II]{202}
$\mathcal{S}$-Module gradué quasi-cohérent $\mathcal{M}$ et tout entier $k$
tel que $0\leqslant k\leqslant d-1$, on a aussi un $h$-isomorphisme canonique
(avec les notations de (\hyperref[II.3.1.1-fr]{3.1.1}))
\begin{equation}
\label{II.8.14.13.1-fr}
  (\operatorname{Proj}(\mathcal{M}))^{(d,k)}
  \xleftarrow{\sim}\operatorname{Proj}(\mathcal{M}^{(d,k)}).
\tag{8.14.13.1}
\end{equation}

Supposons d'abord $Y=\operatorname{Spec}(A)$ affine,
$\mathcal{S}=\widetilde{S}$, $\mathcal{M}=\widetilde{M}$, où $S$ est une
$A$-algèbre graduée à degrés positifs et $M$ un $S$-module gradué. On sait
que pour tout $f\in S_e$ ($e>0$) $h$ applique $D_+(f)$ sur $D_+(f^d)$ et
correspond à l'isomorphisme canonique
$S_{(f^d)}\xrightarrow{\sim}S_{(f)}$
(\hyperref[II.2.2.2-fr]{2.2.2}). La restriction de
(\hyperref[II.8.14.13.1-fr]{8.14.13.1}) à $D_+(f^d)$ correspond alors au
di-isomorphisme canonique $M_{f^d}\xrightarrow{\sim}M_f$, restreint aux
éléments de $M_{f^d}$ de degré congru à $k$ (mod. $d$). Nous laissons au
lecteur le soin de vérifier que ces isomorphismes sont compatibles avec le
passage de $f$ à un de ses multiples homogènes $fg$, puis qu'on a une
compatibilité analogue avec le passage de $S$ à une $A'$-algèbre graduée
$S'=S\otimes_A A'$, lorsque $A'$ est une $A$-algèbre.

En particulier, on a ainsi un $h$-isomorphisme
\begin{equation}
\label{II.8.14.13.2-fr}
  (\mathcal{S}^{(d)})_{X^{(d)}}
  \xrightarrow{\sim}(\mathcal{S}_X)^{(d)}
\tag{8.14.13.2}
\end{equation}
qui respecte les structures multiplicatives des deux membres, et grâce auquel
(\hyperref[II.8.14.13.1-fr]{8.14.13.1}) devient un $h$-di-isomorphisme d'un
$(\mathcal{S}^{(d)})_{X^{(d)}}$-Module gradué sur un
$(\mathcal{S}_X)^{(d)}$-Module gradué. De même, on a un $h$-isomorphisme
\begin{equation}
\label{II.8.14.13.3-fr}
  \operatorname{Proj}_0(\mathcal{S}^{(d,k)}(n))
  \xrightarrow{\sim}\mathcal{O}_X(nd+k)
\tag{8.14.13.3}
\end{equation}
qui complète le résultat de (\hyperref[II.3.2.9-fr]{3.2.9, (ii)}).

On déduit immédiatement de l'isomorphisme
(\hyperref[II.8.14.13.1-fr]{8.14.13.1}) un isomorphisme de
$\mathcal{S}^{(d)}$-Modules gradués
\begin{equation}
\label{II.8.14.13.4-fr}
  \Gamma_*^{(d)}(\operatorname{Proj}(\mathcal{M}^{(d,k)}))
  \xrightarrow{\sim}
  \Gamma_*((\operatorname{Proj}(\mathcal{M}))^{(d,k)})
\tag{8.14.13.4}
\end{equation}
où $\Gamma_*^{(d)}$ correspond au morphisme structural
$q^{(d)}:X^{(d)}\to Y$~; il est immédiat de vérifier que l'homomorphisme
canonique $\alpha$ (\hyperref[II.8.14.5.4-fr]{8.14.5.4}) et l'homomorphisme
analogue $\alpha^{(d)}$ pour $X^{(d)}$ rendent commutatif le diagramme
\begin{equation}
\label{II.8.14.13.5-fr}
  \xymatrix{
    & \mathcal{M}^{(d,k)}\ar[dl]_{\alpha^{(d)}}\ar[dr]^{\alpha} & \\
    \Gamma_*^{(d)}(\operatorname{Proj}(\mathcal{M}^{(d,k)}))
      \ar[rr]^{\sim}
      && \Gamma_*((\operatorname{Proj}(\mathcal{M}))^{(d,k)})
  }
\tag{8.14.13.5}
\end{equation}
la vérification se faisant en supposant $Y$ affine, et calculant les
restrictions des images par $\alpha^{(d)}$ et $\alpha$ d'un même élément de
$M^{(d,k)}$ aux ouverts $D_+(f^d)$ et $D_+(f)$ (avec les mêmes notations que
ci-dessus).
\end{remark}

\begin{proposition}[8.14.14]
\label{II.8.14.14-fr}
Soient $Y$ un préschéma quasi-compact, $\mathcal{S}$ une
$\mathcal{O}_Y$-Algèbre graduée quasi-cohérente de type fini, $\mathcal{M}$
un $\mathcal{S}$-Module gradué quasi-cohérent vérifiant la condition
\textbf{(TF)}~; soit $X=\operatorname{Proj}(\mathcal{S})$. Alors
$\mathcal{S}_X$ est une $\mathcal{O}_X$-Algèbre graduée périodique
(\hyperref[II.8.14.12-fr]{8.14.12}), et il existe une

\oldpage[II]{203}
période $d$ de $\mathcal{S}_X$ telle que les
$(\operatorname{Proj}(\mathcal{M}))^{(d,k)}$
$(0\leqslant k\leqslant d-1)$ soient des
$(\mathcal{S}_X)^{(d)}$-Modules de type fini.
\end{proposition}

\begin{proof}
En effet, (\hyperref[II.3.1.10-fr]{3.1.10}) prouve qu'il existe $d$ tel que
$\mathcal{S}^{(d)}$ soit engendrée par
$\mathcal{S}_d=(\mathcal{S}^{(d)})_1$, ce dernier étant un
$\mathcal{S}_0$-Module de type fini. Pour démontrer la première assertion, on
peut donc, en vertu de (\hyperref[II.8.14.13.2-fr]{8.14.13.2}), se borner au
cas où $d=1$, et alors la proposition résulte de
(\hyperref[II.3.2.7-fr]{3.2.7}). En outre, compte tenu de
(\hyperref[II.8.14.13.1-fr]{8.14.13.1}), la seconde assertion est conséquence
de (\hyperref[II.2.1.6-fr]{2.1.6, (iii)}) et de
(\hyperref[II.3.4.3-fr]{3.4.3}).
\end{proof}

\clearpage
\thispagestyle{plain}
\markboth{PRÉLIMINAIRES}{PRÉLIMINAIRES}
\oldpage[0\textsubscript{III}]{5}
\phantomsection
\label{chapter:0III-fr}

\begin{center}
{\large CHAPITRE 0 (suite)\footnotemark\par}
\vspace{1.5em}
{\Large\bfseries PRÉLIMINAIRES\par}
\end{center}

\footnotetext{Pour faciliter la recherche des références, on renverra
désormais aux paragraphes du chapitre 0 publiés avec le chapitre I par le
signe $0_{\mathrm I}$.}

\setcounter{section}{7}
\section{Foncteurs représentables}
\label{section:0.8-fr}

\setcounter{subsection}{0}
\subsection{Foncteurs représentables}
\label{subsection:0.8.1-fr}

\begin{env}[8.1.1]
\label{0.8.1.1-fr}
Nous désignerons par $\mathrm{Ens}$ la catégorie des ensembles. Soit $C$ une
catégorie~; pour deux objets $X$, $Y$ de $C$, nous poserons
$h_X(Y)=\operatorname{Hom}(Y,X)$~; pour tout morphisme $u:Y\to Y'$ dans $C$,
nous désignerons par $h_X(u)$ l'application $v\mapsto vu$ de
$\operatorname{Hom}(Y',X)$ dans $\operatorname{Hom}(Y,X)$. Il est immédiat
qu'avec ces définitions, $h_X:C\to\mathrm{Ens}$ est un \emph{foncteur
contravariant}, c'est-à-dire un objet de la catégorie, notée
$\operatorname{Hom}(C^0,\mathrm{Ens})$, des foncteurs covariants de la
catégorie $C^0$, duale de la catégorie $C$, dans la catégorie
$\mathrm{Ens}$ (T, 1.7, d) et [29]).
\end{env}

\begin{env}[8.1.2]
\label{0.8.1.2-fr}
Soit maintenant $w:X\to X'$ un morphisme dans $C$~; pour tout $Y\in C$ et
tout $v\in\operatorname{Hom}(Y,X)=h_X(Y)$, on a
$wv\in\operatorname{Hom}(Y,X')=h_{X'}(Y)$~; désignons par $h_w(Y)$
l'application $v\mapsto wv$ de $h_X(Y)$ dans $h_{X'}(Y)$. Il est immédiat
que pour tout morphisme $u:Y\to Y'$ dans $C$, le diagramme
\[
  \xymatrix{
    h_X(Y')\ar[r]^{h_X(u)}\ar[d]_{h_w(Y')} &
    h_X(Y)\ar[d]^{h_w(Y)}\\
    h_{X'}(Y')\ar[r]_{h_{X'}(u)} &
    h_{X'}(Y)
  }
\]
est commutatif~; autrement dit, $h_w$ est un \emph{morphisme fonctoriel}
$h_X\to h_{X'}$ (T, 1.2), ou encore un morphisme dans la catégorie
$\operatorname{Hom}(C^0,\mathrm{Ens})$ (T, 1.7, d)). Les définitions de
$h_X$ et de $h_w$ constituent donc la définition d'un \emph{foncteur
covariant canonique}
\[
  h:C\to\operatorname{Hom}(C^0,\mathrm{Ens}),\qquad X\mapsto h_X.
  \tag{8.1.2.1}
\]
\end{env}

\begin{env}[8.1.3]
\label{0.8.1.3-fr}
Soient $X$ un objet de $C$, $F$ un foncteur contravariant de $C$ dans
$\mathrm{Ens}$ (objet de $\operatorname{Hom}(C^0,\mathrm{Ens})$). Soit
$g:h_X\to F$ un \emph{morphisme fonctoriel}~: pour tout $Y\in C$,
\oldpage[0\textsubscript{III}]{6}
$g(Y)$ est donc une application $h_X(Y)\to F(Y)$ telle que pour tout morphisme
$u:Y\to Y'$ dans $C$, le diagramme
\[
  \xymatrix{
    h_X(Y')\ar[r]^{h_X(u)}\ar[d]_{g(Y')} &
    h_X(Y)\ar[d]^{g(Y)}\\
    F(Y')\ar[r]_{F(u)} & F(Y)
  }
  \tag{8.1.3.1}\label{0.8.1.3.1-fr}
\]
soit commutatif. En particulier, on a une application
$g(X):h_X(X)=\operatorname{Hom}(X,X)\to F(X)$, d'où un élément
\[
  \alpha(g)=(g(X))(1_X)\in F(X)
  \tag{8.1.3.2}
\]
et par suite une application canonique
\[
  \alpha:\operatorname{Hom}(h_X,F)\to F(X).
  \tag{8.1.3.3}
\]

Inversement, considérons un élément $\xi\in F(X)$~; pour tout morphisme
$v:Y\to X$ dans $C$, $F(v)$ est une application $F(X)\to F(Y)$~;
considérons l'application
\[
  v\mapsto(F(v))(\xi)
  \tag{8.1.3.4}
\]
de $h_X(Y)$ dans $F(Y)$~; si on désigne par $(\beta(\xi))(Y)$ cette
application,
\[
  \beta(\xi):h_X\to F
  \tag{8.1.3.5}
\]
est un \emph{morphisme fonctoriel}, car on a pour tout morphisme
$u:Y\to Y'$ dans $C$,
$(F(vu))(\xi)=(F(v)\circ F(u))(\xi)$, ce qui vérifie la commutativité de
(\hyperref[0.8.1.3.1-fr]{8.1.3.1}) pour $g=\beta(\xi)$. On a ainsi défini
une application canonique
\[
  \beta:F(X)\to\operatorname{Hom}(h_X,F).
  \tag{8.1.3.6}
\]
\end{env}

\begin{proposition}[8.1.4]
\label{0.8.1.4-fr}
Les applications $\alpha$ et $\beta$ sont des bijections réciproques l'une de
l'autre.
\end{proposition}

Calculons $\alpha(\beta(\xi))$ pour $\xi\in F(X)$~; pour tout $Y\in C$,
$(\beta(\xi))(Y)$ est l'application
$g_1(Y):v\mapsto(F(v))(\xi)$ de $h_X(Y)$ dans $F(Y)$. On a donc
\[
  \alpha(\beta(\xi))=(g_1(X))(1_X)=(F(1_X))(\xi)
  =1_{F(X)}(\xi)=\xi.
\]
Calculons maintenant $\beta(\alpha(g))$ pour
$g\in\operatorname{Hom}(h_X,F)$~; pour tout $Y\in C$,
$(\beta(\alpha(g)))(Y)$ est l'application
$v\mapsto(F(v))((g(X))(1_X))$~; en vertu de la commutativité de
(\hyperref[0.8.1.3.1-fr]{8.1.3.1}), cette application n'est autre que
$v\mapsto(g(Y))((h_X(v))(1_X))=(g(Y))(v)$ par définition de $h_X(v)$,
autrement dit, elle est égale à $g(Y)$, ce qui démontre la proposition.

\begin{env}[8.1.5]
\label{0.8.1.5-fr}
Rappelons qu'une \emph{sous-catégorie} $C'$ d'une catégorie $C$ est définie
par la condition que ses objets soient des objets de $C$, et que si $X'$,
$Y'$ sont deux objets de $C'$, l'ensemble
$\operatorname{Hom}_{C'}(X',Y')$ des morphismes $X'\to Y'$ \emph{dans
$C'$} est une partie de l'ensemble $\operatorname{Hom}_C(X',Y')$ des
morphismes $X'\to Y'$ \emph{dans $C$}, l'application canonique de
<<~composition des morphismes~>>
\[
  \operatorname{Hom}_{C'}(X',Y')\times\operatorname{Hom}_{C'}(Y',Z')
  \to\operatorname{Hom}_{C'}(X',Z')
\]
\oldpage[0\textsubscript{III}]{7}
étant la restriction de l'application canonique
\[
  \operatorname{Hom}_C(X',Y')\times\operatorname{Hom}_C(Y',Z')
  \to\operatorname{Hom}_C(X',Z').
\]

On dit que $C'$ est une \emph{sous-catégorie pleine} de $C$ si
$\operatorname{Hom}_{C'}(X',Y')=\operatorname{Hom}_C(X',Y')$ pour tout
couple d'objets de $C'$. La sous-catégorie $C''$ de $C$ formée des objets de
$C$ isomorphes aux objets de $C'$ est encore alors une sous-catégorie pleine
de $C$, \emph{équivalente} (T, 1.2) à $C'$ comme on le vérifie sans peine.

Un foncteur covariant $F:C_1\to C_2$ est dit \emph{pleinement fidèle} si,
pour tout couple d'objets $X_1$, $Y_1$ de $C_1$, l'application
$u\mapsto F(u)$ de $\operatorname{Hom}(X_1,Y_1)$ dans
$\operatorname{Hom}(F(X_1),F(Y_1))$ est \emph{bijective}~; cela entraîne que
la sous-catégorie $F(C_1)$ de $C_2$ est \emph{pleine}. En outre, si deux
objets $X_1$, $X'_1$ ont même image $X_2$, il existe un isomorphisme unique
$u:X_1\to X'_1$ tel que $F(u)=1_{X_2}$. Pour tout objet $X_2$ de $F(C_1)$,
soit alors $G(X_2)$ un des objets $X_1$ de $C_1$ tel que $F(X_1)=X_2$
($G$ étant défini au moyen de l'axiome de choix)~; pour tout morphisme
$v:X_2\to Y_2$ dans $F(C_1)$, $G(v)$ sera l'unique morphisme
$u:G(X_2)\to G(Y_2)$ tel que $F(u)=v$~; $G$ est alors un \emph{foncteur} de
$F(C_1)$ dans $C_1$~; $FG$ est le foncteur identique dans $F(C_1)$, et ce
qui précède montre qu'il existe un isomorphisme de foncteurs
$\varphi:1_{C_1}\to GF$ tel que $F$, $G$, $\varphi$ et l'identité
$1_{F(C_1)}\to FG$ définissent une \emph{équivalence} de la catégorie $C_1$
et de la sous-catégorie pleine $F(C_1)$ de $C_2$ (T, 1.2).
\end{env}

\begin{env}[8.1.6]
\label{0.8.1.6-fr}
Appliquons la prop. (\hyperref[0.8.1.4-fr]{8.1.4}) au cas où le foncteur $F$
est $h_{X'}$, $X'$ étant un objet quelconque de $C$~; l'application
$\beta:\operatorname{Hom}(X,X')\to\operatorname{Hom}(h_X,h_{X'})$ n'est
autre ici que l'application $w\mapsto h_w$ définie dans
(\hyperref[0.8.1.2-fr]{8.1.2})~; cette application étant \emph{bijective},
on voit, avec la terminologie de (\hyperref[0.8.1.5-fr]{8.1.5}), que~:
\end{env}

\begin{proposition}[8.1.7]
\label{0.8.1.7-fr}
Le foncteur canonique
$h:C\to\operatorname{Hom}(C^0,\mathrm{Ens})$ est pleinement fidèle.
\end{proposition}

\begin{env}[8.1.8]
\label{0.8.1.8-fr}
Soit $F$ un foncteur contravariant de $C$ dans $\mathrm{Ens}$~; on dit que
$F$ est \emph{représentable} s'il existe un objet $X\in C$ tel que $F$ soit
\emph{isomorphe} à $h_X$~; il résulte de
(\hyperref[0.8.1.7-fr]{8.1.7}) que la donnée d'un $X\in C$ et d'un
isomorphisme de foncteurs $g:h_X\to F$ détermine $X$ à un isomorphisme
unique près. La prop. (\hyperref[0.8.1.7-fr]{8.1.7}) signifie encore que $h$
définit une \emph{équivalence} de $C$ et de la sous-catégorie pleine de
$\operatorname{Hom}(C^0,\mathrm{Ens})$ formée des \emph{foncteurs
contravariants représentables}. Il résulte d'ailleurs de
(\hyperref[0.8.1.4-fr]{8.1.4}) que la donnée d'un morphisme fonctoriel
$g:h_X\to F$ équivaut à celle d'un élément $\xi\in F(X)$~; dire que $g$ est
un \emph{isomorphisme} équivaut pour $\xi$ à la condition suivante~:
\emph{pour tout objet $Y$ de $C$ l'application
$v\mapsto(F(v))(\xi)$ de $\operatorname{Hom}(Y,X)$ dans $F(Y)$ est
bijective}. Lorsque $\xi$ vérifie cette condition, on dira que le couple
$(X,\xi)$ \emph{représente} le foncteur représentable $F$. Par abus de
langage, on dira aussi que l'objet $X\in C$ représente $F$ s'il existe
$\xi\in F(X)$ tel que $(X,\xi)$ représente $F$, autrement dit si $h_X$ est
isomorphe à $F$.

Soient $F$, $F'$ deux foncteurs contravariants représentables de $C$ dans
$\mathrm{Ens}$, $h_X\to F$ et $h_{X'}\to F'$ deux isomorphismes de
foncteurs. Alors il résulte de (\hyperref[0.8.1.6-fr]{8.1.6}) qu'il y a une
correspondance biunivoque canonique entre $\operatorname{Hom}(X,X')$ et
l'ensemble $\operatorname{Hom}(F,F')$ des morphismes fonctoriels $F\to F'$.
\end{env}

\begin{env}[8.1.9]
\label{0.8.1.9-fr}
\textit{Exemples. I : limites projectives.} La notion de foncteur
contravariant représentable couvre en particulier la notion <<~duale~>> de
la notion usuelle de <<~solution d'un
\oldpage[0\textsubscript{III}]{8}
problème universel~>>. Plus généralement, nous allons voir que la notion de
\emph{limite projective} est un cas particulier de celle de foncteur
représentable. Rappelons que dans une catégorie $C$, on définit un
\emph{système projectif} par la donnée d'un ensemble préordonné $I$, d'une
famille $(A_\alpha)_{\alpha\in I}$ d'objets de $C$, et, pour tout couple
d'indices $(\alpha,\beta)$ tel que $\alpha\leq\beta$, d'un morphisme
$u_{\alpha\beta}:A_\beta\to A_\alpha$. Une \emph{limite projective} de ce
système dans $C$ est constituée par un objet $B$ de $C$ (noté
$\varprojlim A_\alpha$), et, pour chaque $\alpha\in I$, un morphisme
$u_\alpha:B\to A_\alpha$, tels que~: 1\textsuperscript{o}
$u_\alpha=u_{\alpha\beta}u_\beta$ pour $\alpha\leq\beta$~;
2\textsuperscript{o} Pour tout objet $X$ de $C$ et toute famille
$(v_\alpha)_{\alpha\in I}$ de morphismes $v_\alpha:X\to A_\alpha$, telle
que $v_\alpha=u_{\alpha\beta}v_\beta$ pour $\alpha\leq\beta$, il existe un
morphisme unique $v:X\to B$ (noté $\varprojlim v_\alpha$) tel que
$v_\alpha=u_\alpha v$ pour tout $\alpha\in I$ (T, 1.8). Ceci s'interprète de
la façon suivante~: les $u_{\alpha\beta}$ définissent canoniquement des
applications
\[
  \overline{u}_{\alpha\beta}:\operatorname{Hom}(X,A_\beta)
  \to\operatorname{Hom}(X,A_\alpha)
\]
qui définissent un \emph{système projectif} d'ensembles
$(\operatorname{Hom}(X,A_\alpha),\overline{u}_{\alpha\beta})$, et
$(v_\alpha)$ est par définition un élément de l'ensemble
$\varprojlim_\alpha\operatorname{Hom}(X,A_\alpha)$~; il est clair que
$X\mapsto\varprojlim_\alpha\operatorname{Hom}(X,A_\alpha)$ est un
\emph{foncteur contravariant} de $C$ dans $\mathrm{Ens}$, et l'existence de
la limite projective $B$ équivaut à dire que
$(v_\alpha)\mapsto\varprojlim v_\alpha$ est un \emph{isomorphisme} de
foncteurs en $X$
\[
  \varprojlim_\alpha\operatorname{Hom}(X,A_\alpha)
  \xrightarrow{\sim}\operatorname{Hom}(X,B)
  \tag{8.1.9.1}
\]
autrement dit que le foncteur
$X\mapsto\varprojlim_\alpha\operatorname{Hom}(X,A_\alpha)$ est
\emph{représentable}.
\end{env}

\begin{env}[8.1.10]
\label{0.8.1.10-fr}
\textit{Exemples. II : Objet final.} Soient $C$ une catégorie, $\{a\}$ un
ensemble réduit à un seul élément. Considérons le foncteur contravariant
$F:C\to\mathrm{Ens}$ qui, à tout objet $X$ de $C$ fait correspondre
l'ensemble $\{a\}$ et à tout morphisme $X\to X'$ dans $C$ l'unique
application $\{a\}\to\{a\}$. Dire que ce foncteur est \emph{représentable}
signifie qu'il existe un objet $e\in C$ tel que pour tout $Y\in C$,
$\operatorname{Hom}(Y,e)=h_e(Y)$ soit \emph{réduit à un élément}~; on dit
que $e$ est un \emph{objet final} de $C$, et il est clair que deux objets
finaux de $C$ sont isomorphes (ce qui permet de définir, en général à l'aide
de l'axiome de choix, un objet final de $C$ qu'on note alors $e_C$). Par
exemple, dans la catégorie $\mathrm{Ens}$, les objets finaux sont les
ensembles réduits à un élément~; dans la catégorie des \emph{algèbres
augmentées} sur un corps $K$ (où les morphismes sont les homomorphismes
d'algèbres compatibles avec les augmentations), $K$ est un objet final~;
dans la catégorie des \emph{$S$-préschémas} (I, 2.5.1), $S$ est un objet
final.
\end{env}

\begin{env}[8.1.11]
\label{0.8.1.11-fr}
Pour deux objets $X$, $Y$ d'une catégorie $C$, posons
$h'_X(Y)=\operatorname{Hom}(X,Y)$ et pour tout morphisme $u:Y\to Y'$, soit
$h'_X(u)$ l'application $v\mapsto uv$ de $\operatorname{Hom}(X,Y)$ dans
$\operatorname{Hom}(X,Y')$~; $h'_X$ est alors un \emph{foncteur covariant}
$C\to\mathrm{Ens}$, d'où l'on déduit comme dans
(\hyperref[0.8.1.2-fr]{8.1.2}) la définition d'un foncteur covariant
canonique $h':C^0\to\operatorname{Hom}(C,\mathrm{Ens})$~; un \emph{foncteur
covariant} $F$ de $C$ dans $\mathrm{Ens}$, autrement dit un objet de
$\operatorname{Hom}(C,\mathrm{Ens})$ est alors dit \emph{représentable}
s'il existe un objet $X\in C$ (nécessairement unique à isomorphisme unique
près) tel que $F$ soit \emph{isomorphe à $h'_X$}~; nous laissons au lecteur
le soin de développer les considérations <<~duales~>> des précédentes pour
cette notion, qui couvre cette fois celle de \emph{limite inductive}, et en
particulier la notion usuelle de <<~solution de problème universel~>>.
\end{env}

\subsection{Structures algébriques dans les catégories}
\label{subsection:0.8.2-fr}

\begin{env}[8.2.1]
\label{0.8.2.1-fr}
\oldpage[0\textsubscript{III}]{9}
Étant donnés deux foncteurs contravariants $F$, $F'$ de $C$ dans
$\mathrm{Ens}$, rappelons que pour tout objet $Y\in C$, on pose
$(F\times F')(Y)=F(Y)\times F'(Y)$, et pour tout morphisme $u:Y\to Y'$ dans
$C$, on pose $(F\times F')(u)=F(u)\times F'(u)$ qui est l'application
$(t,t')\mapsto((F(u))(t),(F'(u))(t'))$ de $F(Y')\times F'(Y')$ dans
$F(Y)\times F'(Y)$~; $F\times F':C\to\mathrm{Ens}$ est donc un
\emph{foncteur contravariant} (qui n'est autre d'ailleurs que le
\emph{produit} des objets $F$, $F'$ dans la catégorie
$\operatorname{Hom}(C^0,\mathrm{Ens})$). Étant donné un objet $X\in C$, nous
appellerons \emph{loi de composition interne sur $X$} un morphisme fonctoriel
\[
  \gamma_X:h_X\times h_X\to h_X.
  \tag{8.2.1.1}
\]

Autrement dit (T, 1.2), pour tout objet $Y\in C$, $\gamma_X(Y)$ est une
application $h_X(Y)\times h_X(Y)\to h_X(Y)$ (donc par définition une
\emph{loi de composition interne} sur l'ensemble $h_X(Y)$) soumise à la
condition que, pour tout morphisme $u:Y\to Y'$ dans $C$, le diagramme
\[
  \xymatrix{
    h_X(Y')\times h_X(Y')
      \ar[rr]^{h_X(u)\times h_X(u)}\ar[dd]_{\gamma_X(Y')} &&
    h_X(Y)\times h_X(Y)\ar[dd]^{\gamma_X(Y)}\\\\
    h_X(Y')\ar[rr]_{h_X(u)} && h_X(Y)
  }
\]
soit commutatif~; cela signifie que pour les lois de composition
$\gamma_X(Y)$ et $\gamma_X(Y')$, $h_X(u)$ est un \emph{homomorphisme} de
$h_X(Y')$ dans $h_X(Y)$.

De la même façon, étant donnés deux objets $Z$, $X$ de $C$, on appelle
\emph{loi de composition externe sur $X$, ayant $Z$ comme domaine
d'opérateurs}, un morphisme fonctoriel
\[
  \omega_{X,Z}:h_Z\times h_X\to h_X.
  \tag{8.2.1.2}
\]

On voit comme ci-dessus que pour tout $Y\in C$, $\omega_{X,Z}(Y)$ est une loi
de composition externe sur $h_X(Y)$, ayant $h_Z(Y)$ comme domaine
d'opérateurs, et telle que pour tout morphisme $u:Y\to Y'$, $h_X(u)$ et
$h_Z(u)$ forment un \emph{dihomomorphisme} de
$(h_Z(Y'),h_X(Y'))$ dans $(h_Z(Y),h_X(Y))$.
\end{env}

\begin{env}[8.2.2]
\label{0.8.2.2-fr}
Soit $X'$ un second objet de $C$, et supposons donnée sur $X'$ une loi de
composition interne $\gamma_{X'}$~; nous dirons qu'un morphisme
$w:X\to X'$ dans $C$ est un \emph{homomorphisme} pour ces lois de
composition, si pour tout $Y\in C$, $h_w(Y):h_X(Y)\to h_{X'}(Y)$ est un
homomorphisme pour les lois de composition $\gamma_X(Y)$ et
$\gamma_{X'}(Y)$. Si $X''$ est un troisième
\oldpage[0\textsubscript{III}]{10}
objet de $C$ muni d'une loi de composition interne $\gamma_{X''}$, et
$w':X'\to X''$ un morphisme dans $C$ qui est un homomorphisme pour
$\gamma_{X'}$ et $\gamma_{X''}$, il est clair que le morphisme
$w'w:X\to X''$ est un homomorphisme pour les lois de composition $\gamma_X$
et $\gamma_{X''}$. Un isomorphisme $w:X\xrightarrow{\sim}X'$ dans $C$ est
appelé \emph{isomorphisme pour les lois de composition $\gamma_X$ et
$\gamma_{X'}$} si $w$ est un homomorphisme pour ces lois de composition, et
si son morphisme réciproque $w^{-1}$ est un homomorphisme pour les lois de
composition $\gamma_{X'}$ et $\gamma_X$.

On définit de la même manière les \emph{dihomomorphismes} pour les couples
d'objets de $C$ munis de lois de composition externes.
\end{env}

\begin{env}[8.2.3]
\label{0.8.2.3-fr}
Lorsqu'une loi de composition interne $\gamma_X$ sur un objet $X\in C$ est
telle que $\gamma_X(Y)$ soit une \emph{loi de groupe} sur $h_X(Y)$ pour
\emph{tout} $Y\in C$, on dit que $X$, muni de cette loi, est un
\emph{$C$-groupe} ou un \emph{$C$-objet en groupes}. On définit de même les
\emph{$C$-anneaux}, \emph{$C$-modules}, etc.
\end{env}

\begin{env}[8.2.4]
\label{0.8.2.4-fr}
Supposons que le \emph{produit} $X\times X$ d'un objet $X\in C$ par lui-même
existe dans $C$~; par définition, on a alors
$h_{X\times X}=h_X\times h_X$ à un isomorphisme canonique près, puisqu'il
s'agit d'un cas particulier de limite projective
(\hyperref[0.8.1.9-fr]{8.1.9})~; une loi de composition interne sur $X$ peut
donc être considérée comme un morphisme fonctoriel
$\gamma_X:h_{X\times X}\to h_X$, et détermine donc canoniquement
(\hyperref[0.8.1.6-fr]{8.1.6}) un élément
$c_X\in\operatorname{Hom}(X\times X,X)$ tel que $h_{c_X}=\gamma_X$~; dans
ce cas, la donnée d'une loi de composition interne sur $X$ est donc
équivalente à celle d'un morphisme $X\times X\to X$~; lorsque $C$ est la
catégorie $\mathrm{Ens}$, on retrouve la notion classique de loi de
composition interne sur un ensemble. On a un résultat analogue pour une loi
de composition externe lorsque le produit $Z\times X$ existe dans $C$.
\end{env}

\begin{env}[8.2.5]
\label{0.8.2.5-fr}
Avec les notations précédentes, supposons en outre que $X\times X\times X$
existe dans $C$~; la caractérisation du produit comme objet représentant un
foncteur (\hyperref[0.8.1.9-fr]{8.1.9}) entraîne l'existence d'isomorphismes
canoniques
\[
  (X\times X)\times X\xrightarrow{\sim}X\times X\times X
  \xrightarrow{\sim}X\times(X\times X)~;
\]
si on identifie canoniquement $X\times X\times X$ à $(X\times X)\times X$,
l'application $\gamma_X(Y)\times 1_{h_X(Y)}$ s'identifie à
$h_{c_X\times 1_X}(Y)$ pour tout $Y\in C$. Il est par suite équivalent de
dire que pour tout $Y\in C$, la loi interne $\gamma_X(Y)$ est
\emph{associative}, ou que le diagramme d'applications
\[
  \xymatrix{
    h_X(Y)\times h_X(Y)\times h_X(Y)
      \ar[rr]^{\gamma_X(Y)\times 1}\ar[dd]_{1\times\gamma_X(Y)} &&
    h_X(Y)\times h_X(Y)\ar[dd]^{\gamma_X(Y)}\\\\
    h_X(Y)\times h_X(Y)\ar[rr]_{\gamma_X(Y)} && h_X(Y)
  }
\]
\oldpage[0\textsubscript{III}]{11}
est commutatif, ou que le diagramme de morphismes
\[
  \xymatrix{
    X\times X\times X
      \ar[rr]^{c_X\times 1_X}\ar[dd]_{1_X\times c_X} &&
    X\times X\ar[dd]^{c_X}\\\\
    X\times X\ar[rr]_{c_X} && X
  }
\]
est commutatif.
\end{env}

\begin{env}[8.2.6]
\label{0.8.2.6-fr}
Sous les hypothèses de (\hyperref[0.8.2.5-fr]{8.2.5}), si on veut exprimer
que pour tout $Y\in C$, la loi interne $\gamma_X(Y)$ est une \emph{loi de
groupe}, il faut d'une part exprimer qu'elle est associative, et de l'autre
qu'il existe une application $\alpha_X(Y):h_X(Y)\to h_X(Y)$ ayant les
propriétés de l'\emph{inverse} dans un groupe~; comme pour tout morphisme
$u:Y\to Y'$ dans $C$, on a vu que $h_X(u)$ doit être un homomorphisme de
groupes $h_X(Y')\to h_X(Y)$, on voit d'abord que $\alpha_X:h_X\to h_X$ doit
être un \emph{morphisme fonctoriel}. On peut d'autre part exprimer les
propriétés caractéristiques de l'inverse $s\mapsto s^{-1}$ dans un groupe
$G$ sans faire intervenir l'élément neutre~: il suffit d'écrire que les deux
applications composées
\[
  (s,t)\mapsto(s,s^{-1},t)\mapsto(s,s^{-1}t)\mapsto s(s^{-1}t),
\]
\[
  (s,t)\mapsto(s,s^{-1},t)\mapsto(s,ts^{-1})\mapsto(ts^{-1})s
\]
sont égales à la seconde projection $(s,t)\mapsto t$ de $G\times G$ dans
$G$. En vertu de (\hyperref[0.8.1.3-fr]{8.1.3}), on a
$\alpha_X=h_{a_X}$, où $a_X\in\operatorname{Hom}(X,X)$~; la première
condition précédente exprime alors que le morphisme composé
\[
  X\times X\xrightarrow{(1_X,a_X)\times 1_X}X\times X\times X
  \xrightarrow{1_X\times c_X}X\times X\xrightarrow{c_X}X
\]
est la seconde projection $X\times X\to X$ dans $C$, et la seconde condition
se traduit de même.
\end{env}

\begin{env}[8.2.7]
\label{0.8.2.7-fr}
Supposons maintenant qu'il existe dans $C$ un \emph{objet final} $e$
(\hyperref[0.8.1.10-fr]{8.1.10}). Supposons toujours que $\gamma_X(Y)$ soit
une loi de groupe sur $h_X(Y)$ pour tout $Y\in C$, et désignons par
$\eta_X(Y)$ l'élément neutre de $\gamma_X(Y)$. Comme, pour tout morphisme
$u:Y\to Y'$ dans $C$, $h_X(u)$ est un homomorphisme de groupes, on a
$\eta_X(Y)=(h_X(u))(\eta_X(Y'))$~; prenant en particulier $Y'=e$, auquel cas
$u$ est l'unique élément $\varepsilon$ de $\operatorname{Hom}(Y,e)$, on voit
que l'élément $\eta_X(e)$ détermine complètement $\eta_X(Y)$ pour tout
$Y\in C$. Posons $e_X=\eta_X(X)$, élément neutre du groupe
$h_X(X)=\operatorname{Hom}(X,X)$~; la commutativité du diagramme
\[
  \xymatrix{
    h_X(e)\ar[r]^{h_X(\varepsilon)}\ar[d]_{h_{e_X}(e)} &
    h_X(Y)\ar[d]^{h_{e_X}(Y)}\\
    h_X(e)\ar[r]_{h_X(\varepsilon)} & h_X(Y)
  }
\]
(cf. \hyperref[0.8.1.2-fr]{8.1.2}) montre que, dans l'ensemble $h_X(Y)$,
l'application $h_{e_X}(Y)$ n'est autre que
\oldpage[0\textsubscript{III}]{12}
$s\mapsto\eta_X(Y)$ transformant tout élément en l'élément neutre. On
vérifie alors que le fait que $\eta_X(Y)$ soit élément neutre de
$\gamma_X(Y)$ pour tout $Y\in C$ équivaut à dire que le morphisme composé
\[
  X\xrightarrow{(1_X,1_X)}X\times X
  \xrightarrow{1_X\times e_X}X\times X\xrightarrow{c_X}X,
\]
et l'analogue où on permute $1_X$ et $e_X$, sont tous deux \emph{égaux} à
$1_X$.
\end{env}

\begin{env}[8.2.8]
\label{0.8.2.8-fr}
On pourrait bien entendu multiplier sans peine les exemples de structures
algébriques dans les catégories. L'exemple des groupes a été traité avec
assez de détails, mais par la suite nous laisserons généralement au lecteur
le soin de développer des considérations analogues dans les exemples de
structures algébriques que nous rencontrerons.
\end{env}

\section{Ensembles constructibles}
\label{section:0.9-fr}

\subsection{Ensembles constructibles}
\label{subsection:0.9.1-fr}

\begin{definition}[9.1.1]
\label{0.9.1.1-fr}
On dit qu'une application continue $f:X\to Y$ est \emph{quasi-compacte} si
pour tout ouvert quasi-compact $U$ de $Y$, $f^{-1}(U)$ est quasi-compact. On
dit qu'une partie $Z$ d'un espace topologique $X$ est \emph{rétrocompacte
dans $X$} si l'injection canonique $Z\to X$ est quasi-compacte, autrement dit
si pour tout ouvert quasi-compact $U$ de $X$, $U\cap Z$ est quasi-compact.
\end{definition}

Une partie \emph{fermée} de $X$ est rétrocompacte dans $X$, mais une partie
quasi-compacte de $X$ n'est pas nécessairement rétrocompacte dans $X$. Si
$X$ est quasi-compact, toute partie ouverte rétrocompacte dans $X$ est
quasi-compacte. Il est clair que toute réunion \emph{finie} d'ensembles
rétrocompacts dans $X$ est rétrocompacte dans $X$, toute réunion finie
d'ensembles quasi-compacts étant quasi-compacte. Toute intersection finie
d'ouverts rétrocompacts dans $X$ est un ouvert rétrocompact dans $X$. Dans un
espace \emph{localement noethérien} $X$, tout ensemble quasi-compact est un
sous-espace noethérien, et par suite \emph{toute partie} de $X$ est
rétrocompacte dans $X$.

\begin{definition}[9.1.2]
\label{0.9.1.2-fr}
Étant donné un espace topologique $X$, on dit qu'une partie de $X$ est
\emph{constructible} si elle appartient au plus petit ensemble de parties
$\mathfrak{F}$ de $X$ contenant toutes les parties ouvertes rétrocompactes de
$\mathfrak{F}$ et stable par intersection finie et passage au complémentaire
(ce qui implique que $\mathfrak{F}$ est aussi stable par réunion finie).
\end{definition}

\begin{proposition}[9.1.3]
\label{0.9.1.3-fr}
Pour qu'une partie de $X$ soit constructible, il faut et il suffit qu'elle
soit réunion finie d'ensembles de la forme $U\cap\complement{V}$, où $U$ et
$V$ sont des ouverts rétrocompacts dans $X$.
\end{proposition}

Il est clair que la condition est suffisante. Pour voir qu'elle est
nécessaire, considérons l'ensemble $\mathfrak{G}$ des réunions finies
d'ensembles de la forme $U\cap\complement{V}$ où $U$ et $V$ sont ouverts
rétrocompacts dans $X$~; il suffit de voir que tout complémentaire d'un
ensemble de $\mathfrak{G}$ appartient à $\mathfrak{G}$. Soit donc
$Z=\bigcup_{i\in I}(U_i\cap\complement{V_i})$, où $I$ est fini, $U_i$ et
$V_i$ ouverts rétrocompacts dans $X$~; on a
$\complement{Z}=\bigcap_{i\in I}(V_i\cup\complement{U_i})$, donc
$\complement{Z}$ est réunion finie d'ensembles qui sont intersections d'un
certain nombre des $V_i$ et d'un certain nombre de $\complement{U_i}$, donc
de la forme
\oldpage[0\textsubscript{III}]{13}
$V\cap\complement{U}$, où $U$ est réunion d'un certain nombre des $U_i$ et
$V$ intersection d'un certain nombre des $V_i$~; mais on a remarqué plus
haut que les réunions et intersections finies d'ouverts rétrocompacts dans
$X$ sont des ouverts rétrocompacts dans $X$, d'où la conclusion.

\begin{corollary}[9.1.4]
\label{0.9.1.4-fr}
Toute partie constructible de $X$ est rétrocompacte dans $X$.
\end{corollary}

Il suffit de montrer que si $U$ et $V$ sont ouverts rétrocompacts dans $X$,
$U\cap\complement{V}$ est rétrocompact dans $X$~; or, si $W$ est ouvert
quasi-compact dans $X$, $W\cap U\cap\complement{V}$ est fermé dans l'espace
quasi-compact $W\cap U$, donc est quasi-compact.

En particulier~:

\begin{corollary}[9.1.5]
\label{0.9.1.5-fr}
Pour qu'une partie ouverte $U$ de $X$ soit constructible, il faut et il
suffit qu'elle soit rétrocompacte dans $X$. Pour qu'une partie fermée $F$ de
$X$ soit constructible, il faut et il suffit que l'ouvert $\complement{F}$
soit rétrocompact.
\end{corollary}

\begin{env}[9.1.6]
\label{0.9.1.6-fr}
Un cas important est celui où toute partie ouverte quasi-compacte de $X$ est
rétrocompacte, autrement dit, où l'intersection de deux parties ouvertes
quasi-compactes de $X$ est quasi-compacte
(cf. \hyperref[I.5.5.6-fr]{I, 5.5.6}). Lorsque $X$ lui-même est
quasi-compact, cela signifie que les parties ouvertes rétrocompactes dans
$X$ sont identiques aux parties ouvertes quasi-compactes de $X$, et les
parties constructibles de $X$ aux réunions finies d'ensembles de la forme
$U\cap\complement{V}$, où $U$ et $V$ sont ouverts quasi-compacts.
\end{env}

\begin{corollary}[9.1.7]
\label{0.9.1.7-fr}
Pour qu'une partie d'un espace noethérien $X$ soit constructible, il faut et
il suffit qu'elle soit réunion finie de parties localement fermées de $X$.
\end{corollary}

\begin{proposition}[9.1.8]
\label{0.9.1.8-fr}
Soient $X$ un espace topologique, $U$ une partie ouverte de $X$.
\begin{enumerate}
  \item[(i)] Si $T$ est une partie constructible de $X$, $T\cap U$ est une
    partie constructible de $U$.
  \item[(ii)] Supposons en outre $U$ rétrocompact dans $X$. Pour qu'une
    partie $Z$ de $U$ soit constructible dans $X$, il faut et il suffit
    qu'elle soit constructible dans $U$.
\end{enumerate}
\end{proposition}

\begin{enumerate}
  \item[(i)] Utilisant (\hyperref[0.9.1.3-fr]{9.1.3}), on est ramené à
    montrer que si $T$ est ouvert rétrocompact dans $X$, $T\cap U$ est ouvert
    rétrocompact dans $U$, autrement dit, pour tout ouvert quasi-compact
    $W\subset U$, $T\cap U\cap W=T\cap W$ est quasi-compact, ce qui résulte
    aussitôt de l'hypothèse.
  \item[(ii)] La condition étant nécessaire en vertu de (i), il reste à
    démontrer qu'elle est suffisante. Compte tenu de
    (\hyperref[0.9.1.3-fr]{9.1.3}), il suffit de considérer le cas où $Z$ est
    ouvert rétrocompact \emph{dans $U$}, car il s'ensuivra alors que
    $U\setminus Z$ est constructible dans $X$, et si $Z$, $Z'$ sont deux
    ouverts rétrocompacts \emph{dans $U$},
    $Z\cap(U\setminus Z')$ sera bien constructible dans $X$. Or, si $W$ est
    ouvert quasi-compact dans $X$ et $Z$ ouvert rétrocompact dans $U$, on a
    $Z\cap W=Z\cap(W\cap U)$ et par hypothèse $W\cap U$ est ouvert
    quasi-compact dans $U$~; donc $W\cap Z$ est bien quasi-compact, et par
    suite $Z$ est ouvert rétrocompact dans $X$, et \emph{a fortiori}
    constructible dans $X$.
\end{enumerate}

\begin{corollary}[9.1.9]
\label{0.9.1.9-fr}
Soient $X$ un espace topologique, $(U_i)_{i\in I}$ un recouvrement fini de
$X$ formé d'ensembles ouverts rétrocompacts dans $X$. Pour qu'une partie $Z$
de $X$ soit constructible dans $X$, il faut et il suffit que pour tout
$i\in I$, $Z\cap U_i$ soit constructible dans $U_i$.
\end{corollary}

\begin{env}[9.1.10]
\label{0.9.1.10-fr}
Supposons en particulier que $X$ soit quasi-compact et que tout point
\oldpage[0\textsubscript{III}]{14}
de $X$ admette un système fondamental de voisinages ouverts rétrocompacts
dans $X$ (et \emph{a fortiori} quasi-compacts)~; alors la condition pour une
partie $Z$ de $X$ d'être constructible dans $X$ est de nature \emph{locale},
autrement dit, il faut et il suffit que pour tout $x\in X$, il existe un
voisinage ouvert $V$ de $x$ tel que $V\cap Z$ soit constructible dans $V$. En
effet, si cette condition est vérifiée, il existe pour tout $x\in X$ un
voisinage ouvert $V$ de $x$ \emph{rétrocompact dans $X$} et tel que
$V\cap Z$ soit constructible dans $V$, en vertu de l'hypothèse sur $X$ et de
(\hyperref[0.9.1.8-fr]{9.1.8, (i)})~; il suffit alors de recouvrir $X$ par un
nombre fini de ces voisinages, et d'appliquer
(\hyperref[0.9.1.9-fr]{9.1.9}).
\end{env}

\begin{definition}[9.1.11]
\label{0.9.1.11-fr}
Soit $X$ un espace topologique. On dit qu'une partie $T$ de $X$ est
\emph{localement constructible dans $X$} si, pour tout $x\in X$, il existe
un voisinage ouvert $V$ de $x$ tel que $T\cap V$ soit constructible dans $V$.
\end{definition}

Il résulte aussitôt de (\hyperref[0.9.1.8-fr]{9.1.8, (i)}) que si $V$ est
tel que $V\cap T$ soit constructible dans $V$, alors pour tout ouvert
$W\subset V$, $W\cap T$ est constructible dans $W$. Si $T$ est localement
constructible dans $X$, alors, pour tout ouvert $U$ de $X$, $T\cap U$ est
localement constructible dans $U$, comme il résulte de la remarque
précédente. Cette même remarque montre que l'ensemble des parties localement
constructibles dans $X$ est stable par réunion finie et intersection finie~;
il est clair, d'autre part, qu'il est aussi stable par passage aux
complémentaires.

\begin{proposition}[9.1.12]
\label{0.9.1.12-fr}
Soit $X$ un espace topologique. Tout ensemble constructible dans $X$ est
localement constructible dans $X$. La réciproque est vraie si $X$ est
quasi-compact et si sa topologie admet une base formée d'ensembles
rétrocompacts dans $X$.
\end{proposition}

La première assertion résulte de la définition
(\hyperref[0.9.1.11-fr]{9.1.11}) et la seconde de
(\hyperref[0.9.1.10-fr]{9.1.10}).

\begin{corollary}[9.1.13]
\label{0.9.1.13-fr}
Soit $X$ un espace topologique dont la topologie admet une base formée
d'ensembles rétrocompacts dans $X$. Alors toute partie $T$ localement
constructible dans $X$ est rétrocompacte dans $X$.
\end{corollary}

En effet, soit $U$ un ensemble ouvert quasi-compact dans $X$~; $T\cap U$ est
localement constructible dans $U$, donc constructible dans $U$ en vertu de
(\hyperref[0.9.1.12-fr]{9.1.12}), et par suite quasi-compact en vertu de
(\hyperref[0.9.1.4-fr]{9.1.4}).

\subsection{Ensembles constructibles dans les espaces noethériens}
\label{subsection:0.9.2-fr}

\begin{env}[9.2.1]
\label{0.9.2.1-fr}
On a vu (\hyperref[0.9.1.7-fr]{9.1.7}) que dans un espace noethérien $X$, les
parties constructibles dans $X$ sont les \emph{réunions finies de parties
localement fermées de $X$}.

L'image réciproque d'un ensemble constructible dans $X$ par une application
continue d'un espace noethérien $X'$ dans $X$ est constructible dans $X'$.
Si $Y$ est une partie constructible d'un espace noethérien $X$, les parties
de $Y$ qui sont constructibles en tant que sous-espaces de $Y$ sont
identiques à celles qui sont constructibles en tant que sous-espaces de $X$.
\end{env}

\begin{proposition}[9.2.2]
\label{0.9.2.2-fr}
Soient $X$ un espace irréductible noethérien, $E$ une partie constructible de
$X$. Pour que $E$ soit partout dense dans $X$, il faut et il suffit que $E$
contienne une partie ouverte non vide de $X$.
\end{proposition}

\oldpage[0\textsubscript{III}]{15}
La condition est évidemment suffisante, tout ensemble ouvert non vide étant
dense dans $X$. Inversement, soit
$E=\bigcup_{i=1}^n(U_i\cap F_i)$ une partie constructible de $X$, les $U_i$
étant ouverts non vides et les $F_i$ fermés dans $X$~; on a donc
$\overline{E}\subset\bigcup_i F_i$. Par suite, si $\overline{E}=X$, $X$ est
égal à l'un des $F_i$, donc $E\supset U_i$, ce qui achève la démonstration.

Lorsque $X$ admet un point générique $x$
(\hyperref[0.2.1.2-fr]{$0_{\mathrm I}$, 2.1.2}), la condition de
(\hyperref[0.9.2.2-fr]{9.2.2}) équivaut à la relation $x\in E$.

\begin{proposition}[9.2.3]
\label{0.9.2.3-fr}
Soit $X$ un espace noethérien. Pour qu'une partie $E$ de $X$ soit
constructible, il faut et il suffit que, pour toute partie fermée
irréductible $Y$ de $X$, $E\cap Y$ soit rare dans $Y$ ou contienne une partie
ouverte non vide de $Y$.
\end{proposition}

La nécessité de la condition provient de ce que $E\cap Y$ doit être une
partie constructible de $Y$ et de (\hyperref[0.9.2.2-fr]{9.2.2}), car une
partie non dense de $Y$ est nécessairement rare dans l'espace irréductible
$Y$ (\hyperref[0.2.1.1-fr]{$0_{\mathrm I}$, 2.1.1}). Pour prouver que la
condition est suffisante, appliquons le principe de récurrence noethérienne
(\hyperref[0.2.2.2-fr]{$0_{\mathrm I}$, 2.2.2}) à l'ensemble
$\mathfrak{F}$ des parties fermées $Y$ de $X$ telles que $Y\cap E$ soit
constructible (par rapport à $Y$ ou par rapport à $X$, ce qui revient au
même)~: on peut donc supposer que pour toute partie fermée $Y\ne X$ de $X$,
$E\cap Y$ est constructible. Supposons d'abord que $X$ ne soit pas
irréductible, et soient $X_i$ ($1\leq i\leq m$) ses composantes
irréductibles, nécessairement en nombre fini
(\hyperref[0.2.2.5-fr]{$0_{\mathrm I}$, 2.2.5})~; par hypothèse, les
$E\cap X_i$ sont constructibles, donc aussi leur réunion $E$. Supposons
ensuite que $X$ soit irréductible~; alors, par hypothèse, ou bien $E$ est
rare, donc $\overline{E}\ne X$ et $E=E\cap\overline{E}$ est constructible~;
ou bien $E$ contient un ouvert non vide $U$ de $E$, donc est réunion de $U$
et de $E\cap(X\setminus U)$~; mais $X\setminus U$ est un ensemble fermé
distinct de $X$, donc $E\cap(X\setminus U)$ est constructible~; $E$
lui-même est par suite constructible, ce qui achève la démonstration.

\begin{corollary}[9.2.4]
\label{0.9.2.4-fr}
Soient $X$ un espace noethérien, $(E_\alpha)$ une famille filtrante
croissante de parties constructibles de $X$, telle que~:
\begin{enumerate}
  \item[$1^\circ$] $X$ est réunion de la famille $(E_\alpha)$.
  \item[$2^\circ$] Toute partie fermée irréductible de $X$ est contenue dans
    l'adhérence d'un des $E_\alpha$.
\end{enumerate}
Alors il existe un indice $\alpha$ tel que $X=E_\alpha$.

Lorsque toute partie fermée irréductible de $X$ admet un point générique,
l'hypothèse $2^\circ$ peut être supprimée.
\end{corollary}

Appliquons le principe de récurrence noethérienne
(\hyperref[0.2.2.2-fr]{$0_{\mathrm I}$, 2.2.2}) à l'ensemble
$\mathfrak{M}$ des parties fermées de $X$ contenues dans l'un des $E_\alpha$
au moins~; on peut donc supposer que toute partie fermée $Y\ne X$ de $X$ est
contenue dans un des $E_\alpha$. La proposition est évidente si $X$ n'est pas
irréductible, car chacune des composantes irréductibles $X_i$ de $X$
($1\leq i\leq m$) est contenue dans un $E_{\alpha_i}$, et il existe un
$E_\alpha$ contenant tous les $E_{\alpha_i}$. Supposons donc $X$
irréductible. Par hypothèse, il existe $\beta$ tel que
$X=\overline{E_\beta}$, donc (\hyperref[0.9.2.2-fr]{9.2.2}) $E_\beta$
contient un ouvert non vide $U$ de $X$. Mais alors l'ensemble fermé
$X\setminus U$ est contenu dans un $E_\gamma$, et il suffit de prendre
$E_\alpha$ contenant $E_\beta$ et $E_\gamma$. Lorsque toute partie fermée
irréductible $Y$ de $X$
\oldpage[0\textsubscript{III}]{16}
admet un point générique $y$, il existe $\alpha$ tel que $y\in E_\alpha$,
donc $Y=\overline{\{y\}}\subset\overline{E_\alpha}$, et la condition
$2^\circ$ est conséquence de $1^\circ$.

\begin{proposition}[9.2.5]
\label{0.9.2.5-fr}
Soient $X$ un espace noethérien, $x$ un point de $X$, $E$ une partie
constructible de $X$. Pour que $E$ soit un voisinage de $x$, il faut et il
suffit que pour toute partie fermée irréductible $Y$ de $X$ contenant $x$,
$E\cap Y$ soit dense dans $Y$ (s'il existe un point générique $y$ de $Y$,
cela signifie aussi (\hyperref[0.9.2.2-fr]{9.2.2}) que $y\in E$).
\end{proposition}

La condition est évidemment nécessaire~; prouvons qu'elle est suffisante.
Appliquant le principe de récurrence noethérienne à l'ensemble $\mathfrak{M}$
des parties fermées $Y$ de $X$ contenant $x$ et telles que $E\cap Y$ soit un
voisinage de $x$ dans $Y$, on peut supposer que toute partie fermée $Y\ne X$
de $X$ contenant $x$ appartient à $\mathfrak{M}$. Si $X$ n'est pas
irréductible, chacune des composantes irréductibles $X_i$ de $X$ contenant
$x$ est distincte de $X$, donc $E\cap X_i$ est un voisinage de $x$ par
rapport à $X_i$~; par suite, $E$ est un voisinage de $x$ dans la réunion des
composantes irréductibles de $X$ contenant $x$, et comme cette réunion est un
voisinage de $x$ dans $X$, il en est de même de $E$. Si $X$ est
irréductible, $E$ est dense dans $X$ par hypothèse, donc contient une partie
ouverte non vide $U$ de $X$ (\hyperref[0.9.2.2-fr]{9.2.2})~; la proposition
est alors évidente si $x\in U$~; sinon, $x$ est par hypothèse intérieur à
$E\cap(X\setminus U)$ par rapport à $X\setminus U$, donc l'adhérence dans
$X$ de $X\setminus E$ ne contient pas $x$, et le complémentaire de cette
adhérence est un voisinage de $x$ contenu dans $E$, ce qui achève la
démonstration.

\begin{corollary}[9.2.6]
\label{0.9.2.6-fr}
Soient $X$ un espace noethérien, $E$ une partie de $X$. Pour que $E$ soit un
ensemble ouvert dans $X$, il faut et il suffit que pour toute partie fermée
irréductible $Y$ de $X$ rencontrant $E$, $E\cap Y$ contienne une partie
ouverte non vide de $Y$.
\end{corollary}

La condition est évidemment nécessaire~; inversement, si elle est vérifiée,
elle implique que $E$ est constructible en vertu de
(\hyperref[0.9.2.3-fr]{9.2.3}). En outre,
(\hyperref[0.9.2.5-fr]{9.2.5}) montre que $E$ est alors voisinage de chacun
de ses points, d'où la conclusion.

\subsection{Fonctions constructibles}
\label{subsection:0.9.3-fr}

\begin{definition}[9.3.1]
\label{0.9.3.1-fr}
Soit $h$ une application d'un espace topologique $X$ dans un ensemble $T$.
On dit que $h$ est \emph{constructible} si $h^{-1}(t)$ est constructible pour
tout $t\in T$, et vide sauf pour un nombre fini de valeurs de $t$~; pour
toute partie $S$ de $T$, $h^{-1}(S)$ est alors constructible. On dit que $h$
est \emph{localement constructible} si tout $x\in X$ possède un voisinage
ouvert $V$ tel que $h|V$ soit constructible.
\end{definition}

Toute fonction constructible est localement constructible~; la réciproque
est vraie quand $X$ est quasi-compact et admet une base formée d'ensembles
ouverts rétrocompacts dans $X$ (en particulier quand $X$ est noethérien).

\begin{proposition}[9.3.2]
\label{0.9.3.2-fr}
Soit $h$ une application d'un espace noethérien $X$ dans un ensemble $T$.
Pour que $h$ soit constructible, il faut et il suffit que pour toute partie
fermée irréductible $Y$ de $X$, il existe une partie non vide $U$ de $Y$,
ouverte par rapport à $Y$, et dans laquelle $h$ soit constante.
\end{proposition}

La condition est nécessaire~: en effet, par hypothèse, $h$ ne prend dans $Y$
qu'un nombre fini de valeurs $t_i$, et chacun des ensembles
$h^{-1}(t_i)\cap Y$ est constructible dans $Y$
(\hyperref[0.9.2.1-fr]{9.2.1})~; comme ils ne peuvent être tous des parties
rares de l'espace $Y$, un d'eux au moins contient un ensemble ouvert non vide
(\hyperref[0.9.2.3-fr]{9.2.3}).

\oldpage[0\textsubscript{III}]{17}
Pour voir que la condition est suffisante, appliquons le principe de
récurrence noethérienne à l'ensemble $\mathfrak{M}$ des parties fermées $Y$
de $X$ telles que la restriction $h|Y$ soit constructible~; on peut donc
supposer que pour toute partie fermée $Y\ne X$ de $X$, $h|Y$ est
constructible. Si $X$ n'est pas irréductible, la restriction de $h$ à
chacune des composantes irréductibles $X_i$ de $X$ (en nombre fini) est donc
constructible, et il résulte alors aussitôt de la définition
(\hyperref[0.9.3.1-fr]{9.3.1}) que $h$ est constructible. Si $X$ est
irréductible, il existe par hypothèse une partie ouverte non vide $U$ de $X$
dans laquelle $h$ est constante~; d'autre part, la restriction de $h$ à
$X\setminus U$ est constructible par hypothèse, et il en résulte aussitôt
que $h$ est constructible.

\begin{corollary}[9.3.3]
\label{0.9.3.3-fr}
Soit $X$ un espace noethérien dans lequel toute partie fermée irréductible
admet un point générique. Si $h$ est une application de $X$ dans un ensemble
$T$ telle que, pour tout $t\in T$, $h^{-1}(t)$ soit constructible, alors $h$
est constructible.
\end{corollary}

En effet, si $Y$ est une partie fermée irréductible de $X$ et $y$ son point
générique, $Y\cap h^{-1}(h(y))$ est constructible et contient $y$, donc
(\hyperref[0.9.2.2-fr]{9.2.2}) cet ensemble contient une partie ouverte non
vide de $Y$, et il suffit d'appliquer (\hyperref[0.9.3.2-fr]{9.3.2}).

\begin{proposition}[9.3.4]
\label{0.9.3.4-fr}
Soient $X$ un espace noethérien dans lequel toute partie fermée irréductible
admet un point générique, $h$ une application constructible de $X$ dans un
ensemble ordonné. Pour que $h$ soit semi-continue supérieurement dans $X$,
il faut et il suffit que pour tout $x\in X$ et toute générisation
(\hyperref[0.2.1.2-fr]{$0_{\mathrm I}$, 2.1.2}) $x'$ de $x$, on ait
$h(x')\leq h(x)$.
\end{proposition}

La fonction $h$ ne prend qu'un nombre fini de valeurs~; dire qu'elle est
semi-continue supérieurement signifie donc que pour tout $x\in X$,
l'ensemble $E$ des $y\in X$ tels que $h(y)\leq h(x)$ est un voisinage de $x$.
Par hypothèse, $E$ est une partie constructible de $X$~; d'autre part, dire
qu'une partie fermée irréductible $Y$ de $X$ contient $x$ signifie que son
point générique $y$ est une générisation de $x$~; la conclusion résulte
alors de (\hyperref[0.9.2.5-fr]{9.2.5}).

\section{Compléments sur les modules plats}
\label{section:0.10-fr}

Pour les démonstrations des propriétés énoncées sans démonstration dans les
n\textsuperscript{os} (\hyperref[subsection:0.10.1-fr]{10.1}) et
(\hyperref[subsection:0.10.2-fr]{10.2}), nous renvoyons le lecteur à Bourbaki,
\emph{Alg. comm.}, chap. II et III.

\subsection{Relations entre modules plats et modules libres}
\label{subsection:0.10.1-fr}

\begin{env}[10.1.1]
\label{0.10.1.1-fr}
Soient $A$ un anneau, $\mathfrak{J}$ un idéal de $A$, $M$ un $A$-module~; pour
tout entier $p\geq0$, on a un homomorphisme canonique de
$(A/\mathfrak{J})$-modules
\[
\label{0.10.1.1.1-fr}
  \varphi_p:(M/\mathfrak{J}M)\otimes_{A/\mathfrak{J}}
  (\mathfrak{J}^p/\mathfrak{J}^{p+1})
  \longrightarrow \mathfrak{J}^pM/\mathfrak{J}^{p+1}M
  \tag{10.1.1.1}
\]
qui est évidemment \emph{surjectif}. Nous désignerons par
$\operatorname{gr}(A)=\bigoplus_{p\geq0}\mathfrak{J}^p/\mathfrak{J}^{p+1}$ l'anneau gradué
associé à $A$ filtré par les $\mathfrak{J}^p$, par
$\operatorname{gr}(M)=\bigoplus_{p\geq0}\mathfrak{J}^pM/\mathfrak{J}^{p+1}M$ le
$\operatorname{gr}(A)$-module gradué associé à $M$ filtré par les $\mathfrak{J}^pM$~; on a
donc $\operatorname{gr}_p(A)=\mathfrak{J}^p/\mathfrak{J}^{p+1}$,
$\operatorname{gr}_p(M)=\mathfrak{J}^pM/\mathfrak{J}^{p+1}M$~; les $\varphi_p$ définissent
un homomorphisme \emph{surjectif} de $\operatorname{gr}(A)$-modules gradués
\[
\label{0.10.1.1.2-fr}
  \varphi:\operatorname{gr}_0(M)\otimes_{\operatorname{gr}_0(A)}
  \operatorname{gr}(A)\longrightarrow\operatorname{gr}(M).
  \tag{10.1.1.2}
\]
\end{env}

\oldpage[0\textsubscript{III}]{18}
\begin{env}[10.1.2]
\label{0.10.1.2-fr}
Supposons vérifiée l'une des hypothèses suivantes~:
\begin{enumerate}
  \item[(i)] $\mathfrak{J}$ est nilpotent~;
  \item[(ii)] $A$ est noethérien, $\mathfrak{J}$ est contenu dans le radical de
    $A$, et $M$ est de type fini.
\end{enumerate}
Alors les propriétés suivantes sont équivalentes~:
\begin{enumerate}
  \item[$a)$] $M$ est un $A$-module libre.
  \item[$b)$] $M/\mathfrak{J}M=M\otimes_A(A/\mathfrak{J})$ est un
    $(A/\mathfrak{J})$-module libre, et
    $\operatorname{Tor}_1^A(M,A/\mathfrak{J})=0$.
  \item[$c)$] $M/\mathfrak{J}M$ est un $(A/\mathfrak{J})$-module libre et
    l'homomorphisme canonique (\hyperref[0.10.1.1.2-fr]{10.1.1.2}) est
    injectif (donc bijectif).
\end{enumerate}
\end{env}

\begin{env}[10.1.3]
\label{0.10.1.3-fr}
Supposons que $A/\mathfrak{J}$ soit un \emph{corps} (autrement dit que
$\mathfrak{J}$ soit maximal), et que l'une des hypothèses (i), (ii) de
(\hyperref[0.10.1.2-fr]{10.1.2}) soit vérifiée. Alors les propriétés suivantes
sont équivalentes~:
\begin{enumerate}
  \item[$a)$] $M$ est un $A$-module libre.
  \item[$b)$] $M$ est un $A$-module projectif.
  \item[$c)$] $M$ est un $A$-module plat.
  \item[$d)$] $\operatorname{Tor}_1^A(M,A/\mathfrak{J})=0$.
  \item[$e)$] L'homomorphisme canonique
    (\hyperref[0.10.1.1.2-fr]{10.1.1.2}) est bijectif.
\end{enumerate}
Ce résultat s'appliquera en particulier dans les deux cas suivants~:
\begin{enumerate}
  \item[(i)] $M$ est un module \emph{quelconque} sur un anneau local $A$ dont
    l'idéal maximal $\mathfrak{J}$ est \emph{nilpotent} (par exemple un anneau
    local artinien).
  \item[(ii)] $M$ est un module \emph{de type fini} sur un anneau local
    \emph{noethérien}.
\end{enumerate}
\end{env}

\subsection{Critères locaux de platitude}
\label{subsection:0.10.2-fr}

\begin{env}[10.2.1]
\label{0.10.2.1-fr}
Les hypothèses et notations étant celles de
(\hyperref[0.10.1.1-fr]{10.1.1}), considérons les conditions suivantes~:
\begin{enumerate}
  \item[$a)$] $M$ est un $A$-module plat.
  \item[$b)$] $M/\mathfrak{J}M$ est un $(A/\mathfrak{J})$-module plat et
    $\operatorname{Tor}_1^A(M,A/\mathfrak{J})=0$.
  \item[$c)$] $M/\mathfrak{J}M$ est un $(A/\mathfrak{J})$-module plat et
    l'homomorphisme canonique (\hyperref[0.10.1.1.2-fr]{10.1.1.2}) est
    bijectif.
  \item[$d)$] Pour tout $n\geq1$, $M/\mathfrak{J}^nM$ est un
    $(A/\mathfrak{J}^n)$-module plat.
\end{enumerate}
On a alors les implications
\[
  a\Rightarrow b\Rightarrow c\Rightarrow d
\]
et si $\mathfrak{J}$ est \emph{nilpotent}, les quatre conditions $a)$, $b)$,
$c)$, $d)$ sont \emph{équivalentes}. Il en est de même si $A$ est noethérien
et en outre si $M$ est \emph{idéalement séparé}, c'est-à-dire que pour tout
idéal $\mathfrak{a}$ de $A$, le $A$-module $\mathfrak{a}\otimes_A M$ est
\emph{séparé} pour la topologie $\mathfrak{J}$-pré-adique.
\end{env}

\begin{env}[10.2.2]
\label{0.10.2.2-fr}
Soient $A$ un anneau noethérien, $B$ une $A$-algèbre commutative noethérienne,
$\mathfrak{J}$ un idéal de $A$ tel que $\mathfrak{J}B$ soit contenu dans le
radical de $B$, $M$ un $B$-module de type fini. Alors, lorsque $M$ est
considéré comme $A$-module, les conditions $a)$, $b)$, $c)$, $d)$ de
(\hyperref[0.10.2.1-fr]{10.2.1}) sont équivalentes.

Ce résultat s'applique surtout lorsque $A$ et $B$ sont des anneaux
\emph{locaux} noethériens, l'homomorphisme $A\to B$ un homomorphisme
\emph{local}. Plus particulièrement, si $\mathfrak{J}$ est alors l'idéal
\emph{maximal} de $A$, on peut, dans les conditions $b)$ et $c)$, supprimer
l'hypothèse que $M/\mathfrak{J}M$ est plat, qui est automatiquement vérifiée,
et la condition $d)$ signifie que les modules $M/\mathfrak{J}^nM$ sont
\emph{libres} sur les $A/\mathfrak{J}^n$.
\end{env}

\oldpage[0\textsubscript{III}]{19}
\begin{env}[10.2.3]
\label{0.10.2.3-fr}
Les hypothèses sur $A$, $B$, $\mathfrak{J}$, $M$ étant celles formulées au
début de (\hyperref[0.10.2.2-fr]{10.2.2}), soient $\widehat A$ le séparé
complété de $A$ pour la topologie $\mathfrak{J}$-pré-adique, $\widehat M$ le
séparé complété de $M$ pour la topologie $\mathfrak{J}B$-pré-adique. Alors,
pour que $M$ soit un $A$-module plat, il faut et il suffit que $\widehat M$
soit un $\widehat A$-module plat.
\end{env}

\begin{env}[10.2.4]
\label{0.10.2.4-fr}
Soient $\rho:A\to B$ un homomorphisme local d'anneaux locaux noethériens, $k$
le corps résiduel de $A$, $M$, $N$ deux $B$-modules de type fini, $N$ étant
supposé être $A$-\emph{plat}. Soit $u:M\to N$ un $B$-homomorphisme. Alors les
conditions suivantes sont équivalentes~:
\begin{enumerate}
  \item[$a)$] $u$ est injectif et $\operatorname{Coker}(u)$ est un $A$-module plat.
  \item[$b)$] $u\otimes1:M\otimes_A k\to N\otimes_A k$ est injectif.
\end{enumerate}
\end{env}

\begin{env}[10.2.5]
\label{0.10.2.5-fr}
Soient $\rho:A\to B$, $\sigma:B\to C$ des homomorphismes locaux d'anneaux
locaux noethériens, $k$ le corps résiduel de $A$, $M$ un $C$-module de type
fini. On suppose que $B$ est un $A$-module \emph{plat} Alors les conditions
suivantes sont équivalentes~:
\begin{enumerate}
  \item[$a)$] $M$ est un $B$-module plat.
  \item[$b)$] $M$ est un $A$-module plat, et $M\otimes_A k$ est un
    $(B\otimes_A k)$-module plat.
\end{enumerate}
\end{env}

\begin{proposition}[10.2.6]
\label{0.10.2.6-fr}
Soient $A$, $B$ deux anneaux locaux noethériens, $\rho:A\to B$ un
homomorphisme local, $\mathfrak{J}$ un idéal de $B$ contenu dans l'idéal
maximal, $M$ un $B$-module de type fini. Supposons que pour tout $n\geq0$,
$M_n=M/\mathfrak{J}^{n+1}M$ soit un $A$-module plat. Alors $M$ est un
$A$-module plat.
\end{proposition}

Il faut prouver que pour tout homomorphisme injectif $u:N'\to N$ de
$A$-modules de type fini, $v=1\otimes u:M\otimes_A N'\to M\otimes_A N$ est
injectif. Or, $M\otimes_A N'$ et $M\otimes_A N$ sont des $B$-modules de type
fini, donc séparés pour la topologie $\mathfrak{J}$-pré-adique
(\hyperref[0.7.3.5-fr]{$0_{\mathrm I}$, 7.3.5})~; il suffit donc de prouver que
l'homomorphisme
$\widehat v:\widehat{(M\otimes_A N')}\to\widehat{(M\otimes_A N)}$ pour les
séparés complétés est injectif. Or, on a $\widehat v=\varprojlim v_n$, où
$v_n$ est l'homomorphisme $1\otimes u:M_n\otimes_A N'\to M_n\otimes_A N$~;
comme par hypothèse $M_n$ est $A$-plat, $v_n$ est injectif pour tout $n$, donc
il en est de même de $v$, le foncteur $\varprojlim$ étant exact à gauche.

\begin{corollary}[10.2.7]
\label{0.10.2.7-fr}
Soient $A$ un anneau noethérien, $B$ un anneau local noethérien,
$\rho:A\to B$ un homomorphisme, $f$ un élément de l'idéal maximal de $B$, $M$
un $B$-module de type fini. Supposons que l'homothétie
$f_M:x\mapsto fx$ de $M$ soit injective et que $M/fM$ soit un $A$-module
plat. Alors $M$ est un $A$-module plat.
\end{corollary}

Posons $M_i=f^iM$ pour $i\geq0$~; comme $f_M$ est injective,
$M_i/M_{i+1}$ est isomorphe à $M/fM$, donc $A$-plat pour tout $i\geq0$~; de
la suite exacte
\[
  0\longrightarrow M_i/M_{i+1}\longrightarrow M/M_{i+1}
  \longrightarrow M/M_i\longrightarrow0
\]
on tire par récurrence sur $i$ que $M/M_i$ est $A$-\emph{plat} pour tout
$i\geq0$ (\hyperref[0.6.1.2-fr]{$0_{\mathrm I}$, 6.1.2})~; on peut donc
appliquer (\hyperref[0.10.2.6-fr]{10.2.6}). On peut aussi raisonner directement
comme suit~: pour tout $A$-module $N$ de type fini, $M\otimes_A N$ est un
$B$-module de type fini~; comme $f$ appartient au radical $\mathfrak n$ de
$B$, la topologie $(f)$-adique sur $M\otimes_A N$ est plus fine que la
\oldpage[0\textsubscript{III}]{20}
topologie $\mathfrak n$-adique, et on sait que cette dernière est \emph{séparée}
(\hyperref[0.7.3.5-fr]{$0_{\mathrm I}$, 7.3.5}). D'ailleurs, comme $M/M_i$ est
$A$-plat, on a
\[
  f^i(M\otimes_A N)=\operatorname{Im}(M_i\otimes_A N\to M\otimes_A N)
  =\operatorname{Ker}(M\otimes_A N\to(M/M_i)\otimes_A N)
\]
(\hyperref[0.6.1.2-fr]{$0_{\mathrm I}$, 6.1.2}). Soient alors $N$ un
$A$-module de type fini, $N'$ un sous-module de $N$, $j:N'\to N$ l'injection
canonique~; dans le diagramme commutatif
\[
\xymatrix{
  M\otimes_A N'\ar[r]\ar[d]_{1_M\otimes j} &
  (M/M_i)\otimes_A N'\ar[d]^{1_{M/M_i}\otimes j}\\
  M\otimes_A N\ar[r] & (M/M_i)\otimes_A N
}
\]
$1_{M/M_i}\otimes j$ est injectif puisque $M/M_i$ est $A$-plat~; on en
conclut que
\[
  \operatorname{Ker}(M\otimes_A N'\to M\otimes_A N)
  \subset\operatorname{Ker}(M\otimes_A N'\to(M/M_i)\otimes_A N')
\]
quel que soit $i$~; puisque l'intersection des seconds membres est réduite à
$0$ comme on l'a vu plus haut, il en est de même du premier membre, et par
suite $M$ est $A$-plat.

\begin{proposition}[10.2.8]
\label{0.10.2.8-fr}
Soient $A$ un anneau noethérien réduit, $M$ un $A$-module de type fini. On
suppose que pour toute $A$-algèbre $B$, qui est un anneau de valuation
discrète, $M\otimes_A B$ soit un $B$-module plat (donc libre
(\hyperref[0.10.1.3-fr]{10.1.3})). Alors $M$ est un $A$-module plat.
\end{proposition}

On sait que pour que $M$ soit plat, il faut et il suffit que pour tout idéal
maximal $\mathfrak m$ de $A$, $M_{\mathfrak m}$ soit un
$A_{\mathfrak m}$-module plat
(\hyperref[0.6.3.3-fr]{$0_{\mathrm I}$, 6.3.3})~; on peut donc se borner au cas
où $A$ est \emph{local} (\hyperref[0.1.2.8-fr]{$0_{\mathrm I}$, 1.2.8}).
Soient alors $\mathfrak m$ l'idéal maximal de $A$, $\mathfrak p_i$
$(1\leq i\leq r)$ ses idéaux premiers minimaux, $k$ le corps résiduel
$A/\mathfrak m$. On sait (\hyperref[II.7.1.7-fr]{II, 7.1.7}) qu'il existe
pour chaque $i$ un anneau de valuation discrète $B_i$, ayant même corps des
fractions $K_i$ que l'anneau intègre $A/\mathfrak p_i$, et dominant ce
dernier. Posons $M_i=M\otimes_A B_i$. Par hypothèse, $M_i$ est libre sur
$B_i$, donc on a, en désignant par $k_i$ le corps résiduel de $B_i$
\[
\label{0.10.2.8.1-fr}
  \operatorname{rg}_{k_i}(M_i\otimes_{B_i}k_i)
  =\operatorname{rg}_{K_i}(M_i\otimes_{B_i}K_i).
  \tag{10.2.8.1}
\]
Mais il est clair que l'homomorphisme composé
$A\to A/\mathfrak p_i\to B_i$ est local, donc $k$ est une extension de $k_i$,
et l'on a
$M_i\otimes_{B_i}k_i=M\otimes_A k_i=(M\otimes_A k)\otimes_k k_i$, et par
ailleurs $M_i\otimes_{B_i}K_i=M\otimes_A K_i$. L'égalité
(\hyperref[0.10.2.8.1-fr]{10.2.8.1}) entraîne donc
\[
  \operatorname{rg}_k(M\otimes_A k)=\operatorname{rg}_{K_i}(M\otimes_A K_i)
  \qquad\text{pour }1\leq i\leq r
\]
et comme $A$ est réduit, on sait que cette condition entraîne que $M$ est un
$A$-module \emph{libre} (Bourbaki, \emph{Alg. comm.}, chap. II, \S~3,
n\textsuperscript{o}~2, prop. 7).

\subsection{Existence d'extensions plates d'anneaux locaux}
\label{subsection:0.10.3-fr}

\begin{proposition}[10.3.1]
\label{0.10.3.1-fr}
Soient $A$ un anneau local noethérien, $\mathfrak J$ son idéal maximal,
$k=A/\mathfrak J$ son corps résiduel. Soit $K$ une extension du corps $k$. Il
existe un homomorphisme local de $A$ dans un anneau local noethérien $B$, tel
que $B/\mathfrak J B$ soit $k$-isomorphe à $K$, et que $B$ soit un
$A$-module plat.
\end{proposition}

Nous démontrerons cette proposition en plusieurs étapes.

\begin{env}[10.3.1.1]
\label{0.10.3.1.1-fr}
Supposons d'abord que $K=k(T)$, où $T$ est une indéterminée. Dans l'anneau de
polynômes $A'=A[T]$, considérons l'idéal premier
$\mathfrak J'=\mathfrak J A'$ formé des
\oldpage[0\textsubscript{III}]{21}
polynômes ayant leurs coefficients dans l'idéal $\mathfrak J$~; il est clair
que $A'/\mathfrak J'$ est canoniquement isomorphe à $k[T]$. Montrons que
l'anneau de fractions $B=A'_{\mathfrak J'}$ répond à la question~; c'est
évidemment un anneau local noethérien dont l'idéal maximal
$\mathfrak L=\mathfrak J B$. En outre,
$B/\mathfrak L=(A'/\mathfrak J')_{\mathfrak J'}=(k[T])_{\mathfrak J'}$
n'est autre que le corps des fractions $K$ de $k[T]$. Enfin, $B$ est un
$A'$-module plat et $A'$ un $A$-module libre, donc $B$ est un $A$-module plat
(\hyperref[0.6.2.1-fr]{$0_{\mathrm I}$, 6.2.1}).
\end{env}

\begin{env}[10.3.1.2]
\label{0.10.3.1.2-fr}
Supposons ensuite que $K=k(t)=k[t]$, où $t$ est algébrique sur $k$~; soit
$f\in k[T]$ le polynôme minimal de $t$~; il existe un polynôme unitaire
$F\in A[T]$ dont l'image canonique dans $k[T]$ soit $f$. Posons encore
$A'=A[T]$, et soit $\mathfrak J'$ l'idéal $\mathfrak J A'+(F)$ dans $A'$.
Nous allons voir que l'anneau quotient $B=A'/(F)$ répond cette fois à la
question. Tout d'abord, il est clair que $B$ est un $A$-module \emph{libre},
donc plat. L'anneau $A'/\mathfrak J'$ est isomorphe à
\[
  (A'/\mathfrak J A')/((\mathfrak J A'+(F))/\mathfrak J A')
  =k[T]/(f)=K~;
\]
l'image $\mathfrak L$ de $\mathfrak J'$ dans $B$ est donc maximal et on a
évidemment $\mathfrak L=\mathfrak J B$. Enfin, $B$ est un anneau semi-local,
puisqu'il est un $A$-module de type fini (Bourbaki, \emph{Alg. comm.}, chap.
IV, \S~2, n\textsuperscript{o}~5, cor. 3 de la prop. 9), et ses idéaux
maximaux sont en correspondance biunivoque avec ceux de $B/\mathfrak J B$
([13], vol. I, p. 259)~; ce qui précède prouve donc que $B$ est un anneau
local.
\end{env}

\begin{lemma}[10.3.1.3]
\label{0.10.3.1.3-fr}
Soit $(A_\lambda,f_{\mu\lambda})$ un système inductif filtrant d'anneaux
locaux, tel que les $f_{\mu\lambda}$ soient des homomorphismes locaux~; soit
$\mathfrak m_\lambda$ l'idéal maximal de $A_\lambda$, et soit
$K_\lambda=A_\lambda/\mathfrak m_\lambda$. Alors
$A'=\varinjlim A_\lambda$ est un anneau local dont
$\mathfrak m'=\varinjlim\mathfrak m_\lambda$ est l'idéal maximal, et
$K=\varinjlim K_\lambda$ le corps résiduel. En outre, si
$\mathfrak m_\mu=\mathfrak m_\lambda A_\mu$ pour $\lambda<\mu$, on a
$\mathfrak m'=\mathfrak m_\lambda A'$ pour tout $\lambda$. Si, de plus, pour
$\lambda<\mu$, $A_\mu$ est un $A_\lambda$-module plat, et si tous les
$A_\lambda$ sont noethériens, alors $A'$ est noethérien et est un
$A_\lambda$-module plat pour tout $\lambda$.
\end{lemma}

Comme $f_{\mu\lambda}(\mathfrak m_\lambda)\subset\mathfrak m_\mu$ pour
$\lambda<\mu$ par hypothèse, les $\mathfrak m_\lambda$ forment un système
inductif, et sa limite $\mathfrak m'$ est évidemment un idéal de $A'$. En
outre, si $x'\notin\mathfrak m'$, il existe $\lambda$ tel que
$x'=f_\lambda(x_\lambda)$ pour un $x_\lambda\in A_\lambda$
($f_\lambda:A_\lambda\to A'$ désignant l'homomorphisme canonique)~; puisque
$x'\notin\mathfrak m'$, on a nécessairement
$x_\lambda\notin\mathfrak m_\lambda$, donc $x_\lambda$ admet un inverse
$y_\lambda$ dans $A_\lambda$, et $y'=f_\lambda(y_\lambda)$ est l'inverse de
$x'$ dans $A'$, ce qui prouve que $A'$ est un anneau local d'idéal maximal
$\mathfrak m'$~; l'assertion relative à $K$ résulte aussitôt du fait que
$\varinjlim$ est un foncteur exact. L'hypothèse
$\mathfrak m_\mu=\mathfrak m_\lambda A_\mu$ signifie que l'application
canonique
$\mathfrak m_\lambda\otimes_{A_\lambda}A_\mu\to\mathfrak m_\mu$ est
surjective~; la relation $\mathfrak m'=\mathfrak m_\lambda A'$ résulte donc
encore de l'exactitude du foncteur $\varinjlim$ et du fait qu'il commute avec
le produit tensoriel.

Supposons maintenant que pour $\lambda<\mu$, on ait
$\mathfrak m_\mu=\mathfrak m_\lambda A_\mu$ et que $A_\mu$ soit un
$A_\lambda$-module plat. Alors $A'$ est un $A_\lambda$-module plat pour tout
$\lambda$, en vertu de (\hyperref[0.6.2.3-fr]{$0_{\mathrm I}$, 6.2.3})~;
comme $A'$ et $A_\lambda$ sont des anneaux locaux et que
$\mathfrak m'=\mathfrak m_\lambda A'$, $A'$ est même un
$A_\lambda$-module \emph{fidèlement plat}
(\hyperref[0.6.6.2-fr]{$0_{\mathrm I}$, 6.6.2}). Supposons enfin, en outre,
les $A_\lambda$ \emph{noethériens}~; les topologies
$\mathfrak m_\lambda$-préadiques sont alors séparées
(\hyperref[0.7.3.5-fr]{$0_{\mathrm I}$, 7.3.5})~; montrons qu'il en résulte
d'abord que sur $A'$ la topologie $\mathfrak m'$-adique est \emph{séparée}.
En effet, si $x'\in A'$ appartient à tous les $\mathfrak m'^{n}$ $(n>0)$,
il est l'image d'un $x_\mu\in A_\mu$ pour un certain indice $\mu$, et comme
l'image réciproque dans $A_\mu$ de
$\mathfrak m'^{n}=\mathfrak m_\mu^n A'$ est $\mathfrak m_\mu^n$
(\hyperref[0.6.6.1-fr]{$0_{\mathrm I}$, 6.6.1}), $x_\mu$ appartient à tous
les $\mathfrak m_\mu^n$, donc $x_\mu=0$ par hypothèse, et par conséquent
$x'=0$. Notons $\widehat A'$ le complété de $A'$ pour la topologie
$\mathfrak m'$-adique~; ce qui précède montre que l'on a
$A'\subset\widehat A'$. Nous allons montrer que $\widehat A'$ est
\emph{noethérien} et $A_\lambda$-\emph{plat} pour tout $\lambda$~; il en
\oldpage[0\textsubscript{III}]{22}
résultera que $\widehat A'$ est $A'$-plat
(\hyperref[0.6.2.3-fr]{$0_{\mathrm I}$, 6.2.3}), et comme
$\mathfrak m'\widehat A'\ne\widehat A'$, $\widehat A'$ est un $A'$-module
fidèlement plat (\hyperref[0.6.6.2-fr]{$0_{\mathrm I}$, 6.6.2}), d'où on
conclura finalement que $A'$ est \emph{noethérien}
(\hyperref[0.6.5.2-fr]{$0_{\mathrm I}$, 6.5.2}), ce qui achèvera la
démonstration du lemme.

On a $\widehat A'=\varprojlim_n A'/\mathfrak m'^{n}$~; en raison de ce que
$A'$ est $A_\lambda$-plat, on a
\[
  \mathfrak m'^{n}/\mathfrak m'^{n+1}
  =(\mathfrak m_\lambda^n/\mathfrak m_\lambda^{n+1})
    \otimes_{A_\lambda}A'
  =(\mathfrak m_\lambda^n/\mathfrak m_\lambda^{n+1})
    \otimes_{K_\lambda}(K_\lambda\otimes_{A_\lambda}A')
  =(\mathfrak m_\lambda^n/\mathfrak m_\lambda^{n+1})
    \otimes_{K_\lambda}K~;
\]
comme $\mathfrak m_\lambda^n/\mathfrak m_\lambda^{n+1}$ est un
$K_\lambda$-espace vectoriel de dimension finie,
$\mathfrak m'^{n}/\mathfrak m'^{n+1}$ est un $K$-espace vectoriel de
dimension finie pour tout $n\geq0$. Il résulte donc de
(\hyperref[0.7.2.12-fr]{$0_{\mathrm I}$, 7.2.12}) et
(\hyperref[0.7.2.8-fr]{$0_{\mathrm I}$, 7.2.8}) que $\widehat A'$ est
noethérien. On sait en outre que l'idéal maximal de $\widehat A'$ est
$\mathfrak m'\widehat A'$ et que
$\widehat A'/\mathfrak m'^{n}\widehat A'$ est isomorphe à
$A'/\mathfrak m'^{n}$~; comme
$A'/\mathfrak m'^{n}=(A_\lambda/\mathfrak m_\lambda^n)
\otimes_{A_\lambda}A'$, $A'/\mathfrak m'^{n}$ est un
$(A_\lambda/\mathfrak m_\lambda^n)$-module plat
(\hyperref[0.6.2.1-fr]{$0_{\mathrm I}$, 6.2.1})~; le critère
(\hyperref[0.10.2.2-fr]{10.2.2}) est donc applicable à la
$A_\lambda$-algèbre noethérienne $\widehat A'$, et montre que $\widehat A'$
est $A_\lambda$-plat. C.Q.F.D.

\begin{env}[10.3.1.4]
\label{0.10.3.1.4-fr}
Abordons maintenant le cas général. Il existe un ordinal $\gamma$ et pour
tout ordinal $\lambda\leq\gamma$ un sous-corps $k_\lambda$ de $K$ contenant
$k$, tels que~: $1^\circ$ Pour tout $\lambda<\gamma$, $k_{\lambda+1}$ soit
une extension de $k_\lambda$ engendrée par un seul élément~; $2^\circ$ Pour
tout ordinal $\mu$ sans prédécesseur,
$k_\mu=\bigcup_{\lambda<\mu}k_\lambda$~; $3^\circ$ $K=k_\gamma$. Il suffit en
effet de considérer une bijection $\xi\mapsto t_\xi$ de l'ensemble des
ordinaux $\xi\leq\beta$ (pour un $\beta$ convenable) sur $K$, de définir
$k_\lambda$ par induction transfinie (pour $\lambda\leq\beta$) comme la
réunion des $k_\mu$ pour $\mu<\lambda$ si $\lambda$ n'a pas de prédécesseur,
et, si $\lambda=\nu+1$, comme $k_\nu(t_\xi)$, où $\xi$ est le plus petit
ordinal tel que $t_\xi\notin k_\nu$~; $\gamma$ est alors par définition le plus
petit ordinal $\leq\beta$ tel que $k_\gamma=K$.

Cela étant, nous allons définir, par récurrence transfinie, une famille
d'anneaux locaux noethériens $A_\lambda$ pour $\lambda\leq\gamma$, et des
homomorphismes locaux $f_{\mu\lambda}:A_\lambda\to A_\mu$ pour
$\lambda\leq\mu$, vérifiant les conditions suivantes~:
\begin{enumerate}
  \item[(i)] $(A_\lambda,f_{\mu\lambda})$ est un système inductif et $A_0=A$.
  \item[(ii)] Pour tout $\lambda$, on a un $k$-isomorphisme
    $A_\lambda/\mathfrak J A_\lambda\xrightarrow{\sim}k_\lambda$.
  \item[(iii)] Pour $\lambda\leq\mu$, $A_\mu$ est un $A_\lambda$-module plat.
\end{enumerate}

Supposons donc les $A_\lambda$ et les $f_{\mu\lambda}$ définis pour
$\lambda<\mu<\xi$, et supposons en premier lieu que $\xi=\zeta+1$, de sorte
que $k_\xi=k_\zeta(t)$. Si $t$ est transcendant sur $k_\zeta$, on définit
$A_\xi$ suivant le procédé de (\hyperref[0.10.3.1.1-fr]{10.3.1.1}) comme
égal à $(A_\zeta[t])_{\mathfrak J A_\zeta[t]}$~;
$f_{\xi\zeta}$ est l'application canonique, et pour $\lambda<\zeta$, on prend
$f_{\xi\lambda}=f_{\xi\zeta}\circ f_{\zeta\lambda}$~; la vérification des
conditions (i) à (iii) est alors immédiate, vu ce qui a été démontré en
(\hyperref[0.10.3.1.1-fr]{10.3.1.1}). Supposons ensuite $t$ algébrique, et
soient $h$ son polynôme minimal dans $k_\zeta[T]$, $H$ un polynôme unitaire de
$A_\zeta[T]$ dont l'image dans $k_\zeta[T]$ est $h$~; on prend alors $A_\xi$
égal à $A_\zeta[T]/(H)$, les $f_{\xi\lambda}$ se définissant comme
précédemment~; la vérification des conditions (i) à (iii) résulte alors de
ce qui a été vu en (\hyperref[0.10.3.1.2-fr]{10.3.1.2}).

Supposons maintenant que $\xi$ n'ait pas de prédécesseur~; on prend alors
pour $A_\xi$ la limite inductive du système inductif d'anneaux locaux
$(A_\lambda,f_{\mu\lambda})$ pour $\lambda<\xi$~; $f_{\xi\lambda}$ est définie
comme l'application canonique pour $\lambda<\xi$. Le fait que $A_\xi$ soit
local noethérien, que les $f_{\xi\lambda}$ soient des homomorphismes locaux,
et les conditions (i) à (iii) pour $\lambda\leq\xi$
\oldpage[0\textsubscript{III}]{23}
résultent alors de l'hypothèse de récurrence et du lemme
(\hyperref[0.10.3.1.3-fr]{10.3.1.3}). Cette construction faite, il est clair
que l'anneau $B=A_\gamma$ vérifie l'énoncé de
(\hyperref[0.10.3.1-fr]{10.3.1}).
\end{env}

On notera qu'en vertu de (\hyperref[0.10.2.1-fr]{10.2.1, $c)$}), on a un
isomorphisme canonique
\[
\label{0.10.3.1.5-fr}
  \operatorname{gr}(A)\otimes_k K\xrightarrow{\sim}\operatorname{gr}(B).
  \tag{10.3.1.5}
\]

D'autre part, on peut remplacer $B$ par son complété $\mathfrak J B$-adique
$\widehat B$ sans changer les conclusions de
(\hyperref[0.10.3.1-fr]{10.3.1}), puisque $\widehat B$ est un $B$-module plat
(\hyperref[0.7.3.3-fr]{$0_{\mathrm I}$, 7.3.3}), donc un $A$-module plat
(\hyperref[0.6.2.1-fr]{$0_{\mathrm I}$, 6.2.1}).

On a en outre démontré le

\begin{corollary}[10.3.2]
\label{0.10.3.2-fr}
Si $K$ est une extension de degré fini, on peut supposer que $B$ est une
$A$-algèbre finie.
\end{corollary}

\section{Compléments d'algèbre homologique}
\label{section:0.11-fr}

\subsection{Rappels sur les suites spectrales}
\label{subsection:0.11.1-fr}

\begin{env}[11.1.1]
\label{0.11.1.1-fr}
Nous utiliserons dans la suite une notion de suite spectrale plus générale que
celle définie dans (T, 2.4)~; gardant les notations de (T, 2.4), nous
appellerons \emph{suite spectrale} dans une catégorie abélienne $\mathcal C$ un
système E formé des éléments suivants~:
\begin{enumerate}
  \item[(a)] Une famille $(E_r^{pq})$ d'objets de $\mathcal C$ définis pour
    $p\in\mathbb Z$, $q\in\mathbb Z$ et $r\geq2$.
  \item[(b)] Une famille de morphismes
    $d_r^{pq}:E_r^{pq}\to E_r^{p+r,q-r+1}$ tels que
    $d_r^{p+r,q-r+1}d_r^{pq}=0$. On pose
    $Z_{r+1}(E_r^{pq})=\operatorname{Ker}(d_r^{pq})$,
    $B_{r+1}(E_r^{pq})=\operatorname{Im}(d_r^{p-r,q+r-1})$, de sorte que
    \[
      B_{r+1}(E_r^{pq})\subset Z_{r+1}(E_r^{pq})\subset E_r^{pq}.
    \]
  \item[(c)] Une famille d'isomorphismes
    $\alpha_r^{pq}:Z_{r+1}(E_r^{pq})/B_{r+1}(E_r^{pq})
    \xrightarrow{\sim}E_{r+1}^{pq}$.
\end{enumerate}

On définit alors pour $k\geq r+1$, par récurrence sur $k$, les sous-objets
$B_k(E_r^{pq})$ et $Z_k(E_r^{pq})$ comme images réciproques, par le morphisme
canonique $E_r^{pq}\to E_r^{pq}/B_{r+1}(E_r^{pq})$, des sous-objets de ce
quotient identifié par $\alpha_r^{pq}$ aux sous-objets $B_k(E_{r+1}^{pq})$ et
$Z_k(E_{r+1}^{pq})$ respectivement. Il est clair que l'on a alors, à un
isomorphisme près
\[
\label{0.11.1.1.1-fr}
  Z_k(E_r^{pq})/B_k(E_r^{pq})=E_k^{pq}
  \qquad\text{pour }k\geq r+1,
  \tag{11.1.1.1}
\]
et, si on pose encore $B_r(E_r^{pq})=0$ et $Z_r(E_r^{pq})=E_r^{pq}$, on a
les relations d'inclusion
\[
\label{0.11.1.1.2-fr}
  0=B_r(E_r^{pq})\subset B_{r+1}(E_r^{pq})\subset B_{r+2}(E_r^{pq})
  \subset\cdots\subset Z_{r+2}(E_r^{pq})\subset Z_{r+1}(E_r^{pq})
  \subset Z_r(E_r^{pq})=E_r^{pq}.
  \tag{11.1.1.2}
\]

Les autres éléments de la donnée de E sont alors~:
\begin{enumerate}
  \item[(d)] Deux sous-objets $B_\infty(E_2^{pq})$ et
    $Z_\infty(E_2^{pq})$ de $E_2^{pq}$ tels que l'on ait
    $B_\infty(E_2^{pq})\subset Z_\infty(E_2^{pq})$ et, pour tout $k\geq2$,
    \[
      B_k(E_2^{pq})\subset B_\infty(E_2^{pq})
      \qquad\text{et}\qquad
      Z_\infty(E_2^{pq})\subset Z_k(E_2^{pq}).
    \]
    On pose
    \[
    \label{0.11.1.1.3-fr}
      E_\infty^{pq}=Z_\infty(E_2^{pq})/B_\infty(E_2^{pq}).
      \tag{11.1.1.3}
    \]
  \item[(e)]
    \oldpage[0\textsubscript{III}]{24}
    Une famille $(E^n)$ d'objets de $\mathcal C$, dont chacun est muni d'une
    \emph{filtration décroissante} $(F^p(E^n))_{p\in\mathbb Z}$. On désigne
    comme d'ordinaire par $\operatorname{gr}(E^n)$ l'objet gradué associé à
    l'objet filtré $E^n$, somme directe des
    $\operatorname{gr}_p(E^n)=F^p(E^n)/F^{p+1}(E^n)$.
  \item[(f)] Pour tout couple $(p,q)\in\mathbb Z\times\mathbb Z$, un
    isomorphisme
    $\beta^{pq}:E_\infty^{pq}\xrightarrow{\sim}
    \operatorname{gr}_p(E^{p+q})$.
\end{enumerate}

La famille $(E^n)$, sans les filtrations, est appelée \emph{l'aboutissement}
de la suite spectrale E.

Supposons que la catégorie $\mathcal C$ admette des sommes directes infinies,
ou que pour tout $r\geq2$ et tout $n\in\mathbb Z$, les couples $(p,q)$ tels
que $p+q=n$ et $E_r^{pq}\ne0$ soient en nombre fini (il suffit qu'il en soit
ainsi pour $r=2$). Alors les
\[
  E_r^{(n)}=\sum_{p+q=n}E_r^{pq}
\]
sont définis, et en désignant par $d_r^{(n)}$ le morphisme
$E_r^{(n)}\to E_r^{(n+1)}$ dont la restriction à $E_r^{pq}$ est $d_r^{pq}$
pour chaque couple $(p,q)$ tel que $p+q=n$, on a
$d_r^{(n+1)}\circ d_r^{(n)}=0$, autrement dit $(E_r^{(n)})_{n\in\mathbb Z}$
est un \emph{complexe $E_r^{(\bullet)}$} dans $\mathcal C$, à opérateur de
dérivation de degré $+1$, et il résulte de c) que
$\mathrm H^n(E_r^{(\bullet)})$ est isomorphe à $E_{r+1}^{(n)}$ pour tout
$r\geq2$.
\end{env}

\begin{env}[11.1.2]
\label{0.11.1.2-fr}
Un \emph{morphisme $u:E\to E'$} d'une suite spectrale E dans une suite
spectrale $E'=(E_r^{\prime pq},E^{\prime n})$ consiste en des systèmes de
morphismes $u_r^{pq}:E_r^{pq}\to E_r^{\prime pq}$,
$u^n:E^n\to E^{\prime n}$, les $u^n$ étant compatibles avec les filtrations
de $E^n$ et $E^{\prime n}$, les diagrammes
\[
  \xymatrix{
    E_r^{pq}\ar[r]^-{d_r^{pq}}\ar[d]_{u_r^{pq}} &
    E_r^{p+r,q-r+1}\ar[d]^{u_r^{p+r,q-r+1}}\\
    E_r^{\prime pq}\ar[r]_{d_r^{\prime pq}} &
    E_r^{\prime p+r,q-r+1}
  }
\]
étant commutatifs~; en outre, par passage aux quotients, $u_r^{pq}$ donne un
morphisme
\[
  \overline u_r^{pq}:
  Z_{r+1}(E_r^{pq})/B_{r+1}(E_r^{pq})\longrightarrow
  Z_{r+1}(E_r^{\prime pq})/B_{r+1}(E_r^{\prime pq})
\]
et on doit avoir
$\alpha_r^{\prime pq}\circ\overline u_r^{pq}
=u_{r+1}^{pq}\circ\alpha_r^{pq}$~; enfin, on doit avoir
$u_2^{pq}(B_\infty(E_2^{pq}))\subset B_\infty(E_2^{\prime pq})$,
$u_2^{pq}(Z_\infty(E_2^{pq}))\subset Z_\infty(E_2^{\prime pq})$~; par
passage aux quotients, $u_2^{pq}$ donne alors un morphisme
$u_\infty^{pq}:E_\infty^{pq}\to E_\infty^{\prime pq}$, et le diagramme
\[
  \xymatrix{
    E_\infty^{pq}\ar[r]^-{u_\infty^{pq}}\ar[d]_{\beta^{pq}} &
    E_\infty^{\prime pq}\ar[d]^{\beta^{\prime pq}}\\
    \operatorname{gr}_p(E^{p+q})
      \ar[r]_{\operatorname{gr}_p(u^{p+q})} &
    \operatorname{gr}_p(E^{\prime p+q})
  }
\]
doit être commutatif.

Les définitions précédentes montrent, par récurrence sur $r$, que si les
$u_2^{pq}$ sont des \emph{isomorphismes}, il en est de même des $u_r^{pq}$
pour $r\geq2$~; \emph{si l'on sait en outre que}
$u_2^{pq}(B_\infty(E_2^{pq}))=B_\infty(E_2^{\prime pq})$ et
$u_2^{pq}(Z_\infty(E_2^{pq}))=Z_\infty(E_2^{\prime pq})$ \emph{et que les
$u^n$ sont des isomorphismes}, alors on peut conclure que $u$ est un
\emph{isomorphisme}.
\end{env}

\begin{env}[11.1.3]
\label{0.11.1.3-fr}
\oldpage[0\textsubscript{III}]{25}
Rappelons que si $(F^p(X))_{p\in\mathbb Z}$ est une filtration (décroissante)
d'un objet $X\in\mathcal C$, on dit que cette filtration est \emph{séparée}
si $\inf(F^p(X))=0$, \emph{discrète} s'il existe un $p$ tel que $F^p(X)=0$,
\emph{exhaustive} (ou \emph{co-séparée}) si $\sup(F^p(X))=X$,
\emph{co-discrète} s'il existe $p$ tel que $F^p(X)=X$.

Nous dirons qu'une suite spectrale $E=(E_r^{pq},E^n)$ est \emph{faiblement
convergente} si l'on a
\[
  B_\infty(E_2^{pq})=\sup_k(B_k(E_2^{pq})),\qquad
  Z_\infty(E_2^{pq})=\inf_k(Z_k(E_2^{pq}))
\]
(autrement dit les objets $B_\infty(E_2^{pq})$ et $Z_\infty(E_2^{pq})$
sont déterminés par les données a) à c) de la suite spectrale E). Nous dirons
que la suite spectrale E est \emph{régulière} si elle est faiblement
convergente et si en outre~:
\begin{enumerate}
  \item[$1^\circ$] Pour tout couple $(p,q)$, la suite décroissante
    $(Z_k(E_2^{pq}))_{k\geq2}$ est \emph{stationnaire}~; l'hypothèse que E
    est faiblement convergente entraîne alors
    $Z_\infty(E_2^{pq})=Z_k(E_2^{pq})$ pour $k$ assez grand (dépendant de
    $p$ et $q$).
  \item[$2^\circ$] Pour tout $n$, la filtration
    $(F^p(E^n))_{p\in\mathbb Z}$ de $E^n$ est \emph{discrète} et
    \emph{exhaustive}.
\end{enumerate}

On dit que la suite spectrale E est \emph{co-régulière} si elle est
faiblement convergente et si en outre~:
\begin{enumerate}
  \item[$3^\circ$] Pour tout couple $(p,q)$, la suite croissante
    $(B_k(E_2^{pq}))_{k\geq2}$ est \emph{stationnaire}, ce qui entraîne
    $B_\infty(E_2^{pq})=B_k(E_2^{pq})$, et par suite
    $E_\infty^{pq}=\inf E_k^{pq}$.
  \item[$4^\circ$] Pour tout $n$, la filtration de $E^n$ est
    \emph{co-discrète}.
\end{enumerate}

Enfin, on dit que E est \emph{birégulière} si elle est à la fois régulière
et co-régulière, autrement dit si on a les conditions suivantes~:
\begin{enumerate}
  \item[(a)] Pour tout couple $(p,q)$, les suites
    $(B_k(E_2^{pq}))_{k\geq2}$ et $(Z_k(E_2^{pq}))_{k\geq2}$ sont
    \emph{stationnaires} et l'on a
    $B_\infty(E_2^{pq})=B_k(E_2^{pq})$ et
    $Z_\infty(E_2^{pq})=Z_k(E_2^{pq})$ pour $k$ assez grand (ce qui entraîne
    $E_\infty^{pq}=E_k^{pq}$).
  \item[(b)] Pour tout $n$, la filtration $(F^p(E^n))_{p\in\mathbb Z}$ est
    \emph{discrète} et \emph{co-discrète} (ce qu'on exprime aussi en disant
    qu'elle est \emph{finie}).
\end{enumerate}

Les suites spectrales définies dans (T, 2.4) sont donc les suites spectrales
birégulières.
\end{env}

\begin{env}[11.1.4]
\label{0.11.1.4-fr}
Supposons que dans la catégorie $\mathcal C$, les limites inductives filtrantes
existent et que le foncteur $\varinjlim$ soit \emph{exact} (ce qui équivaut à
dire que l'axiome AB 5) de (T, 1.5) est vérifié (cf. T, 1.8)). La condition
que la filtration $(F^p(X))_{p\in\mathbb Z}$ d'un objet $X\in\mathcal C$ est
exhaustive s'écrit alors aussi
$\varinjlim_{p\to-\infty}F^p(X)=X$. Si une suite spectrale E est faiblement
convergente, on a
$B_\infty(E_2^{pq})=\varinjlim_{k\to\infty}B_k(E_2^{pq})$~; si en outre
$u:E\to E'$ est un morphisme de E dans une suite spectrale $E'$ de
$\mathcal C$ faiblement convergente, on a
$u_2^{pq}(B_\infty(E_2^{pq}))=B_\infty(E_2^{\prime pq})$ en raison de
l'exactitude de $\varinjlim$. De plus~:
\end{env}

\begin{proposition}[11.1.5]
\label{0.11.1.5-fr}
Soient $\mathcal C$ une catégorie abélienne dans laquelle les limites
inductives filtrantes sont exactes, E, $E'$ deux suites spectrales régulières
de $\mathcal C$, $u:E\to E'$ un morphisme de suites spectrales. Si les
$u_2^{pq}$ sont des isomorphismes, il en est de même de $u$.
\end{proposition}

\begin{proof}
Nous savons déjà (\hyperref[0.11.1.2-fr]{11.1.2}) que les $u_r^{pq}$ sont des
isomorphismes et que
\[
  u_2^{pq}(B_\infty(E_2^{pq}))=B_\infty(E_2^{\prime pq})~;
\]
\oldpage[0\textsubscript{III}]{26}
l'hypothèse que E et $E'$ sont régulières entraîne aussi
$u_2^{pq}(Z_\infty(E_2^{pq}))=Z_\infty(E_2^{\prime pq})$, et comme
$u_2^{pq}$ est un isomorphisme, il en est de même de $u_\infty^{pq}$~; on en
conclut donc que $\operatorname{gr}_p(u^{p+q})$ est aussi un isomorphisme.
Mais comme les filtrations des $E^n$ et des $E^{\prime n}$ sont discrètes et
exhaustives, cela entraîne que les $u^n$ sont aussi des isomorphismes
(Bourbaki, \emph{Alg. comm.}, chap. III, \S~2, n\textsuperscript{o} 8,
th. 1).
\end{proof}

\begin{env}[11.1.6]
\label{0.11.1.6-fr}
Il résulte de (\hyperref[0.11.1.1.2-fr]{11.1.1.2}) et de la définition
(\hyperref[0.11.1.1.3-fr]{11.1.1.3}) que si, dans une suite spectrale E, on a
$E_r^{pq}=0$, on a $E_k^{pq}=0$ pour $k\geq r$ et $E_\infty^{pq}=0$. On dit
qu'une suite spectrale est \emph{dégénérée} s'il existe un entier $r\geq2$ et,
pour tout entier $n\in\mathbb Z$, un entier $q(n)$ tel que
$E_r^{n-q,q}=0$ pour tout $q\ne q(n)$. On déduit d'abord de la remarque
précédente que l'on a aussi $E_k^{n-q,q}=0$ pour $k\geq r$ (y compris
$k=\infty$) et $q\ne q(n)$. En outre, la définition de $E_{r+1}^{pq}$ montre
que l'on a
$E_{r+1}^{n-q(n),q(n)}=E_r^{n-q(n),q(n)}$~; si E est faiblement convergente,
on a donc aussi
$E_\infty^{n-q(n),q(n)}=E_r^{n-q(n),q(n)}$~; autrement dit, pour tout
$n\in\mathbb Z$, $\operatorname{gr}_p(E^n)=0$ pour $p\ne q(n)$ et
$\operatorname{gr}_{q(n)}(E^n)=E_r^{n-q(n),q(n)}$. Si en outre la filtration
de $E^n$ est discrète et exhaustive, la suite E est régulière et on a
$E^n=E_r^{n-q(n),q(n)}$ à un isomorphisme près.
\end{env}

\begin{env}[11.1.7]
\label{0.11.1.7-fr}
Supposons que dans la catégorie $\mathcal C$ les limites inductives filtrantes
existent et soient exactes, et soit $(E_\lambda,u_{\mu\lambda})$ un système
inductif (suivant un ensemble d'indices filtrant) de suites spectrales de
$\mathcal C$. Alors la \emph{limite inductive} de ce système inductif existe
dans la catégorie additive des suites spectrales d'objets de $\mathcal C$~:
il suffit pour le voir de définir $E_r^{pq}$, $d_r^{pq}$, $\alpha_r^{pq}$,
$B_\infty(E_2^{pq})$, $Z_\infty(E_2^{pq})$, $E^n$, $F^p(E^n)$ et
$\beta^{pq}$ comme limites inductives respectives de $E_{r,\lambda}^{pq}$,
$d_{r,\lambda}^{pq}$, $\alpha_{r,\lambda}^{pq}$,
$B_\infty(E_{2,\lambda}^{pq})$, $Z_\infty(E_{2,\lambda}^{pq})$,
$E_\lambda^n$, $F^p(E_\lambda^n)$ et $\beta_\lambda^{pq}$~; la vérification
des conditions de (\hyperref[0.11.1.1-fr]{11.1.1}) résulte de l'exactitude du
foncteur $\varinjlim$ dans $\mathcal C$.
\end{env}

\begin{remark}[11.1.8]
\label{0.11.1.8-fr}
Supposons que la catégorie $\mathcal C$ soit la catégorie des $A$-modules sur
un anneau \emph{noethérien} $A$ (resp. sur un anneau $A$). Alors, les
définitions de (\hyperref[0.11.1.1-fr]{11.1.1}) montrent que si pour un $r$
donné, les $E_r^{pq}$ sont des $A$-modules \emph{de type fini} (resp. de
\emph{longueur finie}), il en est de même de tous les modules $E_s^{pq}$ pour
$s\geq r$, ainsi que des $E_\infty^{pq}$. Si en outre la filtration de
l'aboutissement $(E^n)$ est discrète et co-discrète pour tout $n$, on en
conclut que chacun des $E^n$ est aussi un $A$-module \emph{de type fini}
(resp. de \emph{longueur finie}).
\end{remark}

\begin{env}[11.1.9]
\label{0.11.1.9-fr}
Nous aurons à considérer des conditions assurant qu'une suite spectrale E est
birégulière de façon «~uniforme~» en $p+q=n$. On utilisera alors le lemme
suivant~:
\end{env}

\begin{lemma}[11.1.10]
\label{0.11.1.10-fr}
Soit $(E_r^{pq})$ une famille d'objets de $\mathcal C$ liés par les données
a), b), c) de (\hyperref[0.11.1.1-fr]{11.1.1}). Pour un entier fixé $n$, les
propriétés suivantes sont équivalentes~:
\begin{enumerate}
  \item[(a)] Il existe un entier $r(n)$ tel que pour $r\geq r(n)$,
    $p+q=n$ ou $p+q=n-1$, les morphismes $d_r^{pq}$ soient tous nuls.
  \item[(b)] Il existe un entier $r(n)$ tel que pour $p+q=n$ ou $p+q=n+1$,
    on ait $B_r(E_2^{pq})=B_s(E_2^{pq})$ pour $s\geq r\geq r(n)$.
  \item[(c)] Il existe un entier $r(n)$ tel que pour $p+q=n$ ou $p+q=n-1$,
    on ait $Z_r(E_2^{pq})=Z_s(E_2^{pq})$ pour $s\geq r\geq r(n)$.
  \item[(d)] Il existe un entier $r(n)$ tel que pour $p+q=n$, on ait
    $B_r(E_2^{pq})=B_s(E_2^{pq})$ et $Z_r(E_2^{pq})=Z_s(E_2^{pq})$ pour
    $s\geq r\geq r(n)$.
\end{enumerate}
\end{lemma}

\begin{proof}
\oldpage[0\textsubscript{III}]{27}
En effet, d'après les conditions a), b), c) de
(\hyperref[0.11.1.1-fr]{11.1.1}), dire que
$Z_{r+1}(E_2^{pq})=Z_r(E_2^{pq})$ équivaut à dire que $d_r^{pq}=0$, et dire
que
$B_r(E_2^{p+r,q-r+1})=B_{r+1}(E_2^{p+r,q-r+1})$ équivaut aussi à dire que
$d_r^{pq}=0$~; le lemme résulte aussitôt de cette remarque.
\end{proof}
\subsection{La suite spectrale d'un complexe filtré}
\label{subsection:0.11.2-fr}

\begin{env}[11.2.1]
\label{0.11.2.1-fr}
Étant donnée une catégorie abélienne $\mathcal C$, nous conviendrons de
désigner par des notations telles que $K^\bullet$ les \emph{complexes}
$(K^i)_{i\in\mathbb Z}$ d'objets de $\mathcal C$ dans lesquels l'opérateur de
dérivation est de degré $+1$, par des notations telles que $K_\bullet$ les
complexes $(K_i)_{i\in\mathbb Z}$ d'objets de $\mathcal C$ dans lesquels
l'opérateur de dérivation est de degré $-1$. À tout complexe
$K^\bullet=(K^i)$ dont l'opérateur de dérivation $d$ est de degré $+1$, on
peut associer un complexe $K'_\bullet=(K'_i)$ en posant $K'_i=K^{-i}$,
l'opérateur de dérivation $K'_i\to K'_{i-1}$ étant l'opérateur
$d:K^{-i}\to K^{-i+1}$~; et \emph{vice versa}, ce qui permettra suivant les
circonstances de considérer l'un ou l'autre type de complexes et de traduire
tout résultat pour l'un des types en résultats pour l'autre. Nous désignerons
de même par des notations telles que $K^{\bullet\bullet}=(K^{ij})$ (resp.
$K_{\bullet\bullet}=(K_{ij})$) des \emph{bicomplexes} d'objets de
$\mathcal C$ dans lesquels les \emph{deux} opérateurs de dérivation sont de
degré $+1$ (resp. $-1$)~; on passe encore de l'un à l'autre type en changeant
les signes des indices, et on a des notations et remarques analogues pour des
multicomplexes quelconques. La notation $K^\bullet$ ou $K_\bullet$ sera aussi
utilisée pour des \emph{objets gradués} de $\mathcal C$, de type $\mathbb Z$,
qui ne sont pas nécessairement des complexes (ou que l'on peut considérer
comme tels pour des opérateurs de dérivation nuls)~; par exemple, nous
écrirons
$\mathrm H^\bullet(K^\bullet)=(\mathrm H^i(K^\bullet))_{i\in\mathbb Z}$ la
\emph{cohomologie} d'un complexe $K^\bullet$ dont l'opérateur de dérivation
est de degré $+1$, par
$\mathrm H_\bullet(K_\bullet)=(\mathrm H_i(K_\bullet))_{i\in\mathbb Z}$
l'\emph{homologie} d'un complexe $K_\bullet$ dont l'opérateur de dérivation
est de degré $-1$~; quand on passe de $K^\bullet$ à $K'_\bullet$ par
l'opération décrite ci-dessus, on a
$\mathrm H_i(K'_\bullet)=\mathrm H^{-i}(K^\bullet)$.

Rappelons à ce propos que pour un complexe $K^\bullet$ (resp. $K_\bullet$),
nous écrirons en général
$Z^i(K^\bullet)=\operatorname{Ker}(K^i\to K^{i+1})$ («~objet des cocycles~»)
et $B^i(K^\bullet)=\operatorname{Im}(K^{i-1}\to K^i)$ («~objet des
cobords~») (resp.
$Z_i(K_\bullet)=\operatorname{Ker}(K_i\to K_{i-1})$ («~objet des cycles~»)
et $B_i(K_\bullet)=\operatorname{Im}(K_{i+1}\to K_i)$ («~objet des
bords~»)) de sorte que
$\mathrm H^i(K^\bullet)=Z^i(K^\bullet)/B^i(K^\bullet)$ (resp.
$\mathrm H_i(K_\bullet)=Z_i(K_\bullet)/B_i(K_\bullet)$).

Si $K^\bullet=(K^i)$ (resp. $K_\bullet=(K_i)$) est un complexe dans
$\mathcal C$, et $T:\mathcal C\to\mathcal C'$ un foncteur de $\mathcal C$
dans une catégorie abélienne $\mathcal C'$, nous désignerons par
$T(K^\bullet)$ (resp. $T(K_\bullet)$) le complexe $(T(K^i))$ (resp.
$(T(K_i))$) dans $\mathcal C'$.

Nous ne revenons pas sur la définition des $\partial$-foncteurs (T, 2.1),
sauf pour signaler que nous dirons aussi $\partial$-foncteur au lieu de
$\partial^*$-foncteur lorsque le morphisme $\partial$ diminue le degré d'une
unité, le contexte devant chaque fois préciser ce point s'il peut y avoir
doute.

Enfin, nous dirons qu'un \emph{objet gradué} $(A_i)_{i\in\mathbb Z}$ de
$\mathcal C$ est \emph{limité inférieurement} (resp. \emph{supérieurement})
s'il existe $i_0$ tel que $A_i=0$ pour $i<i_0$ (resp. $i>i_0$).
\end{env}

\begin{env}[11.2.2]
\label{0.11.2.2-fr}
Soit $K^\bullet$ un complexe de $\mathcal C$ dont l'opérateur de dérivation
$d$ est de degré $+1$, et supposons-le muni d'une \emph{filtration}
$F(K^\bullet)=(F^p(K^\bullet))_{p\in\mathbb Z}$ formé de sous-
\oldpage[0\textsubscript{III}]{28}
objets gradués de $K^\bullet$, autrement dit
$F^p(K^\bullet)=(K^i\cap F^p(K^\bullet))_{i\in\mathbb Z}$~; en outre, on
suppose que $d(F^p(K^\bullet))\subset F^p(K^\bullet)$ pour tout
$p\in\mathbb Z$. Rappelons alors rapidement comment on définit
\emph{fonctoriellement} une suite spectrale $E(K^\bullet)$ à partir de
$K^\bullet$ (M, XV, 4 et G, I, 4.3). Pour $r\geq2$, le morphisme canonique
$F^p(K^\bullet)/F^{p+r}(K^\bullet)\to
F^p(K^\bullet)/F^{p+1}(K^\bullet)$ définit un morphisme pour la cohomologie
\[
  \mathrm H^{p+q}(F^p(K^\bullet)/F^{p+r}(K^\bullet))
  \longrightarrow
  \mathrm H^{p+q}(F^p(K^\bullet)/F^{p+1}(K^\bullet)).
\]
On désigne par $Z_r^{pq}(K^\bullet)$ l'image de ce morphisme. De même, de la
suite exacte
\[
  0\to F^p(K^\bullet)/F^{p+1}(K^\bullet)
  \to F^{p-r+1}(K^\bullet)/F^{p+1}(K^\bullet)
  \to F^{p-r+1}(K^\bullet)/F^p(K^\bullet)\to0
\]
on déduit par la suite exacte de cohomologie un morphisme
\[
  \mathrm H^{p+q-1}(F^{p-r+1}(K^\bullet)/F^p(K^\bullet))
  \longrightarrow
  \mathrm H^{p+q}(F^p(K^\bullet)/F^{p+1}(K^\bullet))
\]
et on désigne par $B_r^{pq}(K^\bullet)$ l'image de ce morphisme~; on montre
que $B_r^{pq}(K^\bullet)\subset Z_r^{pq}(K^\bullet)$ et on prend
$E_r^{pq}(K^\bullet)=Z_r^{pq}(K^\bullet)/B_r^{pq}(K^\bullet)$~; nous ne
préciserons pas la définition des $d_r^{pq}$, ni des $\alpha_r^{pq}$.

On notera ici que tous les $Z_r^{pq}(K^\bullet)$ et
$B_r^{pq}(K^\bullet)$, pour $p$, $q$ fixés, sont des sous-objets du même
objet
$\mathrm H^{p+q}(F^p(K^\bullet)/F^{p+1}(K^\bullet))$, que l'on note
$Z_1^{pq}(K^\bullet)$~; on pose $B_1^{pq}(K^\bullet)=0$, de sorte que les
définitions précédentes pour $Z_r^{pq}(K^\bullet)$ et
$B_r^{pq}(K^\bullet)$ s'appliquent aussi pour $r=1$~; on pose encore
$E_1^{pq}(K^\bullet)=Z_1^{pq}(K^\bullet)$. On définit encore les $d_1^{pq}$
et $\alpha_1^{pq}$ de sorte que les conditions de
(\hyperref[0.11.1.1-fr]{11.1.1}) soient vérifiées pour $r=1$. On définit
d'autre part les sous-objets $Z_\infty^{pq}(K^\bullet)$, image du morphisme
\[
  \mathrm H^{p+q}(F^p(K^\bullet))\longrightarrow
  \mathrm H^{p+q}(F^p(K^\bullet)/F^{p+1}(K^\bullet))
  =E_1^{pq}(K^\bullet)
\]
et $B_\infty^{pq}(K^\bullet)$, image du morphisme
\[
  \mathrm H^{p+q-1}(K^\bullet/F^p(K^\bullet))\longrightarrow
  \mathrm H^{p+q}(F^p(K^\bullet)/F^{p+1}(K^\bullet))
  =E_1^{pq}(K^\bullet)
\]
déduit comme ci-dessus d'une suite exacte de cohomologie. On prend pour
$Z_\infty(E_2^{pq}(K^\bullet))$ et $B_\infty(E_2^{pq}(K^\bullet))$ les
images canoniques dans $E_2^{pq}(K^\bullet)$ de
$Z_\infty^{pq}(K^\bullet)$ et $B_\infty^{pq}(K^\bullet)$.

Enfin, on désigne par $F^p(\mathrm H^n(K^\bullet))$ l'image dans
$\mathrm H^n(K^\bullet)$ du morphisme
$\mathrm H^n(F^p(K^\bullet))\to\mathrm H^n(K^\bullet)$ provenant de
l'injection canonique $F^p(K^\bullet)\to K^\bullet$~; par la suite exacte de
cohomologie, c'est aussi le noyau du morphisme
$\mathrm H^n(K^\bullet)\to\mathrm H^n(K^\bullet/F^p(K^\bullet))$. On
définit ainsi une filtration sur
$E^n(K^\bullet)=\mathrm H^n(K^\bullet)$~; nous ne donnerons pas non plus ici
la définition des isomorphismes $\beta^{pq}$.
\end{env}

\begin{env}[11.2.3]
\label{0.11.2.3-fr}
Le caractère \emph{fonctoriel} de $E(K^\bullet)$ doit s'entendre de la façon
suivante~: étant donné deux complexes \emph{filtrés} $K^\bullet$,
$K^{\prime\bullet}$ de $\mathcal C$, et un morphisme de complexes
$u:K^\bullet\to K^{\prime\bullet}$ \emph{compatible avec les filtrations}, on
en déduit de façon évidente les morphismes $u_r^{pq}$ (pour $r\geq1$) et
$u^n$, et on montre que ces morphismes sont compatibles avec les $d_r^{pq}$,
$\alpha_r^{pq}$ et $\beta^{pq}$ au sens de
(\hyperref[0.11.1.2-fr]{11.1.2}), donc définissent bien un morphisme
$E(u):E(K^\bullet)\to E(K^{\prime\bullet})$ de suites spectrales. En outre,
on montre que si $u$ et $v$ sont des morphismes
$K^\bullet\to K^{\prime\bullet}$ du type précédent, \emph{homotopes d'ordre
$\leq k$}, alors $u_r^{pq}=v_r^{pq}$ pour $r>k$ et $u^n=v^n$ pour tout $n$
(M, XV, 3.1).
\end{env}

\begin{env}[11.2.4]
\label{0.11.2.4-fr}
\oldpage[0\textsubscript{III}]{29}
Supposons que dans $\mathcal C$ les limites inductives filtrantes soient
exactes. Alors, si la filtration $(F^p(K^\bullet))$ de $K^\bullet$ est
\emph{exhaustive}, il en est de même de la filtration
$(F^p(\mathrm H^n(K^\bullet)))$ pour tout $n$, car on a par hypothèse
$K^\bullet=\varinjlim_{p\to-\infty}F^p(K^\bullet)$ et l'hypothèse sur
$\mathcal C$ entraîne que la cohomologie commute aux limites inductives. En
outre, on a, pour la même raison,
$B_\infty(E_2^{pq}(K^\bullet))=\sup_k B_k(E_2^{pq}(K^\bullet))$. On dit que
la filtration $(F^p(K^\bullet))$ de $K^\bullet$ est \emph{régulière} si pour
tout $n$ il existe un entier $u(n)$ tel que
$\mathrm H^n(F^p(K^\bullet))=0$ pour $p>u(n)$. Il en est ainsi en particulier
lorsque la filtration de $K^\bullet$ est \emph{discrète}. Lorsque la
filtration de $K^\bullet$ est régulière et exhaustive, et que les limites
inductives filtrantes dans $\mathcal C$ sont exactes, on montre (M, XV, 4)
que la suite spectrale $E(K^\bullet)$ est \emph{régulière}.
\end{env}

\subsection{Les suites spectrales d'un bicomplexe}
\label{subsection:0.11.3-fr}

\begin{env}[11.3.1]
\label{0.11.3.1-fr}
En ce qui concerne les conventions relatives aux bicomplexes, nous suivons
celles de (T, 2.4) plutôt que celles de (M), les deux dérivations $d'$, $d''$
(de degré $+1$) d'un tel bicomplexe
$K^{\bullet\bullet}=(K^{ij})$ étant donc supposées \emph{permutables}.
Supposons vérifiées l'une des deux conditions suivantes~:
\begin{enumerate}
  \item[$1^\circ$] Les sommes directes infinies existent dans $\mathcal C$.
  \item[$2^\circ$] Pour tout $n\in\mathbb Z$, il n'y a qu'un nombre
    \emph{fini} de couples $(p,q)$ tels que $p+q=n$ et $K^{pq}\ne0$.
\end{enumerate}
Alors, le bicomplexe $K^{\bullet\bullet}$ définit un complexe (simple)
$(K^{\prime n})_{n\in\mathbb Z}$, avec
$K^{\prime n}=\sum_{i+j=n}K^{ij}$, l'opérateur de dérivation $d$ (de degré
$+1$) de ce complexe étant donné par
$dx=d'x+(-1)^i d''x$ pour $x\in K^{ij}$. Lorsque nous parlerons par la suite
du complexe (simple) \emph{défini par un bicomplexe $K^{\bullet\bullet}$}, il
sera toujours sous-entendu que l'une des conditions précédentes est satisfaite.
On adoptera des conventions analogues pour les multicomplexes.

On désigne par $K^{i,\bullet}$ (resp. $K^{\bullet,j}$) le complexe simple
$(K^{ij})_{j\in\mathbb Z}$ (resp. $(K^{ij})_{i\in\mathbb Z}$), par
$Z_{\mathrm{II}}^p(K^{i,\bullet})$, $B_{\mathrm{II}}^p(K^{i,\bullet})$,
$\mathrm H_{\mathrm{II}}^p(K^{i,\bullet})$ (resp.
$Z_{\mathrm I}^p(K^{\bullet,j})$, $B_{\mathrm I}^p(K^{\bullet,j})$,
$\mathrm H_{\mathrm I}^p(K^{\bullet,j})$) ses $p$-èmes objets des cocycles,
des cobords et de cohomologie respectivement~; la dérivation
$d':K^{i,\bullet}\to K^{i+1,\bullet}$ est un morphisme de complexes, qui
donne donc un opérateur dans les cocycles, les cobords et la cohomologie,
\[
\begin{aligned}
  d'&:Z_{\mathrm{II}}^p(K^{i,\bullet})
      \longrightarrow Z_{\mathrm{II}}^p(K^{i+1,\bullet}),\\
  d'&:B_{\mathrm{II}}^p(K^{i,\bullet})
      \longrightarrow B_{\mathrm{II}}^p(K^{i+1,\bullet}),\\
  d'&:\mathrm H_{\mathrm{II}}^p(K^{i,\bullet})
      \longrightarrow \mathrm H_{\mathrm{II}}^p(K^{i+1,\bullet}),
\end{aligned}
\]
et il est clair que pour ces opérateurs,
$(Z_{\mathrm{II}}^p(K^{i,\bullet}))_{i\in\mathbb Z}$,
$(B_{\mathrm{II}}^p(K^{i,\bullet}))_{i\in\mathbb Z}$ et
$(\mathrm H_{\mathrm{II}}^p(K^{i,\bullet}))_{i\in\mathbb Z}$ sont des
complexes~; nous désignerons le complexe
$(\mathrm H_{\mathrm{II}}^p(K^{i,\bullet}))_{i\in\mathbb Z}$ par
$\mathrm H_{\mathrm{II}}^p(K^{\bullet\bullet})$, ses $q$-èmes objets de
cocycles, de cobords et de cohomologie par
$Z_{\mathrm I}^q(\mathrm H_{\mathrm{II}}^p(K^{\bullet\bullet}))$,
$B_{\mathrm I}^q(\mathrm H_{\mathrm{II}}^p(K^{\bullet\bullet}))$ et
$\mathrm H_{\mathrm I}^q(\mathrm H_{\mathrm{II}}^p(K^{\bullet\bullet}))$.
On définit de même les complexes
$\mathrm H_{\mathrm I}^p(K^{\bullet\bullet})$ et leurs objets de cohomologie
$\mathrm H_{\mathrm{II}}^q(\mathrm H_{\mathrm I}^p(K^{\bullet\bullet}))$.
Rappelons d'autre part que $\mathrm H^n(K^{\bullet\bullet})$ désigne le
$n$-ème objet de cohomologie du complexe (simple) défini par
$K^{\bullet\bullet}$.
\end{env}

\begin{env}[11.3.2]
\label{0.11.3.2-fr}
Sur le complexe défini par un bicomplexe $K^{\bullet\bullet}$, on peut
considérer deux filtrations canoniques
$(F_{\mathrm I}^p(K^{\bullet\bullet}))$ et
$(F_{\mathrm{II}}^p(K^{\bullet\bullet}))$ données par
\[
\label{0.11.3.2.1-fr}
  F_{\mathrm I}^p(K^{\bullet\bullet})
  =\left(\sum_{\substack{i+j=n\\i\geq p}}K^{ij}\right)_{n\in\mathbb Z},
  \qquad
  F_{\mathrm{II}}^p(K^{\bullet\bullet})
  =\left(\sum_{\substack{i+j=n\\j\geq p}}K^{ij}\right)_{n\in\mathbb Z}.
  \tag{11.3.2.1}
\]
\oldpage[0\textsubscript{III}]{30}
qui, par définition, sont bien des sous-objets gradués du complexe (simple)
défini par $K^{\bullet\bullet}$, et font donc de ce complexe un complexe
filtré~; par ailleurs, il est clair que ces filtrations sont exhaustives et
séparées.

Il correspond à chacune de ces filtrations une suite spectrale
(\hyperref[0.11.2.2-fr]{11.2.2})~; nous désignerons par
$'E(K^{\bullet\bullet})$ et $''E(K^{\bullet\bullet})$ les suites spectrales
correspondant à $(F_{\mathrm I}^p(K^{\bullet\bullet}))$ et
$(F_{\mathrm{II}}^p(K^{\bullet\bullet}))$ respectivement, dites
\emph{suites spectrales du bicomplexe $K^{\bullet\bullet}$}, et ayant toutes
deux pour aboutissement la cohomologie
$(\mathrm H^n(K^{\bullet\bullet}))$. On montre en outre (M, XV, 6) que l'on a
\[
\label{0.11.3.2.2-fr}
  'E_2^{pq}(K^{\bullet\bullet})
  =\mathrm H_{\mathrm I}^p(\mathrm H_{\mathrm{II}}^q(K^{\bullet\bullet})),
  \qquad
  ''E_2^{pq}(K^{\bullet\bullet})
  =\mathrm H_{\mathrm{II}}^q(\mathrm H_{\mathrm I}^p(K^{\bullet\bullet})).
  \tag{11.3.2.2}
\]

Tout morphisme $u:K^{\bullet\bullet}\to K^{\prime\bullet\bullet}$ de
bicomplexes est \emph{ipso facto} compatible avec les filtrations de même
type de $K^{\bullet\bullet}$ et $K^{\prime\bullet\bullet}$, donc définit un
morphisme pour chacune des deux suites spectrales~; en outre, deux morphismes
\emph{homotopes} définissent une homotopie \emph{d'ordre $\leq1$} des
complexes (simples) filtrés correspondant, donc le \emph{même} morphisme pour
chacune des deux suites spectrales (M, XV, 6.1).
\end{env}

\begin{proposition}[11.3.3]
\label{0.11.3.3-fr}
Soit $K^{\bullet\bullet}=(K^{ij})$ un bicomplexe dans une catégorie abélienne
$\mathcal C$.
\begin{enumerate}
  \item[(i)] S'il existe $i_0$ et $j_0$ tels que $K^{ij}=0$ pour $i<i_0$ ou
    $j<j_0$ (resp. $i>i_0$ ou $j>j_0$), les deux suites spectrales
    $'E(K^{\bullet\bullet})$ et $''E(K^{\bullet\bullet})$ sont
    birégulières.
  \item[(ii)] S'il existe $i_0$ et $i_1$ tels que $K^{ij}=0$ pour $i<i_0$ ou
    $i>i_1$ (resp. s'il existe $j_0$ et $j_1$ tels que $K^{ij}=0$ pour
    $j<j_0$ ou $j>j_1$) les deux suites spectrales
    $'E(K^{\bullet\bullet})$ et $''E(K^{\bullet\bullet})$ sont
    birégulières.
\end{enumerate}
Supposons en outre que dans $\mathcal C$ les limites inductives filtrantes
existent et soient exactes. Alors~:
\begin{enumerate}
  \item[(iii)] S'il existe $i_0$ tel que $K^{ij}=0$ pour $i>i_0$ (resp. s'il
    existe $j_0$ tel que $K^{ij}=0$ pour $j<j_0$), la suite
    $'E(K^{\bullet\bullet})$ est régulière.
  \item[(iv)] S'il existe $i_0$ tel que $K^{ij}=0$ pour $i<i_0$ (resp. s'il
    existe $j_0$ tel que $K^{ij}=0$ pour $j>j_0$), la suite
    $''E(K^{\bullet\bullet})$ est régulière.
\end{enumerate}
\end{proposition}

\begin{proof}
La proposition résulte aussitôt des définitions
(\hyperref[0.11.1.3-fr]{11.1.3}) et de
(\hyperref[0.11.2.4-fr]{11.2.4}), ainsi que des observations suivantes
relatives à la filtration $F_{\mathrm I}$ (et des observations analogues qu'on
en déduit pour $F_{\mathrm{II}}$ en échangeant les rôles des deux indices
dans $K^{\bullet\bullet}$)~:
\begin{enumerate}
  \item[$1^\circ$] S'il existe $i_0$ tel que $K^{ij}=0$ pour $i>i_0$, la
    filtration $F_{\mathrm I}(K^{\bullet\bullet})$ est discrète.
  \item[$2^\circ$] S'il existe $i_0$ tel que $K^{ij}=0$ pour $i<i_0$, la
    filtration $F_{\mathrm I}(K^{\bullet\bullet})$ est co-discrète. On en
    déduit aussitôt qu'il en est de même de la filtration correspondante
    $F_{\mathrm I}(\mathrm H^n(K^{\bullet\bullet}))$ pour tout $n$~; en
    outre, la définition de $B_r^{pq}$ correspondant à la filtration
    $F_{\mathrm I}(K^{\bullet\bullet})$
    (\hyperref[0.11.2.2-fr]{11.2.2}) montre que pour tout couple $(p,q)$, la
    suite $(B_r^{pq})_{r\geq2}$ est stationnaire.
  \item[$3^\circ$] S'il existe $j_0$ tel que $K^{ij}=0$ pour $j<j_0$, on a
    \[
      F_{\mathrm I}^{p+r}(K^{\bullet\bullet})
      \cap\left(\sum_{i+j=n}K^{ij}\right)=0
    \]
    dès que $p+r+j_0>n$, donc
    $Z_r^{pq}=Z_\infty(E_2^{pq})$ pour $r>q-j_0+1$~; d'autre part,
    $\mathrm H^n(F_{\mathrm I}^p(K^{\bullet\bullet}))=0$ pour
    $p>n-j_0+1$.
  \item[$4^\circ$] S'il existe $j_0$ tel que $K^{ij}=0$ pour $j>j_0$, on a
    \[
      F_{\mathrm I}^{p-r+1}(K^{\bullet\bullet})
      \cap\left(\sum_{i+j=n}K^{ij}\right)=\sum_{i+j=n}K^{ij}.
    \]
    \oldpage[0\textsubscript{III}]{31}
    dès que $p-r+1+j_0<n$, donc
    $B_r^{pq}=B_\infty(E_2^{pq})$ pour $r<j_0-q+1$~; d'autre part,
    $\mathrm H^n(F_{\mathrm I}^p(K^{\bullet\bullet}))
    =\mathrm H^n(K^{\bullet\bullet})$ pour $p+j_0<n-1$.
\end{enumerate}
\end{proof}

\begin{env}[11.3.4]
\label{0.11.3.4-fr}
Supposons que le bicomplexe $K^{\bullet\bullet}=(K^{ij})$ soit tel que
$K^{ij}=0$ pour $i<0$ ou $j<0$. On sait qu'on peut alors définir pour tout
$p\in\mathbb Z$ un «~edge-homomorphisme~» canonique
\[
\label{0.11.3.4.1-fr}
  'E_2^{p0}(K^{\bullet\bullet})\longrightarrow
  \mathrm H^p(K^{\bullet\bullet})
  \tag{11.3.4.1}
\]
(M, XV, 6). Rappelons rapidement que cela est dû, d'une part à ce que l'on a
dans la suite spectrale $'E(K^{\bullet\bullet})$,
$Z_r^{p0}=Z_{\mathrm I}^p(Z_{\mathrm{II}}^0(K^{\bullet\bullet}))$ pour
$2\leq r\leq+\infty$, et de l'autre au fait que
$\mathrm H^p(F_{\mathrm I}^{p+1}(K^{\bullet\bullet}))=0$, si bien que
l'isomorphisme
$\beta^{p0}:'E_\infty^{p0}\xrightarrow{\sim}
\mathrm H^p(F_{\mathrm I}^p)/\mathrm H^p(F_{\mathrm I}^{p+1})$ donne un
homomorphisme
$'E_\infty^{p0}\to\mathrm H^p(F_{\mathrm I}^p(K^{\bullet\bullet}))
\to\mathrm H^p(K^{\bullet\bullet})$~; l'égalité de tous les $Z_r^{p0}$ permet
alors de définir des homomorphismes canoniques
$'E_r^{p0}\to'E_s^{p0}$ pour $r\leq s$, et en particulier un homomorphisme
$'E_2^{p0}\to'E_\infty^{p0}$, d'où par composition l'edge-homomorphisme
$'E_2^{p0}\to\mathrm H^p(K^{\bullet\bullet})$~; en outre, on vérifie aussitôt
que, à la classe mod. $B_2^{p0}$ d'un élément
$z\in Z_{\mathrm{II}}^0(K^{\bullet\bullet})\subset K^{p0}$ tel que $d'z=0$,
l'edge-homomorphisme ainsi défini fait correspondre dans $'E_\infty^{p0}$ la
classe de $z$ mod. $B_\infty^{p0}$, puis à cette dernière la classe de
cohomologie de $z$ dans $\mathrm H^p(K^{\bullet\bullet})$. On voit donc
finalement que \emph{l'edge-homomorphisme
(\hyperref[0.11.3.4.1-fr]{11.3.4.1}) provient, par passage à la cohomologie,
de l'injection canonique}
$Z_{\mathrm{II}}^0(K^{\bullet\bullet})\to K^{\bullet\bullet}$ (où
$K^{\bullet\bullet}$ est considéré comme complexe simple). On interprète
naturellement de même l'edge-homomorphisme
\[
\label{0.11.3.4.2-fr}
  ''E_2^{p0}(K^{\bullet\bullet})\longrightarrow
  \mathrm H^p(K^{\bullet\bullet})
  \tag{11.3.4.2}
\]
comme provenant de l'injection canonique
$Z_{\mathrm I}^0(K^{\bullet\bullet})\to K^{\bullet\bullet}$.
\end{env}

\begin{env}[11.3.5]
\label{0.11.3.5-fr}
Soit maintenant $K_{\bullet\bullet}=(K_{ij})$ un bicomplexe de
$\mathcal C$ dont les deux opérateurs de dérivation sont de degré $-1$. Nous
écrirons alors $K_{i,\bullet}$ (resp. $K_{\bullet,j}$) le complexe simple
$(K_{ij})_{j\in\mathbb Z}$ (resp. $(K_{ij})_{i\in\mathbb Z}$),
$\mathrm H_p^{\mathrm{II}}(K_{i,\bullet})$ (resp.
$\mathrm H_p^{\mathrm I}(K_{\bullet,j})$) son $p$-ème objet d'homologie,
$\mathrm H_p^{\mathrm{II}}(K_{\bullet\bullet})$ (resp.
$\mathrm H_p^{\mathrm I}(K_{\bullet\bullet})$) le complexe
$(\mathrm H_p^{\mathrm{II}}(K_{i,\bullet}))_{i\in\mathbb Z}$ (resp.
$(\mathrm H_p^{\mathrm I}(K_{\bullet,j}))_{j\in\mathbb Z}$),
$\mathrm H_q^{\mathrm I}(\mathrm H_p^{\mathrm{II}}(K_{\bullet\bullet}))$
(resp.
$\mathrm H_q^{\mathrm{II}}(\mathrm H_p^{\mathrm I}(K_{\bullet\bullet}))$)
son $q$-ème objet d'homologie~; notations analogues pour les objets des cycles
et objets des bords~; enfin, $\mathrm H_n(K_{\bullet\bullet})$ désignera
(lorsqu'il existe) le $n$-ème objet d'homologie du complexe simple (à
opérateur de dérivation de degré $-1$) défini par $K_{\bullet\bullet}$.

Soit $K^{\prime\bullet\bullet}=(K^{\prime ij})$ avec
$K^{\prime ij}=K_{-i,-j}$ le bicomplexe à opérateurs de dérivation de degrés
$+1$ associé à $K_{\bullet\bullet}$~; par définition, les
\emph{suites spectrales} de $K_{\bullet\bullet}$ sont celles de
$K^{\prime\bullet\bullet}$, que l'on écrit $'E(K_{\bullet\bullet})$ et
$''E(K_{\bullet\bullet})$, où l'on change toutefois les notations, en posant
\[
  'E_{pq}^r(K_{\bullet\bullet})
  ='E_r^{-p,-q}(K^{\prime\bullet\bullet}),\qquad
  ''E_{pq}^r(K_{\bullet\bullet})
  =''E_r^{-p,-q}(K^{\prime\bullet\bullet})
\]
pour $2\leq r\leq\infty$. Avec ces notations, on a
\[
  'E_{pq}^2(K_{\bullet\bullet})
  =\mathrm H_p^{\mathrm I}(\mathrm H_q^{\mathrm{II}}(K_{\bullet\bullet})),
  \qquad
  ''E_{pq}^2(K_{\bullet\bullet})
  =\mathrm H_q^{\mathrm{II}}(\mathrm H_p^{\mathrm I}(K_{\bullet\bullet})).
\]
Pour éviter des erreurs de signe, il sera en général préférable, pour les
relations entre ces suites spectrales et leur aboutissement, de revenir au
complexe $K^{\prime\bullet\bullet}$. Notons toutefois les critères
correspondant à (\hyperref[0.11.3.3-fr]{11.3.3})~:
\end{env}

\begin{env}[11.3.6]
\label{0.11.3.6-fr}
Les suites spectrales $'E(K_{\bullet\bullet})$ et
$''E(K_{\bullet\bullet})$ sont \emph{birégulières} dans les cas suivants~:
\begin{enumerate}
  \item[(a)] Il existe $i_0$ et $j_0$ tels que $K_{ij}=0$ pour $i>i_0$ ou
    pour $j>j_0$ (resp. pour $i<i_0$
    \oldpage[0\textsubscript{III}]{32}
    ou pour $j<j_0$).
  \item[(b)] Il existe $i_0$ et $i_1$ tels que $K_{ij}=0$ pour $i<i_0$ et
    $i>i_1$.
  \item[(c)] Il existe $j_0$ et $j_1$ tels que $K_{ij}=0$ pour $j<j_0$ et
    $j>j_1$.
\end{enumerate}

La suite $'E(K_{\bullet\bullet})$ est régulière s'il existe $i_0$ tel que
$K_{ij}=0$ pour $i<i_0$, ou s'il existe $j_0$ tel que $K_{ij}=0$ pour
$j>j_0$.

La suite $''E(K_{\bullet\bullet})$ est régulière s'il existe $i_0$ tel que
$K_{ij}=0$ pour $i>i_0$, ou s'il existe $j_0$ tel que $K_{ij}=0$ pour
$j<j_0$.
\end{env}

\subsection{Hypercohomologie d'un foncteur par rapport à un complexe
$K^\bullet$}
\label{subsection:0.11.4-fr}

\begin{env}[11.4.1]
\label{0.11.4.1-fr}
Soit $\mathcal C$ une catégorie abélienne~; rappelons que l'on appelle
\emph{résolution droite} (ou \emph{cohomologique}) d'un objet $A$ de
$\mathcal C$ un complexe d'objets de $\mathcal C$, dont l'opérateur de
dérivation est de degré $+1$,
\[
  0\longrightarrow L^0\longrightarrow L^1\longrightarrow L^2
  \longrightarrow\cdots
\]
muni d'un morphisme $\varepsilon:A\to L^0$ dit \emph{augmentation} de la
résolution (et que l'on peut considérer comme un morphisme de complexes
\[
\xymatrix{
  0\ar[r] & A\ar[r]\ar[d]_{\varepsilon} & 0\ar[r]\ar[d]
    & 0\ar[r]\ar[d] & \cdots\\
  0\ar[r] & L^0\ar[r] & L^1\ar[r] & L^2\ar[r] & \cdots
})
\]
tel que la suite
\[
  0\longrightarrow A\xrightarrow{\varepsilon}L^0\longrightarrow L^1
  \longrightarrow\cdots
\]
soit exacte~; de même une \emph{résolution gauche} (ou \emph{homologique})
de $A$ est un complexe
$0\leftarrow L_0\leftarrow L_1\leftarrow\cdots$ d'objets de $\mathcal C$
dont l'opérateur de dérivation est de degré $-1$, muni d'une augmentation
$\varepsilon:L_0\to A$, de sorte que la suite
\[
  0\longleftarrow A\xleftarrow{\varepsilon}L_0\longleftarrow L_1
  \longleftarrow\cdots
\]
soit exacte.

Lorsqu'une résolution droite $(L_i)_{i\geq0}$ d'un objet $A$ est telle que
$L_i=0$ pour $i\geq n+1$, on dit que cette résolution est de
\emph{longueur $\leq n$}. On définit de même une résolution gauche de
longueur $\leq n$. Une résolution qui est de longueur $\leq n$ pour un entier
$n$ est dite \emph{finie}.

Une résolution de $A$ est dite \emph{projective} (resp. \emph{injective}) si
les objets de $\mathcal C$ autres que $A$ qui la composent sont
\emph{projectifs} (resp. \emph{injectifs}). Lorsque $\mathcal C$ est la
catégorie des modules (à gauche par exemple) sur un anneau, on dira de même
qu'une résolution de $A$ est \emph{plate} (resp. \emph{libre}) lorsque les
modules autres que $A$ qui la composent sont \emph{plats} (resp.
\emph{libres}).
\end{env}

\begin{env}[11.4.2]
\label{0.11.4.2-fr}
Soit $K^\bullet=(K^i)_{i\in\mathbb Z}$ un complexe d'objets de $\mathcal C$,
dont l'opérateur de dérivation est de degré $+1$.

On appelle \emph{résolution de Cartan-Eilenberg droite} de $K^\bullet$ le
couple formé d'un bicomplexe $L^{\bullet\bullet}=(L^{ij})$ à opérateurs de
dérivation de degré $+1$, avec $L^{ij}=0$ pour $j<0$, et d'un morphisme de
complexes simples $\varepsilon:K^\bullet\to L^{\bullet,0}$, de façon que les
conditions suivantes soient remplies~:
\oldpage[0\textsubscript{III}]{33}
\begin{enumerate}
  \item[(i)] Pour chaque indice $i$, les suites
    \[
    \begin{aligned}
      0\longrightarrow{}&K^i\xrightarrow{\varepsilon}L^{i0}
        \longrightarrow L^{i1}\longrightarrow\cdots\\
      0\longrightarrow{}&\mathrm B^i(K^\bullet)\xrightarrow{\varepsilon}
        \mathrm B_{\mathrm I}^i(L^{\bullet,0})
        \longrightarrow\mathrm B_{\mathrm I}^i(L^{\bullet,1})
        \longrightarrow\cdots\\
      0\longrightarrow{}&\mathrm Z^i(K^\bullet)\xrightarrow{\varepsilon}
        \mathrm Z_{\mathrm I}^i(L^{\bullet,0})
        \longrightarrow\mathrm Z_{\mathrm I}^i(L^{\bullet,1})
        \longrightarrow\cdots\\
      0\longrightarrow{}&\mathrm H^i(K^\bullet)\xrightarrow{\varepsilon}
        \mathrm H_{\mathrm I}^i(L^{\bullet,0})
        \longrightarrow\mathrm H_{\mathrm I}^i(L^{\bullet,1})
        \longrightarrow\cdots
    \end{aligned}
    \]
    sont exactes, autrement dit $(L^{i,\bullet})$,
    $(\mathrm B_{\mathrm I}^i(L^{\bullet,\bullet}))$,
    $(\mathrm Z_{\mathrm I}^i(L^{\bullet,\bullet}))$ et
    $(\mathrm H_{\mathrm I}^i(L^{\bullet,\bullet}))$ sont respectivement
    des \emph{résolutions} de $K^i$, $\mathrm B^i(K^\bullet)$,
    $\mathrm Z^i(K^\bullet)$ et $\mathrm H^i(K^\bullet)$.
  \item[(ii)] Pour chaque $j$, le complexe simple $L^{\bullet,j}$ est
    \emph{scindé}, autrement dit, les suites exactes
    \[
      \label{0.11.4.2.1-fr}
      0\longrightarrow\mathrm B_{\mathrm I}^i(L^{\bullet,j})
      \longrightarrow\mathrm Z_{\mathrm I}^i(L^{\bullet,j})
      \longrightarrow\mathrm H_{\mathrm I}^i(L^{\bullet,j})
      \longrightarrow0
      \tag{11.4.2.1}
    \]
    \[
      \label{0.11.4.2.2-fr}
      0\longrightarrow\mathrm Z_{\mathrm I}^i(L^{\bullet,j})
      \longrightarrow L^{ij}
      \longrightarrow\mathrm B_{\mathrm I}^{i+1}(L^{\bullet,j})
      \longrightarrow0
      \tag{11.4.2.2}
    \]
    sont scindées.
\end{enumerate}

On prouve (M, XVII, 1.2) que si tout objet de $\mathcal C$ est sous-objet
d'un objet injectif, tout complexe $K^\bullet$ de $\mathcal C$ admet une
résolution de Cartan-Eilenberg \emph{injective}, c'est-à-dire formée d'objets
injectifs $L^{ij}$ (la condition (ii) ci-dessus entraîne alors que les
$\mathrm B_{\mathrm I}^i(L^{\bullet,j})$,
$\mathrm Z_{\mathrm I}^i(L^{\bullet,j})$ et
$\mathrm H_{\mathrm I}^i(L^{\bullet,j})$ sont aussi des objets injectifs).
En outre, pour tout morphisme $f:K^\bullet\to K'^\bullet$ de complexes de
$\mathcal C$, toute résolution de Cartan-Eilenberg $L^{\bullet\bullet}$ de
$K^\bullet$ et toute résolution injective de Cartan-Eilenberg
$L'^{\bullet\bullet}$ de $K'^\bullet$, il existe un morphisme de bicomplexes
$F:L^{\bullet\bullet}\to L'^{\bullet\bullet}$ compatible avec $f$ et les
augmentations, et si $f$ et $g$ sont deux morphismes homotopes de $K^\bullet$
dans $K'^\bullet$, les morphismes correspondants de $L^{\bullet\bullet}$
dans $L'^{\bullet\bullet}$ sont homotopes (\emph{loc. cit.}).

Lorsque $K^\bullet$ est \emph{limité inférieurement} (resp.
\emph{supérieurement}), on peut prendre $L^{\bullet\bullet}$ tel que
$L^{ij}=0$ pour $i<i_0$ (resp. $i>i_0$) si $K^i=0$ pour $i<i_0$ (resp.
$i>i_0$) (M, XVII, 1.3).

Supposons d'autre part qu'il existe un entier $n$ tel que tout objet de
$\mathcal C$ admette une \emph{résolution injective de longueur $\leq n$}~;
alors on peut supposer que l'on a $L^{ij}=0$ pour $j>n$ (M, XVII, 1.4).
\end{env}

\begin{env}[11.4.3]
\label{0.11.4.3-fr}
Soit maintenant $T$ un \emph{foncteur covariant additif} de $\mathcal C$
dans une catégorie abélienne $\mathcal C'$. Étant donnés un complexe
$K^\bullet$ de $\mathcal C$ et une résolution de Cartan-Eilenberg injective
$L^{\bullet\bullet}$ de $K^\bullet$, supposons que le complexe (simple) défini
par le bicomplexe $T(L^{\bullet\bullet})$ existe
(cf. \hyperref[0.11.3.1-fr]{11.3.1})~; alors les deux suites spectrales
$'E(T(L^{\bullet\bullet}))$ et $''E(T(L^{\bullet\bullet}))$ de ce bicomplexe
sont dites \emph{suites spectrales d'hypercohomologie de $T$ par rapport au
complexe $K^\bullet$}~; en vertu de
(\hyperref[0.11.4.2-fr]{11.4.2}) et
(\hyperref[0.11.3.2-fr]{11.3.2}), elles ne dépendent effectivement que de
$K^\bullet$ et non de la résolution de Cartan-Eilenberg injective
$L^{\bullet\bullet}$ choisie~; en outre, elles dépendent
\emph{fonctoriellement} de $K^\bullet$. Elles ont un même aboutissement
$\mathrm H^\bullet(T(L^{\bullet\bullet}))$ appelé encore
l'\emph{hypercohomologie de $T$ par rapport à $K^\bullet$}, et noté
$\mathrm R^\bullet T(K^\bullet)$. On montre que les termes $E_2$ des deux
suites spectrales précédentes sont donnés par
\[
  \label{0.11.4.3.1-fr}
  'E_2^{pq}=\mathrm H^p(\mathrm R^qT(K^\bullet))
  \tag{11.4.3.1}
\]
\[
  \label{0.11.4.3.2-fr}
  ''E_2^{pq}=\mathrm R^pT(\mathrm H^q(K^\bullet))
  \tag{11.4.3.2}
\]
\oldpage[0\textsubscript{III}]{34}
$\mathrm R^pT$ désignant comme d'ordinaire le $p$-ième \emph{foncteur
dérivé} de $T$ pour $p\in\mathbb Z$~; $\mathrm R^qT(K^\bullet)$ désigne le
complexe $(\mathrm R^qT(K^i))_{i\in\mathbb Z}$.

\phantomsection\label{0.11.4.4-fr}
Sauf mention expresse du contraire, nous supposerons désormais que tout objet
de $\mathcal C$ est sous-objet d'un objet injectif de $\mathcal C$, de sorte
que les résolutions de Cartan-Eilenberg injectives existent pour tout complexe
de $\mathcal C$. Comme $L^{ij}=0$ pour $j<0$, les critères de
(\hyperref[0.11.3.3-fr]{11.3.3}) montrent que les \emph{deux} suites
spectrales d'hypercohomologie de $T$ par rapport à $K^\bullet$ existent et
sont \emph{birégulières} dans chacun des deux cas suivants~:
\begin{enumerate}
  \item[$1^\circ$] $K^\bullet$ est limité inférieurement~;
  \item[$2^\circ$] Tout objet de $\mathcal C$ admet une résolution injective
    de longueur au plus égale à un entier $n$ (indépendant de l'objet
    considéré).
\end{enumerate}
En effet, dans le premier cas, on peut supposer
(\hyperref[0.11.4.2-fr]{11.4.2}) qu'il existe $i_0$ tel que $L^{ij}=0$ pour
$i<i_0$ et dans le second qu'il existe $j_1$ tel que $L^{ij}=0$ pour $j>j_1$~;
dans chacun des deux cas, il est clair en outre que pour $n$ donné, il n'y a
qu'un nombre fini de couples $(i,j)$ tels que $L^{ij}\neq0$ et $i+j=n$, ce
qui établit nos assertions.

Lorsqu'on suppose que dans $\mathcal C'$ les limites inductives filtrantes
existent et sont exactes (ce qui implique en particulier l'existence dans
$\mathcal C'$ des sommes directes infinies), alors le complexe défini par le
bicomplexe $T(L^{\bullet\bullet})$ existe, et le critère
(\hyperref[0.11.3.3-fr]{11.3.3}) montre que la suite
$'E(T(L^{\bullet\bullet}))$ est toujours \emph{régulière}.

Lorsque $K^\bullet$ est un complexe dont tous les termes $K^i$ sont nuls sauf
un seul $K^{i_0}$, $\mathrm R^nT(K^\bullet)$ est isomorphe à
$\mathrm R^{n-i_0}T(K^{i_0})$, ainsi qu'il résulte aussitôt des définitions en
prenant une résolution de Cartan-Eilenberg $L^{\bullet\bullet}$ telle que
$L^{ij}=0$ pour $i\neq i_0$.

Si $K^\bullet$ et $K'^\bullet$ sont deux complexes de $\mathcal C$, $f$, $g$
deux morphismes \emph{homotopes} de $K^\bullet$ dans $K'^\bullet$, alors les
morphismes $\mathrm R^\bullet T(K^\bullet)\to
\mathrm R^\bullet T(K'^\bullet)$ déduits de $f$ et $g$ sont identiques, et il
en est de même pour les morphismes des suites spectrales de cohomologie.
\end{env}

\begin{proposition}[11.4.5]
\label{0.11.4.5-fr}
Supposons que dans $\mathcal C'$ les limites inductives filtrantes existent et
soient exactes. Si $\mathrm R^nT(K^i)=0$ pour tout $n>0$ et tout
$i\in\mathbb Z$, on a des isomorphismes fonctoriels
\[
  \label{0.11.4.5.1-fr}
  \mathrm R^iT(K^\bullet)\xrightarrow{\sim}
  \mathrm H^i(T(K^\bullet))\qquad(i\in\mathbb Z).
  \tag{11.4.5.1}
\]
\end{proposition}

\begin{proof}
En effet, les seuls termes $E_2$ non nuls de la première suite spectrale
(\hyperref[0.11.4.3.1-fr]{11.4.3.1}) sont alors
$'E_2^{p0}=\mathrm H^p(T(K^\bullet))$~; autrement dit, cette suite est
\emph{dégénérée}~; comme elle est régulière
(\hyperref[0.11.4.4-fr]{11.4.4}), la conclusion résulte de
(\hyperref[0.11.1.6-fr]{11.1.6}).
\end{proof}

\begin{env}[11.4.6]
\label{0.11.4.6-fr}
Considérons maintenant, par exemple, un \emph{bifoncteur covariant}
$(M,N)\mapsto T(M,N)$ de $\mathcal C\times\mathcal C'$ dans $\mathcal C''$,
$\mathcal C$, $\mathcal C'$, $\mathcal C''$ étant trois catégories
abéliennes~; on suppose, pour simplifier, $T$ additif en chacun de ses
arguments, et en outre que tout objet de $\mathcal C$ et tout objet de
$\mathcal C'$ sont sous-objets d'un objet injectif, et que les limites
inductives filtrantes existent dans $\mathcal C''$ et sont exactes. On définit
alors l'\emph{hypercohomologie de $T$ par rapport à deux complexes}
$K^\bullet$, $K'^\bullet$ de $\mathcal C$ et $\mathcal C'$ respectivement, à
opérateurs de dérivation de degré $+1$, en prenant pour $K^\bullet$ (resp.
$K'^\bullet$) une résolution de Cartan-Eilenberg injective
$L^{\bullet\bullet}$ (resp. $L'^{\bullet\bullet}$)~; alors
$T(L^{\bullet\bullet},L'^{\bullet\bullet})$ est un quadricomplexe de
$\mathcal C''$, que l'on considère comme \emph{bicomplexe} de $\mathcal C''$
en prenant pour degrés de $T(L^{ij},L'^{hk})$ les entiers $i+h$ et $j+k$.
L'\emph{hypercohomologie de $T$ par rapport à $K^\bullet$ et $K'^\bullet$}
est par définition la cohomologie
$\mathrm H^\bullet(T(L^{\bullet\bullet},L'^{\bullet\bullet}))$ de ce
bicomplexe (autrement dit, celle du complexe simple associé) notée
$\mathrm R^\bullet T(K^\bullet,K'^\bullet)$~; elle est l'aboutissement de
deux suites spectrales dont les termes $E_2$ sont donnés par
\oldpage[0\textsubscript{III}]{35}
\[
  \label{0.11.4.6.1-fr}
  'E_2^{pq}=\mathrm H^p(\mathrm R^qT(K^\bullet,K'^\bullet))
  \tag{11.4.6.1}
\]
\[
  \label{0.11.4.6.2-fr}
  ''E_2^{pq}=\sum_{q'+q''=q}
  \mathrm R^pT(\mathrm H^{q'}(K^\bullet),\mathrm H^{q''}(K'^\bullet))
  \tag{11.4.6.2}
\]
(cf. M, XVII, 2).

Ici $\mathrm R^qT(K^\bullet,K'^\bullet)$ est le bicomplexe
$(\mathrm R^qT(K^i,K'^j))_{(i,j)\in\mathbb Z\times\mathbb Z}$ et le second
membre de (\hyperref[0.11.4.6.1-fr]{11.4.6.1}) est sa cohomologie quand on le
considère comme complexe simple.

En outre, la première suite spectrale est toujours régulière, et les deux
suites spectrales sont birégulières lorsqu'il existe $n$ tel que tout objet de
$\mathcal C$ et tout objet de $\mathcal C'$ admettent une résolution injective
de longueur $\leq n$, ou lorsque $K^\bullet$ et $K'^\bullet$ sont limités
inférieurement~; dans ce dernier cas, on peut en outre omettre l'hypothèse que
les limites inductives existent dans $\mathcal C$ et $\mathcal C'$.

Si $K_1^\bullet$, $K_1'^\bullet$ sont deux autres complexes de $\mathcal C$ et
$\mathcal C'$ respectivement, $f$, $g$ deux morphismes homotopes de
$K^\bullet$ dans $K_1^\bullet$, $f'$, $g'$ deux morphismes homotopes de
$K'^\bullet$ dans $K_1'^\bullet$, alors les morphismes
$\mathrm R^\bullet T(K^\bullet,K'^\bullet)\to
\mathrm R^\bullet T(K_1^\bullet,K_1'^\bullet)$ déduits de $f$ et $f'$ d'une
part, de $g$ et $g'$ d'autre part, sont identiques, et il en est de même pour
les morphismes des suites spectrales d'hypercohomologie.

On généralise aisément à un multifoncteur covariant additif quelconque.
\end{env}

\begin{proposition}[11.4.7]
\label{0.11.4.7-fr}
Supposons que pour tout objet injectif $I$ de $\mathcal C$ (resp. $I'$ de
$\mathcal C'$), $A'\mapsto T(I,A')$ (resp. $A\mapsto T(A,I')$) soit un
foncteur exact. Alors, avec les notations de
(\hyperref[0.11.4.6-fr]{11.4.6}), on a des isomorphismes canoniques
\[
  \label{0.11.4.7.1-fr}
  \mathrm R^\bullet T(K^\bullet,K'^\bullet)
  \xrightarrow{\sim}\mathrm H^\bullet(T(L^{\bullet\bullet},K'^\bullet))
  \xrightarrow{\sim}\mathrm H^\bullet(T(K^\bullet,L'^{\bullet\bullet}))
  \tag{11.4.7.1}
\]
où les deux derniers termes sont la cohomologie des complexes simples définis
par les tricomplexes $T(L^{\bullet\bullet},K'^\bullet)$ et
$T(K^\bullet,L'^{\bullet\bullet})$ respectivement.
\end{proposition}

\begin{proof}
Définissons par exemple le premier de ces isomorphismes. Le quadricomplexe
$T(L^{\bullet\bullet},L'^{\bullet\bullet})$ peut être considéré comme un
\emph{bicomplexe}, en prenant comme degrés de $T(L^{ij},L'^{hk})$ les nombres
$i+j+h$ et $k$. Comme pour chaque $h$, $L'^{h,\bullet}$ est une
\emph{résolution} de $K'^h$, on a, pour ce bicomplexe, en vertu de l'hypothèse
sur $T$,
$\mathrm H_{\mathrm{II}}^q(T(L^{\bullet\bullet},L'^{\bullet\bullet}))=0$
pour $q\neq0$ et
$\mathrm H_{\mathrm{II}}^0(T(L^{\bullet\bullet},L'^{\bullet\bullet}))
=T(L^{\bullet\bullet},K'^\bullet)$~; la première suite spectrale de ce
bicomplexe est donc \emph{dégénérée}~; comme $L'^{hk}=0$ pour $k<0$, cette
suite est en outre régulière (\hyperref[0.11.3.3-fr]{11.3.3}), et la
conclusion résulte donc de (\hyperref[0.11.1.6-fr]{11.1.6}).
\end{proof}

On a des résultats analogues pour un multifoncteur covariant en un nombre
quelconque $n$ d'arguments~: dans le calcul de l'hypercohomologie, il n'est pas
nécessaire de remplacer \emph{tous} les complexes par une résolution de
Cartan-Eilenberg, mais seulement $n-1$ d'entre eux, pourvu que, lorsqu'on fixe
$n-1$ arguments quelconques en leur donnant pour valeurs des objets
\emph{injectifs}, le foncteur covariant en l'argument restant soit
\emph{exact}.

\subsection{Passage à la limite inductive dans l'hypercohomologie}
\label{subsection:0.11.5-fr}

\begin{lemma}[11.5.1]
\label{0.11.5.1-fr}
Soit $K^\bullet=(K^i)_{i\in\mathbb Z}$ un complexe de $\mathcal C$, et pour
tout entier $r\in\mathbb Z$, soit $K_{(r)}^\bullet$ le complexe tel que
$K_{(r)}^i=0$ pour $i<r$, $K_{(r)}^i=K^i$ pour $i\geq r$. Soit $T$ un
foncteur covariant additif
\oldpage[0\textsubscript{III}]{36}
de $\mathcal C$ dans $\mathcal C'$, permutant aux
limites inductives (on suppose que les limites inductives filtrantes existent
et sont exactes dans $\mathcal C$ et $\mathcal C'$). Alors
$\mathrm R^\bullet T(K^\bullet)$ est isomorphe à la limite inductive
\[
  \varinjlim_{r\to-\infty}\mathrm R^\bullet T(K_{(r)}^\bullet)
\]
lorsque $r$ tend vers $-\infty$.
\end{lemma}

La construction d'une résolution injective de Cartan-Eilenberg de $K^\bullet$
se fait en choisissant arbitrairement, pour chaque $i$, une résolution
injective $(X_{\mathrm B}^{ij})_{j\geq0}$ de $\mathrm B^i(K^\bullet)$ et une
résolution injective $(X_{\mathrm H}^{ij})_{j\geq0}$ de
$\mathrm H^i(K^\bullet)$~; cela fait, la méthode de construction montre que
la résolution injective $(L^{ij})_{j\geq0}$ de $K^i$ et les opérateurs de
dérivation $L^{i,j}\to L^{i+1,j}$ \emph{ne dépendent que des résolutions}
$(X_{\mathrm B}^{i,\bullet})$, $(X_{\mathrm H}^{i,\bullet})$ et
$(X_{\mathrm B}^{i+1,\bullet})$ (M, XVII, 1.2). Or, il est clair que l'on a
$\mathrm B^i(K_{(r)}^\bullet)=\mathrm B^i(K^\bullet)$ et
$\mathrm H^i(K_{(r)}^\bullet)=\mathrm H^i(K^\bullet)$ pour $i\geq r+1$. On
a, d'autre part, pour chaque $r$, une injection canonique
$\varphi_{r-1,r}^\bullet:K_{(r)}^\bullet\to K_{(r-1)}^\bullet$,
$\varphi_{r-1,r}^i$ étant l'identité pour $i\neq r-1$. La remarque précédente
montre que, si $L^{\bullet\bullet}=(L^{ij})$ est une résolution injective de
Cartan-Eilenberg de $K^\bullet$, on peut pour chaque $r$ définir une résolution
injective de Cartan-Eilenberg
$L_{(r)}^{\bullet\bullet}=(L_{(r)}^{ij})$ de $K_{(r)}^\bullet$ telle que
$L_{(r)}^{ij}=0$ pour $i<r$ et $L_{(r)}^{ij}=L^{ij}$ pour $i\geq r+1$. On peut
d'autre part définir un morphisme de bicomplexes
$\Phi_{r-1,r}^{\bullet\bullet}:L_{(r)}^{\bullet\bullet}\to
L_{(r-1)}^{\bullet\bullet}$ correspondant à $\varphi_{r-1,r}^\bullet$, et la
méthode de définition de ce morphisme (\emph{loc. cit.}) montre encore qu'on
peut le construire de sorte que $\Phi_{r-1,r}^{ij}$ soit l'identité pour
$i\neq r$ et $i\neq r-1$. On a ainsi défini un système inductif
$(L_{(r)}^{\bullet\bullet})$ de bicomplexes de $\mathcal C$, dont
$L^{\bullet\bullet}$ est évidemment la limite inductive lorsque $r$ tend vers
$-\infty$~; en raison de la permutabilité des sommes directes et des limites
inductives, le complexe simple associé à $L^{\bullet\bullet}$ est aussi limite
inductive du complexe simple associé à $L_{(r)}^{\bullet\bullet}$. Comme $T$
permute par hypothèse aux limites inductives et qu'il en est de même de la
cohomologie (en raison de l'exactitude du foncteur $\varinjlim$), on a bien
\[
  \mathrm H^\bullet(T(L^{\bullet\bullet}))=
  \varinjlim\mathrm H^\bullet(T(L_{(r)}^{\bullet\bullet}))
\]
à un isomorphisme près.

Le lemme (\hyperref[0.11.5.1-fr]{11.5.1}) permet d'étendre à des complexes
$K^\bullet$ quelconques, par passage à la limite inductive, des propriétés
valables pour les complexes \emph{limités inférieurement}. Comme premier
exemple, nous prouverons~:

\begin{proposition}[11.5.2]
\label{0.11.5.2-fr}
Sous les hypothèses de (\hyperref[0.11.5.1-fr]{11.5.1}) concernant
$\mathcal C$, $\mathcal C'$ et $T$, $\mathrm R^\bullet T(K^\bullet)$ est un
foncteur cohomologique dans la catégorie abélienne des complexes de
$\mathcal C$.
\end{proposition}

\begin{proof}
Montrons qu'on peut se ramener au cas des complexes limités inférieurement~:
si on a une suite exacte
\[
  0\longrightarrow K'^\bullet\longrightarrow K^\bullet
  \longrightarrow K''^\bullet\longrightarrow0
\]
de complexes, on en déduit évidemment pour chaque $r$ une suite exacte
\[
  0\longrightarrow K_{(r)}'^\bullet\longrightarrow K_{(r)}^\bullet
  \longrightarrow K_{(r)}''^\bullet\longrightarrow0,
\]
d'où par hypothèse, une suite exacte
\[
  \cdots\longrightarrow\mathrm R^nT(K_{(r)}'^\bullet)
  \longrightarrow\mathrm R^nT(K_{(r)}^\bullet)
  \longrightarrow\mathrm R^nT(K_{(r)}''^\bullet)
  \xrightarrow{\partial}\mathrm R^{n+1}T(K_{(r)}'^\bullet)
  \longrightarrow\cdots
\]
ces suites exactes formant un système inductif~; le lemme
(\hyperref[0.11.5.1-fr]{11.5.1}) et l'exactitude du foncteur $\varinjlim$
démontrent que l'on a une suite exacte
\[
  \cdots\longrightarrow\mathrm R^nT(K'^\bullet)
  \longrightarrow\mathrm R^nT(K^\bullet)
  \longrightarrow\mathrm R^nT(K''^\bullet)
  \xrightarrow{\partial}\mathrm R^{n+1}T(K'^\bullet)
  \longrightarrow\cdots
\]
\end{proof}

Pour traiter le cas des complexes limités inférieurement, on peut se borner à
ceux tels que $K^i=0$ pour $i<0$~; ils forment évidemment une catégorie
abélienne $\mathcal K$.

\begin{lemma}[11.5.2.1]
\label{0.11.5.2.1-fr}
Dans $\mathcal K$, soit $\mathfrak J$ l'ensemble des complexes
$Q^\bullet=(Q^i)_{i\geq0}$ ayant les
\oldpage[0\textsubscript{III}]{37}
propriétés suivantes~:
\begin{enumerate}
  \item[$1^\circ$] Tout $Q^i$ est un objet injectif de $\mathcal C$~;
  \item[$2^\circ$] Pour tout $i\geq0$, on a
    $\mathrm Z^i(Q^\bullet)=\mathrm B^i(Q^\bullet)$, et
    $\mathrm Z^i(Q^\bullet)$ est facteur direct de $Q^i$.
\end{enumerate}
Alors~:
\begin{enumerate}
  \item[(i)] Tout $Q^\bullet\in\mathfrak J$ est un objet injectif de
    $\mathcal K$.
  \item[(ii)] Tout objet de $\mathcal K$ est isomorphe à un sous-complexe
    d'un complexe appartenant à $\mathfrak J$.
\end{enumerate}
\end{lemma}

\begin{proof}
\textup{(i)} Soient $A^\bullet=(A^i)$ un objet de $\mathcal K$,
$A'^\bullet=(A'^i)$ un sous-objet de $A^\bullet$,
$Q^\bullet=(Q^i)$ un objet de $\mathfrak J$, et supposons donné un morphisme
$f=(f^i):A'^\bullet\to Q^\bullet$, qu'il s'agit de prolonger en un morphisme
$g=(g^i):A^\bullet\to Q^\bullet$. Nous utiliserons le langage de la catégorie
des modules pour simplifier (cf. [27]).

Identifions $Q^i$ à
$\mathrm B^i(Q^\bullet)\oplus\mathrm B^{i+1}(Q^\bullet)$~; procédons par
récurrence sur $i$, et supposons donc les $g^j$ définis pour $j<i$,
compatibles avec les opérateurs de dérivation $d^j:A^j\to A^{j+1}$ et
$d'^j:Q^j\to Q^{j+1}$ pour $j<i-1$ et tels en outre que~:
\begin{enumerate}
  \item[$1^\circ$] $g^{i-1}(\mathrm Z^{i-1}(A^\bullet))\subset
    \mathrm Z^{i-1}(Q^\bullet)$~;
  \item[$2^\circ$] Si on pose $C^j=(d^j)^{-1}(A'^{j+1})$ pour tout $j$, alors
    $d'^{i-1}\circ g^{i-1}$ coïncide avec $f^i\circ d^{i-1}$ sur $C^{i-1}$.
\end{enumerate}
Le morphisme $f^i:A'^i\to Q^i$ donne par composition avec les projections
deux morphismes
$f'^i:A'^i\to\mathrm B^i(Q^\bullet)$ et
$f''^i:A'^i\to\mathrm B^{i+1}(Q^\bullet)$. Comme
$d'^{i-1}\circ g^{i-1}$ applique $A^{i-1}$ dans
$\mathrm B^i(Q^\bullet)$ et s'annule dans $\mathrm Z^{i-1}(A^\bullet)$, il
définit un morphisme
$h^i:\mathrm B^i(A^\bullet)\to\mathrm B^i(Q^\bullet)$, et puisque
$d'^{i-1}\circ g^{i-1}$ coïncide avec $f^i\circ d^{i-1}$ sur $C^{i-1}$,
$h^i$ coïncide avec $f'^i$ dans
$\mathrm B^i(A^\bullet)\cap A'^i$. Comme $\mathrm B^i(Q^\bullet)$, facteur
direct de $Q^i$, est injectif, il y a un morphisme
$g'^i:A^i\to\mathrm B^i(Q^\bullet)$ qui coïncide avec $h^i$ dans
$\mathrm B^i(A^\bullet)$ et avec $f'^i$ dans $A'^i$. Considérons d'autre part
le morphisme
$f''^{i+1}\circ d^i:C^i\to\mathrm B^{i+1}(Q^\bullet)$, qui s'annule dans
$\mathrm Z^i(A^\bullet)$~; comme $\mathrm B^{i+1}(Q^\bullet)$ est injectif,
il y a un morphisme $g''^i:A^i\to\mathrm B^{i+1}(Q^\bullet)$, qui coïncide
avec $f''^{i+1}\circ d^i$ dans $C^i$ et avec $0$ dans
$\mathrm Z^i(A^\bullet)$. Il suffit alors de prendre $g^i=g'^i+g''^i$ pour
pouvoir poursuivre la récurrence.

\textup{(ii)} Pour plonger $A^\bullet=(A^i)$ dans un complexe appartenant à
$\mathfrak J$, on prend pour chaque $i\geq1$ un objet injectif $Q'^i$ de
$\mathcal C$ tel qu'il existe une injection $f'^i:A^i\to Q'^i$. Posons alors
\[
  Q^i=0\quad(i<0),\qquad Q^0=Q'^1,qquad
  Q^i=Q'^i\oplus Q'^{i+1}\quad(i\geq1),
\]
avec l'opérateur de dérivation évident. Posons
$f^i=f'^i+(f'^{i+1}\circ d^i)$ pour tout $i$ (avec $f'^i=0$ pour $i\leq0$)~;
il est immédiat que $f^i$ est injectif pour tout $i$, et que les $f^i$ sont
compatibles avec les opérateurs de dérivation.
\end{proof}

\begin{corollary}[11.5.2.2]
\label{0.11.5.2.2-fr}
Tout objet $K^\bullet$ de $\mathcal K$ admet une résolution droite formée
d'objets de $\mathfrak J$. Si $L^{\bullet\bullet}$ est une telle résolution,
pour toute résolution $L'^{\bullet\bullet}$ de $K^\bullet$, formée d'objets
de $\mathcal K$, il y a un morphisme de bicomplexes
$F:L'^{\bullet\bullet}\to L^{\bullet\bullet}$ compatible avec les
augmentations, et deux tels morphismes $F$, $F'$ sont homotopes.
\end{corollary}

\begin{proof}
Ce n'est autre que (M, V, 1.1 a)) appliqué à la catégorie abélienne
$\mathcal K$.
\end{proof}

\begin{env}[11.5.2.3]
\label{0.11.5.2.3-fr}
Ces préliminaires étant posés, considérons une résolution de
Cartan-Eilenberg injective $L'^{\bullet\bullet}$ de $K^\bullet$ et une
résolution $L^{\bullet\bullet}$ de $K^\bullet$ formée d'objets de
$\mathfrak J$, et montrons que l'on a un isomorphisme
\[
  \mathrm H^\bullet(T(L'^{\bullet\bullet}))\xrightarrow{\sim}
  \mathrm H^\bullet(T(L^{\bullet\bullet})).
\]
On déduit en effet de (\hyperref[0.11.5.2.2-fr]{11.5.2.2}) un morphisme de
bicomplexes $T(L'^{\bullet\bullet})\to T(L^{\bullet\bullet})$, et par suite
un morphisme
$'E(T(L'^{\bullet\bullet}))\to{}'E(T(L^{\bullet\bullet}))$ des premières
suites spectrales de ces bicomplexes. Comme en vertu de
(\hyperref[0.11.3.3-fr]{11.3.3}), ces suites spectrales sont régulières, il
suffit (\hyperref[0.11.1.5-fr]{11.1.5}) de voir que le morphisme précédent est
un isomorphisme pour les termes $E_2$, ou encore que
$\mathrm H_{\mathrm{II}}^q(T(L^{i,\bullet}))$ est égal à
$\mathrm R^qT(K^i)$~; comme $L^{i,\bullet}$ est une résolution droite de
$K^i$, on est ramené à prouver le
\end{env}

\oldpage[0\textsubscript{III}]{38}
\begin{lemma}[11.5.2.4]
\label{0.11.5.2.4-fr}
Si $L^\bullet=(L^i)_{i\in\mathbb Z}$ est une résolution droite d'un objet $A$
de $\mathcal C$ telle que $\mathrm R^nT(L^i)=0$ pour tout $i\in\mathbb Z$ et
tout $n>0$, alors on a
$\mathrm H^\bullet T(L^\bullet)=\mathrm R^\bullet T(A)$.
\end{lemma}

\begin{proof}
C'est un cas particulier de (T, 2.5.2).
\end{proof}

\begin{env}[11.5.2.5]
\label{0.11.5.2.5-fr}
La démonstration de (\hyperref[0.11.5.2-fr]{11.5.2}) se termine maintenant de
façon immédiate, car $L'^{\bullet\bullet}$ est une résolution injective de
$K^\bullet$ dans la catégorie abélienne $\mathcal K$, autrement dit,
$K^\bullet\mapsto\mathrm H^\bullet(T(L'^{\bullet\bullet}))$ n'est autre que
le foncteur cohomologique dérivé droit de $T$ dans la catégorie $\mathcal K$
(T, 2.3).
\end{env}

\begin{proposition}[11.5.3]
\label{0.11.5.3-fr}
Sous les hypothèses de (\hyperref[0.11.5.1-fr]{11.5.1}) concernant
$\mathcal C$, $\mathcal C'$ et $T$, soit
$L^{\bullet\bullet}=(L^{ij})$ un bicomplexe tel que $L^{ij}=0$ pour $j<0$ et
que, pour tout $i$, $L^{i,\bullet}$ soit une résolution de $K^i$~; supposons
enfin que $\mathrm R^nT(L^{ij})=0$ pour tout couple $(i,j)$ et tout $n>0$.
Alors il existe un isomorphisme fonctoriel
\[
  \label{0.11.5.3.1-fr}
  \mathrm R^\bullet T(K^\bullet)\xrightarrow{\sim}
  \mathrm H^\bullet(T(L^{\bullet\bullet})).
  \tag{11.5.3.1}
\]
\end{proposition}

\begin{proof}
Soit $L_{(r)}^{\bullet\bullet}=(L_{(r)}^{ij})$ le bicomplexe tel que
$L_{(r)}^{ij}=0$ pour $i<r$, $L_{(r)}^{ij}=L^{ij}$ pour $i\geq r$~; il est
immédiat que $L^{\bullet\bullet}$ est limite inductive de
$L_{(r)}^{\bullet\bullet}$ pour $r$ tendant vers $-\infty$~; en vertu de
l'hypothèse sur $T$ et de (\hyperref[0.11.5.1-fr]{11.5.1}), il suffit donc de
prouver la proposition lorsque $K^\bullet$ est limité inférieurement, par
exemple $K^i=0$ pour $i<0$, et $L^{ij}=0$ pour $i<0$. Soit alors
$L'^{\bullet\bullet}=(L'^{ij})$ une résolution droite de $K^\bullet$ formée
d'objets de $\mathfrak J$ (\hyperref[0.11.5.2.2-fr]{11.5.2.2})~; il y a un
morphisme de bicomplexes
$L^{\bullet\bullet}\to L'^{\bullet\bullet}$ compatible avec les
augmentations, d'où un morphisme
$'E(T(L^{\bullet\bullet}))\to{}'E(T(L'^{\bullet\bullet}))$ pour les premières
suites spectrales~; le lemme (\hyperref[0.11.5.2.4-fr]{11.5.2.4}) montre,
comme dans (\hyperref[0.11.5.2.3-fr]{11.5.2.3}), que ce morphisme est un
\emph{isomorphisme}, d'où la conclusion.
\end{proof}

\begin{remark}[11.5.4]
\label{0.11.5.4-fr}
Les raisonnements précédents prouvent que les conclusions de
(\hyperref[0.11.5.2-fr]{11.5.2}) et
(\hyperref[0.11.5.3-fr]{11.5.3}) sont valables dans la catégorie des
complexes $K^\bullet$ limités inférieurement, \emph{sans supposer que $T$
permute aux limites inductives filtrantes}. En outre, quand on considère
seulement la catégorie $\mathcal K$ des complexes $K^\bullet$ tels que
$K^i=0$ pour $i<0$, la caractérisation de $\mathrm R^\bullet T(K^\bullet)$
comme le système des \emph{foncteurs dérivés droits} de $T$ dans $\mathcal K$
montre que ce foncteur cohomologique est \emph{universel} (T, 2.3).

Un autre cas où (\hyperref[0.11.5.2-fr]{11.5.2}) est valable sans faire
d'hypothèse supplémentaire sur $T$ est le cas où il existe un entier $m>0$
tel que tout objet de $\mathcal C$ admette une résolution injective de longueur
$\leq m$. En effet, dans la démonstration de
(\hyperref[0.11.5.1-fr]{11.5.1}), toutes les résolutions injectives d'objets
de $\mathcal C$ qui interviennent peuvent être prises de longueur $\leq m$,
d'où il résulte aussitôt que les termes de degré total $n$ du bicomplexe
$T(L_{(r)}^{\bullet\bullet})$ sont égaux à ceux de
$T(L^{\bullet\bullet})$ et en nombre fini, dès que $r$ est assez grand, ce qui
entraîne que pour tout $n$,
$\mathrm H^n(T(L^{\bullet\bullet}))=
\mathrm H^n(T(L_{(r)}^{\bullet\bullet}))$ dès que $r$ est assez grand. Avec
les notations de (\hyperref[0.11.5.2-fr]{11.5.2}), on a donc aussi
$\mathrm R^nT(K_{(r)}^\bullet)=\mathrm R^nT(K^\bullet)$ pour $r$ assez grand
(dépendant de $n$) et de même pour $K'^\bullet$ et $K''^\bullet$, d'où la
conclusion. De la même manière, (\hyperref[0.11.5.3-fr]{11.5.3}) est valable
sans condition supplémentaire sur $T$ lorsque $\mathcal C$ vérifie
l'hypothèse précédente et que l'on suppose que les résolutions
$L^{i,\bullet}$ sont de longueur $\leq m$.
\end{remark}

\begin{env}[11.5.5]
\label{0.11.5.5-fr}
Le résultat de (\hyperref[0.11.5.2-fr]{11.5.2}) se généralise aux
multifoncteurs covariants. Considérons par exemple la situation de
(\hyperref[0.11.4.6-fr]{11.4.6}), où l'on suppose que dans $\mathcal C$,
$\mathcal C'$ et $\mathcal C''$
\oldpage[0\textsubscript{III}]{39}
les limites inductives filtrantes existent et sont exactes, et que $T$ commute
aux limites inductives. Alors $\mathrm R^\bullet T(K^\bullet,K'^\bullet)$ est
un bifoncteur cohomologique en chacun des complexes $K^\bullet$,
$K'^\bullet$~; pour le voir, on se ramène comme dans
(\hyperref[0.11.5.2-fr]{11.5.2}) au cas où $K^\bullet$ et $K'^\bullet$ sont
limités inférieurement~; prenant alors pour $K^\bullet$ et $K'^\bullet$ des
résolutions injectives du type décrit dans
(\hyperref[0.11.5.2.2-fr]{11.5.2.2}), on est ramené à la propriété générale
démontrée dans (M, V, 4.1).
\end{env}

\begin{env}[11.5.6]
\label{0.11.5.6-fr}
De même, les résultats de (\hyperref[0.11.4.7-fr]{11.4.7}) et
(\hyperref[0.11.5.3-fr]{11.5.3}) se généralisent comme suit. Supposons (sous
les hypothèses de (\hyperref[0.11.5.5-fr]{11.5.5})) que l'on ait deux
bicomplexes $L^{\bullet\bullet}=(L^{ij})$,
$L'^{\bullet\bullet}=(L'^{ij})$ tels que $L^{ij}=0$ et $L'^{ij}=0$ pour
$j<0$, que pour tout $i$, $L^{i,\bullet}$ soit une résolution de $K^i$ et
$L'^{i,\bullet}$ une résolution de $K'^i$, et enfin que
$\mathrm R^nT(L^{ij},L'^{hk})=0$ pour $n>0$ et pour tout système d'indices
$(i,j,h,k)$. Alors on a un isomorphisme fonctoriel en $K^\bullet$ et
$K'^\bullet$
\[
  \label{0.11.5.6.1-fr}
  \mathrm R^\bullet T(K^\bullet,K'^\bullet)\xrightarrow{\sim}
  \mathrm H^\bullet(T(L^{\bullet\bullet},L'^{\bullet\bullet})).
  \tag{11.5.6.1}
\]
Cela s'établit comme dans (\hyperref[0.11.5.3-fr]{11.5.3}) en se ramenant au
cas où $K^\bullet$ et $K'^\bullet$ sont limités inférieurement.

Supposons en outre que pour tout couple $(i,j)$ et pour tout couple $(h,k)$,
les foncteurs $A\mapsto T(A,L'^{hk})$ et $A'\mapsto T(L^{ij},A')$ soient
exacts dans $\mathcal C$ et $\mathcal C'$ respectivement. Alors on a aussi des
isomorphismes fonctoriels
\[
  \label{0.11.5.6.2-fr}
  \mathrm R^\bullet T(K^\bullet,K'^\bullet)
  \xrightarrow{\sim}\mathrm H^\bullet(T(L^{\bullet\bullet},K'^\bullet))
  \xrightarrow{\sim}\mathrm H^\bullet(T(K^\bullet,L'^{\bullet\bullet})).
  \tag{11.5.6.2}
\]
La démonstration est semblable à celle de
(\hyperref[0.11.4.7-fr]{11.4.7}).
\end{env}

\begin{env}[11.5.7]
\label{0.11.5.7-fr}
On notera encore que les résultats de (\hyperref[0.11.5.5-fr]{11.5.5}) et
(\hyperref[0.11.5.6-fr]{11.5.6}) sont valables sans supposer que $T$ permute
aux limites inductives, pourvu que l'on se restreigne aux complexes
$K^\bullet$, $K'^\bullet$ limités inférieurement, ou que l'on suppose que tout
objet de $\mathcal C$ (resp. $\mathcal C'$) admet une résolution injective de
longueur bornée, et que dans (\hyperref[0.11.5.6-fr]{11.5.6}) les bicomplexes
$L^{\bullet\bullet}$ et $L'^{\bullet\bullet}$ ont leur second degré limité
supérieurement.
\end{env}

\subsection{Hyperhomologie d'un foncteur par rapport à un complexe
$K_\bullet$}
\label{subsection:0.11.6-fr}

\begin{env}[11.6.1]
\label{0.11.6.1-fr}
Soient $\mathcal C$ une catégorie abélienne,
$K_\bullet=(K_i)_{i\in\mathbb Z}$ un complexe d'objets de $\mathcal C$, dont
l'opérateur de dérivation est de degré $-1$. Une \emph{résolution de
Cartan-Eilenberg gauche} de $K_\bullet$ est formée d'un bicomplexe
$L_{\bullet\bullet}=(L_{ij})$ à opérateurs de dérivation de degré $-1$ avec
$L_{ij}=0$ pour $j<0$, et d'un morphisme de complexes simples
$\varepsilon:L_{\bullet,0}\to K_\bullet$, de façon que les conditions obtenues
à partir de celles de (\hyperref[0.11.4.2-fr]{11.4.2}) par «~renversement des
flèches~» soient remplies. Si tout objet de $\mathcal C$ est quotient d'un
objet projectif, tout complexe $K_\bullet$ de $\mathcal C$ admet une
résolution de Cartan-Eilenberg \emph{projective}, c'est-à-dire formée d'objets
projectifs $L_{ij}$, avec des propriétés fonctorielles semblables à celles de
(\hyperref[0.11.4.2-fr]{11.4.2}). En outre, si $K_\bullet$ est limité
inférieurement (resp. supérieurement), on peut supposer que $L_{ij}=0$ pour
$i<i_0$ (resp. $i>i_0$) si $K_i=0$ pour $i<i_0$ (resp. $i>i_0$). Si tout
objet de $\mathcal C$ admet une résolution projective de longueur $\leq n$,
on peut supposer que $L_{ij}=0$ pour $j>n$.
\end{env}

\begin{env}[11.6.2]
\label{0.11.6.2-fr}
\oldpage[0\textsubscript{III}]{40}
Supposons que $T$ soit un foncteur covariant additif de $\mathcal C$ dans une
catégorie abélienne $\mathcal C'$. La définition de
\emph{l'hyperhomologie} $L_\bullet T(K_\bullet)$ et des \emph{suites
spectrales d'hyperhomologie} de $T$ par rapport à un complexe $K_\bullet$ de
$\mathcal C$ (lorsqu'elles existent) se fait encore à partir de
(\hyperref[0.11.4.3-fr]{11.4.3}) par «~renversement des flèches~», les termes
$E_2$ des deux suites spectrales ainsi obtenues étant
\[
\label{0.11.6.2.1-fr}
\tag{11.6.2.1}
{}'E^2_{pq}=H_p(L_qT(K_\bullet))
\]
\[
\label{0.11.6.2.2-fr}
\tag{11.6.2.2}
{}''E^2_{pq}=L_pT(H_q(K_\bullet)),
\]
où $L_pT$ désigne le $p$-ième dérivé gauche de $T$ pour $p\geq 0$, et $0$
pour $p<0$, $L_qT(K_\bullet)$ le complexe
$(L_qT(K_i))_{i\in\mathbb Z}$.

Les propriétés de l'hyperhomologie ne se déduisent pas toutes par simple
«~renversement des flèches~» de celles de l'hypercohomologie (à moins qu'on
ne fasse des hypothèses supplémentaires du type (AB~5*) de (T, 1.5) sur la
catégorie $\mathcal C'$) en raison des conditions de régularité des deux
suites spectrales précédentes, auxquelles il faut cette fois appliquer les
critères de (\hyperref[0.11.3.4-fr]{11.3.4}). Ces derniers montrent que
lorsqu'on suppose que dans $\mathcal C'$ les limites inductives filtrantes
existent et sont exactes, alors le complexe défini par le bicomplexe
$T(L_{\bullet\bullet})$ existe, et la seconde suite spectrale
$\,{}''E(T(L_{\bullet\bullet}))$ est cette fois régulière. Si l'on suppose,
soit que $K_\bullet$ soit limité inférieurement, soit qu'il existe un entier
$n$ tel que tout objet de $\mathcal C$ admette une résolution projective de
longueur $\leq n$, alors les deux suites spectrales d'hyperhomologie existent
(sans hypothèse sur $\mathcal C'$) et sont birégulières.
\end{env}

\begin{proposition}[11.6.3]
\label{0.11.6.3-fr}
Soient $\mathcal C$, $\mathcal C'$ deux catégories abéliennes, $T$ un foncteur
covariant additif de $\mathcal C$ dans $\mathcal C'$. Alors :
\begin{enumerate}
\item[{\rm(i)}] \emph{L'hyperhomologie $L_\bullet T(K_\bullet)$ est un
foncteur homologique dans la catégorie abélienne des complexes de
$\mathcal C$ limités inférieurement.}
\item[{\rm(ii)}] \emph{Soit $K_\bullet$ un complexe de $\mathcal C$ limité
inférieurement. Si $L_nT(K_i)=0$ pour tout $n>0$ et tout $i\in\mathbb Z$, on
a des isomorphismes fonctoriels}
\[
\label{0.11.6.3.1-fr}
\tag{11.6.3.1}
L_iT(K_\bullet)\simeq H_i(T(K_\bullet))\qquad(i\in\mathbb Z).
\]
\item[{\rm(iii)}] \emph{Soit $K_\bullet$ un complexe de $\mathcal C$ limité
inférieurement. Soit $L_{\bullet\bullet}=(L_{ij})$ un bicomplexe tel que
$L_{ij}=0$ pour $j<0$ et que, pour tout $i$, $L_{i,\bullet}$ soit une
résolution de $K_i$; supposons enfin que $L_nT(L_{ij})=0$ pour tout couple
$(i,j)$ et tout $n>0$. Alors il existe un isomorphisme fonctoriel}
\[
\label{0.11.6.3.2-fr}
\tag{11.6.3.2}
L_\bullet T(K_\bullet)\simeq H_\bullet(T(L_{\bullet\bullet})).
\]
\end{enumerate}
Les démonstrations se font comme celles de
(\hyperref[0.11.5.2-fr]{11.5.2}), (\hyperref[0.11.4.5-fr]{11.4.5}) et
(\hyperref[0.11.5.3-fr]{11.5.3}) dans le cas des complexes limités
inférieurement. Nous laissons les détails de ces raisonnements au lecteur.
\end{proposition}

\begin{env}[11.6.4]
\label{0.11.6.4-fr}
On a des résultats tout à fait analogues pour les multifoncteurs covariants
additifs en chacun des arguments. Par exemple pour un bifoncteur $T$, on a
les deux suites spectrales d'hyperhomologie de termes $E_2$ donnés par
\oldpage[0\textsubscript{III}]{41}
\[
\label{0.11.6.4.1-fr}
\tag{11.6.4.1}
{}'E^2_{pq}=H_p(L_qT(K_\bullet,K'_\bullet))
\]
\[
\label{0.11.6.4.2-fr}
\tag{11.6.4.2}
{}''E^2_{pq}=
\sum_{q'+q''=q}L_pT(H_{q'}(K_\bullet),H_{q''}(K'_\bullet)).
\]
Ici encore, c'est la seconde suite spectrale qui est régulière, les deux
suites étant birégulières lorsqu'il s'agit de complexes $K_\bullet$,
$K'_\bullet$ limités inférieurement, ou quand les objets des catégories
abéliennes que l'on considère ont des résolutions projectives de longueur
fixée.

En outre :
\end{env}

\begin{proposition}[11.6.5]
\label{0.11.6.5-fr}
Soient $\mathcal C$, $\mathcal C'$, $\mathcal C''$ trois catégories
abéliennes, $T$ un bifoncteur covariant biadditif de
$\mathcal C\times\mathcal C'$ dans $\mathcal C''$.
\begin{enumerate}
\item[{\rm(i)}] \emph{$L_\bullet T(K_\bullet,K'_\bullet)$ est un bifoncteur
homologique en chacun des complexes limités inférieurement $K_\bullet$,
$K'_\bullet$ (formés respectivement d'objets de $\mathcal C$ et d'objets de
$\mathcal C'$).}
\item[{\rm(ii)}] \emph{Supposons $K_\bullet$ et $K'_\bullet$ limités
inférieurement. Soient $L_{\bullet\bullet}=(L_{ij})$,
$L'_{\bullet\bullet}=(L'_{ij})$ deux bicomplexes tels que $L_{ij}=0$ et
$L'_{ij}=0$ pour $j<0$, que pour tout $i$, $L_{i,\bullet}$ soit une
résolution de $K_i$ et $L'_{i,\bullet}$ une résolution de $K'_i$ et enfin que
$L_nT(L_{ij},L'_{hk})=0$ pour $n>0$ et tout système $(i,j,h,k)$. Alors on a
un isomorphisme fonctoriel}
\[
\label{0.11.6.5.1-fr}
\tag{11.6.5.1}
L_\bullet T(K_\bullet,K'_\bullet)
\simeq H_\bullet(T(L_{\bullet\bullet},L'_{\bullet\bullet})).
\]
\item[{\rm(iii)}] \emph{Supposons en outre que pour tout couple $(i,j)$ et
tout couple $(h,k)$, les foncteurs
\[
A\longmapsto T(A,L'_{hk}),\qquad
A'\longmapsto T(L_{ij},A')
\]
soient exacts dans $\mathcal C$ et $\mathcal C'$ respectivement. Alors on a
des isomorphismes fonctoriels}
\[
\label{0.11.6.5.2-fr}
\tag{11.6.5.2}
L_\bullet T(K_\bullet,K'_\bullet)
\simeq H_\bullet(T(L_{\bullet\bullet},K'_\bullet))
\simeq H_\bullet(T(K_\bullet,L'_{\bullet\bullet})).
\]
\end{enumerate}
Les démonstrations sont analogues à celles de
(\hyperref[0.11.5.5-fr]{11.5.5}) et (\hyperref[0.11.5.6-fr]{11.5.6}).
\end{proposition}

\subsection{Hyperhomologie d'un foncteur par rapport à un bicomplexe
$K_{\bullet\bullet}$}
\label{subsection:0.11.7-fr}

\begin{env}[11.7.1]
\label{0.11.7.1-fr}
Soit $\mathcal C$ une catégorie abélienne dans laquelle tout objet est
quotient d'un objet projectif. Considérons un bicomplexe
$K_{\bullet\bullet}=(K_{ij})$ formé d'objets de $\mathcal C$, et dont les deux
degrés sont limités inférieurement; on peut toujours se borner au cas où
$K_{ij}=0$ pour $i<0$ ou $j<0$, et c'est ce que nous ferons désormais. On
peut considérer $K_{\bullet\bullet}$ comme un complexe (simple) formé
d'objets $K_{i,\bullet}=(K_{ij})_{j\geq 0}$ de la catégorie abélienne
$\mathcal K$ des complexes à degrés positifs d'objets de $\mathcal C$. Il
résulte du lemme (\hyperref[0.11.5.2.1-fr]{11.5.2.1}) (ou plutôt du lemme
«~dual~» obtenu par «~renversement des flèches~») et de (M, V, 2.2) que
$K_{\bullet\bullet}$ admet une \emph{résolution projective de
Cartan-Eilenberg} dans la catégorie $\mathcal K$; une telle résolution est un
tricomplexe $M_{\bullet\bullet\bullet}=(M_{ijk})$ de $\mathcal C$, à degrés
tous $\geq 0$, formé d'objets projectifs, tel que pour tout $i$,
$M_{i,\bullet\bullet}$, $\mathrm B_i^{\mathrm I}(M_{\bullet\bullet\bullet})$,
$\mathrm Z_i^{\mathrm I}(M_{\bullet\bullet\bullet})$,
$H_i^{\mathrm I}(M_{\bullet\bullet\bullet})$ constituent des résolutions
projectives de $K_{i,\bullet}$,
$\mathrm B_i^{\mathrm I}(K_{\bullet\bullet})$,
$\mathrm Z_i^{\mathrm I}(K_{\bullet\bullet})$,
$H_i^{\mathrm I}(K_{\bullet\bullet})$ respectivement dans la catégorie
$\mathcal K$; en particulier, pour tout couple $(i,j)$,
$M_{ij,\bullet}$ est une résolution projective de $K_{ij}$ dans
$\mathcal C$.
\end{env}

\begin{proposition}[11.7.2]
\label{0.11.7.2-fr}
Soit $T$ un foncteur covariant additif de $\mathcal C$ dans une catégorie
abélienne $\mathcal C'$. Avec les notations de
(\hyperref[0.11.7.1-fr]{11.7.1}), \emph{l'homologie
$H_\bullet(T(M_{\bullet\bullet\bullet}))$ du complexe simple associé au
tricomplexe $T(M_{\bullet\bullet\bullet})$ est canoniquement isomorphe à
l'hyperhomologie $L_\bullet T(K_{\bullet\bullet})$ du complexe
\oldpage[0\textsubscript{III}]{42}
simple associé à $K_{\bullet\bullet}$
(\hyperref[0.11.6.2-fr]{11.6.2}), et est
l'aboutissement commun de six suites spectrales birégulières notées
$\mathop{}^{(t)}E$ (avec $t=a,b,a',b',c$ ou $d$), dont les termes $E_2$ sont
donnés par les formules}
\[
\begin{aligned}
{}^{(a)}E^2_{pq}&=L_pT(H_q(K_{\bullet\bullet})),&
{}^{(b)}E^2_{pq}&=H_p(L_q^{\mathrm{II}}T(K_{\bullet\bullet})),\\
{}^{(a')}E^2_{pq}&=L_pT(H_q^{\mathrm I}(K_{\bullet\bullet})),&
{}^{(b')}E^2_{pq}&=H_p(L_qT(K_{\bullet\bullet})),\\
{}^{(c)}E^2_{pq}&=L_pT(H_q^{\mathrm{II}}(K_{\bullet\bullet})),&
{}^{(d)}E^2_{pq}&=H_p(L_q^{\mathrm I}T(K_{\bullet\bullet})).
\end{aligned}
\]
\end{proposition}

On rappelle que nous utilisons la notation $F(A_\bullet)$ pour désigner le
complexe des objets $F(A_i)$ pour tout complexe $A_\bullet=(A_i)$; par
exemple $L_q^{\mathrm{II}}T(K_{\bullet\bullet})$ désigne le complexe
$(L_q^{\mathrm{II}}T(K_{i,\bullet}))_{i\geq0}$, où
$L_q^{\mathrm{II}}T(K_{i,\bullet})$ est l'hyperhomologie d'indice $q$ du
foncteur $T$ par rapport au complexe simple $K_{i,\bullet}$.

\begin{proof}
Désignons par $L_\bullet$ le complexe simple associé à
$K_{\bullet\bullet}$, de sorte que
\[
L_i=\bigoplus_{r+s=i}K_{rs},
\]
et posons
\[
N_{ij}=\bigoplus_{r+s=i}M_{rsj};
\]
il est clair que pour chaque $i$, $N_{i,\bullet}$ est une résolution
projective de $L_i$ dans $\mathcal C$; il résulte donc de
(\hyperref[0.11.6.3-fr]{11.6.3}) et
(\hyperref[0.11.6.4-fr]{11.6.4}) que l'on a un isomorphisme fonctoriel
\[
L_\bullet T(L_\bullet)\simeq H_\bullet(T(N_{\bullet\bullet}));
\]
comme le complexe simple associé au bicomplexe $T(N_{\bullet\bullet})$ est
aussi associé au tricomplexe $T(M_{\bullet\bullet\bullet})$, cela prouve la
première assertion de l'énoncé.

En outre, $L_\bullet T(L_\bullet)$ est l'aboutissement des deux suites
spectrales d'hyperhomologie (\hyperref[0.11.6.2-fr]{11.6.2}) de $T$
relativement au complexe simple $L_\bullet$, qui ne sont autres que les suites
$\mathop{}^{(b')}E$ et $\mathop{}^{(a)}E$ respectivement.

Considérons maintenant $M_{\bullet\bullet\bullet}$ comme un bicomplexe
$U_{\bullet\bullet}$ avec
\[
U_{ij}=\bigoplus_{r+s=j}M_{irs};
\]
$H_\bullet(T(M_{\bullet\bullet\bullet}))$ est encore l'aboutissement des deux
suites spectrales du bicomplexe $T(U_{\bullet\bullet})$. Or, pour tout $i$,
$M_{i,\bullet\bullet}$ est un bicomplexe vérifiant les conditions de
(\hyperref[0.11.6.3-fr]{11.6.3}) relativement au complexe simple
$K_{i,\bullet}$; donc
$H_q^{\mathrm{II}}(T(U_{i,\bullet}))=
L_q^{\mathrm{II}}T(K_{i,\bullet})$, et la première suite spectrale de
$T(U_{\bullet\bullet})$ n'est autre que $\mathop{}^{(b)}E$. D'autre part, pour
tout $r$, $M_{\bullet,r,\bullet}$ est une résolution de Cartan-Eilenberg du
complexe simple $K_{\bullet,r}$, le calcul fait dans (M, XV, 2) montre que
\[
H_q^{\mathrm I}(T(M_{\bullet,rs}))
=T(H_q^{\mathrm I}(M_{\bullet,rs})),
\]
d'où
\[
H_q^{\mathrm I}(T(U_{\bullet,j}))
=\bigoplus_{r+s=j}T(H_q^{\mathrm I}(M_{\bullet,rs}));
\]
autrement dit, le complexe simple
$H_q^{\mathrm I}(T(U_{\bullet\bullet}))$ n'est autre que le complexe simple
associé au bicomplexe
$T(H_q^{\mathrm I}(M_{\bullet\bullet\bullet}))$. Or, pour tout $q$,
$H_q^{\mathrm I}(M_{\bullet\bullet\bullet})$ est une résolution projective du
complexe simple $H_q^{\mathrm I}(K_{\bullet\bullet})$ dans la catégorie
$\mathcal K$; appliquant (\hyperref[0.11.6.3-fr]{11.6.3}), on voit que l'on a
\[
H_p^{\mathrm{II}}
\bigl(T(H_q^{\mathrm I}(M_{\bullet\bullet\bullet}))\bigr)
=L_pT(H_q^{\mathrm I}(K_{\bullet\bullet})),
\]
et on obtient ainsi la suite $\mathop{}^{(a')}E$. Enfin, les suites
$\mathop{}^{(c)}E$ et $\mathop{}^{(d)}E$ s'obtiennent en intervertissant les
rôles des indices $i$ et $j$ dans la définition du tricomplexe
$M_{\bullet\bullet\bullet}$ et en appliquant à ce nouveau tricomplexe les
raisonnements qui précèdent.
\end{proof}

On dit que $L_\bullet T(K_{\bullet\bullet})$ est \emph{l'hyperhomologie} de
$T$ relative au bicomplexe $K_{\bullet\bullet}$.

\begin{env}[11.7.3]
\label{0.11.7.3-fr}
\emph{Remarques}.
\begin{enumerate}
\item[{\rm(i)}] Il résulte de (\hyperref[0.11.6.3-fr]{11.6.3}) que
$L_\bullet T(K_{\bullet\bullet})$ est un foncteur homologique dans la
catégorie des bicomplexes $K_{\bullet\bullet}$ de $\mathcal C$ limités
inférieurement en chacun de leurs degrés.
\oldpage[0\textsubscript{III}]{43}
\item[{\rm(ii)}] Soit $M_{\bullet\bullet\bullet}$ un tricomplexe de
$\mathcal C$ tel que pour chaque couple $(r,s)$,
$M_{rs,\bullet}$ soit une résolution de $K_{rs}$ et que
$L_nT(M_{ijk})=0$ pour tous les triplets $(i,j,k)$ et tout $n>0$. Alors on a
un isomorphisme
\[
L_\bullet T(K_{\bullet\bullet})
\simeq H_\bullet(T(M_{\bullet\bullet\bullet}));
\]
en effet, avec les notations de la démonstration de
(\hyperref[0.11.7.2-fr]{11.7.2}), $N_{i,\bullet}$ est une résolution de
$L_i$ telle que $L_nT(N_{ij})=0$ pour $n>0$ et tout couple $(i,j)$, et il
suffit d'appliquer (\hyperref[0.11.6.3-fr]{11.6.3}, (iii)).
\item[{\rm(iii)}] On généralise aussitôt les résultats de
(\hyperref[0.11.7.2-fr]{11.7.2}) aux multifoncteurs covariants; par exemple,
soient $\mathcal C'$ une seconde catégorie abélienne dans laquelle tout objet
est quotient d'un objet projectif, $K'_{\bullet\bullet}$ un bicomplexe de
$\mathcal C'$ dont les deux degrés sont limités inférieurement, et $T$ un
bifoncteur covariant additif de $\mathcal C\times\mathcal C'$ dans une
catégorie abélienne $\mathcal C''$. Si $L_\bullet$ et $L'_\bullet$ sont les
complexes simples associés à $K_{\bullet\bullet}$ et
$K'_{\bullet\bullet}$ respectivement, on définit l'hyperhomologie de $T$ par
rapport aux deux bicomplexes $K_{\bullet\bullet}$,
$K'_{\bullet\bullet}$ comme l'hyperhomologie
$L_\bullet T(L_\bullet,L'_\bullet)$; appliquant
(\hyperref[0.11.6.4-fr]{11.6.4}) et
(\hyperref[0.11.6.5-fr]{11.6.5}), on a de même que dans
(\hyperref[0.11.7.2-fr]{11.7.2}) six suites spectrales aboutissant à cette
hyperhomologie, et que nous laissons le soin d'écrire au lecteur.
\end{enumerate}
\end{env}

\subsection{Compléments sur la cohomologie des complexes simpliciaux}
\label{subsection:0.11.8-fr}

\begin{env}[11.8.1]
\label{0.11.8.1-fr}
Soient $A$ un ensemble fini, $\Sigma(A)$ l'ensemble des suites finies
$\sigma=(\alpha_0,\ldots,\alpha_h)$ d'éléments de $A$ («~simplexes~» de
$A$); on pose $|\sigma|=\{\alpha_0,\ldots,\alpha_h\}$; on rappelle que le
complexe de chaînes $\mathrm C_\bullet(A)$ est le groupe abélien libre
gradué engendré par les éléments de $\Sigma(A)$,
$(\alpha_0,\ldots,\alpha_h)$ étant de degré $h$, avec une différentielle
définie par
\[
d(\alpha_0,\ldots,\alpha_h)=
\sum_{j=0}^{h}(-1)^j
(\alpha_0,\ldots,\widehat{\alpha_j},\ldots,\alpha_h).
\]
Le sous-groupe $\mathrm D_\bullet(A)$ de $\mathrm C_\bullet(A)$ engendré par
les chaînes $\sigma=(\alpha_0,\ldots,\alpha_h)$ pour lesquelles deux des
$\alpha_i$ sont égaux, et par les chaînes
$\pi(\sigma)-\varepsilon_\pi\sigma$, où pour toute permutation $\pi$,
\[
\pi(\sigma)=(\alpha_{\pi(0)},\ldots,\alpha_{\pi(h)})
\]
et $\varepsilon_\pi$ est la signature de $\pi$, est un sous-complexe de
$\mathrm C_\bullet(A)$ dont les éléments sont dits \emph{chaînes
dégénérées}; on pose
$\mathrm L_\bullet(A)=\mathrm C_\bullet(A)/\mathrm D_\bullet(A)$, et on a un
homomorphisme naturel de complexes
\[
p:\mathrm C_\bullet(A)\longrightarrow\mathrm L_\bullet(A).
\]
On définit d'autre part un homomorphisme de complexes
\[
j:\mathrm L_\bullet(A)\longrightarrow\mathrm C_\bullet(A)
\]
de la façon suivante : on ordonne totalement $A$; à la classe
mod.~$\mathrm D_\bullet(A)$ d'un simplexe
$\sigma=(\alpha_0,\ldots,\alpha_h)$, on fait correspondre $0$ si deux des
$\alpha_i$ sont égaux, et la suite $(\beta_0,\ldots,\beta_h)$ des $\alpha_i$
rangés dans l'ordre croissant dans le cas contraire. Il est clair que $p\circ
j$ est l'identité de $\mathrm L_\bullet(A)$.
\end{env}

\begin{env}[11.8.2]
\label{0.11.8.2-fr}
Soit $B$ un second ensemble fini; si $d'$, $d''$ sont les différentielles de
$\mathrm C_\bullet(A)$ et $\mathrm C_\bullet(B)$, le complexe produit
tensoriel $\mathrm C_\bullet(A)\otimes\mathrm C_\bullet(B)$ peut être
considéré comme le groupe abélien libre engendré par les éléments de
$\Sigma(A)\times\Sigma(B)$, avec la différentielle
\[
d(\sigma,\tau)=(d'\sigma,\tau)+(-1)^{h+1}(\sigma,d''\tau)
\quad\text{si }\operatorname{Card}|\sigma|=h+1.
\]
Les homomorphismes naturels
$\mathrm C_\bullet(A)\to\mathrm L_\bullet(A)$,
$\mathrm C_\bullet(B)\to\mathrm L_\bullet(B)$ définissent un homomorphisme
\[
p:\mathrm C_\bullet(A)\otimes\mathrm C_\bullet(B)
\longrightarrow\mathrm L_\bullet(A)\otimes\mathrm L_\bullet(B),
\]
ce dernier produit tensoriel étant isomorphe à
\[
\frac{\mathrm C_\bullet(A)\otimes\mathrm C_\bullet(B)}
{\mathrm C_\bullet(A)\otimes\mathrm D_\bullet(B)
+\mathrm D_\bullet(A)\otimes\mathrm C_\bullet(B)}.
\]
De même, à l'aide des homomorphismes
$\mathrm L_\bullet(A)\to\mathrm C_\bullet(A)$ et
$\mathrm L_\bullet(B)\to\mathrm C_\bullet(B)$ définis dans
(\hyperref[0.11.8.1-fr]{11.8.1}) (à l'aide d'ordres totaux sur $A$ et $B$),
on définit un homomorphisme
\[
j:\mathrm L_\bullet(A)\otimes\mathrm L_\bullet(B)
\longrightarrow\mathrm C_\bullet(A)\otimes\mathrm C_\bullet(B)
\]
tel que $p\circ j$ soit l'identité.
\end{env}

\oldpage[0\textsubscript{III}]{44}
Avec ces notations :
\begin{proposition}[11.8.3]
\label{0.11.8.3-fr}
\emph{Il existe une homotopie
\[
h:\mathrm C_\bullet(A)\otimes\mathrm C_\bullet(B)
\longrightarrow\mathrm C_\bullet(A)\otimes\mathrm C_\bullet(B),
\]
telle que $h(\sigma,\tau)$ soit combinaison linéaire de couples de simplexes
$(\sigma_i,\tau_i)$ avec $|\sigma_i|\subset|\sigma|$,
$|\tau_i|\subset|\tau|$, et que l'on ait pour $f=j\circ p$,}
\[
\label{0.11.8.3.1-fr}
\tag{11.8.3.1}
f-1=h\circ d+d\circ h.
\]
\end{proposition}

\begin{proof}
Il suffit de définir $h$ sur chaque couple $(\sigma,\tau)$ de simplexes, en
raisonnant par récurrence sur la somme des degrés de $\sigma$ et $\tau$, car
on peut prendre $h=0$ lorsque cette somme est $0$. Soit
\[
\omega=f(\sigma,\tau)-(\sigma,\tau)-h(d(\sigma,\tau));
\]
en vertu de l'hypothèse de récurrence et de la définition de $d$, on a
\[
\omega\in\mathrm C_\bullet(|\sigma|)\otimes\mathrm C_\bullet(|\tau|).
\]
On a
\[
d\omega=f(d(\sigma,\tau))-d(\sigma,\tau)-d(h(d(\sigma,\tau)))
=h(d(d(\sigma,\tau)))=0
\]
en vertu de (\hyperref[0.11.8.3.1-fr]{11.8.3.1}) et de l'hypothèse de
récurrence. Or, on a
$H_q(\mathrm C_\bullet(A))=0$ pour $q>0$ (G, I, 3.7.4), donc aussi
\[
H_q(\mathrm C_\bullet(A)\otimes\mathrm C_\bullet(B))=0\quad(q>0),
\]
en vertu de la formule de Künneth (G, I, 5.5.2). Appliquant cette remarque en
remplaçant $A$ par $|\sigma|$ et $B$ par $|\tau|$, on voit qu'il existe un
élément
\[
\omega'\in\mathrm C_\bullet(|\sigma|)\otimes\mathrm C_\bullet(|\tau|)
\]
tel que $\omega=d\omega'$; prenant $h(\sigma,\tau)=\omega'$, on vérifie
(\hyperref[0.11.8.3.1-fr]{11.8.3.1}) pour le couple $(\sigma,\tau)$ et la
récurrence peut se poursuivre.
\end{proof}

\begin{env}[11.8.4]
\label{0.11.8.4-fr}
Les notations étant celles de (\hyperref[0.11.8.2-fr]{11.8.2}), nous poserons
$(\sigma,\tau)\leq(\sigma',\tau')$ si $|\sigma|\subset|\sigma'|$ et
$|\tau|\subset|\tau'|$. Un \emph{système de coefficients} $\Gamma$ sur
$\Sigma(A)\times\Sigma(B)$ est constitué par une famille
$(\Gamma_{\sigma,\tau})$ de groupes abéliens, où
$\Gamma_{\sigma,\tau}$ ne dépend que des ensembles $|\sigma|$ et $|\tau|$,
et une famille d'homomorphismes
\[
\Gamma_{\sigma,\tau}\longrightarrow\Gamma_{\sigma',\tau'}
\quad\text{pour }(\sigma,\tau)\leq(\sigma',\tau'),
\]
formant un système inductif pour cette relation de préordre. On définit alors
un complexe de cochaînes $\mathrm C^\bullet(A,B;\Gamma)$ comme l'ensemble
des familles $\lambda=(\lambda(\sigma,\tau))$, où $(\sigma,\tau)$ parcourt
$\Sigma(A)\times\Sigma(B)$, avec
$\lambda(\sigma,\tau)\in\Gamma_{\sigma,\tau}$ pour tout couple
$(\sigma,\tau)$. La différentielle est donnée de la façon suivante : si
\[
d(\sigma,\tau)=\sum_i\pm(\sigma_i,\tau_i),
\]
on a $|\sigma_i|\subset|\sigma|$, $|\tau_i|\subset|\tau|$ pour tout $i$, et
l'on prend
\[
d\lambda(\sigma,\tau)=\sum_i\pm\lambda_i(\sigma_i,\tau_i),
\]
où $\lambda_i(\sigma_i,\tau_i)$ désigne l'image canonique de
$\lambda(\sigma_i,\tau_i)$ dans $\Gamma_{\sigma,\tau}$.

Nous dirons qu'une cochaîne $\lambda\in\mathrm C^\bullet(A,B;\Gamma)$ est
\emph{bi-alternée} si $\lambda(\sigma,\tau)=0$ lorsque l'un des deux
simplexes $\sigma$, $\tau$ a deux termes égaux, et si l'on a
\[
\lambda(\pi(\sigma),\tau)=\varepsilon_\pi\lambda(\sigma,\tau),
\qquad
\lambda(\sigma,\pi'(\tau))=\varepsilon_{\pi'}\lambda(\sigma,\tau)
\]
pour des permutations quelconques $\pi$, $\pi'$ des indices. Il est clair
que ces cochaînes engendrent un sous-complexe
$\mathrm L^\bullet(A,B;\Gamma)$ de $\mathrm C^\bullet(A,B;\Gamma)$.
\end{env}

\begin{proposition}[11.8.5]
\label{0.11.8.5-fr}
\emph{L'injection canonique
\[
\mathrm L^\bullet(A,B;\Gamma)\longrightarrow
\mathrm C^\bullet(A,B;\Gamma)
\]
définit un isomorphisme pour la cohomologie de ces deux complexes.}
\end{proposition}

\begin{proof}
Notons que si $p$ et $j$ ont le sens défini dans
(\hyperref[0.11.8.2-fr]{11.8.2}), les applications
\[
{}'p:\lambda\longmapsto\lambda\circ p,\qquad
{}'j:\lambda\longmapsto\lambda\circ j
\]
sont définies dans $\mathrm L^\bullet(A,B;\Gamma)$ et
$\mathrm C^\bullet(A,B;\Gamma)$ respectivement, la première n'étant autre
que l'injection canonique. Comme $\mathop{}'j\circ{}'p$ est l'identité, il
suffit de démontrer que $\mathop{}'p\circ{}'j$ est homotope à l'identité; or,
d'après (\hyperref[0.11.8.3-fr]{11.8.3}),
$\mathop{}'h:\lambda\mapsto\lambda\circ h$ est définie dans
$\mathrm C^\bullet(A,B;\Gamma)$ et on peut donc transposer l'identité
(\hyperref[0.11.8.3.1-fr]{11.8.3.1}), qui fournit le résultat cherché.
\end{proof}

\begin{env}[11.8.6]
\label{0.11.8.6-fr}
\oldpage[0\textsubscript{III}]{45}
La proposition (\hyperref[0.11.8.5-fr]{11.8.5}) ramène le calcul de la
cohomologie de $\mathrm L^\bullet(A,B;\Gamma)$ à celui de la cohomologie de
$\mathrm C^\bullet(A,B;\Gamma)$. Rappelons d'autre part que cette dernière
est, en vertu du th. d'Eilenberg-Zilber (G, I, 3.10.2), canoniquement
isomorphe à la cohomologie du complexe de cochaînes défini comme suit : on
forme le complexe de chaînes $\mathrm P_\bullet(A,B)$, constitué par les
combinaisons linéaires des $(\sigma,\tau)\in\Sigma(A)\times\Sigma(B)$ tels
que $\sigma$ et $\tau$ aient le même degré; la différentielle de ce complexe
est donnée par
\[
d(\sigma,\tau)=\sum_{j,k}(-1)^{j+k}(\sigma_j,\tau_k)
\]
si
\[
d\sigma=\sum_j(-1)^j\sigma_j,\qquad
d\tau=\sum_k(-1)^k\tau_k;
\]
on a alors deux homomorphismes canoniques de complexes
\[
f:\mathrm P_\bullet(A,B)\longrightarrow
\mathrm C_\bullet(A)\otimes\mathrm C_\bullet(B),\qquad
g:\mathrm C_\bullet(A)\otimes\mathrm C_\bullet(B)
\longrightarrow\mathrm P_\bullet(A,B),
\]
et on démontre (\emph{loc. cit.}) qu'il y a des homotopies $h$, $h'$ telles
que
\[
f\circ g-1=d\circ h+h\circ d,\qquad
g\circ f-1=d\circ h'+h'\circ d.
\]
En outre, on a
\[
f(\sigma,\tau)\in\mathrm C_\bullet(|\sigma|)
\otimes\mathrm C_\bullet(|\tau|),\qquad
g(\sigma,\tau)\in\mathrm P_\bullet(|\sigma|,|\tau|)
\]
et les homotopies $h$, $h'$ peuvent être prises telles que
\[
h(\sigma,\tau)\in\mathrm C_\bullet(|\sigma|)
\otimes\mathrm C_\bullet(|\tau|),\qquad
h'(\sigma,\tau)\in\mathrm P_\bullet(|\sigma|,|\tau|).
\]
Ce point provient de ce que la définition de $h(\sigma,\tau)$ et
$h'(\sigma,\tau)$ peut se faire par récurrence sur la somme des degrés de
$\sigma$ et $\tau$, et de ce que les
\[
H^q(\mathrm C_\bullet(|\sigma|)\otimes\mathrm C_\bullet(|\tau|))
\quad\text{et}\quad
H^q(\mathrm P_\bullet(|\sigma|,|\tau|))
\]
sont nuls pour $q>0$ (\emph{loc. cit.}); on raisonne alors comme dans
(\hyperref[0.11.8.3-fr]{11.8.3}) et la conclusion en résulte.

On définit alors $\mathrm P^\bullet(A,B;\Gamma)$ comme l'ensemble des
familles $\lambda=(\lambda(\sigma,\tau))$ où $(\sigma,\tau)$ parcourt les
couples dont les termes ont même degré, avec
$\lambda(\sigma,\tau)\in\Gamma_{\sigma,\tau}$, et comme on a
\[
d\sigma=\sum_j(-1)^j\sigma_j,\qquad
d\tau=\sum_k(-1)^k\tau_k,
\]
\[
d\lambda(\sigma,\tau)=
\sum_{j,k}(-1)^{j+k}\lambda(\sigma_j,\tau_k)
\]
est défini et donne la différentielle du complexe
$\mathrm P^\bullet(A,B;\Gamma)$. Cela étant, les applications
\[
{}'f:\lambda\mapsto\lambda\circ f,\quad
{}'g:\lambda\mapsto\lambda\circ g,\quad
{}'h:\lambda\mapsto\lambda\circ h,\quad
{}'h':\lambda\mapsto\lambda\circ h'
\]
sont toutes définies en vertu des remarques qui précèdent;
$\mathop{}'f\circ{}'g$ et $\mathop{}'g\circ{}'f$ sont donc homotopes à
l'identité, d'où l'isomorphisme cherché entre la cohomologie de
$\mathrm C^\bullet(A,B;\Gamma)$ et celle de
$\mathrm P^\bullet(A,B;\Gamma)$.
\end{env}

\begin{env}[11.8.7]
\label{0.11.8.7-fr}
\emph{Remarque}. Le même raisonnement que dans
(\hyperref[0.11.8.3-fr]{11.8.3}), mais appliqué à
$\mathrm C_\bullet(A)$ et $\mathrm L_\bullet(A)$, montre que si $j$ et $p$
sont définis comme dans (\hyperref[0.11.8.1-fr]{11.8.1}), $f=j\circ p$
vérifie encore une relation (\hyperref[0.11.8.3.1-fr]{11.8.3.1}), avec
$|h(\sigma)|\subset|\sigma|$, d'où on déduit comme dans
(\hyperref[0.11.8.5-fr]{11.8.5}) un isomorphisme de la cohomologie de
$\mathrm L^\bullet(A;\Gamma)$ sur celle de
$\mathrm C^\bullet(A;\Gamma)$, ces deux complexes étant définis de façon
évidente. C'est le résultat dont la démonstration est esquissée dans
(G, I, 3.8.1).
\end{env}

\begin{env}[11.8.8]
\label{0.11.8.8-fr}
Reprenons maintenant les notations et hypothèses de
(\hyperref[0.11.8.2-fr]{11.8.2}), et considérons un complexe
$\Gamma^\bullet=(\Gamma^k)$ de systèmes de coefficients sur
$\Sigma(A)\times\Sigma(B)$ : pour chaque $(\sigma,\tau)$, les
$\Gamma^k_{\sigma,\tau}$ forment donc un complexe de groupes abéliens
$(k\in\mathbb Z)$, et on a les diagrammes commutatifs
\[
\xymatrix{
\Gamma^k_{\sigma,\tau}\ar[r]\ar[d] &
\Gamma^{k+1}_{\sigma,\tau}\ar[d]\\
\Gamma^k_{\sigma',\tau'}\ar[r] &
\Gamma^{k+1}_{\sigma',\tau'}
}
\]
\oldpage[0\textsubscript{III}]{46}
pour $(\sigma,\tau)\leq(\sigma',\tau')$. Alors on vérifie aussitôt que
\[
\mathrm C^\bullet(A,B;\Gamma^\bullet)
=\bigl(\mathrm C^h(A,B;\Gamma^k)\bigr)_{(h,k)\in\mathbb Z\times\mathbb Z}
\]
est un bicomplexe de groupes abéliens, et
\[
\mathrm L^\bullet(A,B;\Gamma^\bullet)
=\bigl(\mathrm L^h(A,B;\Gamma^k)\bigr)
\]
en est un sous-bicomplexe.
\end{env}

\begin{proposition}[11.8.9]
\label{0.11.8.9-fr}
\emph{L'injection canonique
\[
\mathrm L^\bullet(A,B;\Gamma^\bullet)\longrightarrow
\mathrm C^\bullet(A,B;\Gamma^\bullet)
\]
définit un isomorphisme pour la cohomologie de ces deux bicomplexes.}
\end{proposition}

\begin{proof}
Posons $\mathrm C^{\bullet\bullet}=
\mathrm C^\bullet(A,B;\Gamma^\bullet)$ et
$\mathrm L^{\bullet\bullet}=
\mathrm L^\bullet(A,B;\Gamma^\bullet)$ pour simplifier, et notons que
puisque $\mathrm C^{hk}=\mathrm L^{hk}=0$ pour $h<0$, les secondes suites
spectrales de ces bicomplexes sont régulières
(\hyperref[0.11.3.3-fr]{11.3.3}); l'homomorphisme
$\mathrm L^{\bullet\bullet}\to\mathrm C^{\bullet\bullet}$ fournit donc un
morphisme de suites spectrales
\[
{}''E(\mathrm L^{\bullet\bullet})\longrightarrow
{}''E(\mathrm C^{\bullet\bullet})
\]
qui, pour les termes $E_2$, se réduit à
\[
\label{0.11.8.9.1-fr}
\tag{11.8.9.1}
H^p_{\mathrm{II}}(H^q_{\mathrm I}(\mathrm L^{\bullet\bullet}))
\longrightarrow
H^p_{\mathrm{II}}(H^q_{\mathrm I}(\mathrm C^{\bullet\bullet})).
\]
Mais pour tout $k\in\mathbb Z$, il résulte de
(\hyperref[0.11.8.3-fr]{11.8.3}) que l'homomorphisme
$H^q_{\mathrm I}(\mathrm L^{\bullet,k})\to
H^q_{\mathrm I}(\mathrm C^{\bullet,k})$ est bijectif; la conclusion résulte
donc de (\hyperref[0.11.1.5-fr]{11.1.5}).
\end{proof}

\begin{env}[11.8.10]
\label{0.11.8.10-fr}
De même, avec les notations de (\hyperref[0.11.8.6-fr]{11.8.6}), on a des
homomorphismes canoniques de bicomplexes
\[
\mathrm C^\bullet(A,B;\Gamma^\bullet)\longrightarrow
\mathrm P^\bullet(A,B;\Gamma^\bullet)
\]
(avec des notations évidentes), et le même raisonnement que dans
(\hyperref[0.11.8.9-fr]{11.8.9}), basé cette fois sur
(\hyperref[0.11.8.6-fr]{11.8.6}), montre que cet homomorphisme donne encore
un isomorphisme en cohomologie.
\end{env}

\subsection{Un lemme sur les complexes de type fini}
\label{subsection:0.11.9-fr}

\begin{proposition}[11.9.1]
\label{0.11.9.1-fr}
Soient $\mathcal C$ une catégorie abélienne, $K'$ et $K''$ des parties de
l'ensemble des objets de $\mathcal C$, telles que $K''\subset K'$, et
vérifiant les conditions suivantes :
\begin{enumerate}
\item[{\rm(i)}] \emph{Pour tout objet $A'\in K'$ et tout épimorphisme
$u:A\to A'$ dans $\mathcal C$, il existe un objet $B\in K''$ et un morphisme
$v:B\to A$ tels que $u\circ v$ soit un épimorphisme.}
\item[{\rm(ii)}] \emph{Pour tout couple d'objets $A\in K''$, $B\in K'$ et
tout épimorphisme $u:A\to B$, $\operatorname{Ker}(u)$ appartient à $K'$.}
\item[{\rm(iii)}] \emph{Le produit de deux objets de $K''$ appartient à
$K''$.}
\end{enumerate}
\emph{Soit $P_\bullet=(P_i)_{i\in\mathbb Z}$ un complexe dans $\mathcal C$,
tel que $H_i(P_\bullet)\in K'$ pour tout $i$, et qu'il existe $d$ tel que
$H_i(P_\bullet)=0$ pour $i<d$. Alors il existe un complexe
$Q_\bullet=(Q_i)_{i\in\mathbb Z}$ dans $\mathcal C$, tel que $Q_i\in K''$
pour tout $i$ et $Q_i=0$ pour $i<d$, et un morphisme
$u:Q_\bullet\to P_\bullet$ de complexes tel que le morphisme correspondant
$H_\bullet(Q_\bullet)\to H_\bullet(P_\bullet)$ soit un isomorphisme.}
\end{proposition}

\begin{proof}
Démontrons d'abord la conséquence suivante de la propriété (i) :

\emph{(i bis) Soient $u:C\to B$ un épimorphisme dans $\mathcal C$, $A$ un
objet de $K'$, $v:A\to B$ un morphisme dans $\mathcal C$; il existe alors un
objet $D\in K''$, un épimorphisme $u':D\to A$ et un morphisme $v':D\to C$
tels que le diagramme}
\[
\xymatrix{
D\ar[r]^{u'}\ar[d]_{v'} & A\ar[d]^{v}\\
C\ar[r]_{u} & B
}
\tag{11.9.1.1}\label{0.11.9.1.1-fr}
\]
\emph{soit commutatif.}

Considérons en effet le produit fibré $C\times_B A$ dans $\mathcal C$ et les
projections canoniques $p:C\times_B A\to C$, $q:C\times_B A\to A$, rendant
commutatif le diagramme
\oldpage[0\textsubscript{III}]{47}
\[
\xymatrix{
C\times_B A\ar[r]^{q}\ar[d]_{p} & A\ar[d]^{v}\\
C\ar[r]_{u} & B
}
\tag{11.9.1.2}\label{0.11.9.1.2-fr}
\]
On sait ([27], p.~1-12) que le conoyau de $q$ est le quotient de $A$ par
$v^{-1}(u(C))$; comme $u$ est un épimorphisme, $u(C)=B$ et
$v^{-1}(u(C))=A$, donc $q$ est un épimorphisme; il suffit alors d'appliquer
(i) à l'épimorphisme $q:C\times_B A\to A$ : il y a un objet $D\in K''$ et un
morphisme $w:D\to C\times_B A$ tels que $q\circ w$ soit un épimorphisme; on
prendra $u'=q\circ w$, $v'=p\circ w$.

Cela étant, pour démontrer la proposition, procédons par récurrence.
Supposons, pour un $i\geq d-1$, construits, pour $j\leq i$, les objets $Q_j$,
les morphismes $d_j:Q_j\to Q_{j-1}$ et les morphismes $u_j:Q_j\to P_j$, de
sorte que $Q_j=0$ pour $j<d$, que $d_{j-1}\circ d_j=0$ et
$d_j\circ u_j=u_{j-1}\circ d_j$ pour $j\leq i$; en outre, nous supposons
vérifiées les conditions suivantes :
\begin{enumerate}
\item[$(\mathrm I_i)$] On a $Q_j\in K''$ pour $j\leq i$ et
$B_j(Q_\bullet)\in K'$ pour $j<i$.
\item[$(\mathrm{II}_i)$] Pour $j<i$, l'homomorphisme
$H_j(Q_\bullet)\to H_j(P_\bullet)$ déduit de la famille $(u_k)_{k\leq i}$
est un isomorphisme.
\item[$(\mathrm{III}_i)$] Le morphisme composé
\[
v_i:Z_i(Q_\bullet)\longrightarrow Z_i(P_\bullet)
\longrightarrow H_i(P_\bullet)
\]
(où la flèche de gauche est la restriction de $u_i$ et la flèche de droite
le morphisme canonique) est un épimorphisme.
\end{enumerate}
Notons que, d'après (ii), $Z_i(Q_\bullet)$, noyau de l'épimorphisme
$Q_i\to B_{i-1}(Q_\bullet)$, appartient à $K'$ en vertu de l'hypothèse
$(\mathrm I_i)$. On tire encore de (ii) que
$N_i=\operatorname{Ker}(v_i)$ appartient aussi à $K'$, compte tenu de
l'hypothèse $(\mathrm{III}_i)$. En vertu de (i bis), il existe un
$Q'_{i+1}\in K''$, un épimorphisme $d'_{i+1}:Q'_{i+1}\to N_i$ et un
morphisme $u'_{i+1}:Q'_{i+1}\to P_{i+1}$, tels que le diagramme
\[
\xymatrix{
Q'_{i+1}\ar[r]^{d'_{i+1}}\ar[d]_{u'_{i+1}} &
N_i\ar[d]^{u_i}\\
P_{i+1}\ar[r]_{d_{i+1}} & B_i(P_\bullet)
}
\tag{11.9.1.3}\label{0.11.9.1.3-fr}
\]
soit commutatif.

Comme le morphisme canonique $Z_{i+1}(P_\bullet)\to H_{i+1}(P_\bullet)$ est
un épimorphisme et que $H_{i+1}(P_\bullet)\in K'$ par hypothèse, il résulte
de (i) qu'il existe un objet $Q''_{i+1}\in K''$ et un morphisme
$u''_{i+1}:Q''_{i+1}\to Z_{i+1}(P_\bullet)$ tels que le composé
$Q''_{i+1}\to Z_{i+1}(P_\bullet)\to H_{i+1}(P_\bullet)$ soit un
épimorphisme. Si on prend $d''_{i+1}:Q''_{i+1}\to N_i$ égal à $0$, le
diagramme
\[
\xymatrix{
Q''_{i+1}\ar[r]^{d''_{i+1}}\ar[d]_{u''_{i+1}} &
N_i\ar[d]^{u_i}\\
Z_{i+1}(P_\bullet)\ar[r]_{d_{i+1}} & P_i
}
\tag{11.9.1.4}\label{0.11.9.1.4-fr}
\]
est commutatif, la flèche horizontale inférieure étant $0$.
\oldpage[0\textsubscript{III}]{48}
Prenons alors
\[
Q_{i+1}=Q'_{i+1}\times Q''_{i+1},\qquad
d_{i+1}=d'_{i+1}+d''_{i+1},\qquad
u_{i+1}=u'_{i+1}+u''_{i+1}.
\]
Comme $d_{i+1}(Q_{i+1})=d'_{i+1}(Q'_{i+1})=N_i\subset Z_i(Q_\bullet)$, on a
$d_i\circ d_{i+1}=0$ et, avec les notations usuelles,
$B_i(Q_\bullet)=N_i$, ce qui vérifie $(\mathrm I_{i+1})$. La commutativité
des diagrammes (\hyperref[0.11.9.1.3-fr]{11.9.1.3}) et
(\hyperref[0.11.9.1.4-fr]{11.9.1.4}) montre que
$d_{i+1}\circ u_{i+1}=u_i\circ d_{i+1}$. Par définition de $N_i$, le
morphisme
\[
H_i(Q_\bullet)=Z_i(Q_\bullet)/N_i\longrightarrow H_i(P_\bullet)
\]
déduit du système des $u_k$ ($k\leq i+1$) est le morphisme déduit de $v_i$
par passage aux quotients, donc c'est un isomorphisme puisque $v_i$ est un
épimorphisme, d'où $(\mathrm{II}_{i+1})$. Enfin, on a
$Q''_{i+1}\subset Z_{i+1}(Q_\bullet)$ par définition; le choix de
$u''_{i+1}$ montre que le morphisme
\[
v_{i+1}:Z_{i+1}(Q_\bullet)\longrightarrow Z_{i+1}(P_\bullet)
\longrightarrow H_{i+1}(P_\bullet)
\]
est un épimorphisme, sa restriction à $Q''_{i+1}$ l'étant déjà, d'où
$(\mathrm{III}_{i+1})$. La récurrence peut donc se poursuivre, et la
proposition est démontrée.
\end{proof}

\begin{corollary}[11.9.2]
\label{0.11.9.2-fr}
Soient $A$ un anneau noethérien (non nécessairement commutatif),
$P_\bullet=(P_i)_{i\in\mathbb Z}$ un complexe de $A$-modules à droite. On
suppose que les $H_i(P_\bullet)$ sont des $A$-modules de type fini et qu'il
existe $d$ tel que $H_i(P_\bullet)=0$ pour $i<d$. Alors il existe un complexe
$Q_\bullet=(Q_i)_{i\in\mathbb Z}$ formé de $A$-modules à droite libres de
rang fini, tel que $Q_i=0$ pour $i<d$, et un homomorphisme
$u:Q_\bullet\to P_\bullet$ de complexes, tel que l'homomorphisme
$H_\bullet(Q_\bullet)\to H_\bullet(P_\bullet)$ correspondant à $u$ soit
bijectif.
\end{corollary}

\begin{proof}
On applique (\hyperref[0.11.9.1-fr]{11.9.1}) en prenant pour $\mathcal C$ la
catégorie des $A$-modules à droite, pour $K'$ (resp. $K''$) l'ensemble des
$A$-modules de type fini (resp. l'ensemble des $A$-modules libres de rang
fini); la vérification des conditions (i), (ii) et (iii) de
(\hyperref[0.11.9.1-fr]{11.9.1}) est immédiate, compte tenu de l'hypothèse
que $A$ est noethérien.
\end{proof}

\begin{env}[11.9.3]
\label{0.11.9.3-fr}
\emph{Remarques}.
\begin{enumerate}
\item[{\rm(i)}] Sous les conditions de
(\hyperref[0.11.9.2-fr]{11.9.2}), supposons en outre que les $P_i$ soient des
$A$-modules à droite plats. Alors, pour tout $A$-module à gauche $M$,
l'homomorphisme de complexes
\[
u\otimes 1:Q_\bullet\otimes_A M\longrightarrow P_\bullet\otimes_A M
\]
définit encore un isomorphisme
$H_\bullet(Q_\bullet\otimes_A M)\simeq
H_\bullet(P_\bullet\otimes_A M)$ de l'homologie comme nous le verrons au
chap.~III.
\item[{\rm(ii)}] La conclusion de (\hyperref[0.11.9.2-fr]{11.9.2}) n'est
plus nécessairement exacte lorsqu'on ne suppose pas $A$ noethérien; en
effet, en l'appliquant à un complexe réduit à $0$ sauf pour un seul terme, on
en conclurait que tout $A$-module à gauche de type fini admet une résolution
par des modules libres de type fini, ce qui n'est pas vrai en général
(cf. Bourbaki, \emph{Alg. comm.}, chap.~I, \S~2, exerc.~6).

Toutefois, au lieu de supposer $A$ noethérien, on peut supposer seulement que
les $H_i(P_\bullet)$ ont une $\infty$-présentation finie (cf. chap.~IV).
\end{enumerate}
\end{env}

\subsection{Caractéristique d'Euler-Poincaré d'un complexe de modules de
longueur finie}
\label{subsection:0.11.10-fr}

\begin{env}[11.10.1]
\label{0.11.10.1-fr}
Soient $A$ un anneau (non nécessairement commutatif),
\[
\label{0.11.10.1.1-fr}
\tag{11.10.1.1}
M^\bullet:0\longrightarrow M^0\longrightarrow M^1
\longrightarrow\cdots\longrightarrow M^n\longrightarrow0
\]
un complexe de $A$-modules à gauche de longueur finie. On appelle
\emph{caractéristique d'Euler-Poincaré} de ce complexe le nombre
\oldpage[0\textsubscript{III}]{49}
\[
\label{0.11.10.1.2-fr}
\tag{11.10.1.2}
\chi(M^\bullet)=\sum_{i=0}^{n}(-1)^i\operatorname{long}(M^i).
\]
\end{env}

\begin{proposition}[11.10.2]
\label{0.11.10.2-fr}
\emph{Pour tout complexe fini $M^\bullet$ de $A$-modules à gauche de longueur
finie, on a
\[
\chi(M^\bullet)=\chi(H^\bullet(M^\bullet))
\]
($H^\bullet(M^\bullet)$ étant considéré comme un complexe pour la dérivation
triviale). En particulier, si la suite
(\hyperref[0.11.10.1.1-fr]{11.10.1.1}) est exacte, on a
$\chi(M^\bullet)=0$.}
\end{proposition}

\begin{proof}
Posons pour abréger $B^i=B^i(M^\bullet)$, $Z^i=Z^i(M^\bullet)$,
$H^i=H^i(M^\bullet)=Z^i/B^i$; les $B^i$, $Z^i$, $H^i$ sont de longueur
finie. Des suites exactes
\[
0\longrightarrow B^i\longrightarrow Z^i\longrightarrow H^i
\longrightarrow0,
\qquad
0\longrightarrow Z^i\longrightarrow M^i\longrightarrow B^{i+1}
\longrightarrow0
\]
on tire les relations
\[
\operatorname{long}(Z^i)=\operatorname{long}(H^i)+
\operatorname{long}(B^i),
\qquad
\operatorname{long}(M^i)=\operatorname{long}(Z^i)+
\operatorname{long}(B^{i+1}),
\]
d'où
\[
\operatorname{long}(M^i)-\operatorname{long}(H^i)=
\operatorname{long}(B^{i+1})+\operatorname{long}(B^i).
\]
Multiplions cette relation par $(-1)^i$ et faisons la somme des relations
obtenues pour $0\leq i\leq n$, en notant que
$B^0=B^{n+1}=0$, il vient l'égalité cherchée.
\end{proof}

\begin{corollary}[11.10.3]
\label{0.11.10.3-fr}
Soit $E=(E_r^{pq})$ une suite spectrale dans la catégorie des modules sur un
anneau $A$. On suppose que les $E_2^{pq}$ sont des $A$-modules de longueur
finie et qu'il n'y a qu'un nombre fini de couples $(p,q)$ tels que
$E_2^{pq}\neq0$. Alors les caractéristiques d'Euler-Poincaré
$\chi(E_r^\bullet)$ de tous les complexes
$E_r^\bullet=(E_r^{(n)})_{n\in\mathbb Z}$
(\hyperref[0.11.1.1-fr]{11.1.1}) sont toutes égales. Si, en outre, la suite
$E$ est faiblement convergente et si on pose
\[
E_\infty^{(n)}=\bigoplus_{p+q=n}E_\infty^{pq}\quad(n\in\mathbb Z),
\]
on a aussi $\chi(E_\infty^\bullet)=\chi(E_2^\bullet)$,
$E_\infty^\bullet=(E_\infty^{(n)})_{n\in\mathbb Z}$ étant considéré comme
un complexe à dérivation triviale.
\end{corollary}

\begin{proof}
Notons d'abord que si $E_2^{pq}=0$, on a $E_r^{pq}=0$ pour
$2\leq r<+\infty$, donc tous les complexes $E_r^\bullet$ sont finis et
formés de $A$-modules de longueur finie; la relation
$\chi(E_r^\bullet)=\chi(E_{r+1}^\bullet)$ pour tout $r$ fini résulte donc de
(\hyperref[0.11.10.2-fr]{11.10.2}) et de l'isomorphie entre
$H^\bullet(E_r^\bullet)$ et $E_{r+1}^\bullet$ (en tant que complexes à
dérivation triviale). L'hypothèse que les $E_r^{pq}$ sont de longueur finie
entraîne que pour tout couple $(p,q)$, les suites
$(B_k(E_2^{pq}))_{k\geq2}$ et $(Z_k(E_2^{pq}))_{k\geq2}$ sont stationnaires;
l'hypothèse que $E$ est faiblement convergente et que $E_2^{pq}=0$, sauf
pour un nombre fini de couples $(p,q)$, entraîne donc qu'il existe un entier
$r\geq2$ tel que $E_\infty^{pq}=E_r^{pq}$ pour tout couple $(p,q)$; d'où
l'assertion relative à $\chi(E_\infty^\bullet)$.
\end{proof}

\section{Compléments sur la cohomologie des faisceaux}
\label{section:0.12-fr}

\subsection{Cohomologie des faisceaux de modules sur les espaces annelés}
\label{subsection:0.12.1-fr}

\begin{env}[12.1.1]
\label{0.12.1.1-fr}
Soit $(X,\mathcal O_X)$ un espace annelé. Rappelons que pour tout
$\mathcal O_X$-Module $\mathcal F$, on définit la cohomologie
$H^\bullet(X,\mathcal F)$, qui est un foncteur cohomologique universel
(T, 2.2) de la catégorie $C(X)$ des $\mathcal O_X$-Modules dans la catégorie
des groupes abéliens; c'est le foncteur dérivé du foncteur exact à gauche
$\mathcal F\to\Gamma(X,\mathcal F)$. Le foncteur $H^\bullet(X,\mathcal F)$
est isomorphe
\oldpage[0\textsubscript{III}]{50}
à la restriction à la catégorie $C(X)$ du foncteur cohomologique défini de
même sur la catégorie des \emph{faisceaux de groupes abéliens} sur $X$
(G, II, 7.2.1).
\end{env}

\begin{env}[12.1.2]
\label{0.12.1.2-fr}
Posons $A=\Gamma(X,\mathcal O_X)$. Comme tout élément de $A$ définit un
endomorphisme du groupe abélien $\Gamma(X,\mathcal F)$, il définit par
fonctorialité un endomorphisme de $\delta$-foncteur de
$H^\bullet(X,\mathcal F)$; ces endomorphismes définissent sur chacun des
$H^p(X,\mathcal F)$ une structure de $A$-module, et l'opérateur $\partial$
est $A$-linéaire. En outre, pour deux entiers positifs quelconques $p$, $q$,
et deux $\mathcal O_X$-Modules $\mathcal F$, $\mathcal G$, on a un
homomorphisme de $A$-modules, dit \emph{cup-produit}
\[
\label{0.12.1.2.1-fr}
\tag{12.1.2.1}
H^p(X,\mathcal F)\otimes_A H^q(X,\mathcal G)
\longrightarrow H^{p+q}(X,\mathcal F\otimes_{\mathcal O_X}\mathcal G)
\]
(G, II, 6.6). Ces homomorphismes font de la somme directe $S$ des
$H^p(X,\mathcal O_X)$ (pour $p\geq0$) une $A$-algèbre graduée
anticommutative et de la somme directe des $H^p(X,\mathcal F)$ un $S$-module
gradué.

Pour tout \emph{recouvrement ouvert} $\mathfrak U$ de $X$, nous désignerons
toujours par $\check C^\bullet(\mathfrak U,\mathcal F)$ (contrairement à
(G, II, 5.1)) le complexe des cochaînes \emph{alternées} du nerf de
$\mathfrak U$ à valeurs dans le système de coefficients
$\Gamma(U_\alpha,\mathcal F)$. Il est clair que
$\check C^\bullet(\mathfrak U,\mathcal F)$ est un $A$-module gradué, donc les
groupes de cohomologie $\check H^p(\mathfrak U,\mathcal F)$ de ce complexe
sont munis d'une structure de $A$-module; en outre, les applications
canoniques
\[
\check H^p(\mathfrak U,\mathcal F)\longrightarrow H^p(X,\mathcal F)
\]
(G, II, 5.4) sont $A$-linéaires.
\end{env}

\begin{env}[12.1.3]
\label{0.12.1.3-fr}
Soit $(X',\mathcal O_{X'})$ un second espace annelé, et soit
$f=(\psi,\theta)$ un morphisme de $X'$ dans $X$.

Posons $A'=\Gamma(X',\mathcal O_{X'})$; le $\psi$-morphisme $\theta$ définit
canoniquement un homomorphisme d'anneaux $A\to A'$. Soient $\mathcal F$ un
$\mathcal O_X$-Module, $\mathcal F'$ un $\mathcal O_{X'}$-Module; pour tout
$f$-morphisme $u:\mathcal F\to\mathcal F'$
(\hyperref[0.4.4.1-fr]{0\textsubscript{I}, 4.4.1}), nous allons voir qu'on
peut définir pour tout $p\geq0$, un \emph{di-homomorphisme}
\[
\label{0.12.1.3.1-fr}
\tag{12.1.3.1}
u_p:H^p(X,\mathcal F)\longrightarrow H^p(X',\mathcal F').
\]
En effet, comme $\psi^*$ est exact dans la catégorie des faisceaux de groupes
abéliens sur $X$, $\mathcal F\to H^\bullet(X',\psi^*(\mathcal F))$ est un
$\delta$-foncteur dans cette catégorie, et on sait qu'on a un homomorphisme
canonique de $\delta$-foncteurs
\[
\label{0.12.1.3.2-fr}
\tag{12.1.3.2}
H^\bullet(X,\mathcal F)\longrightarrow H^\bullet(X',\psi^*(\mathcal F))
\]
uniquement déterminé par la condition de se réduire à l'homomorphisme
canonique $\Gamma(X,\mathcal F)\to\Gamma(X',\psi^*(\mathcal F))$ en degré
0 (T, 3.2.2). En outre, tout élément de $A$ détermine un endomorphisme $\mu$
de $\Gamma(X,\mathcal F)$ et un endomorphisme $\mu'$ de
$\Gamma(X',\psi^*(\mathcal F))$ tels que le diagramme
\[
\label{0.12.1.3.3-fr}
\tag{12.1.3.3}
\xymatrix{
\Gamma(X,\mathcal F)\ar[r]\ar[d]_\mu
&\Gamma(X',\psi^*(\mathcal F))\ar[d]^{\mu'}\\
\Gamma(X,\mathcal F)\ar[r]
&\Gamma(X',\psi^*(\mathcal F))
}
\]
soit commutatif; par la propriété d'unicité de prolongement des morphismes
pour les foncteurs cohomologiques universels (T, 2.2), on en déduit des
prolongements uniques de $\mu$ et $\mu'$ à la cohomologie rendant commutatifs
les diagrammes analogues à (\hyperref[0.12.1.3.3-fr]{12.1.3.3}), ce qui
signifie que (\hyperref[0.12.1.3.2-fr]{12.1.3.2}) est un homomorphisme de
$A$-modules. Notons maintenant que l'on a
\[
f^*(\mathcal F)=\psi^*(\mathcal F)
\otimes_{\psi^*(\mathcal O_X)}\mathcal O_{X'},
\]
et que l'on a donc un di-homomorphisme canonique
$\psi^*(\mathcal F)\to f^*(\mathcal F)$ du
$\psi^*(\mathcal O_X)$-Module $\psi^*(\mathcal F)$ dans le
$\mathcal O_{X'}$-Module $f^*(\mathcal F)$. Par fonctorialité, on en déduit
donc un di-homomorphisme fonctoriel
\[
\label{0.12.1.3.4-fr}
\tag{12.1.3.4}
H^p(X',\psi^*(\mathcal F))\longrightarrow H^p(X',f^*(\mathcal F))
\]
les anneaux correspondants étant $A$ et $A'$; en composant ce
di-homomorphisme avec (\hyperref[0.12.1.3.2-fr]{12.1.3.2}), on obtient un
di-homomorphisme canonique fonctoriel en $\mathcal F$
\[
\label{0.12.1.3.5-fr}
\tag{12.1.3.5}
\theta_p:H^p(X,\mathcal F)\longrightarrow H^p(X',f^*(\mathcal F)).
\]
Enfin, par fonctorialité, on déduit de l'homomorphisme
$u^\sharp:f^*(\mathcal F)\to\mathcal F'$ un homomorphisme de $A'$-modules
$H^p(X',f^*(\mathcal F))\to H^p(X',\mathcal F')$, qui, composé avec
(\hyperref[0.12.1.3.5-fr]{12.1.3.5}), donne
(\hyperref[0.12.1.3.1-fr]{12.1.3.1}).

Soit $f'=(\psi',\theta'):X''\to X'$ un second morphisme d'espaces annelés,
$f''=ff'$ le morphisme composé. Compte tenu de la permutabilité du foncteur
$\psi'^*$ et du produit tensoriel
(\hyperref[0.4.3.3-fr]{0\textsubscript{I}, 4.3.3}), on vérifie aussitôt que le
composé du di-homomorphisme
$H^p(X',f^*(\mathcal F))\to H^p(X'',f'^*(f^*(\mathcal F)))$ et de
(\hyperref[0.12.1.3.5-fr]{12.1.3.5}) est le di-homomorphisme correspondant
$H^p(X,\mathcal F)\to H^p(X'',f''^*(\mathcal F))$.
\end{env}

\begin{env}[12.1.4]
\label{0.12.1.4-fr}
Une définition directe de l'homomorphisme
(\hyperref[0.12.1.3.2-fr]{12.1.3.2}) peut s'obtenir de la façon suivante : on
considère une résolution injective $\mathcal L^\bullet=(\mathcal L^i)$ de
$\mathcal F$ formée de faisceaux de groupes abéliens sur $X$; comme le
foncteur $\psi^*$ est exact, $\psi^*(\mathcal L^\bullet)$ est une résolution
de $\psi^*(\mathcal F)$ formée de faisceaux sur $X'$. Si
$\mathcal L'^\bullet=(\mathcal L'^i)$ est une résolution injective de
$\psi^*(\mathcal F)$ dans la catégorie des faisceaux de groupes abéliens sur
$X'$, il y a donc un morphisme
$\psi^*(\mathcal L^\bullet)\to\mathcal L'^\bullet$ de complexes de faisceaux
de groupes abéliens, compatible avec les augmentations (M, V, 1.1.a)), bien
déterminé à une homotopie près. On en déduit des homomorphismes
\[
\Gamma(X,\mathcal L^\bullet)\longrightarrow
\Gamma(X',\psi^*(\mathcal L^\bullet))\longrightarrow
\Gamma(X',\mathcal L'^\bullet)
\]
de complexes de groupes abéliens, dont le composé, par passage à la
cohomologie donne un morphisme de $\delta$-foncteurs
$H^\bullet(X,\mathcal F)\to H^\bullet(X',\psi^*(\mathcal F))$; comme il
coïncide avec (\hyperref[0.12.1.3.2-fr]{12.1.3.2}) en degré 0, il lui est
identique (T, 2.2).

Considérons maintenant un \emph{recouvrement ouvert}
$\mathfrak U=(U_\alpha)$ de $X$, et soit
$\mathfrak U'=(f^{-1}(U_\alpha))$ le recouvrement ouvert de $X'$, image
réciproque de $\mathfrak U$. Les homomorphismes canoniques
$\Gamma(V,\mathcal F)\to\Gamma(f^{-1}(V),f^*(\mathcal F))$ pour tout ouvert
$V$ de $X$ définissent aussitôt (cf. G, II, 5.1) un homomorphisme de complexes
$\check C^\bullet(\mathfrak U,\mathcal F)\to
\check C^\bullet(\mathfrak U',f^*(\mathcal F))$, d'où des homomorphismes
canoniques
\[
\label{0.12.1.4.1-fr}
\tag{12.1.4.1}
\theta_p:\check H^p(\mathfrak U,\mathcal F)
\longrightarrow\check H^p(\mathfrak U',f^*(\mathcal F)).
\]

\oldpage[0\textsubscript{III}]{52}
En outre, on a des diagrammes commutatifs
\[
\label{0.12.1.4.2-fr}
\tag{12.1.4.2}
\xymatrix{
\check H^p(\mathfrak U,\mathcal F)\ar[r]^{\theta_p}\ar[d]
&\check H^p(\mathfrak U',f^*(\mathcal F))\ar[d]\\
H^p(X,\mathcal F)\ar[r]_{\theta_p}
&H^p(X',f^*(\mathcal F))
}
\]
où les flèches verticales sont les homomorphismes canoniques de (G, II, 5.2).
Pour établir la commutativité de
(\hyperref[0.12.1.4.2-fr]{12.1.4.2}), considérons le complexe de faisceaux de
cochaînes (alternées) de $\mathcal F$ relatives à $\mathfrak U$,
$\mathcal C^\bullet(\mathfrak U,\mathcal F)$, tel que
$\Gamma(X,\mathcal C^\bullet(\mathfrak U,\mathcal F))=
\check C^\bullet(\mathfrak U,\mathcal F)$ (G, II, 5.2). Les homomorphismes
canoniques $\Gamma(V,\mathcal F)\to
\Gamma(f^{-1}(V),\psi^*(\mathcal F))$ définissent alors un $\psi$-morphisme
$\mathcal C^\bullet(\mathfrak U,\mathcal F)\to
\mathcal C^\bullet(\mathfrak U',\psi^*(\mathcal F))$, et on a, avec les
notations ci-dessus, un diagramme commutatif
\[
\xymatrix{
\Gamma(X,\mathcal C^\bullet(\mathfrak U,\mathcal F))\ar[r]\ar[d]
&\Gamma(X',\mathcal C^\bullet(\mathfrak U',\psi^*(\mathcal F)))\ar[d]\\
\Gamma(X,\mathcal L^\bullet)\ar[r]
&\Gamma(X',\psi^*(\mathcal L^\bullet))\ar[r]
&\Gamma(X',\mathcal L'^\bullet)
}
\]
qui, en passant à la cohomologie, donne des diagrammes commutatifs
\[
\xymatrix{
\check H^p(\mathfrak U,\mathcal F)\ar[r]\ar[d]
&\check H^p(\mathfrak U',\psi^*(\mathcal F))\ar[d]\\
H^p(X,\mathcal F)\ar[r]
&H^p(X',\psi^*(\mathcal F))
}
\]
où les flèches verticales sont les homomorphismes canoniques de (G, II, 5.2).
Il suffit alors de combiner ces diagrammes avec les diagrammes commutatifs
\[
\xymatrix{
\check H^p(\mathfrak U',\psi^*(\mathcal F))\ar[r]\ar[d]
&\check H^p(\mathfrak U',f^*(\mathcal F))\ar[d]\\
H^p(X',\psi^*(\mathcal F))\ar[r]
&H^p(X',f^*(\mathcal F))
}
\]
\oldpage[0\textsubscript{III}]{53}
qui proviennent de l'homomorphisme $\psi^*(\mathcal F)\to f^*(\mathcal F)$
et du caractère fonctoriel des homomorphismes canoniques de (G, II, 5.2),
pour obtenir la commutativité de
(\hyperref[0.12.1.4.2-fr]{12.1.4.2}).

On notera que si $A=\Gamma(X,\mathcal O_X)$,
$A'=\Gamma(X',\mathcal O_{X'})$, l'homomorphisme (12.1.4.3) est un
di-homomorphisme de modules
correspondant aux anneaux $A$ et $A'$. On a une propriété de
\emph{transitivité} de (\hyperref[0.12.1.4.1-fr]{12.1.4.1}) pour le composé de
deux morphismes, analogue à la transitivité de
(\hyperref[0.12.1.3.5-fr]{12.1.3.5}). Enfin, notons que dans les définitions
précédentes, au lieu d'une résolution injective $\mathcal L^\bullet$ de
$\mathcal F$, on aurait aussi bien pu partir d'une résolution telle que
$H^p(X,\mathcal L^i)=0$ pour tout $i$ et tout $p>0$ (G, II, 4.7.1).
\end{env}

\begin{env}[12.1.5]
\label{0.12.1.5-fr}
Soient $\mathcal F$, $\mathcal G$, $\mathcal H$ trois
$\mathcal O_X$-Modules, et considérons un $\mathcal O_X$-homomorphisme
$u:\mathcal F\otimes_{\mathcal O_X}\mathcal G\to\mathcal H$, qui donne pour
la cohomologie des homomorphismes
\[
\label{0.12.1.5.1-fr}
\tag{12.1.5.1}
H^p(X,\mathcal F)\otimes_A H^q(X,\mathcal G)
\longrightarrow H^{p+q}(X,\mathcal H)
\]
déduits du cup-produit (\hyperref[0.12.1.2.1-fr]{12.1.2.1}). Montrons qu'avec
les hypothèses et notations de (\hyperref[0.12.1.3-fr]{12.1.3}), on a des
diagrammes commutatifs
\[
\label{0.12.1.5.2-fr}
\tag{12.1.5.2}
\xymatrix@C=10pt{
H^p(X,\mathcal F)\otimes_A H^q(X,\mathcal G)\ar[r]\ar[d]
&H^{p+q}(X,\mathcal H)\ar[d]\\
H^p(X',f^*(\mathcal F))\otimes_{A'}H^q(X',f^*(\mathcal G))\ar[r]
&H^{p+q}(X',f^*(\mathcal H))
}
\]
où les flèches verticales proviennent des homomorphismes canoniques
(\hyperref[0.12.1.3.5-fr]{12.1.3.5}). Pour cela, rappelons que
(\hyperref[0.12.1.5.1-fr]{12.1.5.1}) peut s'obtenir en partant des résolutions
\emph{canoniques} (G, II, 4.3) $\mathcal L^\bullet$, $\mathcal M^\bullet$,
$\mathcal N^\bullet$ de $\mathcal F$, $\mathcal G$, $\mathcal H$
respectivement (qui sont formées de $\mathcal O_X$-Modules), de l'application
linéaire $\mathcal L^\bullet\otimes_{\mathcal O_X}\mathcal M^\bullet
\to\mathcal N^\bullet$ de complexes de $\mathcal O_X$-Modules correspondant
à $u$, qui fournit un homomorphisme de complexes de $A$-modules
$\Gamma(X,\mathcal L^\bullet)\otimes_A\Gamma(X,\mathcal M^\bullet)
\to\Gamma(X,\mathcal N^\bullet)$, et en passant à la cohomologie des
homomorphismes
\[
H^p(\Gamma(X,\mathcal L^\bullet))\otimes_A
H^q(\Gamma(X,\mathcal M^\bullet))
\longrightarrow H^{p+q}(\Gamma(X,\mathcal N^\bullet))
\]
(G, II, 6.6). Or, on a évidemment un diagramme commutatif
\[
\label{0.12.1.5.3-fr}
\tag{12.1.5.3}
\xymatrix@C=8pt{
\Gamma(X,\mathcal L^\bullet)\otimes_A\Gamma(X,\mathcal M^\bullet)
\ar[r]\ar[d]
&\Gamma(X,\mathcal N^\bullet)\ar[d]\\
\Gamma(X',\psi^*(\mathcal L^\bullet))
\otimes_{\Gamma(X',\psi^*(\mathcal O_X))}
\Gamma(X',\psi^*(\mathcal M^\bullet))\ar[r]
&\Gamma(X',\psi^*(\mathcal N^\bullet))
}
\]
\oldpage[0\textsubscript{III}]{54}
qui donne en passant à la cohomologie des diagrammes commutatifs
\[
\label{0.12.1.5.4-fr}
\tag{12.1.5.4}
\xymatrix@C=7pt{
H^p(X,\mathcal F)\otimes_A H^q(X,\mathcal G)\ar[r]\ar[d]
&H^{p+q}(X,\mathcal H)\ar[d]\\
H^p(\Gamma(X',\psi^*(\mathcal L^\bullet)))
\otimes_{\Gamma(X',\psi^*(\mathcal O_X))}
H^q(\Gamma(X',\psi^*(\mathcal M^\bullet)))\ar[r]
&H^{p+q}(\Gamma(X',\psi^*(\mathcal N^\bullet)))
}
\]
Mais comme $\psi^*(\mathcal L^\bullet)$, $\psi^*(\mathcal M^\bullet)$ et
$\psi^*(\mathcal N^\bullet)$ sont des \emph{résolutions} de
$\psi^*(\mathcal F)$, $\psi^*(\mathcal G)$, $\psi^*(\mathcal H)$
respectivement, on a un diagramme commutatif (G, II, 6.6.1)
\[
\label{0.12.1.5.5-fr}
\tag{12.1.5.5}
\xymatrix@C=7pt{
H^p(\Gamma(X',\psi^*(\mathcal L^\bullet)))
\otimes_{\Gamma(X',\psi^*(\mathcal O_X))}
H^q(\Gamma(X',\psi^*(\mathcal M^\bullet)))\ar[r]\ar[d]
&H^{p+q}(\Gamma(X',\psi^*(\mathcal N^\bullet)))\ar[d]\\
H^p(X',\psi^*(\mathcal F))
\otimes_{\Gamma(X',\psi^*(\mathcal O_X))}
H^q(X',\psi^*(\mathcal G))\ar[r]
&H^{p+q}(X',\psi^*(\mathcal H))
}
\]
Enfin, par fonctorialité, on a un diagramme commutatif
\[
\label{0.12.1.5.6-fr}
\tag{12.1.5.6}
\xymatrix@C=7pt{
H^p(X',\psi^*(\mathcal F))
\otimes_{\Gamma(X',\psi^*(\mathcal O_X))}
H^q(X',\psi^*(\mathcal G))\ar[r]\ar[d]
&H^{p+q}(X',\psi^*(\mathcal H))\ar[d]\\
H^p(X',f^*(\mathcal F))\otimes_{A'}H^q(X',f^*(\mathcal G))\ar[r]
&H^{p+q}(X',f^*(\mathcal H))
}
\]
et par combinaison des trois diagrammes
(\hyperref[0.12.1.5.4-fr]{12.1.5.4}),
(\hyperref[0.12.1.5.5-fr]{12.1.5.5}) et
(\hyperref[0.12.1.5.6-fr]{12.1.5.6}), on obtient le diagramme commutatif
(\hyperref[0.12.1.5.2-fr]{12.1.5.2}) cherché.
\end{env}

\begin{remark}[12.1.6]
\label{0.12.1.6-fr}
\oldpage[0\textsubscript{III}]{55}
Avec les notations de (\hyperref[0.12.1.3-fr]{12.1.3}), supposons que l'on
ait un diagramme commutatif
\[
\label{0.12.1.6.1-fr}
\tag{12.1.6.1}
\xymatrix{
0\ar[r]&\mathcal F\ar[r]^r\ar[d]_u
&\mathcal G\ar[r]^s\ar[d]_v
&\mathcal H\ar[r]\ar[d]_w&0\\
0\ar[r]&\mathcal F'\ar[r]_{r'}
&\mathcal G'\ar[r]_{s'}
&\mathcal H'\ar[r]&0
}
\]
où $r$, $s$ sont des homomorphismes de $\mathcal O_X$-Modules, $r'$, $s'$ des
homomorphismes de $\mathcal O_{X'}$-Modules, $u$, $v$, $w$ des $f$-morphismes
et les lignes sont \emph{exactes}. On en déduit alors un diagramme commutatif
\[
\label{0.12.1.6.2-fr}
\tag{12.1.6.2}
\xymatrix@C=6pt{
\cdots\ar[r]&H^p(X,\mathcal F)\ar[r]\ar[d]_{u_p}
&H^p(X,\mathcal G)\ar[r]\ar[d]_{v_p}
&H^p(X,\mathcal H)\ar[r]^\partial\ar[d]_{w_p}
&H^{p+1}(X,\mathcal F)\ar[r]\ar[d]_{u_{p+1}}&\cdots\\
\cdots\ar[r]&H^p(X',\mathcal F')\ar[r]
&H^p(X',\mathcal G')\ar[r]
&H^p(X',\mathcal H')\ar[r]_\partial
&H^{p+1}(X',\mathcal F')\ar[r]&\cdots
}
\]
En effet, (\hyperref[0.12.1.6.1-fr]{12.1.6.1}) se factorise en
\[
\xymatrix@R=14pt{
0\ar[r]&\mathcal F\ar[r]\ar[d]
&\mathcal G\ar[r]\ar[d]
&\mathcal H\ar[r]\ar[d]&0\\
0\ar[r]&\psi^*(\mathcal F)\ar[r]\ar[d]
&\psi^*(\mathcal G)\ar[r]\ar[d]
&\psi^*(\mathcal H)\ar[r]\ar[d]&0\\
0\ar[r]&\mathcal F'\ar[r]
&\mathcal G'\ar[r]
&\mathcal H'\ar[r]&0
}
\]
où la ligne du milieu est \emph{exacte}
(\hyperref[0.3.7.2-fr]{0\textsubscript{I}, 3.7.2}) et il suffit d'utiliser le
fait que (\hyperref[0.12.1.3.2-fr]{12.1.3.2}) est un homomorphisme de
$\delta$-foncteurs et que les $H^p(X',\mathcal F')$ forment un
$\delta$-foncteur en $\mathcal F'$.
\end{remark}

\begin{env}[12.1.7]
\label{0.12.1.7-fr}
Les hypothèses et notations étant celles de
(\hyperref[0.12.1.3-fr]{12.1.3}), considérons maintenant le cas où
$\mathcal F=f_*(\mathcal F')=\psi_*(\mathcal F')$; nous allons voir que le
di-homomorphisme défini dans (\hyperref[0.12.1.3-fr]{12.1.3})
\[
\label{0.12.1.7.1-fr}
\tag{12.1.7.1}
H^p(X,f_*(\mathcal F'))\longrightarrow H^p(X',\mathcal F')
\]
peut s'obtenir (à un automorphisme près de $H^p(X',\mathcal F')$) comme
\emph{edge-homomorphisme d'une suite spectrale} du foncteur composé
$\mathcal F'\leadsto\Gamma(X,\psi_*(\mathcal F'))$ (T, 2.4). La description
de l'homomorphisme (\hyperref[0.12.1.7.1-fr]{12.1.7.1}) donnée dans
(\hyperref[0.12.1.4-fr]{12.1.4}) montre ici qu'on peut obtenir cet
homomorphisme de la façon suivante : on considère des résolutions injectives
$\mathcal L^\bullet$ et $\mathcal L'^\bullet$ de $\psi_*(\mathcal F')$ et de
$\mathcal F'$ respectivement, puis on prend un homomorphisme de complexes
$v:\psi^*(\mathcal L^\bullet)\to\mathcal L'^\bullet$ «~au-dessus~» de
l'homomorphisme canonique $\psi^*(\psi_*(\mathcal F'))\to\mathcal F'$;
\oldpage[0\textsubscript{III}]{56}
on note ensuite que l'on a
$\Gamma(X',\mathcal L'^\bullet)=\Gamma(X,\psi_*(\mathcal L'^\bullet))$ et
que l'homomorphisme composé
\[
\Gamma(X,\mathcal L^\bullet)\longrightarrow
\Gamma(X',\psi^*(\mathcal L^\bullet))
\xrightarrow{\Gamma(v)}\Gamma(X',\mathcal L'^\bullet)
\]
n'est autre que
\[
\label{0.12.1.7.2-fr}
\tag{12.1.7.2}
\Gamma(v^\flat):\Gamma(X,\mathcal L^\bullet)
\longrightarrow\Gamma(X,\psi_*(\mathcal L'^\bullet))
\]
(\hyperref[0.3.7.1-fr]{0\textsubscript{I}, 3.7.1}), et
(\hyperref[0.12.1.7.1-fr]{12.1.7.1}) s'obtient par passage à la cohomologie
dans (\hyperref[0.12.1.7.2-fr]{12.1.7.2}). D'autre part, les suites spectrales
du foncteur composé $\mathcal F'\leadsto\Gamma(X,\psi_*(\mathcal F'))$
s'obtiennent en considérant une résolution injective de Cartan-Eilenberg
$\mathcal M^{\bullet\bullet}=(\mathcal M^{ij})$ du complexe
$\psi_*(\mathcal L'^\bullet)$ dans la catégorie des faisceaux de groupes
abéliens sur $X$; les suites spectrales en question sont celles du bicomplexe
$\Gamma(X,\mathcal M^{\bullet\bullet})$ (qui sont birégulières puisque
$\mathcal M^{ij}=0$ pour $i<0$ ou $j<0$). Or, la première suite spectrale de
ce bicomplexe \emph{dégénère}, car les faisceaux
$\psi_*(\mathcal L'^i)$ sont flasques (G, II, 3.1.1), donc
$H_{\mathrm{II}}^q(\Gamma(X,\mathcal M^{i,\bullet}))
=H^q(\psi_*(\mathcal L'^i))=0$ pour $q>0$ (G, II, 4.4.3); on a donc des
edge-homomorphismes \emph{bijectifs} (\hyperref[0.11.1.6-fr]{11.1.6})
\[
\label{0.12.1.7.3-fr}
\tag{12.1.7.3}
'E_2^{i0}=H^i(H_{\mathrm{II}}^0(\Gamma(X,\mathcal M^{\bullet\bullet})))
\longrightarrow H^i(\Gamma(X,\mathcal M^{\bullet\bullet}))
\]
et on sait (\hyperref[0.11.3.4-fr]{11.3.4}) que cet homomorphisme provient,
par passage à la cohomologie, de l'augmentation
\[
\label{0.12.1.7.4-fr}
\tag{12.1.7.4}
\Gamma(X,\psi_*(\mathcal L'^\bullet))
\longrightarrow\Gamma(X,\mathcal M^{\bullet\bullet})
\]
laquelle provient elle-même de l'augmentation
$\eta:\psi_*(\mathcal L'^\bullet)\to\mathcal M^{\bullet\bullet}$. D'autre
part, pour la seconde suite spectrale, on a des edge-homomorphismes
\[
\label{0.12.1.7.5-fr}
\tag{12.1.7.5}
''E_2^{i0}=H^i(H_{\mathrm I}^0(\Gamma(X,\mathcal M^{\bullet\bullet})))
\longrightarrow H^i(\Gamma(X,\mathcal M^{\bullet\bullet}))
\]
provenant (\hyperref[0.11.3.4-fr]{11.3.4}) par passage à la cohomologie, de
l'homomorphisme de complexes
$Z_{\mathrm I}^0(\Gamma(X,\mathcal M^{\bullet\bullet}))
\to\Gamma(X,\mathcal M^{\bullet\bullet})$. Or, comme $\psi_*$ est exact à
gauche, la suite
\[
0\longrightarrow\psi_*(\mathcal F')\longrightarrow
\psi_*(\mathcal L'^0)\longrightarrow\psi_*(\mathcal L'^1)
\]
est exacte; par définition d'une résolution de Cartan-Eilenberg
(\hyperref[0.11.4.2-fr]{11.4.2}), on peut donc prendre
$B_{\mathrm I}^0(\mathcal M^{\bullet\bullet})=0$,
$Z_{\mathrm I}^0(\mathcal M^{\bullet\bullet})=\mathcal L^\bullet$; comme le
diagramme
\[
\xymatrix@C=50pt@R=24pt{
\mathcal L^0\ar[r]^{i^0}&\mathcal M^{00}\\
\psi_*(\mathcal F')\ar[u]^\varepsilon
\ar[r]_{\varepsilon'}\ar[ur]^{\varepsilon''}
&\psi_*(\mathcal L'^0)\ar[u]_{\eta^0}
}
\]
est commutatif, l'injection de complexes
$i:\mathcal L^\bullet\to\mathcal M^{\bullet\bullet}$ est compatible avec les
augmentations $\varepsilon$ et $\varepsilon''$. On a ainsi deux
homomorphismes de complexes de $\mathcal L^\bullet$ dans
$\mathcal M^{\bullet\bullet}$
\[
\xymatrix@C=52pt@R=22pt{
\mathcal L^\bullet\ar[r]^i\ar[dr]_{v^\flat}
&\mathcal M^{\bullet\bullet}\\
&\psi_*(\mathcal L'^\bullet)\ar[u]_\eta
}
\]
\oldpage[0\textsubscript{III}]{57}
compatibles avec les augmentations $\varepsilon$ et $\varepsilon''$; comme
$\mathcal L^\bullet$ est une \emph{résolution} injective et que
$\mathcal M^{\bullet\bullet}$ est formée de faisceaux injectifs, il résulte de
(M, V, 1.1a)) que ces deux homomorphismes sont \emph{homotopes}; il en est
donc de même des deux homomorphismes correspondants
$\Gamma(X,\mathcal L^\bullet)\to\Gamma(X,\mathcal M^{\bullet\bullet})$, et en
passant à la cohomologie, on obtient donc le \emph{même} homomorphisme; en
d'autres termes, on a bien montré que le edge-homomorphisme
(\hyperref[0.12.1.7.5-fr]{12.1.7.5}), qui s'écrit
$H^p(X,\psi_*(\mathcal F'))\to H^p(\Gamma(X,\mathcal M^{\bullet\bullet}))$
est composé de (\hyperref[0.12.1.7.1-fr]{12.1.7.1}) et de
(\hyperref[0.12.1.7.3-fr]{12.1.7.3}), qui s'écrit
$H^p(X',\mathcal F')\to H^p(\Gamma(X,\mathcal M^{\bullet\bullet}))$ et qu'on
a vu être un \emph{isomorphisme}; d'où notre assertion.
\end{env}

\subsection{Images directes supérieures}
\label{subsection:0.12.2-fr}

\begin{env}[12.2.1]
\label{0.12.2.1-fr}
Soient $(X,\mathcal O_X)$, $(Y,\mathcal O_Y)$ deux espaces annelés,
$f=(\psi,\theta)$ un morphisme de $X$ dans $Y$, qui définit le foncteur
image directe
\[
f_*:C(X)\longrightarrow C(Y),
\]
identique d'ailleurs à la restriction à $C(X)$ du foncteur $\psi_*$ défini
dans la catégorie des faisceaux de groupes abéliens sur $X$. Ce dernier
foncteur est additif et exact à gauche, et comme tout faisceau de groupes
abéliens sur $X$ est isomorphe à un sous-faisceau d'un faisceau injectif de
groupes abéliens, on définit les foncteurs dérivés droits
$\mathcal F\leadsto R^p\psi_*(\mathcal F)$ du foncteur $\psi_*$; les
$R^p\psi_*(\mathcal F)$ sont des faisceaux de groupes abéliens sur $Y$, et
les $R^p\psi_*$ forment un foncteur cohomologique universel (T, 2.3).

En outre, le faisceau $R^p\psi_*(\mathcal F)$ est le faisceau associé au
préfaisceau $V\leadsto H^p(f^{-1}(V),\mathcal F)$ (T, 3.7.2). Si maintenant
on suppose que $\mathcal F$ est un $\mathcal O_X$-Module,
$H^p(f^{-1}(V),\mathcal F)$ est naturellement muni d'une structure de
$\Gamma(f^{-1}(V),\mathcal O_X)$-module, donc de
$\Gamma(V,\psi_*(\mathcal O_X))$-module, et la donnée de l'homomorphisme
$\theta:\mathcal O_Y\to\psi_*(\mathcal O_X)$ permet d'en déduire une
structure de $\Gamma(V,\mathcal O_Y)$-module. Pour les structures ainsi
définies, il est clair que la restriction d'un ouvert $V$ à un ouvert
$V'\subset V$ définit un di-homomorphisme, et cela permet donc de définir sur
chacun des $R^p\psi_*(\mathcal F)$ une structure de $\mathcal O_Y$-Module;
c'est ce $\mathcal O_Y$-Module que nous noterons $R^p f_*(\mathcal F)$,
$R^p f_*$ étant ainsi défini comme un foncteur additif de $C(X)$ dans
$C(Y)$. En outre, les $R^p f_*$ forment un $\partial$-foncteur, car si
\[
0\longrightarrow\mathcal F'\longrightarrow\mathcal F
\longrightarrow\mathcal F''\longrightarrow0
\]
est une suite exacte de $\mathcal O_X$-Modules, la description des
$R^p\psi_*$ et de la structure de $\mathcal O_Y$-Module sur
$R^p\psi_*(\mathcal F)$ donnée ci-dessus montre aussitôt que
l'homomorphisme
\[
\partial:R^p\psi_*(\mathcal F'')\longrightarrow
R^{p+1}\psi_*(\mathcal F')
\]
est dans ce cas un homomorphisme de $\mathcal O_Y$-Modules. Enfin, les
$R^p f_*$ s'identifient aux foncteurs dérivés droits de $f_*$ : en effet,
tout $\mathcal O_X$-Module admet une résolution injective formée de
$\mathcal O_X$-Modules, et comme une telle résolution est formée de faisceaux
flasques de groupes abéliens (G, II, 7.1), elle peut servir à calculer les
$R^p\psi_*(\mathcal F)$, puisque $R^n\psi_*(\mathcal G)=0$ pour $n\geq1$ et
pour tout faisceau flasque $\mathcal G$ (T, 2.4.1, Remarque 3, et cor. de la
prop. 3.3.2). On conclut donc que les $R^p f_*$ forment un foncteur
cohomologique universel de $C(X)$ dans $C(Y)$ (T, 2.3).
\end{env}

\begin{env}[12.2.2]
\label{0.12.2.2-fr}
Soient $\mathcal F$ et $\mathcal G$ deux $\mathcal O_X$-Modules. Avec les
notations de (\hyperref[0.12.2.1-fr]{12.2.1}), pour tout ouvert $V$ de $Y$,
on a l'homomorphisme de cup-produit
(\hyperref[0.12.1.2.1-fr]{12.1.2.1})
\[
H^p(f^{-1}(V),\mathcal F)
\otimes_{\Gamma(f^{-1}(V),\mathcal O_X)}
H^q(f^{-1}(V),\mathcal G)
\longrightarrow
H^{p+q}(f^{-1}(V),\mathcal F\otimes_{\mathcal O_X}\mathcal G)
\]
\oldpage[0\textsubscript{III}]{58}
et il résulte aussitôt de la définition du cup-produit (G, II, 6.6) que ces
homomorphismes commutent au passage de $V$ à un sous-espace ouvert $V'$ de
$V$. D'autre part, on a un homomorphisme d'anneaux
\[
\Gamma(V,\mathcal O_Y)\longrightarrow\Gamma(V,\psi_*(\mathcal O_X))
=\Gamma(f^{-1}(V),\mathcal O_X)
\]
provenant de $\theta$, d'où un homomorphisme canonique de produits tensoriels
\[
\begin{aligned}
&H^p(f^{-1}(V),\mathcal F)\otimes_{\Gamma(V,\mathcal O_Y)}
H^q(f^{-1}(V),\mathcal G)\\
&\qquad\longrightarrow
H^p(f^{-1}(V),\mathcal F)
\otimes_{\Gamma(f^{-1}(V),\mathcal O_X)}
H^q(f^{-1}(V),\mathcal G)
\end{aligned}
\]
qui est lui aussi compatible avec la restriction de $V'$ à $V$. Par
composition, on obtient donc un homomorphisme de
$\Gamma(V,\mathcal O_Y)$-modules, qui définit un homomorphisme canonique
fonctoriel en $\mathcal F$ et $\mathcal G$ pour les faisceaux associés aux
préfaisceaux considérés :
\[
\label{0.12.2.2.1-fr}
\tag{12.2.2.1}
R^p f_*(\mathcal F)\otimes_{\mathcal O_Y}R^q f_*(\mathcal G)
\longrightarrow
R^{p+q}f_*(\mathcal F\otimes_{\mathcal O_X}\mathcal G).
\]
On notera que pour $p=q=0$, cet homomorphisme se réduit à
(\hyperref[0.4.2.2.1-fr]{0\textsubscript{I}, 4.2.2.1}).
\end{env}

\begin{proposition}[12.2.3]
\label{0.12.2.3-fr}
\emph{Pour tout $\mathcal O_X$-Module $\mathcal F$ et tout
$\mathcal O_Y$-Module localement libre de rang fini $\mathcal L$, on a des
isomorphismes canoniques fonctoriels
\[
\label{0.12.2.3.1-fr}
\tag{12.2.3.1}
R^p f_*(\mathcal F)\otimes_{\mathcal O_Y}\mathcal L
\xrightarrow{\sim}
R^p f_*(\mathcal F\otimes_{\mathcal O_X}f^*(\mathcal L)).
\]}
\end{proposition}

\begin{proof}
L'homomorphisme (\hyperref[0.12.2.3.1-fr]{12.2.3.1}) s'obtient en composant
l'homomorphisme, cas particulier de
(\hyperref[0.12.2.2.1-fr]{12.2.2.1}) :
\[
\label{0.12.2.3.2-fr}
\tag{12.2.3.2}
R^p f_*(\mathcal F)\otimes_{\mathcal O_Y}f_*(f^*(\mathcal L))
\longrightarrow
R^p f_*(\mathcal F\otimes_{\mathcal O_X}f^*(\mathcal L))
\]
avec l'homomorphisme du premier membre de
(\hyperref[0.12.2.3.1-fr]{12.2.3.1}) dans celui de
(\hyperref[0.12.2.3.2-fr]{12.2.3.2}), provenant de l'homomorphisme canonique
(\hyperref[0.4.4.3.2-fr]{0\textsubscript{I}, 4.4.3.2}). Pour vérifier que
(\hyperref[0.12.2.3.1-fr]{12.2.3.1}) est un isomorphisme lorsque $\mathcal L$
est localement libre, on peut aussitôt se ramener au cas
$\mathcal L=\mathcal O_Y$, la question étant locale sur $Y$, et les foncteurs
envisagés étant additifs en $\mathcal L$. Mais alors, la proposition se
ramène, vu la définition de (\hyperref[0.12.2.2.1-fr]{12.2.2.1}), à la
vérification du fait que l'homomorphisme correspondant de préfaisceaux est
bijectif, ce qui est immédiat en vertu de la relation
$f^*(\mathcal O_Y)=\mathcal O_X$.
\end{proof}

\begin{env}[12.2.4]
\label{0.12.2.4-fr}
Soient $(Z,\mathcal O_Z)$ un troisième espace annelé, $g:Y\to Z$ un
morphisme d'espaces annelés. On sait (G, II, 7.1 et 3.1.1) que pour tout
$\mathcal O_X$-Module injectif $\mathcal G$, $f_*(\mathcal G)$ est un
faisceau flasque de groupes abéliens, et par suite
(\hyperref[0.12.2.1-fr]{12.2.1}) on a
$R^p g_*(f_*(\mathcal G))=0$ pour tout $p>0$. Il s'ensuit (T, 2.4.1) que la
suite spectrale de Leray des foncteurs composés est applicable au foncteur
composé $g_*f_*$ : il y a une suite spectrale birégulière dont
l'aboutissement est le foncteur $R^\bullet h_*$, où $h=g\circ f$, et dont le
terme $E_2$ est donné par
\[
\label{0.12.2.4.1-fr}
\tag{12.2.4.1}
E_2^{pq}=R^p g_*(R^q f_*(\mathcal F)).
\]
\end{env}

\begin{env}[12.2.5]
\label{0.12.2.5-fr}
Sous les conditions de (\hyperref[0.12.2.4-fr]{12.2.4}), nous allons définir
directement des homomorphismes canoniques de $\mathcal O_Z$-Modules
\[
\label{0.12.2.5.1-fr}
\tag{12.2.5.1}
R^n g_*(f_*(\mathcal F))\longrightarrow R^n h_*(\mathcal F)
\]
\[
\label{0.12.2.5.2-fr}
\tag{12.2.5.2}
R^n h_*(\mathcal F)\longrightarrow g_*(R^n f_*(\mathcal F))
\]
\oldpage[0\textsubscript{III}]{59}
qu'on pourrait identifier aux «~edge-homomorphismes~» de la suite spectrale
de Leray (cf. (\hyperref[0.12.1.7-fr]{12.1.7})). Il suffit d'opérer sur les
préfaisceaux auxquels sont associés les faisceaux images supérieures
(\hyperref[0.12.2.1-fr]{12.2.1}). Pour cela, considérons un ouvert quelconque
$W$ de $Z$ et son image réciproque $g^{-1}(W)$ dans $Y$; on a un
di-homomorphisme canonique
\[
\label{0.12.2.5.3-fr}
\tag{12.2.5.3}
H^n(g^{-1}(W),f_*(\mathcal F))\longrightarrow
H^n(f^{-1}(g^{-1}(W)),f^*(f_*(\mathcal F)))
\]
les anneaux correspondants étant $\Gamma(g^{-1}(W),\mathcal O_Y)$ et
$\Gamma(h^{-1}(W),\mathcal O_X)$; d'autre part, l'homomorphisme canonique
(\hyperref[0.4.4.3.3-fr]{0\textsubscript{I}, 4.4.3.3}) fournit par
fonctorialité des homomorphismes canoniques
\[
\label{0.12.2.5.4-fr}
\tag{12.2.5.4}
H^n(h^{-1}(W),f^*(f_*(\mathcal F)))\longrightarrow
H^n(h^{-1}(W),\mathcal F)
\]
qui sont des homomorphismes de $\Gamma(h^{-1}(W),\mathcal O_X)$-modules.
Compte tenu de l'homomorphisme d'anneaux
$\Gamma(W,\mathcal O_Z)\to\Gamma(g^{-1}(W),\mathcal O_Y)$, on voit qu'en
composant (\hyperref[0.12.2.5.4-fr]{12.2.5.4}) et
(\hyperref[0.12.2.5.3-fr]{12.2.5.3}), on obtient un homomorphisme de
préfaisceaux, qui fournit l'homomorphisme de faisceaux
(\hyperref[0.12.2.5.1-fr]{12.2.5.1}).

La définition de (\hyperref[0.12.2.5.2-fr]{12.2.5.2}) est encore plus
simple; par définition, $R^n h_*(\mathcal F)$ est associé au préfaisceau
$W\leadsto H^n(f^{-1}(g^{-1}(W)),\mathcal F)$ et $R^n f_*(\mathcal F)$ au
préfaisceau $V\leadsto H^n(f^{-1}(V),\mathcal F)$; on a donc un homomorphisme
canonique
\[
H^n(f^{-1}(g^{-1}(W)),\mathcal F)\longrightarrow
\Gamma(g^{-1}(W),R^n f_*(\mathcal F)),
\]
et il est immédiat que ces homomorphismes définissent un homomorphisme de
préfaisceaux, qui à son tour définit
(\hyperref[0.12.2.5.2-fr]{12.2.5.2}).
\end{env}

\begin{env}[12.2.6]
\label{0.12.2.6-fr}
Sous les hypothèses de (\hyperref[0.12.2.4-fr]{12.2.4}), soient
$\mathcal F$, $\mathcal G$, $\mathcal K$ trois $\mathcal O_X$-Modules et
$u:\mathcal F\otimes\mathcal G\to\mathcal K$ un
$\mathcal O_X$-homomorphisme. On a alors des diagrammes commutatifs
\[
\xymatrix@C=10pt@R=24pt{
R^p g_*(f_*(\mathcal F))\otimes_{\mathcal O_Z}R^q g_*(f_*(\mathcal G))
\ar[r]\ar[d]&R^{p+q}g_*(f_*(\mathcal K))\ar[d]\\
R^p h_*(\mathcal F)\otimes_{\mathcal O_Z}R^q h_*(\mathcal G)
\ar[r]&R^{p+q}h_*(\mathcal K)
}
\tag{12.2.6.1}\label{0.12.2.6.1-fr}
\]
et
\[
\xymatrix@C=10pt@R=24pt{
R^p h_*(\mathcal F)\otimes_{\mathcal O_Z}R^q h_*(\mathcal G)
\ar[r]\ar[d]&R^{p+q}h_*(\mathcal K)\ar[d]\\
g_*(R^p f_*(\mathcal F))\otimes_{\mathcal O_Z}g_*(R^q f_*(\mathcal G))
\ar[r]&g_*(R^{p+q}f_*(\mathcal K))
}
\tag{12.2.6.2}\label{0.12.2.6.2-fr}
\]
\oldpage[0\textsubscript{III}]{60}
où les flèches horizontales proviennent de
(\hyperref[0.12.2.2.1-fr]{12.2.2.1}) (la dernière combinée avec
(\hyperref[0.4.2.2.1-fr]{0\textsubscript{I}, 4.2.2.1})) et les flèches
verticales des homomorphismes (\hyperref[0.12.2.5.1-fr]{12.2.5.1}) et
(\hyperref[0.12.2.5.2-fr]{12.2.5.2}) respectivement.

Il suffit en effet de le vérifier pour les homomorphismes correspondants de
préfaisceaux; si on revient aux définitions données dans
(\hyperref[0.12.2.2-fr]{12.2.2}) et
(\hyperref[0.12.2.5-fr]{12.2.5}) pour ces homomorphismes, on est aussitôt
ramené, pour (\hyperref[0.12.2.6.1-fr]{12.2.6.1}), aux diagrammes commutatifs
(\hyperref[0.12.1.5.2-fr]{12.1.5.2}); la vérification est encore plus simple
pour (\hyperref[0.12.2.6.2-fr]{12.2.6.2}).
\end{env}

\subsection{Compléments sur les foncteurs Ext de faisceaux}
\label{subsection:0.12.3-fr}

\begin{env}[12.3.1]
\label{0.12.3.1-fr}
Considérons un espace annelé $(X,\mathcal O_X)$; nous ne reviendrons pas sur
la définition et les principales propriétés des bifoncteurs
$\operatorname{Ext}_{\mathcal O_X}^p(X;\mathcal F,\mathcal G)$ de la
catégorie des $\mathcal O_X$-Modules dans celle des
$\Gamma(X,\mathcal O_X)$-modules, et
$\mathcal{E}xt_{\mathcal O_X}^p(\mathcal F,\mathcal G)$ de la catégorie des
$\mathcal O_X$-Modules dans elle-même, ni sur la suite spectrale birégulière
$E(\mathcal F,\mathcal G)$ qui les relie (T, 4.2 et G, II, 7.3).
\end{env}

\begin{env}[12.3.2]
\label{0.12.3.2-fr}
On définit de la même manière que dans (M, XIV, 1) la notion d'extension d'un
$\mathcal O_X$-Module $\mathcal F$ par un $\mathcal O_X$-Module $\mathcal G$
et la loi de composition entre classes d'extensions équivalentes : les
raisonnements faits pour les modules s'adaptent en effet de façon évidente à
une catégorie abélienne quelconque. La seconde démonstration de (M, XIV, 1.1),
qui n'utilise que l'existence de plongements dans les objets injectifs, est
alors encore valable pour la catégorie des $\mathcal O_X$-Modules, et montre
donc que $\operatorname{Ext}_{\mathcal O_X}^1(X;\mathcal F,\mathcal G)$
s'identifie canoniquement au groupe abélien des classes d'extensions de
$\mathcal F$ par $\mathcal G$.
\end{env}

\begin{proposition}[12.3.3]
\label{0.12.3.3-fr}
\emph{Soit $(X,\mathcal O_X)$ un espace annelé tel que le faisceau d'anneaux
$\mathcal O_X$ soit cohérent. Alors, pour tout couple de
$\mathcal O_X$-Modules cohérents $\mathcal F$, $\mathcal G$ et tout
$p\geq0$, $\mathcal{E}xt_{\mathcal O_X}^p(\mathcal F,\mathcal G)$ est un
$\mathcal O_X$-Module cohérent.}
\end{proposition}

\begin{proof}
Notons que les $\mathcal{E}xt_{\mathcal O_X}^p(\mathcal F,\mathcal G)$
forment un foncteur cohomologique contravariant en $\mathcal F$. Puisque
$\mathcal F$ est cohérent, il existe pour tout $p$ et tout point $x\in X$ un
voisinage ouvert $U$ de $X$ et une suite exacte de
$(\mathcal O_X|U)$-Modules
\[
0\longrightarrow\mathcal R\longrightarrow\mathcal L_{p-1}
\longrightarrow\cdots\longrightarrow\mathcal L_0
\longrightarrow\mathcal F|U\longrightarrow0
\]
où chacun des $\mathcal L_i$ ($0\leq i\leq p-1$) est isomorphe à un
$\mathcal O_X^{n_i}|U$ et $\mathcal R$ est cohérent : cela résulte par
récurrence sur $p$ de (\hyperref[0.5.3.2-fr]{0\textsubscript{I}, 5.3.2}) et
(\hyperref[0.5.3.4-fr]{0\textsubscript{I}, 5.3.4}), vu l'hypothèse que
$\mathcal O_X$ est cohérent.

Notons maintenant que, pour $p\geq1$, on a
\[
\operatorname{Ext}_{\mathcal O_X|U}^p(\mathcal L|U,\mathcal G|U)=0
\]
pour tout $\mathcal O_X$-Module $\mathcal L$ tel que $\mathcal L|U$ soit
isomorphe à un $\mathcal O_X^n|U$ (T, 4.2.3); le raisonnement de (M, V, 7.2)
s'applique donc au foncteur cohomologique contravariant
\[
\mathcal F\leadsto
\mathcal{H}om_{\mathcal O_X|U}(\mathcal F|U,\mathcal G|U),
\]
et donne une suite exacte
\[
\mathcal{H}om_{\mathcal O_X|U}(\mathcal L_{p-1},\mathcal G|U)
\longrightarrow
\mathcal{H}om_{\mathcal O_X|U}(\mathcal R,\mathcal G|U)
\longrightarrow
\mathcal{E}xt_{\mathcal O_X|U}^p(\mathcal F|U,\mathcal G|U)
\longrightarrow0
\]
et comme les deux premiers termes de cette suite sont des
$(\mathcal O_X|U)$-Modules cohérents
(\hyperref[0.5.3.5-fr]{0\textsubscript{I}, 5.3.5}), il en est de même du
troisième (\hyperref[0.5.3.4-fr]{0\textsubscript{I}, 5.3.4}).
\end{proof}

\oldpage[0\textsubscript{III}]{61}
\begin{proposition}[12.3.4]
\label{0.12.3.4-fr}
\emph{Soit $f:X\to Y$ un morphisme plat d'espaces annelés, et soient
$\mathcal F$, $\mathcal G$ deux $\mathcal O_Y$-Modules.}
\begin{enumerate}
\item[{\rm(i)}] \emph{Il existe un homomorphisme de bifoncteurs
cohomologiques
\[
\label{0.12.3.4.1-fr}
\tag{12.3.4.1}
f^*(\mathcal{E}xt_{\mathcal O_Y}^{\bullet}(\mathcal F,\mathcal G))
\xrightarrow{\sim}
\mathcal{E}xt_{\mathcal O_X}^{\bullet}(f^*(\mathcal F),f^*(\mathcal G))
\]
se réduisant en degré $0$ à l'homomorphisme canonique
(\hyperref[0.4.4.6-fr]{0\textsubscript{I}, 4.4.6}).}

\item[{\rm(ii)}] \emph{Il existe un morphisme canonique de suites spectrales
\[
\label{0.12.3.4.2-fr}
\tag{12.3.4.2}
E(\mathcal F,\mathcal G)\longrightarrow
E(f^*(\mathcal F),f^*(\mathcal G))
\]
qui, pour les termes $E_2$, se réduit aux homomorphismes
\[
\label{0.12.3.4.3-fr}
\tag{12.3.4.3}
H^p(Y,\mathcal{E}xt_{\mathcal O_Y}^q(\mathcal F,\mathcal G))
\longrightarrow
H^p(X,\mathcal{E}xt_{\mathcal O_X}^q(f^*(\mathcal F),f^*(\mathcal G)))
\]
déduits de (\hyperref[0.12.3.4.1-fr]{12.3.4.1}) et de
(\hyperref[0.12.1.3.1-fr]{12.1.3.1}).}
\end{enumerate}
\end{proposition}

\begin{proof}
\begin{enumerate}
\item[{\rm(i)}]
Comme $f^*$ est un foncteur exact dans la catégorie des
$\mathcal O_Y$-Modules (\hyperref[0.6.7.2-fr]{0\textsubscript{I}, 6.7.2}),
les foncteurs
\[
\mathcal G\leadsto
f^*(\mathcal{H}om_{\mathcal O_Y}(\mathcal F,\mathcal G))
\quad\hbox{et}\quad
\mathcal G\leadsto
\mathcal{H}om_{\mathcal O_X}(f^*(\mathcal F),f^*(\mathcal G))
\]
sont exacts à gauche; on déduit canoniquement de
(\hyperref[0.4.4.6-fr]{0\textsubscript{I}, 4.4.6}) un homomorphisme de leurs
foncteurs dérivés.

Pour calculer ces derniers, on prend une résolution injective
$\mathcal L^\bullet=(\mathcal L^i)$ de $\mathcal G$, et on a donc des
morphismes
\[
\mathcal H^p(f^*(\mathcal{H}om_{\mathcal O_Y}
(\mathcal F,\mathcal L^\bullet)))
\longrightarrow
\mathcal H^p(\mathcal{H}om_{\mathcal O_X}
(f^*(\mathcal F),f^*(\mathcal L^\bullet)))
\]
de cohomologies de complexes de faisceaux. D'ailleurs, en vertu de
l'exactitude de $f^*$, on a
\[
\mathcal H^p(f^*(\mathcal{H}om_{\mathcal O_Y}
(\mathcal F,\mathcal L^\bullet)))
=f^*(\mathcal H^p(\mathcal{H}om_{\mathcal O_Y}
(\mathcal F,\mathcal L^\bullet)))
=f^*(\mathcal{E}xt_{\mathcal O_Y}^p(\mathcal F,\mathcal G))
\]
par définition. D'autre part, l'exactitude de $f^*$ entraîne que
$f^*(\mathcal L^\bullet)$ est une résolution de $f^*(\mathcal G)$; si
$\mathcal L'^\bullet=(\mathcal L'^{i})$ est une résolution injective de
$f^*(\mathcal G)$ dans la catégorie des $\mathcal O_X$-Modules, il y a donc
un homomorphisme de complexes
$f^*(\mathcal L^\bullet)\to\mathcal L'^\bullet$, déterminé à une homotopie
près, et qui définit par suite un homomorphisme bien déterminé en cohomologie;
composant cet homomorphisme avec l'homomorphisme défini plus haut, on obtient
(\hyperref[0.12.3.4.1-fr]{12.3.4.1}).

\item[{\rm(ii)}]
Avec les notations précédentes, on a un homomorphisme de complexes de
faisceaux de $\mathcal O_X$-Modules
\[
f^*(\mathcal{H}om_{\mathcal O_Y}(\mathcal F,\mathcal L^\bullet))
\longrightarrow
\mathcal{H}om_{\mathcal O_X}(f^*(\mathcal F),\mathcal L'^\bullet).
\]
Soit $\mathcal M^{\bullet\bullet}$ une résolution injective de
Cartan-Eilenberg du complexe
$\mathcal{H}om_{\mathcal O_Y}(\mathcal F,\mathcal L^\bullet)$ dans la
catégorie des $\mathcal O_Y$-Modules; alors, en vertu de l'exactitude du
foncteur $f^*$, $f^*(\mathcal M^{\bullet\bullet})$ est une résolution de
Cartan-Eilenberg du complexe
$f^*(\mathcal{H}om_{\mathcal O_Y}(\mathcal F,\mathcal L^\bullet))$; si
$\mathcal M'^{\bullet\bullet}$ est une résolution injective de
Cartan-Eilenberg du complexe
$\mathcal{H}om_{\mathcal O_X}(f^*(\mathcal F),\mathcal L'^\bullet)$, il y a
donc (\hyperref[0.11.4.2-fr]{11.4.2}), un homomorphisme (déterminé à
homotopie près)
\[
f^*(\mathcal M^{\bullet\bullet})\longrightarrow
\mathcal M'^{\bullet\bullet}
\]
compatible avec l'homomorphisme considéré ci-dessus, autrement dit un
$f$-morphisme $\mathcal M^{\bullet\bullet}\to
\mathcal M'^{\bullet\bullet}$ de bicomplexes de faisceaux. On en déduit un
di-homomorphisme
$\Gamma(Y,\mathcal M^{\bullet\bullet})\to
\Gamma(X,\mathcal M'^{\bullet\bullet})$ de bicomplexes de modules, déterminé
à homotopie près, et un morphisme bien déterminé de suites spectrales
(\hyperref[0.11.3.2-fr]{11.3.2}), qui n'est autre que le morphisme
(\hyperref[0.12.3.4.2-fr]{12.3.4.2}) cherché, la caractérisation de
(\hyperref[0.12.3.4.3-fr]{12.3.4.3}) se déduisant aussitôt des définitions.
\end{enumerate}
\end{proof}

\begin{proposition}[12.3.5]
\label{0.12.3.5-fr}
\emph{Sous les hypothèses de (\hyperref[0.12.3.4-fr]{12.3.4}), supposons en
outre le faisceau d'anneaux $\mathcal O_Y$ cohérent; alors, pour tout
$\mathcal O_Y$-Module cohérent $\mathcal F$, les homomorphismes canoniques
(\hyperref[0.12.3.4.1-fr]{12.3.4.1}) sont bijectifs.}
\end{proposition}

\begin{proof}
\oldpage[0\textsubscript{III}]{62}
La question étant locale sur $Y$, on peut supposer qu'il existe une suite
exacte
\[
0\longrightarrow\mathcal R\longrightarrow\mathcal O_Y^n
\longrightarrow\mathcal F\longrightarrow0,
\]
et $\mathcal R$ est alors aussi un $\mathcal O_Y$-Module cohérent
(\hyperref[0.5.3.4-fr]{0\textsubscript{I}, 5.3.4}). Pour prouver que les
homomorphismes
\[
f^*(\mathcal{E}xt_{\mathcal O_Y}^p(\mathcal F,\mathcal G))
\longrightarrow
\mathcal{E}xt_{\mathcal O_X}^p(f^*(\mathcal F),f^*(\mathcal G))
\]
sont bijectifs, raisonnons par récurrence sur $p$, la proposition résultant de
(\hyperref[0.6.7.6.1-fr]{0\textsubscript{I}, 6.7.6.1}) lorsque $p=0$. Or, on
a le diagramme commutatif
\[
{\tiny
\xymatrix@C=2pt@R=24pt{
f^*(\mathcal{E}xt_{\mathcal O_Y}^{p-1}(\mathcal O_Y^n,\mathcal G))
\ar[r]\ar[d]&
f^*(\mathcal{E}xt_{\mathcal O_Y}^{p-1}(\mathcal R,\mathcal G))
\ar[r]^\partial\ar[d]&
f^*(\mathcal{E}xt_{\mathcal O_Y}^{p}(\mathcal F,\mathcal G))
\ar[r]\ar[d]&
f^*(\mathcal{E}xt_{\mathcal O_Y}^{p}(\mathcal O_Y^n,\mathcal G))
\ar[d]\\
\mathcal{E}xt_{\mathcal O_X}^{p-1}(\mathcal O_X^n,f^*(\mathcal G))
\ar[r]&
\mathcal{E}xt_{\mathcal O_X}^{p-1}(f^*(\mathcal R),f^*(\mathcal G))
\ar[r]^\partial&
\mathcal{E}xt_{\mathcal O_X}^{p}(f^*(\mathcal F),f^*(\mathcal G))
\ar[r]&
\mathcal{E}xt_{\mathcal O_X}^{p}(\mathcal O_X^n,f^*(\mathcal G)).
}}
\]
puisque $f^*(\mathcal O_Y)=\mathcal O_X$; comme $f^*$ est exact, les deux
lignes sont exactes. En outre, on a
\[
\mathcal{E}xt_{\mathcal O_Y}^p(\mathcal O_Y^n,\mathcal G)=0
\quad\hbox{pour tout }p>0
\quad\hbox{et de même}\quad
\mathcal{E}xt_{\mathcal O_X}^p(\mathcal O_X^n,f^*(\mathcal G))=0
\quad\hbox{pour tout }p>0
\]
(T, 4.2.3). Vu l'hypothèse de récurrence, les deux premières flèches
verticales du diagramme précédent sont des isomorphismes, et les termes de
droite sont 0, donc
\[
f^*(\mathcal{E}xt_{\mathcal O_Y}^p(\mathcal F,\mathcal G))
\longrightarrow
\mathcal{E}xt_{\mathcal O_X}^p(f^*(\mathcal F),f^*(\mathcal G))
\]
est un isomorphisme.
\end{proof}

\subsection{Hypercohomologie du foncteur image directe}
\label{subsection:0.12.4-fr}

\begin{env}[12.4.1]
\label{0.12.4.1-fr}
Soient $(X,\mathcal O_X)$, $(Y,\mathcal O_Y)$ deux espaces annelés,
$f:X\to Y$ un morphisme d'espaces annelés. On peut prendre
l'\emph{hypercohomologie} de $f_*$ par rapport à un complexe quelconque de
$\mathcal O_X$-Modules
$\mathcal K^\bullet=(\mathcal K^i)_{i\in\mathbb Z}$
(\hyperref[0.11.4.4-fr]{11.4.4}), car dans la catégorie abélienne des
$\mathcal O_Y$-Modules, les limites inductives filtrantes existent et sont
exactes (T, 3.1.1). Les $\mathcal O_Y$-Modules d'hypercohomologie
$R^p f_*(\mathcal K^\bullet)$ se noteront aussi
$\mathcal H^p(f,\mathcal K^\bullet)$ ou
$\mathcal H_f^p(\mathcal K^\bullet)$. Rappelons que
$\mathcal H^\bullet(f,\mathcal K^\bullet)$ est la cohomologie du bicomplexe
de $\mathcal O_Y$-Modules $f_*(\mathcal L^{\bullet\bullet})$, où
$\mathcal L^{\bullet\bullet}$ est une résolution injective de
Cartan-Eilenberg de $\mathcal K^\bullet$ dans la catégorie des
$\mathcal O_X$-Modules~; $\mathcal H^\bullet(f,\mathcal K^\bullet)$ est
l'aboutissement de deux suites spectrales
${}'E(f,\mathcal K^\bullet)$ et ${}''E(f,\mathcal K^\bullet)$ dont les
termes $E_2$ sont donnés par
\begin{equation}
\label{0.12.4.1.1-fr}
\tag{12.4.1.1}
{}'E_2^{pq}=\mathcal H^p(\mathcal H^q(f,\mathcal K^\bullet))
\end{equation}
\begin{equation}
\label{0.12.4.1.2-fr}
\tag{12.4.1.2}
{}''E_2^{pq}=\mathcal H^p(f,\mathcal H^q(\mathcal K^\bullet))
\quad(=R^p f_*(\mathcal H^q(\mathcal K^\bullet))).
\end{equation}
On a, dans ces formules, adopté la notation générale $T(A^\bullet)$ pour
le transformé d'un complexe par un foncteur
(\hyperref[0.11.2.1-fr]{11.2.1}), et on écrit
$\mathcal H^p(f,\mathcal F)$ au lieu de $R^p f_*(\mathcal F)$ pour un
$\mathcal O_X$-Module $\mathcal F$. Rappelons encore que la suite
${}'E(f,\mathcal K^\bullet)$ est toujours \emph{régulière}~; les deux
suites spectrales ${}'E(f,\mathcal K^\bullet)$ et
${}''E(f,\mathcal K^\bullet)$ sont \emph{birégulières} lorsque
$\mathcal K^\bullet$ est limité inférieurement,
\oldpage[0\textsubscript{III}]{63}
ou lorsqu'il existe un entier $m$ tel que tout $\mathcal O_X$-Module
admette une résolution flasque de longueur $\leq m$
(\hyperref[0.11.4.4-fr]{11.4.4}).
\end{env}

\begin{env}[12.4.2]
\label{0.12.4.2-fr}
Nous désignerons de même par $\mathbf H^\bullet(X,\mathcal K^\bullet)$
l'hypercohomologie du foncteur $\Gamma$ par rapport à un complexe
$\mathcal K^\bullet$ de $\mathcal O_X$-Modules~; les
$\mathbf H^p(X,\mathcal K^\bullet)$ sont donc des
$\Gamma(X,\mathcal O_X)$-modules. On peut d'ailleurs considérer
$\mathbf H^\bullet(X,\mathcal K^\bullet)$ comme un cas particulier de
$\mathcal H^\bullet(f,\mathcal K^\bullet)$, où $f$ est un morphisme de
$(X,\mathcal O_X)$ sur un espace annelé réduit à un point muni de l'anneau
$\Gamma(X,\mathcal O_X)$.

Pour tout ouvert $V$ de $X$, nous écrirons
$\mathbf H^\bullet(V,\mathcal K^\bullet)$ au lieu de
$\mathbf H^\bullet(V,\mathcal K^\bullet|V)$.
\end{env}

\begin{proposition}[12.4.3]
\label{0.12.4.3-fr}
\emph{Pour tout entier $p\in\mathbb Z$, le $\mathcal O_Y$-Module
$\mathcal H^p(f,\mathcal K^\bullet)$ est canoniquement isomorphe au
faisceau associé au préfaisceau
$U\mapsto\mathbf H^p(f^{-1}(U),\mathcal K^\bullet)$ sur $Y$.}
\end{proposition}

\begin{proof}
En effet, avec les notations de
(\hyperref[0.12.4.1-fr]{12.4.1}), le faisceau de cohomologie
$\mathcal H^p(f_*(\mathcal L^{\bullet\bullet}))$ est associé au préfaisceau
\[
U\longmapsto H^p(\Gamma(U,f_*(\mathcal L^{\bullet\bullet})))
=H^p(\Gamma(f^{-1}(U),\mathcal L^{\bullet\bullet})).
\]
Mais il est clair que
$\mathcal L^{\bullet\bullet}|f^{-1}(U)$ est une résolution injective de
Cartan-Eilenberg de $\mathcal K^\bullet|f^{-1}(U)$ (T, 3.1.3), donc
\[
H^p(\Gamma(f^{-1}(U),\mathcal L^{\bullet\bullet}))
=\mathbf H^p(f^{-1}(U),\mathcal K^\bullet)
\]
par définition.
\end{proof}

\begin{proposition}[12.4.4]
\label{0.12.4.4-fr}
\emph{L'hypercohomologie $\mathcal H^\bullet(f,\mathcal K^\bullet)$ est
un foncteur cohomologique en $\mathcal K^\bullet$ dans chacun des cas
suivants~:}
\begin{enumerate}
\item[\emph{a)}]
\emph{$\mathcal K^\bullet$ varie dans la catégorie des complexes limités
inférieurement.}
\item[\emph{b)}]
\emph{Il existe un entier $m$ tel que tout $\mathcal O_X$-Module admette
une résolution flasque de longueur $\leq m$.}
\item[\emph{c)}]
\emph{$X$ est un espace noethérien.}
\end{enumerate}
\end{proposition}

\begin{proof}
Les cas a) et b) sont des cas particuliers de
(\hyperref[0.11.5.4-fr]{11.5.4}). D'autre part, le cas c) se déduit de
(\hyperref[0.11.5.2-fr]{11.5.2}), car on sait que dans ce cas, le foncteur
$f_*$ permute aux limites inductives (G, II, 3.10.1).
\end{proof}

\begin{env}[12.4.5]
\label{0.12.4.5-fr}
Considérons maintenant un recouvrement ouvert
$\mathfrak U=(U_\alpha)$ de $X$, et pour tout complexe de
\emph{préfaisceaux} $\mathcal K^\bullet=(\mathcal K^j)$ sur $X$, le
bicomplexe $\check C^\bullet(\mathfrak U,\mathcal K^\bullet)$, dont le
composant d'indices $(i,j)$ est
$\check C^i(\mathfrak U,\mathcal K^j)$, groupe des $i$-cochaînes alternées
du nerf de $\mathfrak U$ à valeurs dans $\mathcal K^j$ (G, II, 5.1). Nous
dirons que la cohomologie de ce bicomplexe est
l'\emph{hypercohomologie du recouvrement $\mathfrak U$ à coefficients dans
$\mathcal K^\bullet$}, et nous la noterons
\[
\mathbf H^\bullet(\mathfrak U,\mathcal K^\bullet)
=H^\bullet(\check C^\bullet(\mathfrak U,\mathcal K^\bullet)).
\]
La suite spectrale de Leray d'un recouvrement (T, 3.8.1 et G, II, 5.9.1)
se généralise comme suit à l'hypercohomologie~:
\end{env}

\begin{proposition}[12.4.6]
\label{0.12.4.6-fr}
\emph{Soit $\mathcal K^\bullet=(\mathcal K^i)$ un complexe de
$\mathcal O_X$-Modules. Il existe un foncteur spectral régulier en
$\mathcal K^\bullet$ ayant pour aboutissement l'hypercohomologie
$\mathbf H^\bullet(X,\mathcal K^\bullet)$, et dont le terme $E_2$ est donné
par}
\begin{equation}
\label{0.12.4.6.1-fr}
\tag{12.4.6.1}
E_2^{pq}=\check H^p(\mathfrak U,h^q(\mathcal K^\bullet)),
\end{equation}
\emph{où $h^q(\mathcal K^\bullet)$ désigne le complexe de préfaisceaux
$V\mapsto H^q(V,\mathcal K^\bullet)$ sur $X$. La suite spectrale précédente
est birégulière si $\mathcal K^\bullet$ est limité inférieurement.}
\end{proposition}

\begin{proof}
Considérons une résolution injective de Cartan-Eilenberg
$\mathcal L^{\bullet\bullet}$ de $\mathcal K^\bullet$, et le tricomplexe
$\check C^\bullet(\mathfrak U,\mathcal L^{\bullet\bullet})
=(\check C^i(\mathfrak U,\mathcal L^{jk}))$~; considérons d'abord ce
tricomplexe comme un bicomplexe pour les degrés $i$ et $j+k$. Comme $i$ ne
prend que des valeurs $\geq0$, la seconde suite spectrale de ce bicomplexe
est régulière (\hyperref[0.11.3.3-fr]{11.3.3}) et dégénérée, car on a
$\check H^q(\mathfrak U,\mathcal L^{jk})=0$ pour tout $q>0$, les
$\mathcal O_X$-Modules $\mathcal L^{jk}$ étant des faisceaux
\oldpage[0\textsubscript{III}]{64}
flasques (G, II, 5.2.3). On a par suite
(\hyperref[0.11.1.6-fr]{11.1.6}) un isomorphisme canonique
\[
H^n(\check C^\bullet(\mathfrak U,\mathcal L^{\bullet\bullet}))
\xrightarrow{\sim}
H^n(\Gamma(X,\mathcal L^{\bullet\bullet}))
\]
(en vertu de (G, II, 5.2.2)), donc par définition
(\hyperref[0.12.4.2-fr]{12.4.2}) un isomorphisme
\[
H^n(\check C^\bullet(\mathfrak U,\mathcal L^{\bullet\bullet}))
\xrightarrow{\sim}\mathbf H^n(X,\mathcal K^\bullet).
\]
Considérons d'autre part le tricomplexe
$\check C^\bullet(\mathfrak U,\mathcal L^{\bullet\bullet})$ comme un
bicomplexe pour les degrés $i+j$ et $k$. Comme $k$ ne prend que des valeurs
$\geq0$, la première suite spectrale de ce bicomplexe est toujours
régulière~; elle est birégulière si $\mathcal L^{jk}=0$ pour $j<j_0$,
c'est-à-dire lorsque $\mathcal K^\bullet$ est limité inférieurement
(\hyperref[0.11.3.3-fr]{11.3.3}). Cette suite spectrale est la suite
cherchée~; en effet, pour tout $j$, $\mathcal L^{j,\bullet}$ est une
résolution injective de $\mathcal K^j$~; par suite,
$H^q(\check C^i(\mathfrak U,\mathcal L^{j,\bullet}))$ n'est autre que le
complexe de cochaînes
$\check C^i(\mathfrak U,h^q(\mathcal K^j))$, ce qui termine la
démonstration.
\end{proof}

\begin{corollary}[12.4.7]
\label{0.12.4.7-fr}
\emph{Si, pour tout simplexe $\sigma$ du nerf de $\mathfrak U$, et pour
tout entier $i$, on a
$H^q(U_\sigma,\mathcal K^i)=0$ pour $q>0$, alors on a un isomorphisme
canonique}
\begin{equation}
\label{0.12.4.7.1-fr}
\tag{12.4.7.1}
\mathbf H^\bullet(\mathfrak U,\mathcal K^\bullet)
\xrightarrow{\sim}\mathbf H^\bullet(X,\mathcal K^\bullet).
\end{equation}
\end{corollary}

\begin{proof}
En effet, l'hypothèse entraîne que
$\check C^\bullet(\mathfrak U,h^q(\mathcal K^i))=0$ pour $q>0$, donc
$E_2^{pq}=0$ pour $q>0$~; la suite
(\hyperref[0.12.4.6.1-fr]{12.4.6.1}) étant dégénérée et régulière, la
conclusion résulte de la définition
(\hyperref[0.12.4.5-fr]{12.4.5}) de
$\mathbf H^\bullet(\mathfrak U,\mathcal K^\bullet)$ et de
(\hyperref[0.11.1.6-fr]{11.1.6}).
\end{proof}

\begin{env}[12.4.8]
\label{0.12.4.8-fr}
Soit $(X',\mathcal O_{X'})$ un second espace annelé, et soit
$f=(\psi,\theta)$ un morphisme de $X'$ dans $X$. Par la même méthode que
dans (\hyperref[0.12.1.3-fr]{12.1.3}) et
(\hyperref[0.12.1.4-fr]{12.1.4}), on définit un di-homomorphisme pour
l'hypercohomologie d'un complexe $\mathcal K^\bullet$ de
$\mathcal O_X$-Modules
\begin{equation}
\label{0.12.4.8.1-fr}
\tag{12.4.8.1}
\mathbf H^p(X,\mathcal K^\bullet)
\longrightarrow\mathbf H^p(X',f^*(\mathcal K^\bullet)).
\end{equation}
On part d'une résolution injective de Cartan-Eilenberg
$\mathcal L^{\bullet\bullet}$ de $\mathcal K^\bullet$, et comme $\psi^*$
est exact, $\psi^*(\mathcal L^{\bullet\bullet})$ est une résolution de
Cartan-Eilenberg de $\psi^*(\mathcal K^\bullet)$ dans la catégorie des
$\psi^*(\mathcal O_X)$-Modules~; il y a alors un morphisme
$\psi^*(\mathcal L^{\bullet\bullet})\to
\mathcal L'^{\bullet\bullet}$, où $\mathcal L'^{\bullet\bullet}$ est une
résolution injective de Cartan-Eilenberg de
$\psi^*(\mathcal K^\bullet)$, et on en déduit un morphisme pour la
cohomologie~:
\[
\mathbf H^\bullet(X,\mathcal K^\bullet)
\longrightarrow\mathbf H^\bullet(X',\psi^*(\mathcal K^\bullet))~;
\]
par composition avec le morphisme déduit par fonctorialité de
$\psi^*(\mathcal K^\bullet)\to f^*(\mathcal K^\bullet)$, on obtient le
morphisme (\hyperref[0.12.4.8.1-fr]{12.4.8.1}) cherché.

Partant de (\hyperref[0.12.4.8.1-fr]{12.4.8.1}) et de
(\hyperref[0.12.4.3-fr]{12.4.3}), on peut alors, en raisonnant comme dans
(\hyperref[0.12.2.5-fr]{12.2.5}), définir, pour deux morphismes
$f:X\to Y$, $g:Y\to Z$ d'espaces annelés, des homomorphismes pour
l'hypercohomologie d'un complexe $\mathcal K^\bullet$ de
$\mathcal O_X$-Modules
\begin{equation}
\label{0.12.4.8.2-fr}
\tag{12.4.8.2}
\mathcal H^n(g,f_*(\mathcal K^\bullet))
\longrightarrow\mathcal H^n(h,\mathcal K^\bullet)
\end{equation}
\begin{equation}
\label{0.12.4.8.3-fr}
\tag{12.4.8.3}
\mathcal H^n(h,\mathcal K^\bullet)
\longrightarrow g_*(\mathcal H^n(f,\mathcal K^\bullet)).
\end{equation}
Nous laissons au lecteur le détail des définitions.
\end{env}

\section{Limites projectives en algèbre homologique}
\label{section:0.13-fr}

\subsection{La condition de Mittag-Leffler}
\label{subsection:0.13.1-fr}

\begin{env}[13.1.1]
\label{0.13.1.1-fr}
Soit $\mathcal C$ une catégorie abélienne dans laquelle les produits infinis
existent (axiome AB~3* de T, 1.5))~; alors la borne inférieure d'une famille
de sous-objets d'un objet de $\mathcal C$ existe, et tout système projectif
d'objets de $\mathcal C$ admet une limite projective, qui est un foncteur
exact à gauche du système projectif considéré (T, 1.8). Soit
$(A_\alpha,f_{\alpha\beta})$
\oldpage[0\textsubscript{III}]{65}
un système projectif d'objets de $\mathcal C$ dont l'ensemble d'indices $I$
est filtrant à droite~; soient $A=\varprojlim A_\alpha$ et, pour tout
$\alpha\in I$, soit $f_\alpha:A\to A_\alpha$ le morphisme canonique. Pour
tout $\alpha\in I$, les $f_{\alpha\beta}(A_\beta)$ pour
$\alpha\leq\beta$ forment une famille filtrante décroissante de sous-objets
de $A_\alpha$~; le sous-objet
\[
A'_\alpha=\inf_{\beta\geq\alpha}f_{\alpha\beta}(A_\beta)
\]
est dit sous-objet des «~images universelles~» dans $A_\alpha$~; il est
clair que $f_\alpha(A)\subset A'_\alpha$ et
$f_{\alpha\beta}(A'_\beta)\subset A'_\alpha$ pour
$\alpha\leq\beta$~; donc
$(A'_\alpha,f_{\alpha\beta}|A'_\beta)$ est un système projectif et
$A=\varprojlim A'_\alpha$.
\end{env}

\begin{env}[13.1.2]
\label{0.13.1.2-fr}
Étant donné un système projectif $(A_\alpha,f_{\alpha\beta})$ dans
$\mathcal C$, on appelle \emph{condition de Mittag-Leffler} la condition
suivante~:
\begin{itemize}
\item[\textbf{(ML)}]
\emph{Pour tout indice $\alpha$, il existe $\beta\geq\alpha$ tel que, pour
tout $\gamma\geq\beta$, on ait
$f_{\alpha\gamma}(A_\gamma)=f_{\alpha\beta}(A_\beta)$.}
\end{itemize}
Il est clair que si les $f_{\alpha\beta}$ sont des \emph{épimorphismes}, la
condition (ML) est vérifiée. Inversement, si (ML) est vérifiée, et si pour
tout $\alpha\in I$, $A'_\alpha$ est le sous-objet des «~images
universelles~» dans $A_\alpha$, la restriction de $f_{\alpha\beta}$ à
$A'_\beta$ est un \emph{épimorphisme}
$A'_\beta\to A'_\alpha$ pour $\alpha\leq\beta$~: en effet, si
$\gamma\geq\beta$ est tel que
$f_{\beta\delta}(A_\delta)=f_{\beta\gamma}(A_\gamma)$ pour
$\delta\geq\gamma$, on a
$A'_\beta=f_{\beta\gamma}(A_\gamma)$ et cela entraîne d'autre part
$f_{\alpha\delta}(A_\delta)=f_{\alpha\gamma}(A_\gamma)$ pour
$\delta\geq\gamma$, donc
$A'_\alpha=f_{\alpha\gamma}(A_\gamma)=f_{\alpha\beta}(A'_\beta)$.

Notons aussi que la condition (ML) est vérifiée lorsque les objets
$A_\alpha$ sont \emph{artiniens} dans $\mathcal C$, c'est-à-dire que toute
famille de sous-objets de $A_\alpha$ admet un élément minimal~: un élément
minimal de la famille filtrante décroissante
$(f_{\alpha\beta}(A_\beta))$ de sous-objets de $A_\alpha$ est en effet
alors nécessairement le plus petit de ces sous-objets.
\end{env}

\begin{remark}[13.1.3]
\label{0.13.1.3-fr}
La condition (ML) peut se formuler également lorsque $\mathcal C$ est par
exemple la catégorie des ensembles~; on peut alors encore définir le
sous-ensemble des «~images universelles~» de $A_\alpha$ et les remarques
faites à ce sujet dans (\hyperref[0.13.1.1-fr]{13.1.1}) et
(\hyperref[0.13.1.2-fr]{13.1.2}) restent valables.
\end{remark}

\subsection{La condition de Mittag-Leffler pour les groupes abéliens}
\label{subsection:0.13.2-fr}

\begin{proposition}[13.2.1]
\label{0.13.2.1-fr}
\emph{Soit}
\[
0\longrightarrow A_\alpha\xrightarrow{u_\alpha}B_\alpha
\xrightarrow{v_\alpha}C_\alpha\longrightarrow0
\]
\emph{une suite exacte de systèmes projectifs de groupes abéliens (relatifs
à un même ensemble filtrant d'indices $I$).}
\begin{enumerate}
\item[\emph{(i)}]
\emph{Si $(B_\alpha)$ vérifie (ML), il en est de même de $(C_\alpha)$.}
\item[\emph{(ii)}]
\emph{Si $(A_\alpha)$ et $(C_\alpha)$ vérifient (ML), il en est de même de
$(B_\alpha)$.}
\end{enumerate}
\end{proposition}

\begin{proof}
Soient $(f_{\alpha\beta})$, $(g_{\alpha\beta})$, $(h_{\alpha\beta})$ les
systèmes d'homomorphismes définissant les systèmes projectifs
$(A_\alpha)$, $(B_\alpha)$, $(C_\alpha)$ respectivement.

(i) Supposons que
$g_{\alpha\beta}(B_\beta)=g_{\alpha\lambda}(B_\lambda)$ pour
$\lambda\geq\beta$~; comme $v_\beta$ et $v_\lambda$ sont surjectives, on a
\[
h_{\alpha\beta}(C_\beta)
=v_\alpha(g_{\alpha\beta}(B_\beta))
=v_\alpha(g_{\alpha\lambda}(B_\lambda))
=h_{\alpha\lambda}(C_\lambda)
\]
pour $\lambda\geq\beta$.

(ii) Soit $\alpha\in I$, et soit $\beta\geq\alpha$ un indice tel que pour
$\lambda\geq\beta$, on ait
$f_{\alpha\beta}(A_\beta)=f_{\alpha\lambda}(A_\lambda)$~; soit d'autre
part $\gamma\geq\beta$ un indice tel que, pour $\lambda\geq\gamma$, on ait
$h_{\beta\gamma}(C_\gamma)=h_{\beta\lambda}(C_\lambda)$. Soit alors
$y_\alpha$ un élément de $g_{\alpha\gamma}(B_\gamma)$~; on a donc
$y_\alpha=g_{\alpha\gamma}(y_\gamma)$ avec $y_\gamma\in B_\gamma$~;
posons $y_\beta=g_{\beta\gamma}(y_\gamma)$, de sorte que
$v_\beta(y_\beta)=h_{\beta\gamma}(v_\gamma(y_\gamma))$. Pour tout
$\lambda\geq\gamma$, il existe par hypothèse $y_\lambda\in B_\lambda$ tel
que
$h_{\beta\gamma}(v_\gamma(y_\gamma))
=h_{\beta\lambda}(v_\lambda(y_\lambda))
=v_\beta(g_{\beta\lambda}(y_\lambda))$, d'où
$v_\beta(y_\beta-g_{\beta\lambda}(y_\lambda))=0$, et par suite il existe
$x_\beta\in A_\beta$
\oldpage[0\textsubscript{III}]{66}
tel que
$y_\beta=g_{\beta\lambda}(y_\lambda)+u_\beta(x_\beta)$. On en déduit
\[
y_\alpha=g_{\alpha\lambda}(y_\lambda)
+u_\alpha(f_{\alpha\beta}(x_\beta))~;
\]
mais comme $\lambda\geq\beta$, il existe $x_\lambda\in A_\lambda$ tel que
$f_{\alpha\beta}(x_\beta)=f_{\alpha\lambda}(x_\lambda)$, et finalement
\[
y_\alpha
=g_{\alpha\lambda}(y_\lambda)+u_\alpha(f_{\alpha\lambda}(x_\lambda))
=g_{\alpha\lambda}(y_\lambda+u_\lambda(x_\lambda))
\in g_{\alpha\lambda}(B_\lambda),
\]
ce qui achève la démonstration.
\end{proof}

\begin{proposition}[13.2.2]
\label{0.13.2.2-fr}
\emph{Soit $I$ un ensemble ordonné filtrant ayant une partie cofinale
dénombrable. Soit}
\[
0\longrightarrow A_\alpha\xrightarrow{u_\alpha}B_\alpha
\xrightarrow{v_\alpha}C_\alpha\longrightarrow0
\]
\emph{une suite exacte de systèmes projectifs de groupes abéliens ayant
$I$ pour ensemble d'indices. Si $(A_\alpha)$ satisfait à la condition
(ML), la suite}
\[
0\longrightarrow\varprojlim A_\alpha
\longrightarrow\varprojlim B_\alpha
\longrightarrow\varprojlim C_\alpha\longrightarrow0
\]
\emph{est exacte.}
\end{proposition}

\begin{proof}
Tout revient à prouver que l'homomorphisme
\[
v=\varprojlim v_\alpha:\varprojlim B_\alpha\longrightarrow
\varprojlim C_\alpha
\]
est surjectif. Soit $z=(z_\alpha)$ un élément de
$\varprojlim C_\alpha$, et posons
$E_\alpha=v_\alpha^{-1}(z_\alpha)$~; il est clair que les $E_\alpha$
forment un système projectif d'ensembles non vides pour les restrictions
des homomorphismes $g_{\alpha\beta}:B_\beta\to B_\alpha$. Montrons que ce
système projectif vérifie la condition (ML)~; identifiant $A_\alpha$ à une
partie de $B_\alpha$ par $u_\alpha$, pour tout $\alpha\in I$, il existe
$\beta\geq\alpha$ tel que
$g_{\alpha\beta}(A_\beta)=g_{\alpha\lambda}(A_\lambda)$ pour
$\lambda\geq\beta$~; montrons que l'on a aussi
$g_{\alpha\beta}(E_\beta)=g_{\alpha\lambda}(E_\lambda)$ pour
$\lambda\geq\beta$. En effet, prenons un $y_\lambda\in E_\lambda$ et
posons $y_\beta=g_{\beta\lambda}(y_\lambda)$,
$y_\alpha=g_{\alpha\lambda}(y_\lambda)$~; soit
$y'_\alpha\in g_{\alpha\beta}(E_\beta)$, de sorte que
$y'_\alpha=g_{\alpha\beta}(y'_\beta)$ pour un $y'_\beta\in E_\beta$~;
on a $y'_\beta-y_\beta=x_\beta\in A_\beta$, et par hypothèse il existe
$x_\lambda\in A_\lambda$ tel que
$g_{\alpha\beta}(x_\beta)=g_{\alpha\lambda}(x_\lambda)$~; donc
\[
y'_\alpha=g_{\alpha\beta}(y_\beta)+g_{\alpha\beta}(x_\beta)
=g_{\alpha\lambda}(y_\lambda)+g_{\alpha\lambda}(x_\lambda)
=g_{\alpha\lambda}(y_\lambda+x_\lambda)
\in g_{\alpha\lambda}(E_\lambda),
\]
ce qui démontre notre assertion. Cela étant, on sait (Bourbaki,
\emph{Top. gén.}, chap.~II, 3\textsuperscript{e} éd., §~3, th.~1) que sous
les hypothèses faites sur $I$, un système projectif d'ensembles non vides
vérifiant (ML) a une limite projective non vide~; par suite, il existe un
point $y=(y_\alpha)\in\varprojlim E_\alpha$, et comme
$v_\alpha(y_\alpha)=z_\alpha$ par définition pour tout $\alpha$, on a
$z=v(y)$, C.Q.F.D.
\end{proof}

\begin{proposition}[13.2.3]
\label{0.13.2.3-fr}
\emph{Les hypothèses sur $I$ étant celles de
(\hyperref[0.13.2.2-fr]{13.2.2}), soit
$(K_\alpha^\bullet)_{\alpha\in I}$ un système projectif de complexes de
groupes abéliens
$K_\alpha^\bullet=(K_\alpha^n)_{n\in\mathbb Z}$ dont l'opérateur de
dérivation est de degré $+1$. Pour chaque $n$, il existe un homomorphisme
canonique fonctoriel}
\begin{equation}
\label{0.13.2.3.1-fr}
\tag{13.2.3.1}
h_n:H^n(\varprojlim K_\alpha^\bullet)
\longrightarrow\varprojlim H^n(K_\alpha^\bullet).
\end{equation}
\emph{Si, pour tout degré $n$, le système projectif de groupes abéliens
$(K_\alpha^n)_{\alpha\in I}$ vérifie (ML), alors tous les homomorphismes
$h_n$ sont surjectifs. Si en outre, pour un degré $n$, le système
projectif $(H^{n-1}(K_\alpha^\bullet))_{\alpha\in I}$ vérifie (ML),
l'homomorphisme $h_n$ est bijectif.}
\end{proposition}

\begin{proof}
Posons, pour tout $n$, $K^n=\varprojlim K_\alpha^n$~; la définition des
homomorphismes $h_n$ provient de la commutativité des diagrammes
\[
\xymatrix{
\cdots\ar[r]&K^{n-1}\ar[r]\ar[d]&K^n\ar[r]\ar[d]
&K^{n+1}\ar[r]\ar[d]&\cdots\\
\cdots\ar[r]&K_\alpha^{n-1}\ar[r]&K_\alpha^n\ar[r]
&K_\alpha^{n+1}\ar[r]&\cdots
}
\]
\oldpage[0\textsubscript{III}]{67}
les opérateurs de dérivation dans $K^\bullet$ étant limites projectives
des opérateurs correspondants dans les $K_\alpha^\bullet$.

Considérons les suites exactes
\[
(*_n)\qquad
0\longrightarrow \mathrm B^n(K_\alpha^\bullet)
\longrightarrow \mathrm Z^n(K_\alpha^\bullet)
\longrightarrow H^n(K_\alpha^\bullet)\longrightarrow0
\]
\[
(**_n)\qquad
0\longrightarrow \mathrm Z^{n-1}(K_\alpha^\bullet)
\longrightarrow K_\alpha^{n-1}
\longrightarrow \mathrm B^n(K_\alpha^\bullet)\longrightarrow0.
\]
L'hypothèse et la prop.
(\hyperref[0.13.2.1-fr]{13.2.1}, (i)) montrent que le système projectif
$(\mathrm B^n(K_\alpha^\bullet))_{\alpha\in I}$ vérifie (ML) pour tout
$n$~; il résulte donc de (\hyperref[0.13.2.2-fr]{13.2.2}) que la suite
\[
(***_n)\qquad
0\longrightarrow\varprojlim_\alpha\mathrm B^n(K_\alpha^\bullet)
\longrightarrow\varprojlim_\alpha\mathrm Z^n(K_\alpha^\bullet)
\longrightarrow\varprojlim_\alpha H^n(K_\alpha^\bullet)
\longrightarrow0
\]
est exacte. Or il est clair que
$\varprojlim_\alpha\mathrm B^n(K_\alpha^\bullet)$ s'identifie à un
sous-groupe de $K^{n+1}$ contenant $\mathrm B^n(K^\bullet)$ et que
$\varprojlim_\alpha\mathrm Z^n(K_\alpha^\bullet)$ s'identifie à un
sous-groupe de $\mathrm Z^n(K^\bullet)$~; par suite, $h_n$ est surjectif.
Si maintenant, on suppose en outre que le système projectif
$(H^{n-1}(K_\alpha^\bullet))_{\alpha\in I}$ vérifie (ML), les suites
exactes $(*_{n-1})$ et la prop.
(\hyperref[0.13.2.1-fr]{13.2.1}, (ii)) montrent que le système projectif
$(\mathrm Z^{n-1}(K_\alpha^\bullet))_{\alpha\in I}$ vérifie (ML)~; mais
alors, (\hyperref[0.13.2.2-fr]{13.2.2}) appliquée aux suites exactes
$(**_n)$ montre que la suite
\[
0\longrightarrow
\varprojlim_\alpha\mathrm Z^{n-1}(K_\alpha^\bullet)
\longrightarrow K^{n-1}\xrightarrow{u}
\varprojlim_\alpha\mathrm B^n(K_\alpha^\bullet)
\longrightarrow0
\]
est exacte~; comme
$\varprojlim_\alpha\mathrm B^n(K_\alpha^\bullet)
\supset\mathrm B^n(K^\bullet)$, et que la composée de l'injection
$\varprojlim_\alpha\mathrm B^n(K_\alpha^\bullet)\to K^n$ et de $u$ est
l'opérateur de dérivation $K^{n-1}\to K^n$, le fait que $u$ soit surjectif
entraîne
$\varprojlim_\alpha\mathrm B^n(K_\alpha^\bullet)
=\mathrm B^n(K^\bullet)$, donc $h_n$ est injectif. C.Q.F.D.
\end{proof}

\begin{env}[13.2.4]
\label{0.13.2.4-fr}
\emph{Remarques} (13.2.4).
\begin{enumerate}
\item[\emph{(i)}]
Le raisonnement de (\hyperref[0.13.2.2-fr]{13.2.2}) (cf. Bourbaki,
\emph{loc. cit.}) montre que la conclusion de cette proposition reste
valable lorsqu'on suppose seulement que les $A_\alpha$ peuvent être munis
de structures d'espaces \emph{métrisables complets}, dans lesquels les
translations sont des homéomorphismes, que les applications
$f_{\alpha\beta}:A_\beta\to A_\alpha$ définissant le système projectif
$(A_\alpha)$ sont \emph{uniformément continues} pour les distances
considérées, et enfin que le système $(A_\alpha)$ vérifie la condition
\begin{itemize}
\item[\textbf{(ML')}]
\emph{Pour tout indice $\alpha$, il existe $\beta\geq\alpha$ tel que pour
tout $\gamma\geq\beta$, $f_{\alpha\gamma}(A_\gamma)$ est dense dans
$f_{\alpha\beta}(A_\beta)$.}
\end{itemize}

Cela permet d'apporter un complément analogue à
(\hyperref[0.13.2.3-fr]{13.2.3})~: supposons que
$K_\alpha^n=0$ pour $n<0$ et pour tout $\alpha$~; supposons de plus que
$(K_\alpha^n)_{\alpha\in I}$ vérifie (ML) pour $n\geq0$ et que les
$A_\alpha=H^0(K_\alpha^\bullet)$ puissent être munis de structures
d'espaces métriques vérifiant les propriétés ci-dessus. Alors les
conclusions de (\hyperref[0.13.2.3-fr]{13.2.3}) sont inchangées pour
$n\geq2$, et en outre $h_1$ est \emph{bijectif}, car le raisonnement de
(\hyperref[0.13.2.2-fr]{13.2.2}) montre encore que
$(\mathrm B^1(K_\alpha^\bullet))_{\alpha\in I}$ vérifie (ML), que la suite
$(***_1)$ est exacte, et enfin, en vertu de ce qui précède, que
$\varprojlim\mathrm B^1(K_\alpha^\bullet)=\mathrm B^1(K^\bullet)$. On a
ainsi établi entre autres les assertions de (T, 3.10.2).

\item[\emph{(ii)}]
Il est possible d'introduire les foncteurs dérivés à droite
$\varinjlim^{(n)}$ du foncteur $\varinjlim$, et d'obtenir des énoncés
plus complets que les précédents [28].
\end{enumerate}
\end{env}

\subsection{Application : cohomologie d'une limite projective de faisceaux}
\label{subsection:0.13.3-fr}

\begin{proposition}[13.3.1]
\label{0.13.3.1-fr}
\oldpage[0\textsubscript{III}]{68}
\emph{Soient $X$ un espace topologique,
$(\mathcal F_k)_{k\in\mathbb N}$ un système projectif de faisceaux de
groupes abéliens sur $X$, et soit
$\mathcal F=\varprojlim_k\mathcal F_k$. On suppose vérifiées les conditions
suivantes :}
\begin{enumerate}
\item[\emph{(i)}]
\emph{Il existe une base $\mathfrak B$ de la topologie de $X$ telle que,
pour tout $U\in\mathfrak B$ et tout $i\geq0$, le système projectif
$(H^i(U,\mathcal F_k))_{k\in\mathbb N}$ vérifie (ML).}
\item[\emph{(ii)}]
\emph{Pour tout $x\in X$ et tout $i>0$, on a}
\[
\varinjlim_U\bigl(\varprojlim_k H^i(U,\mathcal F_k)\bigr)=0
\]
\emph{lorsque $U$ parcourt l'ensemble des voisinages de $x$ appartenant à
$\mathfrak B$.}
\item[\emph{(iii)}]
\emph{Les homomorphismes
$u_{hk}:\mathcal F_k\to\mathcal F_h$ $(h\leq k)$ définissant le système
projectif $(\mathcal F_k)$ sont surjectifs.}
\end{enumerate}
\emph{Dans ces conditions, pour tout $i>0$, l'homomorphisme canonique}
\[
h_i:H^i(X,\mathcal F)\longrightarrow
\varprojlim_k H^i(X,\mathcal F_k)
\]
\emph{est surjectif~; si en outre, pour une valeur de $i$, le système
projectif $(H^{i-1}(X,\mathcal F_k))_{k\in\mathbb N}$ vérifie (ML), $h_i$
est bijectif.}
\end{proposition}

\begin{proof}
\emph{a)} Nous allons d'abord supposer que les $\mathcal F_k$ sont flasques
ainsi que les noyaux $\mathcal N_{hk}$ des $u_{hk}$~; nous allons montrer
alors que la condition (iii) de l'énoncé entraîne que
$H^i(X,\mathcal F)=0$ pour $i>0$. Il suffira de prouver que pour tout ouvert
$U$ de $X$ et tout recouvrement $\mathfrak U$ de $U$ par des ouverts de
$U$, on a $\check H^i(\mathfrak U,\mathcal F)=0$ pour $i>0$. Il en résultera
en effet tout d'abord que pour la cohomologie de Čech, on a
$\check H^i(U,\mathcal F)=0$ pour tout $i>0$, puis (en vertu de
(G, II, 5.9.2) appliqué à l'ensemble de tous les ouverts de $X$) que
$H^i(X,\mathcal F)=0$ pour tout $i>0$. Comme les $\mathcal F_k$ sont
flasques, on a $\check H^i(\mathfrak U,\mathcal F_k)=0$ pour $i>0$
(G, II, 5.2.3)~; considérons pour chaque $k$ le complexe
$\check C^\bullet(\mathfrak U,\mathcal F_k)$ des cochaînes alternées du nerf
du recouvrement $\mathfrak U$ (G, II, 5.1), qui forment évidemment un
système projectif de complexes de groupes abéliens. Montrons que toutes les
applications
$\check C^\bullet(\mathfrak U,\mathcal F_k)\to
\check C^\bullet(\mathfrak U,\mathcal F_h)$ $(h\leq k)$ sont surjectives.
Il suffit évidemment, par définition, de montrer que pour tout ouvert $V$
de $X$, l'application
$\Gamma(V,\mathcal F_k)\to\Gamma(V,\mathcal F_h)$ est surjective~; mais la
suite
\[
0\longrightarrow\mathcal N_{hk}\longrightarrow\mathcal F_k
\longrightarrow\mathcal F_h\longrightarrow0
\]
étant exacte par hypothèse, donne la suite exacte de cohomologie
\[
\Gamma(V,\mathcal F_k)\longrightarrow\Gamma(V,\mathcal F_h)
\longrightarrow H^1(V,\mathcal N_{hk})=0
\]
puisque $\mathcal N_{hk}$ est flasque. Le système projectif
$(\check C^\bullet(\mathfrak U,\mathcal F_k))_{k\in\mathbb N}$ vérifie donc
(ML)~; il en est de même de
$(\check H^i(\mathfrak U,\mathcal F_k))_{k\in\mathbb N}$ pour tout
$i\geq0$, car cela est trivial pour $i>0$, et comme
$\check H^0(\mathfrak U,\mathcal F_k)=\Gamma(U,\mathcal F_k)$
(G, II, 5.2.2), la condition (ML) est aussi remplie pour $i=0$ d'après ce
qui précède. On peut donc appliquer
(\hyperref[0.13.2.3-fr]{13.2.3}), qui montre que
\[
\check H^i(\mathfrak U,\mathcal F)
=\varprojlim_k\check H^i(\mathfrak U,\mathcal F_k)=0
\quad\hbox{pour tout }i>0.
\]

\emph{b)} Passons au cas général, et considérons pour chaque $k\in\mathbb N$,
la résolution canonique
$\mathcal C^\bullet(X,\mathcal F_k)
=(\mathcal C^i(X,\mathcal F_k))_{i\geq0}$ de $\mathcal F_k$ par des
faisceaux flasques (G, II, 4.3). Pour chaque $i\geq0$, il est clair que
$(\mathcal C^i(X,\mathcal F_k))_{k\in\mathbb N}$ est un système projectif
de faisceaux flasques~; montrons qu'il satisfait aux conditions de
\emph{a)}.
\oldpage[0\textsubscript{III}]{69}
En effet, si $\mathcal N_{hk}$ est le noyau de $u_{hk}$ pour $h\leq k$, la
suite
\[
0\longrightarrow\mathcal N_{hk}\longrightarrow\mathcal F_k
\longrightarrow\mathcal F_h\longrightarrow0
\]
est exacte d'après (iii), et notre assertion résulte de ce que le foncteur
$\mathcal A\mapsto\mathcal C^i(X,\mathcal A)$ est exact (G, II, 4.3).
Soit
$\mathcal G^i=\varprojlim_k\mathcal C^i(X,\mathcal F_k)$~; on a donc
$H^j(X,\mathcal G^i)=0$ pour $j>0$ et $i\geq0$ en vertu de \emph{a)}. Nous
allons montrer que $\mathcal G^\bullet=(\mathcal G^i)_{i\geq0}$ est une
résolution du faisceau $\mathcal F$~; comme
$H^j(X,\mathcal G^i)=0$ pour $j>0$, la cohomologie $H^\bullet(X,\mathcal F)$
sera égale à $H^\bullet(\Gamma(X,\mathcal G^\bullet))$ (G, II, 4.7.1).

Il est clair que, par passage à la limite projective, on déduit des suites
exactes
\[
0\longrightarrow\mathcal F_k
\longrightarrow\mathcal C^0(X,\mathcal F_k)
\longrightarrow\mathcal C^1(X,\mathcal F_k)\longrightarrow\cdots
\]
un complexe de faisceaux de groupes abéliens
\[
0\longrightarrow\mathcal F\longrightarrow\mathcal G^0
\longrightarrow\mathcal G^1\longrightarrow\cdots.
\]
Pour prouver notre assertion, il faut établir que
$\mathcal H^i(\mathcal G^\bullet)=0$ pour $i>0$. Ce faisceau est engendré
par le préfaisceau
$U\mapsto H^i(\Gamma(U,\mathcal G^\bullet))$ (G, II, 4.1)~; or, le complexe
$\Gamma(U,\mathcal G^\bullet)$ est la limite projective du système
projectif de complexes de groupes abéliens
\[
(\Gamma(U,\mathcal C^\bullet(X,\mathcal F_k)))_{k\in\mathbb N}
\]
(0\textsubscript{I}, 3.2.6). On a vu dans \emph{a)} que pour chaque
$i\geq0$, les applications
$\Gamma(U,\mathcal C^i(X,\mathcal F_k))\to
\Gamma(U,\mathcal C^i(X,\mathcal F_h))$ $(h\leq k)$ sont surjectives~;
d'autre part, on a
$H^i(U,\mathcal F_k)=
H^i(\Gamma(U,\mathcal C^\bullet(X,\mathcal F_k)))$, la résolution canonique
$\mathcal C^\bullet(U,\mathcal F_k|U)$ étant induite sur $U$ par
$\mathcal C^\bullet(X,\mathcal F_k)$~; en vertu de l'hypothèse (i), pour
tout $U\in\mathfrak B$, on peut appliquer
(\hyperref[0.13.2.3-fr]{13.2.3}) au système projectif de complexes
$(\Gamma(U,\mathcal C^\bullet(X,\mathcal F_k)))_{k\in\mathbb N}$, et on a
donc
\[
H^i(\Gamma(U,\mathcal G^\bullet))
=\varprojlim_k H^i(U,\mathcal F_k)
\quad\hbox{pour tout }i\geq0.
\]
L'hypothèse (ii) prouve bien alors, par définition, que les faisceaux
$\mathcal H^i(\mathcal G^\bullet)$ sont nuls pour $i>0$.

On a alors pour tout $i\geq0$,
\[
H^i(X,\mathcal F)=H^i(\Gamma(X,\mathcal G^\bullet))
\quad\hbox{et}\quad
\Gamma(X,\mathcal G^\bullet)
=\varprojlim_k\Gamma(X,\mathcal C^\bullet(X,\mathcal F_k)).
\]
On vient de remarquer que les applications
$\Gamma(X,\mathcal C^i(X,\mathcal F_k))\to
\Gamma(X,\mathcal C^i(X,\mathcal F_h))$ $(h\leq k)$ sont toutes
surjectives~; la conclusion résulte donc encore de
(\hyperref[0.13.2.3-fr]{13.2.3}).
\end{proof}

\begin{env}[13.3.2]
\label{0.13.3.2-fr}
\emph{Remarques} (13.3.2).
\begin{enumerate}
\item[\emph{(i)}]
L'énoncé (\hyperref[0.13.3.1-fr]{13.3.1}) n'a d'intérêt que pour $i>0$,
puisque pour $i=0$, $h_i$ est toujours un isomorphisme sans hypothèse
(0\textsubscript{I}, 3.2.6).
\item[\emph{(ii)}]
Les conditions (i) et (ii) de (\hyperref[0.13.3.1-fr]{13.3.1}) seront en
particulier vérifiées si $H^i(U,\mathcal F_k)=0$ pour tout $k$, tout $i>0$
et tout $U\in\mathfrak B$, et si pour $U\in\mathfrak B$, les applications
$\Gamma(U,\mathcal F_k)\to\Gamma(U,\mathcal F_h)$ sont surjectives. Ce sera
le cas le plus fréquent d'application de
(\hyperref[0.13.3.1-fr]{13.3.1}).
\end{enumerate}
\end{env}

\subsection{Condition de Mittag-Leffler et objets gradués associés aux systèmes projectifs}
\label{subsection:0.13.4-fr}
\label{0.13.4-fr}

\begin{env}[13.4.1]
\label{0.13.4.1-fr}
Soit $\mathbf A=(A_k,u_{kh})_{k\in\mathbb Z}$ un système projectif dans une
catégorie abélienne $C$~; nous dirons qu'il est \emph{limité inférieurement}
s'il existe $k_0$ tel que $A_k=0$ pour $k<k_0$. Nous définirons sur chaque
$A_k$ une \emph{filtration} $(\mathrm F^p(A_k))_{p\in\mathbb Z}$ par les
formules
\begin{equation}
\label{0.13.4.1.1-fr}
\tag{13.4.1.1}
\mathrm F^p(A_k)=
\begin{cases}
\operatorname{Ker}(A_k\longrightarrow A_{p-1})&\text{pour }p\leq k+1,\\
0&\text{pour }p\geq k+1.
\end{cases}
\end{equation}
\oldpage[0\textsubscript{III}]{70}
On a donc par hypothèse
$\mathrm F^{k_0}(A_k)=A_k$ et $\mathrm F^{k+1}(A_k)=0$, autrement dit la
filtration considérée est \emph{finie}
(\hyperref[0.11.1.3-fr]{11.1.3}). Les objets gradués associés à cette
filtration sont donc
\[
\operatorname{gr}^p(A_k)=
\operatorname{Ker}(A_k\longrightarrow A_{p-1})
/\operatorname{Ker}(A_k\longrightarrow A_p)
\]
et par suite, $\operatorname{gr}^p(A_k)$ est isomorphe à l'image par
$A_k\to A_p$ de $\operatorname{Ker}(A_k\to A_{p-1})$~; en vertu de la
transitivité des morphismes définissant un système projectif, on a donc
\begin{equation}
\label{0.13.4.1.2-fr}
\tag{13.4.1.2}
\operatorname{gr}^p(A_k)=
\operatorname{Ker}(A_p\longrightarrow A_{p-1})
\cap\operatorname{Im}(A_k\longrightarrow A_p)
\end{equation}
mais comme, en vertu de
(\hyperref[0.13.4.1.1-fr]{13.4.1.1}), on a
$\operatorname{Ker}(A_p\to A_{p-1})=\operatorname{gr}^p(A_p)$, on a aussi
\begin{equation}
\label{0.13.4.1.3-fr}
\tag{13.4.1.3}
\operatorname{gr}^p(A_k)=
\operatorname{gr}^p(A_p)\cap\operatorname{Im}(A_k\longrightarrow A_p).
\end{equation}
Les définitions précédentes montrent en outre que l'on a pour $k\leq h$
\[
u_{kh}(\mathrm F^p(A_h))\subset\mathrm F^p(A_k)
\]
et par suite que les $\operatorname{gr}^p(u_{kh})$ définissent un
\emph{système projectif}
$(\operatorname{gr}^p(A_k))_{k\in\mathbb Z}$ pour tout $p\in\mathbb Z$.
\end{env}

\begin{env}[13.4.2]
\label{0.13.4.2-fr}
Nous dirons que le système projectif $\mathbf A$ est
\emph{essentiellement constant} si les morphismes
$A_{k+1}\to A_k$ sont des \emph{isomorphismes} pour $k$ assez grand. Nous
dirons que le système projectif $\mathbf A$ est \emph{strict} si les
morphismes $A_i\to A_j$ $(j\leq i)$ sont des \emph{épimorphismes}.
Lorsque $\mathbf A$ est strict, il résulte de
(\hyperref[0.13.4.1.3-fr]{13.4.1.3}) que pour $p\leq k\leq h$, le morphisme
canonique $\operatorname{gr}^p(A_h)\to\operatorname{gr}^p(A_k)$ est un
\emph{isomorphisme}, autrement dit, le système projectif
$(\operatorname{gr}^p(A_k))_{k\in\mathbb Z}$ est essentiellement constant.
La suite des objets $\operatorname{gr}^p(A_p)$ (identifiés à
$\varprojlim_k\operatorname{gr}^p(A_k)$ pour tout $p$) se note alors
$\operatorname{gr}^\bullet(\mathbf A)$ et s'appelle \emph{l'objet gradué
associé au système projectif strict} $\mathbf A=(A_k)$.

Si l'on suppose maintenant que le système projectif $\mathbf A$ (limité
inférieurement) vérifie (ML), on sait
(\hyperref[0.13.1.2-fr]{13.1.2}) que le système projectif
$\mathbf A'=(A'_k)$ des objets d'«~images universelles~» est \emph{strict},
et il est par ailleurs limité inférieurement~; l'objet gradué
$\operatorname{gr}^\bullet(\mathbf A')$ associé à $\mathbf A'$ est alors
encore appelé l'objet gradué \emph{associé à} $\mathbf A$ et noté
$\operatorname{gr}^\bullet(\mathbf A)$.
\end{env}

\begin{proposition}[13.4.3]
\label{0.13.4.3-fr}
\emph{Soit $\mathbf A=(A_k)_{k\in\mathbb Z}$ un système projectif limité
inférieurement dans une catégorie abélienne $C$. Les deux conditions
suivantes sont équivalentes :}
\begin{enumerate}
\item[\emph{a)}] \emph{$\mathbf A$ vérifie la condition (ML).}
\item[\emph{b)}] \emph{Pour tout $p\in\mathbb Z$, le système projectif
$(\operatorname{gr}^p(A_k))_{k\in\mathbb Z}$ est essentiellement constant.}
\end{enumerate}
\emph{En outre, lorsque ces conditions sont satisfaites, on a pour tout
$p\in\mathbb Z$ un isomorphisme canonique}
\begin{equation}
\label{0.13.4.3.1-fr}
\tag{13.4.3.1}
\operatorname{gr}^p(\mathbf A)\xrightarrow{\sim}
\varprojlim_k\operatorname{gr}^p(A_k).
\end{equation}
\end{proposition}

\begin{proof}
Il résulte aussitôt de
(\hyperref[0.13.4.1.2-fr]{13.4.1.2}) que \emph{a)} implique \emph{b)}~; la
même formule appliquée au système projectif $\mathbf A'$ (notations de
(\hyperref[0.13.4.2-fr]{13.4.2})) donne l'isomorphisme
(\hyperref[0.13.4.3.1-fr]{13.4.3.1}) par définition. Pour $k\leq h$,
posons $A_{kh}=\operatorname{Im}(A_h\to A_k)$~; si $k\leq h\leq j$, on a
$A_{kj}\subset A_{kh}\subset A_k$. Munissons $A_{kh}$ de la filtration
induite par $(\mathrm F^p(A_k))$~; on vérifie aussitôt, en vertu de la
transitivité des morphismes définissant $\mathbf A$, que cette filtration
est aussi la filtration quotient de $(\mathrm F^p(A_h))$~; par suite, on a
\begin{equation}
\label{0.13.4.3.2-fr}
\tag{13.4.3.2}
\operatorname{gr}^p(A_{kh})=
\operatorname{Im}\bigl(\operatorname{gr}^p(A_h)\longrightarrow
\operatorname{gr}^p(A_k)\bigr).
\end{equation}
\oldpage[0\textsubscript{III}]{71}
Cela étant, supposons \emph{b)} vérifiée~; pour tout $p\in\mathbb Z$, et
tout $k\geq p$, il existe un entier $\mathrm L(p,k)$ tel que le second
membre de (\hyperref[0.13.4.3.2-fr]{13.4.3.2}) soit constant pour
$h\geq\mathrm L(p,k)$~; comme $\operatorname{gr}^p(A_k)=0$ pour $p<k_0$,
il n'y a (pour $k$ donné) qu'un nombre fini de $\mathrm L(p,k)$ non nuls
lorsque $p$ parcourt l'ensemble des entiers $\leq k$. Soit
$\mathrm L(k)=m$ le plus grand de ces entiers~; pour tout $h\geq m$, on a
$A_{kh}\subset A_{km}$, et par définition de $m$, l'injection canonique
$A_{kh}\to A_{km}$ définit un \emph{isomorphisme}
$\operatorname{gr}^\bullet(A_{kh})\xrightarrow{\sim}
\operatorname{gr}^\bullet(A_{km})$~; comme les filtrations sont
\emph{finies}, on en conclut que l'injection précédente est elle-même
bijective (Bourbaki, \emph{Alg. comm.}, chap.~III, §~2,
n\textsuperscript{o}~8, th.~1), ce qui prouve que $\mathbf A$ vérifie
(ML).
\end{proof}

\begin{env}[13.4.4]
\label{0.13.4.4-fr}
Supposons que dans $C$ la limite projective
$A=\varprojlim_k A_k$ existe. Dans les définitions de
(\hyperref[0.13.4.1-fr]{13.4.1}), on peut alors remplacer $A_k$ par $A$,
et la filtration ainsi définie sur $A$ est encore telle que
\begin{equation}
\label{0.13.4.4.1-fr}
\tag{13.4.4.1}
\operatorname{gr}^p(A)=\operatorname{gr}^p(A_p)
\cap\operatorname{Im}(A\longrightarrow A_p).
\end{equation}
\end{env}

\begin{corollary}[13.4.5]
\label{0.13.4.5-fr}
\emph{Supposons que $C$ soit la catégorie des groupes abéliens. Si le
système projectif $\mathbf A$ vérifie (ML) et si
$A=\varprojlim_k A_k$, on a pour tout $p\in\mathbb Z$ un isomorphisme
canonique}
\begin{equation}
\label{0.13.4.5.1-fr}
\tag{13.4.5.1}
\operatorname{gr}^p(A)\xrightarrow{\sim}
\varprojlim_k\operatorname{gr}^p(A_k).
\end{equation}
\end{corollary}
\begin{proof}
En effet, on a
$\operatorname{Im}(A_k\to A_p)=\operatorname{Im}(A\to A_p)$ dès que $k$
est assez grand (Bourbaki, \emph{Top. gén.}, chap.~II,
3\textsuperscript{e} éd., §~3, n\textsuperscript{o}~5, th.~1), et la
conclusion résulte de
(\hyperref[0.13.4.1.3-fr]{13.4.1.3}) et
(\hyperref[0.13.4.4.1-fr]{13.4.4.1}).
\end{proof}

\subsection{Limites projectives de suites spectrales de complexes filtrés}
\label{subsection:0.13.5-fr}
\begin{env}[13.5.1]
\label{0.13.5.1-fr}
Soient $C$ une catégorie abélienne, $X^\bullet$ un complexe d'objets de
$C$ muni d'une \emph{filtration}
$(\mathrm F^p(X^\bullet))_{p\in\mathbb Z}$ telle que
$\mathrm F^{p_0}(X^\bullet)=X^\bullet$ pour un indice $p_0$.
Considérons pour chaque $k\in\mathbb Z$ le complexe
$X_k^\bullet=X^\bullet/\mathrm F^{k+1}(X^\bullet)$~; il est canoniquement
muni de la filtration formée des
$\mathrm F^p(X_k^\bullet)=
\mathrm F^p(X^\bullet)/\mathrm F^{k+1}(X^\bullet)$ pour $p\leq k$ et des
$\mathrm F^p(X_k^\bullet)=0$ pour $p\geq k+1$. En outre, on a des
morphismes canoniques $X_{k+1}^\bullet\to X_k^\bullet$, qui font de
$\mathbf X^\bullet=(X_k^\bullet)_{k\in\mathbb Z}$ un \emph{système
projectif de complexes filtrés d'objets de $C$}. On notera que ce système
projectif est \emph{strict} et tel que $X_k^\bullet=0$ pour $k<p_0$.
\end{env}
\begin{env}[13.5.2]
\label{0.13.5.2-fr}
Considérons plus généralement un système projectif \emph{strict}
$\mathbf X^\bullet=(X_k^\bullet)_{k\in\mathbb Z}$ de complexes d'objets de
$C$, \emph{limité inférieurement}~; considérons sur chaque
$X_k^\bullet$ la filtration définie dans
(\hyperref[0.13.4.1-fr]{13.4.1}) (en se plaçant dans la catégorie abélienne
des complexes de $C$ limités inférieurement). Les
$X_k^\bullet\to X_p^\bullet$ $(p\leq k)$ deviennent des morphismes de
complexes \emph{filtrés}, à filtrations finies. Le caractère fonctoriel des
suites spectrales des complexes filtrés
(\hyperref[0.11.2.3-fr]{11.2.3}) montre que les morphismes de définition du
système projectif $\mathbf X^\bullet$ fournissent des morphismes faisant de
$\mathrm E(\mathbf X^\bullet)=(\mathrm E(X_k^\bullet))$ un \emph{système
projectif de suites spectrales}.
\end{env}
\begin{lemma}[13.5.3]
\label{0.13.5.3-fr}
\emph{Supposons que le système projectif
$\mathbf X^\bullet=(X_k^\bullet)_{k\in\mathbb Z}$ de complexes filtrés soit
obtenu comme dans (\hyperref[0.13.5.2-fr]{13.5.2}). Alors :}
\begin{enumerate}
\item[\emph{a)}]\emph{Pour $r\geq p-p_0$, on a
$\mathrm B_r^{pq}(X_k^\bullet)=\mathrm B_\infty^{pq}(X_k^\bullet)$ pour
tout $k\in\mathbb Z$.}
\item[\emph{b)}]\emph{Pour $k+1\leq p+r$, on a
$\mathrm Z_r^{pq}(X_k^\bullet)=\mathrm Z_\infty^{pq}(X_k^\bullet)$.}
\item[\emph{c)}]\emph{Pour $k+1\geq p+r$, les morphismes
$\mathrm Z_r^{pq}(X_h^\bullet)\to\mathrm Z_r^{pq}(X_k^\bullet)$ et
$\mathrm B_r^{pq}(X_h^\bullet)\to\mathrm B_r^{pq}(X_k^\bullet)$ sont des
isomorphismes pour tout $h\geq k$.}
\end{enumerate}
\end{lemma}
\oldpage[0\textsubscript{III}]{72}
Ces trois propriétés résultent aussitôt des définitions de
(\hyperref[0.11.2.2-fr]{11.2.2}), en tenant compte de ce que
$\mathrm F^{p-r+1}(X_k^\bullet)=X_k^\bullet$ pour $p-r<p_0$.

\begin{env}[13.5.4]
\label{0.13.5.4-fr}
Supposons vérifiées les hypothèses de
(\hyperref[0.13.5.3-fr]{13.5.3}). Alors, pour $p,q,r$ fixés ($r$ fini), les
systèmes projectifs
\[
(\mathrm Z_r^{pq}(X_k^\bullet))_{k\in\mathbb Z},\quad
(\mathrm B_r^{pq}(X_k^\bullet))_{k\in\mathbb Z},\quad
(\mathrm E_r^{pq}(X_k^\bullet))_{k\in\mathbb Z}
\]
sont \emph{essentiellement constants}~; on désignera par
$\mathrm Z_r^{pq}(\mathbf X^\bullet)$,
$\mathrm B_r^{pq}(\mathbf X^\bullet)$ et
$\mathrm E_r^{pq}(\mathbf X^\bullet)=
\mathrm Z_r^{pq}(\mathbf X^\bullet)/
\mathrm B_r^{pq}(\mathbf X^\bullet)$ leurs limites projectives
respectives. Les $\mathrm Z_r^{pq}(\mathbf X^\bullet)$ et
$\mathrm B_r^{pq}(\mathbf X^\bullet)$ s'identifient canoniquement à des
sous-objets de $\mathrm E_1^{pq}(\mathbf X^\bullet)$. La définition des
$d_r^{pq}$ (M, XV, 1) montre que ces morphismes (relatifs aux
$X_k^\bullet$) sont aussi essentiellement constants, et par suite
définissent des morphismes
\begin{equation}
\label{0.13.5.4.1-fr}
\tag{13.5.4.1}
d_r^{pq}:\mathrm E_r^{pq}(\mathbf X^\bullet)
\longrightarrow\mathrm E_r^{p+r,q-r+1}(\mathbf X^\bullet)
\end{equation}
tels que $d_r^{p+r,q-r+1}\circ d_r^{pq}=0$~; en outre, on a des
isomorphismes canoniques de $\operatorname{Ker}(d_r^{pq})$ sur
$\mathrm Z_{r+1}^{pq}(\mathbf X^\bullet)/
\mathrm B_r^{pq}(\mathbf X^\bullet)$ et de
$\operatorname{Im}(d_r^{pq})$ sur
$\mathrm B_{r+1}^{p+r,q-r+1}(\mathbf X^\bullet)/
\mathrm B_r^{p+r,q-r+1}(\mathbf X^\bullet)$.
\end{env}

\begin{lemma}[13.5.5]
\label{0.13.5.5-fr}
\emph{Sous les hypothèses de
(\hyperref[0.13.5.3-fr]{13.5.3}), on a, pour $s\geq r>p-p_0$, un
monomorphisme canonique}
\begin{equation}
\label{0.13.5.5.1-fr}
\tag{13.5.5.1}
i:\mathrm E_s^{pq}(\mathbf X^\bullet)
\longrightarrow\mathrm E_r^{pq}(\mathbf X^\bullet)
\end{equation}
\emph{et un isomorphisme canonique}
\begin{equation}
\label{0.13.5.5.2-fr}
\tag{13.5.5.2}
j_r:\mathrm E_r^{pq}(\mathbf X^\bullet)
\xrightarrow{\sim}\mathrm E_\infty^{pq}(X_{p+r-1}^\bullet)
\end{equation}
\emph{tels que le diagramme}
\begin{equation}
\label{0.13.5.5.3-fr}
\tag{13.5.5.3}
\xymatrix{
\mathrm E_s^{pq}(\mathbf X^\bullet)\ar[r]^{j_s}\ar[d]_i
&\mathrm E_\infty^{pq}(X_{p+s-1}^\bullet)\ar[d]\\
\mathrm E_r^{pq}(\mathbf X^\bullet)\ar[r]^{j_r}
&\mathrm E_\infty^{pq}(X_{p+r-1}^\bullet)
}
\end{equation}
\emph{soit commutatif} (la flèche verticale de droite provenant du
morphisme $X_{p+s-1}^\bullet\to X_{p+r-1}^\bullet$).
\end{lemma}
\begin{proof}
L'existence de $i$ provient de ce que
$\mathrm B_r^{pq}(X_k^\bullet)=\mathrm B_\infty^{pq}(X_k^\bullet)$ pour
$r>p-p_0$ (\hyperref[0.13.5.3-fr]{13.5.3}, a)~; on a
$\mathrm Z_r^{pq}(X_k^\bullet)=\mathrm Z_\infty^{pq}(X_k^\bullet)$ pour
$k+1\leq p+r$ (\hyperref[0.13.5.3-fr]{13.5.3}, b), d'où en particulier
$\mathrm Z_r^{pq}(X_{p+r-1}^\bullet)=
\mathrm Z_\infty^{pq}(X_{p+r-1}^\bullet)$ et d'autre part
$\mathrm Z_r^{pq}(X_{p+r-1}^\bullet)$ et
$\mathrm B_r^{pq}(X_{p+r-1}^\bullet)$ s'identifient canoniquement à
$\mathrm Z_r^{pq}(\mathbf X^\bullet)$ et
$\mathrm B_r^{pq}(\mathbf X^\bullet)$ en vertu de
(\hyperref[0.13.5.3-fr]{13.5.3}, c)), d'où l'existence de $j_r$ et la
commutativité de
(\hyperref[0.13.5.5.1-fr]{13.5.5.1}).
\end{proof}

\begin{corollary}[13.5.6]
\label{0.13.5.6-fr}
\emph{Sous les hypothèses de
(\hyperref[0.13.5.3-fr]{13.5.3}), si l'une des limites projectives
$\varprojlim_r\mathrm E_r^{pq}(\mathbf X^\bullet)$,
$\varprojlim_k\mathrm E_\infty^{pq}(X_k^\bullet)$ existe, il en est de
même de l'autre, et on a un isomorphisme canonique}
\begin{equation}
\label{0.13.5.6.1-fr}
\tag{13.5.6.1}
j_\infty:\varprojlim_r\mathrm E_r^{pq}(\mathbf X^\bullet)
\xrightarrow{\sim}
\varprojlim_k\mathrm E_\infty^{pq}(X_k^\bullet).
\end{equation}
\emph{En outre, pour que le système projectif
$(\mathrm E_r^{pq}(\mathbf X^\bullet))_{r\in\mathbb Z}$ soit
essentiellement constant
(\hyperref[0.13.4.2-fr]{13.4.2}), il faut et il suffit que le système
projectif $(\mathrm E_\infty^{pq}(X_k^\bullet))_{k\in\mathbb Z}$ le soit.}
\end{corollary}

\begin{env}[13.5.7]
\label{0.13.5.7-fr}
On note $\mathrm B_\infty^{pq}(\mathbf X^\bullet)$ et
$\mathrm Z_\infty^{pq}(\mathbf X^\bullet)$ les sous-objets de
$\mathrm E_1^{pq}(\mathbf X^\bullet)$ égaux respectivement à
$\mathrm B_r^{pq}(\mathbf X^\bullet)$ pour $r>p-p_0$ et à
$\inf_r\mathrm Z_r^{pq}(\mathbf X^\bullet)$ (quand ce dernier existe), de
sorte que $\varprojlim_r\mathrm E_r^{pq}(\mathbf X^\bullet)$
\oldpage[0\textsubscript{III}]{73}
s'identifie canoniquement à
$\mathrm E_\infty^{pq}(\mathbf X^\bullet)=
\mathrm Z_\infty^{pq}(\mathbf X^\bullet)/
\mathrm B_\infty^{pq}(\mathbf X^\bullet)$. On remarquera que les objets
$\mathrm Z_r^{pq}(\mathbf X^\bullet)$,
$\mathrm B_r^{pq}(\mathbf X^\bullet)$,
$\mathrm E_r^{pq}(\mathbf X^\bullet)$ $(1\leq r\leq+\infty)$ et
$d_r^{pq}$ dépendent \emph{fonctoriellement} du système projectif
$\mathbf X^\bullet$ soumis aux restrictions de
(\hyperref[0.13.5.5-fr]{13.5.5}), et que les morphismes définis dans
(\hyperref[0.13.5.5-fr]{13.5.5}) et
(\hyperref[0.13.5.6-fr]{13.5.6}) sont fonctoriels.
\end{env}

\subsection{Suite spectrale d'un foncteur relative à un objet muni d'une filtration finie}
\label{subsection:0.13.6-fr}

\begin{env}[13.6.1]
\label{0.13.6.1-fr}
Soient $C$, $C'$ deux catégories abéliennes, $\mathrm T:C\to C'$ un
foncteur additif covariant. Supposons que tout objet de $C$ soit isomorphe
à un sous-objet d'un objet injectif, de sorte que les foncteurs dérivés
droits $\mathrm R^p\mathrm T$ $(p\geq0)$ existent.
\end{env}

\begin{lemma}[13.6.2]
\label{0.13.6.2-fr}
\emph{Soit $A$ un objet de $C$, muni d'une filtration finie
$(\mathrm F^i(A))_{i\in\mathbb Z}$. Il existe une résolution injective
$X^\bullet=(X^j)_{j\geq0}$ de $A$ munie d'une filtration finie
$(\mathrm F^i(X^\bullet))_{i\in\mathbb Z}$ telle que la relation
$\mathrm F^i(A)=A$ (resp. $\mathrm F^i(A)=0$) entraîne
$\mathrm F^i(X^\bullet)=X^\bullet$ (resp. $\mathrm F^i(X^\bullet)=0$) et
que, pour tout $i\in\mathbb Z$, $\mathrm F^i(X^\bullet)$ soit une
résolution injective de $\mathrm F^i(A)$.}
\end{lemma}
\begin{proof}
Soit $p$ (resp. $q>p$) le plus grand indice tel que
$\mathrm F^p(A)=A$ (resp. le plus petit indice pour lequel
$\mathrm F^q(A)=0$). On raisonne par récurrence sur $q-p$, le lemme étant
évident pour $q-p=1$. Ayant formé une résolution injective
${X''}^\bullet$ de $A/\mathrm F^{q-1}(A)$ ayant les propriétés voulues,
on considère la suite exacte
\[
0\to\mathrm F^{q-1}(A)\to A\to A/\mathrm F^{q-1}(A)\to0,
\]
on prend une résolution injective ${X'}^\bullet$ de
$\mathrm F^{q-1}(A)$, puis on détermine une résolution injective
$X^\bullet$ de $A$ de façon à avoir une suite exacte
\[
0\to {X'}^\bullet\to X^\bullet\to {X''}^\bullet\to0
\]
compatible avec la précédente (M, V, 2.2)~; il est clair que
$X^\bullet$ répond à la question.
\end{proof}

\begin{corollary}[13.6.3]
\label{0.13.6.3-fr}
\emph{Soit $B$ un second objet de $C$, muni d'une filtration finie
$(\mathrm F^i(B))_{i\in\mathbb Z}$, $s$ un entier, et soit $u:A\to B$
un morphisme tel que $u(\mathrm F^i(A))\subset\mathrm F^{i+s}(B)$ pour
tout $i\in\mathbb Z$. Si $Y^\bullet=(Y^j)_{j\geq0}$ est une résolution
injective de $B$ munie d'une filtration
$(\mathrm F^i(Y^\bullet))_{i\in\mathbb Z}$ ayant les propriétés énoncées
en (\hyperref[0.13.6.2-fr]{13.6.2}), il existe un morphisme
$v:X^\bullet\to Y^\bullet$ compatible avec $u$ et tel que
$v(\mathrm F^i(X^\bullet))\subset\mathrm F^{i+s}(Y^\bullet)$ pour tout
$i\in\mathbb Z$. En outre, deux tels morphismes $v$, $v'$ sont
homotopes.}
\end{corollary}

Cela résulte aussitôt par récurrence sur $q-p$ de la construction
précédente et de (M, V, 2.3).

\begin{env}[13.6.4]
\label{0.13.6.4-fr}
Sous les hypothèses de (\hyperref[0.13.6.2-fr]{13.6.2}), considérons
maintenant le complexe $\mathrm T(X^\bullet)$ dans $C'$, qui est
évidemment filtré par les complexes
$\mathrm T(\mathrm F^i(X^\bullet))$, puisque
$\mathrm F^i(X^\bullet)$ est facteur direct de $X^\bullet$. Il résulte de
(\hyperref[0.13.6.3-fr]{13.6.3}) que la \emph{suite spectrale} de ce
complexe filtré ne dépend que de l'objet filtré $A$, à un isomorphisme
près. Son aboutissement est la cohomologie $\mathrm R^\bullet\mathrm T(A)$,
avec la filtration
\begin{equation}
\label{0.13.6.4.1-fr}
\tag{13.6.4.1}
\begin{split}
\mathrm F^p(\mathrm R^n\mathrm T(A))
&=\operatorname{Im}\bigl(\mathrm R^n\mathrm T(\mathrm F^p(A))
\longrightarrow\mathrm R^n\mathrm T(A)\bigr)\\
&=\operatorname{Ker}\bigl(\mathrm R^n\mathrm T(A)
\longrightarrow\mathrm R^n\mathrm T(A/\mathrm F^p(A))\bigr)
\end{split}
\end{equation}
(\hyperref[0.11.2.2-fr]{11.2.2}), et son terme $\mathrm E_1$ est donné
par
\begin{equation}
\label{0.13.6.4.2-fr}
\tag{13.6.4.2}
\mathrm E_1^{pq}=\mathrm R^{p+q}\mathrm T(\operatorname{gr}^p(A))
\end{equation}
$\operatorname{gr}^p(A)$ désignant comme d'ordinaire
$\mathrm F^p(A)/\mathrm F^{p+1}(A)$. Il est clair, d'après
(\hyperref[0.11.2.2-fr]{11.2.2}), que la filtration de l'aboutissement
est finie, et que pour $p$, $q$ donnés, les suites des
\oldpage[0\textsubscript{III}]{74}
$\mathrm B_r^{pq}(A)=\mathrm B_r^{pq}(\mathrm T(X^\bullet))$ et
$\mathrm Z_r^{pq}(A)=\mathrm Z_r^{pq}(\mathrm T(X^\bullet))$ sont
stationnaires, donc la suite spectrale précédente est \emph{birégulière}
(\hyperref[0.11.1.3-fr]{11.1.3}). Nous noterons cette suite
$\mathrm E(A)=(\mathrm E_r^{pq}(A))$ et nous dirons que c'est la
\emph{suite spectrale du foncteur $\mathrm T$ relative à l'objet filtré
$A$}.
\end{env}

\begin{env}[13.6.5]
\label{0.13.6.5-fr}
Supposons maintenant vérifiées les hypothèses de
(\hyperref[0.13.6.3-fr]{13.6.3}), dont nous conservons les notations.
Comme $\mathrm F^i(X^\bullet)$ (resp. $\mathrm F^i(Y^\bullet)$) est
facteur direct de $X^\bullet$ (resp. $Y^\bullet$), on a
\[
(\mathrm T v)(\mathrm T(\mathrm F^i(X^\bullet)))
\subset\mathrm T(\mathrm F^{i+s}(Y^\bullet))
\]
pour tout $i\in\mathbb Z$~; les définitions de
(\hyperref[0.11.2.2-fr]{11.2.2}) montrent alors que pour
$1\leq r\leq+\infty$, $\mathrm T v$ définit un morphisme
$\mathrm B_r^{pq}(\mathrm T(X^\bullet))\to
\mathrm B_r^{p+s,q-s}(\mathrm T(Y^\bullet))$ et un morphisme
$\mathrm Z_r^{pq}(\mathrm T(X^\bullet))\to
\mathrm Z_r^{p+s,q-s}(\mathrm T(Y^\bullet))$, d'où un morphisme
\[
w_r:\mathrm E_r^{pq}(A)\to\mathrm E_r^{p+s,q-s}(B)~;
\]
de même, on a pour l'aboutissement des morphismes
$u_n:\mathrm R^n\mathrm T(A)\to\mathrm R^n\mathrm T(B)$ tels que
$u_n(\mathrm F^p(\mathrm R^n\mathrm T(A)))\subset
\mathrm F^{p+s}(\mathrm R^n\mathrm T(B))$.

La définition des $d_r^{pq}$ (M, XV, 1) montre en outre que les
diagrammes
\[
\xymatrix{
\mathrm E_r^{pq}(A)\ar[r]^{d_r^{pq}}\ar[d]_{w_r}
&\mathrm E_r^{p+r,q-r+1}(A)\ar[d]^{w_r}\\
\mathrm E_r^{p+s,q-s}(B)\ar[r]_{d_r^{p+s,q-s}}
&\mathrm E_r^{p+r+s,q-r-s+1}(B)
}
\]
sont commutatifs~; on en déduit un diagramme commutatif analogue pour
les isomorphismes $\alpha_r^{pq}$, que nous laisserons au lecteur le
soin d'expliciter. Enfin (\emph{loc. cit.}), on a aussi des diagrammes
commutatifs pour les aboutissements
\[
\xymatrix{
\mathrm E_\infty^{pq}(A)\ar[r]^{\beta^{pq}}\ar[d]_{w_\infty}
&\operatorname{gr}^p(\mathrm R^{p+q}\mathrm T(A))
  \ar[d]^{u_{p+q}}\\
\mathrm E_\infty^{p+s,q-s}(B)\ar[r]_{\beta^{p+s,q-s}}
&\operatorname{gr}^{p+s}(\mathrm R^{p+q}\mathrm T(B)).
}
\]
\end{env}

\begin{env}[13.6.6]
\label{0.13.6.6-fr}
Supposons en particulier qu'il existe un anneau $S$, muni d'une
filtration $(\mathrm F^i(S))_{i\in\mathbb Z}$, et un homomorphisme
d'anneaux
\begin{equation}
\label{0.13.6.6.1-fr}
\tag{13.6.6.1}
h:S\to\operatorname{Hom}_C(A,A)
\end{equation}
\oldpage[0\textsubscript{III}]{75}
tel que pour tout $t\in\mathrm F^i(S)$, on ait
$h_t(\mathrm F^j(A))\subset\mathrm F^{i+j}(A)$ pour tout couple $i$,
$j$. Nous dirons pour abréger que $A$ est alors muni d'une structure de
\emph{$S$-$C$-module filtré} sur l'anneau filtré $S$. Par passage aux
objets gradués associés, tout $h_t$, pour $t\in\mathrm F^j(S)$, définit
un endomorphisme gradué $\overline h_t$ de
$\operatorname{gr}^\bullet(A)$, homogène de degré $j$~; en outre, ce
morphisme ne dépend que de la classe de $t$ dans
$\operatorname{gr}^j(S)$, et on définit ainsi un homomorphisme d'anneaux
gradués
\[
\overline h:\operatorname{gr}^\bullet(S)\to
\operatorname{Hom}_C^g(\operatorname{gr}^\bullet(A),
\operatorname{gr}^\bullet(A))
\]
où le second membre est l'anneau des endomorphismes gradués de
$\operatorname{gr}^\bullet(A)$. Nous dirons que
$\operatorname{gr}^\bullet(A)$ est muni d'une structure de
\emph{$\operatorname{gr}^\bullet(S)$-$C$-module gradué}. Il résulte alors
de (\hyperref[0.13.6.5-fr]{13.6.5}) que pour
$1\leq r\leq+\infty$, tout
$\overline t\in\operatorname{gr}^j(S)$ définit canoniquement dans les
objets bigradués
$(\mathrm B_r^{pq}(A))_{(p,q)\in\mathbb Z\times\mathbb Z}$,
$(\mathrm Z_r^{pq}(A))_{(p,q)\in\mathbb Z\times\mathbb Z}$ et
$\mathrm E_r(A)=(\mathrm E_r^{pq}(A))_{(p,q)\in\mathbb Z\times\mathbb Z}$
des endomorphismes bigradués de degrés $(s,-s)$~; dans
$\mathrm E_r(A)$ (pour $r$ fini), cet endomorphisme commute à
l'endomorphisme bigradué défini par les $d_r^{pq}$. Comme ces
endomorphismes vérifient les conditions habituelles d'associativité et
de distributivité par rapport à l'addition dans
$\operatorname{gr}^\bullet(S)$ et dans les objets bigradués considérés,
nous dirons pour abréger que ces derniers sont des
\emph{$\operatorname{gr}^\bullet(S)$-$C'$-modules bigradués}~; il est
immédiat que les $\alpha_r^{pq}$ définissent un isomorphisme pour ce type
de structures. Pour tout entier $n$, on notera
$\mathrm B_r^{(n)}(A)$ (resp. $\mathrm Z_r^{(n)}(A)$,
$\mathrm E_r^{(n)}(A)$) le sous-objet gradué de
$\mathrm B_r^\bullet(A)$ (resp. $\mathrm Z_r^\bullet(A)$,
$\mathrm E_r^\bullet(A)$) formé des $\mathrm B_r^{pq}(A)$ (resp.
$\mathrm Z_r^{pq}(A)$, $\mathrm E_r^{pq}(A)$) tels que $p+q=n$ (pour
$1\leq r\leq+\infty$)~; il est immédiat que ce sont des
\emph{$\operatorname{gr}^\bullet(S)$-$C'$-modules gradués}. Enfin, tout
$\overline t\in\operatorname{gr}^j(S)$ définit pour tout $n$ un
endomorphisme gradué de degré $j$ dans l'objet gradué
$\operatorname{gr}^\bullet(\mathrm R^n\mathrm T(A))$, qui est ainsi muni
d'une structure de
$\operatorname{gr}^\bullet(S)$-$C'$-module gradué~; les
$\beta^{pq}$ (pour $p+q=n$) définissent un isomorphisme de
$\mathrm E_\infty^{(n)}(A)$ sur
$\operatorname{gr}^\bullet(\mathrm R^n\mathrm T(A))$ pour cette espèce
de structure.

On notera que lorsque $C'$ est la catégorie des groupes abéliens, les
structures de $S$-$C'$-module (resp. de
$\operatorname{gr}^\bullet(S)$-$C'$-module gradué ou bigradué) ne sont
autres que les structures usuelles de $S$-module (resp.
$\operatorname{gr}^\bullet(S)$-module gradué, bigradué).
\end{env}

\subsection{Foncteurs dérivés d'une limite projective d'arguments}
\label{subsection:0.13.7-fr}

\begin{env}[13.7.1]
\label{0.13.7.1-fr}
Soient $C$, $C'$ deux catégories abéliennes, $C$ étant supposée telle
que tout objet de $C$ soit sous-objet d'un objet injectif~; soit
$\mathrm T:C\to C'$ un foncteur additif covariant. Considérons un
système projectif strict $\mathbf A=(A_k)_{k\in\mathbb Z}$ dans $C$,
\emph{limité inférieurement}~; de façon précise, nous supposerons que
$A_k=0$ pour $k<k_0$. Nous associons canoniquement à ce système une
filtration $(\mathrm F^p(A_k))_{p\in\mathbb Z}$ sur chaque $A_k$, par les
formules (\hyperref[0.13.4.1.1-fr]{13.4.1.1}), et comme il s'agit d'un
système projectif strict, les morphismes canoniques
\begin{equation}
\label{0.13.7.1.1-fr}
\tag{13.7.1.1}
\mathrm F^i(A_h)/\mathrm F^j(A_h)\to
\mathrm F^i(A_k)/\mathrm F^j(A_k)\qquad(h\geq k)
\end{equation}
pour $i\leq j\leq k+1$ sont des isomorphismes. Rappelons en outre que
l'on a $\mathrm F^{k_0}(A_k)=A_k$ et $\mathrm F^{k+1}(A_k)=0$ pour tout
$k$.
\end{env}

\begin{env}[13.7.2]
\label{0.13.7.2-fr}
Construisons maintenant pour chaque $k$ une résolution injective
$X_k^\bullet=(X_k^j)_{j\geq0}$ de $A_k$ ayant les propriétés de
(\hyperref[0.13.6.2-fr]{13.6.2}). Les morphismes canoniques
$A_{k+1}\to A_k$ permettent
(\hyperref[0.13.6.3-fr]{13.6.3}) de définir pour chaque $k$ un morphisme
de complexes $X_{k+1}^\bullet\to X_k^\bullet$ compatibles avec les
filtrations, et faisant de
$\mathbf X^\bullet=(X_k^\bullet)_{k\in\mathbb Z}$ un système projectif
de complexes.
\oldpage[0\textsubscript{III}]{76}
On peut en outre supposer que ce système projectif est strict. Pour
cela, on observe qu'en vertu de l'isomorphisme
(\hyperref[0.13.7.1.1-fr]{13.7.1.1}), $A_k$ est isomorphe à
$A_{k+1}/\mathrm F^{k+1}(A_{k+1})$~; on peut donc prendre, dans la
construction de $X_{k+1}^\bullet$, la résolution injective de
$A_{k+1}/\mathrm F^{k+1}(A_{k+1})$ égale à $X_k^\bullet$, et il résulte
de (M, V, 2.3) que la construction du morphisme de complexes
$X_{k+1}^\bullet\to X_k^\bullet$ peut se faire de sorte que ce morphisme
fournisse par passage aux quotients un isomorphisme
\[
X_{k+1}^\bullet/\mathrm F^{k+1}(X_{k+1}^\bullet)
\xrightarrow{\sim}X_k^\bullet
\]
respectant les filtrations, ce qui est la condition de
(\hyperref[0.13.5.1-fr]{13.5.1}).
\end{env}

\begin{env}[13.7.3]
\label{0.13.7.3-fr}
Par construction, le système projectif $\mathrm T(\mathbf X^\bullet)$
des complexes $\mathrm T(X_k^\bullet)$ satisfait aux hypothèses de
(\hyperref[0.13.5.3-fr]{13.5.3}). Les résultats de
(\hyperref[0.13.5.4-fr]{13.5.4}),
(\hyperref[0.13.5.5-fr]{13.5.5}) et
(\hyperref[0.13.5.6-fr]{13.5.6}) sont donc applicables aux suites
spectrales $\mathrm E(\mathrm T(X_k^\bullet))=\mathrm E(A_k)$~; nous
écrirons $\mathrm E_r^{pq}(\mathbf A)$ au lieu de
$\mathrm E_r^{pq}(\mathrm T(\mathbf X^\bullet))$ pour
$1\leq r\leq+\infty$ (cf. (\hyperref[0.13.5.7-fr]{13.5.7}) pour
$r=+\infty$) et de même pour les notations analogues. On notera en
particulier que l'on a
\begin{equation}
\label{0.13.7.3.1-fr}
\tag{13.7.3.1}
\mathrm E_1^{pq}(\mathbf A)=
\mathrm R^{p+q}\mathrm T(\operatorname{gr}^p(\mathbf A))
\end{equation}
en vertu de (\hyperref[0.13.6.4.2-fr]{13.6.4.2}) et du fait que le
système $(\operatorname{gr}^p(A_k))$ est essentiellement constant.

Ces résultats et (\hyperref[0.13.4.3-fr]{13.4.3}) donnent la proposition
suivante, démontrée d'abord par Shih Weishu par une méthode différente
(inédite) :
\end{env}

\begin{proposition}[13.7.4]
\label{0.13.7.4-fr}
\emph{(Shih). --- Soit $n$ un entier. Les deux conditions suivantes sont
équivalentes :}
\begin{enumerate}
\item[\emph{a)}] \emph{Pour tout couple $(p,q)$ tel que $p+q=n$, le
système projectif $(\mathrm E_r^{pq}(\mathbf A))_{r\geq2}$ est
essentiellement constant.}
\item[\emph{b)}] \emph{Le système projectif
$\mathrm R^n\mathrm T(\mathbf A)=
(\mathrm R^n\mathrm T(A_k))_{k\in\mathbb Z}$ vérifie (ML).}
\end{enumerate}
\emph{En outre, lorsque ces conditions sont remplies, on a un
isomorphisme canonique}
\begin{equation}
\label{0.13.7.4.1-fr}
\tag{13.7.4.1}
\operatorname{gr}^p(\mathrm R^n\mathrm T(\mathbf A))
\xrightarrow{\sim}\mathrm E_\infty^{p,n-p}(\mathbf A)
\qquad\emph{pour tout $p\in\mathbb Z$.}
\end{equation}
\end{proposition}
\begin{proof}
En effet, en vertu de (\hyperref[0.13.5.6-fr]{13.5.6}), la condition
\emph{a)} équivaut à dire que le système projectif
$(\mathrm E_\infty^{pq}(A_k))_{k\in\mathbb Z}$ est essentiellement
constant pour $p+q=n$, et d'autre part
$\operatorname{gr}^p(\mathrm R^n\mathrm T(A_k))$ est canoniquement
isomorphe à $\mathrm E_\infty^{p,n-p}(A_k)$, donc il résulte de
(\hyperref[0.13.4.3-fr]{13.4.3}) que \emph{a)} et \emph{b)} sont
équivalentes~; l'isomorphisme
(\hyperref[0.13.7.4.1-fr]{13.7.4.1}) n'est autre que
(\hyperref[0.13.5.6.1-fr]{13.5.6.1}) appliqué au cas envisagé ici.
\end{proof}

\begin{corollary}[13.7.5]
\label{0.13.7.5-fr}
\emph{Soit $\mathcal F=(\mathcal F_k)_{k\in\mathbb N}$ un système
projectif de faisceaux de groupes abéliens satisfaisant aux conditions
\emph{(i)}, \emph{(ii)} et \emph{(iii)} de
(\hyperref[0.13.3.1-fr]{13.3.1}) et soit
$\mathcal F=\varprojlim_k\mathcal F_k$. Supposons que, pour le foncteur
$\mathcal G\mapsto\Gamma(X,\mathcal G)$, le système projectif
$(\mathrm E_r^{pq}(\mathcal F))_{r\in\mathbb Z}$ soit essentiellement
constant pour tout couple $(p,q)$ tel que $p+q=n$ ou $p+q=n+1$.
Considérons sur $H^{n+1}(X,\mathcal F)$ la filtration définie par
$\mathrm F^p(H^{n+1}(X,\mathcal F))=
\operatorname{Ker}(H^{n+1}(X,\mathcal F)\to
H^{n+1}(X,\mathcal F_{p-1}))$. On a alors un isomorphisme canonique}
\begin{equation}
\label{0.13.7.5.1-fr}
\tag{13.7.5.1}
\operatorname{gr}^p(H^{n+1}(X,\mathcal F))
\xrightarrow{\sim}\mathrm E_\infty^{p,n-p+1}(\mathcal F)
\qquad\emph{pour tout $p\in\mathbb Z$.}
\end{equation}
\end{corollary}
\begin{proof}
\oldpage[0\textsubscript{III}]{77}
Il résulte de (\hyperref[0.13.7.4-fr]{13.7.4}) appliqué au cas où $C$
est la catégorie des faisceaux de groupes abéliens sur $X$, $C'$ la
catégorie des groupes abéliens, et $\mathrm T=\Gamma$, que l'on a un
isomorphisme canonique
$\operatorname{gr}^p(\mathrm R^{n+1}\Gamma(\mathcal F))
\xrightarrow{\sim}\mathrm E_\infty^{p,n-p+1}(\mathcal F)$ pour tout
$p\in\mathbb Z$. D'autre part, comme en vertu de
(\hyperref[0.13.7.4-fr]{13.7.4}), le système projectif
$(H^n(X,\mathcal F_k))_{k\in\mathbb Z}$ vérifie (ML), on déduit de
(\hyperref[0.13.3.1-fr]{13.3.1}) un isomorphisme canonique
\begin{equation}
\label{0.13.7.5.1b-fr}
\tag{13.7.5.1}
H^{n+1}(X,\mathcal F)\xrightarrow{\sim}
\varprojlim_k H^{n+1}(X,\mathcal F_k).
\end{equation}

Comme le système projectif $\mathrm R^{n+1}\Gamma(\mathcal F)$ vérifie
(ML) en vertu de (\hyperref[0.13.7.4-fr]{13.7.4}), on a un isomorphisme
canonique
\[
\operatorname{gr}^p(\mathrm R^{n+1}\Gamma(\mathcal F))
\xrightarrow{\sim}
\varprojlim_k\operatorname{gr}^p(H^{n+1}(X,\mathcal F_k))
\]
(\hyperref[0.13.4.3-fr]{13.4.3}), et un isomorphisme canonique
\[
\varprojlim_k\operatorname{gr}^p(H^{n+1}(X,\mathcal F_k))
\xrightarrow{\sim}
\operatorname{gr}^p(\varprojlim_k H^{n+1}(X,\mathcal F_k))
\]
(\hyperref[0.13.4.5-fr]{13.4.5}). Tout revient donc à voir que
l'isomorphisme (\hyperref[0.13.7.5.1-fr]{13.7.5.1}) est compatible avec
les filtrations des deux membres~; mais cela résulte aussitôt des
définitions et de la commutativité du diagramme
\[
\xymatrix{
H^{n+1}(X,\mathcal F)\ar[r]^{\sim}\ar[dr]
&\varprojlim_k H^{n+1}(X,\mathcal F_k)\ar[d]\\
&H^{n+1}(X,\mathcal F_{p-1})
}
\]
pour tout $p$.
\end{proof}

\begin{env}[13.7.6]
\label{0.13.7.6-fr}
Soit $S$ un anneau muni d'une filtration
$(\mathrm F^i(S))_{i\in\mathbb Z}$ telle que $\mathrm F^0(S)=S$ (donc
$\operatorname{gr}^i(S)=0$ pour $i<0$). Supposons que chacun des $A_k$,
muni de la filtration définie dans
(\hyperref[0.13.7.1-fr]{13.7.1}), soit un $S$-$C$-module filtré
(\hyperref[0.13.6.6-fr]{13.6.6}), les morphismes $A_h\to A_k$ pour
$k\leq h$ étant des morphismes pour la structure de $S$-$C$-module
filtré~; nous dirons pour abréger que $\mathbf A$ est un \emph{système
projectif de $S$-$C$-modules filtrés}. Alors il est immédiat que les
morphismes $\mathrm B_r^{pq}(A_h)\to\mathrm B_r^{pq}(A_k)$ et
$\mathrm Z_r^{pq}(A_h)\to\mathrm Z_r^{pq}(A_k)$ pour $k\leq h$,
$1\leq r\leq+\infty$, sont des morphismes pour les structures de
$\operatorname{gr}^\bullet(S)$-$C'$-module bigradué
(\hyperref[0.13.6.5-fr]{13.6.5}), et que les familles
$(\mathrm Z_r^{pq}(\mathbf A))$, $(\mathrm B_r^{pq}(\mathbf A))$ et
$(\mathrm E_r^{pq}(\mathbf A))$ sont des
$\operatorname{gr}^\bullet(S)$-$C'$-modules bigradués pour $r$ fini,
les deux premiers étant des sous-modules de
$(\mathrm E_1^{pq}(\mathbf A))$. On notera encore
$\mathrm Z_r^{(n)}(\mathbf A)$, $\mathrm B_r^{(n)}(\mathbf A)$,
$\mathrm E_r^{(n)}(\mathbf A)$ les sous-objets respectifs des précédents
obtenus en ne prenant que les termes tels que $p+q=n$~; ce sont des
$\operatorname{gr}^\bullet(S)$-$C'$-modules gradués.

Lorsque le système $(\mathrm E_r^{pq}(\mathbf A))_{r\in\mathbb Z}$ est
essentiellement constant, $(\mathrm E_\infty^{pq}(\mathbf A))$ est donc
aussi un $\operatorname{gr}^\bullet(S)$-$C'$-module bigradué, et chaque
$\mathrm E_\infty^{(n)}(\mathbf A)$ un
$\operatorname{gr}^\bullet(S)$-$C'$-module gradué. En outre, les
$\beta^{p,n-p}:\mathrm E_\infty^{p,n-p}(A_k)\xrightarrow{\sim}
\operatorname{gr}^p(\mathrm R^n\mathrm T(A_k))$ constituent pour chaque
$k$ un isomorphisme pour la structure de
$\operatorname{gr}^\bullet(S)$-$C'$-module gradué de
$\mathrm E_\infty^{(n)}(A_k)$ sur
$\operatorname{gr}^\bullet(\mathrm R^n\mathrm T(A_k))$~; si on est dans
les conditions précédentes,
$\beta^{p,n-p}:\mathrm E_\infty^{(n)}(\mathbf A)\xrightarrow{\sim}
\varprojlim_k\operatorname{gr}^\bullet(\mathrm R^n\mathrm T(A_k))$
sera donc aussi un isomorphisme pour ces structures et il en est
évidemment de même de l'isomorphisme canonique
$\operatorname{gr}^\bullet(\mathrm R^n\mathrm T(\mathbf A))
\xrightarrow{\sim}
\varprojlim_k\operatorname{gr}^\bullet(\mathrm R^n\mathrm T(A_k))$,
donc les isomorphismes
(\hyperref[0.13.7.4.1-fr]{13.7.4.1}) constituent un isomorphisme pour les
structures de $\operatorname{gr}^\bullet(S)$-$C'$-module gradué.
\end{env}

\begin{proposition}[13.7.7]
\label{0.13.7.7-fr}
Soit $S$ un anneau noethérien $\mathfrak J$-adique. Supposons que $C$ soit
une catégorie abélienne dont tout objet est isomorphe à un sous-objet d'un
objet injectif, et soit $\mathrm T$ un foncteur covariant additif de $C$ dans
la catégorie des groupes abéliens. Soit
$\mathbf A=(A_k)_{k\in\mathbb Z}$ un
\oldpage[0\textsubscript{III}]{78}
système projectif strict de $S$-$C$-modules filtrés (pour la filtration
$\mathfrak J$-adique sur $S$) limité inférieurement. On suppose que pour un
entier $n$ donné, la condition suivante soit vérifiée~:
\begin{itemize}
\item[$(\mathrm F_n)$] Le $\operatorname{gr}^\bullet(S)$-module gradué
\[
\mathrm E_1^{(m)}(\mathbf A)=
\bigl(\mathrm R^m\mathrm T(\operatorname{gr}^p(\mathbf A))\bigr)_{p\in\mathbb Z}
\]
de (\hyperref[0.13.7.3.1-fr]{13.7.3.1}) est de type fini pour $m=n$ et
$m=n+1$.
\end{itemize}
Dans ces conditions~:
\begin{enumerate}
\item[\rm(i)] Les systèmes projectifs
$(\mathrm R^n\mathrm T(A_k))_{k\in\mathbb Z}$ et
$(\mathrm R^{n+1}\mathrm T(A_k))_{k\in\mathbb Z}$ vérifient (ML).
\item[\rm(ii)] Si on pose
\[
\mathrm R'^{\,n}\mathrm T(\mathbf A)=
\varprojlim_k\mathrm R^n\mathrm T(A_k),
\]
$\mathrm R'^{\,n}\mathrm T(\mathbf A)$ est un $S$-module de type fini.
\item[\rm(iii)] La filtration définie par
\[
\mathrm F^p(\mathrm R'^{\,n}\mathrm T(\mathbf A))=
\operatorname{Ker}\bigl(\mathrm R'^{\,n}\mathrm T(\mathbf A)
\longrightarrow\mathrm R^n\mathrm T(A_{p-1})\bigr)
\quad(p\in\mathbb Z)
\]
sur $\mathrm R'^{\,n}\mathrm T(\mathbf A)$ est $\mathfrak J$-bonne
(c'est-à-dire que
$\mathfrak J\mathrm F^p(\mathrm R'^{\,n}\mathrm T(\mathbf A))
\subset\mathrm F^{p+1}(\mathrm R'^{\,n}\mathrm T(\mathbf A))$ pour tout
$p$, l'égalité des deux membres ayant lieu dès que $p$ est assez grand). En
particulier, la topologie sur $\mathrm R'^{\,n}\mathrm T(\mathbf A)$ définie
par cette filtration est identique à la topologie $\mathfrak J$-adique.
\item[\rm(iv)] Le système projectif
$(\mathrm E_r^{pq}(\mathbf A))_{r\in\mathbb Z}$ est essentiellement constant
pour $p+q=n$ et $p+q=n+1$, $\mathrm E_\infty^{pq}(\mathbf A)$ est donc défini
(\hyperref[0.13.5.7-fr]{13.5.7}) et on a un isomorphisme canonique de
$\operatorname{gr}^\bullet(S)$-modules gradués
\begin{equation}
\label{0.13.7.7.1-fr}
\tag{13.7.7.1}
\operatorname{gr}^p(\mathrm R'^{\,n}\mathrm T(\mathbf A))
\xrightarrow{\sim}\mathrm E_\infty^{p,n-p}(\mathbf A)
\quad(p\in\mathbb Z).
\end{equation}
\end{enumerate}
\end{proposition}

\begin{proof}
On notera que l'isomorphisme (\hyperref[0.13.7.7.1-fr]{13.7.7.1})
permettra de noter $\mathrm R^n\mathrm T(\mathbf A)$, par abus de langage, la
limite projective $\mathrm R'^{\,n}\mathrm T(\mathbf A)$ du système projectif
$\mathrm R^n\mathrm T(A_k)$, compte tenu des isomorphismes
(\hyperref[0.13.7.4.1-fr]{13.7.4.1}).

Comme l'anneau gradué $\operatorname{gr}^\bullet(S)$ est noethérien
(Bourbaki, \emph{Alg. comm.}, chap.~III, \S~2, n\textsuperscript{o}~9,
cor.~5 du th.~2), la suite croissante des sous-$\operatorname{gr}^\bullet(S)$-
modules gradués $\mathrm B_r^{(m)}(\mathbf A)$ de
$\mathrm E_1^{(m)}(\mathbf A)$ (\hyperref[0.13.6.6-fr]{13.6.6}) est
stationnaire pour $m=n$ et $m=n+1$, et par suite la condition b) de
(\hyperref[0.11.1.10-fr]{11.1.10}) est vérifiée. Il s'ensuit que la condition
a) de (\hyperref[0.13.7.4-fr]{13.7.4}) est remplie pour $n$ et pour $n+1$,
et cela prouve déjà (i). En outre, les isomorphismes
(\hyperref[0.13.7.4.1-fr]{13.7.4.1}) (compte tenu des remarques de
(\hyperref[0.13.7.6-fr]{13.7.6})) montrent que
$\operatorname{gr}^\bullet(\mathrm R^n\mathrm T(\mathbf A))$ est un
$\operatorname{gr}^\bullet(S)$-module gradué isomorphe à
\[
\mathrm E_\infty^{(n)}(\mathbf A)=
\mathrm Z_\infty^{(n)}(\mathbf A)/\mathrm B_\infty^{(n)}(\mathbf A)~;
\]
comme $\mathrm Z_\infty^{(n)}(\mathbf A)$ est un sous-module de
$\mathrm E_1^{(n)}(\mathbf A)$, il est de type fini, et il en est donc de
même de $\mathrm E_\infty^{(n)}(\mathbf A)$. En outre, pour la filtration
$(\mathrm F^p(\mathrm R'^{\,n}\mathrm T(\mathbf A)))$, il résulte de
(\hyperref[0.13.4.5-fr]{13.4.5}) que
$\operatorname{gr}^\bullet(\mathrm R^n\mathrm T(\mathbf A))$ et
$\operatorname{gr}^\bullet(\mathrm R'^{\,n}\mathrm T(\mathbf A))$ sont des
$\operatorname{gr}^\bullet(S)$-modules isomorphes, ce qui démontre (iv). Les
assertions (ii) et (iii) seront enfin conséquences des résultats précédents
et du lemme suivant~:
\end{proof}

\begin{lemma}[13.7.7.2]
\label{0.13.7.7.2-fr}
Soient $S$ un anneau noethérien $\mathfrak J$-adique, $M$ un $S$-module muni
d'une filtration co-discrète $(\mathrm F^p(M))_{p\in\mathbb Z}$ telle que
$\mathfrak J\mathrm F^p(M)\subset\mathrm F^{p+1}(M)$ (ce qui exprime que
$M$ est un module filtré sur l'anneau $S$ filtré par la filtration
$\mathfrak J$-adique). Supposons en outre $M$ séparé pour la topologie définie
par la filtration $(\mathrm F^p(M))$. Alors les conditions suivantes sont
équivalentes~:
\begin{enumerate}
\item[\rm a)] $M$ est un $S$-module de type fini et $(\mathrm F^p(M))$ une
filtration $\mathfrak J$-bonne.
\item[\rm b)] $\operatorname{gr}^\bullet(M)$ est un
$\operatorname{gr}^\bullet(S)$-module de type fini.
\item[\rm c)] Les $\operatorname{gr}^p(M)$ sont des $S$-modules de type fini
et pour $p$ assez grand les homomorphismes canoniques
\begin{equation}
\label{0.13.7.7.3-fr}
\tag{13.7.7.3}
\mathfrak J\otimes_S\operatorname{gr}^p(M)\longrightarrow
\operatorname{gr}^{p+1}(M)
\end{equation}
\oldpage[0\textsubscript{III}]{79}
(déduits de
$\mathfrak J\otimes_S\mathrm F^p(M)\to\mathrm F^{p+1}(M)$, compte tenu de ce
que l'image de l'homomorphisme composé
\[
\mathfrak J\otimes_S\mathrm F^{p+1}(M)\longrightarrow
\mathfrak J\otimes_S\mathrm F^p(M)\longrightarrow\mathrm F^{p+1}(M)
\]
est $\mathfrak J\mathrm F^{p+1}(M)\subset\mathrm F^{p+2}(M)$) sont
surjectifs.
\end{enumerate}
Pour la démonstration, voir Bourbaki, \emph{Alg. comm.}, chap.~III, \S~3,
n\textsuperscript{o}~1, prop.~3.
\end{lemma}

\begin{env}[13.7.7.4]
\label{0.13.7.7.4-fr}
Pour appliquer le lemme (\hyperref[0.13.7.7.2-fr]{13.7.7.2}), il reste à
observer que la topologie définie sur
$\mathrm R'^{\,n}\mathrm T(\mathbf A)$ par la filtration considérée fait de
$\mathrm R'^{\,n}\mathrm T(\mathbf A)$ un $S$-module séparé et complet,
cette topologie étant celle de la limite projective des groupes discrets
$\mathrm R^n\mathrm T(A_k)$~; d'autre part, si $A_k=0$ pour $k<k_0$, on a
aussi $\mathrm R^n\mathrm T(A_k)=0$ pour $k<k_0$, donc
$\mathrm F^{k_0}(\mathrm R'^{\,n}\mathrm T(\mathbf A))=
\mathrm R'^{\,n}\mathrm T(\mathbf A)$, et on est bien dans les conditions
d'application du lemme.
\end{env}

\begin{corollary}[13.7.8]
\label{0.13.7.8-fr}
Si l'hypothèse $(\mathrm F_n)$ est vérifiée, on a, pour tout élément $f\in S$,
un isomorphisme canonique
\begin{equation}
\label{0.13.7.8.1-fr}
\tag{13.7.8.1}
\varprojlim_k\bigl((\mathrm R^n\mathrm T(A_k))_f\bigr)
\xrightarrow{\sim}\mathrm R^n\mathrm T(\mathbf A)\otimes_S S_{\{f\}}.
\end{equation}
\end{corollary}

\begin{proof}
En effet, $\mathrm R^n\mathrm T(\mathbf A)$ est un $S$-module de type fini,
$S_{\{f\}}$ une $S$-algèbre adique noethérienne
(\hyperref[0.7.6.11-fr]{0\textsubscript{I}, 7.6.11}), séparée complétée de
$S_f$ pour la topologie $\mathfrak J$-préadique
(\hyperref[0.7.6.2-fr]{0\textsubscript{I}, 7.6.2}). On conclut de
(\hyperref[0.7.7.8-fr]{0\textsubscript{I}, 7.7.8}) et
(\hyperref[0.7.7.1-fr]{0\textsubscript{I}, 7.7.1}) que
$\mathrm R^n\mathrm T(\mathbf A)\otimes_S S_{\{f\}}$ est isomorphe au séparé
complété de $\mathrm R^n\mathrm T(\mathbf A)\otimes_S S_f$ pour la topologie
$\mathfrak J$-préadique~; un système fondamental de voisinages de $0$ pour
cette topologie est
$(\mathfrak J^p\mathrm R^n\mathrm T(\mathbf A))\otimes_S S_f$, donc
$\mathrm F^p(\mathrm R^n\mathrm T(\mathbf A))\otimes_S S_f$ est aussi un tel
système~; ce dernier est le noyau de l'application canonique
$(\mathrm R^n\mathrm T(\mathbf A))_f\to
(\mathrm R^n\mathrm T(A_{p-1}))_f$ et par suite le groupe séparé associé à
$\mathrm R^n\mathrm T(\mathbf A)\otimes_S S_f$ s'identifie à un sous-groupe
$G$ de $\varprojlim_k((\mathrm R^n\mathrm T(A_k))_f)$. Mais le système
projectif $((\mathrm R^n\mathrm T(A_k))_f)$ vérifie évidemment la condition
(ML), et l'image de $(\mathrm R^n\mathrm T(\mathbf A))_f$ dans chacun des
$(\mathrm R^n\mathrm T(A_k))_f$ est égale à l'image commune des
$(\mathrm R^n\mathrm T(A_h))_f$ pour $h\geq k$ assez grand. On en conclut
aussitôt que $G$ est \emph{partout dense} dans
$\varprojlim_k((\mathrm R^n\mathrm T(A_k))_f)$, et comme ce dernier groupe est
séparé et complet, le corollaire est démontré.
\end{proof}

\begin{flushright}
\textit{(A suivre.)}
\end{flushright}

\clearpage
\thispagestyle{empty}
\null
\clearpage
\thispagestyle{plain}
\markboth{ÉTUDE COHOMOLOGIQUE DES FAISCEAUX COHÉRENTS}{CHAPITRE III}
\oldpage[III]{81}
\phantomsection
\label{chapter:III-fr}

\begin{center}
{\large CHAPITRE III\par}
\vspace{1.5em}
{\Large\bfseries ÉTUDE COHOMOLOGIQUE\par}
{\Large\bfseries DES FAISCEAUX COHÉRENTS\par}
\vspace{1.8em}
{\bfseries Sommaire\par}
\end{center}

\medskip
\begin{tabular}{@{}r@{\ }p{0.82\textwidth}@{}}
\S~1. & Cohomologie des schémas affines.\\
\S~2. & Étude cohomologique des morphismes projectifs.\\
\S~3. & Le théorème de finitude pour les morphismes propres.\\
\S~4. & Le théorème fondamental des morphismes propres~; applications.\\
\S~5. & Un théorème d'existence de faisceaux algébriques cohérents.\\
\S~6. & Foncteurs Tor et Ext locaux et globaux~; formule de Künneth.\\
\S~7. & Étude du changement de base dans les foncteurs cohomologiques
covariants de Modules.\\
\S~8. & Le théorème de dualité sur les fibrés projectifs.\\
\S~9. & Cohomologie relative et cohomologie locale~; dualité locale.\\
\S~10. & Relations entre cohomologie projective et cohomologie locale.
Technique de complétion formelle le long d'un diviseur.\\
\S~11. & Groupes de Picard globaux et locaux\footnotemark.
\end{tabular}

\medskip
Ce chapitre donne les théorèmes fondamentaux sur la cohomologie des faisceaux
algébriques cohérents, à l'exception de ceux découlant de la théorie des
résidus (théorèmes de dualité), qui feront l'objet d'un chapitre ultérieur.
Parmi les premiers, il y a essentiellement six théorèmes fondamentaux,
faisant l'objet des six premiers paragraphes du présent chapitre. Ces
résultats seront dans la suite des outils essentiels, même dans des questions
qui ne sont pas de nature proprement cohomologique~; le lecteur en verra les
premiers exemples dès le \S~4. Le \S~7 donne des résultats de nature plus
technique, mais d'un usage constant dans les applications. Enfin, dans les
\S\S~8 à 11, nous développerons certains résultats, liés à la dualité des
faisceaux cohérents, particulièrement importants pour les applications, et
qui peuvent s'exposer antérieurement à la théorie générale des résidus.

\footnotetext{Le chapitre IV ne dépend pas des \S\S~8 à 11, et sera sans doute
publié avant ces derniers.}

\oldpage[III]{82}
Le contenu des \S\S~1 et 2 est dû à J.-P. Serre, et le lecteur constatera que
nous n'avons eu qu'à suivre l'exposé de (FAC). Les \S\S~8 et 9 sont également
inspirés par (FAC) (les transpositions nécessitées par les contextes
différents étant toutefois moins évidentes). Enfin, comme nous l'avons dit
dans l'Introduction, le \S~4 doit être considéré comme la mise en forme, en
langage moderne, du «~théorème d'invariance~» fondamental de la «~théorie des
fonctions holomorphes~» de Zariski.

Signalons enfin que les résultats du n\textsuperscript{o}~3.4 (et les
propositions préliminaires de (0, 13.4 à 13.7)) ne seront pas utilisés dans la
suite du chap.~III et peuvent donc être omis en première lecture.

\setcounter{section}{0}
\section{Cohomologie des schémas affines}
\label{section:III.1-fr}

\subsection{Rappels sur le complexe de l'algèbre extérieure}
\label{subsection:III.1.1-fr}

\begin{env}[1.1.1]
\label{III.1.1.1-fr}
Soient $A$ un anneau, $\mathbf f=(f_i)_{1\leq i\leq r}$ un système de $r$
éléments de $A$. Le \emph{complexe de l'algèbre extérieure
$K_\bullet(\mathbf f)$} correspondant à $\mathbf f$ est un complexe de chaînes
(G, I, 2.2) se définissant de la façon suivante~: le $A$-module gradué
$K_\bullet(\mathbf f)$ est égal à l'\emph{algèbre extérieure
$\bigwedge(A^r)$}, graduée de la façon usuelle, et l'opérateur bord est la
\emph{multiplication intérieure $i_{\mathbf f}$} par $\mathbf f$ considéré
comme élément du dual $(A^r)^\vee$~; on rappelle que $i_{\mathbf f}$ est une
\emph{antidérivation} de degré $-1$ de $\bigwedge(A^r)$, et que si
$(\mathbf e_i)_{1\leq i\leq r}$ est la base canonique de $A^r$, on a
$i_{\mathbf f}(\mathbf e_i)=f_i$~; la vérification de la condition
$i_{\mathbf f}\circ i_{\mathbf f}=0$ est immédiate.

Une définition équivalente est la suivante~: pour chaque $i$, on considère
un complexe de chaînes $K_\bullet(f_i)$ défini comme suit~:
$K_0(f_i)=K_1(f_i)=A$, $K_n(f_i)=0$ pour $n\neq0,1$~; l'opérateur bord est
défini par la condition que $d_1:A\to A$ est la \emph{multiplication par
$f_i$}. On prend alors pour $K_\bullet(\mathbf f)$ le \emph{produit tensoriel
$K_\bullet(f_1)\otimes K_\bullet(f_2)\otimes\cdots\otimes K_\bullet(f_r)$}
(G, I, 2.7) muni de son degré total~; la vérification de l'isomorphisme de ce
complexe et du complexe défini plus haut est immédiate.
\end{env}

\begin{env}[1.1.2]
\label{III.1.1.2-fr}
Pour tout $A$-module $M$, on définit le \emph{complexe de chaînes}
\begin{equation}
\label{III.1.1.2.1-fr}
\tag{1.1.2.1}
K_\bullet(\mathbf f,M)=K_\bullet(\mathbf f)\otimes_A M
\end{equation}
et le \emph{complexe de cochaînes} (G, I, 2.2)
\begin{equation}
\label{III.1.1.2.2-fr}
\tag{1.1.2.2}
K^\bullet(\mathbf f,M)=\operatorname{Hom}_A(K_\bullet(\mathbf f),M).
\end{equation}
Si $g$ est une $k$-cochaîne de ce dernier complexe, et si on pose
\[
g(i_1,\ldots,i_k)=g(\mathbf e_{i_1}\wedge\cdots\wedge\mathbf e_{i_k}),
\]
$g$ s'identifie à une application \emph{alternée} de $[1,r]^k$ dans $M$, et
il résulte des définitions précédentes que l'on a
\begin{equation}
\label{III.1.1.2.3-fr}
\tag{1.1.2.3}
d^kg(i_1,i_2,\ldots,i_{k+1})=
\sum_{h=1}^{k+1}(-1)^{h-1}f_{i_h}
g(i_1,\ldots,\widehat{i_h},\ldots,i_{k+1}).
\end{equation}
\end{env}

\oldpage[III]{83}
\begin{env}[1.1.3]
\label{III.1.1.3-fr}
Des complexes précédents, on déduit comme d'ordinaire les \emph{$A$-modules
d'homologie et de cohomologie} (G, I, 2.2)
\begin{align}
\label{III.1.1.3.1-fr}
\tag{1.1.3.1}
H_\bullet(\mathbf f,M)&=H_\bullet(K_\bullet(\mathbf f,M)),\\
\label{III.1.1.3.2-fr}
\tag{1.1.3.2}
H^\bullet(\mathbf f,M)&=H^\bullet(K^\bullet(\mathbf f,M)).
\end{align}
On définit d'ailleurs un $A$-isomorphisme
$K_i(\mathbf f,M)\xrightarrow{\sim}K^{r-i}(\mathbf f,M)$ en faisant
correspondre à toute chaîne
$z=\sum(\mathbf e_{i_1}\wedge\cdots\wedge\mathbf e_{i_k})\otimes
z_{i_1\ldots i_k}$ la cochaîne $g_z$ telle que
$g_z(j_1,\ldots,j_{r-k})=\varepsilon z_{i_1\ldots i_k}$, où
$(j_h)_{1\leq h\leq r-k}$ est la suite strictement croissante complémentaire
de la suite strictement croissante $(i_h)_{1\leq h\leq k}$ dans $[1,r]$ et
$\varepsilon=(-1)^\nu$, $\nu$ étant le nombre d'inversions de la permutation
$i_1,\ldots,i_k,j_1,\ldots,j_{r-k}$ de $[1,r]$. On vérifie que
$g_{dz}=d(g_z)$, ce qui donne un isomorphisme
\begin{equation}
\label{III.1.1.3.3-fr}
\tag{1.1.3.3}
H^i(\mathbf f,M)\xrightarrow{\sim}H_{r-i}(\mathbf f,M)
\qquad\text{pour }0\leq i\leq r.
\end{equation}
Dans ce chapitre, nous aurons surtout à considérer les modules de cohomologie
$H^\bullet(\mathbf f,M)$.

Pour un $\mathbf f$ donné, il est immédiat (G, I, 2.1) que
$M\mapsto H^\bullet(\mathbf f,M)$ est un \emph{foncteur cohomologique}
(T, II, 2.1), de la catégorie des $A$-modules dans celle des $A$-modules
gradués, nul pour les degrés $<0$ et $>r$. En outre, on a
\begin{equation}
\label{III.1.1.3.4-fr}
\tag{1.1.3.4}
H^0(\mathbf f,M)=\operatorname{Hom}_A(A/(\mathbf f),M)
\end{equation}
en désignant par $(\mathbf f)$ l'idéal de $A$ engendré par
$f_1,\ldots,f_r$~; cela résulte aussitôt de
(\hyperref[III.1.1.2.3-fr]{1.1.2.3}), et il est clair que
$H^0(\mathbf f,M)$ s'identifie au sous-module de $M$ \emph{annulé par
$(\mathbf f)$}. De même, on a d'après
(\hyperref[III.1.1.2.3-fr]{1.1.2.3})
\begin{equation}
\label{III.1.1.3.5-fr}
\tag{1.1.3.5}
H^r(\mathbf f,M)=M/\Bigl(\sum_{i=1}^r f_iM\Bigr)
=(A/(\mathbf f))\otimes_A M.
\end{equation}
Nous utiliserons le résultat connu suivant, dont nous rappellerons une
démonstration pour être complet~:
\end{env}

\begin{proposition}[1.1.4]
\label{III.1.1.4-fr}
Soient $A$ un anneau, $\mathbf f=(f_i)_{1\leq i\leq r}$ une famille finie
d'éléments de $A$, $M$ un $A$-module. Si, pour $1\leq i\leq r$, l'homothétie
$z\mapsto f_i z$ dans $M_{i-1}=M/(f_1M+\cdots+f_{i-1}M)$ est injective, on a
$H^i(\mathbf f,M)=0$ pour $i\neq r$.
\end{proposition}

\begin{proof}
Il suffit en effet de prouver que $H_i(\mathbf f,M)=0$ pour tout $i>0$ en
vertu de (\hyperref[III.1.1.3.3-fr]{1.1.3.3}). Raisonnons par récurrence sur
$r$, le cas $r=0$ étant trivial. Posons
$\mathbf f'=(f_i)_{1\leq i\leq r-1}$~; cette famille vérifie les conditions de
l'énoncé, donc, si on pose $L_\bullet=K_\bullet(\mathbf f',M)$, on a
$H_i(L_\bullet)=0$ pour $i>0$ par hypothèse, et
$H_0(L_\bullet)=M_{r-1}$ en vertu de
(\hyperref[III.1.1.3.3-fr]{1.1.3.3}) et
(\hyperref[III.1.1.3.5-fr]{1.1.3.5}). Posons pour abréger
$K_\bullet=K_\bullet(f_r)=K_0\oplus K_1$, avec $K_0=K_1=A$,
$d_1:K_1\to K_0$ étant la multiplication par $f_r$~; on a par définition
(\hyperref[III.1.1.1-fr]{1.1.1})
$K_\bullet(\mathbf f,M)=K_\bullet\otimes_A L_\bullet$. Or, on a le lemme
suivant~:

\begin{lemma}[1.1.4.1]
\label{III.1.1.4.1-fr}
Soit $K_\bullet$ un complexe de chaînes formé de $A$-modules libres, nuls sauf
en dimensions $0$ et $1$. Pour tout complexe de chaînes $L_\bullet$ formé de
$A$-modules, on a une suite exacte
\[
0\longrightarrow H_0(K_\bullet\otimes H_p(L_\bullet))
\longrightarrow H_p(K_\bullet\otimes L_\bullet)
\longrightarrow H_1(K_\bullet\otimes H_{p-1}(L_\bullet))
\longrightarrow0
\]
pour tout indice $p$.
\end{lemma}

\oldpage[III]{84}
C'est un cas particulier d'une suite exacte de termes de bas degré de la
suite spectrale de Künneth (M, XVII, 5.2 a) et G, I, 5.5.2)~; on peut le
démontrer directement de la façon suivante. Considérons $K_0$ et $K_1$ comme
des complexes de chaînes (nuls en dimensions $\neq0$ et $\neq1$
respectivement)~; on a alors une suite exacte de complexes
\[
0\longrightarrow K_0\otimes L_\bullet
\longrightarrow K_\bullet\otimes L_\bullet
\longrightarrow K_1\otimes L_\bullet\longrightarrow0
\]
à laquelle nous pouvons appliquer la suite exacte d'homologie
\[
\cdots\longrightarrow H_{p+1}(K_1\otimes L_\bullet)
\xrightarrow{\partial}H_p(K_0\otimes L_\bullet)
\longrightarrow H_p(K_\bullet\otimes L_\bullet)
\longrightarrow H_p(K_1\otimes L_\bullet)
\xrightarrow{\partial}H_{p-1}(K_0\otimes L_\bullet)
\longrightarrow\cdots.
\]
Mais il est évident que
$H_p(K_0\otimes L_\bullet)=K_0\otimes H_p(L_\bullet)$ et
$H_p(K_1\otimes L_\bullet)=K_1\otimes H_{p-1}(L_\bullet)$ pour tout $p$~; en
outre, on vérifie aussitôt que l'opérateur
$\partial:K_1\otimes H_p(L_\bullet)\to K_0\otimes H_p(L_\bullet)$ n'est autre
que $d_1\otimes1$~; le lemme résulte donc de la suite exacte précédente et de
la définition de $H_0(K_\bullet\otimes H_p(L_\bullet))$ et de
$H_1(K_\bullet\otimes H_{p-1}(L_\bullet))$.

Ce lemme étant établi, la fin de la démonstration de
(\hyperref[III.1.1.4-fr]{1.1.4}) est immédiate~: l'hypothèse de récurrence et
ce lemme (\hyperref[III.1.1.4.1-fr]{1.1.4.1}) donnent
$H_p(K_\bullet\otimes L_\bullet)=0$ pour $p\geq2$~; en outre, si on prouve que
$H_1(K_\bullet,H_0(L_\bullet))=0$, on déduira aussi
$H_1(K_\bullet\otimes L_\bullet)=0$ du lemme
(\hyperref[III.1.1.4.1-fr]{1.1.4.1})~; mais par définition,
$H_1(K_\bullet,H_0(L_\bullet))$ n'est autre que le noyau de l'homothétie
$z\mapsto f_rz$ dans $M_{r-1}$, et comme par hypothèse ce noyau est nul, cela
achève la démonstration.
\end{proof}

\begin{env}[1.1.5]
\label{III.1.1.5-fr}
Soit $\mathbf g=(g_i)_{1\leq i\leq r}$ une seconde suite de $r$ éléments de
$A$, et posons $\mathbf f\mathbf g=(f_ig_i)_{1\leq i\leq r}$. On peut définir
un homomorphisme canonique de complexes
\begin{equation}
\label{III.1.1.5.1-fr}
\tag{1.1.5.1}
\varphi_{\mathbf g}:K_\bullet(\mathbf f\mathbf g)\longrightarrow
K_\bullet(\mathbf f)
\end{equation}
comme l'extension canonique à l'algèbre extérieure $\bigwedge(A^r)$ de
l'application $A$-linéaire
$(x_1,\ldots,x_r)\mapsto(g_1x_1,\ldots,g_rx_r)$ de $A^r$ dans lui-même. Pour
voir qu'on a bien un homomorphisme de complexes, il suffit de remarquer, de
façon générale, que si $u:E\to F$ est une application $A$-linéaire,
$\mathbf x\in F^\vee$ et $\mathbf y={}^tu(\mathbf x)\in E^\vee$, alors on a
la formule
\begin{equation}
\label{III.1.1.5.2-fr}
\tag{1.1.5.2}
(\bigwedge u)\circ i_{\mathbf y}=i_{\mathbf x}\circ(\bigwedge u)~;
\end{equation}
en effet, les deux membres sont des antidérivations de $\bigwedge F$, et il
suffit de vérifier qu'ils coïncident dans $F$, ce qui résulte aussitôt des
définitions.

Lorsqu'on identifie $K_\bullet(\mathbf f)$ au produit tensoriel des
$K_\bullet(f_i)$ (\hyperref[III.1.1.1-fr]{1.1.1}), $\varphi_{\mathbf g}$ est
le produit tensoriel des $\varphi_{g_i}$, où $\varphi_{g_i}$ se réduit à
l'identité en degré $0$ et à la multiplication par $g_i$ en degré $1$.
\end{env}

\begin{env}[1.1.6]
\label{III.1.1.6-fr}
En particulier, pour tout couple d'entiers $m,n$ tels que $0\leq n\leq m$,
on a des homomorphismes de complexes
\begin{equation}
\label{III.1.1.6.1-fr}
\tag{1.1.6.1}
\varphi_{\mathbf f^{m-n}}:K_\bullet(\mathbf f^m)\longrightarrow
K_\bullet(\mathbf f^n)
\end{equation}
et par suite des homomorphismes
\begin{align}
\label{III.1.1.6.2-fr}
\tag{1.1.6.2}
\varphi_{\mathbf f^{m-n}} &:K^\bullet(\mathbf f^n,M)\longrightarrow
K^\bullet(\mathbf f^m,M),\\
\label{III.1.1.6.3-fr}
\tag{1.1.6.3}
\varphi_{\mathbf f^{m-n}} &:H^\bullet(\mathbf f^n,M)\longrightarrow
H^\bullet(\mathbf f^m,M).
\end{align}

\oldpage[III]{85}
Ces derniers vérifient évidemment la condition de transitivité
$\varphi_{\mathbf f^{m-p}}=
\varphi_{\mathbf f^{m-n}}\circ\varphi_{\mathbf f^{n-p}}$ pour $p\leq n\leq m$~;
ils définissent donc deux \emph{systèmes inductifs} de $A$-modules~; on posera
\begin{align}
\label{III.1.1.6.4-fr}
\tag{1.1.6.4}
C^\bullet((\mathbf f),M)&=\varinjlim_n K^\bullet(\mathbf f^n,M),\\
\label{III.1.1.6.5-fr}
\tag{1.1.6.5}
H^\bullet((\mathbf f),M)&=H^\bullet(C^\bullet((\mathbf f),M))
=\varinjlim_n H^\bullet(\mathbf f^n,M),
\end{align}
la dernière égalité provenant du fait que le passage à la limite inductive
permute au foncteur $H^\bullet$ (G, I, 2.1). On verra ultérieurement
(\hyperref[III.1.4.3-fr]{1.4.3}) que $H^\bullet((\mathbf f),M)$ ne dépend en
fait que de l'\emph{idéal $(\mathbf f)$} dans $A$ (et même de la topologie
$(\mathbf f)$-préadique sur $A$), ce qui justifie les notations.

Il est clair que $M\mapsto C^\bullet((\mathbf f),M)$ est un foncteur
$A$-linéaire exact, et $M\mapsto H^\bullet((\mathbf f),M)$ un foncteur
cohomologique.
\end{env}

\begin{env}[1.1.7]
\label{III.1.1.7-fr}
Soient $\mathbf f=(f_i)\in A^r$, $\mathbf g=(g_i)\in A^r$~; désignons par
$e_{\mathbf g}$ la multiplication à gauche par le vecteur
$\mathbf g\in A^r$ dans l'algèbre extérieure $\bigwedge(A^r)$~; on sait que
l'on a la \emph{formule d'homotopie}
\begin{equation}
\label{III.1.1.7.1-fr}
\tag{1.1.7.1}
i_{\mathbf f}e_{\mathbf g}+e_{\mathbf g}i_{\mathbf f}
=\langle\mathbf g,\mathbf f\rangle 1
\end{equation}
dans le $A$-module $A^r$ ($1$ désignant l'automorphisme identique de $A^r$)~;
cette relation signifie aussi que, \emph{dans le complexe
$K_\bullet(\mathbf f)$}, on a
\begin{equation}
\label{III.1.1.7.2-fr}
\tag{1.1.7.2}
de_{\mathbf g}+e_{\mathbf g}d=\langle\mathbf g,\mathbf f\rangle 1.
\end{equation}
Si l'idéal $(\mathbf f)$ est égal à $A$, il existe $\mathbf g\in A^r$ tel que
$\langle\mathbf g,\mathbf f\rangle=\sum_{i=1}^r g_if_i=1$. Par suite
(G, I, 2.4)~:
\end{env}

\begin{proposition}[1.1.8]
\label{III.1.1.8-fr}
Supposons que l'idéal $(\mathbf f)$ engendré par les $f_i$ soit égal à $A$.
Alors le complexe $K_\bullet(\mathbf f)$ est homotopiquement trivial, et il en
est donc de même des complexes $K_\bullet(\mathbf f,M)$ et
$K^\bullet(\mathbf f,M)$ pour tout $A$-module $M$.
\end{proposition}

\begin{corollary}[1.1.9]
\label{III.1.1.9-fr}
Si $(\mathbf f)=A$, on a $H^\bullet(\mathbf f,M)=0$ et
$H^\bullet((\mathbf f),M)=0$ pour tout $A$-module $M$.
\end{corollary}

\begin{proof}
En effet, on a alors $(\mathbf f^n)=A$ pour tout $n$.
\end{proof}

\begin{remark}[1.1.10]
\label{III.1.1.10-fr}
Avec les mêmes notations que ci-dessus, soient $X=\operatorname{Spec}(A)$,
$Y$ le sous-schéma fermé de $X$ défini par l'idéal $(\mathbf f)$. Nous
prouverons au \S~9 que $H^\bullet((\mathbf f),M)$ est isomorphe à la
cohomologie $H_Y^\bullet(X,\widetilde M)$ correspondant à l'antifiltre $\Phi$
des parties fermées de $Y$ (T, 3.2). Nous montrerons aussi que la prop.
(\hyperref[III.1.2.3-fr]{1.2.3}) appliquée à $X$ et à
$\mathcal F=\widetilde M$, est un cas particulier d'une suite exacte de
cohomologie
\[
\cdots\longrightarrow H_Y^p(X,\mathcal F)\longrightarrow H^p(X,\mathcal F)
\longrightarrow H^p(X-Y,\mathcal F)\longrightarrow H_Y^{p+1}(X,\mathcal F)
\longrightarrow\cdots.
\]
\end{remark}

\subsection{Cohomologie de Čech d'un recouvrement ouvert}
\label{subsection:III.1.2-fr}

\begin{notation}[1.2.1]
\label{III.1.2.1-fr}
Dans ce numéro, on notera~:
\begin{itemize}
\item $X$ un préschéma~;
\item $\mathcal F$ un $\mathcal O_X$-Module quasi-cohérent~;
\oldpage[III]{86}
\item $A=\Gamma(X,\mathcal O_X)$, $M=\Gamma(X,\mathcal F)$~;
\item $\mathbf f=(f_i)_{1\leq i\leq r}$ un système fini d'éléments de $A$~;
\item $U_i=X_{f_i}$ l'ensemble ouvert
(\hyperref[0.5.5.2-fr]{0\textsubscript{I}, 5.5.2}) des $x\in X$ tels que
$f_i(x)\neq0$~;
\item $U=\bigcup_{i=1}^r U_i$~;
\item $\mathfrak U$ le recouvrement $(U_i)_{1\leq i\leq r}$ de $U$.
\end{itemize}
\end{notation}

\begin{env}[1.2.2]
\label{III.1.2.2-fr}
Supposons que $X$ soit, ou bien un préschéma dont l'espace sous-jacent est
\emph{noethérien}, ou bien un \emph{schéma} dont l'espace sous-jacent soit
\emph{quasi-compact}. On sait alors (\hyperref[I.9.3.3-fr]{I, 9.3.3}) que
l'on a $\Gamma(U_i,\mathcal F)=M_{f_i}$. Nous poserons
\[
U_{i_0i_1\ldots i_p}=\bigcap_{k=0}^p U_{i_k}
=X_{f_{i_0}f_{i_1}\cdots f_{i_p}}
\]
(\hyperref[0.5.5.3-fr]{0\textsubscript{I}, 5.5.3})~; on a donc aussi
\begin{equation}
\label{III.1.2.2.1-fr}
\tag{1.2.2.1}
\Gamma(U_{i_0i_1\ldots i_p},\mathcal F)=M_{f_{i_0}f_{i_1}\cdots f_{i_p}}.
\end{equation}
Or (\hyperref[0.1.6.1-fr]{0\textsubscript{I}, 1.6.1})
$M_{f_{i_0}f_{i_1}\cdots f_{i_p}}$ s'identifie à la limite inductive
$\varinjlim_n M_{i_0i_1\ldots i_p}^{(n)}$, où le système inductif est formé
des $M_{i_0i_1\ldots i_p}^{(n)}=M$, l'homomorphisme
$\varphi_{nm}:M_{i_0i_1\ldots i_p}^{(m)}\to
M_{i_0i_1\ldots i_p}^{(n)}$ étant la multiplication par
$(f_{i_0}f_{i_1}\cdots f_{i_p})^{n-m}$ pour $m\leq n$. Désignons par
$C_n^p(M)$ l'ensemble des applications \emph{alternées} de $[1,r]^{p+1}$
dans $M$ (pour tout $n$)~; ces $A$-modules forment encore un système inductif
pour les $\varphi_{nm}$. Si $C^p(\mathfrak U,\mathcal F)$ est le groupe des
$p$-cochaînes \emph{alternées} de Čech relatif au recouvrement $\mathfrak U$,
à coefficients dans $\mathcal F$ (G, II, 5.1), il résulte de ce qui précède
que l'on peut écrire
\begin{equation}
\label{III.1.2.2.2-fr}
\tag{1.2.2.2}
C^p(\mathfrak U,\mathcal F)=\varinjlim_n C_n^p(M).
\end{equation}
Or, avec les notations de (\hyperref[III.1.1.2-fr]{1.1.2}), $C_n^p(M)$
s'identifie à $K^{p+1}(\mathbf f^n,M)$, et l'application $\varphi_{nm}$
s'identifie à l'application $\varphi_{\mathbf f^{n-m}}$ définie dans
(\hyperref[III.1.1.6-fr]{1.1.6}). On a donc, pour tout $p\geq0$, un
isomorphisme canonique fonctoriel en $\mathcal F$
\begin{equation}
\label{III.1.2.2.3-fr}
\tag{1.2.2.3}
C^p(\mathfrak U,\mathcal F)\xrightarrow{\sim}C^{p+1}((\mathbf f),M).
\end{equation}
En outre, la formule (\hyperref[III.1.1.2.3-fr]{1.1.2.3}) et la définition de
la cohomologie d'un recouvrement (G, II, 5.1) montrent que les isomorphismes
(\hyperref[III.1.2.2.3-fr]{1.2.2.3}) sont compatibles avec les opérateurs
cobords.
\end{env}

\begin{proposition}[1.2.3]
\label{III.1.2.3-fr}
Si $X$ est un préschéma dont l'espace sous-jacent est noethérien, ou un
schéma dont l'espace sous-jacent est quasi-compact, il existe un isomorphisme
canonique fonctoriel en $\mathcal F$
\begin{equation}
\label{III.1.2.3.1-fr}
\tag{1.2.3.1}
H^p(\mathfrak U,\mathcal F)\xrightarrow{\sim}H^{p+1}((\mathbf f),M)
\qquad\text{pour tout }p\geq1.
\end{equation}
On a en outre une suite exacte fonctorielle en $\mathcal F$
\begin{equation}
\label{III.1.2.3.2-fr}
\tag{1.2.3.2}
0\longrightarrow H^0((\mathbf f),M)\longrightarrow M
\longrightarrow H^0(\mathfrak U,\mathcal F)
\longrightarrow H^1((\mathbf f),M)\longrightarrow0.
\end{equation}
\end{proposition}

\oldpage[III]{87}
\begin{proof}
Les relations (\hyperref[III.1.2.3.1-fr]{1.2.3.1}) sont en effet conséquences
immédiates de ce qu'on a vu dans (\hyperref[III.1.2.2-fr]{1.2.2}). D'autre
part, on a
$C^0(\mathfrak U,\mathcal F)=C^1((\mathbf f),M)$~;
$H^0(\mathfrak U,\mathcal F)$ s'identifie par suite au sous-groupe des
$1$-cocycles de $C^1((\mathbf f),M)$~; comme
$M=C^0((\mathbf f),M)$, la suite exacte
(\hyperref[III.1.2.3.2-fr]{1.2.3.2}) n'est autre que celle qui résulte de la
définition des groupes de cohomologie $H^0((\mathbf f),M)$ et
$H^1((\mathbf f),M)$.
\end{proof}

\begin{corollary}[1.2.4]
\label{III.1.2.4-fr}
Supposons que les $X_{f_i}$ soient quasi-compacts et qu'il existe des
$g_i\in\Gamma(U,\mathcal F)$ tels que
$\sum_i g_i(f_i\vert U)=1\vert U$. Alors, pour tout
$(\mathcal O_X\vert U)$-Module quasi-cohérent $\mathcal G$, on a
$H^p(\mathfrak U,\mathcal G)=0$ pour $p>0$~; si en outre $U=X$,
l'homomorphisme canonique (\hyperref[III.1.2.3.2-fr]{1.2.3.2})
$M\to H^0(\mathfrak U,\mathcal F)$ est bijectif.
\end{corollary}

\begin{proof}
Comme par hypothèse les $U_i=X_{f_i}$ sont quasi-compacts, il en est de même
de $U$, et on peut donc se borner au cas où $U=X$~; l'hypothèse entraîne alors
$H^p((\mathbf f),M)=0$ pour tout $p\geq0$
(\hyperref[III.1.1.9-fr]{1.1.9}). Le corollaire résulte alors aussitôt de
(\hyperref[III.1.2.3.1-fr]{1.2.3.1}) et
(\hyperref[III.1.2.3.2-fr]{1.2.3.2}).
\end{proof}

On observera que puisque
$H^0(\mathfrak U,\mathcal F)=H^0(U,\mathcal F)$ (G, II, 5.2.2), on a ainsi
démontré à nouveau (I, 1.3.7) comme cas particulier.

\begin{remark}[1.2.5]
\label{III.1.2.5-fr}
Supposons que $X$ soit un \emph{schéma affine}~; alors les
$U_i=X_{f_i}=D(f_i)$ sont des ouverts affines, ainsi que les
$U_{i_0i_1\ldots i_p}$ (mais $U$ n'est pas nécessairement affine). Dans ce
cas, les foncteurs $\Gamma(X,\mathcal F)$ et
$\Gamma(U_{i_0i_1\ldots i_p},\mathcal F)$ sont exacts en $\mathcal F$
(I, 1.3.11). Si on a une suite exacte
$0\to\mathcal F'\to\mathcal F\to\mathcal F''\to0$ de
$\mathcal O_X$-Modules quasi-cohérents, la suite des complexes
\[
0\longrightarrow C^\bullet(\mathfrak U,\mathcal F')
\longrightarrow C^\bullet(\mathfrak U,\mathcal F)
\longrightarrow C^\bullet(\mathfrak U,\mathcal F'')
\longrightarrow0
\]
est exacte, et donne donc une suite exacte de cohomologie
\[
\cdots\longrightarrow H^p(\mathfrak U,\mathcal F')
\longrightarrow H^p(\mathfrak U,\mathcal F)
\longrightarrow H^p(\mathfrak U,\mathcal F'')
\xrightarrow{\partial}H^{p+1}(\mathfrak U,\mathcal F')
\longrightarrow\cdots.
\]
D'autre part, si on pose $M'=\Gamma(X,\mathcal F')$,
$M''=\Gamma(X,\mathcal F'')$, la suite
$0\to M'\to M\to M''\to0$ est exacte~; comme
$C^\bullet((\mathbf f),M)$ est un foncteur exact en $M$, on a aussi la suite
exacte de cohomologie
\[
\cdots\longrightarrow H^p((\mathbf f),M')
\longrightarrow H^p((\mathbf f),M)
\longrightarrow H^p((\mathbf f),M'')
\xrightarrow{\partial}H^{p+1}((\mathbf f),M')
\longrightarrow\cdots.
\]
Cela étant, comme le diagramme
\[
\xymatrix{
0\ar[r]&C^\bullet(\mathfrak U,\mathcal F')\ar[r]\ar[d]
&C^\bullet(\mathfrak U,\mathcal F)\ar[r]\ar[d]
&C^\bullet(\mathfrak U,\mathcal F'')\ar[r]\ar[d]&0\\
0\ar[r]&C^\bullet((\mathbf f),M')\ar[r]
&C^\bullet((\mathbf f),M)\ar[r]
&C^\bullet((\mathbf f),M'')\ar[r]&0
}
\]
\oldpage[III]{88}
est commutatif, on en conclut que les diagrammes
\begin{equation}
\label{III.1.2.5.1-fr}
\tag{1.2.5.1}
\xymatrix{
H^p(\mathfrak U,\mathcal F'')\ar[r]^{\partial}\ar[d]
&H^{p+1}(\mathfrak U,\mathcal F')\ar[d]\\
H^{p+1}((\mathbf f),M'')\ar[r]^{\partial}
&H^{p+2}((\mathbf f),M')
}
\end{equation}
sont commutatifs pour tout $p$ (G, I, 2.1.1).
\end{remark}

\subsection{Cohomologie d'un schéma affine}
\label{subsection:III.1.3-fr}

\begin{theorem}[1.3.1]
\label{III.1.3.1-fr}
Soit $X$ un schéma affine. Pour tout $\mathcal O_X$-Module quasi-cohérent
$\mathcal F$, on a $H^p(X,\mathcal F)=0$ pour tout $p>0$.
\end{theorem}

\begin{proof}
Soit $\mathfrak U$ un recouvrement fini de $X$ par des ouverts affines
$X_{f_i}=D(f_i)$ ($1\leq i\leq r$)~; on sait qu'alors l'idéal engendré dans
$A=\Gamma(X,\mathcal O_X)$ par les $f_i$ est égal à $A$. On conclut donc de
(\hyperref[III.1.2.4-fr]{1.2.4}) que l'on a
$H^p(\mathfrak U,\mathcal F)=0$ pour $p>0$. Comme il y a des recouvrements
finis de $X$ par des ouverts affines qui sont arbitrairement fins (I, 1.1.10),
la définition de la cohomologie de Čech (G, II, 5.8) montre que l'on a aussi
$\check H^p(X,\mathcal F)=0$ pour $p>0$. Mais ceci s'applique aussi à tout
préschéma $X_f$ pour $f\in A$ (I, 1.3.6), donc
$\check H^p(X_f,\mathcal F)=0$ pour $p>0$. Comme l'on a
$X_f\cap X_g=X_{fg}$, on en déduit que l'on a aussi
$H^p(X,\mathcal F)=0$ pour tout $p>0$, en vertu de (G, II, 5.9.2).
\end{proof}

\begin{corollary}[1.3.2]
\label{III.1.3.2-fr}
Soient $Y$ un préschéma, $f:X\to Y$ un morphisme affine (II, 1.6.1). Pour
tout $\mathcal O_X$-Module quasi-cohérent $\mathcal F$, on a
$R^qf_*(\mathcal F)=0$ pour $q>0$.
\end{corollary}

\begin{proof}
En effet, par définition $R^qf_*(\mathcal F)$ est le $\mathcal O_Y$-Module
associé au préfaisceau $U\mapsto H^q(f^{-1}(U),\mathcal F)$, $U$ parcourant
les ouverts de $Y$. Mais les ouverts affines forment une base de $Y$, et pour
un tel ouvert $U$, $f^{-1}(U)$ est affine (II, 1.3.2), donc
$H^q(f^{-1}(U),\mathcal F)=0$ par
(\hyperref[III.1.3.1-fr]{1.3.1}), ce qui démontre le corollaire.
\end{proof}

\begin{corollary}[1.3.3]
\label{III.1.3.3-fr}
Soient $Y$ un préschéma, $f:X\to Y$ un morphisme affine. Pour tout
$\mathcal O_X$-Module quasi-cohérent $\mathcal F$, l'homomorphisme canonique
$H^p(Y,f_*(\mathcal F))\to H^p(X,\mathcal F)$ (0, 12.1.3.1) est bijectif
pour tout $p$.
\end{corollary}

\begin{proof}
En effet, il suffit (en vertu de (0, 12.1.7)) de montrer que les
edge-homomorphismes
$''E_2^{p0}=H^p(Y,f_*(\mathcal F))\to H^p(X,\mathcal F)$ de la seconde suite
spectrale du foncteur composé $\Gamma f_*$ sont bijectifs. Mais le terme
$E_2$ de cette suite est donné par
$''E_2^{pq}=H^p(Y,R^qf_*(\mathcal F))$ (G, II, 4.17.1), donc il résulte de
(\hyperref[III.1.3.2-fr]{1.3.2}) que $''E_2^{pq}=0$ pour $q>0$, et la suite
spectrale dégénère~; d'où notre assertion (0, 11.1.6).
\end{proof}

\begin{corollary}[1.3.4]
\label{III.1.3.4-fr}
Soient $f:X\to Y$ un morphisme affine, $g:Y\to Z$ un morphisme.
\oldpage[III]{89}
Pour tout $\mathcal O_X$-Module quasi-cohérent $\mathcal F$,
l'homomorphisme canonique
$R^pg_*(f_*(\mathcal F))\to R^p(g\circ f)_*(\mathcal F)$ (0, 12.2.5.1)
est bijectif pour tout $p$.
\end{corollary}

\begin{proof}
Il suffit de remarquer que, d'après
(\hyperref[III.1.3.3-fr]{1.3.3}), pour tout ouvert affine $W$ de $Z$,
l'homomorphisme canonique
$H^p(g^{-1}(W),f_*(\mathcal F))\to
H^p(f^{-1}(g^{-1}(W)),\mathcal F)$ est bijectif~; cela prouve que
l'homomorphisme de préfaisceaux définissant l'homomorphisme canonique
$R^pg_*(f_*(\mathcal F))\to R^p(g\circ f)_*(\mathcal F)$ est bijectif
(0, 12.2.5).
\end{proof}

\subsection{Application à la cohomologie des préschémas quelconques}
\label{subsection:III.1.4-fr}

\begin{proposition}[1.4.1]
\label{III.1.4.1-fr}
Soient $X$ un schéma, $\mathfrak U=(U_\alpha)$ un recouvrement de $X$ par
des ouverts affines. Pour tout $\mathcal O_X$-Module quasi-cohérent
$\mathcal F$, les modules de cohomologie $H^\bullet(X,\mathcal F)$ et
$H^\bullet(\mathfrak U,\mathcal F)$ (sur $\Gamma(X,\mathcal O_X)$) sont
canoniquement isomorphes.
\end{proposition}

\begin{proof}
En effet, comme $X$ est un schéma, toute intersection finie $V$ d'ouverts du
recouvrement $\mathfrak U$ est affine (I, 5.5.6), donc
$H^q(V,\mathcal F)=0$ pour $q\geq1$ en vertu de
(\hyperref[III.1.3.1-fr]{1.3.1}). La proposition résulte alors du th. de
Leray (G, II, 5.4.1).
\end{proof}

\begin{remark}[1.4.2]
\label{III.1.4.2-fr}
On notera que la conclusion de (\hyperref[III.1.4.1-fr]{1.4.1}) est encore
valable lorsque les intersections finies des ensembles $U_\alpha$ sont
affines, même lorsque l'on ne suppose pas nécessairement que $X$ soit un
schéma.
\end{remark}

\begin{corollary}[1.4.3]
\label{III.1.4.3-fr}
Soient $X$ un schéma dont l'espace sous-jacent est quasi-compact,
$A=\Gamma(X,\mathcal O_X)$, $\mathbf f=(f_i)_{1\leq i\leq r}$ une suite
finie d'éléments de $A$ tels que les $X_{f_i}$ (notations de
(\hyperref[III.1.2.1-fr]{1.2.1})) soient affines. Alors (avec les notations
de (\hyperref[III.1.2.1-fr]{1.2.1})), pour tout $\mathcal O_X$-Module
quasi-cohérent $\mathcal F$, on a un isomorphisme canonique fonctoriel en
$\mathcal F$
\begin{equation}
\label{III.1.4.3.1-fr}
\tag{1.4.3.1}
H^q(U,\mathcal F)\xrightarrow{\sim}H^{q+1}((\mathbf f),M)
\qquad\text{pour }q\geq1
\end{equation}
et une suite exacte fonctorielle en $\mathcal F$
\begin{equation}
\label{III.1.4.3.2-fr}
\tag{1.4.3.2}
0\longrightarrow H^0((\mathbf f),M)\longrightarrow M
\longrightarrow H^0(U,\mathcal F)
\longrightarrow H^1((\mathbf f),M)\longrightarrow0.
\end{equation}
\end{corollary}

\begin{proof}
Cela résulte en effet aussitôt de
(\hyperref[III.1.4.1-fr]{1.4.1}) et
(\hyperref[III.1.2.3-fr]{1.2.3}).
\end{proof}

\begin{env}[1.4.4]
\label{III.1.4.4-fr}
Si $X$ est un \emph{schéma affine}, il résulte de
(\hyperref[III.1.2.5-fr]{1.2.5}) et
(\hyperref[III.1.4.1-fr]{1.4.1}) que pour tout $q\geq0$, les diagrammes
\begin{equation}
\label{III.1.4.4.1-fr}
\tag{1.4.4.1}
\xymatrix{
H^q(U,\mathcal F'')\ar[r]^{\partial}\ar[d]
&H^{q+1}(U,\mathcal F')\ar[d]\\
H^{q+1}((\mathbf f),M'')\ar[r]^{\partial}
&H^{q+2}((\mathbf f),M')
}
\end{equation}
correspondant à une suite exacte
$0\to\mathcal F'\to\mathcal F\to\mathcal F''\to0$ de
$\mathcal O_X$-Modules quasi-cohérents (avec les notations de
(\hyperref[III.1.2.5-fr]{1.2.5})), sont commutatifs.
\end{env}

\oldpage[III]{90}
\begin{proposition}[1.4.5]
\label{III.1.4.5-fr}
Soient $X$ un schéma quasi-compact, $\mathcal L$ un $\mathcal O_X$-Module
inversible, et considérons l'anneau gradué
$A_\bullet=\Gamma_\bullet(\mathcal L)$ (0\textsubscript{I}, 5.4.6)~; alors
\[
H^\bullet(\mathcal F,\mathcal L)
=\bigoplus_{n\in\mathbb Z}H^\bullet(X,\mathcal F\otimes\mathcal L^{\otimes n})
\]
est un $A_\bullet$-module gradué, et pour tout $f\in A_n$, on a un
isomorphisme canonique
\begin{equation}
\label{III.1.4.5.1-fr}
\tag{1.4.5.1}
H^\bullet(X_f,\mathcal F)
\xrightarrow{\sim}(H^\bullet(\mathcal F,\mathcal L))_{(f)}
\end{equation}
de $(A_\bullet)_{(f)}$-modules.
\end{proposition}

\begin{proof}
Comme $X$ est un schéma quasi-compact, on peut calculer la cohomologie de
tous les $\mathcal O_X$-Modules
$\mathcal F\otimes\mathcal L^{\otimes n}$ à l'aide d'un même recouvrement
fini $\mathfrak U=(U_i)$ par des ouverts affines tels que la restriction
$\mathcal L\vert U_i$ soit isomorphe à $\mathcal O_X\vert U_i$ pour chaque
$i$ (\hyperref[III.1.4.1-fr]{1.4.1}). Il est alors immédiat que les
$U_i\cap X_f$ sont des ouverts affines (I, 1.3.6), et l'on peut donc aussi
calculer la cohomologie
$H^\bullet(X_f,\mathcal F\otimes\mathcal L^{\otimes n})$ à l'aide du
recouvrement $\mathfrak U\vert X_f=(U_i\cap X_f)$
(\hyperref[III.1.4.1-fr]{1.4.1}). Il est immédiat que pour tout $f\in A_n$,
la multiplication par $f$ définit un homomorphisme
\[
C^\bullet(\mathfrak U,\mathcal F\otimes\mathcal L^{\otimes m})
\longrightarrow
C^\bullet(\mathfrak U,\mathcal F\otimes\mathcal L^{\otimes(m+n)}),
\]
d'où un homomorphisme
\[
H^\bullet(\mathfrak U,\mathcal F\otimes\mathcal L^{\otimes m})
\longrightarrow
H^\bullet(\mathfrak U,\mathcal F\otimes\mathcal L^{\otimes(m+n)}),
\]
ce qui établit la première assertion. D'autre part, pour $f\in A_n$ donné,
il résulte de (I, 9.3.2) que l'on a un isomorphisme de complexes de
$(A_\bullet)_{(f)}$-modules
\[
C^\bullet(\mathfrak U\vert X_f,\mathcal F)
\xrightarrow{\sim}
\left(C^\bullet\left(\mathfrak U,
\bigoplus_{n\in\mathbb Z}\mathcal F\otimes\mathcal L^{\otimes n}
\right)\right)_{(f)}
\]
compte tenu de (I, 1.3.9, (ii)). Passant à la cohomologie de ces deux
complexes, on en déduit l'isomorphisme
(\hyperref[III.1.4.5.1-fr]{1.4.5.1}), en se souvenant que le foncteur
$M\mapsto M_{(f)}$ est exact dans la catégorie des $A_\bullet$-modules
gradués.
\end{proof}

\begin{corollary}[1.4.6]
\label{III.1.4.6-fr}
On suppose vérifiées les hypothèses de
(\hyperref[III.1.4.5-fr]{1.4.5}) et on suppose en outre que
$\mathcal L=\mathcal O_X$. Si on pose $A=\Gamma(X,\mathcal O_X)$, alors,
pour tout $f\in A$, on a un isomorphisme canonique
\[
H^\bullet(X_f,\mathcal F)
\xrightarrow{\sim}(H^\bullet(X,\mathcal F))_f
\]
de $A_f$-modules.
\end{corollary}

\begin{corollary}[1.4.7]
\label{III.1.4.7-fr}
Soient $X$ un schéma quasi-compact, $f$ un élément de
$\Gamma(X,\mathcal O_X)$.
\begin{enumerate}
\item[(i)] Supposons l'ouvert $X_f$ affine. Alors, pour tout
$\mathcal O_X$-Module quasi-cohérent $\mathcal F$, tout $i>0$ et tout
$\xi\in H^i(X,\mathcal F)$, il existe un entier $n>0$ tel que $f^n\xi=0$.
\item[(ii)] Inversement, supposons que $X_f$ soit quasi-compact et que pour
tout faisceau quasi-cohérent $\mathcal J$ d'idéaux de $\mathcal O_X$ et tout
$\zeta\in H^1(X,\mathcal J)$, il existe $n>0$ tel que $f^n\zeta=0$. Alors
$X_f$ est affine.
\end{enumerate}
\end{corollary}

\begin{proof}
\begin{enumerate}
\item[(i)] Si $X_f$ est affine, on a $H^i(X_f,\mathcal F)=0$ pour tout
$i>0$ (\hyperref[III.1.3.1-fr]{1.3.1}), donc l'assertion résulte directement
de (\hyperref[III.1.4.6-fr]{1.4.6}).
\item[(ii)] En vertu du critère de Serre (II, 5.2.1), il suffit de prouver
que pour tout Idéal quasi-cohérent $\mathcal K$ de
$\mathcal O_X\vert X_f$, on a $H^1(X_f,\mathcal K)=0$. Comme $X_f$ est un
ouvert quasi-compact dans un schéma quasi-compact $X$, il existe un Idéal
quasi-cohérent $\mathcal J$ de $\mathcal O_X$ tel que
$\mathcal K=\mathcal J\vert X_f$ (I, 9.4.2). En vertu de
(\hyperref[III.1.4.6-fr]{1.4.6}), on a
$H^1(X_f,\mathcal K)=(H^1(X,\mathcal J))_f$, et l'hypothèse entraîne que le
second membre est nul, d'où la conclusion.
\end{enumerate}
\end{proof}

\begin{remark}[1.4.8]
\label{III.1.4.8-fr}
On notera que (\hyperref[III.1.4.7-fr]{1.4.7}, (i)) redonne une
démonstration plus simple de la relation (II, 4.5.13.2).
\end{remark}

\oldpage[III]{91}
\begin{lemma}[1.4.9]
\label{III.1.4.9-fr}
Soient $X$ un schéma quasi-compact,
$\mathfrak U=(U_i)_{1\leq i\leq n}$ un recouvrement fini de $X$ par des
ouverts affines, $\mathcal F$ un $\mathcal O_X$-Module quasi-cohérent. Le
complexe de faisceaux $\mathcal C^\bullet(\mathfrak U,\mathcal F)$ défini
par le recouvrement $\mathfrak U$ (G, II, 5.2) est alors un
$\mathcal O_X$-Module quasi-cohérent.
\end{lemma}

\begin{proof}
Il résulte des définitions (G, II, 5.2) que
$\mathcal C^p(\mathfrak U,\mathcal F)$ est somme directe des faisceaux
images directes des $\mathcal F\vert U_{i_0\ldots i_p}$ par l'injection
canonique $U_{i_0\ldots i_p}\to X$. Or, l'hypothèse que $X$ est un schéma
entraîne que ces injections sont des morphismes affines (I, 5.5.6), donc les
$\mathcal C^p(\mathfrak U,\mathcal F)$ sont quasi-cohérents (II, 1.2.6).
\end{proof}

\begin{proposition}[1.4.10]
\label{III.1.4.10-fr}
Soit $u:X\to Y$ un morphisme séparé et quasi-compact. Pour tout
$\mathcal O_X$-Module quasi-cohérent $\mathcal F$, les
$R^qu_*(\mathcal F)$ sont des $\mathcal O_Y$-Modules quasi-cohérents.
\end{proposition}

\begin{proof}
La question étant locale sur $Y$, on peut supposer $Y$ affine. Alors $X$ est
réunion finie d'ouverts affines $U_i$ ($1\leq i\leq n$)~; soit
$\mathfrak U$ le recouvrement $(U_i)$. En outre, comme $Y$ est un schéma,
il résulte de (I, 5.5.10) que pour tout ouvert affine $V\subset Y$,
l'injection canonique $u^{-1}(V)\to X$ est un morphisme affine~; on en
conclut ((\hyperref[III.1.4.1-fr]{1.4.1}) et (G, II, 5.2)) que l'on a un
isomorphisme canonique
\begin{equation}
\label{III.1.4.10.1-fr}
\tag{1.4.10.1}
H^\bullet(u^{-1}(V),\mathcal F)
\xrightarrow{\sim}H^\bullet(\Gamma(V,\mathcal K^\bullet))
\end{equation}
où l'on a posé
$\mathcal K^\bullet=u_*(\mathcal C^\bullet(\mathfrak U,\mathcal F))$. En
vertu de (\hyperref[III.1.4.9-fr]{1.4.9}) et (I, 9.2.2),
$\mathcal K^\bullet$ est un $\mathcal O_Y$-Module quasi-cohérent~; par
ailleurs, il constitue un \emph{complexe de faisceaux} puisqu'il en est
ainsi de $\mathcal C^\bullet(\mathfrak U,\mathcal F)$. Il résulte alors de la
définition de la cohomologie
$\mathcal H^\bullet(\mathcal K^\bullet)$ (G, II, 4.1) que cette dernière est
formée de $\mathcal O_Y$-Modules quasi-cohérents (I, 4.1.1). Comme (pour $V$
affine dans $Y$) le foncteur $\Gamma(V,\mathcal G)$ est exact en
$\mathcal G$ dans la catégorie des $\mathcal O_Y$-Modules quasi-cohérents,
on a (G, II, 4.1)
\begin{equation}
\label{III.1.4.10.2-fr}
\tag{1.4.10.2}
H^\bullet(\Gamma(V,\mathcal K^\bullet))
=\Gamma(V,\mathcal H^\bullet(\mathcal K^\bullet)).
\end{equation}
Notons enfin qu'il résulte de la définition de l'homomorphisme canonique
\[
H^\bullet(\mathfrak U,\mathcal F)\longrightarrow H^\bullet(X,\mathcal F)
\]
donnée dans (G, II, 5.2), que si $V'\subset V$ est un second ouvert affine
de $Y$, le diagramme
\[
\xymatrix{
H^\bullet(u^{-1}(V),\mathcal F)\ar[r]^{\sim}\ar[d]
&H^\bullet(\Gamma(V,\mathcal K^\bullet))\ar[d]\\
H^\bullet(u^{-1}(V'),\mathcal F)\ar[r]^{\sim}
&H^\bullet(\Gamma(V',\mathcal K^\bullet))
}
\]
est commutatif. On conclut donc de ce qui précède que les isomorphismes
(\hyperref[III.1.4.10.1-fr]{1.4.10.1}) définissent un isomorphisme de
$\mathcal O_Y$-Modules
\begin{equation}
\label{III.1.4.10.3-fr}
\tag{1.4.10.3}
R^\bullet u_*(\mathcal F)
\xrightarrow{\sim}\mathcal H^\bullet(\mathcal K^\bullet)
\end{equation}
\oldpage[III]{92}
et par suite $R^\bullet u_*(\mathcal F)$ est quasi-cohérent.
\end{proof}

Il résulte en outre de (\hyperref[III.1.4.10.3-fr]{1.4.10.3}),
(\hyperref[III.1.4.10.2-fr]{1.4.10.2}) et
(\hyperref[III.1.4.10.1-fr]{1.4.10.1}) que~:

\begin{corollary}[1.4.11]
\label{III.1.4.11-fr}
Sous les hypothèses de (\hyperref[III.1.4.10-fr]{1.4.10}), pour tout ouvert
affine $V$ de $Y$, l'homomorphisme canonique
\begin{equation}
\label{III.1.4.11.1-fr}
\tag{1.4.11.1}
H^q(u^{-1}(V),\mathcal F)\longrightarrow
\Gamma(V,R^qu_*(\mathcal F))
\end{equation}
est un isomorphisme pour tout $q\geq0$.
\end{corollary}

\begin{corollary}[1.4.12]
\label{III.1.4.12-fr}
Supposons vérifiées les hypothèses de
(\hyperref[III.1.4.10-fr]{1.4.10}) et supposons en outre que $Y$ soit
quasi-compact. Alors il existe un entier $r>0$ tel que pour tout
$\mathcal O_X$-Module quasi-cohérent $\mathcal F$ et tout entier $q>r$, on
ait $R^qu_*(\mathcal F)=0$. Si $Y$ est affine, on peut prendre pour $r$ un
entier tel qu'il existe un recouvrement de $X$ formé de $r$ ouverts affines.
\end{corollary}

\begin{proof}
Comme on peut recouvrir $Y$ par un nombre fini d'ouverts affines, on est
ramené à démontrer la seconde assertion, en vertu de
(\hyperref[III.1.4.11-fr]{1.4.11}). Or, si $\mathfrak U$ est un
recouvrement de $X$ par $r$ ouverts affines, on a
$H^q(\mathfrak U,\mathcal F)=0$ pour $q>r$, puisque les cochaînes de
$C^q(\mathfrak U,\mathcal F)$ sont alternées~; la conclusion résulte donc de
(\hyperref[III.1.4.1-fr]{1.4.1}).
\end{proof}

\begin{corollary}[1.4.13]
\label{III.1.4.13-fr}
Supposons vérifiées les hypothèses de
(\hyperref[III.1.4.10-fr]{1.4.10}) et en outre supposons que
$Y=\operatorname{Spec}(A)$ soit affine. Alors, pour tout
$\mathcal O_X$-Module quasi-cohérent $\mathcal F$ et tout $f\in A$, on a
\[
\Gamma(Y_f,R^qu_*(\mathcal F))
=(\Gamma(Y,R^qu_*(\mathcal F)))_f
\]
à un isomorphisme canonique près.
\end{corollary}

\begin{proof}
En effet, cela résulte de ce que $R^qu_*(\mathcal F)$ est un
$\mathcal O_Y$-Module quasi-cohérent (I, 1.3.7).
\end{proof}

\begin{proposition}[1.4.14]
\label{III.1.4.14-fr}
Soient $f:X\to Y$ un morphisme séparé quasi-compact, $g:Y\to Z$ un
morphisme affine. Pour tout $\mathcal O_X$-Module quasi-cohérent
$\mathcal F$, l'homomorphisme canonique
$R^p(g\circ f)_*(\mathcal F)\to g_*(R^pf_*(\mathcal F))$ (0, 12.2.5.2) est
bijectif pour tout $p$.
\end{proposition}

\begin{proof}
En effet, pour tout ouvert affine $W$ de $Z$, $g^{-1}(W)$ est un ouvert
affine de $Y$. L'homomorphisme de préfaisceaux définissant l'homomorphisme
canonique
\[
R^p(g\circ f)_*(\mathcal F)\longrightarrow g_*(R^pf_*(\mathcal F))
\]
(0, 12.2.5) est donc bijectif en vertu de
(\hyperref[III.1.4.11-fr]{1.4.11}).
\end{proof}

\begin{proposition}[1.4.15]
\label{III.1.4.15-fr}
Soient $u:X\to Y$ un morphisme séparé de type fini, $v:Y'\to Y$ un
morphisme plat de préschémas (0\textsubscript{I}, 6.7.1)~; soit
$u'=u_{(Y')}$, de sorte qu'on a le diagramme commutatif
\begin{equation}
\label{III.1.4.15.1-fr}
\tag{1.4.15.1}
\xymatrix{
X\ar[d]_u & X'=X_{(Y')}\ar[l]_{v'}\ar[d]^{u'}\\
Y & Y'\ar[l]_v
}
\end{equation}
Alors, pour tout $\mathcal O_X$-Module quasi-cohérent $\mathcal F$,
$R^qu_*'(\mathcal F')$, où
$\mathcal F'=v'^*(\mathcal F)=\mathcal F\otimes_{\mathcal O_Y}\mathcal O_{Y'}$,
est canoniquement isomorphe à
$R^qu_*(\mathcal F)\otimes_{\mathcal O_Y}\mathcal O_{Y'}
=v^*(R^qu_*(\mathcal F))$ pour tout $q\geq0$.
\end{proposition}

\oldpage[III]{93}
\begin{proof}
L'homomorphisme canonique
$\rho:\mathcal F\to v_*'(v'^*(\mathcal F))$ (0\textsubscript{I}, 4.4.3.2)
définit par fonctorialité un homomorphisme
\begin{equation}
\label{III.1.4.15.2-fr}
\tag{1.4.15.2}
R^qu_*(\mathcal F)\longrightarrow R^qu_*(v_*'(\mathcal F')).
\end{equation}
D'autre part, on a, en posant $w=u\circ v'=v\circ u'$, les homomorphismes
canoniques (0, 12.2.5.1 et 12.2.5.2)
\begin{equation}
\label{III.1.4.15.3-fr}
\tag{1.4.15.3}
R^qu_*(v_*'(\mathcal F'))\longrightarrow R^qw_*(\mathcal F')
\longrightarrow v_*(R^qu_*'(\mathcal F')).
\end{equation}
Composant (\hyperref[III.1.4.15.3-fr]{1.4.15.3}) et
(\hyperref[III.1.4.15.2-fr]{1.4.15.2}), on a un homomorphisme
\[
\psi:R^qu_*(\mathcal F)\longrightarrow v_*(R^qu_*'(\mathcal F'))
\]
et finalement on en déduit l'homomorphisme canonique (dont la définition ne
fait intervenir \emph{aucune hypothèse sur $v$})
\begin{equation}
\label{III.1.4.15.4-fr}
\tag{1.4.15.4}
\psi^\sharp:v^*(R^qu_*(\mathcal F))\longrightarrow R^qu_*'(\mathcal F')
\end{equation}
dont il s'agit de prouver que c'est un isomorphisme lorsque $v$ est
\emph{plat}. Il est clair que la question est locale sur $Y$ et $Y'$ et l'on
peut donc supposer que $Y=\operatorname{Spec}(A)$,
$Y'=\operatorname{Spec}(B)$~; nous utiliserons en outre le

\begin{lemma}[1.4.15.5]
\label{III.1.4.15.5-fr}
Soient $\varphi:A\to B$ un homomorphisme d'anneaux,
$Y=\operatorname{Spec}(A)$, $X=\operatorname{Spec}(B)$,
$f:X\to Y$ le morphisme correspondant à $\varphi$, $M$ un $B$-module. Pour
que le $\mathcal O_X$-Module $\widetilde M$ soit $f$-plat
(0\textsubscript{I}, 6.7.1), il faut et il suffit que $M$ soit un $A$-module
plat. En particulier, pour que le morphisme $f$ soit plat, il faut et il
suffit que $B$ soit un $A$-module plat.
\end{lemma}

En effet, cela résulte de la définition (0\textsubscript{I}, 6.7.1) et de
(0\textsubscript{I}, 6.3.3), compte tenu de (I, 1.3.4).

Cela étant, il résulte de
(\hyperref[III.1.4.11.1-fr]{1.4.11.1}) et des définitions des
homomorphismes (\hyperref[III.1.4.15.3-fr]{1.4.15.3})
(cf. 0, 12.2.5) que $\psi$ correspond alors au morphisme composé
\[
H^q(X,\mathcal F)\xrightarrow{\rho_q}
H^q(X,v_*'(v'^*(\mathcal F)))\xrightarrow{\theta_q}
H^q(X',v'^*(v_*'(v'^*(\mathcal F))))\xrightarrow{\sigma_q}
H^q(X',v'^*(\mathcal F))
\]
où $\rho_q$ et $\sigma_q$ sont les homomorphismes correspondant dans la
cohomologie aux morphismes canoniques $\rho$ et
$\sigma:v'^*(v_*'(\mathcal G'))\to\mathcal G'$, et $\theta_q$ est le
$\varphi$-morphisme (0, 12.1.3.1) relatif au $\mathcal O_X$-Module
$v_*'(v'^*(\mathcal F))$. Mais par fonctorialité de $\theta_q$, on a le
diagramme commutatif
\[
\xymatrix{
H^q(X,\mathcal F)\ar[rr]^{\rho_q}\ar[dd]_{\theta_q}
&&H^q(X,v_*'(v'^*(\mathcal F)))\ar[dd]^{\theta_q}\\\\
H^q(X',v'^*(\mathcal F))\ar[rr]^{v'^*(\rho_q)}
&&H^q(X',v'^*(v_*'(v'^*(\mathcal F))))
}
\]
\oldpage[III]{94}
et comme par définition (0\textsubscript{I}, 4.4.3), $v'^*(\rho)$ est
l'inverse de $\sigma$, on voit que le morphisme composé considéré plus haut
n'est autre finalement que $\theta_q$~; $\psi^\sharp$ est par suite le
$B$-homomorphisme associé
\[
H^q(X,\mathcal F)\otimes_A B\longrightarrow H^q(X',\mathcal F').
\]
Comme $u$ est de type fini, $X$ est réunion finie d'ouverts affines $U_i$
($1\leq i\leq r$)~; soit $\mathfrak U$ le recouvrement $(U_i)$. D'ailleurs
$v$ est un morphisme affine, donc il en est de même de $v'$ (II, 1.6.2,
(iii)), et les $U_i'=v'^{-1}(U_i)$ forment par suite un recouvrement ouvert
affine $\mathfrak U'$ de $X'$. On sait alors (0, 12.1.4.2) que le diagramme
\[
\xymatrix{
H^q(\mathfrak U,\mathcal F)\ar[r]^{\theta_q}\ar[d]
&H^q(\mathfrak U',\mathcal F')\ar[d]\\
H^q(X,\mathcal F)\ar[r]^{\theta_q}&H^q(X',\mathcal F')
}
\]
est commutatif, et les flèches verticales sont des isomorphismes puisque
$X$ et $X'$ sont des schémas (\hyperref[III.1.4.1-fr]{1.4.1}). Il suffit par
suite de prouver que le $\varphi$-morphisme canonique
$\theta_q:H^q(\mathfrak U,\mathcal F)\to H^q(\mathfrak U',\mathcal F')$
est tel que le $B$-homomorphisme associé
\[
H^q(\mathfrak U,\mathcal F)\otimes_A B
\longrightarrow H^q(\mathfrak U',\mathcal F')
\]
soit un isomorphisme. Pour toute suite
$\mathbf s=(i_k)_{0\leq k\leq p}$ de $p+1$ indices de $[1,r]$, posons
\[
U_{\mathbf s}=\bigcap_{k=0}^pU_{i_k},\qquad
U_{\mathbf s}'=\bigcap_{k=0}^pU_{i_k}'=v'^{-1}(U_{\mathbf s}),\qquad
M_{\mathbf s}=\Gamma(U_{\mathbf s},\mathcal F),\qquad
M_{\mathbf s}'=\Gamma(U_{\mathbf s}',\mathcal F').
\]
L'application canonique
$M_{\mathbf s}\otimes_A B\to M_{\mathbf s}'$ est un isomorphisme
(I, 1.6.5), donc l'application canonique
$C^p(\mathfrak U,\mathcal F)\otimes_A B\to
C^p(\mathfrak U',\mathcal F')$ est un isomorphisme, par lequel
$d\otimes1$ s'identifie à l'opérateur cobord
$C^p(\mathfrak U',\mathcal F')\to C^{p+1}(\mathfrak U',\mathcal F')$.
Comme $B$ est un $A$-module \emph{plat}, il résulte aussitôt de la définition
des modules de cohomologie que l'application canonique
$H^q(\mathfrak U,\mathcal F)\otimes_A B\to
H^q(\mathfrak U',\mathcal F')$ est un isomorphisme
(0\textsubscript{I}, 6.1.1). Ce résultat sera généralisé plus tard (\S~6).
\end{proof}

\begin{corollary}[1.4.16]
\label{III.1.4.16-fr}
Soient $A$ un anneau, $X$ un $A$-schéma de type fini, $B$ une $A$-algèbre
fidèlement plate sur $A$. Pour que $X$ soit affine, il faut et il suffit que
$X\otimes_A B$ le soit.
\end{corollary}

\begin{proof}
La condition est évidemment nécessaire (I, 3.2.2)~; montrons qu'elle est
suffisante. Comme $X$ est séparé sur $A$ et que le morphisme
$\operatorname{Spec}(B)\to\operatorname{Spec}(A)$ est plat, il résulte de
(1.4.1) que l'on a
\begin{equation}
\label{III.1.4.16.1-fr}
\tag{1.4.16.1}
H^i(X\otimes_A B,\mathcal F\otimes_A B)
=H^i(X,\mathcal F)\otimes_A B
\end{equation}
pour tout $i\geq0$ et tout $\mathcal O_X$-Module quasi-cohérent $\mathcal F$.
Si $X\otimes_A B$ est affine, le premier membre de
(\hyperref[III.1.4.16.1-fr]{1.4.16.1}) est nul pour $i=1$, donc il en est de
même de $H^1(X,\mathcal F)$ puisque $B$ est un $A$-module fidèlement plat.
Comme $X$ est un schéma quasi-compact, on conclut par le critère de Serre
(II, 5.2.1).
\end{proof}

\begin{proposition}[1.4.17]
\label{III.1.4.17-fr}
Soient $X$ un préschéma,
$0\to\mathcal F\xrightarrow{u}\mathcal G\xrightarrow{v}\mathcal H\to0$
une suite exacte de $\mathcal O_X$-Modules. Si $\mathcal F$ et $\mathcal H$
sont quasi-cohérents, il en est de même de $\mathcal G$.
\end{proposition}

\oldpage[III]{95}
\begin{proof}
La question étant locale sur $X$, on peut supposer
$X=\operatorname{Spec}(A)$ affine et il suffit alors de prouver que
$\mathcal G$ vérifie les conditions d~1) et d~2) de (I, 1.4.1) (avec
$V=X$). La vérification de d~2) est immédiate, car si
$t\in\Gamma(X,\mathcal G)$ a une restriction nulle à $D(f)$, il en est de
même de son image $v(t)\in\Gamma(X,\mathcal H)$~; il y a donc $m>0$ tel que
$f^mv(t)=v(f^mt)=0$ (I, 1.4.1), et comme $\Gamma$ est exact à gauche,
$f^mt=u(s)$, où $s\in\Gamma(X,\mathcal F)$~; comme $u$ est injectif, la
restriction de $s$ à $D(f)$ est nulle, d'où (I, 1.4.1) l'existence d'un
entier $n>0$ tel que $f^ns=0$~; on en déduit finalement
$f^{m+n}t=u(f^ns)=0$.

Vérifions maintenant d~2), et soit $t'\in\Gamma(D(f),\mathcal G)$~; comme
$\mathcal H$ est quasi-cohérent, il existe un entier $m$ tel que
$f^mv(t')=v(f^mt')$ se prolonge en une section
$z\in\Gamma(X,\mathcal H)$ (I, 1.4.1). Mais en vertu de
(\hyperref[III.1.3.1-fr]{1.3.1}) (ou de (I, 5.1.9.2)) appliqué au
$\mathcal O_X$-Module quasi-cohérent $\mathcal F$, la suite
$\Gamma(X,\mathcal G)\to\Gamma(X,\mathcal H)\to0$ est exacte, donc il existe
$t\in\Gamma(X,\mathcal G)$ tel que $z=v(t)$~; on voit donc que
$v(f^mt'-t'')=0$, en désignant par $t''$ la restriction de $t$ à $D(f)$~;
on a donc $f^mt'-t''=u(s')$, où $s'\in\Gamma(D(f),\mathcal F)$. Mais comme
$\mathcal F$ est quasi-cohérent, il existe un entier $n>0$ tel que $f^ns'$
se prolonge en une section $s\in\Gamma(X,\mathcal F)$~; comme
$f^{m+n}t'-f^nt''=u(f^ns')$, on voit que $f^{m+n}t'$ est la restriction à
$D(f)$ de la section $f^nt+u(f^ns)\in\Gamma(X,\mathcal G)$, ce qui achève
la démonstration.
\end{proof}

\section{Étude cohomologique des morphismes projectifs}
\label{section:III.2-fr}

\subsection{Calculs explicites de certains groupes de cohomologie}
\label{subsection:III.2.1-fr}

\begin{env}[2.1.1]
\label{III.2.1.1-fr}
Soient $X$ un préschéma, $\mathcal L$ un $\mathcal O_X$-Module inversible~;
considérons l'anneau gradué (0\textsubscript{I}, 5.4.6)
\begin{equation}
\label{III.2.1.1.1-fr}
\tag{2.1.1.1}
S=\Gamma_\bullet(X,\mathcal L)
=\bigoplus_{n\in\mathbb Z}\Gamma(X,\mathcal L^{\otimes n}).
\end{equation}
Soit $(f_i)_{1\leq i\leq r}$ une famille finie d'éléments \emph{homogènes}
de $S$, soit $f_i\in S_{d_i}$~; posons
$U_i=X_{f_i}$, $U=\bigcup_iU_i$, et désignons par $\mathfrak U$ le
recouvrement $(U_i)$ de $U$. Pour tout $\mathcal O_X$-Module quasi-cohérent
$\mathcal F$, on posera
\begin{align}
\label{III.2.1.1.2-fr}
\tag{2.1.1.2}
H^\bullet(U,\mathcal F(*))
&=\bigoplus_{n\in\mathbb Z}
H^\bullet(U,\mathcal F\otimes\mathcal L^{\otimes n}),\\
\label{III.2.1.1.3-fr}
\tag{2.1.1.3}
H^\bullet(\mathfrak U,\mathcal F(*))
&=\bigoplus_{n\in\mathbb Z}
H^\bullet(\mathfrak U,\mathcal F\otimes\mathcal L^{\otimes n}).
\end{align}
Les groupes abéliens (\hyperref[III.2.1.1.2-fr]{2.1.1.2}) et
(\hyperref[III.2.1.1.3-fr]{2.1.1.3}) sont \emph{bigradués}, en prenant
\[
(H^\bullet(U,\mathcal F(*)))_{mn}
=H^m(U,\mathcal F\otimes\mathcal L^{\otimes n})
\]
et une définition analogue pour
(\hyperref[III.2.1.1.3-fr]{2.1.1.3}). Pour le deuxième degré, il est clair
que ces groupes sont des $S$-modules gradués, comme il résulte par exemple du
fait que $\mathcal F\mapsto H^m(U,\mathcal F)$ et
$\mathcal F\mapsto H^m(\mathfrak U,\mathcal F)$ sont des foncteurs.
\end{env}

\begin{env}[2.1.2]
\label{III.2.1.2-fr}
Considérons maintenant le $S$-module gradué (0\textsubscript{I}, 5.4.6)
\begin{equation}
\label{III.2.1.2.1-fr}
\tag{2.1.2.1}
M=\Gamma_\bullet(\mathcal L,\mathcal F)
=H^0(X,\mathcal F(*))
=\bigoplus_{n\in\mathbb Z}
\Gamma(X,\mathcal F\otimes\mathcal L^{\otimes n}).
\end{equation}
\oldpage[III]{96}
Si $X$ est un préschéma dont l'espace sous-jacent est noethérien, ou un
schéma quasi-compact, il résulte de (I, 9.3.1) qu'en posant comme d'ordinaire
$U_{i_0i_1\ldots i_p}=\bigcap_{k=0}^pU_{i_k}$, on a, à un isomorphisme
canonique près,
\[
\Gamma(U_{i_0i_1\ldots i_p},\mathcal F(*))
=H^0(U_{i_0i_1\ldots i_p},\mathcal F(*))
=M_{f_{i_0}f_{i_1}\cdots f_{i_p}}.
\]
On peut encore, avec les notations de
(\hyperref[III.1.2.2-fr]{1.2.2}), identifier
$M_{f_{i_0}f_{i_1}\cdots f_{i_p}}$ à
$\varinjlim_n M_{i_0i_1\ldots i_p}^{(n)}$. Cette identification est un
isomorphisme de $S$-modules gradués, si l'on définit le degré d'un élément
homogène $z\in\varinjlim_n M_{i_0i_1\ldots i_p}^{(n)}$ de la façon
suivante~: $z$ est l'image canonique d'un élément homogène
$x\in M_{i_0i_1\ldots i_p}^{(n)}=M$, de degré $m$~; on prend alors pour
degré de $z$ le nombre
$m-n(d_{i_0}+d_{i_1}+\cdots+d_{i_p})$. Tenant compte de la définition des
homomorphismes
$\varphi_{kn}:M_{i_0i_1\ldots i_p}^{(n)}\to
M_{i_0i_1\ldots i_p}^{(k)}$ (\hyperref[III.1.2.2-fr]{1.2.2}), on voit
aussitôt que cette définition ne dépend pas du «~représentant~» $x$ de $z$
que l'on a considéré. Désignant comme dans
(\hyperref[III.1.2.2-fr]{1.2.2}) par $C_n^p(M)$ l'ensemble des applications
alternées de $[1,r]^{p+1}$ dans $M$ (pour tout $n$), on définit de la même
manière que ci-dessus une structure de $S$-module gradué sur
$\varinjlim_n C_n^p(M)$~; on a encore comme dans
(\hyperref[III.1.2.2-fr]{1.2.2})
\begin{equation}
\label{III.2.1.2.2-fr}
\tag{2.1.2.2}
C^p(\mathfrak U,\mathcal F(*))=\varinjlim_n C_n^p(M)
\end{equation}
l'isomorphisme des deux membres \emph{respectant les degrés}. On a alors,
comme dans (\hyperref[III.1.2.2-fr]{1.2.2})
\begin{equation}
\label{III.2.1.2.3-fr}
\tag{2.1.2.3}
C^p(\mathfrak U,\mathcal F(*))
=C^{p+1}((\mathbf f),M)
=\varinjlim_n K^{p+1}(\mathbf f^n,M)
\end{equation}
l'isomorphisme conservant les degrés~: le degré d'un élément de
$\varinjlim_n K^{p+1}(\mathbf f^n,M)$, image canonique d'une cochaîne
$g\in K^{p+1}(\mathbf f^n,M)$ dont les valeurs
$g(i_0,\ldots,i_p)$ sont dans une même composante homogène $M_m$ de $M$,
est $m-n(d_{i_0}+\cdots+d_{i_p})$, et il est indépendant du choix de cette
cochaîne comme représentant de l'élément considéré.
\end{env}

Comme les isomorphismes précédents sont compatibles avec les opérateurs
cobords, on en conclut, comme dans
(\hyperref[III.1.2.2-fr]{1.2.2}) que l'on a~:

\begin{proposition}[2.1.3]
\label{III.2.1.3-fr}
Soit $X$ un préschéma dont l'espace sous-jacent est noethérien, ou un schéma
quasi-compact. Il existe un isomorphisme canonique fonctoriel en $\mathcal F$
\begin{equation}
\label{III.2.1.3.1-fr}
\tag{2.1.3.1}
H^p(\mathfrak U,\mathcal F(*))
\xrightarrow{\sim}H^{p+1}((\mathbf f),M)
\qquad\text{pour tout }p\geq1.
\end{equation}
On a en outre une suite exacte fonctorielle en $\mathcal F$
\begin{equation}
\label{III.2.1.3.2-fr}
\tag{2.1.3.2}
0\longrightarrow H^0((\mathbf f),M)\longrightarrow M
\longrightarrow H^0(\mathfrak U,\mathcal F(*))
\longrightarrow H^1((\mathbf f),M)\longrightarrow0.
\end{equation}
De plus, tous les homomorphismes introduits sont de degré $0$ pour les
structures de $S$-modules gradués ($S$ étant l'anneau
(\hyperref[III.2.1.1.1-fr]{2.1.1.1})).
\end{proposition}
\providecommand{\bbGamma}{\boldsymbol{\Gamma}}
\providecommand{\shProj}{\mathcal{P}\mathrm{roj}}

\oldpage[III]{97}
\begin{corollary}[2.1.4]
\label{III.2.1.4-fr}
Si $X$ est un schéma quasi-compact et si les $U_i=X_{f_i}$ sont affines,
il existe un isomorphisme canonique fonctoriel en $\mathcal F$, de degré $0$,
\begin{equation}
\label{III.2.1.4.1-fr}
\tag{2.1.4.1}
H^p(U,\mathcal F(*))\xrightarrow{\sim}H^{p+1}((\mathbf f),M)
\qquad\text{pour }p\geq1,
\end{equation}
et une suite exacte fonctorielle en $\mathcal F$
\begin{equation}
\label{III.2.1.4.2-fr}
\tag{2.1.4.2}
0\longrightarrow H^0((\mathbf f),M)\longrightarrow M
\longrightarrow H^0(U,\mathcal F(*))
\longrightarrow H^1((\mathbf f),M)\longrightarrow0,
\end{equation}
où tous les homomorphismes sont de degré $0$.
\end{corollary}

\begin{proof}
Il suffit en effet d'appliquer (1.4.1) au résultat de
(\hyperref[III.2.1.3-fr]{2.1.3}).
\end{proof}

La proposition «~locale~» analogue à
(\hyperref[III.2.1.3-fr]{2.1.3}) est la suivante~:

\begin{proposition}[2.1.5]
\label{III.2.1.5-fr}
Soient $S$ un anneau gradué à degrés positifs,
$f_i$ $(1\leq i\leq r)$ un élément homogène de $S_+$ de degré $d_i$,
$M$ un $S$-module gradué. Soit $X=\operatorname{Proj}(S)$ le spectre
premier homogène de $S$ et posons $U_i=D_+(f_i)$,
$U=\bigcup_iU_i$,
\[
H^\bullet(U,\widetilde M(*))
=\bigoplus_{n\in\mathbb Z}H^\bullet(U,(M(n))^\sim).
\]
Il existe alors des isomorphismes canoniques fonctoriels en $M$, de degré
$0$ pour les structures de $S$-module gradué
\begin{equation}
\label{III.2.1.5.1-fr}
\tag{2.1.5.1}
H^p(U,\widetilde M(*))\xrightarrow{\sim}H^{p+1}((\mathbf f),M)
\qquad\text{pour }p\geq1,
\end{equation}
et une suite exacte fonctorielle en $M$
\begin{equation}
\label{III.2.1.5.2-fr}
\tag{2.1.5.2}
0\longrightarrow H^0((\mathbf f),M)\longrightarrow M
\longrightarrow H^0(U,\widetilde M(*))
\longrightarrow H^1((\mathbf f),M)\longrightarrow0,
\end{equation}
où tous les homomorphismes sont de degré $0$.
\end{proposition}

\begin{proof}
En effet, on a
\[
\Gamma(U_{i_0i_1\ldots i_p},(M(n))^\sim)
=(M_{f_{i_0}\cdots f_{i_p}})_n
\]
par définition (II, 2.5.2), donc
$\Gamma(U_{i_0i_1\ldots i_p},\widetilde M(*))
=M_{f_{i_0}\cdots f_{i_p}}$. Le reste du raisonnement est alors le même
que pour démontrer (\hyperref[III.2.1.4-fr]{2.1.4}), tenant compte de ce
que $X$ est un schéma.
\end{proof}

\begin{env}[2.1.6]
\label{III.2.1.6-fr}
\emph{Remarques.} ---
(i) Sous les conditions de (\hyperref[III.2.1.5-fr]{2.1.5}), les foncteurs
$\Gamma(U_{i_0i_1\ldots i_p},\widetilde M(*))$ sont exacts en $M$, en
vertu de (0\textsubscript{I}, 1.3.2)~; le même raisonnement que dans
(1.2.5) montre alors que si
$0\to M'\to M\to M''\to0$ est une suite exacte de $S$-modules gradués
(où les homomorphismes sont de degré $0$), on a des diagrammes
commutatifs pour tout $p\geq0$
\begin{equation}
\label{III.2.1.6.1-fr}
\tag{2.1.6.1}
\xymatrix{
H^p(U,\widetilde {M''}(*))\ar[r]^-\partial\ar[d] &
H^{p+1}(U,\widetilde {M'}(*))\ar[d]\\
H^{p+1}((\mathbf f),M'')\ar[r]^-\partial &
H^{p+2}((\mathbf f),M').
}
\end{equation}

(ii) La prop. (\hyperref[III.2.1.5-fr]{2.1.5}) sera surtout intéressante
lorsque $S$ sera une $A$-algèbre engendrée par un nombre fini d'éléments de
degré $1$, $A$ étant supposé noethérien~; en
\oldpage[III]{98}
effet, lorsqu'il en est ainsi, tout $\mathcal O_X$-Module quasi-cohérent est
de la forme $\widetilde M$ (II, 2.7.7).
\end{env}

\begin{env}[2.1.7]
\label{III.2.1.7-fr}
Nous allons appliquer (\hyperref[III.2.1.5-fr]{2.1.5}) dans le cas
$S=A[T_0,\ldots,T_r]$, où $A$ est un anneau quelconque, les $T_i$ sont des
indéterminées, avec $M=S$ et $f_i=T_i$. On est donc essentiellement ramené
à calculer $H^\bullet((\mathbf T),S)$, où
$\mathbf T=(T_i)_{0\leq i\leq r}$.
\end{env}

\begin{lemma}[2.1.8]
\label{III.2.1.8-fr}
Si $S=A[T_0,\ldots,T_r]$, on a, avec
$\mathbf T=(T_i)_{0\leq i\leq r}$,
\begin{align}
\label{III.2.1.8.1-fr}
\tag{2.1.8.1}
H^i(\mathbf T^n,S)&=0 &&\text{si }i\neq r+1,\\
\label{III.2.1.8.2-fr}
\tag{2.1.8.2}
H^{r+1}(\mathbf T^n,S)&=S/(\mathbf T^n).
\end{align}
Le $A$-module $H^{r+1}(\mathbf T^n,S)$ a donc une base sur $A$ formée des
classes mod. $(\mathbf T^n)$ des monômes
\[
T_{\mathbf p}=T_0^{p_0}\cdots T_r^{p_r},
\qquad \mathbf p=(p_0,\ldots,p_r),\quad 0\leq p_i<n
\quad\text{pour tout }i.
\]
\end{lemma}

\begin{proof}
C'est une conséquence immédiate de (1.1.3.5) et de la prop. (1.1.4), dont
les hypothèses sont trivialement vérifiées.
\end{proof}

\begin{env}[2.1.9]
\label{III.2.1.9-fr}
Passons à la limite inductive sur $n$~; les relations
(\hyperref[III.2.1.8.1-fr]{2.1.8.1}) donnent
$H^i((\mathbf T),S)=0$ pour $i\neq r+1$. Pour $i=r+1$, le système inductif
est formé des $S/(\mathbf T^n)$, l'homomorphisme
\[
\varphi_{nm}:S/(\mathbf T^n)\longrightarrow S/(\mathbf T^m)
\qquad(0\leq n\leq m)
\]
étant la multiplication par $(T_0\cdots T_r)^{m-n}$. Pour
$n\geq\sup(p_i)_{0\leq i\leq r}$, désignons par
$\xi_{\mathbf p}^{(n)}=\xi_{p_0,\ldots,p_r}^{(n)}$ la classe de
$T_0^{n-p_0}\cdots T_r^{n-p_r}\pmod{(\mathbf T^n)}$~; on a alors
$\varphi_{nm}(\xi_{\mathbf p}^{(n)})=\xi_{\mathbf p}^{(m)}$, et ces éléments
ont donc même image canonique
$\xi_{\mathbf p}=\xi_{p_0,\ldots,p_r}$ dans la limite inductive
$H^{r+1}((\mathbf T),S)$~; en vertu de la définition du degré donnée dans
(\hyperref[III.2.1.2-fr]{2.1.2}), le degré de $\xi_{\mathbf p}$ est donc
égal à $-|\mathbf p|=-(p_0+p_1+\cdots+p_r)$. Il est clair que les
$\xi_{\mathbf p}^{(n)}$ pour $0<p_i\leq n$ et $0\leq i\leq r$ forment une
base de $S/(\mathbf T^n)$. On déduit donc aussitôt de
(\hyperref[III.2.1.8-fr]{2.1.8})~:
\end{env}

\begin{corollary}[2.1.10]
\label{III.2.1.10-fr}
Avec les notations de (\hyperref[III.2.1.8-fr]{2.1.8}), on a
\begin{equation}
\label{III.2.1.10.1-fr}
\tag{2.1.10.1}
H^i((\mathbf T),S)=0\qquad\text{pour }i\neq r+1,
\end{equation}
et $H^{r+1}((\mathbf T),S)$ est un $A$-module libre dont une base est formée
des éléments $\xi_{p_0,\ldots,p_r}$ tels que $p_i>0$ pour $0\leq i\leq r$.
\end{corollary}

\begin{env}[2.1.11]
\label{III.2.1.11-fr}
\emph{Remarque.} --- Soit $N$ un $A$-module quelconque et soit
$M=S\otimes_A N$~; le raisonnement de
(\hyperref[III.2.1.8-fr]{2.1.8}) montre que l'on a plus généralement
\begin{align}
\label{III.2.1.11.1-fr}
\tag{2.1.11.1}
H^i(\mathbf T^n,M)&=0 &&\text{si }i\neq r+1,\\
\label{III.2.1.11.2-fr}
\tag{2.1.11.2}
H^{r+1}(\mathbf T^n,M)&=(S/(\mathbf T^n))\otimes_A N,
\end{align}
car la dernière formule se déduit directement de (1.1.3.5), et d'autre
part il est clair que
\[
M/(T_0^nM+\cdots+T_{i-1}^nM)
=\bigl(S/(T_0^nS+\cdots+T_{i-1}^nS)\bigr)\otimes_A N,
\]
l'idéal $T_0^nS+\cdots+T_{i-1}^nS$ étant facteur direct dans le
$A$-module $S$~; cela permet d'appliquer (1.1.4) à $M$, et on obtient ainsi
(\hyperref[III.2.1.11.1-fr]{2.1.11.1}).
\end{env}

Combinant (\hyperref[III.2.1.10-fr]{2.1.10}) et
(\hyperref[III.2.1.5-fr]{2.1.5}), on obtient~:

\begin{proposition}[2.1.12]
\label{III.2.1.12-fr}
Soient $A$ un anneau, $r$ un entier $>0$, et
$X=\mathbf P_A^r$ (II, 4.1.1). Alors~:
\begin{enumerate}
\item[{\rm(i)}] On a $H^i(X,\mathcal O_X(*))=0$ pour $i\neq0,r$.
\item[{\rm(ii)}] L'homomorphisme canonique
$\alpha:S\to H^0(X,\mathcal O_X(*))$ (II, 2.6.2) est bijectif.
\oldpage[III]{99}
\item[{\rm(iii)}] $H^r(X,\mathcal O_X(*))$ est un $A$-module libre ayant
une base formée d'éléments $\xi_{p_0,\ldots,p_r}$, où $p_i>0$ pour
$0\leq i\leq r$, $\xi_{p_0,\ldots,p_r}$ étant de degré
$-|\mathbf p|=-(p_0+\cdots+p_r)$, et le produit
$T_i\xi_{p_0,\ldots,p_r}$ étant
$\xi_{p_0,\ldots,p_i-1,\ldots,p_r}$.
\end{enumerate}
\end{proposition}

\begin{proof}
Remarquons en effet que, dans la suite exacte
(\hyperref[III.2.1.5.2-fr]{2.1.5.2}) appliquée à
$M=S=A[T_0,\ldots,T_r]$, on a
$H^0((\mathbf T),S)=0$ et $H^1((\mathbf T),S)=0$ d'après
(\hyperref[III.2.1.10.1-fr]{2.1.10.1}), et que la prop.
(\hyperref[III.2.1.5-fr]{2.1.5}) s'applique à $U=X$, puisque $X$ est
réunion des $D_+(T_i)$ (II, 2.3.14). Il reste à identifier l'application
$S\to H^0(X,\mathcal O_X(*))$ de la suite exacte
(\hyperref[III.2.1.5.2-fr]{2.1.5.2}) et l'application canonique $\alpha$~;
mais cela résulte de l'identification canonique de
$H^0(U,\mathcal O_X(*))$ et de $H^0(\mathfrak U,\mathcal O_X(*))$.
\end{proof}

\begin{corollary}[2.1.13]
\label{III.2.1.13-fr}
Les seules valeurs de $(i,n)$ pour lesquelles on puisse avoir
$H^i(X,\mathcal O_X(n))\neq0$ sont les suivantes~:
$i=0$ et $n\geq0$, $i=r$ et $n\leq-(r+1)$.
\end{corollary}

On notera que si $A\neq0$, on a effectivement
$H^i(X,\mathcal O_X(n))\neq0$ pour les couples énumérés dans
(\hyperref[III.2.1.13-fr]{2.1.13})~; cela résulte de
(\hyperref[III.2.1.12-fr]{2.1.12}), puisque $S_n$ est alors $\neq0$ pour
tous les degrés $n\geq0$.

Dans les applications qui seront faites dans ce chapitre, nous utiliserons
surtout le résultat moins précis~:

\begin{corollary}[2.1.14]
\label{III.2.1.14-fr}
Les $A$-modules $H^i(X,\mathcal O_X(n))$ sont libres de type fini~; si
$i>0$, ils sont nuls pour $n>0$.
\end{corollary}

\begin{proposition}[2.1.15]
\label{III.2.1.15-fr}
Soient $Y$ un préschéma, $\mathcal E$ un $\mathcal O_Y$-Module localement
libre de rang $r+1$, $X=\mathbf P(\mathcal E)$ le fibré projectif défini
par $\mathcal E$, $f:X\to Y$ le morphisme structural. Les seules valeurs de
$i$ et $n$ pour lesquelles $R^if_*(\mathcal O_X(n))\neq0$ sont
$i=0$ et $n\geq0$, $i=r$ et $n\leq-(r+1)$~; en outre, l'homomorphisme
canonique (II, 3.3.2)
\[
\alpha:S_{\mathcal O_Y}(\mathcal E)
\longrightarrow\Gamma_*(\mathcal O_X)
=R^0f_*(\mathcal O_X(*))
=\bigoplus_{n\in\mathbb Z}f_*(\mathcal O_X(n))
\]
est un isomorphisme.
\end{proposition}

\begin{proof}
La question étant locale sur $Y$, on peut supposer $Y$ affine d'anneau
$A$ et $\mathcal E=\widetilde E$, où $E=A^{r+1}$~; on est alors
immédiatement ramené à (\hyperref[III.2.1.12-fr]{2.1.12}), compte tenu de
(1.4.11).
\end{proof}

\begin{env}[2.1.16]
\label{III.2.1.16-fr}
\emph{Remarque.} --- Nous compléterons plus tard les résultats de
(\hyperref[III.2.1.15-fr]{2.1.15}) en démontrant les propositions
suivantes~: posons
\[
\boldsymbol\omega=f^*\left(\bigwedge^{r+1}\mathcal E\right)(-r-1),
\]
qui est un $\mathcal O_X$-Module inversible. Alors~:
\begin{enumerate}
\item[{\rm(i)}] On a un isomorphisme canonique
\begin{equation}
\label{III.2.1.16.1-fr}
\tag{2.1.16.1}
\rho:R^rf_*(\boldsymbol\omega)\xleftarrow{\sim}\mathcal O_Y.
\end{equation}
\item[{\rm(ii)}] L'accouplement par cup-produit (0, 12.2.2)
\begin{equation}
\label{III.2.1.16.2-fr}
\tag{2.1.16.2}
R^rf_*(\mathcal O_X(n))\times R^0f_*(\boldsymbol\omega(-n))
\longrightarrow R^rf_*(\boldsymbol\omega)
\end{equation}
composé avec l'isomorphisme $\rho^{-1}$, définit un \emph{isomorphisme}
de $R^rf_*(\mathcal O_X(n))$ sur le \emph{dual} du $\mathcal O_Y$-Module
localement libre
\[
R^0f_*(\boldsymbol\omega(-n))
=\left(\bigwedge^{r+1}\mathcal E\right)
\otimes_{\mathcal O_Y}\bigl(S_{\mathcal O_Y}(\mathcal E)\bigr)_{-n}.
\]
\end{enumerate}
\end{env}
\subsection{Le théorème fondamental des morphismes projectifs}
\label{subsection:III.2.2-fr}
\oldpage[III]{100}

\begin{theorem}[2.2.1]
\label{III.2.2.1-fr}
(Serre). Soient $Y$ un préschéma noethérien, $f:X\to Y$ un morphisme
propre, $\mathcal L$ un $\mathcal O_X$-Module inversible ample pour $f$.
Pour tout $\mathcal O_X$-Module $\mathcal F$, posons
\[
\mathcal F(n)=\mathcal F\otimes_{\mathcal O_X}\mathcal L^{\otimes n}
\qquad\text{pour tout }n\in\mathbb Z.
\]
Alors, pour tout $\mathcal O_X$-Module cohérent $\mathcal F$~:
\begin{enumerate}
\item[{\rm(i)}] Les $R^qf_*(\mathcal F)$ sont des
$\mathcal O_Y$-Modules cohérents.
\item[{\rm(ii)}] Il existe un entier $N$ tel que pour $n\geq N$, on ait
\[
R^qf_*(\mathcal F(n))=0\qquad\text{pour tout }q>0.
\]
\item[{\rm(iii)}] Il existe un entier $N$ tel que, pour $n\geq N$,
l'homomorphisme canonique
\[
f^*(f_*(\mathcal F(n)))\longrightarrow\mathcal F(n)
\]
soit surjectif.
\end{enumerate}
\end{theorem}

\begin{proof}
Notons d'abord que si le théorème est vrai quand on y remplace
$\mathcal L$ par $\mathcal L^{\otimes d}$ $(d>0)$, il est vrai sous sa
forme initiale. En effet, on peut alors écrire
\[
\mathcal F(n)
=(\mathcal F\otimes\mathcal L^{\otimes r})
\otimes\mathcal L^{\otimes hd}
\]
avec un $h>0$ et $0\leq r<d$, et par hypothèse pour chaque $r$ il y a
un entier $N_r$ tel que pour $h\geq N_r$ les propriétés (ii) et (iii)
aient lieu pour le $\mathcal O_X$-Module
$\mathcal F\otimes\mathcal L^{\otimes r}$~; prenant pour $N$ le plus grand
des $dN_r$, (ii) et (iii) auront lieu pour $n\geq N$. On peut donc
supposer $\mathcal L$ très ample relativement à $f$ (II, 4.6.11)~; il
existe par suite une $Y$-immersion ouverte dominante
$i:X\to P$, où $P=\operatorname{Proj}(\mathcal S)$, où $\mathcal S$ est
une $\mathcal O_Y$-Algèbre quasi-cohérente graduée à degrés positifs, dans
laquelle $\mathcal S_1$ est de type fini et engendre $\mathcal S$~; en
outre, $\mathcal L$ est isomorphe à $i^*(\mathcal O_P(1))$ (II, 4.4.7).
Mais comme $f$ est propre, il en est de même de $i$ (II, 5.4.4), donc
$i$ est un isomorphisme $X\xrightarrow{\sim}P$. On peut donc se borner au
cas où $X=\operatorname{Proj}(\mathcal S)$ et
$\mathcal L=\mathcal O_X(1)$. Le th. (2.2.1) est alors conséquence de la
\end{proof}

\begin{proposition}[2.2.2]
\label{III.2.2.2-fr}
Soient $A$ un anneau noethérien, $S$ une $A$-algèbre graduée à degrés
positifs, dans laquelle $S_1$ est un $A$-module ayant un système de $r+1$
générateurs, et qui engendre l'algèbre $S$. Soit
$X=\operatorname{Proj}(S)$. Pour tout $\mathcal O_X$-Module cohérent
$\mathcal F$~:
\begin{enumerate}
\item[{\rm(i)}] Les $A$-modules $H^q(X,\mathcal F)$ sont de type fini.
\item[{\rm(ii)}] On a $H^q(X,\mathcal F)=0$ pour $q>r$.
\item[{\rm(iii)}] Il existe un entier $N$ tel que pour $n\geq N$, on ait
$H^q(X,\mathcal F(n))=0$ pour tout $q>0$.
\item[{\rm(iv)}] Il existe un entier $N$ tel que pour $n\geq N$,
$\mathcal F(n)$ soit engendré par ses sections au-dessus de $X$.
\end{enumerate}
\end{proposition}

\begin{proof}
Montrons d'abord comment (2.2.2) entraîne (2.2.1)~: dans (2.2.1)
(ramené au cas particulier $X=\operatorname{Proj}(\mathcal S)$ considéré
ci-dessus), $Y$ est quasi-compact, donc peut être recouvert par un nombre
fini d'ouverts affines, d'anneaux noethériens, tels que la restriction de
$\mathcal S_1$ à chacun de ces ouverts $U_\alpha$ soit engendrée par un
nombre fini de sections de $\mathcal S_1$ au-dessus de $U_\alpha$. Si on
suppose (2.2.2) démontré, il suffira alors de prendre pour $N$ dans les
parties (ii) et (iii) de (2.2.1) le plus grand des entiers analogues
correspondant aux $U_\alpha$ (tenant compte de (1.4.11) et de
(II, 3.4.7)).

Pour prouver (2.2.2), remarquons que $X$ s'identifie à un sous-schéma
fermé de $P=\mathbf P_A^r$ (II, 3.6.2)~; en outre, si $j:X\to P$ est
l'injection canonique, $j_*(\mathcal F)$ est un
$\mathcal O_P$-Module cohérent et on a
$j_*(\mathcal F(n))=(j_*(\mathcal F))(n)$ (II, 3.4.5 et 3.5.2).
Compte tenu de (G, II, cor. du th. 4.9.1), on est donc ramené à prouver
(2.2.2) dans le cas particulier où
\[
X=\mathbf P_A^r\quad\text{et}\quad S=A[T_0,\ldots,T_r].
\]

\oldpage[III]{101}
Comme $X$ est recouvert par les ouverts affines $D_+(T_i)$ en nombre
$r+1$, (ii) résulte de (1.4.12). Notons d'autre part que (iv) a déjà été
démontré (II, 2.7.9).

Nous allons prouver simultanément (i) et (iii). Notons que ces assertions
sont vraies pour $\mathcal F=\mathcal O_X(m)$
(\hyperref[III.2.1.13-fr]{2.1.13})~; elles le sont donc aussi lorsque
$\mathcal F$ est somme directe d'un nombre fini de
$\mathcal O_X$-Modules de la forme $\mathcal O_X(m_j)$. D'autre part,
(i) et (iii) sont vraies trivialement pour $q>r$ en vertu de (ii). Nous
allons procéder par \emph{récurrence descendante} sur $q$. On sait que
$\mathcal F$ est isomorphe à un quotient d'une somme directe $\mathcal E$
d'un nombre fini de faisceaux $\mathcal O_X(m_j)$ (II, 2.7.10)~;
autrement dit, on a une suite exacte
\[
0\longrightarrow\mathcal R\longrightarrow\mathcal E
\longrightarrow\mathcal F\longrightarrow0,
\]
où $\mathcal R$ est cohérent (0\textsubscript{I}, 5.3.3) et où
$\mathcal E$ vérifie (i) et (iii). Comme $\mathcal F(n)$ est un foncteur
exact en $\mathcal F$, on a aussi la suite exacte
\[
0\longrightarrow\mathcal R(n)\longrightarrow\mathcal E(n)
\longrightarrow\mathcal F(n)\longrightarrow0
\]
pour tout $n\in\mathbb Z$. On en déduit la suite exacte de cohomologie
\[
H^{q-1}(X,\mathcal E(n))\longrightarrow
H^{q-1}(X,\mathcal F(n))\longrightarrow
H^q(X,\mathcal R(n)).
\]
Comme $\mathcal E(n)$ est somme directe des $\mathcal O_X(n+m_j)$
(II, 2.5.14), $H^{q-1}(X,\mathcal E(n))$ est de type fini, et il en est
de même de $H^q(X,\mathcal R(n))$ par l'hypothèse de récurrence~; comme
$A$ est noethérien, on en conclut que
$H^{q-1}(X,\mathcal F(n))$ est de type fini pour tout $n\in\mathbb Z$, et
en particulier pour $n=0$. D'autre part, par l'hypothèse de récurrence,
il existe un entier $N$ tel que pour $n\geq N$ on ait
$H^q(X,\mathcal R(n))=0$~; par ailleurs, on peut supposer aussi $N$ choisi
tel que $H^{q-1}(X,\mathcal E(n))=0$ pour $n\geq N$, puisque $\mathcal E$
vérifie (iii)~; on en conclut que
$H^{q-1}(X,\mathcal F(n))=0$ pour $n\geq N$, ce qui achève la
démonstration.
\end{proof}

\begin{corollary}[2.2.3]
\label{III.2.2.3-fr}
Sous les hypothèses de (\hyperref[III.2.2.1-fr]{2.2.1}), soit
$\mathcal F\to\mathcal G\to\mathcal H$ une suite exacte de
$\mathcal O_X$-Modules cohérents. Il existe alors un entier $N$ tel que
pour $n\geq N$, la suite
\[
f_*(\mathcal F(n))\longrightarrow f_*(\mathcal G(n))
\longrightarrow f_*(\mathcal H(n))
\]
soit exacte.
\end{corollary}

\begin{proof}
Soient $\mathcal F'$, $\mathcal G'$, $\mathcal G''$ le noyau, l'image et
le conoyau de $\mathcal F\to\mathcal G$~; $\mathcal G'$ est le noyau et
$\mathcal G''$ l'image de $\mathcal G\to\mathcal H$, soit $\mathcal H''$
le conoyau de cet homomorphisme~; tous ces $\mathcal O_X$-Modules sont
cohérents (0\textsubscript{I}, 5.3.4). Comme $\mathcal F(n)$ est un
foncteur exact en $\mathcal F$, il suffit de prouver que pour $n$ assez
grand, chacune des suites
\begin{align*}
0&\longrightarrow f_*(\mathcal F'(n))
\longrightarrow f_*(\mathcal F(n))
\longrightarrow f_*(\mathcal G'(n))\longrightarrow0,\\
0&\longrightarrow f_*(\mathcal G'(n))
\longrightarrow f_*(\mathcal G(n))
\longrightarrow f_*(\mathcal G''(n))\longrightarrow0,\\
0&\longrightarrow f_*(\mathcal G''(n))
\longrightarrow f_*(\mathcal H(n))
\longrightarrow f_*(\mathcal H''(n))\longrightarrow0
\end{align*}
est exacte~; par suite, on peut supposer que
$0\to\mathcal F\to\mathcal G\to\mathcal H\to0$ est exacte. On a alors
la suite exacte de cohomologie
\[
0\longrightarrow f_*(\mathcal F(n))
\longrightarrow f_*(\mathcal G(n))
\longrightarrow f_*(\mathcal H(n))
\longrightarrow R^1f_*(\mathcal F(n))\longrightarrow\cdots
\]
et la conclusion résulte de (\hyperref[III.2.2.1-fr]{2.2.1}, (ii)).
\end{proof}

\begin{corollary}[2.2.4]
\label{III.2.2.4-fr}
Soient $Y$ un préschéma noethérien, $f:X\to Y$ un morphisme de type fini,
$\mathcal L$ un $\mathcal O_X$-Module inversible ample pour $f$~; pour
tout $\mathcal O_X$-Module $\mathcal F$, on pose
$\mathcal F(n)=\mathcal F\otimes_{\mathcal O_X}\mathcal L^{\otimes n}$
(pour $n\in\mathbb Z$). Soit
$\mathcal F\to\mathcal G\to\mathcal H$ une suite exacte de
$\mathcal O_X$-Modules cohérents, telle que les supports de $\mathcal F$
et de $\mathcal H$ soient propres sur $Y$ (II, 5.4.10). Il existe alors
un entier $N$ tel que, pour $n\geq N$, la suite
\[
f_*(\mathcal F(n))\longrightarrow f_*(\mathcal G(n))
\longrightarrow f_*(\mathcal H(n))
\]
soit exacte.
\end{corollary}

\begin{proof}
Le même raisonnement qu'au début de
(\hyperref[III.2.2.1-fr]{2.2.1}) montre que si le corollaire est vrai
pour $\mathcal L^{\otimes d}$ $(d>0)$, il l'est aussi pour $\mathcal L$~;
on peut donc se borner au cas où $\mathcal L$ est très ample pour $f$
(II, 4.6.11), et par suite on peut identifier $X$ à un ouvert dans un
$Y$-schéma $Z=\operatorname{Proj}(\mathcal S)$, où $\mathcal S$ est une
$\mathcal O_Y$-Algèbre quasi-cohérente graduée à degrés positifs, dans
laquelle $\mathcal S_1$ est de type fini et engendre $\mathcal S$, de
sorte que $\mathcal L=i^*(\mathcal O_Z(1))$, où $i$ est l'immersion
canonique $X\to Z$ (II, 4.4.7). Cela étant, comme
$\operatorname{Supp}(\mathcal G)$ est fermé dans $X$ et contenu dans
$\operatorname{Supp}(\mathcal F)\cap\operatorname{Supp}(\mathcal H)$, il
est propre sur $Y$~; les supports de $\mathcal F$, $\mathcal G$,
$\mathcal H$ sont donc \emph{fermés dans $Z$} (II, 5.4.10). Les faisceaux
$\mathcal F'=i_*(\mathcal F)$, $\mathcal G'=i_*(\mathcal G)$,
$\mathcal H'=i_*(\mathcal H)$ sont donc des $\mathcal O_Z$-Modules
cohérents, et la suite $\mathcal F'\to\mathcal G'\to\mathcal H'$ est
exacte~; en outre, si $g:Z\to Y$ est le morphisme structural, on a
$f=g\circ i$, et il est clair que
$\mathcal F'(n)=i_*(\mathcal F(n))$, et de même pour $\mathcal G'$ et
$\mathcal H'$~; la conclusion résulte donc de
(\hyperref[III.2.2.3-fr]{2.2.3}) appliqué à
$\mathcal F'$, $\mathcal G'$, $\mathcal H'$.
\end{proof}

\begin{env}[2.2.5]
\label{III.2.2.5-fr}
\emph{Remarques.} ---
(i) L'assertion (i) de (\hyperref[III.2.2.1-fr]{2.2.1}) est encore vraie
lorsqu'on suppose seulement que $Y$ est \emph{localement noethérien}~; en
effet, la propriété est évidemment locale sur $Y$~; d'autre part, les
hypothèses de (\hyperref[III.2.2.1-fr]{2.2.1}) impliquent que pour tout
ouvert $U\subset Y$, la restriction de $f$ à $f^{-1}(U)$ est un morphisme
projectif $f^{-1}(U)\to U$ (II, 5.5.5, (iii)) et
$\mathcal L|f^{-1}(U)$ est ample pour ce morphisme (II, 4.6.4).

(ii) L'assertion (iii) de (\hyperref[III.2.2.1-fr]{2.2.1}) est encore
valable, comme on l'a vu, lorsque l'on suppose seulement que $X$ est un
schéma quasi-compact ou un préschéma dont l'espace sous-jacent est
noethérien, et $f:X\to Y$ un morphisme quasi-compact (II, 4.6.8). Mais il
faut noter que même lorsqu'on suppose que $Y$ est le \emph{spectre d'un
corps} $K$ et que $f$ est \emph{quasi-projectif}, l'assertion (ii) de
(\hyperref[III.2.2.1-fr]{2.2.1}) n'est plus nécessairement vérifiée. Par
exemple, soit $X'=\operatorname{Spec}(K[T_0,\ldots,T_r])$ et soit $X$ la
réunion des ouverts affines $D(T_i)$ de $X'$ $(0\leq i\leq r)$~; comme
l'immersion $X\to X'$ est quasi-compacte, le morphisme structural
$f:X\to Y$ est quasi-affine (II, 5.1.10), donc $\mathcal O_X$ est très
ample pour $f$ (II, 5.1.6). Mais l'anneau
$\Gamma(X,\mathcal O_X)$ s'identifie à l'intersection des anneaux de
fractions $(K[T_0,\ldots,T_r])_{T_i}$ pour $0\leq i\leq r$
(I, 8.2.1.1), c'est-à-dire à $K[T_0,\ldots,T_r]$. Par suite, il résulte
des formules (1.4.3.1) et (1.1.3.5) que l'on a
$H^r(X,\mathcal O_X^{\otimes n})=H^r(X,\mathcal O_X)=A\neq0$ pour tout
$n$.
\end{env}

\subsection{Application aux faisceaux gradués d'algèbres et de modules}
\label{subsection:III.2.3-fr}

\begin{theorem}[2.3.1]
\label{III.2.3.1-fr}
Soient $Y$ un préschéma noethérien, $\mathcal S$ une
$\mathcal O_Y$-Algèbre graduée à degrés positifs, quasi-cohérente et de
type fini, $X=\operatorname{Proj}(\mathcal S)$,
$q:X\to Y$ le morphisme structural, $\mathcal M$ un $\mathcal S$-Module
gradué quasi-cohérent vérifiant la condition \textbf{(TF)}. Alors il existe
un entier $N$ tel que, pour $n\geq N$, l'homomorphisme canonique
(II, 8.14.5.1)
\[
\alpha_n:\mathcal M_n\longrightarrow
q_*(\shProj_0(\mathcal M(n)))
=q_*((\shProj(\mathcal M))_n)
\]
\oldpage[III]{103}
soit bijectif. En d'autres termes, l'homomorphisme canonique
\[
\alpha:\mathcal M\longrightarrow\bbGamma_*(\shProj(\mathcal M))
\]
est un \textbf{(TN)}-isomorphisme.
\end{theorem}

\begin{proof}
On peut se borner au cas où $\mathcal M$ est un $\mathcal S$-module de
type fini (II, 3.4.2). Comme $Y$ est quasi-compact, il existe un entier
$d>0$ tel que $\mathcal S^{(d)}$ soit engendrée par le
$\mathcal O_Y$-Module quasi-cohérent $\mathcal S_d$, ce dernier étant de
type fini (II, 3.1.10), donc cohérent puisque $Y$ est noethérien.
Remarquons maintenant que $\mathcal M$ est somme directe des
$\mathcal M^{(d,k)}$ pour $0\leq k<d$ et que chacun des
$\mathcal M^{(d,k)}$ est un $\mathcal S^{(d)}$-Module quasi-cohérent de
type fini, ainsi qu'il résulte de (II, 2.1.6, (iii)), la question étant
locale sur $Y$. Or, il suffit évidemment de prouver que chacun des
homomorphismes canoniques
\[
\alpha:\mathcal M^{(d,k)}
\longrightarrow\bbGamma_*((\shProj(\mathcal M))^{(d,k)})
\]
est un \textbf{(TN)}-isomorphisme. Compte tenu de (II, 8.14.13), et
notamment du diagramme (8.14.13.4), on voit qu'on est ramené à prouver le
théorème lorsque $\mathcal S$ est engendrée par $\mathcal S_1$ et que
$\mathcal S_1$ est un $\mathcal O_Y$-Module cohérent.

Comme $Y$ est noethérien, le même raisonnement qu'au début de
(\hyperref[III.2.2.2-fr]{2.2.2}) montre qu'on peut se borner au cas où
$Y=\operatorname{Spec}(A)$, $\mathcal S=\widetilde S$,
$\mathcal M=\widetilde M$, $A$ étant un anneau noethérien, $S_1$ un
$A$-module de type fini et $M$ un $S$-module gradué de type fini.
Montrons qu'il suffit alors de prouver le théorème lorsque $M=S$. En effet,
dans le cas général, on a une suite exacte
\[
L'\longrightarrow L\longrightarrow M\longrightarrow0,
\]
où $L$ et $L'$ sont des sommes directes de modules gradués de la forme
$S(m)$. Si le résultat est vrai pour $M=S$, il l'est aussi pour $M=S(m)$,
donc pour $L$ et $L'$. Considérons alors le diagramme commutatif
\[
\xymatrix{
\widetilde {L'_n}\ar[r]\ar[d]_{\alpha_n} &
\widetilde {L_n}\ar[r]\ar[d]^{\alpha_n} &
\widetilde {M_n}\ar[r]\ar[d]^{\alpha_n} & 0\\
q_*(\widetilde {L'}(n))\ar[r] &
q_*(\widetilde L(n))\ar[r] &
q_*(\widetilde M(n))\ar[r] & 0 .
}
\]
La deuxième ligne est exacte en vertu de
(\hyperref[III.2.2.3-fr]{2.2.3}), dès que $n$ est assez grand~; comme il
en est de même de la première et que les deux flèches verticales de gauche
sont des isomorphismes, il en est de même de la troisième.

Cela étant, pour prouver le théorème lorsque $M=S$, supposons d'abord que
$S=A[T_0,\ldots,T_r]$ ($T_i$ indéterminées)~; dans ce cas, notre assertion
n'est autre que (\hyperref[III.2.1.11-fr]{2.1.11}, (ii)). Dans le cas
général, $S$ s'identifie à un quotient d'un anneau
$S'=A[T_0,\ldots,T_r]$ par un idéal gradué, donc $X$ à un sous-schéma
fermé de $X'=\mathbf P_A^r$ (II, 2.9.2). Si $j$ est l'injection canonique
$X\to X'$, $j_*(\widetilde S(n))$ n'est autre que le
$\mathcal O_{X'}$-Module $(\shProj(\widetilde S))(n)$ où $S$ est considéré
comme un $S'$-module gradué~; cela résulte en effet aussitôt de
(II, 2.8.7). Comme $j_*(\widetilde S(n))$ est un
$\mathcal O_{X'}$-Module vérifiant \textbf{(TF)}, l'homomorphisme canonique
$\alpha_n:S_n\to\Gamma(X',j_*(\widetilde S(n)))$ est bijectif pour $n$
assez grand, en vertu de ce qui précède~; cela achève la démonstration,
puisque
$\Gamma(X',j_*(\widetilde S(n)))=\Gamma(X,\widetilde S(n))$.
\end{proof}

\oldpage[III]{104}
\begin{corollary}[2.3.2]
\label{III.2.3.2-fr}
Sous les hypothèses de (\hyperref[III.2.3.1-fr]{2.3.1}), soit
\[
\mathcal S_X=\bigoplus_{n\in\mathbb Z}\mathcal O_X(n),
\]
et soit $\mathcal F$ un $\mathcal S_X$-Module gradué quasi-cohérent de
type fini. Alors $\bbGamma_*(\mathcal F)$ vérifie la condition
\textbf{(TF)}.
\end{corollary}

\begin{proof}
On a vu dans la démonstration de
(\hyperref[III.2.3.1-fr]{2.3.1}) que $X$, qui est isomorphe à
$\operatorname{Proj}(\mathcal S^{(d)})$ (II, 3.1.8), est de type fini sur
$Y$ (II, 3.4.1). Il résulte alors de (II, 8.14.9) que $\mathcal F$ est
isomorphe à un $\mathcal S_X$-Module gradué de la forme
$\shProj(\mathcal M)$, où $\mathcal M$ est un $\mathcal S$-Module gradué
quasi-cohérent de type fini. En vertu de
(\hyperref[III.2.3.1-fr]{2.3.1}), $\bbGamma_*(\mathcal F)$ est
\textbf{(TN)}-isomorphe à $\mathcal M$ et par suite vérifie
\textbf{(TF)}.
\end{proof}

\begin{env}[2.3.3]
\label{III.2.3.3-fr}
\emph{Scholie.} --- Soient $Y$ un préschéma noethérien, $\mathcal S$ une
$\mathcal O_Y$-Algèbre graduée vérifiant les conditions de
(\hyperref[III.2.3.1-fr]{2.3.1}) et $X=\operatorname{Proj}(\mathcal S)$.
Soient $K_{\mathcal S}$ la catégorie abélienne des $\mathcal S$-Modules
gradués quasi-cohérents vérifiant \textbf{(TF)}, $K'_{\mathcal S}$ la
sous-catégorie de $K_{\mathcal S}$ formée des $\mathcal S$-Modules
vérifiant \textbf{(TN)}~; enfin, soit $K_X$ la catégorie des
$\mathcal S_X$-Modules gradués quasi-cohérents de type fini $\mathcal F$
(ce qui revient à dire, puisque $\mathcal S_X$ est périodique
(II, 8.14.4 et 8.14.12), que les $\mathcal F_i$ sont des
$\mathcal O_X$-Modules cohérents). Alors les foncteurs
\[
\mathcal M\longmapsto\shProj(\mathcal M)
\quad\text{dans }K_{\mathcal S},
\qquad
\mathcal F\longmapsto\bbGamma_*(\mathcal F)
\quad\text{dans }K_X
\]
définissent, en vertu de (II, 8.14.8 et 8.14.10) et de
(\hyperref[III.2.3.2-fr]{2.3.2}), une équivalence (T, I, 1.2) de la
catégorie quotient $K_{\mathcal S}/K'_{\mathcal S}$ (T, I, 1.11) avec la
catégorie $K_X$. Lorsque $\mathcal S$ est engendrée par $\mathcal S_1$,
on peut remplacer $K_X$ par la catégorie des $\mathcal O_X$-Modules
cohérents (II, 8.14.12).
\end{env}

\begin{proposition}[2.3.4]
\label{III.2.3.4-fr}
Soit $Y$ un préschéma noethérien.
\begin{enumerate}
\item[{\rm(i)}] Soit $\mathcal S$ une $\mathcal O_Y$-Algèbre graduée à
degrés positifs, quasi-cohérente et de type fini. Soient
$X=\operatorname{Proj}(\mathcal S)$, et
\[
\mathcal S_X=\shProj(\mathcal S)
=\bigoplus_{n\in\mathbb Z}\mathcal O_X(n).
\]
Alors $\mathcal S_X$ est une $\mathcal O_X$-Algèbre graduée périodique
(II, 8.14.12) dont les composants homogènes
$(\mathcal S_X)_n=\mathcal O_X(n)$ sont des $\mathcal O_X$-Modules
cohérents, et si $d>0$ est une période de $\mathcal S_X$,
$(\mathcal S_X)_d=\mathcal O_X(d)$ est un $\mathcal O_X$-Module
inversible $Y$-ample. En outre, l'homomorphisme canonique
\[
\alpha:\mathcal S\longrightarrow\bbGamma_*(\mathcal S_X)
\]
est un \textbf{(TN)}-isomorphisme.
\item[{\rm(ii)}] Inversement, soit $q:X\to Y$ un morphisme projectif, et
soit $\mathcal S'$ une $\mathcal O_X$-Algèbre graduée, dont les composants
homogènes $\mathcal S'_n$ $(n\in\mathbb Z)$ sont des
$\mathcal O_X$-Modules cohérents, et qui admet une période $d>0$ telle que
$\mathcal S'_d$ soit un $\mathcal O_X$-Module inversible ample pour $q$.
Alors
\[
\mathcal S=\bigoplus_{n\geq0}q_*(\mathcal S'_n)
\]
est une $\mathcal O_Y$-Algèbre graduée à degrés positifs quasi-cohérente
et de type fini, et il existe un $Y$-isomorphisme
\[
r:X\xrightarrow{\sim}\operatorname{Proj}(\mathcal S)
\]
tel que $r^*(\shProj(\mathcal S))$ soit isomorphe (en tant que
$\mathcal O_X$-Algèbre graduée) à $\mathcal S'$.
\end{enumerate}
\end{proposition}

\begin{proof}
(i) Toutes les assertions ont pratiquement déjà été démontrées, la
dernière n'étant autre qu'un cas particulier de
(\hyperref[III.2.3.2-fr]{2.3.2}). Le fait que $\mathcal S_X$ est
périodique a été vu en (II, 8.14.14) et le fait qu'il y a une période
$d>0$ telle que $\mathcal O_X(d)$ soit inversible et $Y$-ample n'est
autre que (II, 4.6.18). Enfin, pour $0\leq k<d$,
$(\mathcal S_X)^{(d,k)}$ est un $(\mathcal S_X)^{(d)}$-Module de type fini
(II, 8.14.14), donc chacun des $(\mathcal S_X)_n$ est un
$\mathcal O_X$-Module quasi-cohérent de type fini, en vertu de
(II, 2.1.6, (ii))~; la question étant locale, comme $\mathcal O_X$ est
cohérent, il en est de même des $(\mathcal S_X)_n$.

(ii) Quitte à remplacer la période $d$ par un de ses multiples, on peut
supposer que $\mathcal L=\mathcal S'_d$ est un $\mathcal O_X$-Module
\emph{très ample} relativement à $q$ (II, 4.6.11). On a en outre
\[
\mathcal S'^{(d)}=\bigoplus_{n\in\mathbb Z}\mathcal L^{\otimes n}
\quad\text{par hypothèse, donc}\quad
\mathcal S^{(d)}=\bigoplus_{n\geq0}q_*(\mathcal L^{\otimes n})~;
\]
on sait (II, 3.1.8 et 3.2.9) qu'il y a un $Y$-isomorphisme $s$ de
$X'=\operatorname{Proj}(\mathcal S)$ sur
$X''=\operatorname{Proj}(\mathcal S^{(d)})$ tel que
\oldpage[III]{105}
\[
s^*(\mathcal O_{X''}(n))=\mathcal O_{X'}(nd).
\]
On établira donc l'existence d'un $Y$-isomorphisme $X\xrightarrow{\sim}X'$
si l'on prouve la

\begin{env}[2.3.4.1]
\label{III.2.3.4.1-fr}
Soient $Y$ un préschéma noethérien, $q:X\to Y$ un morphisme projectif,
$\mathcal L$ un $\mathcal O_X$-Module inversible très ample pour $q$.
Alors $\mathcal S=\bigoplus_{n\geq0}q_*(\mathcal L^{\otimes n})$ est une
$\mathcal O_Y$-Algèbre graduée quasi-cohérente de type fini, telle que
$\mathcal S_n=\mathcal S_1^n$ pour $n$ assez grand, et il existe un
$Y$-isomorphisme $r:X\xrightarrow{\sim}P=\operatorname{Proj}(\mathcal S)$
tel que $\mathcal L=r^*(\mathcal O_P(1))$.
\end{env}

Comme $q$ est un morphisme projectif, il résulte de (II, 5.4.4 et 4.4.7)
qu'il existe un $Y$-isomorphisme
$r':X\xrightarrow{\sim}P'=\operatorname{Proj}(\mathcal T)$, où
$\mathcal T$ est une $\mathcal O_Y$-Algèbre quasi-cohérente,
$\mathcal T_1$ est de type fini et engendre $\mathcal T$, et
$\mathcal L=r'^*(\mathcal O_{P'}(1))$. Si $q':P'\to Y$ est le morphisme
structural, on a
\[
\mathcal S=\bigoplus_{n\geq0}q'_*(\mathcal O_{P'}(n)).
\]
En vertu de (\hyperref[III.2.3.1-fr]{2.3.1}), l'homomorphisme canonique
$\alpha_n:\mathcal T_n\to\mathcal S_n=q'_*(\mathcal O_{P'}(n))$ est
bijectif pour $n$ assez grand~; comme $\mathcal T_n=\mathcal T_1^n$, on a
\emph{a fortiori} $\mathcal S_n=\mathcal S_1^n$ dès que $n$ est assez
grand. En outre, l'homomorphisme canonique
$\alpha:\mathcal T\to\mathcal S$ de $\mathcal O_Y$-Algèbres graduées est
un \textbf{(TN)}-isomorphisme~; par suite
\phantomsection\label{III.2.3.5.1-fr}
\[
\Phi=\operatorname{Proj}(\alpha):
\operatorname{Proj}(\mathcal S)\longrightarrow
\operatorname{Proj}(\mathcal T)
\]
est un isomorphisme (II, 3.6.1). En outre, il résulte de (II, 3.5.2) et
du fait que les $\mathcal T$-Modules gradués obtenus à partir des
$\mathcal S(n)$ par restriction des scalaires sont \textbf{(TN)}-isomorphes
aux $\mathcal T(n)$ que
\[
\Phi_*(\mathcal O_P(n))=\mathcal O_{P'}(n)
\]
pour tout $n$ (II, 3.4.2). Pour achever de prouver
(\hyperref[III.2.3.4.1-fr]{2.3.4.1}), il reste à montrer que $\mathcal S$
est une $\mathcal O_Y$-Algèbre de type fini~; or les
$\mathcal S_n=q'_*(\mathcal O_{P'}(n))$ sont des
$\mathcal O_Y$-Modules cohérents en vertu de
(\hyperref[III.2.2.1-fr]{2.2.1}), et si
$\mathcal S_n=\mathcal S_1^n$ pour $n\geq n_0$, $\mathcal S$ est
engendrée par $\bigoplus_{i\leq n_0}\mathcal S_i$, qui est cohérent~;
d'où notre assertion (I, 9.6.2).

Revenons à la démonstration de
(\hyperref[III.2.3.4-fr]{2.3.4}), dont nous reprenons les notations. Nous
avons démontré l'existence d'un $Y$-isomorphisme $r'':X\xrightarrow{\sim}X''$
tel que $r''_*(\mathcal L^{\otimes n})=\mathcal O_{X''}(n)$ pour tout
$n\in\mathbb Z$~; nous désignerons par $q''$ le morphisme structural
$X''\to Y$. Notons maintenant que $\mathcal S'$ est somme directe des
$\mathcal S'^{(d)}$-Modules gradués $\mathcal S'^{(d,k)}$~; chacun de ces
derniers est un $\mathcal S'^{(d)}$-Module quasi-cohérent de type fini, en
vertu de la périodicité de $\mathcal S'$ et de l'hypothèse que les
$\mathcal S'_n$ sont des $\mathcal O_X$-Modules de type fini
(II, 8.14.12). Posons
$\mathcal F^{(k)}=r''_*(\mathcal S'^{(d,k)})$, de sorte que les
$\mathcal F^{(k)}$ sont des $\mathcal S_{X''}$-Modules gradués
quasi-cohérents de type fini~; par suite (II, 8.14.8), l'homomorphisme
canonique
\[
\beta:\shProj(\bbGamma_*(\mathcal F^{(k)}))
\longrightarrow\mathcal F^{(k)}
\]
est un isomorphisme de $\mathcal S_{X''}$-Modules. Mais on a
$q''_*((\mathcal F^{(k)})_n)=q_*((\mathcal S'^{(d,k)})_n)$, et pour
$n\geq0$ ce dernier $\mathcal O_Y$-Module est par définition égal à
$(\mathcal S^{(d,k)})_n$. Autrement dit, l'injection canonique
$\mathcal S^{(d,k)}\to\bbGamma_*(\mathcal F^{(k)})$ est un
\textbf{(TN)}-isomorphisme~; donc
\[
\shProj(\mathcal S^{(d,k)})
=\shProj(\bbGamma_*(\mathcal F^{(k)})),
\]
et par suite
$r''{}^*(\shProj(\mathcal S^{(d,k)}))\simeq\mathcal S'^{(d,k)}$.
Il reste à remarquer que, à un isomorphisme canonique près,
\[
\shProj(\mathcal S^{(d,k)})
=s_*((\shProj(\mathcal S))^{(d,k)})
\]
(II, 8.14.13.1) pour avoir démontré l'isomorphisme de
$r^*(\shProj(\mathcal S))$ et de $\mathcal S'$.

Enfin, en vertu de (\hyperref[III.2.3.2-fr]{2.3.2}), chacun des
$\bbGamma_*(\mathcal F^{(k)})$ vérifie la condition \textbf{(TF)}, donc il
en est de même de chacun des $\mathcal S^{(d,k)}$~; en outre, les
$\mathcal S_n=q_*(\mathcal S'_n)$ sont cohérents en vertu de
(\hyperref[III.2.2.1-fr]{2.2.1}), puisque les $\mathcal S'_n$ sont
cohérents. On en conclut aussitôt que les $\mathcal S^{(d,k)}$ sont des
$\mathcal S^{(d)}$-Modules de type fini~; comme
(\hyperref[III.2.3.4.1-fr]{2.3.4.1}) montre que $\mathcal S^{(d)}$ est une
$\mathcal O_Y$-Algèbre de type fini, il en est de même de $\mathcal S$.
\end{proof}

\oldpage[III]{106}
\begin{proposition}[2.3.5]
\label{III.2.3.5-fr}
Soient $Y$ un préschéma intègre noethérien, $X$ un préschéma intègre,
$f:X\to Y$ un morphisme projectif birationnel. Il existe alors un Idéal
fractionnaire cohérent $\mathcal J\subset\mathcal R(Y)$ (II, 8.1.2) tel
que $X$ soit $Y$-isomorphe au préschéma obtenu en faisant éclater
$\mathcal J$ (II, 8.1.3). En outre, il existe un ouvert $U$ de $Y$ tel que
la restriction de $f$ à $f^{-1}(U)$ soit un isomorphisme de $f^{-1}(U)$
sur $U$ (cf. I, 6.5.5), et que $\mathcal J|U$ soit inversible.
\end{proposition}

\begin{proof}
Comme il existe un $\mathcal O_X$-Module inversible $\mathcal L$ très
ample pour $f$ (II, 4.4.2 et 5.3.2), on peut appliquer
(\hyperref[III.2.3.4.1-fr]{2.3.4.1}), et on voit que $X$ s'identifie à
$\operatorname{Proj}(\mathcal S)$, où
$\mathcal S=\bigoplus_{n\geq0}f_*(\mathcal L^{\otimes n})$. On sait en
outre que les $f_*(\mathcal L^{\otimes n})$ sont des
$\mathcal O_Y$-Modules sans torsion (I, 7.4.5), donc il en est de même du
$\mathcal O_Y$-Module $\mathcal S$, et par suite $\mathcal S$ s'identifie
canoniquement à un sous-$\mathcal O_Y$-Module de
$\mathcal S\otimes_{\mathcal O_Y}\mathcal R(Y)$ (I, 7.4.1)~; ce dernier
est un faisceau \emph{simple} (I, 7.3.6) qui est connu lorsqu'on connaît
sa restriction à un ouvert non vide, par exemple à un ouvert non vide
$U'\subset U$ tel que $\mathcal L|f^{-1}(U')$ soit isomorphe à
$\mathcal O_X|f^{-1}(U')$. Comme par hypothèse les
$f_*(\mathcal L^{\otimes n})|U'$ sont alors isomorphes à
$\mathcal O_Y|U'$, on voit que
$\mathcal S\otimes_{\mathcal O_Y}\mathcal R(Y)$ est un
$\mathcal R(Y)$-Module isomorphe à $\mathcal R(Y)[T]$, où $T$ est une
indéterminée, et $\mathcal S$ est \textbf{(TN)}-isomorphe à la
sous-$\mathcal O_Y$-Algèbre engendrée par l'image canonique de
$f_*(\mathcal L)$ dans $\mathcal S\otimes_{\mathcal O_Y}\mathcal R(Y)$
(\hyperref[III.2.3.4.1-fr]{2.3.4.1})~; mais si on identifie ce dernier
faisceau à $\mathcal R(Y)[T]$, l'image de $f_*(\mathcal L)$ s'identifie à
$\mathcal J T$, où $\mathcal J$ est un sous-$\mathcal O_Y$-Module
cohérent (\hyperref[III.2.2.1-fr]{2.2.1}) de $\mathcal R(Y)$, dont la
restriction à $U'$ est isomorphe à $\mathcal O_Y|U'$, et qui par suite
est tel que $\mathcal J|U$ soit inversible. On voit alors que $\mathcal S$
est \textbf{(TN)}-isomorphe à $\bigoplus_{n\geq0}\mathcal J^n$, ce qui
achève la démonstration.
\end{proof}

\begin{corollary}[2.3.6]
\label{III.2.3.6-fr}
Sous les hypothèses de (\hyperref[III.2.3.5-fr]{2.3.5}), supposons en
outre que, pour tout sous-$\mathcal O_Y$-Module cohérent
$\mathcal J\neq0$ de $\mathcal R(Y)$, il existe un
$\mathcal O_Y$-Module inversible $\mathcal L$ tel que
\[
\Gamma(Y,\mathcal L\otimes_{\mathcal O_Y}
\mathcal{H}om(\mathcal J,\mathcal O_Y))\neq0~;
\]
alors, dans l'énoncé de (\hyperref[III.2.3.5-fr]{2.3.5}), on peut supposer
que $\mathcal J$ est un Idéal de $\mathcal O_Y$. Cette condition
supplémentaire est toujours vérifiée s'il existe un
$\mathcal O_Y$-Module ample.
\end{corollary}

\begin{proof}
En effet, on a (0\textsubscript{I}, 5.4.2)
\[
\mathcal L\otimes\mathcal{H}om(\mathcal J,\mathcal O_Y)
=\mathcal{H}om(\mathcal L^{-1},\mathcal{H}om(\mathcal J,\mathcal O_Y))
=\mathcal{H}om(\mathcal J\otimes\mathcal L^{-1},\mathcal O_Y)~;
\]
l'hypothèse signifie donc qu'il y a un homomorphisme non nul $u$ de
$\mathcal J\otimes\mathcal L^{-1}$ dans $\mathcal O_Y$. Comme, pour tout
$y\in Y$, $(\mathcal J\otimes\mathcal L^{-1})_y$ s'identifie à un
sous-$\mathcal O_y$-Module du corps des fractions $\mathcal R(Y)_y$ de
$\mathcal O_y$ (I, 7.1.5), $u_y$ est nécessairement injectif, donc $u$
est un isomorphisme de $\mathcal J\otimes\mathcal L^{-1}$ sur un Idéal
$\mathcal J'$ de $\mathcal O_Y$. Mais comme
\[
\operatorname{Proj}\left(\bigoplus_{n\geq0}\mathcal J^n\right)
\quad\text{et}\quad
\operatorname{Proj}\left(
\bigoplus_{n\geq0}(\mathcal J\otimes\mathcal L^{-1})^n\right)
\]
sont $Y$-isomorphes (II, 3.1.8), cela prouve la première assertion du
corollaire. Pour démontrer la seconde, notons que
$\mathcal F=\mathcal{H}om_{\mathcal O_Y}(\mathcal J,\mathcal O_Y)$ est
cohérent et $\neq0$, puisqu'il existe un ouvert $U$ de $Y$ tel que
$\mathcal J|U$ soit inversible. Si $\mathcal L$ est un
$\mathcal O_Y$-Module ample, il existe un entier $n$ tel que
$\mathcal F(n)=\mathcal F\otimes\mathcal L^{\otimes n}$ soit engendré par
ses sections au-dessus de $Y$ (II, 4.5.5)~; \emph{a fortiori}, on a
$\Gamma(Y,\mathcal F(n))\neq0$, d'où la conclusion.
\end{proof}

\begin{corollary}[2.3.7]
\label{III.2.3.7-fr}
Soient $X$ et $Y$ deux schémas intègres, projectifs sur un corps $k$, et
soit $f:X\to Y$ un $k$-morphisme birationnel. Alors $X$ est $k$-isomorphe
à un $Y$-schéma obtenu en faisant éclater un sous-schéma fermé $Y'$
(non nécessairement réduit) de $Y$.
\end{corollary}

\oldpage[III]{107}
\begin{proof}
En effet, $f$ est projectif (II, 5.5.5, (v)) et comme $Y$ est projectif
sur $k$, la condition supplémentaire de
(\hyperref[III.2.3.6-fr]{2.3.6}) est vérifiée~; il suffit alors de
considérer le sous-schéma fermé $Y'$ de $Y$ défini par l'Idéal cohérent
$\mathcal J$ du cor. (\hyperref[III.2.3.6-fr]{2.3.6}).
\end{proof}

\begin{env}[2.3.8]
\label{III.2.3.8-fr}
\emph{Remarque.} Au chap. IV, en étudiant la notion de diviseur, nous
verrons que si, dans l'énoncé de
(\hyperref[III.2.3.5-fr]{2.3.5}), on suppose que les anneaux
$\mathcal O_y$ $(y\in Y)$ sont factoriels (ce qui est le cas par exemple
si $Y$ est non singulier), alors $X$ peut se déduire de $Y$ en faisant
éclater un sous-préschéma fermé $Y'$ de $Y$ dont l'espace sous-jacent est
contenu dans $Y-U$.
\end{env}

\subsection{Une généralisation du théorème fondamental}
\label{subsection:III.2.4-fr}

\begin{theorem}[2.4.1]
\label{III.2.4.1-fr}
Soient $Y$ un préschéma noethérien, $\mathcal S$ une
$\mathcal O_Y$-Algèbre quasi-cohérente de type fini. Soient
$f:X\to Y$ un morphisme projectif, $\mathcal S'=f^*(\mathcal S)$,
$\mathcal M$ un $\mathcal S'$-Module quasi-cohérent de type fini.
Alors~:
\begin{enumerate}
\item[(i)] Pour tout $p\in\mathbb Z$, $R^pf_*(\mathcal M)$ est un
$\mathcal S$-Module de type fini.
\item[(ii)] Soit de plus $\mathcal L$ un $\mathcal O_X$-Module inversible
ample pour $f$, et posons
$\mathcal M(n)=\mathcal M\otimes\mathcal L^{\otimes n}$ pour tout
$n\in\mathbb Z$. Il existe un entier $N$ tel que, pour $n\geq N$, on ait
\[
R^pf_*(\mathcal M(n))=0
\tag{2.4.1.1}\label{III.2.4.1.1-fr}
\]
pour tout $p>0$, et que l'homomorphisme canonique
\[
f^*(f_*(\mathcal M(n)))\longrightarrow\mathcal M(n)
\]
(0\textsubscript{I}, 4.4.4.3) soit surjectif.
\end{enumerate}
\end{theorem}

\begin{proof}
Posons $Y'=\operatorname{Spec}(\mathcal S)$,
$X'=\operatorname{Spec}(\mathcal S')$ de sorte que
$X'=X\times_Y Y'$ (II, 1.5.5)~; soient $g:Y'\to Y$,
$g':X'\to X$ les morphismes structuraux, qui sont affines par
définition, et $f'=f_{(Y')}:X'\to Y'$~; on a donc un diagramme
commutatif
\[
\xymatrix{
X\ar[d]_f & X'\ar[l]_{g'}\ar[d]^{f'}\ar[dl]^{h}\\
Y & Y'\ar[l]^{g}
}
\]
et le morphisme $f'$ est projectif (II, 5.5.5, (iii))~; posons
$h=f\circ g'=g\circ f'$.

(i) Soit $\widetilde{\mathcal M}$ le $\mathcal O_{X'}$-Module associé au
$\mathcal S'$-Module quasi-cohérent $\mathcal M$, quand $X'$ est
considéré comme un $X$-schéma affine (II, 1.4.3), de sorte que l'on a
$\mathcal M=g'_*(\widetilde{\mathcal M})$~; comme $\mathcal M$ est un
$\mathcal S'$-Module de type fini, $\widetilde{\mathcal M}$ est un
$\mathcal O_{X'}$-Module de type fini (II, 1.4.5)~; comme $h$ est de type
fini, puisque $g$ et $f'$ le sont (II, 1.3.7 et I, 6.3.4, (ii)), $X'$ est
noethérien (I, 6.3.7) et $\widetilde{\mathcal M}$ est par suite cohérent.
Cela étant, comme $g'$ est affine, l'homomorphisme canonique
$R^pf_*(\mathcal M)\to R^ph_*(\widetilde{\mathcal M})$ est bijectif
(\hyperref[III.1.3.4-fr]{1.3.4}). En outre, cet homomorphisme est un
homomorphisme de $\mathcal S$-Modules~; en effet, de l'homomorphisme
canonique
\[
g'_*(\mathcal O_{X'})\otimes_{\mathcal O_X}
g'_*(\widetilde{\mathcal M})\longrightarrow
g'_*(\widetilde{\mathcal M})
\tag{2.4.1.2}\label{III.2.4.1.2-fr}
\]
qui définit la structure de $\mathcal S'$-Module de $\mathcal M$ (en se
rappelant que $\mathcal S'=g'_*(\mathcal O_{X'})$), on déduit
canoniquement un homomorphisme
\[
f_*(g'_*(\mathcal O_{X'}))\otimes
R^pf_*(g'_*(\widetilde{\mathcal M}))\longrightarrow
R^pf_*(g'_*(\widetilde{\mathcal M})).
\]

\oldpage[III]{108}
(0, 12.2.2), et comme (\hyperref[III.2.4.1.2-fr]{2.4.1.2}) provient
lui-même (par application de (0\textsubscript{I}, 4.2.2.1)) de
l'homomorphisme
\[
\mathcal O_{X'}\otimes\widetilde{\mathcal M}
\longrightarrow\widetilde{\mathcal M}
\]
définissant la structure de $\mathcal O_{X'}$-Module de
$\widetilde{\mathcal M}$, le diagramme
\[
\xymatrix{
f_*(g'_*(\mathcal O_{X'}))\otimes
R^pf_*(g'_*(\widetilde{\mathcal M}))
\ar[r]\ar[d] &
R^pf_*(g'_*(\widetilde{\mathcal M}))\ar[d]\\
h_*(\mathcal O_{X'})\otimes R^ph_*(\widetilde{\mathcal M})
\ar[r] & R^ph_*(\widetilde{\mathcal M})
}
\]
est commutatif (0, 12.2.6)~; composant les flèches horizontales avec
l'homomorphisme provenant de l'homomorphisme canonique
\[
\mathcal S\longrightarrow f_*(f^*(\mathcal S))=f_*(\mathcal S')
=f_*(g'_*(\mathcal O_{X'}))=h_*(\mathcal O_{X'}),
\]
on obtient notre assertion. D'autre part, puisque $g$ est affine et que
$f'$ est séparé et quasi-compact, l'homomorphisme canonique
$R^ph_*(\widetilde{\mathcal M})\to
g_*(R^pf'_*(\widetilde{\mathcal M}))$ est bijectif
(\hyperref[III.1.4.14-fr]{1.4.14}), et on démontre comme ci-dessus que
c'est un isomorphisme de $\mathcal S$-Modules (en utilisant cette fois la
commutativité de (0, 12.2.6.2)). Or, $f'$ étant projectif et
$\widetilde{\mathcal M}$ cohérent,
$R^pf'_*(\widetilde{\mathcal M})$ est un
$\mathcal O_{Y'}$-Module cohérent en vertu de
(\hyperref[III.2.2.1-fr]{2.2.1})~; on en conclut que
$g_*(R^pf'_*(\widetilde{\mathcal M}))$ est un $\mathcal S$-Module de type
fini (II, 1.4.5).

(ii) Soit $\mathcal L'=g'^*(\mathcal L)$, qui est un
$\mathcal O_{X'}$-Module inversible~; pour tout $n\in\mathbb Z$, on a
\[
g'_*(\widetilde{\mathcal M}\otimes\mathcal L'^{\otimes n})
=g'_*(\widetilde{\mathcal M})\otimes\mathcal L^{\otimes n}
=\mathcal M(n)
\]
(0\textsubscript{I}, 5.4.10) à un isomorphisme près~; on peut appliquer
à $\widetilde{\mathcal M}\otimes\mathcal L'^{\otimes n}$ le
raisonnement fait dans (i) pour $\widetilde{\mathcal M}$, qui prouve que
\[
R^pf_*(g'_*(\widetilde{\mathcal M}\otimes\mathcal L'^{\otimes n}))
\quad\hbox{est isomorphe à}\quad
g_*(R^pf'_*(\widetilde{\mathcal M}\otimes\mathcal L'^{\otimes n})).
\]
Or $\mathcal L'$ est ample pour $f'$ (II, 4.6.13, (iii)), donc il résulte
de (\hyperref[III.2.2.1-fr]{2.2.1}) qu'il existe un entier $N$ tel que
\[
R^pf'_*(\widetilde{\mathcal M}\otimes\mathcal L'^{\otimes n})=0
\]
pour tout $p$ et tout $n\geq N$, ce qui prouve
(\hyperref[III.2.4.1.1-fr]{2.4.1.1}). Enfin, il résulte encore de
(\hyperref[III.2.2.1-fr]{2.2.1}) qu'on peut supposer $N$ tel que pour
$n\geq N$, l'homomorphisme canonique
\[
f'^*(f'_*(\widetilde{\mathcal M}\otimes\mathcal L'^{\otimes n}))
\longrightarrow\widetilde{\mathcal M}\otimes\mathcal L'^{\otimes n}
\]
soit surjectif~; comme $g'_*$ est un foncteur exact (II, 1.4.4),
l'homomorphisme correspondant
\[
g'_*(f'^*(f'_*(\widetilde{\mathcal M}\otimes
\mathcal L'^{\otimes n})))
\longrightarrow
g'_*(\widetilde{\mathcal M}\otimes\mathcal L'^{\otimes n})
=\mathcal M(n)
\]
est surjectif. Or, on a
\[
g'_*(f'^*(f'_*(\widetilde{\mathcal M}\otimes
\mathcal L'^{\otimes n})))
=f^*(g_*(f'_*(\widetilde{\mathcal M}\otimes
\mathcal L'^{\otimes n})))
\]
(II, 1.5.2) et comme $g_*\circ f'_* = f_*\circ g'_*$, on voit
finalement que l'on a
\[
g'_*(f'^*(f'_*(\widetilde{\mathcal M}\otimes\mathcal L'^{\otimes n})))
=f^*(f_*(g'_*(\widetilde{\mathcal M}\otimes
\mathcal L'^{\otimes n})))=f^*(f_*(\mathcal M(n))),
\]
ce qui achève la démonstration.
\end{proof}

\begin{env}[2.4.2]
\label{III.2.4.2-fr}
Nous aurons en particulier à appliquer
(\hyperref[III.2.4.1-fr]{2.4.1}) lorsque $\mathcal S$ est une
$\mathcal O_Y$-Algèbre graduée à degrés positifs,
$\mathcal M=\bigoplus_{k\in\mathbb Z}\mathcal M_k$ un
$\mathcal S'$-Module gradué. Alors (avec les

\oldpage[III]{109}
mêmes hypothèses de finitude sur $\mathcal S$ et $\mathcal M$) on conclut
de (\hyperref[III.2.4.1-fr]{2.4.1}) que
\[
\bigoplus_{k\in\mathbb Z}R^pf_*(\mathcal M_k)
\]
est un $\mathcal S$-Module de type fini pour tout $p$, et (sous les
hypothèses supplémentaires de (\hyperref[III.2.4.1-fr]{2.4.1}), (ii))
qu'il existe $N$ tel que pour $n\geq N$, on ait
$R^pf_*(\mathcal M_k(n))=0$ pour tout $p>0$ et tout $k\in\mathbb Z$, et
que l'homomorphisme canonique
$f^*(f_*(\mathcal M_k(n)))\to\mathcal M_k(n)$ soit surjectif pour tout
$k\in\mathbb Z$.
\end{env}

\subsection{Caractéristique d'Euler-Poincaré et polynôme de Hilbert}
\label{subsection:III.2.5-fr}

\begin{env}[2.5.1]
\label{III.2.5.1-fr}
Soient $A$ un anneau artinien, $X$ un $A$-schéma projectif sur
$Y=\operatorname{Spec}(A)$. Pour tout $\mathcal O_X$-Module cohérent
$\mathcal F$, les $H^i(X,\mathcal F)$ $(i\geq0)$ sont des $A$-modules de
type fini (\hyperref[III.2.2.1-fr]{2.2.1}), donc ici de longueur finie
puisque $A$ est artinien. On sait en outre
(\hyperref[III.2.2.1-fr]{2.2.1}) que $H^i(X,\mathcal F)=0$ sauf pour un
nombre fini de valeurs de $i\geq0$~; le nombre entier
\[
\chi_A(\mathcal F)=\sum_{i=0}^{\infty}(-1)^i
\operatorname{long}(H^i(X,\mathcal F))
\tag{2.5.1.1}\label{III.2.5.1.1-fr}
\]
est donc défini pour tout $\mathcal O_X$-Module cohérent $\mathcal F$.
Lorsque $A$ est un anneau local artinien, on dit que $\chi_A(\mathcal F)$
est la caractéristique d'Euler-Poincaré de $\mathcal F$ (par rapport à
l'anneau $A$). Pour $\mathcal F=\mathcal O_X$, on dit que
$\chi_A(\mathcal O_X)$ est le genre arithmétique de $X$ (par rapport à
$A$).
\end{env}

\begin{proposition}[2.5.2]
\label{III.2.5.2-fr}
Soit
\[
0\longrightarrow\mathcal F'\longrightarrow\mathcal F
\longrightarrow\mathcal F''\longrightarrow0
\]
une suite exacte de $\mathcal O_X$-Modules cohérents~; on a alors
\[
\chi_A(\mathcal F)=\chi_A(\mathcal F')+\chi_A(\mathcal F'').
\tag{2.5.2.1}\label{III.2.5.2.1-fr}
\]
\end{proposition}

\begin{proof}
Comme les modules de cohomologie de $\mathcal F$, $\mathcal F'$,
$\mathcal F''$ sont nuls sauf un nombre fini d'entre eux, il y a un
entier $r>0$ tel que la suite exacte de cohomologie s'écrive
\[
0\to H^0(X,\mathcal F')\to H^0(X,\mathcal F)
\to H^0(X,\mathcal F'')\to H^1(X,\mathcal F')\to\cdots
\]
\[
\cdots\to H^r(X,\mathcal F')\to H^r(X,\mathcal F)
\to H^r(X,\mathcal F'')\to0.
\]
Or, on sait que dans une suite exacte de $A$-modules de longueur finie,
ayant des 0 aux deux extrémités, la somme alternée des longueurs est
nulle (0, 11.10.1)~; appliquant ce résultat, on trouve immédiatement la
formule (\hyperref[III.2.5.2.1-fr]{2.5.2.1}).
\end{proof}

On notera que le résultat de (\hyperref[III.2.5.2-fr]{2.5.2}) s'applique
toutes les fois que l'on sait que $X$ est un $A$-schéma quasi-compact et
que les $A$-modules $H^i(X,\mathcal F)$ sont de type fini pour tout
$\mathcal O_X$-Module cohérent $\mathcal F$
(\hyperref[III.1.4.12-fr]{1.4.12}).

\begin{theorem}[2.5.3]
\label{III.2.5.3-fr}
Soient $A$ un anneau local artinien, $X$ un schéma projectif sur
$Y=\operatorname{Spec}(A)$, $\mathcal L$ un $\mathcal O_X$-Module
inversible très ample relativement à $Y$, $\mathcal F$ un
$\mathcal O_X$-Module cohérent~; on pose
$\mathcal F(n)=\mathcal F\otimes_{\mathcal O_X}\mathcal L^{\otimes n}$
pour tout $n\in\mathbb Z$.
\begin{enumerate}
\item[(i)] Il existe un polynôme unique $P\in\mathbb Q[T]$ tel que
$\chi_A(\mathcal F(n))=P(n)$ pour tout $n\in\mathbb Z$ (on dit que $P$
est le polynôme de Hilbert de $\mathcal F$ par rapport à $A$).
\item[(ii)] Pour $n$ assez grand, on a
$\chi_A(\mathcal F(n))=\operatorname{long}_A
\Gamma(X,\mathcal F(n))$.
\item[(iii)] Le terme de plus haut degré de $\chi_A(\mathcal F(n))$ a un
coefficient $\geq0$.
\end{enumerate}
\end{theorem}

Ajoutons qu'au chap. IV, dans le paragraphe consacré à la notion de
dimension, nous démontrerons en outre que le degré de
$\chi_A(\mathcal F(n))$ est égal à la dimension du support de
$\mathcal F$.

\begin{proof}
Comme on a $H^i(X,\mathcal F(n))=0$ pour tout $i>0$ dès que $n$ est
assez grand (\hyperref[III.2.2.1-fr]{2.2.1}),

\oldpage[III]{110}
on a
\[
\chi_A(\mathcal F(n))=\operatorname{long}H^0(X,\mathcal F(n))
=\operatorname{long}\Gamma(X,\mathcal F(n))
\]
pour $n$ assez grand, d'où (ii)~; cela entraîne
$\chi_A(\mathcal F(n))\geq0$ pour $n$ assez grand, et (iii) résulte donc
de (i)~; comme d'ailleurs l'assertion d'unicité de (i) est immédiate, il
reste à prouver l'existence du polynôme $P$.

Montrons d'abord qu'on peut supposer que $\mathfrak m\mathcal F=0$, où
$\mathfrak m$ est l'idéal maximal de $A$. En effet, il existe un entier
$s>0$ tel que $\mathfrak m^s=0$, et $\mathcal F(n)$ admet donc une
filtration finie
\[
\mathcal F(n)\supset\mathfrak m\mathcal F(n)\supset\cdots
\supset\mathfrak m^{s-1}\mathcal F(n)\supset0.
\]
Par récurrence, on déduit de
(\hyperref[III.2.5.1-fr]{2.5.1}) que
\[
\chi_A(\mathcal F(n))=\sum_{k=1}^{s}
\chi_A(\mathfrak m^{k-1}\mathcal F(n)/\mathfrak m^k\mathcal F(n))~;
\]
comme
$\mathfrak m^{k-1}\mathcal F(n)/\mathfrak m^k\mathcal F(n)
=\mathcal F'_k(n)$, où
$\mathcal F'_k=\mathfrak m^{k-1}\mathcal F/\mathfrak m^k\mathcal F$,
cela prouve notre assertion.

Supposons donc $\mathfrak m\mathcal F=0$~; si $X'$ est le sous-schéma
fermé de $X$, image réciproque par le morphisme structural
$X\to\operatorname{Spec}(A)$ de l'unique point fermé de
$\operatorname{Spec}(A)$, et $j:X'\to X$ l'injection canonique, on a
$\mathcal F=j_*(\mathcal F')$, où $\mathcal F'$ est un
$\mathcal O_{X'}$-Module cohérent~; $X'$ est un schéma projectif sur
$\operatorname{Spec}(K)$, où $K=A/\mathfrak m$. Si
$\mathcal L'=j^*(\mathcal L)$, $\mathcal L'$ est très ample relativement
à $\operatorname{Spec}(K)$ (II, 4.4.10), et on a
$\mathcal F(n)=j_*(\mathcal F'(n))$, où
$\mathcal F'(n)=\mathcal F'\otimes_{\mathcal O_{X'}}
\mathcal L'^{\otimes n}$ (0\textsubscript{I}, 5.4.10). On en conclut que
$\chi_A(\mathcal F(n))=\chi_K(\mathcal F'(n))$ (G, II, 4.9.1), et on est
donc ramené au cas où $A$ est un corps.

Notons maintenant que $X$ peut être considéré comme un sous-schéma
fermé de $P=\mathbf P_A^r$ pour un $r$ convenable (II, 5.5.4, (ii))~;
si $i:X\to P$ est l'injection canonique, on voit comme ci-dessus que l'on
a $\chi_A(\mathcal F(n))=\chi_A(i_*(\mathcal F)(n))$, de sorte que l'on
peut se borner au cas où $X=\mathbf P_A^r=\operatorname{Proj}(S)$ avec
$S=A[T_0,\ldots,T_r]$, $A$ étant un corps.

Cela étant, on a $\mathcal F=\widetilde M$, où $M$ est un $S$-module
gradué de type fini (II, 2.7.8)~; il existe par suite une résolution finie
de $M$ par des $S$-modules gradués libres de type fini
\[
0\to L_q\to L_{q-1}\to\cdots\to L_1\to M\to0
\]
en vertu du th. des syzygies de Hilbert (M, VIII, 6.5)~; comme
$M\mapsto\widetilde M$ est un foncteur exact en $M$ (II, 2.5.4), on a
aussi une suite exacte
\[
0\to\widetilde L_q\to\widetilde L_{q-1}\to\cdots
\to\widetilde L_1\to\widetilde M\to0
\]
et par suite, pour tout $n\in\mathbb Z$, la suite
\[
0\to\widetilde L_q(n)\to\widetilde L_{q-1}(n)\to\cdots
\to\widetilde L_1(n)\to\widetilde M(n)\to0
\]
est exacte~; appliquant par récurrence sur $q$ la prop. (2.5.1), il vient
\[
\chi_A(\widetilde M(n))=\sum_{j=1}^{q}(-1)^{j+1}
\chi_A(\widetilde L_j(n))
\]
et pour prouver (i), on est donc ramené au cas où $M$ est libre et gradué
de type fini, donc au cas où $M=S(h)$ pour un $h\in\mathbb Z$. Comme
alors $\widetilde M(n)=(M(n))^\sim=(S(n+h))^\sim$ (II, 2.5.15), on voit
finalement que le théorème résultera du
\end{proof}

\oldpage[III]{111}
\begin{lemma}[2.5.3.1]
\label{III.2.5.3.1-fr}
Soient $A$ un corps, $r$ un entier $>0$, et $X=\mathbf P_A^r$~; on a
alors
\[
\chi_A(\mathcal O_X(n))=\binom{n+r}{r}
\]
pour tout $n\in\mathbb Z$.
\end{lemma}

\begin{proof}
En effet, pour $n\geq0$, on a
$\chi_A(\mathcal O_X(n))=\operatorname{long}H^0(X,\mathcal O_X(n))$, qui
est le nombre de monômes en les $T_i$ de degré total $n$, c'est-à-dire
$\binom{n+r}{r}$ (\hyperref[III.2.1.12-fr]{2.1.12}). Pour
$n\leq-r-1$, on a de même
\[
\chi_A(\mathcal O_X(n))=(-1)^r
\operatorname{long}H^r(X,\mathcal O_X(n))~;
\]
si $n=-r-h$, la dimension de $H^r(X,\mathcal O_X(n))$ sur $A$ est le
nombre des suites $(p_i)_{0\leq i\leq r}$ d'entiers $p_i>0$ tels que
$\sum_{i=0}^{r}p_i=r+h$ (\hyperref[III.2.1.12-fr]{2.1.12}), ou encore le
nombre des suites d'entiers $q_i\geq0$ $(0\leq i\leq r)$ tels que
$\sum_{i=0}^{r}q_i=h-1$~; c'est donc le nombre
\[
\binom{h+r-1}{r}=(-1)^r\binom{n+r}{r}.
\]
Enfin, pour $-r\leq n<0$, on a $\binom{n+r}{r}=0$ et d'autre part
$H^i(X,\mathcal O_X(n))=0$ pour tout $i\geq0$
(\hyperref[III.2.1.12-fr]{2.1.12}), ce qui prouve le lemme.
\end{proof}

\begin{corollary}[2.5.4]
\label{III.2.5.4-fr}
Soient $A$ un anneau local artinien, $S$ une $A$-algèbre graduée de type
fini engendrée par $S_1$, $M$ un $S$-module gradué de type fini,
$X=\operatorname{Proj}(S)$. On a alors
\[
\chi_A(\widetilde M(n))=\operatorname{long}M_n
\]
pour $n$ assez grand.
\end{corollary}

\begin{proof}
Cela résulte de ce que $M_n$ et $\Gamma(X,\widetilde M(n))$ sont
isomorphes pour $n$ assez grand
(\hyperref[III.2.3.1-fr]{2.3.1}).
\end{proof}

\subsection{Application : critères d'amplitude}
\label{subsection:III.2.6-fr}

\begin{proposition}[2.6.1]
\label{III.2.6.1-fr}
Soient $Y$ un préschéma noethérien, $f:X\to Y$ un morphisme propre,
$\mathcal L$ un $\mathcal O_X$-Module inversible. Les conditions
suivantes sont équivalentes~:
\begin{enumerate}
\item[a)] $\mathcal L$ est ample pour $f$.
\item[b)] Pour tout $\mathcal O_X$-Module cohérent $\mathcal F$, il
existe un entier $N$ tel que pour $n\geq N$, on ait
\[
R^qf_*(\mathcal F\otimes\mathcal L^{\otimes n})=0
\]
pour tout $q>0$.
\item[c)] Pour tout Idéal cohérent $\mathcal J$ de $\mathcal O_X$, il
existe un entier $N$ tel que pour $n\geq N$ on ait
\[
R^1f_*(\mathcal J\otimes\mathcal L^{\otimes n})=0.
\]
\end{enumerate}
\end{proposition}

\begin{proof}
On a vu que a) entraîne b) (\hyperref[III.2.2.1-fr]{2.2.1}, (iii)). Il
est trivial que b) entraîne c), et il reste à prouver que c) implique
a). On peut se borner au cas où $Y$ est affine (II, 4.6.4), et prouver
dans ce cas que $\mathcal L$ est ample~; il suffira de montrer que lorsque
$h$ parcourt l'ensemble des sections des $\mathcal L^{\otimes n}$
$(n>0)$ au-dessus de $X$, ceux des $X_h$ qui sont affines forment un
recouvrement de $X$ (II, 4.5.2). Pour cela, montrons que pour tout point
fermé $x$ de $X$ et tout voisinage ouvert affine $U$ de $x$, il existe un
$n$ et un $h\in\Gamma(X,\mathcal L^{\otimes n})$ tels que
$x\in X_h\subset U$~; $X_h$ sera nécessairement affine (I, 1.3.6) et la
réunion de ces $X_h$ sera un ouvert de $X$ contenant tous les points
fermés de $X$, et par suite $X$ lui-même puisque $X$ est noethérien
(I, 6.3.7 et 0\textsubscript{I}, 2.1.3). Soit $\mathcal J$ (resp.
$\mathcal J'$) le faisceau quasi-cohérent d'idéaux de $\mathcal O_X$
définissant le sous-préschéma fermé réduit de $X$ ayant pour espace
sous-jacent $X-U$ (resp. $(X-U)\cup\{x\}$) (I, 5.2.1)~; il est clair que
$\mathcal J$ et $\mathcal J'$ sont cohérents (I, 6.1.1), que
$\mathcal J'\subset\mathcal J$ et que
$\mathcal J''=\mathcal J/\mathcal J'$ est un $\mathcal O_X$-Module
cohérent (0\textsubscript{I}, 5.3.3) ayant pour support $\{x\}$ et tel
que $\mathcal J''_x=k(x)$. Comme $\mathcal L$ est localement libre, la
suite
\[
0\to\mathcal J'\otimes\mathcal L^{\otimes n}
\to\mathcal J\otimes\mathcal L^{\otimes n}
\to\mathcal J''\otimes\mathcal L^{\otimes n}\to0
\]
est exacte pour tout $n$, et par hypothèse il existe $n$ assez grand tel
que $H^1(X,\mathcal J'\otimes\mathcal L^{\otimes n})=0$~; la suite
exacte de cohomologie

\oldpage[III]{112}
prouve donc que l'homomorphisme
\[
\Gamma(X,\mathcal J\otimes\mathcal L^{\otimes n})
\longrightarrow
\Gamma(X,\mathcal J''\otimes\mathcal L^{\otimes n})
\]
est surjectif. Une section $g$ de
$\mathcal J''\otimes\mathcal L^{\otimes n}$ au-dessus de $X$ telle que
$g(x)\neq0$ est donc l'image d'une section
$h\in\Gamma(X,\mathcal J\otimes\mathcal L^{\otimes n})
\subset\Gamma(X,\mathcal L^{\otimes n})$ (car en vertu de
(0\textsubscript{I}, 5.4.1),
$\mathcal J\otimes\mathcal L^{\otimes n}$ est un sous-$\mathcal O_X$-Module
de $\mathcal L^{\otimes n}$)~; on a par définition $h(x)\neq0$ et
$h(z)=0$ pour $z\notin U$, ce qui achève la démonstration.
\end{proof}

\begin{proposition}[2.6.2]
\label{III.2.6.2-fr}
Soient $Y$ un préschéma noethérien, $f:X\to Y$ un morphisme de type
fini, $g:X'\to X$ un morphisme fini surjectif, $\mathcal L$ un
$\mathcal O_X$-Module inversible et $\mathcal L'=g^*(\mathcal L)$. On
suppose vérifiée la condition suivante~: Il existe une partie $Z$ de
$X$, propre sur $Y$ (II, 5.4.10) telle que pour tout $x\in X-Z$, ou bien
$X$ est normal au point $x$, ou bien $(g_*(\mathcal O_{X'}))_x$ est un
$\mathcal O_x$-module libre. Dans ces conditions, pour que $\mathcal L$
soit ample pour $f$, il faut et il suffit que $\mathcal L'$ soit ample
pour $f\circ g$.
\end{proposition}

\begin{proof}
\begin{env}[2.6.2.1]
\label{III.2.6.2.1-fr}
Puisque $g$ est affine, la condition est nécessaire (II, 5.1.12). Pour
voir qu'elle est suffisante, on peut supposer $Y$ affine (II, 4.6.4).
Montrons en outre qu'on peut se borner au cas où $X$ est réduit. En
effet, soit $j:X_{\mathrm{red}}\to X$ l'injection canonique, et posons
$X_1=X_{\mathrm{red}}$, $X'_1=X'\times_X X_1$, de sorte que l'on a le
diagramme commutatif
\[
\xymatrix{
X'\ar[d]_g & X'_1\ar[l]_{j'}\ar[d]^{g_1\, .}\\
X & X_1\ar[l]^{j}
}
\tag{2.6.2.2}\label{III.2.6.2.2-fr}
\]
Le morphisme $f\circ j$ est alors de type fini (I, 6.3.4) et $g_1$ est
un morphisme fini (II, 6.1.5, (iii))~; si $\mathcal L'$ est ample pour
$f\circ g$, $j'^*(\mathcal L')$ est ample pour $f\circ g\circ j'$ puisque
$j'$ est une immersion fermée (II, 5.1.12 et I, 4.3.2). Si on pose
$Z_1=j^{-1}(Z)$, $Z_1$ est propre sur $Y$ (II, 5.4.10)~; d'autre part,
si $X$ est normal en un point $x$, il en est évidemment de même de
$X_{\mathrm{red}}$~; enfin, si $(g_*(\mathcal O_{X'}))_x$ est un
$\mathcal O_x$-module libre, il résulte aussitôt de (II, 1.5.2) que
$((g_1)_*(\mathcal O_{X'_1}))_x$ est un $\mathcal O_{X_1,x}$-module
libre. Enfin, comme $X$ est noethérien (I, 6.3.7), si
$j^*(\mathcal L)$ est ample, $\mathcal L$ est ample (II, 4.5.14), et
comme $j'^*(\mathcal L')=g_1^*(j^*(\mathcal L))$, cela achève la
réduction annoncée. Nous supposons donc désormais $Y$ affine et $X$
réduit.
\end{env}

Les hypothèses de (II, 6.6.11) étant alors vérifiées, il existe un
$Y$-préschéma réduit $X_2$ et un $Y$-morphisme $h:X_2\to X$ fini et
birationnel tel que la restriction de $h$ à $h^{-1}(X-Z)$ soit un
isomorphisme sur $X-Z$ et que $h^*(\mathcal L)$ soit ample. Remplaçant
$X'$ par $X_2$, on voit qu'on est ramené à démontrer la proposition en
supposant en outre que $g$ possède les propriétés qui viennent d'être
énumérées pour $h$. Nous désignerons encore par $Z$ un sous-préschéma de
$X$ ayant $Z$ pour espace sous-jacent, qui est propre sur $Y$
(II, 5.4.10).

\begin{env}[2.6.2.3]
\label{III.2.6.2.3-fr}
Soient maintenant $X_1$ un sous-préschéma fermé de $X$,
$j:X_1\to X$ l'injection canonique,
$X'_1=g^{-1}(X_1)=X'\times_X X_1$ son image réciproque,
$j':X'_1\to X'$ l'injection canonique, de sorte que l'on a le diagramme
commutatif (\hyperref[III.2.6.2.2-fr]{2.6.2.2})~; posons
$\mathcal L_1=j^*(\mathcal L)$,
$\mathcal L'_1=j'^*(\mathcal L')=g_1^*(\mathcal L_1)$, de sorte que
$\mathcal L'_1$ est ample pour $f\circ g\circ j'$ (II, 5.1.12). Si on
pose $Z_1=j^{-1}(Z)$, le sous-préschéma fermé $Z_1$ de $X_1$ est propre
sur $Y$ (II, 5.4.2, (ii)). En d'autres termes, les hypothèses de
(\hyperref[III.2.6.2-fr]{2.6.2}) sont vérifiées pour $X_1$,
$\mathcal L_1$, $g_1$ et $Z_1$.
\end{env}

\oldpage[III]{113}
Ceci va nous permettre de démontrer
(\hyperref[III.2.6.2-fr]{2.6.2}) par récurrence noethérienne
(0\textsubscript{I}, 2.2.2) dans le cas où la restriction de $g$ à
$g^{-1}(X-Z)$ est un isomorphisme sur $X-Z$ (ce qui est suffisant pour
notre propos, comme on l'a vu en
(\hyperref[III.2.6.2.1-fr]{2.6.2.1}))~: il suffira d'établir que si, pour
tout sous-préschéma fermé $X_1$ de $X$, dont l'espace sous-jacent est
$\neq X$, la conclusion de (\hyperref[III.2.6.2-fr]{2.6.2}) est vraie
pour le faisceau $\mathcal L_1$, alors elle est vraie aussi pour le
faisceau $\mathcal L$.

\begin{env}[2.6.2.4]
\label{III.2.6.2.4-fr}
Soient alors $\mathcal A=\mathcal O_X$, $\mathcal B=g_*(\mathcal O_{X'})$,
de sorte que $\mathcal B$ est une sous-$\mathcal A$-Algèbre de
$\mathcal R(X)$, qui est un $\mathcal A$-Module cohérent~; en outre, la
restriction $\mathcal B|(X-Z)$ est égale à $\mathcal A|(X-Z)$. Soit
$\mathcal K$ le conducteur de $\mathcal B$ sur $\mathcal A$, c'est-à-dire
le plus grand sous-$\mathcal A$-Module quasi-cohérent de $\mathcal A$ tel
que $\mathcal B\mathcal K\subset\mathcal A$ (ou encore l'annulateur du
$\mathcal A$-Module $\mathcal B/\mathcal A$
(0\textsubscript{I}, 5.3.7)), ce qui entraîne
$\mathcal B\mathcal K=\mathcal K$. Il est clair que
$\mathcal K_x=\mathcal A_x$ en tous les points admettant un voisinage
$W_x$ tel que $g$ soit un isomorphisme de $g^{-1}(W_x)$ sur $W_x$, et en
particulier en tous les points de $X-Z$ et dans un voisinage de tout
point générique d'une composante irréductible de $X$. Considérons alors
le sous-préschéma fermé $Z_1=\operatorname{Spec}(\mathcal A/\mathcal K)$
de $X$ défini par $\mathcal K$~; il est encore propre sur $Y$ car le
sous-espace $Z_1$ est fermé dans $Z$ (II, 5.4.10). De plus, la définition
de $\mathcal K$ montre que
$\mathcal B|(X-Z_1)=\mathcal A|(X-Z_1)$~; on voit donc qu'on peut toujours
se ramener au cas où
$Z=\operatorname{Spec}(\mathcal A/\mathcal K)$, et comme on a vu que
$X-Z_1$ est un ouvert non vide de $X$, on peut toujours supposer que
l'espace $Z$ est distinct de $X$.
\end{env}

\begin{env}[2.6.2.5]
\label{III.2.6.2.5-fr}
Considérons $X'$ comme égal à $\operatorname{Spec}(\mathcal B)$ (puisque
$g$ est affine) et soit $\mathcal K'=\widetilde{\mathcal K}$, Idéal
cohérent de $\mathcal O_{X'}$ tel que $g_*(\mathcal K')=\mathcal K$
(II, 1.4.1)~; le sous-préschéma fermé
$Z'=g^{-1}(Z)=Z\times_X X'$ de $X'$ est défini par $\mathcal K'$ et égal
à $\operatorname{Spec}(\mathcal B/\mathcal K)$ (II, 1.4.10)~; comme
$h:Z'\to Z$ est un morphisme fini (II, 6.1.5, (iii)), $Z'$ est propre
sur $Y$ (II, 6.1.11 et 5.4.2, (ii)).

Cela posé, il nous faut prouver que pour tout $x\in X$ et tout voisinage
ouvert $U$ de $x$, il existe une section $s$ d'un
$\mathcal L^{\otimes n}$ $(n>0)$ au-dessus de $X$ telle que
$x\in X_s\subset U$ (II, 4.5.2)~; nous distinguerons deux cas~:

$1^\circ$ On a $x\in X-Z$~; on peut évidemment supposer alors que l'on a
aussi $U\subset X-Z$, donc l'ouvert $U'=g^{-1}(U)$ ne rencontre pas
$Z'$. Comme $\mathcal L'$ est ample par hypothèse, il existe un $n>0$ et
une section $s'$ de $\mathcal L'^{\otimes n}$ au-dessus de $X'$ telle
que $x'=g^{-1}(x)\in X'_{s'}\subset g^{-1}(U)$ (II, 4.5.2). En outre, on
peut supposer que $\mathcal K'\otimes\mathcal L'^{\otimes n}$ soit
engendré par ses sections au-dessus de $X'$ (II, 4.5.5), donc, comme
$\mathcal K'_{x'}=\mathcal O_{x'}$, il y a une de ces sections $s''$
telle que $s''(x')\neq0$~; en la multipliant par $s'$ (ce qui revient à
remplacer $n$ par $2n$), on voit qu'on peut supposer aussi que
$x'\in X'_{s''}\subset g^{-1}(U)$. Cela étant, il résulte de
(0\textsubscript{I}, 5.4.10) que l'on a un isomorphisme canonique
\[
\Gamma(X,\mathcal K\otimes\mathcal L^{\otimes n})
\xrightarrow{\sim}
\Gamma(X',\mathcal K'\otimes\mathcal L'^{\otimes n}).
\]
La section $s$ de $\mathcal K\otimes\mathcal L^{\otimes n}$ qui
correspond à $s''$ par cet isomorphisme a évidemment les propriétés
voulues.

$2^\circ$ On a $x\in Z$. Soit $\mathcal J$ l'Idéal cohérent de
$\mathcal O_X$ définissant le sous-préschéma fermé réduit de $X$ ayant
pour espace sous-jacent $X-U$, et considérons dans $\mathcal B$ les deux

\oldpage[III]{114}
Idéaux cohérents $\mathcal J\mathcal B$ et
$\mathcal J\mathcal K=\mathcal J(\mathcal K\mathcal B)
=\mathcal K(\mathcal J\mathcal B)$, de sorte que l'on a le diagramme
d'inclusions
\[
\xymatrix@R=1.5em{
\mathcal J\mathcal B\ar@{^{(}->}[r] & \mathcal B\\
\mathcal J\ar@{^{(}->}[r]\ar@{^{(}->}[u] &
\mathcal A\ar@{^{(}->}[u]\\
\mathcal J\mathcal K\mathcal B=\mathcal J\mathcal K
\ar@{^{(}->}[r]\ar@{^{(}->}[u] &
\mathcal K\ar@{^{(}->}[u]
}
\tag{2.6.2.6}\label{III.2.6.2.6-fr}
\]
Soit $\mathcal J'$ l'Idéal cohérent $(\mathcal J\mathcal B)^\sim$ de
$\mathcal O_{X'}$, de sorte que $\mathcal J\mathcal B=g_*(\mathcal J')$,
$\mathcal J'\mathcal K'=(\mathcal J\mathcal K\mathcal B)^\sim$, et par
suite
\[
\mathcal J'/\mathcal J'\mathcal K'
=(\mathcal J\mathcal B/\mathcal J\mathcal K\mathcal B)^\sim
\]
(II, 1.4.4). Comme $\mathcal J|V=\mathcal J\mathcal K|V$ pour tout
ouvert $V$ ne rencontrant pas $Z$, on voit que le support de
$\mathcal J'/\mathcal J'\mathcal K'$ est contenu dans $Z'$. Comme $Z'$ est
propre sur $Y$, on peut appliquer (\hyperref[III.2.2.4-fr]{2.2.4}) et on
voit que pour $n$ assez grand, l'application canonique
\[
\Gamma(X',\mathcal J'\otimes\mathcal L'^{\otimes n})
\longrightarrow
\Gamma(X',(\mathcal J'/\mathcal J'\mathcal K')
\otimes\mathcal L'^{\otimes n})
\]
est surjective. Mais en vertu de (0\textsubscript{I}, 5.4.10), on en
conclut que l'application canonique
\[
\Gamma(X,\mathcal J\mathcal B\otimes\mathcal L^{\otimes n})
\longrightarrow
\Gamma(X,(\mathcal J\mathcal B/\mathcal J\mathcal K\mathcal B)
\otimes\mathcal L^{\otimes n})
\]
est surjective.

Cela étant, soient $i:Z\to X$ l'injection canonique,
$i':Z'\to X'$ l'injection canonique, de sorte qu'on a le diagramme
commutatif
\[
\xymatrix{
X'\ar[d]_g & Z'\ar[l]_{i'}\ar[d]^{h}\\
X & Z\ar[l]^{i}
}
\]
Soient $\mathcal M=i^*(\mathcal L)$, $\mathcal M'=i'^*(\mathcal L')$~;
comme $\mathcal L'$ est ample, $\mathcal M'$ est ample (II, 5.1.12), et
d'autre part $\mathcal M'=h^*(\mathcal M)$~; on conclut donc de
l'hypothèse de récurrence noethérienne (puisque $Z\neq X$) que
$\mathcal M$ est ample. Par suite
$i^*(\mathcal J/\mathcal J\mathcal K)\otimes\mathcal M^{\otimes n}$ est
engendré par ses sections au-dessus de $Z$ pour $n$ assez grand
(II, 4.5.5). Comme
$\mathcal J/\mathcal J\mathcal K=i_*(i^*(\mathcal J/\mathcal J\mathcal K))$,
on déduit encore de (0\textsubscript{I}, 5.4.10) qu'il existe une section
$s$ de $(\mathcal J/\mathcal J\mathcal K)\otimes
\mathcal L^{\otimes n}$ au-dessus de $X$ (pour un $n$ assez grand) telle
que $s(x)\neq0$, puisque l'on a $\mathcal J_x=\mathcal A_x$ par
définition de $\mathcal J$ et $\mathcal K_x\neq\mathcal A_x$ par
hypothèse. Le diagramme (\hyperref[III.2.6.2.6-fr]{2.6.2.6}) montre que
$s$ est aussi une section de
$(\mathcal J\mathcal B/\mathcal J\mathcal K\mathcal B)
\otimes\mathcal L^{\otimes n}$ au-dessus de $X$, donc $s$ est l'image
canonique d'une section $t$ de
$(\mathcal J\mathcal B)\otimes\mathcal L^{\otimes n}$ au-dessus de $X$.
Mais par définition, l'image canonique $s$ de $t$ mod.
$(\mathcal J\mathcal K\mathcal B)\otimes\mathcal L^{\otimes n}$ est
dans
\[
(\mathcal J\otimes\mathcal L^{\otimes n})/
(\mathcal J\mathcal K\mathcal B\otimes\mathcal L^{\otimes n})
\]
donc, en vertu de (\hyperref[III.2.6.2.6-fr]{2.6.2.6}), cela implique que
$t$ est une section de $\mathcal J\otimes\mathcal L^{\otimes n}$
au-dessus de $X$, et \emph{a fortiori} une section de
$\mathcal L^{\otimes n}$. On a vu ci-dessus que $t(x)\neq0$, donc
$x\in X_t$, et par définition de $\mathcal J$, $t(y)=0$ dans $X-U$ qui
est le support de $\mathcal O_X/\mathcal J$~; donc $X_t\subset U$, ce qui
achève la démonstration.
\end{env}
\end{proof}

\oldpage[III]{115}
\begin{env}[2.6.3]
\label{III.2.6.3-fr}
\emph{Remarque.} Lorsque $X$ est propre sur $Y$, on peut démontrer
(\hyperref[III.2.6.2-fr]{2.6.2}) plus simplement, en raisonnant comme dans
le th. de Chevalley (II, 6.7.1), à l'aide de
(\hyperref[III.2.6.1-fr]{2.6.1}) et du lemme (II, 6.7.1.1).
\end{env}

\section{Le théorème de finitude pour les morphismes propres}
\label{section:III.3-fr}

\subsection{Le lemme de dévissage}
\label{subsection:III.3.1-fr}

\begin{definition}[3.1.1]
\label{III.3.1.1-fr}
Soit $K$ une catégorie abélienne. On dit qu'un sous-ensemble $K'$ de
l'ensemble des objets de $K$ est \emph{exact} si $0\in K'$ et si, pour toute
suite exacte
\[
0\longrightarrow A'\longrightarrow A\longrightarrow A''\longrightarrow0
\]
dans $K$ telle que deux des objets $A$, $A'$, $A''$ soit dans $K'$, alors le
troisième est aussi dans $K'$.
\end{definition}

\begin{theorem}[3.1.2]
\label{III.3.1.2-fr}
Soit $X$ un préschéma noethérien~; on désigne par $K$ la catégorie abélienne
des $\mathcal O_X$-Modules cohérents. Soient $K'$ un sous-ensemble exact de
$K$, $X'$ une partie fermée de l'espace sous-jacent à $X$. On suppose que
pour toute partie fermée irréductible $Y$ de $X'$, de point générique $y$, il
existe un $\mathcal O_X$-Module $\mathcal G\in K'$ tel que $\mathcal G_y$ soit
un $k(y)$-espace vectoriel de dimension $1$. Alors tout $\mathcal O_X$-Module
cohérent de support contenu dans $X'$ appartient à $K'$ (et en particulier,
si $X'=X$, on a $K'=K$).
\end{theorem}

\begin{proof}
Considérons la propriété suivante $P(Y)$ d'une partie fermée $Y$ de $X'$~:
tout $\mathcal O_X$-Module cohérent de support contenu dans $Y$ appartient à
$K'$. En vertu du principe de récurrence noethérienne (0\textsubscript{I},
2.2.2), on voit qu'on est ramené à démontrer que \emph{si $Y$ est une partie
fermée de $X'$ telle que la propriété $P(Y')$ soit vraie pour toute partie
fermée $Y'$ de $Y$, distincte de $Y$, alors $P(Y)$ est vraie}.

Soit donc $\mathcal F\in K$ à support contenu dans $Y$ et prouvons que
$\mathcal F\in K'$. Désignons encore par $Y$ le sous-préschéma fermé réduit de
$X$ ayant $Y$ pour espace sous-jacent (I, 5.2.1)~; il est défini par un Idéal
cohérent $\mathcal J$ de $\mathcal O_X$. On sait (I, 9.3.4) qu'il existe un
entier $n>0$ tel que $\mathcal J^n\mathcal F=0$~; pour $1\leq k\leq n$, on a
donc une suite exacte
\[
0\longrightarrow
\mathcal J^{k-1}\mathcal F/\mathcal J^k\mathcal F
\longrightarrow \mathcal F/\mathcal J^k\mathcal F
\longrightarrow \mathcal F/\mathcal J^{k-1}\mathcal F
\longrightarrow0
\]
de $\mathcal O_X$-Modules cohérents (0\textsubscript{I}, 5.3.6 et 5.3.3)~;
comme $K'$ est exact, on voit, par récurrence sur $k$, qu'il suffit de montrer
que chacun des
$\mathcal F_k=\mathcal J^{k-1}\mathcal F/\mathcal J^k\mathcal F$ est dans
$K'$. On est donc ramené à prouver que $\mathcal F\in K'$ sous l'hypothèse
supplémentaire que $\mathcal J\mathcal F=0$~; il revient au même de dire que
$\mathcal F=j_*(j^*(\mathcal F))$, où $j$ est l'injection canonique
$Y\to X$. Distinguons maintenant deux cas~:
\begin{enumerate}
\item[a)] $Y$ est \emph{réductible}. Soient $Y=Y'\cup Y''$, $Y'$ et $Y''$
étant des parties fermées de $Y$, distinctes de $Y$~; soient encore $Y'$,
$Y''$ les sous-préschémas fermés réduits de $X$ ayant pour espaces
sous-jacents $Y'$, $Y''$ respectivement, qui sont définis respectivement par
les Idéaux cohérents $\mathcal J'$, $\mathcal J''$ de $\mathcal O_X$. Posons
\[
\mathcal F'=\mathcal F\otimes_{\mathcal O_X}
(\mathcal O_X/\mathcal J'),\qquad
\mathcal F''=\mathcal F\otimes_{\mathcal O_X}
(\mathcal O_X/\mathcal J'').
\]
Les homomorphismes canoniques $\mathcal F\to\mathcal F'$,
$\mathcal F\to\mathcal F''$ définissent donc un homomorphisme
$u:\mathcal F\to\mathcal F'\oplus\mathcal F''$. Montrons que pour tout
$z\notin Y'\cap Y''$, l'homomorphisme
$u_z:\mathcal F_z\to\mathcal F'_z\oplus\mathcal F''_z$ est \emph{bijectif}.
En effet, on a $\mathcal J'\cap\mathcal J''=\mathcal J$, car la question est
locale et

\oldpage[III]{116}
l'égalité précédente résulte de (I, 5.2.1 et 1.1.5)~; si $z\notin Y''$, on a
donc $\mathcal J'_z=\mathcal J_z$, d'où $\mathcal F'_z=\mathcal F_z$ et
$\mathcal F''_z=0$, ce qui établit notre assertion dans ce cas~; on raisonne
de même si $z\notin Y'$. Par suite, le noyau et le conoyau de $u$, qui sont
dans $K$ (0\textsubscript{I}, 5.3.4) ont leurs supports dans $Y'\cap Y''$,
et sont donc dans $K'$ par hypothèse~; pour la même raison, $\mathcal F'$ et
$\mathcal F''$ sont dans $K'$, donc aussi $\mathcal F'\oplus\mathcal F''$,
puisque $K'$ est exact. La conclusion résulte alors de la considération des
deux suites exactes
\[
0\longrightarrow\operatorname{Im}u\longrightarrow
\mathcal F'\oplus\mathcal F''\longrightarrow
\operatorname{Coker}u\longrightarrow0,
\]
\[
0\longrightarrow\operatorname{Ker}u\longrightarrow\mathcal F
\longrightarrow\operatorname{Im}u\longrightarrow0
\]
et de l'hypothèse que $K'$ est exact.

\item[b)] $Y$ est irréductible, et par suite le sous-préschéma $Y$ de $X$ est
\emph{intègre}. Si $y$ est son point générique, on a
$(\mathcal O_Y)_y=k(y)$, et comme $j^*(\mathcal F)$ est un
$\mathcal O_Y$-Module cohérent, $\mathcal F_y=(j^*(\mathcal F))_y$ est un
$k(y)$-espace vectoriel de dimension finie $m$. Par hypothèse, il y a un
$\mathcal O_X$-Module cohérent $\mathcal G$ (nécessairement de support $Y$)
tel que $\mathcal G_y$ soit un $k(y)$-espace vectoriel de dimension $1$. Par
suite, il y a un $k(y)$-isomorphisme
$(\mathcal G_y)^m\xrightarrow{\sim}\mathcal F_y$, qui est aussi un
$\mathcal O_Y$-isomorphisme, et comme $\mathcal G^m$ et $\mathcal F$ sont
cohérents, il existe un voisinage ouvert $W$ de $y$ dans $X$ et un
isomorphisme $\mathcal G^m|W\xrightarrow{\sim}\mathcal F|W$
(0\textsubscript{I}, 5.2.7). Soit $\mathcal H$ le graphe de cet isomorphisme,
qui est un sous-$(\mathcal O_X|W)$-Module cohérent de
$(\mathcal G^m\oplus\mathcal F)|W$, canoniquement isomorphe à
$\mathcal G^m|W$ et à $\mathcal F|W$~; il existe donc un sous-$\mathcal O_X$-
Module cohérent $\mathcal H_0$ de $\mathcal G^m\oplus\mathcal F$, induisant
$\mathcal H$ sur $W$ et $0$ sur $X-Y$ puisque $\mathcal G^m$ et
$\mathcal F$ ont pour support $Y$ (I, 9.4.7). Les restrictions
$v:\mathcal H_0\to\mathcal G^m$ et $w:\mathcal H_0\to\mathcal F$ des
projections canoniques de $\mathcal G^m\oplus\mathcal F$ sont alors des
homomorphismes de $\mathcal O_X$-Modules cohérents, qui, dans $W$ et dans
$X-Y$, se réduisent à des isomorphismes~; autrement dit, les noyaux et
conoyaux de $v$ et $w$ ont leur support dans l'ensemble fermé
$Y-(Y\cap W)$, distinct de $Y$. Ils appartiennent donc à $K'$~; d'autre
part, on a $\mathcal G^m\in K'$ puisque $\mathcal G\in K'$ et que $K'$ est
exact. On en conclut successivement, par l'exactitude de $K'$, que
$\mathcal H_0\in K'$ puis que $\mathcal F\in K'$. C.Q.F.D.
\end{enumerate}
\end{proof}

\begin{corollary}[3.1.3]
\label{III.3.1.3-fr}
Supposons que le sous-ensemble exact $K'$ de $K$ ait en outre la propriété
que tout facteur direct cohérent d'un $\mathcal O_X$-Module cohérent
$\mathcal M\in K'$ appartienne encore à $K'$. Dans ces conditions, la
conclusion de (\hyperref[III.3.1.2-fr]{3.1.2}) subsiste encore lorsque la
condition «~$\mathcal G_y$ est un $k(y)$-espace vectoriel de dimension $1$~»
est remplacée par $\mathcal G_y\neq0$ (ce qui équivaut à
$\operatorname{Supp}(\mathcal G)=Y$).
\end{corollary}

\begin{proof}
En effet, le raisonnement de (\hyperref[III.3.1.2-fr]{3.1.2}) ne doit être
modifié que dans le cas b)~; cette fois $\mathcal G_y$ est un $k(y)$-espace
vectoriel de dimension $q>0$, et on a par suite un $\mathcal O_Y$-
isomorphisme $(\mathcal G_y)^m\xrightarrow{\sim}(\mathcal F_y)^q$~; la fin
du raisonnement de (\hyperref[III.3.1.2-fr]{3.1.2}) prouve alors que
$\mathcal F^q\in K'$, et l'hypothèse additionnelle sur $K'$ implique que
$\mathcal F\in K'$.
\end{proof}

\subsection{Le théorème de finitude : cas des schémas usuels}
\label{subsection:III.3.2-fr}

\begin{theorem}[3.2.1]
\label{III.3.2.1-fr}
Soient $Y$ un préschéma localement noethérien, $f:X\to Y$ un morphisme
propre. Pour tout $\mathcal O_X$-Module cohérent $\mathcal F$, les
$\mathcal O_Y$-Modules $R^qf_*(\mathcal F)$ sont cohérents pour $q\geq0$.
\end{theorem}

\begin{proof}
La question étant locale sur $Y$, on peut supposer $Y$ noethérien, donc $X$
noethérien (I, 6.3.7). Les $\mathcal O_X$-Modules cohérents $\mathcal F$ pour
lesquels la conclusion du th. (\hyperref[III.3.2.1-fr]{3.2.1}) est vraie
forment un sous-ensemble exact $K'$ de la catégorie $K$ des
$\mathcal O_X$-Modules cohérents.

\oldpage[III]{117}
En effet, soit
$0\to\mathcal F'\to\mathcal F\to\mathcal F''\to0$ une suite exacte de
$\mathcal O_X$-Modules cohérents~; supposons par exemple que $\mathcal F'$
et $\mathcal F''$ appartiennent à $K'$~; on a la suite exacte de cohomologie
\[
R^{q-1}f_*(\mathcal F'')\xrightarrow{\partial}R^qf_*(\mathcal F')
\longrightarrow R^qf_*(\mathcal F)\longrightarrow R^qf_*(\mathcal F'')
\xrightarrow{\partial}R^{q+1}f_*(\mathcal F')
\]
dans laquelle par hypothèse les quatre termes extrêmes sont cohérents~; il
en est donc de même du terme médian $R^qf_*(\mathcal F)$ par
(0\textsubscript{I}, 5.3.4 et 5.3.3). On montre de la même manière que
lorsque $\mathcal F$ et $\mathcal F'$ (resp. $\mathcal F$ et
$\mathcal F''$) sont dans $K'$, il en est de même de $\mathcal F''$ (resp.
$\mathcal F'$). De plus, tout \emph{facteur direct} cohérent $\mathcal F'$
d'un $\mathcal O_X$-Module $\mathcal F\in K'$ appartient aussi à $K'$~: en
effet, $R^qf_*(\mathcal F')$ est alors facteur direct de
$R^qf_*(\mathcal F)$ (G, II, 4.4.4), donc est de type fini, et comme il est
quasi-cohérent (\hyperref[III.1.4.10-fr]{1.4.10}), il est cohérent, $Y$
étant noethérien. En vertu de (\hyperref[III.3.1.3-fr]{3.1.3}), on est
ramené à prouver que lorsque $X$ est \emph{irréductible} de point générique
$x$, il existe un $\mathcal O_X$-Module cohérent $\mathcal F$ appartenant à
$K'$, tel que $\mathcal F_x\neq0$~: en effet, si ce point est établi, on
pourra l'appliquer à tout sous-préschéma fermé irréductible $Y$ de $X$, car
si $j:Y\to X$ est l'injection canonique, $f\circ j$ est propre (II, 5.4.2),
et si $\mathcal G$ est un $\mathcal O_Y$-Module cohérent de support $Y$,
$j_*(\mathcal G)$ est un $\mathcal O_X$-Module cohérent tel que
\[
R^q(f\circ j)_*(\mathcal G)=R^qf_*(j_*(\mathcal G))
\]
(G, II, 4.9.1), donc on est bien dans les conditions d'application de
(\hyperref[III.3.1.3-fr]{3.1.3}).

Or, en vertu du lemme de Chow (II, 5.6.2), il existe un préschéma
irréductible $X'$ et un morphisme \emph{projectif} et surjectif
$g:X'\to X$ tel que $f\circ g:X'\to Y$ soit \emph{projectif}. Il existe un
$\mathcal O_{X'}$-Module $\mathcal L$ ample pour $g$ (II, 5.3.1)~;
appliquons le th. fondamental des morphismes projectifs
(\hyperref[III.2.2.1-fr]{2.2.1}) à $g:X'\to X$ et à $\mathcal L$~: il
existe donc un entier $n$ tel que
$\mathcal F=g_*(\mathcal O_{X'}(n))$ soit un $\mathcal O_X$-Module cohérent
et $R^qg_*(\mathcal O_{X'}(n))=0$ pour tout $q>0$~; en outre, comme
\[
g^*(g_*(\mathcal O_{X'}(n)))\longrightarrow\mathcal O_{X'}(n)
\]
est surjectif pour $n$ assez grand (\hyperref[III.2.2.1-fr]{2.2.1}), on
voit qu'on peut supposer que, au point générique $x$ de $X$, on a
$\mathcal F_x\neq0$ (II, 3.4.7). D'autre part, comme $f\circ g$ est
projectif et $Y$ noethérien, les
$R^p(f\circ g)_*(\mathcal O_{X'}(n))$ sont \emph{cohérents}
(\hyperref[III.2.2.1-fr]{2.2.1}). Cela étant,
$R^*(f\circ g)_*(\mathcal O_{X'}(n))$ est l'aboutissement d'une suite
spectrale de Leray, dont le terme $E_2$ est donné par
\[
E_2^{pq}=R^pf_*(R^qg_*(\mathcal O_{X'}(n)))~;
\]
ce qui précède montre que cette suite spectrale est dégénérée, et on sait
alors (0, 11.1.6) que $E_2^{p0}=R^pf_*(\mathcal F)$ est isomorphe à
$R^p(f\circ g)_*(\mathcal O_{X'}(n))$, ce qui achève la démonstration.
\end{proof}

\begin{corollary}[3.2.2]
\label{III.3.2.2-fr}
Soit $Y$ un préschéma localement noethérien. Pour tout morphisme propre
$f:X\to Y$, l'image directe par $f$ de tout $\mathcal O_X$-Module cohérent
est un $\mathcal O_Y$-Module cohérent.
\end{corollary}

\begin{corollary}[3.2.3]
\label{III.3.2.3-fr}
Soient $A$ un anneau noethérien, $X$ un schéma propre sur $A$~; pour tout
$\mathcal O_X$-Module cohérent $\mathcal F$, les $H^p(X,\mathcal F)$ sont
des $A$-modules de type fini, et il existe un entier $r>0$ tel que pour tout
$\mathcal O_X$-Module cohérent $\mathcal F$ et tout $p>r$,
$H^p(X,\mathcal F)=0$.
\end{corollary}

\begin{proof}
La seconde assertion a déjà été démontrée
(\hyperref[III.1.4.12-fr]{1.4.12})~; la première résulte du th. de finitude
(\hyperref[III.3.2.1-fr]{3.2.1}), compte tenu de
(\hyperref[III.1.4.11-fr]{1.4.11}).
\end{proof}

En particulier, si $X$ est un \emph{schéma algébrique propre} sur un corps
$k$, alors, pour tout $\mathcal O_X$-Module cohérent $\mathcal F$, les
$H^p(X,\mathcal F)$ sont des $k$-espaces vectoriels de \emph{dimension
finie}.

\begin{corollary}[3.2.4]
\label{III.3.2.4-fr}
Soient $Y$ un préschéma localement noethérien, $f:X\to Y$ un morphisme de
type fini. Pour tout $\mathcal O_X$-Module cohérent $\mathcal F$ dont le
support est propre sur $Y$ (II, 5.4.10), les $\mathcal O_Y$-Modules
$R^qf_*(\mathcal F)$ sont cohérents.
\end{corollary}

\oldpage[III]{118}
\begin{proof}
La question étant locale sur $Y$, on peut supposer $Y$ noethérien, et il en
est donc de même de $X$ (I, 6.3.7). Par hypothèse, tout sous-préschéma fermé
$Z$ de $X$ dont l'espace sous-jacent est $\operatorname{Supp}(\mathcal F)$
est propre sur $Y$, autrement dit, si $j:Z\to X$ est l'injection canonique,
$f\circ j:Z\to Y$ est propre. Or, on peut supposer $Z$ tel que
$\mathcal F=j_*(\mathcal G)$, où $\mathcal G=j^*(\mathcal F)$ est un
$\mathcal O_Z$-Module cohérent (I, 9.3.5)~; comme on a
$R^qf_*(\mathcal F)=R^q(f\circ j)_*(\mathcal G)$ par
(\hyperref[III.1.3.4-fr]{1.3.4}), la conclusion résulte aussitôt de
(\hyperref[III.3.2.1-fr]{3.2.1}).
\end{proof}

\subsection{Généralisation du théorème de finitude (schémas usuels)}
\label{subsection:III.3.3-fr}

\begin{proposition}[3.3.1]
\label{III.3.3.1-fr}
Soient $Y$ un préschéma noethérien, $\mathcal S$ une
$\mathcal O_Y$-Algèbre graduée à degrés positifs, quasi-cohérente et de type
fini, $Y'=\operatorname{Proj}(\mathcal S)$ et $g:Y'\to Y$ le morphisme
structural. Soient $f:X\to Y$ un morphisme propre,
$\mathcal S'=f^*(\mathcal S)$,
$\mathcal M=\bigoplus_{k\in\mathbb Z}\mathcal M_k$ un $\mathcal S'$-Module
gradué quasi-cohérent et de type fini. Alors les
\[
R^pf_*(\mathcal M)=\bigoplus_{k\in\mathbb Z}R^pf_*(\mathcal M_k)
\]
sont des $\mathcal S$-Modules gradués de type fini pour tout $p$. Supposons
en outre que $\mathcal S$ soit engendrée par $\mathcal S_1$~; alors, pour
chaque $p\in\mathbb Z$, il existe un entier $k_p$ tel que pour tout
$k\geq k_p$ et tout $r\geq0$, on ait
\[
R^pf_*(\mathcal M_{k+r})=\mathcal S_rR^pf_*(\mathcal M_k).
\tag{3.3.1.1}\label{III.3.3.1.1-fr}
\]
\end{proposition}

\begin{proof}
La première assertion est identique à l'énoncé de
(\hyperref[III.2.4.1-fr]{2.4.1}, (i)), où on a simplement remplacé
«~morphisme projectif~» par «~morphisme propre~». Or, dans la démonstration
de (\hyperref[III.2.4.1-fr]{2.4.1}, (i)), l'hypothèse sur $f$ a été utilisée
uniquement pour montrer (avec les notations de cette démonstration) que
$R^pf'_*(\widetilde{\mathcal M})$ est un $\mathcal O_{Y'}$-Module cohérent.
Avec les hypothèses de (\hyperref[III.3.3.1-fr]{3.3.1}), $f'$ est propre
(II, 5.4.2, (iii)), donc on peut reprendre sans changement toute la
démonstration de (\hyperref[III.2.4.1-fr]{2.4.1}, (i)), grâce au théorème de
finitude (\hyperref[III.3.2.1-fr]{3.2.1}).

Quant à la seconde assertion, il suffit de remarquer qu'il y a un
recouvrement ouvert affine fini $(U_i)$ de $Y$ tel que les restrictions à
$U_i$ des deux membres de
(\hyperref[III.3.3.1.1-fr]{3.3.1.1}) soient égales pour tout
$k\geq k_{p,i}$ (II, 2.1.6, (ii))~; il suffit de prendre pour $k_p$ le plus
grand des $k_{p,i}$.
\end{proof}

\begin{corollary}[3.3.2]
\label{III.3.3.2-fr}
Soient $A$ un anneau noethérien, $\mathfrak m$ un idéal de $A$, $X$ un
$A$-schéma propre, $\mathcal F$ un $\mathcal O_X$-Module cohérent. Alors,
pour tout $p\geq0$, la somme directe
$\bigoplus_{k\geq0}H^p(X,\mathfrak m^k\mathcal F)$ est un module de type fini
sur l'anneau $S=\bigoplus_{k\geq0}\mathfrak m^k$~; en particulier, il existe
un entier $k_p\geq0$ tel que pour tout $k\geq k_p$, et tout $r\geq0$, on ait
\[
H^p(X,\mathfrak m^{k+r}\mathcal F)
=\mathfrak m^rH^p(X,\mathfrak m^k\mathcal F).
\tag{3.3.3.1}\label{III.3.3.2.1-fr}
\]
\end{corollary}

\begin{proof}
Il suffit d'appliquer (\hyperref[III.3.3.2-fr]{3.3.2}) avec
$Y=\operatorname{Spec}(A)$, $\mathcal S=\widetilde S$,
$\mathcal M_k=\mathfrak m^k\mathcal F$, compte tenu de
(\hyperref[III.1.4.11-fr]{1.4.11}).
\end{proof}

Il convient de rappeler que la structure de $S$-module de
$\bigoplus_{k\geq0}H^p(X,\mathfrak m^k\mathcal F)$ s'obtient en considérant,
pour tout $a\in\mathfrak m^r$, l'application
\[
H^p(X,\mathfrak m^k\mathcal F)
\longrightarrow H^p(X,\mathfrak m^{k+r}\mathcal F)
\]
qui provient par passage à la cohomologie de la multiplication
$\mathfrak m^k\mathcal F\to\mathfrak m^{k+r}\mathcal F$ définie par $a$
(\hyperref[III.2.4.1-fr]{2.4.1}).

\oldpage[III]{119}
\subsection{Le théorème de finitude : cas des schémas formels}
\label{subsection:III.3.4-fr}

Les résultats de cette section (sauf la définition
(\hyperref[III.3.4.1-fr]{3.4.1})) ne seront pas utilisés dans la suite de ce
chapitre.

\begin{env}[3.4.1]
\label{III.3.4.1-fr}
Soient $\mathfrak X$ et $\mathfrak S$ deux préschémas formels localement
noethériens (I, 10.4.2), $f:\mathfrak X\to\mathfrak S$ un morphisme de
préschémas formels. Nous dirons que $f$ est un morphisme \emph{propre} s'il
vérifie les conditions suivantes~:
\begin{enumerate}
\item[$1^\circ$] $f$ est un morphisme de type fini (I, 10.13.3).
\item[$2^\circ$] Si $\mathcal K$ est un Idéal de définition de $\mathfrak S$
et si l'on pose
$\mathcal J=f^*(\mathcal K)\mathcal O_{\mathfrak X}$,
$X_0=(\mathfrak X,\mathcal O_{\mathfrak X}/\mathcal J)$,
$S_0=(\mathfrak S,\mathcal O_{\mathfrak S}/\mathcal K)$, le morphisme
$f_0:X_0\to Y_0$ déduit de $f$ (I, 10.5.6) est propre.
\end{enumerate}

Il est immédiat que cette définition ne dépend pas de l'Idéal de définition
$\mathcal K$ de $\mathfrak S$ considéré~; en effet, si $\mathcal K'$ est un
second Idéal de définition tel que $\mathcal K'\subset\mathcal K$, et si on
pose $\mathcal J'=f^*(\mathcal K')\mathcal O_{\mathfrak X}$,
$X'_0=(\mathfrak X,\mathcal O_{\mathfrak X}/\mathcal J')$,
$S'_0=(\mathfrak S,\mathcal O_{\mathfrak S}/\mathcal K')$, le morphisme
$f'_0:X'_0\to S'_0$ déduit de $f$ est tel que le diagramme
\[
\xymatrix{
X_0\ar[r]^{f_0}\ar[d]_i & S_0\ar[d]^j\\
X'_0\ar[r]_{f'_0} & S'_0
}
\]
est commutatif, $i$ et $j$ étant des immersions surjectives~; il revient donc
au même de dire que $f_0$ ou $f'_0$ est propre, en vertu de (II, 5.4.5).

On notera que, pour tout $n\geq0$, si on pose
$X_n=(\mathfrak X,\mathcal O_{\mathfrak X}/\mathcal J^{n+1})$,
$S_n=(\mathfrak S,\mathcal O_{\mathfrak S}/\mathcal K^{n+1})$, le morphisme
$f_n:X_n\to Y_n$ déduit de $f$ (I, 10.5.6) est propre pour tout $n$ dès qu'il
l'est pour $n=0$ (II, 5.4.6).

Si $g:Y\to Z$ est un morphisme propre de préschémas usuels, localement
noethériens, $Z'$ une partie fermée de $Z$, $Y'$ une partie fermée de $Y$
telle que $g(Y')\subset Z'$, le prolongement
$\widehat g:Y_{/Y'}\to Z_{/Z'}$ de $g$ aux complétés (I, 10.9.1) est un
morphisme propre de préschémas formels, comme il résulte de la définition et
de (II, 5.4.5).

Soient $\mathfrak X$ et $\mathfrak S$ deux préschémas formels localement
noethériens, $f:\mathfrak X\to\mathfrak S$ un morphisme \emph{de type fini}
(I, 10.13.3)~; les notations étant les mêmes que ci-dessus, on dit qu'une
partie $Z$ de l'espace sous-jacent à $\mathfrak X$ est \emph{propre sur}
$\mathfrak S$ (ou propre pour $f$) si, considérée comme partie de $X_0$, $Z$
est \emph{propre sur $S_0$} (II, 5.4.10). Toutes les propriétés des parties
propres de préschémas usuels énoncées dans (II, 5.4.10) sont encore valables
pour les parties propres de préschémas formels comme il résulte aussitôt des
définitions.
\end{env}

\begin{theorem}[3.4.2]
\label{III.3.4.2-fr}
Soient $\mathfrak X$, $\mathfrak Y$ deux préschémas formels localement
noethériens, $f:\mathfrak X\to\mathfrak Y$ un morphisme propre. Pour tout
$\mathcal O_{\mathfrak X}$-Module cohérent $\mathcal F$, les
$\mathcal O_{\mathfrak Y}$-Modules $R^qf_*(\mathcal F)$ sont cohérents pour
tout $q\geq0$.
\end{theorem}

Soient $\mathcal J$ un Idéal de définition de $\mathfrak Y$,
$\mathcal K=f^*(\mathcal J)\mathcal O_{\mathfrak X}$, et considérons les
$\mathcal O_{\mathfrak X}$-Modules
\[
\mathcal F_k=\mathcal F\otimes_{\mathcal O_{\mathfrak Y}}
(\mathcal O_{\mathfrak Y}/\mathcal J^{k+1})
=\mathcal F/\mathcal K^{k+1}\mathcal F\qquad(k\geq0)
\tag{3.4.2.1}\label{III.3.4.2.1-fr}
\]
qui forment évidemment un \emph{système projectif} de
$\mathcal O_{\mathfrak X}$-Modules topologiques, tel que

\oldpage[III]{120}
\[
\mathcal F=\varprojlim_k\mathcal F_k
\]
par (I, 10.11.3). D'autre part, il résultera de
(\hyperref[III.3.4.2-fr]{3.4.2}) que chacun des $R^qf_*(\mathcal F)$, étant
cohérent, est naturellement muni d'une structure de
$\mathcal O_{\mathfrak Y}$-Module topologique (I, 10.11.6), et il en est de
même des $R^qf_*(\mathcal F_k)$. Aux homomorphismes canoniques
$\mathcal F\to\mathcal F_k=\mathcal F/\mathcal K^{k+1}\mathcal F$
correspondent canoniquement des homomorphismes
\[
R^qf_*(\mathcal F)\longrightarrow R^qf_*(\mathcal F_k)
\]
qui sont nécessairement continus pour les structures de
$\mathcal O_{\mathfrak Y}$-Module topologique précédentes (I, 10.11.6), et
forment un système projectif, donnant à la limite un homomorphisme canonique
fonctoriel
\[
R^qf_*(\mathcal F)\longrightarrow\varprojlim_kR^qf_*(\mathcal F_k)
\tag{3.4.2.2}\label{III.3.4.2.2-fr}
\]
qui sera un homomorphisme continu de $\mathcal O_{\mathfrak Y}$-Modules
topologiques. Nous allons démontrer en même temps que
(\hyperref[III.3.4.2-fr]{3.4.2}) le

\begin{corollary}[3.4.3]
\label{III.3.4.3-fr}
Chacun des homomorphismes (\hyperref[III.3.4.2.2-fr]{3.4.2.2}) est un
isomorphisme topologique. En outre, si $\mathfrak Y$ est noethérien, le
système projectif
$(R^qf_*(\mathcal F/\mathcal K^{k+1}\mathcal F))_{k\geq0}$ satisfait à la
condition (ML) (0, 13.1.1).
\end{corollary}

Nous commencerons par établir (\hyperref[III.3.4.2-fr]{3.4.2}) et
(\hyperref[III.3.4.3-fr]{3.4.3}) lorsque $\mathfrak Y$ est un schéma affine
formel noethérien (I, 10.4.1)~:

\begin{corollary}[3.4.4]
\label{III.3.4.4-fr}
Sous les hypothèses de (\hyperref[III.3.4.2-fr]{3.4.2}), supposons en outre
que $\mathfrak Y=\operatorname{Spf}(A)$, où $A$ est un anneau adique
noethérien. Soit $\mathfrak J$ un idéal de définition de $A$, et posons
$\mathcal F_k=\mathcal F/\mathfrak J^{k+1}\mathcal F$ pour $k\geq0$. Alors
les $H^n(\mathfrak X,\mathcal F)$ sont des $A$-modules de type fini~; le
système projectif $(H^n(\mathfrak X,\mathcal F_k))_{k\geq0}$ satisfait à la
condition (ML) pour tout $n$~; si on pose
\[
N_{n,k}=\operatorname{Ker}\bigl(H^n(\mathfrak X,\mathcal F)
\longrightarrow H^n(\mathfrak X,\mathcal F_k)\bigr)
\tag{3.4.4.1}\label{III.3.4.4.1-fr}
\]
(égal aussi à
$\operatorname{Im}(H^n(\mathfrak X,\mathfrak J^{k+1}\mathcal F)
\to H^n(\mathfrak X,\mathcal F))$ par la suite exacte de cohomologie), les
$N_{k,n}$ définissent sur $H^n(\mathfrak X,\mathcal F)$ une filtration
$\mathfrak J$-bonne (0, 13.7.7)~; enfin, l'homomorphisme canonique
\[
H^n(\mathfrak X,\mathcal F)
\longrightarrow\varprojlim_kH^n(\mathfrak X,\mathcal F_k)
\tag{3.4.4.2}\label{III.3.4.4.2-fr}
\]
est un isomorphisme topologique pour tout $n$ (le premier membre étant muni
de la topologie $\mathfrak J$-adique, les $H^n(\mathfrak X,\mathcal F_k)$ de
la topologie discrète).
\end{corollary}

Posons
\[
S=\operatorname{gr}(A)=\bigoplus_{k\geq0}
\mathfrak J^k/\mathfrak J^{k+1},\qquad
\mathcal M=\operatorname{gr}(\mathcal F)=\bigoplus_{k\geq0}
\mathfrak J^k\mathcal F/\mathfrak J^{k+1}\mathcal F.
\tag{3.4.4.3}\label{III.3.4.4.3-fr}
\]
On sait que $\mathfrak J^\Delta$ est un Idéal de définition de $\mathfrak Y$
(I, 10.3.1)~; soient
$\mathcal K=f^*(\mathfrak J^\Delta)\mathcal O_{\mathfrak X}$,
$X_0=(\mathfrak X,\mathcal O_{\mathfrak X}/\mathcal K)$,
$Y_0=(\mathfrak Y,\mathcal O_{\mathfrak Y}/\mathfrak J^\Delta)
=\operatorname{Spec}(A_0)$ avec $A_0=A/\mathfrak J$. Il est clair que les
$\mathcal M_k=\mathfrak J^k\mathcal F/\mathfrak J^{k+1}\mathcal F$ sont des
$\mathcal O_{X_0}$-Modules cohérents (I, 10.11.3). Considérons d'autre part la
$\mathcal O_{X_0}$-Algèbre graduée quasi-cohérente
\[
\mathcal S=\mathcal O_{X_0}\otimes_{A_0}S
=\operatorname{gr}(\mathcal O_{\mathfrak X})
=\bigoplus_{k\geq0}\mathcal K^k/\mathcal K^{k+1}.
\tag{3.4.4.4}\label{III.3.4.4.4-fr}
\]

L'hypothèse que $\mathcal F$ est un $\mathcal O_{\mathfrak X}$-Module de type
fini entraîne d'abord que $\mathcal M$ est

\oldpage[III]{121}
un $\mathcal S$-Module gradué \emph{de type fini}. En effet, la question est
locale sur $\mathfrak X$, et on peut donc supposer pour la traiter que
$\mathfrak X=\operatorname{Spf}(B)$, où $B$ est un anneau adique noethérien,
et $\mathcal F=N^\Delta$, où $N$ est un $B$-module de type fini
(I, 10.10.5)~; on a en outre $X_0=\operatorname{Spec}(B_0)$ où
$B_0=B/\mathfrak J B$, et les $\mathcal O_{X_0}$-Modules quasi-cohérents
$\mathcal S$ et $\mathcal M$ sont respectivement égaux à $\widetilde{S'}$ et
$\widetilde{M'}$, où
\[
S'=\bigoplus_{k\geq0}
((\mathfrak J^k/\mathfrak J^{k+1})\otimes_{A_0}B_0),\qquad
M'=\bigoplus_{k\geq0}
((\mathfrak J^k/\mathfrak J^{k+1})\otimes_{A_0}N_0),
\]
avec $N_0=N/\mathfrak JN$~; on a donc évidemment
$M'=S'\otimes_{B_0}N_0$, et comme $N_0$ est un $B_0$-module de type fini,
$M'$ est un $S'$-module de type fini, d'où notre assertion (I, 1.3.13).

Comme le morphisme $f_0:X_0\to Y_0$ est \emph{propre} par hypothèse, on peut
appliquer (\hyperref[III.3.3.2-fr]{3.3.2}) à $\mathcal S$, $\mathcal M$ et au
morphisme $f_0$~; compte tenu de (\hyperref[III.1.4.11-fr]{1.4.11}), on en
conclut que pour \emph{tout $n\geq0$},
$\bigoplus_{k\geq0}H^n(X_0,\mathcal M_k)$ est un $S$-module gradué \emph{de
type fini}. Cela prouve que la condition $(\mathrm F_n)$ de (0, 13.7.7) est
vérifiée pour \emph{tout $n\geq0$}, lorsqu'on considère le système projectif
strict $(\mathcal F/\mathfrak J^k\mathcal F)_{k\geq0}$ de faisceaux de
groupes abéliens sur $X_0$, munis chacun de sa structure naturelle de
«~$A$-module filtré~». On peut donc appliquer (0, 13.7.7), qui prouve que~:
\begin{enumerate}
\item[$1^\circ$] Le système projectif
$(H^n(\mathfrak X,\mathcal F_k))_{k\geq0}$ vérifie la condition (ML).
\item[$2^\circ$] Si
$H'^{,n}=\varprojlim_kH^n(\mathfrak X,\mathcal F_k)$, $H'^{,n}$ est un
$A$-module de type fini.
\item[$3^\circ$] La filtration définie sur $H'^{,n}$ par les noyaux des
homomorphismes canoniques $H'^{,n}\to H^n(\mathfrak X,\mathcal F_k)$ est
$\mathfrak J$-bonne.
\end{enumerate}

Notons d'autre part que si on pose
$X_k=(\mathfrak X,\mathcal O_{\mathfrak X}/\mathcal K^{k+1})$,
$\mathcal F_k$ est un $\mathcal O_{X_k}$-Module cohérent (I, 10.11.3) et si
$U$ est un ouvert affine dans $X_0$, $U$ est aussi un ouvert affine dans
chacun des $X_k$ (I, 5.1.9), donc $H^n(U,\mathcal F_k)=0$ pour tout $n>0$ et
tout $k$ (\hyperref[III.1.3.1-fr]{1.3.1}) et
$H^0(U,\mathcal F_k)\to H^0(U,\mathcal F_h)$ est surjectif pour $h\leq k$
(I, 1.3.9). On est donc dans les conditions de (0, 13.3.2) et l'application
de (0, 13.3.1) prouve que $H'^{,n}$ s'identifie canoniquement à
$H^n(\mathfrak X,\varprojlim_k\mathcal F_k)=H^n(\mathfrak X,\mathcal F)$~;
cela achève de prouver (\hyperref[III.3.4.4-fr]{3.4.4}).

\begin{env}[3.4.5]
\label{III.3.4.5-fr}
Revenons maintenant à la démonstration de (\hyperref[III.3.4.2-fr]{3.4.2})
et (\hyperref[III.3.4.3-fr]{3.4.3}). Prouvons d'abord ces propositions dans
le cas $\mathfrak Y=\operatorname{Spf}(A)$ envisagé dans
(\hyperref[III.3.4.4-fr]{3.4.4})~; pour cela, pour tout $g\in A$, appliquons
(\hyperref[III.3.4.4-fr]{3.4.4}) au schéma formel affine noethérien induit
sur l'ouvert $\mathfrak Y_g=\mathfrak D(g)$ de $\mathfrak Y$, qui est égal à
$\operatorname{Spf}(A_{\{g\}})$, et au préschéma formel induit par
$\mathfrak X$ sur $f^{-1}(\mathfrak Y_g)$~; notons que $\mathfrak Y_g$ est
aussi un ouvert affine dans le préschéma
$Y_k=(\mathfrak Y,\mathcal O_{\mathfrak Y}/(\mathfrak J^\Delta)^{k+1})$,
et comme $\mathcal F_k$ est un $\mathcal O_{X_k}$-Module cohérent, on a
\[
H^n(f^{-1}(\mathfrak Y_g),\mathcal F_k)
=\Gamma(\mathfrak Y_g,R^nf_*(\mathcal F_k))
\]
pour tout $k\geq0$ en vertu de
(\hyperref[III.1.4.11-fr]{1.4.11}). L'homomorphisme canonique
\[
H^n(f^{-1}(\mathfrak Y_g),\mathcal F)
\longrightarrow\varprojlim_k\Gamma(\mathfrak Y_g,R^nf_*(\mathcal F_k))
\]
est un isomorphisme~; mais on a (0\textsubscript{I}, 3.2.6)
\[
\varprojlim_k\Gamma(\mathfrak Y_g,R^nf_*(\mathcal F_k))
=\Gamma(\mathfrak Y_g,\varprojlim_kR^nf_*(\mathcal F_k))
\]

\oldpage[III]{122}
et comme le faisceau $R^nf_*(\mathcal F)$ est associé au préfaisceau
$\mathfrak Y_g\mapsto H^n(f^{-1}(\mathfrak Y_g),\mathcal F)$ sur les
$\mathfrak Y_g$ (0\textsubscript{I}, 3.2.1), on a bien montré que
l'homomorphisme (\hyperref[III.3.4.2.2-fr]{3.4.2.2}) est bijectif. Prouvons
ensuite que $R^nf_*(\mathcal F)$ est un $\mathcal O_{\mathfrak Y}$-Module
cohérent, et de façon plus précise que l'on a
\[
R^nf_*(\mathcal F)=(H^n(\mathfrak X,\mathcal F))^\Delta.
\tag{3.4.5.1}\label{III.3.4.5.1-fr}
\]
Avec les notations précédentes, on a, puisque $\mathcal F_k$ est un
$\mathcal O_{X_k}$-Module cohérent (\hyperref[III.1.4.13-fr]{1.4.13})
\[
\Gamma(\mathfrak Y_g,R^nf_*(\mathcal F_k))
=(\Gamma(\mathfrak Y,R^nf_*(\mathcal F_k)))_g
=(H^n(\mathfrak X,\mathcal F_k))_g.
\]
Or, les $H^n(\mathfrak X,\mathcal F_k)$ forment un système projectif
vérifiant (ML) et leur limite projective $H^n(\mathfrak X,\mathcal F)$ est
un $A$-module de type fini. On en conclut (0, 13.7.8) que l'on a
\[
\varprojlim_k(H^n(\mathfrak X,\mathcal F_k))_g
=H^n(\mathfrak X,\mathcal F)\otimes_AA_{\{g\}}
=\Gamma(\mathfrak Y_g,(H^n(\mathfrak X,\mathcal F))^\Delta)
\]
compte tenu de (I, 10.10.8) appliqué à $A$ et $A_{\{g\}}$~; ceci démontre
(\hyperref[III.3.4.5.1-fr]{3.4.5.1}) puisque
\[
\Gamma(\mathfrak Y_g,R^nf_*(\mathcal F))
=\varprojlim_k\Gamma(\mathfrak Y_g,R^nf_*(\mathcal F_k)).
\]
Comme (\hyperref[III.3.4.2.2-fr]{3.4.2.2}) est alors un isomorphisme de
$\mathcal O_{\mathfrak Y}$-Modules cohérents, c'est nécessairement un
isomorphisme \emph{topologique} (I, 10.11.6). Enfin, il résulte des relations
$R^nf_*(\mathcal F_k)=(H^n(\mathfrak X,\mathcal F_k))^\Delta$ que le système
projectif $(R^nf_*(\mathcal F_k))_{k\geq0}$ vérifie (ML) (I, 10.10.2).

Une fois (\hyperref[III.3.4.2-fr]{3.4.2}) et
(\hyperref[III.3.4.3-fr]{3.4.3}) démontrés dans le cas où le préschéma formel
$\mathfrak Y$ est affine noethérien, il est immédiat de passer de là au cas
général pour (\hyperref[III.3.4.2-fr]{3.4.2}) et la première assertion de
(\hyperref[III.3.4.3-fr]{3.4.3}), qui sont locales sur $\mathfrak Y$. Quant à
la seconde assertion de (\hyperref[III.3.4.3-fr]{3.4.3}), il suffit,
$\mathfrak Y$ étant noethérien, de le recouvrir par un nombre fini d'ouverts
affines noethériens $U_i$ et de remarquer que les restrictions du système
projectif $(R^qf_*(\mathcal F_k))$ à chacun des $U_i$ vérifient (ML).
\end{env}

Nous avons en outre démontré en cours de route~:

\begin{corollary}[3.4.6]
\label{III.3.4.6-fr}
Sous les hypothèses de (\hyperref[III.3.4.4-fr]{3.4.4}), l'homomorphisme
canonique
\[
H^q(\mathfrak X,\mathcal F)
\longrightarrow\Gamma(\mathfrak Y,R^qf_*(\mathcal F))
\tag{3.4.6.1}\label{III.3.4.6.1-fr}
\]
est bijectif.
\end{corollary}

\section{Le théorème fondamental des morphismes propres. Applications}
\label{section:III.4-fr}

\subsection{Le théorème fondamental}
\label{subsection:III.4.1-fr}

\begin{env}[4.1.1]
\label{III.4.1.1-fr}
Soient $X$, $Y$ deux préschémas usuels noethériens, $f:X\to Y$ un
morphisme propre, $Y'$ une partie fermée de $Y$, $X'$ son image réciproque
$f^{-1}(Y')$. Nous désignerons par $\widehat X$ et $\widehat Y$ les
préschémas formels $X_{/X'}$ et $Y_{/Y'}$, complétés de $X$ et $Y$ le long
de $X'$ et $Y'$ respectivement (I, 10.8.5), par $\widehat f$ le prolongement
de $f$ à ces complétés (I, 10.9.1) qui est un morphisme
$\widehat X\to\widehat Y$ de préschémas formels. Pour tout
$\mathcal O_X$-Module cohérent $\mathcal F$, nous

\oldpage[III]{123}

désignerons par $\widehat{\mathcal F}$ son complété $\mathcal F_{/X'}$ le
long de $X'$ (I, 10.8.4) qui est un $\mathcal O_{\widehat X}$-Module
cohérent (I, 10.8.8).
\end{env}

\begin{env}[4.1.2]
\label{III.4.1.2-fr}
Soit $\mathcal J$ un Idéal cohérent de $\mathcal O_Y$ tel que
$\operatorname{Supp}(\mathcal O_Y/\mathcal J)=Y'$ (I, 5.2.1)~; on sait
(I, 4.4.5) que
\[
\mathcal K=f^*(\mathcal J)\mathcal O_X
\]
est un Idéal cohérent de $\mathcal O_X$ tel que
\[
\operatorname{Supp}(\mathcal O_X/\mathcal K)=X'.
\]
Nous considérerons pour tout $k\geq0$ les $\mathcal O_X$-Modules cohérents
\[
\mathcal F_k
=\mathcal F\otimes_{\mathcal O_Y}(\mathcal O_Y/\mathcal J^{k+1})
=\mathcal F/\mathcal K^{k+1}\mathcal F.
\]
Les $\mathcal O_Y$-Modules $R^nf_*(\mathcal F)$ et
$R^nf_*(\mathcal F_k)$ sont cohérents pour tout $n$
(\hyperref[III.3.2.1-fr]{3.2.1}). Pour tout $k\geq0$ et tout $n$,
l'homomorphisme canonique $\mathcal F\to\mathcal F_k$ définit par
fonctorialité un homomorphisme
\[
R^nf_*(\mathcal F)\longrightarrow R^nf_*(\mathcal F_k).
\tag{4.1.2.1}\label{III.4.1.2.1-fr}
\]
En outre, comme $\mathcal F_k$ est un
$\mathcal O_X/\mathcal K^{k+1}$-Module, $R^nf_*(\mathcal F_k)$ est un
$\mathcal O_Y/\mathcal J^{k+1}$-Module (0, 12.2.1) et on déduit donc de
(\hyperref[III.4.1.2.1-fr]{4.1.2.1}) un homomorphisme
\[
R^nf_*(\mathcal F)\otimes_{\mathcal O_Y}
(\mathcal O_Y/\mathcal J^{k+1})
\longrightarrow R^nf_*(\mathcal F_k).
\tag{4.1.2.2}\label{III.4.1.2.2-fr}
\]
Les deux membres de (\hyperref[III.4.1.2.2-fr]{4.1.2.2}) forment deux
systèmes projectifs, et la limite projective du premier membre n'est autre
que le complété $(R^nf_*(\mathcal F))_{/Y'}$, que nous noterons
$\widehat{R^nf_*(\mathcal F)}$. En outre, il est immédiat que les
homomorphismes (\hyperref[III.4.1.2.2-fr]{4.1.2.2}) forment un système
projectif, d'où par passage à la limite un homomorphisme canonique
\[
\varphi_n:\widehat{R^nf_*(\mathcal F)}
\longrightarrow\varprojlim_kR^nf_*(\mathcal F_k).
\tag{4.1.2.3}\label{III.4.1.2.3-fr}
\]
D'ailleurs (\hyperref[III.4.1.2.2-fr]{4.1.2.2}) est un homomorphisme de
$(\mathcal O_Y/\mathcal J^{k+1})$-Modules, et par suite (I, 10.8.3) peut
être considéré comme un homomorphisme continu de
$\mathcal O_{\widehat Y}$-Modules topologiques pseudo-discrets
(0\textsubscript{I}, 3.8.1). L'homomorphisme $\varphi_n$ est par suite un
homomorphisme continu de $\mathcal O_{\widehat Y}$-Modules topologiques.
\end{env}

\begin{env}[4.1.3]
\label{III.4.1.3-fr}
Soit $i:\widehat X\to X$ le morphisme canonique d'espaces annelés défini
dans (I, 10.8.7), de sorte que l'on a le diagramme commutatif
\[
\xymatrix{
X_k\ar[r]^{h_k}\ar[dr]_{i_k}&\widehat X\ar[d]^i\\
&X
}
\tag{4.1.3.1}\label{III.4.1.3.1-fr}
\]
où $X_k$ est le sous-préschéma fermé de $X$ défini par l'Idéal
$\mathcal K^{k+1}$, $i_k$ l'injection canonique, $h_k$ le morphisme
d'espaces annelés correspondant à l'identité dans les espaces sous-jacents
et à l'homomorphisme canonique
$\mathcal O_{\widehat X}\to\mathcal O_X/\mathcal K^{k+1}$ (I, 10.5.2).
En outre, on a $\widehat{\mathcal F}=i^*(\mathcal F)$ (I, 10.8.8) à un
isomorphisme canonique près. On sait que l'on a
\[
H^n(X_k,i_k^*(\mathcal F_k))=H^n(X,\mathcal F_k)
\tag{4.1.3.2}\label{III.4.1.3.2-fr}
\]

\oldpage[III]{124}

à un isomorphisme canonique près, puisque
$\mathcal F_k=(i_k)_*(i_k^*(\mathcal F_k))$ (G, II, 4.9.1)~;
l'homomorphisme canonique
$H^n(\widehat X,\widehat{\mathcal F})\to
H^n(X_k,h_k^*(\widehat{\mathcal F}))$ (0, 12.1.3.5) s'écrit donc aussi
\[
H^n(\widehat X,\widehat{\mathcal F})\longrightarrow H^n(X,\mathcal F_k)
\tag{4.1.3.3}\label{III.4.1.3.3-fr}
\]
et ces homomorphismes forment évidemment un système projectif, d'où par
passage à la limite, un homomorphisme canonique
\[
\psi_{n,X}:H^n(\widehat X,\widehat{\mathcal F})
\longrightarrow\varprojlim_kH^n(X,\mathcal F_k).
\tag{4.1.3.4}\label{III.4.1.3.4-fr}
\]
Remplaçant $X$ par un ouvert de la forme $f^{-1}(V)$, où $V$ est un
ouvert affine de $Y$, on a, compte tenu de
(\hyperref[III.1.4.11-fr]{1.4.11}), des homomorphismes
\[
\psi_{n,V}:H^n(\widehat X\cap f^{-1}(V),\widehat{\mathcal F})
\longrightarrow\varprojlim_k\Gamma(V,R^nf_*(\mathcal F_k))~;
\tag{4.1.3.5}\label{III.4.1.3.5-fr}
\]
ces homomorphismes commutent de façon évidente à la restriction de $V$ à
un ouvert affine plus petit, donc définissent finalement un homomorphisme
canonique de faisceaux
\[
\psi_n:R^n\widehat f_*(\widehat{\mathcal F})
\longrightarrow\varprojlim_kR^nf_*(\mathcal F_k).
\tag{4.1.3.6}\label{III.4.1.3.6-fr}
\]
\end{env}

\begin{env}[4.1.4]
\label{III.4.1.4-fr}
Soit enfin $j:\widehat Y\to Y$ le morphisme canonique d'espaces annelés
(I, 10.8.7)~; comme $R^nf_*(\mathcal F)$ est un $\mathcal O_Y$-Module
cohérent (\hyperref[III.3.2.1-fr]{3.2.1}), on a
$j^*(R^nf_*(\mathcal F))=\widehat{R^nf_*(\mathcal F)}$ à un isomorphisme
canonique près (I, 10.8.8) et on a donc un homomorphisme canonique
\[
\rho_n:\widehat{R^nf_*(\mathcal F)}=j^*(R^nf_*(\mathcal F))
\longrightarrow R^n\widehat f_*(i^*(\mathcal F))
=R^n\widehat f_*(\widehat{\mathcal F})
\tag{4.1.4.1}\label{III.4.1.4.1-fr}
\]
défini de façon générale pour les espaces annelés (voir la démonstration
de (\hyperref[III.1.4.15-fr]{1.4.15})). Montrons que le diagramme
\[
\xymatrix{
\widehat{R^nf_*(\mathcal F)}\ar[rr]^{\rho_n}\ar[dr]_{\varphi_n}
&&R^n\widehat f_*(\widehat{\mathcal F})\ar[dl]^{\psi_n}\\
&\displaystyle\varprojlim_kR^nf_*(\mathcal F_k)&
}
\tag{4.1.4.2}\label{III.4.1.4.2-fr}
\]
est commutatif. Il suffit évidemment de prouver la commutativité du
diagramme correspondant d'homomorphismes de préfaisceaux, donc on peut se
borner au cas où $Y$ est affine, et tout revient à prouver que le diagramme
\[
\xymatrix{
\widehat{H^n(X,\mathcal F)}\ar[rr]^{\rho_n}\ar[dr]_{\varphi_n}
&&H^n(\widehat X,\widehat{\mathcal F})\ar[dl]^{\psi_{n,X}}\\
&\displaystyle\varprojlim_kH^n(X,\mathcal F_k)&
}
\tag{4.1.4.3}\label{III.4.1.4.3-fr}
\]
est commutatif. Mais la commutativité de
(\hyperref[III.4.1.3.1-fr]{4.1.3.1}) et les relations vues dans
(\hyperref[III.4.1.3-fr]{4.1.3}) entre les groupes de cohomologie donnent
aussitôt le diagramme commutatif

\oldpage[III]{125}

\[
\xymatrix{
H^n(X,\mathcal F)\ar[r]\ar[dr]&
H^n(\widehat X,\widehat{\mathcal F})
=H^n(\widehat X,i^*(\mathcal F))\ar[d]\\
&H^n(X,\mathcal F_k)=H^n(X_k,i_k^*(\mathcal F))
}
\]
d'où on déduit aussitôt la commutativité de
(\hyperref[III.4.1.4.3-fr]{4.1.4.3}).
\end{env}

\begin{theorem}[4.1.5]
\label{III.4.1.5-fr}
Soient $f:X\to Y$ un morphisme propre de préschémas noethériens, $Y'$ une
partie fermée de $Y$, $X'=f^{-1}(Y')$. Alors, pour tout
$\mathcal O_X$-Module cohérent $\mathcal F$,
$R^n\widehat f_*(\widehat{\mathcal F})$ est un
$\mathcal O_{\widehat X}$-Module cohérent et les homomorphismes
$\varphi_n$, $\psi_n$ et $\rho_n$ du diagramme
(\hyperref[III.4.1.4.2-fr]{4.1.4.2}) sont des isomorphismes topologiques.
\end{theorem}

\begin{proof}
Il suffira évidemment de prouver que $\varphi_n$ et $\psi_n$ sont des
isomorphismes~; comme $R^nf_*(\mathcal F)$ est cohérent
(\hyperref[III.3.2.1-fr]{3.2.1}), il en résultera que
$\widehat{R^nf_*(\mathcal F)}$ est cohérent (I, 10.8.8) et la bicontinuité
de $\varphi_n$, $\psi_n$ et $\rho_n$ est alors automatique (I, 10.11.6).
\end{proof}

\begin{env}[4.1.6]
\label{III.4.1.6-fr}
\emph{Remarques.}
\begin{enumerate}
\item[(i)] Si on pose
$\widehat{\mathcal F}_k=
\widehat{\mathcal F}/\widehat{\mathcal K}^{k+1}
\widehat{\mathcal F}$, il est immédiat que
$\mathcal F_k=i_*(\widehat{\mathcal F}_k)$, et l'homomorphisme canonique
(\hyperref[III.4.1.3.6-fr]{4.1.3.6}) n'est autre que l'homomorphisme déjà
défini en (\hyperref[III.3.4.2.2-fr]{3.4.2.2})
\[
R^n\widehat f_*(\widehat{\mathcal F})
\longrightarrow\varprojlim_kR^n\widehat f_*(\widehat{\mathcal F}_k)~;
\tag{4.1.6.1}\label{III.4.1.6.1-fr}
\]
par suite le fait que $\psi_n$ soit un isomorphisme est un cas particulier
de (\hyperref[III.3.4.3-fr]{3.4.3}). Mais nous en donnerons ci-dessous une
démonstration directe, évitant les considérations délicates sur les limites
projectives de suites spectrales (0, 13.7), sur lesquelles repose le
théorème général (\hyperref[III.3.4.3-fr]{3.4.3}).

\item[(ii)] Compte tenu du fait que les $\psi_n$ sont des isomorphismes, il
est équivalent de dire que les $\varphi_n$ ou les
$\rho_n=\psi_n^{-1}\circ\varphi_n$ sont des isomorphismes. Le th.
(\hyperref[III.4.1.5-fr]{4.1.5}) exprime entre autres que la formation des
$R^nf_*$ commute à la complétion et pourra s'appeler le premier théorème de
comparaison de la théorie «~algébrique~» à la théorie «~formelle~».
\end{enumerate}
Nous commencerons par établir la forme affine de
(\hyperref[III.4.1.5-fr]{4.1.5})~:
\end{env}

\begin{corollary}[4.1.7]
\label{III.4.1.7-fr}
Les hypothèses étant celles de (\hyperref[III.4.1.5-fr]{4.1.5}), supposons
en outre $Y=\operatorname{Spec}(A)$, où $A$ est noethérien, et
$\mathcal J=\widetilde{\mathfrak J}$, où $\mathfrak J$ est un idéal de
$A$, de sorte que
$\mathcal F_k=\mathcal F/\mathfrak J^{k+1}\mathcal F$.
L'homomorphisme canonique
\[
\varphi_n:\widehat{H^n(X,\mathcal F)}
\longrightarrow\varprojlim_kH^n(X,\mathcal F_k)
\tag{4.1.7.1}\label{III.4.1.7.1-fr}
\]
(où le premier membre est le séparé complété de $H^n(X,\mathcal F)$ pour
la topologie $\mathfrak J$-préadique) est un isomorphisme. Le système
projectif $(H^n(X,\mathcal F_k))_{k\geq0}$ vérifie pour tout $n$ la
condition (ML), et l'homomorphisme canonique
\[
\psi_n:H^n(\widehat X,\widehat{\mathcal F})
\longrightarrow\varprojlim_kH^n(X,\mathcal F_k)
\tag{4.1.7.2}\label{III.4.1.7.2-fr}
\]

\oldpage[III]{126}

est un isomorphisme. Enfin, la filtration sur $H^n(X,\mathcal F)$, définie
par les noyaux des homomorphismes canoniques
$H^n(X,\mathcal F)\to H^n(X,\mathcal F_k)$, est $\mathfrak J$-bonne
(0, 13.7.7) et $\varphi_n$ est un isomorphisme topologique.
\footnote{La démonstration qui suit, plus simple que la démonstration
initiale, et le complément relatif à la filtration de
$H^n(X,\mathcal F)$, nous ont été communiqués par J.-P. Serre.}
\end{corollary}

\begin{proof}
L'entier $n\geq0$ étant fixé dans cette démonstration, nous poserons pour
simplifier
\[
\mathrm H=H^n(X,\mathcal F),\qquad
\mathrm H_k=H^n(X,\mathcal F_k)
\tag{4.1.7.3}\label{III.4.1.7.3-fr}
\]
et
\[
\mathrm R_k=\operatorname{Ker}(\mathrm H\longrightarrow\mathrm H_k),
\qquad\text{sous-$A$-module de $\mathrm H$.}
\tag{4.1.7.4}\label{III.4.1.7.4-fr}
\]
La suite exacte de cohomologie
\[
H^n(X,\mathfrak J^{k+1}\mathcal F)\longrightarrow H^n(X,\mathcal F)
\longrightarrow H^n(X,\mathcal F_k)\xrightarrow{\partial}
H^{n+1}(X,\mathfrak J^{k+1}\mathcal F)
\longrightarrow H^{n+1}(X,\mathcal F)
\]
montre que l'on a aussi
\[
\mathrm R_k=\operatorname{Im}\bigl(
H^n(X,\mathfrak J^{k+1}\mathcal F)\longrightarrow H^n(X,\mathcal F)
\bigr)~;
\]
nous poserons
\[
\begin{aligned}
\mathrm Q_k
&=\operatorname{Ker}\bigl(H^{n+1}(X,\mathfrak J^{k+1}\mathcal F)
\longrightarrow H^{n+1}(X,\mathcal F)\bigr)\\
&=\operatorname{Im}\bigl(H^n(X,\mathcal F_k)
\longrightarrow H^{n+1}(X,\mathfrak J^{k+1}\mathcal F)\bigr).
\end{aligned}
\tag{4.1.7.5}\label{III.4.1.7.5-fr}
\]
On a donc la suite exacte
\[
0\longrightarrow\mathrm R_k\longrightarrow\mathrm H
\longrightarrow\mathrm H_k\longrightarrow\mathrm Q_k\longrightarrow0.
\tag{4.1.7.6}\label{III.4.1.7.6-fr}
\]

\begin{env}[4.1.7.7]
\label{III.4.1.7.7-fr}
Soit $x$ un élément de $\mathfrak J^m$ ($m\geq0$)~; la multiplication par
$x$ dans $\mathfrak J^k\mathcal F$ est un homomorphisme
$\mathfrak J^k\mathcal F\to\mathfrak J^{k+m}\mathcal F$, et donne lieu
par suite à un homomorphisme
\[
\mu_{x,m}:H^n(X,\mathfrak J^k\mathcal F)
\longrightarrow H^n(X,\mathfrak J^{k+m}\mathcal F).
\tag{4.1.7.8}\label{III.4.1.7.8-fr}
\]
Si on désigne par $S$ la $A$-algèbre graduée
$\bigoplus_{k\geq0}\mathfrak J^k$, on sait que les multiplications
$\mu_{x,m}$ définissent sur
$\mathrm E=\bigoplus_{k\geq0}H^n(X,\mathfrak J^k\mathcal F)$ une
structure de module gradué de type fini sur l'anneau gradué $S$
(\hyperref[III.3.3.2-fr]{3.3.2}), qui est noethérien (II, 2.1.5).
\end{env}

\begin{lemma}[4.1.7.9]
\label{III.4.1.7.9-fr}
Les sous-modules $\mathrm R_k$ de $\mathrm H$ définissent sur $\mathrm H$
une filtration $\mathfrak J$-bonne.
\end{lemma}

\begin{proof}
En premier lieu, montrons que l'on a
\[
\mathfrak J^m\mathrm R_k\subset\mathrm R_{k+m},
\tag{4.1.7.10}\label{III.4.1.7.10-fr}
\]
la multiplication dans $\mathrm H=H^n(X,\mathcal F)$ par un élément
$x\in\mathfrak J^m$ étant donc l'application $\mu_{x,0}$. Pour tout
$x\in\mathfrak J^m$, le diagramme
\[
\xymatrix{
\mathfrak J^{k+1}\mathcal F\ar[r]^x\ar[d]&
\mathfrak J^{k+m+1}\mathcal F\ar[d]\\
\mathcal F\ar[r]_x&\mathcal F
}
\tag{4.1.7.11}\label{III.4.1.7.11-fr}
\]

\oldpage[III]{127}

(où les flèches horizontales sont la multiplication par $x$, et les
flèches verticales les injections canoniques) est commutatif~; donc le
diagramme correspondant
\[
\xymatrix{
H^n(X,\mathfrak J^{k+1}\mathcal F)\ar[r]^{\mu_{x,m}}\ar[d]&
H^n(X,\mathfrak J^{k+m+1}\mathcal F)\ar[d]\\
H^n(X,\mathcal F)\ar[r]_{\mu_{x,0}}&H^n(X,\mathcal F)
}
\]
est commutatif, ce qui, compte tenu de l'interprétation de $\mathrm R_k$
comme image de
$H^n(X,\mathfrak J^{k+1}\mathcal F)\to H^n(X,\mathcal F)$, prouve
(\hyperref[III.4.1.7.10-fr]{4.1.7.10}) et montre en outre que le $S$-module
gradué $\mathrm R=\bigoplus_{k\geq0}\mathrm R_k$ est un quotient du
sous-$S$-module
$\mathrm M=\bigoplus_{k\geq0}H^n(X,\mathfrak J^{k+1}\mathcal F)$ de
$\mathrm E$~; la remarque faite plus haut montre donc que $\mathrm R$ est
un $S$-module de type fini, ce qui équivaut à l'assertion de
(\hyperref[III.4.1.7.9-fr]{4.1.7.9}) (Bourbaki, \emph{Alg. comm.},
chap. III, §~3, n\textsuperscript{o}~1, th.~1).
\end{proof}

\begin{env}[4.1.7.12]
\label{III.4.1.7.12-fr}
Considérons maintenant le $S$-module gradué
\[
\mathrm N=\bigoplus_{k\geq0}H^{n+1}(X,\mathfrak J^{k+1}\mathcal F)
\]
défini comme dans (4.7.1.8)~; c'est encore un $S$-module de type fini en
vertu de (\hyperref[III.3.3.2-fr]{3.3.2})~; on a
$\mathrm Q_k\subset\mathrm N_k$ pour tout $k$ par
(\hyperref[III.4.1.7.5-fr]{4.1.7.5}) et le diagramme
(\hyperref[III.4.1.7.11-fr]{4.1.7.11}) où on remplace $n$ par $n+1$,
montre que
$S_m\mathrm Q_k=\mathfrak J^m\mathrm Q_k\subset\mathrm Q_{k+m}$.
Autrement dit, $\mathrm Q$ est un sous-$S$-module gradué de $\mathrm N$,
et par suite est de type fini.
\end{env}

\begin{env}[4.1.7.13]
\label{III.4.1.7.13-fr}
Désignons par $\alpha_m$ l'injection canonique
$\mathfrak J^m\to A$, qu'on peut écrire $S_m\to S_0$. Comme
$\mathfrak J^{k+1}\mathcal F_k=0$, le $A$-module
$H^n(X,\mathcal F_k)$ est annulé par $\mathfrak J^{k+1}$~; comme
$\mathrm Q_k$ est l'image du $A$-homomorphisme
$H^n(X,\mathcal F_k)\to H^{n+1}(X,\mathfrak J^{k+1}\mathcal F)$,
$\mathrm Q_k$, en tant que $A$-module, est aussi annulé par
$\mathfrak J^{k+1}$~; cela signifie encore que, dans le $S$-module
$\mathrm Q$, on a
\[
\alpha_{k+1}(S_{k+1})\mathrm Q_k=0.
\tag{4.1.7.14}\label{III.4.1.7.14-fr}
\]
Comme $\mathrm Q$ est un $S$-module de type fini, il existe un entier
$k_0$ et un entier $h$ tels que
$\mathrm Q_{k+h}=S_h\mathrm Q_k$ pour $k\geq k_0$ (II, 2.1.6, (ii))~;
on déduit de cette relation et de
(\hyperref[III.4.1.7.14-fr]{4.1.7.14}) qu'il existe un entier $r>0$ tel
que
\[
\alpha_r(S_r)\mathrm Q=0.
\tag{4.1.7.15}\label{III.4.1.7.15-fr}
\]
\end{env}

\begin{env}[4.1.7.16]
\label{III.4.1.7.16-fr}
Notons maintenant que l'injection canonique
$\mathfrak J^{k+m}\mathcal F\to\mathfrak J^k\mathcal F$ donne en passant
à la cohomologie un $A$-homomorphisme
\[
\nu_m:H^{n+1}(X,\mathfrak J^{k+m}\mathcal F)
\longrightarrow H^{n+1}(X,\mathfrak J^k\mathcal F)
\tag{4.1.7.17}\label{III.4.1.7.17-fr}
\]
et, pour tout $x\in\mathfrak J^m$, on a évidemment la factorisation
\[
\xymatrix@C=4.5em{
H^{n+1}(X,\mathfrak J^k\mathcal F)\ar[r]^{\mu_{x,m}}&
H^{n+1}(X,\mathfrak J^{k+m}\mathcal F)\ar[r]^{\nu_m}&
H^{n+1}(X,\mathfrak J^k\mathcal F)
}
\tag{4.1.7.18}\label{III.4.1.7.18-fr}
\]
dont le composé est $\mu_{x,0}$.
\end{env}

\oldpage[III]{128}

d'où on conclut que, pour tout sous-$A$-module $P$ de
$H^{n+1}(X,\mathfrak J^k\mathcal F)$, on a, dans le $S$-module
$\mathrm N$,
\[
\nu_m(S_mP)=\alpha_m(S_m)P.
\tag{4.1.7.19}\label{III.4.1.7.19-fr}
\]

\begin{lemma}[4.1.7.20]
\label{III.4.1.7.20-fr}
Il existe un entier $m>0$ tel que $\nu_m(\mathrm Q_{k+m})=0$ pour tout
$k\geq k_0$.
\end{lemma}

\begin{proof}
Prenons en effet pour $m$ un multiple de $h$ qui soit $\geq r$~; comme
$\mathrm Q_{k+m}=S_m\mathrm Q_k$ pour $k\geq k_0$, on a en vertu de
(\hyperref[III.4.1.7.19-fr]{4.1.7.19}) et
(\hyperref[III.4.1.7.15-fr]{4.1.7.15})
\[
\nu_m(\mathrm Q_{k+m})=\alpha_m(S_m)\mathrm Q_k
\subset\alpha_r(S_r)\mathrm Q_k=0.
\]
\end{proof}

\begin{env}[4.1.7.21]
\label{III.4.1.7.21-fr}
Remarquons que du diagramme commutatif
\[
\xymatrix@C=2.5em{
H^n(X,\mathcal F)\ar[r]&H^n(X,\mathcal F_k)\ar[r]^-\partial&
H^{n+1}(X,\mathfrak J^{k+1}\mathcal F)\ar[r]&H^{n+1}(X,\mathcal F)\\
H^n(X,\mathcal F)\ar[r]\ar[u]&H^n(X,\mathcal F_{k+m})\ar[r]^-\partial\ar[u]&
H^{n+1}(X,\mathfrak J^{k+m+1}\mathcal F)\ar[r]\ar[u]&
H^{n+1}(X,\mathcal F)\ar[u]
}
\]
provenant lui-même du diagramme commutatif
\[
\xymatrix{
0\ar[r]&\mathfrak J^{k+1}\mathcal F\ar[r]&\mathcal F\ar[r]&
\mathcal F_k\ar[r]&0\\
0\ar[r]&\mathfrak J^{k+m+1}\mathcal F\ar[r]\ar[u]&
\mathcal F\ar[r]\ar[u]&\mathcal F_{k+m}\ar[r]\ar[u]&0
}
\]
où les flèches verticales sont les applications canoniques, on déduit un
diagramme commutatif
\[
\xymatrix@C=2.2em{
0\ar[r]&\mathrm R_k\ar[r]&\mathrm H\ar[r]&\mathrm H_k\ar[r]&
\mathrm Q_k\ar[r]&0\\
0\ar[r]&\mathrm R_{k+m}\ar[r]\ar[u]&\mathrm H\ar[r]\ar@{=}[u]&
\mathrm H_{k+m}\ar[r]\ar[u]&\mathrm Q_{k+m}\ar[r]\ar[u]^{\nu_m}&0
}
\]
où les lignes sont exactes. Comme la dernière flèche verticale est nulle
pour $k\geq k_0$ (\hyperref[III.4.1.7.20-fr]{4.1.7.20}), l'image de
$\mathrm H_{k+m}$ dans $\mathrm H_k$ est contenue dans
$\operatorname{Ker}(\mathrm H_k\to\mathrm Q_k)
=\operatorname{Im}(\mathrm H\to\mathrm H_k)$, mais par ailleurs elle
contient $\operatorname{Im}(\mathrm H\to\mathrm H_k)$ par la commutativité
du diagramme, donc elle lui est égale~; il en est donc de même des images
dans $\mathrm H_k$ des $\mathrm H_{k'}$ pour $k'\geq k+m$, ce qui démontre
la condition (ML) pour le système projectif $(\mathrm H_k)_{k\geq0}$. En
outre, pour tout ouvert affine $U$ de $X$, on a
$H^i(U,\mathcal F_k)=0$ pour $i>0$
(\hyperref[III.1.3.1-fr]{1.3.1}), et pour $m>0$, l'application
$H^0(U,\mathcal F_{k+m})\to H^0(U,\mathcal F_k)$ est surjective
(I, 1.3.9). On peut donc appli-

\oldpage[III]{129}

quer (0, 13.3.1), et l'homomorphisme canonique
$H^n(\widehat X,\widehat{\mathcal F})\to
\varprojlim_kH^n(X,\mathcal F_k)$ est bijectif pour tout $n\geq0$.

Comme le système projectif $(\mathrm H/\mathrm R_k)_{k\geq0}$ est strict,
on peut passer à la limite projective dans les suites exactes
\[
0\longrightarrow\mathrm H/\mathrm R_k\longrightarrow\mathrm H_k
\longrightarrow\mathrm Q_k\longrightarrow0
\tag{4.1.7.22}\label{III.4.1.7.22-fr}
\]
(0, 13.2.2)~; comme $\nu_m(\mathrm Q_{k+m})=0$, on a
$\varprojlim_k\mathrm Q_k=0$, d'où un isomorphisme topologique
\[
\varprojlim_k(\mathrm H/\mathrm R_k)\xrightarrow{\sim}
\varprojlim_k\mathrm H_k.
\]
Mais comme la filtration $(\mathrm R_k)$ de $\mathrm H$ est
$\mathfrak J$-bonne, elle définit sur $\mathrm H$ la topologie
$\mathfrak J$-préadique~; donc $\varprojlim_k(\mathrm H/\mathrm R_k)$ est
le séparé complété de $\mathrm H$ pour la topologie
$\mathfrak J$-préadique, ce qui achève de démontrer
(\hyperref[III.4.1.7-fr]{4.1.7}).
\end{env}
\end{proof}

\begin{env}[4.1.8]
\label{III.4.1.8-fr}
Passons enfin à la démonstration de (\hyperref[III.4.1.5-fr]{4.1.5})~:
pour tout ouvert affine $V$ de $Y$,
$\Gamma(V,\widehat{R^nf_*(\mathcal F)})$ est le séparé complété de
$\Gamma(V,R^nf_*(\mathcal F))$ pour la topologie $\mathfrak J$-préadique
(si $\mathcal J|V=\widetilde{\mathfrak J}$) puisque
$R^nf_*(\mathcal F)$ est un $\mathcal O_Y$-Module cohérent (I, 10.8.4),
et
\[
\Gamma\bigl(V,\varprojlim_kR^nf_*(\mathcal F_k)\bigr)
=\varprojlim_k\Gamma(V,R^nf_*(\mathcal F_k))
\]
(0\textsubscript{I}, 3.2.6)~; le fait que $\varphi_n$ soit un isomorphisme
topologique résulte alors de (\hyperref[III.4.1.7-fr]{4.1.7}) et de
(\hyperref[III.1.4.11-fr]{1.4.11}). D'autre part (toujours en vertu de
(\hyperref[III.1.4.11-fr]{1.4.11})), il résulte de
(\hyperref[III.4.1.7-fr]{4.1.7}) que l'homomorphisme $\psi_{n,V}$ de
(\hyperref[III.4.1.3.3-fr]{4.1.3.3}) est un isomorphisme, donc $\psi_n$
est un isomorphisme par définition de
$R^n\widehat f_*(\widehat{\mathcal F})$.
\end{env}

\begin{corollary}[4.1.9]
\label{III.4.1.9-fr}
Sous les hypothèses de (\hyperref[III.4.1.4-fr]{4.1.4}), pour tout ouvert
affine $V$ de $Y$, l'homomorphisme canonique
\[
H^n(\widehat X\cap f^{-1}(V),\widehat{\mathcal F})
\longrightarrow\Gamma(\widehat Y\cap V,
R^n\widehat f_*(\widehat{\mathcal F}))
\]
est bijectif.
\end{corollary}

\begin{env}[4.1.10]
\label{III.4.1.10-fr}
\emph{Remarque.}
Soit $f:X\to Y$ un morphisme de type fini de préschémas (usuels)
noethériens, et soit $\mathcal F$ un $\mathcal O_X$-Module cohérent dont le
support soit propre sur $Y$ (II, 5.4.10). On sait alors
(\hyperref[III.3.2.4-fr]{3.2.4}) que $R^nf_*(\mathcal F)$ est un
$\mathcal O_Y$-Module cohérent pour tout $n\geq0$. En outre, on peut
toujours supposer que $\mathcal F=u_*(\mathcal G)$, où
$\mathcal G=u^*(\mathcal F)$ est un $\mathcal O_Z$-Module cohérent, $Z$
désignant un sous-préschéma fermé convenable de $X$ dont l'espace
sous-jacent est $\operatorname{Supp}(\mathcal F)$, et $u:Z\to X$
l'injection canonique (I, 9.3.5). Si on pose
\[
\mathcal G_k=\mathcal G\otimes_{\mathcal O_Y}
(\mathcal O_Y/\mathcal J^{k+1}),
\]
on a $\mathcal G_k=u^*(\mathcal F_k)$,
$R^nf_*(\mathcal F_k)=R^n(f\circ u)_*(\mathcal G_k)$ et
$R^nf_*(\mathcal F)=R^n(f\circ u)_*(\mathcal G)$
(\hyperref[III.1.3.4-fr]{1.3.4}), et enfin, compte tenu de (I, 10.9.5),
\[
R^n\widehat f_*(\widehat{\mathcal F})
=R^n\widehat{(f\circ u)}_*(\widehat{\mathcal G}).
\]
On peut alors appliquer (\hyperref[III.4.1.5-fr]{4.1.5}) à $\mathcal G$ et
au morphisme propre $f\circ u$, et on en conclut que sous ces hypothèses,
les résultats de (\hyperref[III.4.1.5-fr]{4.1.5}) sont valables pour
$\mathcal F$ et $f$.
\end{env}

\subsection{Cas particuliers et variantes}
\label{subsection:III.4.2-fr}

La forme la plus utile du th. de comparaison
(\hyperref[III.4.1.5-fr]{4.1.5}) est la suivante~:

\begin{proposition}[4.2.1]
\label{III.4.2.1-fr}
Soient $Y$ un préschéma localement noethérien, $f:X\to Y$ un morphisme
propre, $\mathcal F$ un $\mathcal O_X$-Module cohérent. Alors, pour tout
$y\in Y$ et tout $p$, $(R^pf_*(\mathcal F))_y$ est

\oldpage[III]{130}

un $\mathcal O_y$-module de type fini, donc séparé pour la topologie
$\mathfrak m_y$-préadique, et on a un isomorphisme topologique canonique
\[
\widehat{(R^pf_*(\mathcal F))_y}
\xrightarrow{\sim}
\varprojlim_kH^p\bigl(f^{-1}(y),
\mathcal F\otimes_{\mathcal O_Y}(\mathcal O_y/\mathfrak m_y^k)\bigr)
\tag{4.2.1.1}\label{III.4.2.1.1-fr}
\]
où le premier membre est le complété de $(R^pf_*(\mathcal F))_y$ pour la
topologie $\mathfrak m_y$-préadique, et au second membre $f^{-1}(y)$ est
considéré, pour tout $k\geq0$, comme espace sous-jacent au préschéma
$X\times_Y\operatorname{Spec}(\mathcal O_y/\mathfrak m_y^k)$
(I, 3.6.1).
\end{proposition}

\begin{proof}
Comme $\mathcal O_y$ est un anneau local noethérien et
$(R^pf_*(\mathcal F))_y$ un $\mathcal O_y$-module de type fini
(\hyperref[III.3.2.1-fr]{3.2.1}), la topologie $\mathfrak m_y$-préadique
sur $(R^pf_*(\mathcal F))_y$ est séparée (0\textsubscript{I}, 7.3.5). Les
autres assertions sont conséquences de (\hyperref[III.4.1.7-fr]{4.1.7})
lorsque $Y$ est noethérien et le point $y$ fermé, en remplaçant $Y$ par un
voisinage affine de $y$ et prenant $Y'=\{y\}$, compte tenu de
(G, II, 4.9.1). Dans le cas général, posons
\[
Y_1=\operatorname{Spec}(\mathcal O_y),\qquad
X_1=X\times_YY_1,\qquad
\mathcal F_1=\mathcal F\otimes_{\mathcal O_Y}\mathcal O_{Y_1}
\]
et soit $f_1=f\times1_{Y_1}:X_1\to Y_1$~; $Y_1$ est noethérien et $f_1$
propre (II, 5.4.2, (iii)) et $\mathcal F_1$ est cohérent
(0\textsubscript{I}, 5.3.11). Soit $y_1$ l'unique point fermé de $Y_1$~;
la proposition est valable pour $f_1$, $\mathcal F_1$ et $y_1$~; on a
$\mathcal O_{y_1}=\mathcal O_y$, $f_1^{-1}(y_1)=f^{-1}(y)$ (I, 3.6.5),
les préschémas
$X\times_Y\operatorname{Spec}(\mathcal O_y/\mathfrak m_y^k)$ et
$X_1\times_{Y_1}\operatorname{Spec}(\mathcal O_y/\mathfrak m_y^k)$ étant
canoniquement identifiés (I, 3.3.9)~; en outre,
$\mathcal F_1\otimes_{\mathcal O_{Y_1}}
(\mathcal O_y/\mathfrak m_y^k)$ s'identifie à
$\mathcal F\otimes_{\mathcal O_Y}(\mathcal O_y/\mathfrak m_y^k)$
(I, 9.1.6). Il reste à voir que $R^pf_{1*}(\mathcal F_1)$ est
canoniquement isomorphe à
$R^pf_*(\mathcal F)\otimes_{\mathcal O_Y}\mathcal O_{Y_1}$, ce qui
résulte de (\hyperref[III.1.4.15-fr]{1.4.15}), le morphisme local
$\operatorname{Spec}(\mathcal O_y)\to Y$ étant plat
(0\textsubscript{I}, 6.7.1 et I, 2.4.2).
\end{proof}

Le corollaire suivant utilise la terminologie de la théorie de la dimension
(chap. IV) et ne sera pas appliqué avant le chap. IV.

\begin{corollary}[4.2.2]
\label{III.4.2.2-fr}
Soient $Y$ un préschéma localement noethérien, $f:X\to Y$ un morphisme
propre, $y$ un point de $Y$, $r$ la dimension de $f^{-1}(y)$. Alors, pour
tout $\mathcal O_X$-Module cohérent $\mathcal F$, les faisceaux
$R^pf_*(\mathcal F)$ sont nuls au voisinage de $y$ pour tout $p>r$.
\end{corollary}

\begin{proof}
En effet, on a alors
\[
H^p\bigl(f^{-1}(y),
\mathcal F\otimes_{\mathcal O_Y}(\mathcal O_y/\mathfrak m_y^k)\bigr)=0
\]
(G, II, 4.15.2) pour tout $k$, donc
(\hyperref[III.4.2.1-fr]{4.2.1}) le séparé complété de
$(R^pf_*(\mathcal F))_y$ pour la topologie $\mathfrak m_y$-préadique est
nul, et comme cette topologie est séparée, on a aussi
$(R^pf_*(\mathcal F))_y=0$~; d'où la conclusion, puisque
$R^pf_*(\mathcal F)$ est cohérent (0\textsubscript{I}, 5.2.2).
\end{proof}

\begin{env}[4.2.3]
\label{III.4.2.3-fr}
Le résultat (\hyperref[III.4.2.1-fr]{4.2.1}) est surtout employé pour
$p=0$~; on obtient donc le corollaire suivant~:
\end{env}

\begin{corollary}[4.2.4]
\label{III.4.2.4-fr}
Sous les hypothèses de (\hyperref[III.4.2.1-fr]{4.2.1}), on a un
isomorphisme canonique topologique
\[
\widehat{(f_*(\mathcal F))_y}
\xrightarrow{\sim}
\varprojlim_k\Gamma\bigl(f^{-1}(y),
\mathcal F_y/\mathfrak m_y^k\mathcal F_y\bigr).
\]
\end{corollary}

\subsection{Le théorème de connexion de Zariski}
\label{subsection:III.4.3-fr}
\label{III.4.3-fr}

Les résultats de ce numéro et du suivant généralisent des théorèmes bien
connus de Zariski, et peuvent tous se déduire de
(\hyperref[III.4.2.4-fr]{4.2.4}). Ils sont conséquences du th. suivant~:

\begin{theorem}[4.3.1]
\label{III.4.3.1-fr}
\textnormal{(théorème de connexion).}
Soient $Y$ un préschéma localement noethérien, $f:X\to Y$ un morphisme
propre. Alors $\mathcal A(X)=f_*(\mathcal O_X)$ est une
$\mathcal O_Y$-Algèbre cohérente.

\oldpage[III]{131}

Soit $Y'$ le $Y$-schéma fini au-dessus de $Y$ tel que
$\mathcal A(Y')=\mathcal A(X)$, qui est déterminé à un $Y$-isomorphisme
près (II, 1.3.1 et 6.1.3)~; si $f'=\mathcal A(e)$ est le $Y$-morphisme
$X\to Y'$ déduit de l'isomorphisme identique
$e:\mathcal A(Y')\to\mathcal A(X)$ (II, 1.2.7), $f'$ est propre,
$f'_*(\mathcal O_X)$ est isomorphe à $\mathcal O_{Y'}$, et les fibres
$f'^{-1}(y')$ du morphisme $f'$ sont connexes et non vides pour tout
$y'\in Y'$.
\end{theorem}

\begin{proof}
Soit $g:Y'\to Y$ le morphisme structural. Pour prouver que l'homomorphisme
$\theta:\mathcal O_{Y'}\to f'_*(\mathcal O_X)$ entrant dans la définition
du morphisme $f'$ est bijectif, il suffit, puisque $Y'$ est affine sur $Y$,
de prouver que
\[
g_*(\theta):g_*(\mathcal O_{Y'})\longrightarrow
g_*(f'_*(\mathcal O_X))=f_*(\mathcal O_X)
\]
est l'identité (II, 1.4.2)~; mais cela résulte des définitions puisque
$g_*(\mathcal O_{Y'})=\mathcal A(Y')$ et
$f_*(\mathcal O_X)=\mathcal A(X)$. Le fait que $\mathcal A(X)$ est
cohérent est un cas particulier du th. de finitude
(\hyperref[III.3.2.1-fr]{3.2.1}). Comme $f$ est propre et $g$ séparé,
$f'$ est propre (II, 5.4.3, (i))~; pour terminer la démonstration de
(\hyperref[III.4.3.1-fr]{4.3.1}), il suffit donc d'en prouver le
\end{proof}

\begin{corollary}[4.3.2]
\label{III.4.3.2-fr}
Sous les hypothèses de (\hyperref[III.4.3.1-fr]{4.3.1}), supposons de plus
que $f_*(\mathcal O_X)$ soit isomorphe à $\mathcal O_Y$. Alors les fibres
$f^{-1}(y)$ de $f$ sont connexes et non vides pour tout $y\in Y$.
\end{corollary}

\begin{proof}
L'hypothèse que $f_*(\mathcal O_X)$ est isomorphe à $\mathcal O_Y$
entraîne déjà que $f$ est dominant, donc surjectif puisque $f$ est une
application fermée. On peut se ramener, comme dans
(\hyperref[III.4.2.1-fr]{4.2.1}), au cas où $y$ est fermé dans $Y$~;
$f^{-1}(y)$ étant un espace noethérien, a un nombre fini de composantes
connexes, et c'est l'espace sous-jacent du complété $\widehat X$ le long de
$f^{-1}(y)$. Si $Z_i$ ($1\leq i\leq n$) sont ces composantes connexes, il
est clair que $\Gamma(\widehat X,\mathcal O_{\widehat X})$ est composé
direct des anneaux $\Gamma(Z_i,\mathcal O_{\widehat X})$, et chacun de ces
derniers n'est pas réduit à $0$, puisque la section unité est distincte de
$0$ en chaque point de $\widehat X$. Or, si on applique
(\hyperref[III.4.1.5-fr]{4.1.5}) à $\mathcal F=\mathcal O_X$, dont le
complété le long de $f^{-1}(y)$ est $\mathcal O_{\widehat X}$, on voit que
$\Gamma(\widehat X,\mathcal O_{\widehat X})$ est isomorphe au séparé
complété $\mathfrak m_y$-adique $\widehat{\mathcal O}_y$ de l'anneau local
$\mathcal O_y$~; c'est donc un anneau local qui ne peut être composé
direct de plusieurs anneaux non réduits à $0$ (sans quoi il aurait
plusieurs idéaux maximaux distincts). On a donc $n=1$, ce qui démontre le
corollaire.
\end{proof}

\begin{corollary}[4.3.3]
\label{III.4.3.3-fr}
Sous les hypothèses de (\hyperref[III.4.3.1-fr]{4.3.1}), pour tout $y\in Y$,
l'ensemble des composantes connexes de la fibre $f^{-1}(y)$ est en
correspondance biunivoque avec l'ensemble fini des points de la fibre
$g^{-1}(y)$, où $g:Y'\to Y$ est le morphisme structural (autrement dit,
l'ensemble des idéaux maximaux de $(f_*(\mathcal O_X))_y$).
\end{corollary}

\begin{proof}
Puisque $Y'$ est fini sur $Y$, on sait en effet que $g^{-1}(y)$ est un
espace fini discret (II, 6.1.7). Comme
$f^{-1}(y)=f'^{-1}(g^{-1}(y))$, le corollaire résulte de cette remarque et
de (\hyperref[III.4.3.1-fr]{4.3.1}).
\end{proof}

On a ainsi une interprétation remarquable du $Y$-préschéma $Y'$ défini
dans (\hyperref[III.4.3.1-fr]{4.3.1}). La factorisation $f=g\circ f'$ du
morphisme propre $f$ est analogue à la factorisation obtenue par K.~Stein
pour les applications holomorphes d'espaces analytiques, et nous
l'appellerons par la suite la \emph{factorisation de Stein} de $f$.

\begin{remark}[4.3.4]
\label{III.4.3.4-fr}
Soit $k$ une extension du corps $k(y)$~: si le préschéma
\[
f^{-1}(y)\otimes_{k(y)}k=X\times_Y\operatorname{Spec}(k)
\]
est connexe, il en est de même de $f^{-1}(y)$, qui en est l'image par un
morphisme de projection (I, 3.4.7). Nous dirons que, pour un morphisme
$f:X\to Y$ de préschémas et un point $y\in Y$, la fibre $f^{-1}(y)$ est
\emph{géométriquement connexe}, si pour toute extension $k$ de $k(y)$, le
préschéma
$f^{-1}(y)\otimes_{k(y)}k=X\times_Y\operatorname{Spec}(k)$ est connexe.

\oldpage[III]{132}

Sous les hypothèses de (\hyperref[III.4.3.2-fr]{4.3.2}), on peut alors en
renforcer la conclusion~: les fibres $f^{-1}(y)$ sont en effet
\emph{géométriquement connexes}. Pour le voir, observons que pour toute
extension $k$ de $k(y)$, il existe un anneau local noethérien $A$ et un
homomorphisme local $\varphi:\mathcal O_y\to A$ qui fait de $A$ un
$\mathcal O_y$-module \emph{plat} et tel que le corps résiduel de $A$ soit
$k(y)$-isomorphe à $k$ (0, 10.3.1). Soit alors
$Y_1=\operatorname{Spec}(A)$ et soit $h:Y_1\to Y$ le morphisme local
correspondant à $\varphi$, transformant l'unique point fermée $y_1$ de $Y_1$
en $y$ (I, 2.4.1)~; posons
$X_1=X\times_YY_1$, et $f_1=f\times1_{Y_1}$~; $f_1$ est propre
(II, 5.4.2, (iii)) et $f_1^{-1}(y_1)$ est un $k(y_1)$-préschéma isomorphe à
$X\times_Y\operatorname{Spec}(k)$. Tout revient donc à montrer que
$f_{1*}(\mathcal O_{X_1})=\mathcal O_{Y_1}$ pour pouvoir appliquer
(\hyperref[III.4.3.2-fr]{4.3.2}) à $f_1$. Or $g$ est un morphisme plat,
comme il résulte de (I, 2.4.2) et de
(\hyperref[III.1.4.15-fr]{1.4.15.5})~; on a donc
\[
f_{1*}(\mathcal O_{X_1})
=h^*(f_*(\mathcal O_X))
=h^*(\mathcal O_Y)
=\mathcal O_{Y_1}
\]
en vertu de (\hyperref[III.1.4.15-fr]{1.4.15}) appliqué pour $q=0$.

Dans le cas général (\hyperref[III.4.3.1-fr]{4.3.1}), le même raisonnement
montre que l'on a (avec les notations de
(\hyperref[III.4.3.1-fr]{4.3.1}))
\[
f_{1*}(\mathcal O_{X_1})=h^*(g_*(\mathcal O_{Y'})),
\]
et la factorisation de Stein $f_1=g_1\circ f'_1$ de $f_1$ est telle que
$g_1=g\times1_{Y_1}$ (II, 1.5.2), le $Y_1$-schéma fini correspondant étant
$Y'_1=Y'\times_YY_1$. Tenant compte de la transitivité des fibres
(I, 3.6.4), on voit donc que le nombre des composantes connexes de
$f_1^{-1}(y_1)$ est, en vertu de
(\hyperref[III.4.3.3-fr]{4.3.3}), égal au nombre d'éléments de
\[
g_1^{-1}(y_1)=g^{-1}(y)\otimes_{k(y)}k.
\]
Si on prend pour $k$ une extension algébriquement close de $k(y)$, ce nombre
est indépendant de l'extension algébriquement close considérée et égal au
\emph{nombre géométrique de points} de $g^{-1}(y)$ (I, 6.4.7), ou encore à
la \emph{somme des rangs séparables}
\[
\sum_i[k(y'_i):k(y)]_s
\]
où $y'_i$ parcourt l'ensemble fini $g^{-1}(y)$. On dit encore que ce nombre
est le \emph{nombre géométrique de composantes connexes} de $f^{-1}(y)$. On
notera que les $k(y'_i)$ ne sont autres que les corps résiduels de l'anneau
semi-local $(f_*(\mathcal O_X))_y$.
\end{remark}

\begin{proposition}[4.3.5]
\label{III.4.3.5-fr}
Soient $X$ et $Y$ deux préschémas localement noethériens intègres et
$f:X\to Y$ un morphisme propre dominant. Pour tout $y\in Y$, le nombre de
composantes connexes de $f^{-1}(y)$ est au plus égal au nombre des idéaux
maximaux de la fermeture intégrale $\mathcal O'_y$ de $\mathcal O_y$ dans le
corps des fonctions rationnelles $\mathrm R(X)$.
\end{proposition}

\begin{proof}
En effet, pour tout ouvert $U$ de $Y$,
\[
\Gamma(U,f_*(\mathcal O_X))=\Gamma(f^{-1}(U),\mathcal O_X)
\]
est l'intersection des anneaux locaux $\mathcal O_x$ tels que
$x\in f^{-1}(U)$ (I, 8.2.1.1). On en conclut aussitôt que la fibre
$(f_*(\mathcal O_X))_y$ est un sous-anneau de $\mathrm R(X)$ contenant
$\mathcal O_y$. En outre, comme $f_*(\mathcal O_X)$ est un
$\mathcal O_Y$-Module cohérent, $(f_*(\mathcal O_X))_y$ est un
$\mathcal O_y$-module de type fini, donc contenu dans $\mathcal O'_y$~; on
sait ([13], vol. I, p. 257 et 259) que tout idéal maximal d'un tel anneau
$A$ est l'intersection de $A$ et d'un idéal maximal de $\mathcal O'_y$,
d'où la proposition.
\end{proof}

\begin{definition}[4.3.6]
\label{III.4.3.6-fr}
On dit qu'un anneau local intègre est \emph{unibranche} si sa clôture
intégrale est un anneau local. On dit qu'un point $y$ d'un préschéma intègre
$Y$ est unibranche si l'anneau local $\mathcal O_y$ est unibranche (ce qui
est en particulier le cas lorsque $Y$ est normal au point $y$).
\end{definition}

Soit $A$ un anneau local intègre, et soit $K$ son corps des fractions~; pour
que $A$ soit unibranche, il faut et il suffit que tout sous-anneau $A_1$ de
$K$, contenant $A$ et qui est une $A$-algèbre finie, soit un anneau local.
En effet, soit $A'$ la clôture intégrale de $A$~; il résulte du premier
théorème de Cohen-Seidenberg (Bourbaki, \emph{Alg. comm.}, chap. V, \S~2,
n\textsuperscript{o}~1, th.~1) que tout idéal maximal de $A_1$ est la trace
d'un idéal maximal de $A'$, donc si $A'$ est local, il en est de même de
$A_1$. Réciproquement, $A'$ est limite inductive

\oldpage[III]{133}

de la famille filtrante croissante des sous-$A$-algèbres finies $A_\alpha$
de $A'$, et si chacune des $A_\alpha$ est un anneau local, l'idéal maximal
de $A_\alpha$ est la trace sur $A_\alpha$ de celui de $A_\beta$ pour
$A_\alpha\subset A_\beta$ par le même raisonnement que ci-dessus, donc $A'$
est un anneau local (0, 10.3.1.3).

On notera que si le complété d'un anneau local noethérien $A$ est intègre
(ce qu'on exprime en disant que $A$ est \emph{analytiquement intègre}), $A$
est unibranche. Soient en effet $\mathfrak m$ l'idéal maximal de $A$, $K$
son corps des fractions, $K'$ le corps des fractions de $\widehat A$~; on
a donc
\[
K'=K\otimes_A\widehat A.
\]
Soit $A'_F$ une sous-$A$-algèbre finie de $K$. Le sous-anneau $B_F$ de $K'$
engendré par $\widehat A$ et $A'_F$ est isomorphe à
$A'_F\otimes_A\widehat A$~; c'est un $\widehat A$-module de type fini,
complété de $A'_F$ pour la topologie $\mathfrak m$-adique
(0\textsubscript{I}, 7.3.3 et 7.3.6). Comme $A'_F$ est un anneau semi-local
(Bourbaki, \emph{Alg. comm.}, chap. IV, \S~2, n\textsuperscript{o}~5,
cor. 3 de la prop. 9) et que son complété est intègre, $A'_F$ ne peut avoir
qu'un seul idéal maximal $\mathfrak m'_F$ et on a
$\mathfrak m'_F\cap A=\mathfrak m$~; d'où notre assertion.

\begin{corollary}[4.3.7]
\label{III.4.3.7-fr}
Sous les hypothèses de (\hyperref[III.4.3.5-fr]{4.3.5}), supposons que la
fermeture algébrique de $\mathrm R(Y)$ dans $\mathrm R(X)$ soit de degré
séparable $n$ et que $y\in Y$ soit unibranche. Alors la fibre $f^{-1}(y)$ a
au plus $n$ composantes connexes. En particulier si la fermeture algébrique
de $\mathrm R(Y)$ dans $\mathrm R(X)$ est radicielle sur $\mathrm R(X)$,
$f^{-1}(y)$ est connexe.
\end{corollary}

\begin{proof}
En effet, soit $\mathcal O'_y$ la clôture intégrale de $\mathcal O_y$~; la
fermeture intégrale $\mathcal O''_y$ de $\mathcal O_y$ dans $\mathrm R(X)$
est aussi celle de $\mathcal O'_y$~; mais on sait que si $\mathcal O'_y$ est
un anneau local, $\mathcal O''_y$ est un anneau semi-local dont le nombre
d'idéaux maximaux est au plus égal à $n$ ([13], vol. I, p. 289, th. 22).

Ce corollaire est essentiellement la forme sous laquelle Zariski énonce son
«~théorème de connexion~» pour les schémas algébriques.
\end{proof}

\begin{remark}[4.3.8]
\label{III.4.3.8-fr}
Si on ajoute aux hypothèses de (\hyperref[III.4.3.7-fr]{4.3.7}) l'hypothèse
que $Y$ est normal au point $y$, la fibre $f^{-1}(y)$ est
\emph{géométriquement connexe}, puisque (avec les notations de
(\hyperref[III.4.3.4-fr]{4.3.4})) $g^{-1}(y)$ est réduit à un point $y'$ et
que $k(y')$ est radiciel sur $k(y)$.
\end{remark}

\begin{definition}[4.3.9]
\label{III.4.3.9-fr}
Étant donné un préschéma localement noethérien $Y$, on dit qu'un morphisme de
type fini $f:X\to Y$ est \emph{universellement ouvert} si, pour tout
préschéma irréductible localement noethérien $Y'$, et tout morphisme dominant
$g:Y'\to Y$, toute composante irréductible de $X'=X\times_YY'$ domine $Y'$.
\end{definition}

Si $Y$ est irréductible, cela revient à dire que si $\eta$, $\eta'$ sont les
points génériques de $Y$ et $Y'$ respectivement (de sorte que
$g(\eta')=\eta$), et si on pose $f'=f_{(Y')}$, toute composante irréductible
de $X'$ rencontre $f'^{-1}(\eta')$ (0\textsubscript{I}, 2.1.8)~; cela
implique donc que pour tout ouvert $U$ de $Y$, le morphisme
$f^{-1}(U)\to U$, restriction de $f$, est universellement ouvert.

\begin{corollary}[4.3.10]
\label{III.4.3.10-fr}
Soient $X$, $Y$ deux préschémas localement noethériens intègres.
$f:X\to Y$ un morphisme propre dominant et universellement ouvert. Si la
fermeture algébrique de $\mathrm R(Y)$ dans $\mathrm R(X)$ est radicielle
sur $\mathrm R(Y)$, toute fibre $f^{-1}(y)$ ($y\in Y$) est géométriquement
connexe.
\end{corollary}

\begin{proof}
On peut se borner au cas où $Y=\operatorname{Spec}(B)$, $B$ étant un anneau
intègre noethérien. Il résulte alors de (II, 7.1.7) qu'il existe un anneau
local noethérien intégralement clos $A$ qui domine $\mathcal O_y$ et a
$\mathrm R(Y)$ pour corps des fractions. Soit
$Y'=\operatorname{Spec}(A)$, et soit $h:Y'\to Y$ le morphisme correspondant
à l'injection canonique $B\to A$, qui est birationnel (donc dominant)~; en
outre, si $y'$ est l'unique point fermé de $Y'$, on a $h(y')=y$. Soient
$X'=X\times_YY'$, $f'=f\times1_{Y'}$~; désignons par $\eta$, $\eta'$,
$\xi$ les points génériques de $Y$, $Y'$ et $X$ respectivement, de sorte
que $f(\xi)=\eta$ et $h^{-1}(\eta)=\{\eta'\}$~; en outre,
$k(\eta)=k(\eta')=\mathrm R(Y)$, donc $f'^{-1}(\eta')$ est isomorphe à
$f^{-1}(\eta)$ (I, 3.6.4) et en particulier, puisque $\xi$ est le point
générique de $f^{-1}(\eta)$ (0\textsubscript{I}, 2.1.8),
$f'^{-1}(\eta')$ a un seul point générique. Mais par hypothèse, toute
composante irréductible de $X'$ a son point générique dans
$f'^{-1}(\eta')$, donc $X'$ est nécessairement irréductible, son point
générique $\xi'$ est le point générique de $f'^{-1}(\eta')$ et on a
$k(\xi')=k(\xi)$.

\oldpage[III]{134}

Posons $X''=X'_{\mathrm{red}}$, $f''=f'_{\mathrm{red}}$~; $X''$ est donc
intègre et noethérien, $f''$ est propre (II, 5.4.6) et les espaces
sous-jacents aux fibres $f'^{-1}(y')$ et $f''^{-1}(y')$ sont les mêmes~; en
outre, $\mathrm R(X'')=k(\xi')=\mathrm R(X)$, donc $f''$ vérifie les
hypothèses de (\hyperref[III.4.3.8-fr]{4.3.8}), et $f''^{-1}(y')$ est
géométriquement connexe. Soit maintenant $k$ une extension quelconque de
$k(y)$~; il existe une extension $k_1$ de $k(y)$ dont $k(y')$ et $k$ peuvent
être considérés comme des sous-extensions (Bourbaki, \emph{Alg.}, chap. V,
\S~4, prop. 2). Par hypothèse,
$f''^{-1}(y')\times_{Y'}\operatorname{Spec}(k_1)$ est connexe, et il a même
préschéma réduit que
$f'^{-1}(y')\times_{Y'}\operatorname{Spec}(k_1)$ (I, 5.1.8), donc ce dernier
est connexe, et comme il est isomorphe à
$f^{-1}(y)\times_Y\operatorname{Spec}(k_1)$ (I, 3.6.4), on en conclut que ce
dernier est connexe~; \emph{a fortiori}, il en est de même de
$f^{-1}(y)\times_Y\operatorname{Spec}(k)$ d'après la remarque du début de
(\hyperref[III.4.3.4-fr]{4.3.4}), ce qui achève la démonstration.
\end{proof}

\begin{env}[4.3.11]
\label{III.4.3.11-fr}
\emph{Remarques} (4.3.11). ---
\begin{enumerate}
\item[(i)] Le raisonnement précédent est dû en substance à Zariski [20], à
cela près qu'il peut prendre pour $A$ la clôture intégrale de
$\mathcal O_y$, celle-ci étant un anneau noethérien pour les anneaux locaux
de la géométrie algébrique classique. D'autre part, Zariski prouve que si
$Y$ est la variété de Chow d'un espace projectif $\mathbf P_k^r$ sur un
corps $k$, et si $X$ est la partie fermée de $\mathbf P_k^r\times_kY$ qui
définit la correspondance de Chow entre $\mathbf P_k^r$ et $Y$, alors la
projection $X\to Y$ est un morphisme universellement ouvert
(\emph{loc. cit.}, lemme de la p. 82). Il semble bien que ce soit la seule
propriété formelle des «~coordonnées de Chow~» dont on se soit servi dans
certaines applications~; il y a par suite intérêt dans une telle situation,
à substituer le langage~: fibres d'un morphisme propre (éventuellement
supposé universellement ouvert ou soumis à d'autres restrictions analogues
de régularité locale) au langage~: spécialisation de cycles dans l'espace
projectif.

\item[(ii)] Au chap. IV, nous verrons qu'un morphisme universellement ouvert
$f:X\to Y$ peut encore être défini de la façon suivante (qui justifie la
terminologie)~: pour tout morphisme $Z\to Y$, le morphisme
$f_{(Z)}:X_{(Z)}\to Z$ est ouvert. On peut montrer en outre que si $f$
vérifie les hypothèses de (\hyperref[III.4.3.10-fr]{4.3.10}), alors, si
$y,y'$ sont deux points de $Y$ tels que $y$ soit une spécialisation de $y'$,
le nombre géométrique de composantes connexes de $f^{-1}(y)$ est au plus
égal à celui des composantes connexes de $f^{-1}(y')$.
\end{enumerate}
\end{env}

\begin{corollary}[4.3.12]
\label{III.4.3.12-fr}
Sous les hypothèses de (\hyperref[III.4.3.5-fr]{4.3.5}), supposons en outre
$\mathrm R(Y)$ algébriquement fermé dans $\mathrm R(X)$, et soit $y$ un
point normal de $Y$. Alors $f^{-1}(y)$ est géométriquement connexe, et il
existe un voisinage ouvert $U$ de $y$ dans $Y$ tel que
$f_*(\mathcal O_X|f^{-1}(U))$ soit isomorphe à $\mathcal O_Y|U$. Si plus
particulièrement, on suppose $Y$ normal (et $\mathrm R(Y)$ algébriquement
fermé dans $\mathrm R(X)$), $f_*(\mathcal O_X)$ est isomorphe à
$\mathcal O_Y$.
\end{corollary}

\begin{proof}
La première assertion relative à $f^{-1}(y)$ est un cas particulier de
(\hyperref[III.4.3.8-fr]{4.3.8}). On en

\oldpage[III]{135}

déduit que si $X\xrightarrow{f'}Y'\xrightarrow{g}Y$ est la factorisation
de Stein de $f$ (\hyperref[III.4.3.3-fr]{4.3.3}), $g^{-1}(y)$ se réduit à
un seul point $y'$~; en outre, on a
\[
\mathcal O_y\subset\mathcal O_{y'}=(f_*(\mathcal O_X))_y
\subset\mathrm R(X),
\]
et comme $\mathcal O_{y'}$ est fini sur $\mathcal O_y$ (et \emph{a
fortiori} sur $\mathrm R(Y)$), il est contenu dans $\mathrm R(Y)$ en vertu
de l'hypothèse~; comme $y$ est normal, on a nécessairement
$\mathcal O_{y'}=\mathcal O_y$~; on en conclut que $g$ est un isomorphisme
local au point $y'$ (I, 6.5.4), ce qui achève de prouver la première partie
du corollaire. La seconde résulte de la première, car l'hypothèse
supplémentaire entraîne que $g$ est bijective et un isomorphisme local au
voisinage de tout point de $Y'$, donc un isomorphisme.
\end{proof}

Le fait que (\hyperref[III.4.3.7-fr]{4.3.7}) soit établi dans le cadre des
schémas permet des applications telles que la suivante~:

\begin{proposition}[4.3.13]
\label{III.4.3.13-fr}
Soient $A$ un anneau local noethérien unibranche, $\mathfrak a$ un idéal de
définition de $A$, $A_0=A/\mathfrak a$,
$S=\operatorname{gr}_{\mathfrak a}(A)$ l'anneau gradué associé à $A$ pour
la filtration $\mathfrak a$-préadique~; $S$ est une $A_0$-algèbre graduée
engendrée par $S_1$, $S_1$ étant un $A_0$-module de type fini. Alors
$\operatorname{Proj}(S)$ est un $A_0$-schéma connexe.
\end{proposition}

\begin{proof}
Soit $\mathfrak m$ l'idéal maximal de $A$~; $Y=\operatorname{Spec}(A)$ est
un schéma intègre dont le point $y$ correspondant à $\mathfrak m$ est
l'unique point fermé. Par hypothèse, on a
$\mathfrak m^p\subset\mathfrak a\subset\mathfrak m$ pour un entier $p$,
donc $V(\mathfrak a)=\{\mathfrak m\}$. Soit
$S'=\bigoplus_{n\geq0}\mathfrak a^n$, et soit
$X=\operatorname{Proj}(S')$, qui est le $Y$-schéma obtenu en faisant éclater
l'idéal $\mathfrak a$~; $X$ est intègre et le morphisme structural
$f:X\to Y$ est birationnel (II, 8.1.4) et évidemment projectif. Par suite,
(\hyperref[III.4.3.7-fr]{4.3.7}) est applicable et montre que $f^{-1}(y)$
est connexe~; mais l'espace $f^{-1}(y)$ est sous-jacent à
$\operatorname{Proj}(S'\otimes_AA_0)$ (I, 3.6.1 et II, 2.8.10)~; comme
$S'\otimes_AA_0=S$ par définition, la proposition est démontrée.
\end{proof}


\subsection{Le «~main theorem~» de Zariski}
\label{subsection:III.4.4-fr}

\begin{proposition}[4.4.1]
\label{III.4.4.1-fr}
Soient $Y$ un préschéma localement noethérien, $f:X\to Y$ un morphisme
propre. Soit $X'$ l'ensemble des points $x\in X$ qui sont isolés dans leur
fibre $f^{-1}(f(x))$. Alors l'ensemble $X'$ est ouvert dans $X$, et si
$f=gf'$ est la factorisation de Stein de $f$
(\hyperref[III.4.3.3-fr]{4.3.3}), la restriction de $f'$ à $X'$ est un
isomorphisme de $X'$ sur un sous-préschéma induit sur un ouvert $U$ de $Y'$,
et on a $X'=f'^{-1}(U)$.
\end{proposition}

\begin{proof}
Comme $g^{-1}(f(x))$ est fini et discret
(\hyperref[III.4.3.3-fr]{4.3.3} et II, 6.1.7), pour que $x$ soit isolé dans
$f^{-1}(f(x))$, il faut et il suffit qu'il le soit dans
$f'^{-1}(f'(x))$~; on peut donc se borner au cas où $f'=f$, donc
$f_*(\mathcal O_X)=\mathcal O_Y$. Alors, si $x\in X'$,
$f^{-1}(f(x))$, qui est connexe (\hyperref[III.4.3.2-fr]{4.3.2}), est
nécessairement réduit au point $x$. Comme $f$ est fermé, pour tout voisinage
ouvert $V$ de $x$ dans $X$, $f(X-V)$ est fermé dans $Y$ et ne contient pas
$y=f(x)$, puisque $f^{-1}(y)=\{x\}$~; si $U$ est le complémentaire de
$f(X-V)$ dans $Y$, on a $f^{-1}(U)\subset V$, et on en conclut que les
images réciproques par $f$ d'un système fondamental de voisinages ouverts de
$y$ forment un système fondamental de voisinages ouverts de $x$. L'hypothèse
$f_*(\mathcal O_X)=\mathcal O_Y$ et la définition de l'image directe d'un
faisceau (0\textsubscript{I}, 3.4.1 et 4.2.1) entraînent alors que, si
$f=(\psi,\theta)$, l'homomorphisme
$\theta_y^\#:\mathcal O_y\to\mathcal O_x$ est un isomorphisme. On en conclut
qu'il existe un voisinage ouvert $V$ de $x$ et un voisinage ouvert $U$ de
$y$ tels que la restriction de $f$ à $V$ soit un isomorphisme de $V$ sur
$U$ (I, 6.5.4)~; en outre, d'après ce qu'on vient de voir,

\oldpage[III]{136}

on peut supposer que $f^{-1}(U)=V$, d'où on conclut aussitôt, par définition,
que $V\subset X'$, ce qui achève la démonstration.
\end{proof}

La proposition suivante a été démontrée par Chevalley dans le cas des schémas
algébriques~:

\begin{proposition}[4.4.2]
\label{III.4.4.2-fr}
Soient $Y$ un préschéma localement noethérien, $f:X\to Y$ un morphisme. Les
conditions suivantes sont équivalentes~:
\begin{enumerate}
\item[a)] $f$ est fini.
\item[b)] $f$ est affine et propre.
\item[c)] $f$ est propre et, pour tout $y\in Y$, $f^{-1}(y)$ est un ensemble
fini.
\end{enumerate}
\end{proposition}

\begin{proof}
On sait que a) entraîne b) (II, 6.1.2 et 6.1.11). Si $f$ est propre et affine,
il en est de même du morphisme
$f^{-1}(y)\to\operatorname{Spec}(k(y))$ (II, 1.6.2, (iii) et 5.4.2, (iii)),
et le th. de finitude (\hyperref[III.3.2.1-fr]{3.2.1}) appliqué au faisceau
structural de $f^{-1}(y)$, montre que $f^{-1}(y)=\operatorname{Spec}(A)$, où
$A$ est une $k(y)$-algèbre finie~; donc $f^{-1}(y)$ est un ensemble fini
(II, 6.1.7), et on voit que b) entraîne c). Enfin, comme $f^{-1}(y)$ est un
préschéma algébrique sur $k(y)$, l'hypothèse que l'ensemble $f^{-1}(y)$ est
fini entraîne que l'espace $f^{-1}(y)$ est discret (I, 6.4.4). Avec les
notations de (\hyperref[III.4.4.1-fr]{4.4.1}), on a donc $X'=X$, et
$f':X\to Y'$ est un isomorphisme~; comme $g$ est un morphisme fini, on voit
que c) entraîne a).
\end{proof}

\begin{theorem}[4.4.3]
\label{III.4.4.3-fr}
\textnormal{(«~Main theorem~» de Zariski).}
Soient $Y$ un préschéma noethérien, $f:X\to Y$ un morphisme quasi-projectif,
$X'$ l'ensemble des points $x\in X$ qui sont isolés dans leur fibre
$f^{-1}(f(x))$. Alors $X'$ est une partie ouverte de $X$, et le
sous-préschéma induit $X'$ est isomorphe à un préschéma induit sur une partie
ouverte d'un $Y$-préschéma $Y'$ fini sur $Y$.
\end{theorem}

\begin{proof}
L'hypothèse entraîne qu'il existe un $Y$-préschéma projectif $Z$ tel que $X$
soit $Y$-isomorphe à un sous-préschéma induit sur un ouvert de $Z$
(II, 5.3.2 et 5.5.1). On est donc ramené à démontrer le théorème lorsque $f$
est un morphisme projectif, donc propre (II, 5.5.3), et il résulte alors
aussitôt de (\hyperref[III.4.4.1-fr]{4.4.1}).
\end{proof}

\begin{remark}[4.4.4]
\label{III.4.4.4-fr}
Si $X$ est réduit (resp. irréductible et $X'$ non vide), on peut supposer,
dans l'énoncé de (\hyperref[III.4.4.3-fr]{4.4.3}), que $Y'$ est réduit (resp.
irréductible). En effet, on peut toujours remplacer $Y'$ par le
sous-préschéma adhérence $\overline{X'}$ de $X'$ dans $Y'$ (I, 9.5.11 et
II, 6.1.5, (i) et (ii)), et on sait que si $X'$ est réduit, il en est de même
de $\overline{X'}$ (I, 9.5.9, (i))~; par ailleurs, si $X'$ n'est pas vide, il
est irréductible si $X$ l'est, et $\overline{X'}$ est alors aussi
irréductible.
\end{remark}

\begin{corollary}[4.4.5]
\label{III.4.4.5-fr}
Soient $Y$ un schéma localement noethérien, $f:X\to Y$ un morphisme de type
fini, $x$ un point de $X$ isolé dans sa fibre $f^{-1}(f(x))$. Alors il existe
un voisinage ouvert de $x$ dans $X$ qui est isomorphe à une partie ouverte
d'un $Y$-préschéma fini sur $Y$.
\end{corollary}

\begin{proof}
Soient en effet $y=f(x)$, $U$ un voisinage ouvert affine de $y$ dans $Y$,
$V$ un voisinage ouvert affine de $x$ dans $X$, contenu dans $f^{-1}(U)$.
Comme $Y$ est séparé, l'injection $U\to Y$ est affine (II, 1.6.3), et comme
$V$ est affine sur $U$ (\emph{ibid.}), la restriction de $f$ à $V$ est un
morphisme affine $V\to Y$ (II, 1.6.2, (ii))~; \emph{a fortiori}, cette
restriction est un morphisme quasi-projectif puisqu'il est de type fini
(I, 6.3.5 et II, 5.3.4, (i)). Il suffit alors d'appliquer à cette restriction
le th. (\hyperref[III.4.4.3-fr]{4.4.3}).
\end{proof}

Le cor. (\hyperref[III.4.4.5-fr]{4.4.5}) s'énonce dans le langage de
l'algèbre commutative~:

\oldpage[III]{137}

\begin{corollary}[4.4.6]
\label{III.4.4.6-fr}
Soient $A$ un anneau noethérien, $B$ une $A$-algèbre de type fini,
$\mathfrak q$ un idéal premier de $B$, $\mathfrak p$ son image réciproque
dans $A$. On suppose que $\mathfrak q$ soit à la fois maximal et minimal dans
l'ensemble des idéaux premiers de $B$ dont l'image réciproque est
$\mathfrak p$. Alors il existe $g\in B-\mathfrak q$, une $A$-algèbre finie
$A'$ et un élément $f'\in A'$ tels que les $A$-algèbres $B_g$ et $A'_{f'}$
soient isomorphes.
\end{corollary}

\begin{proof}
Il suffit en effet d'appliquer (\hyperref[III.4.4.5-fr]{4.4.5}) à
$Y=\operatorname{Spec}(A)$ et $X=\operatorname{Spec}(B)$, l'hypothèse sur
$\mathfrak q$ signifiant exactement qu'il est isolé dans sa fibre
(I, 1.1.7).
\end{proof}

On en déduit le résultat suivant, moins général en apparence~:

\begin{corollary}[4.4.7]
\label{III.4.4.7-fr}
Soient $A$ un anneau local noethérien, $B$ une $A$-algèbre de type fini,
$\mathfrak n$ un idéal premier de $B$ dont l'image réciproque dans $A$ est
l'idéal maximal $\mathfrak m$. On suppose que $\mathfrak n$ est maximal dans
$B$ et est minimal dans l'ensemble des idéaux premiers de $B$ dont l'image
réciproque est $\mathfrak m$ (ce qui signifie aussi que
$B\mathfrak m$ est primaire pour $\mathfrak n$). Alors il existe une
$A$-algèbre finie $A'$ et un idéal maximal $\mathfrak m'$ de $A'$ (dont
$\mathfrak m$ est l'image réciproque dans $A$) tels que $B_{\mathfrak n}$
soit isomorphe à la $A$-algèbre $A'_{\mathfrak m'}$.
\end{corollary}

Le cas particulier suivant de (\hyperref[III.4.4.7-fr]{4.4.7}) est aussi
parfois appelé «~Main Theorem~»~:

\begin{corollary}[4.4.8]
\label{III.4.4.8-fr}
Sous les conditions de (\hyperref[III.4.4.7-fr]{4.4.7}), supposons en outre
$A$ et $B$ intègres et ayant même corps des fractions $K$. Alors, si $A$ est
intégralement clos, on a $B=A$.
\end{corollary}

\begin{proof}
En effet, la Remarque (\hyperref[III.4.4.4-fr]{4.4.4}) montre que l'on peut
supposer, dans l'application de (\hyperref[III.4.4.7-fr]{4.4.7}) que $A'$ est
intègre et a $K$ pour corps des fractions~; l'hypothèse sur $A$ entraîne alors
$A'=A$, donc $B_{\mathfrak n}=A$~; comme on a
$A\subset B\subset B_{\mathfrak n}$, on en conclut bien $A=B$.
\end{proof}

L'énoncé (\hyperref[III.4.4.8-fr]{4.4.8}) est la forme donnée par Zariski à
son «~Main theorem~» (étendu aux anneaux locaux noethériens intègres
quelconques).

Les corollaires précédents étaient des variantes de nature locale de
(\hyperref[III.4.4.3-fr]{4.4.3}), qui est un résultat global. Voici une autre
conséquence de nature globale~:

\begin{corollary}[4.4.9]
\label{III.4.4.9-fr}
Soient $Y$ un préschéma intègre localement noethérien, $f:X\to Y$ un
morphisme séparé, de type fini et birationnel. Supposons en outre $Y$ normal
et toutes les fibres $f^{-1}(y)$ finies pour $y\in Y$. Alors $f$ est une
immersion ouverte~; si en outre $f$ est fermé (et en particulier si $f$ est
propre), $f$ est un isomorphisme.
\end{corollary}

\begin{proof}
Soit en effet $x\in X$, et posons $y=f(x)$. Comme $f^{-1}(y)$ est un schéma
algébrique sur $k(y)$, l'hypothèse qu'il est fini entraîne qu'il est discret
(I, 6.4.4)~; en outre $\mathcal O_y$ est intégralement clos et
$\mathcal O_x$ et $\mathcal O_y$ ont même corps des fractions (I, 7.1.5).
On peut donc appliquer (\hyperref[III.4.4.8-fr]{4.4.8}), et si
$f=(\psi,\theta)$, l'homomorphisme
$\theta_y^\#:\mathcal O_y\to\mathcal O_x$ est bijectif~; on en conclut
(I, 6.5.4) que $f$ est un isomorphisme local. Mais comme $f$ est séparé et
$X$ intègre, $f$ est une immersion ouverte (I, 8.2.8). La dernière assertion
résulte de ce que $f$ est dominant.
\end{proof}

\begin{proposition}[4.4.10]
\label{III.4.4.10-fr}
Soient $Y$ un préschéma localement noethérien, $f:X\to Y$ un morphisme
localement de type fini. L'ensemble $X'$ des $x\in X$ isolés dans leur fibre
$f^{-1}(f(x))$ est ouvert dans $X$.
\end{proposition}

\begin{proof}
La question étant locale sur $X$ et $Y$, on peut supposer $X$ et $Y$ affines
noethériens et $f$ de type fini~; $f$ est alors un morphisme affine de type
fini, donc quasi-projectif (II, 5.3.4, (i)), et il suffit d'appliquer
(\hyperref[III.4.4.3-fr]{4.4.3}).
\end{proof}

\oldpage[III]{138}

\begin{corollary}[4.4.11]
\label{III.4.4.11-fr}
Soient $Y$ un préschéma localement noethérien, $f:X\to Y$ un morphisme
propre. L'ensemble $U$ des points $y\in Y$ tels que $f^{-1}(y)$ soit discret
est ouvert dans $Y$, et le morphisme $f^{-1}(U)\to U$ restriction de $f$ est
fini. En particulier, un morphisme propre et quasi-fini $X\to Y$ est fini.
\end{corollary}

\begin{proof}
En effet, le complémentaire de $U$ dans $Y$ est l'image par $f$ de $X-X'$ qui
est fermé dans $X$ en vertu de (\hyperref[III.4.4.10-fr]{4.4.10})~; comme $f$
est une application fermée, $U$ est ouvert. En outre, il résulte de
(II, 6.2.2) que $f^{-1}(y)$ est fini pour tout $y\in U$~; comme le morphisme
$f^{-1}(U)\to U$, restriction de $f$ est propre (II, 5.4.1), il est fini en
vertu de (\hyperref[III.4.4.2-fr]{4.4.2}).
\end{proof}

\begin{env}[4.4.12]
\label{III.4.4.12-fr}
\emph{Remarques} (4.4.12). ---
\begin{enumerate}
\item[(i)] Comme on l'a annoncé dans (II, 6.2.7), nous montrerons au chap. V
que si $Y$ est localement noethérien, tout morphisme quasi-fini et séparé
$f:X\to Y$ est quasi-affine, donc quasi-projectif. Il s'ensuivra que, dans le
Main Theorem (\hyperref[III.4.4.3-fr]{4.4.3}) la conclusion reste valable
lorsqu'on suppose seulement $f$ séparé et de type fini. En effet, il résulte
de (\hyperref[III.4.4.10-fr]{4.4.10}) que $X'$ est ouvert dans $X$, et comme
$X$ est localement noethérien, la restriction de $f$ à $X'$ est encore de
type fini (I, 6.3.5), donc quasi-fini par définition de $X'$, et évidemment
séparé~; on peut donc appliquer (\hyperref[III.4.4.3-fr]{4.4.3}) à cette
restriction, d'où la conclusion.

\item[(ii)] Nous donnerons au chap. IV une démonstration plus élémentaire de
(\hyperref[III.4.4.10-fr]{4.4.10}), utilisant la théorie de la dimension.
\end{enumerate}
\end{env}

\subsection{Complétés de modules d'homomorphismes}
\label{subsection:III.4.5-fr}

\begin{proposition}[4.5.1]
\label{III.4.5.1-fr}
Soient $A$ un anneau noethérien, $\mathfrak J$ un idéal de $A$, $X$ un
$A$-préschéma de type fini, $\mathcal F$, $\mathcal G$ deux
$\mathcal O_X$-Modules cohérents dont les supports ont une intersection
propre sur $Y=\operatorname{Spec}(A)$ (II, 5.4.10). Alors, pour tout entier
$n\geq0$, $\operatorname{Ext}^n_{\mathcal O_X}(X;\mathcal F,\mathcal G)$ est un $A$-module
de type fini, et son séparé complété pour la topologie $\mathfrak J$-préadique
s'identifie canoniquement (avec les notations de
(\hyperref[III.4.1.7-fr]{4.1.7})) à
$\operatorname{Ext}^n_{\mathcal O_{\widehat X}}(\widehat X;\widehat{\mathcal F},
\widehat{\mathcal G})$.
\end{proposition}

\begin{proof}
On sait (T, 4.2) qu'il existe une suite spectrale birégulière
$\mathrm E(\mathcal F,\mathcal G)$ dont l'aboutissement est
$\operatorname{Ext}_{\mathcal O_X}(X;\mathcal F,\mathcal G)$ et dont les termes
$\mathrm E_2$ sont donnés par
\[
\mathrm E_2^{pq}=\mathrm H^p\bigl(X,
\mathcal{E}\!xt^q_{\mathcal O_X}(\mathcal F,\mathcal G)\bigr).
\]
On sait que $\mathcal{E}\!xt^q_{\mathcal O_X}(\mathcal F,\mathcal G)$ est un
$\mathcal O_X$-Module cohérent (0, 12.3.3) dont le support est contenu dans
l'intersection de ceux de $\mathcal F$ et $\mathcal G$ (T, 4.2.2), et est
par suite propre sur $Y$ (II, 5.4.10). On conclut de
(\hyperref[III.3.2.4-fr]{3.2.4}) que les $\mathrm E_2^{pq}$ sont des
$A$-modules de type fini, et par suite (0, 11.1.8) il en est de même de tous
les termes $\mathrm E_r^{pq}$ de la suite spectrale et de son aboutissement.
D'autre part, si $i:\widehat X\to X$ est le morphisme canonique,
$\widehat{\mathcal F}$ et $\widehat{\mathcal G}$ s'identifient canoniquement
à $i^*(\mathcal F)$ et $i^*(\mathcal G)$, et $i$ est plat (I, 10.8.8 et
10.8.9). On sait alors (0, 12.3.4) qu'il existe, pour tout $q\geq0$, un
$i$-morphisme canonique
\[
u_q:\mathcal{E}\!xt^q_{\mathcal O_X}(\mathcal F,\mathcal G)
\longrightarrow
\mathcal{E}\!xt^q_{\mathcal O_{\widehat X}}
(\widehat{\mathcal F},\widehat{\mathcal G}),
\]
et que le $\mathcal O_{\widehat X}$-homomorphisme correspondant
\[
u_q^\#:i^*\bigl(\mathcal{E}\!xt^q_{\mathcal O_X}(\mathcal F,\mathcal G)\bigr)
\longrightarrow
\mathcal{E}\!xt^q_{\mathcal O_{\widehat X}}
(\widehat{\mathcal F},\widehat{\mathcal G})
\]
est un isomorphisme (0, 12.3.5)~; autrement dit (I, 10.8.8),
$\mathcal{E}\!xt^q_{\mathcal O_{\widehat X}}
(\widehat{\mathcal F},\widehat{\mathcal G})$ s'identifie canoniquement au
complété $(\mathcal{E}\!xt^q_{\mathcal O_X}(\mathcal F,\mathcal G))^\wedge$ (avec les
notations de (\hyperref[III.4.1.7-fr]{4.1.7})). On conclut alors du th. de
comparaison (\hyperref[III.4.1.10-fr]{4.1.10}) que pour tout $p\geq0$,
\[
\mathrm H^p\bigl(\widehat X,
\mathcal{E}\!xt^q_{\mathcal O_{\widehat X}}
(\widehat{\mathcal F},\widehat{\mathcal G})\bigr)
\]
s'identifie canoniquement au séparé complété de
$\mathrm H^p\bigl(X,\mathcal{E}\!xt^q_{\mathcal O_X}
(\mathcal F,\mathcal G)\bigr)$ pour la topologie $\mathfrak J$-préadique.

\oldpage[III]{139}

Si on désigne par $\mathrm E(\widehat{\mathcal F},
\widehat{\mathcal G})$ la suite spectrale birégulière définie dans (T, 4.2)
relative à $\widehat{\mathcal F}$ et $\widehat{\mathcal G}$, on voit donc que
si $\widehat A$ désigne le séparé complété de $A$ pour la topologie
$\mathfrak J$-préadique, on a, à un isomorphisme canonique près,
\[
\mathrm E_2^{pq}(\widehat{\mathcal F},\widehat{\mathcal G})
=\mathrm E_2^{pq}(\mathcal F,\mathcal G)\otimes_A\widehat A
\qquad (0\textsubscript{I}, 7.3.3).
\]

Cela étant, on sait que la donnée du morphisme plat $i$ définit un
homomorphisme canonique de suites spectrales
\[
\varphi:\mathrm E(\mathcal F,\mathcal G)\longrightarrow
\mathrm E(\widehat{\mathcal F},\widehat{\mathcal G})
=\mathrm E(i^*(\mathcal F),i^*(\mathcal G))
\]
qui, pour les termes $\mathrm E_2$ (resp. l'aboutissement) se réduit à
l'homomorphisme
\[
\varphi_2^{pq}:\mathrm H^p\bigl(X,
\mathcal{E}\!xt^q_{\mathcal O_X}(\mathcal F,\mathcal G)\bigr)
\longrightarrow
\mathrm H^p\bigl(\widehat X,
\mathcal{E}\!xt^q_{\mathcal O_{\widehat X}}
(\widehat{\mathcal F},\widehat{\mathcal G})\bigr)
\]
\[
\text{(resp. }\varphi^n:\operatorname{Ext}^n_{\mathcal O_X}
(X;\mathcal F,\mathcal G)\longrightarrow
\operatorname{Ext}^n_{\mathcal O_{\widehat X}}
(\widehat X;\widehat{\mathcal F},\widehat{\mathcal G})\text{)}
\]
déduit de $u_q$ (resp. $u_0$) par fonctorialité (0, 12.3.4). Par
tensorisation avec $\widehat A$, les $\varphi_r^{pq}$ et $\varphi^n$ donnent
des homomorphismes de $\widehat A$-modules
\[
\psi_r^{pq}:\mathrm E_r^{pq}(\mathcal F,\mathcal G)
\otimes_A\widehat A\longrightarrow
\mathrm E_r^{pq}(\widehat{\mathcal F},\widehat{\mathcal G})
\]
\[
\psi^n:\operatorname{Ext}^n_{\mathcal O_X}(X;\mathcal F,\mathcal G)
\otimes_A\widehat A\longrightarrow
\operatorname{Ext}^n_{\mathcal O_{\widehat X}}
(\widehat X;\widehat{\mathcal F},\widehat{\mathcal G}).
\]
Comme $\widehat A$ est un $A$-module plat (0\textsubscript{I}, 7.3.3), les
$\widehat A$-modules
$\mathrm E_r^{pq}(\mathcal F,\mathcal G)\otimes_A\widehat A$ forment une
suite spectrale birégulière d'aboutissement les
$\operatorname{Ext}^n_{\mathcal O_X}(X;\mathcal F,\mathcal G)\otimes_A\widehat A$, et les
$\psi_r^{pq}$ et $\psi^n$ un morphisme de suites spectrales. Comme les
$\psi_2^{pq}$ sont des isomorphismes, il en est de même des $\psi^n$
(0, 11.1.5).
\end{proof}

\begin{corollary}[4.5.2]
\label{III.4.5.2-fr}
Sous les hypothèses de (\hyperref[III.4.5.1-fr]{4.5.1}), supposons en outre
que $A$ soit un anneau $\mathfrak J$-adique noethérien. Alors, pour tout
entier $n\geq0$, $\operatorname{Ext}^n_{\mathcal O_X}(X;\mathcal F,\mathcal G)$
s'identifie canoniquement à
$\operatorname{Ext}^n_{\mathcal O_{\widehat X}}
(\widehat X;\widehat{\mathcal F},\widehat{\mathcal G})$.
\end{corollary}

\begin{proof}
Il suffit de remarquer que $\operatorname{Ext}^n_{\mathcal O_X}
(X;\mathcal F,\mathcal G)$, étant un $A$-module de type fini, est séparé et
complet pour la topologie $\mathfrak J$-préadique (0\textsubscript{I}, 7.3.6).
\end{proof}

Le cas particulier $n=0$ de (\hyperref[III.4.5.1-fr]{4.5.1}) s'énonce de la
façon suivante~:

\begin{corollary}[4.5.3]
\label{III.4.5.3-fr}
Sous les hypothèses de (\hyperref[III.4.5.1-fr]{4.5.1}), pour tout
homomorphisme $u:\mathcal F\to\mathcal G$, désignons par
$\widehat u:\widehat{\mathcal F}\to\widehat{\mathcal G}$ l'homomorphisme
complété (I, 10.8.4). Alors on a un isomorphisme canonique
\[
\label{III.4.5.3.1-fr}
\left(\operatorname{Hom}_{\mathcal O_X}(\mathcal F,\mathcal G)\right)^\wedge
\xrightarrow{\sim}
\operatorname{Hom}_{\mathcal O_{\widehat X}}
(\widehat{\mathcal F},\widehat{\mathcal G}),
\tag{4.5.3.1}
\]
où le premier membre est le séparé complété pour la topologie
$\mathfrak J$-préadique du $A$-module
$\operatorname{Hom}_{\mathcal O_X}(\mathcal F,\mathcal G)$, cet isomorphisme étant obtenu
par passage aux séparés complétés à partir de l'homomorphisme
$u\mapsto\widehat u$.
\end{corollary}

\subsection{Relations entre morphismes formels et morphismes usuels}

\begin{proposition}[4.6.1]
\label{III.4.6.1-fr}
Soient $Y$ un préschéma localement noethérien, $f:X\to Y$ un morphisme
propre, $\mathcal F$ un $\mathcal O_X$-Module cohérent et $f$-plat, $y$ un
point de $Y$. Supposons que pour un entier $n$, on ait
\[
H^n\bigl(f^{-1}(y),\mathcal F\otimes_{\mathcal O_Y}k(y)\bigr)=0.
\]
Alors il existe un voisinage $U$ de $y$ dans $Y$ tel que
$R^nf_*(\mathcal F)|U=0$, et pour tout entier $p\geq0$, l'homomorphisme
canonique
\[
\bigl(R^{n-1}f_*(\mathcal F)\bigr)_y
\longrightarrow
H^{n-1}\bigl(f^{-1}(y),
\mathcal F\otimes_{\mathcal O_Y}
(\mathcal O_y/\mathfrak m_y^{p+1})\bigr)
\]
de (\hyperref[III.4.2.1.1-fr]{4.2.1.1}) est surjectif.
\end{proposition}

\begin{proof}
Comme $R^nf_*(\mathcal F)$ est un $\mathcal O_Y$-Module cohérent
(\hyperref[III.3.2.1-fr]{3.2.1}), la première assertion de la proposition
sera établie si l'on prouve que l'on a
$(R^nf_*(\mathcal F))_y=0$ (0\textsubscript{I}, 5.2.2)~; en vertu de
(\hyperref[III.4.2.1-fr]{4.2.1}), il suffira de prouver que
\[
H^n\bigl(f^{-1}(y),\mathcal F\otimes_{\mathcal O_Y}
(\mathcal O_y/\mathfrak m_y^{p+1})\bigr)=0
\]
pour tout $p$. C'est vrai par hypothèse pour $p=0$~; nous allons le démontrer
par récurrence sur $p$. Posons
\[
X_p=X\times_Y\operatorname{Spec}
(\mathcal O_y/\mathfrak m_y^{p+1}),
\]
de sorte que $X_{p-1}$ est un sous-préschéma fermé de $X_p$, ayant même
espace sous-jacent (I, 3.6.1)~; l'hypothèse de récurrence
\[
H^n\bigl(X_{p-1},\mathcal F\otimes_{\mathcal O_Y}
(\mathcal O_y/\mathfrak m_y^p)\bigr)=0
\]
entraîne donc
\[
H^n\bigl(X_p,\mathcal F\otimes_{\mathcal O_Y}
(\mathcal O_y/\mathfrak m_y^p)\bigr)=0~;
\]
d'autre part, la suite exacte de cohomologie donne, à partir de la suite
exacte
\[
0\longrightarrow
\mathfrak m_y^p\mathcal F/\mathfrak m_y^{p+1}\mathcal F
\longrightarrow
\mathcal F/\mathfrak m_y^{p+1}\mathcal F
\longrightarrow
\mathcal F/\mathfrak m_y^p\mathcal F
\longrightarrow0
\]
de $\mathcal O_{X_p}$-Modules, la suite exacte
\[
H^n\bigl(X_p,
\mathfrak m_y^p\mathcal F/\mathfrak m_y^{p+1}\mathcal F\bigr)
\longrightarrow
H^n\bigl(X_p,\mathcal F/\mathfrak m_y^{p+1}\mathcal F\bigr)
\longrightarrow
H^n\bigl(X_p,\mathcal F/\mathfrak m_y^p\mathcal F\bigr)
\]
et il suffira de montrer que l'on a
\[
\label{III.4.6.1.1-fr}
H^n\bigl(X_p,
\mathfrak m_y^p\mathcal F/\mathfrak m_y^{p+1}\mathcal F\bigr)=0.
\tag{4.6.1.1}
\]
car alors $H^n(X_p,\mathcal F/\mathfrak m_y^{p+1}\mathcal F)$ sera un
sous-module de $H^n(X_p,\mathcal F/\mathfrak m_y^p\mathcal F)$, donc $0$ en
vertu de l'hypothèse de récurrence.

\oldpage[III]{140}

Notons maintenant que la fibre
$Z=f^{-1}(y)=X\times_Y\operatorname{Spec}(k(y))$ est un sous-préschéma fermé
de $X_p$, et que
$\mathfrak m_y^p\mathcal F/\mathfrak m_y^{p+1}\mathcal F$ est annulé par
$\mathfrak m_y$, donc peut être considéré comme un $\mathcal O_Z$-Module, de
sorte que
\[
H^n\bigl(Z,\mathfrak m_y^p\mathcal F/
\mathfrak m_y^{p+1}\mathcal F\bigr)
=H^n\bigl(X_p,\mathfrak m_y^p\mathcal F/
\mathfrak m_y^{p+1}\mathcal F\bigr).
\]
Cela étant, nous allons montrer que le $\mathcal O_Z$-homomorphisme canonique
\[
\label{III.4.6.1.2-fr}
(\mathcal F/\mathfrak m_y\mathcal F)
\otimes_{k(y)}(\mathfrak m_y^p/\mathfrak m_y^{p+1})
\longrightarrow
\mathfrak m_y^p\mathcal F/\mathfrak m_y^{p+1}\mathcal F
\tag{4.6.1.2}
\]
est bijectif~; cela établi, il en résultera, puisque
$\mathfrak m_y^p/\mathfrak m_y^{p+1}$ est un $k(y)$-module libre, que l'on a
\[
\begin{split}
H^n\bigl(Z,\mathfrak m_y^p\mathcal F/
\mathfrak m_y^{p+1}\mathcal F\bigr)
&=H^n\bigl(Z,\mathcal F/\mathfrak m_y\mathcal F\bigr)
\otimes_{k(y)}(\mathfrak m_y^p/\mathfrak m_y^{p+1})\\
&=0
\end{split}
\]
(0, 12.2.3), puisque
$H^n(Z,\mathcal F/\mathfrak m_y\mathcal F)=0$ par hypothèse, d'où
(\hyperref[III.4.6.1.1-fr]{4.6.1.1}). Pour établir la première assertion, il
reste donc à prouver que (\hyperref[III.4.6.1.2-fr]{4.6.1.2}) est bijectif~;
comme la question est ponctuelle sur $X$ et que $\mathcal F_x$ est un
$\mathcal O_y$-module plat par hypothèse pour tout $x\in f^{-1}(y)$, il
suffit d'appliquer (0, 10.2.1, c)), car
$\mathcal F_x/\mathfrak m_y\mathcal F_x$ est un module plat sur le corps
$k(y)=\mathcal O_y/\mathfrak m_y$.

Pour démontrer la seconde assertion de (\hyperref[III.4.6.1-fr]{4.6.1}), on
se ramène aussitôt, comme dans (\hyperref[III.4.2.1-fr]{4.2.1}), au cas où
$Y$ est affine et $y$ fermé. Notons que
(\hyperref[III.4.6.1.1-fr]{4.6.1.1}) donne, par un raisonnement analogue,
pour tout $k>0$, la relation
\[
\label{III.4.6.1.3-fr}
H^n\bigl(X_{k+1},
\mathfrak m_y^k\mathcal F/
\mathfrak m_y^{k+p+1}\mathcal F\bigr)=0.
\tag{4.6.1.3}
\]

\oldpage[III]{141}

d'où on déduit, par (\hyperref[III.4.2.1-fr]{4.2.1}), que l'on a aussi
\[
\label{III.4.6.1.4-fr}
\bigl(R^nf_*(\mathfrak m_y^k\mathcal F)\bigr)_y=0.
\tag{4.6.1.4}
\]
Cela étant, on tire de la suite exacte de cohomologie l'exactitude de la
suite
\[
\bigl(R^{n-1}f_*(\mathcal F)\bigr)_y
\longrightarrow
\bigl(R^{n-1}f_*(\mathcal F/\mathfrak m_y^p\mathcal F)\bigr)_y
\longrightarrow
\bigl(R^nf_*(\mathfrak m_y^p\mathcal F)\bigr)_y=0
\]
et comme $y$ est fermé et $Y$ affine, on a (\hyperref[III.1.4.11-fr]{1.4.11})
\[
R^{n-1}f_*(\mathcal F/\mathfrak m_y^p\mathcal F)
=\bigl(H^{n-1}(X,\mathcal F/\mathfrak m_y^p\mathcal F)\bigr)^\sim
=\bigl(H^{n-1}(f^{-1}(y),
\mathcal F/\mathfrak m_y^p\mathcal F)\bigr)^\sim
\]
(G, II, 4.9.1)~; or
$H^{n-1}(f^{-1}(y),\mathcal F/\mathfrak m_y^p\mathcal F)$ est un
$(\mathcal O_y/\mathfrak m_y^p)$-module, d'où
\[
\bigl(R^{n-1}f_*(\mathcal F/\mathfrak m_y^p\mathcal F)\bigr)_y
=H^{n-1}(f^{-1}(y),\mathcal F/\mathfrak m_y^p\mathcal F)
\]
et cela achève de prouver (\hyperref[III.4.6.1-fr]{4.6.1}).
\end{proof}

\begin{corollary}[4.6.2]
\label{III.4.6.2-fr}
Soient $Y$ un préschéma localement noethérien, $f:X\to Y$ un morphisme propre
et plat, $\mathcal F$, $\mathcal G$ deux $\mathcal O_X$-Modules localement
libres, $y$ un point de $Y$. Posons
\[
X_y=f^{-1}(y)=X\otimes_Yk(y),\qquad
\mathcal F_y=\mathcal F\otimes_{\mathcal O_Y}k(y),\qquad
\mathcal G_y=\mathcal G\otimes_{\mathcal O_Y}k(y),
\]
et supposons que l'on ait
\[
\label{III.4.6.2.1-fr}
H^1\bigl(X_y,
\mathcal{H}\!om_{\mathcal O_{X_y}}(\mathcal F_y,\mathcal G_y)\bigr)=0.
\tag{4.6.2.1}
\]
Alors, pour tout homomorphisme $u_0:\mathcal F_y\to\mathcal G_y$, il existe
un voisinage ouvert $U$ de $y$ et un homomorphisme
\[
u:\mathcal F|f^{-1}(U)\longrightarrow\mathcal G|f^{-1}(U)
\]
tels que $u_0$ soit égal à l'homomorphisme $u\otimes1$.
\end{corollary}

\begin{proof}
En effet, l'hypothèse permet d'appliquer (\hyperref[III.4.6.1-fr]{4.6.1}) au
$\mathcal O_X$-Module cohérent
$\mathcal H=\mathcal{H}\!om_{\mathcal O_X}(\mathcal F,\mathcal G)$ pour
$n=1$ et $p=0$, car $\mathcal H$ est localement libre et \emph{a fortiori}
$f$-plat, et le $\mathcal O_{X_y}$-Module
$\mathcal H\otimes_{\mathcal O_Y}k(y)$ s'identifie alors à
$\mathcal{H}\!om_{\mathcal O_{X_y}}(\mathcal F_y,\mathcal G_y)$
(0\textsubscript{I}, 6.2.2). On peut supposer $Y=\operatorname{Spec}(A)$
affine, et alors (\hyperref[III.1.4.11-fr]{1.4.11})
\[
R^0f_*(\mathcal H)=
\bigl(\operatorname{Hom}_{\mathcal O_X}(\mathcal F,\mathcal G)\bigr)^\sim,
\]
donc
\[
\bigl(R^0f_*(\mathcal H)\bigr)_y
=\operatorname{Hom}_{\mathcal O_X}(\mathcal F,\mathcal G)
\otimes_A\mathcal O_y~;
\]
l'homomorphisme canonique
\[
\operatorname{Hom}_{\mathcal O_X}(\mathcal F,\mathcal G)
\otimes_A\mathcal O_y
\longrightarrow
\operatorname{Hom}_{\mathcal O_{X_y}}(\mathcal F_y,\mathcal G_y)
\]
étant surjectif par (\hyperref[III.4.6.1-fr]{4.6.1}), cela établit le
corollaire, tout élément de
$\operatorname{Hom}_{\mathcal O_X}(\mathcal F,\mathcal G)
\otimes_A\mathcal O_y$ pouvant toujours se mettre sous la forme
$u\otimes(1/s)$, où $s\notin\mathfrak m_y$ est un élément de $A$.
\end{proof}

Ce corollaire peut se compléter par le suivant~:

\begin{corollary}[4.6.3]
\label{III.4.6.3-fr}
Sous les hypothèses de (\hyperref[III.4.6.2-fr]{4.6.2}), si $u_0$ est
injectif (resp. surjectif, bijectif), on peut supposer qu'il en est de même de
$u$.
\end{corollary}

\begin{proof}
On peut se borner au cas où $U=Y$. Il suffit de prouver que si $u_0$ est
injectif (resp. surjectif), $\operatorname{Ker}u_x=0$ (resp.
$\operatorname{Coker}u_x=0$) pour tout $x\in f^{-1}(y)$~: en effet,
$\operatorname{Ker}u$ et $\operatorname{Coker}u$ sont des
$\mathcal O_X$-Modules cohérents (0\textsubscript{I}, 5.3.4), donc il
existera un voisinage $V$ de $f^{-1}(y)$ dans $X$ tel que la restriction de
$\operatorname{Ker}u$ (resp. $\operatorname{Coker}u$) à $V$ soit $0$
(0\textsubscript{I}, 5.2.2)~; comme $f$ est fermé, il existera un voisinage
$U'\subset U$ de $y$ tel que $f^{-1}(U')\subset V$, et
(\hyperref[III.4.6.3-fr]{4.6.3}) sera démontré. Par hypothèse,
\[
u_x\otimes1:
\mathcal F_x\otimes_{\mathcal O_y}k(y)
\longrightarrow
\mathcal G_x\otimes_{\mathcal O_y}k(y)
\]
est injectif (resp. surjectif), $\mathcal F_x$ et $\mathcal G_x$ sont des
$\mathcal O_x$-modules libres de type fini et $\mathcal O_x$ est un
$\mathcal O_y$-module plat. Lorsque l'on suppose $u_x\otimes1$ injectif, le
fait que $u_x$ soit injectif résulte de (0, 10.2.4). Lorsque l'on suppose
$u_x\otimes1$ surjectif, \emph{a fortiori} l'homomorphisme
\[
\mathcal F_x/\mathfrak m_x\mathcal F_x
\longrightarrow
\mathcal G_x/\mathfrak m_x\mathcal G_x,
\]
qui s'en déduit par passage aux quotients, est surjectif~; comme
$\mathcal G_x$ est un $\mathcal O_x$-module de type fini et que
$\mathcal O_x$ est un anneau local d'idéal maximal $\mathfrak m_x$, la
conclusion résulte du lemme de Nakayama (Bourbaki, \emph{Alg.}, chap. VIII,
\S\,6, no~3, cor. 4 de la prop. 6).
\end{proof}

On déduit en particulier de (\hyperref[III.4.6.3-fr]{4.6.3})~:

\begin{corollary}[4.6.4]
\label{III.4.6.4-fr}
Soient $Y$ un préschéma localement noethérien, $f:X\to Y$ un morphisme propre
et plat, $y$ un point de $Y$, $X_y=X\otimes_Yk(y)$. Soit $\mathcal E_0$ un
$\mathcal O_{X_y}$-Module localement libre tel que
\[
\label{III.4.6.4.1-fr}
H^1\bigl(X_y,
\mathcal{H}\!om_{\mathcal O_{X_y}}(\mathcal E_0,\mathcal E_0)\bigr)=0.
\tag{4.6.4.1}
\]
Soient $\mathcal F$, $\mathcal G$ deux $\mathcal O_X$-Modules localement
libres tels que $\mathcal F_y$ et $\mathcal G_y$ (avec les notations de
(\hyperref[III.4.6.2-fr]{4.6.2})) soient isomorphes à $\mathcal E_0$. Alors
il existe un voisinage ouvert $U$ de $y$ tel que $\mathcal F|f^{-1}(U)$ et
$\mathcal G|f^{-1}(U)$ soient isomorphes.
\end{corollary}

Plus particulièrement~:

\begin{corollary}[4.6.5]
\label{III.4.6.5-fr}
Sous les hypothèses de (\hyperref[III.4.6.4-fr]{4.6.4}) sur $f$, $X$, $Y$,
supposons que $H^1(X_y,\mathcal O_{X_y})=0$. Si $\mathcal F$ et $\mathcal G$
sont deux $\mathcal O_X$-Modules inversibles tels que $\mathcal F_y$ et
$\mathcal G_y$ soient isomorphes, il existe un voisinage ouvert $U$ de $y$
tel que $\mathcal F|f^{-1}(U)$ et $\mathcal G|f^{-1}(U)$ soient isomorphes.
\end{corollary}

\begin{proof}
Il suffit d'appliquer (\hyperref[III.4.6.4-fr]{4.6.4}) aux modules
$\mathcal F^{-1}\otimes\mathcal G$ et $\mathcal O_X$.
\end{proof}

\begin{remarks}[4.6.6]
\label{III.4.6.6-fr}
\begin{enumerate}
\item[(i)] En utilisant (\hyperref[III.4.6.5-fr]{4.6.5}), nous établirons au
chap. V la classification des faisceaux inversibles sur un fibré projectif,
annoncée dans (II, 4.2.7).
\item[(ii)] Le résultat de (\hyperref[III.4.6.1-fr]{4.6.1}) apparaîtra au
\S\,7 comme conséquence de propositions plus générales.
\end{enumerate}
\end{remarks}

\begin{proposition}[4.6.7]
\label{III.4.6.7-fr}
Soient $Z$ un préschéma localement noethérien, $X$, $Y$ deux $Z$-préschémas
tels que les morphismes structuraux $g:X\to Z$, $h:Y\to Z$ soient propres.
Soient $f:X\to Y$ un $Z$-morphisme, $z$ un point de $Z$, et soit
\[
f_z=f\times_Z1:
X\otimes_Zk(z)\longrightarrow Y\otimes_Zk(z).
\]
\begin{enumerate}
\item[(i)] Si $f_z$ est un morphisme fini (resp. une immersion fermée), il
existe un voisinage ouvert $U$ de $z$ tel que le morphisme
$g^{-1}(U)\to h^{-1}(U)$, restriction de $f$, soit un morphisme fini (resp.
une immersion fermée).
\item[(ii)] On suppose de plus que $g$ soit un morphisme plat. Alors, si
$f_z$ est un isomorphisme, il existe un voisinage ouvert $U$ de $z$ tel que
le morphisme $g^{-1}(U)\to h^{-1}(U)$, restriction de $f$, soit un
isomorphisme.
\end{enumerate}
\end{proposition}

\begin{proof}
Dans les deux cas, il suffira de prouver que pour tout $y\in h^{-1}(z)$ il
existe un voisinage $V_y$ de $y$ tel que la restriction
$f^{-1}(V_y)\to V_y$ de $f$ soit un morphisme fini (resp. une immersion
fermée, un isomorphisme)~; il en résultera alors en effet que si $V$ est la
réunion des $V_y$, la restriction $f^{-1}(V)\to V$ de $f$ est un morphisme
fini (resp. une immersion fermée, un isomorphisme) (II, 6.1.1 et I, 4.2.4).
Comme $h$ est un morphisme fermé, il existera un voisinage ouvert $U$ de $z$
tel que $h^{-1}(U)\subset V$, et la proposition sera démontrée.

\noindent\textit{(i)} Notons tout d'abord que $f$ est un morphisme propre
(II, 5.4.3)~; si on suppose $f_z$ fini, l'existence pour tout
$y\in h^{-1}(z)$ d'un voisinage $V_y$ tel que $f^{-1}(V_y)\to V_y$ soit fini
résulte de (\hyperref[III.4.4.11-fr]{4.4.11}).

\oldpage[III]{143}

Pour traiter le cas où $f_z$ est une immersion fermée, on peut donc déjà
supposer que le morphisme $f$ est fini, donc que
$X=\operatorname{Spec}(\mathcal B)$, où $\mathcal B$ est une
$\mathcal O_Y$-Algèbre cohérente, le morphisme $f$ correspondant (II, 1.2.7)
à l'homomorphisme canonique $u:\mathcal O_Y\to\mathcal B$. Si l'on prouve
que pour tout $y\in h^{-1}(z)$, l'homomorphisme
$u_y:\mathcal O_y\to\mathcal B_y$ est surjectif, il en résultera que pour un
voisinage $V_y$ de $y$, $u|V_y$ sera surjectif, le faisceau
$\operatorname{Coker}u$ étant cohérent (0\textsubscript{I}, 5.3.4 et 5.2.2).
Cela étant, le morphisme fini $f_z$ correspond à l'homomorphisme
\[
v=u\otimes1:
\mathcal O_Y\otimes_{\mathcal O_Z}k(z)
\longrightarrow
\mathcal B\otimes_{\mathcal O_Z}k(z),
\]
et l'hypothèse que $f_z$ est une immersion fermée entraîne que
l'homomorphisme
\[
u_y\otimes1:
\mathcal O_y\otimes_{\mathcal O_z}k(z)
\longrightarrow
\mathcal B_y\otimes_{\mathcal O_z}k(z)
\]
est surjectif. Comme $\mathcal B_y$ est un $\mathcal O_y$-module de type
fini et $\mathcal O_y$ un anneau local noethérien, la conclusion résulte
comme dans (\hyperref[III.4.6.3-fr]{4.6.3}) du lemme de Nakayama.

\noindent\textit{(ii)} Le même raisonnement que ci-dessus montre qu'il
suffit cette fois de prouver que $u_y$ est bijectif, sachant que
$u_y\otimes1$ est bijectif.

Cela résultera du lemme suivant~:
\end{proof}

\begin{lemma}[4.6.7.1]
\label{III.4.6.7.1-fr}
Soient $A$, $B$ deux anneaux locaux noethériens,
$\rho:A\to B$ un homomorphisme local, $u:N\to M$ un homomorphisme de
$B$-modules. On suppose que $M$ est un $A$-module plat, $N$ un $B$-module
de type fini et que
\[
u\otimes1:N\otimes_Ak\longrightarrow M\otimes_Ak
\]
(où $k$ est le corps résiduel de $A$) est injectif. Alors $N$ est un
$A$-module plat et $u$ est injectif.
\end{lemma}

\begin{proof}
Pour établir la première assertion, il faut montrer que pour tout couple de
$A$-modules de type fini $P$, $Q$ et tout $A$-homomorphisme injectif
$v:P\to Q$,
$1_N\otimes v:N\otimes_AP\to N\otimes_AQ$ est injectif. Or, on a le
diagramme commutatif
\[
\xymatrix{
N\otimes_AP\ar[r]^{1_N\otimes v}\ar[d]_{u\otimes1_P}
  &N\otimes_AQ\ar[d]^{u\otimes1_Q}\\
M\otimes_AP\ar[r]_{1_M\otimes v}
  &M\otimes_AQ
}
\]
et comme $1_M\otimes v$ est injectif par hypothèse, il suffira de prouver
qu'il en est de même de $u\otimes1_P$. Soit $\mathfrak m$ l'idéal maximal
de $A$~; la filtration $\mathfrak m$-adique sur le $A$-module
$N\otimes_AP$ est aussi sa filtration $\mathfrak mB$-adique en tant que
$B$-module~; la topologie définie par cette filtration est donc séparée,
puisque $B$ est noethérien, que $\mathfrak mB$ est contenu dans le radical de
$B$, et que $N\otimes_AP$ est un $B$-module de type fini, $N$ étant un
$B$-module de type fini et $P$ un $A$-module de type fini
(0\textsubscript{I}, 7.3.5). Il suffit donc de prouver que l'homomorphisme
\[
\operatorname{gr}(u\otimes1_P):
\operatorname{gr}_\bullet(N\otimes_AP)
\longrightarrow
\operatorname{gr}_\bullet(M\otimes_AP)
\]
(où les modules gradués sont relatifs aux filtrations $\mathfrak m$-adiques)
est injectif (Bourbaki, \emph{Alg. comm.}, chap. III, \S\,2, no~8, cor. 1 du
th. 1). Notons maintenant que puisque $M$ est un $A$-module plat, les
homomorphismes
\[
M\otimes_A(\mathfrak m^nP)
\longrightarrow
\mathfrak m^n(M\otimes_AP)
\]
sont bijectifs~; il en est donc de même de l'homomorphisme canonique
\[
\varphi_M:
\operatorname{gr}_0(M)\otimes_A\operatorname{gr}_\bullet(P)
\longrightarrow
\operatorname{gr}_\bullet(M\otimes_AP).
\]

\oldpage[III]{144}

Or, on a un diagramme commutatif
\[
\xymatrix{
\operatorname{gr}_0(N)\otimes_A\operatorname{gr}_\bullet(P)
  \ar[r]^{\operatorname{gr}_0(u)\otimes1}\ar[d]_{\varphi_N}
  &\operatorname{gr}_0(M)\otimes_A\operatorname{gr}_\bullet(P)
  \ar[d]^{\varphi_M}\\
\operatorname{gr}_\bullet(N\otimes_AP)
  \ar[r]_{\operatorname{gr}(u\otimes1)}
  &\operatorname{gr}_\bullet(M\otimes_AP)
}
\]
dans lequel $\varphi_M$ est bijectif, $\varphi_N$ surjectif~; en outre,
$\operatorname{gr}_0(u)$ est injectif par hypothèse, et comme
\[
\begin{split}
\operatorname{gr}_0(N)\otimes_A\operatorname{gr}_\bullet(P)
&=\operatorname{gr}_0(N)\otimes_k\operatorname{gr}_\bullet(P),\\
\operatorname{gr}_0(M)\otimes_A\operatorname{gr}_\bullet(P)
&=\operatorname{gr}_0(M)\otimes_k\operatorname{gr}_\bullet(P),
\end{split}
\]
$\operatorname{gr}_0(u)\otimes1$ est aussi injectif. On en conclut que
$\operatorname{gr}(u\otimes1)$ est injectif, ce qui achève de démontrer la
première assertion. La seconde se déduit du raisonnement précédent en faisant
$P=A$.
\end{proof}

\begin{proposition}[4.6.8]
\label{III.4.6.8-fr}
Soient $Z$ un préschéma localement noethérien, $X$, $Y$, deux $Z$-préschémas
tels que les morphismes structuraux $g:X\to Z$, $h:Y\to Z$ soient propres,
$Z'$ une partie fermée de $Z$, $X'=g^{-1}(Z')$, $Y'=h^{-1}(Z')$ ses images
réciproques,
\[
\widehat X=X_{/X'},\qquad
\widehat Y=Y_{/Y'},\qquad
\widehat Z=Z_{/Z'}
\]
les complétés formels de $X$, $Y$, $Z$ le long de ces parties fermées,
$f:X\to Y$ un $Z$-morphisme, $\widehat f:\widehat X\to\widehat Y$ son
prolongement aux complétés. Pour que $\widehat f$ soit un isomorphisme (resp.
une immersion fermée), il faut et il suffit qu'il existe un voisinage ouvert
$U$ de $Z'$ tel que le morphisme $g^{-1}(U)\to h^{-1}(U)$, restriction de
$f$, soit un isomorphisme (resp. une immersion fermée).
\end{proposition}

\begin{proof}
La suffisance de la condition est immédiate (I, 10.14.7). Pour en montrer la
nécessité, il suffit encore de prouver que pour tout $y\in Y'$, il existe un
voisinage ouvert $V_y$ de $y$ tel que la restriction
$f^{-1}(V_y)\to V_y$ de $f$ soit un isomorphisme (resp. une immersion
fermée), par le même raisonnement que dans
(\hyperref[III.4.6.7-fr]{4.6.7}). On est ainsi ramené au cas où $Y=Z$,
$Y=\operatorname{Spec}(A)$ étant affine noethérien. Par hypothèse
(I, 10.9.1 et 10.14.2) la fibre $f^{-1}(y)$ est réduite à un point pour
$y\in Y'$, donc comme $f$ est propre (II, 5.4.3), il existe un voisinage
ouvert $U$ de $y$ tel que la restriction $f^{-1}(U)\to U$ de $f$ soit un
morphisme fini (\hyperref[III.4.4.11-fr]{4.4.11}). On peut donc déjà supposer
que $f$ soit un morphisme fini, donc $X=\operatorname{Spec}(B)$, où $B$ est
une $A$-algèbre finie sur $A$. Si $Y'=V(\mathfrak J)$, on a alors
\[
\widehat Y=\operatorname{Spf}(\widehat A),\qquad
\widehat X=\operatorname{Spf}(\widehat B),
\]
$\widehat A$ étant le séparé complété de $A$ pour la topologie
$\mathfrak J$-préadique, $\widehat B$ le séparé complété de $B$ pour la
topologie $\mathfrak JB$-préadique, ou (ce qui revient au même), le séparé
complété du $A$-module $B$ pour la topologie $\mathfrak J$-préadique~; en
outre, $\widehat f$ est le morphisme de schémas formels affines correspondant
au prolongement continu $\widehat\varphi:\widehat A\to\widehat B$ de
l'homomorphisme canonique d'anneaux $\varphi:A\to B$, et l'hypothèse est que
$\widehat\varphi$ est surjectif (resp. bijectif) (I, 10.14.2). Or,
$\widehat\varphi$ est aussi le prolongement continu de $\varphi$ considéré
comme homomorphisme de $A$-modules~; on sait alors (I, 10.8.14) qu'il existe
un voisinage ouvert $U$ de $Y'$ tel que la restriction à $U$ de
l'homomorphisme $\widetilde\varphi:\widetilde A\to\widetilde B$ de
$\mathcal O_Y$-Modules soit surjectif (resp. bijectif), ce qui achève la
démonstration.
\end{proof}

\oldpage[III]{145}

\subsection{Un critère d'amplitude}

\begin{theorem}[4.7.1]
\label{III.4.7.1-fr}
Soient $Y$ un préschéma localement noethérien, $f:X\to Y$ un morphisme
propre, $\mathcal L$ un $\mathcal O_X$-Module inversible, $y$ un point de
$Y$, $X_y=X\otimes_Yk(y)=f^{-1}(y)$, $g$ la projection de $X_y$ dans $X$.
Si $\mathcal L_y=g^*(\mathcal L)=\mathcal L\otimes_{\mathcal O_Y}k(y)$ est
ample sur $X_y$, il existe un voisinage ouvert $U$ de $y$ dans $Y$ tel que
$\mathcal L|f^{-1}(U)$ soit ample pour la restriction de $f$ à $f^{-1}(U)$.
\end{theorem}

\begin{proof}
\noindent\textit{I)} Posons
\[
Y'=\operatorname{Spec}(\mathcal O_y),\qquad
X'=X\times_YY',\qquad
\mathcal L'=\mathcal L\otimes_{\mathcal O_Y}\mathcal O_y~;
\]
nous allons d'abord prouver que $\mathcal L'$ est ample pour
$f'=f_{(Y')}$. On a le diagramme commutatif
\[
\xymatrix{
X\ar[d]_f & X'\ar[l]\ar[d]_{f'} & X_y\ar[l]\ar[d]_{f_y}\\
Y & Y'\ar[l] & \operatorname{Spec}(k(y))\ar[l]
}
\]
Comme $f'$ est propre (II, 5.4.2, (iii)) et $\mathcal O_y$ noethérien, on
voit qu'on peut se borner au cas où
$Y=Y'=\operatorname{Spec}(\mathcal O_y)$, donc $X=X'$, supposer que
$\mathcal L_y$ est ample pour $f_y$ et prouver que $\mathcal L$ est ample
pour $f$ (II, 4.6.6). Nous allons appliquer le critère
(\hyperref[III.2.6.1-fr]{2.6.1, c)}) et montrer en fait que pour tout
$\mathcal O_X$-Module cohérent $\mathcal F$, il existe un entier $N$ tel que
\[
H^1(X,\mathcal F(n))=0\quad\text{pour tout }n\geq N,
\qquad
\mathcal F(n)=\mathcal F\otimes\mathcal L^{\otimes n}.
\]
Notons que $y$ est un point fermé de $Y$ correspondant à l'idéal maximal
$\mathfrak m$ de $\mathcal O_y$~; $X_y$ est donc un sous-préschéma fermé de
$X$ défini par l'Idéal cohérent
$\mathcal J=f^*(\widetilde{\mathfrak m})\mathcal O_X=
\mathfrak m\mathcal O_X$ de $\mathcal O_X$ (I, 4.4.5), et $g$ l'injection
canonique. Considérons alors la $k(y)$-algèbre graduée
\[
S=\operatorname{gr}(\mathcal O_y)
=\bigoplus_{j\geq0}\mathfrak m^j/\mathfrak m^{j+1},
\]
qui est de type fini puisque $\mathcal O_y$ est noethérien~; la
$\mathcal O_X$-Algèbre $\mathcal S=f^*(\widetilde S)$ est donc
quasi-cohérente et de type fini, et elle est évidemment annulée par
$\mathcal J$, donc si on pose $\mathcal S_y=g^*(\mathcal S)$,
$\mathcal S_y$ est une $\mathcal O_{X_y}$-Algèbre quasi-cohérente de type
fini, et $\mathcal S=g_*(\mathcal S_y)$. Posons d'autre part
\[
\mathcal M_j=\mathfrak m^j\mathcal F/\mathfrak m^{j+1}\mathcal F,
\qquad
\mathcal M=\bigoplus_{j\geq0}\mathcal M_j
=\operatorname{gr}(\mathcal F)~;
\]
comme $\mathcal F$ est cohérent, $\mathcal M$ est un $\mathcal S$-Module
quasi-cohérent de type fini (0, 10.1.1) qui est aussi annulé par
$\mathcal J$, de sorte que si l'on pose
\[
\mathcal M'_j=g^*(\mathcal M_j),\qquad
\mathcal M'=g^*(\mathcal M)=\bigoplus_{j\geq0}\mathcal M'_j,
\]
$\mathcal M'$ est un $\mathcal S_y$-Module gradué quasi-cohérent de type
fini tel que $\mathcal M=g_*(\mathcal M')$. En outre, si on pose
$\mathcal M'_j(n)=\mathcal M'_j\otimes\mathcal L_y^{\otimes n}$, on a
$\mathcal M'_j(n)=g^*(\mathcal M_j(n))$. Cela étant, $f_y$ est propre
(II, 5.4.2, (iii)) et $\mathcal L_y$ est ample, donc $f_y$ est projectif
(II, 5.5.4 et 4.6.11), et on peut appliquer à
$\operatorname{Spec}(k(y))$, $f_y$, $\mathcal S_y$, $\mathcal L_y$ et
$\mathcal M'$ le théorème (\hyperref[III.2.4.1-fr]{2.4.1, (ii)})~: il
existe un entier $N$ tel que pour $n\geq N$, on ait
$H^q(X_y,\mathcal M'_j(n))=0$ pour tout $q>0$ et tout $j$~; par suite, on a
aussi $H^q(X,\mathcal M_j(n))=0$ pour tout $q>0$ et tout $j$
(G, II, 4.9.1). Posons alors
\[
\mathcal F(n)_j=\mathcal F(n)/\mathfrak m^{j+1}\mathcal F(n),
\]
de sorte que $\mathcal F(n)_{j-1}=\mathcal F(n)_j/\mathcal M_j(n)$ pour
$j\geq1$ et $\mathcal F(n)_0=\mathcal M_0(n)$. On a
$H^1(X,\mathcal F(n)_0)=0$, et, par la suite exacte de cohomologie,
$H^1(X,\mathcal F(n)_j)=H^1(X,\mathcal F(n)_{j-1})$ pour tout $j\geq1$,
donc $H^1(X,\mathcal F(n)_j)=0$ pour tout $j\geq0$. On conclut donc de
(\hyperref[III.4.2.1-fr]{4.2.1}) que $H^1(X,\mathcal F(n))=0$, ce qui
achève de prouver notre assertion.

\oldpage[III]{146}

\noindent\textit{II)} Revenons aux notations du début de la démonstration,
et remarquons qu'on peut toujours supposer $Y=\operatorname{Spec}(A)$
affine~; comme $f'$ est de type fini et $\mathcal L'$ ample pour $f'$, il
existe un entier $m>0$ tel que $\mathcal L'^{\otimes m}$ soit très ample pour
$f'$ (II, 4.6.11)~; remplaçant au besoin $\mathcal L$ par
$\mathcal L^{\otimes m}$, on peut se borner à considérer le cas où
$\mathcal L'$ est très ample pour $f'$, et à prouver que
$\mathcal L|f^{-1}(U)$ est alors très ample pour $f$. Comme $f'$ est propre,
il existe alors une $Y'$-immersion fermée
$j:X'\to P=\mathbf P_{Y'}^r$ pour un entier $r>0$ convenable, telle que
$\mathcal L'$ soit isomorphe à $j^*(\mathcal O_P(1))$ (II, 5.5.4, (ii))~;
cette immersion correspond canoniquement à un $\mathcal O_{X'}$-homomorphisme
surjectif
\[
u:\mathcal O_{X'}^{r+1}\longrightarrow\mathcal L'
\]
(II, 4.2.3). Ce dernier correspond (0\textsubscript{I}, 5.1.1) à la donnée
de $r+1$ sections $s'_i$ ($0\leq i\leq r$) de $\mathcal L'$ au-dessus de
$X'$ qui engendrent ce $\mathcal O_{X'}$-Module. Ces sections sont aussi par
définition des sections de $f'_*(\mathcal L')$ au-dessus de $Y'$~; on a
\[
f'_*(\mathcal L')=f_*(\mathcal L)\otimes_{\mathcal O_Y}\mathcal O_{Y'}
\]
(0\textsubscript{I}, 5.4.10), $Y$ est affine et $\mathcal O_y$ est l'anneau
local en l'idéal premier $\mathfrak j_y$ de $A$, donc on a
$s'_i=s''_i/t_i$, où les $s''_i$ sont des sections de $f_*(\mathcal L)$
au-dessus de $Y$ et les $t_i$ des éléments de $A$ n'appartenant pas à
$\mathfrak j_y$~; on en conclut qu'il existe un voisinage ouvert affine $V$
de $y$ dans $Y$ et des sections $s_i$ de $f_*(\mathcal L)|V$ telles que
$s'_i=s_i/1$ (on rappelle que l'espace $Y'$ est contenu dans $V$, cf. I,
2.4.2). Les $s_i$ sont alors des sections de $\mathcal L$ au-dessus de
$f^{-1}(V)$, définissant donc un homomorphisme
\[
v:\bigl(\mathcal O_X|f^{-1}(V)\bigr)^{r+1}
\longrightarrow\mathcal L|f^{-1}(V)
\]
qui, par hypothèse, est surjectif en tous les points de $f^{-1}(y)$~; comme
$\operatorname{Coker}(v)$ est cohérent (0\textsubscript{I}, 5.3.4), son
support est fermé (0\textsubscript{I}, 5.2.2) et par suite il existe un
voisinage ouvert $W\subset f^{-1}(V)$ de $f^{-1}(y)$ tel que la restriction
de $v$ à $W$ soit un homomorphisme surjectif. Puisque le morphisme $f$ est
fermé, on peut supposer que $W$ est de la forme $f^{-1}(U)$, où $U$ est un
voisinage ouvert de $y$, et la conclusion résulte alors de (II, 4.2.3).
\end{proof}

\subsection{Morphismes finis de préschémas formels}

\begin{proposition}[4.8.1]
\label{III.4.8.1-fr}
Soient $\mathfrak Y$ un préschéma formel localement noethérien, $\mathcal K$ un
Idéal de définition de $\mathfrak Y$, $f:\mathfrak X\to\mathfrak Y$ un
morphisme de préschémas formels. Les conditions suivantes sont équivalentes~:
\begin{enumerate}
\item[a)] $\mathfrak X$ est localement noethérien, $f$ est un morphisme adique
  (I, 10.12.1) et si l'on pose
  $\mathcal J=f^*(\mathcal K)\mathcal O_{\mathfrak X}$, le morphisme
  \[
  f_0:(\mathfrak X,\mathcal O_{\mathfrak X}/\mathcal J)
  \longrightarrow
  (\mathfrak Y,\mathcal O_{\mathfrak Y}/\mathcal K)
  \]
  déduit de $f$ est fini.
\item[b)] $\mathfrak X$ est localement noethérien et est limite inductive d'un
  $(Y_n)$-système inductif adique $(X_n)$ tel que le morphisme
  $X_0\to Y_0$ soit fini.
\item[c)] Tout point de $\mathfrak Y$ possède un voisinage ouvert formel affine
  noethérien $V$ tel que $f^{-1}(V)$ soit un ouvert formel affine et que
  $\Gamma(f^{-1}(V),\mathcal O_{\mathfrak X})$ soit un
  $\Gamma(V,\mathcal O_{\mathfrak Y})$-module de type fini.
\end{enumerate}
\end{proposition}

\begin{proof}
Il est immédiat que a) entraîne b) en vertu de (I, 10.12.3). Pour voir que b)
entraîne c), on peut supposer que
$\mathfrak Y=\operatorname{Spf}(B)$, où $B$ est adique noethérien et
$\mathcal K=\mathfrak R^\Delta$, où $\mathfrak R$ est un idéal de définition
de $B$. Par hypothèse, $X_0$ est un schéma affine dont l'anneau $A_0$ est un
$B/\mathfrak R$-module de type fini (II, 6.1.3). En vertu de (I, 5.1.9), chacun
des $X_n$ est un schéma affine, et si $A_n$ est son anneau, l'hypothèse b)
entraîne que pour $m\leq n$, $A_m$ est isomorphe à
$A_n/\mathfrak R^{m+1}A_n$. On en déduit que $\mathfrak X$ est isomorphe à
$\operatorname{Spf}(A)$, où $A=\varprojlim A_n$~; on conclut en vertu de
(0\textsubscript{I}, 7.2.9). Enfin, pour prouver que c)

\oldpage[III]{147}

entraîne a), on peut se borner encore au cas où
$\mathfrak Y=\operatorname{Spf}(B)$, $\mathfrak X=\operatorname{Spf}(A)$,
$A$ étant une $B$-algèbre finie~; comme $A/\mathfrak RA$ est alors une
$B/\mathfrak R$-algèbre finie, il résulte de (I, 10.10.9) que les conditions
de a) sont satisfaites.
\end{proof}

\begin{definition}[4.8.2]
\label{III.4.8.2-fr}
Lorsque les propriétés équivalentes a), b), c) de
(\hyperref[III.4.8.1-fr]{4.8.1}) sont vérifiées, on dit que le morphisme $f$
est \emph{fini}, ou que $\mathfrak X$ est un \emph{$\mathfrak Y$-préschéma
formel fini}, ou un préschéma formel \emph{fini au-dessus de $\mathfrak Y$}.
\end{definition}

\begin{proposition}[4.8.3]
\label{III.4.8.3-fr}
\begin{enumerate}
\item[(i)] Une immersion fermée de préschémas formels localement noethériens
  est un morphisme fini.
\item[(ii)] Le composé de deux morphismes finis de préschémas formels
  localement noethériens est un morphisme fini.
\item[(iii)] Soient $\mathfrak X$, $\mathfrak Y$, $\mathfrak S$ trois
  préschémas formels localement noethériens, $f:\mathfrak X\to\mathfrak S$ un
  morphisme fini, $g:\mathfrak Y\to\mathfrak S$ un morphisme~; alors le
  morphisme
  $\mathfrak X\times_{\mathfrak S}\mathfrak Y\to\mathfrak Y$ est fini.
\item[(iv)] Soient $\mathfrak S$ un préschéma formel localement noethérien,
  $\mathfrak X'$, $\mathfrak Y'$ deux préschémas formels localement
  noethériens tels que $\mathfrak X'\times_{\mathfrak S}\mathfrak Y'$ soit
  localement noethérien. Si $\mathfrak X$, $\mathfrak Y$ sont des
  $\mathfrak S$-préschémas formels localement noethériens,
  $f:\mathfrak X\to\mathfrak X'$, $g:\mathfrak Y\to\mathfrak Y'$ deux
  $\mathfrak S$-morphismes finis, alors $f\times_{\mathfrak S}g$ est un
  morphisme fini.
\item[(v)] Soient $f:\mathfrak X\to\mathfrak Y$,
  $g:\mathfrak Y\to\mathfrak Z$ deux morphismes de préschémas formels
  localement noethériens tels que $g$ soit de type fini et séparé~; alors, si
  $g\circ f$ est un morphisme fini, $f$ est un morphisme fini.
\end{enumerate}
\end{proposition}

\begin{proof}
(i) est trivial, et les autres assertions se ramènent aussitôt aux
propositions correspondantes pour les morphismes de préschémas usuels
(II, 6.1.5) à l'aide du critère a) de
(\hyperref[III.4.8.1-fr]{4.8.1})~; nous laissons les détails au lecteur, sur le
modèle de (I, 10.13.5).
\end{proof}

\begin{corollary}[4.8.4]
\label{III.4.8.4-fr}
Sous les hypothèses de (I, 10.9.9), si $f$ est un morphisme fini, il en est de
même de son prolongement $\widehat f$ aux complétés.
\end{corollary}

\begin{corollary}[4.8.5]
\label{III.4.8.5-fr}
Si $\mathfrak X$ est un préschéma formel fini au-dessus de $\mathfrak Y$,
$f:\mathfrak X\to\mathfrak Y$ le morphisme structural, alors, pour tout ouvert
$U\subset\mathfrak Y$, $f^{-1}(U)$ est fini au-dessus de $U$.
\end{corollary}

\begin{proposition}[4.8.6]
\label{III.4.8.6-fr}
Si $f:\mathfrak X\to\mathfrak Y$ est un morphisme fini de préschémas formels
localement noethériens, $f_*(\mathcal O_{\mathfrak X})$ est une
$\mathcal O_{\mathfrak Y}$-Algèbre cohérente.
\end{proposition}

\begin{proof}
On peut considérer $f$ comme limite inductive d'un système inductif $(f_n)$ de
morphismes $f_n:X_n\to Y_n$~; montrons que les $f_n$ sont des morphismes finis
et que $f_*(\mathcal O_{\mathfrak X})$ est isomorphe à la limite projective des
$(f_n)_*(\mathcal O_{X_n})$, ce qui établira notre assertion (I, 10.10.5). Il
suffit de se borner au cas où $\mathfrak Y=\operatorname{Spf}(B)$,
$\mathfrak X=\operatorname{Spf}(A)$, et de remarquer que si $\mathfrak R$ est
un idéal de définition de $B$ et $A$ un $B$-module de type fini,
$A/\mathfrak R^{n+1}A$ est un module de type fini sur
$B/\mathfrak R^{n+1}B$, et que $A$ est limite projective des
$A/\mathfrak R^{n+1}A$.
\end{proof}

Réciproquement~:

\begin{proposition}[4.8.7]
\label{III.4.8.7-fr}
Soient $\mathfrak Y$ un préschéma formel localement noethérien, $\mathcal A$
une $\mathcal O_{\mathfrak Y}$-Algèbre cohérente. Il existe un préschéma
formel $\mathfrak X$ fini au-dessus de $\mathfrak Y$, défini à un
$\mathfrak Y$-isomorphisme unique près, et tel que
$f_*(\mathcal O_{\mathfrak X})=\mathcal A$, $f:\mathfrak X\to\mathfrak Y$
étant le morphisme structural.
\end{proposition}

\begin{proof}
Soit $\mathcal K$ un Idéal de définition de $\mathfrak Y$, et posons
\[
Y_n=(\mathfrak Y,\mathcal O_{\mathfrak Y}/\mathcal K^{n+1}),\qquad
\mathcal A_n=\mathcal A/\mathcal K^{n+1}\mathcal A~;
\]
il est clair que $\mathcal A_n$ est une $\mathcal O_{Y_n}$-Algèbre finie et
définit donc un

\oldpage[III]{148}

$Y_n$-préschéma fini $X_n=\operatorname{Spec}(\mathcal A_n)$ (II, 6.1.3)~;
pour $m\leq n$, l'homomorphisme canonique surjectif
$h_{mn}:\mathcal A_n\to\mathcal A_m$ définit un morphisme
$u_{mn}:X_m\to X_n$ tel que le diagramme
\[
\xymatrix{
X_n\ar[d]_{f_n} & X_m\ar[l]_{u_{mn}}\ar[d]^{f_m}\\
Y_n & Y_m\ar[l]
}
\]
($f_n$ étant le morphisme structural) soit commutatif et identifie $X_m$ au
produit $X_n\times_{Y_n}Y_m$, comme on le voit aussitôt (II, 1.4.6). Le
préschéma formel $\mathfrak X$, limite inductive du système inductif $(X_n)$
est alors localement noethérien et tel que le morphisme structural
$f:\mathfrak X\to\mathfrak Y$, limite inductive du système $(f_n)$, soit fini
(\hyperref[III.4.8.1-fr]{4.8.1} et II, 10.12.3.1)~; on a vu en outre dans la
démonstration de (\hyperref[III.4.8.6-fr]{4.8.6}) que
$f_*(\mathcal O_{\mathfrak X})$ est limite projective des $\mathcal A_n$, donc
égale à $\mathcal A$ (I, 10.10.6). Quant à l'assertion d'unicité, elle est
conséquence du résultat plus général suivant~:
\end{proof}

\begin{proposition}[4.8.8]
\label{III.4.8.8-fr}
Soient $\mathfrak Y$ un préschéma formel localement noethérien,
$\mathfrak X$, $\mathfrak X'$ deux $\mathfrak Y$-préschémas formels finis
au-dessus de $\mathfrak Y$, $f:\mathfrak X\to\mathfrak Y$,
$f':\mathfrak X'\to\mathfrak Y$ les morphismes structuraux. Il existe une
bijection canonique de
\[
\operatorname{Hom}_{\mathfrak Y}(\mathfrak X,\mathfrak X')
\quad\hbox{sur}\quad
\operatorname{Hom}_{\mathcal O_{\mathfrak Y}}
\bigl(f'_*(\mathcal O_{\mathfrak X'}),f_*(\mathcal O_{\mathfrak X})\bigr)
\footnotemark.
\]
\footnotetext{La dernière expression désigne l'ensemble des homomorphismes de
$\mathcal O_{\mathfrak Y}$-Algèbres
$f'_*(\mathcal O_{\mathfrak X'})\to f_*(\mathcal O_{\mathfrak X})$.}
\end{proposition}

\begin{proof}
La définition de cette application $h\to\mathcal A(h)$ est la même que dans
(II, 1.1.2), et pour voir qu'elle est bijective, on est aussitôt ramené au cas
où $\mathfrak Y=\operatorname{Spf}(B)$ est un schéma formel affine noethérien.
Mais alors $\mathfrak X=\operatorname{Spf}(A)$,
$\mathfrak X'=\operatorname{Spf}(A')$, où $A$ et $A'$ sont deux $B$-algèbres
finies et $f_*(\mathcal O_{\mathfrak X})=A^\Delta$,
$f'_*(\mathcal O_{\mathfrak X'})=A'^\Delta$. La conclusion résulte alors de
la correspondance biunivoque, d'une part entre les $\mathfrak Y$-morphismes
$\mathfrak X\to\mathfrak X'$ et les $B$-homomorphismes (nécessairement
continus) $A'\to A$ qui sont des homomorphismes d'algèbres (I, 10.2.2), et
d'autre part entre les homomorphismes de $B$-modules $A'\to A$ et les
homomorphismes de $\mathcal O_{\mathfrak Y}$-Modules
$A'^\Delta\to A^\Delta$ (I, 10.10.2.3).
\end{proof}

\begin{corollary}[4.8.9]
\label{III.4.8.9-fr}
Dans la correspondance biunivoque canonique définie dans
(\hyperref[III.4.8.8-fr]{4.8.8}), les immersions fermées
$\mathfrak X\to\mathfrak X'$ correspondent aux homomorphismes surjectifs de
$\mathcal O_{\mathfrak Y}$-Algèbres
$f'_*(\mathcal O_{\mathfrak X'})\to f_*(\mathcal O_{\mathfrak X})$.
\end{corollary}

\begin{proof}
La question étant encore locale sur $\mathfrak Y$, on est ramené à la
définition des immersions fermées de préschémas formels localement noethériens
(I, 10.14.2).
\end{proof}

\begin{corollary}[4.8.10]
\label{III.4.8.10-fr}
Les notations et hypothèses étant celles de
(\hyperref[III.4.8.1-fr]{4.8.1}), pour qu'un morphisme adique $f$ soit une
immersion fermée, il faut et il suffit que $f_0$ soit une immersion fermée (de
préschémas usuels).
\end{corollary}

\begin{proof}
Cela résulte aussitôt de (\hyperref[III.4.8.9-fr]{4.8.9}) et de la condition de
surjectivité pour un homomorphisme de $\mathcal O_{\mathfrak Y}$-Modules
cohérents (I, 10.11.5).
\end{proof}

\begin{proposition}[4.8.11]
\label{III.4.8.11-fr}
Pour qu'un morphisme $f:\mathfrak X\to\mathfrak Y$ de préschémas formels localement

\oldpage[III]{149}

noethériens soit fini, il faut et il suffit qu'il soit propre et ait ses fibres
$f^{-1}(y)$ finies (pour tout $y\in\mathfrak Y$).
\end{proposition}

\begin{proof}
Grâce aux définitions (3.4.1 et 4.8.2), on est aussitôt ramené à la même proposition
pour $f_0$ (notations de (\hyperref[III.4.8.1-fr]{4.8.1})), ce qui n'est autre que
(4.4.2).
\end{proof}

\section{Un théorème d'existence de faisceaux algébriques cohérents}
\label{section:III.5-fr}

\subsection{Énoncé du théorème}
\label{subsection:III.5.1-fr}

\begin{env}[5.1.1]
\label{III.5.1.1-fr}
Soient $A$ un anneau adique noethérien, $\mathfrak J$ un idéal de définition de
$A$, de sorte que $A$ est séparé et complet pour la topologie $\mathfrak J$-adique.
Si $Y=\operatorname{Spec}(A)$, le schéma formel affine $\operatorname{Spf}(A)$
s'identifie au complété $\widehat Y$ de $Y$ le long de la partie fermée
$Y'=\mathrm V(\mathfrak J)$ (I, 10.10.1). Soient $X$ un $Y$-préschéma (usuel) de
type fini, $f:X\to Y$ le morphisme structural~; nous désignerons par $\widehat X$
le complété de $X$ le long de la partie fermée $X'=f^{-1}(Y')$, ou encore le
$\widehat Y$-préschéma formel $X\times_Y\widehat Y$~; par
$\widehat f:\widehat X\to\widehat Y$ le prolongement de $f$ aux complétés~; enfin,
pour tout $\mathcal O_X$-Module cohérent $\mathcal F$, nous noterons
$\widehat{\mathcal F}$ son complété, $\mathcal F_{/X'}$, qui est un
$\mathcal O_{\widehat X}$-Module cohérent.
\end{env}

\begin{proposition}[5.1.2]
\label{III.5.1.2-fr}
Les hypothèses et notations étant celles de
(\hyperref[III.5.1.1-fr]{5.1.1}), soit $\mathcal F$ un $\mathcal O_X$-Module
cohérent dont le support est propre sur $Y$ (II, 5.4.10). Les homomorphismes
canoniques (4.1.4)
\[
\rho_i:\mathrm H^i(X,\mathcal F)\longrightarrow
\mathrm H^i(\widehat X,\widehat{\mathcal F})
\]
sont alors des isomorphismes.
\end{proposition}

\begin{proof}
Comme $\mathrm H^i(X,\mathcal F)$ est un $A$-module de type fini (3.2.4), donc
identique à son séparé complété pour la topologie $\mathfrak J$-préadique
($0_{\mathrm I}$, 7.3.6), la proposition n'est qu'un cas particulier de (4.1.10).
\end{proof}

Rappelons que les isomorphismes canoniques $\rho_i$ commutent aux cobords pour
toute suite exacte
$0\to\mathcal F'\to\mathcal F\to\mathcal F''\to0$ de $\mathcal O_X$-Modules
cohérents (($0$, 12.1.6) et (I, 10.8.9)).

\begin{corollary}[5.1.3]
\label{III.5.1.3-fr}
Soient $\mathcal F$, $\mathcal G$ deux $\mathcal O_X$-Modules cohérents tels que
l'intersection de leurs supports soit propre sur $Y$. Alors l'homomorphisme canonique
\begin{equation}
\operatorname{Hom}_{\mathcal O_X}(\mathcal F,\mathcal G)
\longrightarrow
\operatorname{Hom}_{\mathcal O_{\widehat X}}
(\widehat{\mathcal F},\widehat{\mathcal G})
\tag{5.1.3.1}
\label{III.5.1.3.1.eq-fr}
\end{equation}
qui, à tout homomorphisme $u:\mathcal F\to\mathcal G$, fait correspondre son
complété $\widehat u:\widehat{\mathcal F}\to\widehat{\mathcal G}$, est un
isomorphisme. De plus, lorsque le morphisme $f$ est fermé, pour que $\widehat u$ soit
injectif (resp. surjectif), il faut et il suffit que $u$ le soit.
\end{corollary}

\begin{proof}
La première assertion est un cas particulier de (4.5.3), dû encore au fait que le
premier membre de (\ref{III.5.1.3.1.eq-fr}) est un $A$-module de type fini, donc
identique à son séparé complété. Pour démontrer la seconde, notons en vertu de
(I, 10.8.14) que $\widehat u$ est injectif (resp. surjectif) si et seulement s'il existe
un voisinage de $X'$ dans lequel $u$ soit injectif (resp. surjectif). La conclusion
résulte donc du lemme suivant~:
\end{proof}

\oldpage[III]{150}

\begin{lemma}[5.1.3.1]
\label{III.5.1.3.1-fr}
Sous les hypothèses de (\hyperref[III.5.1.1-fr]{5.1.1}), si on suppose en outre le
morphisme $f$ fermé, tout voisinage de $X'$ dans $X$ est identique à $X$.
\end{lemma}

\begin{proof}
Tout d'abord, on peut se ramener au cas où $f(X)=Y$. En effet, par hypothèse,
$f(X)$ est une partie fermée $Z$ de $Y$~; on peut en outre remplacer $f$ par
$f_{\mathrm{red}}$ (I, 6.3.4), et supposer par suite $X$ et $Y$ réduits~; on peut alors
remplacer $Y$ par le sous-préschéma fermé réduit de $Y$ ayant $Z$ pour espace
sous-jacent (I, 5.2.2), car tout idéal de $A$ est fermé, et tout anneau quotient de
$A$ est donc adique et noethérien. On a alors $f(X')=Y'$~; si $V$ est un voisinage
ouvert de $X'$ dans $X$, $f(X-V)$ est fermé dans $Y$ par hypothèse, et ne rencontre
pas $Y'$~; mais cela est impossible à moins que $X-V$ ne soit vide, puisque
$\mathfrak J$ est contenu dans le radical de $A$ (I, 1.1.15 et $0_{\mathrm I}$,
7.1.10), d'où la conclusion.
\end{proof}

Lorsqu'on se borne aux $\mathcal O_X$-Modules cohérents dont le support est propre
sur $Y$, (\hyperref[III.5.1.3-fr]{5.1.3}) peut s'énoncer, dans le langage des
catégories, en disant que le foncteur
$\mathcal F\mapsto\widehat{\mathcal F}$ est pleinement fidèle de la catégorie des
$\mathcal O_X$-Modules du type précédent, dans la catégorie des
$\mathcal O_{\widehat X}$-Modules cohérents, et établit par suite une équivalence de
la première de ces catégories avec une sous-catégorie pleine de la seconde
($0$, 8.1.6). Le théorème d'existence va prouver que lorsque $X$ est propre sur $Y$,
cette sous-catégorie est en fait la catégorie de tous les
$\mathcal O_{\widehat X}$-Modules cohérents. De façon précise~:

\begin{theorem}[5.1.4]
\label{III.5.1.4-fr}
Soient $A$ un anneau adique noethérien, $Y=\operatorname{Spec}(A)$,
$\mathfrak J$ un idéal de définition de $A$, $Y'=\mathrm V(\mathfrak J)$,
$f:X\to Y$ un morphisme séparé de type fini, $X'=f^{-1}(Y')$. Soient
$\widehat Y=Y_{/Y'}=\operatorname{Spf}(A)$, $\widehat X=X_{/X'}$ les complétés de
$Y$ et $X$ le long de $Y'$ et $X'$, $\widehat f:\widehat X\to\widehat Y$ le
prolongement de $f$ aux complétés~; alors, le foncteur
$\mathcal F\mapsto\mathcal F_{/X'}=\widehat{\mathcal F}$ est une équivalence de la
catégorie des $\mathcal O_X$-Modules cohérents de support propre sur
$\operatorname{Spec}(A)$, avec la catégorie des $\mathcal O_{\widehat X}$-Modules
cohérents de support propre sur $\operatorname{Spf}(A)$.
\end{theorem}

En d'autres termes, compte tenu de (\hyperref[III.5.1.3-fr]{5.1.3})~:

\begin{corollary}[5.1.5]
\label{III.5.1.5-fr}
Pour qu'un $\mathcal O_{\widehat X}$-Module soit isomorphe au complété d'un
$\mathcal O_X$-Module cohérent et de support propre sur $\operatorname{Spec}(A)$,
il faut et il suffit qu'il soit cohérent et de support propre sur
$\operatorname{Spf}(A)$.
\end{corollary}

Le cas le plus important est le

\begin{corollary}[5.1.6]
\label{III.5.1.6-fr}
Supposons $X$ propre sur $Y=\operatorname{Spec}(A)$. Alors le foncteur
$\mathcal F\mapsto\widehat{\mathcal F}$ est une équivalence de la catégorie des
$\mathcal O_X$-Modules cohérents et de la catégorie des
$\mathcal O_{\widehat X}$-Modules cohérents.
\end{corollary}

\begin{env}[5.1.7]
\label{III.5.1.7-fr}
\emph{Scholie.}
Si on tient compte de la caractérisation des faisceaux cohérents sur les préschémas
formels (I, 10.11.3), on voit que sous les conditions de
(\hyperref[III.5.1.1-fr]{5.1.1}), la donnée d'un $\mathcal O_X$-Module cohérent de
support propre sur $\operatorname{Spec}(A)$ équivaut (en posant
$Y_n=\operatorname{Spec}(A/\mathfrak J^{n+1})$ et $X_n=X\times_Y Y_n$) à la donnée
d'un système projectif de $\mathcal O_{X_n}$-Modules cohérents $(\mathcal F_n)$ tel
que pour $m\leq n$ on ait
\[
\mathcal F_m=\mathcal F_n\otimes_{\mathcal O_{Y_n}}\mathcal O_{Y_m}
\quad\text{(ou encore }\mathcal F_m=\mathcal F_n/\mathfrak J^{m+1}\mathcal F_n\text{)}
\]
et que le support de $\mathcal F_0$ soit une partie de $X_0$ propre sur $Y_0$. Au
moyen de (I, 10.11.4), on interprète de même les homomorphismes de
$\mathcal O_X$-Modules cohérents comme des homomorphismes de systèmes projectifs de
$\mathcal O_{X_n}$-Modules cohérents.

\oldpage[III]{151}

Dans tous les cas d'application connus, $A$ est en fait un anneau adique local
noethérien, donc les $Y_n$ sont des spectres d'anneaux artiniens locaux, et les
résultats de ce paragraphe et des précédents réduisent dans une large mesure la
géométrie algébrique sur un anneau adique local noethérien, à la géométrie algébrique
sur des anneaux locaux artiniens.
\end{env}

\begin{corollary}[5.1.8]
\label{III.5.1.8-fr}
Sous les conditions de (\hyperref[III.5.1.4-fr]{5.1.4}), l'application
$Z\mapsto\widehat Z=Z_{/(Z\cap X')}$ est une bijection de l'ensemble des
sous-préschémas fermés $Z$ de $X$, propres sur $Y$, sur l'ensemble des
sous-préschémas formels fermés de $\widehat X$, propres sur $\widehat Y$.
\end{corollary}

\begin{proof}
En effet, un sous-préschéma formel fermé de $\widehat X$ est de la forme
$(T,(\mathcal O_{\widehat X}/\mathcal A)|T)$, où $\mathcal A$ est un Idéal
cohérent de $\mathcal O_{\widehat X}$ (I, 10.14.2)~; si $T$ est propre sur
$\widehat Y$, il résulte de (\hyperref[III.5.1.4-fr]{5.1.4}) que
$\mathcal O_{\widehat X}/\mathcal A$ est isomorphe à un
$\mathcal O_{\widehat X}$-Module de la forme $\widehat{\mathcal F}$, où
$\mathcal F$ est un $\mathcal O_X$-Module cohérent de support propre sur $Y$~; en
outre, il résulte de (\hyperref[III.5.1.3-fr]{5.1.3}) que l'homomorphisme canonique
$\mathcal O_{\widehat X}\to\mathcal O_{\widehat X}/\mathcal A$ est de la forme
$\widehat u$, où $u:\mathcal O_X\to\mathcal F$ est un homomorphisme surjectif de
$\mathcal O_X$-Modules. Donc $\mathcal F$ est de la forme
$\mathcal O_X/\mathcal N$, où $\mathcal N$ est un Idéal cohérent de
$\mathcal O_X$, et $\mathcal A=\widehat{\mathcal N}$ (I, 10.8.8), d'où la
conclusion (I, 10.14.7).
\end{proof}

\subsection{Démonstration du théorème d'existence : cas projectif et quasi-projectif}
\label{subsection:III.5.2-fr}

\begin{env}[5.2.1]
\label{III.5.2.1-fr}
Sous les conditions de (\hyperref[III.5.1.4-fr]{5.1.4}), nous dirons provisoirement
qu'un $\mathcal O_{\widehat X}$-Module cohérent est algébrisable s'il est isomorphe
à un complété $\widehat{\mathcal F}$ d'un $\mathcal O_X$-Module cohérent
$\mathcal F$ de support propre sur $Y$.
\end{env}

\begin{lemma}[5.2.2]
\label{III.5.2.2-fr}
Soient $\mathcal F'$, $\mathcal G'$ deux $\mathcal O_{\widehat X}$-Modules
algébrisables. Pour tout homomorphisme $u:\mathcal F'\to\mathcal G'$,
$\operatorname{Ker}(u)$, $\operatorname{Im}(u)$ et $\operatorname{Coker}(u)$ sont
algébrisables.
\end{lemma}

\begin{proof}
En effet, si $\mathcal F'=\widehat{\mathcal F}$,
$\mathcal G'=\widehat{\mathcal G}$, où $\mathcal F$ et $\mathcal G$ sont des
$\mathcal O_X$-Modules cohérents de supports propres sur $Y$, on a
$u=\widehat v$, où $v:\mathcal F\to\mathcal G$ est un homomorphisme
(\hyperref[III.5.1.3-fr]{5.1.3}). En vertu de l'exactitude du foncteur
$\mathcal F\mapsto\widehat{\mathcal F}$,
$\operatorname{Ker}(\widehat v)$ est isomorphe à
$\widehat{\operatorname{Ker}(v)}$ et comme le support de
$\operatorname{Ker}(v)$ est contenu dans celui de $\mathcal F$, on voit que
$\operatorname{Ker}(u)$ est algébrisable~; démonstration analogue pour
$\operatorname{Im}(u)$ et $\operatorname{Coker}(u)$.
\end{proof}

\begin{proposition}[5.2.3]
\label{III.5.2.3-fr}
Soient $A$ un anneau adique noethérien, $\mathfrak J$ un idéal de définition de
$A$, $\mathfrak Y=\operatorname{Spf}(A)$,
$f:\mathfrak X\to\mathfrak Y$ un morphisme propre de préschémas formels. On pose
$Y_k=\operatorname{Spec}(A/\mathfrak J^{k+1})$,
$X_k=\mathfrak X\times_{\mathfrak Y}Y_k$, et pour tout
$\mathcal O_{\mathfrak X}$-Module $\mathcal F$,
\[
\mathcal F_k=\mathcal F\otimes_{\mathcal O_{\mathfrak X}}\mathcal O_{X_k}
=\mathcal F/\mathfrak J^{k+1}\mathcal F.
\]
Soit $\mathcal L$ un $\mathcal O_{\mathfrak X}$-Module inversible, et supposons que
$\mathcal L_0=\mathcal L/\mathfrak J\mathcal L$ soit un
$\mathcal O_{X_0}$-Module ample~; pour tout $\mathcal O_{\mathfrak X}$-Module
$\mathcal F$ et tout entier $n$, posons
$\mathcal F(n)=\mathcal F\otimes\mathcal L^{\otimes n}$. Alors, pour tout
$\mathcal O_{\mathfrak X}$-Module cohérent $\mathcal F$, il existe un entier $n_0$
tel que, pour tout $n\geq n_0$, les propriétés suivantes aient lieu~:
\begin{enumerate}
\item[(i)] L'homomorphisme canonique
$\mathrm H^0(\mathfrak X,\mathcal F(n))\to
\mathrm H^0(\mathfrak X,\mathcal F_k(n))$ est surjectif pour tout $k\geq0$.
\item[(ii)] On a $\mathrm H^q(\mathfrak X,\mathcal F(n))=0$ pour tout $q>0$.
\end{enumerate}
\end{proposition}

\begin{proof}
On sait que les espaces sous-jacents à $\mathfrak X$ et à $X_0$ sont les mêmes~;
les faisceaux
\[
\mathcal M_k=\mathfrak J^k\mathcal F/\mathfrak J^{k+1}\mathcal F
\]
étant annulés par $\mathfrak J$, peuvent être considérés comme des
$\mathcal O_{X_0}$-Modules cohérents ($0_{\mathrm I}$, 5.3.10)~; en outre, si on
pose $\mathcal M_k(n)=\mathcal M_k\otimes_{\mathcal O_{X_0}}
\mathcal L_0^{\otimes n}$, on voit aussitôt que
\[
\mathcal M_k(n)=\mathfrak J^k\mathcal F(n)/
\mathfrak J^{k+1}\mathcal F(n).
\]
Notons que, puisque $\mathcal L_0$ est ample pour $f_0:X_0\to Y_0$, et

\oldpage[III]{152}

que $f_0$ est propre, on en conclut que $f_0$ est projectif (II, 5.5.4). Soit
\[
S=\bigoplus_{k\geq0}\mathfrak J^k/\mathfrak J^{k+1}
\]
la $A$-algèbre graduée associée à la filtration $\mathfrak J$-adique de $A$, qui est
de type fini puisque $A$ est noethérien~; si on pose
$\mathcal S'=f_0^*(\widetilde S)$,
$\mathcal M=\bigoplus_{k\geq0}\mathcal M_k$, $\mathcal S'$ est une
$\mathcal O_{X_0}$-Algèbre quasi-cohérente de type fini, et $\mathcal M$ un
$\mathcal S'$-Module gradué quasi-cohérent et de type fini (puisque $\mathcal F_0$
est cohérent et engendre le $\mathcal S'$-Module $\mathcal M$). On est donc dans les
conditions d'application du théorème (2.4.1, (ii)), et on en conclut qu'il existe
$n_0$ tel que, pour $n\geq n_0$ et pour tout $k$, on ait
\begin{equation}
\mathrm H^q(X_0,\mathcal M_k(n))=0\qquad\text{pour tout }q>0.
\tag{5.2.3.1}
\label{III.5.2.3.1-fr}
\end{equation}
On a donc aussi $\mathrm H^q(\mathfrak X,\mathcal M_k(n))=0$ pour $q>0$ et
$n\geq n_0$, $\mathcal M_k(n)$ étant cette fois considéré comme
$\mathcal O_{\mathfrak X}$-Module. Appliquant la suite exacte de cohomologie à
\[
0\to
\frac{\mathfrak J^k\mathcal F(n)}{\mathfrak J^{k+1}\mathcal F(n)}
\to
\frac{\mathfrak J^h\mathcal F(n)}{\mathfrak J^{k+1}\mathcal F(n)}
\to
\frac{\mathfrak J^h\mathcal F(n)}{\mathfrak J^k\mathcal F(n)}
\to0,
\]
on en déduit tout d'abord que pour $0\leq h<k$, $n\geq n_0$ et $q>0$, on a par
récurrence sur $k-h$
\begin{equation}
\mathrm H^q\left(\mathfrak X,
\frac{\mathfrak J^h\mathcal F(n)}{\mathfrak J^k\mathcal F(n)}\right)=0
\tag{5.2.3.2}
\label{III.5.2.3.2-fr}
\end{equation}
et en particulier pour $h=0$
\begin{equation}
\mathrm H^q(\mathfrak X,\mathcal F_k(n))=0
\qquad\text{pour }n\geq n_0,\ k\geq0\text{ et }q>0.
\tag{5.2.3.3}
\label{III.5.2.3.3-fr}
\end{equation}
Une autre portion de la suite exacte de cohomologie, pour $h=0$, donne la suite
exacte
\begin{equation}
\mathrm H^0(\mathfrak X,\mathcal F_{k+1}(n))\to
\mathrm H^0(\mathfrak X,\mathcal F_k(n))\to
\mathrm H^1\left(\mathfrak X,
\frac{\mathfrak J^k\mathcal F(n)}{\mathfrak J^{k+1}\mathcal F(n)}\right)=0
\tag{5.2.3.4}
\label{III.5.2.3.4-fr}
\end{equation}
d'où on déduit que pour $h\leq k$, l'application canonique
\begin{equation}
\mathrm H^0(\mathfrak X,\mathcal F_k(n))\longrightarrow
\mathrm H^0(\mathfrak X,\mathcal F_h(n))
\tag{5.2.3.5}
\label{III.5.2.3.5-fr}
\end{equation}
est surjective. Pour tout $q$, le système projectif
$(\mathrm H^q(\mathfrak X,\mathcal F_k(n)))_{k\geq0}$ vérifie donc la condition
(ML) pour $n\geq n_0$. Par ailleurs, tout ouvert formel affine $U$ de $\mathfrak X$
est aussi un ouvert affine dans chacun des $X_k$ (I, 10.5.2), donc on a
$\mathrm H^q(U,\mathcal F_k(n))=0$ pour tout $q>0$ (I, 1.3.1), et
$\mathrm H^0(U,\mathcal F_k(n))\to\mathrm H^0(U,\mathcal F_h(n))$ est surjective
pour $h\leq k$ (I, 1.3.9). Les conditions d'application de ($0$, 13.3.1) sont par
suite remplies, et on en conclut que, pour $n\geq n_0$~:
\begin{enumerate}
\item[$1^\circ$] Pour tout $q>0$,
\[
\mathrm H^q(\mathfrak X,\mathcal F(n))\longrightarrow
\varprojlim_k\mathrm H^q(\mathfrak X,\mathcal F_k(n))
\]
est bijectif, donc, en vertu de (\ref{III.5.2.3.3-fr}),
$\mathrm H^q(\mathfrak X,\mathcal F(n))=0$.
\item[$2^\circ$] L'homomorphisme
\[
\mathrm H^0(\mathfrak X,\mathcal F(n))\longrightarrow
\varprojlim_k\mathrm H^0(\mathfrak X,\mathcal F_k(n))
\]
est bijectif~; par ailleurs, comme les homomorphismes
(\ref{III.5.2.3.5-fr}) sont surjectifs, il en est de même de chacun des
homomorphismes
\[
\varprojlim_k\mathrm H^0(\mathfrak X,\mathcal F_k(n))
\longrightarrow\mathrm H^0(\mathfrak X,\mathcal F_h(n)),
\]
\end{enumerate}
ce qui achève la démonstration.
\end{proof}

\begin{corollary}[5.2.4]
\label{III.5.2.4-fr}
Les hypothèses étant celles de (\hyperref[III.5.2.3-fr]{5.2.3}), pour tout
$\mathcal O_{\mathfrak X}$-Module cohérent $\mathcal F$, il existe un entier $N$
tel que pour $n\geq N$, $\mathcal F(n)$ soit engendré par ses sections au-dessus

\oldpage[III]{153}

de $\mathfrak X$~; en d'autres termes, $\mathcal F$ est isomorphe au quotient d'un
$\mathcal O_{\mathfrak X}$-Module de la forme
$(\mathcal O_{\mathfrak X}(-n))^k$.
\end{corollary}

\begin{proof}
Comme $X_0$ est noethérien, il résulte de l'hypothèse sur $\mathcal L_0$ et de
(II, 4.5.5) qu'il existe $n_0$ tel que, pour $n\geq n_0$, $\mathcal F_0(n)$ soit
engendré par ses sections au-dessus de $\mathfrak X$~; par ailleurs, on peut supposer
$n_0$ pris assez grand pour que l'homomorphisme
\[
\Gamma(\mathfrak X,\mathcal F(n))\longrightarrow
\Gamma(\mathfrak X,\mathcal F_0(n))
\]
soit surjectif pour $n\geq n_0$ (\hyperref[III.5.2.3-fr]{5.2.3}). Il existe donc
un nombre fini de sections $s_i\in\Gamma(\mathfrak X,\mathcal F(n))$ dont les
images dans $\Gamma(\mathfrak X,\mathcal F_0(n))$ engendrent $\mathcal F_0(n)$
($0_{\mathrm I}$, 5.2.3). Comme $\mathfrak J$ est contenu dans l'idéal maximal de
l'anneau local en tout point de $\mathfrak X$, il résulte du lemme de Nakayama,
appliqué à ces anneaux locaux, que les $s_i$ engendrent $\mathcal F(n)$
($0_{\mathrm I}$, 5.1.1).
\end{proof}

\begin{env}[5.2.5]
\label{III.5.2.5-fr}
\emph{Démonstration du théorème d'existence : cas projectif.}
Les notations étant celles de (\hyperref[III.5.1.4-fr]{5.1.4}), supposons $f$
projectif, de sorte qu'il existe un $\mathcal O_X$-Module ample $\mathcal L$
(I, 5.5.4). Par définition, $X_n=\widehat X\times_{\widehat Y}Y_n$ est égal au
sous-préschéma fermé $X\times_Y Y_n=f^{-1}(Y_n)$ de $X$~; si
$\mathcal L'=\mathcal L\otimes_{\mathcal O_X}\mathcal O_{\widehat X}$ est le
complété de $\mathcal L$, on a donc
$\mathcal L'_0=\mathcal L/\mathfrak J\mathcal L$, considéré comme
$\mathcal O_{X_0}$-Module~; on sait que $\mathcal L'_0$ est ample
(II, 4.6.13, (i bis)). On peut donc appliquer à $\mathcal L'$ et à tout
$\mathcal O_{\widehat X}$-Module cohérent $\mathcal F$ le cor.
(\hyperref[III.5.2.4-fr]{5.2.4})~; on voit donc que $\mathcal F$ est isomorphe à
un quotient de
\[
\mathcal G=(\mathcal L'^{\otimes(-n)})^k
\]
pour des entiers $n>0$ et $k>0$ convenables. Or, il est clair que $\mathcal G$ est
le complété de $(\mathcal L^{\otimes(-n)})^k$ (I, 10.8.10), donc est
algébrisable. Considérons ensuite l'homomorphisme canonique surjectif
$u:\mathcal G\to\mathcal F$, et soit $\mathcal H=\operatorname{Ker}(u)$, qui est
un $\mathcal O_{\widehat X}$-Module cohérent ($0_{\mathrm I}$, 5.3.4). On voit de
même qu'il existe un $\mathcal O_{\widehat X}$-Module algébrisable $\mathcal K$ et
un homomorphisme $v:\mathcal K\to\mathcal G$ tel que
$\mathcal H=\operatorname{Im}(v)$. On a alors
$\mathcal F=\operatorname{Coker}(v)$, et $\mathcal F$ est algébrisable en vertu de
(\hyperref[III.5.2.2-fr]{5.2.2}).
\end{env}

\begin{env}[5.2.6]
\label{III.5.2.6-fr}
\emph{Démonstration du théorème d'existence; cas quasi-projectif.}
Les notations étant toujours celles de (\hyperref[III.5.1.4-fr]{5.1.4}), supposons
maintenant que $f$ soit quasi-projectif. Il existe alors un morphisme projectif
$g:Z\to Y$ tel que $X$ s'identifie au $Y$-préschéma induit sur un ouvert de $Z$
(II, 5.3.2)~; si on pose $Z'=g^{-1}(Y')$, on a $X'=X\cap Z'$. Par suite, le
complété $\widehat X=X_{/X'}$ s'identifie au préschéma formel induit par le complété
$\widehat Z=Z_{/Z'}$ sur l'ouvert $X\cap Z'$ de $\widehat Z$ (I, 10.8.5). Soit
$\mathcal F'$ un $\mathcal O_{\widehat X}$-Module cohérent, dont le support $T'$ est
propre sur $\widehat Y$~; cela signifie par définition qu'il existe un sous-préschéma
fermé de $X'$, ayant $T'\subset X'$ comme espace sous-jacent, tel que la restriction
$T'\to Y'$ de $f$ soit propre~; on en conclut que $T'$ est propre sur $Y$, donc
fermé dans $Z'$ (II, 5.4.10). Il en résulte que $\mathcal F'$ est le faisceau induit
sur $\widehat X$ par le $\mathcal O_{\widehat Z}$-Module $\mathcal G'$ obtenu par
recollement de $\mathcal F'$ (défini sur l'ouvert $\widehat X$ de $\widehat Z$) et
du faisceau $0$ sur l'ouvert $\widehat Z-T'$ de $\widehat Z$, ces deux faisceaux
coïncidant dans l'ouvert intersection $\widehat X-T'$. Il est clair que
$\mathcal G'$ est cohérent~; en vertu de
(\hyperref[III.5.2.5-fr]{5.2.5}), il existe un $\mathcal O_Z$-Module cohérent
$\mathcal G$ tel que $\mathcal G'=\widehat{\mathcal G}$~; soit $T$ le support de
$\mathcal G$, de sorte que $T'=T\cap Z'$ (I, 10.8.12). Si $h$ est la restriction
de $g$ au sous-préschéma fermé réduit de $Z$ ayant $T$ comme espace sous-jacent, on
a donc $T'=h^{-1}(Y')=T\cap g^{-1}(Y')$, et par suite $X\cap T$ est un ouvert de
$T$ contenant $T'$.

\oldpage[III]{154}

Comme $h$ est propre (II, 5.4.2), donc fermé, il résulte de
(\hyperref[III.5.1.3.1-fr]{5.1.3.1}) que $X\cap T=T$, autrement dit $T\subset X$,
et comme $T$ est fermé dans $Z$, $T$ est propre sur $Y$. Si $\mathcal F$ est le
faisceau induit sur $X$ par $\mathcal G$, son complété $\widehat{\mathcal F}$ est
induit sur $\widehat X$ par $\widehat{\mathcal G}$ (I, 10.8.4), donc est égal à
$\mathcal F'$, ce qui achève la démonstration.
\end{env}

\subsection{Démonstration du théorème d'existence : cas général}
\label{subsection:III.5.3-fr}

\begin{lemma}[5.3.1]
\label{III.5.3.1-fr}
Sous les conditions de (\hyperref[III.5.1.4-fr]{5.1.4}), si
\[
0\to\mathcal H\to\mathcal F\to\mathcal G\to0
\]
est une suite exacte de $\mathcal O_{\widehat X}$-Modules cohérents telle que
$\mathcal G$ et $\mathcal H$ soient algébrisables, alors $\mathcal F$ est
algébrisable.
\end{lemma}

\begin{proof}
En effet, supposons que $\mathcal G=\widehat{\mathcal B}$,
$\mathcal H=\widehat{\mathcal C}$, $\mathcal B$ et $\mathcal C$ étant des
$\mathcal O_X$-Modules cohérents à supports propres sur $Y$~; $\mathcal F$ définit
canoniquement un élément du $A$-module
\[
\operatorname{Ext}^1_{\mathcal O_{\widehat X}}
(\widehat X;\widehat{\mathcal B},\widehat{\mathcal C})
\]
($0$, 12.3.2), et les hypothèses entraînent que ce $A$-module est canoniquement
isomorphe à
\[
\operatorname{Ext}^1_{\mathcal O_X}(X;\mathcal B,\mathcal C)
\]
(4.5.2)~; il existe donc une suite exacte
$0\to\mathcal C\to\mathcal A\to\mathcal B\to0$ de $\mathcal O_X$-Modules
cohérents telle que l'image canonique de l'élément de
$\operatorname{Ext}^1_{\mathcal O_X}(X;\mathcal B,\mathcal C)$ correspondant à
$\mathcal A$ soit l'élément de
$\operatorname{Ext}^1_{\mathcal O_{\widehat X}}
(\widehat X;\widehat{\mathcal B},\widehat{\mathcal C})$ correspondant à
$\mathcal F$. Mais par définition (compte tenu de (I, 10.8.8, (iii))), cela signifie
que $\mathcal F$ est isomorphe à $\widehat{\mathcal A}$, d'où le lemme, car
$\operatorname{Supp}(\mathcal A)$ est contenu dans la réunion de
$\operatorname{Supp}(\mathcal B)$ et $\operatorname{Supp}(\mathcal C)$, donc est
propre sur $Y$.
\end{proof}

\begin{corollary}[5.3.2]
\label{III.5.3.2-fr}
Sous les conditions de (\hyperref[III.5.1.1-fr]{5.1.1}), soit
$u:\mathcal F\to\mathcal G$ un homomorphisme de
$\mathcal O_{\widehat X}$-Modules cohérents~; si $\mathcal G$,
$\operatorname{Ker}(u)$ et $\operatorname{Coker}(u)$ sont algébrisables, il en est
de même de $\mathcal F$.
\end{corollary}

\begin{proof}
Le lemme (\hyperref[III.5.2.2-fr]{5.2.2}) appliqué à l'homomorphisme
$\mathcal G\to\operatorname{Coker}(u)$ montre en effet que
$\operatorname{Im}(u)$ est algébrisable, et il suffit ensuite d'appliquer le lemme
(\hyperref[III.5.3.1-fr]{5.3.1}) à la suite exacte
\[
0\to\operatorname{Ker}(u)\to\mathcal F\to\operatorname{Im}(u)\to0.
\]
\end{proof}

\begin{lemma}[5.3.3]
\label{III.5.3.3-fr}
Sous les conditions de (\hyperref[III.5.1.1-fr]{5.1.1}), soient $h:Z\to Y$ un
morphisme de type fini, $\widehat Z$ le complété de $Z$ le long de
$Z'=h^{-1}(Y')$, $g:Z\to X$ un $Y$-morphisme propre,
$\widehat g:\widehat Z\to\widehat X$ son prolongement aux complétés. Pour tout
$\mathcal O_{\widehat Z}$-Module algébrisable $\mathcal F'$,
$\widehat g_*(\mathcal F')$ est un $\mathcal O_{\widehat X}$-Module algébrisable.
\end{lemma}

\begin{proof}
En effet, si $\mathcal F$ est un $\mathcal O_Z$-Module cohérent tel que
$\mathcal F'=\widehat{\mathcal F}$, il résulte du premier th. de comparaison
(4.1.5) que $\widehat g_*(\widehat{\mathcal F})$ est isomorphe au complété de
$g_*(\mathcal F)$.
\end{proof}

\begin{lemma}[5.3.4]
\label{III.5.3.4-fr}
Soient $X$ un schéma (usuel) noethérien, $X'$ une partie fermée de $X$,
$f:Z\to X$ un morphisme propre, $Z'=f^{-1}(X')$,
$\widehat X=X_{/X'}$, $\widehat Z=Z_{/Z'}$,
$\widehat f:\widehat Z\to\widehat X$ le prolongement de $f$ aux complétés. Soit
$\mathcal M$ un Idéal cohérent de $\mathcal O_X$ tel que, si
$U=X-\operatorname{Supp}(\mathcal O_X/\mathcal M)$, la restriction
$f^{-1}(U)\to U$ de $f$ soit un isomorphisme. Alors, pour tout
$\mathcal O_{\widehat X}$-Module cohérent $\mathcal F$, il existe un entier $n>0$
tel que le noyau et le conoyau de l'homomorphisme canonique
\[
\rho_{\mathcal F}:\mathcal F\longrightarrow
\widehat f_*\bigl(\widehat f^{,*}(\mathcal F)\bigr)
\]
($0_{\mathrm I}$, 4.4.3) soient annulés par $\widehat{\mathcal M}^{,n}$.
\end{lemma}

\begin{proof}
On peut se borner au cas où $X=\operatorname{Spec}(B)$ où $B$ est un anneau
noethérien, donc $X'=\mathrm V(\mathfrak K)$, où $\mathfrak K$ est un idéal de
$B$. Nous allons voir qu'on peut se ramener au cas où $B$ est un anneau adique
noethérien et $\mathfrak K$ un idéal de définition de $B$. Soit en effet $B_1$ le
séparé complété de $B$ pour la topologie $\mathfrak K$-préadique~; si
$\mathfrak K_1=\mathfrak K B_1$, $B_1$ est donc un

\oldpage[III]{155}

anneau adique noethérien dont $\mathfrak K_1$ est un idéal de définition. Posons
$X_1=\operatorname{Spec}(B_1)$ et soit $h:X_1\to X$ le morphisme correspondant à
l'homomorphisme canonique $B\to B_1$~; si $X'_1=h^{-1}(X')$, on a donc
$X'_1=\mathrm V(\mathfrak K_1)$. Posons enfin
$Z_1=Z\times_X X_1=Z_{(X_1)}$, $f_1=f_{(X_1)}:Z_1\to X_1$, qui est un morphisme
propre (II, 5.4.2), et désignons par $\widehat X_1$ le complété de $X_1$ le long
de $X'_1$, par
$\widehat Z_1=Z_1\times_{X_1}\widehat X_1$ le complété de $Z_1$ le long de
$Z'_1=f_1^{-1}(X'_1)$, par $\widehat f_1$ le prolongement de $f_1$ aux complétés.
Il est immédiat que le prolongement $\widehat h:\widehat X_1\to\widehat X$ de $h$
aux complétés est un isomorphisme, correspondant à l'application identique de $B_1$
(I, 10.9.1)~; on en conclut que l'homomorphisme correspondant
$\widehat Z_1\to\widehat Z$ est aussi un isomorphisme, ces isomorphismes identifiant
$\widehat f_1$ et $\widehat f$. Enfin,
$\mathcal M_1=h^*(\mathcal M)$ est un Idéal cohérent de $\mathcal O_{X_1}$ et
\[
\operatorname{Supp}(\mathcal O_{X_1}/\mathcal M_1)
=h^{-1}(\operatorname{Supp}(\mathcal O_X/\mathcal M))
\]
(I, 9.1.13), donc, si
$U_1=X_1-\operatorname{Supp}(\mathcal O_{X_1}/\mathcal M_1)$, on a
$U_1=h^{-1}(U)$, d'où résulte aussitôt que la restriction
$f_1^{-1}(U_1)\to U_1$ de $f_1$ est un isomorphisme (I, 3.2.7)~; en outre, les
complétés $\widehat{\mathcal M}$ et $\widehat{\mathcal M}_1$ s'identifient par
$\widehat h$ (I, 10.9.5). Toutes les hypothèses de
(\hyperref[III.5.3.4-fr]{5.3.4}) sont donc remplies par $X_1$, $X'_1$, $f_1$ et
$\mathcal M_1$, et on peut donc désormais supposer $B$ adique noethérien et
$\mathfrak K$ un idéal de définition de $B$. On a alors
$\widehat X=\operatorname{Spf}(B)$, et $\mathcal F=N^\Delta$, où $N$ est un
$B$-module de type fini, d'où $\mathcal F=\widehat{\mathcal G}$, où $\mathcal G$
est le $\mathcal O_X$-Module cohérent $\widetilde N$ (I, 10.10.5), et par suite
\[
\widehat f^{,*}(\mathcal F)=\widehat{f^*(\mathcal G)}
\]
(I, 10.9.5). En outre, en vertu du premier théorème de comparaison (4.1.5),
\[
\widehat f_*\bigl(\widehat{f^*(\mathcal G)}\bigr)
\simeq\widehat{f_*\bigl(f^*(\mathcal G)\bigr)},
\]
et l'homomorphisme canonique $\rho_{\mathcal F}$ n'est autre que
$\widehat{\rho}_{\mathcal G}$ en vertu de
(\hyperref[III.5.1.3-fr]{5.1.3}). Or, le noyau $\mathcal P$ et le conoyau
$\mathcal R$ de
\[
\rho_{\mathcal G}:\mathcal G\longrightarrow f_*\bigl(f^*(\mathcal G)\bigr)
\]
sont des $\mathcal O_X$-Modules cohérents, et par hypothèse leurs restrictions à
$U$ sont évidemment nulles. Il existe par suite un entier $n>0$ tel que
$\mathcal M^n\mathcal P=\mathcal M^n\mathcal R=0$ (I, 9.3.4)~; on en conclut que
$\widehat{\mathcal M}^{,n}\widehat{\mathcal P}=
\widehat{\mathcal M}^{,n}\widehat{\mathcal R}=0$
(I, 10.8.8 et 10.8.10).
\end{proof}

\begin{env}[5.3.5]
\label{III.5.3.5-fr}
\emph{Fin de la démonstration du théorème d'existence.}
Les hypothèses étant celles de (\hyperref[III.5.1.4-fr]{5.1.4}), nous allons
utiliser le principe de récurrence noethérienne ($0_{\mathrm I}$, 2.2.2) en supposant
donc le théorème vrai pour tout sous-préschéma fermé $T$ de $X$ dont l'espace
sous-jacent est distinct de $X$ (le complété $\widehat T$ étant bien entendu le
complété de $T$ le long de $T\cap X'$). On peut supposer $X$ non vide. Comme $f$
est séparé et de type fini, on peut appliquer le lemme de Chow (II, 5.6.1)~: il
existe donc un $Y$-schéma $Z$ et un $Y$-morphisme $g:Z\to X$ tels que le morphisme
structural $h:Z\to Y$ soit quasi-projectif, le morphisme $g$ projectif et surjectif,
et en outre un ouvert non vide $U$ de $X$ tel que la restriction
$g^{-1}(U)\to U$ soit un isomorphisme. Soit $\mathcal M$ un Idéal cohérent de
$\mathcal O_X$ définissant un sous-préschéma fermé d'espace sous-jacent $X-U$
(I, 5.2.2), et soit $\mathcal F$ un $\mathcal O_{\widehat X}$-Module cohérent dont
le support $E$ est propre sur $Y$~; désignons par $\widehat Z$ le complété de $Z$
le long de $h^{-1}(Y')$, par $\widehat g:\widehat Z\to\widehat X$ le prolongement
de $g$ aux complétés. Alors $\widehat g^{,*}(\mathcal F)$ est un
$\mathcal O_{\widehat Z}$-Module cohérent dont le support est contenu dans
$g^{-1}(E)$ et est par suite propre sur $Y$, puisque $g$ est projectif, donc propre
(II, 5.4.6). Comme $h$ est quasi-projectif, $\widehat g^{,*}(\mathcal F)$ est
algébrisable en vertu de (\hyperref[III.5.2.6-fr]{5.2.6}). On en conclut

\oldpage[III]{156}

que $\widehat g_*\bigl(\widehat g^{,*}(\mathcal F)\bigr)$ est un
$\mathcal O_{\widehat X}$-Module algébrisable
(\hyperref[III.5.3.3-fr]{5.3.3}) puisque $g$ est propre. On peut maintenant
appliquer à $\mathcal F$ et à $g$ le résultat de
(\hyperref[III.5.3.4-fr]{5.3.4})~: le noyau $\mathcal P$ et le conoyau
$\mathcal R$ de l'homomorphisme
$\rho_{\mathcal F}:\mathcal F\to
\widehat g_*\bigl(\widehat g^{,*}(\mathcal F)\bigr)$ sont annulés par une
puissance $\widehat{\mathcal M}^{,n}$~; soient $T$ le sous-préschéma fermé de $X$
défini par $\mathcal M^n$, ayant $X-U$ pour espace sous-jacent,
$j:T\to X$ l'injection canonique, de sorte que le prolongement aux complétés
$\widehat j:\widehat T\to\widehat X$ est l'injection canonique (I, 10.14.7). On
peut donc écrire
$\mathcal P=\widehat j_*\bigl(\widehat j^{,*}(\mathcal P)\bigr)$ et
$\mathcal R=\widehat j_*\bigl(\widehat j^{,*}(\mathcal R)\bigr)$ et comme $U$
n'est pas vide, il résulte de l'hypothèse de récurrence que
$\widehat j^{,*}(\mathcal P)$ et $\widehat j^{,*}(\mathcal R)$ sont des
$\mathcal O_{\widehat T}$-Modules algébrisables~; en vertu de
(\hyperref[III.5.3.3-fr]{5.3.3}), $\mathcal P$ et $\mathcal R$ sont
algébrisables, et on peut alors appliquer (\hyperref[III.5.3.2-fr]{5.3.2}), qui
prouve finalement que $\mathcal F$ est algébrisable.

\hfill C.Q.F.D.
\end{env}

\subsection{Application : comparaison de morphismes de schémas usuels et de morphismes de schémas formels. Schémas formels algébrisables}
\label{subsection:III.5.4-fr}

\begin{theorem}[5.4.1]
\label{III.5.4.1-fr}
Soient $A$ un anneau adique noethérien, $\mathfrak J$ un idéal de définition de
$A$, $S=\operatorname{Spec}(A)$, $S'=\mathrm V(\mathfrak J)$. Soient
$u:X\to S$ un morphisme propre, $v:Y\to S$ un morphisme séparé de type fini, et
soient $\widehat S$, $\widehat X$, $\widehat Y$ les complétés de $S$, $X$, $Y$ le
long de $S'$, $u^{-1}(S')$, $v^{-1}(S')$ respectivement. Si, pour tout
$S$-morphisme $f:X\to Y$, $\widehat f:\widehat X\to\widehat Y$ est le
prolongement de $f$ aux complétés, l'application $f\mapsto\widehat f$ est une
bijection
\[
\operatorname{Hom}_S(X,Y)\xrightarrow{\ \sim\ }
\operatorname{Hom}_{\widehat S}(\widehat X,\widehat Y).
\]
\end{theorem}

\begin{proof}
Montrons d'abord que $f\mapsto\widehat f$ est injective. Supposons en effet que deux
$S$-morphismes $f$, $g$ de $X$ dans $Y$ soient tels que $\widehat f=\widehat g$.
On sait alors (I, 10.9.4) qu'il existe un voisinage ouvert $V$ de
$X'=u^{-1}(S')$ dans lequel $f$ et $g$ coïncident. Or, comme $u$ est une application
fermée, on a $V=X$ (\hyperref[III.5.1.3.1-fr]{5.1.3.1}), d'où $f=g$.

Prouvons maintenant que $f\mapsto\widehat f$ est surjective, et soit donc $h$ un
$\widehat S$-morphisme $\widehat X\to\widehat Y$. Soit $Z=X\times_S Y$, et
désignons par $p:Z\to X$ et $q:Z\to Y$ les projections canoniques~; $Z$ est de
type fini sur $S$ (I, 6.3.4), donc noethérien~; désignons par $\widehat Z$ son
complété le long de $Z'=p^{-1}(u^{-1}(S'))$~; on sait que $\widehat Z$ s'identifie
canoniquement à $\widehat X\times_{\widehat S}\widehat Y$, les projections
$\widehat Z\to\widehat X$ et $\widehat Z\to\widehat Y$ s'identifiant aux
prolongements $\widehat p$ et $\widehat q$ (I, 10.9.7). Comme $Y$ est séparé sur
$S$, $\widehat Y$ est séparé sur $\widehat S$ (I, 10.15.7), donc le morphisme
graphe
\[
\Gamma_h=(1_{\widehat X},h):\widehat X\to\widehat Z
\]
est une immersion fermée (I, 10.15.4). Soient $\mathfrak T$ le sous-préschéma
formel fermé de $\widehat Z$ associé à cette immersion, et
$j:\mathfrak T\to\widehat Z$ l'injection canonique, de sorte que
$\Gamma_h=j\circ w$, où $w:\widehat X\to\mathfrak T$ est un isomorphisme
(I, 10.14.3) dont l'isomorphisme réciproque est $\widehat p\circ j$~; en outre,
$\mathfrak T$ est évidemment propre sur $\widehat S$, puisque $\widehat X$ l'est~;
on en conclut (\hyperref[III.5.1.8-fr]{5.1.8}) qu'il existe un sous-préschéma
fermé $T$ de $Z$ tel que
$\mathfrak T=\widehat T=T_{/(T\cap Z')}$, et que $j=\widehat i$, où $i$ est
l'injection canonique $T\to Z$ (I, 10.14.7). Alors $p\circ i:T\to X$ est un
isomorphisme, car il en est ainsi de
$\widehat{p\circ i}=\widehat p\circ\widehat i$ par hypothèse, et il suffit
d'appliquer

\oldpage[III]{157}

(4.6.8) en notant comme ci-dessus que $S$ est le seul voisinage de $S'$ dans $S$.
Soit $g:X\to T$ l'isomorphisme réciproque de $p\circ i$, et posons
$f=q\circ i\circ g$, qui est un morphisme $X\to Y$ dont par définition le graphe
$\Gamma_f=i\circ g$. Comme $\widehat g$ est l'isomorphisme réciproque de
$\widehat{p\circ i}=w$, on a
\[
\widehat{\Gamma_f}=\widehat i\circ\widehat g=j\circ w=\Gamma_h.
\]
Mais on sait que $\widehat{\Gamma_f}=\Gamma_{\widehat f}$ (I, 10.9.8), d'où
finalement $h=\widehat f$, ce qui achève la démonstration.
\end{proof}

On peut donc dire, dans le langage des catégories, que le foncteur
$X\mapsto\widehat X$ est pleinement fidèle ($0$, 8.1.6) de la catégorie des
schémas propres sur $\operatorname{Spec}(A)$ dans la catégorie des schémas formels
propres sur $\operatorname{Spf}(A)$, pour tout anneau adique noethérien $A$~; il
établit par suite une équivalence entre la première de ces catégories et une
sous-catégorie de la seconde~; les objets de cette dernière seront appelés
\emph{schémas formels algébrisables}. Pour un tel schéma $\mathfrak X$, il existe
un schéma usuel $X$, propre sur $\operatorname{Spec}(A)$, déterminé à isomorphisme
unique près, tel que $\mathfrak X$ soit isomorphe à $\widehat X$.

\begin{env}[5.4.2]
\label{III.5.4.2-fr}
\emph{Scholie.}
Avec les notations de (\hyperref[III.5.4.1-fr]{5.4.1}), posons
$S_n=\operatorname{Spec}(A/\mathfrak J^{n+1})$,
$X_n=X\times_S S_n$, $Y_n=Y\times_S S_n$. Il résulte de
(\hyperref[III.5.4.1-fr]{5.4.1}) et de (I, 10.12.3) que se donner un
$S$-morphisme $f:X\to Y$ équivaut à se donner un $(S_n)$-système inductif adique
(I, 10.12.2) de $S_n$-morphismes $f_n:X_n\to Y_n$.
\end{env}

\begin{remark}[5.4.3]
\label{III.5.4.3-fr}
Contrairement à ce que pourrait suggérer le th. d'existence
(\hyperref[III.5.1.6-fr]{5.1.6}), il y a des schémas formels propres sur
$\operatorname{Spf}(A)$ et qui ne sont pas algébrisables (tout comme il y a des
espaces analytiques compacts qui ne proviennent pas de variétés algébriques
complexes). Nous rencontrerons plus tard de tels schémas dans la «~théorie des
modules~», qui traite précisément (lorsque le corps de base est $\mathbb C$) des
variations infinitésimales de la structure complexe d'une variété algébrique
complète, et on sait que de telles variations peuvent donner naissance à des
variétés analytiques qui ne sont pas algébriques.
\end{remark}

\begin{proposition}[5.4.4]
\label{III.5.4.4-fr}
Soient $A$ un anneau adique noethérien, $\mathfrak S=\operatorname{Spf}(A)$,
$g:\mathfrak X\to\mathfrak S$, $h:\mathfrak Y\to\mathfrak S$ deux morphismes
propres de schémas formels, $f:\mathfrak X\to\mathfrak Y$ un
$\mathfrak S$-morphisme. Si $f$ est fini et si $\mathfrak Y$ est algébrisable,
alors $\mathfrak X$ est algébrisable.
\end{proposition}

\begin{proof}
On notera que les hypothèses sur $g$ et $h$ entraînent déjà que $f$ est propre
(3.4.1), et pour que $f$ soit fini, il suffit que pour tout $y\in\mathfrak Y$, la
fibre $f^{-1}(y)$ soit finie (\hyperref[III.4.8.11-fr]{4.8.11}). L'hypothèse
entraîne que $\mathcal B=f_*(\mathcal O_{\mathfrak X})$ est une
$\mathcal O_{\mathfrak Y}$-Algèbre cohérente
(\hyperref[III.4.8.6-fr]{4.8.6}), donc il résulte du th. d'existence que, si
$\mathfrak Y=\widehat Y$ et $h=\widehat w$, où
$w:Y\to\operatorname{Spec}(A)$ est un morphisme propre de schémas usuels, il
existe une $\mathcal O_Y$-Algèbre cohérente $\mathcal C$ telle que
$\mathcal B=\widehat{\mathcal C}$. Soient $X=\operatorname{Spec}(\mathcal C)$, et
$u:X\to Y$ le morphisme structural~; alors, il résulte aussitôt de la définition de
$\mathfrak X$ à partir de $\mathcal B$
(\hyperref[III.4.8.7-fr]{4.8.7}) que $\mathfrak X$ est canoniquement isomorphe à
$\widehat X$ et que $f$ s'identifie à $\widehat u$ (il suffit pour le voir de
considérer le cas où $Y$ est affine).

On notera que (\hyperref[III.5.1.8-fr]{5.1.8}) est un cas particulier de
(\hyperref[III.5.4.4-fr]{5.4.4}).
\end{proof}

\begin{theorem}[5.4.5]
\label{III.5.4.5-fr}
Soient $A$ un anneau adique noethérien, $\mathfrak J$ un idéal de définition de
$A$, $S=\operatorname{Spec}(A)$,
$\mathfrak S=\widehat S=\operatorname{Spf}(A)$,
$f:\mathfrak X\to\mathfrak S$ un morphisme propre de schémas formels. On pose
$S_k=\operatorname{Spec}(A/\mathfrak J^{k+1})$,
$X_k=\mathfrak X\times_{\mathfrak S}S_k$, et pour tout
$\mathcal O_{\mathfrak X}$-Module $\mathcal F$,
$\mathcal F_k=\mathcal F\otimes_{\mathcal O_{\mathfrak X}}\mathcal O_{X_k}
=\mathcal F/\mathfrak J^{k+1}\mathcal F$.

\oldpage[III]{158}

Soit $\mathcal L$ un $\mathcal O_{\mathfrak X}$-Module inversible, et supposons
que $\mathcal L_0=\mathcal L/\mathfrak J\mathcal L$ soit un
$\mathcal O_{X_0}$-Module ample. Alors $\mathfrak X$ est algébrisable, et si
$X$ est un $S$-schéma propre tel que $\mathfrak X$ soit isomorphe à
$\widehat X$, il existe un $\mathcal O_X$-Module ample $\mathcal M$ tel que
$\mathcal L$ soit isomorphe à $\widehat{\mathcal M}$ (ce qui entraîne que
$X$ est projectif sur $S$).
\end{theorem}

\begin{proof}
Appliquons (\hyperref[III.5.2.3-fr]{5.2.3}) à
$\mathcal F=\mathcal O_{\mathfrak X}$~: il existe donc un entier $n_0$ tel que
pour $n\geq n_0$, l'homomorphisme canonique
$\Gamma(\mathfrak X,\mathcal L^{\otimes n})\to
\Gamma(X_0,\mathcal L_0^{\otimes n})$ soit surjectif. On peut supposer
$n\geq n_0$ choisi assez grand pour que $\mathcal L_0^{\otimes n}$ soit très
ample pour $S_0$ (II, 4.5.10). Comme le morphisme $f_0:X_0\to S_0$ est propre,
$\Gamma(X_0,\mathcal L_0^{\otimes n})$ est un $A$-module de type fini (3.2.1),
donc il existe un sous-$A$-module de type fini $E$ de
$\Gamma(\mathfrak X,\mathcal L^{\otimes n})$ dont l'image dans
$\Gamma(X_0,\mathcal L_0^{\otimes n})$ soit ce dernier module tout entier.
Cela étant, pour tout $k\geq0$, considérons l'homomorphisme
\[
u_k:E/\mathfrak J^{k+1}E\longrightarrow
\Gamma(X_k,\mathcal L_k^{\otimes n})
\]
déduit de l'injection canonique
$E\to\Gamma(\mathfrak X,\mathcal L^{\otimes n})$. Notons que
$(f_k)_*(\mathcal L_k^{\otimes n})$ est quasi-cohérent, et comme
$\Gamma(S_k,(f_k)_*(\mathcal L_k^{\otimes n}))
=\Gamma(X_k,\mathcal L_k^{\otimes n})$, $u_k$ définit un homomorphisme
\[
\widetilde u_k:(E/\mathfrak J^{k+1}E)^\sim
\longrightarrow(f_k)_*(\mathcal L_k^{\otimes n}),
\]
et par suite aussi un homomorphisme
\[
\widetilde u_k^\sharp:
f_k^*((E/\mathfrak J^{k+1}E)^\sim)\longrightarrow
\mathcal L_k^{\otimes n}.
\]
D'ailleurs, si on pose
$\mathcal G_k=f_k^*((E/\mathfrak J^{k+1}E)^\sim)$, on a
$\mathcal G_0=\mathcal G_k/\mathfrak J\mathcal G_k$ (I, 9.1.5), donc
$\widetilde u_0^\sharp:\mathcal G_k/\mathfrak J\mathcal G_k
\to\mathcal L_k^{\otimes n}/\mathfrak J\mathcal L_k^{\otimes n}$ se déduit
de $\widetilde u_k^\sharp$ par passage aux quotients. Or, par définition de
$E$, $\widetilde u_0^\sharp$ n'est autre que l'homomorphisme canonique
\[
\sigma:f_0^*((f_0)_*(\mathcal L_0^{\otimes n}))\longrightarrow
\mathcal L_0^{\otimes n},
\]
et l'hypothèse que $\mathcal L_0^{\otimes n}$ est très ample entraîne que
$\widetilde u_0^\sharp$ est surjectif (II, 4.4.3)~; on déduit alors du lemme
de Nakayama que chacun des $\widetilde u_k^\sharp$ est aussi surjectif.
Chacun des $\widetilde u_k^\sharp$ définit donc (II, 4.2.2) un
$S_k$-morphisme
\[
g_k:X_k\longrightarrow P_k=\mathbf P(E/\mathfrak J^{k+1}E),
\]
et comme $P_h=P_k\times_{S_k}S_h$ pour $h\leq k$ en vertu de (II, 4.1.3),
$(g_k)$ est un $(S_k)$-système inductif adique (I, 10.12.2) en vertu des
relations entre les $\widetilde u_k^\sharp$ et de (II, 4.2.10). Les $g_k$
définissent donc un $\mathfrak S$-morphisme de schémas formels
$g:\mathfrak X\to\mathfrak P$, où $\mathfrak P$ est la limite inductive du
système $(P_k)$, ou encore le complété $\widehat P$, où $P=\mathbf P(E)$. De
plus, l'hypothèse que $\mathcal L_0^{\otimes n}$ est très ample entraîne que
$g_0$ est une immersion fermée (II, 4.4.3)~; on en conclut que $g$ est une
immersion fermée de schémas formels (4.8.10), donc $\mathfrak X$ est
algébrisable (5.1.8). Le fait que $\mathcal L$ soit isomorphe au complété
$\widehat{\mathcal M}$ d'un $\mathcal O_X$-Module inversible résulte alors du
th. d'existence (5.1.6). En outre, $\mathcal L^{\otimes n}$ est alors le
complété de $\mathcal M^{\otimes n}$ (I, 10.8.10), et les homomorphismes
$\widetilde u_k^\sharp$ définissent un homomorphisme bien déterminé
\[
v:f^*(\widetilde E)\longrightarrow\mathcal M^{\otimes n}
\]
(5.1.7)~; d'ailleurs, comme $\widetilde u_k^\sharp$ est surjectif, il en est de
même de $\widehat v$ (I, 10.11.5), donc de $v$ (5.1.3)~; en outre, le morphisme
$r:X\to P$ défini par $v$ (II, 4.2.2) a pour prolongement aux complétés $g$,
et comme $g$ est une immersion fermée, il en est de même de $r$, par (5.1.8)
et (5.4.1)~; on en conclut que $\mathcal M^{\otimes n}$ est très ample
(II, 4.4.6) et $\mathcal M$ est ample (II, 4.5.10).
\end{proof}

\begin{remark}[5.4.6]
\label{III.5.4.6-fr}
Soient $A$ un anneau adique noethérien,
$\mathfrak S=\operatorname{Spf}(A)$,
$f:\mathfrak X\to\mathfrak S$ un morphisme propre de schémas formels. Soit
$\mathcal N$ l'Idéal cohérent de $\mathcal O_{\mathfrak X}$ tel que pour tout
ouvert formel affine $U$ de $\mathfrak X$, $\Gamma(U,\mathcal N)$ soit le
nilradical de $\Gamma(U,\mathcal O_{\mathfrak X})$~; l'existence de cet Idéal
résulte facilement de (I, 10.10.2) et du fait que tout homomorphisme d'anneaux
$B\to C$ applique le nilradical de $B$ dans celui de $C$. Soit
$\mathfrak X'$ le sous-schéma formel fermé de $\mathfrak X$ défini par
$\mathcal N$ (I, 10.14.2)~; il serait intéressant de savoir si, lorsque
$\mathfrak X'$ est algébrisable, $\mathfrak X$ lui-même est algébrisable. On
parviendrait sans doute à une solution de ce problème si on savait classifier
(par exemple au moyen d'invariants de nature

\oldpage[III]{159}

cohomologique), les \emph{extensions} d'un faisceau structural
$\mathcal O_{\mathfrak X}$ (pour un préschéma usuel ou un préschéma formel) par
un Idéal de carré nul, autrement dit les $\mathcal O_{\mathfrak X}$-Algèbres
$\mathcal A$ telles que $\mathcal O_{\mathfrak X}$ soit isomorphe à
$\mathcal A/\mathcal J$, où $\mathcal J$ est un Idéal de carré nul de
$\mathcal A$.
\end{remark}

\subsection{Une décomposition de certains schémas}
\label{subsection:III.5.5-fr}

\begin{proposition}[5.5.1]
\label{III.5.5.1-fr}
Soient $A$ un anneau adique noethérien, $\mathfrak J$ un idéal de définition
de $A$, $Y=\operatorname{Spec}(A)$. Soit $f:X\to Y$ un morphisme séparé de
type fini~; on pose $Y_0=\operatorname{Spec}(A/\mathfrak J)$,
$X_0=X\times_Y Y_0=f^{-1}(Y_0)$. Soit $Z_0$ une partie ouverte de $X_0$,
propre sur $Y_0$~; il existe alors dans $X$ une partie ouverte et fermée $Z$,
propre sur $Y$ et telle que $Z\cap X_0=Z_0$.
\end{proposition}

\begin{proof}
Par hypothèse, il y a un ouvert $T$ de $X$ tel que $T\cap X_0=Z_0$~; soit
$\widehat T$ le complété le long de $Z_0$ du schéma induit par $X$ sur
l'ouvert $T$~; le support de $\mathcal O_{\widehat T}$ étant $Z_0$, qui est
propre sur $Y_0$, $\widehat T$ est propre sur
$\widehat Y=\operatorname{Spf}(A)$ (3.4.1). Il résulte de (5.1.8) qu'il existe
un sous-schéma fermé $Z$ de $T$ propre sur $Y$, tel que, si $i:Z\to T$ est
l'injection canonique, $\widehat i:\widehat Z\to\widehat T$ soit un
isomorphisme ($\widehat Z$ étant le complété de $Z$ le long de $Z_0$). On en
conclut (4.6.8) qu'il existe dans $T$ un voisinage ouvert $V$ de $Z_0$ tel que
la restriction $i^{-1}(V)\to V$ de $i$ soit un isomorphisme. Mais
$i^{-1}(V)$ est un voisinage de $Z_0$ dans $Z$, donc est nécessairement
identique à $Z$ (5.1.3.1). On en conclut que $Z$ est ouvert dans $T$, donc
dans $X$, ce qui achève la démonstration.
\end{proof}

\begin{corollary}[5.5.2]
\label{III.5.5.2-fr}
Si $X_0$ est propre sur $Y_0$, $X$ est réunion de deux parties ouvertes
disjointes $Z$ et $Z'$ telles que $Z$ soit propre sur $Y$ et contienne $X_0$~;
en outre, toute partie fermée $P$ de $X$, propre sur $Y$, est contenue dans
$Z$.
\end{corollary}

\begin{proof}
La dernière assertion résulte de ce que $P\cap Z'$ étant fermé dans $P$ est
propre sur $Y$~; si $P\cap Z'$ n'était pas vide, $f(P\cap Z')$ serait fermé
non vide dans $Y$, donc rencontrerait $Y_0$ (5.1.3.1), ce qui contredit la
définition de $Z$.
\end{proof}

\begin{flushright}
\emph{(A suivre.)}
\end{flushright}

% La page imprimée 504 est blanche.
\clearpage
\thispagestyle{empty}
\null
\clearpage
\thispagestyle{plain}

\oldpage[III]{161}
\section*{BIBLIOGRAPHIE (suite)}

\begin{flushleft}
[27] P. Gabriel, \emph{Des catégories abéliennes}, thèse, Paris (1961).

\medskip
[28] J. E. Roos, \emph{Sur les foncteurs dérivés de $\varprojlim$.
Applications}, C. R. Acad. Sci. Paris, t. CCLII, 1961, p. 3702--3704.

\medskip
[29] H. Cartan, \emph{Séminaire de l'École Normale Supérieure},
13\textsuperscript{e} année (1960--1961), exposé n\textsuperscript{o} 11.
\end{flushleft}

\clearpage
\oldpage[III]{162}
\section*{INDEX DES NOTATIONS}

\begin{flushleft}
$\mathbf{Ens}$ : 0, 8.1.1.

$h_X(Y)$, $h_X(u)$ ($X$, $Y$ objets d'une catégorie $C$, $u$ morphisme de
$C$) : 0, 8.1.1.

$h_w(u)$ ($u$, $w$ morphismes de $C$) : 0, 8.1.2.

$Z_k(E_r^{pq})$, $B_k(E_r^{pq})$ ($k\geq r$) : 0, 11.1.1.

$Z_\infty(E_2^{pq})$, $B_\infty(E_2^{pq})$ : 0, 11.1.1.

$K^\bullet$, $K_\bullet$ (complexes) : 0, 11.2.1.

$K^{\bullet\bullet}$, $K_{\bullet\bullet}$ (bicomplexes) : 0, 11.2.1.

$B^i(K^\bullet)$, $Z^i(K^\bullet)$, $H^i(K^\bullet)$,
$H^\bullet(K^\bullet)$ : 0, 11.2.1.

$B_i(K_\bullet)$, $Z_i(K_\bullet)$, $H_i(K_\bullet)$,
$H_\bullet(K_\bullet)$ : 0, 11.2.1.

$T(K^\bullet)$, $T(K_\bullet)$ ($T$ foncteur) : 0, 11.2.1.

$E(K^\bullet)$ ($K^\bullet$ complexe filtré) : 0, 11.2.2.

$K^{i,\bullet}$, $K^{\bullet,j}$, $Z_{\mathrm{II}}^p(K^{i,\bullet})$,
$B_{\mathrm{II}}^p(K^{i,\bullet})$, $H_{\mathrm{II}}^p(K^{i,\bullet})$,
$Z_{\mathrm I}^p(K^{\bullet,j})$, $B_{\mathrm I}^p(K^{\bullet,j})$,
$H_{\mathrm I}^p(K^{\bullet,j})$ : 0, 11.3.1.

$H_{\mathrm{II}}^p(K^{\bullet\bullet})$,
$Z_{\mathrm I}^q(H_{\mathrm{II}}^p(K^{\bullet\bullet}))$,
$B_{\mathrm I}^q(H_{\mathrm{II}}^p(K^{\bullet\bullet}))$,
$H_{\mathrm I}^q(H_{\mathrm{II}}^p(K^{\bullet\bullet}))$,
$H^n(K^{\bullet\bullet})$ : 0, 11.3.1.

$F_{\mathrm I}^p(K^{\bullet\bullet})$,
$F_{\mathrm{II}}^p(K^{\bullet\bullet})$, ${}'E(K^{\bullet\bullet})$,
${}''E(K^{\bullet\bullet})$ : 0, 11.3.2.

$K_{i,\bullet}$, $K_{\bullet,j}$, $H_p^{\mathrm{II}}(K_{i,\bullet})$,
$H_p^{\mathrm I}(K_{\bullet,j})$, $H_p^{\mathrm{II}}(K_{\bullet\bullet})$,
$H_p^{\mathrm I}(K_{\bullet\bullet})$,
$H_q^{\mathrm I}(H_p^{\mathrm{II}}(K_{\bullet\bullet}))$,
$H_q^{\mathrm{II}}(H_p^{\mathrm I}(K_{\bullet\bullet}))$,
$H_n(K_{\bullet\bullet})$ : 0, 11.3.5.

$\mathrm R^\bullet T(K^\bullet)$ ($T$ foncteur covariant additif) : 0, 11.4.3.

$\mathrm R^\bullet T(K^\bullet,K'^{\bullet})$ ($T$ bifoncteur covariant
additif) : 0, 11.4.6.

$\mathrm L_\bullet T(K_\bullet)$ ($T$ foncteur covariant additif) : 0, 11.6.2.

$\mathrm L_\bullet T(K_\bullet,K'_\bullet)$ ($T$ bifoncteur covariant
additif) : 0, 11.6.6.

$\mathrm L_\bullet T(K_{\bullet\bullet})$ ($T$ foncteur covariant additif) :
0, 11.7.2.

$|\sigma|$, $\Sigma(A)$, $\mathrm C_\bullet(A)$, $\mathrm L_\bullet(A)$ ($A$ ensemble fini,
$\sigma$ simplexe de $A$) : 0, 11.8.1.

$\mathrm C^\bullet(A,B;\mathscr S)$, $\mathrm L^\bullet(A,B;\mathscr S)$ ($A$, $B$ ensembles
finis, $\mathscr S$ système de coefficients) : 0, 11.8.4.

$\mathrm P^\bullet(A,B,\mathscr S)$ ($A$, $B$ ensembles finis, $\mathscr S$ système
de coefficients) : 0, 11.8.6.

$\mathrm C^\bullet(A,B;\mathscr S^\bullet)$,
$\mathrm L^\bullet(A,B;\mathscr S^\bullet)$ ($A$, $B$ ensembles finis,
$\mathscr S^\bullet$ complexe de systèmes de coefficients) : 0, 11.8.8.

$\mathrm P^\bullet(A,B;\mathscr S^\bullet)$ ($A$, $B$ ensembles finis,
$\mathscr S^\bullet$ complexe de systèmes de coefficients) : 0, 11.8.10.

$\chi(M^\bullet)$ ($M^\bullet$ complexe fini de modules de longueur finie) :
0, 11.10.1.

$\check C^\bullet(\mathfrak U,\mathcal F)$ ($\mathcal F$ faisceau de groupes abéliens
sur $X$, $\mathfrak U$ recouvrement ouvert de $X$) : 0, 12.1.2.

$\mathrm R^pf_*(\mathcal F)$ ($f:X\to Y$ morphisme d'espaces annelés,
$\mathcal F$ faisceau de $\mathcal O_X$-Modules) : 0, 12.2.1.

$\mathcal H^p(f,\mathcal K^\bullet)$,
$\mathcal H_f^p(\mathcal K^\bullet)$, ${}'E(f,\mathcal K^\bullet)$,
${}''E(f,\mathcal K^\bullet)$ ($f:X\to Y$ morphisme d'espaces
annelés, $\mathcal K^\bullet$ complexe de $\mathcal O_X$-Modules) : 0, 12.4.1.

$\mathbf H^\bullet(X,\mathcal K^\bullet)$,
$\mathbf H^\bullet(V,\mathcal K^\bullet)$ ($X$ espace annelé,
$\mathcal K^\bullet$ complexe de $\mathcal O_X$-Modules, $V$ ouvert de $X$) :
0, 12.4.2.

$\operatorname{gr}^\bullet(A)$ ($A$ système projectif vérifiant (ML)) :
0, 13.4.2.

$E(A)$ ($A$ objet filtré) : 0, 13.6.4.

$K_\bullet(\mathbf f)$, $i_{\mathbf f}$ ($\mathbf f$ famille finie d'éléments
d'un anneau) : III, 1.1.1.

$K^\bullet(\mathbf f,M)$, $K_\bullet(\mathbf f,M)$ ($\mathbf f$ famille
finie d'éléments d'un anneau $A$, $M$ $A$-module) : III, 1.1.2.

$H^\bullet(\mathbf f,M)$, $H_\bullet(\mathbf f,M)$ ($\mathbf f$ famille
finie d'éléments d'un anneau $A$, $M$ $A$-module) : III, 1.1.3.

$C^\bullet((\mathbf f),M)$, $H^\bullet((\mathbf f),M)$ ($\mathbf f$ famille
finie d'éléments d'un anneau $A$, $M$ $A$-module) : III, 1.1.6.

$H^\bullet(U,\mathcal F^{(*)})$,
$H^\bullet(\mathfrak U,\mathcal F^{(*)})$ ($\mathcal F$
$\mathcal O_X$-Module quasi-cohérent, $U$ ouvert de $X$, $\mathfrak U$
recouvrement ouvert de $X$) : III, 2.1.1.
\end{flushleft}

\clearpage
\oldpage[III]{163}
\section*{INDEX TERMINOLOGIQUE}

\begin{flushleft}
Aboutissement d'une suite spectrale : 0, 11.1.1.

Algébrisable ($\mathcal O_X$-Module) : III, 5.2.1.

Algébrisable (schéma formel) : III, 5.4.2.

Analytiquement intègre (anneau) : III, 4.3.6.

Application quasi-compacte : 0, 9.1.1.

Augmentation d'une résolution : 0, 11.4.1.

Caractéristique d'Euler-Poincaré d'un complexe : 0, 11.10.1.

Caractéristique d'Euler-Poincaré d'un $\mathcal O_X$-Module cohérent ($X$
projectif sur $\operatorname{Spec}(A)$, $A$ anneau artinien) : III, 2.5.1.

Cochaîne bi-alternée : 0, 11.8.4.

Complexe défini par un bicomplexe : 0, 11.3.1.

Complexe de l'algèbre extérieure : III, 1.1.1.

Condition (ML) : 0, 13.1.2.

Constructible (partie, ensemble) : 0, 9.1.2.

Constructible (fonction) : 0, 9.3.1.

Cup-produit : 0, 12.1.2.

Dihomomorphisme pour les structures algébriques sur les catégories : 0, 8.2.1.

Exact (sous-ensemble) dans une catégorie abélienne : III, 3.1.1.

Essentiellement constant (système projectif) : 0, 13.4.2.

Factorisation de Stein : III, 4.3.3.

Filtration co-discrète, co-séparée, discrète, exhaustive, finie, séparée :
0, 11.1.3.

Final (objet) d'une catégorie : 0, 8.1.10.

Fini (morphisme) de préschémas formels : III, 4.8.2.

Fini (préschéma formel) au-dessus de $\mathfrak Y$,
$\mathfrak Y$-fini (préschéma formel) : III, 4.8.2.

Foncteur covariant canonique $C\to\operatorname{Hom}(C^o,\mathbf{Ens})$ :
0, 8.1.2.

Genre arithmétique : III, 2.5.1.

Géométriquement connexe (fibre) : III, 4.3.4.

Homomorphisme de structures algébriques sur les catégories : 0, 8.2.1.

Hypercohomologie d'un foncteur par rapport à un complexe $K^\bullet$ :
0, 11.4.3.

Hypercohomologie d'un bifoncteur par rapport à deux complexes
$K^\bullet$, $K'^{\bullet}$ : 0, 11.4.6.

Hyperhomologie d'un foncteur par rapport à un complexe $K_\bullet$ :
0, 11.6.2.

Hyperhomologie d'un bifoncteur par rapport à deux complexes $K_\bullet$,
$K'_\bullet$ : 0, 11.6.6.

Hyperhomologie d'un foncteur par rapport à un bicomplexe : 0, 11.7.2.

Localement constructible (ensemble) : 0, 9.1.11.

Loi de composition externe (interne) sur les objets d'une catégorie :
0, 8.2.1.

\oldpage[III]{164}

$S$-$C$-module filtré, $\operatorname{gr}^\bullet(S)$-$C$-module gradué,
$\operatorname{gr}^{\bullet\bullet}(S)$-$C$-module bigradué : 0, 13.6.6.

Morphisme de suites spectrales : 0, 11.1.2.

Nombre géométrique de composantes connexes d'une fibre : III, 4.3.4.

$C$-objet en groupes, $C$-groupe, $C$-anneau, $C$-module : 0, 8.2.3.

Objet des bords, objet des cobords, objet des cocycles, objet des cycles :
0, 11.2.1.

Objet des images universelles : 0, 13.1.1.

Objet gradué associé à un système projectif satisfaisant à (ML) : 0, 13.4.2.

Objet gradué limité inférieurement (supérieurement) : 0, 11.2.1.

Partie propre (sur $\mathfrak Y$) d'un préschéma formel : III, 3.4.1.

Pleine (sous-catégorie) : 0, 8.1.5.

Pleinement fidèle (foncteur) : 0, 8.1.5.

Polynôme de Hilbert : III, 2.5.3.

Propre (morphisme) de préschémas formels : III, 3.4.1.

Représentable (foncteur) : 0, 8.1.8.

Résolution cohomologique, droite, gauche, homologique, injective, libre,
plate, projective : 0, 11.4.1.

Résolution de Cartan-Eilenberg droite, injective : 0, 11.4.2.

Résolution de Cartan-Eilenberg gauche, projective : 0, 11.6.1.

Rétrocompact (ensemble) : 0, 9.1.1.

Strict (système projectif) : 0, 13.4.2.

Suite spectrale dans une catégorie abélienne : 0, 11.1.1.

Suite spectrale faiblement convergente, régulière, corégulière, birégulière :
0, 11.1.3.

Suite spectrale dégénérée : 0, 11.1.6.

Suite spectrale d'un foncteur relativement à un objet filtré : 0, 13.6.4.

Suites spectrales d'un bicomplexe : 0, 11.3.2 et 11.3.5.

Suites spectrales d'hypercohomologie d'un foncteur par rapport à un complexe :
0, 11.4.3.

Suites spectrales d'hyperhomologie d'un foncteur par rapport à un complexe :
0, 11.6.2.

Suites spectrales d'hyperhomologie d'un foncteur par rapport à un bicomplexe :
0, 11.7.2.

Système de coefficients : 0, 11.8.4.

Unibranche (anneau, point) : III, 4.3.6.

Universellement ouvert (morphisme) : III, 4.3.9.
\end{flushleft}

\clearpage
\oldpage[III]{165}
\section*{TABLE DES MATIÈRES}

\begin{flushleft}
\textsc{Chapitre 0.} --- \textbf{Préliminaires (suite)} \dotfill 5

\medskip
\S~8. Foncteurs représentables \dotfill 5

8.1. Foncteurs représentables \dotfill 5

8.2. Structures algébriques dans les catégories \dotfill 9

\medskip
\S~9. Ensembles constructibles \dotfill 12

9.1. Ensembles constructibles \dotfill 12

9.2. Ensembles constructibles dans les espaces noethériens \dotfill 14

9.3. Fonctions constructibles \dotfill 16

\medskip
\S~10. Compléments sur les modules plats \dotfill 17

10.1. Relations entre modules plats et modules libres \dotfill 17

10.2. Critères locaux de platitude \dotfill 18

10.3. Existence d'extensions plates d'anneaux locaux \dotfill 20

\medskip
\S~11. Compléments d'algèbre homologique \dotfill 23

11.1. Rappels sur les suites spectrales \dotfill 23

11.2. La suite spectrale d'un complexe filtré \dotfill 27

11.3. Les suites spectrales d'un bicomplexe \dotfill 29

11.4. Hypercohomologie d'un foncteur par rapport à un complexe $K^\bullet$
\dotfill 32

11.5. Passage à la limite inductive dans l'hypercohomologie \dotfill 35

11.6. Hyperhomologie d'un foncteur par rapport à un complexe $K_\bullet$
\dotfill 39

11.7. Hyperhomologie d'un foncteur par rapport à un bicomplexe
$K_{\bullet\bullet}$ \dotfill 41

11.8. Compléments sur la cohomologie des complexes simpliciaux \dotfill 43

11.9. Un lemme sur les complexes de type fini \dotfill 46

11.10. Caractéristique d'Euler-Poincaré d'un complexe de modules de longueur
finie \dotfill 48

\medskip
\S~12. Compléments sur la cohomologie des faisceaux \dotfill 49

12.1. Cohomologie des faisceaux de modules sur les espaces annelés \dotfill 49

12.2. Images directes supérieures \dotfill 57

12.3. Compléments sur les foncteurs Ext de faisceaux \dotfill 60

12.4. Hypercohomologie du foncteur image directe \dotfill 62

\oldpage[III]{166}

\S~13. Limites projectives en algèbre homologique \dotfill 64

13.1. La condition de Mittag-Leffler \dotfill 64

13.2. La condition de Mittag-Leffler pour les groupes abéliens \dotfill 65

13.3. Application~: cohomologie d'une limite projective de faisceaux \dotfill 68

13.4. Condition de Mittag-Leffler et objets gradués associés aux systèmes
projectifs \dotfill 69

13.5. Limites projectives de suites spectrales de complexes filtrés \dotfill 71

13.6. Suite spectrale d'un foncteur relative à un objet muni d'une filtration
finie \dotfill 73

13.7. Foncteurs dérivés d'une limite projective d'arguments \dotfill 75

\medskip
\textsc{Chapitre III.} --- \textbf{Étude cohomologique des faisceaux
cohérents} \dotfill 81

\medskip
\S~1. Cohomologie des schémas affines \dotfill 82

1.1. Rappels sur le complexe de l'algèbre extérieure \dotfill 82

1.2. Cohomologie de Čech d'un recouvrement ouvert \dotfill 85

1.3. Cohomologie d'un schéma affine \dotfill 88

1.4. Application à la cohomologie des préschémas quelconques \dotfill 89

\medskip
\S~2. Étude cohomologique des morphismes projectifs \dotfill 95

2.1. Calculs explicites de certains groupes de cohomologie \dotfill 95

2.2. Le théorème fondamental des morphismes projectifs \dotfill 100

2.3. Application aux faisceaux gradués d'algèbres et de modules \dotfill 102

2.4. Une généralisation du théorème fondamental \dotfill 107

2.5. Caractéristique d'Euler-Poincaré et polynôme de Hilbert \dotfill 109

2.6. Application~: critères d'amplitude \dotfill 111

\medskip
\S~3. Le théorème de finitude pour les morphismes propres \dotfill 115

3.1. Le lemme de dévissage \dotfill 115

3.2. Le théorème de finitude~: cas des schémas usuels \dotfill 116

3.3. Généralisation du théorème de finitude (schémas usuels) \dotfill 118

3.4. Le théorème de finitude~: cas des schémas formels \dotfill 119

\medskip
\S~4. Le théorème fondamental des morphismes propres. Applications \dotfill 122

4.1. Le théorème fondamental \dotfill 122

4.2. Cas particuliers et variantes \dotfill 129

4.3. Le théorème de connexion de Zariski \dotfill 130

4.4. Le «~main theorem~» de Zariski \dotfill 135

4.5. Complétés de modules d'homomorphismes \dotfill 138

4.6. Relations entre morphismes formels et morphismes usuels \dotfill 139

4.7. Un critère d'amplitude \dotfill 145

4.8. Morphismes finis de préschémas formels \dotfill 146

\oldpage[III]{167}

\S~5. Un théorème d'existence de faisceaux algébriques cohérents \dotfill 149

5.1. Énoncé du théorème \dotfill 149

5.2. Démonstration du théorème d'existence~: cas projectif et quasi-projectif
\dotfill 151

5.3. Démonstration du théorème d'existence~: cas général \dotfill 154

5.4. Application~: comparaison de morphismes de schémas usuels et de
morphismes de schémas formels. Schémas formels algébrisables \dotfill 156

5.5. Une décomposition de certains schémas \dotfill 159

\medskip
\textsc{Bibliographie (suite)} \dotfill 161

\textsc{Index des notations} \dotfill 162

\textsc{Index terminologique} \dotfill 163

\begin{flushright}
\emph{Reçu le 28 juin 1961.}
\end{flushright}
\end{flushleft}


% EGA III, seconde partie (sections 6--7), admitted from the sealed
% direct-NUMDAM French handoff.  Section 5 remains the live III-1 cursor.
\clearpage
\setcounter{section}{5}
\setcounter{subsection}{0}
\setcounter{equation}{0}
\setcounter{footnote}{0}
\chapter*{CHAPITRE III (suite)\\[1.2ex]
\Large ÉTUDE COHOMOLOGIQUE DES FAISCEAUX COHÉRENTS}

\section{Foncteurs « Tor » locaux et globaux; formule de Künneth}
\label{section:III.6.fr}

\subsection{Introduction}
\label{subsection:III.6.1.fr}

\oldpage[III]{137}

\begin{env}[6.1.1]
Soient $f:X\to Y$ un morphisme de préschémas, $\mathscr F$ un
$\mathscr O_X$-Module quasi-cohérent. Dans l'étude des « images directes
supérieures » $\mathrm R^n f_*(\mathscr F)$, on est amené à considérer le
problème général suivant : étant donné un morphisme « changement de base »
$g:Y'\to Y$, on pose
\[
  X'=X_{(Y')}=X\times_Y Y',\qquad
  \mathscr F'=\mathscr F\otimes_Y\mathscr O_{Y'},\qquad
  f'=f_{(Y')}:X'\to Y',
\]
et l'on se propose d'avoir des renseignements sur les images directes
supérieures $\mathrm R^n f'_*(\mathscr F')$ (en supposant connus les
$\mathrm R^n f_*(\mathscr F)$). On voit aisément (cf. par exemple (7.7.2))
que l'on est ramené à étudier les variations de
$\mathrm R^n f_*(\mathscr F\otimes_{\mathscr O_Y}\mathscr G)$ pour un
$\mathscr O_Y$-Module quasi-cohérent variable $\mathscr G$, autrement dit le
foncteur
\[
  \mathscr G\longmapsto
  \mathrm R^n f_*(\mathscr F\otimes_{\mathscr O_Y}\mathscr G).
\]
Si $\mathscr F$ est plat sur $Y$, le foncteur
$\mathscr G\mapsto\mathscr F\otimes_{\mathscr O_Y}\mathscr G$ est exact
(0\textsubscript{I}, 6.7.4) et par suite le foncteur composé
$\mathscr G\mapsto\mathrm R^n f_*(\mathscr F\otimes_{\mathscr O_Y}\mathscr G)$
est encore un foncteur cohomologique. Mais il n'en est plus de même dans le
cas général; pour pouvoir appliquer les méthodes cohomologiques, on est
conduit à substituer à
$\mathscr G\mapsto\mathrm R^n f_*(\mathscr F\otimes_{\mathscr O_Y}\mathscr G)$
d'autres foncteurs, qui cette fois sont toujours des foncteurs
cohomologiques. Ces foncteurs, qui généralisent les foncteurs « Tor » de la
théorie des modules, sont définis aux n\textsuperscript{os} 6.3 à 6.7; il y a
d'ailleurs deux telles généralisations, l'une « locale » et l'autre
« globale », reliées par des suites spectrales qui seront discutées au
n\textsuperscript{o} 6.7; comme application de ces suites spectrales, on
obtient en particulier, sous certaines conditions, une « formule de Künneth »
exprimant
\[
  \mathrm R^n(f_1\times f_2)_*(\mathscr F_1\otimes_Y\mathscr F_2)
\]
à l'aide des images directes supérieures
$\mathrm R^p f_{1*}(\mathscr F_1)$ et $\mathrm R^q f_{2*}(\mathscr F_2)$.
D'autres suites spectrales (6.8) généralisent les suites spectrales
d'associativité du foncteur « Tor » de modules; enfin, le problème du
changement de base conduit lui aussi à des suites spectrales (6.9).
\end{env}

\begin{env}[6.1.2]
En outre, on constate (6.10) que les foncteurs cohomologiques
$\mathscr G\mapsto\mathscr T_\bullet(\mathscr G)$ ainsi définis sont
localement (sur $Y$) du type
\[
  \mathscr G\longmapsto
  \mathscr H_\bullet(\mathscr L_\bullet\otimes_Y\mathscr G),
\]
où $\mathscr L_\bullet$ est un complexe de $\mathscr O_Y$-modules localement
libres (défini à une homotopie près) et $\mathscr H_\bullet$ l'homologie. Il y
a alors intérêt à oublier la situation particulière qui a donné naissance à
$\mathscr T_\bullet$, et à étudier de

\oldpage[III]{138}

façon générale les foncteurs de la forme précédente
$\mathscr G\mapsto\mathscr H_\bullet(\mathscr L_\bullet\otimes_Y\mathscr G)$
(où l'on fait en outre le cas échéant des hypothèses de finitude appropriées
sur les $\mathscr L_i$ ou les $\mathscr H_i(\mathscr L_\bullet)$) : c'est ce
qui est l'objet du § 7, dont la lecture, pour l'essentiel, est indépendante du
§ 6. Les propriétés les plus importantes de ces foncteurs concernent les
propriétés d'exactitude d'une composante $\mathscr T_i$ de
$\mathscr T_\bullet$; on donnera divers critères permettant d'établir de
telles propriétés; comme application, on obtiendra des conditions permettant
d'affirmer (avec les notations de (6.1.1)) que le foncteur
$\mathscr G\mapsto\mathrm R^n f_*(\mathscr F\otimes_{\mathscr O_Y}\mathscr G)$
est exact (ce qu'on exprimera en disant que $\mathscr F$ est
cohomologiquement plat sur $Y$ en dimension $n$). Une autre propriété
importante pour les composantes $\mathscr T_i$ de
$\mathscr T_\bullet(\mathscr G)=
\mathscr H_\bullet(\mathscr L_\bullet\otimes_Y\mathscr G)$ est une propriété
de semi-continuité de la fonction
$y\mapsto\dim_{k(y)}(T_i(k(y)))$; lorsque $\mathscr T_i$ est exact, cette
propriété est remplacée par une propriété de continuité, la réciproque étant
d'ailleurs vraie d'après Grauert lorsque $Y$ est réduit (7.8.4).
\end{env}

\begin{env}[6.1.3]
Dans les §§ 6 et 7, nous avons systématiquement fait usage de
l'hypercohomologie, en prenant partout comme arguments des complexes de
faisceaux au lieu de faisceaux, bien que la nécessité de ce point de vue
n'apparaisse que dans des chapitres ultérieurs. Le formalisme cohomologique
développé à cette occasion deviendra d'ailleurs plus transparent dans le
chapitre de ce Traité qui sera consacré à la mise au point d'une algèbre des
foncteurs cohomologiques de faisceaux cohérents, incluant le formalisme de la
dualité. Mais cela demandera des développements qui sortent du cadre du
présent chapitre.
\end{env}

\begin{env}[6.1.4]
Pour abréger, étant donnés deux complexes $K^\bullet$, $K'^\bullet$ dans une
catégorie abélienne $C$, nous dirons qu'un morphisme de complexes
$f:K^\bullet\to K'^\bullet$ est un homotopisme s'il existe un morphisme
$g:K'^\bullet\to K^\bullet$ tel que les morphismes composés $f\circ g$ et
$g\circ f$ soient tous deux homotopes à l'identité (par abus de langage,
lorsqu'il existe un tel homotopisme, on dira aussi que $K^\bullet$ et
$K'^\bullet$ sont homotopes). Lorsqu'on peut définir l'hypercohomologie d'un
foncteur covariant additif $T$ de $C$ dans une catégorie abélienne $C'$, par
rapport à un complexe de $C$ (0, 11.4.3), il est immédiat qu'un homotopisme
$K^\bullet\to K'^\bullet$ de complexes de $C$ définit canoniquement un
isomorphisme
\[
  \mathrm R^\bullet T(K^\bullet)\xrightarrow{\sim}
  \mathrm R^\bullet T(K'^\bullet)
\]
pour l'hypercohomologie (loc. cit.).
\end{env}

\subsection{Hypercohomologie des complexes de modules sur un préschéma}
\label{subsection:III.6.2.fr}

\begin{env}[6.2.1]
Soient $X$ un préschéma,
$\mathscr K^\bullet=(\mathscr K^i)_{i\in\mathbf Z}$ un complexe de
$\mathscr O_X$-Modules dont l'opérateur de dérivation est de degré $+1$.
Rappelons que pour tout morphisme $f:X\to Y$ de préschémas, on a défini
(0, 12.4.1) les $\mathscr O_Y$-Modules d'hypercohomologie
$\mathscr H^n(f,\mathscr K^\bullet)$ (aussi notés
$\mathscr H_f^n(\mathscr K^\bullet)$ ou
$\mathrm R^n f_*(\mathscr K^\bullet)$) pour tout $n\in\mathbf Z$;
l'hypercohomologie $\mathscr H^\bullet(f,\mathscr K^\bullet)$ est
l'aboutissement des deux foncteurs spectraux
$'\mathscr E(f,\mathscr K^\bullet)$ et
$''\mathscr E(f,\mathscr K^\bullet)$, dont les termes $\mathscr E_2$ sont
donnés par
\begin{align}
  '\mathscr E_2^{pq}
    &=\mathscr H^p\bigl(\mathscr H^q(f,\mathscr K^\bullet)\bigr),
    \tag{6.2.1.1}\label{III.6.2.1.1.fr}\\
  ''\mathscr E_2^{pq}
    &=\mathscr H^p\bigl(f,\mathscr H^q(\mathscr K^\bullet)\bigr)
      =\mathrm R^p f_*\bigl(\mathscr H^q(\mathscr K^\bullet)\bigr).
    \tag{6.2.1.2}\label{III.6.2.1.2.fr}
\end{align}

\oldpage[III]{139}

où $\mathscr H^q(f,\mathscr K^\bullet)$ est le complexe dont le composant de
degré $i$ est
$\mathscr H^q(f,\mathscr K^i)=\mathrm R^q f_*(\mathscr K^i)$ (loc. cit.).
Rappelons aussi que lorsque $Y$ est réduit à un point, on note
l'hypercohomologie correspondante $\mathbf H^\bullet(X,\mathscr K^\bullet)$
(qui est formée de modules sur $\Gamma(X,\mathscr O_X)$ indépendants du
préschéma ponctuel $Y$ considéré); lorsque $Y=X$ et $f=1_X$, on a
$\mathscr H^n(f,\mathscr K^\bullet)=\mathscr H^n(\mathscr K^\bullet)$
(cohomologie du complexe $\mathscr K^\bullet$); lorsque
$\mathscr K^i=0$, sauf pour $i=i_0$, on a
\[
  \mathscr H^n(f,\mathscr K^\bullet)
    =\mathrm R^{n-i_0}f_*(\mathscr K^{i_0}).
\]
La suite spectrale $'\mathscr E(f,\mathscr K^\bullet)$ est toujours
régulière; les deux suites spectrales sont birégulières lorsque
$\mathscr K^\bullet$ est limité inférieurement (0, 12.4.1).

Tout homotopisme $h:\mathscr K^\bullet\to\mathscr K'^\bullet$ de complexes de
$\mathscr O_X$-Modules (6.1.4) donne un isomorphisme
$\mathscr H^\bullet(f,\mathscr K^\bullet)\xrightarrow{\sim}
\mathscr H^\bullet(f,\mathscr K'^\bullet)$ pour l'hypercohomologie. Il en est
de même lorsqu'on suppose seulement que
$\mathscr H^\bullet(h):\mathscr H^\bullet(\mathscr K^\bullet)\to
\mathscr H^\bullet(\mathscr K'^\bullet)$ est un isomorphisme et que
$\mathscr K^\bullet$ et $\mathscr K'^\bullet$ sont limités inférieurement,
comme il résulte aussitôt de (0, 11.1.5) appliqué à la suite spectrale
(6.2.1.2) et à l'analogue pour $\mathscr K'^\bullet$. Enfin, pour tout
recouvrement ouvert $\mathfrak U=(U_\alpha)$ de $X$, on a aussi défini
(0, 12.4.5) l'hypercohomologie
$\mathbf H^\bullet(\mathfrak U,\mathscr K^\bullet)$ comme la cohomologie du
bicomplexe $C^\bullet(\mathfrak U,\mathscr K^\bullet)$ (dont le composant
d'indices $(i,j)$ est par définition
$C^i(\mathfrak U,\mathscr K^j)$); les
$\mathbf H^n(\mathfrak U,\mathscr K^\bullet)$ sont encore des modules sur
$\Gamma(X,\mathscr O_X)$.
\end{env}

\begin{proposition}[6.2.2]
Soient $X$ un schéma, $\mathfrak U=(U_\alpha)$ un recouvrement de $X$ par des
ouverts affines. Pour tout complexe $\mathscr K^\bullet$ de
$\mathscr O_X$-Modules quasi-cohérents, les modules d'hypercohomologie
$\mathbf H^\bullet(X,\mathscr K^\bullet)$ et
$\mathbf H^\bullet(\mathfrak U,\mathscr K^\bullet)$ sont canoniquement
isomorphes.
\end{proposition}

\begin{proof}
En effet, toute intersection finie $V$ d'ouverts du recouvrement
$\mathfrak U$ est affine (I, 5.5.6), donc
$\mathrm H^q(V,\mathscr K^i)=0$ pour tout $i$ et tout $q>0$ (1.3.1); la
proposition est donc un cas particulier de (0, 12.4.7).
\end{proof}

\begin{proposition}[6.2.3]
Soit $f:X\to Y$ un morphisme quasi-compact et séparé de préschémas. Pour tout
complexe $\mathscr K^\bullet$ de $\mathscr O_X$-Modules quasi-cohérents, les
$\mathscr O_Y$-Modules $\mathscr H^n(f,\mathscr K^\bullet)$ sont
quasi-cohérents.
\end{proposition}

\begin{proof}
Comme les $\mathscr H^q(f,\mathscr K^i)=\mathrm R^q f_*(\mathscr K_i)$ sont
des $\mathscr O_Y$-Modules quasi-cohérents (1.4.10), il en est de même de
$'\mathscr E_2^{pq}$, qui, d'après (6.2.1.1), est quotient d'un noyau
d'homomorphisme de Modules quasi-cohérents par une image d'un tel
homomorphisme (I, 4.1.1). Pour la même raison, tous les
$\mathscr O_Y$-Modules $'\mathscr E_r^{pq}$,
$B_k('\mathscr E_2^{pq})$, $Z_k('\mathscr E_2^{pq})$ de la première suite
spectrale sont quasi-cohérents. La régularité de la suite spectrale
$'\mathscr E(f,\mathscr K^\bullet)$ entraîne que
$Z_\infty('\mathscr E_2^{pq})$ est égal à un des
$Z_k('\mathscr E_2^{pq})$, donc est quasi-cohérent, et il en est de même de
$B_\infty('\mathscr E_2^{pq})=\varinjlim_k B_k('\mathscr E_2^{pq})$
(0, 11.2.4 et I, 4.1.1); les $'\mathscr E_\infty^{pq}$ sont donc aussi
quasi-cohérents. La suite spectrale précédente étant régulière, la filtration
des $F^p(\mathscr H^n(f,\mathscr K^\bullet))$ est discrète et exhaustive;
autrement dit, le $\mathscr O_Y$-Module
$\mathscr H^n(f,\mathscr K^\bullet)$ est réunion d'une suite croissante
$(\mathscr G_k)_{k\geq0}$ de $\mathscr O_Y$-Modules telle que
$\mathscr G_0=0$ et que chaque $\mathscr G_k/\mathscr G_{k-1}$ soit égal à un
des $\mathscr O_Y$-Modules $'\mathscr E_\infty^{pq}$, donc soit
quasi-cohérent. Par récurrence sur $k$, on en déduit que les $\mathscr G_k$
sont quasi-cohérents (1.4.17), et comme
$\mathscr H^n(f,\mathscr K^\bullet)=\varinjlim_k\mathscr G_k$, la proposition
est démontrée (I, 4.1.1).
\end{proof}

\oldpage[III]{140}

\begin{corollary}[6.2.4]
Sous les hypothèses de (6.2.3), pour tout ouvert affine $V$ de $Y$,
l'homomorphisme canonique
\[
  \mathbf H^n(f^{-1}(V),\mathscr K^\bullet)
  \longrightarrow
  \Gamma(V,\mathscr H^n(f,\mathscr K^\bullet))
  \tag{6.2.4.1}\label{III.6.2.4.1.fr}
\]
est bijectif pour tout $n\in\mathbf Z$.
\end{corollary}

\begin{proof}
La démonstration est la même que celle de (1.4.11), en utilisant (6.2.2),
remplaçant $\mathscr F$ par $\mathscr K^\bullet$,
$\mathscr K^\bullet$ par
$f_*(C^\bullet(\mathfrak U,\mathscr K^\bullet))$,
$\mathscr H^\bullet(\mathscr K^\bullet)$ par
$\mathscr H^\bullet(f,\mathscr K^\bullet)$, et notant que ce dernier est un
$\mathscr O_Y$-Module quasi-cohérent par (6.2.3).
\end{proof}

\begin{proposition}[6.2.5]
Soient $Y$ un préschéma localement noethérien, $f:X\to Y$ un morphisme
propre, $\mathscr K^\bullet$ un complexe de $\mathscr O_X$-Modules tel que les
$\mathscr O_X$-Modules $\mathscr H^q(\mathscr K^\bullet)$ soient cohérents.
Alors les $\mathscr O_Y$-Modules $\mathscr H^n(f,\mathscr K^\bullet)$ sont
cohérents.
\end{proposition}

\begin{proof}
La question étant locale sur $Y$, on peut se borner au cas où $Y$ est
noethérien et affine, et il s'agit donc, en vertu de (6.2.4), de prouver que
les $\mathbf H^n(X,\mathscr K^\bullet)$ sont des
$\Gamma(Y,\mathscr O_Y)$-modules de type fini. On a alors
$\mathbf H^\bullet(X,\mathscr K^\bullet)=
\mathbf H^\bullet(\mathfrak U,\mathscr K^\bullet)$ (6.2.2), où l'on peut
supposer que $\mathfrak U$ est fini, puisque $X$ est quasi-compact. Les
cochaînes de chaque complexe $C^\bullet(\mathfrak U,\mathscr K^j)$ étant
alternées par définition, il y a un entier $r>0$ tel que
$C^i(\mathfrak U,\mathscr K^j)=0$ pour $i<0$ et $i>r$; on en conclut
(0, 11.3.3) que les deux suites spectrales du bicomplexe
$C^\bullet(\mathfrak U,\mathscr K^\bullet)$ sont birégulières. Comme les
intersections des ensembles de $\mathfrak U$ sont des ouverts affines
(I, 5.5.6), chaque foncteur
$\mathscr F\mapsto C^i(\mathfrak U,\mathscr F)$ est exact dans la catégorie
des $\mathscr O_X$-Modules quasi-cohérents; donc
$\mathrm H_I^q(C^i(\mathfrak U,\mathscr K^\bullet))
=C^i(\mathfrak U,\mathscr H^q(\mathscr K^\bullet))$, et les termes
$E_2$ de la seconde suite spectrale de
$C^\bullet(\mathfrak U,\mathscr K^\bullet)$ sont donnés (0, 11.3.2) par
\[
  ''E_2^{pq}
  =\mathrm H^p(C^\bullet(\mathfrak U,\mathscr H^q(\mathscr K^\bullet)))
  =\mathbf H^p(\mathfrak U,\mathscr H^q(\mathscr K^\bullet))
  =\mathrm H^p(X,\mathscr H^q(\mathscr K^\bullet))
\]
en vertu de (1.4.1); puisque $f$ est propre, ce sont des
$\Gamma(Y,\mathscr O_Y)$-modules de type fini (3.2.1). La suite spectrale
$''E(C^\bullet(\mathfrak U,\mathscr K^\bullet))$ étant birégulière, on en
déduit bien que les $\mathbf H^n(X,\mathscr K^\bullet)$ sont des
$\Gamma(Y,\mathscr O_Y)$-modules de type fini (0, 11.1.8).

A fortiori, si $\mathscr K^\bullet$ est un complexe de
$\mathscr O_X$-Modules cohérents, les $\mathscr O_Y$-Modules
$\mathscr H^n(f,\mathscr K^\bullet)$ sont cohérents sous les hypothèses de
(6.2.5) relatives à $Y$ et $f$ (0\textsubscript{I}, 5.3.4).
\end{proof}

\begin{env}[6.2.6]
L'hypercohomologie $\mathscr H^\bullet(f,\mathscr K^\bullet)$ est un foncteur
cohomologique dans la catégorie des complexes de $\mathscr O_X$-Modules
limités inférieurement (0, 12.4.4). C'est un foncteur cohomologique dans la
catégorie de tous les complexes de $\mathscr O_X$-Modules lorsque le
morphisme $f$ est quasi-compact et l'espace sous-jacent à $X$ localement
noethérien : en effet, il résulte alors de (G, II, 3.10.1) que $f_*$ permute
aux limites inductives (la question étant locale sur $Y$), et l'on peut
appliquer (0, 11.5.2).

Enfin, si $f$ est séparé, $\mathscr H^\bullet(f,\mathscr K^\bullet)$ est un
foncteur cohomologique dans la catégorie des complexes de
$\mathscr O_X$-Modules quasi-cohérents. C'est immédiat lorsque $Y$ est
affine, car alors $X$ est un schéma, donc, en vertu de l'isomorphisme
canonique (6.2.2), on est ramené à voir que
$\mathscr K^\bullet\mapsto\mathbf H^\bullet(\mathfrak U,\mathscr K^\bullet)$
est un foncteur cohomologique dans la catégorie des complexes de
$\mathscr O_X$-Modules quasi-cohérents, ce qui est immédiat puisque le
foncteur $\mathscr K^\bullet\mapsto
C^\bullet(\mathfrak U,\mathscr K^\bullet)$ est exact dans cette catégorie
(I, 1.3.7). Dans le cas général, pour tout ouvert affine $V$

\oldpage[III]{141}

de $Y$, $f^{-1}(Y)$ est un schéma, et pour appliquer ce qui précède, il
suffit de vérifier que pour une suite exacte
\[
  0\longrightarrow\mathscr K'^\bullet
  \longrightarrow\mathscr K^\bullet
  \longrightarrow\mathscr K''^\bullet
  \longrightarrow0
\]
de complexes de $\mathscr O_X$-Modules quasi-cohérents, l'homomorphisme
\[
  \partial:
  \mathscr H^n(f,\mathscr K''^\bullet)|V
  \longrightarrow
  \mathscr H^{n+1}(f,\mathscr K'^\bullet)|V
\]
ne dépend pas du recouvrement ouvert affine $\mathfrak U$ de $f^{-1}(V)$
utilisé pour le définir. Mais cela résulte de ce que, si $\mathfrak U'$ est
un recouvrement ouvert affine plus fin que $\mathfrak U$, le diagramme
\[
  \xymatrix@C=34pt@R=11pt{
    \mathbf H^\bullet(\mathfrak U,\mathscr K^\bullet)
      \ar[rrd]^-{\sim}\ar[dd]^{\iota}
    && \\
    &&
    \mathbf H^\bullet(f^{-1}(V),\mathscr K^\bullet)
    \\
    \mathbf H^\bullet(\mathfrak U',\mathscr K^\bullet)
      \ar[rru]_-{\sim}
    &&
  }
\]
d'isomorphismes canoniques est commutatif, ainsi que le diagramme
\[
  \xymatrix@C=46pt@R=25pt{
    \mathbf H^n(\mathfrak U,\mathscr K''^\bullet)
      \ar[r]^-{\partial}\ar[d]^{\iota}
    &
    \mathbf H^{n+1}(\mathfrak U,\mathscr K'^\bullet)
      \ar[d]^{\iota}
    \\
    \mathbf H^n(\mathfrak U',\mathscr K''^\bullet)
      \ar[r]_-{\partial}
    &
    \mathbf H^{n+1}(\mathfrak U',\mathscr K'^\bullet).
  }
\]

Lorsque l'une des conditions précédentes est remplie et que
$\mathscr K^\bullet\to\mathscr K'^\bullet$ est un homotopisme (6.1.4),
l'isomorphisme correspondant
$\mathscr H^\bullet(f,\mathscr K^\bullet)\xrightarrow{\sim}
\mathscr H^\bullet(f,\mathscr K'^\bullet)$ est alors un isomorphisme de
$\partial$-foncteurs (0, 11.4.4).
\end{env}

\begin{env}[6.2.7]
Tout ce qui précède s'applique naturellement sans changement (sinon de
notations) à un complexe $\mathscr K_\bullet$ de $\mathscr O_X$-Modules
quasi-cohérents dont l'opérateur de dérivation est de degré $-1$; il suffit
de considérer le complexe $\mathscr K^\bullet=(\mathscr K^i)$ où
$\mathscr K^i=\mathscr K_{-i}$ pour tout $i\in\mathbf Z$.
\end{env}

\subsection{Hypertor de deux complexes de modules}
\label{subsection:III.6.3.fr}

\begin{env}[6.3.1]
Soient $A$ un anneau commutatif, $P_\bullet$, $Q_\bullet$ deux complexes de
$A$-modules dont les opérateurs de dérivation sont de degré $-1$; soit
$L_{\bullet\bullet}$ (resp. $M_{\bullet\bullet}$) une résolution projective
de Cartan-Eilenberg de $P_\bullet$ (resp. $Q_\bullet$) (0, 11.6.1);
$L_{\bullet\bullet}\otimes_A M_{\bullet\bullet}$ est alors (pour la somme
des premiers degrés et la somme des seconds degrés) un bicomplexe (à
opérateurs de dérivation de degré $-1$), dont l'homologie
$H_\bullet(L_{\bullet\bullet}\otimes_A M_{\bullet\bullet})$ ne dépend pas des
résolutions de Cartan-Eilenberg $L_{\bullet\bullet}$,
$M_{\bullet\bullet}$ choisies, et est par définition l'hyperhomologie du
bifoncteur $P_\bullet\otimes_A Q_\bullet$ en $P_\bullet$ et $Q_\bullet$
(0, 11.6.5). Nous poserons par définition
\[
  \mathbf{Tor}_n^A(P_\bullet,Q_\bullet)
  =H_n(L_{\bullet\bullet}\otimes_A M_{\bullet\bullet})
  \tag{6.3.1.1}\label{III.6.3.1.1.fr}
\]

\oldpage[III]{142}

et nous dirons que cet $A$-module est l'hypertor d'indice $n$ des deux
complexes $P_\bullet$, $Q_\bullet$. On sait que dans la catégorie des
complexes de $A$-modules limités inférieurement, les
$\mathbf{Tor}_n^A(P_\bullet,Q_\bullet)$ forment un bifoncteur homologique en
$P_\bullet$, $Q_\bullet$ (0, 11.6.5). En outre :
\end{env}

\begin{proposition}[6.3.2]
Le bifoncteur $\mathbf{Tor}_\bullet^A(P_\bullet,Q_\bullet)$ est
l'aboutissement commun de deux bifoncteurs spectraux
$'E(P_\bullet,Q_\bullet)$, $''E(P_\bullet,Q_\bullet)$, dont les termes $E_2$
sont
\begin{align}
  'E^2_{pq}
    &=H_p\bigl(\operatorname{Tor}_q^A(P_\bullet,Q_\bullet)\bigr),
    \tag{6.3.2.1}\label{III.6.3.2.1.fr}\\
  ''E^2_{pq}
    &=\bigoplus_{q'+q''=q}
      \operatorname{Tor}_p^A\bigl(H_{q'}(P_\bullet),H_{q''}(Q_\bullet)\bigr).
    \tag{6.3.2.2}\label{III.6.3.2.2.fr}
\end{align}
où, dans (6.3.2.1), $\operatorname{Tor}_q^A(P_\bullet,Q_\bullet)$ désigne le
bicomplexe formé des $A$-modules $\operatorname{Tor}_q^A(P_i,Q_j)$. La suite
spectrale (6.3.2.2) est toujours régulière; si $P_\bullet$ et $Q_\bullet$
sont limités inférieurement, ou si $A$ est de dimension cohomologique finie,
les deux suites spectrales (6.3.2.1) et (6.3.2.2) sont birégulières.
\end{proposition}

\begin{proof}
Cela résulte de (0, 11.6.5), car lorsque $A$ est de dimension cohomologique
finie $n$, tout $A$-module admet une résolution projective de longueur $n$
(M, VI, 2.1).
\end{proof}

\begin{corollary}[6.3.3]
Soient $P'_\bullet$, $Q'_\bullet$ deux complexes de $A$-modules,
$u:P_\bullet\to P'_\bullet$, $v:Q_\bullet\to Q'_\bullet$ deux homomorphismes
de complexes. Si les homomorphismes
$H_\bullet(u):H_\bullet(P_\bullet)\to H_\bullet(P'_\bullet)$,
$H_\bullet(v):H_\bullet(Q_\bullet)\to H_\bullet(Q'_\bullet)$ déduits
respectivement de $u$ et $v$ sont bijectifs, alors l'homomorphisme
$\mathbf{Tor}_\bullet^A(P_\bullet,Q_\bullet)\to
\mathbf{Tor}_\bullet^A(P'_\bullet,Q'_\bullet)$ déduit de $u$ et $v$ est
bijectif.
\end{corollary}

\begin{proof}
En effet, l'homomorphisme de suites spectrales
$''E(P_\bullet,Q_\bullet)\to''E(P'_\bullet,Q'_\bullet)$ déduit de $u$ et $v$
est alors un isomorphisme pour les termes $E_2$ et la conclusion résulte de
ce que ces suites sont régulières en vertu de (6.3.2) (0, 11.1.5).
\end{proof}

\begin{proposition}[6.3.4]
Soient $P_\bullet$, $Q_\bullet$ deux complexes de $A$-modules, limités
inférieurement. Soit $L_{\bullet\bullet}$ (resp. $M_{\bullet\bullet}$) un
bicomplexe formé de $A$-modules plats, tel que pour tout $i$,
$L_{i,\bullet}$ (resp. $M_{i,\bullet}$) soit une résolution de $P_i$
(resp. $Q_i$). On a alors des isomorphismes canoniques
\[
  \mathbf{Tor}_\bullet^A(P_\bullet,Q_\bullet)
  \xrightarrow{\sim}H_\bullet(L_{\bullet\bullet}\otimes_AQ_\bullet)
  \xrightarrow{\sim}H_\bullet(P_\bullet\otimes_AM_{\bullet\bullet})
  \xrightarrow{\sim}H_\bullet(L_{\bullet\bullet}\otimes_AM_{\bullet\bullet})
  \tag{6.3.4.1}\label{III.6.3.4.1.fr}
\]
\end{proposition}

\begin{proof}
Cela résulte de (0, 11.6.5, (ii) et (iii)) et de la définition des
$A$-modules plats.
\end{proof}

\begin{namedtheorem}[6.3.5]{Remarques}
(i) Avec les notations de (6.3.1), les bicomplexes
$L_{\bullet\bullet}\otimes_A M_{\bullet\bullet}$ et
$M_{\bullet\bullet}\otimes_A L_{\bullet\bullet}$ sont canoniquement
isomorphes, d'où un isomorphisme canonique
$\mathbf{Tor}_\bullet^A(P_\bullet,Q_\bullet)\xrightarrow{\sim}
\mathbf{Tor}_\bullet^A(Q_\bullet,P_\bullet)$.

(ii) Si $F$ et $G$ sont deux $A$-modules, $P_\bullet$ et $Q_\bullet$ les
complexes de $A$-modules réduits à $F$ et $G$ respectivement en degré 0 et
nuls dans les autres degrés, alors deux résolutions projectives $L_\bullet$,
$M_\bullet$ de $F$ et $G$ respectivement peuvent être considérées comme des
résolutions de Cartan-Eilenberg de $P_\bullet$ et $Q_\bullet$ en les
complétant par des zéros. On a par suite dans ce cas
$\mathbf{Tor}_\bullet^A(P_\bullet,Q_\bullet)
=\operatorname{Tor}_\bullet^A(F,G)$.
\end{namedtheorem}

\begin{proposition}[6.3.6]
Soient $(P_\bullet^\lambda)$, $(Q_\bullet^\mu)$ deux systèmes inductifs
filtrants de complexes de $A$-modules; on a un isomorphisme canonique
\[
  \varinjlim_{\lambda,\mu}
  \mathbf{Tor}_\bullet^A(P_\bullet^\lambda,Q_\bullet^\mu)
  \xrightarrow{\sim}
  \mathbf{Tor}_\bullet^A
  \left(\varinjlim_\lambda P_\bullet^\lambda,
        \varinjlim_\mu Q_\bullet^\mu\right).
  \tag{6.3.6.1}\label{III.6.3.6.1.fr}
\]
\end{proposition}

\begin{proof}
Posons $P_\bullet=\varinjlim_\lambda P_\bullet^\lambda$,
$Q_\bullet=\varinjlim_\mu Q_\bullet^\mu$; par fonctorialité, il est clair que
les $\mathbf{Tor}_\bullet^A(P_\bullet^\lambda,Q_\bullet^\mu)$ forment un
système inductif et que les applications
$\mathbf{Tor}_\bullet^A(P_\bullet^\lambda,Q_\bullet^\mu)\to
\mathbf{Tor}_\bullet^A(P_\bullet,Q_\bullet)$ déduites

\oldpage[III]{143}

des applications canoniques $P_\bullet^\lambda\to P_\bullet$,
$Q_\bullet^\mu\to Q_\bullet$, forment un système inductif d'homomorphismes,
d'où un homomorphisme canonique (6.3.6.1), et plus généralement un
homomorphisme canonique
\[
  \varinjlim_{\lambda,\mu}''E(P_\bullet^\lambda,Q_\bullet^\mu)
  \longrightarrow ''E(P_\bullet,Q_\bullet)
\]
dont (6.3.6.1) est l'homomorphisme des aboutissements. En outre, la suite
spectrale $''E(P_\bullet,Q_\bullet)$ est régulière (6.3.2), et il en est de
même de la suite spectrale
$\varinjlim_{\lambda,\mu}''E(P_\bullet^\lambda,Q_\bullet^\mu)$, comme il
résulte des définitions (0, 11.1.7) et de la démonstration de (0, 11.3.3);
pour démontrer que (6.3.6.1) est bijectif, il suffit donc (0, 11.1.5) de
prouver que l'homomorphisme
\[
  \varinjlim_{\lambda,\mu}''E(P_\bullet^\lambda,Q_\bullet^\mu)
  \longrightarrow ''E(P_\bullet,Q_\bullet)
  \tag{6.3.6.2}\label{III.6.3.6.2.fr}
\]
est bijectif pour les termes $E^2$. Comme le foncteur $H_\bullet$ commute à
la limite inductive des complexes de modules, on est finalement ramené à
prouver que pour deux systèmes inductifs filtrants $(F^\lambda)$,
$(G^\mu)$ de $A$-modules, l'homomorphisme canonique
\[
  \varinjlim_{\lambda,\mu}\operatorname{Tor}_\bullet^A(F^\lambda,G^\mu)
  \longrightarrow
  \operatorname{Tor}_\bullet^A
  \left(\varinjlim_\lambda F^\lambda,
        \varinjlim_\mu G^\mu\right)
\]
est bijectif. Pour cela, considérons pour chaque $F^\lambda$ la résolution
libre canonique
\[
  L_\bullet^\lambda:\quad
  \cdots\longrightarrow L_{i+1}^\lambda
  \longrightarrow L_i^\lambda\longrightarrow\cdots
  \longrightarrow L_1^\lambda\longrightarrow L_0^\lambda\longrightarrow0
\]
où $L_0^\lambda$ est le $A$-module des combinaisons linéaires formelles
d'éléments de $F^\lambda$ et $L_{i+1}^\lambda$ le $A$-module des combinaisons
linéaires formelles d'éléments de
$\operatorname{Ker}(L_i^\lambda\to L_{i-1}^\lambda)$; on vérifie
immédiatement que les $L_\bullet^\lambda$ forment un système inductif de
complexes, et si l'on pose
$F=\varinjlim_\lambda F^\lambda$, $L_i=\varinjlim_\lambda L_i^\lambda$, les
$L_i$ forment une résolution $L_\bullet$ de $F$, le foncteur
$\varinjlim$ étant exact; en outre, les $L_i$, limites inductives de
$A$-modules libres, sont plats (0\textsubscript{I}, 6.1.2). On considère de
même pour chaque $\mu$ la résolution libre canonique $M_\bullet^\mu$ de
$G^\mu$, et $M_\bullet=\varinjlim_\mu M_\bullet^\mu$ est une résolution plate
de $G=\varinjlim_\mu G^\mu$. On a alors
$\operatorname{Tor}_\bullet^A(\varinjlim_\lambda F^\lambda,
\varinjlim_\mu G^\mu)=H_\bullet(L_\bullet\otimes_A M_\bullet)$ en vertu de
(6.3.5) et (6.3.4); mais
\[
  H_\bullet(L_\bullet\otimes_A M_\bullet)
  =\varinjlim_{\lambda,\mu}
   H_\bullet(L_\bullet^\lambda\otimes_A M_\bullet^\mu)
\]
puisque $H_\bullet$ commute aux limites inductives de complexes de modules;
comme $H_\bullet(L_\bullet^\lambda\otimes_A M_\bullet^\mu)
=\operatorname{Tor}_\bullet^A(F^\lambda,G^\mu)$, cela termine la
démonstration.

Lorsqu'on suppose qu'il existe $i_0$ tel que
$P_i^\lambda=Q_i^\mu=0$ pour $i<i_0$ quels que soient $\lambda$ et $\mu$, on
démontre de la même manière que l'homomorphisme canonique
\[
  \varinjlim_{\lambda,\mu}'E(P_\bullet^\lambda,Q_\bullet^\mu)
  \longrightarrow'E(P_\bullet,Q_\bullet)
  \tag{6.3.6.3}\label{III.6.3.6.3.fr}
\]
est bijectif.
\end{proof}

\begin{proposition}[6.3.7]
Supposons $P_\bullet$ et $Q_\bullet$ limités inférieurement. Si le complexe
$P_\bullet$ est formé de $A$-modules plats, on a un $A$-isomorphisme
canonique de $\partial$-foncteurs en $Q_\bullet$
\[
  \mathbf{Tor}_\bullet^A(P_\bullet,Q_\bullet)
  \xrightarrow{\sim}H_\bullet(P_\bullet\otimes_A Q_\bullet).
  \tag{6.3.7.1}\label{III.6.3.7.1.fr}
\]
\end{proposition}

\begin{proof}
En effet, la suite spectrale (6.3.2.1) est birégulière et dégénérée, et
l'existence de l'isomorphisme (6.3.7.1) résulte de (0, 11.1.6). En outre, en
calculant l'hypertor à partir d'une résolution projective de Cartan-Eilenberg
de $P_\bullet$ (6.3.4), on voit aussitôt que l'isomorphisme ainsi défini est
un isomorphisme de $\partial$-foncteurs en $Q_\bullet$.
\end{proof}
\begin{env}[6.3.8]
\oldpage[III]{144}
Soit $\rho:A\to A'$ un homomorphisme d'anneaux. Nous nous proposons de
définir un $A$-homomorphisme factoriel de degré 0 canoniquement associé à
$\rho$ :
\[
  \rho_{P_\bullet,Q_\bullet}:
  \mathbf{Tor}_\bullet^A(P_\bullet,Q_\bullet)
  \longrightarrow
  \mathbf{Tor}_\bullet^{A'}
  (P_\bullet\otimes_A A',Q_\bullet\otimes_A A').
  \tag{6.3.8.1}\label{III.6.3.8.1.fr}
\]

Pour cela, considérons une résolution de Cartan-Eilenberg projective
$L_{\bullet\bullet}$ de $P_\bullet$; considérons d'autre part une résolution
de Cartan-Eilenberg projective $L'_{\bullet\bullet}$ de
$P_\bullet\otimes_A A'$. Nous allons voir qu'on peut définir un
$A'$-homomorphisme de complexes
$L_{\bullet\bullet}\otimes_A A'\to L'_{\bullet\bullet}$, déterminé à
homotopie près. En effet, la construction de $L_{\bullet\bullet}$ est
entièrement déterminée lorsqu'on se donne (arbitrairement) pour chaque $i$,
une résolution projective $(X_{ij}^{\mathrm B})_{j\geqslant0}$ de
$\mathrm B_i(P_\bullet)$ et une résolution projective
$(X_{ij}^{\mathrm H})_{j\geqslant0}$ de $\mathrm H_i(P_\bullet)$, qui sont
respectivement égales à $\mathrm B_i^{\mathrm I}(L_{\bullet\bullet})$ et
$\mathrm H_i^{\mathrm I}(L_{\bullet\bullet})$; on en déduit successivement
\[
  \mathrm Z_i^{\mathrm I}(L_{\bullet\bullet})
  =\mathrm H_i^{\mathrm I}(L_{\bullet\bullet})
   \oplus\mathrm B_i^{\mathrm I}(L_{\bullet\bullet}),
  \qquad
  L_{i,\bullet}=\mathrm Z_i^{\mathrm I}(L_{\bullet\bullet})
   \oplus\mathrm B_{i-1}^{\mathrm I}(L_{\bullet\bullet}).
\]
Cela étant, $X_{i,\bullet}^{\mathrm B}\otimes_A A'$ n'est plus en général
une résolution de $\mathrm B_i(P_\bullet)\otimes_A A'$, mais est encore un
complexe formé de $A'$-modules projectifs, et il y a donc un
$A'$-homomorphisme
\[
  X_{i,\bullet}^{\mathrm B}\otimes_A A'
  \longrightarrow \mathrm B_i^{\mathrm I}(L'_{\bullet\bullet})
\]
compatible avec les augmentations, et déterminé à homotopie près
(M, V, 1.1). On a de même un $A'$-homomorphisme
\[
  X_{i,\bullet}^{\mathrm H}\otimes_A A'
  \longrightarrow \mathrm H_i^{\mathrm I}(L'_{\bullet\bullet})
\]
déterminé à homotopie près, d'où l'on déduit, par la construction rappelée
plus haut, un $A'$-homomorphisme
$L_{i,\bullet}\otimes_A A'\to L'_{i,\bullet}$ pour tout $i$; ces
homomorphismes (pour $i\in\mathbf Z$) sont compatibles avec les opérateurs de
dérivation $L_{i,\bullet}\to L_{i-1,\bullet}$ et les analogues pour
$L'_{\bullet\bullet}$, en vertu de la même construction, et ils constituent
donc le $A'$-homomorphisme
$L_{\bullet\bullet}\otimes_A A'\to L'_{\bullet\bullet}$ cherché.

Pour définir (6.3.8.1), il suffit alors de considérer de même une résolution
de Cartan-Eilenberg projective $M_{\bullet\bullet}$ (resp.
$M'_{\bullet\bullet}$) de $Q_\bullet$ (resp.
$Q_\bullet\otimes_A A'$), et un $A'$-homomorphisme
$M_{\bullet\bullet}\otimes_A A'\to M'_{\bullet\bullet}$. On déduit de ces
homomorphismes un $A'$-homomorphisme
\[
  (L_{\bullet\bullet}\otimes_A A')\otimes_{A'}
  (M_{\bullet\bullet}\otimes_A A')
  \longrightarrow
  L'_{\bullet\bullet}\otimes_{A'}M'_{\bullet\bullet},
\]
puis par composition un $A$-homomorphisme de bicomplexes
\[
  L_{\bullet\bullet}\otimes_A M_{\bullet\bullet}
  \longrightarrow
  L'_{\bullet\bullet}\otimes_{A'}M'_{\bullet\bullet},
\]
et en passant à l'homologie on obtient (6.3.8.1), qui est bien défini
puisqu'il provient d'un morphisme de complexes défini à homotopie près.

Si $\rho':A'\to A''$ est un second homomorphisme d'anneaux, et
$\rho'':A\to A''$ l'homomorphisme composé $\rho'\circ\rho$, il est clair que
\[
  \rho''_{P_\bullet,Q_\bullet}
  =\rho'_{P'_\bullet,Q'_\bullet}\circ\rho_{P_\bullet,Q_\bullet},
  \qquad
  P'_\bullet=P_\bullet\otimes_A A',
  \quad Q'_\bullet=Q_\bullet\otimes_A A'.
\]
Notons encore que le morphisme de bicomplexes
$L_{\bullet\bullet}\otimes_A M_{\bullet\bullet}
\to L'_{\bullet\bullet}\otimes_{A'}M'_{\bullet\bullet}$ considéré
ci-dessus définit des morphismes factoriels (en $P_\bullet$ et $Q_\bullet$)
de suites spectrales
\[
  'E_{pq}^r(P_\bullet,Q_\bullet)
  \longrightarrow
  'E_{pq}^r(P_\bullet\otimes_A A',Q_\bullet\otimes_A A')
\]
et
\[
  ''E_{pq}^r(P_\bullet,Q_\bullet)
  \longrightarrow
  ''E_{pq}^r(P_\bullet\otimes_A A',Q_\bullet\otimes_A A'),
\]
indépendants des résolutions de Cartan-Eilenberg considérées, et ayant aussi
la propriété de transitivité précédente.
\end{env}

\begin{proposition}[6.3.9]
Soit $\rho:A\to A'$ un homomorphisme d'anneaux tel que $A'$ soit un
$A$-module plat. On a alors des isomorphismes canoniques factoriels
\[
  \mathbf{Tor}_\bullet^{A'}
  (P_\bullet\otimes_A A',Q_\bullet\otimes_A A')
  \xrightarrow{\sim}
  \mathbf{Tor}_\bullet^A(P_\bullet,Q_\bullet)\otimes_A A'.
  \tag{6.3.9.1}\label{III.6.3.9.1.fr}
\]
\[
  \left\{
  \begin{aligned}
  'E(P_\bullet\otimes_A A',Q_\bullet\otimes_A A')
    &\xrightarrow{\sim}'E(P_\bullet,Q_\bullet)\otimes_A A',\\
  ''E(P_\bullet\otimes_A A',Q_\bullet\otimes_A A')
    &\xrightarrow{\sim}''E(P_\bullet,Q_\bullet)\otimes_A A'.
  \end{aligned}
  \right.
  \tag{6.3.9.2}\label{III.6.3.9.2.fr}
\]
\end{proposition}

\begin{proof}
\oldpage[III]{145}
En effet, vu l'exactitude du foncteur $M\otimes_A A'$ en $M$,
$L_{\bullet\bullet}\otimes_A A'$ et $M_{\bullet\bullet}\otimes_A A'$ sont
alors des résolutions projectives de Cartan-Eilenberg de
$P_\bullet\otimes_A A'$ et $Q_\bullet\otimes_A A'$ respectivement, d'où la
conclusion.
\end{proof}

\begin{env}[6.3.10]
Soit $\rho:A\to A'$ un homomorphisme d'anneaux; pour tout complexe
$P'_\bullet$ de $A'$-modules, $P'_{\bullet[\rho]}$ est un complexe de
$A$-modules; en outre, l'application identique
$P'_{\bullet[\rho]}\to P'_\bullet$ peut être considérée comme composée des
applications canoniques
\[
  P'_{\bullet[\rho]}\longrightarrow
  P'_{\bullet[\rho]}\otimes_A A'
  \xrightarrow{\mu}P'_\bullet,
\]
où $\mu$ est le $A'$-homomorphisme $\mu(x\otimes a')=a'x$. Si $Q'_\bullet$
est un second complexe de $A'$-modules, on a donc des homomorphismes
canoniques factoriels de degré 0
\[
  \mathbf{Tor}_\bullet^A(P'_{\bullet[\rho]},Q'_{\bullet[\rho]})
  \longrightarrow
  \mathbf{Tor}_\bullet^{A'}
  (P'_{\bullet[\rho]}\otimes_A A',Q'_{\bullet[\rho]}\otimes_A A')
  \longrightarrow
  \mathbf{Tor}_\bullet^{A'}(P'_\bullet,Q'_\bullet),
  \tag{6.3.10.1}\label{III.6.3.10.1.fr}
\]
où la première flèche est le $A$-homomorphisme défini dans (6.3.8) et la
seconde se déduit des $A'$-homomorphismes
$P'_{\bullet[\rho]}\otimes_A A'\to P'_\bullet$ et
$Q'_{\bullet[\rho]}\otimes_A A'\to Q'_\bullet$ par fonctorialité. On a des
homomorphismes analogues pour les suites spectrales de (6.3.2), et des
propriétés évidentes de transitivité, que nous laissons au lecteur le soin
d'énoncer.
\end{env}

\begin{proposition}[6.3.11]
Soit $\rho:A\to A'$ un homomorphisme d'anneaux faisant de $A'$ un
$A$-module plat. Pour tout complexe $P'_\bullet$ de $A'$-modules et tout
complexe $Q_\bullet$ de $A$-modules limités inférieurement, on a un
isomorphisme canonique factoriel
\[
  \mathbf{Tor}_\bullet^A(P'_{\bullet[\rho]},Q_\bullet)
  \xrightarrow{\sim}
  \mathbf{Tor}_\bullet^{A'}(P'_\bullet,Q_\bullet\otimes_A A').
  \tag{6.3.11.1}\label{III.6.3.11.1.fr}
\]
\end{proposition}

\begin{proof}
En effet, si $M_{\bullet\bullet}$ est une résolution projective de
Cartan-Eilenberg de $Q_\bullet$, $M_{\bullet\bullet}\otimes_A A'$ est une
résolution projective de Cartan-Eilenberg de $Q_\bullet\otimes_A A'$, et
l'on a, à un isomorphisme canonique près,
\[
  P'_{\bullet[\rho]}\otimes_A M_{\bullet\bullet}
  =P'_\bullet\otimes_{A'}(M_{\bullet\bullet}\otimes_A A');
\]
la conclusion résulte de (6.3.4).
\end{proof}

\begin{namedtheorem}[6.3.12]{Remarque}
Soit $(A^\lambda)$ un système inductif filtrant d'anneaux, et soient
$(P_\bullet^\lambda)$, $(Q_\bullet^\lambda)$ deux systèmes inductifs de
complexes de $(A^\lambda)$-modules; on a alors un isomorphisme canonique
généralisant (6.3.6.1)
\[
  \varinjlim_\lambda
  \mathbf{Tor}_\bullet^{A^\lambda}
  (P_\bullet^\lambda,Q_\bullet^\lambda)
  \xrightarrow{\sim}
  \mathbf{Tor}_\bullet^A(P_\bullet,Q_\bullet),
  \tag{6.3.12.1}\label{III.6.3.12.1.fr}
\]
où $A=\varinjlim A^\lambda$, $P_\bullet=\varinjlim P_\bullet^\lambda$,
$Q_\bullet=\varinjlim Q_\bullet^\lambda$. Une fois définis les
homomorphismes
\[
  \mathbf{Tor}_\bullet^{A^\lambda}
  (P_\bullet^\lambda,Q_\bullet^\lambda)
  \longrightarrow
  \mathbf{Tor}_\bullet^{A^\mu}
  (P_\bullet^\mu,Q_\bullet^\mu)
\]
pour $\lambda\leqslant\mu$, à l'aide de (6.3.10), la démonstration est celle
de (6.3.6).
\end{namedtheorem}

\begin{proposition}[6.3.13]
Soient $S$ une partie multiplicative de $A$, $P_\bullet$ et $Q_\bullet$ deux
complexes de $A$-modules, dans lesquels les homothéties définies par les
éléments de $S$ soient bijectives, de sorte que, si $A'=S^{-1}A$,
$P_\bullet$ et $Q_\bullet$ sont formés de $A'$-modules. Alors on a un
isomorphisme canonique
\[
  \mathbf{Tor}_\bullet^A(P_\bullet,Q_\bullet)
  \xrightarrow{\sim}
  \mathbf{Tor}_\bullet^{A'}(P_\bullet,Q_\bullet).
\]
\end{proposition}

\begin{proof}
En effet, l'hypothèse entraîne que les homomorphismes canoniques
$P_\bullet\to P_\bullet\otimes_A A'$,
$Q_\bullet\to Q_\bullet\otimes_A A'$ sont bijectifs. D'autre part, la
fonctorialité de l'hypertor montre que tout $s\in S$ définit une homothétie
bijective dans $\mathbf{Tor}_\bullet^A(P_\bullet,Q_\bullet)$, et par suite
\[
  \mathbf{Tor}_\bullet^A(P_\bullet,Q_\bullet)
  \longrightarrow
  \mathbf{Tor}_\bullet^A(P_\bullet,Q_\bullet)\otimes_A A'
\]

\oldpage[III]{146}

est aussi un homomorphisme bijectif. Comme $A'$ est un $A$-module plat, la
conclusion résulte de (6.3.9), et on a de même des isomorphismes canoniques
pour les suites spectrales.
\end{proof}

\subsection{Foncteurs hypertor locaux de complexes de Modules quasi-cohérents : cas des schémas affines}

\begin{env}[6.4.1]
Soient $S$ un schéma affine d'anneau $A$, $X$, $Y$ deux $S$-schémas affines
d'anneaux $B$, $C$ respectivement, de sorte que $B$ et $C$ sont des algèbres
sur $A$. Tout complexe $\mathscr P_\bullet$ (resp. $\mathscr Q_\bullet$) de
$\mathscr O_X$-Modules (resp. $\mathscr O_Y$-Modules) quasi-cohérents est de
la forme $\widetilde{P_\bullet}$ (resp. $\widetilde{Q_\bullet}$), où
$P_\bullet$ (resp. $Q_\bullet$) est un complexe de $B$-modules (resp.
$C$-modules) (I, 1.3.7 et 1.3.8). On peut évidemment considérer $P_\bullet$
et $Q_\bullet$ comme des complexes de $A$-modules et former les
$\mathbf{Tor}_n^A(P_\bullet,Q_\bullet)$; en outre, en vertu du caractère
bifonctoriel de $\mathbf{Tor}_n^A(P_\bullet,Q_\bullet)$, les $A$-algèbres
$B$ et $C$ opèrent dans ce $A$-module, et ces opérations en font un
$(B,C)$-bimodule, ou, ce qui revient au même, un module sur
$B\otimes_A C=A(X\times_S Y)$. On a donc défini de la sorte un
$\mathscr O_{X\times_S Y}$-Module quasi-cohérent
\[
  \mathscr{T}\!or_n^{\mathscr O_S}(\mathscr P_\bullet,\mathscr Q_\bullet)
  =\bigl(\mathbf{Tor}_n^A(P_\bullet,Q_\bullet)\bigr)^{\sim}
  \tag{6.4.1.1}\label{III.6.4.1.1.fr}
\]
que l'on appelle l'\emph{hypertor local d'indice $n$} des complexes
$\mathscr P_\bullet$ et $\mathscr Q_\bullet$ et que l'on note aussi
$\mathscr{T}\!or_n^S(\mathscr P_\bullet,\mathscr Q_\bullet)$.
\end{env}

\begin{namedtheorem}[6.4.2]{Lemme}
Avec les notations de (6.4.1), supposons que l'anneau $A$ soit de la forme
$R^{-1}A'$, où $A'$ est un anneau et $R$ une partie multiplicative de $A'$.
Soit $S'=\operatorname{Spec}(A')$, de sorte que $X$ et $Y$ peuvent être
considérés comme des $S'$-préschémas et que l'on a
$X\times_{S'}Y=X\times_S Y$ (I, 1.6.2 et 3.2.4). On a alors
\[
  \mathscr{T}\!or_\bullet^{S'}(\mathscr P_\bullet,\mathscr Q_\bullet)
  =\mathscr{T}\!or_\bullet^S(\mathscr P_\bullet,\mathscr Q_\bullet).
\]
\end{namedtheorem}

\begin{proof}
Cela résulte de la formule (6.4.1.1) et de (6.3.13).
\end{proof}

\begin{env}[6.4.3]
Avec les notations et hypothèses de (6.4.1), soient
$\mathscr F=\widetilde F$ un $\mathscr O_X$-Module quasi-cohérent,
$\mathscr G=\widetilde G$ un $\mathscr O_Y$-Module quasi-cohérent;
considérant $\mathscr F$ et $\mathscr G$ comme des complexes de Modules, on
notera $\mathscr{T}\!or_n^{\mathscr O_S}(\mathscr F,\mathscr G)$ ou
$\mathscr{T}\!or_n^S(\mathscr F,\mathscr G)$ leur hypertor d'indice $n$; il
résulte de (6.3.5 (ii)) que l'on a
\[
  \mathscr{T}\!or_n^{\mathscr O_S}(\mathscr F,\mathscr G)
  =\bigl(\operatorname{Tor}_n^A(F,G)\bigr)^{\sim}.
  \tag{6.4.3.1}\label{III.6.4.3.1.fr}
\]

Revenons alors au cas général de deux complexes de Modules quasi-cohérents
$\mathscr P_\bullet$, $\mathscr Q_\bullet$. Les formules (6.4.1.1) et
(6.4.3.1) montrent, compte tenu de la prop. (6.3.2), que
$\mathscr{T}\!or_\bullet^S(\mathscr P_\bullet,\mathscr Q_\bullet)$ est
l'aboutissement de deux suites spectrales
$'\mathscr E(\mathscr P_\bullet,\mathscr Q_\bullet)$,
$''\mathscr E(\mathscr P_\bullet,\mathscr Q_\bullet)$, dont les termes
$E_2$ sont donnés par
\[
  '\mathscr E_{pq}^2
  =\mathscr H_p\bigl(\mathscr{T}\!or_q^S
    (\mathscr P_\bullet,\mathscr Q_\bullet)\bigr)
  \tag{6.4.3.2}\label{III.6.4.3.2.fr}
\]
\[
  ''\mathscr E_{pq}^2
  =\bigoplus_{q'+q''=q}
  \mathscr{T}\!or_p^S
  \bigl(\mathscr H_{q'}(\mathscr P_\bullet),
        \mathscr H_{q''}(\mathscr Q_\bullet)\bigr)
  \tag{6.4.3.3}\label{III.6.4.3.3.fr}
\]
où $\mathscr{T}\!or_q^S(\mathscr P_\bullet,\mathscr Q_\bullet)$ est le
bicomplexe de $\mathscr O_{X\times_S Y}$-Modules quasi-cohérents
$\mathscr{T}\!or_q^S(\mathscr P_i,\mathscr Q_j)$.
\end{env}

\begin{env}[6.4.4]
Considérons maintenant deux autres schémas affines
$X^{(1)}=\operatorname{Spec}(B^{(1)})$,
$Y^{(1)}=\operatorname{Spec}(C^{(1)})$, où $B^{(1)}$ et $C^{(1)}$ sont des
$A$-algèbres, et supposons donnés deux

\oldpage[III]{147}
S-morphismes $u:X^{(1)}\to X$, $v:Y^{(1)}\to Y$, correspondant à des
$A$-homomorphismes $\varphi:B\to B^{(1)}$, $\psi:C\to C^{(1)}$.
Considérons les complexes
$u^*(\mathscr P_\bullet)=(P_\bullet\otimes_B B^{(1)})^{\sim}$ de
$\mathscr O_{X^{(1)}}$-Modules,
$v^*(\mathscr Q_\bullet)=(Q_\bullet\otimes_C C^{(1)})^{\sim}$ de
$\mathscr O_{Y^{(1)}}$-Modules. Les $A$-homomorphismes canoniques
\[
  P_\bullet\longrightarrow P_\bullet\otimes_B B^{(1)},
  \qquad
  Q_\bullet\longrightarrow Q_\bullet\otimes_C C^{(1)}
\]
donnent par fonctorialité un $A$-homomorphisme
\[
  \mathbf{Tor}_\bullet^A(P_\bullet,Q_\bullet)
  \longrightarrow
  \mathbf{Tor}_\bullet^A
  (P_\bullet\otimes_B B^{(1)},Q_\bullet\otimes_C C^{(1)});
\]
en outre, toujours par fonctorialité, cet homomorphisme est en fait un
homomorphisme de $(B\otimes_A C)$-modules. D'où l'on conclut que l'on a défini
ainsi un $(u\times_S v)$-morphisme
\[
  \theta:\mathscr{T}\!or_\bullet^S(\mathscr P_\bullet,\mathscr Q_\bullet)
  \longrightarrow
  \mathscr{T}\!or_\bullet^S
  (u^*(\mathscr P_\bullet),v^*(\mathscr Q_\bullet))
  \tag{6.4.4.1}\label{III.6.4.4.1.fr}
\]
et par suite, un homomorphisme de
$\mathscr O_{X^{(1)}\times_S Y^{(1)}}$-Modules
\[
  \theta^\sharp:(u\times_S v)^*
  \bigl(\mathscr{T}\!or_\bullet^S(\mathscr P_\bullet,\mathscr Q_\bullet)\bigr)
  \longrightarrow
  \mathscr{T}\!or_\bullet^S
  (u^*(\mathscr P_\bullet),v^*(\mathscr Q_\bullet))
  \tag{6.4.4.2}\label{III.6.4.4.2.fr}
\]
qui est évidemment un morphisme de bi-$\partial$-foncteurs dans les catégories
de Modules quasi-cohérents limités inférieurement.

L'homomorphisme (6.4.4.2) n'est pas nécessairement bijectif; toutefois :
\end{env}

\begin{namedtheorem}[6.4.5]{Lemme}
Avec les notations de (6.4.4), supposons que $u$ et $v$ soient des immersions
ouvertes; alors l'homomorphisme (6.4.4.2) est bijectif.
\end{namedtheorem}

\begin{proof}
Identifions $X^{(1)}$ (resp. $Y^{(1)}$) à un ouvert de $X$ (resp. $Y$);
$X^{(1)}$ (resp. $Y^{(1)}$) est alors réunion d'ouverts de la forme
$\mathrm D(f)$ (resp. $\mathrm D(g)$), où $f\in B$ (resp. $g\in C$), et les
préschémas induits $\mathrm D(\varphi(f))$ et $\mathrm D(f)$ (resp.
$\mathrm D(\psi(g))$ et $\mathrm D(g)$) sont isomorphes. Il suffira de prouver
le lemme lorsque $X^{(1)}$ (resp. $Y^{(1)}$) est de la forme $\mathrm D(f)$
(resp. $\mathrm D(g)$); en effet, si ce point est établi, et si l'on revient au
cas général, il suffira de prouver que la restriction de $\theta^\sharp$ à
chaque ouvert $\mathrm D(f)\times_S\mathrm D(g)$ est un isomorphisme; or, si
$u_1:\mathrm D(f)\to X^{(1)}$, $v_1:\mathrm D(g)\to Y^{(1)}$ sont les
injections canoniques, la restriction précédente n'est autre que
$(u_1\times_S v_1)^*(\theta^\sharp)$; mais il est immédiat, en vertu des
définitions (6.4.4) et de ($0_{\mathrm I}$, 4.4.8), qu'en la composant avec
l'homomorphisme canonique
\[
  (u_1\times_S v_1)^*
  \bigl(\mathscr{T}\!or_\bullet^S
    (u^*(\mathscr P_\bullet),v^*(\mathscr Q_\bullet))\bigr)
  \longrightarrow
  \mathscr{T}\!or_\bullet^S
  ((u')^*(\mathscr P_\bullet),(v')^*(\mathscr Q_\bullet))
  \tag{6.4.5.1}\label{III.6.4.5.1.fr}
\]
où $u'=u\circ u_1$ et $v'=v\circ v_1$, on obtient l'homomorphisme canonique
\[
  (u'\times_S v')^*
  \bigl(\mathscr{T}\!or_\bullet^S(\mathscr P_\bullet,\mathscr Q_\bullet)\bigr)
  \longrightarrow
  \mathscr{T}\!or_\bullet^S
  ((u')^*(\mathscr P_\bullet),(v')^*(\mathscr Q_\bullet))
  \tag{6.4.5.2}\label{III.6.4.5.2.fr}
\]
et si l'on sait que (6.4.5.1) et (6.4.5.2) sont des isomorphismes, il en
résultera qu'il en est de même de $(u_1\times_S v_1)^*(\theta^\sharp)$.

Supposons donc que $X^{(1)}=\mathrm D(f)$ et $Y^{(1)}=\mathrm D(g)$, de sorte
que $B^{(1)}=B_f$ et $C^{(1)}=C_g$; $u^*(\mathscr P_\bullet)$ (resp.
$v^*(\mathscr Q_\bullet)$) s'identifie alors à
$((P_\bullet)_f)^{\sim}$ (resp. $((Q_\bullet)_g)^{\sim}$); d'autre part,
$X^{(1)}\times_S Y^{(1)}$ s'identifie au sous-schéma ouvert
$\mathrm D(f\otimes g)$ de
$X\times_S Y=\operatorname{Spec}(B\otimes_A C)$ (II, 4.3.2.4); il s'agit de
prouver que l'homomorphisme
\[
  \bigl(\mathbf{Tor}_\bullet^A(P_\bullet,Q_\bullet)\bigr)_{f\otimes g}
  \longrightarrow
  \mathbf{Tor}_\bullet^A((P_\bullet)_f,(Q_\bullet)_g)
  \tag{6.4.5.3}\label{III.6.4.5.3.fr}
\]

\oldpage[III]{148}
déduit par fonctorialité des homomorphismes canoniques
$P_\bullet\to(P_\bullet)_f$, $Q_\bullet\to(Q_\bullet)_g$, est bijectif. Or
($0_{\mathrm I}$, 1.6.1), on peut écrire
$(P_\bullet)_f=\varinjlim P_\bullet^{(n)}$, où les $P_\bullet^{(n)}$ sont tous
des complexes de $B$-modules identiques à $P_\bullet$, l'application
$P_\bullet^{(m)}\to P_\bullet^{(n)}$ pour $m\leqslant n$ étant la
multiplication par $f^{n-m}$; on a un résultat analogue pour $Q_\bullet$ en
remplaçant $f$ par $g$; d'autre part, il est clair que l'homomorphisme
\[
  \mathbf{Tor}_\bullet^A(P_\bullet^{(m)},Q_\bullet^{(m)})
  \longrightarrow
  \mathbf{Tor}_\bullet^A(P_\bullet^{(n)},Q_\bullet^{(n)})
\]
correspondant aux homomorphismes
$P_\bullet^{(m)}\to P_\bullet^{(n)}$ et
$Q_\bullet^{(m)}\to Q_\bullet^{(n)}$ est par définition la multiplication
par $(f\otimes g)^{n-m}$. La conclusion résulte alors de ($0_{\mathrm I}$,
1.6.1) appliqué au premier membre de (6.4.5.3) et de (6.3.6).
\end{proof}

\begin{env}[6.4.6]
Avec les notations de (6.4.4), on définit de même des homomorphismes canoniques
de foncteurs spectraux
\[
  \left\{
  \begin{aligned}
  (u\times_S v)^*\bigl('\mathscr E(\mathscr P_\bullet,\mathscr Q_\bullet)\bigr)
    &\longrightarrow
      '\mathscr E(u^*(\mathscr P_\bullet),v^*(\mathscr Q_\bullet)),\\
  (u\times_S v)^*\bigl(''\mathscr E(\mathscr P_\bullet,\mathscr Q_\bullet)\bigr)
    &\longrightarrow
      ''\mathscr E(u^*(\mathscr P_\bullet),v^*(\mathscr Q_\bullet)).
  \end{aligned}
  \right.
  \tag{6.4.6.1}\label{III.6.4.6.1.fr}
\]
et le raisonnement de (6.4.5) montre que lorsque $u$ et $v$ sont des
immersions ouvertes, les homomorphismes (6.4.6.1) sont bijectifs : en effet,
compte tenu de (6.3.6.2) et (6.3.6.3), il prouve que c'est un isomorphisme pour
les termes $E^2$, et (6.4.5) montre que c'est un isomorphisme pour les
aboutissements; on conclut donc à l'aide de (0, 11.1.2) et (0, 11.2.4).
\end{env}

\subsection{Foncteurs hypertor locaux de complexes de Modules quasi-cohérents : cas général}

\begin{env}[6.5.1]
Considérons maintenant un préschéma $S$ quelconque et deux $S$-préschémas
quelconques $X$, $Y$; soit $\mathscr P_\bullet$ (resp. $\mathscr Q_\bullet$)
un complexe de $\mathscr O_X$-Modules (resp. $\mathscr O_Y$-Modules)
quasi-cohérents. Posons $Z=X\times_S Y$; nous allons définir des
$\mathscr O_Z$-Modules quasi-cohérents
$\mathscr{T}\!or_n^S(\mathscr P_\bullet,\mathscr Q_\bullet)$ dits
\emph{hypertor locaux} de $\mathscr P_\bullet$ et $\mathscr Q_\bullet$ qui se
réduiront à ceux déjà définis dans (6.4) lorsque $S$, $X$ et $Y$ sont affines.

Lorsque $\mathscr P_\bullet$ et $\mathscr Q_\bullet$ se réduisent
respectivement à leurs termes de degré 0, $\mathscr F$ et $\mathscr G$ (les
autres étant nuls), on écrira
$\mathscr{T}\!or_n^S(\mathscr F,\mathscr G)$ au lieu de
$\mathscr{T}\!or_n^S(\mathscr P_\bullet,\mathscr Q_\bullet)$.
\end{env}

\begin{env}[6.5.2]
Supposons d'abord $S$ affine, et soient $(X_\lambda)$, $(Y_\mu)$ des
recouvrements de $X$ et $Y$ respectivement par des ouverts affines; alors les
$Z_{\lambda\mu}=X_\lambda\times_S Y_\mu$ forment un recouvrement ouvert affine
de $Z$. Posons
$\mathscr P_{\lambda\bullet}=\mathscr P_\bullet|X_\lambda$,
$\mathscr Q_{\mu\bullet}=\mathscr Q_\bullet|Y_\mu$; nous avons donc pour tout
couple $(\lambda,\mu)$ un $\mathscr O_{Z_{\lambda\mu}}$-Module quasi-cohérent
\[
  \mathscr F_{\lambda\mu}
  =\mathscr{T}\!or_n^S
   (\mathscr P_{\lambda\bullet},\mathscr Q_{\mu\bullet}),
\]
et il faut montrer que les $\mathscr F_{\lambda\mu}$ vérifient la condition de
recollement ($0_{\mathrm I}$, 3.3.1). Pour cela, il suffit de vérifier que pour
tout ouvert affine $U\subset X_\lambda\cap X_{\lambda'}$ (resp.
$V\subset Y_\mu\cap Y_{\mu'}$), les restrictions de
$\mathscr F_{\lambda\mu}$ et de $\mathscr F_{\lambda'\mu'}$ à $U\times_S V$
sont canoniquement isomorphes; mais cela découle aussitôt de l'existence
d'isomorphismes canoniques de ces restrictions sur
$\mathscr{T}\!or_n^S(\mathscr P_\bullet|U,\mathscr Q_\bullet|V)$ (6.4.5). En
outre, il résulte aussitôt de cette définition et de (6.4.5) que le
$\mathscr O_Z$-Module ainsi défini ne dépend pas (à un isomorphisme près) des
recouvrements ouverts consi\oldpage[III]{149}dérés;
nous le noterons donc
$\mathscr{T}\!or_n^S(\mathscr P_\bullet,\mathscr Q_\bullet)$; il résulte enfin
de (6.4.5) que pour toute partie ouverte $U$ (resp. $V$) de $X$ (resp. $Y$), la
restriction de $\mathscr{T}\!or_n^S(\mathscr P_\bullet,\mathscr Q_\bullet)$ à
$U\times_S V$ est canoniquement isomorphe à
$\mathscr{T}\!or_n^S(\mathscr P_\bullet|U,\mathscr Q_\bullet|V)$.
\end{env}

\begin{env}[6.5.3]
Passons maintenant au cas général où $S$ est quelconque, et soit $(S_\alpha)$
un recouvrement de $S$ formé d'ouverts affines; désignons par $X_\alpha$
(resp. $Y_\alpha$) l'image réciproque de $S_\alpha$ dans $X$ (resp. $Y$); il
faut encore prouver que les faisceaux
\[
  \mathscr{T}\!or_n^{S_\alpha}
  (\mathscr P_\bullet|X_\alpha,\mathscr Q_\bullet|Y_\alpha)
  =\mathscr G_\alpha
\]
vérifient la condition de recollement. Il suffit de définir, pour tout ouvert
affine $T$ contenu dans $S_\alpha\cap S_\beta$, des isomorphismes canoniques des
restrictions de $\mathscr G_\alpha$ et $\mathscr G_\beta$ à $U\times_S V$ (en
désignant par $U$ et $V$ les images réciproques de $T$ dans $X$ et $Y$
respectivement) sur
$\mathscr{T}\!or_n^T(\mathscr P_\bullet|U,\mathscr Q_\bullet|V)$; on peut, en
outre, se borner au cas où $T$ s'écrit à la fois $\mathrm D(f_\alpha)$ et
$\mathrm D(f_\beta)$, $f_\alpha$ (resp. $f_\beta$) étant une section de
$\mathscr O_S$ au-dessus de $S_\alpha$ (resp. $S_\beta$); mais alors
$\mathscr{T}\!or_n^T(\mathscr P_\bullet|U,\mathscr Q_\bullet|V)$ est
canoniquement isomorphe à
$\mathscr{T}\!or_n^{S_\alpha}(\mathscr P_\bullet|U,\mathscr Q_\bullet|V)$
d'une part, à
$\mathscr{T}\!or_n^{S_\beta}(\mathscr P_\bullet|U,\mathscr Q_\bullet|V)$
d'autre part, en vertu de (6.4.2); comme on vient de définir des isomorphismes
canoniques de $\mathscr G_\alpha$ sur
$\mathscr{T}\!or_n^{S_\alpha}(\mathscr P_\bullet|U,\mathscr Q_\bullet|V)$ et
de $\mathscr G_\beta$ sur
$\mathscr{T}\!or_n^{S_\beta}(\mathscr P_\bullet|U,\mathscr Q_\bullet|V)$
(6.5.2), cela achève de définir le $\mathscr O_Z$-Module
$\mathscr{T}\!or_n^S(\mathscr P_\bullet,\mathscr Q_\bullet)$. En outre, pour
toute partie ouverte $U$ (resp. $V$) de $X$ (resp. $Y$),
$\mathscr{T}\!or_n^S(\mathscr P_\bullet|U,\mathscr Q_\bullet|V)$ est
canoniquement isomorphe à la restriction de
$\mathscr{T}\!or_n^S(\mathscr P_\bullet,\mathscr Q_\bullet)$ à $U\times_S V$.

Il est immédiat que l'on a ainsi défini (dans les catégories de complexes de
Modules quasi-cohérents limités inférieurement) un bi-$\partial$-foncteur
$\mathscr{T}\!or_\bullet^S(\mathscr P_\bullet,\mathscr Q_\bullet)$ à valeurs
dans la catégorie des $\mathscr O_Z$-Modules, car il est clair que la question
est locale sur $X$, $Y$ et $S$, en vertu de (6.4.5) et de la remarque que
(6.4.4.2) est un morphisme de bi-$\partial$-foncteurs. On notera que si
$\mathscr P_\bullet$ et $\mathscr Q_\bullet$ sont réduits respectivement à
leurs termes de degré 0, $\mathscr F$ et $\mathscr G$,
$\mathscr{T}\!or_0^S(\mathscr F,\mathscr G)$ n'est autre, en vertu de
(6.4.1.1), que le produit tensoriel externe $\mathscr F\otimes_S\mathscr G$
défini dans (I, 9.1.2); cela résulte en effet de (I, 9.1.3).
\end{env}

\begin{env}[6.5.4]
Il résulte de la construction précédente et des remarques faites dans (6.4.6)
que $\mathscr{T}\!or_\bullet^S(\mathscr P_\bullet,\mathscr Q_\bullet)$ est
l'aboutissement de deux foncteurs spectraux,
$'\mathscr E(\mathscr P_\bullet,\mathscr Q_\bullet)$,
$''\mathscr E(\mathscr P_\bullet,\mathscr Q_\bullet)$, de termes $E_2$ égaux à
\[
  '\mathscr E_{pq}^2
  =\mathscr H_p\bigl(\mathscr{T}\!or_q^S
    (\mathscr P_\bullet,\mathscr Q_\bullet)\bigr)
  \tag{6.5.4.1}\label{III.6.5.4.1.fr}
\]
\[
  ''\mathscr E_{pq}^2
  =\bigoplus_{q'+q''=q}
   \mathscr{T}\!or_p^S
   \bigl(\mathscr H_{q'}(\mathscr P_\bullet),
         \mathscr H_{q''}(\mathscr Q_\bullet)\bigr).
  \tag{6.5.4.2}\label{III.6.5.4.2.fr}
\]
La suite spectrale (6.5.4.2) est toujours régulière; les deux suites spectrales
sont birégulières si $\mathscr P_\bullet$ et $\mathscr Q_\bullet$ sont limités
inférieurement. Un autre cas où les deux suites précédentes sont birégulières
est le suivant :
\end{env}

\begin{env}[6.5.5]
Nous dirons que sur un espace topologique $T$ un faisceau d'anneaux
$\mathscr A$ est \emph{de dimension cohomologique $\leqslant n$} si, pour tout
$t\in T$ l'anneau $\mathscr A_t$ est de dimension cohomologique $\leqslant n$;
on dira alors aussi que l'\emph{espace annelé} $(T,\mathscr A)$ est
\emph{de dimension cohomologique $\leqslant n$}. On dira qu'un faisceau
d'anneaux (resp. un espace annelé) est de dimension cohomologique \emph{finie}
s'il existe un entier $n$ tel qu'il soit de dimension cohomologique
$\leqslant n$. On notera que si les $\mathscr A_t$ sont des anneaux locaux
(commutatifs) noethériens,

\oldpage[III]{150}
dire qu'ils sont de dimension cohomologique $\leqslant n$ signifie qu'ils sont
réguliers et de dimension (de Krull) $\leqslant n$ ($0_{\mathrm{IV}}$, 17.3.1).
Avec la terminologie de la théorie de la dimension que nous introduirons au
chap. IV, il revient au même de dire qu'un préschéma localement noethérien $T$
est de dimension cohomologique $\leqslant n$, ou de dire qu'il est régulier
($0_{\mathrm I}$, 4.1.4) et de dimension $\leqslant n$; cela signifie que pour
tout ouvert affine $U$ de $T$, l'anneau $\Gamma(U,\mathscr O_T)$ est de
dimension cohomologique $\leqslant n$ ($0_{\mathrm{IV}}$, 17.2.6). Cela étant,
cette dernière remarque, jointe à (6.3.2), prouve que si $S$ est localement
noethérien et de dimension cohomologique finie, les suites spectrales
$'\mathscr E(\mathscr P_\bullet,\mathscr Q_\bullet)$ et
$''\mathscr E(\mathscr P_\bullet,\mathscr Q_\bullet)$ sont birégulières.

Il est clair que
$\mathscr{T}\!or_\bullet^S(\mathscr P_\bullet,\mathscr Q_\bullet)$ se
transforme en
$\mathscr{T}\!or_\bullet^S(\mathscr Q_\bullet,\mathscr P_\bullet)$ (à un
isomorphisme près) par l'isomorphisme canonique de $X\times_S Y$ sur
$Y\times_S X$.
\end{env}

\begin{proposition}[6.5.6]
Soit $(\mathscr P_{\alpha\bullet})$ un système inductif filtrant de complexes
de $\mathscr O_X$-Modules quasi-cohérents; il existe alors un isomorphisme
canonique
\[
  \varinjlim_\alpha
  \mathscr{T}\!or_\bullet^S
  (\mathscr P_{\alpha\bullet},\mathscr Q_\bullet)
  \xrightarrow{\sim}
  \mathscr{T}\!or_\bullet^S
  \left(\varinjlim_\alpha\mathscr P_{\alpha\bullet},\mathscr Q_\bullet\right).
  \tag{6.5.6.1}\label{III.6.5.6.1.fr}
\]
\end{proposition}

\begin{proof}
La question étant locale sur $S$, $X$ et $Y$, on peut supposer $S$, $X$, $Y$
affines et la proposition se réduit alors à (6.3.6).
\end{proof}

\begin{namedtheorem}[6.5.7]{Remarques}
\textnormal{(i)} Considérons en particulier le cas où $S=X=Y$,
$\mathscr P_\bullet$ et $\mathscr Q_\bullet$ étant donc deux complexes de
$\mathscr O_S$-Modules quasi-cohérents; alors les
$\mathscr{T}\!or_n^S(\mathscr P_\bullet,\mathscr Q_\bullet)$ sont des
$\mathscr O_S$-Modules quasi-cohérents; en outre, pour tout point $z\in S$, il
résulte de (6.5.6) que l'on a un isomorphisme canonique
\[
  \bigl(\mathscr{T}\!or_n^S
    (\mathscr P_\bullet,\mathscr Q_\bullet)\bigr)_z
  \xrightarrow{\sim}
  \mathbf{Tor}_n^{\mathscr O_z}
  \bigl((\mathscr P_\bullet)_z,(\mathscr Q_\bullet)_z\bigr)
  \tag{6.5.7.1}\label{III.6.5.7.1.fr}
\]
car la question est locale et on est ramené au cas des modules, en vertu de
(6.4.1.1).

\textnormal{(ii)} On peut généraliser la définition des hypertor au cas de
deux complexes de $\mathscr O_X$-Modules $\mathscr P_\bullet$,
$\mathscr Q_\bullet$ sur un même espace annelé $(X,\mathscr O_X)$; pour tout
ouvert $U$ de $X$, posons en effet
$A(U)=\Gamma(U,\mathscr O_X)$,
$P_\bullet(U)=\Gamma(U,\mathscr P_\bullet)$,
$Q_\bullet(U)=\Gamma(U,\mathscr Q_\bullet)$; les $A(U)$-modules
\[
  \mathbf{Tor}_n^{A(U)}(P_\bullet(U),Q_\bullet(U))
\]
forment alors un préfaisceau sur $X$, et l'on désigne par
$\mathscr{T}\!or_n^X(\mathscr P_\bullet,\mathscr Q_\bullet)$ le
$\mathscr O_X$-Module associé à ce préfaisceau. Lorsque $X$ est un préschéma,
il résulte de (6.3.12) que ce $\mathscr O_X$-Module est canoniquement isomorphe
à l'hypertor défini ci-dessus. Nous ne développerons pas davantage cette
généralisation.
\end{namedtheorem}

\begin{proposition}[6.5.8]
Soient $X$, $Y$ deux $S$-préschémas, $\mathscr F$ (resp. $\mathscr G$) un
$\mathscr O_X$-Module (resp. un $\mathscr O_Y$-Module) quasi-cohérent. Si
$\mathscr F$ ou $\mathscr G$ est $S$-plat, on a
$\mathscr{T}\!or_n^S(\mathscr F,\mathscr G)=0$ pour $n\neq0$.
\end{proposition}

\begin{proof}
La question étant locale sur $X$ et $Y$, on peut supposer $X$, $Y$ et $S$
affines, d'anneaux respectifs $B$, $C$, $A$, et
$\mathscr F=\widetilde M$, $\mathscr G=\widetilde N$, $M$ (resp. $N$) étant
un $B$-module (resp. un $C$-module). Supposons par exemple que $\mathscr F$
soit $S$-plat, ce qui signifie que pour tout $s\in S$, $M_s$ est un
$A_s$-module plat ($0_{\mathrm I}$, 6.7.1); par suite $M$ est un $A$-module plat
($0_{\mathrm I}$, 6.3.3), et l'on sait que
$\operatorname{Tor}_n^A(M,N)=0$ pour $n>0$ et pour tout $C$-module $N$
($0_{\mathrm I}$, 6.1.1), d'où la conclusion par (6.4.1.1).
\end{proof}

\begin{namedtheorem}[6.5.9]{Corollaire}
Soient $X$, $Y$ deux $S$-préschémas, $\mathscr P_\bullet$ (resp.
$\mathscr Q_\bullet$) un complexe de

\oldpage[III]{151}
$\mathscr O_X$-Modules (resp. de $\mathscr O_Y$-Modules) quasi-cohérents limité
inférieurement. Supposons que tous les $\mathscr P_i$ soient $S$-plats. Alors
il existe un isomorphisme canonique de $\partial$-foncteurs en
$\mathscr Q_\bullet$
\[
  \mathscr{T}\!or_\bullet^S(\mathscr P_\bullet,\mathscr Q_\bullet)
  \xrightarrow{\sim}
  \mathscr H_\bullet(\mathscr P_\bullet\otimes_S\mathscr Q_\bullet).
  \tag{6.5.9.1}\label{III.6.5.9.1.fr}
\]
\end{namedtheorem}

\begin{proof}
Ce n'est autre que (6.3.7) lorsque $S$, $X$, $Y$ sont affines; on passe de là
au cas général par les raisonnements de (6.5.2) et (6.5.3).
\end{proof}

\begin{namedtheorem}[6.5.10]{Corollaire}
Supposons que $X$ soit plat sur $S$ ($0_{\mathrm I}$, 6.7.1), que
$\mathscr P_\bullet$ et $\mathscr Q_\bullet$ soient limités inférieurement et
que tous les $\mathscr P_i$ soient des $\mathscr O_X$-Modules localement
libres (non nécessairement de type fini). Alors l'homomorphisme (6.5.9.1) est
bijectif.
\end{namedtheorem}

\begin{proof}
En effet, l'hypothèse de (6.5.9) est remplie, la platitude étant une propriété
ponctuelle sur $X$ par définition et toute somme directe de modules plats étant
un module plat ($0_{\mathrm I}$, 6.1.2).
\end{proof}

\begin{proposition}[6.5.11]
Soient $X'$, $Y'$ deux $S$-préschémas, $f:X\to X'$, $g:Y\to Y'$ deux
$S$-morphismes affines. Soit $\mathscr P_\bullet$ (resp.
$\mathscr Q_\bullet$) un complexe de $\mathscr O_X$-Modules (resp. de
$\mathscr O_Y$-Modules) quasi-cohérents; on a alors un isomorphisme canonique
fonctoriel
\[
  (f\times_S g)_*
  \bigl(\mathscr{T}\!or_\bullet^S
    (\mathscr P_\bullet,\mathscr Q_\bullet)\bigr)
  \xrightarrow{\sim}
  \mathscr{T}\!or_\bullet^S
  (f_*(\mathscr P_\bullet),g_*(\mathscr Q_\bullet)).
  \tag{6.5.11.1}\label{III.6.5.11.1.fr}
\]
\end{proposition}

\begin{proof}
Comme $f$ et $g$ sont affines, $f_*(\mathscr P_\bullet)$ et
$g_*(\mathscr Q_\bullet)$ sont des complexes de Modules quasi-cohérents
(II, 1.2.6), et si l'on pose $Z'=X'\times_S Y'$, les deux membres de (6.5.11.1)
sont des $\mathscr O_{Z'}$-Modules quasi-cohérents (6.5.1); on se ramène
aisément au cas où $S$, $X'$ et $Y'$ sont affines; mais alors il en est de même
par hypothèse de $X$ et $Y$ et la vérification résulte aussitôt de (6.4.1.1) et
(I, 1.6.3).
\end{proof}

\begin{namedtheorem}[6.5.12]{Remarque}
Soient $X'$, $Y'$ deux $S$-préschémas et supposons que $X'$ soit $S$-plat;
soient $X$ un sous-préschéma fermé de $X'$, $i:X\to X'$ l'injection canonique,
$\mathscr P_\bullet$ (resp. $\mathscr Q_\bullet$) un complexe de
$\mathscr O_X$-Modules (resp. de $\mathscr O_{Y'}$-Modules) quasi-cohérents,
limité inférieurement. Soit enfin $\mathscr L'_{\bullet\bullet}$ une résolution
de $i_*(\mathscr P_\bullet)$ formé de $\mathscr O_{X'}$-Modules localement
libres, telle que tout point de $X'$ ait un voisinage ouvert affine $U$ pour
lequel $\mathscr L'_{j,\bullet}|U$ soit une résolution libre de
$i_*(\mathscr P_j)|U$ pour tout $j$. On a alors un isomorphisme canonique
\[
  (i\times_S 1)_*
  \bigl(\mathscr{T}\!or_\bullet^S
    (\mathscr P_\bullet,\mathscr Q_\bullet)\bigr)
  \xrightarrow{\sim}
  \mathscr H_\bullet
  (\mathscr L'_{\bullet\bullet}\otimes_S\mathscr Q_\bullet).
  \tag{6.5.12.1}\label{III.6.5.12.1.fr}
\]

Si $S$, $X'$, $Y'$ sont affines et si $\mathscr L'_{j,\bullet}$ est une
résolution libre de $i_*(\mathscr P_j)$ pour tout $j$, on est ramené, en vertu
de (6.5.11), au cas où $X'=X$ est $S$-plat, et il suffit d'appliquer (6.3.4).
Dans le cas général, on définit localement l'isomorphisme (6.5.12.1), et il
s'agit de vérifier que cette définition donne bien un isomorphisme global.
Pour cela, il faut se reporter à la définition du premier isomorphisme
(6.3.4.1) qui provient d'un isomorphisme de suites spectrales (0, 11.6.5 et
11.5.3), obtenu lui-même à partir d'un morphisme de bicomplexes
$L_{\bullet\bullet}\to L''_{\bullet\bullet}$, où $L_{\bullet\bullet}$ est la
résolution donnée de $P_\bullet$, $L''_{\bullet\bullet}$ une résolution
projective de $P_\bullet$ dans la catégorie des complexes de $A$-modules
limités

\oldpage[III]{152}
inférieurement (cf. (0, 11.5.2.2)); notre assertion résulte de ce que
l'isomorphisme (6.3.4.1) ne dépend pas de la résolution projective
$L''_{\bullet\bullet}$ choisie, en vertu de l'existence d'un homotopisme entre
deux telles résolutions (M, V, 1.2).
\end{namedtheorem}

\begin{proposition}[6.5.13]
Soient $X$, $Y$ deux $S$-préschémas, et supposons vérifiée l'une des conditions
suivantes :
\begin{enumerate}
  \item[(i)] $X$ et $Z=X\times_S Y$ sont localement noethériens et $X$ est plat
    sur $S$.
  \item[(ii)] $S$ et $X$ sont localement noethériens et $Y$ est de type fini
    sur $S$.
\end{enumerate}
Soit $\mathscr P_\bullet$ (resp. $\mathscr Q_\bullet$) un complexe de
$\mathscr O_X$-Modules (resp. de $\mathscr O_Y$-Modules) quasi-cohérents,
limité inférieurement. On suppose en outre que, pour tout $n$,
$\mathscr H_n(\mathscr P_\bullet)$ (resp.
$\mathscr H_n(\mathscr Q_\bullet)$) est un $\mathscr O_X$-Module (resp. un
$\mathscr O_Y$-Module) de type fini. Alors les
$\mathscr{T}\!or_n^S(\mathscr P_\bullet,\mathscr Q_\bullet)$ sont des
$\mathscr O_Z$-Modules cohérents.
\end{proposition}

\begin{proof}
Comme $\mathscr P_\bullet$ et $\mathscr Q_\bullet$ sont limités
inférieurement, la suite spectrale
$''\mathscr E(\mathscr P_\bullet,\mathscr Q_\bullet)$ est birégulière (6.5.4),
et en vertu de (0, 11.1.8), il suffit (puisque dans les deux cas (i), (ii), $Z$
est localement noethérien) de prouver que les termes
$''\mathscr E_{pq}^2$ sont cohérents. L'hypothèse sur les
$\mathscr H_n(\mathscr P_\bullet)$ et
$\mathscr H_n(\mathscr Q_\bullet)$ et l'expression (6.5.4.2) des
$''\mathscr E_{pq}^2$ montrent donc que la proposition est équivalente à son
cas particulier correspondant à $\mathscr P_\bullet$ et
$\mathscr Q_\bullet$ réduits à leurs termes de degré 0, autrement dit à son
\end{proof}

\begin{corollary}[6.5.14]
Supposons vérifiée l'une des conditions (i), (ii) de (6.5.13), et soit
$\mathscr F$ (resp. $\mathscr G$) un $\mathscr O_X$-Module (resp. un
$\mathscr O_Y$-Module) quasi-cohérent de type fini; alors les
$\mathscr{T}\!or_n^S(\mathscr F,\mathscr G)$ sont des $\mathscr O_Z$-Modules
cohérents.
\end{corollary}

\begin{proof}
La question étant locale sur $X$ et $Y$, on peut supposer $S$, $X$, $Y$
affines.

\textnormal{(i)} Sous les hypothèses de (i), $S$, $X$ et $Z$ sont noethériens.
Il existe donc une résolution localement libre $\mathscr L_\bullet$ de
$\mathscr F$ formée de $\mathscr O_X$-Modules de type fini (I, 1.3.7); comme
$X$ est plat sur $S$, il résulte de (6.3.4) que l'on a
\[
  \mathscr{T}\!or_n^S(\mathscr F,\mathscr G)
  =\mathscr H_n(\mathscr L_\bullet\otimes_S\mathscr G);
\]
or, les $\mathscr L_i\otimes_S\mathscr G$ sont des $\mathscr O_Z$-Modules
quasi-cohérents de type fini (I, 9.1.1), donc cohérents. On en conclut que
$\mathscr H_n(\mathscr L_\bullet\otimes_S\mathscr G)$ est cohérent
($0_{\mathrm I}$, 5.3.4).

\textnormal{(ii)} Supposons maintenant vérifiées les conditions (ii). Comme
l'anneau $A(Y)$ est quotient d'une $A(S)$-algèbre de polynômes à un nombre fini
d'indéterminées (I, 6.3.3), $Y$ est un sous-$S$-préschéma fermé d'un
$S$-préschéma affine $Y'$, plat et de type fini sur $S$; $Y'$ étant noethérien
(I, 6.3.7), il existe une résolution localement libre $\mathscr M_\bullet$ de
$j_*(\mathscr G)$ par des $\mathscr O_{Y'}$-Modules de type fini
($j:Y\to Y'$ étant l'injection canonique); en vertu de (6.5.12),
$\mathscr{T}\!or_n^S(\mathscr F,\mathscr G)$ est l'image réciproque par
$1\times j$ du $\mathscr O_{Z'}$-Module
$\mathscr H_n(\mathscr F\otimes_S\mathscr M_\bullet)$ (où
$Z'=X\times_S Y'$); on voit comme dans (i) que les
$\mathscr F\otimes_S\mathscr M_i$ sont des $\mathscr O_{Z'}$-Modules
cohérents, et l'on en tire encore la conclusion par ($0_{\mathrm I}$, 5.3.4).
\end{proof}

\begin{env}[6.5.15]
La théorie développée ci-dessus pour deux complexes $\mathscr P_\bullet$,
$\mathscr Q_\bullet$ de faisceaux quasi-cohérents sur deux $S$-préschémas
$X$, $Y$ se généralise sans peine au cas où l'on considère un nombre fini
quelconque de $S$-préschémas $X^{(i)}$ ($1\leqslant i\leqslant m$) et sur
chaque $X^{(i)}$ un complexe $\mathscr P_\bullet^{(i)}$ de
$\mathscr O_{X^{(i)}}$-Modules quasi-cohérents; si
$Z=X^{(1)}\times_S X^{(2)}\times_S\cdots\times_S X^{(m)}$, on définit ainsi un
$\mathscr O_Z$-Module quasi-cohérent
$\mathscr{T}\!or_q^S(\mathscr P_\bullet^{(1)},\mathscr P_\bullet^{(2)},
\ldots,\mathscr P_\bullet^{(m)})$. Nous laissons au lecteur le soin de
développer la théorie dans ce cas général et nous nous bornerons à écrire,
pour référence ultérieure, le terme $E_2$ de la seconde suite spectrale
(régulière) dont l'aboutissement est
$\mathscr{T}\!or_\bullet^S(\mathscr P_\bullet^{(1)},
\mathscr P_\bullet^{(2)},\ldots,\mathscr P_\bullet^{(m)})$ :

\oldpage[III]{153}

\[
  ''\mathscr E_{pq}^2
  =\bigoplus_{q_1+q_2+\cdots+q_m=q}
   \mathscr{T}\!or_p^S
   \bigl(\mathscr H_{q_1}(\mathscr P_\bullet^{(1)}),\ldots,
         \mathscr H_{q_m}(\mathscr P_\bullet^{(m)})\bigr).
  \tag{6.5.15.1}\label{III.6.5.15.1.fr}
\]
Nous étudierons dans (6.8) les suites spectrales d'associativité auxquelles
donnent lieu ces foncteurs hypertor d'un nombre quelconque de complexes.
\end{env}

\subsection{Foncteurs hypertor globaux de complexes de Modules quasi-cohérents
et suites spectrales de Künneth : cas de la base affine}
\label{subsection:III.6.6.fr}

\begin{env}[6.6.1]
Considérons un schéma affine $S=\operatorname{Spec}(A)$ et deux $S$-schémas
quasi-compacts $X^{(i)}$ ($i=1,2$); soit $\mathscr P_\bullet^{(i)}$ un complexe
de $\mathscr O_{X^{(i)}}$-Modules quasi-cohérents, limité inférieurement, dont
l'opérateur de dérivation est de degré $-1$ ($i=1,2$). Considérons d'autre
part un recouvrement fini
$\mathfrak U^{(i)}=(U_\alpha^{(i)})$ de $X^{(i)}$ par des ouverts affines;
soit
\[
  X=X^{(1)}\times_S X^{(2)},
\]
qui est un $S$-schéma quasi-compact (I, 5.5.1 et 6.6.4), et soit
$\mathfrak U=\mathfrak U^{(1)}\times_S\mathfrak U^{(2)}$ le recouvrement de
$X$ formé des ouverts affines
$U_\alpha^{(1)}\times_S U_\beta^{(2)}$. Pour tout couple d'entiers
$p\leqslant0$, $q\in\mathbf Z$, le groupe des $(-p)$-cochaînes alternées
\[
  C^{-p}(\mathfrak U^{(i)},\mathscr P_q^{(i)})
\]
du recouvrement $\mathfrak U^{(i)}$ à coefficients dans le faisceau
$\mathscr P_q^{(i)}$ (G, II, 5.1) est un $A$-module; pour $p>0$, on posera
$C^p(\mathfrak U^{(i)},\mathscr P_q^{(i)})=0$; on a ainsi défini un
bicomplexe
\[
  C^\bullet(\mathfrak U^{(i)},\mathscr P_\bullet^{(i)})
\]
de $A$-modules, dont les deux opérateurs de dérivation sont de degré $-1$. Il
résulte des définitions ($0_{\mathrm I}$, 12.1.2) que le $A$-module
d'homologie
\[
  \mathrm H_n\bigl(C^\bullet
  (\mathfrak U^{(i)},\mathscr P_\bullet^{(i)})\bigr)
\]
de ce bicomplexe (considéré à l'ordinaire comme un complexe simple pour le
degré total) n'est autre que le $A$-module d'hypercohomologie
\[
  \mathbf H^{-n}
  (\mathfrak U^{(i)},\mathscr P^{(i)\bullet}),
\]
où $\mathscr P^{(i)\bullet}$ est le complexe à opérateur de dérivation de
degré $+1$ obtenu en prenant $\mathscr P_{-q}^{(i)}$ pour composante de degré
$q$; par abus de notation, nous l'écrirons
$\mathbf H^{-n}(\mathfrak U^{(i)},\mathscr P_\bullet^{(i)})$. Il résulte alors
de (6.2.2) que ce $A$-module est canoniquement isomorphe au $A$-module
d'hypercohomologie
$\mathbf H^{-n}(X^{(i)},\mathscr P^{(i)\bullet})$, que nous écrirons de même
$\mathbf H^{-n}(X^{(i)},\mathscr P_\bullet^{(i)})$; il ne dépend donc pas du
recouvrement fini $\mathfrak U^{(i)}$ choisi.
\end{env}

\begin{env}[6.6.2]
Nous allons appliquer aux deux bicomplexes de $A$-modules
\[
  L_{\bullet\bullet}^{(i)}
  =C^\bullet(\mathfrak U^{(i)},\mathscr P_\bullet^{(i)})
  \qquad(i=1,2)
\]
et au bifoncteur covariant
$L_{\bullet\bullet}^{(1)}\otimes_A L_{\bullet\bullet}^{(2)}$ en ces deux
bicomplexes, la théorie générale de l'hyperhomologie des foncteurs par rapport
aux bicomplexes ($0_{\mathrm I}$, 11.7.4). Comme les cochaînes considérées sont
alternées et les recouvrements $\mathfrak U^{(i)}$ finis, on notera que les
modules $C^{-p}(\mathfrak U^{(i)},\mathscr P_q^{(i)})$ ne sont $\ne0$ que pour
un nombre fini (indépendant de $q$) de valeurs de $p$, et en particulier les
deux degrés de chacun des $L_{\bullet\bullet}^{(i)}$ sont limités
inférieurement. Nous désignerons par
\[
  \mathbf{Tor}_n^A(L_{\bullet\bullet}^{(1)},L_{\bullet\bullet}^{(2)})
\]
ou
\[
  \mathbf{Tor}_n^S
  (\mathfrak U^{(1)},\mathfrak U^{(2)};
   \mathscr P_\bullet^{(1)},\mathscr P_\bullet^{(2)})
\]
le $n$-ième module d'hyperhomologie de
$L_{\bullet\bullet}^{(1)}\otimes_A L_{\bullet\bullet}^{(2)}$, que nous
appellerons l'\emph{hypertor} d'indice $n$ de
$\mathscr P_\bullet^{(1)}$ et $\mathscr P_\bullet^{(2)}$, relatif aux
recouvrements $\mathfrak U^{(1)}$ et $\mathfrak U^{(2)}$. Lorsque
$\mathscr P_\bullet^{(1)}$ et $\mathscr P_\bullet^{(2)}$ sont réduits à leurs
termes de degré 0, $\mathscr F^{(1)}$ et $\mathscr F^{(2)}$, on écrit
\[
  \mathbf{Tor}_n^S
  (\mathfrak U^{(1)},\mathfrak U^{(2)};
   \mathscr F^{(1)},\mathscr F^{(2)})
\]
leur hypertor. On désignera par
\[
  \mathbf{Tor}_n^S
  (\mathfrak U^{(1)},\mathfrak U^{(2)};
   \mathscr P_\bullet^{(1)},\mathscr P_\bullet^{(2)}),
\]
suivant les conventions générales, le bicomplexe dont le composant d'indices
$(j,k)$ est
\[
  \mathbf{Tor}_n^S
  (\mathfrak U^{(1)},\mathfrak U^{(2)};
   \mathscr P_j^{(1)},\mathscr P_k^{(2)}).
\]

Comme $L_{\bullet\bullet}^{(i)}$ est un foncteur exact en
$\mathscr P_\bullet^{(i)}$, puisque les intersections des ensembles
\oldpage[III]{154}
de $\mathfrak U^{(i)}$ sont affines (I, 5.5.6 et 1.3.11),
$\mathbf{Tor}_\bullet^S
(\mathfrak U^{(1)},\mathfrak U^{(2)};
 \mathscr P_\bullet^{(1)},\mathscr P_\bullet^{(2)})$ est un
bi-$\partial$-foncteur covariant en $\mathscr P_\bullet^{(1)}$,
$\mathscr P_\bullet^{(2)}$, à valeurs dans la catégorie des $A$-modules
($0_{\mathrm I}$, 11.7.3). En outre, on sait ($0_{\mathrm I}$, 11.7.2) que ce
bifoncteur est l'aboutissement commun de six foncteurs spectraux biréguliers,
que nous désignerons par la notation
\[
  {}^{(t)}\mathscr E^S
  (\mathfrak U^{(1)},\mathfrak U^{(2)};
   \mathscr P_\bullet^{(1)},\mathscr P_\bullet^{(2)})
\]
ou
\[
  {}^{(t)}\mathscr E
  (\mathfrak U^{(1)},\mathfrak U^{(2)};
   \mathscr P_\bullet^{(1)},\mathscr P_\bullet^{(2)}),
\]
où $t$ doit être remplacé par une des lettres $a$, $b$, $a'$, $b'$, $c$,
$d$, et dont les termes $E_2$ sont les suivants :
\[
\begin{aligned}
  {}^{(a)}\mathscr E_{pq}^2
  &=\bigoplus_{q_1+q_2=q}
    \mathbf{Tor}_p^A
    \bigl(\mathrm H_{q_1}(L_{\bullet\bullet}^{(1)}),
          \mathrm H_{q_2}(L_{\bullet\bullet}^{(2)})\bigr),\\
  {}^{(b)}\mathscr E_{pq}^2
  &=\mathrm H_p\bigl(
    \mathbf{Tor}_q^{A,\mathrm{II}}
    (L_{\bullet\bullet}^{(1)},L_{\bullet\bullet}^{(2)})\bigr),\\
  {}^{(a')}\mathscr E_{pq}^2
  &=\bigoplus_{q_1+q_2=q}
    \mathbf{Tor}_p^A
    \bigl(\mathrm H_{q_1}^{\mathrm I}(L_{\bullet\bullet}^{(1)}),
          \mathrm H_{q_2}^{\mathrm I}(L_{\bullet\bullet}^{(2)})\bigr),\\
  {}^{(b')}\mathscr E_{pq}^2
  &=\mathrm H_p\bigl(
    \mathbf{Tor}_q^A
    (L_{\bullet\bullet}^{(1)},L_{\bullet\bullet}^{(2)})\bigr),\\
  {}^{(c)}\mathscr E_{pq}^2
  &=\bigoplus_{q_1+q_2=q}
    \mathbf{Tor}_p^A
    \bigl(\mathrm H_{q_1}^{\mathrm{II}}(L_{\bullet\bullet}^{(1)}),
          \mathrm H_{q_2}^{\mathrm{II}}(L_{\bullet\bullet}^{(2)})\bigr),\\
  {}^{(d)}\mathscr E_{pq}^2
  &=\mathrm H_p\bigl(
    \mathbf{Tor}_q^{A,\mathrm I}
    (L_{\bullet\bullet}^{(1)},L_{\bullet\bullet}^{(2)})\bigr),
\end{aligned}
\]
où les notations sont conformes à celles de la théorie générale de
l'hyperhomologie. Nous allons dans ce qui suit expliciter davantage ces
termes initiaux.
\end{env}

\begin{env}[6.6.3]
\emph{Suites spectrales $(a)$ et $(a')$.} Nous avons vu en (6.6.1) que le
module d'homologie $\mathrm H_n(L_{\bullet\bullet}^{(i)})$ du bicomplexe
$L_{\bullet\bullet}^{(i)}$ était égal à
$\mathbf H^{-n}(X^{(i)},\mathscr P_\bullet^{(i)})$; donc
\[
  {}^{(a)}\mathscr E_{pq}^2
  =\bigoplus_{q_1+q_2=q}
   \mathbf{Tor}_p^A
   \bigl(\mathbf H^{-q_1}(X^{(1)},\mathscr P_\bullet^{(1)}),
         \mathbf H^{-q_2}(X^{(2)},\mathscr P_\bullet^{(2)})\bigr).
\]
Par définition, le complexe
$\mathrm H_n^{\mathrm I}(L_{\bullet\bullet}^{(i)})$ a pour terme de degré $k$
le module d'homologie
$\mathrm H_n(C^\bullet(\mathfrak U^{(i)},\mathscr P_k^{(i)}))$, c'est-à-dire,
par définition, le module de cohomologie
$\mathbf H^{-n}(\mathfrak U^{(i)},\mathscr P_k^{(i)})$; on sait (1.4.1) que
ce module est canoniquement isomorphe à
$\mathbf H^{-n}(X^{(i)},\mathscr P_k^{(i)})$; donc
\[
  {}^{(a')}\mathscr E_{pq}^2
  =\bigoplus_{q_1+q_2=q}
   \mathbf{Tor}_p^A
   \bigl(\mathbf H^{-q_1}(X^{(1)},\mathscr P_\bullet^{(1)}),
         \mathbf H^{-q_2}(X^{(2)},\mathscr P_\bullet^{(2)})\bigr).
\]
\end{env}

\begin{env}[6.6.4]
\emph{Suites spectrales $(b)$ et $(b')$.} Par définition,
$\mathbf{Tor}_q^{A,\mathrm{II}}
(L_{\bullet\bullet}^{(1)},L_{\bullet\bullet}^{(2)})$ est un bicomplexe dont
le terme de degré $(h,k)$ est le $A$-module
\[
  \mathbf{Tor}_q^A
  \bigl(C^{-h}(\mathfrak U^{(1)},\mathscr P_\bullet^{(1)}),
        C^{-k}(\mathfrak U^{(2)},\mathscr P_\bullet^{(2)})\bigr).
\]
Soit $\Phi^{(i)}$ l'ensemble d'indices de $\mathfrak U^{(i)}$; par
définition, le complexe de modules
$C^r(\mathfrak U^{(i)},\mathscr P_\bullet^{(i)})$ ($r\geqslant0$) est somme
directe des complexes
$\Gamma(U_\rho^{(i)},\mathscr P_\bullet^{(i)})$, où $U_\rho^{(i)}$ est
l'intersection des $U_\xi^{(i)}$ pour $\xi\in\rho$, et $\rho$ parcourt
$\mathfrak P(\Phi^{(i)})$; donc le $A$-module
\[
  \mathbf{Tor}_q^A
  \bigl(C^{-h}(\mathfrak U^{(1)},\mathscr P_\bullet^{(1)}),
        C^{-k}(\mathfrak U^{(2)},\mathscr P_\bullet^{(2)})\bigr)
\]
est somme directe des $A$-modules
\[
  \mathbf{Tor}_q^A
  \bigl(\Gamma(U_\sigma^{(1)},\mathscr P_\bullet^{(1)}),
        \Gamma(U_\tau^{(2)},\mathscr P_\bullet^{(2)})\bigr),
\]
où $\sigma$ (resp. $\tau$) parcourt les éléments de
$\mathfrak P(\Phi^{(1)})$ (resp. $\mathfrak P(\Phi^{(2)})$) tels que
$\operatorname{Card}(\sigma)=-(h+1)$ (resp.
$\operatorname{Card}(\tau)=-(k+1)$). Comme $X^{(1)}$ et $X^{(2)}$ sont des
schémas, les $U_\rho^{(i)}$ sont affines, donc on a (6.4.1.1)
\[
  \mathbf{Tor}_q^A
  \bigl(\Gamma(U_\sigma^{(1)},\mathscr P_\bullet^{(1)}),
        \Gamma(U_\tau^{(2)},\mathscr P_\bullet^{(2)})\bigr)
  =\Gamma\bigl(U_\sigma^{(1)}\times_S U_\tau^{(2)},
    \mathscr{T}\!or_q^S
    (\mathscr P_\bullet^{(1)},\mathscr P_\bullet^{(2)})\bigr).
\]
\oldpage[III]{155}
On voit donc que ${}^{(b)}\mathscr E_{pq}^2$ est le $(-p)$-ème module de
cohomologie du complexe
\[
  L^\bullet(\Phi^{(1)},\Phi^{(2)};\mathscr S)
\]
des cochaînes bi-alternées sur $\Phi^{(1)}$ et $\Phi^{(2)}$ à valeurs dans le
système de coefficients
\[
  \mathscr S:(\sigma,\tau)\leadsto
  \Gamma\bigl(U_\sigma^{(1)}\times_S U_\tau^{(2)},
    \mathscr{T}\!or_q^S
    (\mathscr P_\bullet^{(1)},\mathscr P_\bullet^{(2)})\bigr)
\]
($0_{\mathrm I}$, 11.8.4). On sait alors ($0_{\mathrm I}$, 11.8.5 et 11.8.6)
que la cohomologie de ce complexe est la même que celle du complexe
$C^\bullet(\Phi^{(1)},\Phi^{(2)};\mathscr S)$ de toutes les cochaînes sur
$\Phi^{(1)}$ et $\Phi^{(2)}$ à valeurs dans $\mathscr S$, et aussi que celle du
complexe $P^\bullet(\Phi^{(1)},\Phi^{(2)};\mathscr S)$, dont les éléments sont
les combinaisons linéaires des
\[
  \lambda(\sigma,\tau)\in
  \Gamma\bigl(U_\sigma^{(1)}\times_S U_\tau^{(2)},
    \mathscr{T}\!or_q^S
    (\mathscr P_\bullet^{(1)},\mathscr P_\bullet^{(2)})\bigr),
\]
où $\sigma=(\alpha_0,\ldots,\alpha_h)$ et
$\tau=(\beta_0,\ldots,\beta_h)$ sont des suites ayant même nombre d'éléments.
Mais on a alors
\[
  U_\sigma^{(1)}\times_S U_\tau^{(2)}
  =(U_{\alpha_0}^{(1)}\times_S U_{\beta_0}^{(2)})\cap\cdots\cap
    (U_{\alpha_h}^{(1)}\times_S U_{\beta_h}^{(2)})
  \qquad\text{(I, 3.2.7)}.
\]
Si l'on désigne par $\mathfrak U$ le recouvrement de
$Z=X^{(1)}\times_S X^{(2)}$ par les ouverts affines
$U_\alpha^{(1)}\times_S U_\beta^{(2)}$, on voit finalement, compte tenu de ce
que $X^{(1)}\times_S X^{(2)}$ est un schéma, que l'on a, en vertu de (1.3.1),
\[
  {}^{(b)}\mathscr E_{pq}^2
  =\mathbf H^{-p}\bigl(X^{(1)}\times_S X^{(2)},
    \mathscr{T}\!or_q^S
    (\mathscr P_\bullet^{(1)},\mathscr P_\bullet^{(2)})\bigr).
\]

En second lieu,
$\mathbf{Tor}_q^A(L_{\bullet\bullet}^{(1)},L_{\bullet\bullet}^{(2)})$ est un
bicomplexe dont le terme de degré $(h,k)$ est la somme directe des $A$-modules
\[
  \mathbf{Tor}_q^A
  \bigl(C^{-h_1}(\mathfrak U^{(1)},\mathscr P_{k_1}^{(1)}),
        C^{-h_2}(\mathfrak U^{(2)},\mathscr P_{k_2}^{(2)})\bigr)
\]
tels que $h_1+h_2=h$ et $k_1+k_2=k$; explicitant les modules
$C^r(\mathfrak U^{(i)},\mathscr P_s^{(i)})$ comme ci-dessus, on voit encore
que ce terme est somme directe des $A$-modules
\[
  \Gamma\bigl(U_\sigma^{(1)}\times_S U_\tau^{(2)},
    \mathscr{T}\!or_q^S
    (\mathscr P_{k_1}^{(1)},\mathscr P_{k_2}^{(2)})\bigr),
\]
où $k_1+k_2=k$, et $\sigma$ (resp. $\tau$) parcourt les éléments de
$\mathfrak P(\Phi^{(1)})$ (resp. $\mathfrak P(\Phi^{(2)})$) tels que
$\operatorname{Card}(\sigma)+\operatorname{Card}(\tau)=-h-2$. Le terme
${}^{(b')}\mathscr E_{pq}^2$ que nous calculons est le $(-p)$-ème module de
cohomologie d'un bicomplexe $N^{\bullet\bullet}=(N^{hk})$, où le complexe
simple $N^{\bullet k}$ est le complexe de cochaînes bi-alternées sur
$\Phi^{(1)}$ et $\Phi^{(2)}$, à valeurs dans le système de coefficients
\[
  \mathscr S_k:(\sigma,\tau)\leadsto
  \Gamma\left(U_\sigma^{(1)}\times_S U_\tau^{(2)},
    \bigoplus_{k_1+k_2=k}
    \mathscr{T}\!or_q^S
    (\mathscr P_{-k_1}^{(1)},\mathscr P_{-k_2}^{(2)})\right),
\]
ces systèmes de coefficients formant un complexe $\mathscr S^\bullet$ où la
différentielle provient de celle du complexe simple associé au bicomplexe
$\mathscr{T}\!or_q^S
(\mathscr P_\bullet^{(1)},\mathscr P_\bullet^{(2)})$. On sait que la
cohomologie de $N^{\bullet\bullet}$ est la même que celle du bicomplexe
$C^\bullet(\Phi^{(1)},\Phi^{(2)};\mathscr S^\bullet)$
($0_{\mathrm I}$, 11.8.9), et aussi la même que celle du bicomplexe
$P^\bullet(\Phi^{(1)},\Phi^{(2)};\mathscr S^\bullet)$, où les éléments de
degrés $(h,k)$ sont les combinaisons linéaires des
\[
  \lambda(\sigma,\tau)\in
  \Gamma\left(U_\sigma^{(1)}\times_S U_\tau^{(2)},
    \bigoplus_{k_1+k_2=k}
    \mathscr{T}\!or_q^S
    (\mathscr P_{-k_1}^{(1)},\mathscr P_{-k_2}^{(2)})\right),
\]
$\sigma=(\alpha_0,\ldots,\alpha_h)$,
$\tau=(\beta_0,\ldots,\beta_h)$ étant des suites ayant même nombre
d'éléments ($0_{\mathrm I}$, 11.8.10). On voit alors comme ci-dessus que
${}^{(b')}\mathscr E_{pq}^2$ est le $(-p)$-ème module de cohomologie du
bicomplexe $C^\bullet(\mathfrak U,\mathscr Q^\bullet)$, où
$\mathscr Q^\bullet$ est le complexe simple associé au bicomplexe
$\mathscr{T}\!or_q^S
(\mathscr P_\bullet^{(1)},\mathscr P_\bullet^{(2)})$ de
$\mathscr O_Z$-Modules. Avec les conventions faites dans (6.6.1), on a donc
\[
  {}^{(b')}\mathscr E_{pq}^2
  =\mathbf H^{-p}\bigl(X^{(1)}\times_S X^{(2)},
    \mathscr{T}\!or_q^S
    (\mathscr P_\bullet^{(1)},\mathscr P_\bullet^{(2)})\bigr).
\]
\end{env}

\oldpage[III]{156}
\begin{env}[6.6.5]
\emph{Suites spectrales $(c)$ et $(d)$.} Par définition, le complexe
$\mathrm H_n^{\mathrm{II}}(L_{\bullet\bullet}^{(i)})$ a pour terme de degré
$h$ le $A$-module
$C^{-h}(\mathfrak U^{(i)},\mathscr H_n(\mathscr P_\bullet^{(i)}))$, en vertu
de l'exactitude du foncteur $C^{-h}$. On a donc, par définition de l'hypertor
de deux Modules relatif à deux recouvrements (6.6.2)
\[
  {}^{(c)}\mathscr E_{pq}^2
  =\bigoplus_{q_1+q_2=q}
   \mathbf{Tor}_p^S
   \bigl(\mathfrak U^{(1)},\mathfrak U^{(2)};
         \mathscr H_{q_1}(\mathscr P_\bullet^{(1)}),
         \mathscr H_{q_2}(\mathscr P_\bullet^{(2)})\bigr).
\]
Enfin, par définition,
$\mathbf{Tor}_q^{A,\mathrm I}
(L_{\bullet\bullet}^{(1)},L_{\bullet\bullet}^{(2)})$ est un bicomplexe dont
le terme de degré $(h,k)$ est le $A$-module
$\mathbf{Tor}_q^S
(\mathfrak U^{(1)},\mathfrak U^{(2)};
 \mathscr P_h^{(1)},\mathscr P_k^{(2)})$. On a donc
\[
  {}^{(d)}\mathscr E_{pq}^2
  =\mathrm H_p\bigl(
    \mathbf{Tor}_q^S
    (\mathfrak U^{(1)},\mathfrak U^{(2)};
     \mathscr P_\bullet^{(1)},\mathscr P_\bullet^{(2)})\bigr).
\]
\end{env}

\begin{env}[6.6.6]
La théorie de l'hyperhomologie des foncteurs de bicomplexes
($0_{\mathrm I}$, 11.7.3) montre, comme dans (6.3.4), que, pour toute
résolution \emph{plate} de Cartan-Eilenberg $M_{\bullet\bullet}^{(i)}$ de
$L_{\bullet\bullet}^{(i)}$ (dans la catégorie des complexes de modules
limités inférieurement) ($i=1,2$), on a des isomorphismes canoniques de
bi-$\partial$-foncteurs
\[
\begin{aligned}
  \mathbf{Tor}_\bullet^S
  (\mathfrak U^{(1)},\mathfrak U^{(2)};
   \mathscr P_\bullet^{(1)},\mathscr P_\bullet^{(2)})
  &\xrightarrow{\sim}
   \mathrm H_\bullet
   (M_{\bullet\bullet}^{(1)}\otimes_A M_{\bullet\bullet}^{(2)})\\
  &\xrightarrow{\sim}
   \mathrm H_\bullet
   (M_{\bullet\bullet}^{(1)}\otimes_A L_{\bullet\bullet}^{(2)})
   \xrightarrow{\sim}
   \mathrm H_\bullet
   (L_{\bullet\bullet}^{(1)}\otimes_A M_{\bullet\bullet}^{(2)}).
\end{aligned}
\tag{6.6.6.1}\label{III.6.6.6.1.fr}
\]
\end{env}

\begin{env}[6.6.7]
Nous allons maintenant montrer que l'hypertor global défini dans (6.6.2), et
les six suites spectrales correspondantes, ne dépendent pas des recouvrements
ouverts affines finis $\mathfrak U^{(i)}$ qui ont servi à les définir (à des
isomorphismes canoniques près). Il suffira pour cela de montrer que si
$\mathfrak V^{(i)}$ sont deux autres recouvrements de même nature, tels que
$\mathfrak V^{(i)}$ soit \emph{plus fin} que $\mathfrak U^{(i)}$ pour
$i=1,2$, alors on a des \emph{isomorphismes canoniques} de foncteurs spectraux
\[
  {}^{(t)}\mathscr E
  (\mathfrak U^{(1)},\mathfrak U^{(2)};
   \mathscr P_\bullet^{(1)},\mathscr P_\bullet^{(2)})
  \xrightarrow{\sim}
  {}^{(t)}\mathscr E
  (\mathfrak V^{(1)},\mathfrak V^{(2)};
   \mathscr P_\bullet^{(1)},\mathscr P_\bullet^{(2)})
  \tag{6.6.7, $t$}\label{III.6.6.7.t.fr}
\]
où $t$ est remplacé par $a$, $b$, $a'$, $b'$, $c$ ou $d$.

Or, on a pour $i=1,2$ des homomorphismes de bicomplexes
\[
  C^\bullet(\mathfrak U^{(i)},\mathscr P_\bullet^{(i)})
  \longrightarrow
  C^\bullet(\mathfrak V^{(i)},\mathscr P_\bullet^{(i)})
\]
bien définis à homotopies près (G, II, 5.7.1); il en résulte déjà des
homomorphismes (6.6.7, $t$) canoniquement définis et compatibles avec les
opérateurs bords dans les aboutissements ($0_{\mathrm I}$, 11.3.2). En outre,
le calcul des termes $E_2$ des suites spectrales $(a)$, $(b)$, $(a')$, $(b')$,
montre que pour ces suites spectrales l'homomorphisme (6.6.7, $t$) est un
isomorphisme sur les termes $E_2$; comme ces suites spectrales sont
birégulières, on voit que (6.6.7, $t$) est un isomorphisme pour ces trois
foncteurs spectraux, donc un \emph{isomorphisme de bi-$\partial$-foncteurs}
pour leur aboutissement commun ($0_{\mathrm I}$, 11.1.5).

En particulier, pour des $\mathscr O_{X^{(i)}}$-Modules quasi-cohérents
$\mathscr F^{(i)}$ ($i=1,2$), l'homomorphisme canonique
\[
  \mathbf{Tor}_\bullet^S
  (\mathfrak U^{(1)},\mathfrak U^{(2)};
   \mathscr F^{(1)},\mathscr F^{(2)})
  \longrightarrow
  \mathbf{Tor}_\bullet^S
  (\mathfrak V^{(1)},\mathfrak V^{(2)};
   \mathscr F^{(1)},\mathscr F^{(2)})
\]
est bijectif; vu le calcul de (6.6.5), on voit que (6.6.7, $t$) est aussi un
isomorphisme des termes $E_2$ pour $t=c$ et $t=d$. On conclut comme ci-dessus
que (6.6.7, $t$) est aussi un isomorphisme de suites spectrales pour $t=c$ et
$t=d$.
\oldpage[III]{157}
On peut considérer que les isomorphismes (6.6.7, $t$) définissent des systèmes
inductifs de foncteurs spectraux sur l'ensemble filtrant des couples
$(\mathfrak U^{(1)},\mathfrak U^{(2)})$ de recouvrements ouverts affines finis
de $X^{(1)}$ et $X^{(2)}$. Nous désignerons par
\[
  {}^{(t)}\mathscr E
  (X^{(1)},X^{(2)};
   \mathscr P_\bullet^{(1)},\mathscr P_\bullet^{(2)})
  \quad\text{ou}\quad
  {}^{(t)}\mathscr E^S
  (X^{(1)},X^{(2)};
   \mathscr P_\bullet^{(1)},\mathscr P_\bullet^{(2)})
\]
la limite inductive de ce système, par
$\mathbf{Tor}_\bullet^S
(X^{(1)},X^{(2)};\mathscr P_\bullet^{(1)},\mathscr P_\bullet^{(2)})$
l'aboutissement de ce foncteur spectral, que nous appellerons l'\emph{hypertor
global} des deux complexes $\mathscr P_\bullet^{(1)}$ et
$\mathscr P_\bullet^{(2)}$; si $\mathscr P_\bullet^{(1)}$ et
$\mathscr P_\bullet^{(2)}$ sont réduits à leurs termes de degré 0,
$\mathscr F^{(1)}$ et $\mathscr F^{(2)}$, nous écrirons
\[
  \mathbf{Tor}_\bullet^S
  (X^{(1)},X^{(2)};\mathscr F^{(1)},\mathscr F^{(2)}),
\]
et conformément aux conventions générales,
\[
  \mathbf{Tor}_q^S
  (X^{(1)},X^{(2)};\mathscr P_\bullet^{(1)},\mathscr P_\bullet^{(2)})
\]
sera donc le bicomplexe des
\[
  \mathbf{Tor}_q^S
  (X^{(1)},X^{(2)};\mathscr P_h^{(1)},\mathscr P_k^{(2)}).
\]
\end{env}

\begin{env}[6.6.8]
Les hypothèses étant celles de (6.6.1), considérons maintenant deux
$S$-morphismes $f_i:X^{(i)}\to Y^{(i)}$, où
$Y^{(i)}=\operatorname{Spec}(B_i)$ est un $S$-schéma \emph{affine}, $B_i$
étant donc une $A$-algèbre ($i=1,2$); cela définit donc un $A$-homomorphisme
$B_i\to\Gamma(X^{(i)},\mathscr O_{X^{(i)}})$ (I, 2.2.4), et par suite chacun
des $L_{\bullet\bullet}^{(i)}$ définis dans (6.6.2) est un bicomplexe de
$B_i$-modules; on en conclut que
$L_{\bullet\bullet}^{(1)}\otimes_A L_{\bullet\bullet}^{(2)}$ est un
quadricomplexe de $(B_1\otimes_A B_2)$-modules, et ses six foncteurs spectraux
d'hyperhomologie peuvent donc être considérés comme prenant leurs valeurs
dans la catégorie des suites spectrales de $(B_1\otimes_A B_2)$-modules. Si
l'on pose
$Y=Y^{(1)}\times_S Y^{(2)}=\operatorname{Spec}(B_1\otimes_A B_2)$, on peut
considérer les $\mathscr O_Y$-Modules quasi-cohérents associés à ces modules
(I, 1.3.4); nous noterons
${}^{(t)}\mathscr E
(f_1,f_2;\mathscr P_\bullet^{(1)},\mathscr P_\bullet^{(2)})$
(pour $t=a$, $b$, $a'$, $b'$, $c$ ou $d$) les six suites spectrales
\[
  \bigl({}^{(t)}\mathscr E
  (X^{(1)},X^{(2)};
   \mathscr P_\bullet^{(1)},\mathscr P_\bullet^{(2)})\bigr)^{\sim}
\]
de $\mathscr O_Y$-Modules, et
$\mathfrak{Tor}_\bullet^S
(f_1,f_2;\mathscr P_\bullet^{(1)},\mathscr P_\bullet^{(2)})$ leur
aboutissement commun
\[
  \bigl(\mathbf{Tor}_\bullet^S
  (X^{(1)},X^{(2)};
   \mathscr P_\bullet^{(1)},\mathscr P_\bullet^{(2)})\bigr)^{\sim}.
\]
On le notera
$\mathfrak{Tor}_\bullet^S(f_1,f_2;\mathscr F^{(1)},\mathscr F^{(2)})$
lorsque $\mathscr P_\bullet^{(i)}$ est réduit à son terme de degré 0,
$\mathscr F^{(i)}$ ($i=1,2$).
\end{env}

\subsection{Foncteurs hypertor globaux de complexes de Modules quasi-cohérents
et suites spectrales de Künneth : cas général.}

\begin{env}[6.7.1]
Nous allons maintenant généraliser les définitions de (6.6.8) au cas où $S$
est un préschéma quelconque, $X^{(i)}$, $Y^{(i)}$ des $S$-préschémas et
$f_i:X^{(i)}\to Y^{(i)}$ des morphismes \emph{séparés et quasi-compacts}. Il
s'agit alors, pour tout couple de complexes $\mathscr P_\bullet^{(i)}$ de
$\mathscr O_{X^{(i)}}$-Modules quasi-cohérents, limités inférieurement
($i=1,2$), de définir pour tout $n$ un $\mathscr O_Y$-Module quasi-cohérent
$\mathfrak{Tor}_n^S
(f_1,f_2;\mathscr P_\bullet^{(1)},\mathscr P_\bullet^{(2)})$ ainsi que 6
foncteurs spectraux, se réduisant aux définitions de (6.6.8) lorsque $S$,
$Y^{(1)}$ et $Y^{(2)}$ sont affines (on pose
$Y=Y^{(1)}\times_S Y^{(2)}$).

Supposons d'abord $S=\operatorname{Spec}(A)$ \emph{affine}, mais $Y^{(1)}$
et $Y^{(2)}$ quelconques; soit $W^{(i)}$ un ouvert affine de $Y^{(i)}$;
$f_i^{-1}(W^{(i)})$ est alors un $S$-schéma quasi-compact,
$W=W^{(1)}\times_S W^{(2)}$ un ouvert affine de $Y$; soit
$f'_i:f_i^{-1}(W^{(i)})\to W^{(i)}$ la restriction de $f_i$, et
$\mathscr P'_\bullet{}^{(i)}$ la restriction
$\mathscr P_\bullet^{(i)}|f_i^{-1}(W^{(i)})$ ($i=1,2$). On a alors d'après
(6.6.8) les suites spectrales
${}^{(t)}\mathscr E
(f'_1,f'_2;\mathscr P'_\bullet{}^{(1)},\mathscr P'_\bullet{}^{(2)})$
de $(\mathscr O_Y|W)$-Modules quasi-cohérents, et il s'agit de vérifier
qu'elles satisfont aux conditions de recollement ($0_{\mathrm I}$, 3.3.1).
On est aussitôt ramené au cas où $Y^{(i)}=\operatorname{Spec}(B_i)$ est affine
et où $W^{(i)}=D(g_i)$, où $g_i\in B_i$, de sorte que
$W=D(g_1\otimes g_2)$
\oldpage[III]{158}
dans $Y=\operatorname{Spec}(B_1\otimes_A B_2)$ (II, 4.3.2.1); si
$X'^{(i)}=f_i^{-1}(W^{(i)})$, il s'agit d'établir un isomorphisme canonique
de foncteurs spectraux
\[
  {}^{(t)}\mathscr E
  (X'^{(1)},X'^{(2)};
   \mathscr P'_\bullet{}^{(1)},\mathscr P'_\bullet{}^{(2)})
  \xrightarrow{\sim}
  {}^{(t)}\mathscr E
  (X^{(1)},X^{(2)};
   \mathscr P_\bullet^{(1)},\mathscr P_\bullet^{(2)})
  \otimes_B B_g
  \tag{6.7.1.1}\label{III.6.7.1.1.fr}
\]
où l'on a posé $B=B_1\otimes_A B_2$ et $g=g_1\otimes g_2$. Pour cela,
partons de recouvrements ouverts affines finis $\mathfrak U^{(i)}$ de
$X^{(i)}$ ($i=1,2$), et soit $\mathfrak U'^{(i)}$ la trace de
$\mathfrak U^{(i)}$ sur $X'^{(i)}$, qui est encore formé d'ouverts affines
(I, 5.5.10); de façon précise, on a
\[
  C^\bullet(\mathfrak U'^{(i)},\mathscr P'_\bullet{}^{(i)})
  =C^\bullet(\mathfrak U^{(i)},\mathscr P_\bullet^{(i)})
   \otimes_{B_i}(B_i)_{g_i}.
\]
Si on pose
$L'_{\bullet\bullet}{}^{(i)}
=C^\bullet(\mathfrak U'^{(i)},\mathscr P'_\bullet{}^{(i)})$, on a donc
\[
  L'_{\bullet\bullet}{}^{(1)}\otimes_A L'_{\bullet\bullet}{}^{(2)}
  =\bigl(L_{\bullet\bullet}^{(1)}\otimes_{B_1}(B_1)_{g_1}\bigr)
   \otimes_A
   \bigl(L_{\bullet\bullet}^{(2)}\otimes_{B_2}(B_2)_{g_2}\bigr);
\]
comme on a $B_g=(B_1)_{g_1}\otimes_A(B_2)_{g_2}$ à un isomorphisme
canonique près, on a, à un isomorphisme canonique près,
\[
  L'_{\bullet\bullet}{}^{(1)}\otimes_A L'_{\bullet\bullet}{}^{(2)}
  =(L_{\bullet\bullet}^{(1)}\otimes_A L_{\bullet\bullet}^{(2)})
    \otimes_B B_g.
\]
Si $M_{\bullet\bullet\bullet}^{(i)}$ est une résolution projective de
Cartan-Eilenberg de $L_{\bullet\bullet}^{(i)}$, qu'on peut supposer formée de
$B_i$-modules, il résulte du fait que $(B_i)_{g_i}$ est \emph{plat} sur $B_i$
que
$M'_{\bullet\bullet\bullet}{}^{(i)}
=M_{\bullet\bullet\bullet}^{(i)}\otimes_{B_i}(B_i)_{g_i}$ est une résolution
projective de Cartan-Eilenberg du bicomplexe
$L'_{\bullet\bullet}{}^{(i)}$; en outre, on a
\[
  M'_{\bullet\bullet\bullet}{}^{(1)}\otimes_A M'_{\bullet\bullet\bullet}{}^{(2)}
  =(M_{\bullet\bullet\bullet}^{(1)}\otimes_A M_{\bullet\bullet\bullet}^{(2)})
    \otimes_B B_g.
\]
L'isomorphisme cherché (6.7.1.1) résulte alors aussitôt des définitions de
l'hyperhomologie d'un bicomplexe, et de l'exactitude du foncteur
$G\otimes_B B_g$ en le $B$-module $G$.
\end{env}

\begin{env}[6.7.2]
Supposons maintenant $S$ quelconque, et soient
$u_i:Y^{(i)}\to S$ les morphismes structuraux ($i=1,2$). Soit $(S_\alpha)$ un
recouvrement ouvert affine de $S$, posons
$Y_\alpha^{(i)}=u_i^{-1}(S_\alpha)$,
$X_\alpha^{(i)}=f_i^{-1}(Y_\alpha^{(i)})$, et soit
$f_{i\alpha}:X_\alpha^{(i)}\to Y_\alpha^{(i)}$ la restriction de $f_i$, qui
est un morphisme \emph{séparé et quasi-compact}. Les
$Y_\alpha=Y_\alpha^{(1)}\times_{S_\alpha}Y_\alpha^{(2)}$ forment un
recouvrement ouvert de $Y$, et sur chaque $Y_\alpha$ sont définis par (6.7.1)
des foncteurs spectraux
\[
  {}^{(t)}\mathscr E^{S_\alpha}
  \bigl(f_{1\alpha},f_{2\alpha};
    \mathscr P_\bullet^{(1)}|X_\alpha^{(1)},
    \mathscr P_\bullet^{(2)}|X_\alpha^{(2)}\bigr);
\]
il s'agit encore de montrer que ces foncteurs vérifient les conditions de
recollement.

On se ramène aussitôt à la situation suivante :
$S=\operatorname{Spec}(A)$ est affine, $S'=D(h)$, où $h\in A$, et
$u_i(Y^{(i)})\subset S'$; on peut en outre supposer
$Y^{(i)}=\operatorname{Spec}(B_i)$ affine; il s'agit de définir des
isomorphismes canoniques
\[
  {}^{(t)}\mathscr E^S
  (X^{(1)},X^{(2)};
   \mathscr P_\bullet^{(1)},\mathscr P_\bullet^{(2)})
  \xrightarrow{\sim}
  {}^{(t)}\mathscr E^{S'}
  (X^{(1)},X^{(2)};
   \mathscr P_\bullet^{(1)},\mathscr P_\bullet^{(2)}).
  \tag{6.7.2.1}\label{III.6.7.2.1.fr}
\]
Or, avec les notations de (6.6.2), les $L_{\bullet\bullet}^{(i)}$ sont formés
de $A_h$-modules, et l'on a donc
$L_{\bullet\bullet}^{(1)}\otimes_{A_h}L_{\bullet\bullet}^{(2)}
=L_{\bullet\bullet}^{(1)}\otimes_A L_{\bullet\bullet}^{(2)}$ à un
isomorphisme canonique près; comme on peut prendre une résolution projective
de Cartan-Eilenberg $M_{\bullet\bullet\bullet}^{(i)}$ de
$L_{\bullet\bullet}^{(i)}$ formée de $A_h$-modules, cela donne aussitôt
l'isomorphisme canonique cherché.
\end{env}

Nous avons en résumé démontré le

\begin{theorem}[6.7.3]
Soient $S$ un préschéma, $f_i:X_i\to Y_i$ un $S$-morphisme séparé et
quasi-compact de $S$-préschémas, $\mathscr P_\bullet^{(i)}$ un complexe de
$\mathscr O_{X^{(i)}}$-Modules quasi-cohérents limité inférieurement
($i=1,2$); on pose $Y=Y^{(1)}\times_S Y^{(2)}$. Il existe un
bi-$\partial$-foncteur
  $\mathfrak{Tor}_\bullet^S
(f_1,f_2;\mathscr P_\bullet^{(1)},\mathscr P_\bullet^{(2)})$
à valeurs dans la catégorie des $\mathscr O_Y$-Modules quasi-cohérents, tel
que si $V^{(i)}$ est un ouvert affine de $Y^{(i)}$ ($i=1,2$) et
$V=V^{(1)}\times_S V^{(2)}$, on ait
\[
  \mathfrak{Tor}_\bullet^S
  (f_1,f_2;\mathscr P_\bullet^{(1)},\mathscr P_\bullet^{(2)})|V
  =\Bigl(\mathbf{Tor}_\bullet^S
    \bigl(f_1^{-1}(V^{(1)}),f_2^{-1}(V^{(2)});
      \mathscr P_\bullet^{(1)}|f_1^{-1}(V^{(1)}),
      \mathscr P_\bullet^{(2)}|f_2^{-1}(V^{(2)})\bigr)\Bigr)^{\sim}.
\]
\oldpage[III]{159}
Ce bifoncteur est l'aboutissement de six foncteurs spectraux biréguliers
\[
  {}^{(t)}\mathscr E
  (f_1,f_2;\mathscr P_\bullet^{(1)},\mathscr P_\bullet^{(2)})
  \qquad(t=a,b,a',b',c,d)
\]
dont les termes $E_2$ sont donnés par
\[
\begin{aligned}
  {}^{(a)}\mathscr E_{pq}^2
  &=\bigoplus_{q_1+q_2=q}
    \mathscr{T}\!or_p^S
    \bigl(\mathscr H^{-q_1}(f_1,\mathscr P_\bullet^{(1)}),
          \mathscr H^{-q_2}(f_2,\mathscr P_\bullet^{(2)})\bigr),\\
  {}^{(b)}\mathscr E_{pq}^2
  &=\mathscr H^{-p}\bigl(f_1\times_S f_2,
    \mathfrak{Tor}_q^S
    (\mathscr P_\bullet^{(1)},\mathscr P_\bullet^{(2)})\bigr),\\
  {}^{(a')}\mathscr E_{pq}^2
  &=\bigoplus_{q_1+q_2=q}
    \mathfrak{Tor}_p^S
    \bigl(\mathscr H^{-q_1}(f_1,\mathscr P_\bullet^{(1)}),
          \mathscr H^{-q_2}(f_2,\mathscr P_\bullet^{(2)})\bigr),\\
  {}^{(b')}\mathscr E_{pq}^2
  &=\mathscr H^{-p}\bigl(f_1\times_S f_2,
    \mathscr{T}\!or_q^S
    (\mathscr P_\bullet^{(1)},\mathscr P_\bullet^{(2)})\bigr),\\
  {}^{(c)}\mathscr E_{pq}^2
  &=\bigoplus_{q_1+q_2=q}
    \mathfrak{Tor}_p^S
    \bigl(f_1,f_2;
      \mathscr H_{q_1}(\mathscr P_\bullet^{(1)}),
      \mathscr H_{q_2}(\mathscr P_\bullet^{(2)})\bigr),\\
  {}^{(d)}\mathscr E_{pq}^2
  &=\mathscr H_p\bigl(\mathfrak{Tor}_q^S
    (f_1,f_2;\mathscr P_\bullet^{(1)},\mathscr P_\bullet^{(2)})\bigr).
\end{aligned}
\]
On dit que les suites spectrales $(a)$ et $(b)$ sont les \emph{suites
spectrales de Künneth}.

On notera que les suites spectrales $(a)$ et $(a')$ (resp. $(b)$ et $(b')$)
sont identiques lorsque $\mathscr P_\bullet^{(1)}$ et
$\mathscr P_\bullet^{(2)}$ se réduisent à leurs termes de degré 0; dans ce cas,
les suites $(c)$ et $(d)$ sont dégénérées et sont donc sans intérêt.
\end{theorem}

\begin{remark}[6.7.4]
Les hypertor globaux que nous avons définis ci-dessus comprennent comme cas
particuliers, à la fois les \emph{Modules d'hypercohomologie} définis dans
(6.2.1) et les \emph{hypertor locaux} définis dans (6.5.3). Montrons que l'on
a, pour tout morphisme $f:X\to Y$ \emph{quasi-compact et séparé} et tout
complexe $\mathscr P_\bullet$ de $\mathscr O_X$-Modules quasi-cohérents,
limité inférieurement, un isomorphisme canonique de $\partial$-foncteurs en
$\mathscr P_\bullet$
\[
  \mathfrak{Tor}_n^Y(f,1_Y;\mathscr P_\bullet,\mathscr O_Y)
  \xrightarrow{\sim}
  \mathscr H^{-n}(f,\mathscr P_\bullet)
  \qquad\text{(pour tout $n\in\mathbf Z$)}.
  \tag{6.7.4.1}\label{III.6.7.4.1.fr}
\]
En effet, les méthodes de recollement de (6.7.2) ramènent aussitôt au cas où
$Y$ est affine; on peut alors, en vertu de (6.2.2), calculer les deux membres
de (6.7.4.1) à l'aide d'un même recouvrement fini $\mathfrak U$ de $X$ par des
ouverts affines, et (pour le premier membre) du recouvrement de $Y$ formé de
$Y$ lui-même; avec les notations de (6.6.2), le bicomplexe
$L_{\bullet\bullet}^{(2)}$ est alors réduit à son terme de degrés $(0,0)$,
égal à $A$, et la conclusion résulte de ($0_{\mathrm I}$, 11.7.5). Pour une
généralisation de ce résultat, voir (6.7.7); mais on notera que lorsque dans
le premier membre de (6.7.4.1), on remplace $\mathscr O_Y$ par un
$\mathscr O_Y$-Module quasi-cohérent \emph{quelconque} $\mathscr F$, on n'a
plus en général un isomorphisme avec
$\mathscr H^{-n}(f,\mathscr P_\bullet\otimes_{\mathscr O_Y}\mathscr F)$, bien
que, dans le calcul précédent, le bicomplexe
$L_{\bullet\bullet}^{(1)}\otimes_A L_{\bullet\bullet}^{(2)}$ s'identifie
encore au bicomplexe
$C^\bullet(\mathfrak U,\mathscr P_\bullet\otimes_Y\mathscr F)$.

D'autre part, on a un isomorphisme canonique de bi-$\partial$-foncteurs
\[
  \mathfrak{Tor}_\bullet^S
  (1_{X^{(1)}},1_{X^{(2)}};
   \mathscr P_\bullet^{(1)},\mathscr P_\bullet^{(2)})
  \xrightarrow{\sim}
  \mathscr{T}\!or_\bullet^S
  (\mathscr P_\bullet^{(1)},\mathscr P_\bullet^{(2)}).
  \tag{6.7.4.2}\label{III.6.7.4.2.fr}
\]
En effet, on se ramène encore, par (6.7.1) et (6.7.2), au cas où $S$ et les
$X^{(i)}$ sont affines; dans le calcul du premier membre de (6.7.4.2), on peut
alors prendre pour
\oldpage[III]{160}
recouvrement $\mathfrak U^{(i)}$ la famille réduite au seul élément $X^{(i)}$,
de sorte qu'avec les notations de (6.6.2), $L_{\bullet\bullet}^{(i)}$ se réduit
à $\Gamma(X^{(i)},\mathscr P_\bullet^{(i)})$ (considéré comme bicomplexe dont
les termes de premier degré $\ne0$ sont nuls), et l'égalité des deux membres
de (6.7.4.2) résulte de (6.4.1.1) et (6.3.1).
\end{remark}

\begin{proposition}[6.7.5]
Soit $u:\mathscr P_\bullet^{(1)}\to\mathscr Q_\bullet^{(1)}$ un
homomorphisme de complexes de $\mathscr O_{X^{(1)}}$-Modules quasi-cohérents,
limités inférieurement, tel que l'homomorphisme
\[
  \mathscr H_\bullet(u):
  \mathscr H_\bullet(\mathscr P_\bullet^{(1)})
  \longrightarrow
  \mathscr H_\bullet(\mathscr Q_\bullet^{(1)})
\]
déduit de $u$ soit un isomorphisme. Alors les homomorphismes
\[
  {}^{(t)}\mathscr E
  (f_1,f_2;\mathscr P_\bullet^{(1)},\mathscr P_\bullet^{(2)})
  \longrightarrow
  {}^{(t)}\mathscr E
  (f_1,f_2;\mathscr Q_\bullet^{(1)},\mathscr P_\bullet^{(2)})
\]
déduits de $u$ sont des isomorphismes pour $t=a$, $t=b$ et $t=c$.
\end{proposition}

\begin{proof}
L'assertion relative à la suite spectrale $(c)$ résulte de ce que cette suite
est birégulière et de ce que l'homomorphisme considéré est un isomorphisme pour
les termes $E_2$ par hypothèse ($0_{\mathrm I}$, 11.1.5). Ceci montre déjà que
\[
  \mathfrak{Tor}_\bullet^S
  (f_1,f_2;\mathscr P_\bullet^{(1)},\mathscr P_\bullet^{(2)})
  \longrightarrow
  \mathfrak{Tor}_\bullet^S
  (f_1,f_2;\mathscr Q_\bullet^{(1)},\mathscr P_\bullet^{(2)})
\]
est un isomorphisme. Appliquant les relations (6.7.4.1) et (6.7.4.2) on voit
d'abord que les homomorphismes
\[
  \mathscr H^{-n}(f_1,\mathscr P_\bullet^{(1)})
  \longrightarrow
  \mathscr H^{-n}(f_1,\mathscr Q_\bullet^{(1)})
  \quad\text{et}\quad
  \mathscr{T}\!or_n^S
  (\mathscr P_\bullet^{(1)},\mathscr P_\bullet^{(2)})
  \longrightarrow
  \mathscr{T}\!or_n^S
  (\mathscr Q_\bullet^{(1)},\mathscr P_\bullet^{(2)})
\]
déduits de $u$ sont des isomorphismes. L'assertion relative aux suites $(a)$
et $(b)$ résulte alors de ce que ces suites sont birégulières (6.7.3) et que
les homomorphismes considérés sont bijectifs pour les termes $E_2$
($0_{\mathrm I}$, 11.1.5).

Notons d'autre part que, si
$u:\mathscr P_\bullet^{(1)}\to\mathscr Q_\bullet^{(1)}$ est un
\emph{homotopisme}, on en déduit des isomorphismes canoniques
\[
  {}^{(t)}\mathscr E
  (f_1,f_2;\mathscr P_\bullet^{(1)},\mathscr P_\bullet^{(2)})
  \longrightarrow
  {}^{(t)}\mathscr E
  (f_1,f_2;\mathscr Q_\bullet^{(1)},\mathscr P_\bullet^{(2)})
\]
pour les six suites spectrales. En effet, si $S$ et les $Y^{(i)}$ sont affines,
on déduit de $u$ un homotopisme de bicomplexes
\[
  C^\bullet(\mathfrak U^{(1)},\mathscr P_\bullet^{(1)})
  \longrightarrow
  C^\bullet(\mathfrak U^{(1)},\mathscr Q_\bullet^{(1)}),
\]
et la proposition résulte de la théorie générale de l'hyperhomologie
($0_{\mathrm I}$, 11.3.2); le passage au cas général se fait par recollement,
en utilisant le fait que, d'un homotopisme de complexes, on déduit un
homotopisme de résolutions projectives de Cartan-Eilenberg de ces complexes
(M, XVII, 1.2).
\end{proof}

\begin{proposition}[6.7.6]
Supposons que le complexe $\mathscr P_\bullet^{(1)}$ ou le complexe
$\mathscr P_\bullet^{(2)}$ soit formé de Modules $S$-plats (les deux complexes
étant limités inférieurement). On a alors un isomorphisme canonique de
bi-$\partial$-foncteurs
\[
  \mathfrak{Tor}_n^S
  (f_1,f_2;\mathscr P_\bullet^{(1)},\mathscr P_\bullet^{(2)})
  \xrightarrow{\sim}
  \mathscr H^{-n}\bigl(f_1\times_S f_2,
    \mathscr P_\bullet^{(1)}\otimes_S\mathscr P_\bullet^{(2)}\bigr).
  \tag{6.7.6.1}\label{III.6.7.6.1.fr}
\]
\end{proposition}

\begin{proof}
Supposons d'abord $S$, $Y^{(1)}$ et $Y^{(2)}$ affines, de sorte qu'on est
dans la situation de (6.6.2), dont nous conservons les notations. Supposons par
exemple que $\mathscr P_\bullet^{(1)}$ soit formé de Modules $S$-plats, et
calculons l'hypertor en utilisant la remarque (6.6.6) : c'est donc l'homologie
de $L_{\bullet\bullet}^{(1)}\otimes_A M_{\bullet\bullet\bullet}^{(2)}$, où
$M_{\bullet\bullet\bullet}^{(2)}$ est une résolution projective de
Cartan-Eilenberg de $L_{\bullet\bullet}^{(2)}$, au sens de
($0_{\mathrm I}$, 11.7.1). D'autre part, les modules
$L_{ij}^{(1)}$ sont \emph{plats} sur $A$ en vertu de l'hypothèse (I, 4.15.1);
on déduit alors de ($0_{\mathrm I}$, 11.7.5) un isomorphisme canonique
\[
  \mathbf{Tor}_\bullet^S
  (\mathfrak U^{(1)},\mathfrak U^{(2)};
   \mathscr P_\bullet^{(1)},\mathscr P_\bullet^{(2)})
  \xrightarrow{\sim}
  \mathscr H_\bullet
  (L_{\bullet\bullet}^{(1)}\otimes_A L_{\bullet\bullet}^{(2)}).
  \tag{6.7.6.2}\label{III.6.7.6.2.fr}
\]
On a d'autre part un homomorphisme naturel de bicomplexes de
$L_{\bullet\bullet}^{(1)}\otimes_A L_{\bullet\bullet}^{(2)}$ dans
$C^\bullet(\mathfrak U,\mathscr Q_\bullet)$, où $\mathfrak U$ est le
recouvrement de $Z=X^{(1)}\times_S X^{(2)}$ par les ouverts affines
\oldpage[III]{161}
$U_\alpha^{(1)}\times_S U_\beta^{(2)}$ et
$\mathscr Q_\bullet
=\mathscr P_\bullet^{(1)}\otimes_S\mathscr P_\bullet^{(2)}$ (considéré comme
complexe simple pour le degré total); en effet, la définition de cet
homomorphisme a en substance été donnée au cours du calcul de la suite $(b')$
dans (6.6.4), pour $q=0$; il suffit simplement (en gardant les notations de
(6.6.4)) de tenir compte de ce qu'il y a d'une part un homomorphisme naturel
du complexe $N^{\bullet k}$ dans le complexe
$C^\bullet(\Phi^{(1)},\Phi^{(2)};\mathscr S_k)$
($0_{\mathrm I}$, 11.8.5), d'autre part un homomorphisme naturel de ce dernier
complexe dans le complexe
$P^\bullet(\Phi^{(1)},\Phi^{(2)};\mathscr S_k)$
($0_{\mathrm I}$, 11.8.6), et enfin un homomorphisme naturel de ce dernier
complexe de cochaînes dans le sous-complexe des cochaînes alternées
($0_{\mathrm I}$, 11.8.7). Par ailleurs, l'homomorphisme de bicomplexes ainsi
défini donne un isomorphisme en homologie, comme on l'a vu en (6.6.4); on a
donc, en composant avec (6.7.6.2), obtenu un isomorphisme
\[
  \mathbf{Tor}_n^S
  (\mathfrak U^{(1)},\mathfrak U^{(2)};
   \mathscr P_\bullet^{(1)},\mathscr P_\bullet^{(2)})
  \xrightarrow{\sim}
  \mathbf H^{-n}
  (\mathfrak U,\mathscr P_\bullet^{(1)}\otimes_S\mathscr P_\bullet^{(2)}).
  \tag{6.7.6.3}\label{III.6.7.6.3.fr}
\]
Il faut ensuite prouver que l'isomorphisme ainsi défini ne dépend pas des
recouvrements ouverts choisis (le second membre de (6.7.6.3) étant
canoniquement isomorphe à
$\mathbf H^{-n}
(X^{(1)}\times_S X^{(2)},
 \mathscr P_\bullet^{(1)}\otimes_S\mathscr P_\bullet^{(2)})$ par (6.2.2));
cela se fait à l'aide de (6.6.7) en remarquant (avec les notations de (6.6.7))
que l'on a un diagramme commutatif à homotopismes près
\[
\xymatrix@C=18pt@R=28pt{
  C^\bullet(\mathfrak U^{(1)},\mathscr P_\bullet^{(1)})
  \otimes_A
  C^\bullet(\mathfrak U^{(2)},\mathscr P_\bullet^{(2)})
  \ar[r]\ar[d]
  & C^\bullet(\mathfrak U,\mathscr Q_\bullet)\ar[d]
  \\
  C^\bullet(\mathfrak V^{(1)},\mathscr P_\bullet^{(1)})
  \otimes_A
  C^\bullet(\mathfrak V^{(2)},\mathscr P_\bullet^{(2)})
  \ar[r]
  & C^\bullet(\mathfrak V,\mathscr Q_\bullet)
}
\]
où les flèches horizontales sont les homomorphismes définis ci-dessus. Enfin,
il faut passer au cas général par recollement, ce qui se fait sans difficulté
comme dans (6.7.1) et (6.7.2); nous laissons les détails au lecteur.
\end{proof}

\begin{proposition}[6.7.7]
Supposons que $\mathscr P_\bullet^{(1)}$ et
$\mathscr P_\bullet^{(2)}$ soient limités inférieurement, et que tous les
Modules $\mathscr H^{-n}(f_1,\mathscr P_\bullet^{(1)})$ ou tous les Modules
$\mathscr H^{-n}(f_2,\mathscr P_\bullet^{(2)})$ soient $S$-plats. On a alors un
isomorphisme canonique de bi-$\partial$-foncteurs ($n$ parcourant $\mathbf Z$)
\[
  \mathfrak{Tor}_n^S
  (f_1,f_2;\mathscr P_\bullet^{(1)},\mathscr P_\bullet^{(2)})
  \xrightarrow{\sim}
  \bigoplus_{q_1+q_2=n}
  \mathscr H^{-q_1}(f_1,\mathscr P_\bullet^{(1)})
  \otimes_S
  \mathscr H^{-q_2}(f_2,\mathscr P_\bullet^{(2)}).
  \tag{6.7.7.1}\label{III.6.7.7.1.fr}
\]
\end{proposition}

\begin{proof}
En effet, vu (6.5.8), la suite spectrale $(a)$ de (6.7.3) est dégénérée, et
la proposition résulte aussitôt de ($0_{\mathrm I}$, 11.1.6), cette suite
étant birégulière (6.7.3).
\end{proof}

\begin{theorem}[6.7.8]
Supposons que : $1^\circ$ les complexes $\mathscr P_\bullet^{(1)}$ et
$\mathscr P_\bullet^{(2)}$ soient limités inférieurement; $2^\circ$ le
complexe $\mathscr P_\bullet^{(1)}$ ou le complexe
$\mathscr P_\bullet^{(2)}$ soit formé de Modules $S$-plats; $3^\circ$ tous les
\oldpage[III]{162}
Modules $\mathscr H^{-n}(f_1,\mathscr P_\bullet^{(1)})$ ou tous les Modules
$\mathscr H^{-n}(f_2,\mathscr P_\bullet^{(2)})$ soient $S$-plats. On a alors un
isomorphisme canonique de bi-$\partial$-foncteurs ($n$ parcourant $\mathbf Z$)
\[
  \mathscr H^n\bigl(f_1\times_S f_2,
    \mathscr P_\bullet^{(1)}\otimes_S\mathscr P_\bullet^{(2)}\bigr)
  \xrightarrow{\sim}
  \bigoplus_{n_1+n_2=n}
  \mathscr H^{n_1}(f_1,\mathscr P_\bullet^{(1)})
  \otimes_S
  \mathscr H^{n_2}(f_2,\mathscr P_\bullet^{(2)}).
  \tag{6.7.8.1}\label{III.6.7.8.1.fr}
\]
(« formule de Künneth »).
\end{theorem}

\begin{proof}
Cela résulte de (6.7.6) et (6.7.7).

Lorsque $S$, $Y^{(1)}$ et $Y^{(2)}$ sont affines, l'isomorphisme réciproque de
(6.7.8.1) se déduit (avec les notations de (6.7.6)) de l'homomorphisme de
bicomplexes
\[
  C^\bullet(\mathfrak U^{(1)},\mathscr P_\bullet^{(1)})
  \otimes_A
  C^\bullet(\mathfrak U^{(2)},\mathscr P_\bullet^{(2)})
  \longrightarrow
  C^\bullet(\mathfrak U,\mathscr Q_\bullet)
\]
par le procédé défini dans (G, I, 2.7), comme il résulte de (G, I, 5.5).
\end{proof}

\begin{proposition}[6.7.9]
Supposons vérifiées les trois conditions suivantes : $1^\circ$ $S$,
$Y^{(1)}$ et $Y^{(2)}$ sont localement noethériens, $f_1$ et $f_2$ sont
propres, $Y^{(1)}$ ou $Y^{(2)}$ de type fini sur $S$; $2^\circ$
$\mathscr P_\bullet^{(1)}$ et $\mathscr P_\bullet^{(2)}$ sont limités
inférieurement; $3^\circ$ pour tout $n\in\mathbf Z$,
$\mathscr H_n(\mathscr P_\bullet^{(i)})$ est un Module cohérent ($i=1,2$).
Dans ces conditions,
$\mathfrak{Tor}_n^S
(f_1,f_2;\mathscr P_\bullet^{(1)},\mathscr P_\bullet^{(2)})$ est un
$\mathscr O_Y$-Module cohérent (avec
$Y=Y^{(1)}\times_S Y^{(2)}$).
\end{proposition}

\begin{proof}
Il résulte de (6.5.13) que les hypertor locaux
$\mathscr{T}\!or_n^S(\mathscr P_\bullet^{(1)},\mathscr P_\bullet^{(2)})$ sont
des $\mathscr O_X$-Modules cohérents
($X=X^{(1)}\times_S X^{(2)}$ étant localement noethérien, car un des
$X^{(i)}$ est par hypothèse de type fini sur $S$ (I, 6.3.4 et 6.3.8)). Comme
$Y$ est localement noethérien et $f_1\times_S f_2$ propre (II, 5.4.2), il
résulte de (6.2.5) que les termes ${}^{(b)}\mathscr E_{pq}^2$ de (6.7.3) sont
des $\mathscr O_Y$-Modules cohérents. Comme toutes les suites spectrales de
(6.7.3) sont birégulières en vertu de l'hypothèse $2^\circ$, on conclut par
($0_{\mathrm I}$, 11.1.8).
\end{proof}

\begin{env}[6.7.10]
Soient maintenant $Y'^{(i)}$ deux $S$-préschémas,
$v_i:Y'^{(i)}\to Y^{(i)}$ deux $S$-morphismes ($i=1,2$),
$v=v_1\times_S v_2$ leur produit, qui est un $S$-morphisme $Y'\to Y$, où
l'on pose $Y'=Y'^{(1)}\times_S Y'^{(2)}$. Considérons d'autre part, pour
$i=1,2$, un $S$-préschéma $X'^{(i)}$, et deux $S$-morphismes
$u_i:X'^{(i)}\to X^{(i)}$, $f'_i:X'^{(i)}\to Y'^{(i)}$ de sorte que les
diagrammes
\[
\xymatrix@C=36pt@R=28pt{
  X'^{(i)}\ar[r]^{u_i}\ar[d]_{f'_i}
  & X^{(i)}\ar[d]^{f_i}
  \\
  Y'^{(i)}\ar[r]_{v_i}
  & Y^{(i)}
}
\tag{6.7.10.1}\label{III.6.7.10.1.fr}
\]
soient commutatifs, les morphismes $f'_i$ étant \emph{séparés et
quasi-compacts}. On a alors des $\mathscr O_{Y'}$-homomorphismes canoniques de
foncteurs spectraux
\[
  v^*\bigl({}^{(t)}\mathscr E
  (f_1,f_2;\mathscr P_\bullet^{(1)},\mathscr P_\bullet^{(2)})\bigr)
  \longrightarrow
  {}^{(t)}\mathscr E
  (f'_1,f'_2;u_1^*(\mathscr P_\bullet^{(1)}),
                 u_2^*(\mathscr P_\bullet^{(2)}))
  \tag{6.7.10.2}\label{III.6.7.10.2.fr}
\]
pour $t=a$, $a'$, $b$, $b'$, $c$, $d$. Pour les définir, supposons d'abord
$S=\operatorname{Spec}(A)$,
$Y^{(i)}=\operatorname{Spec}(B_i)$,
$Y'^{(i)}=\operatorname{Spec}(B'_i)$ affines; les $X^{(i)}$ et $X'^{(i)}$
sont alors des schémas quasi-compacts. Pour calculer les suites spectrales
${}^{(t)}\mathscr E
(f_1,f_2;\mathscr P_\bullet^{(1)},\mathscr P_\bullet^{(2)})$, nous
considérons comme dans (6.6.1) des recouvrements finis
$\mathfrak U^{(i)}$ par des ouverts affines de $X^{(i)}$ ($i=1,2$); pour
calculer
${}^{(t)}\mathscr E
(f'_1,f'_2;u_1^*(\mathscr P_\bullet^{(1)}),
u_2^*(\mathscr P_\bullet^{(2)}))$, nous considérons des recouvrements finis
$\mathfrak U'^{(i)}$

\oldpage[III]{163}
de $X'^{(i)}$ par des ouverts affines, \emph{plus fins} respectivement que les
recouvrements $u_i^{-1}(\mathfrak U^{(i)})$ ($i=1,2$). Il est clair que le
bicomplexe
$C^\bullet(\mathfrak U^{(i)},\mathscr P_\bullet^{(i)})
=L_{\bullet\bullet}^{(i)}$ peut être considéré canoniquement comme un
sous-bicomplexe de
$C^\bullet(u_i^{-1}(\mathfrak U^{(i)}),u_i^*(\mathscr P_\bullet^{(i)}))$
($0_{\mathrm I}$, 4.4.3.2); en outre, en choisissant une application
simpliciale (G, II, 5.7) de $\mathfrak U'^{(i)}$ dans
$u_i^{-1}(\mathfrak U^{(i)})$, on définit un homomorphisme de bicomplexes
\[
  C^\bullet(u_i^{-1}(\mathfrak U^{(i)}),u_i^*(\mathscr P_\bullet^{(i)}))
  \longrightarrow
  C^\bullet(\mathfrak U'^{(i)},u_i^*(\mathscr P_\bullet^{(i)}))
\]
d'où, par composition, un homomorphisme de bicomplexes
\[
  L_{\bullet\bullet}^{(i)}
  \longrightarrow
  L_{\bullet\bullet}^{\prime(i)}
  =C^\bullet(\mathfrak U'^{(i)},u_i^*(\mathscr P_\bullet^{(i)})).
\]
En outre, cet homomorphisme est remplacé par un homomorphisme homotope quand
on change d'application simpliciale (G, II, 5.7.1); on a ainsi un
homomorphisme bien défini de foncteurs spectraux :
\[
  {}^{(t)}\mathscr E
  (\mathfrak U^{(1)},\mathfrak U^{(2)};
   \mathscr P_\bullet^{(1)},\mathscr P_\bullet^{(2)})
  \longrightarrow
  {}^{(t)}\mathscr E
  (\mathfrak U'^{(1)},\mathfrak U'^{(2)};
   u_1^*(\mathscr P_\bullet^{(1)}),u_2^*(\mathscr P_\bullet^{(2)})).
  \tag{6.7.10.3}\label{III.6.7.10.3.fr}
\]
On vérifie aussitôt que si $\mathfrak V^{(i)}$ est un recouvrement affine
fini de $X^{(i)}$ plus fin que $\mathfrak U^{(i)}$,
$\mathfrak V'^{(i)}$ un recouvrement affine fini de $X'^{(i)}$, plus fin que
$u_i^{-1}(\mathfrak V^{(i)})$ et que $\mathfrak U'^{(i)}$, le diagramme
\[
\xymatrix@C=34pt@R=30pt{
  C^\bullet(\mathfrak U^{(i)},\mathscr P_\bullet^{(i)})
  \ar[r]\ar[d]
  & C^\bullet(\mathfrak V^{(i)},\mathscr P_\bullet^{(i)})\ar[d]
  \\
  C^\bullet(\mathfrak U'^{(i)},u_i^*(\mathscr P_\bullet^{(i)}))
  \ar[r]
  & C^\bullet(\mathfrak V'^{(i)},u_i^*(\mathscr P_\bullet^{(i)}))
}
\]
est commutatif, ce qui implique que l'homomorphisme (6.7.10.3) ne dépend pas
essentiellement des recouvrements $\mathfrak U^{(i)}$ et
$\mathfrak U'^{(i)}$ considérés. On a donc en fait défini un homomorphisme de
$A$-modules
\[
  {}^{(t)}E
  (X^{(1)},X^{(2)};\mathscr P_\bullet^{(1)},\mathscr P_\bullet^{(2)})
  \longrightarrow
  {}^{(t)}E
  (X'^{(1)},X'^{(2)};
   u_1^*(\mathscr P_\bullet^{(1)}),u_2^*(\mathscr P_\bullet^{(2)})).
  \tag{6.7.10.4}\label{III.6.7.10.4.fr}
\]
mais il est clair par définition des $u_i^*(\mathscr P_\bullet^{(i)})$ et en
vertu de la commutativité de (6.7.10.1) que cet homomorphisme est aussi un
homomorphisme de $(B_1\otimes_A B_2)$-modules; comme le second membre de
(6.7.10.4) est formé de $(B'_1\otimes_A B'_2)$-modules, on déduit
canoniquement de (6.7.10.4) un homomorphisme de
$(B'_1\otimes_A B'_2)$-modules
\[
  {}^{(t)}E
  (X^{(1)},X^{(2)};\mathscr P_\bullet^{(1)},\mathscr P_\bullet^{(2)})
  \otimes_{B_1\otimes_A B_2}(B'_1\otimes_A B'_2)
  \longrightarrow
  {}^{(t)}E
  (X'^{(1)},X'^{(2)};
   u_1^*(\mathscr P_\bullet^{(1)}),u_2^*(\mathscr P_\bullet^{(2)}))
  \tag{6.7.10.5}\label{III.6.7.10.5.fr}
\]
ce qui, compte tenu de (I, 1.6.5) n'est autre que l'homomorphisme cherché
(6.7.10.2) dans le cas particulier considéré.

Il reste à passer au cas général en suivant les recollements de (6.7.1) et
(6.7.2); le second passage est immédiat; en ce qui concerne le premier, on
considère comme dans (6.7.1) des éléments $g_i\in B_i$, et leurs images
$g'_i\in B'_i$, le produit tensoriel $g=g_1\otimes g_2$ dans
$B=B_1\otimes_A B_2$ et son image $g'=g'_1\otimes g'_2$ dans
$B'=B'_1\otimes_A B'_2$, et tout revient à utiliser

\oldpage[III]{164}
l'isomorphisme canonique
\[
  (M\otimes_B B')_{g'}
  \xrightarrow{\sim}
  M_g\otimes_{B_g}B'_{g'}
\]
($0_{\mathrm I}$, 1.5.4); nous laissons les détails au lecteur.
\end{env}

\begin{env}[6.7.11]
La théorie des hypertor globaux, développée ci-dessus pour deux
$S$-morphismes $X^{(i)}\to Y^{(i)}$ et deux complexes
$\mathscr P_\bullet^{(i)}$ de Modules quasi-cohérents limités inférieurement,
s'étend aussitôt au cas général suivant : on a un préschéma $S$, une famille
finie de $S$-préschémas $Y^{(i)}$ ($i\in I$), une famille finie de
$S$-morphismes séparés et quasi-compacts
$f_i:X^{(i)}\to Y^{(i)}$, et pour chaque $i$ un complexe de
$\mathscr O_{X^{(i)}}$-Modules quasi-cohérents
$\mathscr P_\bullet^{(i)}$ limités inférieurement. Si $Y$ est le produit des
$S$-préschémas $Y^{(i)}$, on définit alors pour chaque entier
$n\in\mathbf Z$, un $\mathscr O_Y$-Module quasi-cohérent
$\mathfrak{Tor}_n^S((f_i)_{i\in I};
(\mathscr P_\bullet^{(i)})_{i\in I})$,
ces Modules formant un $\partial$-foncteur covariant en chacun des complexes
$\mathscr P_\bullet^{(i)}$; en outre, ce foncteur est l'aboutissement commun
de six foncteurs spectraux
${}^{(t)}\mathscr E((f_i)_{i\in I};(\mathscr P_\bullet^{(i)})_{i\in I})$.
Nous laissons au lecteur le soin de répéter pour ce cas général les
définitions et les raisonnements faits ci-dessus pour $I=\{1,2\}$. Notons
simplement que lorsque $I$ se réduit à un seul élément, on retrouve
l'hypercohomologie $\mathscr H^\bullet(f,\mathscr P_\bullet)$ définie dans
(6.2.7) (comme on l'a déjà observé dans (6.7.4)). Lorsque $I$ est l'intervalle
$1\leq i\leq m$ de $\mathbf N$, nous écrirons
\[
  \mathfrak{Tor}_n^S
  (f_1,\ldots,f_m;
   \mathscr P_\bullet^{(1)},\ldots,\mathscr P_\bullet^{(m)})
  \quad\text{pour}\quad
  \mathfrak{Tor}_n^S
  ((f_i)_{i\in I};(\mathscr P_\bullet^{(i)})_{i\in I}).
\]
\end{env}

\begin{proposition}[6.7.12]
Les notations étant celles de (6.7.11), soit $J$ une partie de $I$ telle que,
pour $i\in I-J$, on ait $X^{(i)}=Y^{(i)}=S$, $f_i$ étant réduit à l'identité,
et $\mathscr P_\bullet^{(i)}$ égal au complexe réduit au terme de degré $0$
égal à $\mathscr O_S$. Il y a alors un isomorphisme canonique de
$\partial$-foncteurs
\[
  \mathfrak{Tor}_\bullet^S
  ((f_i)_{i\in I};(\mathscr P_\bullet^{(i)})_{i\in I})
  \xrightarrow{\sim}
  \mathfrak{Tor}_\bullet^S
  ((f_i)_{i\in J};(\mathscr P_\bullet^{(i)})_{i\in J}).
  \tag{6.7.12.1}\label{III.6.7.12.1.fr}
\]
\end{proposition}

\begin{proof}
On peut se borner à définir cet isomorphisme lorsque $S$ et les $Y^{(i)}$
sont affines, le recollement se faisant comme d'ordinaire. Pour $i\in I-J$,
on peut prendre le recouvrement $\mathfrak U^{(i)}$ formé du seul ensemble
$S$, et alors
$L_{\bullet\bullet}^{(i)}=C^\bullet(\mathfrak U^{(i)},
\mathscr P_\bullet^{(i)})$ est réduit à son seul terme de degrés $(0,0)$,
égal à $\Gamma(S,\mathscr O_S)=A(S)$; l'isomorphisme (6.7.12.1) est alors
évident.
\end{proof}

\begin{remark}[6.7.13]
Les notations étant celles de (6.7.3), considérons le $S$-isomorphisme
canonique
$Y^{(1)}\times_S Y^{(2)}\to Y^{(2)}\times_S Y^{(1)}$ (I, 3.3.5); alors
l'image par cet isomorphisme de
$\mathfrak{Tor}_\bullet^S(f_1,f_2;
\mathscr P_\bullet^{(1)},\mathscr P_\bullet^{(2)})$ est
$\mathfrak{Tor}_\bullet^S(f_2,f_1;
\mathscr P_\bullet^{(2)},\mathscr P_\bullet^{(1)})$; la question étant
locale, on est ramené au cas envisagé dans (6.6.2), et si on désigne par
$M_{\bullet\bullet\bullet}^{(i)}$ une résolution projective de
Cartan-Eilenberg de $L_{\bullet\bullet}^{(i)}$ ($i=1,2$), l'isomorphisme
considéré transforme
$M_{\bullet\bullet\bullet}^{(1)}\otimes_A
M_{\bullet\bullet\bullet}^{(2)}$ en
$M_{\bullet\bullet\bullet}^{(2)}\otimes_A
M_{\bullet\bullet\bullet}^{(1)}$, d'où notre assertion en considérant
l'homologie des complexes simples associés à ces tricomplexes.
\end{remark}

\subsection{Les suites spectrales d'associativité des hypertor globaux.}

\begin{env}[6.8.1]
Les hypothèses et notations étant celles de (6.7.11) (et en particulier les
$\mathscr P_\bullet^{(i)}$ étant supposés limités inférieurement, supposons
donnée une partition $(I_j)_{j\in J}$ de l'ensemble d'indices $I$; nous nous
proposons de donner une relation d'« associativité » entre les hypertor
$\mathfrak{Tor}_n^S((f_i)_{i\in I};(\mathscr P_\bullet^{(i)})_{i\in I})$ et
chacun des hypertor « partiels »
\[
  \mathcal T_{\bullet j}
  =\mathfrak{Tor}_\bullet^S
  ((f_i)_{i\in I_j};(\mathscr P_\bullet^{(i)})_{i\in I_j}).
\]

\oldpage[III]{165}
Pour simplifier l'écriture, nous nous bornerons au cas où $I$ est
l'intervalle $1\leq i\leq m$, et où la partition $(I_j)$ se compose des deux
intervalles $\{1,2,\ldots,r\}$ et $\{r+1,\ldots,m\}$.
\end{env}

\begin{proposition}[6.8.2]
Il existe un foncteur spectral canonique birégulier (dit « foncteur spectral
d'associativité ») noté
\[
  {}^{(e)}\mathscr E^S
  (f_1,\ldots,f_m;
   \mathscr P_\bullet^{(1)},\ldots,\mathscr P_\bullet^{(m)})
\]
(ou simplement ${}^{(e)}\mathscr E(f_1,\ldots,f_m;
\mathscr P_\bullet^{(1)},\ldots,\mathscr P_\bullet^{(m)})$) dont
l'aboutissement est
$\mathfrak{Tor}_\bullet^S(f_1,\ldots,f_m;
\mathscr P_\bullet^{(1)},\ldots,\mathscr P_\bullet^{(m)})$, et dont le terme
$E_2$ est donné par
\[
\begin{split}
  {}^{(e)}\mathscr E_{pq}^2
  =\bigoplus_{q_1+q_2=q}
  \mathscr{T}\!or_p^S\Bigl(
  &\mathfrak{Tor}_{q_1}^S
  (f_1,\ldots,f_r;
   \mathscr P_\bullet^{(1)},\ldots,\mathscr P_\bullet^{(r)}),\\
  &\mathfrak{Tor}_{q_2}^S
  (f_{r+1},\ldots,f_m;
   \mathscr P_\bullet^{(r+1)},\ldots,\mathscr P_\bullet^{(m)})
  \Bigr).
\end{split}
\]
\end{proposition}

\begin{proof}
Dans cet énoncé, on a identifié canoniquement $Y$ au produit
$Z^{(1)}\times_S Z^{(2)}$, où
$Z^{(1)}=Y^{(1)}\times_S Y^{(2)}\times_S\cdots\times_S Y^{(r)}$, et
$Z^{(2)}=Y^{(r+1)}\times_S\cdots\times_S Y^{(m)}$. Nous nous bornerons au
cas où $S$ et les $Y^{(i)}$ sont affines; on passe de ce cas particulier au
cas général par les méthodes développées dans (6.7.1) et (6.7.2), et nous
laissons les détails du raisonnement (sans difficulté) au lecteur. Nous
démontrerons donc le corollaire suivant.
\end{proof}

\begin{corollary}[6.8.3]
Soient $A$ un anneau, $S=\operatorname{Spec}(A)$, $X^{(i)}$
($1\leq i\leq m$) des $S$-schémas quasi-compacts et, pour chaque $i$, soit
$\mathscr P_\bullet^{(i)}$ un complexe de
$\mathscr O_{X^{(i)}}$-Modules quasi-cohérents limité inférieurement. Il
existe un foncteur spectral canonique birégulier ayant pour aboutissement
\[
  \mathbf{Tor}_\bullet^S
  (X^{(1)},\ldots,X^{(m)};
   \mathscr P_\bullet^{(1)},\ldots,\mathscr P_\bullet^{(m)})
\]
et dont le terme $E_2$ est donné par
\[
\begin{split}
  {}^{(e)}E_{pq}^2
  =\bigoplus_{q_1+q_2=q}
  \operatorname{Tor}_p^A\Bigl(
  &\mathbf{Tor}_{q_1}^S
  (X^{(1)},\ldots,X^{(r)};
   \mathscr P_\bullet^{(1)},\ldots,\mathscr P_\bullet^{(r)}),\\
  &\mathbf{Tor}_{q_2}^S
  (X^{(r+1)},\ldots,X^{(m)};
   \mathscr P_\bullet^{(r+1)},\ldots,\mathscr P_\bullet^{(m)})
  \Bigr).
\end{split}
\]
\end{corollary}

\begin{proof}
Suivant la définition donnée en (6.6.2), le calcul de l'hypertor considéré se
fait en prenant pour chaque $i$ un recouvrement ouvert affine fini
$\mathfrak U^{(i)}$ de $X^{(i)}$, en considérant les bicomplexes
$L_{\bullet\bullet}^{(i)}=C^\bullet(\mathfrak U^{(i)},
\mathscr P_\bullet^{(i)})$, une résolution projective
$M_{\bullet\bullet\bullet}^{(i)}$ de Cartan-Eilenberg de chacun de ces
bicomplexes (au sens de (0, 11.7.1)), le produit tensoriel
\[
  M_{\bullet\bullet\bullet}
  =\bigotimes_{i=1}^m M_{\bullet\bullet\bullet}^{(i)}
\]
de ces tricomplexes, et en prenant l'homologie de
$M_{\bullet\bullet\bullet}$. Considérons $M_{\bullet\bullet\bullet}$ comme
un complexe simple $N_\bullet$, produit tensoriel des deux complexes simples
\[
  N'_\bullet=\bigotimes_{i=1}^r M_{\bullet\bullet\bullet}^{(i)},
  \qquad
  N''_\bullet=\bigotimes_{i=r+1}^m M_{\bullet\bullet\bullet}^{(i)},
\]
où $N'_\bullet$ et $N''_\bullet$ sont gradués par la somme des degrés totaux
des $M_{\bullet\bullet\bullet}^{(i)}$. En outre, les $A$-modules des
complexes $N'_\bullet$ et $N''_\bullet$ sont projectifs, donc il résulte de
(6.5.9) que l'on a
\[
  H_\bullet(M_{\bullet\bullet\bullet})
  =\mathbf{Tor}_\bullet^A(N'_\bullet,N''_\bullet);
\]
la suite spectrale cherchée n'est autre alors que la suite (6.3.2.2)
appliquée aux complexes $N'_\bullet$ et $N''_\bullet$, compte tenu de
l'interprétation des modules d'homologie de ces complexes qui résulte de ce
qui précède (quand on applique les remarques du début à chacun des produits
partiels
$X^{(1)}\times_S\cdots\times_S X^{(r)}$ et
$X^{(r+1)}\times_S\cdots\times_S X^{(m)}$). Enfin, les propriétés de
régularité résultent de (6.3.2) et du fait que, les
$\mathscr P_\bullet^{(i)}$ étant limités inférieurement, il en est de même
des $M_{\bullet\bullet\bullet}^{(i)}$; par suite, $N'_\bullet$ et
$N''_\bullet$ sont limités inférieurement.
\end{proof}

\oldpage[III]{166}
\subsection{Les suites spectrales de changement de base dans les hypertor
globaux.}

\begin{env}[6.9.1]
Les hypothèses et notations étant toujours celles de (6.7.11) (et en
particulier les $\mathscr P_\bullet^{(i)}$ étant supposés limités
inférieurement), considérons un morphisme $g:S'\to S$ de préschémas, et
posons
$Y'^{(i)}=Y_{(S')}^{(i)}$, $X'^{(i)}=X_{(S')}^{(i)}$ et
$\mathscr P_\bullet'^{(i)}=
\mathscr P_\bullet^{(i)}\otimes_{\mathscr O_S}\mathscr O_{S'}$,
$\mathscr P_\bullet'^{(i)}$ étant donc un complexe de
$\mathscr O_{X'^{(i)}}$-Modules quasi-cohérents; soit
$f'_i=(f_i)_{(S')}:X'^{(i)}\to Y'^{(i)}$, qui est un morphisme séparé et
quasi-compact (I, 5.5.1 et 6.6.4). Nous nous proposons d'étudier les relations
entre les $\mathscr O_{Y'}$-Modules quasi-cohérents
$\mathfrak{Tor}_n^{S'}((f'_i)_{i\in I};
(\mathscr P_\bullet'^{(i)})_{i\in I})$ et
$\mathfrak{Tor}_n^S((f_i)_{i\in I};
(\mathscr P_\bullet^{(i)})_{i\in I})
\otimes_{\mathscr O_S}\mathscr O_{S'}$, où
$Y'=Y\times_S S'=Y_{(S')}$. Un cas particulièrement simple est le suivant,
qui se réduit à (1.4.15) lorsque $I$ se réduit à un seul élément et
$\mathscr P_\bullet$ à un seul module :
\end{env}

\begin{proposition}[6.9.2]
Si le morphisme $g:S'\to S$ est plat, on a un isomorphisme canonique de
$\partial$-foncteurs (en les $\mathscr P_\bullet^{(i)}$) :
\[
  \mathfrak{Tor}_\bullet^S
  ((f_i)_{i\in I};(\mathscr P_\bullet^{(i)})_{i\in I})
  \otimes_{\mathscr O_S}\mathscr O_{S'}
  \xrightarrow{\sim}
  \mathfrak{Tor}_\bullet^{S'}
  ((f'_i)_{i\in I};(\mathscr P_\bullet'^{(i)})_{i\in I}).
  \tag{6.9.2.1}\label{III.6.9.2.1.fr}
\]
\end{proposition}

\begin{proof}
On peut encore se borner au cas où $S$, $S'$ et les $Y^{(i)}$ sont affines,
le recollement se faisant suivant les méthodes de (6.7.1) et (6.7.2). Soient
$S=\operatorname{Spec}(A)$, $S'=\operatorname{Spec}(A')$, et prenons pour
chaque $i$ un recouvrement ouvert affine $\mathfrak U^{(i)}$ de $X^{(i)}$;
si $u_i:X'^{(i)}\to X^{(i)}$ est la projection canonique,
$u_i^{-1}(\mathfrak U^{(i)})$ est un recouvrement ouvert affine de
$X'^{(i)}$ (II, 1.5.5), que nous noterons $\mathfrak U'^{(i)}$; il est clair
alors que
\[
  C^\bullet(\mathfrak U'^{(i)},\mathscr P_\bullet'^{(i)})
  =C^\bullet(\mathfrak U^{(i)},\mathscr P_\bullet^{(i)})\otimes_A A',
\]
et l'existence de l'isomorphisme (6.9.2.1) est immédiate, car si
$M_{\bullet\bullet\bullet}^{(i)}$ est une résolution projective de
Cartan-Eilenberg de
$L_{\bullet\bullet}^{(i)}=C^\bullet(\mathfrak U^{(i)},
\mathscr P_\bullet^{(i)})$ au sens de (0, 11.7.1), formée de $A$-modules,
$M_{\bullet\bullet\bullet}^{(i)}\otimes_A A'$ est une résolution projective
de Cartan-Eilenberg (au même sens) de
$L_{\bullet\bullet}^{(i)}\otimes_A A'$ formée de $A'$-modules, en vertu de
l'hypothèse que $A'$ est un $A$-module plat; cette même hypothèse montre en
outre que
\[
  H_\bullet\left(
    \bigotimes_{i=1}^m
    (M_{\bullet\bullet\bullet}^{(i)}\otimes_A A')
  \right)
  =H_\bullet\left(
    \bigotimes_{i=1}^m M_{\bullet\bullet\bullet}^{(i)}
  \right)\otimes_A A'.
\]
\end{proof}

\begin{env}[6.9.2.2]
On notera que lorsque $I$ est réduit au seul élément $1$, la formule
(6.9.2.1) se déduit directement de (6.7.7), appliqué en prenant
$Y_2=X_2=S'$, $f_2=1_{S'}$, et le complexe $\mathscr P_\bullet^{(2)}$
réduit à son terme de degré $0$, égal à $\mathscr O_{S'}$; on sait alors que
l'hypercohomologie $\mathscr H^n(1_{S'},\mathscr O_{S'})$ est nulle pour tout
$n\ne0$ et se réduit à $\mathscr O_{S'}$ pour $n=0$ (6.2.1).
\end{env}

Dans le cas général, nous allons introduire à la place de $\mathscr O_{S'}$
un complexe $\mathscr Q'_\bullet$ de $\mathscr O_{S'}$-Modules
quasi-cohérents limités inférieurement, de sorte que si, pour simplifier, on
prend $I=\{1,2,\ldots,m\}$, on peut considérer le $\partial$-foncteur
\[
  \mathfrak{Tor}_\bullet^S
  (f_1,\ldots,f_m,1_{S'};
   \mathscr P_\bullet^{(1)},\ldots,\mathscr P_\bullet^{(m)},
   \mathscr Q'_\bullet).
\]

\begin{proposition}[6.9.3]
Il existe trois foncteurs spectraux canoniques biréguliers notés
${}^{(t)}\mathscr E(f_1,\ldots,f_m;
\mathscr P_\bullet^{(1)},\ldots,\mathscr P_\bullet^{(m)},
\mathscr Q'_\bullet)$ (avec $t=e$, $f$ ou $f'$) ayant pour aboutissement
commun
$\mathfrak{Tor}_\bullet^S(f_1,\ldots,f_m,1_{S'};
\mathscr P_\bullet^{(1)},\ldots,\mathscr P_\bullet^{(m)},
\mathscr Q'_\bullet)$ et dont les termes $E_2$ sont respectivement

\oldpage[III]{167}
\[
\begin{split}
  {}^{(e)}\mathscr E_{pq}^2
  =\bigoplus_{q'+q''=q}
  \mathscr{T}\!or_p^S\Bigl(
  &\mathfrak{Tor}_{q'}^S
  (f_1,\ldots,f_m;
   \mathscr P_\bullet^{(1)},\ldots,\mathscr P_\bullet^{(m)}),\\
  &\mathscr H_{q''}(\mathscr Q'_\bullet)
  \Bigr),
\end{split}
\]
\[
\begin{split}
  {}^{(f)}\mathscr E_{pq}^2
  =\bigoplus_{q_1+q_2+\cdots+q_{m+1}=q}
  \mathscr{T}\!or_p^{S'}\Bigl(
  &\mathfrak{Tor}_{q_1}^S
  (f_1,1_{S'};\mathscr P_\bullet^{(1)},\mathscr O_{S'}),\ldots,\\
  &\mathfrak{Tor}_{q_m}^S
  (f_m,1_{S'};\mathscr P_\bullet^{(m)},\mathscr O_{S'}),
  \mathscr H_{q_{m+1}}(\mathscr Q'_\bullet)
  \Bigr),
\end{split}
\]
\[
\begin{split}
  {}^{(f')}\mathscr E_{pq}^2
  =\bigoplus_{q_1+\cdots+q_m=q}
  \mathfrak{Tor}_p^{S'}\Bigl(
  &f'_1,\ldots,f'_m,1_{S'};\\
  &\mathscr{T}\!or_{q_1}^S
  (\mathscr P_\bullet^{(1)},\mathscr O_{S'}),\ldots,
  \mathscr{T}\!or_{q_m}^S
  (\mathscr P_\bullet^{(m)},\mathscr O_{S'}),
  \mathscr Q'_\bullet
  \Bigr).
\end{split}
\]
\end{proposition}

\begin{proof}
La suite $(e)$ n'est autre que la suite d'associativité de (6.8.2) pour
$r=m$. Pour définir les deux autres suites spectrales, on va encore se limiter
au cas où $S$, $S'$ et les $Y^{(i)}$ sont affines, le passage au cas général
se faisant par les méthodes de (6.7.1) et (6.7.2) et étant laissé au lecteur.
Nous démontrerons donc le
\end{proof}

\begin{corollary}[6.9.4]
Soient $A$ un anneau, $A'$ une $A$-algèbre,
$S=\operatorname{Spec}(A)$, $S'=\operatorname{Spec}(A')$,
$X^{(i)}$ ($1\leq i\leq m$) des $S$-schémas quasi-compacts et, pour chaque
$i$, soit $X'^{(i)}=X_{(S')}^{(i)}$, qui est un $S'$-schéma quasi-compact.
Pour chaque $i$, soit $\mathscr P_\bullet^{(i)}$ un complexe de
$\mathscr O_{X^{(i)}}$-Modules quasi-cohérents; soit enfin $Q'_\bullet$ un
complexe de $A'$-modules, ces complexes étant limités inférieurement. Il
existe trois foncteurs spectraux biréguliers en les
$\mathscr P_\bullet^{(i)}$ et en $Q'_\bullet$, ayant pour aboutissement commun
\[
  \mathbf{Tor}_\bullet^S
  (X^{(1)},\ldots,X^{(m)},S';
   \mathscr P_\bullet^{(1)},\ldots,\mathscr P_\bullet^{(m)},
   \widetilde{Q'_\bullet})
\]
et dont les termes $E_2$ sont respectivement
\[
\begin{split}
  {}^{(e)}E_{pq}^2
  =\bigoplus_{q'+q''=q}
  \operatorname{Tor}_p^A\Bigl(
  &\mathbf{Tor}_{q'}^S
  (X^{(1)},\ldots,X^{(m)};
   \mathscr P_\bullet^{(1)},\ldots,\mathscr P_\bullet^{(m)}),\\
  &H_{q''}(Q'_\bullet)
  \Bigr),
\end{split}
\]
\[
\begin{split}
  {}^{(f)}E_{pq}^2
  =\bigoplus_{q_1+q_2+\cdots+q_{m+1}=q}
  \operatorname{Tor}_p^{A'}\Bigl(
  &\mathbf{Tor}_{q_1}^S
  (X^{(1)},S';\mathscr P_\bullet^{(1)},\mathscr O_{S'}),\ldots,\\
  &\mathbf{Tor}_{q_m}^S
  (X^{(m)},S';\mathscr P_\bullet^{(m)},\mathscr O_{S'}),
  H_{q_{m+1}}(Q'_\bullet)
  \Bigr),
\end{split}
\]
\[
\begin{split}
  {}^{(f')}E_{pq}^2
  =\bigoplus_{q_1+\cdots+q_m=q}
  \mathbf{Tor}_p^{S'}\Bigl(
  &X'^{(1)},\ldots,X'^{(m)},S';\\
  &\mathscr{T}\!or_{q_1}^S
  (\mathscr P_\bullet^{(1)},\mathscr O_{S'}),\ldots,
  \mathscr{T}\!or_{q_m}^S
  (\mathscr P_\bullet^{(m)},\mathscr O_{S'}),
  \widetilde{Q'_\bullet}
  \Bigr).
\end{split}
\]
\end{corollary}

\begin{proof}
Nous ne reviendrons pas sur le premier de ces foncteurs spectraux, qui a été
traité dans (6.8.3) et n'est inclus ici que pour mémoire. Pour définir les
autres, considérons pour chaque $i$ un recouvrement ouvert affine fini
$\mathfrak U^{(i)}$ de $X^{(i)}$ et, si
$u_i:X'^{(i)}\to X^{(i)}$ est la projection canonique, le recouvrement ouvert
affine fini correspondant
$\mathfrak U'^{(i)}=u_i^{-1}(\mathfrak U^{(i)})$.

En vertu de (6.6.6),
$\mathbf{Tor}_\bullet^S(X^{(1)},\ldots,X^{(m)},S';
\mathscr P_\bullet^{(1)},\ldots,\mathscr P_\bullet^{(m)},
\widetilde{Q'_\bullet})$ s'obtient en prenant pour $1\leq i\leq m$ une
résolution projective de Cartan-Eilenberg
$M_{\bullet\bullet\bullet}^{(i)}$ de
$L_{\bullet\bullet}^{(i)}=C^\bullet(\mathfrak U^{(i)},
\mathscr P_\bullet^{(i)})$ (au sens de (0, 11.7.1)), considérant le
tricomplexe
\[
  M_{\bullet\bullet\bullet}
  =M_{\bullet\bullet\bullet}^{(1)}\otimes_A
   M_{\bullet\bullet\bullet}^{(2)}\otimes_A\cdots\otimes_A
   M_{\bullet\bullet\bullet}^{(m)}\otimes_A Q'_\bullet
\]
(où $Q'_\bullet$ est considéré comme un tricomplexe dont les deux derniers
degrés se réduisent à $0$), et en en prenant l'homologie. Si l'on pose
$M_{\bullet\bullet\bullet}'^{(i)}=
M_{\bullet\bullet\bullet}^{(i)}\otimes_A A'$, on a (en se rappelant que
$Q'_\bullet$ est un complexe de $A'$-modules)
\[
  M_{\bullet\bullet\bullet}
  =M_{\bullet\bullet\bullet}'^{(1)}\otimes_{A'}
   M_{\bullet\bullet\bullet}'^{(2)}\otimes_{A'}\cdots\otimes_{A'}
   M_{\bullet\bullet\bullet}'^{(m)}\otimes_{A'} Q'_\bullet.
\]
Or, considérons chacun des complexes
$M_{\bullet\bullet\bullet}'^{(i)}$ comme un complexe simple (pour son degré
total) et notons que ce complexe est formé de $A'$-modules projectifs; il
résulte de (6.3.7) (étendu à un nombre quelconque de complexes) que
$H_\bullet(M_{\bullet\bullet\bullet})$ est aussi égal à
\[
  \operatorname{Tor}_\bullet^{A'}
  (M_{\bullet\bullet\bullet}'^{(1)},\ldots,
   M_{\bullet\bullet\bullet}'^{(m)},Q'_\bullet);
\]
c'est donc (6.5.15) l'aboutissement d'une suite spectrale ayant les propriétés
de régularité voulues (les trois degrés de
$M_{\bullet\bullet\bullet}^{(i)}$ étant limités inférieurement lorsque
$\mathscr P_\bullet^{(i)}$ est limité inférieurement) et dont le terme $E_2$
est donné par
\[
  E_{pq}^2
  =\bigoplus_{q_1+\cdots+q_{m+1}=q}
  \operatorname{Tor}_p^{A'}
  (H_{q_1}(M_{\bullet\bullet\bullet}'^{(1)}),\ldots,
   H_{q_m}(M_{\bullet\bullet\bullet}'^{(m)}),H_{q_{m+1}}(Q'_\bullet)).
\]

\oldpage[III]{168}
On a d'ailleurs
\[
  H_{q_i}(M_{\bullet\bullet\bullet}'^{(i)})
  =H_{q_i}(M_{\bullet\bullet\bullet}^{(i)}\otimes_A A')
  =\mathbf{Tor}_{q_i}^S
  (X^{(i)},S';\mathscr P_\bullet^{(i)},\mathscr O_{S'})
\]
en vertu de la définition des hypertor globaux, ce qui donne la suite $(f)$
cherchée. On peut d'autre part considérer
$M_{\bullet\bullet\bullet}'^{(i)}$ comme un bicomplexe dans lequel le premier
degré est la somme du premier et du second degré du tricomplexe
$M_{\bullet\bullet\bullet}^{(i)}$, le second degré étant le troisième degré
de ce tricomplexe; comme les modules formant les
$M_{\bullet\bullet\bullet}'^{(i)}$ sont des $A'$-modules projectifs, la théorie
générale de l'hyperhomologie montre que l'homologie du bicomplexe
\[
  M_{\bullet\bullet\bullet}'^{(1)}\otimes_{A'}
  M_{\bullet\bullet\bullet}'^{(2)}\otimes_{A'}\cdots\otimes_{A'}
  M_{\bullet\bullet\bullet}'^{(m)}\otimes_{A'} Q'_\bullet
\]
est canoniquement isomorphe à son hyperhomologie (0, 11.6.5); c'est donc
l'aboutissement d'une suite spectrale de terme $E_2$ égal à
\[
  E_{pq}^2
  =\bigoplus_{q_1+\cdots+q_{m+1}=q}
  \operatorname{Tor}_p^{A'}
  (H_{q_1}^{\mathrm{II}}(M_{\bullet\bullet\bullet}'^{(1)}),\ldots,
   H_{q_m}^{\mathrm{II}}(M_{\bullet\bullet\bullet}'^{(m)}),
   H_{q_{m+1}}^{\mathrm{II}}(Q'_\bullet)).
\]
Or, comme le second degré de $Q'_\bullet$ se réduit à $0$, on a
$H_n^{\mathrm{II}}(Q'_\bullet)=0$ pour $n\ne0$ et
$H_0^{\mathrm{II}}(Q'_\bullet)=Q'_\bullet$; la formule précédente s'écrit
aussi
\[
  E_{pq}^2
  =\bigoplus_{q_1+\cdots+q_m=q}
  \operatorname{Tor}_p^{A'}
  (H_{q_1}^{\mathrm{II}}(M_{\bullet\bullet\bullet}'^{(1)}),\ldots,
   H_{q_m}^{\mathrm{II}}(M_{\bullet\bullet\bullet}'^{(m)}),Q'_\bullet).
\]
En outre, on a
\[
  H_{q_i}^{\mathrm{II}}(M_{\bullet\bullet\bullet}'^{(i)})
  =H_{q_i}^{\mathrm{II}}(M_{\bullet\bullet\bullet}^{(i)}\otimes_A A')
  =\operatorname{Tor}_{q_i}^A(L_{\bullet\bullet}^{(i)},A')
\]
en vertu de (6.3.4); mais
$L_{-j,k}^{(i)}=C^j(\mathfrak U^{(i)},\mathscr P_k^{(i)})$, somme directe des
$\Gamma(V,\mathscr P_k^{(i)})$, où $V$ parcourt les intersections (affines) de
$j+1$ ensembles du recouvrement $\mathfrak U^{(i)}$; si
$V'=u_i^{-1}(V)$, $V'$ est affine dans $X'^{(i)}$ et il résulte de (6.4.1.1)
que l'on a
\[
  \Gamma\bigl(V',\mathscr{T}\!or_{q_i}^S
  (\mathscr P_k^{(i)},\mathscr O_{S'})\bigr)
  =\operatorname{Tor}_{q_i}^A(\Gamma(V,\mathscr P_k^{(i)}),A')
\]
d'où pour le bicomplexe
$H_{q_i}^{\mathrm{II}}(M_{\bullet\bullet\bullet}'^{(i)})$ l'expression
\[
  C^\bullet\bigl(\mathfrak U'^{(i)},
  \mathscr{T}\!or_{q_i}^S(\mathscr P_\bullet^{(i)},\mathscr O_{S'})\bigr),
\]
ce qui donne finalement l'expression cherchée pour le terme $E_2$ de la
suite $(f')$. Le fait que cette suite soit birégulière sous les conditions
indiquées se vérifie comme d'ordinaire, tenant compte de ce que, si
$\mathscr P_\bullet^{(i)}$ est limité inférieurement, tous les degrés de
$M_{\bullet\bullet\bullet}'^{(i)}$ sont limités inférieurement.
\end{proof}

\begin{remark}[6.9.5]
On voit comme dans (6.7.6) que le remplacement des
$\mathscr P_\bullet^{(i)}$ et de $\mathscr Q'_\bullet$ par des complexes qui
leur sont respectivement homotopes ne change pas les suites $(e)$, $(f)$ et
$(f')$ à un isomorphisme canonique près. En outre, pour la suite $(f)$, des
homomorphismes
$\mathscr P_\bullet^{(i)}\to\mathscr R_\bullet^{(i)}$,
$\mathscr Q'_\bullet\to\mathscr T'_\bullet$ de complexes qui donnent des
isomorphismes en homologie
$\mathscr H_\bullet(\mathscr P_\bullet^{(i)})\xrightarrow{\sim}
\mathscr H_\bullet(\mathscr R_\bullet^{(i)})$,
$\mathscr H_\bullet(\mathscr Q'_\bullet)\xrightarrow{\sim}
\mathscr H_\bullet(\mathscr T'_\bullet)$ fournissent un isomorphisme de suites
spectrales
\[
\begin{split}
  {}^{(f)}\mathscr E
  (f_1,\ldots,f_m;
   \mathscr P_\bullet^{(1)},\ldots,\mathscr P_\bullet^{(m)},
   \mathscr Q'_\bullet)
  \xrightarrow{\sim}\;&
  {}^{(f)}\mathscr E
  (f_1,\ldots,f_m;\\
  &\mathscr R_\bullet^{(1)},\ldots,\mathscr R_\bullet^{(m)},
   \mathscr T'_\bullet);
\end{split}
\]
la démonstration est la même que pour (6.7.6) en tenant compte du résultat
de (6.7.6) et de la régularité de la suite $(f)$.
\end{remark}

\begin{corollary}[6.9.6]
Sous les conditions de (6.9.1), supposons que :
\begin{enumerate}
  \item[$1^\circ$] Les complexes $\mathscr P_\bullet^{(i)}$ sont formés de
  Modules plats sur $S$, et les $\mathscr O_Y$-Modules
  \[
    \mathfrak{Tor}_n^S
    (f_1,\ldots,f_m;
     \mathscr P_\bullet^{(1)},\ldots,\mathscr P_\bullet^{(m)})
  \]
  sont plats sur $S$.
  \item[$2^\circ$] Les $\mathscr P_\bullet^{(i)}$ et
  $\mathscr Q'_\bullet$ sont limités inférieurement.
\end{enumerate}

\oldpage[III]{169}
On a alors, en posant
$\mathscr P_\bullet'^{(i)}=
\mathscr P_\bullet^{(i)}\otimes_{\mathscr O_S}\mathscr O_{S'}$, des
isomorphismes canoniques fonctoriels
\[
\begin{split}
  &\mathfrak{Tor}_n^{S'}
  (f'_1,\ldots,f'_m,1_{S'};
   \mathscr P_\bullet'^{(1)},\ldots,
   \mathscr P_\bullet'^{(m)},\mathscr Q'_\bullet)
  \xrightarrow{\sim}\\
  &\qquad
  \bigoplus_{n'+n''=n}
  \mathfrak{Tor}_{n'}^S
  (f_1,\ldots,f_m;
   \mathscr P_\bullet^{(1)},\ldots,
   \mathscr P_\bullet^{(m)})
  \otimes_{\mathscr O_S}\mathscr H_{n''}(\mathscr Q'_\bullet).
\end{split}
\tag{6.9.6.1}\label{III.6.9.6.1.fr}
\]
En particulier, pour $\mathscr Q'_\bullet$ réduit à un seul terme
$\mathscr F'$ de degré $0$, on a des isomorphismes canoniques fonctoriels
\[
\begin{split}
  &\mathfrak{Tor}_n^{S'}
  (f'_1,\ldots,f'_m,1_{S'};
   \mathscr P_\bullet'^{(1)},\ldots,
   \mathscr P_\bullet'^{(m)},\mathscr F')
  \xrightarrow{\sim}\\
  &\qquad
  \mathfrak{Tor}_n^S
  (f_1,\ldots,f_m;
   \mathscr P_\bullet^{(1)},\ldots,
   \mathscr P_\bullet^{(m)})
  \otimes_{\mathscr O_S}\mathscr F'
\end{split}
\tag{6.9.6.2}\label{III.6.9.6.2.fr}
\]
et plus particulièrement, pour $\mathscr F'=\mathscr O_{S'}$,
\[
\begin{split}
  &\mathfrak{Tor}_n^{S'}
  (f'_1,\ldots,f'_m;
   \mathscr P_\bullet'^{(1)},\ldots,
   \mathscr P_\bullet'^{(m)})
  \xrightarrow{\sim}\\
  &\qquad
  \mathfrak{Tor}_n^S
  (f_1,\ldots,f_m;
   \mathscr P_\bullet^{(1)},\ldots,
   \mathscr P_\bullet^{(m)})
  \otimes_{\mathscr O_S}\mathscr O_{S'}.
\end{split}
\tag{6.9.6.3}\label{III.6.9.6.3.fr}
\]
\end{corollary}

\begin{proof}
L'hypothèse de platitude sur les Modules composant les
$\mathscr P_\bullet^{(i)}$ entraîne que les complexes
$\mathscr{T}\!or_{q_i}^S(\mathscr P_\bullet^{(i)},\mathscr O_{S'})$ sont nuls
pour $q_i\ne0$ (6.5.8). La suite $(f')$ est donc dégénérée; l'hypothèse
$2^\circ$ entraîne d'ailleurs qu'elle est birégulière (6.9.3), donc le
edge-homomorphisme
\[
\begin{split}
  &\mathfrak{Tor}_n^S
  (f_1,\ldots,f_m,1_{S'};
   \mathscr P_\bullet^{(1)},\ldots,
   \mathscr P_\bullet^{(m)},\mathscr Q'_\bullet)
  \longrightarrow {}^{(f')}\mathscr E_{n0}^2\\
  &\qquad=
  \mathfrak{Tor}_n^{S'}
  (f'_1,\ldots,f'_m,1_{S'};
   \mathscr P_\bullet'^{(1)},\ldots,
   \mathscr P_\bullet'^{(m)},\mathscr Q'_\bullet)
\end{split}
\tag{6.9.6.4}\label{III.6.9.6.4.fr}
\]
est bijectif (0, 11.1.6). L'hypothèse de platitude sur les Modules
\[
  \mathfrak{Tor}_n^S
  (f_1,\ldots,f_m;
   \mathscr P_\bullet^{(1)},\ldots,\mathscr P_\bullet^{(m)})
\]
entraîne que ${}^{(e)}\mathscr E_{pq}^2=0$ pour $p\ne0$ (6.5.8). La suite
$(e)$ est donc aussi dégénérée, et comme elle est birégulière, le
edge-homomorphisme
\[
\begin{split}
  &\mathfrak{Tor}_n^S
  (f_1,\ldots,f_m,1_{S'};
   \mathscr P_\bullet^{(1)},\ldots,
   \mathscr P_\bullet^{(m)},\mathscr Q'_\bullet)
  \longrightarrow {}^{(e)}\mathscr E_{0n}^2\\
  &\qquad=
  \bigoplus_{n'+n''=n}
  \mathfrak{Tor}_{n'}^S
  (f_1,\ldots,f_m;
   \mathscr P_\bullet^{(1)},\ldots,
   \mathscr P_\bullet^{(m)})
  \otimes_{\mathscr O_S}\mathscr H_{n''}(\mathscr Q'_\bullet)
\end{split}
\tag{6.9.6.5}\label{III.6.9.6.5.fr}
\]
est bijectif (0, 11.1.6); d'où, en combinant les deux isomorphismes
précédents, l'isomorphisme (6.9.6.1). L'isomorphisme (6.9.6.2) s'en déduit
trivialement, puisque l'on a alors
$\mathscr H_q(\mathscr Q'_\bullet)=0$ si $q\ne0$ et
$\mathscr H_0(\mathscr Q'_\bullet)=\mathscr F'$. Enfin, le cas
$\mathscr F'=\mathscr O_{S'}$ dans (6.9.6.2) donne l'isomorphisme (6.9.6.3),
compte tenu de (6.7.12).
\end{proof}

\begin{corollary}[6.9.7]
Sous les conditions de (6.9.1), supposons $S$ et $S'$ affines, et supposons
donné pour chaque $i$ un entier $d_i$ ($1\leq i\leq m$). Il existe alors un
entier $N$ ne dépendant que de $S$, des $X^{(i)}$ et des $d_i$, ayant la
propriété suivante : pour tout entier $n_0$, on a des isomorphismes canoniques
(6.9.6.3) pour $n\leq n_0$ et pour tout système de complexes
$\mathscr P_\bullet^{(i)}$ vérifiant les conditions suivantes :
\begin{enumerate}
  \item[$1^\circ$] $\mathscr P_k^{(i)}=0$ pour $k<d_i$;
  \item[$2^\circ$] $\mathscr P_k^{(i)}$ est plat sur $S$ pour $k<n_0+N$;
  \item[$3^\circ$]
  $\mathfrak{Tor}_q^S(f_1,\ldots,f_m;
  \mathscr P_\bullet^{(1)},\ldots,\mathscr P_\bullet^{(m)})$ est plat sur
  $S$ pour $q<n_0+N$.
\end{enumerate}
\end{corollary}

\begin{proof}
Supposons $\mathscr P_k^{(i)}$ plat sur $S$ pour $k<r$; alors
$\mathscr{T}\!or_{q_i}^S(\mathscr P_k^{(i)},\mathscr O_{S'})=0$ pour $k<r$
et $q_i\ne0$; calculons
\[
\begin{split}
  \mathfrak{Tor}_p^{S'}\Bigl(
  &f'_1,\ldots,f'_m;
  \mathscr{T}\!or_{q_1}^S(\mathscr P_\bullet^{(1)},\mathscr O_{S'}),
  \ldots,\\
  &\mathscr{T}\!or_{q_m}^S(\mathscr P_\bullet^{(m)},\mathscr O_{S'})
  \Bigr).
\end{split}
\tag{6.9.7.1}\label{III.6.9.7.1.fr}
\]

\oldpage[III]{170}
par la méthode de (6.6.2), à l'aide de l'image réciproque d'un recouvrement
affine fixe $\mathfrak U^{(i)}$ de $X^{(i)}$ ($1\leq i\leq m$)
(indépendant de $S'$ et des $\mathscr P_\bullet^{(i)}$); les termes
$L_{jk}^{(i)}$ sont nuls pour $j<-N_i$ (ne dépendant que de
$\mathfrak U^{(i)}$); si l'un des $q_i$ n'est pas nul, le complexe simple dont
l'homologie de degré $p$ est (6.9.7.1) a ses termes nuls pour tous les degrés
\[
  <r+\sum_{j\ne i}d_j-\sum_{i=1}^m N_i,
\]
donc (6.9.7.1) est nul pour $p<r-N$, en désignant par $N$ le plus grand des
nombres
\[
  \sum_{i=1}^m N_i-\sum_{j\ne i}d_j.
\]
On en conclut que l'on a ${}^{(f')}\mathscr E_{pq}^2=0$ pour $q\ne0$ et
$p<r-N$; comme d'autre part ${}^{(f')}\mathscr E_{pq}^2=0$ pour $q<0$, on
voit que le edge-homomorphisme (6.9.6.4) est bijectif pour $n<r-N$
(M, XV, 5.6) (pour $\mathscr Q'_\bullet=\mathscr O_{S'}$). En second lieu, si
\[
  \mathfrak{Tor}_q^S
  (f_1,\ldots,f_m;
   \mathscr P_\bullet^{(1)},\ldots,\mathscr P_\bullet^{(m)})
\]
est plat sur $S$ pour $q<r$, on a ${}^{(e)}\mathscr E_{pq}^2=0$ pour $p\ne0$
et $q<r$; par ailleurs ${}^{(e)}\mathscr E_{pq}^2=0$ pour $p<0$, donc le
edge-homomorphisme (6.9.6.5) est bijectif pour $n<r$, ce qui achève la
démonstration.
\end{proof}

Le cas le plus important de (6.9.3) dans les applications est celui où $m=1$,
$\mathscr Q'_\bullet$ étant réduit à un seul terme $\mathscr F'$ de degré
$0$; nous l'énoncerons à nouveau dans ce cas en vue de références
ultérieures\footnote{Le cas traité dans (6.9.8), et en particulier les suites
spectrales (6.9.8.3) et (6.9.8.4), nous avaient été signalés en 1957 par
J.-P. Serre.} :

\begin{proposition}[6.9.8]
Soient $S$ un préschéma, $g:S'\to S$ un morphisme,
$f:X\to Y$ un $S$-morphisme séparé et quasi-compact de $S$-préschémas,
$\mathscr P_\bullet$ un complexe de $\mathscr O_X$-Modules quasi-cohérents
limité inférieurement, $\mathscr F'$ un $\mathscr O_{S'}$-Module
quasi-cohérent. Il existe deux foncteurs spectraux biréguliers en
$\mathscr P_\bullet$ et $\mathscr F'$, à valeurs dans la catégorie des
$\mathscr O_{Y_{(S')}}$-Modules quasi-cohérents, ayant même aboutissement
$\mathfrak{Tor}_\bullet^S(f,1_{S'};\mathscr P_\bullet,\mathscr F')$, et dont
les termes $E_2$ sont
\[
  {}'\mathscr E_{pq}^2
  =\mathscr{T}\!or_p^S
  (\mathscr H^{-q}(f,\mathscr P_\bullet),\mathscr F')
  \tag{6.9.8.1}\label{III.6.9.8.1.fr}
\]
et
\[
  {}''\mathscr E_{pq}^2
  =\mathscr H^{-p}
  (f',\mathscr{T}\!or_q^S(\mathscr P_\bullet,\mathscr F')),
  \tag{6.9.8.2}\label{III.6.9.8.2.fr}
\]
où $f'=f_{(S')}:X_{(S')}\to Y_{(S')}$. 
\end{proposition}

\begin{proof}
Les suites en question peuvent aussi s'obtenir, non en partant de (6.9.3),
mais des suites $(a)$ et $(b')$ de (6.7.3) pour
$X^{(1)}=X$, $Y^{(1)}=Y$, $X^{(2)}=Y^{(2)}=S'$, $f_1=f$,
$f_2=1_{S'}$. Lorsque $S=S'=Y$, $Y$ étant affine, on obtient deux suites
spectrales de termes $E_2$ égaux à
\[
  {}'\mathscr E_{pq}^2
  =\mathscr{T}\!or_p^Y
  (\mathscr H^{-q}(f,\mathscr P_\bullet),\mathscr F)
  \tag{6.9.8.3}\label{III.6.9.8.3.fr}
\]
\[
  {}''\mathscr E_{pq}^2
  =\mathscr H^{-p}
  (f,\mathscr{T}\!or_q^Y(\mathscr P_\bullet,\mathscr F))
  \tag{6.9.8.4}\label{III.6.9.8.4.fr}
\]
aboutissant (en vertu de (6.7.6)) à l'hypercohomologie
$\mathscr H^\bullet(f,\mathscr P_\bullet\otimes_Y\mathscr F)$ du foncteur
$f_*$ par rapport au complexe
$\mathscr P_\bullet\otimes_Y\mathscr F$ de $\mathscr O_X$-Modules, pour tout
$\mathscr O_Y$-Module quasi-cohérent et $Y$-plat $\mathscr F$ (ou pour tout
$\mathscr O_Y$-Module quasi-cohérent $\mathscr F$ lorsque
$\mathscr P_\bullet$ est formé de $\mathscr O_X$-Modules $Y$-plats), qui sont
distinctes de celles de (6.2.1).
\end{proof}

\begin{corollary}[6.9.9]
Sous les conditions de (6.9.8), supposons que le complexe
$\mathscr P_\bullet$ soit limité inférieurement, formé de Modules plats sur
$S$, et que les $\mathscr O_Y$-Modules
$\mathscr H^n(f,\mathscr P_\bullet)$ soient plats sur $S$.

\oldpage[III]{171}
On a alors des isomorphismes canoniques fonctoriels
\[
  \mathfrak{Tor}_n^{S'}
  (f',1_{S'};\mathscr P'_\bullet,\mathscr F')
  \xrightarrow{\sim}
  \mathscr H^{-n}(f,\mathscr P_\bullet)
  \otimes_{\mathscr O_S}\mathscr F'
  \tag{6.9.9.1}\label{III.6.9.9.1.fr}
\]
où
$\mathscr P'_\bullet=\mathscr P_\bullet\otimes_{\mathscr O_S}
\mathscr O_{S'}$; en particulier, pour $\mathscr F'=\mathscr O_{S'}$, on a
des isomorphismes canoniques fonctoriels
\[
  \mathscr H^n(f',\mathscr P'_\bullet)
  \xrightarrow{\sim}
  \mathscr H^n(f,\mathscr P_\bullet)
  \otimes_{\mathscr O_S}\mathscr O_{S'}.
  \tag{6.9.9.2}\label{III.6.9.9.2.fr}
\]
C'est le cas particulier $m=1$ de (6.9.6). Plus particulièrement :
\end{corollary}

\begin{corollary}[6.9.10]
Soient $S$ un préschéma, $f:X\to Y$ un $S$-morphisme séparé et
quasi-compact de $S$-préschémas, $\mathscr P_\bullet$ un complexe limité
inférieurement, formé de $\mathscr O_X$-Modules quasi-cohérents, plats sur
$S$. On suppose en outre que les $\mathscr O_Y$-Modules
$\mathscr H^n(f,\mathscr P_\bullet)$ soient plats sur $S$. Pour tout $s\in S$,
notons $X_s$ et $Y_s$ les fibres
$X\otimes_S k(s)$, $Y\otimes_S k(s)$,
$f_s:X_s\to Y_s$ le morphisme $f\times_S1$,
$\mathscr P_\bullet^s$ le complexe
$\mathscr P_\bullet\otimes_{\mathscr O_S}k(s)$ de
$\mathscr O_{X_s}$-Modules. On a alors des isomorphismes canoniques
fonctoriels
\[
  \mathscr H^n(f_s,\mathscr P_\bullet^s)
  \xrightarrow{\sim}
  \mathscr H^n(f,\mathscr P_\bullet)\otimes_{\mathscr O_S}k(s).
  \tag{6.9.10.1}\label{III.6.9.10.1.fr}
\]
\end{corollary}

On a donc, moyennant des hypothèses de platitude convenables, un cas où la
formation des foncteurs dérivés $R^nf_*(\mathscr F)$ « commute au passage aux
fibres », que nous retrouverons par une autre méthode au §~7.

\subsection{Structure locale de certains foncteurs cohomologiques.}

\begin{proposition}[6.10.1]
Soient $S=\operatorname{Spec}(A)$ un schéma affine,
$Y^{(i)}$ ($1\leq i\leq n$) une famille finie de $S$-schémas affines, plats
sur $S$; pour chaque $i$, soit $f_i:X^{(i)}\to Y^{(i)}$ un $S$-morphisme
séparé et quasi-compact, et soit $\mathscr P_\bullet^{(i)}$ un complexe de
$\mathscr O_{X^{(i)}}$-Modules quasi-cohérents limité inférieurement. Soit
$Y$ le produit des $S$-schémas $Y^{(i)}$. Il existe un complexe
$\mathscr R_\bullet$ de $\mathscr O_Y$-Modules quasi-cohérents et plats sur
$S$, ayant la propriété suivante : pour tout $S$-schéma affine $S'$ et tout
complexe de $\mathscr O_{S'}$-Modules quasi-cohérents
$\mathscr Q'_\bullet$ limité inférieurement, il y a un isomorphisme
\[
  \mathbf{Tor}_\bullet^S
  (f_1,\ldots,f_n,1_{S'};
   \mathscr P_\bullet^{(1)},\ldots,\mathscr P_\bullet^{(n)},
   \mathscr Q'_\bullet)
  \xrightarrow{\sim}
  \mathscr H_\bullet
  (\mathscr R_\bullet\otimes_{\mathscr O_S}\mathscr Q'_\bullet)
  \tag{6.10.1.1}\label{III.6.10.1.1.fr}
\]
qui est un isomorphisme de $\partial$-foncteurs en $\mathscr Q'_\bullet$. En
outre, pour tout $S$-morphisme $u:S''\to S'$ de $S$-schémas affines, le
diagramme
\[
\xymatrix@C=12pt@R=30pt{
  {\begin{gathered}
    \mathbf{Tor}_\bullet^S
    (f_1,\ldots,f_n,1_{S'};\\
    \mathscr P_\bullet^{(1)},\ldots,\mathscr P_\bullet^{(n)},
    \mathscr Q'_\bullet)
  \end{gathered}}
  \ar[r]^-{\sim}\ar[d]
  &
  {\mathscr H_\bullet
    (\mathscr R_\bullet\otimes_{\mathscr O_S}\mathscr Q'_\bullet)}
  \ar[d]
  \\
  {\begin{gathered}
    \mathbf{Tor}_\bullet^S
    (f_1,\ldots,f_n,1_{S''};\\
    \mathscr P_\bullet^{(1)},\ldots,\mathscr P_\bullet^{(n)},
    u^*(\mathscr Q'_\bullet))
  \end{gathered}}
  \ar[r]^-{\sim}
  &
  {\mathscr H_\bullet
    (\mathscr R_\bullet\otimes_{\mathscr O_S}u^*(\mathscr Q'_\bullet))}
}
\tag{6.10.1.2}\label{III.6.10.1.2.fr}
\]
(où les flèches verticales sont les $(1_Y\times_Su)$-morphismes canoniques
définis dans (6.7.10)) est commutatif.
\end{proposition}

\begin{proof}
\oldpage[III]{172}
Calculons les hypertor par la méthode de (6.6.2), en tenant compte de la remarque
(6.6.6); avec les notations de (6.6.2), chaque
$L_{\bullet\bullet}^{(i)}$ est un bicomplexe de $A_i$-modules, en désignant par
$A_i$ l'anneau de $Y^{(i)}$; il admet donc une résolution de Cartan--Eilenberg
projective $M_{\bullet\bullet\bullet}^{(i)}$ (au sens de (0, 11.7.1)) formée de
$A_i$-modules, et en vertu de (6.6.6), le premier membre de (6.10.1.1) est
canoniquement isomorphe à
\[
  \mathscr H_\bullet
  (M_{\bullet\bullet\bullet}\otimes_A Q'_\bullet),
\]
où
\[
  M_{\bullet\bullet\bullet}
  =
  M_{\bullet\bullet\bullet}^{(1)}
  \otimes_A M_{\bullet\bullet\bullet}^{(2)}
  \otimes_A\cdots\otimes_A M_{\bullet\bullet\bullet}^{(n)}
\]
et $\mathscr Q'_\bullet=\widetilde{Q'_\bullet}$. Comme, par hypothèse, les
anneaux $A_i$ sont des $A$-modules plats, les
$M_{\bullet\bullet\bullet}^{(i)}$ sont des tricomplexes de $A$-modules plats
($0_{\mathrm I}$, 6.2.1), et il en est de même de
$M_{\bullet\bullet\bullet}$; en outre, si $B$ est l'anneau de $Y$, produit
tensoriel des $A_i$, $M_{\bullet\bullet\bullet}$ est un tricomplexe de
$B$-modules; le complexe
\[
  \mathscr R_\bullet=(M_{\bullet\bullet\bullet})^\sim
\]
de $\mathscr O_Y$-Modules (où $M_{\bullet\bullet\bullet}$ est considéré comme
complexe simple) répond donc à la question, comme il résulte aisément de
(6.7.10).
\end{proof}

\begin{corollary}[6.10.2]
Dans l'énoncé de (6.10.1), on peut supposer $\mathscr R_\bullet$ limité
inférieurement. Lorsque les $\mathscr P_\bullet^{(i)}$ sont limités
supérieurement et les $Y^{(i)}$ de dimension cohomologique finie, on peut
supposer $\mathscr R_\bullet$ limité supérieurement.
\end{corollary}

\begin{proof}
La première assertion résulte de ce que les trois degrés de chacun des
$M_{\bullet\bullet\bullet}^{(i)}$ sont limités inférieurement; d'autre part, si
les anneaux $A_i$ sont de dimension cohomologique finie, le troisième degré de
chacun des $M_{\bullet\bullet\bullet}^{(i)}$ ne prend qu'un nombre fini de
valeurs, et il en est de même par construction de son premier degré (6.6.2);
comme son second degré est limité supérieurement s'il en est ainsi du degré de
$\mathscr P_\bullet^{(i)}$ (6.6.2), la seconde assertion en résulte aussitôt.
\end{proof}

\begin{remarks}[6.10.3]
\textnormal{(i)} Avec les notations de (6.10.1),
$\mathscr H_\bullet(\mathscr R_\bullet\otimes_S\mathscr Q'_\bullet)$ est
isomorphe à
$\mathbf{Tor}_\bullet(\mathscr R_\bullet,\mathscr Q'_\bullet)$ puisque
$\mathscr R_\bullet$ est formé de $\mathscr O_Y$-Modules $S$-plats (6.5.9); il
est donc (6.5.4) l'aboutissement d'une suite spectrale régulière de terme
$E_2$ donné par
\[
  \mathscr E_{pq}^2
  =
  \bigoplus_{q'+q''=q}
  \mathbf{Tor}_p^S
  \bigl(
    \mathscr H_{q'}(\mathscr R_\bullet),
    \mathscr H_{q''}(\mathscr Q'_\bullet)
  \bigr)
  \tag{6.10.3.1}\label{III.6.10.3.1.fr}
\]
qui n'est autre que la suite spectrale $(e)$ du changement de base (6.9.3).

\medskip
\textnormal{(ii)} Soit $\mathscr R'_\bullet$ un second complexe de
$\mathscr O_Y$-Modules quasi-cohérents, plat sur $S$, et soit
$g:\mathscr R'_\bullet\to\mathscr R_\bullet$ un homomorphisme de complexes tel
que
\[
  \mathscr H_\bullet(g):
  \mathscr H_\bullet(\mathscr R'_\bullet)
  \longrightarrow
  \mathscr H_\bullet(\mathscr R_\bullet)
\]
soit un isomorphisme. Alors, en vertu de (6.3.3) et (6.5.9), on déduit de $g$
un isomorphisme de $\partial$-foncteurs en $\mathscr Q'_\bullet$
\[
  \mathscr H_\bullet
  (\mathscr R'_\bullet\otimes_S\mathscr Q'_\bullet)
  \xrightarrow{\sim}
  \mathscr H_\bullet
  (\mathscr R_\bullet\otimes_S\mathscr Q'_\bullet)
\]
tel que le diagramme
\[
\xymatrix@C=28pt@R=30pt{
  \mathscr H_\bullet
  (\mathscr R'_\bullet\otimes_S\mathscr Q'_\bullet)
  \ar[r]^-{\sim}\ar[d]
  &
  \mathscr H_\bullet
  (\mathscr R_\bullet\otimes_S\mathscr Q'_\bullet)
  \ar[d]
  \\
  \mathscr H_\bullet
  (\mathscr R'_\bullet\otimes_Su^*(\mathscr Q'_\bullet))
  \ar[r]^-{\sim}
  &
  \mathscr H_\bullet
  (\mathscr R_\bullet\otimes_Su^*(\mathscr Q'_\bullet))
}
\]
\par\medskip\noindent
soit commutatif. Cela prouve donc que le complexe $\mathscr R_\bullet$ n'est
pas entièrement déterminé par les propriétés de (6.10.1).

\medskip
\oldpage[III]{173}
\textnormal{(iii)} Dans la démonstration de (6.10.1), on peut supposer les
$M_{\bullet\bullet\bullet}^{(i)}$ formés de $A_i$-modules libres (comme il
résulte aisément de la démonstration de (0, 11.5.2.1) « dualisée »); les
\[
  M_{\bullet\bullet\bullet}^{(i)}\otimes_{A_i}B
\]
sont alors formés de $B$-modules libres, et comme
$M_{\bullet\bullet\bullet}$ est égal à leur produit tensoriel sur $B$, on voit
qu'on peut supposer en outre dans (6.10.1) que $\mathscr R_\bullet$ est associé
à un complexe de $B$-modules libres. En outre, en vertu de (M, XVII, 1.2), le
tricomplexe $M_{\bullet\bullet\bullet}$ dépend fonctoriellement de chacun des
bicomplexes $L_{\bullet\bullet}^{(i)}$ (donc de chacun des
$\mathscr P_\bullet^{(i)}$, lorsqu'on fixe un recouvrement fini de chacun des
$X^{(i)}$), les « morphismes » de tricomplexes devant être ici entendus comme
les classes d'homomorphismes pour la relation d'homotopie; d'ailleurs, le
remplacement d'un recouvrement de $X^{(i)}$ par un recouvrement plus fin
donnant lieu pour les $L_{\bullet\bullet}^{(i)}$ à des homomorphismes définis
précisément à homotopie près (6.6.8), on voit finalement qu'avec la convention
précédente pour les morphismes, le tricomplexe
$M_{\bullet\bullet\bullet}$ est un foncteur en chacun des
$\mathscr P_\bullet^{(i)}$. Nous préciserons cette dépendance fonctorielle, et
notamment le comportement de $\mathscr R_\bullet$ relativement à des suites
exactes
\[
  0\longrightarrow\mathscr P_\bullet^{\prime(i)}
  \longrightarrow\mathscr P_\bullet^{(i)}
  \longrightarrow\mathscr P_\bullet^{\prime\prime(i)}
  \longrightarrow0
\]
de complexes, dans le chapitre consacré à une algèbre générale de foncteurs
cohomologiques, mentionné dans (6.1.3).
\end{remarks}

\begin{scholium}[6.10.4]
Le fait que $\mathscr R_\bullet$ est formé de $\mathscr O_Y$-Modules
$S$-plats entraîne aisément que
\[
  \mathscr H_\bullet
  (\mathscr R_\bullet\otimes_S\mathscr Q'_\bullet)
\]
est un foncteur homologique en $\mathscr Q'_\bullet$ (voir le raisonnement de
(7.7.1)). C'est cette propriété qui, ainsi qu'on l'a mentionné en (6.1.1), est
la motivation de l'introduction de l'hypertor. Posons en effet
\[
  X=X^{(1)}\times_SX^{(2)}\times_S\cdots\times_SX^{(n)},
  \qquad
  f=f_1\times_Sf_2\times_S\cdots\times_Sf_n,
\]
\[
  \mathscr P_\bullet
  =
  \mathscr P_\bullet^{(1)}
  \otimes_S\mathscr P_\bullet^{(2)}
  \otimes_S\cdots\otimes_S\mathscr P_\bullet^{(n)},
\]
\[
  X'=X\times_SS',
  \qquad
  Y'=Y\times_SS',
  \qquad
  f'=f\times_S1_{S'}.
\]
Les problèmes de changement de base amènent à étudier l'hypercohomologie
\[
  \mathscr H^\bullet_{f'}
  (\mathscr P_\bullet\otimes_S\mathscr N')
\]
en tant que foncteur par rapport au $\mathscr O_{S'}$-Module quasi-cohérent
$\mathscr N'$, ou encore l'hypercohomologie
\[
  \mathscr H^\bullet_f
  (\mathscr P_\bullet\otimes_S\mathscr N)
\]
comme foncteur en le $\mathscr O_S$-Module quasi-cohérent $\mathscr N$. Lorsque
les $\mathscr P_\bullet^{(i)}$ (donc aussi $\mathscr P_\bullet$) sont
$S$-plats, il résulte de ce qui précède et de (6.7.6) que ce foncteur est bien
un foncteur cohomologique en $\mathscr N$; mais il n'en est plus de même
lorsqu'on ne fait plus l'hypothèse de platitude sur les
$\mathscr P_\bullet^{(i)}$, et l'on ne peut plus alors aborder l'étude de
\[
  \mathscr H^\bullet_f
  (\mathscr P_\bullet\otimes_S\mathscr N)
\]
par les méthodes usuelles de l'Algèbre homologique.

Nous aurons toutefois surtout à utiliser le cas où $n=1$, $Y=S$ et où
$\mathscr P_\bullet$ est formé de $\mathscr O_X$-Modules $Y$-plats. On a dans
ce cas le
\end{scholium}

\begin{theorem}[6.10.5]
Soient $Y=\operatorname{Spec}(A)$ un schéma affine noethérien,
$f:X\to Y$ un morphisme propre, $\mathscr P_\bullet$ un complexe de
$\mathscr O_X$-Modules cohérents, plats sur $Y$, limité inférieurement. Il
existe alors un complexe $\mathscr L_\bullet$ de $\mathscr O_Y$-Modules,
limité inférieurement, dont les termes $\mathscr L_i$ sont des
$\mathscr O_Y$-Modules de la forme $\mathscr O_Y^{m_i}$, et un isomorphisme
\[
  \mathscr H^\bullet
  (f,\mathscr P_\bullet\otimes_Y\mathscr Q_\bullet)
  \xrightarrow{\sim}
  \mathscr H_\bullet
  (\mathscr L_\bullet\otimes_Y\mathscr Q_\bullet)
  \tag{6.10.5.1}\label{III.6.10.5.1.fr}
\]
de $\partial$-foncteurs en le complexe $\mathscr Q_\bullet$ de
$\mathscr O_Y$-Modules quasi-cohérents, limité inférieurement. En outre, pour
tout morphisme $u:Y'\to Y$, on a, en posant
\[
  X'=X_{(Y')},
  \qquad
  f'=f_{(Y')},
  \qquad
  \mathscr P'_\bullet
  =
  \mathscr P_\bullet\otimes_Y\mathscr O_{Y'},
  \qquad
  \mathscr L'_\bullet=u^*(\mathscr L_\bullet),
\]
\oldpage[III]{174}
(qui est un complexe de $\mathscr O_{Y'}$-Modules localement libres de type
fini), un isomorphisme
\[
  \mathscr H^\bullet
  (f',\mathscr P'_\bullet\otimes_{Y'}\mathscr Q'_\bullet)
  \xrightarrow{\sim}
  \mathscr H_\bullet
  (\mathscr L'_\bullet\otimes_{Y'}\mathscr Q'_\bullet)
  \tag{6.10.5.2}\label{III.6.10.5.2.fr}
\]
de $\partial$-foncteurs en le complexe $\mathscr Q'_\bullet$ de
$\mathscr O_{Y'}$-Modules quasi-cohérents, limité inférieurement, de sorte que
le diagramme
\[
\xymatrix@C=25pt@R=30pt{
  \mathscr H^\bullet
  (f,\mathscr P_\bullet\otimes_Y\mathscr Q_\bullet)
  \ar[r]^-{\sim}\ar[d]
  &
  \mathscr H_\bullet
  (\mathscr L_\bullet\otimes_Y\mathscr Q_\bullet)
  \ar[d]
  \\
  \mathscr H^\bullet
  (f',\mathscr P'_\bullet\otimes_{Y'}u^*(\mathscr Q_\bullet))
  \ar[r]^-{\sim}
  &
  \mathscr H_\bullet
  (\mathscr L'_\bullet\otimes_{Y'}u^*(\mathscr Q_\bullet))
}
\tag{6.10.5.3}\label{III.6.10.5.3.fr}
\]
soit commutatif.
\end{theorem}

\begin{proof}
L'application de (6.10.1) donne d'abord un complexe limité inférieurement
(6.10.2) $\mathscr R_\bullet$ de $\mathscr O_Y$-Modules quasi-cohérents et
$Y$-libres (6.10.3, (iii)) et un isomorphisme
\[
  \mathscr H^\bullet
  (f,\mathscr P_\bullet\otimes_Y\mathscr Q_\bullet)
  \xrightarrow{\sim}
  \mathscr H_\bullet
  (\mathscr R_\bullet\otimes_Y\mathscr Q_\bullet)
  \tag{6.10.5.4}\label{III.6.10.5.4.fr}
\]
de $\partial$-foncteurs en $\mathscr Q_\bullet$, mais \emph{a priori} les
termes de $\mathscr R_\bullet$ ne sont pas nécessairement des
$\mathscr O_Y$-Modules de type fini. Mais si l'on applique (6.10.5.4) au cas
où $\mathscr Q_\bullet$ est un complexe réduit à un seul terme
$\mathscr O_Y$, on voit que $\mathscr H_\bullet(\mathscr R_\bullet)$ est
isomorphe à $\mathscr H^\bullet(f,\mathscr P_\bullet)$, et est par suite formé
de $\mathscr O_Y$-Modules cohérents (6.2.5). On sait alors (0, 11.9.2) qu'il
existe un complexe $\mathscr L_\bullet$ limité inférieurement, formé de
$\mathscr O_Y$-Modules associés à des $A$-modules libres de type fini, et un
homomorphisme
\[
  \mathscr L_\bullet\longrightarrow\mathscr R_\bullet,
\]
tels que l'homomorphisme correspondant pour l'homologie
\[
  \mathscr H_\bullet(\mathscr L_\bullet)
  \longrightarrow
  \mathscr H_\bullet(\mathscr R_\bullet)
\]
soit bijectif; d'où l'isomorphisme (6.10.5.1), en vertu de (6.10.3, (ii)).
Les autres assertions de (6.10.5) résultent de (6.10.1) et (6.10.3, (ii))
lorsque $Y'$ est affine; dans le cas général, il suffit de vérifier que
lorsqu'on considère un recouvrement $(V_\alpha)$ de $Y'$ par des ouverts
affines, et l'isomorphisme correspondant (6.10.5.2) relatif à chacun des
$V_\alpha$, les restrictions à un ouvert affine
$W\subset V_\alpha\cap V_\beta$ des isomorphismes correspondant à
$V_\alpha$ et à $V_\beta$ coïncident avec l'isomorphisme correspondant à
$W$, ce qui résulte de la commutativité du diagramme (6.10.1.2) appliqué aux
injections canoniques $W\to V_\alpha$ et $W\to V_\beta$.
\end{proof}

\begin{remark}[6.10.6]
Dans les chapitres suivants, nous appliquerons surtout (6.10.5) au cas où
$\mathscr P_\bullet$ est réduit à un seul $\mathscr O_X$-Module cohérent
$\mathscr F$, plat sur $Y$. Comme on a alors
\[
  \mathscr H_n(\mathscr L_\bullet)
  =
  \mathscr H^{-n}(f,\mathscr F)
  =
  R^{-n}f_*(\mathscr F)
  \qquad(6.2.1),
\]
on voit que les $\mathscr H_n(\mathscr L_\bullet)$ sont nuls pour $n>0$;
nous verrons plus loin (7.7.12, (i)) qu'on peut alors supposer que
$\mathscr L_\bullet$ n'a que des termes de degrés $\leqslant0$ (donc en
nombre fini), à condition de remplacer l'hypothèse que les $\mathscr L_i$
sont associés à des $A$-modules libres de type fini par celle que les
$\mathscr L_i$ sont localement libres de type fini.

Le complexe $\mathscr L_\bullet$ correspondant à un tel
$\mathscr O_X$-Module $\mathscr F$ ne paraît posséder aucune
\oldpage[III]{175}
propriété particulière, en dehors de la restriction précédente sur les
degrés. On peut alors se demander si inversement, étant donné un complexe
$\mathscr L_\bullet$ formé de $\mathscr O_Y$-Modules associés à des
$A$-modules projectifs de type fini, limité inférieurement et dont les termes
de degré $>0$ sont nuls, il existe un $Y$-schéma $X$, projectif et plat sur
$Y$ et un $\mathscr O_X$-Module localement libre $\mathscr F$, tels qu'il y
ait un isomorphisme
\[
  \mathscr H^\bullet
  (f,\mathscr F\otimes_Y\mathscr Q_\bullet)
  \xrightarrow{\sim}
  \mathscr H_\bullet
  (\mathscr L_\bullet\otimes_Y\mathscr Q_\bullet)
\]
fonctoriel en $\mathscr Q_\bullet$. L'intérêt d'un tel résultat serait de
réduire complètement la théorie cohomologique des Modules cohérents et
$Y$-plats sur des $Y$-schémas propres, à la théorie « à homotopie près » des
complexes de $A$-modules projectifs de type fini sur un anneau noethérien
$A$.
\end{remark}

\section{Étude du changement de base dans les foncteurs homologiques covariants de Modules}
\label{section:III.7.fr}

\subsection{Foncteurs de $A$-modules}
\label{subsection:III.7.1.fr}

\begin{env}[7.1.1]
Étant donné un anneau $A$ (non nécessairement commutatif), nous noterons
$\mathbf{Ab}_A$ la catégorie des $A$-modules à gauche, et nous noterons
simplement $\mathbf{Ab}$ la catégorie des $\mathbf Z$-modules, identiques aux
groupes commutatifs. Soit
\[
  T:\mathbf{Ab}_A\longrightarrow\mathbf{Ab}
\]
un foncteur covariant additif, et soit $M$ un $(A,A)$-bimodule; $T(M)$ est
alors muni de façon naturelle d'une structure de $A$-module à droite. En
effet, pour tout $a\in A$, notons $h_{a,M}$ (ou simplement $h_a$)
l'endomorphisme $x\mapsto xa$ du $A$-module à gauche $M$. Par hypothèse,
$T(h_a)$ est un endomorphisme du $\mathbf Z$-module $T(M)$; en outre, comme
$T$ est un foncteur covariant additif, on a, pour $a\in A$, $b\in A$,
\[
  T(h_{ab})
  =
  T(h_b\circ h_a)
  =
  T(h_b)\circ T(h_a),
  \qquad\text{et}\qquad
  T(h_{a+b})
  =
  T(h_a+h_b)
  =
  T(h_a)+T(h_b);
\]
cela prouve que l'application
\[
  (a,y)\longmapsto T(h_a)(y)
\]
est une loi externe de $A$-module à droite sur $T(M)$. En particulier,
$T(A_s)$ est un $A$-module à droite.
\end{env}

\begin{env}[7.1.2]
Lorsque $A$ est un anneau commutatif, il résulte de (7.1.1) que pour tout
$A$-module $M$, $T(M)$ est naturellement muni d'une structure de $A$-module;
en outre, si $u:M\to N$ est un homomorphisme de $A$-modules, on a, pour tout
$a\in A$,
\[
  u\circ h_{a,M}=h_{a,N}\circ u,
\]
d'où
\[
  T(u)\circ T(h_{a,M})
  =
  T(h_{a,N})\circ T(u),
\]
ce qui prouve que $T(u):T(M)\to T(N)$ est un homomorphisme de $A$-modules; on
voit donc que $T$ peut être considéré comme un foncteur covariant additif de
la catégorie $\mathbf{Ab}_A$ dans elle-même. De façon précise, on a ainsi
défini une équivalence canonique entre la catégorie des foncteurs covariants
additifs $\mathbf{Ab}_A\to\mathbf{Ab}$ et la catégorie des foncteurs
covariants $A$-linéaires
\[
  T:\mathbf{Ab}_A\longrightarrow\mathbf{Ab}_A,
\]
c'est-à-dire tels que
\[
  T(h_{a,M})=h_{a,T(M)}
  \qquad\text{pour tout }a\in A.
\]
Comme le foncteur d'inclusion
\[
  I:\mathbf{Ab}_A\longrightarrow\mathbf{Ab},
\]
faisant correspondre à tout $A$-module le $\mathbf Z$-module sous-jacent, est
exact et fidèle, les propriétés d'exactitude des deux foncteurs associés par
l'équivalence précédente sont les mêmes.
\end{env}

\begin{env}[7.1.3]
L'anneau $A$ étant toujours supposé commutatif, soit $B$ une $A$-algèbre (non
nécessairement commutative), et soit $\rho:A\to B$ l'homomorphisme d'anneaux
correspondant à cette structure d'algèbre; cet homomorphisme définit un
foncteur covariant
\oldpage[III]{176}
additif
\[
  \rho_*:M\longmapsto M_{[\rho]}
\]
de la catégorie $\mathbf{Ab}_B$ des $B$-modules à gauche dans la catégorie
$\mathbf{Ab}_A$ des $A$-modules. Par composition, on en déduit un foncteur
\[
  T_{(B)}:
  \mathbf{Ab}_B
  \xrightarrow{\rho_*}
  \mathbf{Ab}_A
  \xrightarrow{T}
  \mathbf{Ab},
\]
évidemment covariant et additif, que nous noterons aussi $T^{(B)}$ (pour des
raisons typographiques) ou $T\otimes_A B$, et que nous dirons obtenu à partir
de $T$ par extension des scalaires de $A$ à $B$. Bien entendu, si $B$ est
commutative, on peut considérer $T_{(B)}$ comme un foncteur de
$\mathbf{Ab}_B$ dans elle-même (7.1.2). Lorsque $B$ est commutative, et $C$
une $B$-algèbre, on voit aussitôt que
\[
  T_{(C)}=(T_{(B)})_{(C)};
\]
il est immédiat que l'extension des scalaires est fonctorielle et additive en
$T$; en outre, lorsque $T$ commute aux limites inductives ou aux sommes
directes (resp. est exact à gauche, exact à droite, exact), il en est de même
de $T_{(B)}$: en effet, $\rho_*$ est exact et commute aux limites inductives
et aux sommes directes.
\end{env}

\begin{env}[7.1.4]
Supposons toujours $A$ commutatif, et soit
$T:\mathbf{Ab}_A\to\mathbf{Ab}_A$ un foncteur covariant additif
$A$-linéaire, commutant aux limites inductives. Alors, pour toute partie
multiplicative $S$ de $A$ et tout $A$-module $M$, on a un isomorphisme
canonique fonctoriel de $A$-modules
\[
  T(S^{-1}M)\xrightarrow{\sim}S^{-1}T(M).
  \tag{7.1.4.1}\label{III.7.1.4.1.fr}
\]
Supposons en effet d'abord que $S$ soit l'ensemble des puissances $f^n$
$(n\geqslant0)$ d'un élément $f\in A$. On sait alors que
\[
  M_f=\varinjlim M_n,
\]
où $(M_n,\varphi_{nm})$ est le système inductif des $A$-modules $M_n=M$,
avec
\[
  \varphi_{nm}:z\longmapsto f^{\,n-m}z
  \qquad(0_{\mathrm I}, 1.6.1);
\]
d'où dans ce cas l'isomorphisme (7.1.4.1) en vertu de l'hypothèse sur $T$; si
ensuite $S$ est quelconque, $S^{-1}M$ est limite inductive des $M_f$, pour
$f\in S$ ($0_{\mathrm I}$, 1.4.5) et l'on conclut de la même façon. En outre,
la fonctorialité de l'isomorphisme (7.1.4.1) montre que c'est un isomorphisme
de $S^{-1}A$-modules, et qu'on peut donc écrire, à un isomorphisme canonique
près
\[
  T_{(S^{-1}A)}(S^{-1}M)
  =
  S^{-1}T(M)
  =
  T(S^{-1}M).
  \tag{7.1.4.2}\label{III.7.1.4.2.fr}
\]
Lorsque $S=A-\mathfrak p$ est le complémentaire d'un idéal premier
$\mathfrak p$ de $A$, on écrit $T_{\mathfrak p}$ au lieu de
$T_{(A_{\mathfrak p})}$.
\end{env}

\begin{proposition}[7.1.5]
Sous les hypothèses de (7.1.4), si $T_{\mathfrak m}$ est exact à gauche
(resp. exact à droite, exact) pour tout idéal maximal $\mathfrak m$ de $A$,
$T$ est exact à gauche (resp. exact à droite, exact).
\end{proposition}

\begin{proof}
On sait en effet que lorsque deux sous-modules $N$, $P$ d'un $A$-module $M$
sont tels que $N_{\mathfrak m}=P_{\mathfrak m}$ pour tout idéal maximal
$\mathfrak m$ de $A$, on a $N=P$ (Bourbaki, \emph{Alg. comm.}, chap. II,
\S~3, no~3, th.~1).
\end{proof}

\subsection{Caractérisations du foncteur produit tensoriel}
\label{subsection:III.7.2.fr}

\begin{env}[7.2.1]
Soient $A$ un anneau (non nécessairement commutatif), $M$ (resp. $N$) un
$A$-module à gauche (resp. à droite), $P$ un $\mathbf Z$-module. Rappelons
que la donnée d'un $\mathbf Z$-homomorphisme
\[
  v:N\otimes_A M\longrightarrow P
\]
équivaut à la donnée d'une application $\mathbf Z$-bilinéaire
\[
  u:N\times M\longrightarrow P
\]
telle que
\[
  u(ta,x)=u(t,ax)
  \qquad\text{pour }a\in A,\ t\in N,\ x\in M,
\]
les deux applications étant reliées par $v(t\otimes x)=u(t,x)$. D'autre part,
la donnée de $u$ équivaut à celle d'un
\oldpage[III]{177}
$\mathbf Z$-homomorphisme $x\mapsto f_x$ de $M$ dans
$\operatorname{Hom}_{\mathbf Z}(N,P)$ tel que
\[
  f_{ax}(t)=f_x(ta)
  \qquad\text{pour }a\in A,\ t\in N,\ x\in M,
\]
les deux applications étant reliées par $u(t,x)=f_x(t)$.
\end{env}

\begin{env}[7.2.2]
Soit $T:\mathbf{Ab}_A\to\mathbf{Ab}$ un foncteur covariant additif. Nous
allons définir pour tout $A$-module à gauche $M$, un homomorphisme canonique
fonctoriel en $M$, de $\mathbf Z$-modules
\[
  t_M:T(A_s)\otimes_A M\longrightarrow T(M).
  \tag{7.2.2.1}\label{III.7.2.2.1.fr}
\]
Il suffira pour cela, en vertu de (7.2.1), de définir un
$\mathbf Z$-homomorphisme $x\mapsto t'_M(x)$ de $M$ dans
$\operatorname{Hom}_{\mathbf Z}(T(A_s),T(M))$, tel que l'on ait
\[
  t'_M(ax)(y)=t'_M(x)(ya)
  \qquad\text{pour }a\in A,\ x\in M,\ y\in T(A_s).
\]
Notons pour cela que
$\operatorname{Hom}_{\mathbf Z}(T(A_s),T(M))$ est canoniquement muni d'une
structure de $A$-module à gauche provenant de la structure de $A$-module à
droite de $T(A_s)$, la loi externe étant telle que si
$a\in A$, $v\in\operatorname{Hom}_{\mathbf Z}(T(A_s),T(M))$,
\[
  (a\cdot v)(y)=v(ya)
  \qquad\text{pour }y\in T(A_s).
\]
Cela étant, nous définirons $t'_M$ comme composé des deux homomorphismes
canoniques
\[
  M\xrightarrow{\sim}\operatorname{Hom}_A(A_s,M)
  \xrightarrow{T}
  \operatorname{Hom}_{\mathbf Z}(T(A_s),T(M)),
\]
la seconde flèche étant l'application $u\mapsto T(u)$, la première
l'isomorphisme canonique de $A$-modules $x\mapsto\theta_x$ tel que
$\theta_x(\xi)=\xi x$ pour $\xi\in A$, $x\in M$. On a
$\theta_{ax}=\theta_x\circ h_a$, donc
\[
  T(\theta_{ax})
  =
  T(\theta_x\circ h_a)
  =
  T(\theta_x)\circ T(h_a)
\]
et par suite, pour $y\in T(A_s)$,
\[
  T(\theta_{ax})(y)
  =
  T(\theta_x)(T(h_a)(y))
  =
  T(\theta_x)(ya)
\]
par définition de la loi externe sur $T(A_s)$, ce qui prouve l'existence de
$t_M$; il est immédiat de vérifier que cet homomorphisme est fonctoriel en
$M$, c'est-à-dire que pour tout homomorphisme $w:M\to M'$ de $A$-modules à
gauche, le diagramme
\[
\xymatrix@C=30pt@R=30pt{
  T(A_s)\otimes_A M
  \ar[r]^-{t_M}\ar[d]_{1\otimes w}
  &
  T(M)\ar[d]^{T(w)}
  \\
  T(A_s)\otimes_A M'
  \ar[r]_-{t_{M'}}
  &
  T(M')
}
\tag{7.2.2.2}\label{III.7.2.2.2.fr}
\]
est commutatif.

La fonctorialité de l'homomorphisme (7.2.2.1) montre que lorsque $A$ est
commutatif, c'est un homomorphisme de $A$-modules (cf. (7.1.2)).
\end{env}

\begin{env}[7.2.3]
Lorsque $A$ est commutatif, on peut plus généralement définir un
homomorphisme canonique de $A$-modules
\[
  T(N)\otimes_A M\longrightarrow T(N\otimes_A M)
  \tag{7.2.3.1}\label{III.7.2.3.1.fr}
\]
pour tout $A$-module $N$; il suffit dans la construction de (7.2.2) de
remplacer l'homomorphisme $\theta_x$ par l'homomorphisme de $A$-modules
\[
  N\longrightarrow N\otimes_A M
\]
qui à tout $y\in N$ fait correspondre $y\otimes x$. Il est immédiat que cet
homomorphisme est fonctoriel en $M$ et $N$.

\oldpage[III]{178}
En particulier, si $B$ est une $A$-algèbre (non nécessairement commutative),
on a un homomorphisme fonctoriel en $M$
\[
  (T(M))_{(B)}
  =
  T(M)\otimes_A B
  \longrightarrow
  T(M\otimes_A B)
  =
  T_{(B)}(M_{(B)})
  \tag{7.2.3.2}\label{III.7.2.3.2.fr}
\]
qui, en vertu de la fonctorialité de (7.2.3.1) en $M$, est un homomorphisme
de $B$-modules.

On a en outre le diagramme commutatif
\[
\xymatrix@C=30pt@R=30pt{
  T(A)\otimes_A M
  \ar[r]^-{t_M}\ar[d]
  &
  T(M)\ar[d]
  \\
  T_{(B)}(B_s)\otimes_B M_{(B)}
  \ar[r]_-{t_{M_{(B)}}}
  &
  T_{(B)}(M_{(B)})
}
\tag{7.2.3.3}\label{III.7.2.3.3.fr}
\]
où la flèche verticale de droite est l'homomorphisme composé
\[
  T(M)
  \longrightarrow
  T(M)\otimes_A B
  \longrightarrow
  T(M\otimes_A B)
  =
  T_{(B)}(M_{(B)})
\]
de (7.2.3.2) et de l'homomorphisme canonique; quant à la flèche verticale de
gauche de (7.2.3.3), c'est l'homomorphisme
\[
  T(A)\otimes_A M
  \longrightarrow
  T_{(B)}(B_s)\otimes_B(B\otimes_A M)
  =
  T_{(B)}(B_s)\otimes_A M
\]
où
\[
  T(A)\longrightarrow T_{(B)}(B_s)=T(B)
\]
est $T(\rho)$, $\rho$ étant considéré comme homomorphisme de $A$-modules
$A\to B_{[\rho]}$.
\end{env}

\begin{lemma}[7.2.4]
Si $T$ est un foncteur covariant additif de $\mathbf{Ab}_A$ dans
$\mathbf{Ab}$, commutant avec les sommes directes, l'homomorphisme canonique
$t_L$ (7.2.2.1) est un isomorphisme pour tout $A$-module libre $L$.
\end{lemma}

\begin{proof}
En effet, on a
\[
  L=\bigoplus_{\alpha\in I}L_\alpha
\]
où $L_\alpha$ est isomorphe à $A_s$ pour tout $\alpha\in I$; la définition de
$t_M$ donnée dans (7.2.2) montre que
\[
  t_L=\bigoplus_{\alpha\in I}t_{L_\alpha},
\]
puisque
\[
  T:
  \operatorname{Hom}_A(A_s,L)
  \longrightarrow
  \operatorname{Hom}_{\mathbf Z}(T(A_s),T(L))
\]
est somme directe des applications $\mathbf Z$-linéaires
\[
  T_\alpha:
  \operatorname{Hom}_A(A_s,L_\alpha)
  \longrightarrow
  \operatorname{Hom}_{\mathbf Z}(T(A_s),T(L_\alpha))
\]
en vertu de l'hypothèse sur $T$. On est donc ramené à prouver le lemme pour
$L=A_s$; mais $t_L$ n'est autre alors que l'isomorphisme canonique
\[
  T(A_s)\otimes_A A_s\xrightarrow{\sim}T(A_s)
\]
valable pour tout $A$-module à droite.
\end{proof}

\begin{proposition}[7.2.5]
Soit $T$ un foncteur covariant additif de $\mathbf{Ab}_A$ dans
$\mathbf{Ab}$, commutant aux sommes directes. Les conditions suivantes sont
équivalentes :
\begin{enumerate}
  \item[a)] $T$ est exact à droite.
  \item[b)] L'homomorphisme canonique $t_M$ (7.2.2.1) est un isomorphisme pour
    tout $A$-module à gauche $M$.
  \item[$b'$)] $T$ est semi-exact et l'homomorphisme $t_M$ est surjectif pour
    tout $A$-module à gauche $M$.
  \item[c)] $T$ est isomorphe à un foncteur en $M$ de la forme
    $N\otimes_A M$, où $N$ est un $A$-module à droite.
\end{enumerate}
\end{proposition}

\begin{proof}
Il est clair que $b)$ implique $c)$ et que $c)$ implique $a)$; montrons que
$a)$ implique $b)$.
\oldpage[III]{179}
Posons
\[
  T'(M)=T(A_s)\otimes_A M
\]
pour tout $A$-module à gauche $M$. Il existe une suite exacte
\[
  L'\longrightarrow L\longrightarrow M\longrightarrow0,
\]
où $L$ et $L'$ sont deux $A$-modules à gauche libres; comme $T$ et $T'$ sont
exacts à droite, on a donc le diagramme commutatif
\[
\xymatrix@C=22pt@R=25pt{
  T'(L')\ar[r]\ar[d]^{t_{L'}}
  &
  T'(L)\ar[r]\ar[d]^{t_L}
  &
  T'(M)\ar[r]\ar[d]^{t_M}
  &
  0
  \\
  T(L')\ar[r]
  &
  T(L)\ar[r]
  &
  T(M)\ar[r]
  &
  0
}
\]
où les deux lignes sont exactes; comme $t_L$ et $t_{L'}$ sont des
isomorphismes en vertu de (7.2.4), il en est de même de $t_M$ par le lemme des
cinq. Enfin, il est clair que $b)$ entraîne $b')$. Pour montrer que $b')$
entraîne $a)$, il suffit de prouver le lemme suivant.
\end{proof}

\begin{lemma}[7.2.5.1]
Soient $\mathcal K$, $\mathcal K'$ deux catégories abéliennes, $F$, $G$ deux
foncteurs covariants additifs de $\mathcal K$ dans $\mathcal K'$,
$f:F\to G$ un morphisme fonctoriel (T, I, 1.2) tel que, pour tout objet $E$
de la catégorie $\mathcal K$,
\[
  f_E:F(E)\longrightarrow G(E)
\]
soit un épimorphisme. Alors, si $F$ est exact à droite et $G$ semi-exact,
$G$ est exact à droite.
\end{lemma}

\begin{proof}
Tout revient en effet à démontrer que pour tout épimorphisme
$v:E'\to E$ dans $\mathcal K$,
\[
  G(v):G(E')\longrightarrow G(E)
\]
est un épimorphisme; or, on a le diagramme commutatif
\[
\xymatrix@C=35pt@R=30pt{
  F(E')\ar[r]^{F(v)}\ar[d]_{f_{E'}}
  &
  F(E)\ar[d]^{f_E}
  \\
  G(E')\ar[r]_{G(v)}
  &
  G(E)
}
\]
dans lequel $F(v)$, $f_{E'}$ et $f_E$ sont des épimorphismes; il en est donc
de même de $G(v)$.
\end{proof}

\begin{remark}[7.2.6]
Pour tout $A$-module à droite $N$, posons
\[
  T_N(M)=N\otimes_A M
\]
pour tout $A$-module à gauche $M$, de sorte que $T_N$ est un foncteur
covariant additif de $\mathbf{Ab}_A$ dans $\mathbf{Ab}$, exact à droite et
commutant aux sommes directes. Si on identifie canoniquement $T_N(A_s)$ à
$N$, on vérifie aussitôt que l'homomorphisme correspondant (7.2.2.1) devient
l'identité. On en conclut que le $A$-module à droite $N$ de l'énoncé de
(7.2.5, $c)$) est déterminé à un isomorphisme unique près et est
canoniquement isomorphe à $T(A_s)$. On peut encore dire que les morphismes
fonctoriels
\[
  T\longmapsto T(A_s),
  \qquad
  N\longmapsto T_N
\]
constituent une équivalence (T, I, 1.2) de la catégorie des $A$-modules à
droite et de la catégorie des foncteurs covariants additifs
$\mathbf{Ab}_A\to\mathbf{Ab}$ qui sont exacts à droite et commutent aux
sommes directes.
\end{remark}

\begin{proposition}[7.2.7]
Soient $A$ un anneau artinien à gauche, dont le quotient par son radical
$\mathfrak m$ est un corps $k$. Soit $T$ un foncteur covariant additif de
$\mathbf{Ab}_A$ dans $\mathbf{Ab}$, commutant aux sommes directes. Les
conditions de (7.2.5) sont alors aussi équivalentes à
\begin{enumerate}
  \item[d)] $T$ est semi-exact et l'homomorphisme
    \[
      T(\varepsilon):T(A_s)\longrightarrow T(k)
    \]
    déduit de l'homomorphisme canonique $\varepsilon:A_s\to k$ est surjectif.
\end{enumerate}
\end{proposition}

\begin{proof}
\oldpage[III]{180}
Il est clair que la condition $b')$ de (7.2.5) entraîne $d)$; prouvons que
$d)$ entraîne $b')$.

Il existe un entier $n$ tel que $\mathfrak m^n=0$; posons, pour tout
$A$-module $M$, $M_h=\mathfrak m^hM$; nous prouverons par récurrence
descendante sur $h$ que $t_{M_h}$ est surjectif. La proposition est évidente
pour $h=n$; pour $h<n$, on a une suite exacte
\[
  0\longrightarrow M_{h+1}
  \longrightarrow M_h
  \longrightarrow M_h/M_{h+1}
  \longrightarrow 0
\]
et l'hypothèse de récurrence entraîne que $t_{M_{h+1}}$ est surjectif.
D'autre part, $M_h/M_{h+1}$ est annulé par $\mathfrak m$ et est donc un
$(A/\mathfrak m)$-module, autrement dit est somme directe de $A$-modules
isomorphes à $k$. Pour prouver que $t_{M_h/M_{h+1}}$ est surjectif, il suffit
donc de prouver que $t_k$ l'est, puisque $T$ commute aux sommes directes. Or,
en vertu de la commutativité du diagramme
\[
\xymatrix@C=35pt@R=30pt{
  T(A_s)\otimes_A A_s
  \ar[r]^-{t_{A_s}}\ar[d]_{1\otimes\varepsilon}
  &
  T(A_s)\ar[d]^{T(\varepsilon)}
  \\
  T(A_s)\otimes_A k
  \ar[r]_-{t_k}
  &
  T(k)
}
\]
et de (7.2.4), l'hypothèse $d)$ entraîne que $t_k$ est bien surjectif. Pour
terminer la démonstration, il suffira de prouver que si l'on a une suite
exacte
\[
  0\longrightarrow M'
  \longrightarrow M
  \xrightarrow{v}M''
  \longrightarrow 0
\]
de $A$-modules, telle que $t_{M'}$ et $t_{M''}$ sont surjectifs, alors $t_M$
est surjectif. Or, on a un diagramme commutatif
\[
\xymatrix@C=27pt@R=30pt{
  T'(M')\ar[r]\ar[d]_{t_{M'}}
  &
  T'(M)\ar[r]\ar[d]_{t_M}
  &
  T'(M'')\ar[r]\ar[d]_{t_{M''}}
  &
  0\ar[d]
  \\
  T(M')\ar[r]
  &
  T(M)\ar[r]_-{T(v)}
  &
  T(M'')\ar[r]
  &
  \operatorname{Coker}(T(v))
}
\]
dans lequel les deux lignes sont exactes, en vertu de l'hypothèse que $T$ est
semi-exact. Comme par l'hypothèse de récurrence $t_{M'}$ et $t_{M''}$ sont
des épimorphismes et que la dernière flèche verticale est un monomorphisme,
le lemme des cinq (M, I, 1.1) montre que $t_M$ est un épimorphisme.
\end{proof}

\subsection{Critères d'exactitude des foncteurs homologiques de modules}
\label{subsection:III.7.3.fr}

\begin{proposition}[7.3.1]
Soient $A$ un anneau (non nécessairement commutatif), $T_\bullet$ un foncteur
homologique covariant (T, II, 2.1) de la catégorie $\mathbf{Ab}_A$ dans la
catégorie $\mathbf{Ab}$, commutant aux sommes directes. Soit $p$ un entier tel
que $T_p$ et $T_{p-1}$ soient définis. Les conditions suivantes sont
équivalentes :
\begin{enumerate}
  \item[a)] $T_p$ est exact à droite.
  \item[b)] $T_{p-1}$ est exact à gauche.
  \oldpage[III]{181}
  \item[c)] Pour tout $A$-module à gauche $M$, l'homomorphisme fonctoriel
    canonique (7.2.2.1)
    \[
      T_p(A_s)\otimes_A M\longrightarrow T_p(M)
      \tag{7.3.1.1}\label{III.7.3.1.1.fr}
    \]
    est un isomorphisme.
  \item[d)] Pour tout $A$-module à gauche $M$, l'homomorphisme (7.3.1.1) est
    un épimorphisme.
  \item[e)] $T_p$ est isomorphe à un foncteur
    $M\mapsto N\otimes_A M$, où $N$ est un $A$-module à droite.
\end{enumerate}
Si en outre les conditions de (7.2.7) sur $A$ et $\mathfrak m$ sont vérifiées,
les conditions précédentes sont aussi équivalentes à
\begin{enumerate}
  \item[f)] L'homomorphisme canonique
    \[
      T_p(\varepsilon):T_p(A_s)\longrightarrow T_p(k)
    \]
    est un épimorphisme.
\end{enumerate}
\end{proposition}

\begin{proof}
Comme par définition d'un foncteur homologique, $T_i$ est semi-exact pour tous
les $i$ tels que $T_i$ soit défini, et que, pour toute suite exacte
\[
  0\longrightarrow M'\xrightarrow{u}M\xrightarrow{v}M''\longrightarrow0,
\]
on a
\[
  \operatorname{Ker}(T_{i-1}(u))
  =
  \operatorname{Coker}(T_i(v)),
\]
il est clair que $a)$ et $b)$ sont équivalentes et les autres assertions
résultent trivialement de (7.2.5) et (7.2.7).
\end{proof}

\begin{corollary}[7.3.2]
Soit $A$ un anneau commutatif. Avec les notations de (7.3.1), supposons $T_p$
exact à droite. Si $f\in A$ n'appartient à l'annulateur d'aucun élément
$\ne0$ d'un $A$-module $M$, alors $f$ n'appartient à l'annulateur d'aucun
élément $\ne0$ de $T_{p-1}(M)$. En particulier, si $A$ est intègre, le
$A$-module $T_{p-1}(A)$ est sans torsion.
\end{corollary}

\begin{proof}
En effet, si $h_f$ désigne l'homothétie $x\mapsto fx$ de $M$, l'hypothèse
signifie que $h_f$ est injectif; il en est donc de même de $T_{p-1}(h_f)$ par
la condition $b)$ de (7.3.1).
\end{proof}

\begin{proposition}[7.3.3]
Soient $A$ un anneau, $T_\bullet$ un foncteur homologique covariant de
$\mathbf{Ab}_A$ dans $\mathbf{Ab}$, commutant aux sommes directes. Soit $p$ un
entier tel que $T_{p-1}$, $T_p$ et $T_{p+1}$ soient définis. Les conditions
suivantes sont équivalentes :
\begin{enumerate}
  \item[a)] $T_p$ est exact.
  \item[b)] $T_{p+1}$ et $T_p$ sont exacts à droite.
  \item[c)] $T_p$ et $T_{p-1}$ sont exacts à gauche.
  \item[d)] $T_{p+1}$ est exact à droite et $T_{p-1}$ est exact à gauche.
  \item[e)] Pour tout $A$-module $M$, les homomorphismes canoniques
    \[
      T_i(A_s)\otimes_A M\longrightarrow T_i(M)
      \tag{7.3.3.1}\label{III.7.3.3.1.fr}
    \]
    sont des isomorphismes pour $i=p$ et $i=p+1$.
  \item[$e'$)] Pour tout $A$-module $M$, les homomorphismes canoniques
    (7.3.3.1) sont des épimorphismes pour $i=p$ et $i=p+1$.
  \item[f)] Pour tout $A$-module $M$, l'homomorphisme (7.3.3.1) est un
    isomorphisme pour $i=p$ et $T_p(A_s)$ est un $A$-module à droite plat.
  \item[$f'$)] Pour tout $A$-module $M$, l'homomorphisme (7.3.3.1) est un
    épimorphisme pour $i=p$ et $T_p(A_s)$ est un $A$-module à droite plat.
\end{enumerate}
\end{proposition}

\begin{proof}
L'équivalence des conditions $a)$, $b)$, $c)$, $d)$ résulte de l'équivalence
des conditions $a)$ et $b)$ de (7.3.1). L'équivalence de $b)$, $e)$ et $e')$
résulte de l'équivalence de $a)$, $c)$ et $d)$ dans (7.3.1). Enfin, dire que
$T_p(A_s)$ est plat signifie que le foncteur
\[
  M\mapsto T_p(A_s)\otimes_A M
\]
est exact à gauche; l'équivalence de $a)$, $f)$ et $f')$ résulte encore de
l'équivalence de $a)$, $c)$, $d)$ dans (7.3.1).
\end{proof}

\oldpage[III]{182}

\begin{corollary}[7.3.4]
Supposons $A$ commutatif, $T_p$ exact et supposons en outre que $T_p(A)$ soit
un $A$-module de présentation finie. Alors la fonction
\[
  x\mapsto\operatorname{rang}_{\kappa(x)}(T_p(\kappa(x)))
\]
est localement constante dans $X=\operatorname{Spec}(A)$, donc constante si
$\operatorname{Spec}(A)$ est connexe.
\end{corollary}

\begin{proof}
En effet, comme $T_p(A)$ est un $A$-module plat en vertu de (7.3.3, $f)$), il
est projectif de type fini, et $\bigl(T_p(A)\bigr)^{\sim}$ est donc un
$\mathscr O_X$-Module localement libre (Bourbaki, \emph{Alg. comm.}, chap. II,
\S 5, n\textsuperscript{o} 2, th. 1); on a en outre
\[
  T_p(\kappa(x))=T_p(A)\otimes_A\kappa(x)
\]
(7.3.3, $e)$) et l'on sait que le rang au point $x$ du $A$-module $T_p(A)$ est
localement constant (\emph{loc. cit.}), d'où le corollaire.
\end{proof}

\begin{proposition}[7.3.5]
Supposons que $A$ soit un anneau artinien à gauche dont le quotient par son
radical $\mathfrak m$ est un corps $k$. Alors les conditions de (7.3.3) sont
encore équivalentes à chacune des suivantes :
\begin{enumerate}
  \item[g)] L'homomorphisme canonique
    \[
      T_i(\varepsilon):T_i(A_s)\longrightarrow T_i(k)
    \]
    est un épimorphisme pour $i=p$ et $i=p+1$.
  \item[h)] $T_p(\varepsilon)$ est un épimorphisme et $T_p(A_s)$ est un
    $A$-module à droite plat (ou, ce qui revient au même (Bourbaki,
    \emph{Alg. comm.}, chap. II, \S 3, n\textsuperscript{o} 2, cor. 2 de la
    prop. 5) un $A$-module libre).
\end{enumerate}
Supposons en outre que $A$ soit commutatif et le $A$-module $T_p(k)$ de
longueur finie $d$. Alors les conditions précédentes équivalent aussi à
chacune des suivantes :
\begin{enumerate}
  \item[i)] Pour tout $A$-module $M$ de longueur finie, on a
    \[
      \operatorname{long}(T_p(M))=d\mathbin{.}\operatorname{long}(M).
      \tag{7.3.5.1}\label{III.7.3.5.1.fr}
    \]
  \item[j)] On a
    \[
      \operatorname{long}(T_p(A))=d\mathbin{.}\operatorname{long}(A).
      \tag{7.3.5.2}\label{III.7.3.5.2.fr}
    \]
\end{enumerate}
\end{proposition}

L'équivalence de $g)$ et $h)$ avec les conditions de (7.3.3) découle aussitôt
de (7.2.7). Pour démontrer les autres assertions, nous utiliserons le lemme
suivant :

\begin{lemma}[7.3.5.3]
Soient $\mathcal K$, $\mathcal K'$ deux catégories abéliennes,
$F:\mathcal K\to\mathcal K'$ un foncteur
covariant additif; on suppose que $F$ soit semi-exact, et que, pour tout objet
simple $S$ de $\mathcal K$, $F(S)$ soit un objet de longueur finie dans
$\mathcal K'$. Alors, pour tout objet $E$ de longueur finie dans $\mathcal K$,
$F(E)$ est de longueur finie dans $\mathcal K'$. Pour toute suite exacte
\[
  0\longrightarrow E'\xrightarrow{u}E\xrightarrow{v}E''\longrightarrow0
\]
d'objets de longueur finie dans $\mathcal K$, on a
\[
  \operatorname{long}F(E)
  \leq
  \operatorname{long}F(E')+\operatorname{long}F(E'')
  \tag{7.3.5.4}\label{III.7.3.5.4.fr}
\]
et pour que les deux membres de (7.3.5.4) soient égaux, il faut et il suffit
que la suite
\[
  0\longrightarrow F(E')\longrightarrow F(E)\longrightarrow F(E'')
  \longrightarrow0
\]
soit exacte.
\end{lemma}

\begin{proof}
En effet, la suite
\[
  F(E')\xrightarrow{F(u)}F(E)\xrightarrow{F(v)}F(E'')
\]
est exacte par hypothèse; si l'on suppose $F(E')$ et $F(E'')$ de longueur
finie, il en est de même de $\operatorname{Im}(F(u))$ et de
$\operatorname{Im}(F(v))$ et comme
$\operatorname{Ker}(F(v))=\operatorname{Im}(F(u))$, $F(E)$ est de longueur
finie et l'on a
\[
  \operatorname{long}F(E)
  =
  \operatorname{long}\operatorname{Im}(F(u))
  +
  \operatorname{long}\operatorname{Im}(F(v))
  \leq
  \operatorname{long}F(E')+\operatorname{long}F(E'').
  \tag{7.3.5.5}\label{III.7.3.5.5.fr}
\]
\oldpage[III]{183}
Par récurrence sur la longueur de $E$, cela prouve déjà la première assertion;
en outre, les deux membres de (7.3.5.5) ne peuvent être égaux que si
\[
  \operatorname{long}\operatorname{Im}(F(u))=\operatorname{long}F(E')
\]
(ce qui équivaut à
$\operatorname{long}\operatorname{Ker}(F(u))=0$, ou
$\operatorname{Ker}(F(u))=0$) et
\[
  \operatorname{long}\operatorname{Im}(F(v))=\operatorname{long}F(E'')
\]
(ce qui équivaut à
$\operatorname{long}\operatorname{Coker}(F(v))=0$, ou
$\operatorname{Coker}(F(v))=0$).
\end{proof}

Notons maintenant que si $M$ est un $A$-module de longueur finie ($A$ étant
commutatif), les quotients d'une suite de Jordan-Hölder de $M$ sont
nécessairement isomorphes au $A$-module $k$, donc on déduit de (7.3.5.4), par
récurrence sur la longueur de $M$
\[
  \operatorname{long}T_p(M)\leq d\mathbin{.}\operatorname{long}(M).
  \tag{7.3.5.6}\label{III.7.3.5.6.fr}
\]
En outre, il résulte de (7.3.5.3) que si $T_p$ est exact, on a l'égalité
(7.3.5.1); donc la condition $a)$ de (7.3.3) entraîne $i)$; il est clair que
$i)$ entraîne $j)$, et il reste à prouver le

\begin{lemma}[7.3.5.7]
La relation
\[
  \operatorname{long}T_p(A)=d\mathbin{.}\operatorname{long}A
\]
entraîne que $T_p(\varepsilon)$ est un épimorphisme et que $T_p(A)$ est un
$A$-module plat.
\end{lemma}

\begin{proof}
En effet, partant de la suite exacte
\[
  0\longrightarrow\mathfrak m\longrightarrow A\longrightarrow k
  \longrightarrow0,
\]
il résulte de (7.3.5.4) et (7.3.5.6) que l'on a
\[
  \operatorname{long}T_p(A)
  \leq
  \operatorname{long}T_p(\mathfrak m)+\operatorname{long}T_p(k)
  \leq
  d\bigl(\operatorname{long}\mathfrak m+\operatorname{long}k\bigr)
  =d\mathbin{.}\operatorname{long}A
\]
et que l'égalité ne peut avoir lieu (7.3.5.3) que si la suite
\[
  0\longrightarrow T_p(\mathfrak m)\longrightarrow T_p(A)
  \longrightarrow T_p(k)\longrightarrow0
  \tag{7.3.5.8}\label{III.7.3.5.8.fr}
\]
est exacte. En vertu de (7.2.7) et (7.2.5), $T_p$ est isomorphe à un foncteur
$M\mapsto N\otimes_A M$, et l'exactitude de la suite (7.3.5.8) montre, en vertu
de la suite exacte des $\operatorname{Tor}$, que l'on a
$\operatorname{Tor}_1^A(N,k)=0$. On en conclut que $N=T_p(A)$ est un
$A$-module plat (0, 10.1.3).
\end{proof}

\begin{lemma}[7.3.6]
Soient $A$ un anneau, $T_\bullet$ un foncteur homologique covariant de
$\mathbf{Ab}_A$ dans $\mathbf{Ab}$, commutant aux sommes directes. Supposons
$T_p$ et $T_{p+1}$ définis, et $T_p$ exact à gauche. Pour que $T_{p+1}$ soit
exact, il faut et il suffit que $T_{p+1}(A_s)$ soit un $A$-module à droite
plat.
\end{lemma}

\begin{proof}
En effet, on sait d'après (7.3.1) que l'homomorphisme canonique
\[
  T_{p+1}(A_s)\otimes_A M\longrightarrow T_{p+1}(M)
\]
est un isomorphisme de foncteurs; il suffit d'appliquer la définition d'un
$A$-module plat.
\end{proof}

\begin{proposition}[7.3.7]
Soient $A$ un anneau, $T_\bullet$ un foncteur homologique covariant de
$\mathbf{Ab}_A$ dans $\mathbf{Ab}$, commutant aux sommes directes. On suppose
qu'il existe $i_0$ tel que $T_i$ soit exact pour $i\leq i_0$. Alors, pour tout
entier $p>i_0$, les conditions suivantes sont équivalentes :
\begin{enumerate}
  \item[a)] $T_q$ est exact pour $q\leq p$.
  \item[b)] $T_q(A_s)$ est un $A$-module à droite plat pour $q\leq p$.
  \item[c)] Pour tout $A$-module $M$, l'homomorphisme canonique
    \[
      T_q(A_s)\otimes_A M\longrightarrow T_q(M)
    \]
    est surjectif pour $q\leq p+1$.
\end{enumerate}
\end{proposition}

\begin{proof}
L'équivalence de $a)$ et $b)$ résulte de (7.3.6) par récurrence sur $q$,
puisque $T_{i_0}$ est exact par hypothèse; l'équivalence de $a)$ et $c)$
résulte de l'équivalence des conditions $a)$ et $e')$ dans (7.3.3).
\end{proof}

\oldpage[III]{184}

\begin{env}[7.3.8]
Si $A$ est un anneau commutatif, $B$ une $A$-algèbre (non nécessairement
commutative), $T_\bullet$ un foncteur homologique covariant de
$\mathbf{Ab}_A$ dans $\mathbf{Ab}$, il résulte des définitions (7.1.3) que le
foncteur de $\mathbf{Ab}_B$ dans $\mathbf{Ab}$ obtenu par extension des
scalaires de $A$ à $B$, et que nous noterons
\[
  T_\bullet^{(B)}=(T_i^{(B)}),
\]
est encore un foncteur homologique.
\end{env}

\begin{corollary}[7.3.9]
Supposons que $T_\bullet$ vérifie les conditions générales de (7.3.7) et
commute aux limites inductives, et en outre que $A$ soit un anneau intègre et
tous les $T_n(A)$ des $A$-modules de présentation finie. Alors, pour tout
entier $N$, il existe un $f\in A-\{0\}$ tel que le foncteur
\[
  T_p^{(A_f)}:\mathbf{Ab}_{A_f}\longrightarrow\mathbf{Ab}
\]
soit exact pour $p\leq N$.
\end{corollary}

\begin{proof}
Par hypothèse, $T_i$ est exact pour $i\leq i_0$, donc $T_i(A)$ est plat pour
ces valeurs de $i$. En vertu de (7.3.7, $b)$), il suffit de prendre $f$ tel que
\[
  T_p^{(A_f)}(A_f)=T_p(A_f)
\]
soit un $A_f$-module libre pour $i_0<p\leq N$. Or, on a
$T_p(A_f)=(T_p(A))_f$ puisque $T_p$ commute aux limites inductives (7.1.4). Si
$x$ est le point générique de $\operatorname{Spec}(A)$, $(T_p(A))_x$ est un
espace vectoriel de dimension finie sur le corps des fractions de $A$. Comme
chaque $T_p(A)$ est de présentation finie, il existe bien un $f$ ayant la
propriété voulue (Bourbaki, \emph{Alg. comm.}, chap. II, \S 5,
n\textsuperscript{o} 1, cor. de la prop. 2).

On notera que s'il n'y a qu'un nombre fini d'indices $i$ tels que $T_i\ne0$,
il existe $f\in A-\{0\}$ tel que tous les $T_p^{(A_f)}$ soient exacts.
\end{proof}

\begin{corollary}[7.3.10]
Supposons que $T_\bullet$ vérifie les conditions générales de (7.3.7) et
commute aux limites inductives, que $A$ soit commutatif et noethérien et les
$T_n(A)$ des $A$-modules de type fini. Alors, pour tout entier $N$, il existe
un ouvert dense $U$ de $\operatorname{Spec}(A)$ tel que, pour tout $p\leq N$,
la fonction
\[
  x\mapsto\operatorname{rang}_{\kappa(x)}(T_p(\kappa(x)))
\]
soit constante dans $U$.
\end{corollary}

\begin{proof}
Soit $\mathfrak p$ un idéal premier minimal de $A$; par hypothèse, l'anneau
$B=A/\mathfrak p$ est intègre et $\operatorname{Spec}(B)$ s'identifie à une
composante irréductible de l'espace topologique $\operatorname{Spec}(A)$. Nous
allons montrer, par récurrence sur $p\leq N$ qu'il existe
$f_p\in B-\{0\}$ tel que si on pose $B'=B_{f_p}$, $T_i^{(B')}$ soit exact et
les $T_i(B')$ des $B'$-modules de type fini pour $i\leq p$. La proposition est
vraie pour $p\leq i_0$ en vertu de l'hypothèse, en prenant $f_p=1$ (donc
$B'=B=A/\mathfrak p$) car $T_p$ étant alors exact, $T_p(B)$ est isomorphe à
$T_p(A)/T_p(\mathfrak p)$, donc est un $A$-module (et \emph{a fortiori} un
$B$-module) de type fini. Raisonnons par récurrence sur $p$; $f_p$ est l'image
canonique dans $B$ d'un élément $g_p\in A$, et si l'on pose $A'=A_{g_p}$, on a
$B'=A'/\mathfrak p'$ où $\mathfrak p'$ est un idéal premier minimal de $A'$,
égal à $\mathfrak p_{g_p}$. Comme on a
\[
  T_i(A_{g_p})=(T_i(A))_{g_p},
\]
les $T_i(A')$ sont des $A'$-modules de type fini, donc le foncteur
$T_\bullet^{(A')}$ vérifie les mêmes hypothèses que $T_\bullet$, mais en
remplaçant $i_0$ par $p$. On peut donc se borner au cas où $A'=A$, et où $T_p$
est exact; la suite exacte
\[
  0\longrightarrow\mathfrak p\longrightarrow A\longrightarrow A/\mathfrak p
  \longrightarrow0
\]
donne alors la suite exacte
\[
  T_{p+1}(A)\longrightarrow T_{p+1}(A/\mathfrak p)
  \xrightarrow{\partial}T_p(\mathfrak p)\longrightarrow T_p(A),
\]
et comme $T_p$ est exact, la dernière flèche de droite est injective, donc
$T_{p+1}(A/\mathfrak p)$ est un quotient de $T_{p+1}(A)$ et est par suite de
type fini. Notons maintenant que le raisonnement de (7.3.9) n'a utilisé le fait
que les $T_p(A)$ sont de type fini que pour $p\leq N$; on peut donc l'appliquer
à l'anneau intègre $B$, au foncteur $T_\bullet^{(B)}$ et à $N=p+1$, ce qui
achève le raisonnement par récurrence. Cela étant, il y a un
$f_N\in B-\{0\}$ tel que $\operatorname{Spec}(B_{f_N})$ soit un ouvert $V$
partout dense dans $\operatorname{Spec}(B)$ ne rencontrant aucune autre
composante irréductible de $\operatorname{Spec}(A)$. Si la proposition est
\oldpage[III]{185}
démontrée pour $B_{f_N}$, on aura un ouvert $W$ partout dense dans $V$ dans
lequel les fonctions de l'énoncé seront constantes, puisque
$A_x=(B_{f_N})_x$ pour tout $x\in W$. En faisant le même raisonnement pour
toute composante irréductible de $\operatorname{Spec}(A)$, le corollaire sera
démontré. On peut donc se borner au cas où $A$ est intègre; le raisonnement de
(7.3.9) prouve alors l'existence d'un $f\in A-\{0\}$ tel que les $T_p(A_f)$
soient des $A_f$-modules libres de type fini pour $p\leq N$, ce qui entraîne
la conclusion de (7.3.10) en vertu de (7.3.4).
\end{proof}

\begin{proposition}[7.3.11]
Soient $A$ un anneau commutatif local, $k$ son corps résiduel, $T_\bullet$ un
foncteur homologique covariant de $\mathbf{Ab}_A$ dans $\mathbf{Ab}$,
commutant aux sommes directes. On suppose qu'il existe $i_0$ tel que $T_i$ soit
exact pour $i\leq i_0$, et que tous les $T_n(A)$ soient des $A$-modules de
présentation finie. Alors les conditions équivalentes $a)$, $b)$, $c)$ de
(7.3.7) impliquent les deux suivantes, et leur sont équivalentes lorsque
l'anneau est en outre réduit :
\begin{enumerate}
  \item[d)] Pour tout $x\in\operatorname{Spec}(A)$, on a
    \[
      \operatorname{rang}_{\kappa(x)}T_q(\kappa(x))
      =
      \operatorname{rang}_kT_q(k)
      \quad\text{pour }q\leq p.
    \]
  \item[$d'$)] Pour tout point générique $x_j$ d'une composante irréductible
    de $\operatorname{Spec}(A)$, on a
    \[
      \operatorname{rang}_{\kappa(x_j)}T_q(\kappa(x_j))
      =
      \operatorname{rang}_kT_q(k)
      \quad\text{pour }q\leq p.
    \]
\end{enumerate}
\end{proposition}

\begin{proof}
Comme $T_q(A)$ est un $A$-module de présentation finie, la condition $b)$ de
(7.3.7) équivaut à dire que $T_q(A)$ est un $A$-module libre pour $q\leq p$
(Bourbaki, \emph{Alg. comm.}, chap. II, \S 3, n\textsuperscript{o} 2, cor. 2
de la prop. 5); la condition $c)$ implique que
\[
  T_q(\kappa(x))=T_q(A)\otimes_A\kappa(x)
  \quad\text{pour }q\leq p,
\]
donc les conditions équivalentes de (7.3.7) impliquent $d)$, et il est trivial
que $d)$ entraîne $d')$. Reste à prouver que $d')$ implique $a)$ lorsque $A$
est réduit. Raisonnons par récurrence sur $q\leq p$, puisque $T_q$ est exact
pour $q\leq i_0$. Supposons donc $T_k$ exact pour $k\leq q<p$ et montrons que
$T_{q+1}(A)$ est un $A$-module libre. En vertu de l'hypothèse de récurrence,
$T_{q+1}(A)\otimes_A M$ est isomorphe à $T_{q+1}(M)$ pour tout $A$-module $M$,
par la condition $c)$ de (7.3.7) et (7.3.3); appliquant cette propriété à
$M=\kappa(x_j)$ et $M=k$, on trouve, en vertu de l'hypothèse $d')$, que
\[
  \operatorname{rang}_{\kappa(x_j)}
  \bigl(T_{q+1}(A)\otimes_A\kappa(x_j)\bigr)
  =
  \operatorname{rang}_kT_{q+1}(k)
\]
pour tout $i$; mais cela implique que $T_{q+1}(A)$ est libre (Bourbaki,
\emph{Alg. comm.}, chap. II, \S 3, n\textsuperscript{o} 2, prop. 7), ce qui
achève la démonstration.
\end{proof}

Les résultats précédents seront considérablement améliorés pour les foncteurs
homologiques de type particulier que nous allons étudier dans (7.4); on
obtiendra en effet des critères d'exactitude ne faisant intervenir qu'un seul
des $T_p$.

\subsection{Critères d'exactitude pour les foncteurs
$H_\bullet(P_\bullet\otimes_A M)$}
\label{subsection:III.7.4.fr}

\begin{env}[7.4.1]
Soient $A$ un anneau (non nécessairement commutatif), $P_\bullet$ un complexe
de $A$-modules à droite plats. Comme le foncteur
$M\mapsto P_k\otimes_A M$ est alors exact dans $\mathbf{Ab}_A$ pour tout $k$,
le $\partial$-foncteur
\[
  T_\bullet(M)=H_\bullet(P_\bullet\otimes_A M)
  \tag{7.4.1.1}\label{III.7.4.1.1.fr}
\]
\oldpage[III]{186}
est un foncteur homologique de $\mathbf{Ab}_A$ dans $\mathbf{Ab}$,
évidemment $A$-linéaire lorsque $A$ est commutatif (7.1.2), et commutant aux
limites inductives.

Si $A$ est commutatif, alors, pour toute $A$-algèbre $B$, le foncteur
homologique $T_\bullet^{(B)}$ (7.3.8) est donné par définition par
\[
  T_\bullet^{(B)}(N)
  =
  H_\bullet(P_\bullet\otimes_A N_{[\rho]})
  \tag{7.4.1.2}\label{III.7.4.1.2.fr}
\]
où $\rho:A\to B$ est l'homomorphisme définissant la structure d'algèbre de
$B$; comme on peut aussi écrire
\[
  P_\bullet\otimes_A N_{[\rho]}
  =P_\bullet\otimes_A(B\otimes_B N)_{[\rho]}
  =(P_\bullet\otimes_A B)\otimes_B N,
\]
on voit que l'on a
\[
  T_\bullet^{(B)}(N)=H_\bullet(P'_\bullet\otimes_B N)
  \tag{7.4.1.3}\label{III.7.4.1.3.fr}
\]
pour tout $B$-module $N$, $P'_\bullet$ étant le complexe
$P_\bullet\otimes_A B$ de $B$-modules plats (0\textsubscript{I}, 6.2.1).
\end{env}

\begin{proposition}[7.4.2]
Sous les conditions générales de (7.4.1), et pour un entier $p\in\mathbf Z$
donné, les propriétés suivantes sont équivalentes :
\begin{enumerate}
  \item[a)] $T_p$ est exact à gauche (ou, ce qui revient au même, $T_{p+1}$
    est exact à droite).
  \item[b)]
    \[
      Z'_p(P_\bullet)
      =
      \operatorname{Coker}(P_{p+1}\longrightarrow P_p)
    \]
    est un $A$-module à droite plat.
  \item[c)] Il existe un complexe $P'_\bullet$ de $A$-modules à droite plats
    tel que la différentielle
    \[
      d'_{p+1}:P'_{p+1}\longrightarrow P'_p
    \]
    soit nulle, et un isomorphisme de foncteurs homologiques de
    $H_\bullet(P_\bullet\otimes_A M)$ sur
    $H_\bullet(P'_\bullet\otimes_A M)$.
\end{enumerate}
\end{proposition}

\begin{proof}
Par définition, on a une suite exacte fonctorielle en $M$
\[
  0\longrightarrow T_p(M)
  \longrightarrow Z'_p(P_\bullet\otimes M)
  \longrightarrow P_{p-1}\otimes M
\]
où
\[
  Z'_p(P_\bullet\otimes M)
  =
  \operatorname{Coker}(P_{p+1}\otimes M\longrightarrow P_p\otimes M)
  =
  Z'_p(P_\bullet)\otimes M
\]
en vertu de l'exactitude à droite du produit tensoriel. Pour tout homomorphisme
$f:M\to N$, on a donc un diagramme commutatif
\[
\xymatrix@C=30pt@R=38pt{
  0\ar[r]
  &T_p(M)\ar[r]\ar[d]_u
  &Z'_p(P_\bullet)\otimes M\ar[r]\ar[d]_v
  &P_{p-1}\otimes M\ar[d]_w
  \\
  0\ar[r]
  &T_p(N)\ar[r]
  &Z'_p(P_\bullet)\otimes N\ar[r]
  &P_{p-1}\otimes N
}
\tag{7.4.2.1}\label{III.7.4.2.1.fr}
\]
dont les lignes sont exactes. Si $f$ est un monomorphisme, il en est de même
de $w$ puisque $P_{p-1}$ est plat; si $T_p$ est exact à gauche, $u$ est lui
aussi un monomorphisme; on en conclut que $v$ est un monomorphisme, ce qui
entraîne que $Z'_p(P_\bullet)$ est plat. Inversement, s'il en est ainsi, $v$
est un monomorphisme pour tout monomorphisme $f:M\to N$, donc le diagramme
(7.4.2.1) montre que $u$ est un monomorphisme, et par suite $T_p$ (qui est déjà
semi-exact) est exact à gauche. Ainsi $a)$ et $b)$ sont équivalentes. Il est
immédiat que $c)$ entraîne $a)$, car si
$d'_{p+1}:P'_{p+1}\to P'_p$ est nulle, et
$0\to M'\to M\to M''\to0$ est une suite exacte de $A$-modules, l'opérateur
bord dans la suite exacte
\[
  H_{p+1}(P'_\bullet\otimes M'')
  \xrightarrow{\partial}H_p(P'_\bullet\otimes M')
  \longrightarrow H_p(P'_\bullet\otimes M)
\]
\oldpage[III]{187}
est nul par définition (M, IV, 1), donc $T_p$ est exact à gauche. Montrons
inversement que $b)$ implique $c)$. Si
\[
  Z_{p+1}(P_\bullet)=\operatorname{Ker}(P_{p+1}\longrightarrow P_p),
\]
on a une suite exacte
\[
  0\longrightarrow Z_{p+1}(P_\bullet)\longrightarrow P_{p+1}
  \longrightarrow Z'_p(P_\bullet)\longrightarrow0
\]
dans laquelle $P_{p+1}$ et $Z'_p(P_\bullet)$ sont plats, donc
$Z_{p+1}(P_\bullet)$ est plat ($0_{\mathrm I}$, 6.1.2). On va prendre
\[
  P'_i=P_i\quad\text{pour }i\ne p\text{ et }i\ne p+1,
  \qquad
  P'_p=Z'_p(P_\bullet),
  \qquad
  P'_{p+1}=Z_{p+1}(P_\bullet);
\]
pour différentielle $d'_i:P'_i\to P'_{i-1}$, on prendra celle du complexe
$P_\bullet$ pour $i\ne p$ et $i\ne p+1$, $0$ pour $i=p+1$ et pour $i=p$
l'homomorphisme $Z'_p(P_\bullet)\to P_{p-1}$ déduit de $d_p$ par passage au
quotient. Comme les $P_i$ sont plats, on a
\[
  Z'_i(P_\bullet\otimes M)=Z'_i(P_\bullet)\otimes M,
  \qquad
  Z_i(P_\bullet\otimes M)=Z_i(P_\bullet)\otimes M
\]
et
\[
  B_i(P_\bullet\otimes M)=B_i(P_\bullet)\otimes M
\]
(en posant $B_i(P_\bullet)=\operatorname{Im}(P_{i+1}\to P_i)$); on en conclut
aussitôt pour tout $M$ des isomorphismes fonctoriels
\[
  H_i(P_\bullet\otimes M)\xrightarrow{\sim}H_i(P'_\bullet\otimes M)
\]
pour tout $i$, et la vérification du fait qu'il s'agit d'un isomorphisme de
$\partial$-foncteurs découle sans peine de la définition de $\partial$
(M, IV, 1).
\end{proof}

On remarquera que les conditions de (7.4.2) impliquent aussi que
$B_p(P_\bullet)$ est plat, car on a une suite exacte
\[
  0\longrightarrow B_p(P_\bullet)\longrightarrow P_p
  \longrightarrow Z'_p(P_\bullet)\longrightarrow0,
\]
dans laquelle $P_p$ et $Z'_p(P_\bullet)$ sont plats ($0_{\mathrm I}$, 6.1.2).

\begin{corollary}[7.4.3]
Supposons que $A$ soit un anneau noethérien régulier de dimension $1$
(autrement dit un produit d'anneaux de Dedekind
($0_{\mathrm{IV}}$, 17.1.3 et 17.3.7), par exemple un anneau principal). Alors,
pour que $T_p$ soit exact à gauche, il faut et il suffit que $T_p(A)$ soit un
$A$-module plat. Pour que $T_p$ soit exact, il faut et il suffit que $T_p(A)$
et $T_{p-1}(A)$ soient des $A$-modules plats.
\end{corollary}

\begin{proof}
Rappelons que pour un module $M$ sur un anneau de Dedekind, il revient au même
de dire que $M$ est plat ou qu'il est sans torsion ($0_{\mathrm I}$, 6.3.3 et
6.3.4); sous les hypothèses de (7.4.3), tout sous-module d'un $A$-module plat
est donc plat.

La seconde assertion de (7.4.3) résulte de la première, puisque dire que $T_p$
est exact signifie que $T_p$ et $T_{p-1}$ sont exacts à gauche. Pour démontrer
la première assertion, notons que l'on a une suite exacte
\[
  0\longrightarrow H_p(P_\bullet)\longrightarrow Z'_p(P_\bullet)
  \longrightarrow B_{p-1}(P_\bullet)\longrightarrow0
\]
dans laquelle $B_{p-1}(P_\bullet)$ est un $A$-module plat, en tant que
sous-module du $A$-module plat $P_{p-1}$. Il est donc équivalent de dire que
$H_p(P_\bullet)$ est plat ou que $Z'_p(P_\bullet)$ est plat
($0_{\mathrm I}$, 6.1.2).
\end{proof}

Les applications les plus importantes de (7.4.2) sont les suivantes :

\begin{proposition}[7.4.4]
Soient $A$ un anneau noethérien, $P_\bullet$ un complexe de $A$-modules plats;
on suppose, soit que les $P_i$ sont de type fini, soit que les
$H_i(P_\bullet)$ sont des $A$-modules de type fini et qu'il existe $i_0$ tel
que $H_i(P_\bullet)=0$ pour $i<i_0$. Soit $T$ le foncteur homologique défini
par (7.4.1.1). Alors l'ensemble $U$ des
$y\in\operatorname{Spec}(A)$ tels que $(T_p)_y$ (7.1.4) soit exact à droite
(resp. exact à gauche, exact) est ouvert dans $\operatorname{Spec}(A)$.
\end{proposition}

\begin{proof}
Dans la seconde hypothèse sur $P_\bullet$, on peut remplacer $P_\bullet$ par
un complexe $P'_\bullet$ de
\oldpage[III]{188}
$A$-modules libres de type fini pour lequel le foncteur
$H_\bullet(P'_\bullet\otimes_A M)$ est isomorphe (en tant que
$\partial$-foncteur) à $T_\bullet(M)$ (0, 11.9.3). On peut donc toujours se
ramener à la première hypothèse et dans ce cas les $Z'_i(P_\bullet)$ sont de
type fini; en outre, on peut se borner à démontrer les assertions relatives à
l'exactitude à gauche (cf. (7.4.2, $a)$). Cela étant, soit $x\in U$; comme le
foncteur $M\mapsto M_x$ est exact, on a
\[
  (Z'_p(P_\bullet))_x=Z'_p((P_\bullet)_x),
\]
et (en tenant compte de (7.4.1.3)) l'hypothèse entraîne, en vertu de
(7.4.2, $b)$) que $(Z'_p(P_\bullet))_x$ est un $A_x$-module plat, donc libre
puisqu'il est de type fini et que $A_x$ est un anneau local noethérien
(Bourbaki, \emph{Alg. comm.}, chap. II, \S 3, n\textsuperscript{o} 2, cor. 2
de la prop. 5). On en conclut qu'il y a un $f\in A$ tel que
$(Z'_p(P_\bullet))_f$ soit libre sur $A_f$ (Bourbaki, \emph{Alg. comm.},
chap. II, \S 5, n\textsuperscript{o} 1, cor. de la prop. 2), et \emph{a
fortiori} $(Z'_p(P_\bullet))_y$ est libre sur $A_y$ pour tout $y\in D(f)$, ce
qui achève la démonstration en vertu de (7.4.2, $b)$).
\end{proof}

\begin{corollary}[7.4.5]
Sous les hypothèses de (7.4.4), supposons en outre que $A$ soit intègre. Alors
l'ensemble $U$ des $x\in\operatorname{Spec}(A)$ tels que $(T_p)_x$ soit exact
est ouvert non vide.
\end{corollary}

\begin{proof}
Il suffit en effet de prouver que $(T_p)_x$ est exact pour le point générique
$x$ de $\operatorname{Spec}(A)$, ce qui est immédiat puisque $A_x$ est un
corps, donc tout foncteur additif sur $\mathbf{Ab}_{A_x}$ est exact.
\end{proof}

\begin{proposition}[7.4.6]
Sous les hypothèses générales de (7.4.4), les conditions $a)$, $b)$ et $c)$ de
(7.4.2) sont aussi équivalentes à :
\begin{enumerate}
  \item[d)] Il existe un $A$-module $Q$ et un isomorphisme fonctoriel
    \[
      T_p(M)\xrightarrow{\sim}\operatorname{Hom}_A(Q,M).
      \tag{7.4.6.1}\label{III.7.4.6.1.fr}
    \]
    En outre, le $A$-module $Q$ est déterminé à isomorphisme unique près par
    cette propriété, et il est de type fini.
\end{enumerate}
\end{proposition}

\begin{proof}
L'unicité de $Q$ est un cas particulier de l'unicité d'un objet représentatif
d'un foncteur représentable (0, 8.1.5). Il est clair que le second membre de
(7.4.6.1) est exact à gauche. Inversement, pour démontrer l'existence de $Q$,
lorsque $T_p$ est exact à gauche, on peut d'abord, comme dans (7.4.4), se
ramener au cas où les $P_i$ sont plats et de type fini, donc (puisque $A$ est
noethérien) projectifs de type fini (Bourbaki, \emph{Alg. comm.}, chap. II,
\S 5, n\textsuperscript{o} 2, cor. du th. 1). Le dual $\check P_i$ de $P_i$
est alors aussi un $A$-module projectif de type fini, $P_i$ est canoniquement
isomorphe au dual de $\check P_i$, et l'homomorphisme canonique
\[
  P_i\otimes_A M\longrightarrow\operatorname{Hom}_A(\check P_i,M)
\]
est bijectif (Bourbaki, \emph{Alg.}, chap. II, 3\textsuperscript{e} éd.,
\S 4, n\textsuperscript{o} 2, prop. 2). On sait d'autre part (7.4.2, $c)$)
qu'on peut supposer que $d_{p+1}:P_{p+1}\to P_p$ est nulle, donc que l'on a une
suite exacte
\[
  0\longrightarrow T_p(M)\xrightarrow{u}P_p\otimes M
  \xrightarrow{v}P_{p-1}\otimes M
\]
où $v=d_p\otimes1$. Posons alors $Q'=\operatorname{Ker}(d_p)$, de sorte qu'on
a la suite exacte
\[
  0\longrightarrow Q'\xrightarrow{w}P_p\xrightarrow{d_p}P_{p-1},
\]
d'où, par transposition, la suite exacte
\[
  \check P_{p-1}\xrightarrow{{}^t d_p}\check P_p
  \xrightarrow{{}^t w}\check Q'\longrightarrow0.
\]
Nous allons voir que
\[
  Q=\check Q'=\operatorname{Coker}({}^t d_p)
\]
répond à la question. En effet, on a la suite exacte
\[
  0\longrightarrow\operatorname{Hom}_A(Q,M)
  \longrightarrow\operatorname{Hom}_A(\check P_p,M)
  \xrightarrow{v'}\operatorname{Hom}_A(\check P_{p-1},M)
\]
\oldpage[III]{189}
où $v'=\operatorname{Hom}({}^t d_p,1)$; lorsqu'on identifie canoniquement
$P_i\otimes M$ à $\operatorname{Hom}(\check P_i,M)$, $v'$ s'identifie donc à
$v=d_p\otimes1$, et l'on a par suite l'isomorphisme fonctoriel
\[
  T_p(M)\xrightarrow{\sim}\operatorname{Hom}_A(Q,M)
\]
cherché. En outre, $Q$, étant un quotient de $\check P_p$, est de type fini.
\end{proof}

\begin{proposition}[7.4.7]
Supposons vérifiées les conditions générales de (7.4.4). Alors, pour tout
$A$-module $M$ de type fini :
\begin{enumerate}
  \item[(i)] Les $T_i(M)$ sont des $A$-modules de type fini.
  \item[(ii)] Pour tout idéal $\mathfrak m$ de $A$, l'homomorphisme canonique
    \[
      \widehat{T_i(M)}
      \longrightarrow
      \varprojlim_n T_i\bigl(M\otimes_A(A/\mathfrak m^{n+1})\bigr)
      \tag{7.4.7.1}\label{III.7.4.7.1.fr}
    \]
    (où le premier membre est le séparé complété de $T_i(M)$ pour la topologie
    $\mathfrak m$-préadique) est bijectif.
\end{enumerate}
\end{proposition}

\begin{proof}
Comme dans (7.4.4), on se ramène d'abord au cas où les $P_i$ sont de type fini;
$A$ étant noethérien, les sous-modules des $P_i\otimes_A M$ sont de type fini,
d'où trivialement l'assertion (i). Quant à l'assertion (ii), elle résulte plus
généralement du lemme suivant :
\end{proof}

\begin{lemma}[7.4.7.2]
Soient $A$ un anneau noethérien, $u:E\to F$ un homomorphisme de $A$-modules de
type fini. Pour tout $A$-module de type fini, posons
\[
  K(M)=\operatorname{Ker}(u\otimes1_M),
  \qquad
  C(M)=\operatorname{Coker}(u\otimes1_M);
\]
alors les homomorphismes canoniques
\[
  \widehat{K(M)}\longrightarrow\varprojlim_n K(M_n),
  \qquad
  \widehat{C(M)}\longrightarrow\varprojlim_n C(M_n)
  \tag{7.4.7.3}\label{III.7.4.7.3.fr}
\]
(où l'on a posé
$M_n=M\otimes_A(A/\mathfrak m^{n+1})=M/\mathfrak m^{n+1}M$) sont bijectifs
pour tout idéal $\mathfrak m$ de $A$.
\end{lemma}

\begin{proof}
Comme $E\otimes M$ et $F\otimes M$ sont de type fini, et que le foncteur
$M\mapsto\widehat M$ est exact dans la catégorie des $A$-modules de type fini
($0_{\mathrm I}$, 7.3.3), $\widehat{K(M)}$ et $\widehat{C(M)}$ sont
respectivement le noyau et le conoyau de
\[
  \widehat{(u\otimes1)}:
  \widehat{(E\otimes M)}\longrightarrow\widehat{(F\otimes M)}.
\]
L'exactitude à gauche du foncteur $\varprojlim$ montre donc que
$\widehat{K(M)}=\varprojlim_n K(M_n)$; d'autre part, l'exactitude à droite du
produit tensoriel prouve que
$C(M_n)=C(M)\otimes_A(A/\mathfrak m^{n+1})$, donc
$\widehat{C(M)}=\varprojlim_n C(M_n)$ par définition.
\end{proof}

\begin{remark}[7.4.8]
Compte tenu de (6.10.5) et (6.10.6), on voit que, moyennant une hypothèse de
platitude supplémentaire, (7.4.7) redonne le fait que (4.1.7.1) est un
isomorphisme, c'est-à-dire l'essentiel du « premier théorème de comparaison »
pour les morphismes propres; en outre, l'énoncé s'applique non plus seulement
à un $\mathscr O_X$-Module cohérent, mais à un complexe de tels Modules. Il
serait intéressant d'obtenir un énoncé comprenant à la fois (7.4.7) et
(4.1.7.1) comme cas particuliers. On notera que lorsque les $P_i$ sont de type
fini, la démonstration de (7.4.7) n'utilise pas le fait que ce sont des modules
plats; il y aurait lieu d'examiner si la conclusion de (7.4.7) est encore
valable lorsque les $P_i$ ne sont pas supposés plats ni de type fini, mais que
les $H_i(P_\bullet)$ sont supposés de type fini pour tout $i$ et nuls pour
$i<i_0$. Est-il alors possible de remplacer $P_\bullet$ par un complexe
$P'_\bullet$ de $A$-modules de type fini tel que les foncteurs
$H_\bullet(P_\bullet\otimes M)$ et $H_\bullet(P'_\bullet\otimes M)$ (qui ne
sont plus homologiques) soient encore isomorphes ?
\end{remark}

\oldpage[III]{190}
\subsection{Cas des anneaux locaux noethériens}
\label{subsection:III.7.5.fr}

\begin{env}[7.5.1]
Soient $A$ un anneau local noethérien, $\mathfrak m$ son idéal maximal, et pour
tout $A$-module $M$, désignons par $\widehat M$ son séparé complété pour la
topologie $\mathfrak m$-préadique, isomorphe à
\[
  \varprojlim_n\bigl(M\otimes_A(A/\mathfrak m^{n+1})\bigr)
  =
  \varprojlim_n\bigl(M/\mathfrak m^{n+1}M\bigr).
\]
Soit $T$ un foncteur covariant additif de $\mathbf{Ab}_A$ dans $\mathbf{Ab}$;
les homomorphismes canoniques (7.2.3.1)
\[
  T(M)\otimes_A(A/\mathfrak m^{n+1})
  \longrightarrow
  T\bigl(M\otimes_A(A/\mathfrak m^{n+1})\bigr)
\]
forment évidemment un système projectif de $A$-homomorphismes, qui donnent donc
à la limite un $\widehat A$-homomorphisme fonctoriel en $M$
\[
  \widehat{T(M)}
  \longrightarrow
  \varprojlim_n T(M_n)
  \tag{7.5.1.1}\label{III.7.5.1.1.fr}
\]
où l'on a posé
\[
  M_n=M\otimes_A(A/\mathfrak m^{n+1}),
  \qquad
  A_n=A/\mathfrak m^{n+1}.
\]
\end{env}

\begin{proposition}[7.5.2]
Soient $A$ un anneau local noethérien d'idéal maximal $\mathfrak m$,
$k=A/\mathfrak m$ son corps résiduel, $T$ un foncteur covariant additif de
$\mathbf{Ab}_A$ dans $\mathbf{Ab}$, semi-exact et commutant aux limites
inductives. On suppose en outre que pour tout $A$-module de type fini $M$,
$T(M)$ est un $A$-module de type fini et que l'homomorphisme canonique
(7.5.1.1) est un isomorphisme. Sous ces conditions, les propriétés suivantes
sont équivalentes :
\begin{enumerate}
  \item[a)] $T$ est exact à droite.
  \item[b)] Pour tout $n$, le foncteur $N\mapsto T(N)$ est exact à droite dans
    la catégorie des $A_n$-modules de type fini (ce qui revient à dire que $T$
    est exact à droite dans la catégorie des $A$-modules de longueur finie).
  \item[c)] L'homomorphisme canonique
    \[
      T(\varepsilon):T(A)\longrightarrow T(k)
    \]
    est surjectif.
  \item[d)] Pour tout $n$ assez grand, l'homomorphisme canonique
    \[
      T(A_n)\longrightarrow T(k)
    \]
    est surjectif.
\end{enumerate}
\end{proposition}

\begin{proof}
Il est clair que a) entraîne b). Montrons que b) entraîne a), c'est-à-dire que
si $u:M\to N$ est un épimorphisme de $A$-modules, $T(u)$ est un épimorphisme.
Comme $T$ commute aux limites inductives et que le foncteur $\varinjlim$ est
exact dans la catégorie des modules (pour des ensembles d'indices filtrants),
on peut se borner au cas où $M$ et $N$ sont de type fini. Comme $T(M)$ et
$T(N)$ sont alors de type fini, et que $A$ est un anneau local noethérien, il
suffit de montrer que
\[
  \widehat{T(u)}:
  \widehat{T(M)}
  \longrightarrow
  \widehat{T(N)}
\]
est surjectif ($0_{\mathrm I}$, 7.3.5 et $0_{\mathrm I}$, 6.4.1). Par hypothèse,
$\widehat{T(M)}$ et $\widehat{T(N)}$ sont respectivement $\varprojlim_n T(M_n)$
et $\varprojlim_n T(N_n)$, donc $\widehat{T(u)}$ est limite du système projectif
d'homomorphismes
\[
  T(u\otimes1_{A_n}):T(M_n)\longrightarrow T(N_n).
\]
Or, b) signifie que ces homomorphismes sont surjectifs; en outre, $T(M_n)$ est
un $A_n$-module de type fini, et $A_n$ est un anneau artinien par hypothèse; on
en conclut que $\widehat{T(u)}$ est surjectif ($0_{\mathrm I}$, 13.1.2 et
13.2.2). Il est clair que a) entraîne c), et comme $T(\varepsilon)$ se factorise
en
\[
  T(A)\longrightarrow T(A_n)\longrightarrow T(k),
\]
c) implique d); enfin, il résulte de (7.2.7) que b) et d) sont équivalentes
puisque $T$ est semi-exact dans $\mathbf{Ab}_{A_n}$, ce qui achève la
démonstration.
\end{proof}

\begin{corollary}[7.5.3]
Sous les hypothèses générales de (7.5.2), si $T(k)=0$, alors $T(M)=0$ pour tout
$A$-module $M$.
\end{corollary}

\begin{proof}
\oldpage[III]{191}
Comme $k$ est le seul $A$-module simple, on déduit de (7.3.5.4) que $T(E)=0$
pour tout $A$-module de longueur finie $E$. Si maintenant $M$ est de type fini,
$\widehat{T(M)}$ est isomorphe à $\varprojlim_n T(M_n)$, et comme les $M_n$ sont
de longueur finie, on a $\widehat{T(M)}=0$; comme $T(M)$ est de type fini par
hypothèse, il est isomorphe à un sous-module de $\widehat{T(M)}$
($0_{\mathrm I}$, 7.3.5), donc on a $T(M)=0$. Enfin, pour un $A$-module $M$
quelconque, $T(M)$ est limite inductive des $T(N_\alpha)$ pour les sous-modules
de type fini $N_\alpha$ de $M$, ce qui achève la démonstration.
\end{proof}

\begin{proposition}[7.5.4]
Soient $A$ un anneau local noethérien d'idéal maximal $\mathfrak m$,
$k=A/\mathfrak m$ son corps résiduel, $T_\bullet$ un foncteur homologique de
$\mathbf{Ab}_A$ dans $\mathbf{Ab}$, commutant aux limites inductives. On suppose
en outre que pour tout $i$ et tout $A$-module $M$ de type fini, $T_i(M)$ est de
type fini et l'homomorphisme canonique
\[
  \widehat{T_i(M)}
  \longrightarrow
  \varprojlim_n T_i(M_n)
\]
est bijectif. Pour un entier $p$ donné, les conditions suivantes sont alors
équivalentes :
\begin{enumerate}
  \item[a)] $T_p$ est exact.
  \item[b)] $T_p$ est exact à droite, et $T_p(A)$ est un $A$-module libre.
  \item[c)] Les homomorphismes canoniques
    \[
      T_{p+1}(A)\longrightarrow T_{p+1}(k)
      \quad\text{et}\quad
      T_p(A)\longrightarrow T_p(k)
    \]
    sont surjectifs.
  \item[d)] Pour tout $n$, les homomorphismes canoniques
    \[
      T_{p+1}(A_n)\longrightarrow T_{p+1}(k)
      \quad\text{et}\quad
      T_p(A_n)\longrightarrow T_p(k)
    \]
    sont surjectifs.
  \item[e)] Pour tout $n$, le foncteur $N\mapsto T_p(N)$ est exact dans la
    catégorie des $A_n$-modules de type fini.
\end{enumerate}
\end{proposition}

\begin{proof}
On sait (7.3.3) que a) équivaut à dire que $T_{p+1}$ et $T_p$ sont exacts à
droite; comme $T_\bullet$ est un foncteur homologique dans la catégorie
$\mathbf{Ab}_{A_n}$, le même raisonnement que dans (7.3.1) montre que e)
équivaut à dire que $T_p$ et $T_{p+1}$ sont exacts à droite dans la catégorie
des $A_n$-modules de type fini. On déduit donc de (7.5.2) que a) et e) sont
équivalentes; l'équivalence de a), c) et d) résulte aussi de (7.5.2). Enfin, on
sait que tout $A$-module plat de type fini est libre (Bourbaki,
\emph{Alg. comm.}, chap. II, \S 3, n\textsuperscript{o} 2, cor. 2 de la prop.
5); l'équivalence de a) et b) résulte alors de (7.3.1) et (7.3.3).
\end{proof}

\begin{corollary}[7.5.5]
Supposons vérifiées les conditions générales de (7.5.4).
\begin{enumerate}
  \item[(i)] Si $T_p(k)=0$, on a $T_p=0$, $T_{p+1}$ est exact à droite et
    $T_{p-1}$ est exact à gauche.
  \item[(ii)] Si $T_{p-1}(k)=T_{p+1}(k)=0$, $T_p$ est exact,
    l'homomorphisme canonique
    \[
      T_p(A)\otimes_A M\longrightarrow T_p(M)
    \]
    est bijectif et $T_p(A)$ est un $A$-module libre.
\end{enumerate}
\end{corollary}

\begin{proof}
(i) découle aussitôt de (7.5.3) puisque $T_p$ est semi-exact, la dernière
assertion résultant de la définition d'un foncteur homologique. On conclut
aussitôt de (i) les deux premières assertions de (ii), compte tenu de (7.3.3);
le fait que $T_p(A)$ soit libre résulte de (7.5.4).
\end{proof}

\begin{corollary}[7.5.6]
Supposons vérifiées les hypothèses générales de (7.5.4), et supposons de plus
que $A$ soit un anneau de valuation discrète.
\begin{enumerate}
  \item[(i)] Pour que $T_p$ soit exact à droite, il faut et il suffit que
    $T_{p-1}(A)$ soit un $A$-module libre.
  \item[(ii)] Pour que $T_p$ soit exact, il faut et il suffit que $T_p(A)$ et
    $T_{p-1}(A)$ soient des $A$-modules libres.
\end{enumerate}
\end{corollary}

\begin{proof}
\oldpage[III]{192}
Il est clair que (i) implique (ii) (cf. (7.3.3)). Pour prouver (i), remarquons
que si $f$ est un générateur de l'idéal maximal de $A$ (« uniformisante » de
$A$), pour qu'un $A$-module de type fini $M$ soit libre (ou plat, ce qui revient
au même), il faut et il suffit que l'homothétie $h_f:x\mapsto fx$ de $M$ soit
injective, car cela équivaut ici à dire que $M$ est sans torsion
($0_{\mathrm I}$, 6.3.4). Considérons alors la suite exacte
\[
  0\longrightarrow A\xrightarrow{h_f}A\longrightarrow k\longrightarrow0,
\]
qui fournit la suite exacte d'homologie
\[
  T_p(A)
  \longrightarrow T_p(k)
  \longrightarrow T_{p-1}(A)
  \xrightarrow{h_f}T_{p-1}(A).
\]
On voit que $T_{p-1}(A)$ est libre si et seulement si
$T_p(A)\to T_p(k)$ est surjectif; la conclusion résulte alors de (7.5.2).
\end{proof}

\begin{remark}[7.5.7]
On notera que, en vertu de (7.4.7), les hypothèses générales de (7.4.4)
entraînent que le foncteur homologique $T_\bullet$ défini par (7.4.1.1)
satisfait aux hypothèses générales de (7.5.4). Dans ce cas, (7.5.6) est donc
contenu dans (7.4.3).
\end{remark}

\subsection{Descente des propriétés d'exactitude. Théorème de semi-continuité
et critère d'exactitude de Grauert}
\label{subsection:III.7.6.fr}

\begin{proposition}[7.6.1]
Sous les conditions de (7.4.1), soit $B$ une $A$-algèbre commutative. Si $T_p$
est exact à droite (resp. exact à gauche, exact) il en est de même de
$T_p^{(B)}$; la réciproque est vraie lorsque $B$ est un $A$-module fidèlement
plat.
\end{proposition}

\begin{proof}
La première assertion est un cas particulier d'une assertion triviale de
(7.1.3). Inversement, supposons d'abord que $B$ soit un $A$-module plat. On a
alors, pour tout $A$-module $M$,
\[
  H_\bullet\bigl(P_\bullet\otimes_A(M\otimes_A B)\bigr)
  =
  \bigl(H_\bullet(P_\bullet\otimes_A M)\bigr)\otimes_A B,
\]
ce qui s'écrit aussi, pour tout $p$,
\[
  T_p(M)\otimes_A B
  =
  T_p^{(B)}(M_{(B)})
  \tag{7.6.1.1}\label{III.7.6.1.1.fr}
\]
à un isomorphisme canonique près. Supposons $T_p^{(B)}$ exact à droite (resp.
exact à gauche, exact); comme $M\mapsto M_{(B)}$ est un foncteur exact, le
premier membre de (7.6.1.1) est un foncteur exact à droite (resp. exact à
gauche, exact) en $M$; si maintenant $B$ est fidèlement plat sur $A$, on en
déduit que $T_p$ a la même propriété d'exactitude ($0_{\mathrm I}$, 6.4.1).
\end{proof}

\begin{proposition}[7.6.2]
Sous les conditions de (7.4.1), on suppose en outre que $A$ soit un anneau
noethérien réduit et que les $P_i$ soient des $A$-modules de type fini. Pour que
$T_p$ soit exact à droite (resp. exact à gauche, exact), il faut et il suffit
que, pour toute $A$-algèbre $B$ qui est un anneau de valuation discrète,
$T_p^{(B)}$ le soit.
\end{proposition}

\begin{proof}
En vertu de (7.3.1) et (7.3.3), on peut se borner à considérer l'exactitude à
droite, et il n'y a naturellement qu'à prouver la suffisance de la condition
(7.6.1). En vertu de (7.4.2), il suffit de montrer que
$Z'_{p-1}(P_\bullet)$ est un $A$-module plat; comme $P_{p-1}$ est de type fini,
$Z'_{p-1}(P_\bullet)$ est aussi de type fini; le critère ($0$, 10.2.8) montre
qu'il suffit alors que $Z'_{p-1}(P_\bullet)\otimes_A B$ soit un $B$-module plat
pour toute $A$-algèbre $B$ qui est un anneau de valuation discrète. Or, comme
$P_\bullet$ est un complexe de $A$-modules plats, on a
\[
  Z'_{p-1}(P_\bullet)\otimes_A B
  =
  Z'_{p-1}(P_\bullet\otimes_A B);
\]
\oldpage[III]{193}
$P_\bullet\otimes_A B$ est un complexe de $B$-modules plats
($0_{\mathrm I}$, 6.2.1), et pour tout $B$-module $N$, on a
\[
  H_\bullet(P_\bullet\otimes_A N)
  =
  H_\bullet\bigl((P_\bullet\otimes_A B)\otimes_B N\bigr),
\]
donc
\[
  T_p^{(B)}(N)
  =
  H_p\bigl((P_\bullet\otimes_A B)\otimes_B N\bigr);
\]
appliquant (7.4.2) à $T_p^{(B)}$, on voit que l'hypothèse que $T_p^{(B)}$ est
exact à droite est équivalente au fait que $Z'_{p-1}(P_\bullet\otimes_A B)$ est
un $B$-module plat.

Le critère précédent amène à étudier de plus près le cas des anneaux de
valuation discrète :
\end{proof}

\begin{proposition}[7.6.3]
Sous les conditions de (7.4.1), supposons que $A$ soit un anneau noethérien
régulier de dimension 1 (autrement dit, que $A$ est noethérien et que, pour tout
$x\in\operatorname{Spec}(A)$, $A_x$ est un corps ou un anneau de valuation
discrète). Alors, pour tout entier $p$ et tout $A$-module $M$, on a une suite
exacte canonique fonctorielle en $M$
\[
  0
  \longrightarrow T_p(A)\otimes_A M
  \xrightarrow{t_M}T_p(M)
  \longrightarrow \operatorname{Tor}_1^A\bigl(T_{p-1}(A),M\bigr)
  \longrightarrow0.
  \tag{7.6.3.1}\label{III.7.6.3.1.fr}
\]
\end{proposition}

\begin{proof}
Dans ce qui suit, nous supprimerons pour simplifier la mention du complexe
$P_\bullet$ dans les notations homologiques usuelles $H_p(P_\bullet)$,
$B_p(P_\bullet)$, $Z_p(P_\bullet)$ et $Z'_p(P_\bullet)$. On a les trois suites
exactes
\begin{gather*}
  0\longrightarrow H_p\longrightarrow Z'_p\longrightarrow B_{p-1}
  \longrightarrow0,\\
  0\longrightarrow B_{p-1}\longrightarrow Z_{p-1}\longrightarrow H_{p-1}
  \longrightarrow0,\\
  0\longrightarrow Z_{p-1}\longrightarrow P_{p-1}\longrightarrow B_{p-2}
  \longrightarrow0.
\end{gather*}
Comme $P_{p-1}$ et $P_{p-2}$ sont plats, il en est de même de leurs
sous-modules respectifs $B_{p-1}$, $Z_{p-1}$ et $B_{p-2}$ puisqu'il y a
identité entre $A_x$-modules plats et $A_x$-modules sans torsion (pour tout
$x\in\operatorname{Spec}(A)$); par tensorisation avec $M$, on a donc les suites
exactes
\[
  0=\operatorname{Tor}_1^A(B_{p-1},M)
  \longrightarrow H_p\otimes M
  \longrightarrow Z'_p\otimes M
  \xrightarrow{u}B_{p-1}\otimes M
  \longrightarrow0
  \tag{7.6.3.2}\label{III.7.6.3.2.fr}
\]
\[
  0=\operatorname{Tor}_1^A(Z_{p-1},M)
  \longrightarrow\operatorname{Tor}_1^A(H_{p-1},M)
  \longrightarrow B_{p-1}\otimes M
  \xrightarrow{v}Z_{p-1}\otimes M
  \tag{7.6.3.3}\label{III.7.6.3.3.fr}
\]
\[
  0=\operatorname{Tor}_1^A(B_{p-2},M)
  \longrightarrow Z_{p-1}\otimes M
  \xrightarrow{w}P_{p-1}\otimes M.
  \tag{7.6.3.4}\label{III.7.6.3.4.fr}
\]
Par définition,
$T_p(M)=\operatorname{Ker}(d_p\otimes1)/\operatorname{Im}(d_{p+1}\otimes1)$;
c'est donc le noyau de l'homomorphisme
\[
  (P_p\otimes M)/\operatorname{Im}(d_{p+1}\otimes1)
  \longrightarrow P_{p-1}\otimes M
\]
obtenu à partir de $d_p\otimes1$ par passage au quotient, homomorphisme qui
s'écrit aussi $Z'_p\otimes M\to P_{p-1}\otimes M$ par définition de
$Z'_p=P_p/B_p$; or, cet homomorphisme peut être considéré comme le composé
\[
  Z'_p\otimes M
  \xrightarrow{u}B_{p-1}\otimes M
  \xrightarrow{v}Z_{p-1}\otimes M
  \xrightarrow{w}P_{p-1}\otimes M.
\]
Comme $w$ est injectif d'après (7.6.3.4), on a une suite exacte
\[
  0\longrightarrow\operatorname{Ker}u
  \longrightarrow T_p(M)
  \longrightarrow\operatorname{Ker}v
  \longrightarrow0,
\]
qui n'est autre que (7.6.3.1), compte tenu de (7.6.3.2) et (7.6.3.3) et de ce
que $H_p=T_p(A)$ par définition.
\end{proof}

\begin{remarks}[7.6.4]
\begin{enumerate}
  \item[(i)] $H_\bullet(P_\bullet\otimes_A M)$ est l'homologie du bicomplexe
    $P_\bullet\otimes_A M$, où $M$ est considéré comme un complexe réduit à son
    terme de degré 0; elle est par suite (6.3.6 et 6.3.2) l'aboutissement de la
    suite spectrale régulière dont le terme $E_2$ est
    \[
      E^2_{pq}
      =
      \operatorname{Tor}_p^A\bigl(H_q(P_\bullet),M\bigr)
      =
      \operatorname{Tor}_p^A\bigl(T_q(A),M\bigr).
    \]
    \oldpage[III]{194}
    Or, l'hypothèse sur l'anneau $A$ entraîne que
    $\operatorname{Tor}_p^A(E,F)=0$ pour $p\geq2$ et pour des $A$-modules
    quelconques ($0_{\mathrm{IV}}$, 17.2.2); on sait (M, XV) que cela entraîne
    l'exactitude de la suite
    \[
      0\longrightarrow E^2_{0,q}
      \longrightarrow H_q(P_\bullet\otimes_A M)
      \longrightarrow E^2_{1,q-1}
      \longrightarrow0
    \]
    qui n'est autre que (7.6.3.1).
  \item[(ii)] Compte tenu de (7.3.1), la suite exacte (7.6.3.1) redonne comme
    cas particulier le résultat de (7.4.3).
\end{enumerate}
\end{remarks}

\begin{corollary}[7.6.5]
Sous les conditions de (7.4.1), supposons que $A$ soit un anneau de valuation
discrète, de corps des fractions $K$, de corps résiduel $k$, et que les $T_i(A)$
soient des $A$-modules de type fini. On a alors
\[
  \operatorname{rang}_k T_p(k)
  \geq
  \operatorname{rang}_k\bigl(T_p(A)\otimes_A k\bigr)
  \geq
  \operatorname{rang}_A T_p(A)
  =
  \operatorname{rang}_K T_p(K).
  \tag{7.6.5.1}\label{III.7.6.5.1.fr}
\]
En outre, pour que les termes extrêmes de cette inégalité soient égaux, il faut
et il suffit que $T_p$ soit exact, ou encore que $T_p(A)$ et $T_{p-1}(A)$ soient
des $A$-modules libres.
\end{corollary}

\begin{proof}
Faisons en effet $M=k$ dans la suite exacte (7.6.3.1); il vient, puisqu'il
s'agit d'espaces vectoriels sur $k$
\[
  \operatorname{rang}_k T_p(k)
  =
  \operatorname{rang}_k\bigl(T_p(A)\otimes_A k\bigr)
  +
  \operatorname{rang}_k\bigl(\operatorname{Tor}_1^A(T_{p-1}(A),k)\bigr).
\]
D'autre part, comme $T_p(A)$ est un module de type fini sur l'anneau local
intègre $A$, on a (Bourbaki, \emph{Alg. comm.}, chap. II, \S 3,
n\textsuperscript{o} 2, cor. 1 de la prop. 4)
\[
  \operatorname{rang}_k\bigl(T_p(A)\otimes_A k\bigr)
  \geq
  \operatorname{rang}_A T_p(A)
  =
  \operatorname{rang}_K\bigl(T_p(A)\otimes_A K\bigr)
  \tag{7.6.5.2}\label{III.7.6.5.2.fr}
\]
et en outre les deux membres de (7.6.5.2) sont égaux si et seulement si $T_p(A)$
est un $A$-module libre (\emph{loc. cit.}, prop. 7). On notera d'ailleurs que
puisque $K$ est un $A$-module plat, on a par définition
\[
  T_p(A)\otimes_A K
  =
  H_p(P_\bullet)\otimes_A K
  =
  H_p(P_\bullet\otimes_A K)
  =
  T_p(K).
\]
On a donc bien l'inégalité (7.6.5.1) et on voit en outre que l'égalité n'est
possible que si : 1\textsuperscript{o} $T_p(A)$ est libre;
2\textsuperscript{o} $\operatorname{Tor}_1^A(T_{p-1}(A),k)=0$, condition qui
équivaut, comme on sait ($0$, 10.1.3), au fait que $T_{p-1}(A)$ est un
$A$-module libre. Enfin, comme les $T_i(A)$ sont des $A$-modules de type fini,
il revient au même de dire qu'ils sont plats ou libres (Bourbaki,
\emph{Alg. comm.}, chap. II, \S 3, n\textsuperscript{o} 2, cor. 2 de la prop.
5), et on conclut par (7.4.3).
\end{proof}

\begin{env}[7.6.6]
Les hypothèses étant toujours celles de (7.4.1), nous poserons, pour tout
$x\in\operatorname{Spec}(A)$,
\[
  d_p(x)=d_p^T(x)=\operatorname{rang}_{\kappa(x)}T_p(\kappa(x)).
  \tag{7.6.6.1}\label{III.7.6.6.1.fr}
\]
\end{env}

\begin{lemma}[7.6.7]
Soit $\varphi:A\to A'$ un homomorphisme d'anneaux, et soit
\[
  f={}^a\!\varphi:
  \operatorname{Spec}(A')\longrightarrow\operatorname{Spec}(A)
\]
l'application correspondante (I, 1.2.1). Si l'on pose
$T'_\bullet=T_\bullet^{(A')}$ (7.1.3), on a
\[
  d_p^{T'}=d_p^T\circ f.
  \tag{7.6.7.1}\label{III.7.6.7.1.fr}
\]
\end{lemma}

\begin{proof}
\oldpage[III]{195}
En effet, pour tout $x'\in\operatorname{Spec}(A')$, on a, en posant
$x=f(x')$,
\[
  H_\bullet(P_\bullet\otimes_A\kappa(x'))
  =
  H_\bullet\bigl((P_\bullet\otimes_A\kappa(x))
    \otimes_{\kappa(x)}\kappa(x')\bigr)
  =
  H_\bullet(P_\bullet\otimes_A\kappa(x))
    \otimes_{\kappa(x)}\kappa(x'),
\]
puisque $\kappa(x')$ est plat sur $\kappa(x)$, d'où la relation (7.6.7.1).
\end{proof}

\begin{lemma}[7.6.8]
Si l'anneau $A$ est noethérien et le complexe $P_\bullet$ formé de $A$-modules
de type fini, la fonction $x\mapsto d_p^T(x)$ sur $\operatorname{Spec}(A)$ est
constructible.
\end{lemma}

\begin{proof}
Il faut prouver que pour toute partie fermée irréductible $Y$ de
$X=\operatorname{Spec}(A)$, il existe un ouvert non vide $U$ de $Y$ dans lequel
$d_p$ est constante ($0$, 9.2.2); comme
$Y=\operatorname{Spec}(A/\mathfrak a)$, où $\mathfrak a$ est un idéal de $A$ tel
que $A/\mathfrak a$ soit réduit, on peut, en vertu de (7.6.7), se borner au cas
où $Y=X$ et où $A$ est un anneau noethérien intègre; mais alors l'assertion
résulte de (7.4.5).
\end{proof}

\begin{theorem}[7.6.9]
Soient $A$ un anneau noethérien, $P_\bullet$ un complexe de $A$-modules plats de
type fini, $T_\bullet(M)=H_\bullet(P_\bullet\otimes_A M)$ le foncteur homologique
défini par $P_\bullet$; pour tout $x\in\operatorname{Spec}(A)$, soit
\[
  d_p(x)=\operatorname{rang}_{\kappa(x)}T_p(\kappa(x)).
\]
Alors :
\begin{enumerate}
  \item[(i)] La fonction $d_p$ est constructible et semi-continue
    supérieurement dans $\operatorname{Spec}(A)$.
  \item[(ii)] Si $T_p$ est exact, $d_p$ est continue (donc localement constante)
    dans $\operatorname{Spec}(A)$; la réciproque est vraie lorsque l'anneau $A$
    est réduit.
\end{enumerate}
\end{theorem}

\begin{proof}
\begin{enumerate}
  \item[(i)] La première assertion a été démontrée dans (7.6.8). Pour prouver
    la seconde, il suffit donc ($0$, 9.2.4) de montrer que si $x'\ne x$ est une
    générisation de $x$ dans $\operatorname{Spec}(A)$, on a
    $d_p(x')\leq d_p(x)$. Or, il existe alors un anneau de valuation discrète
    $B$ et un morphisme
    \[
      f:\operatorname{Spec}(B)\longrightarrow\operatorname{Spec}(A)
    \]
    tels que, si $a$ désigne le point fermé de $\operatorname{Spec}(B)$ et $b$
    son point générique, on ait $f(a)=x$ et $f(b)=y$ (II, 7.1.9). En vertu de
    la formule (7.6.7.1), on voit qu'on est ramené à démontrer l'inégalité
    $d_p(a)\geq d_p(b)$ dans $\operatorname{Spec}(B)$; mais cela n'est autre que
    l'inégalité (7.6.5.1)\footnote{Le principe de démonstration de (i) par
    réduction au cas d'un anneau de valuation discrète nous a été communiqué
    oralement par Hironaka.}.
  \item[(ii)] La première assertion a déjà été démontrée (7.3.4). Pour démontrer
    la réciproque, utilisons le critère valuatif (7.6.2); compte tenu de la
    formule (7.6.7.1), on est donc ramené au cas où $A$ est un anneau de
    valuation discrète; mais comme $\operatorname{Spec}(A)$ ne comporte alors
    que deux points, l'hypothèse que $d_p$ est constante implique bien que $T_p$
    est exact, en vertu de (7.6.5).
\end{enumerate}
\end{proof}

\begin{corollary}[7.6.10]
Soient $A$ un anneau noethérien, $\mathfrak p_i$ $(1\leq i\leq r)$ ses idéaux
premiers minimaux, $k_i$ le corps résiduel de $A_{\mathfrak p_i}$
$(1\leq i\leq r)$.
\begin{enumerate}
  \item[(i)] Pour tout $x\in\operatorname{Spec}(A)$, il existe un indice $i$
    tel que
    \[
      d_p(x)\geq\operatorname{rang}_{k_i}T_p(k_i).
      \tag{7.6.10.1}\label{III.7.6.10.1.fr}
    \]
\end{enumerate}
En particulier, si $A$ est intègre et si $K$ est son corps des fractions, on a
\[
  d_p(x)\geq\operatorname{rang}_K T_p(K)
  \tag{7.6.10.2}\label{III.7.6.10.2.fr}
\]
pour tout $x\in\operatorname{Spec}(A)$.
\oldpage[III]{196}
\begin{enumerate}
  \item[(ii)] Supposons en outre que $A$ soit local et réduit, et soit $k$ son
    corps résiduel. Alors, pour que $T_p$ soit exact, il faut et il suffit que
    l'on ait
    \[
      \operatorname{rang}_k T_p(k)
      =
      \operatorname{rang}_{k_i}T_p(k_i)
      \qquad\text{pour }1\leq i\leq r.
      \tag{7.6.10.3}\label{III.7.6.10.3.fr}
    \]
\end{enumerate}
\end{corollary}

\begin{proof}
(i) est immédiat puisque tout voisinage de $x$ contient un des $\mathfrak p_i$,
et il suffit d'appliquer la définition de la semi-continuité. D'autre part, si
$A$ est local, le seul voisinage dans $\operatorname{Spec}(A)$ de l'idéal
maximal $\mathfrak m$ est $\operatorname{Spec}(A)$ tout entier, donc on a
\[
  d_p(x)\leq\operatorname{rang}_k T_p(k)
\]
pour tout $x\in\operatorname{Spec}(A)$; cette relation, jointe à (i), montre que
la condition (7.6.10.3) entraîne que $d_p(x)$ est constante dans
$\operatorname{Spec}(A)$, et par suite que $T_p$ est exact en vertu de
(7.6.9, (ii)); la réciproque est évidente en vertu de (7.6.9, (ii)).
\end{proof}

\begin{remark}[7.6.11]
On peut se demander si l'assertion de (7.6.9, (i)) ne peut être renforcée par
l'inégalité
\[
  \operatorname{rang}_{\kappa(x)}T_p(\kappa(x))
  \geq
  \operatorname{rang}_{\kappa(x)}
    \bigl(T_p(A)\otimes_A\kappa(x)\bigr)
  \tag{7.6.11.1}\label{III.7.6.11.1.fr}
\]
pour tout $x\in\operatorname{Spec}(A)$, qui a lieu effectivement lorsque $A$
est un anneau de valuation discrète et $x$ son idéal maximal (7.6.5).
Bornons-nous au cas où $A$ est un anneau local noethérien, d'idéal maximal
$\mathfrak m$ et de corps résiduel $k$. Alors, les conditions suivantes sont
équivalentes :
\begin{enumerate}
  \item[a)] Pour tout complexe $P_\bullet$ de $A$-modules plats de type fini,
    on a
    \[
      \operatorname{rang}_k T_p(k)
      \geq
      \operatorname{rang}_k\bigl(T_p(A)\otimes_A k\bigr)
      \quad\text{pour tout entier }p.
      \tag{7.6.11.2}\label{III.7.6.11.2.fr}
    \]
  \item[b)] Pour tout $A$-module $M$ de type fini, on a
    \[
      \operatorname{rang}_k(M\otimes_A k)
      \geq
      \operatorname{rang}_k(\check M\otimes_A k).
      \tag{7.6.11.3}\label{III.7.6.11.3.fr}
    \]
  \item[c)] Pour tout $A$-module $N$ de type fini, on a
    \[
      \operatorname{rang}_k\bigl(\operatorname{Tor}_1^A(N,k)\bigr)
      \geq
      \operatorname{rang}_k\bigl(\operatorname{Tor}_2^A(N,k)\bigr).
      \tag{7.6.11.4}\label{III.7.6.11.4.fr}
    \]
\end{enumerate}
On notera qu'il revient au même, par décalage (M, V, 7.2), de dire que l'on a,
pour tout $i\geq1$,
\[
  \operatorname{rang}_k\bigl(\operatorname{Tor}_i^A(N,k)\bigr)
  \geq
  \operatorname{rang}_k\bigl(\operatorname{Tor}_{i+1}^A(N,k)\bigr).
  \tag{7.6.11.5}\label{III.7.6.11.5.fr}
\]
Donnons rapidement quelques indications sur la démonstration. Pour voir que a)
entraîne b), on considère une suite exacte
\[
  L_1\xrightarrow{d}L_0\longrightarrow M\longrightarrow0
\]
où $L_0$ et $L_1$ sont libres de type fini, et on applique a) au complexe
\[
  P_1\xrightarrow{{}^t d}P_0,
  \qquad
  P_0=\check L_1,
  \quad
  P_1=\check L_0,
\]
les autres termes étant nuls; on a alors $T_1(A)=\check M$ et
\[
  T_1(k)
  =
  \operatorname{Hom}_A(M,k)
  =
  \operatorname{Hom}_k(M/\mathfrak mM,k),
\]
autrement dit $T_1(k)$ est le dual de l'espace vectoriel $M\otimes_A k$, et a
donc même rang que ce dernier. Pour prouver que b) entraîne c), nous établissons
d'abord le lemme suivant :

\begin{lemma}[7.6.11.6]
Étant donné un complexe
\[
  \cdots\longrightarrow0\longrightarrow
  P_1\xrightarrow{d}P_0
  \longrightarrow0\longrightarrow\cdots
\]
de $A$-modules plats, on a une suite exacte
\[
  0\longrightarrow\operatorname{Tor}_2^A(Z'_0,k)
  \longrightarrow T_1(A)\otimes_A k
  \longrightarrow T_1(k)
  \longrightarrow\operatorname{Tor}_1^A(Z'_0,k)
  \longrightarrow0.
  \tag{7.6.11.7}\label{III.7.6.11.7.fr}
\]
\end{lemma}

\begin{proof}
\oldpage[III]{197}
En effet, partant de la suite exacte
$0\to Z_1\to P_1\to B_0\to0$, on en tire la suite exacte
\[
  0\longrightarrow\operatorname{Tor}_1^A(B_0,k)
  \longrightarrow Z_1\otimes k
  \longrightarrow P_1\otimes k
  \xrightarrow{u}B_0\otimes k
  \longrightarrow0.
\]
De la suite exacte $0\to B_0\to Z_0\to Z'_0\to0$, on tire, puisque $Z_0=P_0$
est plat,
$\operatorname{Tor}_1^A(B_0,k)=\operatorname{Tor}_2^A(Z'_0,k)$; par définition,
on a $Z_1=T_1(A)$; enfin, on a $T_1(k)=\operatorname{Ker}(d\otimes1)$, et
$d\otimes1$ se factorise en
\[
  P_1\otimes k
  \xrightarrow{u}B_0\otimes k
  \xrightarrow{v}Z_0\otimes k
  =P_0\otimes k;
\]
on a $T_1(k)=u^{-1}(R)$, où $R=\operatorname{Ker}v$, et comme $u$ est
surjectif, $R=u(T_1(k))$; enfin,
$R=\operatorname{Tor}_1^A(Z'_0,k)$ par définition de $v$, ce qui achève
d'établir la suite exacte (7.6.11.7).
\end{proof}

Pour déduire alors c) de b), on considère une suite exacte
\[
  L_1\xrightarrow{d}L_0\longrightarrow N\longrightarrow0,
\]
où $L_0$ et $L_1$ sont des modules libres de type fini; considérons le foncteur
$T$ associé au complexe formé de $L_1$ et $L_0$; comme $L_0$ et $L_1$ sont
libres, ils s'identifient à leurs biduals; donc si
\[
  M=\operatorname{Coker}({}^t d),
  \qquad
  T_1(A)=\operatorname{Ker}(d)=\check M;
\]
par ailleurs, $M\otimes_A k=\operatorname{Coker}({}^t d\otimes1_k)$ a même rang
sur $k$ que $\operatorname{Ker}(d\otimes1_k)$. L'hypothèse b) entraîne par suite
que
\[
  \operatorname{rang}_k\bigl(T_1(A)\otimes_A k\bigr)
  \leq
  \operatorname{rang}_k\bigl(T_1(k)\bigr);
\]
comme $Z'_0=N$, l'inégalité (7.6.11.4) résulte donc de la suite exacte
(7.6.11.7).

Enfin, pour prouver que c) entraîne a), appliquons (7.6.11.6) en remplaçant
$P_0$ et $P_1$ par $P_p$ et $P_{p-1}$; l'hypothèse c) appliquée au module
$Z'_p$ donne $\operatorname{rang}_kR\geq\operatorname{rang}_kS$, où
\[
  R=\operatorname{Ker}(P_p\otimes k\longrightarrow P_{p-1}\otimes k)
  \quad\text{et}\quad
  S=Z_p\otimes k.
\]
Or, si l'on factorise $d_{p+1}:P_{p+1}\to P_p$ en
\[
  P_{p+1}\xrightarrow{v}Z_p\xrightarrow{j}P_p,
\]
on a
\[
  \operatorname{Im}(d_{p+1}\otimes1)
  =
  (j\otimes1)\bigl(\operatorname{Im}(v\otimes1)\bigr).
\]
Comme
\[
  T_p(A)\otimes_A k
  =(Z_p/B_p)\otimes_A k
  =(Z_p\otimes k)/\operatorname{Im}(v\otimes1),
\]
et
\[
  T_p(k)=R/\operatorname{Im}(d_{p+1}\otimes1),
\]
on en conclut bien l'inégalité (7.6.11.2).

Cela étant, supposons que l'anneau local $A$ soit régulier de dimension $n$;
on sait alors [17] que le $A$-module
$\operatorname{Tor}_i^A(k,k)$ est isomorphe à la puissance extérieure
\[
  \bigwedge\nolimits^i(\mathfrak m/\mathfrak m^2);
\]
on voit donc que la condition (7.6.11.4) n'est pas vérifiée pour $N=k$, dès que
$n\geq4$.

Par contre, si l'anneau local intègre $A$ est tel que tout $A$-module réflexif
de type fini soit libre (ce qui est le cas lorsque $A$ est un anneau régulier
de dimension 2), la condition (7.6.11.3) est vérifiée : en effet, on sait que le
dual $\check M$ d'un $A$-module $M$ de type fini est réflexif, donc libre, et
par suite
\[
  \operatorname{rang}_k(\check M\otimes_A k)
  =
  \operatorname{rang}_K(\check M)
  =
  \operatorname{rang}_K(M),
\]
($K$ corps des fractions de $A$); par ailleurs, on sait que toute base sur $k$
de $M\otimes_A k$ est formée d'images d'un système de générateurs de $M$
(Bourbaki, \emph{Alg. comm.}, chap. II, \S 3, n\textsuperscript{o} 2, cor. 2 de
la prop. 4), donc
$\operatorname{rang}_K(M)\leq\operatorname{rang}_k(M\otimes_A k)$, ce qui
prouve notre assertion.
\end{remark}

\subsection{Application aux morphismes propres : I. La propriété d'échange}
\label{subsection:III.7.7.fr}

Les trois numéros qui suivent sont, pour l'essentiel, des traductions, dans le
langage des morphismes de préschémas, des résultats des numéros précédents.

\begin{env}[7.7.1]
Soit $f:X\to Y$ un morphisme quasi-compact et séparé de préschémas, et soit
$\mathscr P_\bullet$ un complexe de $\mathscr O_X$-Modules quasi-cohérents, dont
l'opérateur de dérivation est de degré $-1$; supposons en outre que les
$\mathscr O_X$-Modules $\mathscr P_i$ sont $Y$-plats ($0_{\mathrm I}$, 6.7.1).
\oldpage[III]{198}
Nous allons considérer le $\partial$-foncteur
$\mathscr M\mapsto\mathscr T_\bullet(\mathscr M)$ (aussi noté
$\mathscr T_\bullet(\mathscr P_\bullet,\mathscr M)$) dans la catégorie des
$\mathscr O_Y$-Modules quasi-cohérents, à valeurs dans la catégorie des
$\mathscr O_Y$-Modules quasi-cohérents (en vertu de (6.2.3)), défini par
\[
  \mathscr T_n(\mathscr P_\bullet,\mathscr M)
  =\mathscr T_n(\mathscr M)
  =\mathscr H^{-n}
    \bigl(f,\mathscr P^\bullet\otimes_{\mathscr O_Y}\mathscr M\bigr)
  \qquad\text{pour }n\in\mathbf Z,
  \tag{7.7.1.1}\label{III.7.7.1.1.fr}
\]
où $\mathscr P^\bullet$ est le complexe dont le terme de degré $j$ est
$\mathscr P_{-j}$, l'opérateur de dérivation étant donc alors de degré $+1$. Le foncteur
$\mathscr T_\bullet$ ainsi défini est un \emph{foncteur homologique} en
$\mathscr M$ (6.2.6).
\end{env}

\begin{env}[7.7.2]
Soit $g:Y'\to Y$ un morphisme, et posons
\[
  X'=X_{(Y')}=X\times_Y Y',
  \qquad
  f'=f_{(Y')}:X'\longrightarrow Y',
\]
qui est un morphisme quasi-compact et séparé; soit d'autre part
\[
  \mathscr P'_\bullet
  =\mathscr P_\bullet\otimes_{\mathscr O_Y}\mathscr O_{Y'};
\]
c'est un complexe de $\mathscr O_{X'}$-Modules quasi-cohérents qui sont
$Y'$-plats en vertu de (I, 9.1.12) et ($0_{\mathrm I}$, 6.2.1). Nous poserons
(avec les mêmes conventions sur les degrés)
\[
  \mathscr T^{Y'}_\bullet(\mathscr M')
  =\mathscr H^\bullet
    \bigl(f',\mathscr P'^\bullet\otimes_{\mathscr O_{Y'}}\mathscr M'\bigr)
  =\mathscr H^\bullet
    \bigl(f',\mathscr P^\bullet\otimes_{\mathscr O_Y}\mathscr M'\bigr)
  \tag{7.7.2.1}\label{III.7.7.2.1.fr}
\]
qui est un foncteur homologique en le $\mathscr O_{Y'}$-Module quasi-cohérent
$\mathscr M'$. Lorsque $Y'$ est un schéma affine d'anneau $A'$, on écrira
$\mathscr T^{A'}_\bullet$ au lieu de $\mathscr T^{Y'}_\bullet$; pour tout
$A'$-module $M'$, on a alors
\[
  \mathscr T^{A'}_\bullet(\widetilde{M'})
  =\bigl(\Gamma(Y',\mathscr T^{Y'}_\bullet(\widetilde{M'}))\bigr)^\sim;
\]
on posera
\[
  T^{A'}_\bullet(M')
  =\Gamma(Y',\mathscr T^{Y'}_\bullet(\widetilde{M'})),
\]
qui est un foncteur homologique de $A'$-modules, à valeurs dans la catégorie des
$A'$-modules.

On observera que si $Y=\operatorname{Spec}(A)$ est aussi affine, le foncteur de
$A'$-modules $T^{A'}_\bullet$ coïncide avec le foncteur obtenu par extension des
scalaires de $A$ à $A'$ à partir du foncteur homologique de $A$-modules
$T^A_\bullet$ (7.1.3) : en effet, soit $g:Y'\to Y$ le morphisme correspondant à
l'homomorphisme d'anneaux $A\to A'$, et soit $g':X'\to X$ le morphisme
correspondant qui est affine (II, 1.6.2); si $\mathfrak U$ est un recouvrement
ouvert affine de $X$, $\mathfrak U'=g'^{-1}(\mathfrak U)$ est un recouvrement
ouvert affine de $X'$; en vertu de (6.2.2), tout revient à voir que
\[
  C^\bullet
    \bigl(\mathfrak U,
      \mathscr P_\bullet\otimes_{\mathscr O_Y}g_*(\mathscr M')\bigr)
  =C^\bullet
    \bigl(\mathfrak U',
      \mathscr P_\bullet\otimes_{\mathscr O_Y}\mathscr M'\bigr),
\]
et finalement, que pour tout ouvert affine $U$ de $X$, en posant
$U'=g'^{-1}(U)$, on a
\[
  \Gamma
    \bigl(U,\mathscr P_\bullet\otimes_{\mathscr O_Y}g_*(\mathscr M')\bigr)
  =\Gamma
    \bigl(U',\mathscr P_\bullet\otimes_{\mathscr O_Y}\mathscr M'\bigr),
\]
ce qui est trivial (I, 1.3 et 3.2).

En particulier, si $U$ est un ouvert de $Y$, on a, pour tout
$\mathscr O_Y$-Module quasi-cohérent $\mathscr M$
\[
  \mathscr T^U_\bullet(\mathscr M|U)
  =\bigl(\mathscr T_\bullet(\mathscr M)\bigr)|U.
  \tag{7.7.2.2}\label{III.7.7.2.2.fr}
\]
\end{env}

\begin{env}[7.7.3]
Pour tout $\mathscr O_Y$-Module quasi-cohérent $\mathscr M$, on a un
homomorphisme canonique, fonctoriel en $\mathscr M$ :
\[
  \mathscr T_p(\mathscr O_Y)\otimes_{\mathscr O_Y}\mathscr M
  \longrightarrow
  \mathscr T_p(\mathscr M).
  \tag{7.7.3.1}\label{III.7.7.3.1.fr}
\]
En effet, si $Y$ est affine, cet homomorphisme a été défini en (7.2.2); cette
définition s'étend sans peine au cas général, en remarquant que si $U$, $V$ sont
deux ouverts affines de $Y$ tels que $V\subset U$, le diagramme
\oldpage[III]{199}
\[
\xymatrix@C=18pt@R=32pt{
  {\begin{gathered}
    \bigl(\mathscr T_p(\mathscr O_Y)
      \otimes_{\mathscr O_Y}\mathscr M\bigr)|U\\
    =\mathscr T^U_p(\mathscr O_Y|U)
      \otimes_{\mathscr O_Y|U}(\mathscr M|U)
  \end{gathered}}
  \ar[r]\ar[d]
  &
  {\begin{gathered}
    \mathscr T^U_p(\mathscr M|U)\\
    =\bigl(\mathscr T_p(\mathscr M)\bigr)|U
  \end{gathered}}
  \ar[d]
  \\
  {\begin{gathered}
    \bigl(\mathscr T_p(\mathscr O_Y)
      \otimes_{\mathscr O_Y}\mathscr M\bigr)|V\\
    =\mathscr T^V_p(\mathscr O_Y|V)
      \otimes_{\mathscr O_Y|V}(\mathscr M|V)
  \end{gathered}}
  \ar[r]
  &
  {\begin{gathered}
    \mathscr T^V_p(\mathscr M|V)\\
    =\bigl(\mathscr T_p(\mathscr M)\bigr)|V
  \end{gathered}}
}
\]
est commutatif par (7.2.3.3).

Pour tout morphisme $g:Y'\to Y$ on a un homomorphisme canonique
\[
  \mathscr T_p(\mathscr O_Y)\otimes_{\mathscr O_Y}\mathscr O_{Y'}
  \longrightarrow
  \mathscr T^{Y'}_p(\mathscr O_{Y'}),
  \tag{7.7.3.2}\label{III.7.7.3.2.fr}
\]
qui n'est autre que le cas particulier de (6.7.11.2) (pour les aboutissements)
dans le cas où $S=Y$, $v_1=f$, $v_2=1_Y$, $\mathscr P^{(2)}_\bullet$ réduit au
seul terme $\mathscr M$ de degré 0.

Lorsque $Y=\operatorname{Spec}(A)$, $Y'=\operatorname{Spec}(A')$ sont affines,
(7.7.3.2) n'est autre que l'homomorphisme de faisceaux correspondant à
l'homomorphisme canonique de $A'$-modules défini dans (7.2.2)
\[
  T^A_p(A)\otimes_A A'
  \longrightarrow
  T^{A'}_p(A')=T^A_p(A')
\]
comme il résulte aisément de (6.7.11) (car dans le cas envisagé, on peut prendre
$\mathfrak U'^{(i)}=u_i^{-1}(\mathfrak U^{(i)})$ dans (6.7.11)).
\end{env}

\begin{env}[7.7.4]
Lorsque $f$ est un morphisme \emph{propre}, $Y=\operatorname{Spec}(A)$ un schéma
affine noethérien et $\mathscr P_\bullet$ un complexe de
$\mathscr O_X$-Modules cohérents et $Y$-plats \emph{limité inférieurement}, on a
vu (6.10.5) que l'on peut écrire à un isomorphisme près,
\[
  \mathscr T_p(\mathscr M)
  =\mathscr H_p
    \bigl(\mathscr L_\bullet\otimes_{\mathscr O_Y}\mathscr M\bigr),
  \qquad
  \mathscr L_\bullet=\widetilde{L_\bullet},
\]
où $L_\bullet$ est un complexe de $A$-modules \emph{libres de type fini limité
inférieurement}; le foncteur $\mathscr T_\bullet$ est donc du type qui a été
étudié en détail dans (7.4) et (7.6). Nous allons traduire les résultats de cette
étude :
\end{env}

\begin{theorem}[7.7.5]
Soient $Y$ un préschéma localement noethérien, $(U_\alpha)$ un recouvrement de
$Y$ formé d'ouverts affines, $f:X\to Y$ un morphisme propre,
$\mathscr P_\bullet$ un complexe de $\mathscr O_X$-Modules cohérents et $Y$-plats
limité inférieurement. Le foncteur homologique $\mathscr T_\bullet(\mathscr M)$
défini par (7.7.1.1) possède alors les propriétés suivantes :

\medskip
\noindent
\textnormal{I) (La propriété de semi-continuité).}\footnote{Un cas particulier de
ce théorème se trouve déjà dans la note [3] de Chow-Igusa. La propriété de
semi-continuité a été découverte, dans le cadre des espaces analytiques (et sous
des hypothèses assez particulières), par Kodaira-Spencer (« On the variations of
almost-complex structures », \emph{Algebraic Geometry and Topology, A Symposium in
honor of S. Lefschetz}, Princeton Series n\textsuperscript{o} 12, p. 139--150,
Princeton, 1957) et la version générale démontrée par Grauert [5].}
\emph{La fonction}
\[
  y\longmapsto
  d_p(y)
  =
  \operatorname{rang}_{\kappa(y)}
    \mathscr T_p^{\kappa(y)}(\kappa(y))
  \tag{7.7.5.1}\label{III.7.7.5.1.fr}
\]
\emph{est semi-continue supérieurement.}

\oldpage[III]{200}
\medskip
\noindent
\textnormal{II) (La propriété d'échange).}
\emph{Pour un entier $p$ donné, les conditions suivantes sont équivalentes :}
\begin{enumerate}
  \item[a)] $\mathscr T_p$ est exact à droite.

  \item[a$'$)] $\mathscr T_p(\mathscr M)$ est isomorphe à un foncteur de la forme
    $\mathscr N\otimes_{\mathscr O_Y}\mathscr M$, $\mathscr N$ étant nécessairement
    isomorphe à
    \[
      \mathscr T_p(\mathscr O_Y)
      =
      \mathscr H^{-p}(f,\mathscr P^\bullet).
    \]

  \item[a$''$)] L'homomorphisme canonique fonctoriel (7.7.3.1) est un
    isomorphisme.

  \item[b)] $\mathscr T_{p-1}$ est exact à gauche.

  \item[b$'$)] Il existe un $\mathscr O_Y$-Module $\mathscr Q$ (nécessairement
    cohérent, et déterminé à un isomorphisme unique près) et un isomorphisme de
    foncteurs
    \[
      \mathscr T_{p-1}(\mathscr M)
      \xrightarrow{\sim}
      \mathscr{H}\!om_{\mathscr O_Y}(\mathscr Q,\mathscr M).
      \tag{7.7.5.2}\label{III.7.7.5.2.fr}
    \]

  \item[c)] En désignant par $A_\alpha$ l'anneau de l'ouvert affine $U_\alpha$,
    pour tout indice $\alpha$ le foncteur de $A_\alpha$-modules
    $T_p^{A_\alpha}$ est exact à droite.

  \item[d)] Pour tout morphisme $g:Y'\to Y$, l'homomorphisme canonique
    \[
      \mathscr T_p(\mathscr O_Y)
      \otimes_{\mathscr O_Y}\mathscr O_{Y'}
      \longrightarrow
      \mathscr T_p^{Y'}(\mathscr O_{Y'})
      \tag{7.7.5.3}\label{III.7.7.5.3.fr}
    \]
    est un isomorphisme.
\end{enumerate}
\end{theorem}

\begin{proof}
La propriété de semi-continuité est locale sur $Y$ et résulte donc de la
remarque (7.7.4) et de (7.6.9). Il est clair que a$''$) entraîne a$'$) et que
a$'$) entraîne a). L'équivalence de a), a$''$), b) et b$'$) a été démontrée
dans (7.3.1) et (7.4.6), compte tenu de la remarque (7.7.4), lorsque $Y$ est
affine. Pour passer au cas général, prouvons d'abord que a) est équivalent à c),
ce qui prouvera le caractère local sur $Y$ de la propriété a) ; la démonstration
s'appliquera également pour prouver le caractère local de a$''$) et b). Comme il
est clair que c) entraîne a), tout revient à prouver la réciproque. Il suffit
évidemment de montrer que pour tout ouvert affine $U$ de $Y$ et toute suite
exacte
\[
  0\longrightarrow\mathscr F'
  \longrightarrow\mathscr F
  \longrightarrow\mathscr F''
  \longrightarrow0
\]
de $(\mathscr O_Y|U)$-Modules quasi-cohérents, il existe une suite exacte
\[
  0\longrightarrow\mathscr G'
  \longrightarrow\mathscr G
  \longrightarrow\mathscr G''
  \longrightarrow0
\]
de $\mathscr O_Y$-Modules quasi-cohérents telle que
\[
  \mathscr F'=\mathscr G'|U,
  \qquad
  \mathscr F=\mathscr G|U,
  \qquad
  \mathscr F''=\mathscr G''|U;
\]
or, cela résulte aussitôt de l'hypothèse que $Y$ est localement noethérien, et
de (I, 9.4.2) : on prolonge en effet $\mathscr F$ en un $\mathscr O_Y$-Module
quasi-cohérent $\mathscr G$, $\mathscr F'$ en un sous-$\mathscr O_Y$-Module
$\mathscr G'$ de $\mathscr G$, et il suffit de prendre
$\mathscr G''=\mathscr G/\mathscr G'$.

Pour démontrer l'équivalence de b) et b$'$) dans le cas général, remarquons que
lorsque $Y$ est affine, on sait que $\mathscr Q$ est déterminé à un isomorphisme
unique près ; si alors $U$ est un ouvert affine du schéma affine $Y$, on en
déduit qu'il existe un isomorphisme fonctoriel
\[
  \mathscr T_{p-1}^{U}(\mathscr M|U)
  \xrightarrow{\sim}
  \mathscr{H}\!om_{\mathscr O_Y|U}
    (\mathscr Q|U,\mathscr M|U).
\]
Dans le cas général, pour tout ouvert affine $U$ de $Y$, il y a un
$(\mathscr O_Y|U)$-Module cohérent $\mathscr Q_U$ et un isomorphisme fonctoriel
\[
  \mathscr T_{p-1}^{U}(\mathscr M|U)
  \xrightarrow{\sim}
  \mathscr{H}\!om_{\mathscr O_Y|U}
    (\mathscr Q_U,\mathscr M|U);
\]
la remarque précédente montre que si $V$ est un ouvert affine contenu dans $U$,
on a $\mathscr Q_U|V=\mathscr Q_V$ ; d'où l'existence et l'unicité du
$\mathscr O_Y$-Module $\mathscr Q$ vérifiant (7.7.5.2).

Reste enfin à montrer l'équivalence de a) et d) ; il est clair que d) est de
caractère local sur $Y$, et l'on a vu ci-dessus qu'il en est de même de a) ; en
outre, d) est aussi local sur $Y'$. Or, lorsque $Y=\operatorname{Spec}(A)$,
$Y'=\operatorname{Spec}(A')$, on a vu que $T_\bullet^{A'}$ est le foncteur obtenu

\oldpage[III]{201}
à partir de $T_\bullet^A$ par extension des scalaires à $A'$, et il est clair
alors que a$'$) entraîne que (7.7.5.3) est un isomorphisme. Inversement, supposons
toujours $Y=\operatorname{Spec}(A)$ affine et soit $A'$ la $A$-algèbre
$A\oplus M$, où $M$ est un $A$-module quelconque, la multiplication dans $A'$
étant donnée par
\[
  (a_1,m_1)(a_2,m_2)
  =
  (a_1a_2,a_1m_2+a_2m_1);
\]
alors
\[
  T_p^{A'}(A')
  =
  T_p(A\oplus M)
  =
  T_p(A)\oplus T_p(M),
\]
et l'hypothèse que (7.7.5.3) soit bijectif entraîne qu'il en est de même de
l'application canonique
\[
  T_p(A)\otimes_A M\longrightarrow T_p(M),
\]
autrement dit d) entraîne a$''$), ce qui achève la démonstration.
\end{proof}

\begin{theorem}[7.7.6]
Soient $Y$ un préschéma localement noethérien, $f:X\to Y$ un morphisme propre,
$\mathscr F$ un $\mathscr O_X$-Module cohérent et $Y$-plat. Il existe alors un
$\mathscr O_Y$-Module cohérent $\mathscr Q$ (déterminé à isomorphisme unique près)
et un isomorphisme de foncteurs en le $\mathscr O_Y$-Module quasi-cohérent
$\mathscr M$ :
\[
  f_*(\mathscr F\otimes_{\mathscr O_Y}\mathscr M)
  \xrightarrow{\sim}
  \mathscr{H}\!om_{\mathscr O_Y}(\mathscr Q,\mathscr M)
  \tag{7.7.6.1}\label{III.7.7.6.1.fr}
\]
(d'où un isomorphisme de foncteurs
\[
  \Gamma(X,\mathscr F\otimes_{\mathscr O_Y}\mathscr M)
  \xrightarrow{\sim}
  \operatorname{Hom}_{\mathscr O_Y}(\mathscr Q,\mathscr M)).
  \tag{7.7.6.2}\label{III.7.7.6.2.fr}
\]
\end{theorem}

\begin{proof}
En effet, comme $\mathscr M\mapsto\mathscr F\otimes_{\mathscr O_Y}\mathscr M$ est
exact (0\textsubscript{I}, 6.7.4) et $f_*$ exact à gauche, le foncteur
$\mathscr M\mapsto f_*(\mathscr F\otimes_{\mathscr O_Y}\mathscr M)$ est exact à
gauche. Il suffit alors d'appliquer l'équivalence de (7.7.5, b)) et (7.7.5,
b$'$)) pour $p=1$.
\end{proof}

\begin{corollary}[7.7.7]
Soient $Y$ un préschéma localement noethérien, $f:X\to Y$ un morphisme propre,
$\mathscr F$, $\mathscr F'$ deux $\mathscr O_X$-Modules cohérents et $Y$-plats,
$u:\mathscr F\to\mathscr F'$ un homomorphisme. Considérons les deux foncteurs en
le $\mathscr O_Y$-Module quasi-cohérent $\mathscr M$ :
\begin{align*}
  \mathscr T(\mathscr M)
  &=
  \operatorname{Ker}
  \bigl(
    f_*(\mathscr F\otimes_{\mathscr O_Y}\mathscr M)
    \longrightarrow
    f_*(\mathscr F'\otimes_{\mathscr O_Y}\mathscr M)
  \bigr),\\
  T(\mathscr M)
  &=
  \Gamma(Y,\mathscr T(\mathscr M))\\
  &=
  \operatorname{Ker}
  \bigl(
    \Gamma(X,\mathscr F\otimes_{\mathscr O_Y}\mathscr M)
    \longrightarrow
    \Gamma(X,\mathscr F'\otimes_{\mathscr O_Y}\mathscr M)
  \bigr).
\end{align*}
Alors il existe un $\mathscr O_Y$-Module cohérent $\mathscr R$ (déterminé à
isomorphisme unique près) et des isomorphismes de foncteurs
\[
  \mathscr T(\mathscr M)
  \xrightarrow{\sim}
  \mathscr{H}\!om_{\mathscr O_Y}(\mathscr R,\mathscr M)
  \tag{7.7.7.1}\label{III.7.7.7.1.fr}
\]
\[
  T(\mathscr M)
  \xrightarrow{\sim}
  \operatorname{Hom}_{\mathscr O_Y}(\mathscr R,\mathscr M).
  \tag{7.7.7.2}\label{III.7.7.7.2.fr}
\]
\end{corollary}

\begin{proof}
On peut se borner à démontrer (7.7.7.2) ; cela prouvera en effet (7.7.7.1)
dans le cas où $Y$ est affine, et on passera de là au cas général en raisonnant
comme dans la démonstration de l'équivalence de (7.7.5, b) et b$'$)), grâce à
l'unicité à isomorphisme unique près d'un représentant d'un foncteur
représentable (0, 8.1.8). Il résulte de (7.7.6) qu'il existe deux
$\mathscr O_Y$-Modules cohérents $\mathscr Q$, $\mathscr Q'$ définissant des
isomorphismes fonctoriels
\[
  \Gamma(X,\mathscr F\otimes_{\mathscr O_Y}\mathscr M)
  \xrightarrow{\sim}
  \operatorname{Hom}_{\mathscr O_Y}(\mathscr Q,\mathscr M),
  \qquad
  \Gamma(X,\mathscr F'\otimes_{\mathscr O_Y}\mathscr M)
  \xrightarrow{\sim}
  \operatorname{Hom}_{\mathscr O_Y}(\mathscr Q',\mathscr M).
\]
Or, $u:\mathscr F\to\mathscr F'$ définit canoniquement un morphisme de foncteurs
\[
  \Gamma(X,\mathscr F\otimes_{\mathscr O_Y}\mathscr M)
  \longrightarrow
  \Gamma(X,\mathscr F'\otimes_{\mathscr O_Y}\mathscr M);
\]

\oldpage[III]{202}
il correspond à ce dernier un homomorphisme unique
$v:\mathscr Q'\to\mathscr Q$ de $\mathscr O_Y$-Modules tel que le diagramme
\[
\xymatrix@C=42pt@R=28pt{
  \Gamma(X,\mathscr F\otimes_{\mathscr O_Y}\mathscr M)
    \ar[r]\ar[d]_-{\sim}
  &
  \Gamma(X,\mathscr F'\otimes_{\mathscr O_Y}\mathscr M)
    \ar[d]^-{\sim}
  \\
  \operatorname{Hom}_{\mathscr O_Y}(\mathscr Q,\mathscr M)
    \ar[r]
  &
  \operatorname{Hom}_{\mathscr O_Y}(\mathscr Q',\mathscr M)
}
\]
soit commutatif (0, 8.1.4). Comme le foncteur contravariant
$\mathscr N\mapsto\operatorname{Hom}_{\mathscr O_Y}(\mathscr N,\mathscr M)$ est
exact à gauche dans la catégorie des $\mathscr O_Y$-Modules, il suffit de prendre
$\mathscr R=\operatorname{Coker}(v)$ pour obtenir l'isomorphisme (7.7.7.2)
cherché.
\end{proof}

\begin{corollary}[7.7.8]
Sous les hypothèses de (7.7.6) relatives à $X$, $Y$ et $f$, soient $\mathscr F$,
$\mathscr G$ deux $\mathscr O_X$-Modules cohérents vérifiant les conditions
suivantes : (i) $\mathscr F$ est $Y$-plat ; (ii) $\mathscr G$ est isomorphe au
conoyau d'un homomorphisme de $\mathscr O_X$-Modules localement libres de type
fini $\mathscr E_1\to\mathscr E_0$. Considérons les deux foncteurs en le
$\mathscr O_Y$-Module quasi-cohérent $\mathscr M$ :
\begin{align*}
  \mathscr T(\mathscr M)
  &=
  f_*
  \bigl(
    \mathscr{H}\!om_{\mathscr O_X}
      (\mathscr G,\mathscr F\otimes_{\mathscr O_Y}\mathscr M)
  \bigr),\\
  T(\mathscr M)
  &=
  \Gamma(Y,\mathscr T(\mathscr M))
  =
  \operatorname{Hom}_{\mathscr O_X}
    (\mathscr G,\mathscr F\otimes_{\mathscr O_Y}\mathscr M).
\end{align*}
Alors il existe un $\mathscr O_Y$-Module cohérent $\mathscr N$ (déterminé à
isomorphisme unique près) et des isomorphismes de foncteurs
\[
  \mathscr T(\mathscr M)
  \xrightarrow{\sim}
  \mathscr{H}\!om_{\mathscr O_Y}(\mathscr N,\mathscr M)
  \tag{7.7.8.1}\label{III.7.7.8.1.fr}
\]
\[
  T(\mathscr M)
  \xrightarrow{\sim}
  \operatorname{Hom}_{\mathscr O_Y}(\mathscr N,\mathscr M).
  \tag{7.7.8.2}\label{III.7.7.8.2.fr}
\]
\end{corollary}

\begin{proof}
En vertu de l'isomorphisme fonctoriel ($0_{\mathrm I}$, 5.4.2.1), on a des
isomorphismes fonctoriels en $\mathscr M$
\begin{align*}
  \mathscr{H}\!om_{\mathscr O_X}
    (\mathscr E_i,\mathscr F\otimes_{\mathscr O_Y}\mathscr M)
  &\xrightarrow{\sim}
  \check{\mathscr E}_i
    \otimes_{\mathscr O_X}
    (\mathscr F\otimes_{\mathscr O_Y}\mathscr M)\\
  &\xrightarrow{\sim}
  (\check{\mathscr E}_i\otimes_{\mathscr O_X}\mathscr F)
    \otimes_{\mathscr O_Y}\mathscr M\\
  &\xrightarrow{\sim}
  \mathscr{H}\!om_{\mathscr O_X}(\mathscr E_i,\mathscr F)
    \otimes_{\mathscr O_Y}\mathscr M
\end{align*}
pour $i=0,1$. Posons
$\mathscr F_i=\mathscr{H}\!om_{\mathscr O_X}(\mathscr E_i,\mathscr F)$ pour
$i=0,1$ ; ce sont des $\mathscr O_X$-Modules cohérents ($0_{\mathrm I}$, 5.3.5)
et $Y$-plats ($0_{\mathrm I}$, 5.4.2) ; soit
$u=\mathscr{H}\!om(v,1_{\mathscr F}):\mathscr F_0\to\mathscr F_1$. En vertu de
l'exactitude à gauche du foncteur
$\mathscr H\mapsto\mathscr{H}\!om_{\mathscr O_X}
(\mathscr H,\mathscr F\otimes_{\mathscr O_Y}\mathscr M)$, on a des isomorphismes
fonctoriels en $\mathscr M$
\begin{align*}
  \mathscr{H}\!om_{\mathscr O_X}
    (\mathscr G,\mathscr F\otimes_{\mathscr O_Y}\mathscr M)
  &\xrightarrow{\sim}
  \operatorname{Ker}
  \bigl(
    \mathscr{H}\!om_{\mathscr O_X}
      (\mathscr E_0,\mathscr F\otimes_{\mathscr O_Y}\mathscr M)
    \longrightarrow
    \mathscr{H}\!om_{\mathscr O_X}
      (\mathscr E_1,\mathscr F\otimes_{\mathscr O_Y}\mathscr M)
  \bigr)\\
  &\xrightarrow{\sim}
  \operatorname{Ker}
  \bigl(
    \mathscr F_0\otimes_{\mathscr O_Y}\mathscr M
    \longrightarrow
    \mathscr F_1\otimes_{\mathscr O_Y}\mathscr M
  \bigr).
\end{align*}
Puisque $f_*$ est exact à gauche, on en déduit un isomorphisme fonctoriel
\[
  f_*
  \bigl(
    \mathscr{H}\!om_{\mathscr O_X}
      (\mathscr G,\mathscr F\otimes_{\mathscr O_Y}\mathscr M)
  \bigr)
  \xrightarrow{\sim}
  \operatorname{Ker}
  \bigl(
    f_*(\mathscr F_0\otimes_{\mathscr O_Y}\mathscr M)
    \longrightarrow
    f_*(\mathscr F_1\otimes_{\mathscr O_Y}\mathscr M)
  \bigr)
\]
et il suffit alors d'appliquer (7.7.7).
\end{proof}

\oldpage[III]{203}
\begin{remarks}[7.7.9]
\textnormal{(i)} Dans (7.7.6), (7.7.7), (7.7.8), la formation des
$\mathscr O_Y$-Modules $\mathscr Q$, $\mathscr R$, $\mathscr N$ \emph{permute aux
changements de base}. Par exemple (gardant les notations de (7.7.2)), dans le
cas (7.7.6), on a, pour tout $\mathscr O_{Y'}$-Module quasi-cohérent
$\mathscr M'$, l'isomorphisme
\[
  f'_*(\mathscr F\otimes_{\mathscr O_Y}\mathscr M')
  \xrightarrow{\sim}
  \mathscr{H}\!om_{\mathscr O_{Y'}}(g^*(\mathscr Q),\mathscr M')
\]
car en vertu de la remarque faite dans (7.7.2), tout revient à voir que l'on a
\[
  \operatorname{Hom}_{\mathscr O_Y}(\mathscr Q,g_*(\mathscr M'))
  =
  \operatorname{Hom}_{\mathscr O_{Y'}}(g^*(\mathscr Q),\mathscr M')
\]
ce qui n'est autre que ($0_{\mathrm I}$, 4.4.3.1). De même, quand dans (7.7.7)
on remplace $Y$, $f$, $\mathscr M$, $\mathscr F$, $\mathscr F'$ par $Y'$, $f'$,
$\mathscr M'$, $\mathscr F\otimes_{\mathscr O_X}\mathscr O_{X'}$,
$\mathscr F'\otimes_{\mathscr O_X}\mathscr O_{X'}$, il faut remplacer
$\mathscr R$ par $g^*(\mathscr R)$. Enfin, dans (7.7.8), lorsqu'on remplace $X$,
$Y$, $f$, $\mathscr M$, $\mathscr F$ par $X'$, $Y'$, $f'$, $\mathscr M'$,
$\mathscr F\otimes_{\mathscr O_X}\mathscr O_{X'}$, et $\mathscr G$ par
$\mathscr G'=\mathscr G\otimes_{\mathscr O_Y}\mathscr O_{X'}$, il faut remplacer
$\mathscr N$ par $g^*(\mathscr N)$ : cela résulte de ce que l'on a encore une
suite exacte
\[
  \mathscr E'_1\longrightarrow
  \mathscr E'_0\longrightarrow
  \mathscr G'\longrightarrow0
  \qquad\text{avec}\qquad
  \mathscr E'_i=\mathscr E_i\otimes_{\mathscr O_Y}\mathscr O_{X'}
\]
et de ce que
\[
  \check{\mathscr E}'_i
  \otimes_{\mathscr O_{X'}}
  (\mathscr F\otimes_{\mathscr O_Y}\mathscr M')
  =
  \check{\mathscr E}_i
  \otimes_{\mathscr O_X}
  (\mathscr F\otimes_{\mathscr O_Y}\mathscr M')
  \qquad(i=0,1).
\]

\medskip
\textnormal{(ii)} La condition (ii) de l'énoncé de (7.7.8) relative à
$\mathscr G$ est toujours satisfaite pour un $\mathscr O_X$-Module cohérent
$\mathscr G$ quelconque lorsqu'il existe un $\mathscr O_X$-Module inversible
$Y$-ample, par exemple lorsque $Y$ est affine et $f:X\to Y$ est un morphisme
projectif. Il suffit alors de noter (compte tenu de (II, 5.5.1)) qu'il existe un
$\mathscr O_X$-Module localement libre de type fini $\mathscr E_0$ tel que
$\mathscr G$ soit isomorphe à un quotient de $\mathscr E_0$ (II, 2.7.10) ; comme
$\mathscr E_0$ et $\mathscr G$ sont cohérents, il en est de même du noyau
$\mathscr G_1$ de $\mathscr E_0\to\mathscr G$, et en appliquant le même résultat
à $\mathscr G_1$, on obtient bien une suite exacte
$\mathscr E_1\to\mathscr E_0\to\mathscr G\to0$ où $\mathscr E_0$ et
$\mathscr E_1$ sont localement libres de type fini.

\medskip
\textnormal{(iii)} Nous prouverons au chap. V que, dans (7.7.8), l'hypothèse
restrictive (ii) est \emph{surabondante}.
\end{remarks}

\begin{proposition}[7.7.10]
\textnormal{(Critères locaux pour la propriété d'échange).} Sous les conditions
générales de (7.7.5), soient $y$ un point de $Y$, $p$ un entier. Les propriétés
suivantes sont équivalentes :
\begin{enumerate}
  \item[a)] Le foncteur $T_p^{\mathscr O_y}$ est exact à droite.

  \item[b)] L'homomorphisme canonique
    \[
      T_p^{\mathscr O_y}(\mathscr O_y)
      \longrightarrow
      T_p^{\mathscr O_y}(\kappa(y))
    \]
    est surjectif.

  \item[c)] Pour tout entier $n$, l'homomorphisme canonique
    \[
      T_p^{\mathscr O_y}
      (\mathscr O_y/\mathfrak m_y^{n+1})
      \longrightarrow
      T_p^{\mathscr O_y}(\kappa(y))
    \]
    est surjectif.
\end{enumerate}
De plus, l'ensemble des $y\in Y$ vérifiant ces conditions est le plus grand
ensemble ouvert $U$ de $Y$ tel que $\mathscr T_p^U$ soit exact à droite.
\end{proposition}

\begin{proof}
Compte tenu de (7.7.4), l'équivalence de a), b) et c) résulte de (7.4.7) et
(7.5.2). Le fait que l'ensemble $U$ où $T_p^{\mathscr O_y}$ est exact à droite
soit ouvert est aussi conséquence de (7.4.4), et inversement si
$\mathscr T_p^V$ est exact à droite, il en est de même de
$T_p^{\mathscr O_y}$ pour tout $y\in V$, par la condition c) de (7.7.5) et
(7.6.1).
\end{proof}

\begin{corollary}[7.7.11]
Si $\mathscr T_p$ est exact à droite (resp. à gauche), alors, pour tout morphisme
$g:Y'\to Y$, $\mathscr T_p^{Y'}$ est exact à droite (resp. à gauche). La
réciproque est vraie lorsque le morphisme $g$ est fidèlement plat.
\end{corollary}

\oldpage[III]{204}
\begin{proof}
La première assertion est conséquence immédiate de (7.6.1) et du fait que la
question est locale sur $Y$ et $Y'$, d'après (7.7.5, c) et b)). Pour démontrer
la seconde assertion, il suffit de voir que pour tout $y\in Y$,
$T_p^{\mathscr O_y}$ est exact à droite (resp. à gauche), en vertu de (7.7.10).
Mais par hypothèse, il existe $y'\in Y'$ tel que $g(y')=y$, et
$\mathscr O_{y'}$ est un $\mathscr O_y$-module fidèlement plat ; la conclusion
résulte alors de l'hypothèse et de (7.6.1).
\end{proof}

\begin{remarks}[7.7.12]
\textnormal{(i)} Sous les hypothèses de (7.7.4), supposons en outre que
$\mathscr P_\bullet$ soit un complexe \emph{fini} ; alors il résulte de (7.7.1)
(puisqu'on peut prendre pour $\mathfrak U$ un recouvrement ouvert affine
\emph{fini} de $X$) que le bicomplexe
$C^\bullet(\mathfrak U,\mathscr P^\bullet\otimes_{\mathscr O_Y}\mathscr M)$ est
aussi fini, et de façon précise qu'il existe un ensemble fini $E$ de couples,
\emph{indépendant de $\mathscr M$}, tel que
\[
  C^h(\mathfrak U,\mathscr P^k\otimes_{\mathscr O_Y}\mathscr M)=0
  \qquad\text{pour tous les couples }(h,k)\notin E.
\]
On en conclut qu'il existe $i_1$ tel que, pour $i\geq i_1$, on ait
$\mathscr T_i(\mathscr M)=0$ pour \emph{tout} $\mathscr O_Y$-Module
quasi-cohérent $\mathscr M$. En particulier, $\mathscr T_i$ est trivialement un
foncteur exact en $\mathscr M$ pour ces valeurs de $i$, et par suite (7.4.1),
$Z'_i(L_\bullet)$ est un $A$-module \emph{plat} de type fini (donc
\emph{projectif} de type fini, puisque $A$ est noethérien) pour ces valeurs de
$i$. Considérons alors le complexe $(L'_\bullet)$ de $A$-modules tel que
$L'_i=L_i$ pour $i<i_1$, $L'_{i_1}=Z'_{i_1}(L_\bullet)$ et $L'_i=0$ pour
$i>i_1$ et soit $\mathscr L'_\bullet=\widetilde{L'_\bullet}$. Il est clair que
\[
  \mathscr H_i(\mathscr L'_\bullet\otimes_{\mathscr O_Y}\mathscr M)
  =
  \mathscr H_i(\mathscr L_\bullet\otimes_{\mathscr O_Y}\mathscr M)
\]
pour $i<i_1-1$ et aussi pour $i\geq i_1$ (les deux membres étant alors nuls) ;
enfin, comme
$\operatorname{Im}(Z'_{i_1}\otimes_A M)=\operatorname{Im}(L_{i_1}\otimes_A M)$
par définition, on a aussi
$\mathscr H_i(\mathscr L'_\bullet\otimes_{\mathscr O_Y}\mathscr M)
=\mathscr H_i(\mathscr L_\bullet\otimes_{\mathscr O_Y}\mathscr M)$ pour
$i=i_1-1$. On voit donc qu'on peut supposer dans (7.7.4) que
$\mathscr L_\bullet$ est aussi un complexe \emph{fini}, à condition d'exiger
seulement que les $\mathscr L_i$ soient des $\mathscr O_Y$-Modules
\emph{localement libres} (associés à des $A$-modules projectifs de type fini).

Ce raisonnement s'applique en particulier au cas où $\mathscr P_\bullet$ est
réduit à \emph{un seul} terme $\mathscr F\ne0$, de degré 0 (auquel cas
$\mathscr T_n(\mathscr M)=\mathrm R^{-n}f_*
(\mathscr F\otimes_{\mathscr O_Y}\mathscr M)$) ; on peut alors supposer que les
$\mathscr L_i$ sont nuls pour $i>0$ ; on utilisera de préférence dans ce cas les
notations cohomologiques, écrivant donc $\mathscr T^{-p}$ au lieu de
$\mathscr T_p$.

\medskip
\textnormal{(ii)} Lorsque dans l'énoncé de (7.7.5), on ne suppose plus que les
$\mathscr P_i$ sont $Y$-plats, les conclusions restent valables à condition de
poser cette fois
\[
  \mathscr T_p(\mathscr M)
  =
  \operatorname{Tor}^Y_p
    (f,1_Y;\mathscr P_\bullet,\mathscr M).
  \tag{7.7.12.1}\label{III.7.7.12.1.fr}
\]
En effet, $\operatorname{Tor}^Y_n(f,1_Y;\mathscr P_\bullet,\mathscr O_Y)$ est
alors un $\mathscr O_Y$-Module \emph{cohérent} en vertu de (6.7.9). La
démonstration de (6.10.5) s'applique sans changement, compte tenu de (6.10.1)
et montre encore que lorsque $Y=\operatorname{Spec}(A)$ est affine, l'on a
\[
  \mathscr T_p(\mathscr M)
  =
  \mathscr H_p(\mathscr L_\bullet\otimes_{\mathscr O_Y}\mathscr M),
  \qquad
  \mathscr L_\bullet=\widetilde{L_\bullet},
\]
où $L_\bullet$ est un complexe de $A$-modules libres de type fini ; cela prouve
notre assertion.
\end{remarks}

\subsection{Application aux morphismes propres : II. Critères de platitude
cohomologique}
\label{subsection:III.7.8.fr}

\begin{definition}[7.8.1]
Soient $X$, $Y$ deux préschémas, $f:X\to Y$ un morphisme quasi-compact et
séparé, $\mathscr P_\bullet$ un complexe de $\mathscr O_X$-Modules
quasi-cohérents et $Y$-plats, $\mathscr T_\bullet$ le foncteur homologique de
$\mathscr O_Y$-Modules quasi-cohérents défini par (7.7.1.1), $y$ un point de $Y$.
On dit que $\mathscr P_\bullet$ est

\oldpage[III]{205}
\emph{homologiquement plat sur $Y$ au point $y$, en dimension $p$} (ou
\emph{cohomologiquement plat sur $Y$ au point $y$, en dimension $-p$}) s'il
existe un voisinage ouvert $U$ de $y$ dans $Y$ tel que
$\mathscr T_p^U=\mathscr T_U^{-p}$ soit exact. On dit que $\mathscr P_\bullet$ est
\emph{homologiquement plat en dimension $p$ sur $Y$} (ou
\emph{cohomologiquement plat en dimension $-p$ sur $Y$}) s'il est
homologiquement plat sur $Y$ en tout point $y\in Y$, en dimension $p$.

Lorsque $\mathscr P_\bullet$ est homologiquement plat sur $Y$ (resp. sur $Y$ au
point $y$) pour \emph{toute dimension $p$}, on dit simplement que
$\mathscr P_\bullet$ est \emph{homologiquement plat sur $Y$} (resp. sur $Y$ au
point $y$) ou \emph{cohomologiquement plat sur $Y$} (resp. sur $Y$ au point
$y$).
\end{definition}

\begin{env}[7.8.2]
Par définition, la notion de platitude homologique sur $Y$ est \emph{locale sur
$Y$}. Si $Y$ est localement noethérien, ou un \emph{schéma}, pour que
$\mathscr P_\bullet$ soit homologiquement plat sur $Y$ en dimension $p$, il faut
et il suffit que le foncteur $\mathscr T_p$ soit \emph{exact} : la démonstration
a été faite dans le cas où $Y$ est localement noethérien au cours de la
démonstration de (7.7.5) ; le raisonnement est le même (basé sur (I, 9.4.2)
appliqué à un ouvert affine dans un schéma quasi-compact) lorsque $Y$ est un
schéma.
\end{env}

\begin{proposition}[7.8.3]
Les notations et hypothèses étant celles de (7.8.1), les conditions suivantes
sont équivalentes :
\begin{enumerate}
  \item[a)] $\mathscr P_\bullet$ est homologiquement plat sur $Y$ au point $y$
    en dimension $p$.
  \item[b)] Il existe un voisinage ouvert $U$ de $y$ dans $Y$ tel que
    $\mathscr T_p^U$ et $\mathscr T_{p+1}^U$ soient exacts à droite.
  \item[c)] Il existe un voisinage ouvert $U$ de $y$ dans $Y$ tel que
    $\mathscr T_p^U$ et $\mathscr T_{p-1}^U$ soient exacts à gauche.
  \item[d)] Il existe un voisinage ouvert $U$ de $y$ dans $Y$ tel que
    $\mathscr T_{p+1}^U$ soit exact à droite et $\mathscr T_{p-1}^U$ exact à
    gauche.
\end{enumerate}
\end{proposition}

\begin{proof}
Compte tenu de l'interprétation de $\mathscr T_p$ lorsque $Y$ est affine, cela
n'est qu'une traduction d'une partie de (7.3.3).
\end{proof}

\begin{proposition}[7.8.4]
Soient $Y$ un préschéma localement noethérien, $f:X\to Y$ un morphisme propre,
$\mathscr P_\bullet$ un complexe de $\mathscr O_X$-Modules cohérents et $Y$-plats
limité inférieurement, $\mathscr T_\bullet$ le foncteur défini par (7.7.1.1).
Pour tout $y\in Y$, les conditions suivantes sont équivalentes :
\begin{enumerate}
  \item[a)] $\mathscr P_\bullet$ est homologiquement plat sur $Y$ en $y$ en
    dimension $p$.
  \item[b)] Le foncteur $T_p^{\mathscr O_y}$ est exact.
  \item[c)] Il existe un entier $n_0$ tel que pour $n\geq n_0$, on ait
    \[
      \operatorname{long}T_p^{\mathscr O_y}
        (\mathscr O_y/\mathfrak m_y^{n+1})
      =
      \operatorname{long}T_p^{\mathscr O_y}(\kappa(y))
      \cdot
      \operatorname{long}(\mathscr O_y/\mathfrak m_y^{n+1})
      \tag{7.8.4.1}\label{III.7.8.4.1.fr}
    \]
    (où il s'agit de longueurs de $\mathscr O_y$-modules).
  \item[d)] Il y a un voisinage ouvert $U$ de $y$ tel que
    $(\mathscr H^{-p}(f,\mathscr P^\bullet))|U$ soit isomorphe à un
    $(\mathscr O_Y|U)$-Module de la forme $(\mathscr O_Y|U)^m$ et que, pour tout
    $(\mathscr O_Y|U)$-Module quasi-cohérent $\mathscr M$, l'homomorphisme
    canonique
    \[
      ((\mathscr H^{-p}(f,\mathscr P^\bullet))|U)
      \otimes_{\mathscr O_Y|U}\mathscr M
      \longrightarrow
      \mathscr H^{-p}
      \bigl(
        f,(\mathscr P^\bullet|U)
          \otimes_{\mathscr O_Y|U}\mathscr M
      \bigr)
      \tag{7.8.4.2}\label{III.7.8.4.2.fr}
    \]
    soit bijectif.
\end{enumerate}
Lorsque ces conditions sont vérifiées, on a en outre la propriété suivante :
\begin{enumerate}
  \item[e)] Il existe un voisinage de $y$ dans lequel la fonction
    $z\mapsto d_p(z)$ (définie dans (7.7.5.1)) est constante.
\end{enumerate}
De plus, si $Y$ est réduit au point $y$ ($0_{\mathrm I}$, 4.1.4), e) est
équivalente aux autres conditions.
\end{proposition}

\oldpage[III]{206}
\begin{proof}
En effet, la condition b) équivaut à dire que $T_p^{\mathscr O_y}$ et
$T_{p+1}^{\mathscr O_y}$ sont exacts à droite (7.3.3). L'équivalence de a) et b)
résulte alors de (7.7.10) et (7.8.3). Comme
$\mathscr O_y/\mathfrak m_y^{n+1}$ est artinien, et que
$T_p^{\mathscr O_y}(\mathscr O_y/\mathfrak m_y^{n+1})$ et
$T_p^{\mathscr O_y}(\kappa(y))$ sont des
$(\mathscr O_y/\mathfrak m_y^{n+1})$-modules de type fini (7.7.4), donc de
longueur finie, l'équivalence de b) et c) résulte encore de (7.7.10) et de
(7.3.5.7). Le fait que a) entraîne e), et lui est équivalente lorsque
$\mathscr O_y$ est réduit, est conséquence de (7.6.9). Enfin, a) entraîne que
(7.8.4.2) est bijectif en vertu de la définition (7.8.1) et de (7.7.5) ; d'autre
part, a) entraîne que $(\mathscr H^{-p}(f,\mathscr P^\bullet))_y$ est un
$\mathscr O_y$-module plat (7.3.3, f)), donc libre ($0_{\mathrm I}$, 10.1.3),
puisqu'il s'agit d'un $\mathscr O_y$-module de type fini en vertu de (7.7.4) ;
puisque $\mathscr H^{-p}(f,\mathscr P^\bullet)$ est un $\mathscr O_Y$-Module
cohérent (7.7.4), il est localement libre dans un voisinage de $y$
($0_{\mathrm I}$, 5.2.7). Inversement, il est clair que d) entraîne a) par
définition du foncteur $\mathscr T_p^U$ (7.7.2.2).
\end{proof}

\begin{proposition}[7.8.5]
Sous les hypothèses de (7.8.4), les conditions suivantes sont équivalentes :
\begin{enumerate}
  \item[a)] $\mathscr P_\bullet$ est homologiquement plat sur $Y$ en toute
    dimension $i\leq p$.
  \item[b)] Pour $i\leq p+1$ les foncteurs $\mathscr T_i$ sont exacts à droite.
  \item[c)] Pour $i\geq-p$, les $\mathscr O_Y$-Modules
    $\mathscr H^i(f,\mathscr P^\bullet)$ sont localement libres.
\end{enumerate}
\end{proposition}

\begin{proof}
L'équivalence de a) et b) est triviale (7.8.3) et a) entraîne c) en vertu de
(7.8.4). Inversement, supposons c) vérifiée ; notons d'autre part qu'on a
$\mathscr L_i=0$ pour $i<i_0$ (7.7.4), donc aussi $\mathscr T_i=0$ pour
$i\leq i_0$. Tout point $y\in Y$ a donc un voisinage affine
$U=\operatorname{Spec}(A)$ tel que $T_i^A(A)$ soit un $A$-module libre pour
$i\leq p$ ; en vertu de (7.3.7), on en conclut que
$\mathscr T_i^U=T_i^A$ est exact pour $i\leq p$.
\end{proof}

Nous allons surtout appliquer les critères de platitude cohomologique au cas où
le complexe $\mathscr P_\bullet$ est réduit à un \emph{seul}
$\mathscr O_X$-Module cohérent $\mathscr F$ \emph{plat} sur $Y$, pris égal à
$\mathscr P_0$ ; rappelons que l'on a alors
$\mathscr T_p(\mathscr M)=\mathrm R^{-p}f_*
(\mathscr F\otimes_{\mathscr O_Y}\mathscr M)$.

\begin{proposition}[7.8.6]
Soient $Y$ un préschéma localement noethérien, $f:X\to Y$ un morphisme propre
et plat, $y$ un point de $Y$ ; désignons par $X_y$ la fibre
$f^{-1}(y)=X\otimes_Y\kappa(y)$. Supposons que
$\Gamma(X_y,\mathscr O_{X_y})=R$ soit une $\kappa(y)$-algèbre séparable
(Bourbaki, \emph{Alg.}, chap. VIII, § 7, n\textsuperscript{o} 5) autrement dit
composée d'un nombre fini d'extensions séparables de degré fini de $\kappa(y)$.
Alors $\mathscr O_X$ est cohomologiquement plat sur $Y$ au point $y$ en
dimension 0.
\end{proposition}

\begin{proof}
En vertu de (7.8.4), on peut se borner au cas où $Y$ est le spectre de l'anneau
local $A=\mathscr O_y$ ; l'hypothèse que $f$ est plat entraîne
$\mathscr T_{-1}=0$, donc on voit déjà que $T_0^A$ est exact à gauche et tout
revient à voir qu'il est exact à droite ; en vertu de (7.7.10, c)), on est même
ramené au cas où $A=\mathscr O_y$ est \emph{artinien}. Soit $k'$ une extension
finie de $\kappa(y)$ qui soit un corps neutralisant de $R$, de sorte que
$R\otimes_{\kappa(y)}k'$ est composée directe d'un nombre fini de corps
isomorphes à $k'$. On sait qu'il existe un homomorphisme local de $A$ dans un
anneau local $A'$, faisant de $A'$ une $A$-algèbre \emph{libre} finie sur $A$,
et tel que le corps résiduel de $A'$ soit isomorphe à $k'$ ($0_{\mathrm I}$,
10.3.2). En vertu de (7.6.1), on est ramené à prouver que $T_0^{A'}$ est exact
à droite, autrement dit on peut supposer que $R$ est composé direct de $m$
corps isomorphes à $\kappa(y)$. Notons maintenant le lemme élémentaire
suivant :

\begin{lemma}[7.8.6.1]
Soit $Z$ un espace annelé en anneaux locaux ; pour que $Z$ soit connexe,

\oldpage[III]{207}
il faut et il suffit que l'anneau $\Gamma(Z,\mathscr O_Z)$ ne soit pas un produit
de deux anneaux non réduits à 0.
\end{lemma}

\begin{proof}
Il est clair en effet que si $Z$ est réunion de deux ouverts non vides disjoints,
$\Gamma(Z,\mathscr O_Z)$ est isomorphe au produit des deux anneaux
$\Gamma(Z_1,\mathscr O_Z)$ et $\Gamma(Z_2,\mathscr O_Z)$ non réduits à 0.
Inversement, dire que $\Gamma(Z,\mathscr O_Z)$ est un tel produit équivaut à
dire qu'il y a dans $\Gamma(Z,\mathscr O_Z)$ un idempotent $s$ distinct de 0 et
de 1 ; pour tout $z\in Z$, $s_z$ est alors un idempotent dans $\mathscr O_z$,
donc égal à 0 ou 1. Mais il est clair que l'ensemble des $z$ tels que $s_z=0$
est ouvert ; d'autre part, si $s_z=1$, on a par définition $s(z)\ne0$, donc
l'ensemble des $z$ où $s_z=1$ est aussi ouvert ($0_{\mathrm I}$, 5.5.2) ; d'où
la conclusion.
\end{proof}

Il résulte de ce lemme que $X_y$ a exactement $m$ composantes connexes $X'_i$
et que $\Gamma(X'_i,\mathscr O_{X'_i})=\kappa(y)$ pour tout $i$. Comme $A$ a été
supposé local et artinien, son spectre est réduit à un point, donc $X$ et $X_y$
ont même espace sous-jacent ; $X$ a donc $m$ composantes connexes $X_i$ telles
que $X'_i=X_i\otimes_Y\kappa(y)$. On est ainsi finalement ramené au cas où
$R=\kappa(y)$ ; en vertu de (7.7.10, b)), on est ramené à prouver que
l'homomorphisme canonique
\[
  \Gamma(X,\mathscr O_X)
  \longrightarrow
  \Gamma(X_y,\mathscr O_{X_y})
\]
est surjectif ; mais cela est trivial, car le composé
\[
  \Gamma(Y,\mathscr O_Y)=A
  \longrightarrow
  \Gamma(X,\mathscr O_X)
  \longrightarrow
  \Gamma(X_y,\mathscr O_{X_y})=\kappa(y)
\]
est déjà surjectif.
\end{proof}

\begin{corollary}[7.8.7]
Sous les hypothèses de (7.8.6), il existe un voisinage ouvert $U$ de $y$ tel
que :
\begin{enumerate}
  \item[(i)] $f_*(\mathscr O_X)|U$ soit isomorphe à un
    $(\mathscr O_Y|U)$-Module de la forme $(\mathscr O_Y|U)^m$.
  \item[(ii)] Pour tout $z\in U$, l'homomorphisme canonique
    \[
      (f_*(\mathscr O_X))_z
      \otimes_{\mathscr O_z}\kappa(z)
      \longrightarrow
      \Gamma(X_z,\mathscr O_{X_z})
    \]
    est bijectif.
\end{enumerate}
\end{corollary}

\begin{proof}
(i) résulte de (7.8.6) et (7.8.4).

(ii) résulte de ce que $\mathscr T_0^U$ est exact (pour $U$ convenablement
choisi), et de (7.7.5.3).
\end{proof}

\begin{corollary}[7.8.8]
Supposons vérifiées les conditions de (7.8.6) et en outre que
$\Gamma(X_y,\mathscr O_{X_y})=\kappa(y)$. Alors il existe un voisinage ouvert
$U$ de $y$ tel que l'homomorphisme canonique
$\mathscr O_Y|U\to f_*(\mathscr O_X)|U$ soit bijectif.
\end{corollary}

\begin{proof}
En effet, il résulte de (7.8.7, (ii)) que l'entier $m$ figurant dans (7.8.7,
(i)) est nécessairement égal à 1.
\end{proof}

\begin{corollary}[7.8.9]
Sous les hypothèses de (7.8.6), il existe un voisinage ouvert $U$ de $y$, un
$\mathscr O_U$-Module cohérent $\mathscr Q$ (déterminé à isomorphisme unique
près) et un isomorphisme de foncteurs en le $\mathscr O_U$-Module quasi-cohérent
$\mathscr M$ :
\[
  \mathrm R^1f_*(f^*(\mathscr M))
  \xrightarrow{\sim}
  \mathscr{H}\!om_{\mathscr O_U}(\mathscr Q,\mathscr M).
  \tag{7.8.9.1}\label{III.7.8.9.1.fr}
\]
\end{corollary}

\begin{proof}
En effet, l'hypothèse entraîne que $\mathscr T_0^U$ est exact pour un $U$
convenable ; il suffit donc d'appliquer l'équivalence de (7.7.5, a)) et
(7.7.5, b$'$)) au cas $p=0$ et en prenant pour $\mathscr P_\bullet$ le complexe
réduit à son terme de degré 0 égal à $\mathscr O_X$.
\end{proof}

\begin{remarks}[7.8.10]
\textnormal{(i)} Sous les conditions de (7.8.6), considérons la factorisation de
Stein de $f$ (4.3.3)
\[
\xymatrix{
  X \ar[r]^{f'} & Y' \ar[r]^{g} & Y .
}
\]

\oldpage[III]{208}
avec $Y'=\operatorname{Spec}(f_*(\mathscr O_X))$ ; le morphisme fini $g$ est alors
tel que $g_*(\mathscr O_{Y'})=f_*(\mathscr O_X)$ soit localement libre au
voisinage de $y$, et sa fibre en $y$ est le spectre d'une algèbre séparable sur
$\kappa(y)$ (II, 1.5.1). Nous en déduirons au chap. IV qu'il y a un voisinage
ouvert $U$ de $y$ dans $Y$ tel que pour la restriction $g^{-1}(U)\to U$ de $g$,
\emph{toute} fibre $g^{-1}(z)$ (où $z\in U$) soit spectre d'une algèbre séparable
sur $\kappa(z)$ (c'est ce que nous appellerons un \emph{revêtement étale} de
$U$) ; il résultera alors de (7.8.7, (ii)) que l'hypothèse faite sur le point
$y$ dans (7.8.6) est vérifiée aussi en tous les points d'un voisinage de $y$.

\medskip
\textnormal{(ii)} Nous verrons au chap. V que, même si $X$ est projectif sur $Y$
(et même s'il est en outre « simple » sur $Y$, propriété qui sera définie au
chap. IV), le $\mathscr O_U$-Module $\mathscr Q$ de (7.8.9) n'est pas
nécessairement localement libre ; en d'autres termes, $\mathscr O_X$ (sous ces
conditions) n'est pas nécessairement cohomologiquement plat en \emph{dimension
1} sur $Y$ au point $y$. Au chap. V, nous interpréterons $\mathscr Q$ comme le
faisceau des 1-différentielles du schéma de Picard de $X$ par rapport à $Y$ le
long de la section unité.
\end{remarks}

\subsection{Application aux morphismes propres : III. Invariance de la
caractéristique d'Euler-Poincaré et du polynôme de Hilbert}
\label{subsection:III.7.9.fr}
\label{III.7.9.fr}

\begin{env}[7.9.1]
\label{III.7.9.1.fr}
Soient $A$ un anneau, $M$ un $A$-module projectif de type fini ; rappelons
(Bourbaki, \emph{Alg. comm.}, chap. II, \S~5, n\textsuperscript{o}~2) qu'il
revient au même de dire que le $\mathscr O_X$-Module associé $\widetilde M$ sur
$X=\operatorname{Spec}(A)$ est \emph{localement libre de type fini}. Pour tout
$\mathfrak p\in\operatorname{Spec}(A)$ on appelle \emph{rang de $M$ en
$\mathfrak p$} et on note $\operatorname{rang}_{\mathfrak p}(M)$ le rang du
$A_{\mathfrak p}$-module libre $M_{\mathfrak p}$ (ou encore le rang en
$\mathfrak p$ du $\mathscr O_X$-Module localement libre $\widetilde M$). On a
donc
\[
  \operatorname{rang}_{\mathfrak p}M
  =\operatorname{rang}_{\mathfrak p}(M_{\mathfrak p})
  =\operatorname{rang}_{\kappa(\mathfrak p)}
  \bigl(M\otimes_A\kappa(\mathfrak p)\bigr).
  \tag{7.9.1.1}\label{III.7.9.1.1.fr}
\]
\end{env}

\begin{proposition}[7.9.2]
\label{III.7.9.2.fr}
Soit $P_\bullet$ un complexe fini de $A$-modules projectifs de type fini, et
pour tout $A$-module $M$, soit
$T_\bullet(M)=H_\bullet(P_\bullet\otimes_A M)$. Alors, pour tout
$\mathfrak p\in\operatorname{Spec}(A)$, on a
\[
  \sum_i(-1)^i\operatorname{rang}_{\kappa(\mathfrak p)}
  T_i\bigl(\kappa(\mathfrak p)\bigr)
  =\sum_i(-1)^i\operatorname{rang}_{\mathfrak p}(P_i).
  \tag{7.9.2.1}\label{III.7.9.2.1.fr}
\]
\end{proposition}

\begin{proof}
En effet, on a par définition
$T_i(\kappa(\mathfrak p))=H_i(P_\bullet\otimes_A\kappa(\mathfrak p))$ et,
compte tenu de (7.9.1.1), la formule (7.9.1.2) n'est autre que l'invariance de
la caractéristique d'Euler-Poincaré d'un complexe fini d'espaces vectoriels de
dimension finie par passage à l'homologie ($0_{\mathrm I}$, 11.10.2).
\end{proof}

\begin{corollary}[7.9.3]
\label{III.7.9.3.fr}
La fonction
\[
  \mathfrak p\longmapsto
  \sum_i(-1)^i\operatorname{rang}_{\kappa(\mathfrak p)}
  T_i\bigl(\kappa(\mathfrak p)\bigr)
\]
est localement constante dans $\operatorname{Spec}(A)$.
\end{corollary}

\begin{theorem}[7.9.4]
\label{III.7.9.4.fr}
Soient $Y$ un préschéma localement noethérien, $f:X\to Y$ un morphisme propre,
$\mathscr P_\bullet$ un complexe fini de $\mathscr O_X$-Modules cohérents et
$Y$-plats. Si l'on pose
\[
  \mathscr T_\bullet(\mathscr M)
  =\mathscr H^{-\bullet}
  \bigl(f,\mathscr P^\bullet\otimes_{\mathscr O_Y}\mathscr M\bigr)
  \qquad\text{(cf. (7.7.1.1))}
\]
la fonction

\oldpage[III]{209}
\[
  y\longmapsto
  \sum_i(-1)^i\operatorname{rang}_{\kappa(y)}
  \mathscr T_i\bigl(\kappa(y)\bigr)
  \tag{7.9.4.1}\label{III.7.9.4.1.fr}
\]
est localement constante dans $Y$.
\end{theorem}

\begin{proof}
On peut se borner au cas où $Y=\operatorname{Spec}(A)$ est affine d'anneau $A$
noethérien. Comme le complexe $\mathscr P_\bullet$ est fini, on sait (7.7.12,
(i)) qu'on a
\[
  \mathscr T_p(\mathscr M)
  =\mathscr H_p(\mathscr L_\bullet\otimes_{\mathscr O_Y}\mathscr M),
  \qquad \mathscr L_\bullet=\widetilde{L_\bullet},
\]
$L_\bullet$ étant un complexe \emph{fini} de $A$-modules projectifs de type
fini. Le théorème résulte alors de (7.9.3).
\end{proof}

\begin{env}[7.9.5]
\label{III.7.9.5.fr}
Sous les conditions de (7.9.4), la fonction (7.9.4.1) est \emph{constante}
lorsque $Y$ est \emph{connexe}. Lorsque $Y$ est connexe et non vide, on désigne
la valeur unique (entière) de (7.9.4.1) par
\[
  \operatorname{EP}(f,\mathscr P_\bullet)
  \quad\text{ou}\quad
  \operatorname{EP}(Y,\mathscr P_\bullet),
  \quad\text{ou simplement}\quad
  \operatorname{EP}(\mathscr P_\bullet)
\]
s'il ne peut en résulter de confusion, et l'on dit que cet entier est la
\emph{caractéristique d'Euler-Poincaré de $\mathscr P_\bullet$} relativement à
$f$ (ou à $Y$). Dans le cas général, on notera aussi
\[
  \operatorname{EP}(f,\mathscr P_\bullet;y)
  \quad\text{ou}\quad
  \operatorname{EP}(Y,\mathscr P_\bullet;y)
  \quad\text{ou}\quad
  \operatorname{EP}(\mathscr P_\bullet;y)
\]
le second membre de (7.9.4.1).
\end{env}

\begin{env}[7.9.6]
\label{III.7.9.6.fr}
Sous les hypothèses de (7.9.4) relativement à $X$, $Y$ et $f$, soit
\[
  0\longrightarrow\mathscr P'_\bullet
  \xrightarrow{u}\mathscr P_\bullet
  \xrightarrow{v}\mathscr P''_\bullet
  \longrightarrow0
\]
une suite exacte de complexes finis de $\mathscr O_X$-Modules cohérents et
$Y$-plats, les homomorphismes $u$ et $v$ étant de degrés \emph{pairs} $2d$,
$2d'$ respectivement. Comme $\mathscr T_\bullet$ est un foncteur homologique
(7.7.1), on a une suite exacte d'homologie
\[
  \cdots\longrightarrow
  \mathscr T_i(\mathscr P'_\bullet,\kappa(y))
  \longrightarrow
  \mathscr T_{i+2d}(\mathscr P_\bullet,\kappa(y))
  \longrightarrow
  \mathscr T_{i+2d+2d'}(\mathscr P''_\bullet,\kappa(y))
  \longrightarrow
  \mathscr T_{i-1}(\mathscr P'_\bullet,\kappa(y))
  \longrightarrow\cdots
\]
n'ayant d'ailleurs qu'un nombre fini de termes. En écrivant que la
caractéristique d'Euler-Poincaré de ce complexe est nulle ($0_{\mathrm I}$,
11.10.1), il vient aussitôt
\[
  \operatorname{EP}(\mathscr P_\bullet;y)
  =\operatorname{EP}(\mathscr P'_\bullet;y)
  +\operatorname{EP}(\mathscr P''_\bullet;y)
  \tag{7.9.6.1}\label{III.7.9.6.1.fr}
\]
pour tout $y\in Y$. Or, si par exemple $\mathscr P_\bullet=(\mathscr P_i)$ avec
$\mathscr P_i=0$ pour $i<0$, on a la suite exacte de complexes
\[
\xymatrix@C=2.4em{
  \cdots \ar[r] & 0 \ar[r]\ar[d] & 0 \ar[r]\ar[d]
    & 0 \ar[r]\ar[d] & 0 \ar[r] & \cdots \\
  \cdots \ar[r] & 0 \ar[r]\ar[d] & 0 \ar[r]\ar[d]
    & \mathscr P_1 \ar[r]\ar[d] & \mathscr P_2 \ar[r] & \cdots \\
  \cdots \ar[r] & 0 \ar[r]\ar[d] & \mathscr P_0 \ar[r]\ar[d]
    & \mathscr P_1 \ar[r]\ar[d] & \mathscr P_2 \ar[r] & \cdots \\
  \cdots \ar[r] & 0 \ar[r]\ar[d] & \mathscr P_0 \ar[r]\ar[d]
    & 0 \ar[r]\ar[d] & 0 \ar[r] & \cdots \\
  \cdots \ar[r] & 0 \ar[r] & 0 \ar[r]
    & 0 \ar[r] & 0 \ar[r] & \cdots
}
\]
les flèches verticales non nulles étant les automorphismes identiques ; on peut
appliquer (7.9.6.1) à cette suite exacte, d'où, par récurrence sur la longueur de
$\mathscr P_\bullet$, la formule
\[
  \operatorname{EP}(\mathscr P_\bullet;y)
  =\sum_i(-1)^i\operatorname{EP}(\mathscr P_i;y).
  \tag{7.9.6.2}\label{III.7.9.6.2.fr}
\]

\oldpage[III]{210}
où, pour tout $\mathscr O_X$-Module cohérent $\mathscr F$, plat sur $Y$, on
désigne par $\operatorname{EP}(\mathscr F;y)$ (ou
$\operatorname{EP}(f,\mathscr F;y)$ ou $\operatorname{EP}(Y,\mathscr F;y)$) la
fonction $\operatorname{EP}(\mathscr Q_\bullet;y)$ correspondant au complexe
$\mathscr Q_\bullet$ dont le seul terme $\ne0$ est de degré $0$ et égal à
$\mathscr F$. On voit donc qu'on peut se ramener à étudier les caractéristiques
d'Euler-Poincaré de complexes réduits à un seul terme.
\end{env}

\begin{proposition}[7.9.7]
\label{III.7.9.7.fr}
Sous les hypothèses de (7.9.4), soient $Y'$ un préschéma localement noethérien,
$g:Y'\to Y$ un morphisme,
\[
  X'=X\times_Y Y',
  \qquad f'=f_{(Y')}:X'\longrightarrow Y',
\]
$\mathscr P'_\bullet$ le complexe fini
\[
  \mathscr P_\bullet\otimes_{\mathscr O_Y}\mathscr O_{Y'}
\]
de $\mathscr O_{X'}$-Modules ; $\mathscr P'_\bullet$ est formé de
$\mathscr O_{X'}$-Modules cohérents et $Y'$-plats, et pour tout $y'\in Y'$, on a
\[
  \operatorname{EP}(\mathscr P'_\bullet;y')
  =\operatorname{EP}(\mathscr P_\bullet;g(y')).
  \tag{7.9.7.1}\label{III.7.9.7.1.fr}
\]
\end{proposition}

\begin{proof}
Les $\mathscr O_{X'}$-Modules $\mathscr P'_i$ étant images réciproques des
$\mathscr P_i$ par la projection $X'\to X$ sont cohérents, ils sont $Y'$-plats en
vertu de ($0_{\mathrm I}$, 6.2.1) et (I, 4.14.5), la question étant locale sur
$X$, $Y$ et $Y'$ ; enfin, on sait que $f'$ est propre (II, 5.4.2), donc le
premier membre de (7.9.7.1) est défini. La formule (7.9.7.1) résulte alors de
(6.10.4.2), (7.7.2) et du lemme (7.6.7), en se ramenant, comme on peut toujours
le faire, au cas où $Y$ et $Y'$ sont affines.
\end{proof}

\begin{proposition}[7.9.8]
\label{III.7.9.8.fr}
Supposons vérifiées les hypothèses de (7.9.4) et en outre qu'il existe un entier
$i_0$ tel que
\[
  \mathscr T_i(\kappa(y))=0
  \qquad\text{pour $i\ne i_0$ et tout $y\in Y$.}
\]
Alors
\[
  \mathscr T_{i_0}(\mathscr O_Y)
  =\mathscr H^{-i_0}(f,\mathscr P_\bullet)
\]
est un $\mathscr O_Y$-Module localement libre, dont le rang en $y\in Y$ est égal
à $(-1)^{i_0}\operatorname{EP}(f,\mathscr P_\bullet;y)$.
\end{proposition}

\begin{proof}
Remarquons d'abord que les hypothèses de (7.4.4) sont vérifiées par les
$\mathscr T_j^{\mathscr O_y}$, donc (7.4.7) leur est applicable, et l'hypothèse
entraîne que $\mathscr T_i^{\mathscr O_y}$ est \emph{nul} pour $i\ne i_0$ en
vertu de (7.5.3) ; en raison de (7.3.3), $\mathscr T_{i_0}$ est donc aussi exact,
et par suite (7.8.4), $\mathscr H^{-i_0}(f,\mathscr P_\bullet)$ est localement
libre et son rang en un point $y\in Y$ est
\[
  \operatorname{rang}_{\kappa(y)}
  \mathscr T_{i_0}(\kappa(y))
  =\operatorname{EP}(f,\mathscr P_\bullet;y)
\]
par définition, puisque $\mathscr T_i(\kappa(y))=0$ pour $i\ne i_0$.
\end{proof}

\begin{corollary}[7.9.9]
\label{III.7.9.9.fr}
Soient $Y$ un préschéma localement noethérien, $f:X\to Y$ un morphisme propre,
$\mathscr F$ un $\mathscr O_X$-Module cohérent et $Y$-plat ; on suppose qu'il
existe un entier $i_0$ tel que
\[
  H^i\bigl(f^{-1}(y),
  \mathscr F\otimes_{\mathscr O_Y}\kappa(y)\bigr)=0
  \qquad\text{pour tout $i\ne i_0$ et tout $y\in Y$.}
\]
Alors $\mathrm R^{i_0}f_*(\mathscr F)$ est un $\mathscr O_Y$-Module localement
libre, dont le rang en $y$ est égal à
$(-1)^{i_0}\operatorname{EP}(f,\mathscr F;y)$.
\end{corollary}

En particulier :

\begin{corollary}[7.9.10]
\label{III.7.9.10.fr}
Sous les conditions préliminaires de (7.9.9) pour $X$, $Y$ et $\mathscr F$,
supposons que l'on ait $\mathrm R^if_*(\mathscr F)=0$ pour tout $i>0$. Alors
$f_*(\mathscr F)$ est un $\mathscr O_Y$-Module localement libre, dont le rang en
$y$ est égal à $\operatorname{EP}(f,\mathscr F;y)$.
\end{corollary}

\begin{proof}
Il suffira, en vertu de (7.9.9) de prouver le lemme suivant :
\end{proof}

\begin{lemma}[7.9.10.1]
\label{III.7.9.10.1.fr}
Sous les hypothèses de (7.9.10), on a
\[
  H^i\bigl(f^{-1}(y),
  \mathscr F\otimes_{\mathscr O_Y}\kappa(y)\bigr)=0
\]
pour tout $i>0$ et tout $y\in Y$.
\end{lemma}

\begin{proof}
En effet, on peut se borner au cas où $Y=\operatorname{Spec}(A)$ est affine.
Avec les notations de (7.9.4), et $\mathscr P_\bullet$ étant réduit à son terme
de degré $0$ égal à $\mathscr F$, on a en effet
$\mathscr T_p(\mathscr O_Y)=0$ pour $p<0$ par hypothèse ; on conclut de (7.3.7)
que $\mathscr T_p$ est exact pour $p<0$, et le lemme résulte alors de
l'équivalence de (7.7.5, a)) et (7.7.5, d)).
\end{proof}

\oldpage[III]{211}
\begin{proposition}[7.9.11]
\label{III.7.9.11.fr}
Les hypothèses étant celles de (7.9.4), soit $\mathscr L$ un
$\mathscr O_X$-Module inversible très ample pour $Y$, et posons
$\mathscr P_\bullet(n)=\mathscr P_\bullet\otimes_{\mathscr O_X}
\mathscr L^{\otimes n}$ pour tout $n\in\mathbf Z$. Alors, pour tout $y\in Y$,
la fonction
\[
  n\longmapsto\operatorname{EP}(f,\mathscr P_\bullet(n);y)
  \tag{7.9.11.1}\label{III.7.9.11.1.fr}
\]
est un polynôme à coefficients dans $\mathbf Q$, qui est le même pour tous les
points d'une même composante connexe de $Y$.
\end{proposition}

\begin{proof}
Il est clair que $\mathscr P_\bullet(n)$ est un complexe de
$\mathscr O_X$-Modules $Y$-plats. En vertu de (7.9.6.2), on peut se borner au cas
où $\mathscr P_\bullet$ est réduit à un seul terme $\mathscr F\ne0$ de degré 0 ;
en outre, comme il s'agit de questions locales sur $Y$, on peut supposer $Y$
affine et $f$ projectif (II, 5.5.3) ; posons $X_y=f^{-1}(y)$, et soit
$\mathscr L_y=\mathscr L\otimes_{\mathscr O_Y}\kappa(y)$, qui est un
$\mathscr O_{X_y}$-Module très ample (II, 4.4.10) ; en vertu de (7.7.2), on a,
pour le foncteur $\mathscr T_\bullet$ relatif au complexe
$\mathscr P_\bullet(n)$,
\[
  \mathscr T_i(\kappa(y))
  =H^{-i}(X_y,\mathscr F_y\otimes\mathscr L_y^{\otimes n})
  \qquad
  \bigl(\text{où }\mathscr F_y=\mathscr F\otimes_{\mathscr O_Y}\kappa(y)\bigr) ;
\]
d'où résulte que $\operatorname{EP}(f,\mathscr F(n);y)$ n'est autre que la
caractéristique d'Euler-Poincaré $\chi_{\kappa(y)}(\mathscr F_y(n))$ définie dans
(2.5.1) ; le fait que (7.9.11.1) soit un polynôme résulte alors de (2.5.3) ; en
outre, pour chaque $n$, sa valeur est constante dans une composante connexe de
$Y$ (7.9.4), ce qui achève la démonstration.
\end{proof}

\begin{env}[7.9.12]
\label{III.7.9.12.fr}
Nous désignerons par $\operatorname{PH}(f,\mathscr P_\bullet;y)$ ou
$\operatorname{PH}(\mathscr P_\bullet;y)$ le polynôme (7.9.11.1), à coefficients
rationnels, et nous dirons que c'est le \emph{polynôme de Hilbert en $y$
relatif à $\mathscr P_\bullet$, $f$ et $\mathscr L$} (ou simplement le
\emph{polynôme de Hilbert en $y$ de $\mathscr P_\bullet$}, ou de $f$, s'il n'en
résulte pas de confusion) ; lorsque $Y$ est connexe non vide, on supprime la
mention de $y$ dans la notation et la terminologie. L'invariant ainsi obtenu
jouera un rôle essentiel au chap. V, dans la théorie des « modules » des
faisceaux cohérents quotients d'un faisceau cohérent donné.

Avec les notations de (7.9.6) et (7.9.11), on a
\[
  \operatorname{PH}(\mathscr P_\bullet;y)
  =\operatorname{PH}(\mathscr P'_\bullet;y)
  +\operatorname{PH}(\mathscr P''_\bullet;y)
  \tag{7.9.12.1}\label{III.7.9.12.1.fr}
\]
et en particulier
\[
  \operatorname{PH}(\mathscr P_\bullet;y)
  =\sum_i(-1)^i\operatorname{PH}(\mathscr P_i;y) ;
  \tag{7.9.12.2}\label{III.7.9.12.2.fr}
\]
cela résulte trivialement de (7.9.6.1) et (7.9.6.2). De même, avec les notations
et hypothèses de (7.9.7), on a
\[
  \operatorname{PH}(\mathscr P'_\bullet;y')
  =\operatorname{PH}(\mathscr P_\bullet;g(y')).
  \tag{7.9.12.3}\label{III.7.9.12.3.fr}
\]
\end{env}

La formule (7.9.12.2) ramène l'étude des polynômes de Hilbert d'un complexe à
celle des polynômes de Hilbert d'un seul $\mathscr O_X$-Module $Y$-plat. Ces
derniers admettent une interprétation remarquable indépendante de considérations
homologiques :

\begin{corollary}[7.9.13]
\label{III.7.9.13.fr}
Soient $Y$ un préschéma noethérien, $f:X\to Y$ un morphisme propre,
$\mathscr L$ un $\mathscr O_X$-Module inversible très ample pour $Y$,
$\mathscr F$ un $\mathscr O_X$-Module cohérent et $Y$-plat. Il existe un entier
$n_0$ tel que pour $n\geq n_0$, $f_*(\mathscr F(n))$ soit un
$\mathscr O_Y$-Module localement libre, de rang en $y\in Y$ égal à
$\operatorname{PH}(f,\mathscr F;y)(n)$.
\end{corollary}

\begin{proof}
Comme le morphisme $f$ est projectif (II, 5.5.3), il existe $n_0$ tel que pour
$n\geq n_0$

\oldpage[III]{212}
on ait $\mathrm R^if_*(\mathscr F(n))=0$ pour tout $i>0$ (2.2.1) ; la conclusion
résulte donc de (7.9.10).
\end{proof}

Le critère de platitude suivant sera important dans la théorie des « modules »
des faisceaux cohérents du chap. V :

\begin{proposition}[7.9.14]
\label{III.7.9.14.fr}
Soient $Y$ un préschéma noethérien, $f:X\to Y$ un morphisme projectif,
$\mathscr L$ un $\mathscr O_X$-Module inversible ample pour $f$, et posons
$\mathscr F(n)=\mathscr F\otimes\mathscr L^{\otimes n}$ pour tout
$\mathscr O_X$-Module $\mathscr F$ et tout $n\in\mathbf Z$. Pour qu'un
$\mathscr O_X$-Module cohérent $\mathscr F$ soit $Y$-plat, il faut et il suffit
qu'il existe un entier $n_0$ tel que, pour tout $n\geq n_0$,
$f_*(\mathscr F(n))$ soit un $\mathscr O_Y$-Module localement libre.
\end{proposition}

\begin{proof}
La nécessité de la condition se démontre comme dans (7.9.13) (le résultat de
(2.2.1) s'appliquant à un faisceau ample $\mathscr L$, puisque $f$ est
projectif). Pour démontrer la réciproque, on peut se borner au cas où $Y$ est
affine d'anneau $A$ ; en vertu de l'hypothèse et de (2.2.2, (i)), les
$A$-modules $\Gamma(X,\mathscr F(n))$ sont de type fini et projectifs (Bourbaki,
\emph{Alg. comm.}, chap. II, \S~5, n\textsuperscript{o}~2, th. 1). Soit $S$
l'anneau gradué
\[
  S=\bigoplus_{n\geq0}\Gamma(X,\mathscr L^{\otimes n}) ;
\]
on sait que $X$ s'identifie canoniquement à $\operatorname{Proj}(S)$ (II,
4.5.2, (b) et 5.4.4). Soit
\[
  M=\bigoplus_{n\geq n_0}\Gamma(X,\mathscr F(n)) ;
\]
remplaçant au besoin $\mathscr L$ par une puissance $\mathscr L^{\otimes d}$,
on peut supposer que $S$ est engendré par un nombre fini d'éléments de degré 1
(2.3.5.1), et il résulte alors de (II, 2.7.5 et 2.7.2) que $\mathscr F$
s'identifie à $\mathscr P\!roj_0(M)$. Pour tout élément homogène $g\in S$ de
degré $>0$, on a donc $\Gamma(X_g,\mathscr F)=M_{(g)}$ ; or, $M$, somme directe
de $A$-modules projectifs, est un $A$-module plat, donc il en est de même de
$M_g$ ($0_{\mathrm I}$, 6.3.2), et par suite aussi de $M_{(g)}$, qui est un
facteur direct de $M_g$ ($0_{\mathrm I}$, 6.1.2). On en conclut (I, 4.14.5) que
$\mathscr F$ est $Y$-plat en tout point de $X_g$, et comme les $X_g$ recouvrent
$X$, la proposition est démontrée.
\end{proof}

\bigskip
\noindent\emph{(A suivre.)}

\clearpage
\oldpage[III]{213}
\section*{INDEX DES NOTATIONS}

\begin{flushleft}
$\mathbf{Tor}_n^A(P_\bullet,Q_\bullet)$ : 6.3.1.

\medskip
$\mathscr{T}\!or_n^{\mathscr O_S}(\mathscr P_\bullet,\mathscr Q_\bullet)$,
$\mathscr{T}\!or_n^S(\mathscr P_\bullet,\mathscr Q_\bullet)$ : 6.4.1 et 6.5.1.

\medskip
$\mathscr{T}\!or_n^S(\mathscr F,\mathscr G)$ : 6.5.1.

\medskip
$\mathscr{T}\!or_q^S(\mathscr P_\bullet^{(1)},\mathscr P_\bullet^{(2)},\ldots,
\mathscr P_\bullet^{(m)})$ : 6.5.15.

\medskip
$\mathbf H^{-n}(\mathfrak U^{(i)},\mathscr P_\bullet^{(i)})$,
$\mathbf H^{-n}(X^{(i)},\mathscr P_\bullet^{(i)})$ : 6.6.1.

\medskip
$\mathbf{Tor}_n^A(L_{\bullet\bullet}^{(1)},L_{\bullet\bullet}^{(2)})$,
$\mathbf{Tor}_n^S(\mathfrak U^{(1)},\mathfrak U^{(2)};
\mathscr P_\bullet^{(1)},\mathscr P_\bullet^{(2)})$ : 6.6.2.

\medskip
$\mathscr{T}\!or_n^S(\mathfrak U^{(1)},\mathfrak U^{(2)};
\mathscr F^{(1)},\mathscr F^{(2)})$,
$\mathscr{T}\!or_n^S(\mathfrak U^{(1)},\mathfrak U^{(2)};
\mathscr P_\bullet^{(1)},\mathscr P_\bullet^{(2)})$ : 6.6.2.

\medskip
$\mathscr{T}\!or_n^S(f_1,f_2;\mathscr P_\bullet^{(1)},
\mathscr P_\bullet^{(2)})$ : 6.6.8, 6.7.1 et 6.7.3.

\medskip
$\mathscr{T}\!or_n^S((f_i)_{i\in I};(\mathscr P_\bullet^{(i)})_{i\in I})$,
$\mathscr{T}\!or_n^S(f_1,\ldots,f_m;\mathscr P_\bullet^{(1)},\ldots,
\mathscr P_\bullet^{(m)})$ : 6.7.11.

\medskip
$\mathbf{Ab}$, $\mathbf{Ab}_A$ : 7.1.1.

\medskip
$T_{(B)}$, $T^{(B)}$, $T\otimes_A B$
($T$ foncteur de $\mathbf{Ab}_A$ dans $\mathbf{Ab}$) : 7.1.3.

\medskip
$t_M$ : 7.2.2.

\medskip
$T_{\mathfrak p}$ : 7.1.4.

\medskip
$\mathscr T_\bullet^{Y'}(\mathscr M')$,
$\mathscr T_\bullet^U(\mathscr M|U)$ : 7.7.2.

\medskip
$\operatorname{EP}(f,\mathscr P_\bullet;y)$,
$\operatorname{EP}(Y,\mathscr P_\bullet;y)$,
$\operatorname{EP}(\mathscr P_\bullet;y)$ : 7.9.5.

\medskip
$\operatorname{PH}(f,\mathscr P_\bullet;y)$,
$\operatorname{PH}(\mathscr P_\bullet;y)$ : 7.9.12.
\end{flushleft}

\clearpage
\oldpage[III]{214}
\section*{INDEX TERMINOLOGIQUE}

\begin{flushleft}
Caractéristique d'Euler-Poincaré d'un complexe de Modules : 7.9.5.

\medskip
Cohomologiquement plat en un point $y\in Y$ en dimension $p$,
cohomologiquement plat sur $Y$ en dimension $p$, cohomologiquement plat sur
$Y$ au point $y$, cohomologiquement plat sur $Y$ (complexe de
$\mathscr O_X$-Modules $Y$-plats) : 7.8.1.

\medskip
Complexes homotopes : 6.1.4.

\medskip
Espace annelé de dimension cohomologique $\leq n$ : 6.5.5.

\medskip
Extension des scalaires dans un foncteur covariant additif : 7.1.3.

\medskip
Faisceau d'anneaux de dimension cohomologique $\leq n$ : 6.5.5.

\medskip
Formule de Künneth : 6.7.8.

\medskip
Homologiquement plat en un point $y\in Y$ en dimension $p$,
homologiquement plat sur $Y$ en dimension $p$, homologiquement plat sur $Y$
au point $y$, homologiquement plat sur $Y$ (complexe de $\mathscr O_X$-Modules
$Y$-plats) : 7.8.1.

\medskip
Homotopisme : 6.1.4.

\medskip
Hypertor de deux complexes de $A$-modules : 6.3.1.

\medskip
Hypertor local de deux complexes de Modules : 6.4.1.

\medskip
Hypertor global de deux complexes de Modules : 6.6.2, 6.6.8 et 6.7.3.

\medskip
Polynôme de Hilbert relatif à un complexe de Modules : 7.9.12.

\medskip
Suites spectrales de Künneth : 6.7.3.
\end{flushleft}

\clearpage
\oldpage[III]{215}
\section*{TABLE DES MATIÈRES}

\begin{flushleft}
\textsc{Chapitre III.} --- \textbf{Étude cohomologique des faisceaux cohérents
(suite)} \dotfill 5

\medskip
\S~6. Foncteurs Tor locaux et globaux ; formule de Künneth \dotfill 5

6.1. Introduction \dotfill 5

6.2. Hypercohomologie des complexes de Modules sur un préschéma \dotfill 6

6.3. Hypertor de deux complexes de modules \dotfill 9

6.4. Foncteurs hypertor locaux de complexes de Modules quasi-cohérents ; cas
des schémas affines \dotfill 14

6.5. Foncteurs hypertor locaux de complexes de Modules quasi-cohérents : cas
général \dotfill 16

6.6. Foncteurs hypertor globaux de complexes de Modules quasi-cohérents et
suites spectrales de Künneth : cas de la base affine \dotfill 21

6.7. Foncteurs hypertor globaux de complexes de Modules quasi-cohérents et
suites spectrales de Künneth : cas général \dotfill 25

6.8. Les suites spectrales d'associativité des hypertor globaux \dotfill 32

6.9. Les suites spectrales de changement de base dans les hypertor globaux
\dotfill 34

6.10. Structure locale de certains foncteurs cohomologiques \dotfill 39

\medskip
\S~7. Étude du changement de base dans les foncteurs homologiques covariants
de Modules \dotfill 43

7.1. Foncteurs de $A$-modules \dotfill 43

7.2. Caractérisation du foncteur produit tensoriel \dotfill 44

7.3. Critères d'exactitude des foncteurs homologiques de modules \dotfill 48

7.4. Critères d'exactitude pour les foncteurs
$\mathrm H_i(P\otimes_A M)$ \dotfill 53

7.5. Cas des anneaux locaux noethériens \dotfill 58

7.6. Descente des propriétés d'exactitude. Théorème de semi-continuité et
critère d'exactitude de Grauert \dotfill 60

7.7. Application aux morphismes propres : I. La propriété d'échange
\dotfill 65

7.8. Application aux morphismes propres : II. Critères de platitude
cohomologique \dotfill 72

7.9. Application aux morphismes propres : III. Invariance de la
caractéristique d'Euler-Poincaré et du polynôme de Hilbert \dotfill 76

\begin{flushright}
\emph{Reçu le 15 décembre 1962.}
\end{flushright}
\end{flushleft}

% La page imprimée 216 est blanche.
\clearpage
\thispagestyle{empty}
\null
\clearpage

\oldpage[III]{217}
\section*{ERRATA ET ADDENDA}
\begin{center}
(Liste 2)

\medskip
A) \emph{Erreurs typographiques}
\end{center}

\noindent\textbf{$(0_{\mathrm I},\,3.2.1)$} Ligne 2 de la p. 27, remplacer $X$
par $U$.

\noindent\textbf{$(0_{\mathrm I},\,3.2.6)$} Ligne 1 du bas de la p. 27 et ligne
1 de la p. 28, remplacer (3 fois) $\mathfrak S_\lambda$ par
$\mathscr H_\lambda$.

\noindent\textbf{$(0_{\mathrm I},\,3.4.5)$} Ajouter une parenthèse devant
3.4.5).

\noindent\textbf{$(0_{\mathrm I},\,4.2.1)$} Ligne 10 de la p. 39, remplacer
$\psi_*(A)$ par $\psi_*(\mathscr A)$.

\noindent\textbf{$(0_{\mathrm I},\,5.1.3)$} Ligne 16 de la p. 45, remplacer
$\mathscr F|V$ par $\mathscr F|U$.

\noindent\textbf{$(0_{\mathrm I},\,5.3.9)$} Ligne 12 de la p. 48, avant
« telle que », ajouter : au-dessus d'un voisinage ouvert $U$ de $x$.

\noindent\textbf{$(0_{\mathrm I},\,7.6.9)$} Ligne 3 du bas de la p. 73,
remplacer $I_\lambda$ par $J_\lambda$.

\noindent\textbf{(I, 1.1.15)} Ligne 3 de la p. 83, remplacer $\mathfrak a$ par
$\mathfrak a\ne A$.

\noindent\textbf{(I, 1.4.1)} Ligne 2 du bas de la p. 90, remplacer
$\widetilde N$ par $N$.

\noindent\textbf{(I, 2.3.1)} Ligne 3 du bas de la p. 101, remplacer (0, 4.1.6)
par (0, 4.1.7).

\noindent\textbf{(I, 2.3.2)} Ligne 11 de la p. 101, remplacer $B$ par $B_s$ et
$C$ par $C_t$.

\noindent\textbf{(I, 2.5.5)} Ligne 19 de la p. 104, remplacer $S$-morphisme
par $S$-préschéma.

\noindent\textbf{(I, 3.3.7)} Ligne 17 de la p. 109, remplacer $Y_{(S)}$ par
$Y_{(S')}$.

\noindent\textbf{(I, 3.4.5)} Ligne 9 de la p. 113, remplacer $\mathrm s$ par
$s$.

\noindent\textbf{(I, 3.4.8)} Ligne 4 de la p. 114, remplacer $p(x)$ par
$f(x)$.

\noindent\textbf{(I, 3.5.1)} Ligne 3 de la p. 115, remplacer $1_X\times g$
par $1_{X'}\times g$.

\noindent\textbf{(I, 3.7.2)} Ligne 8 de la p. 119, remplacer le premier $X$
par $X'$.

\noindent\textbf{(I, 4.2.2)} Ligne 8 de la p. 123, dans la flèche verticale de
gauche du diagramme, remplacer $\alpha_{\psi(y)}$ par $\varphi_{\psi(y)}$.

\noindent\textbf{(I, 4.4.3)} Ligne 16 de la p. 126, remplacer isomorphée par
isomorphes.

\noindent\textbf{(I, 4.5.5)} Ligne 17 du bas de la p. 127, remplacer (4.2.4)
par (4.2.5).

\noindent\textbf{(I, 5.1.4)} Ligne 14 du bas de la p. 128, remplacer (2.1.7)
par (2.1.8).

\noindent\textbf{(I, 5.3.13)} Ligne 20 de la p. 134, remplacer (4.2.4) par
(4.2.5).

\noindent\textbf{(I, 6.4.2)} Ligne 1 du bas de la p. 147, remplacer
recouvrement par recouvrement fini.

\noindent\textbf{(I, 9.1.13)} Ligne 15 de la p. 171, remplacer
$p^{-1}(\mathscr F)$ par $p^*(\mathscr F)$.

\noindent\textbf{(I, 9.5.11)} Ligne 7 de la p. 179, remplacer $Y$ par $Y'$.

\noindent\textbf{(I, 9.6.5)} Ligne 17 du bas de la p. 180, remplacer
$(0,\,4.1.4)$ par $(0,\,4.1.3)$.

\noindent\textbf{(I, 10.12.2)} Ligne 14 du bas de la p. 206, remplacer
$(X_n)_{(S_n)}$ par $(X_n)_{(S_m)}$.

\clearpage
\oldpage[III]{218}
\noindent\textbf{(I), Index terminologique :} Ligne 6 du bas de la p. 219,
remplacer 0, 4.1.4 par 0, 4.1.5. Lignes 16, 17, 18 du bas de la p. 222,
remplacer I, 2.1.7 par I, 2.1.8 ; ligne 19 du bas de la p. 222, remplacer I,
2.1.8 par I, 2.1.7.

\noindent\textbf{(II, 1.5.2)} Ligne 17 de la p. 12, insérer un $f$ à côté de
la flèche verticale de gauche du diagramme.

\noindent\textbf{(II, 4.2.7)} Ligne 16 du bas de la p. 75, remplacer (III,
2.1.14) par (III, 2.1.13).

\noindent\textbf{(II, 5.2.1)} Ligne 16 du bas de la p. 97, remplacer $X-U$ par
$U$. Ligne 8 du bas de la p. 97, remplacer $X'$ par $X_f$.

\noindent\textbf{(II, 5.3.6)} Ligne 16 de la p. 100, remplacer (4.6.18) par
(4.6.17).

\noindent\textbf{(II, 6.2.7)} Ligne 9 du bas de la p. 116, remplacer chap. V
par chap. IV.

\noindent\textbf{(II, 6.6.5)} Ligne 15 du bas de la p. 132, remplacer chap. V
par chap. IV.

\noindent\textbf{(II, 7.4.12)} Ligne 1 de la p. 151, remplacer chap. V par
chap. III.

\noindent\textbf{(II, 8.1.4)} Ligne 15 de la p. 154, remplacer (III, 2.3.8)
par (III, 2.3.7).

\noindent\textbf{(II, 8.10.5)} Ligne 17 du bas de la p. 187, remplacer
$g_{(j)x}$ par $g_{j(x)}$.

\noindent\textbf{$(0_{\mathrm{III}},\,10.2.7)$} Ligne 1 de la p. 20,
remplacer $u$-adique par $\mathfrak m$-adique. Ligne 13 du bas de la p. 20,
remplacer $k$ par $k_i$.

\noindent\textbf{$(0_{\mathrm{III}},\,11.1.1)$} Ligne 18 de la p. 23,
remplacer $\operatorname{Im}(d_r^{p+r,q-r+1})$ par
$\operatorname{Im}(d_r^{p-r,q+r-1})$.

\noindent\textbf{$(0_{\mathrm{III}},\,11.8.6)$} Ligne 14 du bas de la p. 45,
remplacer les flèches $\to$ par $\rightsquigarrow$ (4 fois).

\noindent\textbf{$(0_{\mathrm{III}},\,12.3.4)$} Ligne 4 de la p. 61,
remplacer $\rightsquigarrow$ par $\to$.

\noindent\textbf{$(0_{\mathrm{III}},\,13.2.4)$} Ligne 2 du bas de la p. 67,
remplacer $\varinjlim$ par $\varprojlim$ (2 fois).

\noindent\textbf{(III, 1.1.7)} Ligne 15 de la p. 85, remplacer $A^r$ par
$\bigwedge(A^r)$ (2 fois).

\noindent\textbf{(III, 1.2.4)} Ligne 7 de la p. 87, remplacer
$\Gamma(U,\mathscr F)$ par $\Gamma(U,\mathscr O_X)$ et remplacer $\mathscr G$
par $\mathscr F$. Ligne 8 de la p. 87, remplacer
$\mathrm H^p(\mathfrak U,\mathscr G)$ par $\mathrm H^p(\mathfrak U,\mathscr F)$.

\noindent\textbf{(III, 2.5.3.1)} Ligne 3 de la p. 111, remplacer $n>0$ par
$n\geq0$. Ligne 8 de la p. 111, remplacer $-r\leq n\leq0$ par
$-r\leq n<0$.

\noindent\textbf{(III, 3.1.2)} Ligne 13 de la p. 116, remplacer $\mathscr G$
par $\mathscr G\in\mathbf K'$.

\noindent\textbf{(III, 4.3.6)} Ligne 7 de la p. 133, remplacer
$K'=K\otimes_A\widehat A$ par $K'\supset K\otimes_A\widehat A$. Ligne 16 de
la p. 133, remplacer $R(X)$ par $R(Y)$.

\noindent\textbf{(III, 4.4.12)} Ligne 10 de la p. 138, remplacer chap. V par
chap. IV.

\noindent\textbf{(III, 4.6.6)} Ligne 19 de la p. 142, remplacer « chap. V » par
« dans un paragraphe ultérieur ».

\begin{center}
B) \emph{Modifications de texte}\footnote{Pour faciliter la recherche des
références, les modifications appartenant à la liste d'errata insérée dans le
chapitre N seront désormais désignées par $\mathbf{Err}_{\mathrm N}$ suivi d'un
numéro.}
\end{center}

\noindent\textbf{$(\mathbf{Err}_{\mathrm{III}},\,1)$} Dans
$(0_{\mathrm I},\,5.1.3)$, ligne 18 de la p. 45, remplacer « somme directe »
par « somme directe finie ».

\noindent\textbf{$(\mathbf{Err}_{\mathrm{III}},\,2)$} Dans
$(0_{\mathrm I},\,5.4.3)$, lignes 8 à 4 du bas de la p. 49, remplacer depuis
« Dans le cas général\ldots{} » par le texte suivant :

Dans le cas général où $X$ est un espace annelé tel que $\mathscr O_x$ soit un
anneau \emph{local} pour
\oldpage[III]{219}
tout $x\in X$, si $\mathscr L$ est un $\mathscr O_X$-Module \emph{de type fini}
tel qu'il existe un $\mathscr O_X$-Module $\mathscr F$ pour lequel
$\mathscr L\otimes_{\mathscr O_X}\mathscr F$ soit \emph{isomorphe à}
$\mathscr O_X$, le raisonnement précédent montre que pour tout $x\in X$,
$\mathscr L_x$ est isomorphe à $\mathscr O_x$. On en déduit que $\mathscr L$ est
alors \emph{inversible}. En effet, pour tout $x\in X$, soient $U$ un voisinage
ouvert de $x$ tel que $\mathscr L|U$ soit engendré par $n$ sections
$s_i\ (1\leq i\leq n)$ au-dessus de $U$ (5.2.1). On peut supposer par exemple
que $s=s_1$ est telle que $s_x\ne0$, et comme $\mathscr L_x$ est isomorphe à
$\mathscr O_x$, il existe pour chaque $i$ une section $t_i$ de $\mathscr O_X$
au-dessus d'un voisinage ouvert $V_i\subset U$ de $x$ telle que
$(t_i)_x s_x=(s_i)_x$. Il y a par suite un voisinage ouvert
$V\subset\bigcap_{i=1}^nV_i$ de $x$ tel que $t_i s=s_i$ dans $V$, autrement
dit $\mathscr L|V$ est engendré par \emph{l'unique} section $s|V$. En outre,
si $z$ est une section au-dessus d'un ouvert $W\subset V$ du noyau de
l'homomorphisme $\mathscr O_X|V\to\mathscr L|V$ défini par $s$ (5.1.1), $z_y$
annule $\mathscr L_y$ pour tout $y\in W$, donc $z_y=0$ par hypothèse et par
suite $z=0$, ce qui achève de prouver notre assertion. De plus, la
considération du produit tensoriel
$\mathscr L^{-1}\otimes\mathscr L\otimes\mathscr F$ montre aussitôt que
$\mathscr F$ est isomorphe à $\mathscr L^{-1}$.

\noindent\textbf{$(\mathbf{Err}_{\mathrm{III}},\,3)$} Dans
$(0_{\mathrm I},\,7.2.4)$, il faut imposer la condition que les puissances
$\mathfrak J^n$ sont des idéaux \emph{fermés} dans l'anneau admissible $A$ pour
que la proposition soit exacte. La démonstration donnée est incorrecte, et
l'énoncé rectifié résulte de Bourbaki, \emph{Top. gén.}, chap. III,
3\textsuperscript{e} éd., \S~3, n\textsuperscript{o}~5, cor. 1 de la prop. 9.
Il faut de même supposer les $\mathfrak J^n$ \emph{fermés} dans $A$ dans les
énoncés $(0_{\mathrm I},\,7.2.5)$ et $(0_{\mathrm I},\,7.2.6)$.

\noindent\textbf{$(\mathbf{Err}_{\mathrm{III}},\,4)$} Après la proposition
(I, 2.2.5), ajouter : avec les notations de (I, 2.2.5), si $\mathfrak J$ est un
idéal de $A$, on note $\widetilde{\mathfrak J\mathscr F}$ le sous-$\mathscr O_X$-Module
$\mathfrak J\widetilde{\mathscr F}$ défini dans $(0,\,4.3.5)$.

\noindent\textbf{$(\mathbf{Err}_{\mathrm{III}},\,5)$} Dans (I, 3.4.5), ligne 1
du bas de la p. 112, remplacer « \emph{un corps} $K$ » par « \emph{un corps
algébriquement clos} $K$ ». Ligne 1 de la p. 113, remplacer « \emph{extension} »
par « \emph{extension algébriquement close} ». Lignes 1 à 4 de la p. 113,
remplacer depuis « $K$ sera appelé\ldots{} » par le texte suivant :

Pour tout point de $X$ à valeurs dans un corps $K$, $K$ sera appelé le
\emph{corps des valeurs} du point correspondant, et si $x$ est la localité de
ce point on dit encore que ce dernier est \emph{localisé en} $x$. On définit
ainsi une application $X(K)\to X$, faisant correspondre à un point à valeurs
dans $K$ sa localité.

Lignes 7 et 16 de la p. 113, supprimer « géométrique » ; lignes 10 et 13 de
la p. 113, supprimer « géométriques ». Lignes 10 et 11 du bas de la p. 115,
supprimer « \emph{géométrique} ».

\noindent\textbf{$(\mathbf{Err}_{\mathrm{III}},\,6)$} À la fin de la
démonstration de (I, 3.6.1), ligne 12 du bas de la p. 117, ajouter : Si l'on
pose
$X'=X\times_Y\operatorname{Spec}(\mathscr O_y/\mathfrak a_y)$, alors, pour tout
point $x\in X'$, identifié par $p$ à un point de $X$, on a
$\mathscr O_{X',x}=\mathscr O_{X,x}/\mathfrak a_y\mathscr O_{X,x}$. La question
étant en effet locale sur $X$ et $Y$, on peut supposer que
$X=\operatorname{Spec}(B)$, $Y=\operatorname{Spec}(A)$, et l'on a
$(B/\mathfrak a_yB)_x=B_x/\mathfrak a_yB_x$ par platitude $(0,\,1.3.2)$.

\noindent\textbf{$(\mathbf{Err}_{\mathrm{III}},\,7)$} Dans (I, 5.1.1), ligne 3
du bas de la p. 127, remplacer « $\mathscr O_X$-Module quasi-cohérent » par
« \emph{Idéal quasi-cohérent de} $\mathscr B$ ».

\noindent\textbf{$(\mathbf{Err}_{\mathrm{III}},\,8)$} Après (I, 5.1.10),
ligne 16 de la p. 131, ajouter :

\emph{Remarque (5.1.11).} --- Une légère adaptation du raisonnement fait dans
(5.1.9) montre que la conclusion reste valable \emph{sans supposer a priori
que $X$ soit un préschéma}, mais

\oldpage[III]{220}
en supposant seulement que $(X,\mathscr O_X)$ soit un espace annelé en anneaux
locaux, $\mathscr J$ un Idéal de $\mathscr O_X$ tel que $\mathscr J^n=0$, que
l'espace annelé $X_0=(X,\mathscr O_X/\mathscr J)$ soit un schéma affine et enfin
que les Idéaux $\mathscr J^k/\mathscr J^{k+1}$ soient des
$\mathscr O_{X_0}$-Modules quasi-cohérents. Il suffit en effet d'utiliser
(1.8.1) (cf. $(\mathbf{Err}_{\mathrm{II}})$) au lieu de (2.2.4) dans le
raisonnement.

\noindent\textbf{$(\mathbf{Err}_{\mathrm{III}},\,9)$} Dans (I, 5.3.1), ligne
11 de la p. 132, insérer, après « ou $\Delta_X$ », « ou $\Delta_\varphi$ si
$\varphi:X\to S$ est le morphisme structural ».

\noindent\textbf{$(\mathbf{Err}_{\mathrm{III}},\,10)$} Dans (I, 5.3.9), la
démonstration est insuffisante, car elle ne prouve pas que $\Delta_X(X)$ soit
localement fermé dans $X\times_SX$. Pour avoir une démonstration correcte, il
suffit d'utiliser (4.2.4, $a$)) : pour tout $x\in X$ et tout voisinage affine
$U$ de $x$ dans $X$, $U\times_SU$ est un voisinage affine de $\Delta_X(x)$ ;
compte tenu de (5.3.16) (dont la démonstration n'utilise que la définition
(5.3.1.1)), on est ramené à démontrer (5.3.9) lorsque
$S=\operatorname{Spec}(B)$ et $X=\operatorname{Spec}(A)$ sont des schémas
affines ; il est clair alors en vertu de (5.3.1.1) que $\Delta_X$ correspond à
l'homomorphisme canonique $A\otimes_BA\to A$ qui transforme $x\otimes y$ en
$xy$ ; cet homomorphisme étant surjectif, $\Delta_X$ est dans ce cas une
immersion fermée (4.2.3), ce qui achève la démonstration.

\noindent\textbf{$(\mathbf{Err}_{\mathrm{III}},\,11)$} Dans (I, 7.1.15) et
(I, 7.1.16), lignes 7 et 15 du bas de la p. 158, supprimer
« \emph{géométriques} ».

\noindent\textbf{$(\mathbf{Err}_{\mathrm{III}},\,12)$} Remplacer (I, 7.4.7)
par : On étend (par abus de langage) les définitions de (7.4.1) au cas où
$X$ est un préschéma \emph{réduit} dont tout point admet un voisinage ouvert
n'ayant qu'un nombre \emph{fini} de composantes irréductibles ; il résulte
alors de (7.3.4) et (7.4.6) que, pour un $\mathscr O_X$-Module quasi-cohérent de
type fini $\mathscr F$, dire que $\mathscr F$ est un \emph{faisceau de torsion}
équivaut à dire que $\operatorname{Supp}(\mathscr F)$ \emph{ne contient aucune
composante irréductible de $X$}.

\noindent\textbf{$(\mathbf{Err}_{\mathrm{III}},\,13)$} Dans (I, 10.11.7),
lignes 17 à 23 de la p. 205, la fin de la démonstration depuis « Reste à
prouver\ldots{} » est inutile en vertu de (10.11.6), les faisceaux considérés
étant cohérents par $(0,\,5.3.5)$.

\noindent\textbf{$(\mathbf{Err}_{\mathrm{III}},\,14)$} Dans (II, 1.7.8),
après la ligne 10 du bas de la p. 16, ajouter : lorsque
$\mathscr E=\mathscr O_S^n$, on écrit aussi $\mathbf V_S^n$ au lieu de
$\mathbf V(\mathscr E)$ ; si en outre $S=\operatorname{Spec}(A)$ est affine,
on écrit $\mathbf V_A^n$ au lieu de $\mathbf V_S^n$ ; on a
$\mathbf V_A^n=\operatorname{Spec}(A[T_1,\ldots,T_n])$ où les $T_i$ sont des
indéterminées.

\noindent\textbf{$(\mathbf{Err}_{\mathrm{III}},\,15)$} Dans (II, 1.7.10),
ligne 10 de la p. 17, supprimer « \emph{géométriques} » ; lignes 12 et 16--17
de la p. 17, supprimer « géométrique ».

\noindent\textbf{$(\mathbf{Err}_{\mathrm{III}},\,16)$} Dans (II, 1.7.12),
ajouter : On pose de même :
\[
  \mathbf V(\mathscr O_S^n)=\mathbf V_S^n=S[T_1,\ldots,T_n].
\]

\noindent\textbf{$(\mathbf{Err}_{\mathrm{III}},\,17)$} Dans (II, 4.2.6),
ligne 7 de la p. 75, supprimer « \emph{géométriques} » ; ligne 8 de la p.
75, supprimer « \emph{géométrique} ».

\noindent\textbf{$(\mathbf{Err}_{\mathrm{III}},\,18)$} Dans (II, 4.6.13),
ligne 7 de la p. 92, le raisonnement doit être précisé, car avec les notations
de (4.4.10), il faut ici prouver que l'immersion $\Gamma_f$ est
\emph{quasi-compacte}, afin d'appliquer (i bis). En vertu de (4.6.4), pour
prouver que $\mathscr L$ est ample relativement à $f$, on peut se borner au
cas où $Y$ est affine. Notons d'autre part que dans chacune des hypothèses de
(v), $f$ est quasi-compact (I, 6.6.4). D'autre part, si $g$ est
\oldpage[III]{221}
séparé, $\Gamma_f$ est une immersion fermée (I, 5.4.3) donc quasi-compacte
(I, 6.6.4). Si au contraire $X$ est localement noethérien, comme $Y$ est affine
et $f$ quasi-compact, l'espace sous-jacent à $X$ est quasi-compact, donc
noethérien, et on peut de nouveau appliquer (I, 6.6.4) pour prouver que
$\Gamma_f$ est quasi-compacte.

\noindent\textbf{$(\mathbf{Err}_{\mathrm{III}},\,19)$} L'énoncé de (II,
5.2.2) est incorrect, les conditions $b)$, $c)$ et $c')$ n'étant pas locales
sur $Y$ lorsqu'on ne sait pas si pour un ouvert $U$ de $Y$, un
$(\mathscr O_X|f^{-1}(U))$-Module quasi-cohérent est restriction à $f^{-1}(U)$
d'un $\mathscr O_X$-Module quasi-cohérent. Il faut donc remplacer dans la ligne
16 de la p. 98 « \emph{Soit $f:X\to Y$ un morphisme séparé quasi-compact} »
par : « \emph{Soient $X$, $Y$ deux préschémas tels que $X$ soit un schéma ou
que l'espace sous-jacent à $X$ soit localement noethérien, et soit
$f:X\to Y$ un morphisme quasi-compact.} » Dans la démonstration, on observera
que les hypothèses entraînent que $f$ est \emph{séparé} lorsqu'on suppose que
$X$ est un schéma, en vertu de (I, 5.5.5 et 5.5.8) ; on peut donc appliquer
(I, 9.2.2) dans les deux cas.

\noindent\textbf{$(\mathbf{Err}_{\mathrm{III}},\,20)$} Dans (II, 6.2.3),
ajouter, après la ligne 18 du bas de la p. 115 :

\emph{On dit qu'un morphisme $f:X\to Y$ est quasi-fini en un point $x\in X$
s'il existe un voisinage ouvert affine $V$ de $y=f(x)$ et un voisinage ouvert
affine $U$ de $x$ tels que $f(U)\subset V$ et que le morphisme $U\to V$
restriction de $f$ soit quasi-fini. On dit qu'un morphisme $f:X\to Y$ est
localement quasi-fini s'il est quasi-fini en tout point de $X$.}

\noindent\textbf{$(\mathbf{Err}_{\mathrm{III}},\,21)$} Dans (II, 6.4.3), la
démonstration est incorrecte, les éléments d'un sous-$A$-module de type fini de
$K$ n'étant pas nécessairement entiers sur $A$. Remplacer les 9 dernières lignes
de la p. 121 par le texte suivant : Comme on peut supposer que $E$ est sans
torsion, donc fidèle, $E$ est aussi un $A[u]$-module fidèle ($A$ étant plongé
dans l'anneau des endomorphismes de $E$). Comme $E$ est un $A$-module de type
fini, il en résulte que $u$ est \emph{entier} sur $A$ (Bourbaki,
\emph{Alg. comm.}, chap. V, \S~1, n\textsuperscript{o}~1, lemme 1). On en
conclut immédiatement que les valeurs propres de $u\otimes1$ (dans une clôture
algébrique de $K$) sont des éléments \emph{entiers} sur $A$, et il en est donc
de même des $\sigma_i(u)$.

\noindent\textbf{$(\mathbf{Err}_{\mathrm{III}},\,22)$} Dans
$(0_{\mathrm{III}},\,9.1.1)$, ligne 18 du bas de la p. 12, supprimer
« ouverte ».

\noindent\textbf{$(\mathbf{Err}_{\mathrm{III}},\,23)$} Dans
$(0_{\mathrm{III}},\,11.7.3)$, ligne 10 de la p. 43, après $\mathsf C''$,
ajouter : tel que pour tout objet projectif $P$ de $\mathsf C'$ (resp. tout
objet projectif $P'$ de $\mathsf C''$) le foncteur
$A'\rightsquigarrow T(P,A')$ (resp. $A\rightsquigarrow T(A,P')$) soit exact
dans $\mathsf C'$ (resp. $\mathsf C$).

\noindent\textbf{$(\mathbf{Err}_{\mathrm{III}},\,24)$} L'énoncé de
$(0_{\mathrm{III}},\,13.7.7)$ est inexact, et doit être modifié comme suit :
ligne 6 de la p. 78, supprimer « et
$(\mathrm R^{n+1}T(A_k))_{k\in\mathbf Z}$ » ; ligne 7 de la p. 78, remplacer
« $\mathrm R'^nT(A)$ est un $S$-module de type fini » par
« $\mathrm R'^nT(A)$ et $\mathrm R'^{n+1}T(A)$ sont des $S$-modules de type
fini » ; lignes 12--13 de la p. 78, supprimer « et $p+q=n+1$ ». Dans la
démonstration, ligne 23 de la p. 78, supprimer « et pour $n+1$ ».

\noindent\textbf{$(\mathbf{Err}_{\mathrm{III}},\,25)$} Dans (III, 1.4.15),
ligne 6 du bas de la p. 92, remplacer « \emph{de type fini} » par
« \emph{quasi-compact} ».

\noindent\textbf{$(\mathbf{Err}_{\mathrm{III}},\,26)$} Dans (III, 2.2.4),
ligne 1 de la p. 102, remplacer « \emph{que les supports de $\mathscr F$ et de
$\mathscr H$ soient propres sur $Y$} » par « \emph{que le support de
$\operatorname{Im}(\mathscr F\to\mathscr G)$ soit propre sur $Y$} ». Dans la
démonstration, remplacer le texte des lignes 10 à 16, depuis « Cela
étant\ldots{} », par :

\oldpage[III]{222}
Posons
$\mathscr G_1=\operatorname{Im}(\mathscr F\to\mathscr G)
=\operatorname{Ker}(\mathscr G\to\mathscr H)$, qui est cohérent. Comme on a la
suite exacte
$0\to\mathscr G_1(n)\to\mathscr G(n)\to\mathscr H(n)$ et que le foncteur
$f_*$ est exact à gauche, la suite
$0\to f_*(\mathscr G_1(n))\to f_*(\mathscr G(n))\to f_*(\mathscr H(n))$ est
exacte pour tout $n$, et il suffit donc de montrer que, pour $n$ assez grand,
la suite $f_*(\mathscr F(n))\to f_*(\mathscr G_1(n))\to0$ est exacte. Autrement
dit, on peut se borner au cas où $\mathscr H=0$ et où le support de
$\mathscr G$ est propre sur $Y$, donc \emph{fermé dans $Z$} (II, 5.4.10) ;
$\mathscr G'=i_*(\mathscr G)$ est par suite un $\mathscr O_Z$-Module cohérent tel
que $\mathscr G'|X=\mathscr G$. On sait (I, 9.4.3) qu'il existe un
$\mathscr O_Z$-Module cohérent $\mathscr F'$ tel que
$\mathscr F'|X=\mathscr F$. En outre, comme $X$ est ouvert dans $Z$ et
$\operatorname{Supp}(\mathscr G)\subset X$ fermé dans $Z$, il est immédiat que
l'on définit un homomorphisme surjectif $u':\mathscr F'\to\mathscr G'$ de
faisceaux tel que $u'|X$ soit l'homomorphisme surjectif donné
$u:\mathscr F\to\mathscr G$, en prenant $u'|U=0$ pour tout ouvert $U$ de $Z$
ne rencontrant pas $\operatorname{Supp}(\mathscr G)$ et $u'|U=u|U$ pour tout
ouvert $U\subset X$. Cela étant, on voit comme au début de la démonstration de
(2.2.2) que l'on peut se borner au cas où $Y$ est affine, et il s'agit donc de
montrer que, pour $n$ assez grand, l'homomorphisme
$\Gamma(u):\Gamma(X,\mathscr F(n))\to\Gamma(X,\mathscr G(n))$ est surjectif. Or,
on a le diagramme commutatif
\[
\xymatrix@C=35mm@R=12mm{
\Gamma(Z,\mathscr F'(n)) \ar[r]^{\Gamma(u')} \ar[d]_v &
\Gamma(Z,\mathscr G'(n)) \ar[d]^w \\
\Gamma(X,\mathscr F(n)) \ar[r]_{\Gamma(u)} & \Gamma(X,\mathscr G(n))
}
\]
où $v$ et $w$ sont les homomorphismes de restriction. La définition de
$\mathscr G'$ montre en outre que $w$ est \emph{bijectif} ; par ailleurs, en
vertu de (2.2.3), $\Gamma(u')$ est \emph{surjectif} pour $n$ assez grand, donc
il en est de même de $\Gamma(u)$.

\noindent\textbf{$(\mathbf{Err}_{\mathrm{III}},\,27)$} Dans (III, 2.2.5),
après la ligne 7 du bas de la p. 102, ajouter :

(iii) Sous les hypothèses de (2.2.4) concernant $X$, $Y$, $f$ et $\mathscr L$,
il est immédiat que si $\mathscr H$ est un $\mathscr O_X$-Module cohérent dont
le support est \emph{propre sur $Y$}, un raisonnement analogue à celui de
(2.2.4) montre qu'il existe un entier $N$ tel que pour $n\geq N$, on ait
$\mathrm R^if_*(\mathscr H(n))=0$ pour tout $i>0$. On en conclut, par la suite
exacte de cohomologie, que si $u:\mathscr F\to\mathscr G$ est un homomorphisme
de $\mathscr O_X$-Modules cohérents tel que $\operatorname{Ker}(u)$ et
$\operatorname{Coker}(u)$ aient leurs supports \emph{propres sur $Y$}, alors il
existe $N$ tel que pour $n\geq N$, l'homomorphisme correspondant
$\mathrm R^if_*(\mathscr F(n))\to\mathrm R^if_*(\mathscr G(n))$ soit
\emph{bijectif} pour tout $i>0$.

\noindent\textbf{$(\mathbf{Err}_{\mathrm{III}},\,28)$} Dans (III, 4.4.9),
ligne 17 du bas de la p. 137, remplacer « \emph{Soient $Y$ un préschéma
intègre localement noethérien} » par « \emph{Soient $X$ et $Y$ deux
préschémas intègres localement noethériens} » ; ligne 16 du bas de la p.
137, remplacer « \emph{de type fini} » par « \emph{localement de type
fini} ». La démonstration est essentiellement inchangée,
$f^{-1}(y)$ étant localement de type fini sur $k(y)$, donc encore discret.

\noindent\textbf{$(\mathbf{Err}_{\mathrm{III}},\,29)$} dans (I, 9.3.4), ligne
19 du bas de la p. 173, remplacer « \emph{noethérien} » par
« \emph{quasi-compact} » ; lignes 18 et 19 du bas de la p. 173, remplacer
« \emph{cohérent} » par « \emph{quasi-cohérent de type fini} ». Dans la
démonstration, on se ramène au cas où

\oldpage[III]{223}
$X=\operatorname{Spec}(A)$, $\mathscr F=\widetilde M$, où $M$ est un $A$-module
de type fini, $\mathscr J=\widetilde{\mathfrak J}$, où $\mathfrak J$ est un
idéal de type fini de $A$ ; le reste du raisonnement est alors inchangé.

\noindent\textbf{$(\mathbf{Err}_{\mathrm{III}},\,30)$} Dans (I, 9.3.5),
remplacer les lignes 1 à 7 du bas de la p. 173 par le texte suivant :

\emph{Proposition (9.3.5). --- Soient $X$ un préschéma, $\mathscr F$ un
$\mathscr O_X$-Module quasi-cohérent de type fini. Alors il existe un
sous-préschéma fermé $Y$ de $X$, dont l'espace sous-jacent est égal à
$\operatorname{Supp}(\mathscr F)$, et un $\mathscr O_Y$-Module quasi-cohérent de
type fini $\mathscr G$ tels que, si $j:Y\to X$ est l'injection canonique,
$\mathscr F$ soit isomorphe à $j_*(\mathscr G)$.}

Il suffira de montrer que l'Idéal $\mathscr J$ de $\mathscr O_X$,
\emph{annulateur de $\mathscr F$}, est \emph{quasi-cohérent} ; on prendra alors
pour $Y$ le sous-préschéma fermé de $X$ défini par $\mathscr J$ (I, 4.1.2), et
comme $\mathscr J\mathscr F=0$, $\mathscr F$ est un
$(\mathscr O_X/\mathscr J)$-Module et on répondra à la question en prenant
$\mathscr G=j^*(\mathscr F)$. Pour voir que $\mathscr J$ est quasi-cohérent, on
peut (la question étant locale) se borner au cas où
$X=\operatorname{Spec}(A)$, $\mathscr F=\widetilde M$, où $M$ est un $A$-module
engendré par un nombre fini d'éléments $x_i\ (1\leq i\leq r)$ ; l'Idéal
$\mathscr J$ est alors l'intersection des annulateurs des $x_i$. Mais
l'annulateur de $x_i$ est le noyau de l'homomorphisme
$\mathscr O_X\to\mathscr F$ correspondant à l'homomorphisme $s\mapsto sx_i$ de
$A$ dans $M$ ; c'est donc bien un Idéal quasi-cohérent (I, 4.1.1) et toute
intersection finie de tels Idéaux est aussi un Idéal quasi-cohérent (I,
1.3.10).

\end{document}
